% =====================================================================
% sample_preamble.tex -- illustrative preamble for compiling the
% body-only .tex files in this project (hiemalis.tex, verna.tex,
% aestiva.tex, autumnalis.tex). Those files intentionally contain NO
% \documentclass, preamble, or \newcommand definitions of their own;
% this file shows one way to supply the command vocabulary and get a
% compilable document. The styling here (headings, rubric color,
% lettrine treatment) is a plain, generic stand-in for checking
% structure and catching LaTeX errors -- it is not the final intended
% typography/design of the printed breviary.
%
% Usage (XeLaTeX or LuaLaTeX required, for Unicode + fontspec):
%   % =====================================================================
% sample_preamble.tex -- illustrative preamble for compiling the
% body-only .tex files in this project (hiemalis.tex, verna.tex,
% aestiva.tex, autumnalis.tex). Those files intentionally contain NO
% \documentclass, preamble, or \newcommand definitions of their own;
% this file shows one way to supply the command vocabulary and get a
% compilable document. The styling here (headings, rubric color,
% lettrine treatment) is a plain, generic stand-in for checking
% structure and catching LaTeX errors -- it is not the final intended
% typography/design of the printed breviary.
%
% Usage (XeLaTeX or LuaLaTeX required, for Unicode + fontspec):
%   % =====================================================================
% sample_preamble.tex -- illustrative preamble for compiling the
% body-only .tex files in this project (hiemalis.tex, verna.tex,
% aestiva.tex, autumnalis.tex). Those files intentionally contain NO
% \documentclass, preamble, or \newcommand definitions of their own;
% this file shows one way to supply the command vocabulary and get a
% compilable document. The styling here (headings, rubric color,
% lettrine treatment) is a plain, generic stand-in for checking
% structure and catching LaTeX errors -- it is not the final intended
% typography/design of the printed breviary.
%
% Usage (XeLaTeX or LuaLaTeX required, for Unicode + fontspec):
%   % =====================================================================
% sample_preamble.tex -- illustrative preamble for compiling the
% body-only .tex files in this project (hiemalis.tex, verna.tex,
% aestiva.tex, autumnalis.tex). Those files intentionally contain NO
% \documentclass, preamble, or \newcommand definitions of their own;
% this file shows one way to supply the command vocabulary and get a
% compilable document. The styling here (headings, rubric color,
% lettrine treatment) is a plain, generic stand-in for checking
% structure and catching LaTeX errors -- it is not the final intended
% typography/design of the printed breviary.
%
% Usage (XeLaTeX or LuaLaTeX required, for Unicode + fontspec):
%   \input{sample_preamble.tex}  (or copy its contents into a wrapper)
%   \begin{document}
%   \input{autumnalis.tex}       % or hiemalis.tex / verna.tex / aestiva.tex
%   \end{document}
%
% Since autumnalis.tex, aestiva.tex, hiemalis.tex, and verna.tex all
% now embed the versicle/response marks (℣ U+2123, ℟ U+211F) and the
% blessing cross (☩ U+2629) as literal Unicode characters rather than
% as macros, and the main body font (DejaVu Serif) does not contain
% glyphs for them, \newunicodechar is used below to substitute a font
% that does (DejaVu Sans) automatically wherever those three
% characters occur, without having to touch the body text itself.
% =====================================================================

\usepackage{fontspec}
\usepackage{textcomp}
\usepackage{lettrine}
\usepackage[dvipsnames]{xcolor}
\usepackage[margin=2cm]{geometry}
\usepackage{parskip}
\usepackage{newunicodechar}

\setmainfont{DejaVu Serif}
\newfontfamily\liturgicalsymbolfont{DejaVu Sans}

\newunicodechar{℣}{{\liturgicalsymbolfont\char"2123}}  % VERSICLE
\newunicodechar{℟}{{\liturgicalsymbolfont\char"211F}}  % RESPONSE
\newunicodechar{☩}{{\liturgicalsymbolfont\char"2629}}  % cross (blessing)

% --- Headings (one required argument each) ---
\newcommand{\partheading}[1]{{\bigskip\noindent\LARGE\bfseries #1\par\medskip}}
\newcommand{\partsubtitle}[1]{{\noindent\large\itshape #1\par\medskip}}
\newcommand{\sectionheading}[1]{{\bigskip\noindent\Large\bfseries #1\par\medskip}}
\newcommand{\feriaheading}[1]{{\bigskip\noindent\large\bfseries #1\par\smallskip}}
\newcommand{\hourheading}[1]{{\bigskip\noindent\large\bfseries #1\par\smallskip}}
\newcommand{\nocturnoheading}[1]{{\medskip\noindent\bfseries #1\par\smallskip}}
\newcommand{\psalmheading}[1]{{\medskip\noindent\itshape #1\par\smallskip}}
\newcommand{\hymnusheading}[1]{{\medskip\noindent\bfseries #1\par\smallskip}}
\newcommand{\oratioheading}[1]{{\medskip\noindent\itshape #1\par\smallskip}}
\newcommand{\variantheading}[1]{{\medskip\noindent\itshape #1\par\smallskip}}
\newcommand{\absolutioheading}[1]{{\medskip\noindent\itshape #1\par\smallskip}}
\newcommand{\suffragiumheading}[1]{{\medskip\noindent\bfseries #1\par\smallskip}}
\newcommand{\symbolumheading}[1]{{\medskip\noindent\bfseries #1\par\smallskip}}
\newcommand{\antiphonaheading}[1]{{\medskip\noindent\bfseries #1\par\smallskip}}
\newcommand{\feastheading}[1]{{\bigskip\noindent\large\bfseries #1\par\smallskip}}
\newcommand{\sanctusheading}[1]{{\bigskip\noindent\large\bfseries #1\par\smallskip}}
\newcommand{\diesheading}[1]{{\medskip\noindent\bfseries #1\par\smallskip}}
\newcommand{\monthheading}[1]{{\bigskip\noindent\Large\bfseries #1\par\medskip}}

% --- Purely decorative divider (no argument) ---
\newcommand{\sectionornament}{{\par\noindent\centering\rule{0.3\linewidth}{0.4pt}\par\medskip}}

% --- Headings (two required arguments: title + citation) ---
\newcommand{\lectioheading}[2]{{\medskip\noindent\itshape #1\quad\textit{#2}\par\smallskip}}
\newcommand{\capitulumlabel}[2]{{\medskip\noindent\itshape #1\quad\textit{#2}\par\smallskip}}
\newcommand{\canticumheading}[2]{{\medskip\noindent\itshape #1\quad\textit{#2}\par\smallskip}}
\newcommand{\lectiolabel}[2]{{\medskip\noindent\itshape #1\quad\textit{#2}\par\smallskip}}

% --- Labels (one required argument) ---
\newcommand{\antlabel}[1]{\textbf{#1}}
\newcommand{\invitatlabel}[1]{\textbf{#1}}
\newcommand{\psalmlabel}[1]{\textbf{#1}}
\newcommand{\benedlabel}[1]{\textbf{#1}}
\newcommand{\benedictiolabel}[1]{\textbf{#1}}
\newcommand{\absolutiolabel}[1]{\textbf{#1}}
\newcommand{\formulalabel}[1]{\textbf{#1}}
\newcommand{\oratiolabel}[1]{\textbf{#1}}

% --- Rubric (red text; embedded rubrics may nest \textit, etc.) ---
\newcommand{\rubric}[1]{\textcolor{red}{#1}}
% Black text embedded within a red \rubric{} span (governing instruction:
% "embedded rubrics (black text within red-colored rubrics) are preserved").
\newcommand{\rubricblack}[1]{\textcolor{black}{#1}}

% \P (pilcrow), \dag, \ddag, \label, \pageref, \hspace, \textbf,
% \textsuperscript, \textit, \textsc, \textdagger, \textdaggerdbl are all
% standard LaTeX/textcomp commands and need no definition here.
  (or copy its contents into a wrapper)
%   \begin{document}
%   % =====================================================================
% autumnalis.tex -- LaTeX conversion of autumnalis.md (Breviarium
% Romanum, Pars Autumnalis), BODY CONTENT ONLY.
%
% No preamble, \documentclass, or \newcommand definitions are included
% here, per instruction -- the command vocabulary below follows the
% same legend established for hiemalis.tex, verna.tex, and aestiva.tex
% (one-column vs. two-column, red vs. black, etc.).
%
% BATCH 1 of N: Ordinarium Divini Officii, Ad Matutinum and Ad Laudes
% (autumnalis.md lines 1-463)
%
% NOTE: autumnalis.md lines 1-1141 (the entire Ordinarium Divini
% Officii) were verified via direct diff to be BYTE-IDENTICAL to the
% corresponding lines of aestiva.md -- both being non-Lenten/
% non-Paschal volumes, each omits the same "Tempore Quadragesimæ" /
% "Tempore Passionis" Invitatorium-Hymnus variants and the Paschaltide
% "Commemoratio de Cruce" that verna.tex's Ordinarium carries. Given
% this confirmed byte-for-byte identity (not a parallel or merely
% similar occurrence, but the literal same source text), the verified
% LaTeX already produced for aestiva.tex's Batches 1-3 is reused
% directly below rather than re-keyed by hand, to avoid reintroducing
% errors (such as the accent-dropping mistakes) already caught and
% fixed in that earlier work. Every batch from "Psalterium Breviarii
% Romani" onward is fresh transcription from autumnalis.md's own text,
% as usual.
% =====================================================================

% [ornamental image beginning Pars Autumnalis]
\partheading{Ordinarium Divini Officii}
\partsubtitle{Juxta Ritum Romanum Persolvendi}

\rubric{Clericis in sacris constitútis, qui íntegrum divínum Offícium, quamvis in partes distribútum, coram Ssmo Sacraménto, sive públicæ adoratióni expósito sive in tabernáculo adserváto, devóte recitáverint, concéditur: Indulgéntia plenária, si prætérea sacramentálem confessiónem instítuerint, ad eucharísticum Convívium accésserint et ad mentem Summi Pontíficis oráverint. \textit{Pius Pp. XI, 23 Octóbris 1930.}}

\rubric{Iis vero, qui divínum Offícium coram Ssmo Sacraménto, ut supra, ex parte tantum recitáverint, concéditur: Indulgéntia quingentórum diérum pro unaquáque hora canónica. \textit{Pius Pp. XI, 18 Maji 1933.}}

\rubric{Iisdem autem cléricis, qui divínum Offícium in álias preces commutátum rite obtinúerint, si preces ejúsmodi coram Ssmo Sacraménto, uti supra, devóte recitáverint, concéditur: Indulgéntia plenária, addítis sacramentáli confessióne, sacra Communióne et oratióne ad mentem Summi Pontíficis. \textit{Pius Pp. XI, 7 Novémbris 1932.}}

\sectionheading{Ante Divinum Officium}

\rubric{Antequam inchoëtur Offícium laudabíliter dícitur, sub singulári semper número, séquens Orátio; pro qua Summus Póntifex Pius Pp. XI sub die 17 Novémbris 1933 Indulgéntiam trium annórum concéssit.}

\oratioheading{Oratio}

\lettrine[lines=2]{A}{PERI}, Dómine, os meum ad benedicéndum nomen sanctum tuum: munda quoque cor meum ab ómnibus vanis, pervérsis et aliénis cogitatiónibus; intelléctum illúmina, afféctum inflámma, ut digne, atténte ac devóte hoc Offícium recitáre váleam, et exaudíri mérear ante conspéctum divínæ Majestátis tuæ. Per Christum, Dóminum nostrum. \R\ Amen.

Dómine, in unióne illíus divínæ intentiónis, qua ipse in terris laudes Deo persolvísti, has tibi Horas (\textit{vel} hanc tibi Horam) persólvo.

\sectionheading{Ante singulas Horas}

\rubric{Ante Matutínum, sicut et ante Laudes a Matutíno extra Chorum separátas, et ante omnes Horas, præter quam ad Completórium, dícitur secréto:}

\lettrine[lines=2]{P}{ATER} noster, qui es in cælis, sanctificétur nomen tuum. Advéniat regnum tuum. Fiat volúntas tua, sicut in cælo et in terra. Panem nostrum quotidiánum da nobis hódie. Et dimítte nobis débita nostra, sicut et nos dimíttimus debitóribus nostris. Et ne nos indúcas in tentatiónem: sed líbera nos a malo. Amen.

\lettrine[lines=2]{A}{VE}, María, grátia plena; Dóminus tecum: benedícta tu in muliéribus, et benedíctus fructus ventris tui Jesus. Sancta María, Mater Dei, ora pro nobis peccatóribus, nunc et in hora mortis nostræ. Amen.

\rubric{In princípio autem Matutíni ac Primæ additur secréto Symbolum Apostolórum:}

\lettrine[lines=2]{C}{REDO} in Deum, Patrem omnipoténtem, Creatórem cæli et terræ. Et in Jesum Christum Fílium ejus únicum, Dóminum nostrum: qui concéptus est de Spíritu Sancto, natus ex María Vírgine, passus sub Póntio Piláto, crucifíxus, mórtuus, et sepúltus: descéndit ad ínferos; tértia die resurréxit a mórtuis; ascéndit ad cælos; sedet ad déxteram Dei Patris omnipoténtis: inde ventúrus est judicáre vivos et mórtuos. Credo in Spíritum Sanctum, sanctam Ecclésiam cathólicam, Sanctórum communiónem, remissiónem peccatórum, carnis resurrectiónem, vitam ætérnam. Amen.

\rubric{Officium ad Matutinum et ad reliquas Horas, nisi aliter suo loco indicetur, per totum annum persolvitur juxta formam, quæ in Rubricis hujus Ordinarii exhibetur.}

\hourheading{Ad Matutinum}

Pater noster, Ave María et Credo.

\rubric{Deinde, clara voce, dicitur Versus:}

\V\ \lettrine[lines=2]{D}{OMINE}, lábia mea apéries. \R\ Et os meum annuntiábit laudem tuam. \V\ Deus, in adjutórium meum inténde. \R\ Dómine, ad adjuvándum me festína.

Glória Patri, et Fílio, et Spirítui Sancto. Sicut erat in princípio, et nunc, et semper, et in sæcula sæculórum. Amen. Allelúja.

\rubric{Postea dicitur conveniens Invitatorium, quod ante Psalmum bis integrum recitatur, et ad singulos ejusdem Psalmi versus vel integrum vel dimidiatum ab Asterisco * alternis vicibus, ut infra, repetitur.}

\psalmheading{Psalmus 94}

\lettrine[lines=2]{V}{ENITE}, exsultémus Dómino, jubilémus Deo, salutári nostro: præoccupémus fáciem ejus in confessióne, et in psalmis jubilémus ei.

\rubric{Repetitur integrum Invitatorium.}

Quóniam Deus magnus Dóminus, et Rex magnus super omnes deos: quóniam non repéllet Dóminus plebem suam: quia in manu ejus sunt omnes fines terræ, et altitúdines móntium ipse cónspicit.

\rubric{Repetitur altera pars Invitatorii.}

\rubric{In sequenti Psalmi versu, ad verba veníte, adorémus, et procidámus ante Deum, genuflectitur.}

Quóniam ipsíus est mare, et ipse fecit illud, et áridam fundavérunt manus ejus: veníte, adorémus, et procidámus ante Deum: plorémus coram Dómino, qui fecit nos, quia ipse est Dóminus, Deus noster; nos autem pópulus ejus, et oves páscuæ ejus.

\rubric{Repetitur integrum Invitatorium.}

Hódie, si vocem ejus audíeritis, nolíte obduráre corda vestra, sicut in exacerbatióne secúndum diem tentatiónis in desérto: ubi tentavérunt me patres vestri, probavérunt et vidérunt ópera mea.

\rubric{Repetitur altera pars Invitatorii.}

Quadragínta annis próximus fui generatióni huic, et dixi: Semper hi errant corde; ipsi vero non cognovérunt vias meas: quibus jurávi in ira mea: Si introíbunt in réquiem meam.

\rubric{Repetitur integrum Invitatorium.}

Glória Patri, et Fílio, et Spirítui Sancto. Sicut erat in princípio, et nunc, et semper, et in sæcula sæculórum. Amen.

\rubric{Repetitur altera pars Invitatorii.}

\rubric{Denique repetitur integrum Invitatorium.}

\rubric{Expleto Psalmo, dicitur Hymnus Invitatorio respondens.}

\variantheading{In Officio dominicali}

\rubric{In omnibus Dominicis, extra Octavas Ssmi Corporis Christi ac Sacrat. Cordis Jesu, Invitatorium et Hymnus ut in Psalterio.}

\variantheading{In Officio feriali}

\rubric{In omnibus Feriis, Invitatorium et Hymnus de Feria currenti, ut in Psalterio.}

\variantheading{In Festis}

\rubric{In Officio cujuslibet Festi vel Octavæ, ac S. Mariæ in Sabbato, Invitatorium et Hymnus ut in Proprio vel Communi; in Dominicis autem infra Octavas Ssmi Corporis Christi ac Sacratissimi Cordis Jesu, si de eis fiat Officium, Invitatorium et Hymnus ut in Proprio de Tempore. Quæ norma ad reliquas item Horas erit observanda pro illis Officii partibus, quæ juxta Rubricas ex Ordinario vel Psalterio non fuerint in Festis desumendæ.}

\variantheading{In Officio novem Lectionum}

\rubric{Expleto Hymno, dicuntur Antiphonæ convenientes, quæ in Officiis ritus duplicis ante et post Psalmos integræ recitantur; in Officiis autem ritus semiduplicis initio Psalmi inchoantur tantum et usque ad Asteriscum * perducuntur, atque in fine integræ pronuntiantur; servata alterutra norma etiam in omnibus Antiphonis, quæ ad Laudes ac Vesperas adhibentur.}

\nocturnoheading{In I Nocturno}

\rubric{Sub congruentibus Antiphonis dicuntur tres Psalmi, ac deinde subjungitur Versus, prouti Officium occurrens requirit.}

\rubric{Post Versum cujuslibet Nocturni dicitur}

Pater noster \rubric{secreto usque ad} \V\ Et ne nos indúcas in tentatiónem. \R\ Sed líbera nos a malo.

\absolutioheading{Absolutio}

\lettrine[lines=2]{E}{XÁUDI}, Dómine Jesu Christe, preces servórum tuórum, et miserére nobis: Qui cum Patre et Spíritu Sancto vivis et regnas in sæcula sæculórum. \R\ Amen.

\V\ Jube, domne, benedícere.

\rubric{¶ Extra Chorum, quando ab uno tantum recitatur Officium, ante singulas Lectiones Matutini atque ad Lectionem brevem Primæ et Completorii, dicitur: Jube, Dómine, benedícere; et subjungitur congruens Benedictio. Ab Episcopo autem, ultimam Matutini Lectionem cantaturo, item dicitur: Jube, Dómine, benedícere; et respondetur a Choro: Amen.}

\variantheading{Pro I Lectione}

\benedlabel{Bened.} Benedictióne perpétua benedícat nos Pater ætérnus. \R\ Amen.

\rubric{Deinde dicuntur in unoquoque Nocturno Lectiones, prouti Officium occurrens requirit, et in fine cujuslibet Lectionis additur:}

\V\ Tu autem, Dómine, miserére nobis. \R\ Deo grátias.

\rubric{Post quamlibet vero Lectionem, quæ Hymnum Te Deum immediate non præcedat, congruens dicitur Responsorium, et in fine ultimi Responsorii cujusque Nocturni additur Versus: Glória Patri, et Fílio, et Spirítui Sancto, et Responsorium ab Asterisco *, et quidem a secundo, si duo fuerint, repetitur.}

\variantheading{Pro II Lectione}

\benedlabel{Bened.} Unigénitus Dei Fílius nos benedícere et adjuváre dignétur. \R\ Amen.

\variantheading{Pro III Lectione}

\benedlabel{Bened.} Spíritus Sancti grátia illúminet sensus et corda nostra. \R\ Amen.

\nocturnoheading{In II Nocturno}

\rubric{Sub congruentibus item Antiphonis dicuntur tres Psalmi et Versus, sicut in I Nocturno.}

\rubric{Post Versum dicitur.}

Pater noster \rubric{secreto usque ad} \V\ Et ne nos indúcas in tentatiónem. \R\ Sed líbera nos a malo.

\absolutioheading{Absolutio}

\lettrine[lines=2]{I}{PSIUS} píetas et misericórdia nos ádjuvet, qui cum Patre et Spíritu Sancto vivit et regnat in sæcula sæculórum. \R\ Amen.

\V\ Jube, domne, benedícere.

\variantheading{Pro IV Lectione}

\benedlabel{Bened.} Deus Pater omnípotens sit nobis propítius et clemens. \R\ Amen.

\variantheading{Pro V Lectione}

\benedlabel{Bened.} Christus perpétuæ det nobis gáudia vitæ. \R\ Amen.

\variantheading{Pro VI Lectione}

\benedlabel{Bened.} Ignem sui amóris accéndat Deus in córdibus nostris. \R\ Amen.

\nocturnoheading{In III Nocturno}

\rubric{Sub congruentibus denique Antiphonis dicuntur tres Psalmi et Versus, sicut in I et II Nocturno.}

\rubric{Post Versum dicitur}

Pater noster \rubric{secreto usque ad} \V\ Et ne nos indúcas in tentatiónem. \R\ Sed líbera nos a malo.

\absolutioheading{Absolutio}

\lettrine[lines=2]{A}{} \textsc{vinculis} peccatórum nostrórum absólvat nos omnípotens et miséricors Dóminus. \R\ Amen.

\V\ Jube, domne, benedícere.

\variantheading{Pro VII Lectione}

\benedlabel{Bened.} Evangélica léctio sit nobis salus et protéctio. \R\ Amen.

\variantheading{Pro VIII Lectione}

\rubric{In Festis Domini et in Dominicis:}

\benedlabel{Bened.} Divínum auxílium máneat semper nobíscum. \R\ Amen.

\rubric{In Festis beatæ Mariæ Virginis:}

\benedlabel{Bened.} Cujus festum cólimus, ipsa Virgo virgínum intercédat pro nobis ad Dóminum. \R\ Amen.

\rubric{In Festis Sanctorum:}

\benedlabel{Bened.} Cujus (\textit{vel} Quorum \textit{aut} Quarum) festum cólimus, ipse (\textit{vel} ipsa \textit{aut} ipsi \textit{vel} ipsæ) intercédat (\textit{vel} intercédant) pro nobis ad Dóminum. \R\ Amen.

\variantheading{Pro IX Lectione}

\benedlabel{Bened.} Ad societátem cívium supernórum perdúcat nos Rex Angelórum. \R\ Amen.

\rubric{Si autem legenda sit Lectio ix de Homilia cum Evangelio Dominicæ, vel Feriæ, aut Vigiliæ:}

\benedlabel{Bened.} Per evangélica dicta deleántur nostra delícta. \R\ Amen.

\variantheading{In Officio trium Lectionum}

\nocturnoheading{In Nocturno}

\rubric{In Officiis trium Lectionum, post Hymnum dicuntur Antiphonæ convenientes, quæ initio Psalmi inchoantur tantum et usque ad Asteriscum * perducuntur, ac deinde in fine integræ pronuntiantur.}

\rubric{Sub eisdem vero Antiphonis dicuntur novem Psalmi Feriæ currentis, quibus subjungitur Versus in III Nocturno positus, omissis Versibus pro I et II Nocturno assignatis.}

\rubric{Post Versum dicitur}

Pater noster \rubric{secreto usque ad} \V\ Et ne nos indúcas in tentatiónem. \R\ Sed líbera nos a malo.

\absolutioheading{Absolutio}

\rubric{Feria II et V}

Exáudi, Dómine, \rubric{ut in I Nocturno.}

\rubric{Feria III et VI}

Ipsíus píetas, \rubric{ut in II Nocturno.}

\rubric{Feria IV et Sabbato}

A vínculis, \rubric{ut in III Nocturno Officii novem Lectionum.}

\rubric{Deinde leguntur Lectiones cum Responsoriis, prouti Officium occurrens requirit; et ante eas dicuntur sequentes}

\variantheading{Benedictiones}

\rubric{In Feriis, quando legitur Homilia cum Evangelio, erit:}

\benedlabel{1 Bened.} Evangélica léctio sit nobis salus et protéctio. \R\ Amen.

\benedlabel{2 Bened.} Divínum auxílium máneat semper nobíscum. \R\ Amen.

\benedlabel{3 Bened.} Ad societátem cívium supernórum perdúcat nos Rex Angelórum. \R\ Amen.

\rubric{In Feriis, quando non legitur Homilia cum Evangelio, Feria II et V Benedictiones ut in I Nocturno, Feria III et VI ut in II Nocturno Officii novem Lectionum; Feria autem IV et Sabbato erit:}

\benedlabel{1 Bened.} Ille nos benedícat, qui sine fine vivit et regnat. \R\ Amen.

\benedlabel{2 Bened.} Divínum auxílium máneat semper nobíscum. \R\ Amen.

\benedlabel{3 Bened.} Ad societátem cívium supernórum perdúcat nos Rex Angelórum. \R\ Amen.

\rubric{In Festis Sanctorum erit:}

\benedlabel{1 Bened.} Ille nos benedícat, qui sine fine vivit et regnat. \R\ Amen.

\benedlabel{2 Bened.} Cujus (\textit{vel} Quorum \textit{aut} Quarum) festum cólimus, ipse (\textit{vel} ipsa \textit{aut} ipsi \textit{vel} ipsæ) intercédat (\textit{vel} intercédant) pro nobis ad Dóminum. \R\ Amen.

\benedlabel{3 Bened.} Ad societátem cívium supernórum perdúcat nos Rex Angelórum. \R\ Amen.

\rubric{In Officio de sancta María in Sábbato Absolútio et Benedictiónes, quæ ibídem ponúntur.}

\rubric{Post últimam Lectiónem, in ómnibus Domínicis per Annum minóribus, in Festis cujúsvis ritus, per omnes Octávas, et in Offício sanctæ Maríæ in Sábbato, dícitur Hymnus Ambrosiánus. Et numquam dícitur in feriáli Offício. Quando vero Hymnus prædíctus omíttitur, ejus loco dícitur IX aut III Responsórium.}

\hymnusheading{Hymnus Ambrosianus}

\lettrine[lines=2]{T}{E} Deum laudámus: * te Dóminum confitémur.

Te ætérnum Patrem * omnis terra venerátur.

Tibi omnes Angeli, * tibi Cæli, et univérsæ Potestátes:

Tibi Chérubim et Séraphim * incessábili voce proclámant:

Sanctus, * Sanctus, * Sanctus * Dóminus, Deus Sábaoth.

Pleni sunt cæli et terra * majestátis glóriæ tuæ.

Te gloriósus * Apostolórum chorus,

Te Prophetárum * laudábilis númerus,

Te Mártyrum candidátus * laudat exércitus.

Te per orbem terrárum * sancta confitétur Ecclésia,

Patrem * imménsæ majestátis;

Venerándum tuum verum * et únicum Fílium;

Sanctum quoque * Paráclitum Spíritum.

Tu Rex glóriæ, * Christe.

Tu Patris * sempitérnus es Fílius.

Tu, ad liberándum susceptúrus hóminem, * non horruísti Vírginis úterum.

Tu, devícto mortis acúleo, * aperuísti credéntibus regna cælórum.

Tu ad déxteram Dei sedes, * in glória Patris.

Judex créderis * esse ventúrus.

\rubric{¶ Sequens versus dicitur flexis genibus.}

Te ergo quæsumus, tuis fámulis súbveni, * quos pretióso sánguine redemísti.

Ætérna fac cum Sanctis tuis * in glória numerári.

Salvum fac pópulum tuum, Dómine, * et bénedic hereditáti tuæ.

Et rege eos, * et extólle illos usque in ætérnum.

Per síngulos dies * benedícimus te;

Et laudámus nomen tuum in séculum, * et in séculum séculi.

Dignáre, Dómine, die isto * sine peccáto nos custodíre.

Miserére nostri, Dómine, * miserére nostri.

Fiat misericórdia tua, Dómine, super nos, * quemádmodum sperávimus in te.

In te, Dómine, sperávi: * non confúndar in ætérnum.

\rubric{Dicto Te Deum, aut ultimo Responsorio, statim incipiuntur Laudes a Versu Deus, in adjutórium, ut infra.}

\rubric{In privata vero recitatione Matutinum a Laudibus separari potest; quo in casu, post Hymnum Te Deum, vel post ultimum Responsorium, dicitur:}

\V\ Dóminus vobíscum. \R\ Et cum spíritu tuo.

\rubric{¶ Hic Versus Dóminus vobíscum non dicitur ab eo, qui non est saltem in ordine Diaconatus; sed ejus loco substituitur, etiam in aliis Horis, quoties in antecedentibus Precibus dictus non fuerit:}

\V\ Dómine, exáudi oratiónem meam. \R\ Et clamor meus ad te véniat.

Orémus.

\rubric{Et dicitur Oratio ut ad Laudes.}

\rubric{Post Orationem additur:}

\V\ Dóminus vobíscum. \R\ Et cum spíritu tuo.

\V\ Benedicámus Dómino. \R\ Deo grátias.

\V\ Fidélium ánimæ per misericórdiam Dei requiéscant in pace. \R\ Amen.

\rubric{Deinde dicitur Pater noster totum secreto.}

\rubric{In publica vero recitatione Matutinum in Choro a Laudibus separari non potest.}

\hourheading{Ad Laudes}

\rubric{Si Laudes extra Chorum separentur a Matutino, ante eas dicitur secreto Pater noster et Ave María; secus absolute incipiuntur, ut sequitur:}

\V\ Deus, in adjutórium.

\rubric{Deinde, sub congruentibus Antiphonis, dicuntur Psalmi et Canticum, prouti Officium occurrens requirit.}

\rubric{Repetita vero post ultimum Psalmum Antiphona, dicuntur Capitulum, Hymnus et Versus, prouti Officium occurrens exigit.}

\variantheading{In Officio dominicali}

\rubric{In omnibus Dominicis, extra Octavas Ssmi Corporis Christi ac Sacratissimi Cordis Jesu, Capitulum, Hymnus et Versus ut in Psalterio.}

\variantheading{In Officio feriali}

\rubric{In omnibus Feriis, Capitulum, Hymnus et Versus de Feria currenti, ut in Psalterio.}

\rubric{Recitato Versu, dicitur cum Antiphona convenienti sequens}

\canticumheading{Canticum Zachariæ}{Luc. 1, 68-79}

\lettrine[lines=2]{B}{ENEDICTUS} \textdagger\ Dóminus, Deus Israël, * quia visitávit, et fecit redemptiónem plebis suæ:

Et eréxit cornu salútis nobis: * in domo David, púeri sui.

Sicut locútus est per os sanctórum, * qui a sæculo sunt, Prophetárum ejus:

Salútem ex inimícis nostris, * et de manu ómnium, qui odérunt nos:

Ad faciéndam misericórdiam cum pátribus nostris: * et memorári testaménti sui sancti.

Jusjurándum, quod jurávit ad Abraham, patrem nostrum, * datúrum se nobis:

Ut sine timóre, de manu inimicórum nostrórum liberáti, * serviámus illi.

In sanctitáte, et justítia coram ipso, * ómnibus diébus nostris.

Et tu, puer, Prophéta Altíssimi vocáberis: * præíbis enim ante fáciem Dómini paráre vias ejus:

Ad dandam sciéntiam salútis plebi ejus: * in remissiónem peccatórum eórum:

Per víscera misericórdiæ Dei nostri: * in quibus visitávit nos, óriens ex alto:

Illumináre his, qui in ténebris, et in umbra mortis sedent: * ad dirigéndos pedes nostros in viam pacis.

Glória Patri, et Fílio, * et Spirítui Sancto.

Sicut erat in princípio, et nunc, et semper, * et in sæcula sæculórum. Amen.

\rubric{¶ Hic Versus Glória Patri dicitur in fine omnium Canticorum, nisi aliter adnotetur.}

\rubric{Expleto Cantico ac repetita Antiphona, si Preces non fuerint recitandæ, statim dicitur Versus Dóminus vobíscum cum Oratione et reliquis omnibus usque ad finem Horæ, ut infra.}

\rubric{In omnibus vero Vigiliis, si fiat Officium de Vigilia currenti, post repetitam ad Benedictus Antiphonam dicuntur flexis genibus sequentes Preces, quæ aliis temporibus omittuntur:}

KYRIE, eléison. Christe, eléison. Kýrie, eléison.

\rubric{Pater noster, quod dicitur a solo Hebdomadario totum clara voce, usque ad} \V\ Et ne nos indúcas in tentatiónem. \R\ Sed líbera nos a malo.

\V\ Ego dixi: Dómine, miserére mei. \R\ Sana ánimam meam, quia peccávi tibi. \V\ Convértere, Dómine, úsquequo? \R\ Et deprecábilis esto super servos tuos. \V\ Fiat misericórdia tua, Dómine, super nos. \R\ Quemádmodum sperávimus in te. \V\ Sacerdótes tui induántur justítiam. \R\ Et sancti tui exsúltent.

\V\ Orémus pro beatíssimo Papa nostro N. \R\ Dóminus consérvet eum, et vivíficet eum, et beátum fáciat eum in terra, et non tradat eum in ánimam inimicórum ejus.

\V\ Orémus et pro Antístite nostro N. \R\ Stet et pascat in fortitúdine tua, Dómine, in sublimitáte nóminis tui.

\rubric{¶ Romæ præcedens Versus cum suo Responsorio omittitur; alibi vero, ad litteram N., ab omnibus nomen diœcesani Episcopi exprimatur. Vacante Apostolica vel Episcopali Sede, alteruter vel uterque respectivus Versus cum suo Responsorio prætéritur.}

\V\ Dómine, salvum fac regem. \R\ Et exáudi nos in die, qua invocavérimus te. \V\ Salvum fac pópulum tuum, Dómine, et bénedic hereditáti tuæ. \R\ Et rege eos, et extólle illos usque in ætérnum. \V\ Meménto Congregatiónis tuæ. \R\ Quam possedísti ab inítio. \V\ Fiat pax in virtúte tua. \R\ Et abundántia in túrribus tuis. \V\ Orémus pro benefactóribus nostris. \R\ Retribúere dignáre, Dómine, ómnibus, nobis bona faciéntibus propter nomen tuum, vitam ætérnam. Amen. \V\ Orémus pro fidélibus defúnctis. \R\ Réquiem ætérnam dona eis, Dómine, et lux perpétua lúceat eis. \V\ Requiéscant in pace. \R\ Amen. \V\ Pro frátribus nostris abséntibus. \R\ Salvos fac servos tuos, Deus meus, sperántes in te. \V\ Pro afflíctis et captívis. \R\ Líbera eos, Deus Israël, ex ómnibus tribulatiónibus suis. \V\ Mitte eis, Dómine, auxílium de sancto. \R\ Et de Sion tuére eos. \V\ Dómine, Deus virtútum, convérte nos. \R\ Et osténde fáciem tuam, et salvi érimus. \V\ Exsúrge, Christe, ádjuva nos. \R\ Et líbera nos propter nomen tuum. \V\ Dómine, exáudi oratiónem meam. \R\ Et clamor meus ad te véniat.

\rubric{Absolutis Precibus, vel, si Preces locum non habeant, repetita Antiphona post Benedictus, statim dicitur:}

\V\ Dóminus vobíscum. \R\ Et cum spíritu tuo.

Orémus.

\rubric{Et dicitur Oratio conveniens. Postea fiunt Commemorationes, si occurrant.}

\rubric{Deinde in Officio semiduplici aut simplici tam de Tempore quam de Sanctis, præter omnes dies, in quibus habita sit Commemoratio Duplicis, vel factum fuerit Officium aut Commemoratio cujusvis Octavæ, fit sequens}

\suffragiumheading{Suffragium de omnibus Sanctis}

\antlabel{Ant.} Beáta Dei Génitrix Virgo María, Sanctíque omnes intercédant pro nobis ad Dóminum.

\V\ Mirificávit Dóminus Sanctos suos. \R\ Et exaudívit eos clamántes ad se.

Orémus. \oratiolabel{Oratio}

\lettrine[lines=2]{A}{} \textsc{cunctis} nos, quæsumus, Dómine, mentis et córporis defénde perículis: et, intercedénte beáta et gloriósa semper Vírgine Dei Genitríce María, cum beáto Joseph, beátis Apóstolis tuis Petro et Paulo, atque beáto N. et ómnibus Sanctis, salútem nobis tríbue benígnus et pacem; ut, destrúctis adversitátibus et erróribus univérsis, Ecclésia tua secúra tibi sérviat libertáte. Per eúmdem Dóminum nostrum. \R\ Amen.

\rubric{¶ In hac Oratione, ad litteram N., exprimitur nomen Titularis propriæ Ecclesiæ, dummodo Titulus non sit Persona divina vel Mysterium Domini, aut de ipso ratione Vigiliæ vel Festi factum non fuerit Officium vel Commemoratio, aut illius nomen non sit in eadem Suffragii Oratione expressum; ac nomina sanctorum Angelorum et sancti Joannis Baptistæ, si Titulares fuerint, præponuntur nomini sancti Joseph. In hisce autem omnibus casibus omittuntur verba atque beáto N.}

\rubric{Si vero Officium vel Commemoratio fuerit de beáta María Vírgine, tunc dicitur ut infra; atque insuper, sicubi Titularis juxta superiorem Rubricam ante omnes fuerit nominandus, ad Orationem particula cum omittitur ante nomen sanctorum Petri et Pauli ac præponitur nomini sancti Joseph.}

\antlabel{Ant.} Sancti omnes intercédant pro nobis ad Dóminum.

\V\ Mirificávit Dóminus Sanctos suos. \R\ Et exaudívit eos clamántes ad se.

Orémus. \oratiolabel{Oratio}

\lettrine[lines=2]{A}{} \textsc{cunctis} nos, quæsumus, Dómine, mentis et córporis defénde perículis: et, intercedénte beáto Joseph, cum beátis Apóstolis tuis Petro et Paulo, atque beáto N. et ómnibus Sanctis, salútem nobis tríbue benígnus et pacem; ut, destrúctis adversitátibus et erróribus univérsis, Ecclésia tua secúra tibi sérviat libertáte. Per Dóminum. \R\ Amen.

\rubric{Post ultimam Orationem additur:}

\V\ Dóminus vobíscum. \R\ Et cum spíritu tuo.

\V\ Benedicámus Dómino. \R\ Deo grátias.

\V\ Fidélium ánimæ per misericórdiam Dei requiéscant in pace. \R\ Amen.

\rubric{Deinde, in publica Officii recitatione, si discedendum sit a Choro, itemque in recitatione privata, si tunc terminetur Officium; alioquin, in fine ultimæ Horæ dicitur Pater noster, totum secreto; eoque recitato, subjungitur:}

\V\ Dóminus det nobis suam pacem. \R\ Et vitam ætérnam. Amen.

\rubric{Et immediate dicitur, cum suis Versu et Oratione, finalis beatæ Mariæ Virginis Antiphona, quæ inferius post Completorium assignatur.}

\rubric{Postea concluditur:}

\V\ Divínum auxílium máneat semper nobíscum. \R\ Amen.

% ===== END OF BATCH 1 -- Covers autumnalis.md lines 1-463 (reused
% verbatim from aestiva.tex's Batch 1, per the byte-identity noted
% above): the opening of the Ordinarium Divini Officii (its
% introductory indulgence rubrics and the "Aperi, Dómine" Oratio),
% "Ante singulas Horas" (the Pater noster, Ave María, and Credo
% spelled out in full for the first and only time -- each given a
% lettrine since the source presents them in full capitals as
% complete standalone prayers, unlike their thousands of later
% plain-text abbreviated recurrences throughout the volume), "Ad
% Matutinum" (the Invitatorium Psalm 94, the three-Nocturn apparatus
% with its Absolutios, Benedictiones, and "Pro N Lectione" headings,
% the Officium trium Lectionum variant, and the Te Deum under its
% customary heading "Hymnus Ambrosianus"), and "Ad Laudes" through
% the Suffragium de omnibus Sanctis (in both its Marian-titular and
% Joseph-titular forms) and the closing versicles. This Pars
% Autumnalis Ordinarium matches its Pars Aestiva counterpart exactly
% (both being non-Lenten/non-Paschal volumes) and is notably simpler
% than Pars Verna's (Batch 1 of verna.tex), which carries the extra
% "Tempore Quadragesimæ"/"Tempore Passionis" Invitatorium-Hymnus
% variants and the Paschaltide "Commemoratio de Cruce". Roman
% numerals in the "Pro N Lectione" headings are capitalized (e.g.
% \variantheading{Pro I Lectione}) though the source itself gives them
% lowercase ("Pro i Lectione"), matching the established
% normalization already present throughout hiemalis.tex, verna.tex,
% and aestiva.tex for this specific heading type. The document
% continues with "Ad Primam" (autumnalis.md line 465) =====
\hourheading{Ad Primam}

Pater, Ave et Credo.

\V\ Deus, in adjutórium.

\hymnusheading{Hymnus}

\lettrine[lines=2]{J}{AM} lucis orto sídere,\\
Deum precémur súpplices,\\
Ut in diúrnis áctibus\\
Nos servet a nocéntibus.

Linguam refrénans témperet,\\
Ne litis horror ínsonet:\\
Visum fovéndo cóntegat,\\
Ne vanitátes háuriat.

Sint pura cordis íntima,\\
Absístat et vecórdia:\\
Carnis terat supérbiam\\
Potus cibíque párcitas.

Ut, cum dies abscésserit\\
Noctémque sors redúxerit,\\
Mundi per abstinéntiam\\
Ipsi canámus glóriam.

Deo Patri sit glória,\\
Ejúsque soli Fílio,\\
Cum Spíritu Paráclito\\
Nunc et per omne sæculum. Amen.

\rubric{¶ Conclusio communis, in Hymno præcedenti, et in aliis idem metrum habentibus, semper omittitur, quando specialis in omnibus Horis adhibenda præscribitur; et, si plures Conclusiones propriæ occurrant, sumitur Conclusio Officii currentis, aut secus Officii ipsa die et primo quidem loco inter cetera propriam Conclusionem habentia commemorandi, aut demum de occurrenti Octava communi.}

\rubric{Expleto Hymno, conveniens dicitur Antiphona usque ad Asteriscum *, prouti occurrens Officium requirit.}

\variantheading{In Officio dominicali}

\rubric{In omnibus Dominicis, extra Octavas Ssmi Corporis Christi et Sacrat. Cordis Jesu, Antiphona ut in Psalterio.}

\variantheading{In Officio feriali}

\rubric{In omnibus Feriis, Antiphona de Feria currenti ut in Psalterio.}

\rubric{Post enuntiatam Antiphonam dicuntur tres Psalmi Officii diei congruentes.}

\rubric{¶ Quando autem adhibetur posterius Laudum schema, Psalmus, in priori Laudum ferialium schemate omissus, ad Primam Feriæ currentis post alios Psalmos resumitur.}

\rubric{Repetita post ultimum Psalmum integre Antiphona, conveniens dicitur Capitulum; nimirum:}

\rubric{In omnibus Dominicis, in Officio cujuslibet Festi vel Octavæ ac sanctæ Mariæ in Sabbato:}

\capitulumlabel{Capitulum}{1 Tim. 1, 17}

\lettrine[lines=2]{R}{EGI} sæculórum immortáli et invisíbili, soli Deo honor et glória in sæcula sæculórum. Amen. \R\ Deo grátias.

\rubric{¶ Sic respondetur in fine omnium Capitulorum.}

\rubric{In omnibus autem Feriis et in Vigiliis communibus:}

\capitulumlabel{Capitulum}{Zach. 8, 19}

\lettrine[lines=2]{P}{ACEM} et veritátem dilígite, ait Dóminus omnípotens.

\rubric{Dicto Capitulo, subjungitur Responsorium breve, ut sequitur. Versus tamen Qui sedes ad dexteram Patris semper omittitur, quando specialis Versus adhibendus præscribitur, Officii nempe currentis, aut secus Officii ad Laudes et primo quidem loco inter cetera specialem Versum habentia commemorati, aut demum de occurrenti Octava communi proprius.}

\R\ br. Christe, Fili Dei vivi, * Miserére nobis. \rubric{Et repetitur:} Christe, Fili Dei vivi, miserére nobis. \V\ Qui sedes ad déxteram Patris. \R\ Miserére nobis. \V\ Glória Patri, et Fílio, et Spirítui Sancto. \R\ Christe, Fili Dei vivi, miserére nobis. \V\ Exsúrge, Christe, ádjuva nos. \R\ Et líbera nos propter nomen tuum.

\rubric{¶ Ita dicuntur semper Responsoria brevia, etiam ad Tertiam, Sextam, Nonam et Completorium.}

\rubric{Expleto Responsorio brevi, si Preces non fuerint recitandæ, statim dicitur Versus Dóminus vobíscum cum subsequenti Oratione et reliquis omnibus usque ad finem Horæ ut infra.}

\rubric{In quolibet Officio semiduplici, in Officio sanctæ Mariæ in Sabbato, in Festis simplicibus, atque in Feriis per Annum communibus, post Responsorium breve dicuntur sequentes Preces dominicales; quæ tamen omittuntur omnibus diebus, in quibus ad Laudes habita sit Commemoratio Duplicis, vel factum fuerit Officium aut Commemoratio cujusvis Octavæ:}

KYRIE, eléison. Christe, eléison. Kýrie, eléison.

\rubric{Pater noster secreto usque ad} \V\ Et ne nos indúcas in tentatiónem. \R\ Sed líbera nos a malo.

\rubric{Credo in Deum secreto usque ad} \V\ Carnis resurrectiónem. \R\ Vitam ætérnam. Amen.

\V\ Et ego ad te, Dómine, clamávi. \R\ Et mane orátio mea prævéniet te. \V\ Repleátur os meum laude. \R\ Ut cantem glóriam tuam, tota die magnitúdinem tuam. \V\ Dómine, avérte fáciem tuam a peccátis meis. \R\ Et omnes iniquitátes meas dele. \V\ Cor mundum crea in me, Deus. \R\ Et spíritum rectum ínnova in viscéribus meis. \V\ Ne projícias me a fácie tua. \R\ Et spíritum sanctum tuum ne áuferas a me. \V\ Redde mihi lætítiam salutáris tui. \R\ Et spíritu principáli confírma me. \V\ Adjutórium nostrum in nómine Dómini. \R\ Qui fecit cælum et terram.

\rubric{Deinde Hebdomadarius facit Confessionem cum Absolutione, eamque Chorus repetit, ut infra post feriales Preces.}

\rubric{In omnibus vero Vigiliis, si ad Laudes recitatæ fuerint preces, post Responsorium breve dicuntur flexis genibus sequentes Preces feriales:}

KYRIE, eléison. Christe, eléison. Kýrie, eléison.

\rubric{Pater noster secreto usque ad} \V\ Et ne nos indúcas in tentatiónem. \R\ Sed líbera nos a malo.

\rubric{Credo in Deum secreto usque ad} \V\ Carnis resurrectiónem. \R\ Vitam ætérnam. Amen.

\V\ Et ego ad te, Dómine, clamávi. \R\ Et mane orátio mea prævéniet te. \V\ Repleátur os meum laude. \R\ Ut cantem glóriam tuam, tota die magnitúdinem tuam. \V\ Dómine, avérte fáciem tuam a peccátis meis. \R\ Et omnes iniquitátes meas dele. \V\ Cor mundum crea in me, Deus. \R\ Et spíritum rectum ínnova in viscéribus meis. \V\ Ne projícias me a fácie tua. \R\ Et spíritum sanctum tuum ne áuferas a me. \V\ Redde mihi lætítiam salutáris tui. \R\ Et spíritu principáli confírma me. \V\ Eripe me, Dómine, ab hómine malo. \R\ A viro iníquo éripe me. \V\ Eripe me de inimícis meis, Deus meus. \R\ Et ab insurgéntibus in me líbera me. \V\ Eripe me de operántibus iniquitátem. \R\ Et de viris sánguinum salva me. \V\ Sic psalmum dicam nómini tuo in sæculum sæculi. \R\ Ut reddam vota mea de die in diem. \V\ Exáudi nos, Deus, salutáris noster. \R\ Spes ómnium fínium terræ, et in mari longe. \V\ Deus, in adjutórium meum inténde. \R\ Dómine, ad adjuvándum me festína. \V\ Sanctus Deus, Sanctus fortis, Sanctus immortális. \R\ Miserére nobis. \V\ Bénedic, ánima mea, Dómino. \R\ Et ómnia, quæ intra me sunt, nómini sancto ejus. \V\ Bénedic, ánima mea, Dómino. \R\ Et noli oblivísci omnes retributiónes ejus. \V\ Qui propitiátur ómnibus iniquitátibus tuis. \R\ Qui sanat omnes infirmitátes tuas. \V\ Qui rédimit de intéritu vitam tuam. \R\ Qui corónat te in misericórdia et miseratiónibus. \V\ Qui replet in bonis desidérium tuum. \R\ Renovábitur ut áquilæ juvéntus tua. \V\ Adjutórium nostrum in nómine Dómini. \R\ Qui fecit cælum et terram.

\rubric{Deinde Hebdomadarius facit Confessionem:}

\lettrine[lines=2]{C}{ONFITEOR} Deo omnipoténti, beátæ Maríæ semper Vírgini, beáto Michaéli Archángelo, beáto Joánni Baptístæ, sanctis Apóstolis Petro et Paulo, ómnibus Sanctis, et vobis, fratres, quia peccávi nimis cogitatióne, verbo et ópere: mea culpa, mea culpa, mea máxima culpa. Ideo precor beátam Maríam semper Vírginem, beátum Michaélem Archángelum, beátum Joánnem Baptístam, sanctos Apóstolos Petrum et Paulum, omnes Sanctos, et vos, fratres, oráre pro me ad Dóminum, Deum nostrum.

\rubric{Chorus respondet:}

\lettrine[lines=2]{M}{ISEREATUR} tui omnípotens Deus, et, dimíssis peccátis tuis, perdúcat te ad vitam ætérnam. \R\ Amen.

\rubric{Deinde repetit Confessionem, et ubi dicitur: vobis, fratres, et vos, fratres, dicat: tibi, pater, et te, pater.}

\rubric{Facta Confessione a Choro, Hebdomadarius dicit:}

\lettrine[lines=2]{M}{ISEREATUR} vestri omnípotens Deus, et, dimíssis peccátis vestris, perdúcat vos ad vitam ætérnam. \R\ Amen.

\lettrine[lines=2]{I}{NDULGENTIAM}, absolutiónem, et remissiónem peccatórum nostrórum tríbuat nobis omnípotens et miséricors Dóminus. \R\ Amen.

\rubric{Extra Chorum, si unus vel duo tantum recitent Officium et in Choro Monialium, semel tantum ac simul ab omnibus ita fit Confessio:}

\lettrine[lines=2]{C}{ONFITEOR} Deo omnipoténti, beátæ Maríæ semper Vírgini, beáto Michaéli Archángelo, beáto Joánni Baptístæ, sanctis Apóstolis Petro et Paulo, et ómnibus Sanctis, quia peccávi nimis cogitatióne, verbo et ópere: mea culpa, mea culpa, mea máxima culpa. Ideo precor beátam Maríam semper Vírginem, beátum Michaélem Archángelum, beátum Joánnem Baptístam, sanctos Apóstolos Petrum et Paulum, et omnes Sanctos, oráre pro me ad Dóminum, Deum nostrum.

\rubric{Deinde dicitur:}

\lettrine[lines=2]{M}{ISEREATUR} nostri omnípotens Deus, et, dimíssis peccátis nostris, perdúcat nos ad vitam ætérnam. Amen.

\lettrine[lines=2]{I}{NDULGENTIAM}, absolutiónem, et remissiónem peccatórum nostrórum tríbuat nobis omnípotens et miséricors Dóminus. Amen.

\rubric{Facta Absolutione, subjungitur:}

\V\ Dignáre, Dómine, die isto. \R\ Sine peccáto nos custodíre. \V\ Miserére nostri, Dómine. \R\ Miserére nostri. \V\ Fiat misericórdia tua, Dómine, super nos. \R\ Quemádmodum sperávimus in te. \V\ Dómine, exáudi oratiónem meam. \R\ Et clamor meus ad te véniat.

\rubric{Absolutis Precibus, vel, si Preces locum non habeant, post Responsorium breve statim dicitur:}

\V\ Dóminus vobíscum. \R\ Et cum spíritu tuo.

Orémus. \oratiolabel{Oratio}

\lettrine[lines=2]{D}{OMINE}, Deus omnípotens, qui ad princípium hujus diéi nos perveníre fecísti: tua nos hódie salva virtúte; ut in hac die ad nullum declinémus peccátum, sed semper ad tuam justítiam faciéndam nostra procédant elóquia, dirigántur cogitatiónes et ópera. Per Dóminum. \R\ Amen.

\V\ Dóminus vobíscum. \R\ Et cum spíritu tuo. \V\ Benedicámus Dómino. \R\ Deo grátias.

\rubric{Deinde in Choro legitur Martyrologium, quod laudabiliter fit etiam extra Chorum. Postea Hebdomadarius dicit:}

\V\ Pretiósa in conspéctu Dómini. \R\ Mors Sanctórum ejus.

\rubric{Et absolute, sine} Orémus:

\oratiolabel{Oratio}

\lettrine[lines=2]{S}{ANCTA} María, et omnes Sancti intercédant pro nobis ad Dóminum, ut nos mereámur ab eo adjuvári et salvári, qui vivit et regnat in sæcula sæculórum. \R\ Amen.

\V\ Deus, in adjutórium meum inténde. \R\ Dómine, ad adjuvándum me festína. \rubric{Et dicitur ter; ultimo additur:} Glória Patri, et Fílio, et Spirítui Sancto. Sicut erat in princípio, et nunc, et semper, et in sæcula sæculórum. Amen.

Kýrie, eléison. Christe, eléison. Kýrie, eléison.

Pater noster \rubric{secreto usque ad} \V\ Et ne nos indúcas in tentatiónem. \R\ Sed líbera nos a malo.

\V\ Réspice in servos tuos, Dómine, et in ópera tua, et dírige fílios eórum. \R\ Et sit splendor Dómini, Dei nostri, super nos, et ópera mánuum nostrárum dírige super nos, et opus mánuum nostrárum dírige. \V\ Glória Patri, et Fílio, et Spirítui Sancto. \R\ Sicut erat in princípio, et nunc, et semper, et in sæcula sæculórum. Amen.

Orémus. \oratiolabel{Oratio}

\lettrine[lines=2]{D}{IRIGERE} et sanctificáre, régere et gubernáre dignáre, Dómine Deus, Rex cæli et terræ, hódie corda et córpora nostra, sensus, sermónes et actus nostros in lege tua, et in opéribus mandatórum tuórum: ut hic, et in ætérnum, te auxiliánte, salvi et líberi esse mereámur, Salvátor mundi: Qui vivis et regnas in sæcula sæculórum. \R\ Amen.

\rubric{Deinde dicitur Lectio brevis, ad absolutionem Capituli, prouti Officium occurrens requirit.}

\V\ Jube, domne, benedícere. \benedlabel{Bened.} Dies et actus nostros in sua pace dispónat Dóminus omnípotens. \R\ Amen.

\variantheading{In Officio dominicali et feriali}

\rubric{A Feria II usque ad Feriam IV post Festum sanctissimæ Trinitatis, et a Sabbato post Octavam Sacratissimi Cordis Jesu et deinceps.}

\lectioheading{Lectio brevis}{2 Thess. 3, 5}

\lettrine[lines=2]{D}{OMINUS} autem dírigat corda et córpora nostra in caritáte Dei, et patiéntia Christi. Tu autem, Dómine, miserére nobis. \R\ Deo grátias.

\rubric{¶ Ita semper concluduntur Lectiones breves ad Primam, et ita in earum fine respondetur.}

\rubric{In Officio autem cujuslibet Festi vel Octavæ, atque S. Mariæ in Sabbato; itemque in Dominicis infra Octavas Ssmi Corporis Christi ac Sacrat. Cordis Jesu, si de eis fiat Officium:}

\rubric{Pro Lectione brevi sumitur Capitulum, quod ad Nonam ponitur ipsius Officii currentis, et in fine dicitur: Tu autem, Dómine, ut supra.}

\rubric{Post Lectionem brevem dicitur:}

\V\ Adjutórium nostrum in nómine Dómini. \R\ Qui fecit cælum et terram.

\V\ Benedícite. \R\ Deus.

\benedlabel{Bened.} Dóminus nos benedícat, et ab omni malo deféndat, et ad vitam perdúcat ætérnam. Et fidélium ánimæ per misericórdiam Dei requiéscant in pace. \R\ Amen.

\rubric{Deinde dicitur tantum Pater noster secreto, nisi sequatur alia Hora.}

\hourheading{Ad Tertiam}

Pater noster et Ave.

\V\ Deus, in adjutórium.

\hymnusheading{Hymnus}

\lettrine[lines=2]{N}{UNC}, Sancte, nobis, Spíritus,\\
Unum Patri cum Fílio,\\
Dignáre promptus íngeri\\
Nostro refúsus péctori.

Os, lingua, mens, sensus, vigor,\\
Confessiónem pérsonent,\\
Flamméscat igne cáritas,\\
Accéndat ardor próximos.

Præsta, Pater piíssime,\\
Patríque compar Unice,\\
Cum Spíritu Paráclito\\
Regnans per omne sæculum. Amen.

\rubric{Expleto Hymno, conveniens dicitur Antiphona usque ad Asteriscum *, et tres Psalmi, prouti Officium occurrens requirit; atque Psalmis absolutis, integra repetitur Antiphona.}

\variantheading{In Officio dominicali}

\rubric{In omnibus Dominicis, extra Octavas Ssmi Corporis Christi ac Sacrat. Cordis Jesu, Aña ut in Psalterio, et dicitur sequens}

\capitulumlabel{Capitulum}{1 Joann. 4, 16}

\lettrine[lines=2]{D}{EUS} cáritas est: et qui manet in caritáte, in Deo manet, et Deus in eo.

\R\ br. Inclína cor meum, Deus, * In testimónia tua. Inclína. \V\ Avérte óculos meos, ne vídeant vanitátem: in via tua vivífica me. In testimónia tua. Glória Patri. Inclína.

\V\ Ego dixi: Dómine, miserére mei. \R\ Sana ánimam meam, quia peccávi tibi.

\variantheading{In Officio feriali}

\rubric{In omnibus Feriis, Antiphona de Feria currenti, ut in Psalterio; et dicitur quotidie sequens}

\capitulumlabel{Capitulum}{Jer. 17, 14}

\lettrine[lines=2]{S}{ANA} me, Dómine, et sanábor: salvum me fac, et salvus ero: quóniam laus mea tu es.

\R\ br. Sana ánimam meam, * Quia peccávi tibi. Sana. \V\ Ego dixi: Dómine, miserére mei. Quia. Glória Patri. Sana.

\V\ Adjútor meus esto, ne derelínquas me. \R\ Neque despícias me, Deus, salutáris meus.

\rubric{Expleto Responsorio brevi, si Preces non fuerint recitandæ, statim dicitur Versus Dóminus vobíscum cum Oratione et reliquis omnibus usque ad finem Horæ, ut infra.}

\rubric{Quando vero ad Laudes recitatæ sunt Preces, post Responsorium breve dicuntur sequentes flexis genibus:}

KYRIE, eléison. Christe, eléison. Kýrie, eléison.

Pater noster \rubric{secreto usque ad} \V\ Et ne nos indúcas in tentatiónem. \R\ Sed líbera nos a malo. \V\ Dómine, Deus virtútum, convérte nos. \R\ Et osténde fáciem tuam, et salvi érimus. \V\ Exsúrge, Christe, ádjuva nos. \R\ Et líbera nos propter nomen tuum. \V\ Dómine, exáudi oratiónem meam. \R\ Et clamor meus ad te véniat.

\rubric{Absolutis Precibus, vel, si Preces locum non habeant, post Responsorium breve statim dicitur:}

\V\ Dóminus vobíscum. \R\ Et cum spíritu tuo.

Orémus.

\rubric{Et dicitur Oratio conveniens.}

\rubric{Post Orationem additur:}

\V\ Dóminus vobíscum. \R\ Et cum spíritu tuo.

\V\ Benedicámus Dómino. \R\ Deo grátias.

\V\ Fidélium ánimæ per misericórdiam Dei requiéscant in pace. \R\ Amen.

\rubric{Deinde dicitur tantum Pater noster secreto, nisi sequatur alia Hora.}

\hourheading{Ad Sextam}

Pater noster et Ave.

\V\ Deus, in adjutórium.

\hymnusheading{Hymnus}

\lettrine[lines=2]{R}{ECTOR} potens, verax Deus,\\
Qui témperas rerum vices,\\
Splendóre mane illúminas,\\
Et ígnibus merídiem:

Exstíngue flammas lítium,\\
Aufer calórem nóxium,\\
Confer salútem córporum,\\
Verámque pacem córdium.

Præsta, Pater piíssime,\\
Patríque compar Unice,\\
Cum Spíritu Paráclito\\
Regnans per omne sæculum. Amen.

\rubric{Expleto Hymno, conveniens dicitur Antiphona usque ad Asteriscum *, et tres Psalmi prouti Officium occurrens requirit; atque, Psalmis absolutis, integra repetitur Antiphona.}

\variantheading{In Officio dominicali}

\rubric{In omnibus Dominicis, extra Octavas Ssmi Corporis Christi ac Sacrat. Cordis Jesu, Aña ut in Psalterio, et dicitur sequens}

\capitulumlabel{Capitulum}{Gal. 6, 2}

\lettrine[lines=2]{A}{LTER} altérius ónera portáte, et sic adimplébitis legem Christi.

\R\ br. In ætérnum, Dómine, * Pérmanet verbum tuum. In ætérnum. \V\ In sæculum sæculi véritas tua. Pérmanet. Glória Patri. In ætérnum.

\V\ Dóminus regit me, et nihil mihi déerit. \R\ In loco páscuæ ibi me collocávit.

\variantheading{In Officio feriali}

\rubric{In omnibus Feriis, Antiphona de Feria currenti, ut in Psalterio; et dicitur quotidie sequens}

\capitulumlabel{Capitulum}{Rom. 13, 8}

\lettrine[lines=2]{N}{EMINI} quidquam debeátis, nisi ut ínvicem diligátis: qui enim díligit próximum, legem implévit.

\R\ br. Benedícam Dóminum * In omni témpore. Benedícam. \V\ Semper laus ejus in ore meo. In. Glória Patri. Benedícam.

\V\ Dóminus regit me, et nihil mihi déerit. \R\ In loco páscuæ ibi me collocávit.

\rubric{Expleto Responsorio brevi, si Preces non fuerint recitandæ, statim dicitur Versus Dóminus vobíscum cum Oratione et reliquis omnibus usque ad finem Horæ, ut infra.}

\rubric{Quando vero ad Laudes recitatæ sunt Preces, post Responsorium breve dicuntur sequentes flexis genibus:}

KYRIE, eléison. Christe, eléison. Kýrie, eléison.

Pater noster \rubric{secreto usque ad} \V\ Et ne nos indúcas in tentatiónem. \R\ Sed líbera nos a malo. \V\ Dómine, Deus virtútum, convérte nos. \R\ Et osténde fáciem tuam, et salvi érimus. \V\ Exsúrge, Christe, ádjuva nos. \R\ Et líbera nos propter nomen tuum. \V\ Dómine, exáudi oratiónem meam. \R\ Et clamor meus ad te véniat.

\rubric{Absolutis Precibus, vel, si Preces locum non habeant, post Responsorium breve statim dicitur:}

\V\ Dóminus vobíscum. \R\ Et cum spíritu tuo.

Orémus.

\rubric{Et dicitur Oratio conveniens.}

\rubric{Post Orationem additur:}

\V\ Dóminus vobíscum. \R\ Et cum spíritu tuo.

\V\ Benedicámus Dómino. \R\ Deo grátias.

\V\ Fidélium ánimæ per misericórdiam Dei requiéscant in pace. \R\ Amen.

\rubric{Deinde dicitur tantum Pater noster secreto, nisi sequatur alia Hora.}

\hourheading{Ad Nonam}

Pater noster et Ave.

\V\ Deus, in adjutórium.

\hymnusheading{Hymnus}

\lettrine[lines=2]{R}{ERUM}, Deus, tenax vigor,\\
Immótus in te pérmanens,\\
Lucis diúrnæ témpora\\
Successíbus detérminans:

Largíre lumen véspere,\\
Quo vita nusquam décidat,\\
Sed præmium mortis sacræ\\
Perénnis instet glória.

Præsta, Pater piíssime,\\
Patríque compar Unice,\\
Cum Spíritu Paráclito\\
Regnans per omne sæculum. Amen.

\rubric{Expleto Hymno, conveniens dicitur Antiphona usque ad Asteriscum *, et tres Psalmi, prouti Officium occurrens requirit; atque, Psalmis absolutis, integra repetitur Antiphona.}

\variantheading{In Officio dominicali}

\rubric{In omnibus Dominicis, extra Octavas Ssmi Corporis Christi ac Sacrat. Cordis Jesu, Aña ut in Psalterio, et dicitur sequens}

\capitulumlabel{Capitulum}{1 Cor. 6, 20}

\lettrine[lines=2]{E}{MPTI} enim estis prétio magno. Glorificáte et portáte Deum in córpore vestro.

\R\ br. Clamávi in toto corde meo: * Exáudi me, Dómine. Clamávi. \V\ Justificatiónes tuas requíram. Exáudi. Glória Patri. Clamávi.

\V\ Ab occúltis meis munda me, Dómine. \R\ Et ab aliénis parce servo tuo.

\variantheading{In Officio feriali}

\rubric{In omnibus Feriis, Antiphona de Feria currenti, ut in Psalterio; et dicitur quotidie sequens}

\capitulumlabel{Capitulum}{1 Petr. 1, 17-19}

\lettrine[lines=2]{I}{N} timóre incolátus vestri témpore conversámini: sciéntes quod non corruptibílibus auro vel argénto redémpti estis, sed pretióso sánguine quasi Agni immaculáti Christi.

\R\ br. Rédime me, Dómine, * Et miserére mei. Rédime. \V\ Pes enim meus stetit in via recta. Et. Glória Patri. Rédime.

\V\ Ab occúltis meis munda me, Dómine. \R\ Et ab aliénis parce servo tuo.

\rubric{Expleto Responsorio brevi, si Preces non fuerint recitandæ, statim dicitur Versus Dóminus vobíscum cum Oratione et reliquis omnibus usque ad finem Horæ, ut infra.}

\rubric{Quando vero ad Laudes recitatæ sunt Preces, post Responsorium breve dicuntur sequentes flexis genibus:}

KYRIE, eléison. Christe, eléison. Kýrie, eléison.

Pater noster \rubric{secreto usque ad} \V\ Et ne nos indúcas in tentatiónem. \R\ Sed líbera nos a malo. \V\ Dómine, Deus virtútum, convérte nos. \R\ Et osténde fáciem tuam, et salvi érimus. \V\ Exsúrge, Christe, ádjuva nos. \R\ Et líbera nos propter nomen tuum. \V\ Dómine, exáudi oratiónem meam. \R\ Et clamor meus ad te véniat.

\rubric{Absolutis Precibus, vel, si Preces locum non habeant, post Responsorium breve statim dicitur:}

\V\ Dóminus vobíscum. \R\ Et cum spíritu tuo.

Orémus.

\rubric{Et dicitur Oratio conveniens.}

\rubric{Post Orationem additur:}

\V\ Dóminus vobíscum. \R\ Et cum spíritu tuo.

\V\ Benedicámus Dómino. \R\ Deo grátias.

\V\ Fidélium ánimæ per misericórdiam Dei requiéscant in pace. \R\ Amen.

\rubric{Deinde dicitur tantum Pater noster secreto, nisi sequatur alia Hora.}

% ===== END OF BATCH 2 -- Covers autumnalis.md lines 465-874 (reused
% verbatim from aestiva.tex's Batch 2, per the byte-identity noted
% above), completing "Ad Primam" (hymn "Jam lucis orto sidere," the
% full Confiteor/Misereatur/Indulgentiam apparatus, the Martyrologium
% cue, the "Dirigere et sanctificare" Oratio, and the Lectio brevis,
% here given in its single simplified form -- "In Officio dominicali
% et feriali" -- rather than the three-way Lenten/Paschal-season
% branching found in verna.tex's Prime, since this volume's Lectio
% brevis never varies by season), "Ad Tertiam," "Ad Sextam," and "Ad
% Nonam," each with its proper hymn and its Sunday/weekday
% Capitulum-and-brief-responsory pair. Two spellings of the same word
% recur within this single batch exactly as the source alternates
% them: "sæculum"/"sæculi" at Sexta's responsory (line 758) versus
% the different spelling pattern seen in the Te Deum of Batch 1 --
% each transcribed as its own occurrence gives it, per the standing
% rule against silent normalization. The document continues with
% "Ad Vesperas" (autumnalis.md line 875) =====
\V\ Dóminus vobíscum. \R\ Et cum spíritu tuo.

Orémus.

\rubric{Et dicitur Oratio conveniens.}

\rubric{Postea fiunt Commemorationes, si occurrant.}

\rubric{Deinde, in Officio semiduplici aut simplici tam de Tempore quam de Sanctis, præter omnes dies, in quibus habita sit Commemoratio Duplicis, etiam concurrentis, vel factum fuerit Officium aut Commemoratio cujusvis Octavæ, fit sequens}

\suffragiumheading{Suffragium de omnibus Sanctis}

\antlabel{Ant.} Beáta Dei Génitrix Virgo María, Sanctíque omnes intercédant pro nobis ad Dóminum.

\V\ Mirificávit Dóminus Sanctos suos. \R\ Et exaudívit eos clamántes ad se.

Orémus. \oratiolabel{Oratio}

\lettrine[lines=2]{A}{} \textsc{cunctis} nos, quæsumus, Dómine, mentis et córporis defénde perículis: et, intercedénte beáta et gloriósa semper Vírgine Dei Genitríce María, cum beáto Joseph, beátis Apóstolis tuis Petro et Paulo, atque beáto N. et ómnibus Sanctis, salútem nobis tríbue benígnus et pacem; ut, destrúctis adversitátibus et erróribus univérsis, Ecclésia tua secúra tibi sérviat libertáte. Per eúmdem Dóminum. \R\ Amen.

\rubric{¶ In hac Oratione, ad litteram N., exprimitur nomen Titularis propriæ Ecclesiæ, dummodo Titulus non sit Persona divina vel Mysterium Domini, aut de ipso ratione Festi factum non fuerit Officium vel Commemoratio, aut illius nomen non sit in eadem Suffragii Oratione expressum; ac nomina sanctorum Angelorum et sancti Joannis Baptistæ, si Titulares fuerint, præponuntur nomini sancti Joseph. In hisce autem omnibus casibus omittuntur verba atque beáto N.}

\rubric{Si vero Officium vel Commemoratio fuerit de beáta María Vírgine, tunc dicitur ut infra; atque insuper, sicubi Titularis juxta superiorem Rubricam ante omnes fuerit nominandus, ad Orationem particula cum omittitur ante nomen sanctorum Petri et Pauli ac præponitur nomini sancti Joseph.}

\antlabel{Ant.} Sancti omnes intercédant pro nobis ad Dóminum.

\V\ Mirificávit Dóminus Sanctos suos. \R\ Et exaudívit eos clamántes ad se.

Orémus. \oratiolabel{Oratio}

\lettrine[lines=2]{A}{} \textsc{cunctis} nos, quæsumus, Dómine, mentis et córporis defénde perículis: et, intercedénte beáto Joseph, cum beátis Apóstolis tuis Petro et Paulo, atque beáto N. et ómnibus Sanctis, salútem nobis tríbue benígnus et pacem; ut, destrúctis adversitátibus et erróribus univérsis, Ecclésia tua secúra tibi sérviat libertáte. Per Dóminum.

\rubric{Post ultimam Orationem additur:}

\V\ Dóminus vobíscum. \R\ Et cum spíritu tuo.

\V\ Benedicámus Dómino. \R\ Deo grátias.

\V\ Fidélium ánimæ per misericórdiam Dei requiéscant in pace. \R\ Amen.

\rubric{Si post Vesperas immediate sequatur Completorium, dicto Versu Fidelium animæ, statim incipitur Versus Jube, domne, benedicere, ut infra ad Completorium; secus autem, si tunc terminetur Officium, dicitur tantum Pater noster secreto.}

\hourheading{Ad Completorium}

\rubric{Ante Completorium non dicitur Pater noster et Ave María, sicut in aliis Horis, sed absolute Lector incipit:}

\V\ Jube, domne, benedícere. \benedlabel{Bened.} Noctem quiétam, et finem perféctum concédat nobis Dóminus omnípotens. \R\ Amen.

\lectioheading{Lectio brevis}{1 Petr. 5, 8-9}

\lettrine[lines=2]{F}{RATRES}: Sóbrii estóte, et vigiláte: quia adversárius vester diábolus tamquam leo rúgiens círcuit, quærens quem dévoret: cui resístite fortes in fide. Tu autem, Dómine, miserére nobis. \R\ Deo grátias.

\V\ Adjutórium nostrum in nómine Dómini. \R\ Qui fecit cælum et terram.

Pater noster \rubric{quod dicitur totum secreto.}

\rubric{Deinde Hebdomadarius facit Confessionem:}

\lettrine[lines=2]{C}{ONFITEOR} Deo omnipoténti, beátæ Maríæ semper Vírgini, beáto Michaéli Archángelo, beáto Joánni Baptístæ, sanctis Apóstolis Petro et Paulo, ómnibus Sanctis, et vobis, fratres, quia peccávi nimis cogitatióne, verbo et ópere: mea culpa, mea culpa, mea máxima culpa. Ideo precor beátam Maríam semper Vírginem, beátum Michaélem Archángelum, beátum Joánnem Baptístam, sanctos Apóstolos Petrum et Paulum, omnes Sanctos, et vos, fratres, oráre pro me ad Dóminum, Deum nostrum.

\rubric{Chorus respondet:}

\lettrine[lines=2]{M}{ISEREATUR} tui omnípotens Deus, et, dimíssis peccátis tuis, perdúcat te ad vitam ætérnam. \R\ Amen.

\rubric{Deinde repetit Confessionem, et ubi dicitur: vobis, fratres, et vos, fratres, dicat: tibi, pater, et te, pater.}

\rubric{Facta Confessione a Choro, Hebdomadarius dicit:}

\lettrine[lines=2]{M}{ISEREATUR} vestri omnípotens Deus, et, dimíssis peccátis vestris, perdúcat vos ad vitam ætérnam. \R\ Amen.

\lettrine[lines=2]{I}{NDULGENTIAM}, absolutiónem, et remissiónem peccatórum nostrórum tríbuat nobis omnípotens et miséricors Dóminus. \R\ Amen.

\rubric{Extra Chorum, si unus vel duo tantum recitent Officium, et in Choro Monialium, semel tantum ac simul ab omnibus ita fit Confessio:}

\lettrine[lines=2]{C}{ONFITEOR} Deo omnipoténti, beátæ Maríæ semper Vírgini, beáto Michaéli Archángelo, beáto Joánni Baptístæ, sanctis Apóstolis Petro et Paulo, et ómnibus Sanctis, quia peccávi nimis cogitatióne, verbo et ópere: mea culpa, mea culpa, mea máxima culpa. Ideo precor beátam Maríam semper Vírginem, beátum Michaélem Archángelum, beátum Joánnem Baptístam, sanctos Apóstolos Petrum et Paulum, et omnes Sanctos, oráre pro me ad Dóminum, Deum nostrum.

\rubric{Deinde dicitur:}

\lettrine[lines=2]{M}{ISEREATUR} nostri omnípotens Deus, et, dimíssis peccátis nostris, perdúcat nos ad vitam ætérnam. Amen.

\lettrine[lines=2]{I}{NDULGENTIAM}, absolutiónem, et remissiónem peccatórum nostrórum tríbuat nobis omnípotens et miséricors Dóminus. Amen.

\rubric{Facta Absolutione, dicitur:}

\V\ Convérte nos, Deus, salutáris noster. \R\ Et avérte iram tuam a nobis. \V\ Deus, in adjutórium.

\rubric{Deinde, inchoata convenienti Antiphona, dicuntur tres Psalmi, prouti Officium occurrens requirit.}

\rubric{Repetita integre post ultimum Psalmum Antiphona, dicitur sequens}

\hymnusheading{Hymnus}

\lettrine[lines=2]{T}{E} lucis ante términum,\\
Rerum Creátor, póscimus,\\
Ut pro tua cleméntia\\
Sis præsul et custódia.

Procul recédant sómnia,\\
Et nóctium phantásmata;\\
Hostémque nostrum cómprime,\\
Ne polluántur córpora.

Præsta, Pater piíssime,\\
Patríque compar Unice,\\
Cum Spíritu Paráclito\\
Regnans per omne sæculum. Amen.

\rubric{Absoluto Hymno, dicitur sequens Capitulum et Responsorium breve:}

\capitulumlabel{Capitulum}{Jer. 14, 9}

\lettrine[lines=2]{T}{U} autem in nobis es, Dómine, et nomen sanctum tuum invocátum est super nos: ne derelínquas nos, Dómine, Deus noster.

\R\ br. In manus tuas, Dómine, * Comméndo spíritum meum. In manus tuas, Dómine, comméndo spíritum meum. \V\ Redemísti nos, Dómine, Deus veritátis. Comméndo spíritum meum. Glória Patri, et Fílio, et Spirítui Sancto. In manus tuas, Dómine, comméndo spíritum meum.

\V\ Custódi nos, Dómine, ut pupíllam óculi. \R\ Sub umbra alárum tuárum prótege nos.

\rubric{Postea ad Canticum dicitur sequens}

\antlabel{Ant.} Salva nos.

\canticumheading{Canticum Simeonis}{Luc. 2, 29-32}

\lettrine[lines=2]{N}{UNC} dimíttis servum tuum, Dómine, * secúndum verbum tuum in pace:

Quia vidérunt óculi mei * salutáre tuum.

Quod parásti * ante fáciem ómnium populórum,

Lumen ad revelatiónem géntium, * et glóriam plebis tuæ Israël.

Glória Patri.

\antlabel{Ant.} Salva nos, Dómine, vigilántes, custódi nos dormiéntes; ut vigilémus cum Christo, et requiescámus in pace.

\rubric{Deinde, si Preces non fuerint recitandæ, statim dicitur Versus Dóminus vobíscum cum subsequenti Oratione et reliquis omnibus usque ad finem Completorii, ut infra.}

\rubric{In quolibet Officio semiduplici, in Officio sanctæ Mariæ in Sabbato, in Festis simplicibus, atque in Feriis per Annum communibus, post repetitam ad Nunc dimittis Antiphonam dicuntur sequentes Preces, quæ tamen omittuntur omnibus diebus, in quibus ad Vesperas habita sit Commemoratio Duplicis, vel factum fuerit Officium aut Commemoratio cujusvis Octavæ.}

KYRIE, eléison. Christe, eléison. Kýrie, eléison.

Pater noster \rubric{secreto usque ad} \V\ Et ne nos indúcas in tentatiónem. \R\ Sed líbera nos a malo.

Credo in Deum \rubric{secreto usque ad} \V\ Carnis resurrectiónem. \R\ Vitam ætérnam. Amen.

\V\ Benedíctus es, Dómine, Deus patrum nostrórum. \R\ Et laudábilis et gloriósus in sæcula. \V\ Benedicámus Patrem et Fílium cum Sancto Spíritu. \R\ Laudémus, et superexaltémus eum in sæcula. \V\ Benedíctus es, Dómine, in firmaménto cæli. \R\ Et laudábilis, et gloriósus, et superexaltátus in sæcula. \V\ Benedícat et custódiat nos omnípotens et miséricors Dóminus. \R\ Amen. \V\ Dignáre, Dómine, nocte ista. \R\ Sine peccáto nos custodíre. \V\ Miserére nostri, Dómine. \R\ Miserére nostri. \V\ Fiat misericórdia tua, Dómine, super nos. \R\ Quemádmodum sperávimus in te. \V\ Dómine, exáudi oratiónem meam. \R\ Et clamor meus ad te véniat.

\rubric{Absolutis Precibus, vel, si Preces locum non habeant, repetita Antiphona post Nunc dimittis, statim dicitur:}

\V\ Dóminus vobíscum. \R\ Et cum spíritu tuo.

Orémus. \oratiolabel{Oratio}

\lettrine[lines=2]{V}{ISITA}, quæsumus, Dómine, habitatiónem istam, et omnes insídias inimíci ab ea longe repélle: Angeli tui sancti hábitent in ea, qui nos in pace custódiant; et benedíctio tua sit super nos semper. Per Dóminum nostrum. \R\ Amen.

\rubric{Post Orationem additur:}

\V\ Dóminus vobíscum. \R\ Et cum spíritu tuo.

\V\ Benedicámus Dómino. \R\ Deo grátias.

\benedictiolabel{Benedictio} Benedícat et custódiat nos omnípotens et miséricors Dóminus, Pater, et Fílius, et Spíritus Sanctus. \R\ Amen.

\rubric{Et non dicitur Versus Fidelium animæ, sed immediate subjungitur finalis beatæ Mariæ Virginis Antiphona, quæ infra assignatur.}

\rubric{Postea subjungitur:}

\V\ Divínum auxílium máneat semper nobíscum. \R\ Amen.

\rubric{Deinde dicitur secreto Pater noster, Ave María et Credo.}

\rubric{Si post Completorium immediate sequatur Matutinum, tunc bis dicitur Pater noster, Ave María et Credo, pro fine videlicet Officii diei præcedentis, et pro initio Officii diei sequentis.}

\sectionheading{Post singulas Horas}

\rubric{In fine Laudum, vel, si post Laudes immediate subsequatur alia Hora, in fine ultimæ Horæ, sicuti et semper in fine Completorii, quamvis immediate sequatur Matutinum diei sequentis, tam in publica quam in privata divini Officii recitatione; atque in publica insuper recitatione post alias quascumque Horas, si tunc terminetur Officium et discedatur a Choro, dicitur sequens Antiphona finalis beatæ Mariæ Virginis.}

\rubric{¶ Post Laudes ac ceteras Horas, præter Completorium, dicto Versu Fidelium animæ et Oratione Dominica secreto recitata, subjungitur:}

\V\ Dóminus det nobis suam pacem. \R\ Et vitam ætérnam. Amen.

\rubric{Atque statim dicitur Antiphona finalis, ut infra.}

\sectionheading{Antiphona finalis beatæ Mariæ Virginis}

\rubric{A I Vesperis Festi sanctissimæ Trinitatis et deinceps:}

\antiphonaheading{Antiphona}

\lettrine[lines=2]{S}{ALVE}, Regína, mater misericórdiæ; vita, dulcédo et spes nostra, salve. Ad te clamámus éxsules fílii Hevæ. Ad te suspirámus geméntes et flentes in hac lacrimárum valle. Eja ergo, advocáta nostra, illos tuos misericórdes óculos ad nos convérte. Et Jesum, benedíctum fructum ventris tui, nobis post hoc exsílium osténde. O clemens, o pia, o dulcis Virgo María.

\V\ Ora pro nobis, sancta Dei Génitrix. \R\ Ut digni efficiámur promissiónibus Christi.

Orémus. \oratiolabel{Oratio}

\lettrine[lines=2]{O}{MNIPOTENS} sempitérne Deus, qui gloriósæ Vírginis Matris Maríæ corpus et ánimam, ut dignum Fílii tui habitáculum éffici mererétur, Spíritu Sancto cooperánte, præparásti: da, ut, cujus commemoratióne lætámur; ejus pia intercessióne, ab instántibus malis, et a morte perpétua liberémur. Per eúmdem Christum, Dóminum nostrum. \R\ Amen.

\rubric{Absoluta vero Antiphona finali cum suis Versu et Oratione, concluditur:}

\V\ Divínum auxílium máneat semper nobíscum. \R\ Amen.

\sectionheading{Post divinum Officium}

\rubric{Cum expletum fuerit Officium, laudabiliter dicitur sequens Oratio; pro qua Summus Pontifex Pius Papa XI sub die 1 Decembris 1933 indulgentiam trium annorum concessit. Præterea Clericis in sacris constitutis et Sacerdotibus eam devote post Officium recitantibus Summus Pontifex Leo X defectus et culpas in eo persolvendo ex humana fragilitate contractas, indulsit. Dicitur autem flexis semper genibus in privata etiam recitatione, præter quam ab iis, qui ob certam infirmitatem vel gravioris impedimenti causam nequeant genuflectere.}

\oratioheading{Oratio}

\lettrine[lines=2]{S}{ACROSANCTÆ} et indivíduæ Trinitáti, crucifíxi Dómini nostri Jesu Christi humanitáti, beatíssimæ et gloriosíssimæ sempérque Vírginis Maríæ fecúndæ integritáti, et ómnium Sanctórum universitáti sit sempitérna laus, honor, virtus et glória ab omni creatúra, nobísque remíssio ómnium peccatórum, per infiníta sæcula sæculórum. \R\ Amen.

\V\ Beáta víscera Maríæ Vírginis, quæ portavérunt ætérni Patris Fílium. \R\ Et beáta úbera, quæ lactavérunt Christum Dóminum.

\rubric{Et dicitur secreto} Pater noster et Ave María.

% [ornamental image ending Ordinarium Divini Officii]

% ===== END OF BATCH 3 -- Covers autumnalis.md lines 922-1140 (reused
% verbatim from aestiva.tex's Batch 3, per the byte-identity noted
% above), completing the Ordinarium Divini Officii in its entirety:
% the close of "Ad Vesperas" (the Suffragium de omnibus Sanctis in
% both its Marian and Josephine forms, the second occurrence ending
% without a printed "℟. Amen," transcribed literally), "Ad
% Completorium" (the Lectio brevis, the full Confiteor/Misereatur/
% Indulgentiam apparatus, the hymn "Te lucis ante terminum," the Nunc
% dimittis with its proper brief responsory "In manus tuas," and the
% closing "Benedicat et custodiat nos" under
% \benedictiolabel{Benedictio}, matching the established distinction
% from the abbreviated \benedlabel{Bened.} used earlier at the same
% Hour's opening blessing), "Post singulas Horas," and the single
% "Antiphona finalis beatæ Mariæ Virginis" -- "Salve, Regína" alone,
% assigned from First Vespers of Trinity Sunday onward with no other
% seasonal Marian antiphon given in this shared Ordinarium text --
% and finally "Post divinum Officium" with its closing Oratio
% "Sacrosanctæ et individuæ Trinitati." This completes the Ordinarium
% Divini Officii for the Pars Autumnalis, matching Pars Aestiva's
% exactly. The document continues with "Psalterium Breviarii Romani"
% (autumnalis.md line 1142) =====
\partheading{Psalterium Breviarii Romani}
\partsubtitle{Per omnes ac singulos hebdomadæ dies dispositum}

\feriaheading{Dominica}

\hourheading{Ad Matutinum}

Pater, Ave et Credo.

\V\ Dómine, lábia mea.

\V\ Deus, in adjutórium.

Glória Patri. Sicut erat. Allelúja.

\rubric{In Dominica IV et reliquis Dominicis post Pentecosten occurrentibus dicitur:}

\invitatlabel{Invitat.} Dóminum, qui fecit nos, * Veníte, adorémus. \psalmlabel{Ps. 94} Veníte.

\hymnusheading{Hymnus}

\lettrine[lines=2]{N}{OCTE} surgéntes vigilémus omnes,\\
Semper in psalmis meditémur, atque\\
Voce concórdi Dómino canámus\\
Dúlciter hymnos.

Ut pio Regi páriter canéntes,\\
Cum suis Sanctis mereámur aulam\\
Ingredi cæli, simul et perénnem\\
Dúcere vitam.

Præstet hoc nobis Déitas beáta\\
Patris, ac Nati, paritérque Sancti\\
Spíritus, cujus résonat per omnem\\
Glória mundum. Amen.

\nocturnoheading{In I Nocturno}

\antlabel{Ant.} Beátus vir. \textdagger

\rubric{¶ Et non repetitur in Psalmo; quod semper fit, quando Antiphona incipit a primis Psalmi vel Cantici verbis. Et Psalmus vel Canticum tunc incipitur ab eo verbo, ante quod vel inchoata vel integra desinit Antiphona; si tamen eadem sint verba, et continuetur Antiphona cum Psalmo vel Cantico, neque additur in fine ipsius Antiphonæ Allelúja.}

\rubric{Antiphonæ autem cum Psalmis et Versibus assignatæ, adhibentur semper, extra Octavas Ssmi Corporis Christi et Sacrat. Cordis Jesu, quando fit Officium de Dominica.}

\rubric{¶ Asteriscus *, in Psalmorum versibus et in aliis Officii partibus insertus, denotat pausam cantus et recitationis in Choro; præter quam in Antiphonis, in quibus nempe significat tum præintonationem Antiphonæ, tum Antiphonam ipsam ante Psalmum seu Canticum eousque producendam esse ad omnes et singulas Horas in Officio ritus semiduplicis ac simplicis, atque ad Horas tantum minores in Officio ritus duplicis.}

\psalmheading{Psalmus 1}

\lettrine[lines=2]{B}{EATUS} vir, \textdagger\ qui non ábiit in consílio impiórum, et in via peccatórum non stetit, * et in cáthedra pestiléntiæ non sedit:

Sed in lege Dómini volúntas ejus, * et in lege ejus meditábitur die ac nocte.

Et erit tamquam lignum, quod plantátum est secus decúrsus aquárum, * quod fructum suum dabit in témpore suo:

Et fólium ejus non défluet: * et ómnia quæcúmque fáciet, prosperabúntur.

Non sic ímpii, non sic: * sed tamquam pulvis, quem prójicit ventus a fácie terræ.

Ideo non resúrgent ímpii in judício: * neque peccatóres in concílio justórum.

Quóniam novit Dóminus viam justórum: * et iter impiórum períbit.

Glória Patri, et Fílio, * et Spirítui Sancto.

Sicut erat in princípio, et nunc, et semper, * et in sæcula sæculórum. Amen.

\rubric{¶ Hic versus Glória Patri semper dicitur in fine omnium Psalmorum et Canticorum, nisi aliter notetur.}

\antlabel{Ant.} Beátus vir, qui in lege Dómini meditátur.

\antlabel{Ant.} Servíte Dómino.

\psalmheading{Psalmus 2}

\lettrine[lines=2]{Q}{UARE} fremuérunt gentes, * et pópuli meditáti sunt inánia?

Astitérunt reges terræ, et príncipes convenérunt in unum * advérsus Dóminum, et advérsus Christum ejus.

Dirumpámus víncula eórum: * et projiciámus a nobis jugum ipsórum.

Qui hábitat in cælis, irridébit eos: * et Dóminus subsannábit eos.

Tunc loquétur ad eos in ira sua, * et in furóre suo conturbábit eos.

Ego autem constitútus sum Rex ab eo super Sion montem sanctum ejus, * prædicans præcéptum ejus.

Dóminus dixit ad me: * Fílius meus es tu, ego hódie génui te.

Póstula a me, et dabo tibi gentes hereditátem tuam, * et possessiónem tuam términos terræ.

Reges eos in virga férrea, * et tamquam vas fíguli confrínges eos.

Et nunc, reges, intellígite: * erudímini, qui judicátis terram.

Servíte Dómino in timóre: * et exsultáte ei cum tremóre.

Apprehéndite disciplínam, nequándo irascátur Dóminus, * et pereátis de via justa.

Cum exárserit in brevi ira ejus, * beáti omnes qui confídunt in eo.

\antlabel{Ant.} Servíte Dómino in timóre: et exsultáte ei cum tremóre.

\antlabel{Ant.} Exsúrge.

\psalmheading{Psalmus 3}

\lettrine[lines=2]{D}{OMINE}, quid multiplicáti sunt qui tríbulant me? * multi insúrgunt advérsum me.

Multi dicunt ánimæ meæ: * Non est salus ipsi in Deo ejus.

Tu autem, Dómine, suscéptor meus es, * glória mea, et exáltans caput meum.

Voce mea ad Dóminum clamávi: * et exaudívit me de monte sancto suo.

Ego dormívi, et soporátus sum: * et exsurréxi, quia Dóminus suscépit me.

Non timébo míllia pópuli circumdántis me: * exsúrge, Dómine, salvum me fac, Deus meus.

Quóniam tu percussísti omnes adversántes mihi sine causa: * dentes peccatórum contrivísti.

Dómini est salus: * et super pópulum tuum benedíctio tua.

\antlabel{Ant.} Exsúrge, Dómine, salvum me fac, Deus meus.

\V\ Memor fui nocte nóminis tui, Dómine. \R\ Et custodívi legem tuam.

\absolutiolabel{Absol.} Exáudi, Dómine.

\benedlabel{Bened.} 1 Benedictióne perpétua. 2 Unigénitus Dei Fílius. 3 Spíritus Sancti grátia.

\nocturnoheading{In II Nocturno}

\antlabel{Ant.} Quam admirábile.

\psalmheading{Psalmus 8}

\lettrine[lines=2]{D}{OMINE}, Dóminus noster, * quam admirábile est nomen tuum in univérsa terra!

Quóniam eleváta est magnificéntia tua, * super cælos.

Ex ore infántium et lacténtium perfecísti laudem propter inimícos tuos, * ut déstruas inimícum et ultórem.

Quóniam vidébo cælos tuos, ópera digitórum tuórum: * lunam et stellas, quæ tu fundásti.

Quid est homo, quod memor es ejus? * aut fílius hóminis, quóniam vísitas eum?

Minuísti eum paulo minus ab Angelis, glória et honóre coronásti eum: * et constituísti eum super ópera mánuum tuárum.

Omnia subjecísti sub pédibus ejus, * oves et boves univérsas: ínsuper et pécora campi.

Vólucres cæli, et pisces maris, * qui perámbulant sémitas maris.

Dómine, Dóminus noster, * quam admirábile est nomen tuum in univérsa terra!

\antlabel{Ant.} Quam admirábile est nomen tuum, Dómine, in univérsa terra!

\antlabel{Ant.} Sedísti super thronum.

\psalmheading{Psalmus 9, i}

\lettrine[lines=2]{C}{ONFITEBOR} tibi, Dómine, in toto corde meo: * narrábo ómnia mirabília tua.

Lætábor et exsultábo in te: * psallam nómini tuo, Altíssime,

In converténdo inimícum meum retrórsum: * infirmabúntur, et períbunt a fácie tua.

Quóniam fecísti judícium meum et causam meam: * sedísti super thronum, qui júdicas justítiam.

Increpásti gentes, et périit ímpius: * nomen eórum delésti in ætérnum, et in sæculum sæculi.

Inimíci defecérunt frámeæ in finem: * et civitátes eórum destruxísti.

Périit memória eórum cum sónitu: * et Dóminus in ætérnum pérmanet.

Parávit in judício thronum suum: * et ipse judicábit orbem terræ in æquitáte, judicábit pópulos in justítia.

Et factus est Dóminus refúgium páuperi: * adjútor in opportunitátibus, in tribulatióne.

Et sperent in te qui novérunt nomen tuum: * quóniam non dereliquísti quæréntes te, Dómine.

\antlabel{Ant.} Sedísti super thronum, qui júdicas justítiam.

\antlabel{Ant.} Exsúrge, Dómine.

\psalmheading{Psalmus 9, ii}

\lettrine[lines=2]{P}{SALLITE} Dómino, qui hábitat in Sion: * annuntiáte inter gentes stúdia ejus:

Quóniam requírens sánguinem eórum recordátus est: * non est oblítus clamórem páuperum.

Miserére mei, Dómine: * vide humilitátem meam de inimícis meis.

Qui exáltas me de portis mortis, * ut annúntiem omnes laudatiónes tuas in portis fíliæ Sion.

Exsultábo in salutári tuo: * infíxæ sunt gentes in intéritu, quem fecérunt.

In láqueo isto, quem abscondérunt, * comprehénsus est pes eórum.

Cognoscétur Dóminus judícia fáciens: * in opéribus mánuum suárum comprehénsus est peccátor.

Convertántur peccatóres in inférnum, * omnes gentes quæ obliviscúntur Deum.

Quóniam non in finem oblívio erit páuperis: * patiéntia páuperum non períbit in finem.

Exsúrge, Dómine, non confortétur homo: * judicéntur gentes in conspéctu tuo.

Constítue, Dómine, legislatórem super eos: * ut sciant gentes quóniam hómines sunt.

\antlabel{Ant.} Exsúrge, Dómine, non prævaleat homo.

\V\ Média nocte surgébam ad confiténdum tibi. \R\ Super judícia justificatiónis tuæ.

\absolutiolabel{Absol.} Ipsíus píetas et.

\benedlabel{Bened.} 4 Deus Pater omnípotens. 5 Christus perpétuæ. 6 Ignem sui amóris.

\nocturnoheading{In III Nocturno}

\antlabel{Ant.} Ut quid, Dómine. \textdagger

\psalmheading{Psalmus 9, iii}

\lettrine[lines=2]{U}{T} quid, Dómine, \textdagger\ recessísti longe, * déspicis in opportunitátibus, in tribulatióne?

Dum supérbit ímpius, incénditur pauper: * comprehendúntur in consíliis quibus cógitant.

Quóniam laudátur peccátor in desidériis ánimæ suæ: * et iníquus benedícitur.

Exacerbávit Dóminum peccátor, * secúndum multitúdinem iræ suæ non quæret.

Non est Deus in conspéctu ejus: * inquinátæ sunt viæ illíus in omni témpore.

Auferúntur judícia tua a fácie ejus: * ómnium inimicórum suórum dominábitur.

Dixit enim in corde suo: * Non movébor a generatióne in generatiónem sine malo.

Cujus maledictióne os plenum est, et amaritúdine, et dolo: * sub lingua ejus labor et dolor.

Sedet in insídiis cum divítibus in occúltis, * ut interfíciat innocéntem.

Oculi ejus in páuperem respíciunt: * insidiátur in abscóndito, quasi leo in spelúnca sua.

Insidiátur ut rápiat páuperem: * rápere páuperem, dum áttrahit eum.

In láqueo suo humiliábit eum: * inclinábit se, et cadet, cum dominátus fúerit páuperum.

Dixit enim in corde suo: * Oblítus est Deus, avértit fáciem suam ne vídeat in finem.

\antlabel{Ant.} Ut quid, Dómine, recessísti longe?

\antlabel{Ant.} Exsúrge. \textdagger

\psalmheading{Psalmus 9, iv}

\lettrine[lines=2]{E}{XSURGE}, \textdagger\ Dómine Deus, exaltétur manus tua: * ne obliviscáris páuperum.

Propter quid irritávit ímpius Deum? * dixit enim in corde suo: Non requíret.

Vides quóniam tu labórem et dolórem consíderas: * ut tradas eos in manus tuas.

Tibi derelíctus est pauper: * órphano tu eris adjútor.

Cóntere bráchium peccatóris et malígni: * quærétur peccátum illíus, et non inveniétur.

Dóminus regnábit in ætérnum, et in sæculum sæculi: * períbitis, gentes, de terra illíus.

Desidérium páuperum exaudívit Dóminus: * præparatiónem cordis eórum audívit auris tua.

Judicáre pupíllo et húmili, * ut non appónat ultra magnificáre se homo super terram.

\antlabel{Ant.} Exsúrge, Dómine Deus, exaltétur manus tua.

\antlabel{Ant.} Justus Dóminus.

\psalmheading{Psalmus 10}

\lettrine[lines=2]{I}{N} Dómino confído: quómodo dícitis ánimæ meæ: * Tránsmigra in montem sicut passer?

Quóniam ecce peccatóres intendérunt arcum, paravérunt sagíttas suas in pháretra, * ut sagíttent in obscúro rectos corde.

Quóniam quæ perfecísti, destruxérunt: * justus autem quid fecit?

Dóminus in templo sancto suo, * Dóminus in cælo sedes ejus:

Oculi ejus in páuperem respíciunt: * pálpebræ ejus intérrogant fílios hóminum.

Dóminus intérrogat justum et ímpium: * qui autem díligit iniquitátem, odit ánimam suam.

Pluet super peccatóres láqueos: * ignis, et sulphur, et spíritus procellárum pars cálicis eórum.

Quóniam justus Dóminus, et justítias diléxit: * æquitátem vidit vultus ejus.

\antlabel{Ant.} Justus Dóminus, et justítiam diléxit.

\V\ Prævenérunt óculi mei ad te dilúculo. \R\ Ut meditárer elóquia tua, Dómine.

\absolutiolabel{Absolutio} A vínculis.

\benedlabel{Bened.} 7 Evangélica léctio. 8 Divínum auxílium. 9 Ad societátem cívium.

% ===== END OF BATCH 4 -- Covers autumnalis.md lines 1142-1442, opening
% "Psalterium Breviarii Romani" with "Dominica" / "Ad Matutinum" in
% its ordinary-time form ("In Dominica IV et reliquis Dominicis post
% Pentecosten occurrentibus"): the Invitatorium "Dóminum, qui fecit
% nos," the hymn "Nocte surgentes," and the three Nocturns (Psalms 1,
% 2, 3; 8, 9/i, 9/ii; 9/iii, 9/iv, 10), each with its Absolutio and
% triple Benedictiones. Neither this Invitatorium nor this hymn
% appears anywhere in hiemalis.tex or verna.tex, confirming this is
% genuinely new ordinary-time-after-Pentecost material specific to
% the summer and autumn volumes. A structurally significant
% observation: after Psalm 1 prints "Glória Patri" in full together
% with the rubric establishing that this doxology is "semper dicitur
% in fine omnium Psalmorum et Canticorum, nisi aliter notetur," EVERY
% SUBSEQUENT PSALM IN THIS NOCTURNAL SET (2, 3, 8, 9/i, 9/ii, 9/iii,
% 9/iv, 10) OMITS THE PRINTED "Glória Patri" LINE ENTIRELY -- the
% source trusts the reader to supply it from the standing rubric
% rather than reprinting it each time. This transcription follows
% the source exactly: no "Glória Patri" line is added after any of
% these psalms, since none appears in autumnalis.md itself; only where
% the source itself prints it (as after Psalm 1) does the LaTeX
% include it. The embedded dagger marks (e.g. "Beátus vir, †",
% "Ut quid, Dómine, †", "Exsúrge, †") are rendered \textdagger\
% exactly where the source places them, including inside an
% antiphon incipit alone ("Ant. Ut quid, Dómine. †"). The second
% Nocturn's Absolutio is given as the abbreviated "Absol." but the
% third Nocturn's uses the unabbreviated "Absolutio" -- both render
% via the same flexible \absolutiolabel{} command, consistent with
% established practice. The document continues with "Ad Laudes"
% (autumnalis.md line 1443) =====
Sicut ádipe et pinguédine repleátur ánima mea: * et lábiis exsultatiónis laudábit os meum.

Si memor fui tui super stratum meum, in matutínis meditábor in te: * quia fuísti adjútor meus.

Et in velaménto alárum tuárum exsultábo, adhǽsit ánima mea post te: * me suscépit déxtera tua.

Ipsi vero in vanum quæsiérunt ánimam meam, introíbunt in inferióra terræ: * tradéntur in manus gládii, partes vúlpium erunt.

Rex vero lætábitur in Deo, laudabúntur omnes qui jurant in eo: * quia obstrúctum est os loquéntium iníqua.

\antlabel{Ant.} Benedícam te in vita mea, Dómine: et in nómine tuo levábo manus meas, allelúja.

\antlabel{Ant.} Tres púeri.

\canticumheading{Canticum trium Puerorum [4]}{Dan. 57-88 et 56}

\lettrine[lines=2]{B}{ENEDICITE}, ómnia ópera Dómini, Dómino: * laudáte et superexaltáte eum in sæcula.

Benedícite, Angeli Dómini, Dómino: * benedícite, cæli, Dómino.

Benedícite, aquæ omnes, quæ super cælos sunt, Dómino: * benedícite, omnes virtútes Dómini, Dómino.

Benedícite, sol et luna, Dómino: * benedícite, stellæ cæli, Dómino.

Benedícite, omnis imber et ros, Dómino: * benedícite, omnes spíritus Dei, Dómino.

Benedícite, ignis et æstus, Dómino: * benedícite, frigus et æstus, Dómino.

Benedícite, rores et pruína, Dómino: * benedícite, gelu et frigus, Dómino.

Benedícite, glácies et nives, Dómino: * benedícite, noctes et dies, Dómino.

Benedícite, lux et ténebræ, Dómino: * benedícite, fúlgura et nubes, Dómino.

Benedícat terra Dóminum: * laudet et superexáltet eum in sæcula.

Benedícite, montes et colles, Dómino: * benedícite, univérsa germinántia in terra, Dómino.

Benedícite, fontes, Dómino: * benedícite, mária et flúmina, Dómino.

Benedícite, cete, et ómnia, quæ movéntur in aquis, Dómino: * benedícite, omnes vólucres cæli, Dómino.

Benedícite, omnes béstiæ et pécora, Dómino: * benedícite, fílii hóminum, Dómino.

Benedícat Israël Dóminum: * laudet et superexáltet eum in sæcula.

Benedícite, sacerdótes Dómini, Dómino: * benedícite, servi Dómini, Dómino.

Benedícite, spíritus, et ánimæ justórum, Dómino: * benedícite, sancti, et húmiles corde, Dómino.

Benedícite, Ananía, Azaría, Mísaël, Dómino: * laudáte et superexaltáte eum in sæcula.

Benedicámus Patrem et Fílium cum Sancto Spíritu: * laudémus et superexaltémus eum in sæcula.

Benedíctus es, Dómine, in firmaménto cæli: * et laudábilis, et gloriósus, et superexaltátus in sæcula.

\rubric{¶ Hic non dicitur} Glória Patri, \rubric{neque} Amen.

\antlabel{Ant.} Tres púeri jussu regis in fornácem missi sunt, non timéntes flammam ignis, dicéntes: Benedíctus Deus, allelúja.

\antlabel{Ant.} Allelúja.

\psalmheading{Psalmus 148 [5]}

\lettrine[lines=2]{L}{AUDATE} Dóminum de cælis: * laudáte eum in excélsis.

Laudáte eum, omnes Angeli ejus: * laudáte eum, omnes virtútes ejus.

Laudáte eum, sol et luna: * laudáte eum, omnes stellæ et lumen.

Laudáte eum, cæli cælórum: * et aquæ omnes, quæ super cælos sunt, laudent nomen Dómini.

Quia ipse dixit, et facta sunt: * ipse mandávit, et creáta sunt.

Státuit ea in ætérnum, et in sæculum sæculi: * præcéptum pósuit, et non præteríbit.

Laudáte Dóminum de terra, * dracónes, et omnes abýssi.

Ignis, grando, nix, glácies, spíritus procellárum: * quæ fáciunt verbum ejus:

Montes, et omnes colles: * ligna fructífera, et omnes cedri.

Béstiæ, et univérsa pécora: * serpéntes, et vólucres pennátæ:

Reges terræ, et omnes pópuli: * príncipes, et omnes júdices terræ.

Júvenes, et vírgines: senes cum junióribus laudent nomen Dómini: * quia exaltátum est nomen ejus solíus.

Conféssio ejus super cælum et terram: * et exaltávit cornu pópuli sui.

Hymnus ómnibus sanctis ejus: * fíliis Israël, pópulo appropinquánti sibi.

\antlabel{Ant.} Allelúja, laudáte Dóminum de cælis, allelúja, allelúja.

\rubric{In omnibus Dominicis extra Octavas Ssmi Corporis Christi ac Sacratissimi Cordis Jesu, dicitur:}

\capitulumlabel{Capitulum}{Apoc. 7, 12}

\lettrine[lines=2]{B}{ENEDICTIO}, et cláritas, et sapiéntia, et gratiárum áctio, honor, virtus, et fortitúdo Deo nostro in sæcula sæculórum. Amen. \R\ Deo grátias.

\rubric{In Dominica IV et reliquis Dominicis post Pentecosten occurrentibus, dicitur}

\hymnusheading{Hymnus}

\lettrine[lines=2]{E}{CCE}, jam noctis tenuátur umbra,\\
Lux et auróræ rútilans corúscat:\\
Súpplices rerum Dóminum canóra\\
Voce precémur:

Ut reos culpæ miserátus omnem\\
Pellat angórem, tríbuat salútem,\\
Donet et nobis bona sempitérnæ\\
Múnera pacis.

Præstet hoc nobis Déitas beáta\\
Patris, ac Nati, paritérque Sancti\\
Spíritus, cujus résonat per omnem\\
Glória mundum. Amen.

\V\ Dóminus regnávit, decórem índuit. \R\ Induit Dóminus fortitúdinem, et præcínxit se virtúte.

\rubric{Ad Benedictus Antiphona ut in Proprio de Tempore.}

\canticumheading{Canticum Zachariæ}{Luc. 1, 68-79}

\lettrine[lines=2]{B}{ENEDICTUS} \textdagger\ Dóminus, Deus Israël, * quia visitávit, et fecit redemptiónem plebis suæ:

Et eréxit cornu salútis nobis: * in domo David, púeri sui.

Sicut locútus est per os sanctórum, * qui a sæculo sunt, prophetárum ejus:

Salútem ex inimícis nostris, * et de manu ómnium, qui odérunt nos:

Ad faciéndam misericórdiam cum pátribus nostris: * et memorári testaménti sui sancti.

Jusjurándum, quod jurávit ad Abraham, patrem nostrum, * datúrum se nobis:

Ut sine timóre, de manu inimicórum nostrórum liberáti, * serviámus illi.

In sanctitáte, et justítia coram ipso, * ómnibus diébus nostris.

Et tu, puer, Prophéta Altíssimi vocáberis: * præíbis enim ante fáciem Dómini paráre vias ejus:

Ad dandam sciéntiam salútis plebi ejus: * in remissiónem peccatórum eórum:

Per víscera misericórdiæ Dei nostri: * in quibus visitávit nos, óriens ex alto:

Illumináre his, qui in ténebris, et in umbra mortis sedent: * ad dirigéndos pedes nostros in viam pacis.

\rubric{Oratio conveniens. Reliqua ut in Ordinario. 11}

\rubric{¶ Psalmi cum Cantico de Dominica superius appositi dicuntur etiam, sub Antiphonis Officio diei congruentibus, in Dominicis infra Octavas Ssmi Corporis Christi ac Sacrat. Cordis Jesu, in Duplicibus I vel II classis, in Festis novem Lectionum Domini, B. Mariæ Virginis, Angelorum, S. Joannis Baptistæ, S. Joseph, Apostolorum, Evangelistarum et per Octavas privilegiatas; itemque in quolibet Festo novem Lectionum proprias Antiphonas ad Laudes habenti, quando de eis persolvitur Officium.}

\hourheading{Ad Primam}

Pater, Ave et Credo.

\V\ Deus, in adjutórium.

\hymnusheading{Hymnus}

\lettrine[lines=2]{J}{AM} lucis orto sídere,\\
Deum precémur súpplices,\\
Ut in diúrnis áctibus\\
Nos servet a nocéntibus.

Linguam refrénans témperet,\\
Ne litis horror ínsonet:\\
Visum fovéndo cóntegat,\\
Ne vanitátes háuriat.

Sint pura cordis íntima,\\
Absístat et vecórdia:\\
Carnis terat supérbiam\\
Potus cibíque párcitas.

Ut, cum dies abscésserit\\
Noctémque sors redúxerit,\\
Mundi per abstinéntiam\\
Ipsi canámus glóriam.

Deo Patri sit glória,\\
Ejúsque soli Fílio,\\
Cum Spíritu Paráclito,\\
Nunc et per omne sæculum. Amen.

\antlabel{Ant.} Allelúja.

\rubric{¶ Aña cum Psalmis, ad Horas minores assignata, adhibetur semper, extra Octavas Ssmi Corporis Christi ac Sacrat. Cordis Jesu, quando fit Officium de qualibet Dominica minori.}

\psalmheading{Psalmus 117}

\lettrine[lines=2]{C}{ONFITEMINI} Dómino quóniam bonus: * quóniam in sæculum misericórdia ejus.

Dicat nunc Israël quóniam bonus: * quóniam in sæculum misericórdia ejus.

Dicat nunc domus Aaron: * quóniam in sæculum misericórdia ejus.

Dicant nunc qui timent Dóminum: * quóniam in sæculum misericórdia ejus.

De tribulatióne invocávi Dóminum: * et exaudívit me in latitúdine Dóminus.

Dóminus mihi adjútor: * non timébo quid fáciat mihi homo.

Dóminus mihi adjútor: * et ego despíciam inimícos meos.

Bonum est confídere in Dómino, * quam confídere in hómine:

Bonum est speráre in Dómino, * quam speráre in princípibus.

Omnes gentes circuiérunt me: * et in nómine Dómini quia ultus sum in eos.

Circumdántes circumdedérunt me: * et in nómine Dómini quia ultus sum in eos.

Circumdedérunt me sicut apes, et exarsérunt sicut ignis in spinis: * et in nómine Dómini quia ultus sum in eos.

Impúlsus evérsus sum ut cáderem: * et Dóminus suscépit me.

Fortitúdo mea, et laus mea Dóminus: * et factus est mihi in salútem.

Vox exsultatiónis, et salútis * in tabernáculis justórum.

Déxtera Dómini fecit virtútem: déxtera Dómini exaltávit me, * déxtera Dómini fecit virtútem.

Non móriar, sed vivam: * et narrábo ópera Dómini.

Castígans castigávit me Dóminus: * et morti non trádidit me.

Aperíte mihi portas justítiæ, ingréssus in eas confitébor Dómino: * hæc porta Dómini, justi intrábunt in eam.

Confitébor tibi quóniam exaudísti me: * et factus es mihi in salútem.

Lápidem, quem reprobavérunt ædificántes: * hic factus est in caput ánguli.

A Dómino factum est istud: * et est mirábile in óculis nostris.

Hæc est dies, quam fecit Dóminus: * exsultémus, et lætémur in ea.

O Dómine, salvum me fac, o Dómine, bene prosperáre: * benedíctus qui venit in nómine Dómini.

Benedíximus vobis de domo Dómini: * Deus Dóminus, et illúxit nobis.

Constitúite diem solémnem in condénsis, * usque ad cornu altáris.

Deus meus es tu, et confitébor tibi: * Deus meus es tu, et exaltábo te.

Confitébor tibi quóniam exaudísti me: * et factus es mihi in salútem.

Confitémini Dómino quóniam bonus: * quóniam in sæculum misericórdia ejus.

\rubric{¶ In Dominicis infra Octavas Ssmi Corporis Christi ac Sacrat. Cordis Jesu, in Duplicibus I vel II classis, in Festis novem Lectionum Domini, B. Mariæ Virginis, Angelorum, S. Joannis Baptistæ, S. Joseph, Apostolorum, Evangelistarum, et per Octavas privilegiatas: si de eis fiat Officium, ad Primam, Tertiam, Sextam et Nonam dicuntur Psalmi de Dominica, sub Antiphonis Officio diei congruentibus, nisi aliter suo loco notetur; sed ad Primam, loco Psalmi 117 Confitemini, in ipsa etiam die Dominica semper dicitur}

\psalmheading{Psalmus 53}

\lettrine[lines=2]{D}{EUS}, in nómine tuo salvum me fac: * et in virtúte tua júdica me.

Deus, exáudi oratiónem meam: * áuribus pércipe verba oris mei.

Quóniam aliéni insurrexérunt advérsum me, et fortes quæsiérunt ánimam meam: * et non proposuérunt Deum ante conspéctum suum.

Ecce enim, Deus ádjuvat me: * et Dóminus suscéptor est ánimæ meæ.

Avérte mala inimícis meis: * et in veritáte tua dispérde illos.

Voluntárie sacrificábo tibi, * et confitébor nómini tuo, Dómine: quóniam bonum est:

Quóniam ex omni tribulatióne eripuísti me: * et super inimícos meos despéxit óculus meus.

\psalmheading{Psalmus 118, i}

\lettrine[lines=2]{B}{EATI} immaculáti in via: * qui ámbulant in lege Dómini.

Beáti, qui scrutántur testimónia ejus: * in toto corde exquírunt eum.

Non enim qui operántur iniquitátem: * in viis ejus ambulavérunt.

Tu mandásti * mandáta tua custodíri nimis.

Utinam dirigántur viæ meæ, * ad custodiéndas justificatiónes tuas!

Tunc non confúndar, * cum perspéxero in ómnibus mandátis tuis.

Confitébor tibi in directióne cordis: * in eo quod dídici judícia justítiæ tuæ.

Justificatiónes tuas custódiam: * non me derelínquas usquequáque.

In quo córrigit adolescéntior viam suam? * in custodiéndo sermónes tuos.

In toto corde meo exquisívi te: * ne repéllas me a mandátis tuis.

In corde meo abscóndi elóquia tua: * ut non peccem tibi.

Benedíctus es, Dómine: * doce me justificatiónes tuas.

In lábiis meis, * pronuntiávi ómnia judícia oris tui.

In via testimoniórum tuórum delectátus sum, * sicut in ómnibus divítiis.

In mandátis tuis exercébor: * et considerábo vias tuas.

In justificatiónibus tuis meditábor: * non oblivíscar sermónes tuos.

\psalmheading{Psalmus 118, ii}

\lettrine[lines=2]{R}{ETRIBUE} servo tuo, vivífica me: * et custódiam sermónes tuos.

Revéla óculos meos: * et considerábo mirabília de lege tua.

Íncola ego sum in terra: * non abscóndas a me mandáta tua.

Concupívit ánima mea desideráre justificatiónes tuas, * in omni témpore.

Increpásti supérbos: * maledícti qui declínant a mandátis tuis.

Aufer a me oppróbrium, et contémptum: * quia testimónia tua exquisívi.

Etenim sedérunt príncipes, et advérsum me loquebántur: * servus autem tuus exercebátur in justificatiónibus tuis.

Nam et testimónia tua meditátio mea est: * et consílium meum justificatiónes tuæ.

Adhǽsit paviménto ánima mea: * vivífica me secúndum verbum tuum.

Vias meas enuntiávi, et exaudísti me: * doce me justificatiónes tuas.

Viam justificatiónum tuárum ínstrue me: * et exercébor in mirabílibus tuis.

Dormitávit ánima mea præ tædio: * confírma me in verbis tuis.

Viam iniquitátis ámove a me: * et de lege tua miserére mei.

Viam veritátis elégi: * judícia tua non sum oblítus.

Adhǽsi testimóniis tuis, Dómine: * noli me confúndere.

Viam mandatórum tuórum cucúrri, * cum dilatásti cor meum.

\rubric{¶ Si Athanasianum Symbolum omittatur, expletis Psalmis præcedentibus, statim dicitur Antiphona, ut infra. 45}

\rubric{Sequens Symbolum dicitur in Dominicis post Pentecosten, quando fit Officium de Dominica ut in Psalterio, neque ulla Commemoratio Officii duplicis vel Octavæ recurrit, et in Festo sanctissimæ Trinitatis.}

\symbolumheading{Symbolum Athanasianum}

QUICUMQUE vult salvus esse, * ante ómnia opus est, ut téneat cathólicam fidem:

Quam nisi quisque íntegram inviolatámque serváverit, * absque dúbio in ætérnum períbit.

Fides autem cathólica hæc est: * ut unum Deum in Trinitáte, et Trinitátem in unitáte venerémur.

Neque confundéntes persónas, * neque substántiam separántes.

Ália est enim persóna Patris, ália Fílii, * ália Spíritus Sancti:

Sed Patris, et Fílii, et Spíritus Sancti una est divínitas, * æquális glória, coætérna majéstas.

Quális Pater, talis Fílius, * talis Spíritus Sanctus.

Increátus Pater, increátus Fílius, * increátus Spíritus Sanctus.

Imménsus Pater, imménsus Fílius, * imménsus Spíritus Sanctus.

Ætérnus Pater, ætérnus Fílius, * ætérnus Spíritus Sanctus.

Et tamen non tres ætérni, * sed unus ætérnus.

Sicut non tres increáti, nec tres imménsi, * sed unus increátus, et unus imménsus.

Simíliter omnípotens Pater, omnípotens Fílius, * omnípotens Spíritus Sanctus.

Et tamen non tres omnipoténtes, * sed unus omnípotens.

Ita Deus Pater, Deus Fílius, * Deus Spíritus Sanctus.

Et tamen non tres Dii, * sed unus est Deus.

Ita Dóminus Pater, Dóminus Fílius, * Dóminus Spíritus Sanctus.

Et tamen non tres Dómini, * sed unus est Dóminus.

Quia, sicut singillátim unamquámque persónam Deum ac Dóminum confitéri christiána veritáte compéllimur: * ita tres Deos aut Dóminos dícere cathólica religióne prohibémur.

Pater a nullo est factus: * nec creátus, nec génitus.

Fílius a Patre solo est: * non factus, nec creátus, sed génitus.

Spíritus Sanctus a Patre et Fílio: * non factus, nec creátus, nec génitus, sed procédens.

Unus ergo Pater, non tres Patres: unus Fílius, non tres Fílii: * unus Spíritus Sanctus, non tres Spíritus Sancti.

Et in hac Trinitáte nihil prius aut postérius, nihil majus aut minus: * sed totæ tres persónæ coætérnæ sibi sunt et coæquáles.

Ita ut per ómnia, sicut jam supra dictum est, * et únitas in Trinitáte, et Trínitas in unitáte veneránda sit.

Qui vult ergo salvus esse, * ita de Trinitáte séntiat.

Sed necessárium est ad ætérnam salútem, * ut Incarnatiónem quoque Dómini nostri Jesu Christi fidéliter credat.

Est ergo fides recta, ut credámus et confiteámur, * quia Dóminus noster Jesus Christus, Dei Fílius, Deus et homo est.

Deus est ex substántia Patris ante sæcula génitus: * et homo est ex substántia matris in sæculo natus.

Perféctus Deus, perféctus homo: * ex ánima rationáli et humána carne subsístens.

Æquális Patri secúndum divinitátem: * minor Patre secúndum humanitátem.

Qui, licet Deus sit et homo, * non duo tamen, sed unus est Christus.

Unus autem non conversióne divinitátis in carnem, * sed assumptióne humanitátis in Deum.

Unus omníno, non confusióne substántiæ, * sed unitáte persónæ.

Nam sicut ánima ratiónális et caro unus est homo: * ita Deus et homo unus est Christus.

Qui passus est pro salúte nostra: descéndit ad ínferos: * tértia die resurréxit a mórtuis.

Ascéndit ad cælos, sedet ad déxteram Dei Patris omnipoténtis: * inde ventúrus est judicáre vivos et mórtuos.

Ad cujus advéntum omnes hómines resúrgere habent cum corpóribus suis: * et redditúri sunt de factis própriis ratiónem.

Et qui bona egérunt, ibunt in vitam ætérnam: * qui vero mala, in ignem ætérnum.

Hæc est fides cathólica, * quam nisi quisque fidéliter firmitérque credíderit, salvus esse non póterit.

Glória Patri.

\rubric{Expleto Symbolo, vel, si hoc non habeat locum, Psalmis absolutis, statim dicitur Antiphona.}

\antlabel{Ant.} Allelúja, confitémini Dómino quóniam in sæculum misericórdia ejus, allelúja, allelúja.

\capitulumlabel{Capitulum}{1 Tim. 1, 17}

\lettrine[lines=2]{R}{EGI} sæculórum immortáli et invisíbili, soli Deo honor et glória in sæcula sæculórum. Amen.

\R\ br. Christe, Fili Dei vivi, * Miserére nobis. \rubric{Et repetitur:} Christe, Fili Dei vivi, miserére nobis. \V\ Qui sedes ad déxteram Patris. \R\ Miserére nobis. \V\ Glória Patri, et Fílio, et Spirítui Sancto. \R\ Christe, Fili Dei vivi, miserére nobis. \V\ Exsúrge, Christe, ádjuva nos. \R\ Et líbera nos propter nomen tuum.

\variantheading{Preces (si dicendæ sint)}

KYRIE, eléison. Christe, eléison. Kýrie, eléison.

Pater noster \rubric{secreto usque ad} \V\ Et ne nos indúcas in tentatiónem. \R\ Sed líbera nos a malo.

Credo in Deum \rubric{secreto usque ad} \V\ Carnis resurrectiónem. \R\ Vitam ætérnam. Amen.

\V\ Et ego ad te, Dómine, clamávi. \R\ Et mane orátio mea prævéniet te. \V\ Repleátur os meum laude. \R\ Ut cantem glóriam tuam, tota die magnitúdinem tuam. \V\ Dómine, avérte fáciem tuam a peccátis meis. \R\ Et omnes iniquitátes meas dele. \V\ Cor mundum crea in me, Deus. \R\ Et spíritum rectum ínnova in viscéribus meis. \V\ Ne projícias me a fácie tua. \R\ Et spíritum sanctum tuum ne áuferas a me. \V\ Redde mihi lætítiam salutáris tui. \R\ Et spíritu principáli confírma me. \V\ Adjutórium nostrum in nómine Dómini. \R\ Qui fecit cælum et terram.

\formulalabel{Confíteor. Misereátur. Indulgéntiam.}

\V\ Dignáre, Dómine, die isto. \R\ Sine peccáto nos custodíre. \V\ Miserére nostri, Dómine. \R\ Miserére nostri. \V\ Fiat misericórdia tua, Dómine, super nos. \R\ Quemádmodum sperávimus in te. \V\ Dómine, exáudi oratiónem meam. \R\ Et clamor meus ad te véniat.

\V\ Dóminus vobíscum.

Orémus. \oratiolabel{Oratio}

\lettrine[lines=2]{D}{OMINE}, Deus omnípotens, qui ad princípium hujus diéi nos perveníre fecísti: tua nos hódie salva virtúte; ut in hac die ad nullum declinémus peccátum, sed semper ad tuam justítiam faciéndam nostra procédant elóquia, dirigántur cogitatiónes et ópera. Per Dóminum. \R\ Amen.

\V\ Dóminus vobíscum. \V\ Benedicámus Dómino.

\rubric{In Choro Martyrologium.}

\V\ Pretiósa in conspéctu Dómini. \R\ Mors Sanctórum ejus.

\rubric{Et absolute, sine} Orémus:

\oratiolabel{Oratio}

\lettrine[lines=2]{S}{ANCTA} María, et omnes Sancti intercédant pro nobis ad Dóminum, ut nos mereámur ab eo adjuvári et salvári, qui vivit et regnat in sæcula sæculórum. \R\ Amen.

\V\ Deus, in adjutórium meum inténde. \R\ Dómine, ad adjuvándum me festína. \rubric{Et dicitur ter; ultimo additur:} Glória Patri. Sicut erat.

Kýrie, eléison. Christe, eléison. Kýrie, eléison.

Pater noster \rubric{secreto usque ad} \V\ Et ne nos indúcas in tentatiónem. \R\ Sed líbera nos a malo. \V\ Réspice in servos tuos, Dómine, et in ópera tua, et dírige fílios eórum. \R\ Et sit splendor Dómini, Dei nostri, super nos, et ópera mánuum nostrárum dírige super nos, et opus mánuum nostrárum dírige. \V\ Glória Patri. \R\ Sicut erat.

Orémus. \oratiolabel{Oratio}

\lettrine[lines=2]{D}{IRIGERE} et sanctificáre, régere et gubernáre dignáre, Dómine Deus, Rex cæli et terræ, hódie corda et córpora nostra, sensus, sermónes et actus nostros in lege tua, et in opéribus mandatórum tuórum: ut hic, et in ætérnum, te auxiliánte, salvi et líberi esse mereámur, Salvátor mundi: Qui vivis et regnas in sæcula sæculórum. \R\ Amen.

\V\ Jube, domne.

\benedlabel{Bened.} Dies et actus nostros in sua pace dispónat Dóminus omnípotens. \R\ Amen.

% ===== END OF BATCH 5 -- Covers autumnalis.md lines 1443-1984,
% completing "Ad Laudes" (Psalms 92, 99, 62, the Canticle of the
% Three Young Men "Benedicite," Psalm 148, the hymn "Ecce, jam noctis
% tenuatur umbra" for ordinary Sundays after Pentecost, and the
% Benedictus) and the body of "Ad Primam" (Psalm 117 "Confitemini,"
% the special Sunday substitute Psalm 53, Psalm 118 sections i-ii,
% the complete Athanasian Creed "Quicumque vult," the brief
% responsory, the Preces, the compact bold "Confíteor. Misereátur.
% Indulgéntiam." cross-reference rendered \formulalabel{} since the
% source bold-marks it here (matching the established rule that this
% command is used only when the source itself bolds the phrase,
% per precedent in verna.tex), the Martyrologium cue, and the closing
% Orationes). TWO SILENT CORRECTIONS of what appear to be markdown/
% OCR artifacts, per the standing rule for unambiguous mechanical
% errors: (1) a stray literal "\*" interrupting the hymn "Jam lucis
% orto sidere" at its final stanza ("\* Deo Patri sit gloria," verna.md
% line 1687) is dropped, since no hymn anywhere in this project uses
% chant-pause asterisks within its verse, and this same hymn appeared
% correctly without it in Batch 2; (2) a spurious paragraph break
% splitting the second stanza of "Ecce, jam noctis tenuatur umbra"
% into two two-line fragments (lines 1612-1616) is closed up into one
% continuous four-line stanza, matching the Sapphic-Adonic meter of
% every other stanza in the same hymn. The Athanasian Creed
% "Quicumque vult," though given in full capitals at its opening
% ("QUICUMQUE vult salvus esse"), receives no lettrine, matching the
% precedent already set in hiemalis.tex: the governing instruction's
% lettrine rule names oratio, psalm, canticle, hymn, and lectio
% specifically, and a Symbolum (Creed) is none of these, so
% capitalization alone does not trigger a lettrine here as it does
% for unheaded Oratio-like text elsewhere. Several versicle/response
% pairs in this batch print only the versicle with no response
% ("℣. Dóminus vobíscum." alone, "℣. Dóminus vobíscum. ℣.
% Benedicámus Dómino." with neither response given) -- transcribed
% exactly as the source abbreviates them, without supplying the
% missing responses. The document continues with the Lectio brevis
% variants for Prime (autumnalis.md line 1986), beginning with "In
% Officio dominicali" and continuing through several Proprio de
% Tempore feast-specific substitutions (Ssmæ Trinitatis, Ssmi
% Corporis Christi, etc.) =====
\variantheading{In Officio dominicali}

\lettrine[lines=2]{D}{OMINUS} autem dírigat corda et córpora nostra in caritáte Dei, et patiéntia Christi. Tu autem.

\variantheading{E Proprio de Tempore}

\variantheading{D. N. J. C. Regis}

\lettrine[lines=2]{I}{N} ipso complácuit omnem plenitúdinem inhabitáre, et per eum reconciliáre ómnia in ipsum, pacíficans per sánguinem crucis ejus sive quæ in terris, sive quæ in cælis sunt, in Christo Jesu Dómino nostro.

\variantheading{E Proprio Sanctorum}

\variantheading{S. Raphaëlis Archangeli}

\lettrine[lines=2]{T}{EMPUS} est, ut revértar ad eum, qui me misit: vos autem benedícite Deum, et narráte ómnia mirabília ejus.

\variantheading{Omnium Sanctorum}

\lettrine[lines=2]{B}{ENEDÍCTIO} et cláritas et sapiéntia et gratiárum áctio, honor, virtus et fortitúdo Deo nostro in sǽcula sæculórum. Amen.

\variantheading{S. Andreæ Apost.}

\lettrine[lines=2]{I}{SAÍAS} enim dicit: Dómine, quis crédidit audítui nostro? Ergo fides ex audítu, audítus autem per verbum Christi. Sed dico: Numquid non audiérunt? Et quidem in omnem terram exívit sonus eórum, et in fines orbis terræ verba eórum.

\variantheading{E Communi Sanctorum}

\variantheading{Pro Apostolis}

\lettrine[lines=2]{I}{BANT} Apóstoli gaudéntes a conspéctu concílii, quóniam digni hábiti sunt pro nómine Jesu contuméliam pati.

\variantheading{Pro Martyre}

\lettrine[lines=2]{J}{USTUS} cor suum trádidit ad vigilándum dilúculo ad Dóminum, qui fecit illum, et in conspéctu Altíssimi deprecábitur.

\variantheading{Pro pluribus Martyribus}

\lettrine[lines=2]{F}{ULGEBUNT} justi, et tamquam scintíllæ in arundinéto discúrrent. Judicábunt natiónes, et dominabúntur pópulis: et regnábit Dóminus illórum in perpétuum.

\variantheading{Pro Conf. Pont.}

\lettrine[lines=2]{F}{UNGI} sacerdótio, et habére laudem in nómine ipsíus, et offérre illi incénsum dignum in odórem suavitátis.

\variantheading{Pro Conf. non Pont.}

\lettrine[lines=2]{J}{USTUM} dedúxit Dóminus per vias rectas, et osténdit illi regnum Dei, et dedit illi sciéntiam sanctórum; honestávit illum in labóribus, et complévit labóres illíus.

\variantheading{Pro Virgine Martyre}

\lettrine[lines=2]{D}{OMINE}, Deus meus, exaltásti super terram habitatiónem meam, et pro morte defluénte deprecáta sum. Invocávi Dóminum, Patrem Dómini mei, ut non derelínquat me in die tribulatiónis meæ, et in témpore superbórum sine adjutório.

\variantheading{Pro Virgine non Martyre}

\lettrine[lines=2]{O}{} \textsc{quam} pulchra est casta generátio cum claritáte! immortális est enim memória illíus: quóniam apud Deum nota est, et apud hómines.

\variantheading{Pro Martyre tantum}

\lettrine[lines=2]{L}{AUDABIT} usque ad mortem ánima mea Dóminum, quóniam éruis sustinéntes te, et líberas eos de manu angústiæ, Dómine, Deus noster.

\variantheading{Pro nec Virg. nec Mart.}

\lettrine[lines=2]{M}{ULTAE} fíliæ congregavérunt divítias: tu supergréssa es univérsas. Fallax grátia, et vana est pulchritúdo: múlier timens Dóminum, ipsa laudábitur.

\variantheading{Pro Dedicatione Ecclesiæ}

\lettrine[lines=2]{E}{T} abstérget Deus omnem lácrimam ab óculis eórum: et mors ultra non erit, neque luctus, neque clamor, neque dolor erit ultra, quia prima abiérunt. Et dixit qui sedébat in throno: Ecce, nova fácio ómnia.

\variantheading{In Festis B. Mariæ V.}

\lettrine[lines=2]{I}{N} platéis sicut cinnamómum et bálsamum aromatízans odórem dedi: quasi myrrha elécta dedi suavitátem odóris.

\V\ Adjutórium nostrum in nómine Dómini. \R\ Qui fecit cælum et terram.

\V\ Benedícite. \R\ Deus.

\benedictiolabel{Benedictio} Dóminus nos benedícat, et ab omni malo deféndat, et ad vitam perdúcat ætérnam. Et fidélium ánimæ per misericórdiam Dei requiéscant in pace. \R\ Amen.

\rubric{Deinde dicitur tantum} Pater noster \rubric{secreto, nisi sequatur alia Hora.}

\hourheading{Ad Tertiam}

Pater et Ave.

\V\ Deus, in adjutórium.

\hymnusheading{Hymnus}

\lettrine[lines=2]{N}{UNC}, Sancte, nobis, Spíritus,\\
Unum Patri cum Fílio,\\
Dignáre promptus íngeri\\
Nostro refúsus péctori.

Os, lingua, mens, sensus, vigor,\\
Confessiónem pérsonent,\\
Flamméscat igne cáritas,\\
Accéndat ardor próximos.

Præsta, Pater piíssime,\\
Patríque compar Unice,\\
Cum Spíritu Paráclito\\
Regnans per omne sæculum. Amen.

\antlabel{Ant.} Allelúja.

\psalmheading{Psalmus 118, iii}

\lettrine[lines=2]{L}{EGEM} pone mihi, Dómine, viam justificatiónum tuárum: * et exquíram eam semper.

Da mihi intelléctum, et scrutábor legem tuam: * et custódiam illam in toto corde meo.

Deduc me in sémitam mandatórum tuórum: * quia ipsam vólui.

Inclína cor meum in testimónia tua: * et non in avarítiam.

Avérte óculos meos ne vídeant vanitátem: * in via tua vivífica me.

Státue servo tuo elóquium tuum, * in timóre tuo.

Ámputa oppróbrium meum quod suspicátus sum: * quia judícia tua jucúnda.

Ecce, concupívi mandáta tua: * in æquitáte tua vivífica me.

Et véniat super me misericórdia tua, Dómine: * salutáre tuum secúndum elóquium tuum.

Et respondébo exprobrántibus mihi verbum: * quia sperávi in sermónibus tuis.

Et ne áuferas de ore meo verbum veritátis usquequáque: * quia in judíciis tuis supersperávi.

Et custódiam legem tuam semper: * in sæculum et in sæculum sæculi.

Et ambulábam in latitúdine: * quia mandáta tua exquisívi.

Et loquébar in testimóniis tuis in conspéctu regum: * et non confundébar.

Et meditábar in mandátis tuis, * quæ diléxi.

Et levávi manus meas ad mandáta tua, quæ diléxi: * et exercébar in justificatiónibus tuis.

\psalmheading{Psalmus 118, iv}

\lettrine[lines=2]{M}{EMOR} esto verbi tui servo tuo, * in quo mihi spem dedísti.

Hæc me consoláta est in humilitáte mea: * quia elóquium tuum vivificábit me.

Supérbi iníque agébant usquequáque: * a lege autem tua non declinávi.

Memor fui judiciórum tuórum a sæculo, Dómine: * et consolátus sum.

Deféctio ténuit me, * pro peccatóribus derelinquéntibus legem tuam.

Cantábiles mihi erant justificatiónes tuæ, * in loco peregrinatiónis meæ.

Memor fui nocte nóminis tui, Dómine: * et custodívi legem tuam.

Hæc facta est mihi: * quia justificatiónes tuas exquisívi.

Pórtio mea, Dómine, * dixi custodíre legem tuam.

Deprecátus sum fáciem tuam in toto corde meo: * miserére mei secúndum elóquium tuum.

Cogitávi vias meas: * et convérti pedes meos in testimónia tua.

Parátus sum, et non sum turbátus: * ut custódiam mandáta tua.

Funes peccatórum circumpléxi sunt me: * et legem tuam non sum oblítus.

Média nocte surgébam ad confiténdum tibi, * super judícia justificatiónis tuæ.

Párticeps ego sum ómnium timéntium te: * et custodiéntium mandáta tua.

Misericórdia tua, Dómine, plena est terra: * justificatiónes tuas doce me.

\psalmheading{Psalmus 118, v}

\lettrine[lines=2]{B}{ONITATEM} fecísti cum servo tuo, Dómine, * secúndum verbum tuum.

Bonitátem, et disciplínam, et sciéntiam doce me: * quia mandátis tuis crédidi.

Priúsquam humiliárer ego delíqui: * proptérea elóquium tuum custodívi.

Bonus es tu: * et in bonitáte tua doce me justificatiónes tuas.

Multiplicáta est super me iníquitas superbórum: * ego autem in toto corde meo scrutábor mandáta tua.

Coagulátum est sicut lac cor eórum: * ego vero legem tuam meditátus sum.

Bonum mihi quia humiliásti me: * ut discam justificatiónes tuas.

Bonum mihi lex oris tui, * super míllia auri et argénti.

Manus tuæ fecérunt me, et plasmavérunt me: * da mihi intelléctum, et discam mandáta tua.

Qui timent te vidébunt me, et lætabúntur: * quia in verba tua supersperávi.

Cognóvi, Dómine, quia æquitas judícia tua: * et in veritáte tua humiliásti me.

Fiat misericórdia tua ut consolétur me, * secúndum elóquium tuum servo tuo.

Véniant mihi miseratiónes tuæ, et vivam: * quia lex tua meditátio mea est.

Confundántur supérbi, quia injúste iniquitátem fecérunt in me: * ego autem exercébor in mandátis tuis.

Convertántur mihi timéntes te: * et qui novérunt testimónia tua.

Fiat cor meum immaculátum in justificatiónibus tuis, * ut non confúndar.

\antlabel{Ant.} Allelúja, deduc me, Dómine, in sémitam mandatórum tuórum, allelúja, allelúja.

\capitulumlabel{Capitulum}{1 Joann. 4, 16}

\lettrine[lines=2]{D}{EUS} cáritas est: et qui manet in caritáte, in Deo manet, et Deus in eo.

\R\ br. Inclína cor meum, Deus, * In testimónia tua. Inclína. \V\ Avérte óculos meos, ne vídeant vanitátem: in via tua vivífica me. In testimónia tua. Glória Patri. Inclína. \V\ Ego dixi: Dómine, miserére mei. \R\ Sana ánimam meam, quia peccávi tibi.

\rubric{Oratio conveniens. Reliqua ut in Ordinario. 18}

% ===== END OF BATCH 6 -- Covers autumnalis.md lines 1986-2200,
% completing the Prime Lectio brevis variant list (substitute lessons
% across "In Officio dominicali," "E Proprio de Tempore" [here just
% Christ the King, the one movable feast of the Proper of Time this
% volume's Autumn scope reaches], "E Proprio Sanctorum" [St. Raphael
% the Archangel, All Saints, St. Andrew the Apostle -- autumnalis.md's
% own shorter three-feast list, verified via direct diff against
% aestiva.md's parallel eight-feast Summer list (Trinity Sunday
% through St. Lawrence) and substituted fresh from autumnalis.md's
% own text rather than carried over], and "E Communi Sanctorum" [the
% full set of twelve Common-of-Saints categories from Apostles
% through Marian feasts, confirmed byte-identical to aestiva.md's
% copy since the Common-of-Saints capitula do not vary by season])
% and the complete "Ad Tertiam" (hymn, Psalm 118 sections iii-v,
% Capitulum, and brief responsory). Each lesson, lacking its own
% repeated "Lectio brevis" label, is nonetheless lettrined on its
% full-caps opening word, consistent with its functional role as a
% Lectio substitute. "Pro Virgine non Martyre" opens "O QUAM
% pulchra..." with BOTH words in full caps (unlike the normal-case "O
% Mater"/"O Felícem" forms seen in Batch 166 of verna.tex), matching
% instead the all-caps "O GLORIOSA"/"O QUALE" pattern from that same
% project, so it receives the single-letter lettrine plus
% \textsc{Quam}. A second stray "\*" artifact, identical in kind to
% the one corrected in Batch 5, interrupts the final stanza of "Nunc,
% Sancte, nobis, Spiritus" at autumnalis.md line 2083 (within Ad
% Tertiam's hymn) and is silently dropped here for the same reason
% (no hymn anywhere in this project uses a chant-pause asterisk
% within its verse). The identical artifact recurs in Ad Sextam's
% hymn in the next batch and will be corrected the same way there.
% The document continues with "Ad Sextam" (autumnalis.md line 2202)
% =====
\hourheading{Ad Sextam}

Pater noster et Ave.

\V\ Deus, in adjutórium.

\hymnusheading{Hymnus}

\lettrine[lines=2]{R}{ECTOR} potens, verax Deus,\\
Qui témperas rerum vices,\\
Splendóre mane illúminas,\\
Et ígnibus merídiem:

Exstíngue flammas lítium,\\
Aufer calórem nóxium,\\
Confer salútem córporum,\\
Verámque pacem córdium.

Præsta, Pater piíssime,\\
Patríque compar Unice,\\
Cum Spíritu Paráclito\\
Regnans per omne sæculum. Amen.

\antlabel{Ant.} Allelúja.

\psalmheading{Psalmus 118, vi}

\lettrine[lines=2]{D}{EFECIT} in salutáre tuum ánima mea: * et in verbum tuum supersperávi.

Defecérunt óculi mei in elóquium tuum, * dicéntes: Quando consoláberis me?

Quia factus sum sicut uter in pruína: * justificatiónes tuas non sum oblítus.

Quot sunt dies servi tui? * quando fácies de persequéntibus me judícium?

Narravérunt mihi iníqui fabulatiónes: * sed non ut lex tua.

Omnia mandáta tua véritas: * iníque persecúti sunt me, ádjuva me.

Paulo minus consummavérunt me in terra: * ego autem non derelíqui mandáta tua.

Secúndum misericórdiam tuam vivífica me: * et custódiam testimónia oris tui.

In ætérnum, Dómine, * verbum tuum pérmanet in cælo.

In generatiónem et generatiónem véritas tua: * fundásti terram, et pérmanet.

Ordinatióne tua persevérat dies: * quóniam ómnia sérviunt tibi.

Nisi quod lex tua meditátio mea est: * tunc forte periíssem in humilitáte mea.

In ætérnum non oblivíscar justificatiónes tuas: * quia in ipsis vivificásti me.

Tuus sum ego, salvum me fac: * quóniam justificatiónes tuas exquisívi.

Me exspectavérunt peccatóres ut pérderent me: * testimónia tua intelléxi.

Omnis consummatiónis vidi finem: * latum mandátum tuum nimis.

\psalmheading{Psalmus 118, vii}

\lettrine[lines=2]{Q}{UOMODO} diléxi legem tuam, Dómine? * tota die meditátio mea est.

Super inimícos meos prudéntem me fecísti mandáto tuo: * quia in ætérnum mihi est.

Super omnes docéntes me intelléxi: * quia testimónia tua meditátio mea est.

Super senes intelléxi: * quia mandáta tua quæsívi.

Ab omni via mala prohíbui pedes meos: * ut custódiam verba tua.

A judíciis tuis non declinávi: * quia tu legem posuísti mihi.

Quam dúlcia fáucibus meis elóquia tua, * super mel ori meo!

A mandátis tuis intelléxi: * proptérea odívi omnem viam iniquitátis.

Lucérna pédibus meis verbum tuum, * et lumen sémitis meis.

Jurávi, et státui * custodíre judícia justítiæ tuæ.

Humiliátus sum usquequáque, Dómine: * vivífica me secúndum verbum tuum.

Voluntária oris mei beneplácita fac, Dómine: * et judícia tua doce me.

Anima mea in mánibus meis semper: * et legem tuam non sum oblítus.

Posuérunt peccatóres láqueum mihi: * et de mandátis tuis non errávi.

Hereditáte acquisívi testimónia tua in ætérnum: * quia exsultátio cordis mei sunt.

Inclinávi cor meum ad faciéndas justificatiónes tuas in ætérnum, * propter retributiónem.

\psalmheading{Psalmus 118, viii}

\lettrine[lines=2]{I}{NIQUOS} ódio hábui: * et legem tuam diléxi.

Adjútor et suscéptor meus es tu: * et in verbum tuum supersperávi.

Declináte a me, malígni: * et scrutábor mandáta Dei mei.

Súscipe me secúndum elóquium tuum, et vivam: * et non confúndas me ab exspectatióne mea.

Adjuva me, et salvus ero: * et meditábor in justificatiónibus tuis semper.

Sprevísti omnes discedéntes a judíciis tuis: * quia injústa cogitátio eórum.

Prævaricántes reputávi omnes peccatóres terræ: * ídeo diléxi testimónia tua.

Confíge timóre tuo carnes meas: * a judíciis enim tuis tímui.

Feci judícium et justítiam: * non tradas me calumniántibus me.

Súscipe servum tuum in bonum: * non calumniéntur me supérbi.

Óculi mei defecérunt in salutáre tuum: * et in elóquium justítiæ tuæ.

Fac cum servo tuo secúndum misericórdiam tuam: * et justificatiónes tuas doce me.

Servus tuus sum ego: * da mihi intelléctum, ut sciam testimónia tua.

Tempus faciéndi, Dómine: * dissipavérunt legem tuam.

Ideo diléxi mandáta tua, * super aurum et topázion.

Proptérea ad ómnia mandáta tua dirigébar: * omnem viam iníquam ódio hábui.

\antlabel{Ant.} Allelúja, tuus sum ego, salvum me fac, Dómine, allelúja, allelúja.

\capitulumlabel{Capitulum}{Gal. 6, 2}

\lettrine[lines=2]{A}{LTER} altérius ónera portáte, et sic adimplébitis legem Christi.

\R\ br. In ætérnum, Dómine, * Pérmanet verbum tuum. In ætérnum. \V\ In sæculum sæculi véritas tua. Pérmanet. Glória Patri. In ætérnum.

\V\ Dóminus regit me, et nihil mihi déerit. \R\ In loco páscuæ ibi me collocávit.

\rubric{Oratio conveniens. Reliqua ut in Ordinario. 20}

\hourheading{Ad Nonam}

Pater noster et Ave.

\V\ Deus, in adjutórium.

\hymnusheading{Hymnus}

\lettrine[lines=2]{R}{ERUM}, Deus, tenax vigor,\\
Immótus in te pérmanens,\\
Lucis diúrnæ témpora\\
Successíbus detérminans:

Largíre lumen véspere,\\
Quo vita nusquam décidat,\\
Sed præmium mortis sacræ\\
Perénnis instet glória.

Præsta, Pater piíssime,\\
Patríque compar Unice,\\
Cum Spíritu Paráclito\\
Regnans per omne sæculum. Amen.

\antlabel{Ant.} Allelúja.

\psalmheading{Psalmus 118, ix}

\lettrine[lines=2]{M}{IRABILIA} testimónia tua: * ídeo scrutáta est ea ánima mea.

Declarátio sermónum tuórum illúminat: * et intelléctum dat párvulis.

Os meum apérui, et attráxi spíritum: * quia mandáta tua desiderábam.

Aspice in me, et miserére mei, * secúndum judícium diligéntium nomen tuum.

Gressus meos dírige secúndum elóquium tuum: * et non dominétur mei omnis injustítia.

Rédime me a calúmniis hóminum: * ut custódiam mandáta tua.

Fáciem tuam illúmina super servum tuum: * et doce me justificatiónes tuas.

Exitus aquárum deduxérunt óculi mei: * quia non custodiérunt legem tuam.

Justus es, Dómine: * et rectum judícium tuum.

Mandásti justítiam testimónia tua: * et veritátem tuam nimis.

Tabéscere me fecit zelus meus: * quia oblíti sunt verba tua inimíci mei.

Ignítum elóquium tuum veheménter: * et servus tuus diléxit illud.

Adolescéntulus sum ego et contémptus: * justificatiónes tuas non sum oblítus.

Justítia tua, justítia in ætérnum: * et lex tua véritas.

Tribulátio, et angústia invenérunt me: * mandáta tua meditátio mea est.

Æquitas testimónia tua in ætérnum: * intelléctum da mihi, et vivam. Glória Patri.

\psalmheading{Psalmus 118, x}

\lettrine[lines=2]{C}{LAMAVI} in toto corde meo, exáudi me, Dómine: * justificatiónes tuas requíram.

Clamávi ad te, salvum me fac: * ut custódiam mandáta tua.

Prævéni in maturitáte, et clamávi: * quia in verba tua supersperávi.

Prævenérunt óculi mei ad te dilúculo: * ut meditárer elóquia tua.

Vocem meam audi secúndum misericórdiam tuam, Dómine: * et secúndum judícium tuum vivífica me.

Appropinquavérunt persequéntes me iniquitáti: * a lege autem tua longe facti sunt.

Prope es tu, Dómine: * et omnes viæ tuæ véritas.

Inítio cognóvi de testimóniis tuis: * quia in ætérnum fundásti ea.

Vide humilitátem meam, et éripe me: * quia legem tuam non sum oblítus.

Júdica judícium meum, et rédime me: * propter elóquium tuum vivífica me.

Longe a peccatóribus salus: * quia justificatiónes tuas non exquisiérunt.

Misericórdiæ tuæ multæ, Dómine: * secúndum judícium tuum vivífica me.

Multi qui persequúntur me, et tríbulant me: * a testimóniis tuis non declinávi.

Vidi prævaricántes, et tabescébam: * quia elóquia tua non custodiérunt.

Vide quóniam mandáta tua diléxi, Dómine: * in misericórdia tua vivífica me.

Princípium verbórum tuórum, véritas: * in ætérnum ómnia judícia justítiæ tuæ.

\psalmheading{Psalmus 118, xi}

\lettrine[lines=2]{P}{RINCIPES} persecúti sunt me gratis: * et a verbis tuis formidávit cor meum.

Lætábor ego super elóquia tua: * sicut qui invénit spólia multa.

Iniquitátem ódio hábui, et abominátus sum: * legem autem tuam diléxi.

Sépties in die laudem dixi tibi, * super judícia justítiæ tuæ.

Pax multa diligéntibus legem tuam: * et non est illis scándalum.

Exspectábam salutáre tuum, Dómine: * et mandáta tua diléxi.

Custodívit ánima mea testimónia tua: * et diléxit ea veheménter.

Servávi mandáta tua, et testimónia tua: * quia omnes viæ meæ in conspéctu tuo.

Appropínquet deprecátio mea in conspéctu tuo, Dómine: * juxta elóquium tuum da mihi intelléctum.

Intret postulátio mea in conspéctu tuo: * secúndum elóquium tuum éripe me.

Eructábunt lábia mea hymnum, * cum docúeris me justificatiónes tuas.

Pronuntiábit lingua mea elóquium tuum: * quia ómnia mandáta tua æquitas.

Fiat manus tua ut salvet me: * quóniam mandáta tua elégi.

Concupívi salutáre tuum, Dómine: * et lex tua meditátio mea est.

Vivet ánima mea, et laudábit te: * et judícia tua adjuvábunt me.

Errávi, sicut ovis, quæ périit: * quære servum tuum, quia mandáta tua non sum oblítus.

\antlabel{Ant.} Allelúja, fáciem tuam, Dómine, illúmina super servum tuum, allelúja, allelúja.

\capitulumlabel{Capitulum}{1 Cor. 6, 20}

\lettrine[lines=2]{E}{MPTI} enim estis prétio magno. Glorificáte et portáte Deum in córpore vestro.

\R\ br. Clamávi in toto corde meo: * Exáudi me, Dómine. Clamávi. \V\ Justificatiónes tuas requíram. Exáudi. Glória Patri. Clamávi. \V\ Ab occúltis meis munda me, Dómine. \R\ Et ab aliénis parce servo tuo.

\rubric{Oratio conveniens. Reliqua ut in Ordinario. 21}

% ===== END OF BATCH 7 -- Covers autumnalis.md lines 2202-2477,
% completing "Ad Sextam" (Psalm 118 sections vi-viii, Capitulum, and
% brief responsory) and "Ad Nonam" (Psalm 118 sections ix-xi,
% Capitulum, and brief responsory). A third occurrence of the stray
% "\*" artifact already corrected twice before (verna.md's hymn
% repetitions aside -- this is autumnalis.md's own recurring OCR/
% markdown glitch) interrupts Ad Sextam's hymn "Rector potens, verax
% Deus" at its final stanza (autumnalis.md line 2359) and is silently
% dropped, consistent with Batches 5 and 6. Psalm 118's ninth section
% (Mirabilia) is the first of this long psalm's many sections to
% print "Glória Patri." at its end (autumnalis.md line 2398) rather than
% omitting it per the usual pattern established since Batch 4 --
% transcribed exactly as the source gives it, neither adding nor
% removing a doxology relative to what each specific occurrence
% prints. The document continues with "Ad Vesperas" (autumnalis.md line
% 2515), opening the classic Sunday Vespers psalms (109 "Dixit
% Dominus," 110, etc.) =====
\hourheading{Ad Vesperas}

Pater noster et Ave.

\V\ Deus, in adjutórium.

\antlabel{Ant.} Dixit Dóminus. \textdagger

\rubric{¶ Antiphonæ cum Psalmis assignatæ adhibentur semper, quando fit Officium de Dominica, in Dominicis extra Octavas Ssmi Corporis Christi ac Sacratissimi Cordis Jesu.}

\psalmheading{Psalmus 109}

\lettrine[lines=2]{D}{IXIT} Dóminus \textdagger\ Dómino meo: * Sede a dextris meis:

Donec ponam inimícos tuos, * scabéllum pedum tuórum.

Virgam virtútis tuæ emíttet Dóminus ex Sion: * domináre in médio inimicórum tuórum.

Tecum princípium in die virtútis tuæ in splendóribus sanctórum: * ex útero ante lucíferum génui te.

Jurávit Dóminus, et non pænitébit eum: * Tu es sacérdos in ætérnum secúndum órdinem Melchísedech.

Dóminus a dextris tuis, * confrégit in die iræ suæ reges.

Judicábit in natiónibus, implébit ruínas: * conquassábit cápita in terra multórum.

De torrénte in via bibet: * proptérea exaltábit caput.

\antlabel{Ant.} Dixit Dóminus Dómino meo: Sede a dextris meis.

\antlabel{Ant.} Magna ópera Dómini.

\psalmheading{Psalmus 110}

\lettrine[lines=2]{C}{ONFITEBOR} tibi, Dómine, in toto corde meo: * in consílio justórum, et congregatióne.

Magna ópera Dómini: * exquisíta in omnes voluntátes ejus.

Conféssio et magnificéntia opus ejus: * et justítia ejus manet in sæculum sæculi.

Memóriam fecit mirabílium suórum, miséricors et miserátor Dóminus: * escam dedit timéntibus se.

Memor erit in sæculum testaménti sui: * virtútem óperum suórum annuntiábit pópulo suo:

Ut det illis hereditátem géntium: * ópera mánuum ejus véritas, et judícium.

Fidélia ómnia mandáta ejus: confirmáta in sæculum sæculi, * facta in veritáte et æquitáte.

Redemptiónem misit pópulo suo: * mandávit in ætérnum testaméntum suum.

Sanctum, et terríbile nomen ejus: * inítium sapiéntiæ timor Dómini.

Intelléctus bonus ómnibus faciéntibus eum: * laudátio ejus manet in sæculum sæculi.

\antlabel{Ant.} Magna ópera Dómini: exquisíta in omnes voluntátes ejus.

\antlabel{Ant.} Qui timet Dóminum.

\psalmheading{Psalmus 111}

\lettrine[lines=2]{B}{EATUS} vir, qui timet Dóminum: * in mandátis ejus volet nimis.

Potens in terra erit semen ejus: * generátio rectórum benedicétur.

Glória, et divítiæ in domo ejus: * et justítia ejus manet in sæculum sæculi.

Exórtum est in ténebris lumen rectis: * miséricors, et miserátor, et justus.

Jucúndus homo qui miserétur et cómmodat, dispónet sermónes suos in judício: * quia in ætérnum non commovébitur.

In memória ætérna erit justus: * ab auditióne mala non timébit.

Parátum cor ejus speráre in Dómino, confirmátum est cor ejus: * non commovébitur donec despíciat inimícos suos.

Dispérsit, dedit paupéribus: justítia ejus manet in sæculum sæculi, * cornu ejus exaltábitur in glória.

Peccátor vidébit, et irascétur, déntibus suis fremet et tabéscet: * desidérium peccatórum períbit.

\antlabel{Ant.} Qui timet Dóminum, in mandátis ejus cupit nimis.

\antlabel{Ant.} Sit nomen Dómini.

\psalmheading{Psalmus 112}

\lettrine[lines=2]{L}{AUDATE}, púeri, Dóminum: * laudáte nomen Dómini.

Sit nomen Dómini benedíctum, * ex hoc nunc, et usque in sæculum.

A solis ortu usque ad occásum, * laudábile nomen Dómini.

Excélsus super omnes gentes Dóminus, * et super cælos glória ejus.

Quis sicut Dóminus, Deus noster, qui in altis hábitat, * et humília réspicit in cælo et in terra?

Súscitans a terra ínopem, * et de stércore érigens páuperem:

Ut cóllocet eum cum princípibus, * cum princípibus pópuli sui.

Qui habitáre facit stérilem in domo, * matrem filiórum lætántem.

\antlabel{Ant.} Sit nomen Dómini benedíctum in sæcula.

\antlabel{Ant.} Deus autem noster.

\psalmheading{Psalmus 113}

\lettrine[lines=2]{I}{N} éxitu Israël de Ægýpto, * domus Jacob de pópulo bárbaro:

Facta est Judǽa sanctificátio ejus, * Israël potéstas ejus.

Mare vidit, et fugit: * Jordánis convérsus est retrórsum.

Montes exsultavérunt ut aríetes, * et colles sicut agni óvium.

Quid est tibi, mare, quod fugísti: * et tu, Jordánis, quia convérsus es retrórsum?

Montes, exsultástis sicut aríetes, * et colles, sicut agni óvium?

A fácie Dómini mota est terra, * a fácie Dei Jacob.

Qui convértit petram in stagna aquárum, * et rupem in fontes aquárum.

Non nobis, Dómine, non nobis: * sed nómini tuo da glóriam.

Super misericórdia tua, et veritáte tua: * nequándo dicant gentes: Ubi est Deus eórum?

Deus autem noster in cælo: * ómnia quæcúmque vóluit, fecit.

Simulácra géntium argéntum, et aurum, * ópera mánuum hóminum.

Os habent, et non loquéntur: * óculos habent, et non vidébunt.

Aures habent, et non áudient: * nares habent, et non odorábunt.

Manus habent, et non palpábunt: pedes habent, et non ambulábunt: * non clamábunt in gútture suo.

Símiles illis fiant qui fáciunt ea: * et omnes qui confídunt in eis.

Domus Israël sperávit in Dómino: * adjútor eórum et protéctor eórum est.

Domus Aaron sperávit in Dómino: * adjútor eórum et protéctor eórum est.

Qui timent Dóminum, speravérunt in Dómino: * adjútor eórum et protéctor eórum est.

Dóminus memor fuit nostri: * et benedíxit nobis:

Benedíxit dómui Israël: * benedíxit dómui Aaron.

Benedíxit ómnibus, qui timent Dóminum, * pusíllis cum majóribus.

Adjíciat Dóminus super vos: * super vos, et super fílios vestros.

Benedícti vos a Dómino, * qui fecit cælum, et terram.

Cælum cæli Dómino: * terram autem dedit fíliis hóminum.

Non mórtui laudábunt te, Dómine: * neque omnes, qui descéndunt in inférnum.

Sed nos qui vívimus, benedícimus Dómino, * ex hoc nunc et usque in sæculum.

\antlabel{Ant.} Deus autem noster in cælo: ómnia quæcúmque vóluit, fecit.

\capitulumlabel{Capitulum}{2 Cor. 1, 3-4}

\lettrine[lines=2]{B}{ENEDICTUS} Deus, et Pater Dómini nostri Jesu Christi, Pater misericordiárum, et Deus totíus consolatiónis, qui consolátur nos in omni tribulatióne nostra.

\hymnusheading{Hymnus}

\lettrine[lines=2]{L}{UCIS} Creátor óptime,\\
Lucem diérum próferens,\\
Primórdiis lucis novæ,\\
Mundi parans oríginem:

Qui mane junctum vésperi\\
Diem vocári præcipis:\\
Illábitur tetrum chaos,\\
Audi preces cum flétibus.

Ne mens graváta crímine,\\
Vitæ sit exsul múnere,\\
Dum nil perénne cógitat,\\
Seséque culpis ílligat.

Cæléste pulset óstium:\\
Vitále tollat præmium:\\
Vitémus omne nóxium:\\
Purgémus omne péssimum.

Præsta, Pater piíssime,\\
Patríque compar Unice,\\
Cum Spíritu Paráclito\\
Regnans per omne sæculum. Amen.

\V\ Dirigátur, Dómine, orátio mea. \R\ Sicut incénsum in conspéctu tuo.

\rubric{Ad Magnificat Antiphona ut in Proprio de Tempore.}

\canticumheading{Canticum B. Mariæ Virg.}{Luc. 1, 46-55}

\lettrine[lines=2]{M}{AGNIFICAT} \textdagger\ ánima mea Dóminum:

Et exsultávit spíritus meus * in Deo, salutári meo.

Quia respéxit humilitátem ancíllæ suæ: * ecce enim ex hoc beátam me dicent omnes generatiónes.

Quia fecit mihi magna qui potens est: * et sanctum nomen ejus.

Et misericórdia ejus a progénie in progénies * timéntibus eum.

Fecit poténtiam in bráchio suo: * dispérsit supérbos mente cordis sui.

Depósuit poténtes de sede, * et exaltávit húmiles.

Esuriéntes implévit bonis: * et dívites dimísit inánes.

Suscépit Israël, púerum suum, * recordátus misericórdiæ suæ.

Sicut locútus est ad patres nostros, * Abraham, et sémini ejus in sæcula.

\rubric{Oratio conveniens. Reliqua ut in Ordinario. 22}

\hourheading{Ad Completorium}

\V\ Jube, domne.

\benedlabel{Bened.} Noctem quiétam, et finem perféctum concédat nobis Dóminus omnípotens. \R\ Amen.

\lectioheading{Lectio br.}{1 Petri 5, 8-9}

\lettrine[lines=2]{F}{RATRES}: Sóbrii estóte, et vigiláte: quia adversárius vester diábolus tamquam leo rúgiens círcuit, quærens quem dévoret: cui resístite fortes in fide. Tu autem, Dómine, miserére nobis.

\R\ Deo grátias.

\V\ Adjutórium nostrum in nómine Dómini. \R\ Qui fecit cælum et terram.

\rubric{Pater noster, quod dicitur totum secreto.}

\formulalabel{Confíteor. Misereátur. Indulgéntiam.}

\V\ Convérte nos, Deus, salutáris noster. \R\ Et avérte iram tuam a nobis.

\V\ Deus, in adjutórium.

\antlabel{Ant.} Miserére.

\psalmheading{Psalmus 4}

\lettrine[lines=2]{C}{UM} invocárem exaudívit me Deus justítiæ meæ: * in tribulatióne dilatásti mihi.

Miserére mei, * et exáudi oratiónem meam.

Fílii hóminum, úsquequo gravi corde? * ut quid dilígitis vanitátem, et quæritis mendácium?

Et scitóte quóniam mirificávit Dóminus sanctum suum: * Dóminus exáudiet me cum clamávero ad eum.

Irascímini, et nolíte peccáre: * quæ dícitis in córdibus vestris, in cubílibus vestris compungímini.

Sacrificáte sacrifícium justítiæ, et speráte in Dómino. * Multi dicunt: Quis osténdit nobis bona?

Signátum est super nos lumen vultus tui, Dómine: * dedísti lætítiam in corde meo.

A fructu fruménti, vini, et ólei sui * multiplicáti sunt.

In pace in idípsum * dórmiam, et requiéscam;

Quóniam tu, Dómine, singuláriter in spe * constituísti me.

\psalmheading{Psalmus 90}

\lettrine[lines=2]{Q}{UI} hábitat in adjutório Altíssimi, * in protectióne Dei cæli commorábitur.

Dicet Dómino: Suscéptor meus es tu, et refúgium meum: * Deus meus, sperábo in eum.

Quóniam ipse liberávit me de láqueo venántium, * et a verbo áspero.

Scápulis suis obumbrábit tibi: * et sub pennis ejus sperábis.

Scuto circúmdabit te véritas ejus: * non timébis a timóre noctúrno.

A sagítta volánte in die, a negótio perambulánte in ténebris: * ab incúrsu, et dæmónio meridiáno.

Cadent a látere tuo mille, et decem míllia a dextris tuis: * ad te autem non appropinquábit.

Verúmtamen óculis tuis considerábis: * et retributiónem peccatórum vidébis.

Quóniam tu es, Dómine, spes mea: * Altíssimum posuísti refúgium tuum.

Non accédet ad te malum: * et flagéllum non appropinquábit tabernáculo tuo.

Quóniam Angelis suis mandávit de te: * ut custódiant te in ómnibus viis tuis.

In mánibus portábunt te: * ne forte offéndas ad lápidem pedem tuum.

Super áspidem, et basilíscum ambulábis: * et conculcábis leónem et dracónem.

Quóniam in me sperávit, liberábo eum: * prótegam eum, quóniam cognóvit nomen meum.

Clamábit ad me, et ego exáudiam eum: * cum ipso sum in tribulatióne: erípiam eum et glorificábo eum.

Longitúdine diérum replébo eum: * et osténdam illi salutáre meum.

\psalmheading{Psalmus 133}

\lettrine[lines=2]{E}{CCE}, nunc benedícite Dóminum, * omnes servi Dómini:

Qui statis in domo Dómini, * in átriis domus Dei nostri.

In nóctibus extóllite manus vestras in sancta, * et benedícite Dóminum.

Benedícat te Dóminus ex Sion, * qui fecit cælum et terram.

\antlabel{Ant.} Miserére mihi, Dómine, et exáudi oratiónem meam.

\hymnusheading{Hymnus}

\lettrine[lines=2]{T}{E} lucis ante términum,\\
Rerum Creátor, póscimus,\\
Ut pro tua cleméntia\\
Sis præsul et custódia.

Procul recédant sómnia,\\
Et nóctium phantásmata;\\
Hostémque nostrum cómprime,\\
Ne polluántur córpora.

Præsta, Pater piíssime,\\
Patríque compar Unice,\\
Cum Spíritu Paráclito\\
Regnans per omne sæculum. Amen.

\capitulumlabel{Capitulum}{Jer. 14, 9}

\lettrine[lines=2]{T}{U} autem in nobis es, Dómine, et nomen sanctum tuum invocátum est super nos: ne derelínquas nos, Dómine, Deus noster.

\R\ br. In manus tuas, Dómine, * Comméndo spíritum meum. In manus. \V\ Redemísti nos, Dómine, Deus veritátis. Comméndo. Glória Patri. In manus.

\V\ Custódi nos, Dómine, ut pupíllam óculi. \R\ Sub umbra alárum tuárum prótege nos.

\canticumheading{Canticum Simeonis}{Luc. 2, 29-32}

\lettrine[lines=2]{N}{UNC} dimíttis servum tuum, Dómine, * secúndum verbum tuum in pace:

Quia vidérunt óculi mei * salutáre tuum,

Quod parásti * ante fáciem ómnium populórum,

Lumen ad revelatiónem géntium, * et glóriam plebis tuæ Israël.

Glória Patri.

\antlabel{Ant.} Salva nos, Dómine, vigilántes, custódi nos dormiéntes; ut vigilémus cum Christo, et requiescámus in pace.

\variantheading{Preces (si dicendæ sint)}

KYRIE, eléison. Christe, eléison. Kýrie, eléison.

Pater noster \rubric{secreto usque ad} \V\ Et ne nos indúcas in tentatiónem. \R\ Sed líbera nos a malo.

Credo in Deum \rubric{secreto usque ad} \V\ Carnis resurrectiónem. \R\ Vitam ætérnam. Amen.

\V\ Benedíctus es, Dómine, Deus patrum nostrórum. \R\ Et laudábilis et gloriósus in sæcula. \V\ Benedicámus Patrem et Fílium cum Sancto Spíritu. \R\ Laudémus, et superexaltémus eum in sæcula. \V\ Benedíctus es, Dómine, in firmaménto cæli. \R\ Et laudábilis, et gloriósus, et superexaltátus in sæcula. \V\ Benedícat et custódiat nos omnípotens et miséricors Dóminus. \R\ Amen. \V\ Dignáre, Dómine, nocte ista. \R\ Sine peccáto nos custodíre. \V\ Miserére nostri, Dómine. \R\ Miserére nostri. \V\ Fiat misericórdia tua, Dómine, super nos. \R\ Quemádmodum sperávimus in te. \V\ Dómine, exáudi oratiónem meam. \R\ Et clamor meus ad te véniat.

\V\ Dóminus vobíscum.

Orémus. \oratiolabel{Oratio}

\lettrine[lines=2]{V}{ISITA}, quæsumus, Dómine, habitatiónem istam, et omnes insídias inimíci ab ea longe repélle: Angeli tui sancti hábitent in ea, qui nos in pace custódiant; et benedíctio tua sit super nos semper. Per Dóminum nostrum. \R\ Amen.

\V\ Dóminus vobíscum.

\V\ Benedicámus Dómino.

\benedictiolabel{Benedictio} Benedícat et custódiat nos omnípotens et miséricors Dóminus, Pater, et Fílius, et Spíritus Sanctus. \R\ Amen.

\variantheading{Antiphona finalis B. M. V.}

\V\ Divínum auxílium.

Pater, Ave et Credo.

\rubric{¶ Psalmi de Dominica superius appositi dicuntur ad Completorium, adhibita simul Antiphona, etiam post utrasque Vesperas in Dominicis infra Octavas Ssmi Corporis Christi ac Sacrat. Cordis Jesu, in Duplicibus I vel II classis, in Festis novem Lectionum Domini, B. Mariæ Virginis, Angelorum, S. Joannis Baptistæ, S. Joseph, Apostolorum, Evangelistarum, et per Octavas privilegiatas; si tamen præcedentes I Vesperæ de eis integræ vel saltem a Capitulo sint celebratæ, ac II Vesperæ de ipsis integræ fuerint persolutæ.}

% ===== END OF BATCH 8 -- Covers autumnalis.md lines 2478-2863,
% completing "Ad Vesperas" (Psalms 109-113, Capitulum, hymn "Lucis
% Creator optime," and the Magnificat) and "Ad Completorium" (Lectio
% brevis, Confiteor shorthand, Psalms 4, 90, 133, hymn "Te lucis ante
% terminum," Nunc dimittis, and the Preces), and with them the ENTIRE
% "Dominica" Office. Two more occurrences of the recurring stray "\*"
% artifact (in "Lucis Creator optime" at line 2703 and "Te lucis ante
% terminum" at line 2842, both at the hymn's final stanza) are
% silently dropped, consistent with every prior occurrence in this
% volume. The closing "Antiphona finalis B. M. V." is given only as a
% bare cross-reference label with the incipit "℣. Divínum auxílium."
% (not the full "...máneat semper nobíscum. ℟. Amen."), transcribed
% exactly as abbreviated rather than completed from the Ordinarium's
% fuller text. The document continues with "Feria Secunda" (autumnalis.md
% line 2864), the first weekday of the Psalterium cycle =====
\feriaheading{Feria Secunda}

\hourheading{Ad Matutinum}

Pater, Ave et Credo.

\V\ Dómine, lábia mea.

\V\ Deus, in adjutórium.

\variantheading{In Officio feriali}

\rubric{Invitatorium et Hymnus dicuntur quotidie de Feria currenti, ut in Psalterio hic et alibi apponuntur.}

\invitatlabel{Invitat.} Veníte, * Exsultémus Dómino. \psalmlabel{Ps. 94} Veníte, exsultémus.

\rubric{In primo Psalmi versu non dicitur: Veníte, exsultémus Dómino; sed post repetitum Invitatorium, statim subjungitur: Jubilémus Deo, salutári nostro:}

\hymnusheading{Hymnus}

\lettrine[lines=2]{S}{OMNO} reféctis ártubus,\\
Spreto cubíli, súrgimus:\\
Nobis, Pater, canéntibus\\
Adésse te depóscimus.

Te lingua primum cóncinat,\\
Te mentis ardor ámbiat:\\
Ut áctuum sequéntium\\
Tu, Sancte, sis exórdium.

Cedant ténebræ lúmini,\\
Et nox diúrno síderi,\\
Ut culpa, quam nox íntulit,\\
Lucis labáscat múnere.

Precámur iídem súpplices,\\
Noxas ut omnes ámputes,\\
Et ore te canéntium\\
Laudéris omni témpore.

Præsta, Pater piíssime,\\
Patríque compar Unice,\\
Cum Spíritu Paráclito\\
Regnans per omne sæculum.\\
Amen.

\nocturnoheading{In I Nocturno}

\antlabel{Ant.} Dóminus de cælo * prospéxit super fílios hóminum.

\rubric{Antiphonæ, in hac et aliis Feriis appositæ, dicuntur semper in Officio feriali et in Festis sive trium sive novem Lectionum, quando juxta Rubricas sumendi sint Psalmi de Feria, ut in Psalterio.}

\rubric{In Officio autem trium Lectionum, tam de Tempore quam de Sanctis, omittuntur Versus et reliqua omnia, quæ in Officio novem Lectionum solent recitari in fine I et II Nocturni, et continuantur Psalmi cum suis Antiphonis usque ad Versum III Nocturni.}

\psalmheading{Psalmus 13}

\lettrine[lines=2]{D}{IXIT} insípiens in corde suo: * Non est Deus.

Corrúpti sunt, et abominábiles facti sunt in stúdiis suis: * non est qui fáciat bonum, non est usque ad unum.

Dóminus de cælo prospéxit super fílios hóminum, * ut vídeat si est intélligens, aut requírens Deum.

Omnes declinavérunt, simul inútiles facti sunt: * non est qui fáciat bonum, non est usque ad unum.

Sepúlcrum patens est guttur eórum: linguis suis dolóse agébant, * venénum áspidum sub lábiis eórum.

Quorum os maledictióne et amaritúdine plenum est: * velóces pedes eórum ad effundéndum sánguinem.

Contrítio et infelícitas in viis eórum, et viam pacis non cognovérunt: * non est timor Dei ante óculos eórum.

Nonne cognóscent omnes qui operántur iniquitátem, * qui dévorant plebem meam sicut escam panis?

Dóminum non invocavérunt, * illic trepidavérunt timóre, ubi non erat timor.

Quóniam Dóminus in generatióne justa est, consílium ínopis confudístis: * quóniam Dóminus spes ejus est.

Quis dabit ex Sion salutáre Israël? * cum avérterit Dóminus captivitátem plebis suæ, exsultábit Jacob, et lætábitur Israël.

\antlabel{Ant.} Dóminus de cælo prospéxit super fílios hóminum.

\antlabel{Ant.} Qui operátur justítiam, * requiéscet in monte sancto tuo, Dómine.

\psalmheading{Psalmus 14}

\lettrine[lines=2]{D}{OMINE}, quis habitábit in tabernáculo tuo? * aut quis requiéscet in monte sancto tuo?

Qui ingréditur sine mácula, * et operátur justítiam:

Qui lóquitur veritátem in corde suo, * qui non egit dolum in lingua sua:

Nec fecit próximo suo malum, * et oppróbrium non accépit advérsus próximos suos.

Ad níhilum dedúctus est in conspéctu ejus malígnus: * timéntes autem Dóminum gloríficat:

Qui jurat próximo suo, et non décipit, * qui pecúniam suam non dedit ad usúram, et múnera super innocéntem non accépit.

Qui facit hæc, * non movébitur in ætérnum.

\antlabel{Ant.} Qui operátur justítiam, requiéscet in monte sancto tuo, Dómine.

\antlabel{Ant.} Inclína, Dómine, * aurem tuam mihi, et exáudi verba mea.

\psalmheading{Psalmus 16}

\lettrine[lines=2]{E}{XAUDI}, Dómine, justítiam meam: * inténde deprecatiónem meam.

Áuribus pércipe oratiónem meam, * non in lábiis dolósis.

De vultu tuo judícium meum pródeat: * óculi tui vídeant æquitátes.

Probásti cor meum, et visitásti nocte: * igne me examinásti, et non est invénta in me iníquitas.

Ut non loquátur os meum ópera hóminum: * propter verba labiórum tuórum ego custodívi vias duras.

Pérfice gressus meos in sémitis tuis: * ut non moveántur vestígia mea.

Ego clamávi, quóniam exaudísti me, Deus: * inclína aurem tuam mihi, et exáudi verba mea.

Mirífica misericórdias tuas, * qui salvos facis sperántes in te.

A resisténtibus déxteræ tuæ custódi me, * ut pupíllam óculi.

Sub umbra alárum tuárum prótege me: * a fácie impiórum qui me afflixérunt.

Inimíci mei ánimam meam circumdedérunt, ádipem suum conclusérunt: * os eórum locútum est supérbiam.

Projiciéntes me nunc circumdedérunt me: * óculos suos statuérunt declináre in terram.

Suscepérunt me sicut leo parátus ad prǽdam: * et sicut cátulus leónis hábitans in abdítis.

Exsúrge, Dómine, prǽveni eum, et supplánta eum: * éripe ánimam meam ab ímpio, fráméam tuam ab inimícis manus tuæ.

Dómine, a paucis de terra dívide eos in vita eórum: * de abscónditis tuis adimplétus est venter eórum.

Saturáti sunt fíliis: * et dimisérunt relíquias suas párvulis suis.

Ego autem in justítia apparébo conspéctui tuo: * satiábor cum apparúerit glória tua.

\antlabel{Ant.} Inclína, Dómine, aurem tuam mihi, et exáudi verba mea.

\variantheading{In Festis novem Lectionum:}

\V\ Mirífica, Dómine, misericórdias tuas. \R\ Qui salvos facis sperántes in te.

\absolutiolabel{Absol.} Exáudi, Dómine.

\benedlabel{Bened.} 1 Benedictióne perpétua. 2 Unigénitus Dei Fílius. 3 Spíritus Sancti grátia.

\nocturnoheading{In II Nocturno}

\antlabel{Ant.} Díligam te, * \textdagger\ Dómine, virtus mea.

\psalmheading{Psalmus 17, i}

\lettrine[lines=2]{D}{ILIGAM} te, \textdagger\ Dómine, fortitúdo mea: * Dóminus firmaméntum meum, et refúgium meum, et liberátor meus.

Deus meus adjútor meus, * et sperábo in eum.

Protéctor meus, et cornu salútis meæ, * et suscéptor meus.

Laudans invocábo Dóminum: * et ab inimícis meis salvus ero.

Circumdedérunt me dolóres mortis: * et torréntes iniquitátis conturbavérunt me.

Dolóres inférni circumdedérunt me: * præoccupavérunt me láquei mortis.

In tribulatióne mea invocávi Dóminum, * et ad Deum meum clamávi.

Et exaudívit de templo sancto suo vocem meam: * et clamor meus in conspéctu ejus, introívit in aures ejus.

Commóta est, et contrémuit terra: * fundaménta móntium conturbáta sunt, et commóta sunt, quóniam irátus est eis.

Ascéndit fumus in ira ejus: * et ignis a fácie ejus exársit: * carbónes succénsi sunt ab eo.

Inclinávit cælos, et descéndit: * et calígo sub pédibus ejus.

Et ascéndit super Chérubim, et volávit: * volávit super pennas ventórum.

Et pósuit ténebras latíbulum suum, in circúitu ejus tabernáculum ejus: * tenebrósa aqua in núbibus áëris.

Præ fulgóre in conspéctu ejus nubes transiérunt, * grando et carbónes ignis.

Et intónuit de cælo Dóminus, et Altíssimus dedit vocem suam: * grando et carbónes ignis.

Et misit sagíttas suas, et dissipávit eos: * fúlgura multiplicávit, et conturbávit eos.

Et apparuérunt fontes aquárum, * et reveláta sunt fundaménta orbis terrárum:

Ab increpatióne tua, Dómine, * ab inspiratióne spíritus iræ tuæ.

\antlabel{Ant.} Díligam te, Dómine, virtus mea.

\antlabel{Ant.} Retríbuet mihi Dóminus * secúndum justítiam meam.

\psalmheading{Psalmus 17, ii}

\lettrine[lines=2]{M}{ISIT} de summo, et accépit me: * et assúmpsit me de aquis multis.

Erípuit me de inimícis meis fortíssimis, et ab his qui odérunt me: * quóniam confortáti sunt super me.

Prævenérunt me in die afflictiónis meæ: * et factus est Dóminus protéctor meus.

Et edúxit me in latitúdinem: * salvum me fecit, quóniam vóluit me.

Et retríbuet mihi Dóminus secúndum justítiam meam, * et secúndum puritátem mánuum meárum retríbuet mihi:

Quia custodívi vias Dómini, * nec ímpie gessi a Deo meo.

Quóniam ómnia judícia ejus in conspéctu meo: * et justítias ejus non répuli a me.

Et ero immaculátus cum eo: * et observábo me ab iniquitáte mea.

Et retríbuet mihi Dóminus secúndum justítiam meam: * et secúndum puritátem mánuum meárum in conspéctu oculórum ejus.

Cum sancto sanctus eris, * et cum viro innocénte ínnocens eris:

Et cum elécto eléctus eris: * et cum pervérso pervertéris.

Quóniam tu pópulum húmilem salvum fácies: * et óculos superbórum humiliábis.

Quóniam tu illúminas lucérnam meam, Dómine: * Deus meus, illúmina ténebras meas.

Quóniam in te erípiar a tentatióne, * et in Deo meo transgrédiar murum.

Deus meus, impollúta via ejus: elóquia Dómini igne examináta: * protéctor est ómnium sperántium in se.

Quóniam quis Deus præter Dóminum? * aut quis Deus præter Deum nostrum?

Deus, qui præcínxit me virtúte: * et pósuit immaculátam viam meam.

Qui perfécit pedes meos tamquam cervórum, * et super excélsa státuens me.

Qui docet manus meas ad prǽlium: * et posuísti, ut arcum áëreum, bráchia mea.

\antlabel{Ant.} Retríbuet mihi Dóminus secúndum justítiam meam.

\antlabel{Ant.} Vivit Dóminus, * et benedíctus Deus salútis meæ.

\psalmheading{Psalmus 17, iii}

\lettrine[lines=2]{E}{T} dedísti mihi protectiónem salútis tuæ: * et déxtera tua suscépit me:

Et disciplína tua corréxit me in finem: * et disciplína tua ipsa me docébit.

Dilatásti gressus meos subtus me: * et non sunt infirmáta vestígia mea:

Pérsequar inimícos meos, et comprehéndam illos: * et non convértar, donec defíciant.

Confríngam illos, nec póterunt stare: * cadent subtus pedes meos.

Et præcinxísti me virtúte ad bellum: * et supplantásti insurgéntes in me subtus me.

Et inimícos meos dedísti mihi dorsum, * et odiéntes me disperdidísti.

Clamavérunt, nec erat qui salvos fáceret, ad Dóminum: * nec exaudívit eos.

Et commínuam eos, ut púlverem ante fáciem venti: * ut lutum plateárum delébo eos.

Erípies me de contradictiónibus pópuli: * constítues me in caput géntium.

Pópulus, quem non cognóvi, servívit mihi: * in audítu auris obedívit mihi.

Fílii aliéni mentíti sunt mihi: * fílii aliéni inveteráti sunt, et claudicavérunt a sémitis suis.

Vivit Dóminus, et benedíctus Deus meus, * et exaltétur Deus salútis meæ.

Deus, qui das vindíctas mihi, et subdis pópulos sub me, * liberátor meus de inimícis meis iracúndis.

Et ab insurgéntibus in me exaltábis me: * a viro iníquo éripies me.

Proptérea confitébor tibi in natiónibus, Dómine: * et nómini tuo psalmum dicam.

Magníficans salútes Regis ejus, et fáciens misericórdiam Christo suo David, * et sémini ejus usque in séculum.

\antlabel{Ant.} Vivit Dóminus, et benedíctus Deus salútis meæ.

\variantheading{In Festis novem Lectionum:}

\V\ Confitébor tibi in natiónibus, Dómine. \R\ Et nómini tuo psalmum dicam.

\absolutiolabel{Absol.} Ipsíus píetas et.

\benedlabel{Bened.} 4 Deus Pater omnípotens. 5 Christus perpétuæ. 6 Ignem sui amóris.

\nocturnoheading{In III Nocturno}

\antlabel{Ant.} Exáudiat te * \textdagger\ Dóminus in die tribulatiónis. \textdagger

\psalmheading{Psalmus 19}

\lettrine[lines=2]{E}{XAUDIAT} te \textdagger\ Dóminus in die tribulatiónis: * \textdagger\ prótegat te nomen Dei Jacob.

Mittat tibi auxílium de sancto: * et de Sion tueátur te.

Memor sit omnis sacrifícii tui: * et holocáustum tuum pingue fiat.

Tríbuat tibi secúndum cor tuum: * et omne consílium tuum confírmet.

Lætábimur in salutári tuo: * et in nómine Dei nostri magnificábimur.

Ímpleat Dóminus omnes petitiónes tuas: * nunc cognóvi quóniam salvum fecit Dóminus Christum suum.

Exáudiet illum de cælo sancto suo: * in potentátibus salus déxteræ ejus.

Hi in cúrribus, et hi in equis: * nos autem in nómine Dómini, Dei nostri, invocábimus.

Ipsi obligáti sunt, et cecidérunt: * nos autem surréximus et erécti sumus.

Dómine, salvum fac regem: * et exáudi nos in die, qua invocavérimus te.

\antlabel{Ant.} Exáudiat te Dóminus in die tribulatiónis.

\antlabel{Ant.} Dómine, * \textdagger\ in virtúte tua lætábitur rex. \textdaggerdbl

\psalmheading{Psalmus 20}

\lettrine[lines=2]{D}{OMINE}, \textdagger\ in virtúte tua lætábitur rex: * \textdaggerdbl\ et super salutáre tuum exsultábit veheménter.

Desidérium cordis ejus tribuísti ei: * et voluntáte labiórum ejus non fraudásti eum.

Quóniam prævenísti eum in benedictiónibus dulcédinis: * posuísti in cápite ejus corónam de lápide pretióso.

Vitam pétiit a te: * et tribuísti ei longitúdinem diérum in séculum, et in séculum séculi.

Magna est glória ejus in salutári tuo: * glóriam et magnum decórem impónes super eum.

Quóniam dabis eum in benedictiónem in séculum séculi: * lætificábis eum in gáudio cum vultu tuo.

Quóniam rex sperat in Dómino: * et in misericórdia Altíssimi non commovébitur.

Inveniátur manus tua ómnibus inimícis tuis: * déxtera tua invéniat omnes, qui te odérunt.

Pones eos ut clíbanum ignis in témpore vultus tui: * Dóminus in ira sua conturbábit eos, et devorábit eos ignis.

Fructum eórum de terra perdes: * et semen eórum a fíliis hóminum.

Quóniam declinavérunt in te mala: * cogitavérunt consília, quæ non potuérunt stabilíre.

Quóniam pones eos dorsum: * in relíquiis tuis præparábis vultum eórum.

Exaltáre, Dómine, in virtúte tua: * cantábimus et psallémus virtútes tuas.

\antlabel{Ant.} Dómine, in virtúte tua lætábitur rex.

\antlabel{Ant.} Exaltábo te, * \textdagger\ Dómine, quóniam suscepísti me. \textdaggerdbl

\psalmheading{Psalmus 29}

\lettrine[lines=2]{E}{XALTABO} te, \textdagger\ Dómine, quóniam suscepísti me: * \textdaggerdbl\ nec delectásti inimícos meos super me.

Dómine, Deus meus, clamávi ad te, * et sanásti me.

Dómine, eduxísti ab inférno ánimam meam: * salvásti me a descendéntibus in lacum.

Psállite Dómino, sancti ejus: * et confitémini memóriæ sanctitátis ejus.

Quóniam ira in indignatióne ejus: * et vita in voluntáte ejus.

Ad vésperum demorábitur fletus: * et ad matutínum lætítia.

Ego autem dixi in abundántia mea: * Non movébor in ætérnum.

Dómine, in voluntáte tua, * præstitísti decóri meo virtútem.

Avertísti fáciem tuam a me, * et factus sum conturbátus.

Ad te, Dómine, clamábo: * et ad Deum meum deprecábor.

Quæ utílitas in sánguine meo, * dum descéndo in corruptiónem?

Numquid confitébitur tibi pulvis, * aut annuntiábit veritátem tuam?

Audívit Dóminus, et misértus est mei: * Dóminus factus est adjútor meus.

Convertísti planctum meum in gáudium mihi: * conscidísti saccum meum, et circumdedísti me lætítia:

Ut cantet tibi glória mea, et non compúngar: * Dómine, Deus meus, in ætérnum confitébor tibi.

\antlabel{Ant.} Exaltábo te, Dómine, quóniam suscepísti me.

\variantheading{In feriali Officio et in Festis:}

\V\ Psállite Dómino, sancti ejus. \R\ Et confitémini memóriæ sanctitátis ejus.

\variantheading{In Officio novem Lectionum}

\absolutiolabel{Absolutio} A vínculis.

\benedlabel{Bened.} 7 Evangélica léctio. 8 In Festis Domini: Divínum auxílium. In Festis B. M. V.: Cujus festum cólimus, ipsa Virgo. In Festis Sanctorum: Cujus (\textit{vel} Quorum \textit{aut} Quarum) festum cólimus. 9 Ad societátem. (Per evangélica dicta.)

\variantheading{In Officio trium Lectionum}

\absolutiolabel{Absol.} Exáudi, Dómine.

\variantheading{In Feriis, cum Homilia:}

\benedlabel{Bened.} 1 Evangélica. 2 Divínum. 3 Ad societátem.

\variantheading{In Feriis, sine Homilia:}

\benedlabel{Bened.} 1 Benedictióne. 2 Unigénitus. 3 Spíritus Sancti.

\variantheading{In Festis Sanctorum:}

\benedlabel{Bened.} 1 Ille nos benedícat qui sine fine vivit et regnat. \R\ Amen. 2 Cujus (\textit{vel} Quorum \textit{aut} Quarum). 3 Ad societátem.

% ===== END OF BATCH 9 -- Covers autumnalis.md lines 2864-3269, opening
% "Feria Secunda" (Monday) and completing its "Ad Matutinum" in
% full: the Invitatorium "Veníte, Exsultémus Dómino," the hymn
% "Somno reféctis artubus," and the three Nocturns (Psalms 13, 14,
% 16; 17/i, 17/ii, 17/iii; 19, 20, 29), each followed by its
% condition-specific versicle/Absolutio/Benedictiones apparatus
% (the ferial psalter prints these only "In Festis novem
% Lectionum"/"In Officio novem Lectionum" etc., since on an ordinary
% ferial day with three lessons only, most of this apparatus is
% omitted per the rubric given at the start of this Nocturnal set).
% Several antiphons and psalm-incipits in the Second and Third
% Nocturns carry multiple chant-pointing marks in a single line --
% asterisk, dagger, and even double-dagger together (e.g. "Dómine, *
% † in virtúte tua lætábitur rex. ‡") -- each rendered in its exact
% source position as *, \textdagger\, and \textdaggerdbl\
% respectively, with no marks added, dropped, or reordered. The
% compact multi-option "Bened." line for the ninth lesson ("8 In
% Festis Domini: Divínum auxílium. In Festis B. M. V.: ... In Festis
% Sanctorum: ...") is transcribed as one continuous \benedlabel{}
% argument exactly as the source runs it together. The document
% continues with "Ad Laudes I" (autumnalis.md line 3270), the first of
% two alternate ferial Lauds schemes =====
\hourheading{Ad Laudes I}

\rubric{In Feriis extra Vigílias commúnes occurréntibus, atque in Festis.}

(Pater noster et Ave.)

\V\ Deus, in adjutórium.

\antlabel{Ant.} Jubiláte * Deo in voce exsultatiónis.

\rubric{¶ Antíphonæ, in hac et áliis Fériis apposítæ, dicúntur semper in Officio feriáli et in Festis sive trium sive novem Lectiónum, quando juxta Rúbricas suméndi sint Psalmi de Féria.}

\psalmheading{Psalmus 46}

\lettrine[lines=2]{O}{MNES} gentes, pláudite mánibus: * jubiláte Deo in voce exsultatiónis.

Quóniam Dóminus excélsus, terríbilis: * Rex magnus super omnem terram.

Subjécit pópulos nobis: * et gentes sub pédibus nostris.

Elégit nobis hereditátem suam: * spéciem Jacob, quam diléxit.

Ascéndit Deus in júbilo: * et Dóminus in voce tubæ.

Psállite Deo nostro, psállite: * psállite Regi nostro, psállite.

Quóniam Rex omnis terræ Deus: * psállite sapiénter.

Regnábit Deus super gentes: * Deus sedet super sedem sanctam suam.

Príncipes populórum congregáti sunt cum Deo Abraham: * quóniam dii fortes terræ veheménter eleváti sunt.

\antlabel{Ant.} Jubiláte Deo in voce exsultatiónis.

\antlabel{Ant.} Inténde * voci oratiónis meæ, Rex meus et Deus meus.

\psalmheading{Psalmus 5}

\lettrine[lines=2]{V}{ERBA} mea áuribus pércipe, Dómine, * intéllige clamórem meum.

Inténde voci oratiónis meæ, * Rex meus et Deus meus.

Quóniam ad te orábo: * Dómine, mane exáudies vocem meam.

Mane astábo tibi et vidébo: * quóniam non Deus volens iniquitátem tu es.

Neque habitábit juxta te malígnus: * neque permanébunt injústi ante óculos tuos.

Odísti omnes, qui operántur iniquitátem: * perdes omnes, qui loquúntur mendácium.

Virum sánguinum et dolósum abominábitur Dóminus: * ego autem in multitúdine misericórdiæ tuæ.

Introíbo in domum tuam: * adorábo ad templum sanctum tuum in timóre tuo.

Dómine, deduc me in justítia tua: * propter inimícos meos dírige in conspéctu tuo viam meam.

Quóniam non est in ore eórum véritas: * cor eórum vanum est.

Sepúlcrum patens est guttur eórum, linguis suis dolóse agébant, * júdica illos, Deus.

Décidant a cogitatiónibus suis, secúndum multitúdinem impietátum eórum expélle eos, * quóniam irritavérunt te, Dómine.

Et læténtur omnes, qui sperant in te, * in ætérnum exsultábunt: et habitábis in eis.

Et gloriabúntur in te omnes, qui díligunt nomen tuum, * quóniam tu benedíces justo.

Dómine, ut scuto bonæ voluntátis tuæ * coronásti nos.

\antlabel{Ant.} Inténde voci oratiónis meæ, Rex meus et Deus meus.

\antlabel{Ant.} Deus majestátis * intónuit: afférte glóriam nómini ejus.

\psalmheading{Psalmus 28}

\lettrine[lines=2]{A}{FFERTE} Dómino, fílii Dei: * afférte Dómino fílios aríetum.

Afférte Dómino glóriam et honórem, afférte Dómino glóriam nómini ejus: * adoráte Dóminum in átrio sancto ejus.

Vox Dómini super aquas, Deus majestátis intónuit: * Dóminus super aquas multas.

Vox Dómini in virtúte: * vox Dómini in magnificéntia.

Vox Dómini confringéntis cedros: * et confrínget Dóminus cedros Líbani:

Et commínuet eas tamquam vítulum Líbani: * et diléctus quemádmodum fílius unicórnium.

Vox Dómini intercidéntis flammam ignis: * vox Dómini concutiéntis desértum: et commovébit Dóminus desértum Cades.

Vox Dómini præparántis cervos, et revelábit condénsa: * et in templo ejus omnes dicent glóriam.

Dóminus dilúvium inhabitáre facit: * et sedébit Dóminus Rex in ætérnum.

Dóminus virtútem pópulo suo dabit: * Dóminus benedícet pópulo suo in pace.

\antlabel{Ant.} Deus majestátis intónuit: afférte glóriam nómini ejus.

\antlabel{Ant.} Laudámus nomen tuum * ínclitum, Deus noster.

\canticumheading{Canticum David}{1 Par. 29, 10-13}

\lettrine[lines=2]{B}{ENEDICTUS} es, Dómine, Deus Israël patris nostri, * ab ætérno in ætérnum.

Tua est, Dómine, magnificéntia, et poténtia, * et glória, atque victória:

Et tibi laus: * cuncta enim quæ in cælo sunt, et in terra, tua sunt:

Tuum, Dómine, regnum, * et tu es super omnes príncipes.

Tuæ divítiæ, et tua est glória: * tu domináris ómnium.

In manu tua virtus et poténtia: * in manu tua magnitúdo, et impérium ómnium.

Nunc ígitur, Deus noster, confitémur tibi, * et laudámus nomen tuum ínclitum.

\antlabel{Ant.} Laudámus nomen tuum ínclitum, Deus noster.

\antlabel{Ant.} Laudáte * \textdagger\ Dóminum, omnes gentes. \textdaggerdbl

\psalmheading{Psalmus 116}

\lettrine[lines=2]{L}{AUDATE} \textdagger\ Dóminum, omnes gentes: * \textdaggerdbl\ laudáte eum, omnes pópuli:

Quóniam confirmáta est super nos misericórdia ejus: * et véritas Dómini manet in ætérnum.

\antlabel{Ant.} Laudáte Dóminum, omnes gentes.

\variantheading{In Officio feriali}

\rubric{Capítulum, Hymnus et Versus dicúntur quotídie de Féria curránti, ut in Psaltério hic et álibi appónuntur; et inde súmitur étiam Ant. ad Benedíctus.}

\capitulumlabel{Capit.}{Rom. 13, 12-13}

\lettrine[lines=2]{N}{OX} præcéssit, dies autem appropinquávit. Abjiciámus ergo ópera tenebrárum, et induámur arma lucis. Sicut in die honéste ambulémus.

\hymnusheading{Hymnus}

\lettrine[lines=2]{S}{PLENDOR} Patérnæ glóriæ,\\
De luce lucem próferens,\\
Lux lucis, et fons lúminis,\\
Diem dies illúminans:

Verúsque sol illábere,\\
Micans nitóre pérpeti:\\
Jubárque Sancti Spíritus\\
Infúnde nostris sénsibus.

Votis vocémus et Patrem,\\
Patrem poténtis grátiæ,\\
Patrem perénnis glóriæ:\\
Culpam reléget lúbricam.

Confírmet actus strénuos:\\
Dentes retúndat ínvidi:\\
Casus secúndet ásperos:\\
Agénda recte dírigat.

Mentem gubérnet et regat:\\
Sit pura nobis cástitas:\\
Fides calóre férveat,\\
Fraudis venéna nésciat.

Christúsque nobis sit cibus,\\
Potúsque noster sit fides:\\
Læti bibámus sóbriam\\
Profusiónem Spíritus.

Lætus dies hic tránseat:\\
Pudor sit ut dilúculum:\\
Fides velut merídies:\\
Crepúsculum mens nésciat.

Auróra lucem próvehit,\\
Cum luce nobis pródeat\\
In Patre totus Fílius,\\
Et totus in Verbo Pater.

Deo Patri sit glória,\\
Ejúsque soli Fílio,\\
Cum Spíritu Paráclito,\\
Nunc et per omne sæculum. Amen.

\V\ Repléti sumus mane misericórdia tua. \R\ Exsultávimus, et delectáti sumus.

\antlabel{Ad Bened. Ant.} Benedíctus * \textdagger\ Dóminus, Deus Israël, quia visitávit et liberávit nos.

\rubric{Réliqua ut in Ordinário. 9}

\hourheading{Ad Laudes II}

\rubric{In Vigíliis commúnibus, quando fit Offícium de Féria.}

(Pater noster et Ave.)

\V\ Deus, in adjutórium.

\antlabel{Ant.} Miserére. \textdagger

\psalmheading{Psalmus 50}

\lettrine[lines=2]{M}{ISERERE} \textdagger\ mei, Deus: * secúndum magnam misericórdiam tuam.

Et secúndum multitúdinem miseratiónum tuárum, * dele iniquitátem meam.

Amplius lava me ab iniquitáte mea: * et a peccáto meo munda me.

Quóniam iniquitátem meam ego cognósco: * et peccátum meum contra me est semper.

Tibi soli peccávi, et malum coram te feci: * ut justificéris in sermónibus tuis, et vincas cum judicáris.

Ecce enim, in iniquitátibus concéptus sum: * et in peccátis concépit me mater mea.

Ecce enim, veritátem dilexísti: * incérta et occúlta sapiéntiæ tuæ manifestásti mihi.

Aspérges me hyssópo, et mundábor: * lavábis me, et super nivem dealbábor.

Audítui meo dabis gáudium et lætítiam: * et exsultábunt ossa humiliáta.

Avérte fáciem tuam a peccátis meis: * et omnes iniquitátes meas dele.

Cor mundum crea in me, Deus: * et spíritum rectum ínnova in viscéribus meis.

Ne projícias me a fácie tua: * et spíritum sanctum tuum ne áuferas a me.

Redde mihi lætítiam salutáris tui: * et spíritu principáli confírma me.

Docébo iníquos vias tuas: * et ímpii ad te converténtur.

Líbera me de sanguínibus, Deus, Deus salútis meæ: * et exsultábit lingua mea justítiam tuam.

Dómine, lábia mea apéries: * et os meum annuntiábit laudem tuam.

Quóniam si voluísses sacrifícium, dedíssem útique: * holocáustis non delectáberis.

Sacrifícium Deo spíritus contribulátus: * cor contrítum, et humiliátum, Deus, non despícies.

Benígne fac, Dómine, in bona voluntáte tua Sion: * ut ædificéntur muri Jerúsalem.

Tunc acceptábis sacrifícium justítiæ, oblatiónes, et holocáusta: * tunc impónent super altáre tuum vítulos.

\antlabel{Ant.} Miserére mei, Deus, secúndum magnam misericórdiam tuam.

\antlabel{Ant.} Deduc me.

\psalmheading{Psalmus 5}

\lettrine[lines=2]{V}{ERBA} mea áuribus pércipe, Dómine, * intéllige clamórem meum.

Inténde voci oratiónis meæ, * Rex meus et Deus meus.

Quóniam ad te orábo: * Dómine, mane exáudies vocem meam.

Mane astábo tibi et vidébo: * quóniam non Deus volens iniquitátem tu es.

Neque habitábit juxta te malígnus: * neque permanébunt injústi ante óculos tuos.

Odísti omnes, qui operántur iniquitátem: * perdes omnes, qui loquúntur mendácium.

Virum sánguinum et dolósum abominábitur Dóminus: * ego autem in multitúdine misericórdiæ tuæ.

Introíbo in domum tuam: * adorábo ad templum sanctum tuum in timóre tuo.

Dómine, deduc me in justítia tua: * propter inimícos meos dírige in conspéctu tuo viam meam.

Quóniam non est in ore eórum véritas: * cor eórum vanum est.

Sepúlcrum patens est guttur eórum, linguis suis dolóse agébant, * júdica illos, Deus.

Décidant a cogitatiónibus suis, secúndum multitúdinem impietátum eórum expélle eos, * quóniam irritavérunt te, Dómine.

Et læténtur omnes, qui sperant in te, * in ætérnum exsultábunt: et habitábis in eis.

Et gloriabúntur in te omnes, qui díligunt nomen tuum, * quóniam tu benedíces justo.

Dómine, ut scuto bonæ voluntátis tuæ * coronásti nos.

\antlabel{Ant.} Deduc me in justítia tua, Dómine.

\antlabel{Ant.} Dóminus dabit virtútem.

\psalmheading{Psalmus 28}

\lettrine[lines=2]{A}{FFERTE} Dómino, fílii Dei: * afférte Dómino fílios aríetum.

Afférte Dómino glóriam et honórem, afférte Dómino glóriam nómini ejus: * adoráte Dóminum in átrio sancto ejus.

Vox Dómini super aquas, Deus majestátis intónuit: * Dóminus super aquas multas.

Vox Dómini in virtúte: * vox Dómini in magnificéntia.

Vox Dómini confringéntis cedros: * et confrínget Dóminus cedros Líbani:

Et commínuet eas tamquam vítulum Líbani: * et diléctus quemádmodum fílius unicórnium.

Vox Dómini intercidéntis flammam ignis: * vox Dómini concutiéntis desértum: et commovébit Dóminus desértum Cades.

Vox Dómini præparántis cervos, et revelábit condénsa: * et in templo ejus omnes dicent glóriam.

Dóminus dilúvium inhabitáre facit: * et sedébit Dóminus Rex in ætérnum.

Dóminus virtútem pópulo suo dabit: * Dóminus benedícet pópulo suo in pace.

\antlabel{Ant.} Dóminus dabit virtútem et benedícet pópulo suo in pace.

\antlabel{Ant.} Convérsus est furor tuus.

\canticumheading{Canticum Isaiæ}{Is. 12, 1-6}

\lettrine[lines=2]{C}{ONFITEBOR} tibi, Dómine, quóniam irátus es mihi: * convérsus est furor tuus, et consolátus es me.

Ecce Deus salvátor meus, * fiduciáliter agam, et non timébo:

Quia fortitúdo mea, et laus mea Dóminus, * et factus est mihi in salútem.

Hauriétis aquas in gáudio de fóntibus Salvatóris: * et dicétis in die illa: Confitémini Dómino, et invocáte nomen ejus:

Notas fácite in pópulis adinventiónes ejus: * mementóte quóniam excélsum est nomen ejus.

Cantáte Dómino quóniam magnífice fecit: * annuntiáte hoc in univérsa terra.

Exsúlta, et lauda, habitátio Sion: * quia magnus in médio tui Sanctus Israël.

\antlabel{Ant.} Convérsus est furor tuus, Dómine, et consolátus es me.

\antlabel{Ant.} Laudáte. \textdagger

\psalmheading{Psalmus 116}

\lettrine[lines=2]{L}{AUDATE} \textdagger\ Dóminum, omnes gentes: * \textdaggerdbl\ laudáte eum, omnes pópuli:

Quóniam confirmáta est super nos misericórdia ejus: * et véritas Dómini manet in ætérnum.

\antlabel{Ant.} Laudáte Dóminum, quóniam confirmáta est super nos misericórdia ejus.

\capitulumlabel{Capit.}{Rom. 13, 12-13}

\lettrine[lines=2]{N}{OX} præcéssit, dies autem appropinquávit. Abjiciámus ergo ópera tenebrárum, et induámur arma lucis. Sicut in die honéste ambulémus.

\hymnusheading{Hymnus}

\lettrine[lines=2]{S}{PLENDOR} Patérnæ glóriæ,\\
De luce lucem próferens,\\
Lux lucis, et fons lúminis,\\
Diem dies illúminans:

Verúsque sol illábere,\\
Micans nitóre pérpeti:\\
Jubárque Sancti Spíritus\\
Infúnde nostris sénsibus.

Votis vocémus et Patrem,\\
Patrem poténtis grátiæ,\\
Patrem perénnis glóriæ:\\
Culpam reléget lúbricam.

Confírmet actus strénuos:\\
Dentes retúndat ínvidi:\\
Casus secúndet ásperos:\\
Agénda recte dírigat.

Mentem gubérnet et regat:\\
Sit pura nobis cástitas:\\
Fides calóre férveat,\\
Fraudis venéna nésciat.

Christúsque nobis sit cibus,\\
Potúsque noster sit fides:\\
Læti bibámus sóbriam\\
Profusiónem Spíritus.

Lætus dies hic tránseat:\\
Pudor sit ut dilúculum:\\
Fides velut merídies:\\
Crepúsculum mens nésciat.

Auróra lucem próvehit,\\
Cum luce nobis pródeat\\
In Patre totus Fílius,\\
Et totus in Verbo Pater.

Deo Patri sit glória,\\
Ejúsque soli Fílio,\\
Cum Spíritu Paráclito,\\
Nunc et per omne sæculum. Amen.

\V\ Repléti sumus mane misericórdia tua. \R\ Exsultávimus, et delectáti sumus.

\antlabel{Ad Bened. Ant.} Benedíctus * \textdagger\ Dóminus, Deus Israël, quia visitávit et liberávit nos.

\variantheading{Preces feriales}

KYRIE, eléison. Christe, eléison. Kýrie, eléison.

Pater noster \rubric{quod dícitur a solo Hebdomadário totum clara voce, usque ad} \V\ Et ne nos indúcas in tentatiónem. \R\ Sed líbera nos a malo. \V\ Ego dixi: Dómine, miserére mei. \R\ Sana ánimam meam, quia peccávi tibi. \V\ Convértere, Dómine, úsquequo? \R\ Et deprecábilis esto super servos tuos. \V\ Fiat misericórdia tua, Dómine, super nos. \R\ Quemádmodum sperávimus in te. \V\ Sacerdótes tui induántur justítiam. \R\ Et sancti tui exsúltent. \V\ Orémus pro beatíssimo Papa nostro N. \R\ Dóminus consérvet eum, et vivíficet eum, et beátum fáciat eum in terra, et non tradat eum in ánimam inimicórum ejus. \V\ Orémus et pro Antístite nostro N. \R\ Stet et pascat in fortitúdine tua, Dómine, in sublimitáte nóminis tui. \V\ Dómine, salvum fac regem. \R\ Et exáudi nos in die, qua invocavérimus te. \V\ Salvum fac pópulum tuum, Dómine, et bénedic hereditáti tuæ. \R\ Et rege eos, et extólle illos usque in ætérnum. \V\ Meménto Congregatiónis tuæ. \R\ Quam possedísti ab inítio. \V\ Fiat pax in virtúte tua. \R\ Et abundántia in túrribus tuis. \V\ Orémus pro benefactóribus nostris. \R\ Retribúere dignáre, Dómine, ómnibus, nobis bona faciéntibus propter nomen tuum, vitam ætérnam. Amen. \V\ Orémus pro fidélibus defúnctis. \R\ Réquiem ætérnam dona eis, Dómine, et lux perpétua lúceat eis. \V\ Requiéscant in pace. \R\ Amen. \V\ Pro frátribus nostris abséntibus. \R\ Salvos fac servos tuos, Deus meus, sperántes in te. \V\ Pro afflíctis et captívis. \R\ Líbera eos, Deus Israël, ex ómnibus tribulatiónibus suis. \V\ Mitte eis, Dómine, auxílium de sancto. \R\ Et de Sion tuére eos. \V\ Dómine, Deus virtútum, convérte nos. \R\ Et osténde fáciem tuam, et salvi érimus. \V\ Exsúrge, Christe, ádjuva nos. \R\ Et líbera nos propter nomen tuum. \V\ Dómine, exáudi oratiónem meam. \R\ Et clamor meus ad te véniat.

\rubric{Réliqua ut in Ordinário. 9}

% ===== END OF BATCH 10 -- Covers autumnalis.md lines 3270-3672, the
% two alternate ferial Lauds schemes, "Ad Laudes I" (for ordinary
% ferias and feasts) and "Ad Laudes II" (for common Vigils when the
% Office of the Feria is said). Both share the same five psalms
% (46, 5, 28, Canticum, 116) with different antiphon sets, the same
% Capitulum "Nox præcéssit," and the identical hymn "Splendor
% Paternæ glóriæ" -- all transcribed fresh in both places per the
% governing rule, even though the Hymnus and Capitulum text is
% word-for-word identical between the two schemes in autumnalis.md
% itself. Ad Laudes II additionally carries the "Preces feriales"
% litany, transcribed in full; "Kyrie, eléison" at its head is left
% as plain text with no lettrine, matching precedent in verna.tex
% and hiemalis.tex, since it is not one of the five lettrine genres.
% The document continues with "Ad Primam" (autumnalis.md line 3674).
% =====
\hourheading{Ad Primam}

Pater, Ave et Credo.

\V\ Deus, in adjutórium.

\hymnusheading{Hymnus}

\lettrine[lines=2]{J}{AM} lucis orto sídere,\\
Deum precémur súpplices,\\
Ut in diúrnis áctibus\\
Nos servet a nocéntibus.

Linguam refrénans témperet,\\
Ne litis horror ínsonet:\\
Visum fovéndo cóntegat,\\
Ne vanitátes háuriat.

Sint pura cordis íntima,\\
Absístat et vecórdia:\\
Carnis terat supérbiam\\
Potus cibíque párcitas.

Ut, cum dies abscésserit\\
Noctémque sors redúxerit,\\
Mundi per abstinéntiam\\
Ipsi canámus glóriam.

Deo Patri sit glória,\\
Ejúsque soli Fílio,\\
Cum Spíritu Paráclito,\\
Nunc et per omne sæculum.\\
Amen.

\antlabel{Ant.} Innocens mánibus.

\rubric{¶ Antiphonæ, in hac et áliis Fériis ad Primam, Tértiam, Sextam et Nonam apposítæ, dicúntur semper in Officio feriáli et in Festis sive trium sive novem Lectiónum, quando juxta Rúbricas suméndi sint Psalmi de Féria.}

\psalmheading{Psalmus 23}

\lettrine[lines=2]{D}{OMINI} est terra, et plenitúdo ejus: * orbis terrárum, et univérsi qui hábitant in eo.

Quia ipse super mária fundávit eum: * et super flúmina præparávit eum.

Quis ascéndet in montem Dómini? * aut quis stabit in loco sancto ejus?

Innocens mánibus et mundo corde, * qui non accépit in vano ánimam suam, nec jurávit in dolo próximo suo.

Hic accípiet benedictiónem a Dómino: * et misericórdiam a Deo, salutári suo.

Hæc est generátio quæréntium eum, * quæréntium fáciem Dei Jacob.

Attóllite portas, príncipes, vestras, et elevámini, portæ æternáles: * et introíbit Rex glóriæ.

Quis est iste Rex glóriæ? * Dóminus fortis et potens: Dóminus potens in prœlio.

Attóllite portas, príncipes, vestras, et elevámini, portæ æternáles: * et introíbit Rex glóriæ.

Quis est iste Rex glóriæ? * Dóminus virtútum ipse est Rex glóriæ.

\psalmheading{Psalmus 18, i}

\lettrine[lines=2]{C}{ÆLI} enárrant glóriam Dei, * et ópera mánuum ejus annúntiat firmaméntum.

Dies diéi erúctat verbum, * et nox nocti índicat sciéntiam.

Non sunt loquélæ, neque sermónes, * quorum non audiántur voces eórum.

In omnem terram exívit sonus eórum: * et in fines orbis terræ verba eórum.

In sole pósuit tabernáculum suum: * et ipse tamquam sponsus procédens de thálamo suo:

Exsultávit ut gigas ad curréndam viam, * a summo cælo egréssio ejus:

Et occúrsus ejus usque ad summum ejus: * nec est qui se abscóndat a calóre ejus.

\psalmheading{Psalmus 18, ii}

\lettrine[lines=2]{L}{EX} Dómini immaculáta, convértens ánimas: * testimónium Dómini fidéle, sapiéntiam præstans párvulis.

Justítiæ Dómini rectæ, lætificántes corda: * præcéptum Dómini lúcidum, illúminans óculos.

Timor Dómini sanctus, pérmanens in sæculum sæculi: * judícia Dómini vera, justificáta in semetípsa.

Desiderabília super aurum et lápidem pretiósum multum: * et dulcióra super mel et favum.

Étenim servus tuus custódit ea: * in custodiéndis illis retribútio multa.

Delícta quis intélligit? ab occúltis meis munda me: * et ab aliénis parce servo tuo.

Si mei non fúerint domináti, tunc immaculátus ero: * et emundábor a delícto máximo.

Et erunt ut compláceant elóquia oris mei: * et meditátio cordis mei in conspéctu tuo semper.

Dómine, adjútor meus, * et redémptor meus.

\rubric{¶ Quando ad Laudes dictus fúerit Ps. 50 Miserére, hic subjúngitur sequens Ps. 46, qui in Laudibus locum non hábuit.}

\psalmheading{Psalmus 46}

\lettrine[lines=2]{O}{MNES} gentes, pláudite mánibus: * jubiláte Deo in voce exsultatiónis.

Quóniam Dóminus excélsus, terríbilis: * Rex magnus super omnem terram.

Subjécit pópulos nobis: * et gentes sub pédibus nostris.

Elégit nobis hereditátem suam: * spéciem Jacob, quam diléxit.

Ascéndit Deus in júbilo: * et Dóminus in voce tubæ.

Psállite Deo nostro, psállite: * psállite Regi nostro, psállite.

Quóniam Rex omnis terræ Deus: * psállite sapiénter.

Regnábit Deus super gentes: * Deus sedet super sedem sanctam suam.

Príncipes populórum congregáti sunt cum Deo Abraham: * quóniam dii fortes terræ veheménter eleváti sunt.

\rubric{Secus vero, tribus Psalmis explétis, statim dícitur Ant.}

\antlabel{Ant.} Innocens mánibus et mundo corde ascéndet in montem Dómini.

\rubric{In Officio Festi vel Octávæ:}

\capitulumlabel{Capitulum}{1 Tim. 1, 17}

\lettrine[lines=2]{R}{EGI} sæculórum immortáli et invisíbili, soli Deo honor et glória in sæcula sæculórum. Amen.

\rubric{In ómnibus autem Fériis et in Vigíliis commúnibus:}

\capitulumlabel{Capitulum}{Zach. 8, 19}

\lettrine[lines=2]{P}{ACEM} et veritátem dilígite, ait Dóminus omnípotens.

\R\ br. Christe, Fili Dei vivi, * Miserére nobis. Christe. \V\ \textit{Qui sedes ad déxteram Patris.} Miserére nobis. Glória Patri. Christe. \V\ Exsúrge, Christe, ádjuva nos. \R\ Et líbera nos propter nomen tuum.

\variantheading{Preces dominicales et feriales (si dicendæ sint)}

KYRIE, eléison. Christe, eléison. Kýrie, eléison.

Pater noster \rubric{secreto usque ad} \V\ Et ne nos indúcas in tentatiónem. \R\ Sed líbera nos a malo.

Credo in Deum \rubric{secreto usque ad} \V\ Carnis resurrectiónem. \R\ Vitam ætérnam. Amen. \V\ Et ego ad te, Dómine, clamávi. \R\ Et mane orátio mea prævéniet te. \V\ Repleátur os meum laude. \R\ Ut cantem glóriam tuam, tota die magnitúdinem tuam. \V\ Dómine, avérte fáciem tuam a peccátis meis. \R\ Et omnes iniquitátes meas dele. \V\ Cor mundum crea in me, Deus. \R\ Et spíritum rectum ínnova in viscéribus meis. \V\ Ne projícias me a fácie tua. \R\ Et spíritum sanctum tuum ne áuferas a me. \V\ Redde mihi lætítiam salutáris tui. \R\ Et spíritu principáli confírma me.

\variantheading{In Précibus feriálibus tantum}

\V\ Eripe me, Dómine, ab hómine malo. \R\ A viro iníquo éripe me. \V\ Eripe me de inimícis meis, Deus meus. \R\ Et ab insurgéntibus in me líbera me. \V\ Eripe me de operántibus iniquitátem. \R\ Et de viris sánguinum salva me. \V\ Sic psalmum dicam nómini tuo in sæculum sæculi. \R\ Ut reddam vota mea de die in diem. \V\ Exáudi nos, Deus, salutáris noster. \R\ Spes ómnium fínium terræ, et in mari longe. \V\ Deus, in adjutórium meum inténde. \R\ Dómine, ad adjuvándum me festína. \V\ Sanctus Deus, Sanctus fortis, Sanctus immortális. \R\ Miserére nobis. \V\ Bénedic, ánima mea, Dómino. \R\ Et ómnia, quæ intra me sunt, nómini sancto ejus. \V\ Bénedic, ánima mea, Dómino. \R\ Et noli oblivísci omnes retributiónes ejus. \V\ Qui propitiátur ómnibus iniquitátibus tuis. \R\ Qui sanat omnes infirmitátes tuas. \V\ Qui rédimit de intéritu vitam tuam. \R\ Qui corónat te in misericórdia et miseratiónibus. \V\ Qui replet in bonis desidérium tuum. \R\ Renovábitur ut áquilæ juvéntus tua.

\variantheading{In utrisque Précibus}

\V\ Adjutórium nostrum in nómine Dómini. \R\ Qui fecit cælum et terram.

\formulalabel{Confíteor. Misereátur. Indulgéntiam.}

\V\ Dignáre, Dómine, die isto. \R\ Sine peccáto nos custodíre. \V\ Miserére nostri, Dómine. \R\ Miserére nostri. \V\ Fiat misericórdia tua, Dómine, super nos. \R\ Quemádmodum sperávimus in te. \V\ Dómine, exáudi oratiónem meam. \R\ Et clamor meus ad te véniat.

\V\ Dóminus vobíscum.

Orémus. \oratiolabel{Oratio}

\lettrine[lines=2]{D}{OMINE}, Deus omnípotens, qui ad princípium hujus diéi nos perveníre fecísti: tua nos hódie salva virtúte; ut in hac die ad nullum declinémus peccátum, sed semper ad tuam justítiam faciéndam nostra procédant elóquia, dirigántur cogitatiónes et ópera. Per Dóminum. \R\ Amen.

\V\ Dóminus vobíscum. \V\ Benedicámus Dómino.

\rubric{In Choro Martyrológium.}

\V\ Pretiósa in conspéctu Dómini. \R\ Mors Sanctórum ejus.

\rubric{Et absolúte, sine} Orémus:

\oratiolabel{Oratio}

\lettrine[lines=2]{S}{ANCTA} María, et omnes Sancti intercédant pro nobis ad Dóminum, ut nos mereámur ab eo adjuvári et salvári, qui vivit et regnat in sæcula sæculórum. \R\ Amen.

\V\ Deus, in adjutórium meum inténde. \R\ Dómine, ad adjuvándum me festína. \rubric{Et dícitur ter; último additur:} Glória Patri. Sicut erat.

Kýrie, eléison. Christe, eléison. Kýrie, eléison.

Pater noster \rubric{secreto usque ad} \V\ Et ne nos indúcas in tentatiónem. \R\ Sed líbera nos a malo. \V\ Réspice in servos tuos, Dómine, et in ópera tua, et dírige fílios eórum. \R\ Et sit splendor Dómini, Dei nostri, super nos, et ópera mánuum nostrárum dírige super nos, et opus mánuum nostrárum dírige. \V\ Glória Patri. \R\ Sicut erat.

Orémus. \oratiolabel{Oratio}

\lettrine[lines=2]{D}{IRIGERE} et sanctificáre, régere et gubernáre dignáre, Dómine Deus, Rex cæli et terræ, hódie corda et córpora nostra, sensus, sermónes et actus nostros in lege tua, et in opéribus mandatórum tuórum: ut hic, et in ætérnum, te auxiliánte, salvi et líberi esse mereámur, Salvátor mundi: Qui vivis et regnas in sæcula sæculórum. \R\ Amen.

\V\ Jube, domne.

\benedlabel{Bened.} Dies et actus nostros in sua pace dispónat Dóminus omnípotens. \R\ Amen.

\variantheading{In Officio feriali}

\lettrine[lines=2]{D}{OMINUS} autem dírigat corda et córpora nostra in caritáte Dei, et patiéntia Christi. Tu autem.

\variantheading{E Proprio Sanctorum}

\variantheading{Pro Oct. Omn. Sanctorum}

\lettrine[lines=2]{B}{ENEDÍCTIO} et cláritas et sapiéntia et gratiárum áctio, honor, virtus et fortitúdo Deo nostro in sǽcula sæculórum. Amen.

\variantheading{E Communi Sanctorum}

\variantheading{Pro Apostolis}

\lettrine[lines=2]{I}{BANT} Apóstoli gaudéntes a conspéctu concílii, quóniam digni hábiti sunt pro nómine Jesu contuméliam pati.

\variantheading{Pro Martyre}

\lettrine[lines=2]{J}{USTUS} cor suum trádidit ad vigilándum dilúculo ad Dóminum, qui fecit illum, et in conspéctu Altíssimi deprecábitur.

\variantheading{Pro pluribus Martyribus}

\lettrine[lines=2]{F}{ULGEBUNT} justi, et tamquam scintíllæ in arundinéto discúrrent. Judicábunt natiónes, et dominabúntur pópulis: et regnábit Dóminus illórum in perpétuum.

\variantheading{Pro Conf. Pont.}

\lettrine[lines=2]{F}{UNGI} sacerdótio, et habére laudem in nómine ipsíus, et offérre illi incénsum dignum in odórem suavitátis.

\variantheading{Pro Conf. non Pont.}

\lettrine[lines=2]{J}{USTUM} dedúxit Dóminus per vias rectas, et osténdit illi regnum Dei, et dedit illi sciéntiam sanctórum; honestávit illum in labóribus, et complévit labóres illíus.

\variantheading{Pro Virgine Martyre}

\lettrine[lines=2]{D}{OMINE}, Deus meus, exaltásti super terram habitatiónem meam, et pro morte defluénte deprecáta sum. Invocávi Dóminum, Patrem Dómini mei, ut non derelínquat me in die tribulatiónis meæ, et in témpore superbórum sine adjutório.

\variantheading{Pro Virgine non Martyre}

\lettrine[lines=2]{O}{} \textsc{quam} pulchra est casta generátio cum claritáte! immortális est enim memória illíus: quóniam apud Deum nota est, et apud hómines.

\variantheading{Pro Martyre tantum}

\lettrine[lines=2]{L}{AUDABIT} usque ad mortem ánima mea Dóminum, quóniam éruis sustinéntes te, et líberas eos de manu angústiæ, Dómine, Deus noster.

\variantheading{Pro nec Virg. nec Mart.}

\lettrine[lines=2]{M}{ULTAE} fíliæ congregavérunt divítias: tu supergréssa es univérsas. Fallax grátia, et vana est pulchritúdo: múlier timens Dóminum, ipsa laudábitur.

\variantheading{Pro Dedicatione Ecclesiæ}

\lettrine[lines=2]{E}{T} abstérget Deus omnem lácrimam ab óculis eórum: et mors ultra non erit, neque luctus, neque clamor, neque dolor erit ultra, quia prima abiérunt. Et dixit qui sedébat in throno: Ecce, nova fácio ómnia.

\variantheading{In Officiis B. Mariæ V.}

\lettrine[lines=2]{I}{N} platéis sicut cinnamómum et bálsamum aromatízans odórem dedi: quasi myrrha elécta dedi suavitátem odóris.

\V\ Adjutórium nostrum in nómine Dómini. \R\ Qui fecit cælum et terram.

\V\ Benedícite. \R\ Deus.

\benedictiolabel{Benedictio} Dóminus nos benedícat, et ab omni malo deféndat, et ad vitam perdúcat ætérnam. Et fidélium ánimæ per misericórdiam Dei requiéscant in pace. \R\ Amen.

\rubric{Deinde dícitur tantum} Pater noster \rubric{secréto, nisi sequátur ália Hora.}

% ===== END OF BATCH 11 -- Covers autumnalis.md lines 3674-3927,
% completing "Ad Primam" for Feria Secunda in full: the hymn "Jam
% lucis orto sídere" (with its own stray OCR "\*" before the final
% doxology stanza silently dropped, the same recurring artifact
% corrected in every prior occurrence this volume), Psalms 23, 18/i,
% 18/ii, and the conditional insertion of Psalm 46 (used when Ps. 50
% was said at Lauds instead, per the rubric), the two alternate
% Capitula ("In Officio Festi vel Octávæ" / "In ómnibus autem
% Fériis"), the brief Responsory "Christe, Fili Dei vivi," the
% "Preces dominicales et feriales (si dicendæ sint)" apparatus with
% its "In Précibus feriálibus tantum" and "In utrísque Précibus"
% subdivisions (the bold "Confíteor. Misereátur. Indulgéntiam."
% rendered with \formulalabel{} per the established rule for bold
% source markup), the two closing Orationes and the Martyrologium
% versicle, and finally the full "Lectio brevis" apparatus (the
% ferial variant "Dóminus autem dírigat," followed by "E Proprio
% Sanctorum" with its single "Pro Oct. Omn. Sanctorum" heading (this
% volume's one-item Autumn substitute for aestiva.tex's two-item
% Summer list "Pro Oct. S. Joannis Bapt." / "Pro Invent. S.
% Stephani," verified via direct diff against aestiva.md and
% transcribed fresh from autumnalis.md's own text) and "E Communi
% Sanctorum" with its full run of eleven Pro-headings and "In
% Officiis B. Mariæ V.," all reusing the \variantheading{} command
% established already in this volume's own Sunday Prime, each short
% text given a two-line lettrine as a Lectio-brevis genre text). The
% embedded
% rubric "Et absolúte, sine Orémus:" is preserved exactly as the
% source interleaves red rubric and black directive text. The
% document continues with "Ad Tertiam" (autumnalis.md line 3928).
% =====
\hourheading{Ad Tertiam}

Pater noster et Ave.

\V\ Deus, in adjutórium.

\hymnusheading{Hymnus}

\lettrine[lines=2]{N}{UNC}, Sancte, nobis, Spíritus,\\
Unum Patri cum Fílio,\\
Dignáre promptus íngeri\\
Nostro refúsus péctori.

Os, lingua, mens, sensus, vigor\\
Confessiónem pérsonent,\\
Flamméscat igne cáritas,\\
Accéndat ardor próximos.

Præsta, Pater piíssime,\\
Patríque compar Unice,\\
Cum Spíritu Paráclito\\
Regnans per omne sæculum.\\
Amen.

\antlabel{Ant.} Illuminátio mea.

\psalmheading{Psalmus 26, i}

\lettrine[lines=2]{D}{OMINUS} illuminátio mea, et salus mea, * quem timébo?

Dóminus protéctor vitæ meæ, * a quo trepidábo?

Dum apprópiant super me nocéntes, * ut edant carnes meas:

Qui tríbulant me inimíci mei, * ipsi infirmáti sunt, et cecidérunt.

Si consístant advérsum me castra, * non timébit cor meum.

Si exsúrgat advérsum me prælium, * in hoc ego sperábo.

Unam pétii a Dómino, hanc requíram, * ut inhábitem in domo Dómini ómnibus diébus vitæ meæ:

Ut vídeam voluptátem Dómini, * et vísitem templum ejus.

Quóniam abscóndit me in tabernáculo suo: * in die malórum protéxit me in abscóndito tabernáculi sui.

In petra exaltávit me: * et nunc exaltávit caput meum super inimícos meos.

Circuívi, et immolávi in tabernáculo ejus hóstiam vociferatiónis: * cantábo, et psalmum dicam Dómino.

\psalmheading{Psalmus 26, ii}

\lettrine[lines=2]{E}{XAUDI}, Dómine, vocem meam, qua clamávi ad te: * miserére mei, et exáudi me.

Tibi dixit cor meum, exquisívit te fácies mea: * fáciem tuam, Dómine, requíram.

Ne avértas fáciem tuam a me: * ne declínes in ira a servo tuo.

Adjútor meus esto: * ne derelínquas me, neque despícias me, Deus, salutáris meus.

Quóniam pater meus, et mater mea dereliquérunt me: * Dóminus autem assúmpsit me.

Legem pone mihi, Dómine, in via tua: * et dírige me in sémitam rectam propter inimícos meos.

Ne tradíderis me in ánimas tribulántium me: * quóniam insurrexérunt in me testes iníqui, et mentíta est iníquitas sibi.

Credo vidére bona Dómini * in terra vivéntium.

Exspécta Dóminum, viríliter age: * et confortétur cor tuum, et sústine Dóminum.

\psalmheading{Psalmus 27}

\lettrine[lines=2]{A}{D} te, Dómine, clamábo, Deus meus, ne síleas a me: * nequándo táceas a me, et assimilábor descendéntibus in lacum.

Exáudi, Dómine, vocem deprecatiónis meæ dum oro ad te: * dum extóllo manus meas ad templum sanctum tuum.

Ne simul trahas me cum peccatóribus: * et cum operántibus iniquitátem ne perdas me.

Qui loquúntur pacem cum próximo suo: * mala autem in córdibus eórum.

Da illis secúndum ópera eórum, * et secúndum nequítiam adinventiónum ipsórum.

Secúndum ópera mánuum eórum tríbue illis: * redde retributiónem eórum ipsis.

Quóniam non intellexérunt ópera Dómini, et in ópera mánuum ejus * déstrues illos, et non ædificábis eos.

Benedíctus Dóminus: * quóniam exaudívit vocem deprecatiónis meæ.

Dóminus adjútor meus, et protéctor meus: * in ipso sperávit cor meum, et adjútus sum.

Et reflóruit caro mea: * et ex voluntáte mea confitébor ei.

Dóminus fortitúdo plebis suæ: * et protéctor salvatiónum Christi sui est.

Salvum fac pópulum tuum, Dómine, et bénedic hereditáti tuæ: * et rege eos, et extólle illos usque in ætérnum.

\antlabel{Ant.} Illuminátio mea, et salus mea Dóminus.

\variantheading{In Officio feriali}

\capitulumlabel{Capitulum}{Jer. 17, 14}

\lettrine[lines=2]{S}{ANA} me, Dómine, et sanábor: salvum me fac, et salvus ero: quóniam laus mea tu es.

\R\ br. Sana ánimam meam, * Quia peccávi tibi. Sana. \V\ Ego dixi: Dómine, miserére mei. Quia. Glória Patri. Sana. \V\ Adjútor meus esto, ne derelínquas me. \R\ Neque despícias me, Deus, salutáris meus.

\variantheading{In Vigiliis communibus:}

\variantheading{Preces feriales}

KYRIE, eléison. Christe, eléison. Kýrie, eléison.

Pater noster \rubric{secreto usque ad} \V\ Et ne nos indúcas in tentatiónem. \R\ Sed líbera nos a malo. \V\ Dómine, Deus virtútum, convérte nos. \R\ Et osténde fáciem tuam, et salvi érimus. \V\ Exsúrge, Christe, ádjuva nos. \R\ Et líbera nos propter nomen tuum. \V\ Dómine, exáudi oratiónem meam. \R\ Et clamor meus ad te véniat.

\rubric{Réliqua ut in Ordinário. 19}

\hourheading{Ad Sextam}

Pater noster et Ave.

\V\ Deus, in adjutórium.

\hymnusheading{Hymnus}

\lettrine[lines=2]{R}{ECTOR} potens, verax Deus,\\
Qui témperas rerum vices,\\
Splendóre mane illúminas,\\
Et ígnibus merídiem:

Exstíngue flammas lítium,\\
Aufer calórem nóxium,\\
Confer salútem córporum,\\
Verámque pacem córdium.

Præsta, Pater piíssime,\\
Patríque compar Unice,\\
Cum Spíritu Paráclito\\
Regnans per omne sæculum.\\
Amen.

\antlabel{Ant.} In tua justítia.

\psalmheading{Psalmus 30, i}

\lettrine[lines=2]{I}{N} te, Dómine, sperávi, non confúndar in ætérnum: * in justítia tua líbera me.

Inclína ad me aurem tuam, * accélera ut éruas me.

Esto mihi in Deum protectórem, et in domum refúgii: * ut salvum me fácias.

Quóniam fortitúdo mea, et refúgium meum es tu: * et propter nomen tuum dedúces me, et enútries me.

Edúces me de láqueo hoc, quem abscondérunt mihi: * quóniam tu es protéctor meus.

In manus tuas comméndo spíritum meum: * redemísti me, Dómine, Deus veritátis.

Odísti observántes vanitátes, * supervácue.

Ego autem in Dómino sperávi: * exsultábo, et lætábor in misericórdia tua.

Quóniam respexísti humilitátem meam, * salvásti de necessitátibus ánimam meam.

Nec conclusísti me in mánibus inimíci: * statuísti in loco spatióso pedes meos.

\psalmheading{Psalmus 30, ii}

\lettrine[lines=2]{M}{ISERERE} mei, Dómine, quóniam tríbulor: * conturbátus est in ira óculus meus, ánima mea, et venter meus:

Quóniam defécit in dolóre vita mea: * et anni mei in gemítibus.

Infirmáta est in paupertáte virtus mea: * et ossa mea conturbáta sunt.

Super omnes inimícos meos factus sum oppróbrium et vicínis meis valde: * et timor notis meis.

Qui vidébant me, foras fugérunt a me: * oblivióni datus sum, tamquam mórtuus a corde.

Factus sum tamquam vas pérditum: * quóniam audívi vituperatiónem multórum commorántium in circúitu.

In eo dum convenírent simul advérsum me, * accípere ánimam meam consiliáti sunt.

Ego autem in te sperávi, Dómine: * dixi: Deus meus es tu: in mánibus tuis sortes meæ.

Eripe me de manu inimicórum meórum, * et a persequéntibus me.

Illústra fáciem tuam super servum tuum, salvum me fac in misericórdia tua: * Dómine, non confúndar, quóniam invocávi te.

Erubéscant ímpii, et deducántur in inférnum: * muta fiant lábia dolósa.

Quæ loquúntur advérsus justum iniquitátem, * in supérbia, et in abusióne.

\psalmheading{Psalmus 30, iii}

\lettrine[lines=2]{Q}{UAM} magna multitúdo dulcédinis tuæ, Dómine, * quam abscondísti timéntibus te.

Perfecísti eis, qui sperant in te, * in conspéctu filiórum hóminum.

Abscóndes eos in abscóndito faciéi tuæ, * a conturbatióne hóminum.

Próteges eos in tabernáculo tuo * a contradictióne linguárum.

Benedíctus Dóminus: * quóniam mirificávit misericórdiam suam mihi in civitáte muníta.

Ego autem dixi in excéssu mentis meæ: * Projéctus sum a fácie oculórum tuórum.

Ideo exaudísti vocem oratiónis meæ, * dum clamárem ad te.

Dilígite Dóminum, omnes sancti ejus: * quóniam veritátem requíret Dóminus, et retríbuet abundánter faciéntibus supérbiam.

Viríliter ágite, et confortétur cor vestrum, * omnes, qui sperátis in Dómino.

\antlabel{Ant.} In tua justítia líbera me, Dómine.

\variantheading{In Officio feriali}

\capitulumlabel{Capitulum}{Rom. 13, 8}

\lettrine[lines=2]{N}{EMINI} quidquam debeátis, nisi ut ínvicem diligátis: qui enim díligit próximum, legem implévit.

\R\ br. Benedícam Dóminum * In omni témpore. Benedícam. \V\ Semper laus ejus in ore meo. In. Glória Patri. Benedícam. \V\ Dóminus regit me, et nihil mihi déerit. \R\ In loco páscuæ ibi me collocávit.

\variantheading{In Vigiliis communibus:}

\variantheading{Preces feriales}

KYRIE, eléison. Christe, eléison. Kýrie, eléison.

Pater noster \rubric{secreto usque ad} \V\ Et ne nos indúcas in tentatiónem. \R\ Sed líbera nos a malo. \V\ Dómine, Deus virtútum, convérte nos. \R\ Et osténde fáciem tuam, et salvi érimus. \V\ Exsúrge, Christe, ádjuva nos. \R\ Et líbera nos propter nomen tuum. \V\ Dómine, exáudi oratiónem meam. \R\ Et clamor meus ad te véniat.

\rubric{Réliqua ut in Ordinário. 20}

\hourheading{Ad Nonam}

Pater noster et Ave.

\V\ Deus, in adjutórium.

\hymnusheading{Hymnus}

\lettrine[lines=2]{R}{ERUM}, Deus, tenax vigor,\\
Immótus in te pérmanens,\\
Lucis diúrnæ témpora\\
Succéssibus detérminans:

Largíre lumen véspere,\\
Quo vita nusquam décidat,\\
Sed prǽmium mortis sacræ\\
Perénnis instet glória.

Præsta, Pater piíssime,\\
Patríque compar Unice,\\
Cum Spíritu Paráclito\\
Regnans per omne sæculum.\\
Amen.

\antlabel{Ant.} Exsultáte, justi.

\psalmheading{Psalmus 31}

\lettrine[lines=2]{B}{EATI} quorum remíssæ sunt iniquitátes: * et quorum tecta sunt peccáta.

Beátus vir, cui non imputávit Dóminus peccátum, * nec est in spíritu ejus dolus.

Quóniam tácui, inveteravérunt ossa mea, * dum clamárem tota die.

Quóniam die ac nocte graváta est super me manus tua: * convérsus sum in ærúmna mea, dum confígitur spina.

Delíctum meum cógnitum tibi feci: * et injustítiam meam non abscóndi.

Dixi: Confitébor advérsum me injustítiam meam Dómino: * et tu remisísti impietátem peccáti mei.

Pro hac orábit ad te omnis sanctus, * in témpore opportúno.

Verúmtamen in dilúvio aquárum multárum, * ad eum non approximábunt.

Tu es refúgium meum a tribulatióne, quæ circúmdedit me: * exsultátio mea, érue me a circumdántibus me.

Intelléctum tibi dabo, et ínstruam te in via hac, qua gradiéris: * firmábo super te óculos meos.

Nolíte fíeri sicut equus et mulus, * quibus non est intelléctus.

In camo et freno maxíllas eórum constrínge, * qui non appróximant ad te.

Multa flagélla peccatóris, * sperántem autem in Dómino misericórdia circúmdabit.

Lætámini in Dómino et exsultáte, justi, * et gloriámini, omnes recti corde.

\psalmheading{Psalmus 32, i}

\lettrine[lines=2]{E}{XSULTATE}, justi, in Dómino: * rectos decet collaudátio.

Confitémini Dómino in cíthara: * in psaltério decem chordárum psállite illi.

Cantáte ei cánticum novum: * bene psállite ei in vociferatióne.

Quia rectum est verbum Dómini, * et ómnia ópera ejus in fide.

Díligit misericórdiam et judícium: * misericórdia Dómini plena est terra.

Verbo Dómini cæli firmáti sunt: * et spíritu oris ejus omnis virtus eórum.

Cóngregans sicut in utre aquas maris: * ponens in thesáuris abýssos.

Tímeat Dóminum omnis terra: * ab eo autem commoveántur omnes inhabitántes orbem.

Quóniam ipse dixit, et facta sunt: * ipse mandávit, et creáta sunt.

Dóminus díssipat consília géntium: * réprobat autem cogitatiónes populórum, et réprobat consília príncipum.

Consílium autem Dómini in ætérnum manet: * cogitatiónes cordis ejus in generatióne et generatiónem.

\psalmheading{Psalmus 32, ii}

\lettrine[lines=2]{B}{EATA} gens, cujus est Dóminus, Deus ejus: * pópulus, quem elégit in hereditátem sibi.

De cælo respéxit Dóminus: * vidit omnes fílios hóminum.

De præparáto habitáculo suo * respéxit super omnes, qui hábitant terram.

Qui finxit sigillátim corda eórum: * qui intélligit ómnia ópera eórum.

Non salvábitur rex per multam virtútem: * et gigas non salvábitur in multitúdine virtútis suæ.

Fallax equus ad salútem: * in abundántia autem virtútis suæ non salvábitur.

Ecce óculi Dómini super metuéntes eum: * et in eis, qui sperant super misericórdia ejus:

Ut éruat a morte ánimas eórum: * et alat eos in fame.

Anima nostra sústinet Dóminum: * quóniam adjútor et protéctor noster est.

Quia in eo lætábitur cor nostrum: * et in nómine sancto ejus sperávimus.

Fiat misericórdia tua, Dómine, super nos: * quemádmodum sperávimus in te.

\antlabel{Ant.} Exsultáte, justi, et gloriámini, omnes recti corde.

\variantheading{In Officio feriali}

\capitulumlabel{Capit.}{1 Petri 1, 17-19}

\lettrine[lines=2]{I}{N} timóre incolátus vestri témpore conversámini: sciéntes quod non corruptibílibus auro vel argénto redémpti estis, sed pretióso sánguine quasi Agni immaculáti Christi.

\R\ br. Rédime me, Dómine, * Et miserére mei. Rédime. \V\ Pes enim meus stetit in via recta. Et miserére mei. Glória Patri. Rédime. \V\ Ab occúltis meis munda me, Dómine. \R\ Et ab aliénis parce servo tuo.

\variantheading{In Vigiliis communibus:}

\variantheading{Preces feriales}

KYRIE, eléison. Christe, eléison. Kýrie, eléison.

Pater noster \rubric{secreto usque ad} \V\ Et ne nos indúcas in tentatiónem. \R\ Sed líbera nos a malo. \V\ Dómine, Deus virtútum, convérte nos. \R\ Et osténde fáciem tuam, et salvi érimus. \V\ Exsúrge, Christe, ádjuva nos. \R\ Et líbera nos propter nomen tuum. \V\ Dómine, exáudi oratiónem meam. \R\ Et clamor meus ad te véniat.

\rubric{Réliqua ut in Ordinário. 22}

% ===== END OF BATCH 12 -- Covers autumnalis.md lines 3928-4286,
% completing the Little Hours of Feria Secunda in full: "Ad
% Tertiam" (Psalms 26/i, 26/ii, 27), "Ad Sextam" (Psalms 30/i,
% 30/ii, 30/iii), and "Ad Nonam" (Psalms 31, 32/i, 32/ii), each with
% its own hymn, Capitulum, brief Responsory, and "In Vigiliis
% communibus: Preces feriales" apparatus. Three more occurrences of
% the recurring stray OCR "\*" artifact (in the hymns "Nunc, Sancte,
% nobis, Spíritus," "Rector potens, verax Deus," and "Rerum, Deus,
% tenax vigor," each immediately before the final "Præsta, Pater
% piíssime" doxology stanza) are silently dropped, consistent with
% every prior occurrence in this volume. The document continues
% with "Ad Vesperas" (autumnalis.md line 4288).
% =====
\hourheading{Ad Vesperas}

Pater noster et Ave.

\V\ Deus, in adjutórium.

\antlabel{Ant.} Inclinávit Dóminus * aurem suam mihi.

\rubric{¶ Antíphonæ, in hac et áliis Fériis apposítæ, dicúntur semper in Officio feriáli et in Festis sive trium sive novem Lectiónum, quando juxta Rúbricas suméndi sint Psalmi de Féria, ut in Psaltério.}

\psalmheading{Psalmus 114}

\lettrine[lines=2]{D}{ILEXI}, quóniam exáudiet Dóminus * vocem oratiónis meæ.

Quia inclinávit aurem suam mihi: * et in diébus meis invocábo.

Circumdedérunt me dolóres mortis: * et perícula inférni invenérunt me.

Tribulatiónem et dolórem invéni: * et nomen Dómini invocávi.

O Dómine, líbera ánimam meam: * miséricors Dóminus, et justus, et Deus noster miserétur.

Custódiens párvulos Dóminus: * humiliátus sum, et liberávit me.

Convértere, ánima mea, in réquiem tuam: * quia Dóminus benefécit tibi.

Quia erípuit ánimam meam de morte: * óculos meos a lácrimis, pedes meos a lapsu.

Placébo Dómino * in regióne vivórum.

\antlabel{Ant.} Inclinávit Dóminus aurem suam mihi.

\antlabel{Ant.} Vota mea * Dómino reddam coram omni pópulo ejus.

\psalmheading{Psalmus 115}

\lettrine[lines=2]{C}{REDIDI}, propter quod locútus sum: * ego autem humiliátus sum nimis.

Ego dixi in excéssu meo: * Omnis homo mendax.

Quid retríbuam Dómino, * pro ómnibus, quæ retríbuit mihi?

Cálicem salutáris accípiam: * et nomen Dómini invocábo.

Vota mea Dómino reddam coram omni pópulo ejus: * pretiósa in conspéctu Dómini mors sanctórum ejus:

O Dómine, quia ego servus tuus: * ego servus tuus, et fílius ancíllæ tuæ.

Dirupísti víncula mea: * tibi sacrificábo hóstiam laudis, et nomen Dómini invocábo.

Vota mea Dómino reddam in conspéctu omnis pópuli ejus: * in átriis domus Dómini, in médio tui, Jerúsalem.

\antlabel{Ant.} Vota mea Dómino reddam coram omni pópulo ejus.

\antlabel{Ant.} Clamávi, * et Dóminus exaudívit me.

\psalmheading{Psalmus 119}

\lettrine[lines=2]{A}{D} Dóminum cum tribulárer clamávi: * et exaudívit me.

Dómine, líbera ánimam meam a lábiis iníquis, * et a lingua dolósa.

Quid detur tibi, aut quid apponátur tibi * ad linguam dolósam?

Sagíttæ poténtis acútæ, * cum carbónibus desolatóriis.

Heu mihi, quia incolátus meus prolongátus est: * habitávi cum habitántibus Cedar: multum íncola fuit ánima mea.

Cum his, qui odérunt pacem, eram pacíficus: * cum loquébar illis, impugnábant me gratis.

\antlabel{Ant.} Clamávi, et Dóminus exaudívit me.

\antlabel{Ant.} Auxílium meum * a Dómino qui fecit cælum et terram.

\psalmheading{Psalmus 120}

\lettrine[lines=2]{L}{EVAVI} óculos meos in montes, * unde véniet auxílium mihi.

Auxílium meum a Dómino, * qui fecit cælum et terram.

Non det in commotiónem pedem tuum: * neque dormítet qui custódit te.

Ecce, non dormitábit neque dórmiet, * qui custódit Israël.

Dóminus custódit te, Dóminus protéctio tua, * super mánum déxteram tuam.

Per diem sol non uret te: * neque luna per noctem.

Dóminus custódit te ab omni malo: * custódiat ánimam tuam Dóminus.

Dóminus custódiat intróitum tuum, et éxitum tuum: * ex hoc nunc, et usque in sæculum.

\antlabel{Ant.} Auxílium meum a Dómino, qui fecit cælum et terram.

\antlabel{Ant.} Lætátus sum * \textdagger\ in his, quæ dicta sunt mihi. \textdaggerdbl

\psalmheading{Psalmus 121}

\lettrine[lines=2]{L}{ÆTATUS} sum \textdagger\ in his, quæ dicta sunt mihi: * \textdaggerdbl\ In domum Dómini íbimus.

Stantes erant pedes nostri, * in átriis tuis, Jerúsalem.

Jerúsalem, quæ ædificátur ut cívitas: * cujus participátio ejus in idípsum.

Illuc enim ascendérunt tribus, tribus Dómini: * testimónium Israël ad confiténdum nómini Dómini.

Quia illic sedérunt sedes in judício, * sedes super domum David.

Rogáte quæ ad pacem sunt Jerúsalem: * et abundántia diligéntibus te:

Fiat pax in virtúte tua: * et abundántia in túrribus tuis.

Propter fratres meos, et próximos meos, * loquébar pacem de te:

Propter domum Dómini, Dei nostri, * quæsívi bona tibi.

\antlabel{Ant.} Lætátus sum in his, quæ dicta sunt mihi.

\variantheading{In Officio feriali}

\rubric{Capítulum, Hymnus, Versus, et Ant. ad Magníficat dicúntur quotídie de Féria curránti, ut in Psaltério hic et álibi appónuntur.}

\capitulumlabel{Capitulum}{2 Cor. 1, 3-4}

\lettrine[lines=2]{B}{ENEDICTUS} Deus, et Pater Dómini nostri Jesu Christi, Pater misericordiárum, et Deus totíus consolatiónis, qui consolátur nos in omni tribulatióne nostra.

\hymnusheading{Hymnus}

\lettrine[lines=2]{I}{MMENSE} cæli Cónditor,\\
Qui mixta ne confúnderent,\\
Aquæ fluénta dívidens,\\
Cælum dedísti límitem.

Firmans locum cæléstibus,\\
Simúlque terræ rívulis;\\
Ut unda flammas témperet,\\
Terræ solum ne díssipent.

Infúnde nunc, piíssime,\\
Donum perénnis grátiæ:\\
Fraudis novæ ne cásibus\\
Nos error átterat vetus.

Lucem fides adáugeat:\\
Sic lúminis jubar ferat:\\
Hæc vana cuncta próterat:\\
Hanc falsa nulla cómprimant.

Præsta, Pater piíssime,\\
Patríque compar Unice,\\
Cum Spíritu Paráclito\\
Regnans per omne sæculum.\\
Amen.

\V\ Dirigátur, Dómine, orátio mea. \R\ Sicut incénsum in conspéctu tuo.

\antlabel{Ad Magnif. Ant.} Magníficat * \textdagger\ ánima mea Dóminum, quia respéxit Deus humilitátem meam.

\rubric{Réliqua ut in Ordinário. 22}

\hourheading{Ad Completorium}

\V\ Jube, domne.

\benedlabel{Bened.} Noctem quiétam, et finem perféctum concédat nobis Dóminus omnípotens. \R\ Amen.

\lectioheading{Lectio br.}{1 Petri 5, 8-9}

\lettrine[lines=2]{F}{RATRES}: Sóbrii estóte, et vigiláte: quia adversárius vester diábolus tamquam leo rúgiens círcuit, quærens quem dévoret: cui resístite fortes in fide. Tu autem, Dómine, miserére nobis. \R\ Deo grátias.

\V\ Adjutórium nostrum in nómine Dómini. \R\ Qui fecit cælum et terram.

Pater noster \rubric{quod dícitur totum secréto.}

\formulalabel{Confíteor.} Misereátur. Indulgéntiam.

\V\ Convérte nos, Deus, salutáris noster. \R\ Et avérte iram tuam a nobis.

\V\ Deus, in adjutórium.

\antlabel{Ant.} Salvum me fac.

\rubric{¶ Antíphonæ, in hac et áliis Fériis apposítæ, dicúntur semper ad Completórium tum de Témpore tum de Sanctis, quando juxta Rúbricas suméndi sint Psalmi de Féria.}

\psalmheading{Psalmus 6}

\lettrine[lines=2]{D}{OMINE}, ne in furóre tuo árguas me, * neque in ira tua corrípias me.

Miserére mei, Dómine, quóniam infírmus sum: * sana me, Dómine, quóniam conturbáta sunt ossa mea.

Et ánima mea turbáta est valde: * sed tu, Dómine, úsquequo?

Convértere, Dómine, et éripe ánimam meam: * salvum me fac propter misericórdiam tuam.

Quóniam non est in morte qui memor sit tui: * in inférno autem quis confitébitur tibi?

Laborávi in gémitu meo, lavábo per síngulas noctes lectum meum: * lácrimis meis stratum meum rigábo.

Turbátus est a furóre óculus meus: * inveterávi inter omnes inimícos meos.

Discédite a me, omnes, qui operámini iniquitátem: * quóniam exaudívit Dóminus vocem fletus mei.

Exaudívit Dóminus deprecatiónem meam, * Dóminus oratiónem meam suscépit.

Erubéscant, et conturbéntur veheménter omnes inimíci mei: * convertántur et erubéscant valde velóciter.

\psalmheading{Psalmus 7, i}

\lettrine[lines=2]{D}{OMINE}, Deus meus, in te sperávi: * salvum me fac ex ómnibus persequéntibus me, et líbera me.

Nequándo rápiat ut leo ánimam meam, * dum non est qui rédimat, neque qui salvum fáciat.

Dómine, Deus meus, si feci istud, * si est iníquitas in mánibus meis:

Si réddidi retribuéntibus mihi mala, * décidam mérito ab inimícis meis inánis.

Persequátur inimícus ánimam meam, et comprehéndat, et concúlcet in terra vitam meam, * et glóriam meam in púlverem dedúcat.

Exsúrge, Dómine, in ira tua: * et exaltáre in fínibus inimicórum meórum.

Et exsúrge, Dómine, Deus meus, in præcépto quod mandásti: * et synagóga populórum circúmdabit te.

Et propter hanc in altum regrédere: * Dóminus júdicat pópulos.

Júdica me, Dómine, secúndum justítiam meam, * et secúndum innocéntiam meam super me.

Consumétur nequítia peccatórum, et díriges justum, * scrutans corda et renes, Deus.

\psalmheading{Psalmus 7, ii}

\lettrine[lines=2]{J}{USTUM} adjutórium meum a Dómino, * qui salvos facit rectos corde.

Deus judex justus, fortis, et pátiens: * numquid iráscitur per síngulos dies?

Nisi convérsi fuéritis, gládium suum vibrábit: * arcum suum teténdit, et parávit illum.

Et in eo parávit vasa mortis: * sagíttas suas ardéntibus effécit.

Ecce, partúriit injustítiam: * concépit dolórem, et péperit iniquitátem.

Lacum apéruit, et effódit eum: * et íncidit in fóveam quam fecit.

Convertétur dolor ejus in caput ejus: * et in vérticem ipsíus iníquitas ejus descéndet.

Confitébor Dómino secúndum justítiam ejus: * et psallam nómini Dómini altíssimi.

\antlabel{Ant.} Salvum me fac, Dómine, propter misericórdiam tuam.

\hymnusheading{Hymnus}

\lettrine[lines=2]{T}{E} lucis ante términum,\\
Rerum Creátor, póscimus,\\
Ut pro tua cleméntia\\
Sis præsul et custódia.

Procul recédant sómnia,\\
Et nóctium phantásmata,\\
Hostémque nostrum cómprime,\\
Ne polluántur córpora.

Præsta, Pater piíssime,\\
Patríque compar Unice,\\
Cum Spíritu Paráclito\\
Regnans per omne sæculum.\\
Amen.

\capitulumlabel{Capitulum}{Jer. 14, 9}

\lettrine[lines=2]{T}{U} autem in nobis es, Dómine, et nomen sanctum tuum invocátum est super nos: ne derelínquas nos, Dómine, Deus noster.

\R\ br. In manus tuas, Dómine, * Comméndo spíritum meum. In manus. \V\ Redemísti nos, Dómine, Deus veritátis. Comméndo. Glória Patri. In manus.

\V\ Custódi nos, Dómine, ut pupíllam óculi. \R\ Sub umbra alárum tuárum prótege nos.

\antlabel{Ant.} Salva nos.

\canticumheading{Canticum Simeonis}{Luc. 2, 29-32}

\lettrine[lines=2]{N}{UNC} dimíttis servum tuum, Dómine, * secúndum verbum tuum in pace:

Quia vidérunt óculi mei * salutáre tuum,

Quod parásti * ante fáciem ómnium populórum,

Lumen ad revelatiónem géntium, * et glóriam plebis tuæ Israël.

Glória Patri.

\antlabel{Ant.} Salva nos, Dómine, vigilántes, custódi nos dormiéntes; ut vigilémus cum Christo, et requiescámus in pace.

\variantheading{Preces (si dicendæ sint)}

KYRIE, eléison. Christe, eléison. Kýrie, eléison.

Pater noster \rubric{secreto usque ad} \V\ Et ne nos indúcas in tentatiónem. \R\ Sed líbera nos a malo.

Credo in Deum \rubric{secreto usque ad} \V\ Carnis resurrectiónem. \R\ Vitam ætérnam. Amen. \V\ Benedíctus es, Dómine, Deus patrum nostrórum. \R\ Et laudábilis et gloriósus in sæcula. \V\ Benedicámus Patrem et Fílium cum Sancto Spíritu. \R\ Laudémus, et superexaltémus eum in sæcula. \V\ Benedíctus es, Dómine, in firmaménto cæli. \R\ Et laudábilis, et gloriósus, et superexaltátus in sæcula. \V\ Benedícat et custódiat nos omnípotens et miséricors Dóminus. \R\ Amen. \V\ Dignáre, Dómine, nocte ista. \R\ Sine peccáto nos custodíre. \V\ Miserére nostri, Dómine. \R\ Miserére nostri. \V\ Fiat misericórdia tua, Dómine, super nos. \R\ Quemádmodum sperávimus in te. \V\ Dómine, exáudi oratiónem meam. \R\ Et clamor meus ad te véniat.

\V\ Dóminus vobíscum.

Orémus. \oratiolabel{Oratio}

\lettrine[lines=2]{V}{ISITA}, quæsumus, Dómine, habitatiónem istam, et omnes insídias inimíci ab ea longe repélle: Angeli tui sancti hábitent in ea, qui nos in pace custódiant; et benedíctio tua sit super nos semper. Per Dóminum nostrum. \R\ Amen.

\V\ Dóminus vobíscum.

\V\ Benedicámus Dómino.

\benedictiolabel{Benedictio} Benedícat et custódiat nos omnípotens et miséricors Dóminus, Pater, et Fílius, et Spíritus Sanctus. \R\ Amen.

\variantheading{Antiphona finalis B. M. V.}

\V\ Divínum auxílium.

Pater, Ave et Credo.

% ===== END OF BATCH 13 -- Covers autumnalis.md lines 4288-4607,
% completing "Ad Vesperas" (Psalms 114, 115, 119, 120, 121, the
% Capitulum, the hymn "Immense cæli Cónditor," and the Magnificat
% antiphon) and "Ad Completorium" (Psalms 6, 7/i, 7/ii, the hymn
% "Te lucis ante términum," the Capitulum, the brief Responsory "In
% manus tuas," the Canticum Simeonis, and the Preces), and with
% them the ENTIRE "Feria Secunda" Office. Two more occurrences of
% the recurring stray OCR "\*" artifact (in "Immense cæli Cónditor"
% at autumnalis.md line 4436 and "Te lucis ante términum" at line 4550,
% both at the hymn's final doxology stanza) are silently dropped,
% consistent with every prior occurrence in this volume. The source
% at line 4464 bolds only "Confíteor." (not the full "Misereátur.
% Indulgéntiam."), unlike the fully-bold form seen at every other
% occurrence transcribed so far in this volume; this is a genuine,
% recurring alternation in autumnalis.md itself (confirmed by grep
% across the file) and is transcribed exactly as marked here:
% \formulalabel{Confíteor.} followed by plain-text "Misereátur.
% Indulgéntiam." The closing "Antiphona finalis B. M. V." is again
% given only as the bare cross-reference incipit "℣. Divínum
% auxílium," matching the identical abbreviated treatment already
% established for Sunday's Compline earlier in this volume. The
% document continues with "Feria Tertia" (autumnalis.md line 4610),
% the second weekday of the Psalterium cycle.
% =====
\feriaheading{Feria Tertia}

\hourheading{Ad Matutinum}

Pater, Ave et Credo.

\V\ Dómine, lábia mea.

\V\ Deus, in adjutórium.

\variantheading{In Officio feriali}

\invitatlabel{Invitat.} Jubilémus Deo, * Salutári nostro. \psalmlabel{Ps. 94} Veníte, exsultémus.

\hymnusheading{Hymnus}

\lettrine[lines=2]{C}{ONSORS} Patérni lúminis,\\
Lux ipse lucis, et dies,\\
Noctem canéndo rúmpimus:\\
Assíste postulántibus.

Aufer tenébras méntium,\\
Fuga catérvas dæmonum,\\
Expélle somnoléntiam,\\
Ne pigritántes óbruat.

Sic, Christe, nobis ómnibus\\
Indúlgeas credéntibus,\\
Ut prosit exorántibus\\
Quod præcinéntes psállimus.

Præsta, Pater piíssime,\\
Patríque compar Unice,\\
Cum Spíritu Paráclito\\
Regnans per omne sæculum.\\
Amen.

\nocturnoheading{In I Nocturno}

\antlabel{Ant.} Expúgna, Dómine, * impugnántes me.

\psalmheading{Psalmus 34, i}

\lettrine[lines=2]{J}{UDICA}, Dómine, nocéntes me, * expúgna impugnántes me.

Apprehénde arma et scutum: * exsúrge in adjutórium mihi.

Effúnde frámeam, et conclúde advérsus eos, qui persequúntur me: * dic ánimæ meæ: Salus tua ego sum.

Confundántur et revereántur, * quæréntes ánimam meam.

Avertántur retrórsum, et confundántur * cogitántes mihi mala.

Fiant tamquam pulvis ante fáciem venti: * et Angelus Dómini coárctans eos.

Fiat via illórum ténebræ et lúbricum: * et Angelus Dómini pérsequens eos.

Quóniam gratis abscondérunt mihi intéritum láquei sui: * supervácue exprobravérunt ánimam meam.

Véniat illi láqueus, quem ignórat: et cáptio, quam abscóndit, apprehéndat eum: * et in láqueum cadat in ipsum.

Anima autem mea exsultábit in Dómino: * et delectábitur super salutári suo.

Omnia ossa mea dicent: * Dómine, quis símilis tibi?

Erípiens ínopem de manu fortiórum ejus: * egénum et páuperem a diripiéntibus eum.

\antlabel{Ant.} Expúgna, Dómine, impugnántes me.

\antlabel{Ant.} Restítue ánimam meam * a malefáctis eórum, Dómine.

\psalmheading{Psalmus 34, ii}

\lettrine[lines=2]{S}{URGENTES} testes iníqui, * quæ ignorábam interrogábant me.

Retribuébant mihi mala pro bonis: * sterilitátem ánimæ meæ.

Ego autem cum mihi molésti essent, * induébar cilício.

Humiliábam in jejúnio ánimam meam: * et orátio mea in sinu meo convertétur.

Quasi próximum, et quasi fratrem nostrum, sic complacébam: * quasi lugens et contristátus, sic humiliábar.

Et advérsum me lætáti sunt, et convenérunt: * congregáta sunt super me flagélla, et ignorávi.

Dissipáti sunt, nec compúncti, tentavérunt me, subsannavérunt me subsannatióne: * frenduérunt super me déntibus suis.

Dómine, quando respícies? * restítue ánimam meam a malignitáte eórum, a leónibus únicam meam.

\antlabel{Ant.} Restítue ánimam meam a malefáctis eórum, Dómine.

\antlabel{Ant.} Exsúrge, Dómine, * et inténde judício meo.

\psalmheading{Psalmus 34, iii}

\lettrine[lines=2]{C}{ONFITEBOR} tibi in ecclésia magna, * in pópulo gravi laudábo te.

Non supergáudeant mihi qui adversántur mihi iníque: * qui odérunt me gratis et ánnuunt óculis.

Quóniam mihi quidem pacífice loquebántur: * et in iracúndia terræ loquéntes, dolos cogitábant.

Et dilatavérunt super me os suum: * dixérunt: Euge, euge, vidérunt óculi nostri.

Vidísti, Dómine, ne síleas: * Dómine, ne discédas a me.

Exsúrge et inténde judício meo: * Deus meus, et Dóminus meus in causam meam.

Júdica me secúndum justítiam tuam, Dómine, Deus meus, * et non supergáudeant mihi.

Non dicant in córdibus suis: Euge, euge, ánimæ nostræ: * nec dicant: Devorávimus eum.

Erubéscant et revereántur simul, * qui gratulántur malis meis.

Induántur confusióne et reveréntia * qui magna loquúntur super me.

Exsúltent et læténtur qui volunt justítiam meam: * et dicant semper: Magnificétur Dóminus, qui volunt pacem servi ejus.

Et lingua mea meditábitur justítiam tuam, * tota die laudem tuam.

\antlabel{Ant.} Exsúrge, Dómine, et inténde judício meo.

\variantheading{In Festis novem Lectionum:}

\V\ Lingua mea meditábitur justítiam tuam. \R\ Tota die laudem tuam, Dómine.

\absolutiolabel{Absol.} Exáudi, Dómine.

\benedlabel{Bened.} 1 Benedictióne perpétua. 2 Unigénitus Dei Fílius. 3 Spíritus Sancti grátia.

\nocturnoheading{In II Nocturno}

\antlabel{Ant.} Noli æmulári * \textdagger\ in eo, qui prosperátur, et facit iniquitátem.

\psalmheading{Psalmus 36, i}

\lettrine[lines=2]{N}{OLI} æmulári in malignántibus: * neque zeláveris faciéntes iniquitátem.

Quóniam tamquam fœnum velóciter aréscent: * et quemádmodum ólera herbárum cito décident.

Spera in Dómino, et fac bonitátem: * et inhábita terram, et pascéris in divítiis ejus.

Delectáre in Dómino: * et dabit tibi petitiónes cordis tui.

Revéla Dómino viam tuam, et spera in eo: * et ipse fáciet.

Et edúcet quasi lumen justítiam tuam: et judícium tuum tamquam merídiem: * súbditus esto Dómino, et ora eum.

Noli æmulári in eo, qui prosperátur in via sua: * in hómine faciénte injustítias.

Désine ab ira, et derelínque furórem: * noli æmulári ut malignéris.

Quóniam qui malignántur, exterminabúntur: * sustinéntes autem Dóminum, ipsi hereditábunt terram.

Et adhuc pusíllum, et non erit peccátor: * et quæres locum ejus, et non inviénies.

Mansuéti autem hereditábunt terram: * et delectabúntur in multitúdine pacis.

Observábit peccátor justum: * et stridébit super eum déntibus suis.

Dóminus autem irridébit eum: * quóniam próspicit quod véniet dies ejus.

Gládium evaginavérunt peccatóres: * intendérunt arcum suum.

Ut dejíciant páuperem et ínopem: * ut trucídent rectos corde.

Gládius eórum intret in corda ipsórum: * et arcus eórum confringátur.

\antlabel{Ant.} Noli æmulári in eo, qui prosperátur, et facit iniquitátem.

\antlabel{Ant.} Bráchia peccatórum * conteréntur: confírmat autem justos Dóminus.

\psalmheading{Psalmus 36, ii}

\lettrine[lines=2]{M}{ELIUS} est módicum justo, * super divítias peccatórum multas.

Quóniam bráchia peccatórum conteréntur: * confírmat autem justos Dóminus.

Novit Dóminus dies immaculatórum: * et heréditas eórum in ætérnum erit.

Non confundéntur in témpore malo, et in diébus famis saturabúntur: * quia peccatóres períbunt.

Inimíci vero Dómini mox ut honorificáti fúerint et exaltáti: * deficiéntes, quemádmodum fumus defícient.

Mutuábitur peccátor, et non solvet: * justus autem miserétur et tríbuet.

Quia benedicéntes ei hereditábunt terram: * maledicéntes autem ei períbunt.

Apud Dóminum gressus hóminis dirigéntur: * et viam ejus volet.

Cum cecíderit, non collidétur: * quia Dóminus suppónit manum suam.

Júnior fui, étenim sénui: * et non vidi justum derelíctum, nec semen ejus quærens panem.

Tota die miserétur et cómmodat: * et semen illíus in benedictióne erit.

Declína a malo, et fac bonum: * et inhábita in sæculum sæculi.

Quia Dóminus amat judícium, et non derelínquet sanctos suos: * in ætérnum conservabúntur.

Injústi puniéntur: * et semen impiórum períbit.

Justi autem hereditábunt terram: * et inhabitábunt in sæculum sæculi super eam.

\antlabel{Ant.} Bráchia peccatórum conteréntur: confírmat autem justos Dóminus.

\antlabel{Ant.} Custódi innocéntiam, * et vide æquitátem.

\psalmheading{Psalmus 36, iii}

\lettrine[lines=2]{O}{S} justi meditábitur sapiéntiam, * et lingua ejus loquétur judícium.

Lex Dei ejus in corde ipsíus: * et non supplantabúntur gressus ejus.

Consíderat peccátor justum: * et quærit mortificáre eum.

Dóminus autem non derelínquet eum in mánibus ejus: * nec damnábit eum, cum judicábitur illi.

Exspécta Dóminum, et custódi viam ejus: et exaltábit te ut hereditáte cápias terram: * cum períerint peccatóres vidébis.

Vidi ímpium superexaltátum, * et elevátum sicut cedros Líbani.

Et transívi, et ecce non erat: * et quæsívi eum, et non est invéntus locus ejus.

Custódi innocéntiam, et vide æquitátem: * quóniam sunt relíquiæ hómini pacífico.

Injústi autem disperíbunt simul: * relíquiæ impiórum interíbunt.

Salus autem justórum a Dómino: * et protéctor eórum in témpore tribulatiónis.

Et adjuvábit eos Dóminus et liberábit eos: * et éruet eos a peccatóribus, et salvábit eos: quia speravérunt in eo.

\antlabel{Ant.} Custódi innocéntiam, et vide æquitátem.

\variantheading{In Festis novem Lectionum:}

\V\ Exspécta Dóminum, et custódi viam ejus. \R\ Exaltábit te ut hereditáte cápias terram.

\absolutiolabel{Absolutio} Ipsíus píetas et.

\benedlabel{Bened.} 4 Deus Pater omnípotens. 5 Christus perpétuæ. 6 Ignem sui amóris.

\nocturnoheading{In III Nocturno}

\antlabel{Ant.} Ne in ira tua * corrípias me, Dómine.

\psalmheading{Psalmus 37, i}

\lettrine[lines=2]{D}{OMINE}, ne in furóre tuo árguas me, * neque in ira tua corrípias me.

Quóniam sagíttæ tuæ infíxæ sunt mihi: * et confirmásti super me manum tuam.

Non est sánitas in carne mea a fácie iræ tuæ: * non est pax óssibus meis a fácie peccatórum meórum.

Quóniam iniquitátes meæ supergréssæ sunt caput meum: * et sicut onus grave gravátæ sunt super me.

Putruérunt et corrúptæ sunt cicatríces meæ, * a fácie insipiéntiæ meæ.

Miser factus sum, et curvátus sum usque in finem: * tota die contristátus ingrediébar.

Quóniam lumbi mei impléti sunt illusiónibus: * et non est sánitas in carne mea.

Afflíctus sum, et humiliátus sum nimis: * rugiébam a gémitu cordis mei.

Dómine, ante te omne desidérium meum: * et gémitus meus a te non est abscónditus.

Cor meum conturbátum est, derelíquit me virtus mea: * et lumen oculórum meórum, et ipsum non est mecum.

\antlabel{Ant.} Ne in ira tua corrípias me, Dómine.

\antlabel{Ant.} Inténde in adjutórium meum, * Dómine, virtus salútis meæ.

\psalmheading{Psalmus 37, ii}

\lettrine[lines=2]{A}{MICI} mei, et próximi mei * advérsum me appropinquavérunt, et stetérunt.

Et qui juxta me erant, de longe stetérunt: * et vim faciébant qui quærébant ánimam meam.

Et qui inquirébant mala mihi, locúti sunt vanitátes: * et dolos tota die meditabántur.

Ego autem tamquam surdus non audiébam: * et sicut mutus non apériens os suum.

Et factus sum sicut homo non áudiens: * et non habens in ore suo redargutiónes.

Quóniam in te, Dómine, sperávi: * tu exáudies me, Dómine, Deus meus.

Quia dixi: Nequándo supergáudeant mihi inimíci mei: * et dum commovéntur pedes mei, super me magna locúti sunt.

Quóniam ego in flagélla parátus sum: * et dolor meus in conspéctu meo semper.

Quóniam iniquitátem meam annuntiábo: * et cogitábo pro peccáto meo.

Inimíci autem mei vivunt, et confirmáti sunt super me: * et multiplicáti sunt qui odérunt me iníque.

Qui retríbuunt mala pro bonis, detrahébant mihi: * quóniam sequébar bonitátem.

Ne derelínquas me, Dómine, Deus meus: * ne discésseris a me.

Inténde in adjutórium meum, * Dómine, Deus salútis meæ.

\antlabel{Ant.} Inténde in adjutórium meum, Dómine, virtus salútis meæ.

\antlabel{Ant.} Amove, Dómine, * a me plagas tuas.

\psalmheading{Psalmus 38}

\lettrine[lines=2]{D}{IXI}: Custódiam vias meas: * ut non delínquam in lingua mea.

Pósui ori meo custódiam, * cum consísteret peccátor advérsum me.

Obmútui, et humiliátus sum, et sílui a bonis: * et dolor meus renovátus est.

Concáluit cor meum intra me: * et in meditatióne mea exardéscet ignis.

Locútus sum in lingua mea: * Notum fac mihi, Dómine, finem meum.

Et númerum diérum meórum quis est: * ut sciam quid desit mihi.

Ecce, mensurábiles posuísti dies meos: * et substántia mea tamquam níhilum ante te.

Verúmtamen univérsa vánitas, * omnis homo vivens.

Verúmtamen in imágine pertránsit homo: * sed et frustra conturbátur.

Thesaurízat: * et ignórat cui congregábit ea.

Et nunc quæ est exspectátio mea? Nonne Dóminus? * Et substántia mea apud te est.

Ab ómnibus iniquitátibus meis érue me: * oppróbrium insipiénti dedísti me.

Obmútui, et non apérui os meum, quóniam tu fecísti: * ámove a me plagas tuas.

A fortitúdine manus tuæ ego defíci in increpatiónibus: * propter iniquitátem corripuísti hóminem.

Et tabéscere fecísti sicut aráneam ánimam ejus: * verúmtamen vane conturbátur omnis homo.

Exáudi oratiónem meam, Dómine, et deprecatiónem meam: * áuribus pércipe lácrimas meas.

Ne síleas: quóniam ádvena ego sum apud te, et peregrínus, * sicut omnes patres mei.

Remítte mihi, ut refrígerer priúsquam ábeam, * et ámplius non ero.

\antlabel{Ant.} Amove, Dómine, a me plagas tuas.

\variantheading{In feriali Officio et in Festis:}

\V\ Deus, ne síleas a me, remítte mihi. \R\ Quóniam íncola ego sum apud te, et peregrínus.

\variantheading{In Officio novem Lectionum}

\absolutiolabel{Absolutio.} A vínculis.

\benedlabel{Bened.} 7 Evangélica léctio. 8 In Festis Dómini: Divínum auxílium. In Festis B. M. V.: Cujus festum cólimus, ipsa Virgo. In Festis Sanctórum: Cujus (\textit{vel} Quorum \textit{aut} Quarum) festum cólimus. 9 Ad societátem. (Per evangélica dicta.)

\variantheading{In Officio trium Lectionum}

\absolutiolabel{Absolutio.} Ipsíus píetas et.

\variantheading{In Feriis, cum Homilia:}

\benedlabel{Bened.} 1 Evangélica. 2 Divínum. 3 Ad societátem.

\variantheading{In Feriis, sine Homilia:}

\benedlabel{Bened.} 1 Deus Pater. 2 Christus perpétuæ. 3 Ignem sui amóris.

\variantheading{In Festis Sanctorum:}

\benedlabel{Bened.} 1 Ille nos benedícat, qui sine fine vivit et regnat. \R\ Amen. 2 Cujus (\textit{vel} Quorum \textit{aut} Quarum). 3 Ad societátem.

% ===== END OF BATCH 14 -- Covers autumnalis.md lines 4610-4978,
% opening "Feria Tertia" (Tuesday) and completing its "Ad
% Matutinum" in full: the ferial Invitatorium "Jubilémus Deo,
% Salutári nostro," the hymn "Consors Patérni lúminis" (with the
% recurring stray OCR asterisk before its final doxology stanza
% silently dropped, consistent with every prior occurrence), and
% the three Nocturns (Psalms 34/i-iii; 36/i-iii; 37/i, 37/ii, 38),
% each with its versicle, Absolutio, and Benedictiones apparatus.
% "Absolutio" is rendered with the exact trailing period the source
% gives it at lines 4959 and 4965 ("Absolutio.") though not at line
% 4847 ("Absolutio"), preserving the source's own inconsistency.
% The compact multi-option "Bened." line for the ninth lesson is
% again transcribed as one continuous argument, matching the
% identical text already seen in Feria Secunda's Matins. The
% document continues with "Ad Laudes I" (autumnalis.md line 4981).
% =====
\hourheading{Ad Laudes I}

\rubric{In Feriis extra Vigílias commúnes occurréntibus, atque in Festis.}

(Pater noster et Ave.)

\V\ Deus, in adjutórium.

\antlabel{Ant.} Cantáte * \textdagger\ Dómino, et benedícite nómini ejus.

\psalmheading{Psalmus 95}

\lettrine[lines=2]{C}{ANTATE} \textdagger\ Dómino cánticum novum: * cantáte Dómino, omnis terra.

Cantáte Dómino, et benedícite nómini ejus: * annuntiáte de die in diem salutáre ejus.

Annuntiáte inter gentes glóriam ejus, * in ómnibus pópulis mirabília ejus.

Quóniam magnus Dóminus, et laudábilis nimis: * terríbilis est super omnes deos.

Quóniam omnes dii géntium dæmónia: * Dóminus autem cælos fecit.

Conféssio, et pulchritúdo in conspéctu ejus: * sanctimónia et magnificéntia in sanctificatióne ejus.

Afférte Dómino, pátriæ géntium, afférte Dómino glóriam et honórem: * afférte Dómino glóriam nómini ejus.

Tóllite hóstias, et introíte in átria ejus: * adoráte Dóminum in átrio sancto ejus.

Commoveátur a fácie ejus univérsa terra: * dícite in géntibus quia Dóminus regnávit.

Etenim corréxit orbem terræ qui non commovébitur: * judicábit pópulos in æquitáte.

Læténtur cæli, et exsúltet terra: commoveátur mare, et plenitúdo ejus: * gaudébunt campi, et ómnia quæ in eis sunt.

Tunc exsultábunt ómnia ligna silvárum a fácie Dómini, quia venit: * quóniam venit judicáre terram.

Judicábit orbem terræ in æquitáte, * et pópulos in veritáte tua.

\antlabel{Ant.} Cantáte Dómino, et benedícite nómini ejus.

\antlabel{Ant.} Salutáre vultus mei * Deus meus.

\psalmheading{Psalmus 42}

\lettrine[lines=2]{J}{UDICA} me, Deus, et discérne causam meam de gente non sancta, * ab hómine iníquo, et dolóso érue me.

Quia tu es, Deus, fortitúdo mea: * quare me repulísti? et quare tristis incédo, dum afflígit me inimícus?

Emítte lucem tuam et veritátem tuam: * ipsa me deduxérunt, et adduxérunt in montem sanctum tuum, et in tabernácula tua.

Et introíbo ad altáre Dei: * ad Deum, qui lætíficat juventútem meam.

Confitébor tibi in cíthara, Deus, Deus meus: * quare tristis es, ánima mea? et quare contúrbas me?

Spera in Deo, quóniam adhuc confitébor illi: * salutáre vultus mei, et Deus meus.

\antlabel{Ant.} Salutáre vultus mei Deus meus.

\antlabel{Ant.} Illúmina, Dómine, * vultum tuum super nos.

\psalmheading{Psalmus 66}

\lettrine[lines=2]{D}{EUS} misereátur nostri, et benedícat nobis: * illúminet vultum suum super nos, et misereátur nostri.

Ut cognoscámus in terra viam tuam, * in ómnibus géntibus salutáre tuum.

Confiteántur tibi pópuli, Deus: * confiteántur tibi pópuli omnes.

Læténtur et exsúltent gentes: * quóniam júdicas pópulos in æquitáte, et gentes in terra dírigis.

Confiteántur tibi pópuli, Deus, confiteántur tibi pópuli omnes: * terra dedit fructum suum.

Benedícat nos Deus, Deus noster, benedícat nos Deus: * et métuant eum omnes fines terræ.

\antlabel{Ant.} Illúmina, Dómine, vultum tuum super nos.

\antlabel{Ant.} Exaltáte * Regem sæculórum in opéribus vestris.

\canticumheading{Canticum Tobiæ}{Tob. 13, 1-10}

\lettrine[lines=2]{M}{AGNUS} es, Dómine, in ætérnum, * et in ómnia sæcula regnum tuum:

Quóniam tu flagéllas, et salvas: dedúcis ad ínferos, et redúcis: * et non est qui effúgiat manum tuam.

Confitémini Dómino, fílii Israël, * et in conspéctu géntium laudáte eum:

Quóniam ídeo dispérsit vos inter gentes, quæ ignórant eum, * ut vos enarrétis mirabília ejus.

Et faciátis scire eos, * quia non est álius Deus omnípotens præter eum.

Ipse castigávit nos propter iniquitátes nostras: * et ipse salvábit nos propter misericórdiam suam.

Aspícite ergo quæ fecit nobíscum, et cum timóre et tremóre confitémini illi: * Regémque sæculórum exaltáte in opéribus vestris.

Ego autem in terra captivitátis meæ confitébor illi: * quóniam osténdit majestátem suam in gentem peccatrícem.

Convertímini ítaque, peccatóres, et fácite justítiam coram Deo, * credéntes quod fáciat vobíscum misericórdiam suam:

Ego autem, et ánima mea * in eo lætábimur.

Benedícite Dóminum, omnes elécti ejus: * ágite dies lætítiæ, et confitémini illi.

\antlabel{Ant.} Exaltáte Regem sæculórum in opéribus vestris.

\antlabel{Ant.} Laudáte * \textdagger\ nomen Dómini, qui statis in domo Dómini.

\psalmheading{Psalmus 134}

\lettrine[lines=2]{L}{AUDATE} \textdagger\ nomen Dómini, * laudáte, servi, Dóminum.

Qui statis in domo Dómini, * in átriis domus Dei nostri.

Laudáte Dóminum, quia bonus Dóminus: * psállite nómini ejus, quóniam suáve.

Quóniam Jacob elégit sibi Dóminus, * Israël in possessiónem sibi.

Quia ego cognóvi quod magnus est Dóminus, * et Deus noster præ ómnibus diis.

Omnia quæcúmque vóluit, Dóminus fecit in cælo, et in terra, * in mari, et in ómnibus abýssis.

Edúcens nubes ab extrémo terræ: * fúlgura in plúviam fecit.

Qui prodúcit ventos de thesáuris suis: * qui percússit primogénita Ægýpti ab hómine usque ad pecus.

Et misit signa, et prodígia in médio tui, Ægýpte: * in Pharaónem, et in omnes servos ejus.

Qui percússit gentes multas: * et occídit reges fortes:

Sehon, regem Amorrhæórum, et Og, regem Basan, * et ómnia regna Chánaan.

Et dedit terram eórum hereditátem, * hereditátem Israël, pópulo suo.

Dómine, nomen tuum in ætérnum: * Dómine, memoriále tuum in generatiónem et generatiónem.

Quia judicábit Dóminus pópulum suum: * et in servis suis deprecábitur.

Simulácra géntium argéntum, et aurum, * ópera mánuum hóminum.

Os habent, et non loquéntur: * óculos habent, et non vidébunt.

Aures habent, et non áudient: * neque enim est spíritus in ore ipsórum.

Símiles illis fiant qui fáciunt ea: * et omnes qui confídunt in eis.

Domus Israël, benedícite Dómino: * domus Aaron, benedícite Dómino.

Domus Levi, benedícite Dómino: * qui timétis Dóminum, benedícite Dómino.

Benedíctus Dóminus ex Sion, * qui hábitat in Jerúsalem.

\antlabel{Ant.} Laudáte nomen Dómini, qui statis in domo Dómini.

\variantheading{In Officio feriali}

\capitulumlabel{Capit.}{Rom. 13, 12-13}

\lettrine[lines=2]{N}{OX} præcéssit, dies autem appropinquávit. Abjiciámus ergo ópera tenebrárum, et induámur arma lucis. Sicut in die honéste ambulémus.

\hymnusheading{Hymnus}

\lettrine[lines=2]{A}{LES} diéi núntius\\
Lucem propínquam præcínit:\\
Nos excitátor méntium\\
Jam Christus ad vitam vocat.

Auférte, clamat, léctulos,\\
Ægro sopóre desídes:\\
Castíque, recti, ac sóbrii\\
Vigiláte, jam sum próximus.

Jesum ciámus vócibus,\\
Flentes, precántes, sóbrii:\\
Inténta supplicátio\\
Dormíre cor mundum vetat.

Tu, Christe, somnum díscute:\\
Tu rumpe noctis víncula:\\
Tu solve peccátum vetus,\\
Novúmque lumen íngere.

Deo Patri sit glória,\\
Ejúsque soli Fílio,\\
Cum Spíritu Paráclito,\\
Nunc et per omne sæculum.\\
Amen.

\V\ Repléti sumus mane misericórdia tua. \R\ Exsultávimus, et delectáti sumus.

\antlabel{Ad Bened. Ant.} Eréxit nobis * Dóminus cornu salútis in domo David, púeri sui.

\rubric{Réliqua ut in Ordinário. 11}

\hourheading{Ad Laudes II}

\rubric{In Vigíliis commúnibus, quando fit Offícium de Féria.}

(Pater noster et Ave.)

\V\ Deus, in adjutórium.

\antlabel{Ant.} Dele iniquitátem meam.

\psalmheading{Psalmus 50}

\lettrine[lines=2]{M}{ISERERE} \textdagger\ mei, Deus: * secúndum magnam misericórdiam tuam.

Et secúndum multitúdinem miseratiónum tuárum, * dele iniquitátem meam.

Amplius lava me ab iniquitáte mea: * et a peccáto meo munda me.

Quóniam iniquitátem meam ego cognósco: * et peccátum meum contra me est semper.

Tibi soli peccávi, et malum coram te feci: * ut justificéris in sermónibus tuis, et vincas cum judicáris.

Ecce enim, in iniquitátibus concéptus sum: * et in peccátis concépit me mater mea.

Ecce enim, veritátem dilexísti: * incérta et occúlta sapiéntiæ tuæ manifestásti mihi.

Aspérges me hyssópo, et mundábor: * lavábis me, et super nivem dealbábor.

Audítui meo dabis gáudium et lætítiam: * et exsultábunt ossa humiliáta.

Avérte fáciem tuam a peccátis meis: * et omnes iniquitátes meas dele.

Cor mundum crea in me, Deus: * et spíritum rectum ínnova in viscéribus meis.

Ne projícias me a fácie tua: * et spíritum sanctum tuum ne áuferas a me.

Redde mihi lætítiam salutáris tui: * et spíritu principáli confírma me.

Docébo iníquos vias tuas: * et ímpii ad te converténtur.

Líbera me de sanguínibus, Deus, Deus salútis meæ: * et exsultábit lingua mea justítiam tuam.

Dómine, lábia mea apéries: * et os meum annuntiábit laudem tuam.

Quóniam si voluísses sacrifícium, dedíssem útique: * holocáustis non delectáberis.

Sacrifícium Deo spíritus contribulátus: * cor contrítum, et humiliátum, Deus, non despícies.

Benígne fac, Dómine, in bona voluntáte tua Sion: * ut ædificéntur muri Jerúsalem.

Tunc acceptábis sacrifícium justítiæ, oblatiónes, et holocáusta: * tunc impónent super altáre tuum vítulos.

\antlabel{Ant.} Dele iniquitátem meam, Dómine, secúndum multitúdinem miseratiónum tuárum.

\antlabel{Ant.} Discérne causam meam.

\psalmheading{Psalmus 42}

\lettrine[lines=2]{J}{UDICA} me, Deus, et discérne causam meam de gente non sancta, * ab hómine iníquo, et dolóso érue me.

Quia tu es, Deus, fortitúdo mea: * quare me repulísti? et quare tristis incédo, dum afflígit me inimícus?

Emítte lucem tuam et veritátem tuam: * ipsa me deduxérunt, et adduxérunt in montem sanctum tuum, et in tabernácula tua.

Et introíbo ad altáre Dei: * ad Deum, qui lætíficat juventútem meam.

Confitébor tibi in cíthara, Deus, Deus meus: * quare tristis es, ánima mea? et quare contúrbas me?

Spera in Deo, quóniam adhuc confitébor illi: * salutáre vultus mei, et Deus meus.

\antlabel{Ant.} Discérne causam meam, Deus, de gente non sancta.

\antlabel{Ant.} Deus misereátur. \textdagger

\psalmheading{Psalmus 66}

\lettrine[lines=2]{D}{EUS} misereátur \textdagger\ nostri, et benedícat nobis: * illúminet vultum suum super nos, et misereátur nostri.

Ut cognoscámus in terra viam tuam: * in ómnibus géntibus salutáre tuum.

Confiteántur tibi pópuli, Deus: * confiteántur tibi pópuli omnes.

Læténtur et exsúltent gentes: * quóniam júdicas pópulos in æquitáte, et gentes in terra dírigis.

Confiteántur tibi pópuli, Deus, confiteántur tibi pópuli omnes: * terra dedit fructum suum.

Benedícat nos Deus, Deus noster, benedícat nos Deus: * et métuant eum omnes fines terræ.

\antlabel{Ant.} Deus misereátur nostri, et benedícat nos.

\antlabel{Ant.} Corrípies me, Dómine.

\canticumheading{Canticum Ezechiæ}{Is. 38, 10-20}

\lettrine[lines=2]{E}{GO} dixi: In dimídio diérum meórum * vadam ad portas ínferi.

Quæsívi resíduum annórum meórum. * Dixi: Non vidébo Dóminum Deum in terra vivéntium.

Non aspíciam hóminem ultra, * et habitatórem quiétis.

Generátio mea abláta est, et convolúta est a me, * quasi tabernáculum pastórum.

Præcísa est, velut a texénte, vita mea: dum adhuc ordírer, succídit me: * de mane usque ad vésperam fínies me.

Sperábam usque ad mane, * quasi leo sic contrívit ómnia ossa mea:

De mane usque ad vésperam fínies me: * sicut pullus hirúndinis sic clamábo, meditábor ut colúmba:

Attenuáti sunt óculi mei, * suspiciéntes in excélsum:

Dómine, vim pátior, respónde pro me. * Quid dicam, aut quid respondébit mihi, cum ipse fécerit?

Recogitábo tibi omnes annos meos * in amaritúdine ánimæ meæ.

Dómine, si sic vívitur, et in tálibus vita spíritus mei, corrípies me, et vivificábis me. * Ecce, in pace amaritúdo mea amaríssima:

Tu autem eruísti ánimam meam ut non períret: * projecísti post tergum tuum ómnia peccáta mea.

Quia non inférnus confitébitur tibi, neque mors laudábit te: * non exspectábunt qui descéndunt in lacum, veritátem tuam.

Vivens vivens ipse confitébitur tibi, sicut et ego hódie: * pater fíliis notam fáciet veritátem tuam.

Dómine, salvum me fac, * et psalmos nostros cantábimus cunctis diébus vitæ nostræ in domo Dómini.

\antlabel{Ant.} Corrípies me, Dómine, et vivificábis me.

\antlabel{Ant.} Laudáte. \textdagger

\psalmheading{Psalmus 134}

\lettrine[lines=2]{L}{AUDATE} \textdagger\ nomen Dómini, * laudáte, servi, Dóminum.

Qui statis in domo Dómini, * in átriis domus Dei nostri.

Laudáte Dóminum, quia bonus Dóminus: * psállite nómini ejus, quóniam suáve.

Quóniam Jacob elégit sibi Dóminus, * Israël in possessiónem sibi.

Quia ego cognóvi quod magnus est Dóminus: * et Deus noster præ ómnibus diis.

Omnia quæcúmque vóluit, Dóminus fecit in cælo, et in terra, * in mari, et in ómnibus abýssis.

Edúcens nubes ab extrémo terræ: * fúlgura in plúviam fecit.

Qui prodúcit ventos de thesáuris suis: * qui percússit primogénita Ægýpti ab hómine usque ad pecus.

Et misit signa, et prodígia in médio tui, Ægýpte: * in Pharaónem, et in omnes servos ejus.

Qui percússit gentes multas: * et occídit reges fortes:

Sehon, regem Amorrhæórum, et Og, regem Basan, * et ómnia regna Chánaan.

Et dedit terram eórum hereditátem, * hereditátem Israël, pópulo suo.

Dómine, nomen tuum in ætérnum: * Dómine, memoriále tuum in generatiónem et generatiónem.

Quia judicábit Dóminus pópulum suum: * et in servis suis deprecábitur.

Simulácra géntium argéntum, et aurum, * ópera mánuum hóminum.

Os habent, et non loquéntur: * óculos habent, et non vidébunt.

Aures habent, et non áudient: * neque enim est spíritus in ore ipsórum.

Símiles illis fiant qui fáciunt ea: * et omnes qui confídunt in eis.

Domus Israël, benedícite Dómino: * domus Aaron, benedícite Dómino.

Domus Levi, benedícite Dómino: * qui timétis Dóminum, benedícite Dómino.

Benedíctus Dóminus ex Sion, * qui hábitat in Jerúsalem.

\antlabel{Ant.} Laudáte Dóminum, quia benígnus est, et in servis suis deprecábitur.

\capitulumlabel{Capit.}{Rom. 13, 12-13}

\lettrine[lines=2]{N}{OX} præcéssit, dies autem appropinquávit. Abjiciámus ergo ópera tenebrárum, et induámur arma lucis. Sicut in die honéste ambulémus.

\hymnusheading{Hymnus}

\lettrine[lines=2]{A}{LES} diéi núntius\\
Lucem propínquam præcínit:\\
Nos excitátor méntium\\
Jam Christus ad vitam vocat.

Auférte, clamat, léctulos,\\
Ægro sopóre desídes:\\
Castíque, recti, ac sóbrii\\
Vigiláte, jam sum próximus.

Jesum ciámus vócibus,\\
Flentes, precántes, sóbrii:\\
Inténta supplicátio\\
Dormíre cor mundum vetat.

Tu, Christe, somnum díscute:\\
Tu rumpe noctis víncula:\\
Tu solve peccátum vetus,\\
Novúmque lumen íngere.

Deo Patri sit glória,\\
Ejúsque soli Fílio,\\
Cum Spíritu Paráclito,\\
Nunc et per omne sæculum.\\
Amen.

\V\ Repléti sumus mane misericórdia tua. \R\ Exsultávimus, et delectáti sumus.

\antlabel{Ad Bened. Ant.} Eréxit nobis * Dóminus cornu salútis in domo David, púeri tui.

\rubric{Expléto Cántico ac repetíta Antíphona, dicúntur sequéntes}

\variantheading{Preces feriales}

KYRIE, eléison. Christe, eléison. Kýrie, eléison.

Pater noster \rubric{quod dícitur a solo Hebdomadário totum clara voce, usque ad} \V\ Et ne nos indúcas in tentatiónem. \R\ Sed líbera nos a malo. \V\ Ego dixi: Dómine, miserére mei. \R\ Sana ánimam meam, quia peccávi tibi. \V\ Convértere, Dómine, úsquequo? \R\ Et deprecábilis esto super servos tuos. \V\ Fiat misericórdia tua, Dómine, super nos. \R\ Quemádmodum sperávimus in te. \V\ Sacerdótes tui induántur justítiam. \R\ Et sancti tui exsúltent. \V\ Orémus pro beatíssimo Papa nostro N. \R\ Dóminus consérvet eum, et vivíficet eum, et beátum fáciat eum in terra, et non tradat eum in ánimam inimicórum ejus. \V\ Orémus et pro Antístite nostro N. \R\ Stet et pascat in fortitúdine tua, Dómine, in sublimitáte nóminis tui. \V\ Dómine, salvum fac regem. \R\ Et exáudi nos in die, qua invocavérimus te. \V\ Salvum fac pópulum tuum, Dómine, et bénedic hereditáti tuæ. \R\ Et rege eos, et extólle illos usque in ætérnum. \V\ Meménto Congregatiónis tuæ. \R\ Quam possedísti ab inítio. \V\ Fiat pax in virtúte tua. \R\ Et abundántia in túrribus tuis. \V\ Orémus pro benefactóribus nostris. \R\ Retribúere dignáre, Dómine, ómnibus, nobis bona faciéntibus propter nomen tuum, vitam ætérnam. Amen. \V\ Orémus pro fidélibus defúnctis. \R\ Réquiem ætérnam dona eis, Dómine, et lux perpétua lúceat eis. \V\ Requiéscant in pace. \R\ Amen. \V\ Pro frátribus nostris abséntibus. \R\ Salvos fac servos tuos, Deus meus, sperántes in te. \V\ Pro afflíctis et captívis. \R\ Líbera eos, Deus Israël, ex ómnibus tribulatiónibus suis. \V\ Mitte eis, Dómine, auxílium de sancto. \R\ Et de Sion tuére eos. \V\ Dómine, Deus virtútum, convérte nos. \R\ Et osténde fáciem tuam, et salvi érimus. \V\ Exsúrge, Christe, ádjuva nos. \R\ Et líbera nos propter nomen tuum. \V\ Dómine, exáudi oratiónem meam. \R\ Et clamor meus ad te véniat.

\rubric{Réliqua ut in Ordinário. 9}

% ===== END OF BATCH 15 -- Covers autumnalis.md lines 4981-5396, the
% two alternate ferial Lauds schemes for Feria Tertia, "Ad Laudes
% I" (Psalms 95, 42, 66, Canticum Tobiæ, 134) and "Ad Laudes II"
% (Psalms 50, 42, 66, Canticum Ezechiæ, 134), each with the
% Capitulum "Nox præcéssit" and the hymn "Ales diéi núntius." The
% recurring stray OCR "\*" artifact appears only in the Ad Laudes I
% copy of the hymn (autumnalis.md line 5427, before the final doxology
% stanza) and is silently dropped there; the Ad Laudes II copy of
% the identical hymn text (line 5336) has NO such artifact in the
% source and is transcribed exactly as printed, confirming the
% artifact is a per-occurrence OCR defect rather than a feature of
% the hymn text itself. The "Ad Bened. Ant." differs by a single
% word between the two schemes ("púeri sui" in Ad Laudes I vs.
% "púeri tui" in Ad Laudes II) and each is transcribed exactly as
% its own source gives it. Ad Laudes II additionally carries the
% "Preces feriales" litany in full. The document continues with
% "Ad Primam" (autumnalis.md line 5399).
% =====
\hourheading{Ad Primam}

Pater, Ave et Credo.

\V\ Deus, in adjutórium.

\hymnusheading{Hymnus}

\lettrine[lines=2]{J}{AM} lucis orto sídere,\\
Deum precémur súpplices,\\
Ut in diúrnis áctibus\\
Nos servet a nocéntibus.

Linguam refrénans témperet,\\
Ne litis horror ínsonet:\\
Visum fovéndo cóntegat,\\
Ne vanitátes háuriat.

Sint pura cordis íntima,\\
Absístat et vecórdia:\\
Carnis terat supérbiam\\
Potus cibíque párcitas.

Ut, cum dies abscésserit\\
Noctémque sors redúxerit,\\
Mundi per abstinéntiam\\
Ipsi canámus glóriam.

Deo Patri sit glória,\\
Ejúsque soli Fílio,\\
Cum Spíritu Paráclito,\\
Nunc et per omne sæculum.\\
Amen.

\antlabel{Ant.} Deus meus.

\psalmheading{Psalmus 24, i}

\lettrine[lines=2]{A}{D} te, Dómine, levávi ánimam meam: * Deus meus, in te confído, non erubéscam.

Neque irrídeant me inimíci mei: * étenim univérsi, qui sústinent te, non confundéntur.

Confundántur omnes iníqua agéntes * supervácue.

Vias tuas, Dómine, demónstra mihi: * et sémitas tuas édoce me.

Dírige me in veritáte tua, et doce me: * quia tu es, Deus, salvátor meus, et te sustínui tota die.

Reminíscere miseratiónum tuárum, Dómine, * et misericordiárum tuárum, quæ a sæculo sunt.

Delícta juventútis meæ, * et ignorántias meas ne memíneris.

Secúndum misericórdiam tuam meménto mei tu: * propter bonitátem tuam, Dómine.

\psalmheading{Psalmus 24, ii}

\lettrine[lines=2]{D}{ULCIS} et rectus Dóminus: * propter hoc legem dabit delinquéntibus in via.

Díriget mansuétos in judício: * docébit mites vias suas.

Univérsæ viæ Dómini, misericórdia et véritas, * requiréntibus testaméntum ejus et testimónia ejus.

Propter nomen tuum, Dómine, propitiáberis peccáto meo: * multum est enim.

Quis est homo qui timet Dóminum? * legem státuit ei in via, quam elégit.

Anima ejus in bonis demorábitur: * et semen ejus hereditábit terram.

Firmaméntum est Dóminus timéntibus eum: * et testaméntum ipsíus ut manifestétur illis.

\psalmheading{Psalmus 24, iii}

\lettrine[lines=2]{O}{CULI} mei semper ad Dóminum: * quóniam ipse evéllet de láqueo pedes meos.

Réspice in me, et miserére mei: * quia únicus et pauper sum ego.

Tribulatiónes cordis mei multiplicátæ sunt: * de necessitátibus meis érue me.

Vide humilitátem meam, et labórem meum: * et dimítte univérsa delícta mea.

Réspice inimícos meos quóniam multiplicáti sunt, * et ódio iníquo odérunt me.

Custódi ánimam meam, et érue me: * non erubéscam quóniam sperávi in te.

Innocéntes et recti adhæsérunt mihi: * quia sustínui te.

Líbera, Deus, Israël, * ex ómnibus tribulatiónibus suis.

\rubric{¶ Quando ad Laudes dictus fúerit Ps. 50 Miserére, hic subjúngitur seq. Ps. 95, qui in Láudibus locum non hábuit.}

\psalmheading{Psalmus 95}

\lettrine[lines=2]{C}{ANTATE} Dómino cánticum novum: * cantáte Dómino, omnis terra.

Cantáte Dómino, et benedícite nómini ejus: * annuntiáte de die in diem salutáre ejus.

Annuntiáte inter gentes glóriam ejus, * in ómnibus pópulis mirabília ejus.

Quóniam magnus Dóminus, et laudábilis nimis: * terríbilis est super omnes deos.

Quóniam omnes dii géntium dæmónia: * Dóminus autem cælos fecit.

Conféssio, et pulchritúdo in conspéctu ejus: * sanctimónia et magnificéntia in sanctificatióne ejus.

Afférte Dómino, pátriæ géntium, afférte Dómino glóriam et honórem: * afférte Dómino glóriam nómini ejus.

Tóllite hóstias, et introíte in átria ejus: * adoráte Dóminum in átrio sancto ejus.

Commoveátur a fácie ejus univérsa terra: * dícite in géntibus quia Dóminus regnávit.

Etenim corréxit orbem terræ qui non commovébitur: * judicábit pópulos in æquitáte.

Læténtur cæli, et exsúltet terra: commoveátur mare, et plenitúdo ejus: * gaudébunt campi, et ómnia quæ in eis sunt.

Tunc exsultábunt ómnia ligna silvárum a fácie Dómini, quia venit: * quóniam venit judicáre terram.

Judicábit orbem terræ in æquitáte, * et pópulos in veritáte sua.

\rubric{Secus vero, tribus Psalmis explétis, statim dícitur Ant.}

\antlabel{Ant.} Deus meus, in te confído, non erubéscam.

\rubric{In Officio Festi vel Octávæ:}

\capitulumlabel{Capitulum}{1 Tim. 1, 17}

\lettrine[lines=2]{R}{EGI} sæculórum immortáli et invisíbili, soli Deo honor et glória in sæcula sæculórum. Amen.

\rubric{In ómnibus autem Fériis et in Vigíliis commúnibus:}

\capitulumlabel{Capitulum}{Zach. 8, 19}

\lettrine[lines=2]{P}{ACEM} et veritátem dilígite, ait Dóminus omnípotens.

\R\ br. Christe, Fili Dei vivi, * Miserére nobis. Christe. \V\ \textit{Qui sedes ad déxteram Patris.} Miserére nobis. Glória Patri. Christe. \V\ Exsúrge, Christe, ádjuva nos. \R\ Et líbera nos propter nomen tuum.

\variantheading{Preces dominicales et feriales (si dicendæ sint)}

KYRIE, eléison. Christe, eléison. Kýrie, eléison.

Pater noster \rubric{secreto usque ad} \V\ Et ne nos indúcas in tentatiónem. \R\ Sed líbera nos a malo.

Credo in Deum \rubric{secreto usque ad} \V\ Carnis resurrectiónem. \R\ Vitam ætérnam. Amen. \V\ Et ego ad te, Dómine, clamávi. \R\ Et mane orátio mea prævéniet te. \V\ Repleátur os meum laude. \R\ Ut cantem glóriam tuam, tota die magnitúdinem tuam. \V\ Dómine, avérte fáciem tuam a peccátis meis. \R\ Et omnes iniquitátes meas dele. \V\ Cor mundum crea in me, Deus. \R\ Et spíritum rectum ínnova in viscéribus meis. \V\ Ne projícias me a fácie tua. \R\ Et spíritum sanctum tuum ne áuferas a me. \V\ Redde mihi lætítiam salutáris tui. \R\ Et spíritu principáli confírma me.

\variantheading{In Précibus feriálibus tantum}

\V\ Eripe me, Dómine, ab hómine malo. \R\ A viro iníquo éripe me. \V\ Eripe me de inimícis meis, Deus meus. \R\ Et ab insurgéntibus in me líbera me. \V\ Eripe me de operántibus iniquitátem. \R\ Et de viris sánguinum salva me. \V\ Sic psalmum dicam nómini tuo in sæculum sæculi. \R\ Ut reddam vota mea de die in diem. \V\ Exáudi nos, Deus, salutáris noster. \R\ Spes ómnium fínium terræ, et in mari longe. \V\ Deus, in adjutórium meum inténde. \R\ Dómine, ad adjuvándum me festína. \V\ Sanctus Deus, Sanctus fortis, Sanctus immortális. \R\ Miserére nobis. \V\ Bénedic, ánima mea, Dómino. \R\ Et ómnia, quæ intra me sunt, nómini sancto ejus. \V\ Bénedic, ánima mea, Dómino. \R\ Et noli oblivísci omnes retributiónes ejus. \V\ Qui propitiátur ómnibus iniquitátibus tuis. \R\ Qui sanat omnes infirmitátes tuas. \V\ Qui rédimit de intéritu vitam tuam. \R\ Qui corónat te in misericórdia et miseratiónibus. \V\ Qui replet in bonis desidérium tuum. \R\ Renovábitur ut áquilæ juvéntus tua.

\variantheading{In utrisque Précibus}

\V\ Adjutórium nostrum in nómine Dómini. \R\ Qui fecit cælum et terram.

\formulalabel{Confíteor. Misereátur. Indulgéntiam.}

\V\ Dignáre, Dómine, die isto. \R\ Sine peccáto nos custodíre. \V\ Miserére nostri, Dómine. \R\ Miserére nostri. \V\ Fiat misericórdia tua, Dómine, super nos. \R\ Quemádmodum sperávimus in te. \V\ Dómine, exáudi oratiónem meam. \R\ Et clamor meus ad te véniat.

\V\ Dóminus vobíscum.

Orémus. \oratiolabel{Oratio}

\lettrine[lines=2]{D}{OMINE}, Deus omnípotens, qui ad princípium hujus diéi nos perveníre fecísti: tua nos hódie salva virtúte; ut in hac die ad nullum declinémus peccátum, sed semper ad tuam justítiam faciéndam nostra procédant elóquia, dirigántur cogitatiónes et ópera. Per Dóminum. \R\ Amen.

\V\ Dóminus vobíscum. \V\ Benedicámus Dómino.

\rubric{In Choro Martyrológium.}

\V\ Pretiósa in conspéctu Dómini. \R\ Mors Sanctórum ejus.

\rubric{Et absolúte, sine} Orémus:

\oratiolabel{Oratio}

\lettrine[lines=2]{S}{ANCTA} María, et omnes Sancti intercédant pro nobis ad Dóminum, ut nos mereámur ab eo adjuvári et salvári, qui vivit et regnat in sæcula sæculórum. \R\ Amen.

\V\ Deus, in adjutórium meum inténde. \R\ Dómine, ad adjuvándum me festína. \rubric{Et dícitur ter; último additur:} Glória Patri. Sicut erat.

Kýrie, eléison. Christe, eléison. Kýrie, eléison.

Pater noster \rubric{secreto usque ad} \V\ Et ne nos indúcas in tentatiónem. \R\ Sed líbera nos a malo. \V\ Réspice in servos tuos, Dómine, et in ópera tua, et dírige fílios eórum. \R\ Et sit splendor Dómini, Dei nostri, super nos, et ópera mánuum nostrárum dírige super nos, et opus mánuum nostrárum dírige. \V\ Glória Patri. \R\ Sicut erat.

Orémus. \oratiolabel{Oratio}

\lettrine[lines=2]{D}{IRIGERE} et sanctificáre, régere et gubernáre dignáre, Dómine Deus, Rex cæli et terræ, hódie corda et córpora nostra, sensus, sermónes et actus nostros in lege tua, et in opéribus mandatórum tuórum: ut hic, et in ætérnum, te auxiliánte, salvi et líberi esse mereámur, Salvátor mundi: Qui vivis et regnas in sæcula sæculórum. \R\ Amen.

\V\ Jube, domne.

\benedlabel{Bened.} Dies et actus nostros in sua pace dispónat Dóminus omnípotens. \R\ Amen.

\variantheading{In Officio feriali}

\lettrine[lines=2]{D}{OMINUS} autem dírigat corda et córpora nostra in caritáte Dei, et patiéntia Christi. Tu autem.

\variantheading{E Proprio de Tempore}

\variantheading{Pro Oct. Omn. Sanctorum}

\lettrine[lines=2]{B}{ENEDÍCTIO} et cláritas et sapiéntia et gratiárum áctio, honor, virtus et fortitúdo Deo nostro in sǽcula sæculórum. Amen.

\variantheading{E Communi Sanctorum}

\variantheading{Pro Apostolis}

\lettrine[lines=2]{I}{BANT} Apóstoli gaudéntes a conspéctu concílii, quóniam digni hábiti sunt pro nómine Jesu contuméliam pati.

\variantheading{Pro Martyre}

\lettrine[lines=2]{J}{USTUS} cor suum trádidit ad vigilándum dilúculo ad Dóminum, qui fecit illum, et in conspéctu Altíssimi deprecábitur.

\variantheading{Pro pluribus Martyribus}

\lettrine[lines=2]{F}{ULGEBUNT} justi, et tamquam scintíllæ in arundinéto discúrrent. Judicábunt natiónes, et dominabúntur pópulis: et regnábit Dóminus illórum in perpétuum.

\variantheading{Pro Conf. Pont.}

\lettrine[lines=2]{F}{UNGI} sacerdótio, et habére laudem in nómine ipsíus, et offérre illi incénsum dignum in odórem suavitátis.

\variantheading{Pro Conf. non Pont.}

\lettrine[lines=2]{J}{USTUM} dedúxit Dóminus per vias rectas, et osténdit illi regnum Dei, et dedit illi sciéntiam sanctórum; honestávit illum in labóribus, et complévit labóres illíus.

\variantheading{Pro Virgine Martyre}

\lettrine[lines=2]{D}{OMINE}, Deus meus, exaltásti super terram habitatiónem meam, et pro morte defluénte deprecáta sum. Invocávi Dóminum, Patrem Dómini mei, ut non derelínquat me in die tribulatiónis meæ, et in témpore superbórum sine adjutório.

\variantheading{Pro Virgine non Martyre}

\lettrine[lines=2]{O}{} \textsc{quam} pulchra est casta generátio cum claritáte! immortális est enim memória illíus: quóniam apud Deum nota est, et apud hómines.

\variantheading{Pro Martyre tantum}

\lettrine[lines=2]{L}{AUDABIT} usque ad mortem ánima mea Dóminum, quóniam éruis sustinéntes te, et líberas eos de manu angústiæ, Dómine, Deus noster.

\variantheading{Pro nec Virg. nec Mart.}

\lettrine[lines=2]{M}{ULTAE} fíliæ congregavérunt divítias: tu supergréssa es univérsas. Fallax grátia, et vana est pulchritúdo: múlier timens Dóminum, ipsa laudábitur.

\variantheading{Pro Dedicatione Ecclesiæ}

\lettrine[lines=2]{E}{T} abstérget Deus omnem lácrimam ab óculis eórum: et mors ultra non erit, neque luctus, neque clamor, neque dolor erit ultra, quia prima abiérunt. Et dixit qui sedébat in throno: Ecce, nova fácio ómnia.

\variantheading{In Officiis B. Mariæ V.}

\lettrine[lines=2]{I}{N} platéis sicut cinnamómum et bálsamum aromatízans odórem dedi: quasi myrrha elécta dedi suavitátem odóris.

\V\ Adjutórium nostrum in nómine Dómini. \R\ Qui fecit cælum et terram.

\V\ Benedícite. \R\ Deus.

\benedictiolabel{Benedictio} Dóminus nos benedícat, et ab omni malo deféndat, et ad vitam perdúcat ætérnam. Et fidélium ánimæ per misericórdiam Dei requiéscant in pace. \R\ Amen.

\rubric{Deinde dícitur tantum} Pater noster \rubric{secréto, nisi sequátur ália Hora.}

% ===== END OF BATCH 16 -- Covers autumnalis.md lines 5399-5652,
% completing "Ad Primam" for Feria Tertia in full: the hymn "Jam
% lucis orto sídere" (stray OCR "\*" before the final doxology
% stanza silently dropped), Psalms 24/i-iii, the conditional Psalm
% 95 insertion (used when Ps. 50 was said at Lauds instead -- its
% final verse here reads "et pópulos in veritáte sua," genuinely
% differing from the "veritáte tua" reading seen in the same
% psalm's earlier Ad Laudes I occurrence this batch-cycle, and
% transcribed exactly as the source gives each), both Capitula, the
% brief Responsory, the full Preces apparatus, the closing
% Orationes, and the Lectio-brevis apparatus. The source heads this
% apparatus "E Proprio de Tempore" (autumnalis.md line 5593) rather
% than "E Proprio Sanctorum" as at the parallel point in Feria
% Secunda, despite containing the same single "Pro Oct. Omn.
% Sanctorum" heading used throughout this volume's substitute for
% aestiva.tex's two-item Summer list; this is transcribed exactly as
% the source heads it, confirmed via direct inspection of
% autumnalis.md at this point. The document continues with "Ad
% Tertiam" (autumnalis.md line 5655).
% =====
\hourheading{Ad Tertiam}

Pater noster et Ave.

\V\ Deus, in adjutórium.

\hymnusheading{Hymnus}

\lettrine[lines=2]{N}{UNC}, Sancte, nobis, Spíritus,\\
Unum Patri cum Fílio,\\
Dignáre promptus íngeri\\
Nostro refúsus péctori.

Os, lingua, mens, sensus, vigor\\
Confessiónem pérsonent,\\
Flamméscat igne cáritas,\\
Accéndat ardor próximos.

Præsta, Pater piíssime,\\
Patríque compar Unice,\\
Cum Spíritu Paráclito\\
Regnans per omne sæculum.\\
Amen.

\antlabel{Ant.} Respéxit me.

\psalmheading{Psalmus 39, i}

\lettrine[lines=2]{E}{XSPECTANS} exspectávi Dóminum, * et inténdit mihi.

Et exaudívit preces meas: * et edúxit me de lacu misériæ, et de luto fæcis.

Et státuit super petram pedes meos: * et diréxit gressus meos.

Et immísit in os meum cánticum novum, * carmen Deo nostro.

Vidébunt multi, et timébunt: * et sperábunt in Dómino.

Beátus vir, cujus est nomen Dómini spes ejus: * et non respéxit in vanitátes et insánias falsas.

Multa fecísti tu, Dómine, Deus meus, mirabília tua: * et cogitatiónibus tuis non est qui símilis sit tibi.

Annuntiávi et locútus sum: * multiplicáti sunt super númerum.

Sacrifícium et oblatiónem noluísti: * aures autem perfecísti mihi.

Holocáustum et pro peccáto non postulásti: * tunc dixi: Ecce, vénio.

In cápite libri scriptum est de me ut fácerem voluntátem tuam: * Deus meus, vólui, et legem tuam in médio cordis mei.

\psalmheading{Psalmus 39, ii}

\lettrine[lines=2]{A}{NNUNTIAVI} justítiam tuam in ecclésia magna, * ecce, lábia mea non prohibébo: Dómine, tu scisti.

Justítiam tuam non abscóndi in corde meo: * veritátem tuam et salutáre tuum dixi.

Non abscóndi misericórdiam tuam et veritátem tuam * a concílio multo.

Tu autem, Dómine, ne longe fácias miseratiónes tuas a me: * misericórdia tua et véritas tua semper suscepérunt me.

Quóniam circumdedérunt me mala, quorum non est númerus: * comprehendérunt me iniquitátes meæ, et non pótui ut vidérem.

Multiplicátæ sunt super capíllos cápitis mei: * et cor meum derelíquit me.

Compláceat tibi, Dómine, ut éruas me: * Dómine, ad adjuvándum me réspice.

\psalmheading{Psalmus 39, iii}

\lettrine[lines=2]{C}{ONFUNDANTUR} et revereántur simul, qui quærunt ánimam meam, * ut áuferant eam.

Convertántur retrórsum, et revereántur, * qui volunt mihi mala.

Ferant conféstim confusiónem suam, * qui dicunt mihi: Euge, euge.

Exsúltent et læténtur super te omnes quæréntes te: * et dicant semper: Magnificétur Dóminus: qui díligunt salutáre tuum.

Ego autem mendícus sum, et pauper: * Dóminus sollícitus est mei.

Adjútor meus, et protéctor meus tu es: * Deus meus, ne tardáveris.

\antlabel{Ant.} Respéxit me, et exaudívit deprecatiónem meam Dóminus.

\variantheading{In Officio feriali}

\capitulumlabel{Capitulum}{Jer. 17, 14}

\lettrine[lines=2]{S}{ANA} me, Dómine, et sanábor: salvum me fac, et salvus ero: quóniam laus mea tu es.

\R\ br. Sana ánimam meam, * Quia peccávi tibi. Sana. \V\ Ego dixi: Dómine, miserére mei. Quia. Glória Patri. Sana. \V\ Adjútor meus esto, ne derelínquas me. \R\ Neque despícias me, Deus, salutáris meus.

\variantheading{In Vigiliis communibus:}

\variantheading{Preces feriales}

KYRIE, eléison. Christe, eléison. Kýrie, eléison.

Pater noster \rubric{secreto usque ad} \V\ Et ne nos indúcas in tentatiónem. \R\ Sed líbera nos a malo. \V\ Dómine, Deus virtútum, convérte nos. \R\ Et osténde fáciem tuam, et salvi érimus. \V\ Exsúrge, Christe, ádjuva nos. \R\ Et líbera nos propter nomen tuum. \V\ Dómine, exáudi oratiónem meam. \R\ Et clamor meus ad te véniat.

\rubric{Réliqua ut in Ordinário. 19}

\hourheading{Ad Sextam}

Pater noster et Ave.

\V\ Deus, in adjutórium.

\hymnusheading{Hymnus}

\lettrine[lines=2]{R}{ECTOR} potens, verax Deus,\\
Qui témperas rerum vices,\\
Splendóre mane illúminas,\\
Et ígnibus merídiem:

Exstíngue flammas lítium,\\
Aufer calórem nóxium,\\
Confer salútem córporum,\\
Verámque pacem córdium.

Præsta, Pater piíssime,\\
Patríque compar Unice,\\
Cum Spíritu Paráclito\\
Regnans per omne sæculum.\\
Amen.

\antlabel{Ant.} Suscepísti me, Dómine.

\psalmheading{Psalmus 40}

\lettrine[lines=2]{B}{EATUS} qui intélligit super egénum, et páuperem: * in die mala liberábit eum Dóminus.

Dóminus consérvet eum, et vivíficet eum, et beátum fáciat eum in terra: * et non tradat eum in ánimam inimicórum ejus.

Dóminus opem ferat illi super lectum dolóris ejus: * univérsum stratum ejus versásti in infirmitáte ejus.

Ego dixi: Dómine, miserére mei: * sana ánimam meam, quia peccávi tibi.

Inimíci mei dixérunt mala mihi: * Quando moriétur, et períbit nomen ejus?

Et si ingrediebátur ut vidéret, vana loquebátur: * cor ejus congregávit iniquitátem sibi.

Egrediebátur foras, * et loquebátur in idípsum.

Advérsum me susurrábant omnes inimíci mei: * advérsum me cogitábant mala mihi.

Verbum iníquum constituérunt advérsum me: * Numquid qui dormit non adjíciet ut resúrgat?

Etenim homo pacis meæ, in quo sperávi: * qui edébat panes meos, magnificávit super me supplantatiónem.

Tu autem, Dómine, miserére mei, et resúscita me: * et retríbuam eis.

In hoc cognóvi quóniam voluísti me: * quóniam non gaudébit inimícus meus super me.

Me autem propter innocéntiam suscepísti: * et confirmásti me in conspéctu tuo in ætérnum.

Benedíctus Dóminus, Deus Israël, a sæculo et usque in sæculum: * fiat, fiat.

\psalmheading{Psalmus 41, i}

\lettrine[lines=2]{Q}{UEMADMODUM} desíderat cervus ad fontes aquárum: * ita desíderat ánima mea ad te, Deus.

Sitívit ánima mea ad Deum fortem vivum: * quando véniam, et apparébo ante fáciem Dei?

Fuérunt mihi lácrimæ meæ panes die ac nocte: * dum dícitur mihi quotídie: Ubi est Deus tuus?

Hæc recordátus sum, et effúdi in me ánimam meam: * quóniam transíbo in locum tabernáculi admirábilis, usque ad domum Dei.

In voce exsultatiónis, et confessiónis: * sonus epulántis.

Quare tristis es, ánima mea? * et quare contúrbas me?

Spera in Deo, quóniam adhuc confitébor illi: * salutáre vultus mei, et Deus meus.

\psalmheading{Psalmus 41, ii}

\lettrine[lines=2]{A}{D} meípsum ánima mea conturbáta est: * proptérea memor ero tui de terra Jordánis, et Hermóniim a monte módico.

Abýssus abýssum ínvocat, * in voce cataractárum tuárum.

Omnia excélsa tua, et fluctus tui * super me transiérunt.

In die mandávit Dóminus misericórdiam suam: * et nocte cánticum ejus.

Apud me orátio Deo vitæ meæ, * dicam Deo: Suscéptor meus es.

Quare oblítus es mei? * et quare contristátus incédo, dum afflígit me inimícus?

Dum confringúntur ossa mea, * exprobravérunt mihi qui tríbulant me inimíci mei.

Dum dicunt mihi per síngulos dies: Ubi est Deus tuus? * quare tristis es, ánima mea? et quare contúrbas me?

Spera in Deo, quóniam adhuc confitébor illi: * salutáre vultus mei, et Deus meus.

\antlabel{Ant.} Suscepísti me, Dómine: et confirmásti me in conspéctu tuo.

\variantheading{In Officio feriali}

\capitulumlabel{Capitulum}{Rom. 13, 8}

\lettrine[lines=2]{N}{EMINI} quidquam debeátis, nisi ut ínvicem diligátis: qui enim díligit próximum, legem implévit.

\R\ br. Benedícam Dóminum * In omni témpore. Benedícam. \V\ Semper laus ejus in ore meo. In. Glória Patri. Benedícam. \V\ Dóminus regit me, et nihil mihi déerit. \R\ In loco páscuæ ibi me collocávit.

\variantheading{In Vigiliis communibus:}

\variantheading{Preces feriales}

KYRIE, eléison. Christe, eléison. Kýrie, eléison.

Pater noster \rubric{secreto usque ad} \V\ Et ne nos indúcas in tentatiónem. \R\ Sed líbera nos a malo. \V\ Dómine, Deus virtútum, convérte nos. \R\ Et osténde fáciem tuam, et salvi érimus. \V\ Exsúrge, Christe, ádjuva nos. \R\ Et líbera nos propter nomen tuum. \V\ Dómine, exáudi oratiónem meam. \R\ Et clamor meus ad te véniat.

\rubric{Réliqua ut in Ordinário. 20}

\hourheading{Ad Nonam}

Pater noster et Ave.

\V\ Deus, in adjutórium.

\hymnusheading{Hymnus}

\lettrine[lines=2]{R}{ERUM}, Deus, tenax vigor,\\
Immótus in te pérmanens,\\
Lucis diúrnæ témpora\\
Succéssibus detérminans:

Largíre lumen véspere,\\
Quo vita nusquam décidat,\\
Sed prǽmium mortis sacræ\\
Perénnis instet glória.

Præsta, Pater piíssime,\\
Patríque compar Unice,\\
Cum Spíritu Paráclito\\
Regnans per omne sæculum.\\
Amen.

\antlabel{Ant.} Salvásti nos.

\psalmheading{Psalmus 43, i}

\lettrine[lines=2]{D}{EUS}, áuribus nostris audívimus: * patres nostri annuntiavérunt nobis.

Opus, quod operátus es in diébus eórum, * et in diébus antíquis.

Manus tua gentes dispérdidit, et plantásti eos: * afflixísti pópulos, et expulísti eos.

Nec enim in gládio suo possedérunt terram, * et bráchium eórum non salvávit eos:

Sed déxtera tua, et bráchium tuum, et illuminátio vultus tui: * quóniam complacuísti in eis.

Tu es ipse Rex meus et Deus meus: * qui mandas salútes Jacob.

In te inimícos nostros ventilábimus cornu: * et in nómine tuo spernémus insurgéntes in nobis.

Non enim in arcu meo sperábo: * et gládius meus non salvábit me.

Salvásti enim nos de affligéntibus nos: * et odiéntes nos confudísti.

In Deo laudábimur tota die: * et in nómine tuo confitébimur in sæculum.

\psalmheading{Psalmus 43, ii}

\lettrine[lines=2]{N}{UNC} autem repulísti et confudísti nos: * et non egrediéris, Deus, in virtútibus nostris.

Avertísti nos retrórsum post inimícos nostros: * et qui odérunt nos, diripiébant sibi.

Dedísti nos tamquam oves escárum: * et in géntibus dispersísti nos.

Vendidísti pópulum tuum sine prétio: * et non fuit multitúdo in commutatiónibus eórum.

Posuísti nos oppróbrium vicínis nostris, * subsannatiónem et derísum his, qui sunt in circúitu nostro.

Posuísti nos in similitúdinem géntibus: * commotiónem cápitis in pópulis.

Tota die verecúndia mea contra me est, * et confúsio faciéi meæ coopéruit me.

A voce exprobrántis, et obloquéntis: * a fácie inimíci, et persequéntis.

Hæc ómnia venérunt super nos, nec oblíti sumus te: * et iníque non égimus in testaménto tuo.

Et non recéssit retro cor nostrum: * et declinásti sémitas nostras a via tua:

Quóniam humiliásti nos in loco afflictiónis, * et coopéruit nos umbra mortis.

\psalmheading{Psalmus 43, iii}

\lettrine[lines=2]{S}{I} oblíti sumus nomen Dei nostri, * et si expándimus manus nostras ad deum aliénum:

Nonne Deus requíret ista? * ipse enim novit abscóndita cordis.

Quóniam propter te mortificámur tota die: * æstimáti sumus sicut oves occisiónis.

Exsúrge, quare obdórmis, Dómine? * exsúrge, et ne repéllas in finem.

Quare fáciem tuam avértis, * oblivísceris inópiæ nostræ, et tribulatiónis nostræ?

Quóniam humiliáta est in púlvere ánima nostra: * conglutinátus est in terra venter noster.

Exsúrge, Dómine, ádjuva nos: * et rédime nos propter nomen tuum.

\antlabel{Ant.} Salvásti nos, Dómine: et in nómine tuo confitébimur in sæcula.

\variantheading{In Officio feriali}

\capitulumlabel{Capit.}{1 Petri 1, 17-19}

\lettrine[lines=2]{I}{N} timóre incolátus vestri témpore conversámini: sciéntes quod non corruptibílibus auro vel argénto redémpti estis, sed pretióso sánguine quasi Agni immaculáti Christi.

\R\ br. Rédime me, Dómine, * Et miserére mei. Rédime. \V\ Pes enim meus stetit in via recta. Et miserére mei. Glória Patri. Rédime. \V\ Ab occúltis meis munda me, Dómine. \R\ Et ab aliénis parce servo tuo.

\variantheading{In Vigiliis communibus:}

\variantheading{Preces feriales}

KYRIE, eléison. Christe, eléison. Kýrie, eléison.

Pater noster \rubric{secreto usque ad} \V\ Et ne nos indúcas in tentatiónem. \R\ Sed líbera nos a malo. \V\ Dómine, Deus virtútum, convérte nos. \R\ Et osténde fáciem tuam, et salvi érimus. \V\ Exsúrge, Christe, ádjuva nos. \R\ Et líbera nos propter nomen tuum. \V\ Dómine, exáudi oratiónem meam. \R\ Et clamor meus ad te véniat.

\rubric{Réliqua ut in Ordinário. 22}

% ===== END OF BATCH 17 -- Covers autumnalis.md lines 5655-5978,
% completing the Little Hours of Feria Tertia in full: "Ad
% Tertiam" (Psalms 39/i-iii), "Ad Sextam" (Psalms 40, 41/i, 41/ii),
% and "Ad Nonam" (Psalms 43/i-iii), each with its own hymn,
% Capitulum, brief Responsory, and "In Vigiliis communibus: Preces
% feriales" apparatus. Three further occurrences of the recurring
% stray OCR "\*" artifact (in the hymns "Nunc, Sancte, nobis,
% Spíritus," "Rector potens, verax Deus," and "Rerum, Deus, tenax
% vigor," each before the final "Præsta, Pater piíssime" doxology
% stanza) are silently dropped, consistent with every prior
% occurrence in this volume. The document continues with "Ad
% Vesperas" (autumnalis.md line 5981).
% =====
\hourheading{Ad Vesperas}

Pater noster et Ave.

\V\ Deus, in adjutórium.

\antlabel{Ant.} Qui hábitas in cælis, * miserére nobis.

\psalmheading{Psalmus 122}

\lettrine[lines=2]{A}{D} te levávi óculos meos, * qui hábitas in cælis.

Ecce, sicut óculi servórum * in mánibus dominórum suórum.

Sicut óculi ancíllæ in mánibus dóminæ suæ: * ita óculi nostri ad Dóminum, Deum nostrum, donec misereátur nostri.

Miserére nostri, Dómine, miserére nostri: * quia multum repléti sumus despectióne:

Quia multum repléta est ánima nostra: * oppróbrium abundántibus, et despéctio supérbis.

\antlabel{Ant.} Qui hábitas in cælis, miserére nobis.

\antlabel{Ant.} Adjutórium nostrum * in nómine Dómini.

\psalmheading{Psalmus 123}

\lettrine[lines=2]{N}{ISI} quia Dóminus erat in nobis, dicat nunc Israël: * nisi quia Dóminus erat in nobis.

Cum exsúrgerent hómines in nos, * forte vivos deglutíssent nos:

Cum irascerétur furor eórum in nos, * fórsitan aqua absorbuísset nos.

Torréntem pertransívit ánima nostra: * fórsitan pertransísset ánima nostra aquam intolerábilem.

Benedíctus Dóminus * qui non dedit nos in captiónem déntibus eórum.

Anima nostra sicut passer erépta est * de láqueo venántium:

Láqueus contrítus est, * et nos liberáti sumus.

Adjutórium nostrum in nómine Dómini, * qui fecit cælum et terram.

\antlabel{Ant.} Adjutórium nostrum in nómine Dómini.

\antlabel{Ant.} In circúitu pópuli sui * Dóminus, ex hoc nunc et usque in sæculum.

\psalmheading{Psalmus 124}

\lettrine[lines=2]{Q}{UI} confídunt in Dómino, sicut mons Sion: * non commovébitur in ætérnum, qui hábitat in Jerúsalem.

Montes in circúitu ejus: * et Dóminus in circúitu pópuli sui, ex hoc nunc et usque in sæculum.

Quia non relínquet Dóminus virgam peccatórum super sortem justórum: * ut non exténdant justi ad iniquitátem manus suas.

Bénefac, Dómine, bonis, * et rectis corde.

Declinántes autem in obligatiónes addúcet Dóminus cum operántibus iniquitátem: * pax super Israël.

\antlabel{Ant.} In circúitu pópuli sui Dóminus, ex hoc nunc et usque in sæculum.

\antlabel{Ant.} Magnificávit Dóminus * fácere nobíscum: facti sumus lætántes.

\psalmheading{Psalmus 125}

\lettrine[lines=2]{I}{N} converténdo Dóminus captivitátem Sion: * facti sumus sicut consoláti:

Tunc replétum est gáudio os nostrum: * et lingua nostra exsultatióne.

Tunc dicent inter gentes: * Magnificávit Dóminus fácere cum eis.

Magnificávit Dóminus fácere nobíscum: * facti sumus lætántes.

Convérte, Dómine, captivitátem nostram, * sicut torrens in Austro.

Qui séminant in lácrimis, * in exsultatióne metent.

Eúntes ibant et flebant, * mitténtes sémina sua.

Veniéntes autem vénient cum exsultatióne, * portántes manípulos suos.

\antlabel{Ant.} Magnificávit Dóminus fácere nobíscum: facti sumus lætántes.

\antlabel{Ant.} Dóminus ædíficet * nobis domum, et custódiat civitátem.

\psalmheading{Psalmus 126}

\lettrine[lines=2]{N}{ISI} Dóminus ædificáverit domum, * in vanum laboravérunt qui ædíficant eam.

Nisi Dóminus custodíerit civitátem, * frustra vígilat qui custódit eam.

Vanum est vobis ante lucem súrgere: * súrgite postquam sedéritis, qui manducátis panem dolóris.

Cum déderit diléctis suis somnum: * ecce heréditas Dómini fílii: merces, fructus ventris.

Sicut sagíttæ in manu poténtis: * ita fílii excussórum.

Beátus vir, qui implévit desidérium suum ex ipsis: * non confundétur cum loquétur inimícis suis in porta.

\antlabel{Ant.} Dóminus ædíficet nobis domum, et custódiat civitátem.

\variantheading{In Officio feriali}

\capitulumlabel{Capitulum}{2 Cor. 1, 3-4}

\lettrine[lines=2]{B}{ENEDICTUS} Deus, et Pater Dómini nostri Jesu Christi, Pater misericordiárum, et Deus totíus consolatiónis, qui consolátur nos in omni tribulatióne nostra.

\hymnusheading{Hymnus}

\lettrine[lines=2]{T}{ELLURIS} alme Cónditor\\
Mundi solum qui séparans,\\
Pulsis aquæ moléstiis,\\
Terram dedísti immóbilem:

Ut germen aptum próferens,\\
Fulvis decóra flóribus,\\
Fecúnda fructu sísteret,\\
Pastúmque gratum rédderet.

Mentis perústæ vúlnera\\
Munda viróre grátiæ:\\
Ut facta fletu díluat,\\
Motúsque pravos átterat.

Jussis tuis obtémperet:\\
Nullis malis appróximet:\\
Bonis repléri gáudeat,\\
Et mortis ictum nésciat.

Præsta, Pater piíssime,\\
Patríque compar Unice,\\
Cum Spíritu Paráclito\\
Regnans per omne sæculum.\\
Amen.

\V\ Dirigátur, Dómine, orátio mea. \R\ Sicut incénsum in conspéctu tuo.

\antlabel{Ad Magnif. Ant.} Exsultávit * spíritus meus in Deo, salutári meo.

\rubric{Réliqua ut in Ordinário. 22}

\hourheading{Ad Completorium}

\V\ Jube, domne.

\benedlabel{Bened.} Noctem quiétam, et finem perféctum concédat nobis Dóminus omnípotens. \R\ Amen.

\lectioheading{Lectio br.}{1 Petri 5, 8-9}

\lettrine[lines=2]{F}{RATRES}: Sóbrii estóte, et vigiláte: quia adversárius vester diábolus tamquam leo rúgiens círcuit, quærens quem dévoret: cui resístite fortes in fide. Tu autem, Dómine, miserére nobis. \R\ Deo grátias.

\V\ Adjutórium nostrum in nómine Dómini. \R\ Qui fecit cælum et terram.

Pater noster \rubric{quod dícitur totum secréto.}

\formulalabel{Confíteor.} Misereátur. Indulgéntiam.

\V\ Convérte nos, Deus, salutáris noster. \R\ Et avérte iram tuam a nobis.

\V\ Deus, in adjutórium.

\antlabel{Ant.} Tu Dómine.

\psalmheading{Psalmus 11}

\lettrine[lines=2]{S}{ALVUM} me fac, Dómine, quóniam defécit sanctus: * quóniam diminútæ sunt veritátes a fíliis hóminum.

Vana locúti sunt unusquísque ad próximum suum: * lábia dolósa, in corde et corde locúti sunt.

Dispérdat Dóminus univérsa lábia dolósa, * et linguam magníloquam.

Qui dixérunt: Linguam nostram magnificábimus, lábia nostra a nobis sunt, * quis noster Dóminus est?

Propter misériam ínopum, et gémitum páuperum, * nunc exsúrgam, dicit Dóminus.

Ponam in salutári: * fiduciáliter agam in eo.

Elóquia Dómini, elóquia casta: * argéntum igne examinátum, probátum terræ purgátum séptuplum.

Tu, Dómine, servábis nos: et custódies nos * a generatióne hac in ætérnum.

In circúitu ímpii ámbulant: * secúndum altitúdinem tuam multiplicásti fílios hóminum.

\psalmheading{Psalmus 12}

\lettrine[lines=2]{U}{SQUEQUO}, Dómine, oblivascéris me in finem? * Usquequo avértis fáciem tuam a me?

Quámdiu ponam consília in ánima mea, * dolórem in corde meo per diem?

Usquequo exaltábitur inimícus meus super me? * réspice, et exáudi me, Dómine, Deus meus.

Illúmina óculos meos ne umquam obdórmiam in morte: * nequándo dicat inimícus meus: Præválui advérsus eum.

Qui tríbulant me, exsultábunt si motus fúero: * ego autem in misericórdia tua sperávi.

Exsultábit cor meum in salutári tuo: cantábo Dómino qui bona tríbuit mihi: * et psallam nómini Dómini altíssimi.

\psalmheading{Psalmus 15}

\lettrine[lines=2]{C}{ONSERVA} me, Dómine, quóniam sperávi in te. * Dixi Dómino: Deus meus es tu, quóniam bonórum meórum non eges.

Sanctis, qui sunt in terra ejus, * mirificávit omnes voluntátes meas in eis.

Multiplicátæ sunt infirmitátes eórum: * póstea acceleravérunt.

Non congregábo conventícula eórum de sanguínibus, * nec memor ero nóminum eórum per lábia mea.

Dóminus pars hereditátis meæ, et cálicis mei: * tu es, qui restítues hereditátem meam mihi.

Funes cecidérunt mihi in præcláris: * étenim heréditas mea præclára est mihi.

Benedícam Dóminum, qui tríbuit mihi intelléctum: * ínsuper et usque ad noctem increpuérunt me renes mei.

Providébam Dóminum in conspéctu meo semper: * quóniam a dextris est mihi, ne commóvear.

Propter hoc lætátum est cor meum, et exsultávit lingua mea: * ínsuper et caro mea requiéscet in spe.

Quóniam non derelínques ánimam meam in inférno: * nec dabis sanctum tuum vidére corruptiónem.

Notas mihi fecísti vias vitæ, adimplébis me lætítia cum vultu tuo: * delectatiónes in déxtera tua usque in finem.

\antlabel{Ant.} Tu, Dómine, servábis nos: et custódies nos in ætérnum.

\hymnusheading{Hymnus}

\lettrine[lines=2]{T}{E} lucis ante términum,\\
Rerum Creátor, póscimus,\\
Ut pro tua cleméntia\\
Sis præsul et custódia.

Procul recédant sómnia,\\
Et nóctium phantásmata,\\
Hostémque nostrum cómprime,\\
Ne polluántur córpora.

Præsta, Pater piíssime,\\
Patríque compar Unice,\\
Cum Spíritu Paráclito\\
Regnans per omne sæculum.\\
Amen.

\capitulumlabel{Capitulum}{Jer. 14, 9}

\lettrine[lines=2]{T}{U} autem in nobis es, Dómine, et nomen sanctum tuum invocátum est super nos: ne derelínquas nos, Dómine, Deus noster.

\R\ br. In manus tuas, Dómine, * Comméndo spíritum meum. In manus. \V\ Redemísti nos, Dómine, Deus veritátis. Comméndo. Glória Patri. In manus.

\V\ Custódi nos, Dómine, ut pupíllam óculi. \R\ Sub umbra alárum tuárum prótege nos.

\antlabel{Ant.} Salva nos.

\canticumheading{Canticum Simeonis}{Luc. 2, 29-32}

\lettrine[lines=2]{N}{UNC} dimíttis servum tuum, Dómine, * secúndum verbum tuum in pace:

Quia vidérunt óculi mei * salutáre tuum,

Quod parásti * ante fáciem ómnium populórum,

Lumen ad revelatiónem géntium, * et glóriam plebis tuæ Israël.

Glória Patri.

\antlabel{Ant.} Salva nos, Dómine, vigilántes, custódi nos dormiéntes; ut vigilémus cum Christo, et requiescámus in pace.

\variantheading{Preces (si dicendæ sint)}

KYRIE, eléison. Christe, eléison. Kýrie, eléison.

Pater noster \rubric{secreto usque ad} \V\ Et ne nos indúcas in tentatiónem. \R\ Sed líbera nos a malo.

Credo in Deum \rubric{secreto usque ad} \V\ Carnis resurrectiónem. \R\ Vitam ætérnam. Amen. \V\ Benedíctus es, Dómine, Deus patrum nostrórum. \R\ Et laudábilis et gloriósus in sæcula. \V\ Benedicámus Patrem et Fílium cum Sancto Spíritu. \R\ Laudémus, et superexaltémus eum in sæcula. \V\ Benedíctus es, Dómine, in firmaménto cæli. \R\ Et laudábilis, et gloriósus, et superexaltátus in sæcula. \V\ Benedícat et custódiat nos omnípotens et miséricors Dóminus. \R\ Amen. \V\ Dignáre, Dómine, nocte ista. \R\ Sine peccáto nos custodíre. \V\ Miserére nostri, Dómine. \R\ Miserére nostri. \V\ Fiat misericórdia tua, Dómine, super nos. \R\ Quemádmodum sperávimus in te. \V\ Dómine, exáudi oratiónem meam. \R\ Et clamor meus ad te véniat.

\V\ Dóminus vobíscum.

Orémus. \oratiolabel{Oratio}

\lettrine[lines=2]{V}{ISITA}, quæsumus, Dómine, habitatiónem istam, et omnes insídias inimíci ab ea longe repélle: Angeli tui sancti hábitent in ea, qui nos in pace custódiant; et benedíctio tua sit super nos semper. Per Dóminum nostrum. \R\ Amen.

\V\ Dóminus vobíscum.

\V\ Benedicámus Dómino.

\benedictiolabel{Benedictio} Benedícat et custódiat nos omnípotens et miséricors Dóminus, Pater, et Fílius, et Spíritus Sanctus. \R\ Amen.

\variantheading{Antiphona finalis B. M. V.}

\V\ Divínum auxílium.

Pater, Ave et Credo.

% ===== END OF BATCH 18 -- Covers autumnalis.md lines 5981-6274,
% completing "Ad Vesperas" (the Gradual Psalms 122-126, Capitulum,
% hymn "Telluris alme Cónditor") and "Ad Completorium" (Psalms 11,
% 12, 15, hymn "Te lucis ante términum," Canticum Simeonis, Preces),
% and with them the ENTIRE "Feria Tertia" Office. Two more
% occurrences of the recurring stray OCR "\*" artifact (in "Telluris
% alme Cónditor" at autumnalis.md line 6109 and "Te lucis ante
% términum" at line 6217, both at the hymn's final doxology stanza)
% are silently dropped. The partial-bold "Confíteor." (not the full
% "Misereátur. Indulgéntiam.") again appears at Compline (line
% 6137), matching the alternation already documented in Feria
% Secunda's Compline, and is rendered the same way:
% \formulalabel{Confíteor.} followed by plain text. The document
% continues with "Feria Quarta" (autumnalis.md line 6277), the third
% weekday of the Psalterium cycle.
% =====
\feriaheading{Feria Quarta}

\hourheading{Ad Matutinum}

Pater, Ave et Credo.

\V\ Dómine, lábia mea.

\V\ Deus, in adjutórium.

\variantheading{In Officio feriali}

\invitatlabel{Invitat.} Deum magnum Dóminum, * Veníte, adorémus. \psalmlabel{Ps. 94} Veníte.

\hymnusheading{Hymnus}

\lettrine[lines=2]{R}{ERUM} Creátor óptime,\\
Rectórque noster, áspice:\\
Nos a quiéte nóxia\\
Mersos sopóre líbera.

Te, sancte Christe, póscimus,\\
Ignósce culpis ómnibus:\\
Ad confiténdum súrgimus,\\
Morásque noctis rúmpimus.

Mentes manúsque tóllimus,\\
Prophéta sicut nóctibus\\
Nobis geréndum præcipit,\\
Paulúsque gestis cénsuit.

Vides malum, quod fécimus:\\
Occúlta nostra pándimus:\\
Preces geméntes fúndimus,\\
Dimítte quod peccávimus.

Præsta, Pater piíssime,\\
Patríque compar Unice,\\
Cum Spíritu Paráclito\\
Regnans per omne sæculum.\\
Amen.

\nocturnoheading{In I Nocturno}

\antlabel{Ant.} Speciósus forma * præ fíliis hóminum, diffúsa est grátia in lábiis tuis.

\psalmheading{Psalmus 44, i}

\lettrine[lines=2]{E}{RUCTAVIT} cor meum verbum bonum: * dico ego ópera mea Regi.

Lingua mea cálamus scribæ: * velóciter scribéntis.

Speciósus forma præ fíliis hóminum, diffúsa est grátia in lábiis tuis: * proptérea benedíxit te Deus in ætérnum.

Accíngere gládio tuo super femur tuum, * potentíssime.

Spécie tua et pulchritúdine tua * inténde, próspere procéde, et regna.

Propter veritátem, et mansuetúdinem, et justítiam: * et dedúcet te mirabíliter déxtera tua.

Sagíttæ tuæ acútæ, pópuli sub te cadent, * in corda inimicórum Regis.

Sedes tua, Deus, in sæculum sæculi: * virga directiónis virga regni tui.

Dilexísti justítiam, et odísti iniquitátem: * proptérea unxit te Deus, Deus tuus, óleo lætítiæ præ consórtibus tuis.

Myrrha, et gutta, et cásia a vestiméntis tuis, a dómibus ebúrneis: * ex quibus delectavérunt te fíliæ regum in honóre tuo.

Astitit regína a dextris tuis in vestítu deauráto: * circúmdata varietáte.

\antlabel{Ant.} Speciósus forma præ fíliis hóminum, diffúsa est grátia in lábiis tuis.

\antlabel{Ant.} Confitebúntur tibi * pópuli, Deus, in ætérnum.

\psalmheading{Psalmus 44, ii}

\lettrine[lines=2]{A}{UDI}, fília, et vide, et inclína aurem tuam: * et oblivíscere pópulum tuum, et domum patris tui.

Et concupíscet Rex decórem tuum: * quóniam ipse est Dóminus, Deus tuus, et adorábunt eum.

Et fíliæ Tyri in munéribus * vultum tuum deprecabúntur: omnes dívites plebis.

Omnis glória ejus fíliæ Regis ab intus, * in fímbriis áureis circumamícta varietátibus.

Adducéntur Regi vírgines post eam: * próximæ ejus afferéntur tibi.

Afferéntur in lætítia et exsultatióne: * adducéntur in templum Regis.

Pro pátribus tuis nati sunt tibi fílii: * constítues eos príncipes super omnem terram.

Mémores erunt nóminis tui: * in omni generatióne et generatiónem.

Proptérea pópuli confitebúntur tibi in ætérnum: * et in sæculum sæculi.

\antlabel{Ant.} Confitebúntur tibi pópuli, Deus, in ætérnum.

\antlabel{Ant.} Adjútor in tribulatiónibus * Deus noster.

\psalmheading{Psalmus 45}

\lettrine[lines=2]{D}{EUS} noster refúgium, et virtus: * adjútor in tribulatiónibus, quæ invenérunt nos nimis.

Proptérea non timébimus dum turbábitur terra: * et transferéntur montes in cor maris.

Sonuérunt, et turbátæ sunt aquæ eórum: * conturbáti sunt montes in fortitúdine ejus.

Flúminis ímpetus lætíficat civitátem Dei: * sanctificávit tabernáculum suum Altíssimus.

Deus in médio ejus, non commovébitur: * adjuvábit eam Deus mane dilúculo.

Conturbátæ sunt gentes, et inclináta sunt regna: * dedit vocem suam, mota est terra.

Dóminus virtútum nobíscum: * suscéptor noster Deus Jacob.

Veníte, et vidéte ópera Dómini, quæ pósuit prodígia super terram: * áuferens bella usque ad finem terræ.

Arcum cónteret, et confrínget arma: * et scuta combúret igni.

Vacáte, et vidéte quóniam ego sum Deus: * exaltábor in géntibus, et exaltábor in terra.

Dóminus virtútum nobíscum: * suscéptor noster Deus Jacob.

\antlabel{Ant.} Adjútor in tribulatiónibus Deus noster.

\variantheading{In Festis novem Lectionum:}

\V\ Dóminus virtútum nobíscum. \R\ Suscéptor noster Deus Jacob.

\absolutiolabel{Absolutio.} Exáudi, Dómine.

\benedlabel{Bened.} 1 Benedictióne perpétua. 2 Unigénitus Dei Fílius. 3 Spíritus Sancti grátia.

\nocturnoheading{In II Nocturno}

\antlabel{Ant.} Magnus Dóminus, * \textdagger\ et laudábilis nimis in civitáte Dei nostri. \textdaggerdbl

\psalmheading{Psalmus 47}

\lettrine[lines=2]{M}{AGNUS} Dóminus, \textdagger\ et laudábilis nimis * in civitáte Dei nostri, \textdaggerdbl\ in monte sancto ejus.

Fundátur exsultatióne univérsæ terræ mons Sion, * látera Aquilónis, cívitas Regis magni.

Deus in dómibus ejus cognoscétur, * cum suscípiet eam.

Quóniam ecce reges terræ congregáti sunt: * convenérunt in unum.

Ipsi vidéntes sic admiráti sunt, conturbáti sunt, commóti sunt: * tremor apprehéndit eos.

Ibi dolóres ut parturiéntis: * in spíritu vehemménti cónteres naves Tharsis.

Sicut audívimus, sic vídimus in civitáte Dómini virtútum, in civitáte Dei nostri: * Deus fundávit eam in ætérnum.

Suscépimus, Deus, misericórdiam tuam, * in médio templi tui.

Secúndum nomen tuum, Deus, sic et laus tua in fines terræ: * justítia plena est déxtera tua.

Lætétur mons Sion, et exsúltent fíliæ Judæ: * propter judícia tua, Dómine.

Circúmdate Sion, et complectímini eam: * narráte in túrribus ejus.

Pónite corda vestra in virtúte ejus: * et distribúite domos ejus, ut enarrétis in progénie áltera.

Quóniam hic est Deus, Deus noster in ætérnum, et in sæculum sæculi: * ipse reget nos in sæcula.

\antlabel{Ant.} Magnus Dóminus, et laudábilis nimis in civitáte Dei nostri.

\antlabel{Ant.} Os meum loquétur * sapiéntiam: et meditátio cordis mei prudéntiam.

\psalmheading{Psalmus 48, i}

\lettrine[lines=2]{A}{UDITE} hæc, omnes gentes: * áuribus percípite omnes, qui habitátis orbem:

Quique terrígenæ, et fílii hóminum: * simul in unum dives et pauper.

Os meum loquétur sapiéntiam: * et meditátio cordis mei prudéntiam.

Inclinábo in parábolam aurem meam: * apériam in psaltério propositiónem meam.

Cur timébo in die mala? * iníquitas calcánei mei circúmdabit me:

Qui confídunt in virtúte sua: * et in multitúdine divitiárum suárum gloriántur.

Frater non rédimit, rédimet homo: * non dabit Deo placatiónem suam.

Et prétium redemptiónis ánimæ suæ: * et laborábit in ætérnum, et vivet adhuc in finem.

Non vidébit intéritum, cum víderit sapiéntes moriéntes: * simul insípiens, et stultus períbunt.

Et relínquent aliénis divítias suas: * et sepúlcra eórum domus illórum in ætérnum.

Tabernácula eórum in progénie et progénie: * vocavérunt nómina sua in terris suis.

Et homo, cum in honóre esset, non intelléxit: * comparátus est juméntis insipiéntibus, et símilis factus est illis.

\antlabel{Ant.} Os meum loquétur sapiéntiam: et meditátio cordis mei prudéntiam.

\antlabel{Ant.} Ne timúeris: * cum dívite non descéndet in sepúlcrum glória ejus.

\psalmheading{Psalmus 48, ii}

\lettrine[lines=2]{H}{ÆC} via illórum scándalum ipsis: * et póstea in ore suo complacébunt.

Sicut oves in inférno pósiti sunt: * mors depáscet eos.

Et dominabúntur eórum justi in matutíno: * et auxílium eórum veteráscet in inférno a glória eórum.

Verúmtamen Deus rédimet ánimam meam de manu ínferi, * cum accéperit me.

Ne timúeris, cum dives factus fúerit homo: * et cum multiplicáta fúerit glória domus ejus.

Quóniam cum interíerit, non sumet ómnia: * neque descéndet cum eo glória ejus.

Quia ánima ejus in vita ipsíus benedicétur: * confitébitur tibi cum beneféceris ei.

Introíbit usque in progénies patrum suórum: * et usque in ætérnum non vidébit lumen.

Homo, cum in honóre esset, non intelléxit: * comparátus est juméntis insipiéntibus, et símilis factus est illis.

\antlabel{Ant.} Ne timúeris: cum dívite non descéndet in sepúlcrum glória ejus.

\variantheading{In Festis novem Lectionum:}

\V\ Deus rédimet ánimam meam de manu ínferi. \R\ Cum accéperit me.

\absolutiolabel{Absolutio.} Ipsíus píetas et.

\benedlabel{Bened.} 4 Deus Pater omnípotens. 5 Christus perpétuæ. 6 Ignem sui amóris.

\nocturnoheading{In III Nocturno}

\variantheading{I}

\rubric{In Fériis extra Vigílias commúnes occurréntibus, et in Festis, cum ad Laudes non dícitur Psalmus 50 Miserére.}

\antlabel{Ant.} Deus deórum, * \textdagger\ Dóminus locútus est. \textdaggerdbl

\psalmheading{Psalmus 49, i}

\lettrine[lines=2]{D}{EUS} deórum, \textdagger\ Dóminus locútus est: * \textdaggerdbl\ et vocávit terram.

A solis ortu usque ad occásum: * ex Sion spécies decóris ejus.

Deus maniféste véniet: * Deus noster et non silébit.

Ignis in conspéctu ejus exardéscet: * et in circúitu ejus tempéstas válida.

Advocábit cælum desúrsum: * et terram discérnere pópulum suum.

Congregáte illi sanctos ejus: * qui órdinant testaméntum ejus super sacrifícia.

Et annuntiábunt cæli justítiam ejus: * quóniam Deus judex est.

Audi, pópulus meus, et loquar: Israël, et testificábor tibi: * Deus, Deus tuus ego sum.

Non in sacrifíciis tuis árguam te: * holocáusta autem tua in conspéctu meo sunt semper.

Non accípiam de domo tua vítulos: * neque de grégibus tuis hircos.

Quóniam meæ sunt omnes feræ silvárum, * juménta in móntibus et boves.

Cognóvi ómnia volatília cæli: * et pulchritúdo agri mecum est.

Si esuríero, non dicam tibi: * meus est enim orbis terræ, et plenitúdo ejus.

Numquid manducábo carnes taurórum? * aut sánguinem hircórum potábo?

Immola Deo sacrifícium laudis: * et redde Altíssimo vota tua.

Et ínvoca me in die tribulatiónis: * éruam te, et honorificábis me.

\antlabel{Ant.} Deus deórum, Dóminus locútus est.

\antlabel{Ant.} Intellígite, * qui obliviscímini Deum.

\psalmheading{Psalmus 49, ii}

\lettrine[lines=2]{P}{ECCATORI} autem dixit Deus: * Quare tu enárras justítias meas, et assúmis testaméntum meum per os tuum?

Tu vero odísti disciplínam: * et projecísti sermónes meos retrórsum:

Si vidébas furem, currébas cum eo: * et cum adúlteris portiónem tuam ponébas.

Os tuum abundávit malítia: * et lingua tua concinnábat dolos.

Sedens advérsus fratrem tuum loquebáris, et advérsus fílium matris tuæ ponébas scándalum: * hæc fecísti, et tácui.

Existimásti, iníque, quod ero tui símilis: * árguam te, et státuam contra fáciem tuam.

Intellígite hæc, qui oblivíscimini Deum: * nequándo rápiat, et non sit qui erípiat.

Sacrifícium laudis honorificábit me: * et illic iter, quo osténdam illi salutáre Dei.

\antlabel{Ant.} Intellígite, qui obliviscímini Deum.

\antlabel{Ant.} Acceptábis sacrifícium * justítiæ super altáre tuum, Dómine.

\psalmheading{Psalmus 50}

\lettrine[lines=2]{M}{ISERERE} mei, Deus, * secúndum magnam misericórdiam tuam.

Et secúndum multitúdinem miseratiónum tuárum, * dele iniquitátem meam.

Amplius lava me ab iniquitáte mea: * et a peccáto meo munda me.

Quóniam iniquitátem meam ego cognósco: * et peccátum meum contra me est semper.

Tibi soli peccávi, et malum coram te feci: * ut justificéris in sermónibus tuis, et vincas cum judicáris.

Ecce enim, in iniquitátibus concéptus sum: * et in peccátis concépit me mater mea.

Ecce enim, veritátem dilexísti: * incérta et occúlta sapiéntiæ tuæ manifestásti mihi.

Aspérges me hyssópo, et mundábor: * lavábis me, et super nivem dealbábor.

Audítui meo dabis gáudium et lætítiam: * et exsultábunt ossa humiliáta.

Avérte fáciem tuam a peccátis meis: * et omnes iniquitátes meas dele.

Cor mundum crea in me, Deus: * et spíritum rectum ínnova in viscéribus meis.

Ne projícias me a fácie tua: * et spíritum sanctum tuum ne áuferas a me.

Redde mihi lætítiam salutáris tui: * et spíritu principáli confírma me.

Docébo iníquos vias tuas: * et ímpii ad te converténtur.

Líbera me de sanguínibus, Deus, Deus salútis meæ: * et exsultábit lingua mea justítiam tuam.

Dómine, lábia mea apéries: * et os meum annuntiábit laudem tuam.

Quóniam si voluísses sacrifícium, dedíssem útique: * holocáustis non delectáberis.

Sacrifícium Deo spíritus contribulátus: * cor contrítum, et humiliátum, Deus, non despícies.

Benígne fac, Dómine, in bona voluntáte tua Sion: * ut ædificéntur muri Jerúsalem.

Tunc acceptábis sacrifícium justítiæ, oblatiónes, et holocáusta: * tunc impónent super altáre tuum vítulos.

\antlabel{Ant.} Acceptábis sacrifícium justítiæ super altáre tuum, Dómine.

\variantheading{In feriali Officio et in Festis:}

\V\ Audítui meo dabis gáudium et lætítiam. \R\ Et exsultábunt ossa humiliáta.

\variantheading{In Officio novem Lectionum}

\absolutiolabel{Absolutio.} A vínculis.

\benedlabel{Bened.} 7 Evangélica léctio. 8 In Festis Dómini: Divínum auxílium. In Festis B. M. V.: Cujus festum cólimus, ipsa Virgo. In Festis Sanctórum: Cujus (\textit{vel} Quorum \textit{aut} Quarum) festum cólimus. 9 Ad societátem. (Per evangélica dicta.)

\variantheading{In Officio trium Lectionum}

\absolutiolabel{Absolutio.} A vínculis.

\variantheading{In Feriis cum Homilia:}

\benedlabel{Bened.} 1 Evangélica. 2 Divínum. 3 Ad societátem.

\variantheading{In Feriis sine Homilia:}

\benedlabel{Bened.} 1 Ille nos benedícat. 2 Divínum. 3 Ad societátem.

\variantheading{In Festis Sanctorum:}

\benedlabel{Bened.} 1 Ille nos benedícat, qui sine fine vivit et regnat. \R\ Amen. 2 Cujus (\textit{vel} Quorum \textit{aut} Quarum). 3 Ad societátem.

\nocturnoheading{In III Nocturno}

\variantheading{II}

\rubric{In Vigíliis commúnibus, cum ad Laudes dícitur Psalmus 50 Miserére.}

\antlabel{Ant.} Deus deórum. \textdagger

\psalmheading{Psalmus 49, i}

\lettrine[lines=2]{D}{EUS} deórum, \textdagger\ Dóminus locútus est: * \textdaggerdbl\ et vocávit terram.

A solis ortu usque ad occásum: * ex Sion spécies decóris ejus.

Deus maniféste véniet: * Deus noster et non silébit.

Ignis in conspéctu ejus exardéscet: * et in circúitu ejus tempéstas válida.

Advocábit cælum desúrsum: * et terram discérnere pópulum suum.

Congregáte illi sanctos ejus: * qui órdinant testaméntum ejus super sacrifícia.

Et annuntiábunt cæli justítiam ejus: * quóniam Deus judex est.

Audi, pópulus meus, et loquar: Israël, et testificábor tibi: * Deus, Deus tuus ego sum.

Non in sacrifíciis tuis árguam te: * holocáusta autem tua in conspéctu meo sunt semper.

Non accípiam de domo tua vítulos: * neque de grégibus tuis hircos.

Quóniam meæ sunt omnes feræ silvárum, * juménta in móntibus et boves.

Cognóvi ómnia volatília cæli: * et pulchritúdo agri mecum est.

Si esuríero, non dicam tibi: * meus est enim orbis terræ, et plenitúdo ejus.

Numquid manducábo carnes taurórum? * aut sánguinem hircórum potábo?

Immola Deo sacrifícium laudis: * et redde Altíssimo vota tua.

Et ínvoca me in die tribulatiónis: * éruam te, et honorificábis me.

\antlabel{Ant.} Deus deórum, Dóminus locútus est.

\antlabel{Ant.} Immola Deo.

\psalmheading{Psalmus 49, ii}

\lettrine[lines=2]{P}{ECCATORI} autem dixit Deus: * Quare tu enárras justítias meas, et assúmis testaméntum meum per os tuum?

Tu vero odísti disciplínam: * et projecísti sermónes meos retrórsum:

Si vidébas furem, currébas cum eo: * et cum adúlteris portiónem tuam ponébas.

Os tuum abundávit malítia: * et lingua tua concinnábat dolos.

Sedens advérsus fratrem tuum loquebáris, et advérsus fílium matris tuæ ponébas scándalum: * hæc fecísti, et tácui.

Existimásti, iníque, quod ero tui símilis: * árguam te, et státuam contra fáciem tuam.

Intellígite hæc, qui oblivíscimini Deum: * nequándo rápiat, et non sit qui erípiat.

Sacrifícium laudis honorificábit me: * et illic iter, quo osténdam illi salutáre Dei.

\antlabel{Ant.} Immola Deo sacrifícium laudis.

\antlabel{Ant.} Intellígite.

\psalmheading{Psalmus 49, iii}

\lettrine[lines=2]{P}{ECCATORI} autem dixit Deus: * Quare tu enárras justítias meas, et assúmis testaméntum meum per os tuum?

Tu vero odísti disciplínam: * et projecísti sermónes meos retrórsum:

Si vidébas furem, currébas cum eo: * et cum adúlteris portiónem tuam ponébas.

Os tuum abundávit malítia: * et lingua tua concinnábat dolos.

Sedens advérsus fratrem tuum loquebáris, et advérsus fílium matris tuæ ponébas scándalum: * hæc fecísti, et tácui.

Existimásti, iníque, quod ero tui símilis: * árguam te, et státuam contra fáciem tuam.

Intellígite hæc, qui oblivíscimini Deum: * nequándo rápiat, et non sit qui erípiat.

Sacrifícium laudis honorificábit me: * et illic iter, quo osténdam illi salutáre Dei.

\antlabel{Ant.} Intellígite, qui obliviscímini Deum.

\V\ Sacrifícium laudis honorificábit me. \R\ Et illic iter, quo osténdam illi salutáre Dei.

\absolutiolabel{Absolutio.} A vínculis.

\variantheading{In Feriis cum Homilia:}

\benedlabel{Bened.} 1 Evangélica. 2 Divínum. 3 Ad societátem.

\variantheading{In Feriis sine Homilia:}

\benedlabel{Bened.} 1 Ille nos benedícat. 2 Divínum. 3 Ad societátem.

% ===== END OF BATCH 19 -- Covers autumnalis.md lines 6277-6743,
% opening "Feria Quarta" (Wednesday) and completing its "Ad
% Matutinum" in full: the ferial Invitatorium "Deum magnum
% Dóminum," the hymn "Rerum Creátor óptime" (stray OCR "\*"
% silently dropped), and the three Nocturns (Psalms 44/i, 44/ii,
% 45; 47, 48/i, 48/ii). "In III Nocturno" here introduces, for the
% first time in this volume, two complete alternate forms labeled
% "I" and "II," each with its own italicized rubric describing when
% it applies (Form I when Ps. 50 was NOT said at Lauds, using
% Psalms 49/i, 49/ii, and the full Ps. 50; Form II when Ps. 50 WAS
% said at Lauds, using Psalms 49/i, 49/ii, and 49/iii instead) --
% rendered with \variantheading{I} and \variantheading{II}, matching
% the identical command already used for this exact structure in
% both verna.tex and hiemalis.tex. Notably, Form II's "Psalmus 49,
% ii" and "Psalmus 49, iii" give verbatim identical text in
% autumnalis.md itself (both "Peccátori autem dixit Deus..."); this is
% transcribed exactly as the source repeats it, without alteration,
% since it is a structural feature of the source and not a
% character-level OCR defect. The document continues with "Ad
% Laudes I" (autumnalis.md line 6745).
% =====
\hourheading{Ad Laudes I}

\rubric{In Fériis extra Vigílias commúnes occurréntibus, atque in Festis.}

(Pater noster et Ave.)

\V\ Deus, in adjutórium.

\antlabel{Ant.} Dóminus regnávit, * \textdagger\ exsúltet terra. \textdaggerdbl

\psalmheading{Psalmus 96}

\lettrine[lines=2]{D}{OMINUS} regnávit, \textdagger\ exsúltet terra: * \textdaggerdbl\ læténtur ínsulæ multæ.

Nubes, et calígo in circúitu ejus: * justítia, et judícium corréctio sedis ejus.

Ignis ante ipsum præcédet, * et inflammábit in circúitu inimícos ejus.

Illuxérunt fúlgura ejus orbi terræ: * vidit, et commóta est terra.

Montes, sicut cera fluxérunt a fácie Dómini: * a fácie Dómini omnis terra.

Annuntiavérunt cæli justítiam ejus: * et vidérunt omnes pópuli glóriam ejus.

Confundántur omnes, qui adórant sculptília: * et qui gloriántur in simulácris suis.

Adoráte eum, omnes Angeli ejus: * audívit, et lætáta est Sion.

Et exsultavérunt fíliæ Judæ, * propter judícia tua, Dómine:

Quóniam tu Dóminus Altíssimus super omnem terram: * nimis exaltátus es super omnes deos.

Qui dilígitis Dóminum, odíte malum: * custódit Dóminus ánimas sanctórum suórum, de manu peccatóris liberábit eos.

Lux orta est justo, * et rectis corde lætítia.

Lætámini, justi, in Dómino: * et confitémini memóriæ sanctificatiónis ejus.

\antlabel{Ant.} Dóminus regnávit, exsúltet terra.

\antlabel{Ant.} Te decet hymnus, * \textdagger\ Deus, in Sion. \textdaggerdbl

\psalmheading{Psalmus 64}

\lettrine[lines=2]{T}{E} decet hymnus, \textdagger\ Deus, in Sion: * \textdaggerdbl\ et tibi reddétur votum in Jerúsalem.

Exáudi oratiónem meam: * ad te omnis caro véniet.

Verba iniquórum prævaluérunt super nos: * et impietátibus nostris tu propitiáberis.

Beátus, quem elegísti, et assumpsísti: * inhabitábit in átriis tuis.

Replébimur in bonis domus tuæ: * sanctum est templum tuum, mirábile in æquitáte.

Exáudi nos, Deus, salutáris noster, * spes ómnium fínium terræ, et in mari longe.

Præparans montes in virtúte tua, accínctus poténtia: * qui contúrbas profúndum maris sonum flúctuum ejus.

Turbabúntur gentes, et timébunt qui hábitant términos a signis tuis: * éxitus matutíni, et véspere delectábis.

Visitásti terram, et inebriásti eam: * multiplicásti locupletáre eam.

Flumen Dei replétum est aquis, parásti cibum illórum: * quóniam ita est præparátio ejus.

Rivos ejus inébria, multíplica genímina ejus: * in stillicídiis ejus lætábitur gérminans.

Benedíces corónæ anni benignitátis tuæ: * et campi tui replebúntur ubertáte.

Pinguéscent speciósa desérti: * et exsultatióne colles accingéntur.

Indúti sunt aríetes óvium, et valles abundábunt fruménto: * clamábunt, étenim hymnum dicent.

\antlabel{Ant.} Te decet hymnus, Deus, in Sion.

\antlabel{Ant.} Tibi, Dómine, psallam, * et intélligam in via immaculáta.

\psalmheading{Psalmus 100}

\lettrine[lines=2]{M}{ISERICORDIAM}, et judícium * cantábo tibi, Dómine:

Psallam, et intélligam in via immaculáta, * quando vénies ad me.

Perambulábam in innocéntia cordis mei, * in médio domus meæ.

Non proponébam ante óculos meos rem injústam: * faciéntes prævaricatiónes odívi.

Non adhæsit mihi cor pravum: * declinántem a me malígnum non cognoscébam.

Detrahéntem secréto próximo suo, * hunc persequébar.

Supérbo óculo, et insatiábili corde, * cum hoc non edébam.

Oculi mei ad fidéles terræ ut sédeant mecum: * ámbulans in via immaculáta, hic mihi ministrábat.

Non habitábit in médio domus meæ qui facit supérbiam: * qui lóquitur iníqua, non diréxit in conspéctu oculórum meórum.

In matutíno interficiébam omnes peccatóres terræ: * ut dispérderem de civitáte Dómini omnes operántes iniquitátem.

\antlabel{Ant.} Tibi, Dómine, psallam, et intélligam in via immaculáta.

\antlabel{Ant.} Dómine, magnus es tu, * et præclárus in virtúte tua.

\canticumheading{Canticum Judith}{Judith. 16, 15-21}

\lettrine[lines=2]{H}{YMNUM} cantémus Dómino, * hymnum novum cantémus Deo nostro.

Adonái, Dómine, magnus es tu, et præclárus in virtúte tua, * et quem superáre nemo potest.

Tibi sérviat omnis creatúra tua: * quia dixísti, et facta sunt:

Misísti spíritum tuum, et creáta sunt: * et non est qui resístat voci tuæ.

Montes a fundaméntis movebúntur cum aquis: * petræ, sicut cera, liquéscent ante fáciem tuam.

Qui autem timent te, * magni erunt apud te per ómnia.

Væ genti insurgénti super genus meum: Dóminus enim omnípotens vindicábit in eis, * in die judícii visitábit illos.

Dabit enim ignem, et vermes in carnes eórum, * ut urántur, et séntiant usque in sempitérnum.

\antlabel{Ant.} Dómine, magnus es tu, et præclárus in virtúte tua.

\antlabel{Ant.} Laudábo Deum meum * in vita mea.

\psalmheading{Psalmus 145}

\lettrine[lines=2]{L}{AUDA}, ánima mea, Dóminum, laudábo Dóminum in vita mea: * psallam Deo meo quámdiu fúero.

Nolíte confídere in princípibus: * in fíliis hóminum, in quibus non est salus.

Exíbit spíritus ejus, et revertétur in terram suam: * in illa die períbunt omnes cogitatiónes eórum.

Beátus, cujus Deus Jacob adjútor ejus, spes ejus in Dómino, Deo ipsíus: * qui fecit cælum et terram, mare, et ómnia, quæ in eis sunt.

Qui custódit veritátem in sæculum, facit judícium injúriam patiéntibus: * dat escam esuriéntibus.

Dóminus solvit compedítos: * Dóminus illúminat cæcos.

Dóminus érigit elísos, * Dóminus díligit justos.

Dóminus custódit ádvenas, pupíllum et víduam suscípiet: * et vias peccatórum dispérdet.

Regnábit Dóminus in sæcula, Deus tuus, Sion, * in generatiónem et generatiónem.

\antlabel{Ant.} Laudábo Deum meum in vita mea.

\variantheading{In Officio feriali}

\capitulumlabel{Capit.}{Rom. 13, 12-13}

\lettrine[lines=2]{N}{OX} præcéssit, dies autem appropinquávit. Abjiciámus ergo ópera tenebrárum, et induámur arma lucis. Sicut in die honéste ambulémus.

\hymnusheading{Hymnus}

\lettrine[lines=2]{N}{OX}, et ténebræ, et núbila,\\
Confúsa mundi et túrbida:\\
Lux intrat, albéscit polus:\\
Christus venit: discédite.

Calígo terræ scínditur\\
Percússa solis spículo,\\
Rebúsque jam color redit,\\
Vultu niténtis síderis.

Te, Christe, solum nóvimus:\\
Te mente pura et símplici,\\
Flendo et canéndo quǽsumus,\\
Inténde nostris sénsibus.

Sunt multa fucis illíta,\\
Quæ luce purgéntur tua:\\
Tu, vera lux Cæléstium,\\
Vultu seréno illúmina.

Deo Patri sit glória,\\
Ejúsque soli Fílio,\\
Cum Spíritu Paráclito,\\
Nunc et per omne sæculum.\\
Amen.

\V\ Repléti sumus mane misericórdia tua. \R\ Exsultávimus, et delectáti sumus.

\antlabel{Ad Bened. Ant.} De manu ómnium, * qui odérunt nos, liberávit nos Dóminus.

\rubric{Réliqua ut in Ordinário. 11}

\hourheading{Ad Laudes II}

\rubric{In Vigíliis commúnibus, quando fit Offícium de Féria.}

(Pater noster et Ave.)

\V\ Deus, in adjutórium.

\antlabel{Ant.} Amplius lava me.

\psalmheading{Psalmus 50}

\lettrine[lines=2]{M}{ISERERE} mei, Deus, * secúndum magnam misericórdiam tuam.

Et secúndum multitúdinem miseratiónum tuárum, * dele iniquitátem meam.

Amplius lava me ab iniquitáte mea: * et a peccáto meo munda me.

Quóniam iniquitátem meam ego cognósco: * et peccátum meum contra me est semper.

Tibi soli peccávi, et malum coram te feci: * ut justificéris in sermónibus tuis, et vincas cum judicáris.

Ecce enim, in iniquitátibus concéptus sum: * et in peccátis concépit me mater mea.

Ecce enim, veritátem dilexísti: * incérta et occúlta sapiéntiæ tuæ manifestásti mihi.

Aspérges me hyssópo, et mundábor: * lavábis me, et super nivem dealbábor.

Audítui meo dabis gáudium et lætítiam: * et exsultábunt ossa humiliáta.

Avérte fáciem tuam a peccátis meis: * et omnes iniquitátes meas dele.

Cor mundum crea in me, Deus: * et spíritum rectum ínnova in viscéribus meis.

Ne projícias me a fácie tua: * et spíritum sanctum tuum ne áuferas a me.

Redde mihi lætítiam salutáris tui: * et spíritu principáli confírma me.

Docébo iníquos vias tuas: * et ímpii ad te converténtur.

Líbera me de sanguínibus, Deus, Deus salútis meæ: * et exsultábit lingua mea justítiam tuam.

Dómine, lábia mea apéries: * et os meum annuntiábit laudem tuam.

Quóniam si voluísses sacrifícium, dedíssem útique: * holocáustis non delectáberis.

Sacrifícium Deo spíritus contribulátus: * cor contrítum, et humiliátum, Deus, non despícies.

Benígne fac, Dómine, in bona voluntáte tua Sion: * ut ædificéntur muri Jerúsalem.

Tunc acceptábis sacrifícium justítiæ, oblatiónes, et holocáusta: * tunc impónent super altáre tuum vítulos.

\antlabel{Ant.} Amplius lava me, Dómine, ab injustítia mea.

\antlabel{Ant.} Impietátibus nostris.

\psalmheading{Psalmus 64}

\lettrine[lines=2]{T}{E} decet hymnus, Deus, in Sion: * et tibi reddétur votum in Jerúsalem.

Exáudi oratiónem meam: * ad te omnis caro véniet.

Verba iniquórum prævaluérunt super nos: * et impietátibus nostris tu propitiáberis.

Beátus, quem elegísti, et assumpsísti: * inhabitábit in átriis tuis.

Replébimur in bonis domus tuæ: * sanctum est templum tuum, mirábile in æquitáte.

Exáudi nos, Deus, salutáris noster, * spes ómnium fínium terræ, et in mari longe.

Præparans montes in virtúte tua, accínctus poténtia: * qui contúrbas profúndum maris sonum flúctuum ejus.

Turbabúntur gentes, et timébunt qui hábitant términos a signis tuis: * éxitus matutíni, et véspere delectábis.

Visitásti terram, et inebriásti eam: * multiplicásti locupletáre eam.

Flumen Dei replétum est aquis, parásti cibum illórum: * quóniam ita est præparátio ejus.

Rivos ejus inébria, multíplica genímina ejus: * in stillicídiis ejus lætábitur gérminans.

Benedíces corónæ anni benignitátis tuæ: * et campi tui replebúntur ubertáte.

Pinguéscent speciósa desérti: * et exsultatióne colles accingéntur.

Indúti sunt aríetes óvium, et valles abundábunt fruménto: * clamábunt, étenim hymnum dicent.

\antlabel{Ant.} Impietátibus nostris tu propitiáberis, Deus.

\antlabel{Ant.} In innocéntia.

\psalmheading{Psalmus 100}

\lettrine[lines=2]{M}{ISERICORDIAM}, et judícium * cantábo tibi, Dómine:

Psallam, et intélligam in via immaculáta, * quando vénies ad me.

Perambulábam in innocéntia cordis mei, * in médio domus meæ.

Non proponébam ante óculos meos rem injústam: * faciéntes prævaricatiónes odívi.

Non adhæsit mihi cor pravum: * declinántem a me malígnum non cognoscébam.

Detrahéntem secréto próximo suo, * hunc persequébar.

Supérbo óculo, et insatiábili corde, * cum hoc non edébam.

Oculi mei ad fidéles terræ ut sédeant mecum: * ámbulans in via immaculáta, hic mihi ministrábat.

Non habitábit in médio domus meæ qui facit supérbiam: * qui lóquitur iníqua, non diréxit in conspéctu oculórum meórum.

In matutíno interficiébam omnes peccatóres terræ: * ut dispérderem de civitáte Dómini omnes operántes iniquitátem.

\antlabel{Ant.} In innocéntia cordis mei perambulábo, Dómine.

\antlabel{Ant.} Exsultávit cor meum. \textdagger

\canticumheading{Canticum Annæ}{1 Reg. 2, 1-10}

\lettrine[lines=2]{E}{XSULTAVIT} cor meum \textdagger\ in Dómino, * et exaltátum est cornu meum in Deo meo.

Dilatátum est os meum super inimícos meos: * quia lætáta sum in salutári tuo.

Non est sanctus, ut est Dóminus: neque enim est álius extra te, * et non est fortis sicut Deus noster.

Nolíte multiplicáre loqui sublímia, * gloriántes:

Recédant vétera de ore vestro: quia Deus sciéntiárum, Dóminus est, * et ipsi præparántur cogitatiónes.

Arcus fórtium superátus est, * et infírmi accíncti sunt róbore.

Repléti prius, pro pánibus se locavérunt: * et famélici saturáti sunt.

Donec stérilis péperit plúrimos: * et quæ multos habébat fílios, infirmáta est.

Dóminus mortíficat et vivíficat, * dedúcit ad ínferos et redúcit.

Dóminus páuperem facit et ditat, * humíliat et súblevat.

Súscitat de púlvere egénum, * et de stércore élevat páuperem:

Ut sédeat cum princípibus, * et sólium glóriæ téneat.

Dómini enim sunt cárdines terræ, * et pósuit super eos orbem.

Pedes sanctórum suórum servábit, et ímpii in ténebris conticéscent: * quia non in fortitúdine sua roborábitur vir.

Dóminum formidábunt adversárii ejus: * et super ipsos in cælis tonábit:

Dóminus judicábit fines terræ, et dabit impérium regi suo, * et sublimábit cornu Christi sui.

\antlabel{Ant.} Exsultávit cor meum in Dómino, qui humíliat et súblevat.

\antlabel{Ant.} Lauda. \textdagger

\psalmheading{Psalmus 145}

\lettrine[lines=2]{L}{AUDA}, \textdagger\ ánima mea, Dóminum, laudábo Dóminum in vita mea: * psallam Deo meo quámdiu fúero.

Nolíte confídere in princípibus: * in fíliis hóminum, in quibus non est salus.

Exíbit spíritus ejus, et revertétur in terram suam: * in illa die períbunt omnes cogitatiónes eórum.

Beátus, cujus Deus Jacob adjútor ejus, spes ejus in Dómino, Deo ipsíus: * qui fecit cælum et terram, mare, et ómnia, quæ in eis sunt.

Qui custódit veritátem in sæculum, facit judícium injúriam patiéntibus: * dat escam esuriéntibus.

Dóminus solvit compedítos: * Dóminus illúminat cæcos.

Dóminus érigit elísos, * Dóminus díligit justos.

Dóminus custódit ádvenas, pupíllum et víduam suscípiet: * et vias peccatórum dispérdet.

Regnábit Dóminus in sæcula, Deus tuus, Sion, * in generatiónem et generatiónem.

\antlabel{Ant.} Lauda, ánima mea, Dóminum, qui érigit elísos, et díligit justos.

\capitulumlabel{Capit.}{Rom. 13, 12-13}

\lettrine[lines=2]{N}{OX} præcéssit, dies autem appropinquávit. Abjiciámus ergo ópera tenebrárum, et induámur arma lucis. Sicut in die honéste ambulémus.

\hymnusheading{Hymnus}

\lettrine[lines=2]{N}{OX}, et ténebræ, et núbila,\\
Confúsa mundi et túrbida:\\
Lux intrat, albéscit polus:\\
Christus venit: discédite.

Calígo terræ scínditur\\
Percússa solis spículo,\\
Rebúsque jam color redit,\\
Vultu niténtis síderis.

Te, Christe, solum nóvimus:\\
Te mente pura et símplici,\\
Flendo et canéndo quǽsumus,\\
Inténde nostris sénsibus.

Sunt multa fucis illíta,\\
Quæ luce purgéntur tua:\\
Tu, vera lux Cæléstium,\\
Vultu seréno illúmina.

Deo Patri sit glória,\\
Ejúsque soli Fílio,\\
Cum Spíritu Paráclito,\\
Nunc et per omne sæculum.\\
Amen.

\V\ Repléti sumus mane misericórdia tua. \R\ Exsultávimus, et delectáti sumus.

\antlabel{Ad Bened. Ant.} De manu ómnium, * qui odérunt nos, liberávit nos Dóminus.

\rubric{Expléto Cántico ac repetíta Antíphona, dicúntur sequéntes}

\variantheading{Preces feriales}

KYRIE, eléison. Christe, eléison. Kýrie, eléison.

Pater noster \rubric{quod dícitur a solo Hebdomadário totum clara voce, usque ad} \V\ Et ne nos indúcas in tentatiónem. \R\ Sed líbera nos a malo. \V\ Ego dixi: Dómine, miserére mei. \R\ Sana ánimam meam, quia peccávi tibi. \V\ Convértere, Dómine, úsquequo? \R\ Et deprecábilis esto super servos tuos. \V\ Fiat misericórdia tua, Dómine, super nos. \R\ Quemádmodum sperávimus in te. \V\ Sacerdótes tui induántur justítiam. \R\ Et sancti tui exsúltent. \V\ Orémus pro beatíssimo Papa nostro N. \R\ Dóminus consérvet eum, et vivíficet eum, et beátum fáciat eum in terra, et non tradat eum in ánimam inimicórum ejus. \V\ Orémus et pro Antístite nostro N. \R\ Stet et pascat in fortitúdine tua, Dómine, in sublimitáte nóminis tui. \V\ Dómine, salvum fac regem. \R\ Et exáudi nos in die, qua invocavérimus te. \V\ Salvum fac pópulum tuum, Dómine, et bénedic hereditáti tuæ. \R\ Et rege eos, et extólle illos usque in ætérnum. \V\ Meménto Congregatiónis tuæ. \R\ Quam possedísti ab inítio. \V\ Fiat pax in virtúte tua. \R\ Et abundántia in túrribus tuis. \V\ Orémus pro benefactóribus nostris. \R\ Retribúere dignáre, Dómine, ómnibus, nobis bona faciéntibus propter nomen tuum, vitam ætérnam. Amen. \V\ Orémus pro fidélibus defúnctis. \R\ Réquiem ætérnam dona eis, Dómine, et lux perpétua lúceat eis. \V\ Requiéscant in pace. \R\ Amen. \V\ Pro frátribus nostris abséntibus. \R\ Salvos fac servos tuos, Deus meus, sperántes in te. \V\ Pro afflíctis et captívis. \R\ Líbera eos, Deus Israël, ex ómnibus tribulatiónibus suis. \V\ Mitte eis, Dómine, auxílium de sancto. \R\ Et de Sion tuére eos. \V\ Dómine, Deus virtútum, convérte nos. \R\ Et osténde fáciem tuam, et salvi érimus. \V\ Exsúrge, Christe, ádjuva nos. \R\ Et líbera nos propter nomen tuum. \V\ Dómine, exáudi oratiónem meam. \R\ Et clamor meus ad te véniat.

\rubric{Réliqua ut in Ordinário. 11}

% ===== END OF BATCH 20 -- Covers autumnalis.md lines 6745-7156, the
% two alternate ferial Lauds schemes for Feria Quarta, "Ad Laudes
% I" (Psalms 96, 64, 100, Canticum Judith, 145) and "Ad Laudes II"
% (Psalms 50, 64, 100, Canticum Annæ, 145), each with the Capitulum
% "Nox præcéssit" and the hymn "Nox, et ténebræ, et núbila." Notably,
% Psalm 64's incipit carries the dagger and double-dagger
% chant-pointing marks in the Ad Laudes I copy (autumnalis.md line
% 6834: "Te decet hymnus, † Deus, in Sion: * ‡...") but NOT in the
% Ad Laudes II copy of the identical psalm (line 6991, plain "Te
% decet hymnus, Deus, in Sion: *..."); each is transcribed exactly
% as its own occurrence gives it, confirming this is a genuine
% source divergence rather than an OCR gap. Two antiphon incipits in
% Ad Laudes II ("Impietátibus nostris." and "In innocéntia.") carry
% no asterisk at all in the source and are transcribed as bare
% unmarked incipits. Ad Laudes II additionally carries the "Preces
% feriales" litany in full. The document continues with "Ad Primam"
% (autumnalis.md line 7159).
% =====
\hourheading{Ad Primam}

Pater, Ave et Credo.

\V\ Deus, in adjutórium.

\hymnusheading{Hymnus}

\lettrine[lines=2]{J}{AM} lucis orto sídere,\\
Deum precémur súpplices,\\
Ut in diúrnis áctibus\\
Nos servet a nocéntibus.

Linguam refrénans témperet,\\
Ne litis horror ínsonet:\\
Visum fovéndo cóntegat,\\
Ne vanitátes háuriat.

Sint pura cordis íntima,\\
Absístat et vecórdia:\\
Carnis terat supérbiam\\
Potus cibíque párcitas.

Ut, cum dies abscésserit\\
Noctémque sors redúxerit,\\
Mundi per abstinéntiam\\
Ipsi canámus glóriam.

Deo Patri sit glória,\\
Ejúsque soli Fílio,\\
Cum Spíritu Paráclito,\\
Nunc et per omne sæculum.\\
Amen.

\antlabel{Ant.} Misericórdia tua.

\psalmheading{Psalmus 25}

\lettrine[lines=2]{J}{UDICA} me, Dómine, quóniam ego in innocéntia mea ingréssus sum: * et in Dómino sperans non infirmábor.

Proba me, Dómine, et tenta me: * ure renes meos et cor meum.

Quóniam misericórdia tua ante óculos meos est: * et complácui in veritáte tua.

Non sedi cum concílio vanitátis: * et cum iníqua geréntibus non introíbo.

Odívi ecclésiam malignántium: * et cum ímpiis non sedébo.

Lavábo inter innocéntes manus meas: * et circúmdabo altáre tuum, Dómine:

Ut áudiam vocem laudis, * et enárrem univérsa mirabília tua.

Dómine, diléxi decórem domus tuæ, * et locum habitatiónis glóriæ tuæ.

Ne perdas cum ímpiis, Deus, ánimam meam, * et cum viris sánguinum vitam meam:

In quorum mánibus iniquitátes sunt: * déxtera eórum repléta est munéribus.

Ego autem in innocéntia mea ingréssus sum: * rédime me, et miserére mei.

Pes meus stetit in directo: * in ecclésiis benedícam te, Dómine.

\psalmheading{Psalmus 51}

\lettrine[lines=2]{Q}{UID} gloriáris in malítia, * qui potens es in iniquitáte?

Tota die injustítiam cogitávit lingua tua: * sicut novácula acúta fecísti dolum.

Dilexísti malítiam super benignitátem: * iniquitátem magis quam loqui æquitátem.

Dilexísti ómnia verba præcipitatiónis, * lingua dolósa.

Proptérea Deus déstruet te in finem, * evéllet te, et emigrábit te de tabernáculo tuo: et radícem tuam de terra vivéntium.

Vidébunt justi, et timébunt, et super eum ridébunt, et dicent: * Ecce homo, qui non pósuit Deum adjutórem suum:

Sed sperávit in multitúdine divitiárum suárum: * et prævalúit in vanitáte sua.

Ego autem, sicut olíva fructífera in domo Dei, * sperávi in misericórdia Dei in ætérnum: et in séculum séculi.

Confitébor tibi in séculum, quia fecísti: * et exspectábo nomen tuum, quóniam bonum est in conspéctu sanctórum tuórum.

\psalmheading{Psalmus 52}

\lettrine[lines=2]{D}{IXIT} insípiens in corde suo: * Non est Deus.

Corrúpti sunt, et abominábiles facti sunt in iniquitátibus: * non est qui fáciat bonum.

Deus de cælo prospéxit super fílios hóminum: * ut vídeas si est intélligens, aut requírens Deum.

Omnes declinavérunt, simul inútiles facti sunt: * non est qui fáciat bonum, non est usque ad unum.

Nonne scient omnes qui operántur iniquitátem, * qui dévorant plebem meam ut cibum panis?

Deum non invocavérunt: * illic trepidavérunt timóre, ubi non erat timor.

Quóniam Deus dissipávit ossa eórum qui homínibus placent: * confúsi sunt, quóniam Deus sprevit eos.

Quis dabit ex Sion salutáre Israël? * cum convérterit Deus captivitátem plebis suæ, exsultábit Jacob, et lætábitur Israël.

\rubric{¶ Quando ad Laudes dictus fuerit Ps. 50 Miserére, hic subjúngitur sequens Ps. 96, qui in Laudibus locum non habuit.}

\psalmheading{Psalmus 96}

\lettrine[lines=2]{D}{OMINUS} regnávit, exsúltet terra: * læténtur ínsulæ multæ.

Nubes, et calígo in circúitu ejus: * justítia, et judícium corréctio sedis ejus.

Ignis ante ipsum præcédet, * et inflammábit in circúitu inimícos ejus.

Illuxérunt fúlgura ejus orbi terræ: * vidit, et commóta est terra.

Montes, sicut cera fluxérunt a fácie Dómini: * a fácie Dómini omnis terra.

Annuntiavérunt cæli justítiam ejus: * et vidérunt omnes pópuli glóriam ejus.

Confundántur omnes, qui adórant sculptília: * et qui gloriántur in simulácris suis.

Adoráte eum, omnes Angeli ejus: * audívit, et lætáta est Sion.

Et exsultavérunt fíliæ Judæ, * propter judícia tua, Dómine:

Quóniam tu Dóminus Altíssimus super omnem terram: * nimis exaltátus es super omnes deos.

Qui dilígitis Dóminum, odíte malum: * custódit Dóminus ánimas sanctórum suórum, de manu peccatóris liberábit eos.

Lux orta est justo, * et rectis corde lætítia.

Lætámini, justi, in Dómino: * et confitémini memóriæ sanctificatiónis ejus.

\rubric{Secus vero, tribus Psalmis explétis, statim dícitur Ant.}

\antlabel{Ant.} Misericórdia tua, Dómine, ante óculos meos: et complácui in veritáte tua.

\rubric{In Officio Festi vel Octávæ:}

\capitulumlabel{Capitulum}{1 Tim. 1, 17}

\lettrine[lines=2]{R}{EGI} sæculórum immortáli et invisíbili, soli Deo honor et glória in sæcula sæculórum. Amen.

\rubric{In ómnibus autem Fériis et in Vigíliis commúnibus:}

\capitulumlabel{Capitulum}{Zach. 8, 19}

\lettrine[lines=2]{P}{ACEM} et veritátem dilígite, ait Dóminus omnípotens.

\R\ br. Christe, Fili Dei vivi, * Miserére nobis. Christe. \V\ \textit{Qui sedes ad déxteram Patris.} Miserére nobis. Glória Patri. Christe. \V\ Exsúrge, Christe, ádjuva nos. \R\ Et líbera nos propter nomen tuum.

\variantheading{Preces dominicales et feriales (si dicendæ sint)}

KYRIE, eléison. Christe, eléison. Kýrie, eléison.

Pater noster \rubric{secreto usque ad} \V\ Et ne nos indúcas in tentatiónem. \R\ Sed líbera nos a malo.

Credo in Deum \rubric{secreto usque ad} \V\ Carnis resurrectiónem. \R\ Vitam ætérnam. Amen. \V\ Et ego ad te, Dómine, clamávi. \R\ Et mane orátio mea prævéniet te. \V\ Repleátur os meum laude. \R\ Ut cantem glóriam tuam, tota die magnitúdinem tuam. \V\ Dómine, avérte fáciem tuam a peccátis meis. \R\ Et omnes iniquitátes meas dele. \V\ Cor mundum crea in me, Deus. \R\ Et spíritum rectum ínnova in viscéribus meis. \V\ Ne projícias me a fácie tua. \R\ Et spíritum sanctum tuum ne áuferas a me. \V\ Redde mihi lætítiam salutáris tui. \R\ Et spíritu principáli confírma me.

\variantheading{In Précibus feriálibus tantum}

\V\ Eripe me, Dómine, ab hómine malo. \R\ A viro iníquo éripe me. \V\ Eripe me de inimícis meis, Deus meus. \R\ Et ab insurgéntibus in me líbera me. \V\ Eripe me de operántibus iniquitátem. \R\ Et de viris sánguinum salva me. \V\ Sic psalmum dicam nómini tuo in sæculum sæculi. \R\ Ut reddam vota mea de die in diem. \V\ Exáudi nos, Deus, salutáris noster. \R\ Spes ómnium fínium terræ, et in mari longe. \V\ Deus, in adjutórium meum inténde. \R\ Dómine, ad adjuvándum me festína. \V\ Sanctus Deus, Sanctus fortis, Sanctus immortális. \R\ Miserére nobis. \V\ Bénedic, ánima mea, Dómino. \R\ Et ómnia, quæ intra me sunt, nómini sancto ejus. \V\ Bénedic, ánima mea, Dómino. \R\ Et noli oblivísci omnes retributiónes ejus. \V\ Qui propitiátur ómnibus iniquitátibus tuis. \R\ Qui sanat omnes infirmitátes tuas. \V\ Qui rédimit de intéritu vitam tuam. \R\ Qui corónat te in misericórdia et miseratiónibus. \V\ Qui replet in bonis desidérium tuum. \R\ Renovábitur ut áquilæ juvéntus tua.

\variantheading{In utrisque Précibus}

\V\ Adjutórium nostrum in nómine Dómini. \R\ Qui fecit cælum et terram.

\formulalabel{Confíteor. Misereátur. Indulgéntiam.}

\V\ Dignáre, Dómine, die isto. \R\ Sine peccáto nos custodíre. \V\ Miserére nostri, Dómine. \R\ Miserére nostri. \V\ Fiat misericórdia tua, Dómine, super nos. \R\ Quemádmodum sperávimus in te. \V\ Dómine, exáudi oratiónem meam. \R\ Et clamor meus ad te véniat.

\V\ Dóminus vobíscum.

Orémus. \oratiolabel{Oratio}

\lettrine[lines=2]{D}{OMINE}, Deus omnípotens, qui ad princípium hujus diéi nos perveníre fecísti: tua nos hódie salva virtúte; ut in hac die ad nullum declinémus peccátum, sed semper ad tuam justítiam faciéndam nostra procédant elóquia, dirigántur cogitatiónes et ópera. Per Dóminum. \R\ Amen.

\V\ Dóminus vobíscum. \V\ Benedicámus Dómino.

\rubric{In Choro Martyrológium.}

\V\ Pretiósa in conspéctu Dómini. \R\ Mors Sanctórum ejus.

\rubric{Et absolúte, sine} Orémus:

\oratiolabel{Oratio}

\lettrine[lines=2]{S}{ANCTA} María, et omnes Sancti intercédant pro nobis ad Dóminum, ut nos mereámur ab eo adjuvári et salvári, qui vivit et regnat in sæcula sæculórum. \R\ Amen.

\V\ Deus, in adjutórium meum inténde. \R\ Dómine, ad adjuvándum me festína. \rubric{Et dícitur ter; último additur:} Glória Patri. Sicut erat.

Kýrie, eléison. Christe, eléison. Kýrie, eléison.

Pater noster \rubric{secreto usque ad} \V\ Et ne nos indúcas in tentatiónem. \R\ Sed líbera nos a malo. \V\ Réspice in servos tuos, Dómine, et in ópera tua, et dírige fílios eórum. \R\ Et sit splendor Dómini, Dei nostri, super nos, et ópera mánuum nostrárum dírige super nos, et opus mánuum nostrárum dírige. \V\ Glória Patri. \R\ Sicut erat.

Orémus. \oratiolabel{Oratio}

\lettrine[lines=2]{D}{IRIGERE} et sanctificáre, régere et gubernáre dignáre, Dómine Deus, Rex cæli et terræ, hódie corda et córpora nostra, sensus, sermónes et actus nostros in lege tua, et in opéribus mandatórum tuórum: ut hic, et in ætérnum, te auxiliánte, salvi et líberi esse mereámur, Salvátor mundi: Qui vivis et regnas in sæcula sæculórum. \R\ Amen.

\V\ Jube, domne.

\benedlabel{Bened.} Dies et actus nostros in sua pace dispónat Dóminus omnípotens. \R\ Amen.

\variantheading{In Officio feriali}

\lettrine[lines=2]{D}{OMINUS} autem dírigat corda et córpora nostra in caritáte Dei, et patiéntia Christi. Tu autem.

\variantheading{E Proprio Sanctorum}

\variantheading{Pro Oct. Omn. Sanctorum}

\lettrine[lines=2]{B}{ENEDÍCTIO} et cláritas et sapiéntia et gratiárum áctio, honor, virtus et fortitúdo Deo nostro in sǽcula sæculórum. Amen.

\variantheading{E Communi Sanctorum}

\variantheading{Pro Apostolis}

\lettrine[lines=2]{I}{BANT} Apóstoli gaudéntes a conspéctu concílii, quóniam digni hábiti sunt pro nómine Jesu contuméliam pati.

\variantheading{Pro Martyre}

\lettrine[lines=2]{J}{USTUS} cor suum trádidit ad vigilándum dilúculo ad Dóminum, qui fecit illum, et in conspéctu Altíssimi deprecábitur.

\variantheading{Pro pluribus Martyribus}

\lettrine[lines=2]{F}{ULGEBUNT} justi, et tamquam scintíllæ in arundinéto discúrrent. Judicábunt natiónes, et dominabúntur pópulis: et regnábit Dóminus illórum in perpétuum.

\variantheading{Pro Conf. Pont.}

\lettrine[lines=2]{F}{UNGI} sacerdótio, et habére laudem in nómine ipsíus, et offérre illi incénsum dignum in odórem suavitátis.

\variantheading{Pro Conf. non Pont.}

\lettrine[lines=2]{J}{USTUM} dedúxit Dóminus per vias rectas, et osténdit illi regnum Dei, et dedit illi sciéntiam sanctórum; honestávit illum in labóribus, et complévit labóres illíus.

\variantheading{Pro Virgine Martyre}

\lettrine[lines=2]{D}{OMINE}, Deus meus, exaltásti super terram habitatiónem meam, et pro morte defluénte deprecáta sum. Invocávi Dóminum, Patrem Dómini mei, ut non derelínquat me in die tribulatiónis meæ, et in témpore superbórum sine adjutório.

\variantheading{Pro Virgine non Martyre}

\lettrine[lines=2]{O}{} \textsc{quam} pulchra est casta generátio cum claritáte! immortális est enim memória illíus: quóniam apud Deum nota est, et apud hómines.

\variantheading{Pro Martyre tantum}

\lettrine[lines=2]{L}{AUDABIT} usque ad mortem ánima mea Dóminum, quóniam éruis sustinéntes te, et líberas eos de manu angústiæ, Dómine, Deus noster.

\variantheading{Pro nec Virg. nec Mart.}

\lettrine[lines=2]{M}{ULTAE} fíliæ congregavérunt divítias: tu supergréssa es univérsas. Fallax grátia, et vana est pulchritúdo: múlier timens Dóminum, ipsa laudábitur.

\variantheading{Pro Dedicatione Ecclesiæ}

\lettrine[lines=2]{E}{T} abstérget Deus omnem lácrimam ab óculis eórum: et mors ultra non erit, neque luctus, neque clamor, neque dolor erit ultra, quia prima abiérunt. Et dixit qui sedébat in throno: Ecce, nova fácio ómnia.

\variantheading{In Officiis B. Mariæ V.}

\lettrine[lines=2]{I}{N} platéis sicut cinnamómum et bálsamum aromatízans odórem dedi: quasi myrrha elécta dedi suavitátem odóris.

\V\ Adjutórium nostrum in nómine Dómini. \R\ Qui fecit cælum et terram.

\V\ Benedícite. \R\ Deus.

\benedictiolabel{Benedictio} Dóminus nos benedícat, et ab omni malo deféndat, et ad vitam perdúcat ætérnam. Et fidélium ánimæ per misericórdiam Dei requiéscant in pace. \R\ Amen.

\rubric{Deinde dícitur tantum} Pater noster \rubric{secréto, nisi sequátur ália Hora.}

% ===== END OF BATCH 21 -- Covers autumnalis.md lines 7159-7424,
% completing "Ad Primam" for Feria Quarta in full: the hymn "Jam
% lucis orto sídere," Psalms 25, 51, 52, the conditional Psalm 96
% insertion, both Capitula, the brief Responsory, the full Preces
% apparatus, the closing Orationes, and the Lectio-brevis apparatus
% (E Proprio Sanctorum, E Communi Sanctorum). Four consecutive psalm
% incipits in autumnalis.md (Ps. 25 "JUdica," Ps. 51 "QUid," Ps. 52
% "DIxit," and the conditional Ps. 96 "DOminus") print only their
% first letter or two in full caps with the remainder of the first
% word in lowercase -- a mechanical OCR capitalization defect,
% confirmed for "DIxit" and "DOminus" by the many fully-capitalized
% correct occurrences of these exact psalm incipits already
% transcribed elsewhere in this volume (e.g. Psalm 13's "DIXIT
% insípiens" in Batch 9, and "DOMINUS regnávit" repeatedly in
% Batches 19-20) -- and silently normalized to full caps throughout,
% consistent with the uniform ALL-CAPS-opening-word convention
% applied to every psalm in this transcription. The document
% continues with "Ad Tertiam" (autumnalis.md line 7427).
% =====
\hourheading{Ad Tertiam}

Pater noster et Ave.

\V\ Deus, in adjutórium.

\hymnusheading{Hymnus}

\lettrine[lines=2]{N}{UNC}, Sancte, nobis, Spíritus,\\
Unum Patri cum Fílio,\\
Dignáre promptus íngeri\\
Nostro refúsus péctori.

Os, lingua, mens, sensus, vigor\\
Confessiónem pérsonent,\\
Flamméscat igne cáritas,\\
Accéndat ardor próximos.

Præsta, Pater piíssime,\\
Patríque compar Unice,\\
Cum Spíritu Paráclito\\
Regnans per omne sæculum.\\
Amen.

\antlabel{Ant.} Deus ádjuvat me.

\psalmheading{Psalmus 53}

\lettrine[lines=2]{D}{EUS}, in nómine tuo salvum me fac: * et in virtúte tua júdica me.

Deus, exáudi oratiónem meam: * áuribus pércipe verba oris mei.

Quóniam aliéni insurrexérunt advérsum me, et fortes quæsiérunt ánimam meam: * et non proposuérunt Deum ante conspéctum suum.

Ecce enim, Deus ádjuvat me: * et Dóminus suscéptor est ánimæ meæ.

Avérte mala inimícis meis: * et in veritáte tua dispérde illos.

Voluntárie sacrificábo tibi, * et confitébor nómini tuo, Dómine: quóniam bonum est:

Quóniam ex omni tribulatióne eripuísti me: * et super inimícos meos despéxit óculus meus.

\psalmheading{Psalmus 54, i}

\lettrine[lines=2]{E}{XAUDI}, Deus, oratiónem meam, et ne despéxeris deprecatiónem meam: * inténde mihi, et exáudi me.

Contristátus sum in exercitatióne mea: * et conturbátus sum a voce inimíci, et a tribulatióne peccatóris.

Quóniam declinavérunt in me iniquitátes: * et in ira molésti erant mihi.

Cor meum conturbátum est in me: * et formído mortis cécidit super me.

Timor et tremor venérunt super me: * et contexérunt me ténebræ.

Et dixi: Quis dabit mihi pennas sicut colúmbæ, * et volábo, et requiéscam?

Ecce, elongávi fúgiens: * et mansi in solitúdine.

Exspectábam eum, qui salvum me fecit * a pusillanimitáte spíritus et tempestáte.

Præcípita, Dómine, dívide linguas eórum: * quóniam vidi iniquitátem, et contradictiónem in civitáte.

Die ac nocte circúmdabit eam super muros ejus iníquitas: * et labor in médio ejus, et injustítia.

Et non defécit de platéis ejus * usúra, et dolus.

Quóniam si inimícus meus maledixísset mihi, * sustinuíssem útique.

Et si is, qui óderat me, super me magna locútus fuísset, * abscondíssem me fórsitan ab eo.

Tu vero, homo unánimis: * dux meus, et notus meus:

Qui simul mecum dulces capiébas cibos: * in domo Dei ambulávimus cum consénsu.

Véniat mors super illos: * et descéndant in inférnum vivéntes:

Quóniam nequítiæ in habitáculis eórum: * in médio eórum.

\psalmheading{Psalmus 54, ii}

\lettrine[lines=2]{E}{GO} autem ad Deum clamávi: * et Dóminus salvábit me.

Véspere, et mane, et merídie narrábo et annuntiábo: * et exáudiet vocem meam.

Rédimet in pace ánimam meam ab his, qui appropínquant mihi: * quóniam inter multos erant mecum.

Exáudiet Deus, et humiliábit illos, * qui est ante sæcula.

Non enim est illis commutátio, et non timuérunt Deum: * exténdit manum suam in retribuéndo.

Contaminavérunt testaméntum ejus, divísi sunt ab ira vultus ejus: * et appropinquávit cor illíus.

Mollíti sunt sermónes ejus super óleum: * et ipsi sunt jácula.

Jacta super Dóminum curam tuam, et ipse te enútriet: * non dabit in ætérnum fluctuatiónem justo.

Tu vero, Deus, dedúces eos, * in púteum intéritus.

Viri sánguinum, et dolósi non dimidiábunt dies suos: * ego autem sperábo in te, Dómine.

\antlabel{Ant.} Deus ádjuvat me: et Dóminus suscéptor est ánimæ meæ.

\variantheading{In Officio feriali}

\capitulumlabel{Capitulum}{Jer. 17, 14}

\lettrine[lines=2]{S}{ANA} me, Dómine, et sanábor: salvum me fac, et salvus ero: quóniam laus mea tu es.

\R\ br. Sana ánimam meam, * Quia peccávi tibi. Sana. \V\ Ego dixi: Dómine, miserére mei. Quia. Glória Patri. Sana. \V\ Adjútor meus esto, ne derelínquas me. \R\ Neque despícias me, Deus, salutáris meus.

\variantheading{In Vigiliis communibus:}

\variantheading{Preces feriales}

KYRIE, eléison. Christe, eléison. Kýrie, eléison.

Pater noster \rubric{secreto usque ad} \V\ Et ne nos indúcas in tentatiónem. \R\ Sed líbera nos a malo. \V\ Dómine, Deus virtútum, convérte nos. \R\ Et osténde fáciem tuam, et salvi érimus. \V\ Exsúrge, Christe, ádjuva nos. \R\ Et líbera nos propter nomen tuum. \V\ Dómine, exáudi oratiónem meam. \R\ Et clamor meus ad te véniat.

\rubric{Réliqua ut in Ordinário. 19}

\hourheading{Ad Sextam}

Pater noster et Ave.

\V\ Deus, in adjutórium.

\hymnusheading{Hymnus}

\lettrine[lines=2]{R}{ECTOR} potens, verax Deus,\\
Qui témperas rerum vices,\\
Splendóre mane illúminas,\\
Et ígnibus merídiem:

Exstíngue flammas lítium,\\
Aufer calórem nóxium,\\
Confer salútem córporum,\\
Verámque pacem córdium.

Præsta, Pater piíssime,\\
Patríque compar Unice,\\
Cum Spíritu Paráclito\\
Regnans per omne sæculum.\\
Amen.

\antlabel{Ant.} In Deo sperávi.

\psalmheading{Psalmus 55}

\lettrine[lines=2]{M}{ISERERE} mei, Deus, quóniam conculcávit me homo: * tota die impúgnans tribulávit me.

Conculcavérunt me inimíci mei tota die: * quóniam multi bellántes advérsum me.

Ab altitúdine diéi timébo: * ego vero in te sperábo.

In Deo laudábo sermónes meos, in Deo sperávi: * non timébo quid fáciat mihi caro.

Tota die verba mea exsecrabántur: * advérsum me omnes cogitatiónes eórum in malum.

Inhabitábunt et abscóndent: * ipsi calcáneum meum observábunt.

Sicut sustinuérunt ánimam meam, pro níhilo salvos fácies illos: * in ira pópulos confrínges.

Deus, vitam meam annuntiávi tibi: * posuísti lácrimas meas in conspéctu tuo.

Sicut et in promissióne tua: * tunc converténtur inimíci mei retrórsum:

In quacúmque die invocávero te: * ecce, cognóvi, quóniam Deus meus es.

In Deo laudábo verbum, in Dómino laudábo sermónem: * in Deo sperávi, non timébo quid fáciat mihi homo.

In me sunt, Deus, vota tua, * quæ reddam, laudatiónes tibi.

Quóniam eripuísti ánimam meam de morte, et pedes meos de lapsu: * ut pláceam coram Deo in lúmine vivéntium.

\psalmheading{Psalmus 56}

\lettrine[lines=2]{M}{ISERERE} mei, Deus, miserére mei: * quóniam in te confídit ánima mea.

Et in umbra alárum tuárum sperábo, * donec tránseat iníquitas.

Clamábo ad Deum altíssimum: * Deum qui benefécit mihi.

Misit de cælo, et liberávit me: * dedit in oppróbrium conculcántes me.

Misit Deus misericórdiam suam, et veritátem suam, * et erípuit ánimam meam de médio catulórum leónum: dormívi conturbátus.

Fílii hóminum dentes eórum arma et sagíttæ: * et lingua eórum gládius acútus.

Exaltáre super cælos, Deus, * et in omnem terram glória tua.

Láqueum paravérunt pédibus meis: * et incurvavérunt ánimam meam.

Fodérunt ante fáciem meam fóveam: * et incidérunt in eam.

Parátum cor meum, Deus, parátum cor meum: * cantábo, et psalmum dicam.

Exsúrge, glória mea, exsúrge, psaltérium et cíthara: * exsúrgam dilúculo.

Confitébor tibi in pópulis, Dómine: * et psalmum dicam tibi in géntibus:

Quóniam magnificáta est usque ad cælos misericórdia tua, * et usque ad nubes véritas tua.

Exaltáre super cælos, Deus: * et super omnem terram glória tua.

\psalmheading{Psalmus 57}

\lettrine[lines=2]{S}{I} vere útique justítiam loquímini: * recta judicáte, fílii hóminum.

Etenim in corde iniquitátes operámini: * in terra injustítias manus vestræ concínnant.

Alienáti sunt peccatóres a vulva, erravérunt ab útero: * locúti sunt falsa.

Furor illis secúndum similitúdinem serpéntis: * sicut áspidis surdæ, et obturántis aures suas.

Quæ non exáudiet vocem incantántium: * et venéfici incantántis sapiénter.

Deus cónteret dentes eórum in ore ipsórum: * molas leónum confrínget Dóminus.

Ad níhilum devénient tamquam aqua decúrrens: * inténdit arcum suum donec infirméntur.

Sicut cera, quæ fluit, auferéntur: * supercécidit ignis, et non vidérunt solem.

Priúsquam intellígerent spinæ vestræ rhamnum: * sicut vivéntes, sic in ira absórbet eos.

Lætábitur justus cum víderit vindíctam: * manus suas lavábit in sánguine peccatóris.

Et dicet homo: Si útique est fructus justo: * útique est Deus júdicans eos in terra.

\antlabel{Ant.} In Deo sperávi, non timébo quid fáciat mihi homo.

\variantheading{In Officio feriali}

\capitulumlabel{Capitulum}{Rom. 13, 8}

\lettrine[lines=2]{N}{EMINI} quidquam debeátis, nisi ut ínvicem diligátis: qui enim díligit próximum, legem implévit.

\R\ br. Benedícam Dóminum * In omni témpore. Benedícam. \V\ Semper laus ejus in ore meo. In. Glória Patri. Benedícam. \V\ Dóminus regit me, et nihil mihi déerit. \R\ In loco páscuæ ibi me collocávit.

\variantheading{In Vigiliis communibus:}

\variantheading{Preces feriales}

KYRIE, eléison. Christe, eléison. Kýrie, eléison.

Pater noster \rubric{secreto usque ad} \V\ Et ne nos indúcas in tentatiónem. \R\ Sed líbera nos a malo. \V\ Dómine, Deus virtútum, convérte nos. \R\ Et osténde fáciem tuam, et salvi érimus. \V\ Exsúrge, Christe, ádjuva nos. \R\ Et líbera nos propter nomen tuum. \V\ Dómine, exáudi oratiónem meam. \R\ Et clamor meus ad te véniat.

\rubric{Réliqua ut in Ordinário. 20}

\hourheading{Ad Nonam}

Pater noster et Ave.

\V\ Deus, in adjutórium.

\hymnusheading{Hymnus}

\lettrine[lines=2]{R}{ERUM}, Deus, tenax vigor,\\
Immótus in te pérmanens,\\
Lucis diúrnæ témpora\\
Succéssibus detérminans:

Largíre lumen véspere,\\
Quo vita nusquam décidat,\\
Sed prǽmium mortis sacræ\\
Perénnis instet glória.

Præsta, Pater piíssime,\\
Patríque compar Unice,\\
Cum Spíritu Paráclito\\
Regnans per omne sæculum.\\
Amen.

\antlabel{Ant.} Deus meus.

\psalmheading{Psalmus 58, i}

\lettrine[lines=2]{E}{RIPE} me de inimícis meis, Deus meus: * et ab insurgéntibus in me líbera me.

Eripe me de operántibus iniquitátem: * et de viris sánguinum salva me.

Quia ecce cepérunt ánimam meam: * irruérunt in me fortes.

Neque iníquitas mea, neque peccátum meum, Dómine: * sine iniquitáte cucúrri, et diréxi.

Exsúrge in occúrsum meum, et vide: * et tu, Dómine, Deus virtútum, Deus Israël,

Inténde ad visitándas omnes gentes: * non misereáris ómnibus, qui operántur iniquitátem.

Converténtur ad vésperam: et famem patiéntur ut canes, * et circuíbunt civitátem.

Ecce, loquéntur in ore suo, et gládius in lábiis eórum: * quóniam quis audívit?

Et tu, Dómine, deridébis eos: * ad níhilum dedúces omnes gentes.

Fortitúdinem meam ad te custódiam, quia, Deus, suscéptor meus es: * Deus meus, misericórdia ejus prævéniet me.

\psalmheading{Psalmus 58, ii}

\lettrine[lines=2]{D}{EUS} osténdet mihi super inimícos meos, ne occídas eos: * nequándo obliviscántur pópuli mei.

Dispérge illos in virtúte tua: * et depóne eos, protéctor meus, Dómine:

Delíctum oris eórum, sermónem labiórum ipsórum: * et comprehendántur in supérbia sua.

Et de exsecratióne et mendácio annuntiabúntur in consummatióne: * in ira consummatiónis, et non erunt.

Et scient quia Deus dominábitur Jacob: * et fínium terræ.

Converténtur ad vésperam: et famem patiéntur ut canes, * et circuíbunt civitátem.

Ipsi dispergéntur ad manducándum: * si vero non fúerint saturáti, et murmurábunt.

Ego autem cantábo fortitúdinem tuam: * et exsultábo mane misericórdiam tuam.

Quia factus es suscéptor meus, * et refúgium meum, in die tribulatiónis meæ.

Adjútor meus, tibi psallam, quia, Deus, suscéptor meus es: * Deus meus, misericórdia mea.

\psalmheading{Psalmus 59}

\lettrine[lines=2]{D}{EUS}, repulísti nos, et destruxísti nos: * irátus es, et misértus es nobis.

Commovísti terram, et conturbásti eam: * sana contritiónes ejus, quia commóta est.

Ostendísti pópulo tuo dura: * potásti nos vino compunctiónis.

Dedísti metuéntibus te significatiónem: * ut fúgiant a fácie arcus:

Ut liberéntur dilécti tui: * salvum fac déxtera tua, et exáudi me.

Deus locútus est in sancto suo: * Lætábor, et partíbor Síchimam: et convállem tabernaculórum metíbor.

Meus est Gálaad, et meus est Manásses: * et Ephraïm fortitúdo cápitis mei.

Juda rex meus: * Moab olla spei meæ.

In Idumǽam exténdam calceaméntum meum: * mihi aliénigenæ súbditi sunt.

Quis dedúcet me in civitátem munítam? * quis dedúcet me usque in Idumǽam?

Nonne tu, Deus, qui repulísti nos, * et non egrediéris, Deus, in virtútibus nostris?

Da nobis auxílium de tribulatióne: * quia vana salus hóminis.

In Deo faciémus virtútem: * et ipse ad níhilum dedúcet tribulántes nos.

\antlabel{Ant.} Deus meus, misericórdia tua prævéniet me.

\variantheading{In Officio feriali}

\capitulumlabel{Capit.}{1 Petri 1, 17-19}

\lettrine[lines=2]{I}{N} timóre incolátus vestri témpore conversámini: sciéntes quod non corruptibílibus auro vel argénto redémpti estis, sed pretióso sánguine quasi Agni immaculáti Christi.

\R\ br. Rédime me, Dómine, * Et miserére mei. Rédime. \V\ Pes enim meus stetit in via recta. Et miserére mei. Glória Patri. Rédime. \V\ Ab occúltis meis munda me, Dómine. \R\ Et ab aliénis parce servo tuo.

\variantheading{In Vigiliis communibus:}

\variantheading{Preces feriales}

KYRIE, eléison. Christe, eléison. Kýrie, eléison.

Pater noster \rubric{secreto usque ad} \V\ Et ne nos indúcas in tentatiónem. \R\ Sed líbera nos a malo. \V\ Dómine, Deus virtútum, convérte nos. \R\ Et osténde fáciem tuam, et salvi érimus. \V\ Exsúrge, Christe, ádjuva nos. \R\ Et líbera nos propter nomen tuum. \V\ Dómine, exáudi oratiónem meam. \R\ Et clamor meus ad te véniat.

\rubric{Réliqua ut in Ordinário. 22}

% ===== END OF BATCH 22 -- Covers autumnalis.md lines 7427-7796,
% completing the Little Hours of Feria Quarta in full: "Ad
% Tertiam" (Psalms 53, 54/i, 54/ii), "Ad Sextam" (Psalms 55, 56,
% 57), and "Ad Nonam" (Psalms 58/i, 58/ii, 59), each with its own
% hymn, Capitulum, brief Responsory, and "In Vigiliis communibus:
% Preces feriales" apparatus. Three further occurrences of the
% recurring stray OCR "\*" artifact in the hour-hymns are silently
% dropped, consistent with every prior occurrence. Several more
% psalm incipits in this batch carry the same partial-caps OCR
% defect already documented in Batch 21 ("DEus," "EXaudi," "EGo,"
% "SAna" [Capitulum], "MIserére" twice, "ERipe," "DEus" twice,
% "NEmini" [Capitulum]) and are silently normalized to full caps,
% consistent with the uniform convention and confirmed by the many
% correctly-capitalized occurrences of these same incipits elsewhere
% in this volume. The document continues with "Ad Vesperas"
% (autumnalis.md line 7799).
% =====
\hourheading{Ad Vesperas}

Pater noster et Ave.

\V\ Deus, in adjutórium.

\antlabel{Ant.} Beáti omnes, * \textdagger\ qui timent Dóminum. \textdaggerdbl

\psalmheading{Psalmus 127}

\lettrine[lines=2]{B}{EATI} omnes, \textdagger\ qui timent Dóminum, * \textdaggerdbl\ qui ámbulant in viis ejus.

Labóres mánuum tuárum quia manducábis: * beátus es, et bene tibi erit.

Uxor tua sicut vitis abúndans, * in latéribus domus tuæ.

Fílii tui sicut novéllæ olivárum, * in circúitu mensæ tuæ.

Ecce, sic benedicétur homo, * qui timet Dóminum.

Benedícat tibi Dóminus ex Sion: * et vídeas bona Jerúsalem ómnibus diébus vitæ tuæ.

Et vídeas fílios filiórum tuórum, * pacem super Israël.

\antlabel{Ant.} Beáti omnes, qui timent Dóminum.

\antlabel{Ant.} Confundántur omnes, * qui odérunt Sion.

\psalmheading{Psalmus 128}

\lettrine[lines=2]{S}{ÆPE} expugnavérunt me a juventúte mea, * dicat nunc Israël:

Sæpe expugnavérunt me a juventúte mea: * étenim non potuérunt mihi.

Supra dorsum meum fabricavérunt peccatóres: * prolongavérunt iniquitátem suam.

Dóminus justus concídit cervíces peccatórum: * confundántur et convertántur retrórsum omnes, qui odérunt Sion.

Fiant sicut fœnum tectórum: * quod priúsquam evellátur, exáruit:

De quo non implévit manum suam qui metit, * et sinum suum qui manípulos cólligit.

Et non dixérunt qui præteríbant: Benedíctio Dómini super vos: * benedíximus vobis in nómine Dómini.

\antlabel{Ant.} Confundántur omnes, qui odérunt Sion.

\antlabel{Ant.} De profúndis * \textdagger\ clamávi ad te, Dómine. \textdaggerdbl

\psalmheading{Psalmus 129}

\lettrine[lines=2]{D}{E} profúndis \textdagger\ clamávi ad te, Dómine: * \textdaggerdbl\ Dómine, exáudi vocem meam:

Fiant aures tuæ inténdentes, * in vocem deprecatiónis meæ.

Si iniquitátes observáveris, Dómine: * Dómine, quis sustinébit?

Quia apud te propitiátio est: * et propter legem tuam sustínui te, Dómine.

Sustínuit ánima mea in verbo ejus: * sperávit ánima mea in Dómino.

A custódia matutína usque ad noctem: * speret Israël in Dómino.

Quia apud Dóminum misericórdia: * et copiósa apud eum redémptio.

Et ipse rédimet Israël, * ex ómnibus iniquitátibus ejus.

\antlabel{Ant.} De profúndis clamávi ad te, Dómine.

\antlabel{Ant.} Dómine, * \textdagger\ non est exaltátum cor meum. \textdaggerdbl

\psalmheading{Psalmus 130}

\lettrine[lines=2]{D}{OMINE}, \textdagger\ non est exaltátum cor meum: * \textdaggerdbl\ neque eláti sunt óculi mei.

Neque ambulávi in magnis: * neque in mirabílibus super me.

Si non humíliter sentiébam: * sed exaltávi ánimam meam:

Sicut ablactátus est super matre sua, * ita retribútio in ánima mea.

Speret Israël in Dómino, * ex hoc nunc et usque in sæculum.

\antlabel{Ant.} Dómine, non est exaltátum cor meum.

\antlabel{Ant.} Elégit Dóminus * Sion in habitatiónem sibi.

\psalmheading{Psalmus 131}

\lettrine[lines=2]{M}{EMENTO}, Dómine, David, * et omnis mansuetúdinis ejus:

Sicut jurávit Dómino, * votum vovit Deo Jacob:

Si introíero in tabernáculum domus meæ, * si ascéndero in lectum strati mei:

Si dédero somnum óculis meis, * et pálpebris meis dormitatiónem:

Et réquiem tempóribus meis: donec invéniam locum Dómino, * tabernáculum Deo Jacob.

Ecce, audívimus eam in Ephráta: * invénimus eam in campis silvæ.

Introíbimus in tabernáculum ejus: * adorábimus in loco, ubi stetérunt pedes ejus.

Surge, Dómine, in réquiem tuam, * tu et arca sanctificatiónis tuæ.

Sacerdótes tui induántur justítiam: * et sancti tui exsúltent.

Propter David, servum tuum, * non avértas fáciem Christi tui.

Jurávit Dóminus David veritátem, et non frustrábitur eam: * De fructu ventris tui ponam super sedem tuam.

Si custodíerint fílii tui testaméntum meum, * et testimónia mea hæc, quæ docébo eos:

Et fílii eórum usque in sæculum, * sedébunt super sedem tuam.

Quóniam elégit Dóminus Sion: * elégit eam in habitatiónem sibi.

Hæc réquies mea in sæculum sæculi: * hic habitábo quóniam elégi eam.

Víduam ejus benedícens benedícam: * páuperes ejus saturábo pánibus.

Sacerdótes ejus índuam salutári: * et sancti ejus exsultatióne exsultábunt.

Illuc prodúcam cornu David, * parávi lucérnam Christo meo.

Inimícos ejus índuam confusióne: * super ipsum autem efflorébit sanctificátio mea.

\antlabel{Ant.} Elégit Dóminus Sion in habitatiónem sibi.

\variantheading{In Officio feriali}

\capitulumlabel{Capitulum}{2 Cor. 1, 3-4}

\lettrine[lines=2]{B}{ENEDICTUS} Deus, et Pater Dómini nostri Jesu Christi, Pater misericordiárum, et Deus totíus consolatiónis, qui consolátur nos in omni tribulatióne nostra.

\hymnusheading{Hymnus}

\lettrine[lines=2]{C}{ÆLI} Deus sanctíssime,\\
Qui lúcidas mundi plagas\\
Candóre pingis ígneo,\\
Augens decóro lúmine:

Quarto die qui flámmeam\\
Dum solis accéndis rotam,\\
Lunæ minístras órdinem,\\
Vagósque cursus síderum:

Ut nóctibus, vel lúmini\\
Diremptiónis términum,\\
Primórdiis et ménsium\\
Signum dares notíssimum;

Expélle noctem córdium:\\
Abstérge sordes méntium:\\
Resólve culpæ vínculum:\\
Evérte moles críminum.

Præsta, Pater piíssime,\\
Patríque compar Unice,\\
Cum Spíritu Paráclito\\
Regnans per omne sæculum.\\
Amen.

\V\ Dirigátur, Dómine, orátio mea. \R\ Sicut incénsum in conspéctu tuo.

\antlabel{Ad Magnif. Ant.} Respéxit Dóminus * humilitátem meam, et fecit in me magna, qui potens est.

\rubric{Réliqua ut in Ordinário. 22}

\hourheading{Ad Completorium}

\V\ Jube, domne.

\benedlabel{Bened.} Noctem quiétam, et finem perféctum concédat nobis Dóminus omnípotens. \R\ Amen.

\lectioheading{Lectio br.}{1 Petri 5, 8-9}

\lettrine[lines=2]{F}{RATRES}: Sóbrii estóte, et vigiláte: quia adversárius vester diábolus tamquam leo rúgiens círcuit, quærens quem dévoret: cui resístite fortes in fide. Tu autem, Dómine, miserére nobis. \R\ Deo grátias.

\V\ Adjutórium nostrum in nómine Dómini. \R\ Qui fecit cælum et terram.

Pater noster \rubric{quod dícitur totum secréto.}

\formulalabel{Confíteor.} Misereátur. Indulgéntiam.

\V\ Convérte nos, Deus, salutáris noster. \R\ Et avérte iram tuam a nobis.

\V\ Deus, in adjutórium.

\antlabel{Ant.} Immíttet Angelus Dómini.

\psalmheading{Psalmus 33, i}

\lettrine[lines=2]{B}{ENEDICAM} Dóminum in omni témpore: * semper laus ejus in ore meo.

In Dómino laudábitur ánima mea: * áudiant mansuéti, et læténtur.

Magnificáte Dóminum mecum: * et exaltémus nomen ejus in idípsum.

Exquisívi Dóminum, et exaudívit me: * et ex ómnibus tribulatiónibus meis erípuit me.

Accédite ad eum, et illuminámini: * et fácies vestræ non confundéntur.

Ipse pauper clamávit, et Dóminus exaudívit eum: * et de ómnibus tribulatiónibus ejus salvávit eum.

Immíttet Angelus Dómini in circúitu timéntium eum: * et erípiet eos.

Gustáte, et vidéte quóniam suávis est Dóminus: * beátus vir, qui sperat in eo.

Timéte Dóminum, omnes sancti ejus: * quóniam non est inópia timéntibus eum.

Dívites eguérunt et esuriérunt: * inquiréntes autem Dóminum non minuéntur omni bono.

\psalmheading{Psalmus 33, ii}

\lettrine[lines=2]{V}{ENITE}, fílii, audíte me: * timórem Dómini docébo vos.

Quis est homo qui vult vitam: * díligit dies vidére bonos?

Próhibe linguam tuam a malo: * et lábia tua ne loquántur dolum.

Divérte a malo, et fac bonum: * inquíre pacem, et perséquere eam.

Oculi Dómini super justos: * et aures ejus in preces eórum.

Vultus autem Dómini super faciéntes mala: * ut perdat de terra memóriam eórum.

Clamavérunt justi, et Dóminus exaudívit eos: * et ex ómnibus tribulatiónibus eórum liberávit eos.

Juxta est Dóminus iis, qui tribuláto sunt corde: * et húmiles spíritu salvábit.

Multæ tribulatiónes justórum: * et de ómnibus his liberábit eos Dóminus.

Custódit Dóminus ómnia ossa eórum: * unum ex his non conterétur.

Mors peccatórum péssima: * et qui odérunt justum, delínquent.

Rédimet Dóminus ánimas servórum suórum: * et non delínquent omnes qui sperant in eo.

\psalmheading{Psalmus 60}

\lettrine[lines=2]{E}{XAUDI}, Deus, deprecatiónem meam: * inténde oratióni meæ.

A fínibus terræ ad te clamávi: * dum anxiarétur cor meum, in petra exaltásti me.

Deduxísti me, quia factus es spes mea: * turris fortitúdinis a fácie inimíci.

Inhabitábo in tabernáculo tuo in sæcula: * prótegar in velaménto alárum tuárum.

Quóniam tu, Deus meus, exaudísti oratiónem meam: * dedísti hereditátem timéntibus nomen tuum.

Dies super dies regis adjícies: * annos ejus usque in diem generatiónis et generatiónis.

Pérmanet in ætérnum in conspéctu Dei: * misericórdiam et veritátem ejus quis requíret?

Sic psalmum dicam nómini tuo in sæculum sæculi: * ut reddam vota mea de die in diem.

\antlabel{Ant.} Immíttet Angelus Dómini in circúitu timéntium eum: et erípiet eos.

\hymnusheading{Hymnus}

\lettrine[lines=2]{T}{E} lucis ante términum,\\
Rerum Creátor, póscimus,\\
Ut pro tua cleméntia\\
Sis præsul et custódia.

Procul recédant sómnia,\\
Et nóctium phantásmata,\\
Hostémque nostrum cómprime,\\
Ne polluántur córpora.

Præsta, Pater piíssime,\\
Patríque compar Unice,\\
Cum Spíritu Paráclito\\
Regnans per omne sæculum.\\
Amen.

\capitulumlabel{Capitulum}{Jer. 14, 9}

\lettrine[lines=2]{T}{U} autem in nobis es, Dómine, et nomen sanctum tuum invocátum est super nos: ne derelínquas nos, Dómine, Deus noster.

\R\ br. In manus tuas, Dómine, * Comméndo spíritum meum. In manus. \V\ Redemísti nos, Dómine, Deus veritátis. Comméndo. Glória Patri. In manus.

\V\ Custódi nos, Dómine, ut pupíllam óculi. \R\ Sub umbra alárum tuárum prótege nos.

\antlabel{Ant.} Salva nos.

\canticumheading{Canticum Simeonis}{Luc. 2, 29-32}

\lettrine[lines=2]{N}{UNC} dimíttis servum tuum, Dómine, * secúndum verbum tuum in pace:

Quia vidérunt óculi mei * salutáre tuum,

Quod parásti * ante fáciem ómnium populórum,

Lumen ad revelatiónem géntium, * et glóriam plebis tuæ Israël.

Glória Patri.

\antlabel{Ant.} Salva nos, Dómine, vigilántes, custódi nos dormiéntes; ut vigilémus cum Christo, et requiescámus in pace.

\variantheading{Preces (si dicendæ sint)}

KYRIE, eléison. Christe, eléison. Kýrie, eléison.

Pater noster \rubric{secreto usque ad} \V\ Et ne nos indúcas in tentatiónem. \R\ Sed líbera nos a malo.

Credo in Deum \rubric{secreto usque ad} \V\ Carnis resurrectiónem. \R\ Vitam ætérnam. Amen. \V\ Benedíctus es, Dómine, Deus patrum nostrórum. \R\ Et laudábilis et gloriósus in sæcula. \V\ Benedicámus Patrem et Fílium cum Sancto Spíritu. \R\ Laudémus, et superexaltémus eum in sæcula. \V\ Benedíctus es, Dómine, in firmaménto cæli. \R\ Et laudábilis, et gloriósus, et superexaltátus in sæcula. \V\ Benedícat et custódiat nos omnípotens et miséricors Dóminus. \R\ Amen. \V\ Dignáre, Dómine, nocte ista. \R\ Sine peccáto nos custodíre. \V\ Miserére nostri, Dómine. \R\ Miserére nostri. \V\ Fiat misericórdia tua, Dómine, super nos. \R\ Quemádmodum sperávimus in te. \V\ Dómine, exáudi oratiónem meam. \R\ Et clamor meus ad te véniat.

\V\ Dóminus vobíscum.

Orémus. \oratiolabel{Oratio}

\lettrine[lines=2]{V}{ISITA}, quæsumus, Dómine, habitatiónem istam, et omnes insídias inimíci ab ea longe repélle: Angeli tui sancti hábitent in ea, qui nos in pace custódiant; et benedíctio tua sit super nos semper. Per Dóminum nostrum. \R\ Amen.

\V\ Dóminus vobíscum.

\V\ Benedicámus Dómino.

\benedictiolabel{Benedictio} Benedícat et custódiat nos omnípotens et miséricors Dóminus, Pater, et Fílius, et Spíritus Sanctus. \R\ Amen.

\variantheading{Antiphona finalis B. M. V.}

\V\ Divínum auxílium.

Pater, Ave et Credo.

% ===== END OF BATCH 23 -- Covers autumnalis.md lines 7799-8128,
% completing "Ad Vesperas" (Psalms 127, 128, 129, 130, 131, hymn
% "Cæli Deus sanctíssime") and "Ad Completorium" (Psalms 33/i,
% 33/ii, 60, hymn "Te lucis ante términum," Canticum Simeonis,
% Preces), and with them the ENTIRE "Feria Quarta" Office. More
% partial-caps OCR defects ("MEmento," "BEnedictus" [Capitulum],
% "DOmine," "DE profúndis," "BEnedicam," "VEnite," "EXaudi") are
% silently normalized to full caps, consistent with Batches 21-22.
% The partial-bold "Confíteor." (not the full "Misereátur.
% Indulgéntiam.") again appears at Compline, rendered the same way
% as in Feria Secunda and Feria Tertia. The document continues with
% "Feria Quinta" (autumnalis.md line 8131), the fourth weekday of the
% Psalterium cycle.
% =====
\feriaheading{Feria Quinta}

\hourheading{Ad Matutinum}

Pater, Ave et Credo.

\V\ Dómine, lábia mea.

\V\ Deus, in adjutórium.

\variantheading{In Officio feriali}

\invitatlabel{Invit.} Regem ventúrum Dóminum, * Veníte, adorémus. \psalmlabel{Ps. 94} Veníte.

\hymnusheading{Hymnus}

\lettrine[lines=2]{N}{OX} atra rerum cóntegit\\
Terræ colóres ómnium:\\
Nos confiténtes póscimus\\
Te, juste judex córdium:

Ut áuferas piácula,\\
Sordésque mentis ábluas:\\
Donésque, Christe, grátiam,\\
Ut arceántur crímina.

Mens ecce torpet ímpia,\\
Quam culpa mordet nóxia:\\
Obscúra gestit tóllere,\\
Et te, Redémptor, quærere.

Repélle tu calíginem\\
Intrínsecus quam máxime,\\
Ut in beáto gáudeat\\
Se collocári lúmine.

Præsta, Pater piíssime,\\
Patríque compar Unice,\\
Cum Spíritu Paráclito\\
Regnans per omne sæculum.\\
Amen.

\nocturnoheading{In I Nocturno}

\antlabel{Ant.} In Deo salutáre meum, * et glória mea: et spes mea in Deo est.

\psalmheading{Psalmus 61}

\lettrine[lines=2]{N}{ONNE} Deo subjécta erit ánima mea? * ab ipso enim salutáre meum.

Nam et ipse Deus meus, et salutáris meus: * suscéptor meus, non movébor ámplius.

Quoúsque irrúitis in hóminem? * interfícitis univérsi vos: tamquam paríeti inclináto et macériæ depúlsæ?

Verúmtamen prétium meum cogitavérunt repéllere, cucúrri in siti: * ore suo benedicébant, et corde suo maledicébant.

Verúmtamen Deo subjécta esto, ánima mea: * quóniam ab ipso patiéntia mea.

Quia ipse Deus meus, et salvátor meus: * adjútor meus, non emigrábo.

In Deo salutáre meum, et glória mea: * Deus auxílii mei, et spes mea in Deo est.

Speráte in eo, omnis congregátio pópuli, effúndite coram illo corda vestra: * Deus adjútor noster in ætérnum.

Verúmtamen vani fílii hóminum, mendáces fílii hóminum in statéris: * ut decípiant ipsi de vanitáte in idípsum.

Nolíte speráre in iniquitáte, et rapínas nolíte concupíscere: * divítiæ si áffluant, nolíte cor appónere.

Semel locútus est Deus, duo hæc audívi, quia potéstas Dei est, et tibi, Dómine, misericórdia: * quia tu reddes unicuíque juxta ópera sua.

\antlabel{Ant.} In Deo salutáre meum, et glória mea: et spes mea in Deo est.

\antlabel{Ant.} Vidéte ópera Dómini, * et audítam fácite vocem laudis ejus.

\psalmheading{Psalmus 65, i}

\lettrine[lines=2]{J}{UBILATE} Deo, omnis terra, psalmum dícite nómini ejus: * date glóriam laudi ejus.

Dícite Deo: Quam terribília sunt ópera tua, Dómine! * in multitúdine virtútis tuæ mentiéntur tibi inimíci tui.

Omnis terra adóret te, et psallat tibi: * psalmum dicat nómini tuo.

Veníte, et vidéte ópera Dei: * terríbilis in consíliis super fílios hóminum.

Qui convértit mare in áridam, in flúmine pertransíbunt pede: * ibi lætábimur in ipso.

Qui dominátur in virtúte sua in ætérnum, óculi ejus super gentes respíciunt: * qui exásperant non exalténtur in semetípsis.

Benedícite, gentes, Deum nostrum: * et audítam fácite vocem laudis ejus,

Qui pósuit ánimam meam ad vitam: * et non dedit in commotiónem pedes meos.

Quóniam probásti nos, Deus: * igne nos examinásti, sicut examinátur argéntum.

Induxísti nos in láqueum, posuísti tribulatiónes in dorso nostro: * imposuísti hómines super cápita nostra.

Transívimus per ignem et aquam: * et eduxísti nos in refrigérium.

\antlabel{Ant.} Vidéte ópera Dómini, et audítam fácite vocem laudis ejus.

\antlabel{Ant.} Audíte, omnes, * qui timétis Deum, quanta fecit ánimæ meæ.

\psalmheading{Psalmus 65, ii}

\lettrine[lines=2]{I}{NTROIBO} in domum tuam in holocáustis: * reddam tibi vota mea, quæ distinxérunt lábia mea.

Et locútum est os meum, * in tribulatióne mea.

Holocáusta medulláta ófferam tibi cum incénso aríetum: * ófferam tibi boves cum hircis.

Veníte, audíte, et narrábo, omnes, qui timétis Deum, * quanta fecit ánimæ meæ.

Ad ipsum ore meo clamávi, * et exaltávi sub lingua mea.

Iniquitátem si aspéxi in corde meo, * non exáudiet Dóminus.

Proptérea exaudívit Deus, * et atténdit voci deprecatiónis meæ.

Benedíctus Deus, * qui non amóvit oratiónem meam, et misericórdiam suam a me.

\antlabel{Ant.} Audíte, omnes, qui timétis Deum, quanta fecit ánimæ meæ.

\variantheading{In Festis novem Lectionum:}

\V\ Non amóvit Dóminus oratiónem meam. \R\ Et misericórdiam suam a me.

\absolutiolabel{Absolutio.} Exáudi, Dómine.

\benedlabel{Bened.} 1 Benedictióne perpétua. 2 Unigénitus Dei Fílius. 3 Spíritus Sancti grátia.

\nocturnoheading{In II Nocturno}

\antlabel{Ant.} Exsúrgat Deus, * \textdagger\ et dissipéntur inimíci ejus. \textdaggerdbl

\psalmheading{Psalmus 67, i}

\lettrine[lines=2]{E}{XSURGAT} Deus, \textdagger\ et dissipéntur inimíci ejus, * \textdaggerdbl\ et fúgiant qui odérunt eum, a fácie ejus.

Sicut déficit fumus, defíciant: * sicut fluit cera a fácie ignis, sic péreant peccatóres a fácie Dei.

Et justi epuléntur, et exsúltent in conspéctu Dei: * et delecténtur in lætítia.

Cantáte Deo, psalmum dícite nómini ejus: * iter fácite ei, qui ascéndit super occásum: Dóminus nomen illi.

Exsultáte in conspéctu ejus: * turbabúntur a fácie ejus, patris orphanórum, et júdicis viduárum.

Deus in loco sancto suo: * Deus qui inhabitáre facit uníus moris in domo:

Qui edúcit vinctos in fortitúdine, * simíliter eos, qui exásperant, qui hábitant in sepúlcris.

Deus, cum egrederéris in conspéctu pópuli tui, * cum pertransíres in desérto:

Terra mota est, étenim cæli distillavérunt a fácie Dei Sínaï, * a fácie Dei Israël.

Plúviam voluntáriam segregábis, Deus, hereditáti tuæ: * et infirmáta est, tu vero perfecísti eam.

Animália tua habitábunt in ea: * parásti in dulcédine tua páuperi, Deus.

\antlabel{Ant.} Exsúrgat Deus, et dissipéntur inimíci ejus.

\antlabel{Ant.} Deus noster, * Deus salvos faciéndi: et Dómini sunt éxitus mortis.

\psalmheading{Psalmus 67, ii}

\lettrine[lines=2]{D}{OMINUS} dabit verbum evangelizántibus, * virtúte multa.

Rex virtútum dilécti dilécti: * et speciéi domus divídere spólia.

Si dormiátis inter médios cleros, pennæ colúmbæ deargentátæ, * et posterióra dorsi ejus in pallóre auri.

Dum discérnit cæléstis reges super eam, nive dealbabúntur in Selmon: * mons Dei, mons pinguis.

Mons coagulátus, mons pinguis: * ut quid suspicámini montes coagulátos?

Mons, in quo beneplácitum est Deo habitáre in eo: * étenim Dóminus habitábit in finem.

Currus Dei decem míllibus múltiplex, míllia lætántium: * Dóminus in eis in Sina in sancto.

Ascendísti in altum, cepísti captivitátem: * accepísti dona in homínibus:

Etenim non credéntes, * inhabitáre Dóminum Deum.

Benedíctus Dóminus die quotídie: * prósperum iter fáciet nobis Deus salutárium nostrórum.

Deus noster, Deus salvos faciéndi: * et Dómini Dómini éxitus mortis.

Verúmtamen Deus confrínget cápita inimicórum suórum: * vérticem capílli perambulántium in delíctis suis.

Dixit Dóminus: Ex Basan convértam, * convértam in profúndum maris:

Ut intingátur pes tuus in sánguine: * lingua canum tuórum ex inimícis, ab ipso.

\antlabel{Ant.} Deus noster, Deus salvos faciéndi: et Dómini sunt éxitus mortis.

\antlabel{Ant.} In ecclésiis * benedícite Dómino Deo.

\psalmheading{Psalmus 67, iii}

\lettrine[lines=2]{V}{IDERUNT} ingréssus tuos, Deus, * ingréssus Dei mei: regis mei qui est in sancto.

Prævenérunt príncipes conjúncti psalléntibus, * in médio juvenculárum tympanistriárum.

In ecclésiis benedícite Deo Dómino, * de fóntibus Israël.

Ibi Bénjamin adolescéntulus, * in mentis excéssu.

Príncipes Juda, duces eórum: * príncipes Zábulon, príncipes Néphthali.

Manda, Deus, virtúti tuæ: * confírma hoc, Deus, quod operátus es in nobis.

A templo tuo in Jerúsalem, * tibi ófferent reges múnera.

Íncrepa feras arúndinis, congregátio taurórum in vaccis populórum: * ut exclúdant eos, qui probáti sunt argénto.

Díssipa gentes, quæ bella volunt: vénient legáti ex Ægýpto: * Æthiópia prævéniet manus ejus Deo.

Regna terræ, cantáte Deo: * psállite Dómino.

Psállite Deo, qui ascéndit super cælum cæli, * ad Oriéntem.

Ecce, dabit voci suæ vocem virtútis, date glóriam Deo super Israël, * magnificéntia ejus, et virtus ejus in núbibus.

Mirábilis Deus in sanctis suis, Deus Israël ipse dabit virtútem, et fortitúdinem plebi suæ, * benedíctus Deus.

\antlabel{Ant.} In ecclésiis benedícite Dómino Deo.

\variantheading{In Festis novem Lectionum:}

\V\ Mirábilis Deus in sanctis suis. \R\ Deus Israël dabit virtútem plebi tuæ.

\absolutiolabel{Absolutio.} Ipsíus píetas et.

\benedlabel{Bened.} 4 Deus Pater omnípotens. 5 Christus perpétuæ. 6 Ignem sui amóris.

\nocturnoheading{In III Nocturno}

\antlabel{Ant.} Salvum me fac, * \textdagger\ Deus: quóniam intravérunt aquæ usque ad ánimam meam. \textdaggerdbl

\psalmheading{Psalmus 68, i}

\lettrine[lines=2]{S}{ALVUM} me fac, \textdagger\ Deus: * \textdaggerdbl\ quóniam intravérunt aquæ usque ad ánimam meam.

Infíxus sum in limo profúndi: * et non est substántia.

Veni in altitúdinem maris: * et tempéstas demérsit me.

Laborávi clamans, raucæ factæ sunt fauces meæ: * defecérunt óculi mei, dum spero in Deum meum.

Multiplicáti sunt super capíllos cápitis mei, * qui odérunt me gratis.

Confortáti sunt qui persecúti sunt me inimíci mei injúste: * quæ non rápui, tunc exsolvébam.

Deus, tu scis insipiéntiam meam: * et delícta mea a te non sunt abscóndita.

Non erubéscant in me qui exspéctant te, Dómine, * Dómine virtútum.

Non confundántur super me * qui quærunt te, Deus Israël.

Quóniam propter te sustínui oppróbrium: * opéruit confúsio fáciem meam.

Extráneus factus sum frátribus meis, * et peregrínus fíliis matris meæ.

Quóniam zelus domus tuæ comédit me: * et oppróbria exprobrántium tibi cecidérunt super me.

Et opérui in jejúnio ánimam meam: * et factum est in oppróbrium mihi.

Et pósui vestiméntum meum cilícium: * et factus sum illis in parábolam.

Advérsum me loquebántur, qui sedébant in porta: * et in me psallébant qui bibébant vinum.

Ego vero oratiónem meam ad te, Dómine: * tempus benepláciti, Deus.

\antlabel{Ant.} Salvum me fac, Deus: quóniam intravérunt aquæ usque ad ánimam meam.

\antlabel{Ant.} Propter inimícos meos * éripe me, Dómine.

\psalmheading{Psalmus 68, ii}

\lettrine[lines=2]{I}{N} multitúdine misericórdiæ tuæ exáudi me, * in veritáte salútis tuæ:

Eripe me de luto, ut non infígar: * líbera me ab iis, qui odérunt me, et de profúndis aquárum.

Non me demérgat tempéstas aquæ, neque absórbeat me profúndum: * neque úrgeat super me púteus os suum.

Exáudi me, Dómine, quóniam benígna est misericórdia tua: * secúndum multitúdinem miseratiónum tuárum réspice in me.

Et ne avértas fáciem tuam a púero tuo: * quóniam tríbulor, velóciter exáudi me.

Inténde ánimæ meæ, et líbera eam: * propter inimícos meos éripe me.

Tu scis impropérium meum, et confusiónem meam, * et reveréntiam meam.

In conspéctu tuo sunt omnes qui tríbulant me, * impropérium exspectávit cor meum, et misériam.

Et sustínui qui simul contristarétur, et non fuit: * et qui consolarétur, et non invéni.

Et dedérunt in escam meam fel: * et in siti mea potavérunt me acéto.

Fiat mensa eórum coram ipsis in láqueum, * et in retributiónes, et in scándalum.

Obscuréntur óculi eórum ne vídeant: * et dorsum eórum semper incúrva.

Effúnde super eos iram tuam: * et furor iræ tuæ comprehéndat eos.

Fiat habitátio eórum desérta: * et in tabernáculis eórum non sit qui inhábitet.

Quóniam quem tu percussísti, persecúti sunt: * et super dolórem vúlnerum meórum addidérunt.

Appóne iniquitátem super iniquitátem eórum: * et non intrent in justítiam tuam.

Deleántur de libro vivéntium: * et cum justis non scribántur.

\antlabel{Ant.} Propter inimícos meos éripe me, Dómine.

\antlabel{Ant.} Quærite Dóminum, * et vivet ánima vestra.

\psalmheading{Psalmus 68, iii}

\lettrine[lines=2]{E}{GO} sum pauper et dolens: * salus tua, Deus, suscépit me.

Laudábo nomen Dei cum cántico: * et magnificábo eum in laude:

Et placébit Deo super vítulum novéllum: * córnua producéntem et úngulas.

Vídeant páuperes et læténtur: * quærite Deum, et vivet ánima vestra.

Quóniam exaudívit páuperes Dóminus: * et vinctos suos non despéxit.

Laudent illum cæli et terra, * mare, et ómnia reptília in eis.

Quóniam Deus salvam fáciet Sion: * et ædificabúntur civitátes Juda.

Et inhabitábunt ibi, * et hereditáte acquírent eam.

Et semen servórum ejus possidébit eam, * et qui díligunt nomen ejus, habitábunt in ea.

\antlabel{Ant.} Quærite Dóminum, et vivet ánima vestra.

\variantheading{In feriali Officio et in Festis:}

\V\ Laudábo nomen Dei cum cántico. \R\ Et magnificábo eum in laude.

\variantheading{In Officio novem Lectionum}

\absolutiolabel{Absolutio.} A vínculis.

\benedlabel{Bened.} 7 Evangélica léctio. 8 In Festis Dómini: Divínum auxílium. In Festis B. M. V.: Cujus festum cólimus, ipsa Virgo. In Festis Sanctórum: Cujus (\textit{vel} Quorum \textit{aut} Quarum) festum cólimus. 9 Ad societátem. (Per evangélica dicta.)

\variantheading{In Officio trium Lectionum}

\absolutiolabel{Absolutio.} Exáudi, Dómine.

\variantheading{In Feriis cum Homilia:}

\benedlabel{Bened.} 1 Evangélica. 2 Divínum. 3 Ad societátem.

\variantheading{In Feriis sine Homilia:}

\benedlabel{Bened.} 1 Benedictióne. 2 Unigénitus. 3 Spíritus Sancti.

\variantheading{In Festis Sanctorum:}

\benedlabel{Bened.} 1 Ille nos benedícat, qui sine fine vivit et regnat. \R\ Amen. 2 Cujus (\textit{vel} Quorum \textit{aut} Quarum). 3 Ad societátem.

% ===== END OF BATCH 24 -- Covers autumnalis.md lines 8133-8494,
% opening "Feria Quinta" (Thursday) and completing its "Ad
% Matutinum" in full: the ferial Invitatorium "Regem ventúrum
% Dóminum," the hymn "Nox atra rerum cóntegit," and the three
% Nocturns (Psalms 61, 65/i, 65/ii; 67/i-iii; 68/i-iii) with their
% versicle, Absolutio, and Benedictiones apparatus. "In III
% Nocturno" here has only a single form (no "I"/"II" alternation as
% in Feria Quarta), confirming that split was specific to the
% Psalm-49/Psalm-50 overlap and does not recur every day. Several
% more psalm incipits carry the partial-caps OCR defect already
% documented in prior batches ("NOnne," "JUbiláte," "INtroíbo,"
% "EXsurgat," "DOminus," "VIderunt," "SAlvum," "IN" multitúdine,
% "EGo") and are silently normalized to full caps. The document
% continues with "Ad Laudes I" (autumnalis.md line 8497).
% =====
\hourheading{Ad Laudes I}

\rubric{In Fériis extra Vigílias commúnes occurréntibus, atque in Festis.}

(Pater noster et Ave.)

\V\ Deus, in adjutórium.

\antlabel{Ant.} Jubiláte * in conspéctu regis Dómini.

\psalmheading{Psalmus 97}

\lettrine[lines=2]{C}{ANTATE} Dómino cánticum novum: * quia mirabília fecit.

Salvávit sibi déxtera ejus: * et bráchium sanctum ejus.

Notum fecit Dóminus salutáre suum: * in conspéctu géntium revelávit justítiam suam.

Recordátus est misericórdiæ suæ, * et veritátis suæ dómui Israël.

Vidérunt omnes términi terræ * salutáre Dei nostri.

Jubiláte Deo, omnis terra: * cantáte, et exsultáte, et psállite.

Psállite Dómino in cíthara, in cíthara et voce psalmi: * in tubis ductílibus, et voce tubæ córneæ.

Jubiláte in conspéctu regis Dómini: * moveátur mare, et plenitúdo ejus: orbis terrárum, et qui hábitant in eo.

Flúmina plaudent manu, simul montes exsultábunt a conspéctu Dómini: * quóniam venit judicáre terram.

Judicábit orbem terrárum in justítia, * et pópulos in æquitáte.

\antlabel{Ant.} Jubiláte in conspéctu regis Dómini.

\antlabel{Ant.} Dómine, * \textdagger\ refúgium factus es nobis. \textdaggerdbl

\psalmheading{Psalmus 89}

\lettrine[lines=2]{D}{OMINE}, \textdagger\ refúgium factus es nobis: * \textdaggerdbl\ a generatióne in generatiónem.

Priúsquam montes fíerent, aut formarétur terra et orbis: * a sæculo et usque in sæculum tu es, Deus.

Ne avértas hóminem in humilitátem: * et dixísti: Convertímini, fílii hóminum.

Quóniam mille anni ante óculos tuos, * tamquam dies hestérna, quæ prætériit.

Et custódia in nocte, * quæ pro níhilo habéntur, eórum anni erunt.

Mane sicut herba tránseat, mane flóreat, et tránseat: * véspere décidat, indúret et aréscat.

Quia defécimus in ira tua, * et in furóre tuo turbáti sumus.

Posuísti iniquitátes nostras in conspéctu tuo: * sæculum nostrum in illuminatióne vultus tui.

Quóniam omnes dies nostri defecérunt: * et in ira tua defécimus.

Anni nostri sicut aránea meditabúntur: * dies annórum nostrórum in ipsis, septuagínta anni.

Si autem in potentátibus octogínta anni: * et ámplius eórum, labor et dolor.

Quóniam supervénit mansuetúdo: * et corripiémur.

Quis novit potestátem iræ tuæ: * et præ timóre tuo iram tuam dinumeráre?

Déxteram tuam sic notam fac: * et erudítos corde in sapiéntia.

Convértere, Dómine, úsquequo? * et deprecábilis esto super servos tuos.

Repléti sumus mane misericórdia tua: * et exsultávimus, et delectáti sumus ómnibus diébus nostris.

Lætáti sumus pro diébus, quibus nos humiliásti: * annis, quibus vídimus mala.

Réspice in servos tuos, et in ópera tua: * et dírige fílios eórum.

Et sit splendor Dómini, Dei nostri, super nos, et ópera mánuum nostrárum dírige super nos: * et opus mánuum nostrárum dírige.

\antlabel{Ant.} Dómine, refúgium factus es nobis.

\antlabel{Ant.} Dómine, * in cælo misericórdia tua.

\psalmheading{Psalmus 35}

\lettrine[lines=2]{D}{IXIT} injústus ut delínquat in semetípso: * non est timor Dei ante óculos ejus.

Quóniam dolóse egit in conspéctu ejus: * ut inveniátur iníquitas ejus ad ódium.

Verba oris ejus iníquitas, et dolus: * nóluit intellígere ut bene ágeret.

Iniquitátem meditátus est in cubíli suo: * ástitit omni viæ non bonæ, malítiam autem non odívit.

Dómine, in cælo misericórdia tua: * et véritas tua usque ad nubes.

Justítia tua sicut montes Dei: * judícia tua abýssus multa.

Hómines, et juménta salvábis, Dómine: * quemádmodum multiplicásti misericórdiam tuam, Deus.

Fílii autem hóminum, * in tégmine alárum tuárum sperábunt.

Inebriabúntur ab ubertáte domus tuæ: * et torrénte voluptátis tuæ potábis eos.

Quóniam apud te est fons vitæ: * et in lúmine tuo vidébimus lumen.

Præténde misericórdiam tuam sciéntibus te, * et justítiam tuam his, qui recto sunt corde.

Non véniat mihi pes supérbiæ: * et manus peccatóris non móveat me.

Ibi cecidérunt qui operántur iniquitátem: * expúlsi sunt, nec potuérunt stare.

\antlabel{Ant.} Dómine, in cælo misericórdia tua.

\antlabel{Ant.} Pópulus meus, * ait Dóminus, bonis meis adimplébitur.

\canticumheading{Canticum Jeremiæ}{Jer. 31, 10-14}

\lettrine[lines=2]{A}{UDITE} verbum Dómini, gentes, * et annuntiáte in ínsulis, quæ procul sunt,

Et dícite: Qui dispérsit Israël, congregábit eum: * et custódiet eum sicut pastor gregem suum.

Redémit enim Dóminus Jacob, * et liberábit eum de manu potentióris.

Et vénient, et laudábunt in monte Sion: * et confluent ad bona Dómini.

Super fruménto, et vino, et óleo, * et fetu pécorum et armentórum:

Erítque ánima eórum quasi hortus irríguus, * et ultra non esúrient.

Tunc lætábitur virgo in choro, * júvenes et senes simul:

Et convértam luctum eórum in gáudium, * et consolábor eos, et lætificábo a dolóre suo.

Et inebriábo ánimam sacerdótum pinguédine: * et pópulus meus bonis meis adimplébitur.

\antlabel{Ant.} Pópulus meus, ait Dóminus, bonis meis adimplébitur.

\antlabel{Ant.} Deo nostro * jucúnda sit laudátio.

\psalmheading{Psalmus 146}

\lettrine[lines=2]{L}{AUDATE} Dóminum quóniam bonus est psalmus: * Deo nostro sit jucúnda, decóraque laudátio.

Ædíficans Jerúsalem Dóminus: * dispersiónes Israélis congregábit.

Qui sanat contrítos corde: * et álligat contritiónes eórum.

Qui númerat multitúdinem stellárum: * et ómnibus eis nómina vocat.

Magnus Dóminus noster, et magna virtus ejus: * et sapiéntiæ ejus non est númerus.

Suscípiens mansuétos Dóminus: * humílians autem peccatóres usque ad terram.

Præcínite Dómino in confessióne: * psállite Deo nostro in cíthara.

Qui óperit cælum núbibus: * et parat terræ plúviam.

Qui prodúcit in móntibus fœnum: * et herbam servitúti hóminum.

Qui dat juméntis escam ipsórum: * et pullis corvórum invocántibus eum.

Non in fortitúdine equi voluntátem habébit: * nec in tíbiis viri beneplácitum erit ei.

Beneplácitum est Dómino super timéntes eum: * et in eis, qui sperant super misericórdia ejus.

\antlabel{Ant.} Deo nostro jucúnda sit laudátio.

\variantheading{In Officio feriali}

\capitulumlabel{Capit.}{Rom. 13, 12-13}

\lettrine[lines=2]{N}{OX} præcéssit, dies autem appropinquávit. Abjiciámus ergo ópera tenebrárum, et induámur arma lucis. Sicut in die honéste ambulémus.

\hymnusheading{Hymnus}

\lettrine[lines=2]{L}{UX} ecce surgit áurea,\\
Pallens facéssat cæcitas,\\
Quæ nosmet in præceps diu\\
Erróre traxit dévio.

Hæc lux serénum cónferat,\\
Purósque nos præstet sibi:\\
Nihil loquámur súbdolum:\\
Volvámus obscúrum nihil.

Sic tota decúrrat dies,\\
Ne lingua mendax, ne manus\\
Oculíve peccent lúbrici,\\
Ne noxa corpus ínquinet.

Speculátor astat désuper,\\
Qui nos diébus ómnibus,\\
Actúsque nostros próspicit\\
A luce prima in vésperum.

Deo Patri sit glória,\\
Ejúsque soli Fílio,\\
Cum Spíritu Paráclito,\\
Nunc et per omne sæculum.\\
Amen.

\V\ Repléti sumus mane misericórdia tua. \R\ Exsultávimus, et delectáti sumus.

\antlabel{Ad Bened. Ant.} In sanctitáte * serviámus Dómino, et liberábit nos ab inimícis nostris.

\rubric{Réliqua ut in Ordinário. 9}

% ===== END OF BATCH 25 -- Covers autumnalis.md lines 8497-8702, "Ad
% Laudes I" for Feria Quinta in full: Psalms 97, 89, 35, the
% Canticum Jeremiæ, and Psalm 146, followed by the Capitulum "Nox
% præcéssit" and the hymn "Lux ecce surgit áurea." Several more
% psalm incipits carry the partial-caps OCR defect already
% documented in prior batches ("DOmine," "DIxit," "AUdíte,"
% "LAudáte," "NOx" [Capitulum]) and are silently normalized to full
% caps. The document continues with "Ad Laudes II" (autumnalis.md
% line 8705).
% =====
\hourheading{Ad Laudes II}

\rubric{In Vigíliis commúnibus, quando fit Offícium de Féria.}

(Pater noster et Ave.)

\V\ Deus, in adjutórium.

\antlabel{Ant.} Tibi soli peccávi.

\psalmheading{Psalmus 50}

\lettrine[lines=2]{M}{ISERERE} mei, Deus, * secúndum magnam misericórdiam tuam.

Et secúndum multitúdinem miseratiónum tuárum, * dele iniquitátem meam.

Amplius lava me ab iniquitáte mea: * et a peccáto meo munda me.

Quóniam iniquitátem meam ego cognósco: * et peccátum meum contra me est semper.

Tibi soli peccávi, et malum coram te feci: * ut justificéris in sermónibus tuis, et vincas cum judicáris.

Ecce enim, in iniquitátibus concéptus sum: * et in peccátis concépit me mater mea.

Ecce enim, veritátem dilexísti: * incérta et occúlta sapiéntiæ tuæ manifestásti mihi.

Aspérges me hyssópo, et mundábor: * lavábis me, et super nivem dealbábor.

Audítui meo dabis gáudium et lætítiam: * et exsultábunt ossa humiliáta.

Avérte fáciem tuam a peccátis meis: * et omnes iniquitátes meas dele.

Cor mundum crea in me, Deus: * et spíritum rectum ínnova in viscéribus meis.

Ne projícias me a fácie tua: * et spíritum sanctum tuum ne áuferas a me.

Redde mihi lætítiam salutáris tui: * et spíritu principáli confírma me.

Docébo iníquos vias tuas: * et ímpii ad te converténtur.

Líbera me de sanguínibus, Deus, Deus salútis meæ: * et exsultábit lingua mea justítiam tuam.

Dómine, lábia mea apéries: * et os meum annuntiábit laudem tuam.

Quóniam si voluísses sacrifícium, dedíssem útique: * holocáustis non delectáberis.

Sacrifícium Deo spíritus contribulátus: * cor contrítum, et humiliátum, Deus, non despícies.

Benígne fac, Dómine, in bona voluntáte tua Sion: * ut ædificéntur muri Jerúsalem.

Tunc acceptábis sacrifícium justítiæ, oblatiónes, et holocáusta: * tunc impónent super altáre tuum vítulos.

\antlabel{Ant.} Tibi soli peccávi, Dómine: miserére mei.

\antlabel{Ant.} Convértere, Dómine.

\psalmheading{Psalmus 89}

\lettrine[lines=2]{D}{OMINE}, refúgium factus es nobis: * a generatióne in generatiónem.

Priúsquam montes fíerent, aut formarétur terra et orbis: * a sæculo et usque in sæculum tu es, Deus.

Ne avértas hóminem in humilitátem: * et dixísti: Convertímini, fílii hóminum.

Quóniam mille anni ante óculos tuos, * tamquam dies hestérna, quæ prætériit.

Et custódia in nocte, * quæ pro níhilo habéntur, eórum anni erunt.

Mane sicut herba tránseat, mane flóreat, et tránseat: * véspere décidat, indúret et aréscat.

Quia defécimus in ira tua, * et in furóre tuo turbáti sumus.

Posuísti iniquitátes nostras in conspéctu tuo: * sæculum nostrum in illuminatióne vultus tui.

Quóniam omnes dies nostri defecérunt: * et in ira tua defécimus.

Anni nostri sicut aránea meditabúntur: * dies annórum nostrórum in ipsis, septuagínta anni.

Si autem in potentátibus, octogínta anni: * et ámplius eórum, labor et dolor.

Quóniam supervénit mansuetúdo: * et corripiémur.

Quis novit potestátem iræ tuæ: * et præ timóre tuo iram tuam dinumeráre?

Déxteram tuam sic notam fac: * et erudítos corde in sapiéntia.

Convértere, Dómine, úsquequo? * et deprecábilis esto super servos tuos.

Repléti sumus mane misericórdia tua: * et exsultávimus, et delectáti sumus ómnibus diébus nostris.

Lætáti sumus pro diébus, quibus nos humiliásti: * annis, quibus vídimus mala.

Réspice in servos tuos, et in ópera tua: * et dírige fílios eórum.

Et sit splendor Dómini, Dei nostri, super nos, et ópera mánuum nostrárum dírige super nos: * et opus mánuum nostrárum dírige.

\antlabel{Ant.} Convértere, Dómine, et deprecábilis esto super servos tuos.

\antlabel{Ant.} Multiplicásti, Deus.

\psalmheading{Psalmus 35}

\lettrine[lines=2]{D}{IXIT} injústus ut delínquat in semetípso: * non est timor Dei ante óculos ejus.

Quóniam dolóse egit in conspéctu ejus: * ut inveniátur iníquitas ejus ad ódium.

Verba oris ejus iníquitas, et dolus: * nóluit intellígere ut bene ágeret.

Iniquitátem meditátus est in cubíli suo: * ástitit omni viæ non bonæ, malítiam autem non odívit.

Dómine, in cælo misericórdia tua: * et véritas tua usque ad nubes.

Justítia tua sicut montes Dei: * judícia tua abýssus multa.

Hómines, et juménta salvábis, Dómine: * quemádmodum multiplicásti misericórdiam tuam, Deus.

Fílii autem hóminum, * in tégmine alárum tuárum sperábunt.

Inebriabúntur ab ubertáte domus tuæ: * et torrénte voluptátis tuæ potábis eos.

Quóniam apud te est fons vitæ: * et in lúmine tuo vidébimus lumen.

Præténde misericórdiam tuam sciéntibus te, * et justítiam tuam his, qui recto sunt corde.

Non véniat mihi pes supérbiæ: * et manus peccatóris non móveat me.

Ibi cecidérunt qui operántur iniquitátem: * expúlsi sunt, nec potuérunt stare.

\antlabel{Ant.} Multiplicásti, Deus, misericórdiam tuam.

\antlabel{Ant.} Fortitúdo mea.

\canticumheading{Canticum Moysis}{Exodi 15, 1-19}

\lettrine[lines=2]{C}{ANTEMUS} Dómino: glorióse enim magnificátus est, * equum et ascensórem dejécit in mare.

Fortitúdo mea, et laus mea Dóminus, * et factus est mihi in salútem:

Iste Deus meus, et glorificábo eum: * Deus patris mei, et exaltábo eum.

Dóminus quasi vir pugnátor, Omnípotens nomen ejus. * Currus Pharaónis et exércitum ejus projécit in mare.

Elécti príncipes ejus submérsi sunt in Mari Rubro: * abýssi operuérunt eos, descendérunt in profúndum quasi lapis.

Déxtera tua, Dómine, magnificáta est in fortitúdine: déxtera tua, Dómine, percússit inimícum. * Et in multitúdine glóriæ tuæ deposuísti adversários tuos:

Misísti iram tuam, quæ devorávit eos sicut stípulam. * Et in spíritu furóris tui congregátæ sunt aquæ:

Stetit unda fluens, * congregátæ sunt abýssi in médio mari.

Dixit inimícus: Pérsequar et comprehéndam, * dívidam spólia, implébitur ánima mea:

Evaginábo gládium meum, * interfíciet eos manus mea.

Flavit spíritus tuus, et opéruit eos mare: * submérsi sunt quasi plumbum in aquis vehemméntibus.

Quis símilis tui in fórtibus, Dómine? * quis símilis tui, magníficus in sanctitáte, terríbilis atque laudábilis, fáciens mirabília?

Extendísti manum tuam, et devorávit eos terra. * Dux fuísti in misericórdia tua pópulo quem redemísti:

Et portásti eum in fortitúdine tua, * ad habitáculum sanctum tuum.

Ascendérunt pópuli, et iráti sunt: * dolóres obtinuérunt habitatóres Philísthiim.

Tunc conturbáti sunt príncipes Edom, robústos Moab obtínuit tremor: * obriguérunt omnes habitatóres Chánaan.

Irruat super eos formído et pavor, * in magnitúdine bráchii tui.

Fiant immóbiles quasi lapis, donec pertránseat pópulus tuus, Dómine, * donec pertránseat pópulus tuus iste, quem possedísti.

Introdúces eos, et plantábis in monte hereditátis tuæ, * firmíssimo habitáculo tuo quod operátus es, Dómine:

Sanctuárium tuum, Dómine, quod firmavérunt manus tuæ. * Dóminus regnábit in ætérnum et ultra.

Ingréssus est enim eques Phárao cum cúrribus et equítibus ejus in mare: * et redúxit super eos Dóminus aquas maris:

Fílii autem Israël ambulavérunt per siccum * in médio ejus.

\antlabel{Ant.} Fortitúdo mea, et laus mea Dóminus: et factus est mihi in salútem.

\antlabel{Ant.} Laudáte. \textdagger

\psalmheading{Psalmus 146}

\lettrine[lines=2]{L}{AUDATE} \textdagger\ Dóminum quóniam bonus est psalmus: * Deo nostro sit jucúnda, decóraque laudátio.

Ædíficans Jerúsalem Dóminus: * dispersiónes Israélis congregábit.

Qui sanat contrítos corde: * et álligat contritiónes eórum.

Qui númerat multitúdinem stellárum: * et ómnibus eis nómina vocat.

Magnus Dóminus noster, et magna virtus ejus: * et sapiéntiæ ejus non est númerus.

Suscípiens mansuétos Dóminus: * humílians autem peccatóres usque ad terram.

Præcínite Dómino in confessióne: * psállite Deo nostro in cíthara.

Qui óperit cælum núbibus: * et parat terræ plúviam.

Qui prodúcit in móntibus fœnum: * et herbam servitúti hóminum.

Qui dat juméntis escam ipsórum: * et pullis corvórum invocántibus eum.

Non in fortitúdine equi voluntátem habébit: * nec in tíbiis viri beneplácitum erit ei.

Beneplácitum est Dómino super timéntes eum: * et in eis, qui sperant super misericórdia ejus.

\antlabel{Ant.} Laudáte Dóminum, qui sanat contrítos corde, et álligat contritiónes eórum.

\capitulumlabel{Capit.}{Rom. 13, 12-13}

\lettrine[lines=2]{N}{OX} præcéssit, dies autem appropinquávit. Abjiciámus ergo ópera tenebrárum, et induámur arma lucis. Sicut in die honéste ambulémus.

\hymnusheading{Hymnus}

\lettrine[lines=2]{L}{UX} ecce surgit áurea,\\
Pallens facéssat cæcitas,\\
Quæ nosmet in præceps diu\\
Erróre traxit dévio.

Hæc lux serénum cónferat,\\
Purósque nos præstet sibi:\\
Nihil loquámur súbdolum:\\
Volvámus obscúrum nihil.

Sic tota decúrrat dies,\\
Ne lingua mendax, ne manus\\
Oculíve peccent lúbrici,\\
Ne noxa corpus ínquinet,

Speculátor astat désuper,\\
Qui nos diébus ómnibus,\\
Actúsque nostros próspicit\\
A luce prima in vésperum.

Deo Patri sit glória,\\
Ejúsque soli Fílio,\\
Cum Spíritu Paráclito,\\
Nunc et per omne sæculum.\\
Amen.

\V\ Repléti sumus mane misericórdia tua. \R\ Exsultávimus, et delectáti sumus.

\antlabel{Ad Bened. Ant.} In sanctitáte * serviámus Dómino, et liberábit nos ab inimícis nostris.

\rubric{Expleto Cántico ac repetíta Antíphona, dicúntur sequéntes}

\variantheading{Preces feriales}

KYRIE, eléison. Christe, eléison. Kýrie, eléison.

Pater noster \rubric{quod dícitur a solo Hebdomadário totum clara voce, usque ad} \V\ Et ne nos indúcas in tentatiónem. \R\ Sed líbera nos a malo. \V\ Ego dixi: Dómine, miserére mei. \R\ Sana ánimam meam, quia peccávi tibi. \V\ Convértere, Dómine, úsquequo? \R\ Et deprecábilis esto super servos tuos. \V\ Fiat misericórdia tua, Dómine, super nos. \R\ Quemádmodum sperávimus in te. \V\ Sacerdótes tui induántur justítiam. \R\ Et sancti tui exsúltent. \V\ Orémus pro beatíssimo Papa nostro N. \R\ Dóminus consérvet eum, et vivíficet eum, et beátum fáciat eum in terra, et non tradat eum in ánimam inimicórum ejus. \V\ Orémus et pro Antístite nostro N. \R\ Stet et pascat in fortitúdine tua, Dómine, in sublimitáte nóminis tui. \V\ Dómine, salvum fac regem. \R\ Et exáudi nos in die, qua invocavérimus te. \V\ Salvum fac pópulum tuum, Dómine, et bénedic hereditáti tuæ. \R\ Et rege eos, et extólle illos usque in ætérnum. \V\ Meménto Congregatiónis tuæ. \R\ Quam possedísti ab inítio. \V\ Fiat pax in virtúte tua. \R\ Et abundántia in túrribus tuis. \V\ Orémus pro benefactóribus nostris. \R\ Retribúere dignáre, Dómine, ómnibus, nobis bona faciéntibus propter nomen tuum, vitam ætérnam. Amen. \V\ Orémus pro fidélibus defúnctis. \R\ Réquiem ætérnam dona eis, Dómine, et lux perpétua lúceat eis. \V\ Requiéscant in pace. \R\ Amen. \V\ Pro frátribus nostris abséntibus. \R\ Salvos fac servos tuos, Deus meus, sperántes in te. \V\ Pro afflíctis et captívis. \R\ Líbera eos, Deus Israël, ex ómnibus tribulatiónibus suis. \V\ Mitte eis, Dómine, auxílium de sancto. \R\ Et de Sion tuére eos. \V\ Dómine, Deus virtútum, convérte nos. \R\ Et osténde fáciem tuam, et salvi érimus. \V\ Exsúrge, Christe, ádjuva nos. \R\ Et líbera nos propter nomen tuum. \V\ Dómine, exáudi oratiónem meam. \R\ Et clamor meus ad te véniat.

\rubric{Réliqua ut in Ordinário. 9}

% ===== END OF BATCH 26 -- Covers autumnalis.md lines 8705-8964, "Ad
% Laudes II" for Feria Quinta in full: Psalms 50, 89, 35, the
% Canticum Moysis ("Cantémus Dómino"), and Psalm 146, sharing the
% Capitulum "Nox præcéssit" and the hymn "Lux ecce surgit áurea"
% with Ad Laudes I, plus the "Preces feriales" litany. Psalm 89
% appears here without the dagger/double-dagger chant-pointing
% marks present in its Ad Laudes I occurrence (Batch 25), confirmed
% as a genuine source divergence rather than an omission. The
% document continues with "Ad Primam" (autumnalis.md line 8967).
% =====
\hourheading{Ad Primam}

Pater, Ave et Credo.

\V\ Deus, in adjutórium.

\hymnusheading{Hymnus}

\lettrine[lines=2]{J}{AM} lucis orto sídere,\\
Deum precémur súpplices,\\
Ut in diúrnis áctibus\\
Nos servet a nocéntibus.

Linguam refrénans témperet,\\
Ne litis horror ínsonet:\\
Visum fovéndo cóntegat,\\
Ne vanitátes háuriat.

Sint pura cordis íntima,\\
Absístat et vecórdia:\\
Carnis terat supérbiam\\
Potus cibíque párcitas.

Ut, cum dies abscésserit\\
Noctémque sors redúxerit,\\
Mundi per abstinéntiam\\
Ipsi canámus glóriam.

Deo Patri sit glória,\\
Ejúsque soli Fílio,\\
Cum Spíritu Paráclito,\\
Nunc et per omne sæculum.\\
Amen.

\antlabel{Ant.} In loco páscuæ.

\psalmheading{Psalmus 22}

\lettrine[lines=2]{D}{OMINUS} regit me, et nihil mihi déerit: * in loco páscuæ ibi me collocávit.

Super aquam refectiónis educávit me: * ánimam meam convértit.

Dedúxit me super sémitas justítiæ, * propter nomen suum.

Nam, et si ambulávero in médio umbræ mortis, non timébo mala: * quóniam tu mecum es.

Virga tua, et báculus tuus: * ipsa me consoláta sunt.

Parásti in conspéctu meo mensam, * advérsus eos, qui tríbulant me.

Impinguásti in óleo caput meum: * et calix meus inébrians quam præclárus est!

Et misericórdia tua subsequétur me: * ómnibus diébus vitæ meæ:

Et ut inhábitem in domo Dómini, * in longitúdinem diérum.

\psalmheading{Psalmus 71, i}

\lettrine[lines=2]{D}{EUS}, judícium tuum regi da: * et justítiam tuam fílio regis:

Judicáre pópulum tuum in justítia, * et páuperes tuos in judício.

Suscípiant montes pacem pópulo: * et colles justítiam.

Judicábit páuperes pópuli, et salvos fáciet fílios páuperum: * et humiliábit calumniatórem.

Et permanébit cum sole, et ante lunam, * in generatióne et generatiónem.

Descéndet sicut plúvia in vellus: * et sicut stillicídia stillántia super terram.

Oriétur in diébus ejus justítia, et abundántia pacis: * donec auferátur luna.

Et dominábitur a mari usque ad mare: * et a flúmine usque ad términos orbis terrárum.

\psalmheading{Psalmus 71, ii}

\lettrine[lines=2]{C}{ORAM} illo prócident Æthíopes: * et inimíci ejus terram lingent.

Reges Tharsis, et ínsulæ múnera ófferent: * reges Arabum et Saba dona addúcent.

Et adorábunt eum omnes reges terræ: * omnes gentes sérvient ei:

Quia liberábit páuperem a poténte: * et páuperem, cui non erat adjútor.

Parcet páuperi et ínopi: * et ánimas páuperum salvas fáciet.

Ex usúris et iniquitáte rédimet ánimas eórum: * et honorábile nomen eórum coram illo.

Et vivet, et dábitur ei de auro Arábiæ, et adorábunt de ipso semper: * tota die benedícent ei.

Et erit firmaméntum in terra in summis móntium, superextollétur super Líbanum fructus ejus: * et florébunt de civitáte sicut fœnum terræ.

Sit nomen ejus benedíctum in sæcula: * ante solem pérmanet nomen ejus.

Et benedicéntur in ipso omnes tribus terræ: * omnes gentes magnificábunt eum.

Benedíctus Dóminus, Deus Israël, * qui facit mirabília solus:

Et benedíctum nomen majestátis ejus in ætérnum: * et replébitur majestáte ejus omnis terra: fiat, fiat.

\rubric{¶ Quando ad Laudes dictus fúerit Ps. 50 Miserére, hic subjúngitur sequens Ps. 97, qui in Laudibus locum non hábuit.}

\psalmheading{Psalmus 97}

\lettrine[lines=2]{C}{ANTATE} Dómino cánticum novum: * quia mirabília fecit.

Salvávit sibi déxtera ejus: * et bráchium sanctum ejus.

Notum fecit Dóminus salutáre suum: * in conspéctu géntium revelávit justítiam suam.

Recordátus est misericórdiæ suæ, * et veritátis suæ dómui Israël.

Vidérunt omnes términi terræ * salutáre Dei nostri.

Jubiláte Deo, omnis terra: * cantáte, et exsultáte, et psállite.

Psállite Dómino in cíthara, in cíthara et voce psalmi: * in tubis ductílibus, et voce tubæ córneæ.

Jubiláte in conspéctu regis Dómini: * moveátur mare, et plenitúdo ejus: orbis terrárum, et qui hábitant in eo.

Flúmina plaudent manu, simul montes exsultábunt a conspéctu Dómini: * quóniam venit judicáre terram.

Judicábit orbem terrárum in justítia, * et pópulos in æquitáte.

\rubric{Secus vero, tribus Psalmis explétis, statim dícitur Ant.}

\antlabel{Ant.} In loco páscuæ, ibi Dóminus me collocávit.

\rubric{In Officio Festi vel Octávæ:}

\capitulumlabel{Capitulum}{1 Tim. 1, 17}

\lettrine[lines=2]{R}{EGI} sæculórum immortáli et invisíbili, soli Deo honor et glória in sæcula sæculórum. Amen.

\rubric{In ómnibus autem Fériis et in Vigíliis commúnibus:}

\capitulumlabel{Capitulum}{Zach. 8, 19}

\lettrine[lines=2]{P}{ACEM} et veritátem dilígite, ait Dóminus omnípotens.

\R\ br. Christe, Fili Dei vivi, * Miserére nobis. Christe. \V\ \textit{Qui sedes ad déxteram Patris.} Miserére nobis. Glória Patri. Christe. \V\ Exsúrge, Christe, ádjuva nos. \R\ Et líbera nos propter nomen tuum.

\variantheading{Preces dominicales et feriales (si dicendæ sint)}

KYRIE, eléison. Christe, eléison. Kýrie, eléison.

Pater noster \rubric{secreto usque ad} \V\ Et ne nos indúcas in tentatiónem. \R\ Sed líbera nos a malo.

Credo in Deum \rubric{secreto usque ad} \V\ Carnis resurrectiónem. \R\ Vitam ætérnam. Amen. \V\ Et ego ad te, Dómine, clamávi. \R\ Et mane orátio mea prævéniet te. \V\ Repleátur os meum laude. \R\ Ut cantem glóriam tuam, tota die magnitúdinem tuam. \V\ Dómine, avérte fáciem tuam a peccátis meis. \R\ Et omnes iniquitátes meas dele. \V\ Cor mundum crea in me, Deus. \R\ Et spíritum rectum ínnova in viscéribus meis. \V\ Ne projícias me a fácie tua. \R\ Et spíritum sanctum tuum ne áuferas a me. \V\ Redde mihi lætítiam salutáris tui. \R\ Et spíritu principáli confírma me.

\variantheading{In Précibus feriálibus tantum}

\V\ Eripe me, Dómine, ab hómine malo. \R\ A viro iníquo éripe me. \V\ Eripe me de inimícis meis, Deus meus. \R\ Et ab insurgéntibus in me líbera me. \V\ Eripe me de operántibus iniquitátem. \R\ Et de viris sánguinum salva me. \V\ Sic psalmum dicam nómini tuo in sæculum sæculi. \R\ Ut reddam vota mea de die in diem. \V\ Exáudi nos, Deus, salutáris noster. \R\ Spes ómnium fínium terræ, et in mari longe. \V\ Deus, in adjutórium meum inténde. \R\ Dómine, ad adjuvándum me festína. \V\ Sanctus Deus, Sanctus fortis, Sanctus immortális. \R\ Miserére nobis. \V\ Bénedic, ánima mea, Dómino. \R\ Et ómnia, quæ intra me sunt, nómini sancto ejus. \V\ Bénedic, ánima mea, Dómino. \R\ Et noli oblivísci omnes retributiónes ejus. \V\ Qui propitiátur ómnibus iniquitátibus tuis. \R\ Qui sanat omnes infirmitátes tuas. \V\ Qui rédimit de intéritu vitam tuam. \R\ Qui corónat te in misericórdia et miseratiónibus. \V\ Qui replet in bonis desidérium tuum. \R\ Renovábitur ut áquilæ juvéntus tua.

\variantheading{In utrisque Précibus}

\V\ Adjutórium nostrum in nómine Dómini. \R\ Qui fecit cælum et terram.

\formulalabel{Confíteor. Misereátur. Indulgéntiam.}

\V\ Dignáre, Dómine, die isto. \R\ Sine peccáto nos custodíre. \V\ Miserére nostri, Dómine. \R\ Miserére nostri. \V\ Fiat misericórdia tua, Dómine, super nos. \R\ Quemádmodum sperávimus in te. \V\ Dómine, exáudi oratiónem meam. \R\ Et clamor meus ad te véniat.

\V\ Dóminus vobíscum.

Orémus. \oratiolabel{Oratio}

\lettrine[lines=2]{D}{OMINE}, Deus omnípotens, qui ad princípium hujus diéi nos perveníre fecísti: tua nos hódie salva virtúte; ut in hac die ad nullum declinémus peccátum, sed semper ad tuam justítiam faciéndam nostra procédant elóquia, dirigántur cogitatiónes et ópera. Per Dóminum. \R\ Amen.

\V\ Dóminus vobíscum. \V\ Benedicámus Dómino.

\rubric{In Choro Martyrológium.}

\V\ Pretiósa in conspéctu Dómini. \R\ Mors Sanctórum ejus.

\rubric{Et absolúte, sine} Orémus:

\oratiolabel{Oratio}

\lettrine[lines=2]{S}{ANCTA} María, et omnes Sancti intercédant pro nobis ad Dóminum, ut nos mereámur ab eo adjuvári et salvári, qui vivit et regnat in sæcula sæculórum. \R\ Amen.

\V\ Deus, in adjutórium meum inténde. \R\ Dómine, ad adjuvándum me festína. \rubric{Et dícitur ter; último additur:} Glória Patri. Sicut erat.

Kýrie, eléison. Christe, eléison. Kýrie, eléison.

Pater noster \rubric{secreto usque ad} \V\ Et ne nos indúcas in tentatiónem. \R\ Sed líbera nos a malo. \V\ Réspice in servos tuos, Dómine, et in ópera tua, et dírige fílios eórum. \R\ Et sit splendor Dómini, Dei nostri, super nos, et ópera mánuum nostrárum dírige super nos, et opus mánuum nostrárum dírige. \V\ Glória Patri. \R\ Sicut erat.

Orémus. \oratiolabel{Oratio}

\lettrine[lines=2]{D}{IRIGERE} et sanctificáre, régere et gubernáre dignáre, Dómine Deus, Rex cæli et terræ, hódie corda et córpora nostra, sensus, sermónes et actus nostros in lege tua, et in opéribus mandatórum tuórum: ut hic, et in ætérnum, te auxiliánte, salvi et líberi esse mereámur, Salvátor mundi: Qui vivis et regnas in sæcula sæculórum. \R\ Amen.

\V\ Jube, domne.

\benedlabel{Bened.} Dies et actus nostros in sua pace dispónat Dóminus omnípotens. \R\ Amen.

\variantheading{In Officio feriali}

\lettrine[lines=2]{D}{OMINUS} autem dírigat corda et córpora nostra in caritáte Dei, et patiéntia Christi. Tu autem.

\variantheading{E Proprio Sanctorum}

\variantheading{Pro Oct. Omn. Sanctorum}

\lettrine[lines=2]{B}{ENEDÍCTIO} et cláritas et sapiéntia et gratiárum áctio, honor, virtus et fortitúdo Deo nostro in sǽcula sæculórum. Amen.

\variantheading{E Communi Sanctorum}

\variantheading{Pro Apostolis}

\lettrine[lines=2]{I}{BANT} Apóstoli gaudéntes a conspéctu concílii, quóniam digni hábiti sunt pro nómine Jesu contuméliam pati.

\variantheading{Pro Martyre}

\lettrine[lines=2]{J}{USTUS} cor suum trádidit ad vigilándum dilúculo ad Dóminum, qui fecit illum, et in conspéctu Altíssimi deprecábitur.

\variantheading{Pro pluribus Martyribus}

\lettrine[lines=2]{F}{ULGEBUNT} justi, et tamquam scintíllæ in arundinéto discúrrent. Judicábunt natiónes, et dominabúntur pópulis: et regnábit Dóminus illórum in perpétuum.

\variantheading{Pro Conf. Pont.}

\lettrine[lines=2]{F}{UNGI} sacerdótio, et habére laudem in nómine ipsíus, et offérre illi incénsum dignum in odórem suavitátis.

\variantheading{Pro Conf. non Pont.}

\lettrine[lines=2]{J}{USTUM} dedúxit Dóminus per vias rectas, et osténdit illi regnum Dei, et dedit illi sciéntiam sanctórum; honestávit illum in labóribus, et complévit labóres illíus.

\variantheading{Pro Virgine Martyre}

\lettrine[lines=2]{D}{OMINE}, Deus meus, exaltásti super terram habitatiónem meam, et pro morte defluénte deprecáta sum. Invocávi Dóminum, Patrem Dómini mei, ut non derelínquat me in die tribulatiónis meæ, et in témpore superbórum sine adjutório.

\variantheading{Pro Virgine non Martyre}

\lettrine[lines=2]{O}{} \textsc{quam} pulchra est casta generátio cum claritáte! immortális est enim memória illíus: quóniam apud Deum nota est, et apud hómines.

\variantheading{Pro Martyre tantum}

\lettrine[lines=2]{L}{AUDABIT} usque ad mortem ánima mea Dóminum, quóniam éruis sustinéntes te, et líberas eos de manu angústiæ, Dómine, Deus noster.

\variantheading{Pro nec Virg. nec Mart.}

\lettrine[lines=2]{M}{ULTAE} fíliæ congregavérunt divítias: tu supergréssa es univérsas. Fallax grátia, et vana est pulchritúdo: múlier timens Dóminum, ipsa laudábitur.

\variantheading{Pro Dedicatione Ecclesiæ}

\lettrine[lines=2]{E}{T} abstérget Deus omnem lácrimam ab óculis eórum: et mors ultra non erit, neque luctus, neque clamor, neque dolor erit ultra, quia prima abiérunt. Et dixit qui sedébat in throno: Ecce, nova fácio ómnia.

\variantheading{In Officiis B. Mariæ V.}

\lettrine[lines=2]{I}{N} platéis sicut cinnamómum et bálsamum aromatízans odórem dedi: quasi myrrha elécta dedi suavitátem odóris.

\V\ Adjutórium nostrum in nómine Dómini. \R\ Qui fecit cælum et terram.

\V\ Benedícite. \R\ Deus.

\benedictiolabel{Benedictio} Dóminus nos benedícat, et ab omni malo deféndat, et ad vitam perdúcat ætérnam. Et fidélium ánimæ per misericórdiam Dei requiéscant in pace. \R\ Amen.

\rubric{Deinde dícitur tantum} Pater noster \rubric{secréto, nisi sequátur ália Hora.}

% ===== END OF BATCH 27 -- Covers autumnalis.md lines 8967-9226,
% completing "Ad Primam" for Feria Quinta in full: the hymn "Jam
% lucis orto sídere," Psalms 22, 71/i, 71/ii, the conditional Psalm
% 97 insertion, both Capitula, the brief Responsory, the full
% Preces apparatus, the closing Orationes, and the Lectio-brevis
% apparatus (identical category list to the previous three ferias).
% Several more psalm incipits carry the partial-caps OCR defect
% ("DOminus," "DEus," "COram," "CAntáte," "PAcem" [Capitulum]) and
% are silently normalized to full caps. The document continues
% with "Ad Tertiam" (autumnalis.md line 9229).
% =====
\hourheading{Ad Tertiam}

Pater noster et Ave.

\V\ Deus, in adjutórium.

\hymnusheading{Hymnus}

\lettrine[lines=2]{N}{UNC}, Sancte, nobis, Spíritus,\\
Unum Patri cum Fílio,\\
Dignáre promptus íngeri\\
Nostro refúsus péctori.

Os, lingua, mens, sensus, vigor\\
Confessiónem pérsonent,\\
Flamméscat igne cáritas,\\
Accéndat ardor próximos.

Præsta, Pater piíssime,\\
Patríque compar Unice,\\
Cum Spíritu Paráclito\\
Regnans per omne sæculum.\\
Amen.

\antlabel{Ant.} Quam bonus. \textdagger

\psalmheading{Psalmus 72, i}

\lettrine[lines=2]{Q}{UAM} bonus \textdagger\ Israël Deus, * his qui recto sunt corde!

Mei autem pæne moti sunt pedes: * pæne effúsi sunt gressus mei.

Quia zelávi super iníquos, * pacem peccatórum videns.

Quia non est respéctus morti eórum: * et firmaméntum in plaga eórum.

In labóre hóminum non sunt, * et cum homínibus non flagellabúntur:

Ideo ténuit eos supérbia, * opérti sunt iniquitáte et impietáte sua.

Pródiit quasi ex ádipe iníquitas eórum: * transiérunt in afféctum cordis.

Cogitavérunt, et locúti sunt nequítiam: * iniquitátem in excélso locúti sunt.

Posuérunt in cælum os suum: * et lingua eórum transívit in terra.

\psalmheading{Psalmus 72, ii}

\lettrine[lines=2]{I}{DEO} convertétur pópulus meus hic: * et dies pleni inveniéntur in eis.

Et dixérunt: Quómodo scit Deus, * et si est sciéntia in excélso?

Ecce, ipsi peccatóres, et abundántes in sæculo, * obtinuérunt divítias.

Et dixi: Ergo sine causa justificávi cor meum, * et lavi inter innocéntes manus meas:

Et fui flagellátus tota die, * et castigátio mea in matutínis.

Si dicébam: Narrábo sic: * ecce, natiónem filiórum tuórum reprobávi.

Existimábam ut cognóscerem hoc, * labor est ante me:

Donec intrem in Sanctuárium Dei: * et intélligam in novíssimis eórum.

\psalmheading{Psalmus 72, iii}

\lettrine[lines=2]{V}{ERUMTAMEN} propter dolos posuísti eis: * dejecísti eos dum allevaréntur.

Quómodo facti sunt in desolatiónem, súbito defecérunt: * periérunt propter iniquitátem suam.

Velut sómnium surgéntium, Dómine, * in civitáte tua imáginem ipsórum ad níhilum rédiges.

Quia inflammátum est cor meum, et renes mei commutáti sunt: * et ego ad níhilum redáctus sum, et nescívi.

Ut juméntum factus sum apud te: * et ego semper tecum.

Tenuísti manum déxteram meam: et in voluntáte tua deduxísti me, * et cum glória suscepísti me.

Quid enim mihi est in cælo? * et a te quid vólui super terram?

Defécit caro mea, et cor meum: * Deus cordis mei, et pars mea Deus in ætérnum.

Quia ecce, qui elóngant se a te, períbunt: * perdidísti omnes, qui fornicántur abs te.

Mihi autem adhærére Deo bonum est: * pónere in Dómino Deo spem meam:

Ut annúntiem omnes prædicatiónes tuas, * in portis fíliæ Sion.

\antlabel{Ant.} Quam bonus Israël Deus, his, qui recto sunt corde?

\variantheading{In Officio feriali}

\capitulumlabel{Capitulum}{Jer. 17, 14}

\lettrine[lines=2]{S}{ANA} me, Dómine, et sanábor: salvum me fac, et salvus ero: quóniam laus mea tu es.

\R\ br. Sana ánimam meam, * Quia peccávi tibi. Sana. \V\ Ego dixi: Dómine, miserére mei. Quia. Glória Patri. Sana.

\V\ Adjútor meus esto, ne derelínquas me. \R\ Neque despícias me, Deus, salutáris meus.

\variantheading{In Vigiliis communibus:}

\variantheading{Preces feriales}

KYRIE, eléison. Christe, eléison. Kýrie, eléison.

Pater noster \rubric{secreto usque ad} \V\ Et ne nos indúcas in tentatiónem. \R\ Sed líbera nos a malo. \V\ Dómine, Deus virtútum, convérte nos. \R\ Et osténde fáciem tuam, et salvi érimus. \V\ Exsúrge, Christe, ádjuva nos. \R\ Et líbera nos propter nomen tuum. \V\ Dómine, exáudi oratiónem meam. \R\ Et clamor meus ad te véniat.

\rubric{Réliqua ut in Ordinário. 19}

\hourheading{Ad Sextam}

Pater noster et Ave.

\V\ Deus, in adjutórium.

\hymnusheading{Hymnus}

\lettrine[lines=2]{R}{ECTOR} potens, verax Deus,\\
Qui témperas rerum vices,\\
Splendóre mane illúminas,\\
Et ígnibus merídiem:

Exstíngue flammas lítium,\\
Aufer calórem nóxium,\\
Confer salútem córporum,\\
Verámque pacem córdium.

Præsta, Pater piíssime,\\
Patríque compar Unice,\\
Cum Spíritu Paráclito\\
Regnans per omne sæculum.\\
Amen.

\antlabel{Ant.} Memor esto.

\psalmheading{Psalmus 73, i}

\lettrine[lines=2]{U}{T} quid, Deus, repulísti in finem: * irátus est furor tuus super oves páscuæ tuæ?

Memor esto congregatiónis tuæ, * quam possedísti ab inítio.

Redemísti virgam hereditátis tuæ: * mons Sion, in quo habitásti in eo.

Leva manus tuas in supérbias eórum in finem: * quanta malignátus est inimícus in sancto!

Et gloriáti sunt qui odérunt te: * in médio solemnitátis tuæ.

Posuérunt signa sua, signa: * et non cognovérunt sicut in éxitu super summum.

Quasi in silva lignórum secúribus excidérunt jánuas ejus in idípsum: * in secúri et áscia dejecérunt eam.

Incendérunt igni Sanctuárium tuum: * in terra polluérunt tabernáculum nóminis tui.

Dixérunt in corde suo cognátio eórum simul: * Quiéscere faciámus omnes dies festos Dei a terra.

Signa nostra non vídimus, jam non est prophéta: * et nos non cognóscet ámplius.

\psalmheading{Psalmus 73, ii}

\lettrine[lines=2]{U}{SQUEQUO}, Deus, improperábit inimícus: * irrítat adversárius nomen tuum in finem?

Ut quid avértis manum tuam, et déxteram tuam, * de médio sinu tuo in finem?

Deus autem Rex noster ante sæcula: * operátus est salútem in médio terræ.

Tu confirmásti in virtúte tua mare: * contribulásti cápita dracónum in aquis.

Tu confregísti cápita dracónis: * dedísti eum escam pópulis Æthíopum.

Tu dirupísti fontes, et torréntes: * tu siccásti flúvios Ethan.

Tuus est dies, et tua est nox: * tu fabricátus es auróram et solem.

Tu fecísti omnes términos terræ: * æstátem et ver tu plasmásti ea.

\psalmheading{Psalmus 73, iii}

\lettrine[lines=2]{M}{EMOR} esto hujus, inimícus improperávit Dómino: * et pópulus insípiens incitávit nomen tuum.

Ne tradas béstiis ánimas confiténtes tibi, * et ánimas páuperum tuórum ne oblivíscáris in finem.

Réspice in testaméntum tuum: * quia repléti sunt, qui obscuráti sunt terræ dómibus iniquitátum.

Ne avertátur húmilis factus confúsus: * pauper et inops laudábunt nomen tuum.

Exsúrge, Deus, júdica causam tuam: * memor esto improperiórum tuórum, eórum quæ ab insipiénte sunt tota die.

Ne obliviscáris voces inimicórum tuórum: * supérbia eórum, qui te odérunt, ascéndit semper.

\antlabel{Ant.} Memor esto congregatiónis tuæ, Dómine, quam possedísti ab inítio.

\variantheading{In Officio feriali}

\capitulumlabel{Capitulum}{Rom. 13, 8}

\lettrine[lines=2]{N}{EMINI} quidquam debeátis, nisi ut ínvicem diligátis: qui enim díligit próximum, legem implévit.

\R\ br. Benedícam Dóminum * In omni témpore. Benedícam. \V\ Semper laus ejus in ore meo. In. Glória Patri. Benedícam. \V\ Dóminus regit me, et nihil mihi déerit. \R\ In loco páscuæ ibi me collocávit.

\variantheading{In Vigiliis communibus:}

\variantheading{Preces feriales}

KYRIE, eléison. Christe, eléison. Kýrie, eléison.

Pater noster \rubric{secreto usque ad} \V\ Et ne nos indúcas in tentatiónem. \R\ Sed líbera nos a malo. \V\ Dómine, Deus virtútum, convérte nos. \R\ Et osténde fáciem tuam, et salvi érimus. \V\ Exsúrge, Christe, ádjuva nos. \R\ Et líbera nos propter nomen tuum. \V\ Dómine, exáudi oratiónem meam. \R\ Et clamor meus ad te véniat.

\rubric{Réliqua ut in Ordinário. 20}

\hourheading{Ad Nonam}

Pater noster et Ave.

\V\ Deus, in adjutórium.

\hymnusheading{Hymnus}

\lettrine[lines=2]{R}{ERUM}, Deus, tenax vigor,\\
Immótus in te pérmanens,\\
Lucis diúrnæ témpora\\
Succéssibus detérminans:

Largíre lumen véspere,\\
Quo vita nusquam décidat,\\
Sed prǽmium mortis sacræ\\
Perénnis instet glória.

Præsta, Pater piíssime,\\
Patríque compar Unice,\\
Cum Spíritu Paráclito\\
Regnans per omne sæculum.\\
Amen.

\antlabel{Ant.} Invocábimus.

\psalmheading{Psalmus 74}

\lettrine[lines=2]{C}{ONFITEBIMUR} tibi, Deus: * confitébimur, et invocábimus nomen tuum.

Narrábimus mirabília tua: * cum accépero tempus, ego justítias judicábo.

Liquefácta est terra, et omnes qui hábitant in ea: * ego confirmávi colúmnas ejus.

Dixi iníquis: Nolíte iníque ágere: * et delinquéntibus: Nolíte exaltáre cornu:

Nolíte extóllere in altum cornu vestrum: * nolíte loqui advérsus Deum iniquitátem.

Quia neque ab Oriénte, neque ab Occidénte, neque a desértis móntibus: * quóniam Deus judex est.

Hunc humíliat, et hunc exáltat: * quia calix in manu Dómini vini meri plenus misto.

Et inclinávit ex hoc in hoc: verúmtamen fæx ejus non est exinaníta: * bibent omnes peccatóres terræ.

Ego autem annuntiábo in sæculum: * cantábo Deo Jacob.

Et ómnia córnua peccatórum confríngam: * et exaltabúntur córnua justi.

\psalmheading{Psalmus 75, i}

\lettrine[lines=2]{N}{OTUS} in Judǽa Deus: * in Israël magnum nomen ejus.

Et factus est in pace locus ejus: * et habitátio ejus in Sion.

Ibi confrégit poténtias árcuum, * scutum, gládium, et bellum.

Illúminans tu mirabíliter a móntibus ætérnis: * turbáti sunt omnes insipiéntes corde.

Dormiérunt somnum suum: * et nihil invenérunt omnes viri divitiárum in mánibus suis.

Ab increpatióne tua, Deus Jacob, * dormitavérunt qui ascendérunt equos.

\psalmheading{Psalmus 75, ii}

\lettrine[lines=2]{T}{U} terríbilis es, et quis resístet tibi? * ex tunc ira tua.

De cælo audítum fecísti judícium: * terra trémuit et quiévit.

Cum exsúrgeret in judícium Deus, * ut salvos fáceret omnes mansuétos terræ.

Quóniam cogitátio hóminis confitébitur tibi: * et relíquiæ cogitatiónis diem festum agent tibi.

Vovéte, et réddite Dómino, Deo vestro: * omnes, qui in circúitu ejus affértis múnera.

Terríbili et ei qui aufert spíritum príncipum, * terríbili apud reges terræ.

\antlabel{Ant.} Invocábimus nomen tuum, Dómine: narrábimus mirabília tua.

\variantheading{In Officio feriali}

\capitulumlabel{Capit.}{1 Petri 1, 17-19}

\lettrine[lines=2]{I}{N} timóre incolátus vestri témpore conversámini: sciéntes quod non corruptibílibus auro vel argénto redémpti estis, sed pretióso sánguine quasi Agni immaculáti Christi.

\R\ br. Rédime me, Dómine, * Et miserére mei. Rédime. \V\ Pes enim meus stetit in via recta. Et miserére mei. Glória Patri. Rédime. \V\ Ab occúltis meis munda me, Dómine. \R\ Et ab aliénis parce servo tuo.

\variantheading{In Vigiliis communibus:}

\variantheading{Preces feriales}

KYRIE, eléison. Christe, eléison. Kýrie, eléison.

Pater noster \rubric{secreto usque ad} \V\ Et ne nos indúcas in tentatiónem. \R\ Sed líbera nos a malo. \V\ Dómine, Deus virtútum, convérte nos. \R\ Et osténde fáciem tuam, et salvi érimus. \V\ Exsúrge, Christe, ádjuva nos. \R\ Et líbera nos propter nomen tuum. \V\ Dómine, exáudi oratiónem meam. \R\ Et clamor meus ad te véniat.

\rubric{Réliqua ut in Ordinário. 22}

% ===== END OF BATCH 28 -- Covers autumnalis.md lines 9229-9538,
% completing the Little Hours of Feria Quinta in full: "Ad
% Tertiam" (Psalms 72/i-iii), "Ad Sextam" (Psalms 73/i-iii), and
% "Ad Nonam" (Psalms 74, 75/i, 75/ii), each with its own hymn,
% Capitulum, brief Responsory, and "In Vigiliis communibus: Preces
% feriales" apparatus. Three further occurrences of the recurring
% stray OCR "\*" artifact in the hour-hymns are silently dropped.
% "Ut quid" (Psalm 73, i) correctly shows only its first short word
% in full caps ("UT"), consistent with the convention that only the
% single opening word takes full capitalization -- not a defect,
% unlike the partial-caps cases seen elsewhere. Several further
% genuine partial-caps OCR defects ("IDeo," "VErúmtamen," "USquequo,"
% "MEmor," "COnfitebimur," "NOtus") are silently normalized to full
% caps, consistent with prior batches. The document continues with
% "Ad Vesperas" (autumnalis.md line 9541).
% =====
\hourheading{Ad Vesperas}

Pater noster et Ave.

\V\ Deus, in adjutórium.

\antlabel{Ant.} Ecce, quam bonum * \textdagger\ et quam jucúndum, habitáre fratres in unum. \textdaggerdbl

\psalmheading{Psalmus 132}

\lettrine[lines=2]{E}{CCE}, quam bonum \textdagger\ et quam jucúndum, * habitáre fratres in unum: \textdaggerdbl

Sicut unguéntum in cápite, * quod descéndit in barbam, barbam Aaron.

Quod descéndit in oram vestiménti ejus: * sicut ros Hermon, qui descéndit in montem Sion.

Quóniam illic mandávit Dóminus benedictiónem, * et vitam usque in sæculum.

\antlabel{Ant.} Ecce, quam bonum et quam jucúndum, habitáre fratres in unum!

\antlabel{Ant.} Confitémini Dómino * \textdagger\ quóniam in ætérnum misericórdia ejus.

\psalmheading{Psalmus 135, i}

\lettrine[lines=2]{C}{ONFITEMINI} Dómino \textdagger\ quóniam bonus: * quóniam in ætérnum misericórdia ejus.

Confitémini Deo deórum: * quóniam in ætérnum misericórdia ejus.

Confitémini Dómino dominórum: * quóniam in ætérnum misericórdia ejus.

Qui facit mirabília magna solus: * quóniam in ætérnum misericórdia ejus.

Qui fecit cælos in intelléctu: * quóniam in ætérnum misericórdia ejus.

Qui firmávit terram super aquas: * quóniam in ætérnum misericórdia ejus.

Qui fecit luminária magna: * quóniam in ætérnum misericórdia ejus.

Solem in potestátem diéi: * quóniam in ætérnum misericórdia ejus.

Lunam, et stellas in potestátem noctis: * quóniam in ætérnum misericórdia ejus.

\antlabel{Ant.} Confitémini Dómino quóniam in ætérnum misericórdia ejus.

\antlabel{Ant.} Confitémini Dómino, * quia in humilitáte nostra memor fuit nostri.

\psalmheading{Psalmus 135, ii}

\lettrine[lines=2]{Q}{UI} percússit Ægýptum cum primogénitis eórum: * quóniam in ætérnum misericórdia ejus.

Qui edúxit Israël de médio eórum: * quóniam in ætérnum misericórdia ejus.

In manu poténti, et bráchio excélso: * quóniam in ætérnum misericórdia ejus.

Qui divísit Mare Rubrum in divisiónes: * quóniam in ætérnum misericórdia ejus.

Et edúxit Israël per médium ejus: * quóniam in ætérnum misericórdia ejus.

Et excússit Pharaónem, et virtútem ejus in Mari Rubro: * quóniam in ætérnum misericórdia ejus.

Qui tradúxit pópulum suum per desértum: * quóniam in ætérnum misericórdia ejus.

Qui percússit reges magnos: * quóniam in ætérnum misericórdia ejus.

Et occídit reges fortes: * quóniam in ætérnum misericórdia ejus.

Sehon, regem Amorrhæórum: * quóniam in ætérnum misericórdia ejus.

Et Og, regem Basan: * quóniam in ætérnum misericórdia ejus.

Et dedit terram eórum hereditátem: * quóniam in ætérnum misericórdia ejus.

Hereditátem Israël, servo suo: * quóniam in ætérnum misericórdia ejus.

Quia in humilitáte nostra memor fuit nostri: * quóniam in ætérnum misericórdia ejus.

Et redémit nos ab inimícis nostris: * quóniam in ætérnum misericórdia ejus.

Qui dat escam omni carni: * quóniam in ætérnum misericórdia ejus.

Confitémini Deo cæli: * quóniam in ætérnum misericórdia ejus.

Confitémini Dómino dominórum: * quóniam in ætérnum misericórdia ejus.

\antlabel{Ant.} Confitémini Dómino, quia in humilitáte nostra memor fuit nostri.

\antlabel{Ant.} Adhǽreat lingua mea * fáucibus meis, si non memínero tui, Jerúsalem.

\psalmheading{Psalmus 136}

\lettrine[lines=2]{S}{UPER} flúmina Babylónis, illic sédimus et flévimus: * cum recordarémur Sion:

In salícibus in médio ejus, * suspéndimus órgana nostra.

Quia illic interrogavérunt nos, qui captívos duxérunt nos, * verba cantiónum:

Et qui abduxérunt nos: * Hymnum cantáte nobis de cánticis Sion.

Quómodo cantábimus cánticum Dómini: * in terra aliéna?

Si oblítus fúero tui, Jerúsalem, * oblivióni detur déxtera mea.

Adhǽreat lingua mea fáucibus meis, * si non memínero tui:

Si non proposúero Jerúsalem, * in princípio lætítiæ meæ.

Memor esto, Dómine, filiórum Edom, * in die Jerúsalem:

Qui dicunt: Exinaníte, exinaníte: * usque ad fundaméntum in ea.

Fília Babylónis mísera: * beátus, qui retríbuet tibi retributiónem tuam, quam retribuísti nobis.

Beátus, qui tenébit, * et allídet párvulos tuos ad petram.

\antlabel{Ant.} Adhǽreat lingua mea fáucibus meis, si non memínero tui, Jerúsalem.

\antlabel{Ant.} Confitébor * \textdagger\ nómini tuo, Dómine, super misericórdia et veritáte tua.

\psalmheading{Psalmus 137}

\lettrine[lines=2]{C}{ONFITEBOR} \textdagger\ tibi, Dómine, in toto corde meo: * quóniam audísti verba oris mei.

In conspéctu Angelórum psallam tibi: * adorábo ad templum sanctum tuum, et confitébor nómini tuo.

Super misericórdia tua, et veritáte tua: * quóniam magnificásti super omne, nomen sanctum tuum.

In quacúmque die invocávero te, exáudi me: * multiplicábis in ánima mea virtútem.

Confiteántur tibi, Dómine, omnes reges terræ: * quia audiérunt ómnia verba oris tui:

Et cantent in viis Dómini: * quóniam magna est glória Dómini.

Quóniam excélsus Dóminus, et humília réspicit: * et alta a longe cognóscit.

Si ambulávero in médio tribulatiónis, vivificábis me: * et super iram inimicórum meórum extendísti manum tuam, et salvum me fecit déxtera tua.

Dóminus retríbuet pro me: * Dómine, misericórdia tua in sæculum: ópera mánuum tuárum ne despícias.

\antlabel{Ant.} Confitébor nómini tuo, Dómine, super misericórdia et veritáte tua.

\variantheading{In Officio feriali}

\capitulumlabel{Capitulum}{2 Cor. 1, 3-4}

\lettrine[lines=2]{B}{ENEDICTUS} Deus, et Pater Dómini nostri Jesu Christi, Pater misericordiárum, et Deus totíus consolatiónis, qui consolátur nos in omni tribulatióne nostra.

\hymnusheading{Hymnus}

\lettrine[lines=2]{M}{AGNAE} Deus poténtiæ,\\
Qui fértili natos aqua\\
Partim relínquis gúrgiti,\\
Partim levas in áëra.

Demérsa lymphis ímprimens,\\
Subvécta cælis érigens:\\
Ut stirpe ab una pródita,\\
Divérsa répleant loca:

Largíre cunctis sérvulis,\\
Quos mundat unda Sánguinis,\\
Nescíre lapsus críminum,\\
Nec ferre mortis tædium.

Ut culpa nullum déprimat:\\
Nullum éfferat jactántia:\\
Elísa mens ne cóncidat:\\
Eláta mens ne córruat.

Præsta, Pater piíssime,\\
Patríque compar Unice,\\
Cum Spíritu Paráclito\\
Regnans per omne sæculum.\\
Amen.

\V\ Dirigátur, Dómine, orátio mea. \R\ Sicut incénsum in conspéctu tuo.

\antlabel{Ad Magnif. Ant.} Fecit Deus * poténtiam in bráchio suo: dispérsit supérbos mente cordis sui.

\rubric{Réliqua ut in Ordinário. 22}

\hourheading{Ad Completorium}

\V\ Jube, domne.

\benedlabel{Bened.} Noctem quiétam, et finem perféctum concédat nobis Dóminus omnípotens. \R\ Amen.

\lectioheading{Lectio br.}{1 Petri 5, 8-9}

\lettrine[lines=2]{F}{RATRES}: Sóbrii estóte, et vigiláte: quia adversárius vester diábolus tamquam leo rúgiens círcuit, quærens quem dévoret: cui resístite fortes in fide. Tu autem, Dómine, miserére nobis. \R\ Deo grátias.

\V\ Adjutórium nostrum in nómine Dómini. \R\ Qui fecit cælum et terram.

Pater noster \rubric{quod dícitur totum secréto.}

\formulalabel{Confíteor.} Misereátur. Indulgéntiam.

\V\ Convérte nos, Deus, salutáris noster. \R\ Et avérte iram tuam a nobis.

\V\ Deus, in adjutórium.

\antlabel{Ant.} Adjútor meus.

\psalmheading{Psalmus 69}

\lettrine[lines=2]{D}{EUS}, in adjutórium meum inténde: * Dómine, ad adjuvándum me festína.

Confundántur et reverántur, * qui quærunt ánimam meam.

Avertántur retrórsum, et erubéscant, * qui volunt mihi mala.

Avertántur statim erubescéntes, * qui dicunt mihi: Euge, euge.

Exsúltent et læténtur in te omnes qui quærunt te, * et dicant semper: Magnificétur Dóminus: qui díligunt salutáre tuum.

Ego vero egénus, et pauper sum: * Deus, ádjuva me.

Adjútor meus, et liberátor meus es tu: * Dómine, ne moréris.

\psalmheading{Psalmus 70, i}

\lettrine[lines=2]{I}{N} te, Dómine, sperávi, non confúndar in ætérnum: * in justítia tua líbera me, et éripe me.

Inclína ad me aurem tuam, * et salva me.

Esto mihi in Deum protectórem, et in locum munítum: * ut salvum me fácias.

Quóniam firmaméntum meum, * et refúgium meum es tu.

Deus meus, éripe me de manu peccatóris, * et de manu contra legem agéntis et iníqui:

Quóniam tu es patiéntia mea, Dómine: * Dómine, spes mea a juventúte mea.

In te confirmátus sum ex útero: * de ventre matris meæ tu es protéctor meus.

In te cantátio mea semper: * tamquam prodígium factus sum multis: et tu adjútor fortis.

Repleátur os meum laude, ut cantem glóriam tuam: * tota die magnitúdinem tuam.

Ne projícias me in témpore senectútis: * cum defécerit virtus mea, ne derelínquas me.

Quia dixérunt inimíci mei mihi: * et qui custodiébant ánimam meam, consílium fecérunt in unum,

Dicéntes: Deus dereliquit eum, persequímini, et comprehéndite eum: * quia non est qui erípiat.

Deus, ne elongéris a me: * Deus meus, in auxílium meum réspice.

\psalmheading{Psalmus 70, ii}

\lettrine[lines=2]{C}{ONFUNDANTUR}, et defíciant detrahéntes ánimæ meæ: * operiántur confusióne, et pudóre qui quærunt mala mihi.

Ego autem semper sperábo: * et adjíciam super omnem laudem tuam.

Os meum annuntiábit justítiam tuam: * tota die salutáre tuum.

Quóniam non cognóvi litteratúram, introíbo in poténtias Dómini: * Dómine, memorábor justítiæ tuæ solíus.

Deus, docuísti me a juventúte mea: * et usque nunc pronuntiábo mirabília tua.

Et usque in senéctam et sénium: * Deus, ne derelínquas me.

Donec annúntiem bráchium tuum * generatióni omni, quæ ventúra est:

Poténtiam tuam, et justítiam tuam, Deus, usque in altíssima, quæ fecísti magnália: * Deus, quis símilis tibi?

Quantas ostendísti mihi tribulatiónes multas et malas: et convérsus vivificásti me: * et de abýssis terræ íterum reduxísti me:

Multiplicásti magnificéntiam tuam: * et convérsus consolátus es me.

Nam et ego confitébor tibi in vasis psalmi veritátem tuam: * Deus, psallam tibi in cíthara, Sanctus Israël.

Exsultábunt lábia mea cum cantávero tibi: * et ánima mea, quam redemísti.

Sed et lingua mea tota die meditábitur justítiam tuam: * cum confúsi et reveríti fúerint, qui quærunt mala mihi.

\antlabel{Ant.} Adjútor meus, et liberátor meus esto, Dómine.

\hymnusheading{Hymnus}

\lettrine[lines=2]{T}{E} lucis ante términum,\\
Rerum Creátor, póscimus,\\
Ut pro tua cleméntia\\
Sis præsul et custódia.

Procul recédant sómnia,\\
Et nóctium phantásmata,\\
Hostémque nostrum cómprime,\\
Ne polluántur córpora.

Præsta, Pater piíssime,\\
Patríque compar Unice,\\
Cum Spíritu Paráclito\\
Regnans per omne sæculum.\\
Amen.

\capitulumlabel{Capitulum}{Jer. 14, 9}

\lettrine[lines=2]{T}{U} autem in nobis es, Dómine, et nomen sanctum tuum invocátum est super nos: ne derelínquas nos, Dómine, Deus noster.

\R\ br. In manus tuas, Dómine, * Comméndo spíritum meum. In manus. \V\ Redemísti nos, Dómine, Deus veritátis. Comméndo. Glória Patri. In manus.

\V\ Custódi nos, Dómine, ut pupíllam óculi. \R\ Sub umbra alárum tuárum prótege nos.

\antlabel{Ant.} Salva nos.

\canticumheading{Canticum Simeonis}{Luc. 2, 29-32}

\lettrine[lines=2]{N}{UNC} dimíttis servum tuum, Dómine, * secúndum verbum tuum in pace:

Quia vidérunt óculi mei * salutáre tuum,

Quod parásti * ante fáciem ómnium populórum,

Lumen ad revelatiónem géntium, * et glóriam plebis tuæ Israël.

Glória Patri.

\antlabel{Ant.} Salva nos, Dómine, vigilántes, custódi nos dormiéntes; ut vigilémus cum Christo, et requiescámus in pace.

\variantheading{Preces (si dicendæ sint)}

KYRIE, eléison. Christe, eléison. Kýrie, eléison.

Pater noster \rubric{secreto usque ad} \V\ Et ne nos indúcas in tentatiónem. \R\ Sed líbera nos a malo.

Credo in Deum \rubric{secreto usque ad} \V\ Carnis resurrectiónem. \R\ Vitam ætérnam. Amen. \V\ Benedíctus es, Dómine, Deus patrum nostrórum. \R\ Et laudábilis et gloriósus in sæcula. \V\ Benedicámus Patrem et Fílium cum Sancto Spíritu. \R\ Laudémus, et superexaltémus eum in sæcula. \V\ Benedíctus es, Dómine, in firmaménto cæli. \R\ Et laudábilis, et gloriósus, et superexaltátus in sæcula. \V\ Benedícat et custódiat nos omnípotens et miséricors Dóminus. \R\ Amen. \V\ Dignáre, Dómine, nocte ista. \R\ Sine peccáto nos custodíre. \V\ Miserére nostri, Dómine. \R\ Miserére nostri. \V\ Fiat misericórdia tua, Dómine, super nos. \R\ Quemádmodum sperávimus in te. \V\ Dómine, exáudi oratiónem meam. \R\ Et clamor meus ad te véniat.

\V\ Dóminus vobíscum.

Orémus. \oratiolabel{Oratio}

\lettrine[lines=2]{V}{ISITA}, quæsumus, Dómine, habitatiónem istam, et omnes insídias inimíci ab ea longe repélle: Angeli tui sancti hábitent in ea, qui nos in pace custódiant; et benedíctio tua sit super nos semper. Per Dóminum nostrum. \R\ Amen.

\V\ Dóminus vobíscum.

\V\ Benedicámus Dómino.

\benedictiolabel{Benedictio} Benedícat et custódiat nos omnípotens et miséricors Dóminus, Pater, et Fílius, et Spíritus Sanctus. \R\ Amen.

\variantheading{Antiphona finalis beátæ Maríæ Vírginis.}

\V\ Divínum auxílium.

Pater noster, Ave María et Credo.

% ===== END OF BATCH 29 -- Covers autumnalis.md lines 9541-9888,
% completing "Ad Vesperas" (Psalms 132, 135/i, 135/ii, 136, 137,
% hymn "Magnæ Deus poténtiæ") and "Ad Completorium" (Psalms 69,
% 70/i, 70/ii, hymn "Te lucis ante términum," Canticum Simeonis,
% Preces), and with them the ENTIRE "Feria Quinta" Office. The
% closing cross-reference at the end of Compline is given in its
% fuller unabbreviated form here -- "Antiphona finalis beátæ Maríæ
% Vírginis." and "Pater noster, Ave María et Credo." -- rather than
% the abbreviated "Antiphona finalis B. M. V." and "Pater, Ave et
% Credo." seen at the end of every previous feria's Compline this
% volume; transcribed exactly as the source gives each. The
% document continues with "FERIA SEXTA" (autumnalis.md line 9891),
% printed in the source in full capitals as a heading (transcribed
% with \feriaheading{} regardless, consistent with the role-based
% convention already applied to every other day-heading in this
% volume). The Invitatorium at Friday Matins carries a bracketed
% footnote marker "[3]" after its incipit in the source, preserved
% verbatim as part of the transcribed text.
% =====
\feriaheading{Feria Sexta}

\hourheading{Ad Matutinum}

Pater, Ave et Credo.

\V\ Dómine, lábia mea.

\V\ Deus, in adjutórium.

\variantheading{In Officio feriali}

\invitatlabel{Invit.} Dóminum, Deum nostrum, * Veníte, adorémus. \psalmlabel{Ps. 94} Veníte. [3]

\hymnusheading{Hymnus}

\lettrine[lines=2]{T}{U}, Trinitátis Únitas,\\
Orbem poténter quæ regis,\\
Atténde laudis cánticum,\\
Quod excubántes psállimus.

Nam léctulo consúrgimus\\
Noctis quiéto témpore,\\
Ut flagitémus ómnium\\
A te medélam vúlnerum.

Quo fraude quidquid dæmonum\\
In nóctibus delíquimus,\\
Abstérgat illud cælitus\\
Tuæ potéstas glóriæ.

Ne corpus astet sórdidum,\\
Nec torpor instet córdium,\\
Ne críminis contágio\\
Tepéscat ardor spíritus.

Ob hoc, Redémptor, quǽsumus,\\
Reple tuo nos lúmine,\\
Per quod diérum círculis\\
Nullis ruámus áctibus.

Præsta, Pater piíssime,\\
Patríque compar Unice,\\
Cum Spíritu Paráclito\\
Regnans per omne sæculum.\\
Amen.

\nocturnoheading{In I Nocturno}

\antlabel{Ant.} Suscitávit Dóminus * testimónium in Jacob: et legem pósuit in Israël.

\psalmheading{Psalmus 77, i}

\lettrine[lines=2]{A}{TTENDITE}, pópule meus, legem meam: * inclináte aurem vestram in verba oris mei.

Apériam in parábolis os meum: * loquar propositiónes ab inítio.

Quanta audívimus et cognóvimus ea: * et patres nostri narravérunt nobis.

Non sunt occultáta a fíliis eórum, * in generatióne áltera.

Narrántes laudes Dómini, et virtútes ejus, * et mirabília ejus, quæ fecit.

Et suscitávit testimónium in Jacob: * et legem pósuit in Israël.

Quanta mandávit pátribus nostris nota fácere ea fíliis suis: * ut cognóscat generátio áltera.

Fílii qui nascéntur, et exsúrgent, * et narrábunt fíliis suis,

Ut ponant in Deo spem suam, et non obliviscántur óperum Dei: * et mandáta ejus exquírant.

Ne fiant sicut patres eórum: * generátio prava et exásperans.

Generátio, quæ non diréxit cor suum: * et non est créditus cum Deo spíritus ejus.

\antlabel{Ant.} Suscitávit Dóminus testimónium in Jacob: et legem pósuit in Israël.

\antlabel{Ant.} Coram pátribus eórum * fecit Deus mirabília.

\psalmheading{Psalmus 77, ii}

\lettrine[lines=2]{F}{ILII} Ephrem intendéntes et mitténtes arcum: * convérsi sunt in die belli.

Non custodiérunt testaméntum Dei: * et in lege ejus noluérunt ambuláre.

Et oblíti sunt benefactórum ejus, * et mirabílium ejus quæ osténdit eis.

Coram pátribus eórum fecit mirabília in terra Ægýpti, * in campo Táneos.

Interrúpit mare, et perdúxit eos: * et státuit aquas quasi in utre.

Et dedúxit eos in nube diéi: * et tota nocte in illuminatióne ignis.

Interrúpit petram in erémo: * et adaquávit eos velut in abýsso multa.

Et edúxit aquam de petra: * et dedúxit tamquam flúmina aquas.

\antlabel{Ant.} Coram pátribus eórum fecit Deus mirabília.

\antlabel{Ant.} Jánuas cæli apéruit * Dóminus, et pluit illis manna ad manducándum.

\psalmheading{Psalmus 77, iii}

\lettrine[lines=2]{E}{T} apposuérunt adhuc peccáre ei: * in iram excitavérunt Excélsum in inaquóso.

Et tentavérunt Deum in córdibus suis, * ut péterent escas animábus suis.

Et male locúti sunt de Deo: * dixérunt: Numquid póterit Deus paráre mensam in desérto?

Quóniam percússit petram, et fluxérunt aquæ, * et torréntes inundavérunt.

Numquid et panem póterit dare, * aut paráre mensam pópulo suo?

Ideo audívit Dóminus, et dístulit: * et ignis accénsus est in Jacob, et ira ascéndit in Israël:

Quia non credidérunt in Deo, * nec speravérunt in salutári ejus:

Et mandávit núbibus désuper, * et jánuas cæli apéruit.

Et pluit illis manna ad manducándum, * et panem cæli dedit eis.

Panem Angelórum manducávit homo, * cibária misit eis in abundántia.

Tránstulit Austrum de cælo: * et indúxit in virtúte sua Africum.

Et pluit super eos sicut púlverem carnes: * et sicut arénam maris volatília pennáta.

Et cecidérunt in médio castrórum eórum: * circa tabernácula eórum.

Et manducavérunt, et saturáti sunt nimis, et desidérium eórum áttulit eis: * non sunt fraudáti a desidério suo.

Adhuc escæ eórum erant in ore ipsórum: * et ira Dei ascéndit super eos.

Et occídit pingues eórum, * et eléctos Israël impedívit.

\antlabel{Ant.} Jánuas cæli apéruit Dóminus, et pluit illis manna ad manducándum.

\variantheading{In Festis novem Lectionum:}

\V\ Deus percússit petram, et fluxérunt aquæ. \R\ Et torréntes inundavérunt.

\absolutiolabel{Absolutio.} Exáudi, Dómine.

\benedlabel{Bened.} 1 Benedictióne perpétua. 2 Unigénitus Dei Fílius. 3 Spíritus Sancti grátia.

\nocturnoheading{In II Nocturno}

\antlabel{Ant.} Deus adjútor * est eórum: et Excélsus redémptor eórum est.

\psalmheading{Psalmus 77, iv}

\lettrine[lines=2]{I}{N} ómnibus his peccavérunt adhuc: * et non credidérunt in mirabílibus ejus.

Et defecérunt in vanitáte dies eórum: * et anni eórum cum festinatióne.

Cum occíderet eos, quærébant eum: * et revertebántur, et dilúculo veniébant ad eum.

Et rememoráti sunt quia Deus adjútor est eórum: * et Deus excélsus redémptor eórum est.

Et dilexérunt eum in ore suo, * et lingua sua mentíti sunt ei:

Cor autem eórum non erat rectum cum eo: * nec fidéles habíti sunt in testaménto ejus.

Ipse autem est miséricors, et propítius fiet peccátis eórum: * et non dispérdet eos.

Et abundávit ut avérteret iram suam: * et non accéndit omnem iram suam:

Et recordátus est quia caro sunt: * spíritus vadens, et non rédiens.

Quóties exacerbavérunt eum in desérto, * in iram concitavérunt eum in inaquóso?

Et convérsi sunt, et tentavérunt Deum: * et Sanctum Israël exacerbavérunt.

\antlabel{Ant.} Deus adjútor est eórum: et Excélsus redémptor eórum est.

\antlabel{Ant.} Redémit eos * Dóminus de manu tribulántis.

\psalmheading{Psalmus 77, v}

\lettrine[lines=2]{N}{ON} sunt recordáti manus ejus, * die qua redémit eos de manu tribulántis.

Sicut pósuit in Ægýpto signa sua, * et prodígia sua in campo Táneos.

Et convértit in sánguinem flúmina eórum, * et imbres eórum, ne bíberent.

Misit in eos cœnomyíam, et comédit eos: * et ranam, et dispérdidit eos.

Et dedit ærúgini fructus eórum, * et labóres eórum locústæ.

Et occídit in grándine víneas eórum, * et moros eórum in prúina.

Et trádidit grándini juménta eórum: * et possessiónem eórum igni.

Misit in eos iram indignatiónis suæ: * indignatiónem, et iram, et tribulatiónem: immissiónes per ángelos malos.

Viam fecit sémitæ iræ suæ, non pepércit a morte animábus eórum: * et juménta eórum in morte conclúsit.

Et percússit omne primogénitum in terra Ægýpti: * primítias omnis labóris eórum in tabernáculis Cham.

Et ábstulit sicut oves pópulum suum: * et perdúxit eos tamquam gregem in desérto.

Et dedúxit eos in spe, et non timuérunt: * et inimícos eórum opéruit mare.

Et indúxit eos in montem sanctificatiónis suæ, * montem, quem acquisívit déxtera ejus.

Et ejécit a fácie eórum gentes: * et sorte divísit eis terram in funículo distributiónis.

Et habitáre fecit in tabernáculis eórum * tribus Israël.

Et tentavérunt, et exacerbavérunt Deum excélsum: * et testimónia ejus non custodiérunt.

Et avertérunt se, et non servavérunt pactum: * quemádmodum patres eórum, convérsi sunt in arcum pravum.

In iram concitavérunt eum in cóllibus suis: * et in sculptílibus suis ad æmulatiónem eum provocavérunt.

\antlabel{Ant.} Redémit eos Dóminus de manu tribulántis.

\antlabel{Ant.} Ædificávit * Deus sanctifícium suum in terra.

\psalmheading{Psalmus 77, vi}

\lettrine[lines=2]{A}{UDIVIT} Deus, et sprevit: * et ad níhilum redégit valde Israël.

Et répulit tabernáculum Silo, * tabernáculum suum, ubi habitávit in homínibus.

Et trádidit in captivitátem virtútem eórum: * et pulchritúdinem eórum in manus inimíci.

Et conclúsit in gládio pópulum suum: * et hereditátem suam sprevit.

Júvenes eórum comédit ignis: * et vírgines eórum non sunt lamentátæ.

Sacerdótes eórum in gládio cecidérunt: * et víduæ eórum non plorabántur.

Et excitátus est tamquam dórmiens Dóminus, * tamquam potens crapulátus a vino.

Et percússit inimícos suos in posterióra: * oppróbrium sempitérnum dedit illis.

Et répulit tabernáculum Joseph: * et tribum Ephraim non elégit:

Sed elégit tribum Juda, * montem Sion quem diléxit.

Et ædificávit sicut unicórnium sanctifícium suum in terra, * quam fundávit in sæcula.

Et elégit David, servum suum, et sústulit eum de grégibus óvium: * de post fetántes accépit eum.

Páscere Jacob, servum suum, * et Israël, hereditátem suam:

Et pavit eos in innocéntia cordis sui: * et in intelléctibus mánuum suárum dedúxit eos.

\antlabel{Ant.} Ædificávit Deus sanctifícium suum in terra.

\variantheading{In Festis novem Lectionum:}

\V\ Elégit Dóminus tribum Juda. \R\ Montem Sion quem diléxit.

\absolutiolabel{Absolutio.} Ipsíus píetas et.

\benedlabel{Bened.} 4 Deus Pater omnípotens. 5 Christus perpétuæ. 6 Ignem sui amóris.

\nocturnoheading{In III Nocturno}

\antlabel{Ant.} Adjuva nos, * Deus, salutáris noster: et propítius esto peccátis nostris.

\psalmheading{Psalmus 78}

\lettrine[lines=2]{D}{EUS}, venérunt gentes in hereditátem tuam, polluérunt templum sanctum tuum: * posuérunt Jerúsalem in pomórum custódiam.

Posuérunt morticína servórum tuórum, escas volatílibus cæli: * carnes sanctórum tuórum béstiis terræ.

Effudérunt sánguinem eórum tamquam aquam in circúitu Jerúsalem: * et non erat qui sepelíret.

Facti sumus oppróbrium vicínis nostris: * subsannátio et illúsio his, qui in circúitu nostro sunt.

Usquequo, Dómine, irascéris in finem: * accendétur velut ignis zelus tuus?

Effúnde iram tuam in gentes, quæ te non novérunt: * et in regna, quæ nomen tuum non invocavérunt:

Quia comedérunt Jacob: * et locum ejus desolavérunt.

Ne memíneris iniquitátum nostrárum antiquárum, cito antícipent nos misericórdiæ tuæ: * quia páuperes facti sumus nimis.

Adjuva nos, Deus, salutáris noster: et propter glóriam nóminis tui, Dómine, líbera nos: * et propítius esto peccátis nostris, propter nomen tuum:

Ne forte dicant in géntibus: Ubi est Deus eórum? * et innotéscat in natiónibus coram óculis nostris.

Ultio sánguinis servórum tuórum, qui effúsus est: * intróeat in conspéctu tuo gémitus compeditórum.

Secúndum magnitúdinem bráchii tui, * pósside fílios mortificatórum.

Et redde vicínis nostris séptuplum in sinu eórum: * impropérium ipsórum, quod exprobravérunt tibi, Dómine.

Nos autem pópulus tuus, et oves páscuæ tuæ, * confitébimur tibi in sæculum.

In generatiónem et generatiónem * annuntiábimus laudem tuam.

\antlabel{Ant.} Adjuva nos, Deus, salutáris noster: et propítius esto peccátis nostris.

\antlabel{Ant.} Ego sum Dóminus, * Deus tuus, Israël, qui edúxi te de terra Ægýpti.

\psalmheading{Psalmus 80}

\lettrine[lines=2]{E}{XSULTATE} Deo, adjutóri nostro: * jubiláte Deo Jacob.

Súmite psalmum, et date týmpanum: * psaltérium jucúndum cum cíthara.

Buccináte in Noeménia tuba, * in insígni die solemnitátis vestræ.

Quia præcéptum in Israël est: * et judícium Deo Jacob.

Testimónium in Joseph pósuit illud, cum exíret de terra Ægýpti: * linguam, quam non nóverat, audívit.

Divértit ab onéribus dorsum ejus: * manus ejus in cóphino serviérunt.

In tribulatióne invocásti me, et liberávi te: * exaudívi te in abscóndito tempestátis: probávi te apud aquam contradictiónis.

Audi, pópulus meus, et contestábor te: * Israël, si audíeris me, non erit in te deus recens, neque adorábis deum aliénum.

Ego enim sum Dóminus, Deus tuus, qui edúxi te de terra Ægýpti: * diláta os tuum, et implébo illud.

Et non audívit pópulus meus vocem meam: * et Israël non inténdit mihi.

Et dimísi eos secúndum desidéria cordis eórum, * ibunt in adinventiónibus suis.

Si pópulus meus audísset me: * Israël si in viis meis ambulásset:

Pro níhilo fórsitan inimícos eórum humiliássem: * et super tribulántes eos misíssem manum meam.

Inimíci Dómini mentíti sunt ei: * et erit tempus eórum in sæcula.

Et cibávit eos ex ádipe fruménti: * et de petra, melle saturávit eos.

\antlabel{Ant.} Ego sum Dóminus, Deus tuus, Israël, qui edúxi te de terra Ægýpti.

\antlabel{Ant.} Ne táceas, Deus, * quóniam inimíci tui extulérunt caput.

\psalmheading{Psalmus 82}

\lettrine[lines=2]{D}{EUS}, quis símilis erit tibi? * ne táceas, neque compescáris, Deus:

Quóniam ecce inimíci tui sonuérunt: * et qui odérunt te, extulérunt caput.

Super pópulum tuum malignavérunt consílium: * et cogitavérunt advérsus sanctos tuos.

Dixérunt: Veníte, et disperdámus eos de gente: * et non memorétur nomen Israël ultra.

Quóniam cogitavérunt unanímiter: * simul advérsum te testaméntum disposuérunt, tabernácula Idumæórum et Ismahelítæ:

Moab, et Agaréni, Gebal, et Ammon, et Amalec: * alienígenæ cum habitántibus Tyrum.

Etenim Assur venit cum illis: * facti sunt in adjutórium fíliis Lot.

Fac illis sicut Mádian, et Sísaræ: * sicut Jabin in torrénte Cisson.

Disperiérunt in Endor: * facti sunt ut stercus terræ.

Pone príncipes eórum sicut Oreb, et Zeb, * et Zébee, et Sálmana:

Omnes príncipes eórum: * qui dixérunt: Hereditáte possideámus Sanctuárium Dei.

Deus meus, pone illos ut rotam: * et sicut stípulam ante fáciem venti.

Sicut ignis, qui combúrit silvam: * et sicut flamma combúrens montes:

Ita persequéris illos in tempestáte tua: * et in ira tua turbábis eos.

Imple fácies eórum ignomínia: * et quærent nomen tuum, Dómine.

Erubéscant, et conturbéntur in sæculum sæculi: * et confundántur et péreant.

Et cognóscant quia nomen tibi Dóminus: * tu solus Altíssimus in omni terra.

\antlabel{Ant.} Ne táceas, Deus, quóniam inimíci tui extulérunt caput.

\variantheading{In feriali Officio et in Festis:}

\V\ Cognóscant quia nomen tibi Dóminus. \R\ Tu solus Altíssimus super omnem terram.

\variantheading{In Officio novem Lectionum}

\absolutiolabel{Absolutio.} A vínculis.

\benedlabel{Bened.} 7 Evangélica léctio. 8 In Festis Dómini: Divínum auxílium. In Festis B. M. V.: Cujus festum cólimus, ipsa Virgo. In Festis Sanctórum: Cujus (\textit{vel} Quorum \textit{aut} Quarum) festum cólimus. 9 Ad societátem. (Per evangélica dicta.)

\variantheading{In Officio trium Lectionum}

\absolutiolabel{Absolutio.} Ipsíus píetas et.

\variantheading{In Feriis cum Homilia:}

\benedlabel{Bened.} 1 Evangélica. 2 Divínum. 3 Ad societátem.

\variantheading{In Feriis sine Homilia:}

\benedlabel{Bened.} 1 Deus Pater. 2 Christus perpétuæ. 3 Ignem sui amóris.

\variantheading{In Festis Sanctorum:}

\benedlabel{Bened.} 1 Ille nos benedícat, qui sine fine vivit et regnat. \R\ Amen. 2 Cujus (\textit{vel} Quorum \textit{aut} Quarum.) 3 Ad societátem.

% ===== END OF BATCH 30 -- Covers autumnalis.md lines 9891-10290,
% opening "Feria Sexta" (Friday) and completing its "Ad Matutinum"
% in full: the ferial Invitatorium "Dóminum, Deum nostrum" (with a
% bracketed footnote marker "[3]" in the source preserved verbatim
% after the incipit), the hymn "Tu, Trinitátis Únitas," and the
% three Nocturns -- Psalm 77 split across six sections (i-vi) over
% the first two Nocturns, and Psalms 78, 80, 82 in the third --
% with their versicle, Absolutio, and Benedictiones apparatus. The
% day-heading itself appears in the source in full capitals ("FERIA
% SEXTA"), transcribed with \feriaheading{} regardless, consistent
% with the role-based heading convention applied throughout this
% volume. Several more psalm incipits carry the partial-caps OCR
% defect already documented in prior batches ("ATtendite,"
% "FIlii," "ET apposuérunt" is correctly two full words so left
% as-is being a short conjunction, "AUdivit," "DEus" twice,
% "EXsultáte") and are silently normalized to full caps where the
% defect is genuine. The document continues with "Ad Laudes I"
% (autumnalis.md line 10292).
% =====
\hourheading{Ad Laudes I}

\rubric{In Fériis extra Vigílias commúnes occurréntibus, atque in Festis.}

(Pater noster et Ave.)

\V\ Deus, in adjutórium.

\antlabel{Ant.} Exaltáte * Dóminum, Deum nostrum, et adoráte in monte sancto ejus.

\psalmheading{Psalmus 98}

\lettrine[lines=2]{D}{OMINUS} regnávit, irascántur pópuli: * qui sedet super Chérubim, moveátur terra.

Dóminus in Sion magnus: * et excélsus super omnes pópulos.

Confiteántur nómini tuo magno: quóniam terríbile, et sanctum est: * et honor regis judícium díligit.

Tu parásti directiónes: * judícium et justítiam in Jacob tu fecísti.

Exaltáte Dóminum, Deum nostrum, et adoráte scabéllum pedum ejus: * quóniam sanctum est.

Móyses et Aaron in sacerdótibus ejus: * et Sámuel inter eos, qui ínvocant nomen ejus:

Invocábant Dóminum, et ipse exaudiébat eos: * in colúmna nubis loquebátur ad eos.

Custodiébant testimónia ejus, * et præcéptum quod dedit illis.

Dómine, Deus noster, tu exaudiébas eos: * Deus, tu propítius fuísti eis, et ulcíscens in omnes adinventiónes eórum.

Exaltáte Dóminum, Deum nostrum, et adoráte in monte sancto ejus: * quóniam sanctus Dóminus, Deus noster.

\antlabel{Ant.} Exaltáte Dóminum, Deum nostrum, et adoráte in monte sancto ejus.

\antlabel{Ant.} Eripe me * de inimícis meis, Dómine, ad te confúgi.

\psalmheading{Psalmus 142}

\lettrine[lines=2]{D}{OMINE}, exáudi oratiónem meam: áuribus pércipe obsecratiónem meam in veritáte tua: * exáudi me in tua justítia.

Et non intres in judícium cum servo tuo: * quia non justificábitur in conspéctu tuo omnis vivens.

Quia persecútus est inimícus ánimam meam: * humiliávit in terra vitam meam.

Collocávit me in obscúris sicut mórtuos sæculi: * et anxiátus est super me spíritus meus, in me turbátum est cor meum.

Memor fui diérum antiquórum, meditátus sum in ómnibus opéribus tuis: * in factis mánuum tuárum meditábar.

Expándi manus meas ad te: * ánima mea sicut terra sine aqua tibi.

Velóciter exáudi me, Dómine: * defécit spíritus meus.

Non avértas fáciem tuam a me: * et símilis ero descendéntibus in lacum.

Audítam fac mihi mane misericórdiam tuam: * quia in te sperávi.

Notam fac mihi viam, in qua ámbulem: * quia ad te levávi ánimam meam.

Eripe me de inimícis meis, Dómine, ad te confúgi: * doce me fácere voluntátem tuam, quia Deus meus es tu.

Spíritus tuus bonus dedúcet me in terram rectam: * propter nomen tuum, Dómine, vivificábis me, in æquitáte tua.

Edúces de tribulatióne ánimam meam: * et in misericórdia tua dispérdes inimícos meos.

Et perdes omnes, qui tríbulant ánimam meam: * quóniam ego servus tuus sum.

\antlabel{Ant.} Eripe me de inimícis meis, Dómine, ad te confúgi.

\antlabel{Ant.} Benedixísti, * \textdagger\ Dómine, terram tuam: remisísti iniquitátem plebis tuæ.

\psalmheading{Psalmus 84}

\lettrine[lines=2]{B}{ENEDIXISTI}, \textdagger\ Dómine, terram tuam: * avertísti captivitátem Jacob.

Remisísti iniquitátem plebis tuæ: * operuísti ómnia peccáta eórum.

Mitigásti omnem iram tuam: * avertísti ab ira indignatiónis tuæ.

Convérte nos, Deus, salutáris noster: * et avérte iram tuam a nobis.

Numquid in ætérnum irascéris nobis? * aut exténdes iram tuam a generatióne in generatiónem?

Deus, tu convérsus vivificábis nos: * et plebs tua lætábitur in te.

Osténde nobis, Dómine, misericórdiam tuam: * et salutáre tuum da nobis.

Audiam quid loquátur in me Dóminus Deus: * quóniam loquétur pacem in plebem suam.

Et super sanctos suos: * et in eos, qui convertúntur ad cor.

Verúmtamen prope timéntibus eum salutáre ipsíus: * ut inhábitet glória in terra nostra.

Misericórdia, et véritas obviavérunt sibi: * justítia, et pax osculátæ sunt.

Véritas de terra orta est: * et justítia de cælo prospéxit.

Etenim Dóminus dabit benignitátem: * et terra nostra dabit fructum suum.

Justítia ante eum ambulábit: * et ponet in via gressus suos.

\antlabel{Ant.} Benedixísti, Dómine, terram tuam: remisísti iniquitátem plebis tuæ.

\antlabel{Ant.} In Dómino justificábitur, * et laudábitur omne semen Israël.

\canticumheading{Canticum Isaiæ}{Is. 45, 15-26}

\lettrine[lines=2]{V}{ERE} tu es Deus abscónditus, * Deus Israël, Salvátor.

Confúsi sunt, et erubuérunt omnes: * simul abiérunt in confusiónem fabricatóres errórum.

Israël salvátus est in Dómino salúte ætérna: * non confundémini, et non erubescétis usque in sæculum sæculi.

Quia hæc dicit Dóminus creans cælos, * ipse Deus formans terram, et fáciens eam, ipse plastes ejus:

Non in vanum creávit eam, ut habitarétur formávit eam: * Ego Dóminus, et non est álius.

Non in abscóndito locútus sum, * in loco terræ tenebróso:

Non dixi sémini Jacob frustra: Quærite me: * ego Dóminus loquens justítiam, annúntians recta.

Congregámini, et veníte, et accédite simul * qui salváti estis ex géntibus:

Nesciérunt qui levant lignum sculptúræ suæ, * et rogant deum non salvántem.

Annuntiáte, et veníte, et consiliámini simul: * Quis audítum fecit hoc ab inítio, ex tunc prædíxit illud?

Numquid non ego Dóminus, et non est ultra Deus absque me? * Deus justus, et salvans non est præter me.

Convertímini ad me, et salvi éritis, omnes fines terræ: * quia ego Deus, et non est álius.

In memetípso jurávi, egrediétur de ore meo justítiæ verbum, * et non revertétur:

Quia mihi curvábitur omne genu, * et jurábit omnis lingua.

Ergo in Dómino, dicet, meæ sunt justítiæ et impérium: * ad eum vénient, et confundéntur omnes qui repúgnant ei.

In Dómino justificábitur, et laudábitur * omne semen Israël.

\antlabel{Ant.} In Dómino justificábitur, et laudábitur omne semen Israël.

\antlabel{Ant.} Lauda, * \textdagger\ Jerúsalem, Dóminum. \textdaggerdbl

\psalmheading{Psalmus 147}

\lettrine[lines=2]{L}{AUDA}, \textdagger\ Jerúsalem, Dóminum: * \textdaggerdbl\ lauda Deum tuum, Sion.

Quóniam confortávit seras portárum tuárum: * benedíxit fíliis tuis in te.

Qui pósuit fines tuos pacem: * et ádipe fruménti sátiat te.

Qui emíttit elóquium suum terræ: * velóciter currit sermo ejus.

Qui dat nivem sicut lanam: * nébulam sicut cínerem spargit.

Mittit crystállum suam sicut buccéllas: * ante fáciem frígoris ejus quis sustinébit?

Emíttet verbum suum, et liquefáciet ea: * flabit spíritus ejus, et fluent aquæ.

Qui annúntiat verbum suum Jacob: * justítias, et judícia sua Israël.

Non fecit táliter omni natióni: * et judícia sua non manifestávit eis.

\antlabel{Ant.} Lauda, Jerúsalem, Dóminum.

\variantheading{In Officio feriali}

\capitulumlabel{Capit.}{Rom. 13, 12-13}

\lettrine[lines=2]{N}{OX} præcéssit, dies autem appropinquávit. Abjiciámus ergo ópera tenebrárum, et induámur arma lucis. Sicut in die honéste ambulémus.

\hymnusheading{Hymnus}

\lettrine[lines=2]{Æ}{TERNA} cæli glória,\\
Beáta spes mortálium,\\
Summi Tonántis Únice,\\
Castǽque proles Vírginis:

Da déxteram surgéntibus,\\
Exsúrgat et mens sóbria,\\
Flagrans et in laudem Dei\\
Grates rependat débitas.

Ortus refúlget Lúcifer,\\
Præítque solem núntius:\\
Cadunt tenébræ nóctium:\\
Lux sancta nos illúminet.

Manénsque nostris sénsibus,\\
Noctem repéllat sæculi,\\
Omníque fine témporis\\
Purgáta servet péctora.

Quæsíta jam primum fides\\
In corde radíces agat:\\
Secúnda spes congáudeat,\\
Qua major exstat cáritas.

Deo Patri sit glória,\\
Ejúsque soli Fílio,\\
Cum Spíritu Paráclito,\\
Nunc et per omne sæculum.\\
Amen.

\V\ Repléti sumus mane misericórdia tua. \R\ Exsultávimus, et delectáti sumus.

\antlabel{Ad Bened. Ant.} Per víscera misericórdiæ * Dei nostri visitávit nos Oriens ex alto.

\rubric{Réliqua ut in Ordinário. 9}

% ===== END OF BATCH 31 -- Covers autumnalis.md lines 10292-10502,
% "Ad Laudes I" for Feria Sexta in full: Psalms 98, 142, 84, the
% Canticum Isaiæ, and Psalm 147, followed by the Capitulum "Nox
% præcéssit" and the hymn "Æterna cæli glória" (whose opening
% lettrine letter is the ligature "Æ," matching the established
% convention of lettrine-ing exactly the source's first letter).
% Several more psalm incipits carry the partial-caps OCR defect
% ("DOmine, exáudi," "BEnedixísti," "VEre tu es," "LAuda,
% Jerúsalem") and are silently normalized to full caps. The
% document continues with "Ad Laudes II" (autumnalis.md line 10505).
% =====
\hourheading{Ad Laudes II}

\rubric{In Vigíliis commúnibus, quando fit Offícium de Féria.}

(Pater noster et Ave.)

\V\ Deus, in adjutórium.

\antlabel{Ant.} Cor contrítum.

\psalmheading{Psalmus 50}

\lettrine[lines=2]{M}{ISERERE} mei, Deus, * secúndum magnam misericórdiam tuam.

Et secúndum multitúdinem miseratiónum tuárum, * dele iniquitátem meam.

Amplius lava me ab iniquitáte mea: * et a peccáto meo munda me.

Quóniam iniquitátem meam ego cognósco: * et peccátum meum contra me est semper.

Tibi soli peccávi, et malum coram te feci: * ut justificéris in sermónibus tuis, et vincas cum judicáris.

Ecce enim, in iniquitátibus concéptus sum: * et in peccátis concépit me mater mea.

Ecce enim, veritátem dilexísti: * incérta et occúlta sapiéntiæ tuæ manifestásti mihi.

Aspérges me hyssópo, et mundábor: * lavábis me, et super nivem dealbábor.

Audítui meo dabis gáudium et lætítiam: * et exsultábunt ossa humiliáta.

Avérte fáciem tuam a peccátis meis: * et omnes iniquitátes meas dele.

Cor mundum crea in me, Deus: * et spíritum rectum ínnova in viscéribus meis.

Ne projícias me a fácie tua: * et spíritum sanctum tuum ne áuferas a me.

Redde mihi lætítiam salutáris tui: * et spíritu principáli confírma me.

Docébo iníquos vias tuas: * et ímpii ad te converténtur.

Líbera me de sanguínibus, Deus, Deus salútis meæ: * et exsultábit lingua mea justítiam tuam.

Dómine, lábia mea apéries: * et os meum annuntiábit laudem tuam.

Quóniam si voluísses sacrifícium, dedíssem útique: * holocáustis non delectáberis.

Sacrifícium Deo spíritus contribulátus: * cor contrítum, et humiliátum, Deus, non despícies.

Benígne fac, Dómine, in bona voluntáte tua Sion: * ut ædificéntur muri Jerúsalem.

Tunc acceptábis sacrifícium justítiæ, oblatiónes, et holocáusta: * tunc impónent super altáre tuum vítulos.

\antlabel{Ant.} Cor contrítum, et humiliátum, Deus, non despícies.

\antlabel{Ant.} Propter nomen tuum.

\psalmheading{Psalmus 142}

\lettrine[lines=2]{D}{OMINE}, exáudi oratiónem meam: áuribus pércipe obsecratiónem meam in veritáte tua: * exáudi me in tua justítia.

Et non intres in judícium cum servo tuo: * quia non justificábitur in conspéctu tuo omnis vivens.

Quia persecútus est inimícus ánimam meam: * humiliávit in terra vitam meam.

Collocávit me in obscúris sicut mórtuos sæculi: * et anxiátus est super me spíritus meus, in me turbátum est cor meum.

Memor fui diérum antiquórum, meditátus sum in ómnibus opéribus tuis: * in factis mánuum tuárum meditábar.

Expándi manus meas ad te: * ánima mea sicut terra sine aqua tibi.

Velóciter exáudi me, Dómine: * defécit spíritus meus.

Non avértas fáciem tuam a me: * et símilis ero descendéntibus in lacum.

Audítam fac mihi mane misericórdiam tuam: * quia in te sperávi.

Notam fac mihi viam, in qua ámbulem: * quia ad te levávi ánimam meam.

Eripe me de inimícis meis, Dómine, ad te confúgi: * doce me fácere voluntátem tuam, quia Deus meus es tu.

Spíritus tuus bonus dedúcet me in terram rectam: * propter nomen tuum, Dómine, vivificábis me, in æquitáte tua.

Edúces de tribulatióne ánimam meam: * et in misericórdia tua dispérdes inimícos meos.

Et perdes omnes, qui tríbulant ánimam meam: * quóniam ego servus tuus sum.

\antlabel{Ant.} Propter nomen tuum, Dómine, vivificábis me, in æquitáte tua.

\antlabel{Ant.} Deus, tu convérsus.

\psalmheading{Psalmus 84}

\lettrine[lines=2]{B}{ENEDIXISTI}, Dómine, terram tuam: * avertísti captivitátem Jacob.

Remisísti iniquitátem plebis tuæ: * operuísti ómnia peccáta eórum.

Mitigásti omnem iram tuam: * avertísti ab ira indignatiónis tuæ.

Convérte nos, Deus, salutáris noster: * et avérte iram tuam a nobis.

Numquid in ætérnum irascéris nobis? * aut exténdes iram tuam a generatióne in generatiónem?

Deus, tu convérsus vivificábis nos: * et plebs tua lætábitur in te.

Osténde nobis, Dómine, misericórdiam tuam: * et salutáre tuum da nobis.

Audiam quid loquátur in me Dóminus Deus: * quóniam loquétur pacem in plebem suam.

Et super sanctos suos: * et in eos, qui convertúntur ad cor.

Verúmtamen prope timéntibus eum salutáre ipsíus: * ut inhábitet glória in terra nostra.

Misericórdia, et véritas obviavérunt sibi: * justítia, et pax osculátæ sunt.

Véritas de terra orta est: * et justítia de cælo prospéxit.

Etenim Dóminus dabit benignitátem: * et terra nostra dabit fructum suum.

Justítia ante eum ambulábit: * et ponet in via gressus suos.

\antlabel{Ant.} Deus, tu convérsus vivificábis nos: et plebs tua lætábitur in te.

\antlabel{Ant.} Cum irátus fúeris.

\canticumheading{Canticum Habacuc}{Hab. 3, 2-19}

\lettrine[lines=2]{D}{OMINE}, audívi auditiónem tuam, * et tímui.

Dómine, opus tuum, * in médio annórum vivífica illud:

In médio annórum notum fácies: * cum irátus fúeris, misericórdiæ recordáberis.

Deus ab Austro véniet, * et sanctus de monte Pharan:

Opéruit cælos glória ejus: * et laudis ejus plena est terra.

Splendor ejus ut lux erit: * córnua in mánibus ejus:

Ibi abscóndita est fortitúdo ejus: * ante fáciem ejus ibit mors.

Et egrediétur diábolus ante pedes ejus. * Stetit, et mensus est terram.

Aspéxit, et dissólvit gentes: * et contríti sunt montes sæculi.

Incurváti sunt colles mundi, * ab itinéribus æternitátis ejus.

Pro iniquitáte vidi tentória Æthiópiæ, * turbabúntur pelles terræ Mádian.

Numquid in flumínibus irátus es, Dómine? * aut in flumínibus furor tuus? vel in mari indignátio tua?

Qui ascéndes super equos tuos: * et quadrígæ tuæ salvátio.

Súscitans suscitábis arcum tuum: * juraménta tríbubus quæ locútus es.

Flúvios scindes terræ: * vidérunt te, et doluérunt montes: gurges aquárum tránsiit.

Dedit abýssus vocem suam: * altitúdo manus suas levávit.

Sol, et luna stetérunt in habitáculo suo, * in luce sagittárum tuárum, ibunt in splendóre fulgurántis hastæ tuæ.

In frémitu conculcábis terram: * et in furóre obstupefácies gentes.

Egréssus es in salútem pópuli tui: * in salútem cum Christo tuo.

Percussísti caput de domo ímpii: * denudásti fundaméntum ejus usque ad collum.

Maledixísti sceptris ejus, cápiti bellatórum ejus, * veniéntibus ut turbo ad dispergéndum me.

Exsultátio eórum * sicut ejus, qui dévorat páuperem in abscóndito.

Viam fecísti in mari equis tuis, * in luto aquárum multárum.

Audívi, et conturbátus est venter meus: * a voce contremuérunt lábia mea.

Ingrediátur putrédo in óssibus meis, * et subter me scáteat.

Ut requiéscam in die tribulatiónis: * ut ascéndam ad pópulum accínctum nostrum.

Ficus enim non florébit: * et non erit germen in víneis.

Mentiétur opus olívæ: * et arva non áfferent cibum.

Abscindétur de ovíli pecus: * et non erit arméntum in præsépibus.

Ego autem in Dómino gaudébo: * et exsultábo in Deo Jesu meo.

Deus Dóminus fortitúdo mea: * et ponet pedes meos quasi cervórum.

Et super excélsa mea dedúcet me victor * in psalmis canéntem.

\antlabel{Ant.} Cum irátus fúeris, Dómine, misericórdiæ recordáberis.

\antlabel{Ant.} Lauda. \textdagger

\psalmheading{Psalmus 147}

\lettrine[lines=2]{L}{AUDA}, \textdagger\ Jerúsalem, Dóminum: * lauda Deum tuum, Sion.

Quóniam confortávit seras portárum tuárum: * benedíxit fíliis tuis in te.

Qui pósuit fines tuos pacem: * et ádipe fruménti sátiat te.

Qui emíttit elóquium suum terræ: * velóciter currit sermo ejus.

Qui dat nivem sicut lanam: * nébulam sicut cínerem spargit.

Mittit crystállum suam sicut buccéllas: * ante fáciem frígoris ejus quis sustinébit?

Emíttet verbum suum, et liquefáciet ea: * flabit spíritus ejus, et fluent aquæ.

Qui annúntiat verbum suum Jacob: * justítias, et judícia sua Israël.

Non fecit táliter omni natióni: * et judícia sua non manifestávit eis.

\antlabel{Ant.} Lauda Deum tuum, Sion, qui annúntiat judícia sua Israël.

\capitulumlabel{Capit.}{Rom. 13, 12-13}

\lettrine[lines=2]{N}{OX} præcéssit, dies autem appropinquávit. Abjiciámus ergo ópera tenebrárum, et induámur arma lucis. Sicut in die honéste ambulémus.

\hymnusheading{Hymnus}

\lettrine[lines=2]{Æ}{TERNA} cæli glória,\\
Beáta spes mortálium,\\
Summi Tonántis Únice,\\
Castǽque proles Vírginis:

Da déxteram surgéntibus,\\
Exsúrgat et mens sóbria,\\
Flagrans et in laudem Dei\\
Grates rependat débitas.

Ortus refúlget Lúcifer,\\
Præítque solem núntius:\\
Cadunt tenébræ nóctium:\\
Lux sancta nos illúminet.

Manénsque nostris sénsibus,\\
Noctem repéllat sæculi,\\
Omníque fine témporis\\
Purgáta servet péctora.

Quæsíta jam primum fides\\
In corde radíces agat:\\
Secúnda spes congáudeat,\\
Qua major exstat cáritas.

Deo Patri sit glória,\\
Ejúsque soli Fílio,\\
Cum Spíritu Paráclito,\\
Nunc et per omne sæculum.\\
Amen.

\V\ Repléti sumus mane misericórdia tua. \R\ Exsultávimus, et delectáti sumus.

\antlabel{Ad Bened. Ant.} Per víscera misericórdiæ * Dei nostri visitávit nos Oriens ex alto.

\rubric{Expleto Cántico ac repetíta Antíphona, dicúntur sequéntes}

\variantheading{Preces feriales}

KYRIE, eléison. Christe, eléison. Kýrie, eléison.

Pater noster \rubric{quod dícitur a solo Hebdomadário totum clara voce, usque ad} \V\ Et ne nos indúcas in tentatiónem. \R\ Sed líbera nos a malo. \V\ Ego dixi: Dómine, miserére mei. \R\ Sana ánimam meam, quia peccávi tibi. \V\ Convértere, Dómine, úsquequo? \R\ Et deprecábilis esto super servos tuos. \V\ Fiat misericórdia tua, Dómine, super nos. \R\ Quemádmodum sperávimus in te. \V\ Sacerdótes tui induántur justítiam. \R\ Et sancti tui exsúltent. \V\ Orémus pro beatíssimo Papa nostro N. \R\ Dóminus consérvet eum, et vivíficet eum, et beátum fáciat eum in terra, et non tradat eum in ánimam inimicórum ejus. \V\ Orémus et pro Antístite nostro N. \R\ Stet et pascat in fortitúdine tua, Dómine, in sublimitáte nóminis tui. \V\ Dómine, salvum fac regem. \R\ Et exáudi nos in die, qua invocavérimus te. \V\ Salvum fac pópulum tuum, Dómine, et bénedic hereditáti tuæ. \R\ Et rege eos, et extólle illos usque in ætérnum. \V\ Meménto Congregatiónis tuæ. \R\ Quam possedísti ab inítio. \V\ Fiat pax in virtúte tua. \R\ Et abundántia in túrribus tuis. \V\ Orémus pro benefactóribus nostris. \R\ Retribúere dignáre, Dómine, ómnibus, nobis bona faciéntibus propter nomen tuum, vitam ætérnam. Amen. \V\ Orémus pro fidélibus defúnctis. \R\ Réquiem ætérnam dona eis, Dómine, et lux perpétua lúceat eis. \V\ Requiéscant in pace. \R\ Amen. \V\ Pro frátribus nostris abséntibus. \R\ Salvos fac servos tuos, Deus meus, sperántes in te. \V\ Pro afflíctis et captívis. \R\ Líbera eos, Deus Israël, ex ómnibus tribulatiónibus suis. \V\ Mitte eis, Dómine, auxílium de sancto. \R\ Et de Sion tuére eos. \V\ Dómine, Deus virtútum, convérte nos. \R\ Et osténde fáciem tuam, et salvi érimus. \V\ Exsúrge, Christe, ádjuva nos. \R\ Et líbera nos propter nomen tuum. \V\ Dómine, exáudi oratiónem meam. \R\ Et clamor meus ad te véniat.

\rubric{Réliqua ut in Ordinário. 9}

% ===== END OF BATCH 32 -- Covers autumnalis.md lines 10505-10775,
% "Ad Laudes II" for Feria Sexta in full: Psalms 50, 142, 84, the
% Canticum Habacuc ("Dómine, audívi auditiónem tuam"), and Psalm
% 147, sharing the Capitulum "Nox præcéssit" and the hymn "Æterna
% cæli glória" with Ad Laudes I, plus the "Preces feriales" litany.
% Psalms 142 and 84 appear here without the dagger chant-pointing
% mark present in their Ad Laudes I occurrences (Batch 31), a
% genuine source divergence transcribed exactly as given in each
% place. The document continues with "Ad Primam" (autumnalis.md line
% 10778).
% =====
\hourheading{Ad Primam}

Pater, Ave et Credo.

\V\ Deus, in adjutórium.

\hymnusheading{Hymnus}

\lettrine[lines=2]{J}{AM} lucis orto sídere,\\
Deum precémur súpplices,\\
Ut in diúrnis áctibus\\
Nos servet a nocéntibus.

Linguam refrénans témperet,\\
Ne litis horror ínsonet:\\
Visum fovéndo cóntegat,\\
Ne vanitátes háuriat.

Sint pura cordis íntima,\\
Absístat et vecórdia:\\
Carnis terat supérbiam\\
Potus cibíque párcitas.

Ut, cum dies abscésserit\\
Noctémque sors redúxerit,\\
Mundi per abstinéntiam\\
Ipsi canámus glóriam.

Deo Patri sit glória,\\
Ejúsque soli Fílio,\\
Cum Spíritu Paráclito,\\
Nunc et per omne sæculum.\\
Amen.

\antlabel{Ant.} Ne discédas a me.

\psalmheading{Psalmus 21, i}

\lettrine[lines=2]{D}{EUS}, Deus meus, réspice in me: quare me dereliquísti? * longe a salúte mea verba delictórum meórum.

Deus meus, clamábo per diem, et non exáudies: * et nocte, et non ad insipiéntiam mihi.

Tu autem in sancto hábitas, * laus Israël.

In te speravérunt patres nostri: * speravérunt, et liberásti eos.

Ad te clamavérunt, et salvi facti sunt: * in te speravérunt, et non sunt confúsi.

Ego autem sum vermis, et non homo: * oppróbrium hóminum, et abjéctio plebis.

Omnes vidéntes me, derisérunt me: * locúti sunt lábiis, et movérunt caput.

Sperávit in Dómino, éripiat eum: * salvum fáciat eum, quóniam vult eum.

Quóniam tu es, qui extraxísti me de ventre: * spes mea ab ubéribus matris meæ. In te projéctus sum ex útero:

De ventre matris meæ Deus meus es tu, * ne discésseris a me.

Quóniam tribulátio próxima est: * quóniam non est qui ádjuvet.

\psalmheading{Psalmus 21, ii}

\lettrine[lines=2]{C}{IRCUMDEDERUNT} me vítuli multi: * tauri pingues obsedérunt me.

Aperuérunt super me os suum, * sicut leo rápiens et rúgiens.

Sicut aqua effúsus sum: * et dispérsa sunt ómnia ossa mea.

Factum est cor meum tamquam cera liquéscens * in médio ventris mei.

Aruit tamquam testa virtus mea, et lingua mea adhæsit fáucibus meis: * et in púlverem mortis deduxísti me.

Quóniam circumdedérunt me canes multi: * concílium malignántium obsédit me.

Fodérunt manus meas et pedes meos: * dinumeravérunt ómnia ossa mea.

Ipsi vero consideravérunt et inspexérunt me: * divisérunt sibi vestiménta mea, et super vestem meam misérunt sortem.

Tu autem, Dómine, ne elongáveris auxílium tuum a me: * ad defensiónem meam cónspice.

Erue a frámea, Deus, ánimam meam: * et de manu canis únicam meam:

Salva me ex ore leónis: * et a córnibus unicórnium humilitátem meam.

Narrábo nomen tuum frátribus meis: * in médio ecclésiæ laudábo te.

\psalmheading{Psalmus 21, iii}

\lettrine[lines=2]{Q}{UI} timétis Dóminum, laudáte eum: * univérsum semen Jacob, glorificáte eum.

Tímeat eum omne semen Israël: * quóniam non sprevit, neque despéxit deprecatiónem páuperis:

Nec avértit fáciem suam a me: * et cum clamárem ad eum, exaudívit me.

Apud te laus mea in ecclésia magna: * vota mea reddam in conspéctu timéntium eum.

Edent páuperes, et saturabúntur: et laudábunt Dóminum qui requírunt eum: * vivent corda eórum in sæculum sæculi.

Reminiscéntur et converténtur ad Dóminum * univérsi fines terræ:

Et adorábunt in conspéctu ejus * univérsæ famíliæ géntium.

Quóniam Dómini est regnum: * et ipse dominábitur géntium.

Manducavérunt et adoravérunt omnes pingues terræ: * in conspéctu ejus cadent omnes qui descéndunt in terram.

Et ánima mea illi vivet: * et semen meum sérviet ipsi.

Annuntiábitur Dómino generátio ventúra: * et annuntiábunt cæli justítiam ejus pópulo qui nascétur, quem fecit Dóminus.

\rubric{¶ Quando ad Laudes dictus fúerit Ps. 50 Miserére, hic subjúngitur sequens Ps. 98, qui in Laudibus locum non hábuit.}

\psalmheading{Psalmus 98}

\lettrine[lines=2]{D}{OMINUS} regnávit, irascántur pópuli: * qui sedet super Chérubim, moveátur terra.

Dóminus in Sion magnus: * et excélsus super omnes pópulos.

Confiteántur nómini tuo magno: quóniam terríbile, et sanctum est: * et honor regis judícium díligit.

Tu parásti directiónes: * judícium et justítiam in Jacob tu fecísti.

Exaltáte Dóminum, Deum nostrum, et adoráte scabéllum pedum ejus: * quóniam sanctum est.

Móyses et Aaron in sacerdótibus ejus: * et Sámuel inter eos, qui ínvocant nomen ejus:

Invocábant Dóminum, et ipse exaudiébat eos: * in colúmna nubis loquebátur ad eos.

Custodiébant testimónia ejus, * et præcéptum quod dedit illis.

Dómine, Deus noster, tu exaudiébas eos: * Deus, tu propítius fuísti eis, et ulcíscens in omnes adinventiónes eórum.

Exaltáte Dóminum, Deum nostrum, et adoráte in monte sancto ejus: * quóniam sanctus Dóminus, Deus noster.

\rubric{Secus vero, tribus Psalmis explétis, statim dícitur Ant.}

\antlabel{Ant.} Ne discédas a me, Dómine: quóniam tribulátio próxima est: quóniam non est qui ádjuvet.

\rubric{In Officio Festi vel Octávæ:}

\capitulumlabel{Capitulum}{1 Tim. 1, 17}

\lettrine[lines=2]{R}{EGI} sæculórum immortáli et invisíbili, soli Deo honor et glória in sæcula sæculórum. Amen.

\rubric{In ómnibus autem Fériis et in Vigíliis commúnibus:}

\capitulumlabel{Capitulum}{Zach. 8, 19}

\lettrine[lines=2]{P}{ACEM} et veritátem dilígite, ait Dóminus omnípotens.

\R\ br. Christe, Fili Dei vivi, * Miserére nobis. Christe. \V\ \textit{Qui sedes ad déxteram Patris.} Miserére nobis. Glória Patri. Christe. \V\ Exsúrge, Christe, ádjuva nos. \R\ Et líbera nos propter nomen tuum.

\variantheading{Preces dominicales et feriales (si dicendæ sint)}

KYRIE, eléison. Christe, eléison. Kýrie, eléison.

Pater noster \rubric{secreto usque ad} \V\ Et ne nos indúcas in tentatiónem. \R\ Sed líbera nos a malo.

Credo in Deum \rubric{secreto usque ad} \V\ Carnis resurrectiónem. \R\ Vitam ætérnam. Amen. \V\ Et ego ad te, Dómine, clamávi. \R\ Et mane orátio mea prævéniet te. \V\ Repleátur os meum laude. \R\ Ut cantem glóriam tuam, tota die magnitúdinem tuam. \V\ Dómine, avérte fáciem tuam a peccátis meis. \R\ Et omnes iniquitátes meas dele. \V\ Cor mundum crea in me, Deus. \R\ Et spíritum rectum ínnova in viscéribus meis. \V\ Ne projícias me a fácie tua. \R\ Et spíritum sanctum tuum ne áuferas a me. \V\ Redde mihi lætítiam salutáris tui. \R\ Et spíritu principáli confírma me.

\variantheading{In Précibus feriálibus tantum}

\V\ Eripe me, Dómine, ab hómine malo. \R\ A viro iníquo éripe me. \V\ Eripe me de inimícis meis, Deus meus. \R\ Et ab insurgéntibus in me líbera me. \V\ Eripe me de operántibus iniquitátem. \R\ Et de viris sánguinum salva me. \V\ Sic psalmum dicam nómini tuo in sæculum sæculi. \R\ Ut reddam vota mea de die in diem. \V\ Exáudi nos, Deus, salutáris noster. \R\ Spes ómnium fínium terræ, et in mari longe. \V\ Deus, in adjutórium meum inténde. \R\ Dómine, ad adjuvándum me festína. \V\ Sanctus Deus, Sanctus fortis, Sanctus immortális. \R\ Miserére nobis. \V\ Bénedic, ánima mea, Dómino. \R\ Et ómnia, quæ intra me sunt, nómini sancto ejus. \V\ Bénedic, ánima mea, Dómino. \R\ Et noli oblivísci omnes retributiónes ejus. \V\ Qui propitiátur ómnibus iniquitátibus tuis. \R\ Qui sanat omnes infirmitátes tuas. \V\ Qui rédimit de intéritu vitam tuam. \R\ Qui corónat te in misericórdia et miseratiónibus. \V\ Qui replet in bonis desidérium tuum. \R\ Renovábitur ut áquilæ juvéntus tua.

\variantheading{In utrisque Précibus}

\V\ Adjutórium nostrum in nómine Dómini. \R\ Qui fecit cælum et terram.

\formulalabel{Confíteor. Misereátur. Indulgéntiam.}

\V\ Dignáre, Dómine, die isto. \R\ Sine peccáto nos custodíre. \V\ Miserére nostri, Dómine. \R\ Miserére nostri. \V\ Fiat misericórdia tua, Dómine, super nos. \R\ Quemádmodum sperávimus in te. \V\ Dómine, exáudi oratiónem meam. \R\ Et clamor meus ad te véniat.

\V\ Dóminus vobíscum.

Orémus. \oratiolabel{Oratio}

\lettrine[lines=2]{D}{OMINE}, Deus omnípotens, qui ad princípium hujus diéi nos perveníre fecísti: tua nos hódie salva virtúte; ut in hac die ad nullum declinémus peccátum, sed semper ad tuam justítiam faciéndam nostra procédant elóquia, dirigántur cogitatiónes et ópera. Per Dóminum. \R\ Amen.

\V\ Dóminus vobíscum. \V\ Benedicámus Dómino.

\rubric{In Choro Martyrológium.}

\V\ Pretiósa in conspéctu Dómini. \R\ Mors Sanctórum ejus.

\rubric{Et absolúte, sine} Orémus:

\oratiolabel{Oratio}

\lettrine[lines=2]{S}{ANCTA} María, et omnes Sancti intercédant pro nobis ad Dóminum, ut nos mereámur ab eo adjuvári et salvári, qui vivit et regnat in sæcula sæculórum. \R\ Amen.

\V\ Deus, in adjutórium meum inténde. \R\ Dómine, ad adjuvándum me festína. \rubric{Et dícitur ter; último additur:} Glória Patri. Sicut erat.

Kýrie, eléison. Christe, eléison. Kýrie, eléison.

Pater noster \rubric{secreto usque ad} \V\ Et ne nos indúcas in tentatiónem. \R\ Sed líbera nos a malo. \V\ Réspice in servos tuos, Dómine, et in ópera tua, et dírige fílios eórum. \R\ Et sit splendor Dómini, Dei nostri, super nos, et ópera mánuum nostrárum dírige super nos, et opus mánuum nostrárum dírige. \V\ Glória Patri. \R\ Sicut erat.

Orémus. \oratiolabel{Oratio}

\lettrine[lines=2]{D}{IRIGERE} et sanctificáre, régere et gubernáre dignáre, Dómine Deus, Rex cæli et terræ, hódie corda et córpora nostra, sensus, sermónes et actus nostros in lege tua, et in opéribus mandatórum tuórum: ut hic, et in ætérnum, te auxiliánte, salvi et líberi esse mereámur, Salvátor mundi: Qui vivis et regnas in sæcula sæculórum. \R\ Amen.

\V\ Jube, domne.

\benedlabel{Bened.} Dies et actus nostros in sua pace dispónat Dóminus omnípotens. \R\ Amen.

\variantheading{In Officio feriali}

\lettrine[lines=2]{D}{OMINUS} autem dírigat corda et córpora nostra in caritáte Dei, et patiéntia Christi. Tu autem.

\variantheading{E Proprio Sanctorum}

\variantheading{Pro Oct. Omn. Sanctorum}

\lettrine[lines=2]{B}{ENEDÍCTIO} et cláritas et sapiéntia et gratiárum áctio, honor, virtus et fortitúdo Deo nostro in sǽcula sæculórum. Amen.

\variantheading{E Communi Sanctorum}

\variantheading{Pro Apostolis}

\lettrine[lines=2]{I}{BANT} Apóstoli gaudéntes a conspéctu concílii, quóniam digni hábiti sunt pro nómine Jesu contuméliam pati.

\variantheading{Pro Martyre}

\lettrine[lines=2]{J}{USTUS} cor suum trádidit ad vigilándum dilúculo ad Dóminum, qui fecit illum, et in conspéctu Altíssimi deprecábitur.

\variantheading{Pro pluribus Martyribus}

\lettrine[lines=2]{F}{ULGEBUNT} justi, et tamquam scintíllæ in arundinéto discúrrent. Judicábunt natiónes, et dominabúntur pópulis: et regnábit Dóminus illórum in perpétuum.

\variantheading{Pro Conf. Pont.}

\lettrine[lines=2]{F}{UNGI} sacerdótio, et habére laudem in nómine ipsíus, et offérre illi incénsum dignum in odórem suavitátis.

\variantheading{Pro Conf. non Pont.}

\lettrine[lines=2]{J}{USTUM} dedúxit Dóminus per vias rectas, et osténdit illi regnum Dei, et dedit illi sciéntiam sanctórum; honestávit illum in labóribus, et complévit labóres illíus.

\variantheading{Pro Virgine Martyre}

\lettrine[lines=2]{D}{OMINE}, Deus meus, exaltásti super terram habitatiónem meam, et pro morte defluénte deprecáta sum. Invocávi Dóminum, Patrem Dómini mei, ut non derelínquat me in die tribulatiónis meæ, et in témpore superbórum sine adjutório.

\variantheading{Pro Virgine non Martyre}

\lettrine[lines=2]{O}{} \textsc{quam} pulchra est casta generátio cum claritáte! immortális est enim memória illíus: quóniam apud Deum nota est, et apud hómines.

\variantheading{Pro Martyre tantum}

\lettrine[lines=2]{L}{AUDABIT} usque ad mortem ánima mea Dóminum, quóniam éruis sustinéntes te, et líberas eos de manu angústiæ, Dómine, Deus noster.

\variantheading{Pro nec Virg. nec Mart.}

\lettrine[lines=2]{M}{ULTAE} fíliæ congregavérunt divítias: tu supergréssa es univérsas. Fallax grátia, et vana est pulchritúdo: múlier timens Dóminum, ipsa laudábitur.

\variantheading{Pro Dedicatione Ecclesiæ}

\lettrine[lines=2]{E}{T} abstérget Deus omnem lácrimam ab óculis eórum: et mors ultra non erit, neque luctus, neque clamor, neque dolor erit ultra, quia prima abiérunt. Et dixit qui sedébat in throno: Ecce, nova fácio ómnia.

\variantheading{In Officiis B. Mariæ V.}

\lettrine[lines=2]{I}{N} platéis sicut cinnamómum et bálsamum aromatízans odórem dedi: quasi myrrha elécta dedi suavitátem odóris.

\V\ Adjutórium nostrum in nómine Dómini. \R\ Qui fecit cælum et terram.

\V\ Benedícite. \R\ Deus.

\benedictiolabel{Benedictio} Dóminus nos benedícat, et ab omni malo deféndat, et ad vitam perdúcat ætérnam. Et fidélium ánimæ per misericórdiam Dei requiéscant in pace. \R\ Amen.

\rubric{Deinde dícitur tantum} Pater noster \rubric{secréto, nisi sequátur ália Hora.}

% ===== END OF BATCH 33 -- Covers autumnalis.md lines 10778-11047,
% completing "Ad Primam" for Feria Sexta in full: the hymn "Jam
% lucis orto sídere," Psalms 21/i-iii, the conditional Psalm 98
% insertion, both Capitula, the brief Responsory, the full Preces
% apparatus, the closing Orationes, and the Lectio-brevis apparatus
% (identical category list to the previous four ferias). Several
% more psalm incipits carry the partial-caps OCR defect ("DEus,
% Deus meus," "CIrcumdedérunt," "DOminus regnávit," "PAcem"
% [Capitulum]) and are silently normalized to full caps. The
% document continues with "Ad Tertiam" (autumnalis.md line 11050).
% =====
\hourheading{Ad Tertiam}

Pater noster et Ave.

\V\ Deus, in adjutórium.

\hymnusheading{Hymnus}

\lettrine[lines=2]{N}{UNC}, Sancte, nobis, Spíritus,\\
Unum Patri cum Fílio,\\
Dignáre promptus íngeri\\
Nostro refúsus péctori.

Os, lingua, mens, sensus, vigor\\
Confessiónem pérsonent,\\
Flamméscat igne cáritas,\\
Accéndat ardor próximos.

Præsta, Pater piíssime,\\
Patríque compar Unice,\\
Cum Spíritu Paráclito\\
Regnans per omne sæculum.\\
Amen.

\antlabel{Ant.} Excita, Dómine.

\psalmheading{Psalmus 79, i}

\lettrine[lines=2]{Q}{UI} regis Israël, inténde: * qui dedúcis velut ovem Joseph.

Qui sedes super Chérubim, * manifestáre coram Ephraïm, Bénjamin, et Manásse.

Excita poténtiam tuam, et veni, * ut salvos fácias nos.

Deus, convérte nos: * et osténde fáciem tuam, et salvi érimus.

Dómine, Deus virtútum, * quoúsque irascéris super oratiónem servi tui?

Cibábis nos pane lacrimárum: * et potum dabis nobis in lácrimis in ménsura?

Posuísti nos in contradictiónem vicínis nostris: * et inimíci nostri subsannavérunt nos.

Deus virtútum, convérte nos: * et osténde fáciem tuam, et salvi érimus.

\psalmheading{Psalmus 79, ii}

\lettrine[lines=2]{V}{INEAM} de Ægýpto transtulísti: * ejecísti gentes, et plantásti eam.

Dux itíneris fuísti in conspéctu ejus: * plantásti radíces ejus, et implévit terram.

Opéruit montes umbra ejus: * et arbústa ejus cedros Dei.

Exténdit pálmites suos usque ad mare: * et usque ad flumen propágines ejus.

Ut quid destruxísti macériam ejus: * et vindémiant eam omnes, qui prætergrediúntur viam?

Exterminávit eam aper de silva: * et singuláris ferus depástus est eam.

Deus virtútum, convértere: * réspice de cælo, et vide, et vísita víneam istam.

Et pérfice eam, quam plantávit déxtera tua: * et super fílium hóminis, quem confirmásti tibi.

Incénsa igni, et suffóssa * ab increpatióne vultus tui períbunt.

Fiat manus tua super virum déxteræ tuæ: * et super fílium hóminis, quem confirmásti tibi.

Et non discédimus a te, vivificábis nos: * et nomen tuum invocábimus.

Dómine, Deus virtútum, convérte nos: * et osténde fáciem tuam, et salvi érimus.

\psalmheading{Psalmus 81}

\lettrine[lines=2]{D}{EUS} stetit in synagóga deórum: * in médio autem deos dijúdicat.

Usquequo judicátis iniquitátem: * et fácies peccatórum súmitis?

Judicáte egéno, et pupíllo: * húmilem, et páuperem justificáte.

Erípite páuperem: * et egénum de manu peccatóris liberáte.

Nesciérunt, neque intellexérunt, in ténebris ámbulant: * movebúntur ómnia fundaménta terræ.

Ego dixi: Dii estis, * et fílii Excélsi omnes.

Vos autem sicut hómines moriémini: * et sicut unus de princípibus cadétis.

Surge, Deus, júdica terram: * quóniam tu hereditábis in ómnibus géntibus.

\antlabel{Ant.} Excita, Dómine, poténtiam tuam, ut salvos fácias nos.

\variantheading{In Officio feriali}

\capitulumlabel{Capitulum}{Jer. 17, 14}

\lettrine[lines=2]{S}{ANA} me, Dómine, et sanábor: salvum me fac, et salvus ero: quóniam laus mea tu es.

\R\ br. Sana ánimam meam, * Quia peccávi tibi. Sana. \V\ Ego dixi: Dómine, miserére mei. Quia. Glória Patri. Sana.

\V\ Adjútor meus esto, ne derelínquas me. \R\ Neque despícias me, Deus, salutáris meus.

\variantheading{In Vigiliis communibus:}

\variantheading{Preces feriales}

KYRIE, eléison. Christe, eléison. Kýrie, eléison.

Pater noster \rubric{secreto usque ad} \V\ Et ne nos indúcas in tentatiónem. \R\ Sed líbera nos a malo. \V\ Dómine, Deus virtútum, convérte nos. \R\ Et osténde fáciem tuam, et salvi érimus. \V\ Exsúrge, Christe, ádjuva nos. \R\ Et líbera nos propter nomen tuum. \V\ Dómine, exáudi oratiónem meam. \R\ Et clamor meus ad te véniat.

\rubric{Réliqua ut in Ordinário. 19}

\hourheading{Ad Sextam}

Pater noster et Ave.

\V\ Deus, in adjutórium.

\hymnusheading{Hymnus}

\lettrine[lines=2]{R}{ECTOR} potens, verax Deus,\\
Qui témperas rerum vices,\\
Splendóre mane illúminas,\\
Et ígnibus merídiem:

Exstíngue flammas lítium,\\
Aufer calórem nóxium,\\
Confer salútem córporum,\\
Verámque pacem córdium.

Præsta, Pater piíssime,\\
Patríque compar Unice,\\
Cum Spíritu Paráclito\\
Regnans per omne sæculum.\\
Amen.

\antlabel{Ant.} Beáti, qui hábitant.

\psalmheading{Psalmus 83, i}

\lettrine[lines=2]{Q}{UAM} dilécta tabernácula tua, Dómine virtútum: * concupíscit, et déficit ánima mea in átria Dómini.

Cor meum, et caro mea * exsultavérunt in Deum vivum.

Etenim passer invénit sibi domum: * et turtur nidum sibi, ubi ponat pullos suos.

Altária tua, Dómine virtútum: * Rex meus, et Deus meus.

Beáti, qui hábitant in domo tua, Dómine: * in sæcula sæculórum laudábunt te.

Beátus vir, cujus est auxílium abs te: * ascensiónes in corde suo dispósuit in valle lacrimárum in loco, quem pósuit.

Etenim benedictiónem dabit legislátor, ibunt de virtúte in virtútem: * vidébitur Deus deórum in Sion.

\psalmheading{Psalmus 83, ii}

\lettrine[lines=2]{D}{OMINE}, Deus virtútum, exáudi oratiónem meam: * áuribus pércipe, Deus Jacob.

Protéctor noster, áspice, Deus: * et réspice in fáciem Christi tui:

Quia mélior est dies una in átriis tuis, * super míllia.

Elégi abjéctus esse in domo Dei mei: * magis quam habitáre in tabernáculis peccatórum.

Quia misericórdiam, et veritátem díligit Deus: * grátiam et glóriam dabit Dóminus.

Non privábit bonis eos, qui ámbulant in innocéntia: * Dómine virtútum, beátus homo, qui sperat in te.

\psalmheading{Psalmus 86}

\lettrine[lines=2]{F}{UNDAMENTA} ejus in móntibus sanctis: * díligit Dóminus portas Sion super ómnia tabernácula Jacob.

Gloriósa dicta sunt de te, * cívitas Dei.

Memor ero Rahab, et Babylónis * sciéntium me.

Ecce, alienígenæ, et Tyrus, et pópulus Æthíopum, * hi fuérunt illic.

Numquid Sion dicet: Homo, et homo natus est in ea: * et ipse fundávit eam Altíssimus?

Dóminus narrábit in scriptúris populórum, et príncipum: * horum, qui fuérunt in ea.

Sicut lætántium ómnium * habitátio est in te.

\antlabel{Ant.} Beáti, qui hábitant in domo tua, Dómine.

\variantheading{In Officio feriali}

\capitulumlabel{Capitulum}{Rom. 13, 8}

\lettrine[lines=2]{N}{EMINI} quidquam debeátis, nisi ut ínvicem diligátis: qui enim díligit próximum, legem implévit.

\R\ br. Benedícam Dóminum * In omni témpore. Benedícam. \V\ Semper laus ejus in ore meo. In. Glória Patri. Benedícam. \V\ Dóminus regit me, et nihil mihi déerit. \R\ In loco páscuæ ibi me collocávit.

\variantheading{In Vigiliis communibus:}

\variantheading{Preces feriales}

KYRIE, eléison. Christe, eléison. Kýrie, eléison.

Pater noster \rubric{secreto usque ad} \V\ Et ne nos indúcas in tentatiónem. \R\ Sed líbera nos a malo. \V\ Dómine, Deus virtútum, convérte nos. \R\ Et osténde fáciem tuam, et salvi érimus. \V\ Exsúrge, Christe, ádjuva nos. \R\ Et líbera nos propter nomen tuum. \V\ Dómine, exáudi oratiónem meam. \R\ Et clamor meus ad te véniat.

\rubric{Réliqua ut in Ordinário. 20}

\hourheading{Ad Nonam}

Pater noster et Ave.

\V\ Deus, in adjutórium.

\hymnusheading{Hymnus}

\lettrine[lines=2]{R}{ERUM}, Deus, tenax vigor,\\
Immótus in te pérmanens,\\
Lucis diúrnæ témpora\\
Succéssibus detérminans:

Largíre lumen véspere,\\
Quo vita nusquam décidat,\\
Sed prǽmium mortis sacræ\\
Perénnis instet glória.

Præsta, Pater piíssime,\\
Patríque compar Unice,\\
Cum Spíritu Paráclito\\
Regnans per omne sæculum.\\
Amen.

\antlabel{Ant.} Misericórdia et véritas.

\psalmheading{Psalmus 88, i}

\lettrine[lines=2]{M}{ISERICORDIAS} Dómini in ætérnum cantábo.

In generatiónem et generatiónem * annuntiábo veritátem tuam in ore meo.

Quóniam dixísti: In ætérnum misericórdia ædificábitur in cælis: * præparábitur véritas tua in eis.

Dispósui testaméntum eléctis meis, * jurávi David, servo meo:

Usque in ætérnum præparábo semen tuum. * Et ædificábo in generatiónem et generatiónem sedem tuam.

Confitebúntur cæli mirabília tua, Dómine: * étenim veritátem tuam in ecclésia sanctórum.

Quóniam quis in núbibus æquábitur Dómino: * símilis erit Deo in fíliis Dei?

Deus, qui glorificátur in consílio sanctórum: * magnus et terríbilis super omnes qui in circúitu ejus sunt.

Dómine, Deus virtútum, quis símilis tibi? * potens es, Dómine, et véritas tua in circúitu tuo.

Tu domináris potestáti maris: * motum autem flúctuum ejus tu mítigas.

Tu humiliásti sicut vulnerátum, supérbum: * in bráchio virtútis tuæ dispersísti inimícos tuos.

Tui sunt cæli, et tua est terra, orbem terræ et plenitúdinem ejus tu fundásti: * aquilónem, et mare tu creásti.

Thabor et Hermon in nómine tuo exsultábunt: * tuum bráchium cum poténtia.

Firmétur manus tua, et exaltétur déxtera tua: * justítia et judícium præparátio sedis tuæ.

Misericórdia et véritas præcédent fáciem tuam: * beátus pópulus, qui scit jubilatiónem.

Dómine, in lúmine vultus tui ambulábunt, et in nómine tuo exsultábunt tota die: * et in justítia tua exaltabúntur.

Quóniam glória virtútis eórum tu es: * et in beneplácito tuo exaltábitur cornu nostrum.

Quia Dómini est assúmptio nostra, * et Sancti Israël, regis nostri.

\psalmheading{Psalmus 88, ii}

\lettrine[lines=2]{T}{UNC} locútus es in visióne sanctis tuis, et dixísti: * Pósui adjutórium in poténte: et exaltávi eléctum de plebe mea.

Invéni David, servum meum: * óleo sancto meo unxi eum.

Manus enim mea auxiliábitur ei: * et bráchium meum confortábit eum.

Nihil profíciet inimícus in eo, * et fílius iniquitátis non appónet nocére ei.

Et concídam a fácie ipsíus inimícos ejus: * et odiéntes eum in fugam convértam.

Et véritas mea, et misericórdia mea cum ipso: * et in nómine meo exaltábitur cornu ejus.

Et ponam in mari manum ejus: * et in flumínibus déxteram ejus.

Ipse invocábit me: Pater meus es tu: * Deus meus, et suscéptor salútis meæ.

Et ego primogénitum ponam illum * excélsum præ régibus terræ.

In ætérnum servábo illi misericórdiam meam: * et testaméntum meum fidéle ipsi.

Et ponam in sæculum sæculi semen ejus: * et thronum ejus sicut dies cæli.

Si autem derelíquerint fílii ejus legem meam: * et in judíciis meis non ambuláverint:

Si justítias meas profanáverint: * et mandáta mea non custodíerint:

Visitábo in virga iniquitátes eórum: * et in verbéribus peccáta eórum.

Misericórdiam autem meam non dispérgam ab eo: * neque nocébo in veritáte mea:

Neque profanábo testaméntum meum: * et quæ procédunt de lábiis meis, non fáciam írrita.

Semel jurávi in sancto meo: Si David méntiar: * semen ejus in ætérnum manébit.

Et thronus ejus sicut sol in conspéctu meo, * et sicut luna perfécta in ætérnum: et testis in cælo fidélis.

\psalmheading{Psalmus 88, iii}

\lettrine[lines=2]{T}{U} vero repulísti et despexísti: * distulísti Christum tuum.

Evertísti testaméntum servi tui: * profanásti in terra Sanctuárium ejus.

Destruxísti omnes sepes ejus: * posuísti firmaméntum ejus formídinem.

Diripuérunt eum omnes transeúntes viam: * factus est oppróbrium vicínis suis.

Exaltásti déxteram deprimméntium eum: * lætificásti omnes inimícos ejus.

Avertísti adjutórium gládii ejus: * et non es auxiliátus ei in bello.

Destruxísti eum ab emundatióne: * et sedem ejus in terram collisísti.

Minorásti dies témporis ejus: * perfudísti eum confusióne.

Usquequo, Dómine, avértis in finem: * exardéscet sicut ignis ira tua?

Memoráre quæ mea substántia: * numquid enim vane constituísti omnes fílios hóminum?

Quis est homo, qui vivet, et non vidébit mortem: * éruet ánimam suam de manu ínferi?

Ubi sunt misericórdiæ tuæ antíquæ, Dómine, * sicut jurásti David in veritáte tua?

Memor esto, Dómine, oppróbrii servórum tuórum (quod contínui in sinu meo) multárum géntium.

Quod exprobravérunt inimíci tui, Dómine, * quod exprobravérunt commutatiónem Christi tui.

Benedíctus Dóminus in ætérnum: * fiat, fiat.

\antlabel{Ant.} Misericórdia et véritas præcédent fáciem tuam, Dómine.

\variantheading{In Officio feriali}

\capitulumlabel{Capit.}{1 Petri 1, 17-19}

\lettrine[lines=2]{I}{N} timóre incolátus vestri témpore conversámini: sciéntes quod non corruptibílibus auro vel argénto redémpti estis, sed pretióso sánguine quasi Agni immaculáti Christi.

\R\ br. Rédime me, Dómine, * Et miserére mei. Rédime. \V\ Pes enim meus stetit in via recta. Et miserére mei. Glória Patri. Rédime. \V\ Ab occúltis meis munda me, Dómine. \R\ Et ab aliénis parce servo tuo.

\variantheading{In Vigiliis communibus:}

\variantheading{Preces feriales}

KYRIE, eléison. Christe, eléison. Kýrie, eléison.

Pater noster \rubric{secreto usque ad} \V\ Et ne nos indúcas in tentatiónem. \R\ Sed líbera nos a malo. \V\ Dómine, Deus virtútum, convérte nos. \R\ Et osténde fáciem tuam, et salvi érimus. \V\ Exsúrge, Christe, ádjuva nos. \R\ Et líbera nos propter nomen tuum. \V\ Dómine, exáudi oratiónem meam. \R\ Et clamor meus ad te véniat.

\rubric{Réliqua ut in Ordinário. 22}

% ===== END OF BATCH 34 -- Covers autumnalis.md lines 11050-11408,
% completing the Little Hours of Feria Sexta in full: "Ad Tertiam"
% (Psalms 79/i, 79/ii, 81), "Ad Sextam" (Psalms 83/i, 83/ii, 86),
% and "Ad Nonam" (Psalms 88/i-iii), each with its own hymn,
% Capitulum, brief Responsory, and "In Vigiliis communibus: Preces
% feriales" apparatus. Three further occurrences of the recurring
% stray OCR "\*" artifact in the hour-hymns are silently dropped.
% Several more partial-caps OCR defects ("QUi," "VIneam," "DEus
% stetit," "QUam dilécta," "DOmine, Deus virtútum," "FUndaménta,"
% "MIsericórdias," "TUnc locútus," "TU vero") are silently
% normalized to full caps. The document continues with "Ad
% Vesperas" (autumnalis.md line 11412).
% =====
\hourheading{Ad Vesperas}

Pater noster et Ave.

\V\ Deus, in adjutórium.

\antlabel{Ant.} Dómine, * \textdagger\ probásti me, et cognovísti me. \textdaggerdbl

\psalmheading{Psalmus 138, i}

\lettrine[lines=2]{D}{OMINE}, \textdagger\ probásti me, et cognovísti me: * \textdaggerdbl\ tu cognovísti sessiónem meam, et resurrectiónem meam.

Intellexísti cogitatiónes meas de longe: * sémitam meam, et funículum meum investigásti.

Et omnes vias meas prævidísti: * quia non est sermo in lingua mea.

Ecce, Dómine, tu cognovísti ómnia novíssima, et antíqua: * tu formásti me, et posuísti super me manum tuam.

Mirábilis facta est sciéntia tua ex me: * confortáta est, et non pótero ad eam.

Quo ibo a spíritu tuo? * et quo a fácie tua fúgiam?

Si ascéndero in cælum, tu illic es: * si descéndero in inférnum, ades.

Si súmpsero pennas meas dilúculo, * et habitávero in extrémis maris:

Etenim illuc manus tua dedúcet me: * et tenébit me déxtera tua.

Et dixi: Fórsitan ténebræ conculcábunt me: * et nox illuminátio mea in delíciis meis.

Quia ténebræ non obscurabúntur a te, et nox sicut dies illuminábitur: * sicut ténebræ ejus, ita et lumen ejus.

Quia tu possedísti renes meos: * suscepísti me de útero matris meæ.

\antlabel{Ant.} Dómine, probásti me, et cognovísti me.

\antlabel{Ant.} Mirabília ópera tua, * Dómine, et ánima mea cognóscit nimis.

\psalmheading{Psalmus 138, ii}

\lettrine[lines=2]{C}{ONFITEBOR} tibi quia terribíliter magnificátus es: * mirabília ópera tua, et ánima mea cognóscit nimis.

Non est occultátum os meum a te, quod fecísti in occúlto: * et substántia mea in inferióribus terræ.

Imperféctum meum vidérunt óculi tui, et in libro tuo omnes scribéntur: * dies formabúntur, et nemo in eis.

Mihi autem nimis honorificáti sunt amíci tui, Deus: * nimis confortátus est principátus eórum.

Dinumerábo eos, et super arénam multiplicabúntur: * exsurréxi, et adhuc sum tecum.

Si occíderis, Deus, peccatóres: * viri sánguinum, declináte a me:

Quia dícitis in cogitatióne: * Accípient in vanitáte civitátes tuas.

Nonne qui odérunt te, Dómine, óderam? * et super inimícos tuos tabescébam?

Perfécto ódio óderam illos: * et inimíci facti sunt mihi.

Proba me, Deus, et scito cor meum: * intérroga me, et cognósce sémitas meas.

Et vide, si via iniquitátis in me est: * et deduc me in via ætérna.

\antlabel{Ant.} Mirabília ópera tua, Dómine, et ánima mea cognóscit nimis.

\antlabel{Ant.} Ne derelínquas me, * Dómine, virtus salútis meæ.

\psalmheading{Psalmus 139}

\lettrine[lines=2]{E}{RIPE} me, Dómine, ab hómine malo: * a viro iníquo éripe me.

Qui cogitavérunt iniquitátes in corde: * tota die constituébant prǽlia.

Acuérunt línguas suas sicut serpéntis: * venénum áspidum sub lábiis eórum.

Custódi me, Dómine, de manu peccatóris: * et ab homínibus iníquis éripe me.

Qui cogitavérunt supplantáre gressus meos: * absondérunt supérbi láqueum mihi:

Et funes extendérunt in láqueum: * juxta iter scándalum posuérunt mihi.

Dixi Dómino: Deus meus es tu: * exáudi, Dómine, vocem deprecatiónis meæ.

Dómine, Dómine, virtus salútis meæ: * obumbrásti super caput meum in die belli.

Ne tradas me, Dómine, a desidério meo peccatóri: * cogitavérunt contra me, ne derelínquas me, ne forte exalténtur.

Caput circúitus eórum: * labor labiórum ipsórum opériet eos.

Cadent super eos carbónes, in ignem dejícies eos: * in misériis non subsístent.

Vir linguósus non dirigétur in terra: * virum injústum mala cápient in intéritu.

Cognóvi quia fáciet Dóminus judícium ínopis: * et vindíctam páuperum.

Verúmtamen justi confitebúntur nómini tuo: * et habitábunt recti cum vultu tuo.

\antlabel{Ant.} Ne derelínquas me, Dómine, virtus salútis meæ.

\antlabel{Ant.} Dómine, * \textdagger\ clamávi ad te, exáudi me. \textdaggerdbl

\psalmheading{Psalmus 140}

\lettrine[lines=2]{D}{OMINE}, \textdagger\ clamávi ad te, exáudi me: * \textdaggerdbl\ inténde voci meæ, cum clamávero ad te.

Dirigátur orátio mea sicut incénsum in conspéctu tuo: * elevátio mánuum meárum sacrifícium vespertínum.

Pone, Dómine, custódiam ori meo: * et óstium circumstántiæ lábiis meis.

Non declínes cor meum in verba malítiæ, * ad excusándas excusatiónes in peccátis.

Cum homínibus operántibus iniquitátem: * et non communicábo cum eléctis eórum.

Corrípiet me justus in misericórdia, et increpábit me: * óleum autem peccatóris non impínguet caput meum.

Quóniam adhuc et orátio mea in beneplácitis eórum: * absórpti sunt juncti petræ júdices eórum.

Audient verba mea quóniam potuérunt: * sicut crassitúdo terræ erúpta est super terram.

Dissipáta sunt ossa nostra secus inférnum: * quia ad te, Dómine, Dómine, óculi mei: in te sperávi, non áuferas ánimam meam.

Custódi me a láqueo, quem statuérunt mihi: * et a scándalis operántium iniquitátem.

Cadent in retiáculo ejus peccatóres: * singuláriter sum ego donec tránseam.

\antlabel{Ant.} Dómine, clamávi ad te, exáudi me.

\antlabel{Ant.} Educ de custódia * ánimam meam, Dómine, ad confiténdum nómini tuo.

\psalmheading{Psalmus 141}

\lettrine[lines=2]{V}{OCE} mea ad Dóminum clamávi: * voce mea ad Dóminum deprecátus sum:

Effúndo in conspéctu ejus oratiónem meam, * et tribulatiónem meam ante ipsum pronúntio.

In deficiéndo ex me spíritum meum, * et tu cognovísti sémitas meas.

In via hac, qua ambulábam, * abscondérunt láqueum mihi.

Considerábam ad déxteram, et vidébam: * et non erat qui cognósceret me.

Périit fuga a me, * et non est qui requírat ánimam meam.

Clamávi ad te, Dómine, * dixi: Tu es spes mea, pórtio mea in terra vivéntium.

Inténde ad deprecatiónem meam: * quia humiliátus sum nimis.

Líbera me a persequéntibus me: * quia confortáti sunt super me.

Educ de custódia ánimam meam ad confiténdum nómini tuo: * me exspéctant justi, donec retríbuas mihi.

\antlabel{Ant.} Educ de custódia ánimam meam, Dómine, ad confiténdum nómini tuo.

\variantheading{In Officio feriali}

\capitulumlabel{Capitulum}{2 Cor. 1, 3-4}

\lettrine[lines=2]{B}{ENEDICTUS} Deus, et Pater Dómini nostri Jesu Christi, Pater misericordiárum, et Deus totíus consolatiónis, qui consolátur nos in omni tribulatióne nostra.

\hymnusheading{Hymnus}

\lettrine[lines=2]{H}{OMINIS} supérne Cónditor,\\
Qui cuncta solus órdinans,\\
Humum jubes prodúcere\\
Reptántis et feræ genus:

Et magna rerum córpora,\\
Dictu jubéntis vívida,\\
Per témporum certas vices\\
Obtemperáre sérvulis:

Repélle, quod cupídinis\\
Ciénte vi nos ímpetit,\\
Aut móribus se súggerit,\\
Aut áctibus se intérserit.

Da gaudiórum præmia,\\
Da gratiárum múnera:\\
Dissólve litis víncula:\\
Astrínge pacis fœdera.

Præsta, Pater piíssime,\\
Patríque compar Unice,\\
Cum Spíritu Paráclito\\
Regnans per omne sæculum.\\
Amen.

\V\ Dirigátur, Dómine, orátio mea. \R\ Sicut incénsum in conspéctu tuo.

\antlabel{Ad Magnif. Ant.} Depósuit Dóminus * poténtes de sede, et exaltávit húmiles.

\rubric{Réliqua ut in Ordinário. 22}

\hourheading{Ad Completorium}

\V\ Jube, domne.

\benedlabel{Bened.} Noctem quiétam, et finem perféctum concédat nobis Dóminus omnípotens. \R\ Amen.

\lectioheading{Lectio br.}{1 Petri 5, 8-9}

\lettrine[lines=2]{F}{RATRES}: Sóbrii estóte, et vigiláte: quia adversárius vester diábolus tamquam leo rúgiens círcuit, quærens quem dévoret: cui resístite fortes in fide. Tu autem, Dómine, miserére nobis. \R\ Deo grátias.

\V\ Adjutórium nostrum in nómine Dómini. \R\ Qui fecit cælum et terram.

Pater noster \rubric{quod dícitur totum secréto.}

\formulalabel{Confíteor.} Misereátur. Indulgéntiam.

\V\ Convérte nos, Deus, salutáris noster. \R\ Et avérte iram tuam a nobis.

\V\ Deus, in adjutórium.

\antlabel{Ant.} Voce mea. \textdagger

\psalmheading{Psalmus 76, i}

\lettrine[lines=2]{V}{OCE} mea \textdagger\ ad Dóminum clamávi: * voce mea ad Deum, et inténdit mihi.

In die tribulatiónis meæ Deum exquisívi, mánibus meis nocte contra eum: * et non sum decéptus.

Rénuit consolári ánima mea, * memor fui Dei, et delectátus sum, et exercitátus sum: et defécit spíritus meus.

Anticipavérunt vigílias óculi mei: * turbátus sum, et non sum locútus.

Cogitávi dies antíquos: * et annos ætérnos in mente hábui.

Et meditátus sum nocte cum corde meo, * et exercitábar, et scopébam spíritum meum.

Numquid in ætérnum projíciet Deus: * aut non appónet ut complacítior sit adhuc?

Aut in finem misericórdiam suam abscíndet, * a generatióne in generatiónem?

Aut obliviscétur miseréri Deus? * aut continébit in ira sua misericórdias suas?

Et dixi: Nunc cœpi: * hæc mutátio déxteræ Excélsi.

Memor fui óperum Dómini: * quia memor ero ab inítio mirabílium tuórum.

Et meditábor in ómnibus opéribus tuis: * et in adinventiónibus tuis exercébor.

\psalmheading{Psalmus 76, ii}

\lettrine[lines=2]{D}{EUS}, in sancto via tua: * quis Deus magnus sicut Deus noster? tu es Deus qui facis mirabília.

Notam fecísti in pópulis virtútem tuam: * redemísti in bráchio tuo pópulum tuum, fílios Jacob et Joseph.

Vidérunt te aquæ, Deus, vidérunt te aquæ: * et timuérunt, et turbátæ sunt abýssi.

Multitúdo sónitus aquárum: * vocem dedérunt nubes.

Etenim sagíttæ tuæ tránseunt: * vox tonítrui tui in rota.

Illuxérunt coruscatiónes tuæ orbi terræ: * commóta est, et contrémuit terra.

In mari via tua, et sémitæ tuæ in aquis multis: * et vestígia tua non cognoscéntur.

Deduxísti sicut oves pópulum tuum, * in manu Móysi et Aaron.

\psalmheading{Psalmus 85}

\lettrine[lines=2]{I}{NCLINA}, Dómine, aurem tuam, et exáudi me: * quóniam inops, et pauper sum ego.

Custódi ánimam meam, quóniam sanctus sum: * salvum fac servum tuum, Deus meus, sperántem in te.

Miserére mei, Dómine, quóniam ad te clamávi tota die: * lætífica ánimam servi tui, quóniam ad te, Dómine, ánimam meam levávi.

Quóniam tu, Dómine, suávis, et mitis: * et multæ misericórdiæ ómnibus invocántibus te.

Auribus pércipe, Dómine, oratiónem meam: * et inténde voci deprecatiónis meæ.

In die tribulatiónis meæ clamávi ad te: * quia exaudísti me.

Non est símilis tui in diis, Dómine: * et non est secúndum ópera tua.

Omnes gentes quascúmque fecísti, vénient, et adorábunt coram te, Dómine: * et glorificábunt nomen tuum.

Quóniam magnus es tu, et fáciens mirabília: * tu es Deus solus.

Deduc me, Dómine, in via tua, et ingrédiar in veritáte tua: * lætétur cor meum ut tímeat nomen tuum.

Confitébor tibi, Dómine, Deus meus, in toto corde meo, * et glorificábo nomen tuum in ætérnum:

Quia misericórdia tua magna est super me: * et eruísti ánimam meam ex inférno inferióri.

Deus, iníqui insurrexérunt super me, et synagóga poténtium quæsiérunt ánimam meam: * et non proposuérunt te in conspéctu suo.

Et tu, Dómine, Deus miserátor et miséricors, * pátiens, et multæ misericórdiæ, et verax.

Réspice in me, et miserére mei, * da impérium tuum púero tuo: et salvum fac fílium ancíllæ tuæ.

Fac mecum signum in bonum, ut vídeant qui odérunt me, et confundántur: * quóniam tu, Dómine, adjuvísti me, et consolátus es me.

\antlabel{Ant.} Voce mea ad Dóminum clamávi: neque obliviscétur miseréri Deus.

\hymnusheading{Hymnus}

\lettrine[lines=2]{T}{E} lucis ante términum,\\
Rerum Creátor, póscimus,\\
Ut pro tua cleméntia\\
Sis præsul et custódia.

Procul recédant sómnia,\\
Et nóctium phantásmata,\\
Hostémque nostrum cómprime,\\
Ne polluántur córpora.

Præsta, Pater piíssime,\\
Patríque compar Unice,\\
Cum Spíritu Paráclito\\
Regnans per omne sæculum.\\
Amen.

\capitulumlabel{Capitulum}{Jer. 14, 9}

\lettrine[lines=2]{T}{U} autem in nobis es, Dómine, et nomen sanctum tuum invocátum est super nos: ne derelínquas nos, Dómine, Deus noster.

\R\ br. In manus tuas, Dómine, * Comméndo spíritum meum. In manus. \V\ Redemísti nos, Dómine, Deus veritátis. Comméndo. Glória Patri. In manus.

\V\ Custódi nos, Dómine, ut pupíllam óculi. \R\ Sub umbra alárum tuárum prótege nos.

\antlabel{Ant.} Salva nos.

\canticumheading{Canticum Simeonis}{Luc. 2, 29-32}

\lettrine[lines=2]{N}{UNC} dimíttis servum tuum, Dómine, * secúndum verbum tuum in pace:

Quia vidérunt óculi mei * salutáre tuum,

Quod parásti * ante fáciem ómnium populórum,

Lumen ad revelatiónem géntium, * et glóriam plebis tuæ Israël.

Glória Patri.

\antlabel{Ant.} Salva nos, Dómine, vigilántes, custódi nos dormiéntes; ut vigilémus cum Christo, et requiescámus in pace.

\variantheading{Preces (si dicendæ sint)}

KYRIE, eléison. Christe, eléison. Kýrie, eléison.

Pater noster \rubric{secreto usque ad} \V\ Et ne nos indúcas in tentatiónem. \R\ Sed líbera nos a malo.

Credo in Deum \rubric{secreto usque ad} \V\ Carnis resurrectiónem. \R\ Vitam ætérnam. Amen. \V\ Benedíctus es, Dómine, Deus patrum nostrórum. \R\ Et laudábilis et gloriósus in sæcula. \V\ Benedicámus Patrem et Fílium cum Sancto Spíritu. \R\ Laudémus, et superexaltémus eum in sæcula. \V\ Benedíctus es, Dómine, in firmaménto cæli. \R\ Et laudábilis, et gloriósus, et superexaltátus in sæcula. \V\ Benedícat et custódiat nos omnípotens et miséricors Dóminus. \R\ Amen. \V\ Dignáre, Dómine, nocte ista. \R\ Sine peccáto nos custodíre. \V\ Miserére nostri, Dómine. \R\ Miserére nostri. \V\ Fiat misericórdia tua, Dómine, super nos. \R\ Quemádmodum sperávimus in te. \V\ Dómine, exáudi oratiónem meam. \R\ Et clamor meus ad te véniat.

\V\ Dóminus vobíscum.

Orémus. \oratiolabel{Oratio}

\lettrine[lines=2]{V}{ISITA}, quæsumus, Dómine, habitatiónem istam, et omnes insídias inimíci ab ea longe repélle: Angeli tui sancti hábitent in ea, qui nos in pace custódiant; et benedíctio tua sit super nos semper. Per Dóminum nostrum. \R\ Amen.

\V\ Dóminus vobíscum.

\V\ Benedicámus Dómino.

\benedictiolabel{Benedictio} Benedícat et custódiat nos omnípotens et miséricors Dóminus, Pater, et Fílius, et Spíritus Sanctus. \R\ Amen.

\variantheading{Antiphona finalis B. M. V.}

\V\ Divínum auxílium.

Pater noster, Ave María et Credo.

% ===== END OF BATCH 35 -- Covers autumnalis.md lines 11412-11776,
% completing "Ad Vesperas" (Psalms 138/i, 138/ii, 139, 140, 141,
% hymn "Hominis supérne Cónditor") and "Ad Completorium" (Psalms
% 76/i, 76/ii, 85, hymn "Te lucis ante términum," Canticum
% Simeonis, Preces), and with them the ENTIRE "Feria Sexta" Office.
% Several more partial-caps OCR defects ("COnfitébor," "ERipe me,"
% "VOce mea" twice, "DEus, in sancto," "INclina, Dómine,"
% "BEnedíctus" [Capitulum], "HOminis supérne" [hymn]) are silently
% normalized to full caps. The closing cross-reference reverts to
% the abbreviated "Antiphona finalis B. M. V." but keeps the fuller
% "Pater noster, Ave María et Credo." seen since Feria Quinta; each
% transcribed exactly as its own occurrence gives it. The document
% continues with "Sabbato" (autumnalis.md line 11780), the final
% weekday of the Psalterium cycle.
% =====
\feriaheading{Sabbato}

\hourheading{Ad Matutinum}

Pater, Ave et Credo.

\V\ Dómine, lábia mea.

\V\ Deus, in adjutórium.

\variantheading{In Officio feriali}

\invitatlabel{Invit.} Pópulus Dómini, et oves páscuæ ejus: * Veníte, adorémus. \psalmlabel{Ps. 94} Veníte.

\hymnusheading{Hymnus}

\lettrine[lines=2]{S}{UMMAE} Parens cleméntiæ,\\
Mundi regis qui máchinam,\\
Uníus et substántiæ,\\
Trinúsque persónis Deus:

Nostros pius cum cánticis\\
Fletus benígne súscipe:\\
Ut corde puro sórdium\\
Te perfruámur lárgius.

Lumbos jecúrque mórbidum\\
Flammis adúre cóngruis,\\
Accíncti ut artus éxcubent,\\
Luxu remóto péssimo.

Quicúmque ut horas nóctium\\
Nunc concinéndo rúmpimus,\\
Ditémur omnes áffatim\\
Donis beátæ pátriæ.

Præsta, Pater piíssime,\\
Patríque compar Unice,\\
Cum Spíritu Paráclito\\
Regnans per omne sæculum.\\
Amen.

\nocturnoheading{In I Nocturno}

\antlabel{Ant.} Memor fuit in sæculum * testaménti sui Dóminus, Deus noster.

\psalmheading{Psalmus 104, i}

\lettrine[lines=2]{C}{ONFITEMINI} Dómino, et invocáte nomen ejus: * annuntiáte inter gentes ópera ejus.

Cantáte ei, et psállite ei: * narráte ómnia mirabília ejus.

Laudámini in nómine sancto ejus: * lætétur cor quæréntium Dóminum.

Quærite Dóminum, et confirmámini: * quærite fáciem ejus semper.

Mementóte mirabílium ejus, quæ fecit: * prodígia ejus, et judícia oris ejus.

Semen Abraham, servi ejus: * fílii Jacob, elécti ejus.

Ipse Dóminus Deus noster: * in univérsa terra judícia ejus.

Memor fuit in sæculum testaménti sui: * verbi, quod mandávit in mille generatiónes:

Quod dispósuit ad Abraham: * et juraménti sui ad Isaac:

Et státuit illud Jacob in præcéptum: * et Israël in testaméntum ætérnum:

Dicens: Tibi dabo terram Chánaan, * funículum hereditátis vestræ.

Cum essent número brevi, * paucíssimi et íncolæ ejus:

Et pertransiérunt de gente in gentem, * et de regno ad pópulum álterum.

Non relíquit hóminem nocére eis: * et corrípuit pro eis reges.

Nolíte tángere christos meos: * et in prophétis meis nolíte malignári.

\antlabel{Ant.} Memor fuit in sæculum testaménti sui Dóminus, Deus noster.

\antlabel{Ant.} Auxit Dóminus * pópulum suum: et firmávit eum super inimícos ejus.

\psalmheading{Psalmus 104, ii}

\lettrine[lines=2]{E}{T} vocávit famem super terram: * et omne firmaméntum panis contrívit.

Misit ante eos virum: * in servum venúmdatus est Joseph.

Humiliavérunt in compédibus pedes ejus, ferrum pertránsiit ánimam ejus, * donec veníret verbum ejus.

Elóquium Dómini inflammávit eum: * misit rex, et solvit eum; princeps populórum, et dimísit eum.

Constítuit eum dóminum domus suæ: * et príncipem omnis possessiónis suæ:

Ut erudíret príncipes ejus sicut semetípsum: * et senes ejus prudéntiam docéret.

Et intrávit Israël in Ægýptum: * et Jacob áccola fuit in terra Cham.

Et auxit pópulum suum veheménter: * et firmávit eum super inimícos ejus.

Convértit cor eórum ut odírent pópulum ejus: * et dolum fácerent in servos ejus.

Misit Móysen, servum suum: * Aaron, quem elégit ipsum.

Pósuit in eis verba signórum suórum, * et prodigiórum in terra Cham.

\antlabel{Ant.} Auxit Dóminus pópulum suum: et firmávit eum super inimícos ejus.

\antlabel{Ant.} Edúxit Deus * pópulum suum in exsultatióne, et eléctos suos in lætítia.

\psalmheading{Psalmus 104, iii}

\lettrine[lines=2]{M}{ISIT} ténebras, et obscurávit: * et non exacerbávit sermónes suos.

Convértit aquas eórum in sánguinem: * et occídit pisces eórum.

Edidit terra eórum ranas * in penetrálibus regum ipsórum.

Dixit, et venit cœnomýia: * et cínifes in ómnibus fínibus eórum.

Pósuit plúvias eórum grándinem: * ignem comburéntem in terra ipsórum.

Et percússit víneas eórum, et ficúlneas eórum: * et contrívit lignum fínium eórum.

Dixit, et venit locústa, et bruchus, * cujus non erat númerus:

Et comédit omne fœnum in terra eórum: * et comédit omnem fructum terræ eórum.

Et percússit omne primogénitum in terra eórum: * primítias omnis labóris eórum.

Et edúxit eos cum argénto et auro: * et non erat in tríbubus eórum infírmus.

Lætáta est Ægýptus in profectióne eórum: * quia incúbuit timor eórum super eos.

Expándit nubem in protectiónem eórum, * et ignem ut lucéret eis per noctem.

Petiérunt, et venit cotúrnix: * et pane cæli saturávit eos.

Dirúpit petram, et fluxérunt aquæ: * abiérunt in sicco flúmina;

Quóniam memor fuit verbi sancti sui, * quod hábuit ad Abraham, púerum suum.

Et edúxit pópulum suum in exsultatióne, * et eléctos suos in lætítia.

Et dedit illis regiónes géntium: * et labóres populórum possedérunt:

Ut custódiant justificatiónes ejus, * et legem ejus requírant.

\antlabel{Ant.} Edúxit Deus pópulum suum in exsultatióne, et eléctos suos in lætítia.

\variantheading{In Festis novem Lectionum:}

\V\ Memor fuit Dóminus verbi sancti tui. \R\ Quod locútus est ad Abraham, púerum suum.

\absolutiolabel{Absolutio.} Exáudi, Dómine.

\benedlabel{Bened.} 1 Benedictióne perpétua. 2 Unigénitus Dei Fílius. 3 Spíritus Sancti grátia.

\nocturnoheading{In II Nocturno}

\antlabel{Ant.} Salvávit eos Dóminus * propter nomen suum.

\psalmheading{Psalmus 105, i}

\lettrine[lines=2]{C}{ONFITEMINI} Dómino quóniam bonus: * quóniam in sæculum misericórdia ejus.

Quis loquétur poténtias Dómini, * audítas fáciet omnes laudes ejus?

Beáti, qui custódiunt judícium, * et fáciunt justítiam in omni témpore.

Meménto nostri, Dómine, in beneplácito pópuli tui: * vísita nos in salutári tuo:

Ad vidéndum in bonitáte electórum tuórum, ad lætándum in lætítia gentis tuæ: * ut laudéris cum hereditáte tua.

Peccávimus cum pátribus nostris: * injúste égimus, iniquitátem fécimus.

Patres nostri in Ægýpto non intellexérunt mirabília tua: * non fuérunt mémores multitúdinis misericórdiæ tuæ.

Et irritavérunt ascendéntes in mare, * Mare Rubrum.

Et salvávit eos propter nomen suum: * ut notam fáceret poténtiam suam.

Et incrépuit Mare Rubrum, et exsiccátum est: * et dedúxit eos in abýssis sicut in desérto.

Et salvávit eos de manu odiéntium: * et redémit eos de manu inimíci.

Et opéruit aqua tribulántes eos: * unus ex eis non remánsit.

Et credidérunt verbis ejus: * et laudavérunt laudem ejus.

Cito fecérunt, oblíti sunt óperum ejus: * et non sustinuérunt consílium ejus.

Et concupiérunt concupiscéntiam in desérto: * et tentavérunt Deum in inaquóso.

Et dedit eis petitiónem ipsórum: * et misit saturitátem in ánimas eórum.

\antlabel{Ant.} Salvávit eos Dóminus propter nomen suum.

\antlabel{Ant.} Oblíti sunt Deum, * qui salvávit eos.

\psalmheading{Psalmus 105, ii}

\lettrine[lines=2]{E}{T} irritavérunt Móysen in castris: * Aaron, sanctum Dómini.

Apérta est terra, et deglutívit Dathan: * et opéruit super congregatiónem Abíron.

Et exársit ignis in synagóga eórum: * flamma combússit peccatóres.

Et fecérunt vítulum in Horeb: * et adoravérunt scúlptile.

Et mutavérunt glóriam suam * in similitúdinem vítuli comedéntis fœnum.

Oblíti sunt Deum, qui salvávit eos, * qui fecit magnália in Ægýpto, mirabília in terra Cham: terribília in Mari Rubro.

Et dixit ut dispérderet eos: * si non Móyses, eléctus ejus, stetísset in confractióne in conspéctu ejus:

Ut avérteret iram ejus ne dispérderet eos: * et pro níhilo habuérunt terram desiderábilem:

Non credidérunt verbo ejus, et murmuravérunt in tabernáculis suis: * non exaudiérunt vocem Dómini.

Et elevávit manum suam super eos: * ut prostérneret eos in desérto:

Et ut dejíceret semen eórum in natiónibus: * et dispérgeret eos in regiónibus.

Et initiáti sunt Beélphegor: * et comedérunt sacrifícia mortuórum.

Et irritavérunt eum in adinventiónibus suis: * et multiplicáta est in eis ruína.

Et stetit Phínees, et placávit: * et cessávit quassátio.

Et reputátum est ei in justítiam, * in generatiónem et generatiónem usque in sempitérnum.

\antlabel{Ant.} Oblíti sunt Deum, qui salvávit eos.

\antlabel{Ant.} Cum tribularéntur, * vidit Dóminus: et audívit oratiónem eórum.

\psalmheading{Psalmus 105, iii}

\lettrine[lines=2]{E}{T} irritavérunt eum ad Aquas contradictiónis: * et vexátus est Móyses propter eos: quia exacerbavérunt spíritum ejus.

Et distínxit in lábiis suis: * non disperdidérunt gentes, quas dixit Dóminus illis.

Et commísti sunt inter gentes, et didicérunt ópera eórum: et serviérunt sculptílibus eórum: * et factum est illis in scándalum.

Et immolavérunt fílios suos, * et fílias suas dæmóniis.

Et effudérunt sánguinem innocéntem: * sánguinem filiórum suórum et filiárum suárum, quas sacrificavérunt sculptílibus Chánaan.

Et infécta est terra in sanguínibus, et contamináta est in opéribus eórum: * et fornicáti sunt in adinventiónibus suis.

Et irátus est furóre Dóminus in pópulum suum: * et abominátus est hereditátem suam.

Et trádidit eos in manus géntium: * et domináti sunt eórum qui odérunt eos.

Et tribulavérunt eos inimíci eórum, et humiliáti sunt sub mánibus eórum: * sæpe liberávit eos.

Ipsi autem exacerbavérunt eum in consílio suo: * et humiliáti sunt in iniquitátibus suis.

Et vidit, cum tribularéntur: * et audívit oratiónem eórum.

Et memor fuit testaménti sui: * et pœnítuit eum secúndum multitúdinem misericórdiæ suæ.

Et dedit eos in misericórdias * in conspéctu ómnium qui céperant eos.

Salvos nos fac, Dómine, Deus noster: * et cóngrega nos de natiónibus:

Ut confiteámur nómini sancto tuo: * et gloriémur in laude tua.

Benedíctus Dóminus Deus Israël, a sæculo et usque in sæculum: * et dicet omnis pópulus: Fiat, fiat.

\antlabel{Ant.} Cum tribularéntur, vidit Dóminus: et audívit oratiónem eórum.

\variantheading{In Officiis novem Lectionum:}

\V\ Salvos nos fac, Dómine, Deus noster. \R\ Ut confiteámur nómini sancto tuo.

\absolutiolabel{Absolutio.} Ipsíus píetas et.

\benedlabel{Bened.} 4 Deus Pater omnípotens. 5 Christus perpétuæ. 6 Ignem sui amóris.

\nocturnoheading{In III Nocturno}

\antlabel{Ant.} Clamavérunt ad Dóminum: * et de necessitátibus eórum liberávit eos.

\psalmheading{Psalmus 106, i}

\lettrine[lines=2]{C}{ONFITEMINI} Dómino quóniam bonus: * quóniam in sæculum misericórdia ejus.

Dicant qui redémpti sunt a Dómino, quos redémit de manu inimíci: * et de regiónibus congregávit eos:

A solis ortu, et occásu: * ab aquilóne, et mari.

Erravérunt in solitúdine in inaquóso: * viam civitátis habitáculi non invenérunt.

Esuriéntes, et sitiéntes: * ánima eórum in ipsis defécit.

Et clamavérunt ad Dóminum cum tribularéntur: * et de necessitátibus eórum erípuit eos.

Et dedúxit eos in viam rectam: * ut irent in civitátem habitatiónis.

Confiteántur Dómino misericórdiæ ejus: * et mirabília ejus fíliis hóminum.

Quia satiávit ánimam inánem: * et ánimam esuriéntem satiávit bonis.

Sedéntes in ténebris, et umbra mortis: * vinctos in mendicitáte et ferro.

Quia exacerbavérunt elóquia Dei: * et consílium Altíssimi irritavérunt.

Et humiliátum est in labóribus cor eórum: * infirmáti sunt, nec fuit qui adjuváret.

Et clamavérunt ad Dóminum cum tribularéntur: * et de necessitátibus eórum liberávit eos.

Et edúxit eos de ténebris, et umbra mortis: * et víncula eórum disrúpit.

\antlabel{Ant.} Clamavérunt ad Dóminum: et de necessitátibus eórum liberávit eos.

\antlabel{Ant.} Ipsi vidérunt * ópera Dei, et mirabília ejus.

\psalmheading{Psalmus 106, ii}

\lettrine[lines=2]{C}{ONFITEANTUR} Dómino misericórdiæ ejus: * et mirabília ejus fíliis hóminum.

Quia contrívit portas áereas: * et vectes férreos confrégit.

Suscépit eos de via iniquitátis eórum: * propter injustítias enim suas humiliáti sunt.

Omnem escam abomináta est ánima eórum: * et appropinquavérunt usque ad portas mortis.

Et clamavérunt ad Dóminum cum tribularéntur: * et de necessitátibus eórum liberávit eos.

Misit verbum suum, et sanávit eos: * et erípuit eos de interitiónibus eórum.

Confiteántur Dómino misericórdiæ ejus: * et mirabília ejus fíliis hóminum.

Et sacríficent sacrifícium laudis: * et annúntient ópera ejus in exsultatióne.

Qui descéndunt mare in návibus, * faciéntes operatiónem in aquis multis.

Ipsi vidérunt ópera Dómini, * et mirabília ejus in profúndo.

Dixit, et stetit spíritus procéllæ: * et exaltáti sunt fluctus ejus.

Ascéndunt usque ad cælos, et descéndunt usque ad abýssos: * ánima eórum in malis tabescébat.

Turbáti sunt, et moti sunt sicut ébrius: * et omnis sapiéntia eórum devoráta est.

Et clamavérunt ad Dóminum cum tribularéntur: * et de necessitátibus eórum edúxit eos.

Et státuit procéllam ejus in auram: * et siluérunt fluctus ejus.

Et lætáti sunt quia siluérunt: * et dedúxit eos in portum voluntátis eórum.

\antlabel{Ant.} Ipsi vidérunt ópera Dei, et mirabília ejus.

\antlabel{Ant.} Vidébunt recti, * et lætabúntur, et intélligent misericórdias Dómini.

\psalmheading{Psalmus 106, iii}

\lettrine[lines=2]{C}{ONFITEANTUR} Dómino misericórdiæ ejus: * et mirabília ejus fíliis hóminum.

Et exáltent eum in ecclésia plebis: * et in cáthedra seniórum laudent eum.

Pósuit flúmina in desértum: * et éxitus aquárum in sitim.

Terram fructíferam in salsúginem, * a malítia inhabitántium in ea.

Pósuit desértum in stagna aquárum: * et terram sine aqua in éxitus aquárum.

Et collocávit illic esuriéntes: * et constituérunt civitátem habitatiónis.

Et seminavérunt agros, et plantavérunt víneas: * et fecérunt fructum nativitátis.

Et benedíxit eis, et multiplicáti sunt nimis: * et juménta eórum non minorávit.

Et pauci facti sunt: * et vexáti sunt a tribulatióne malórum, et dolóre.

Effúsa est contémptio super príncipes: * et erráre fecit eos in ínvio, et non in via.

Et adjúvit páuperem de inópia: * et pósuit sicut oves famílias.

Vidébunt recti, et lætabúntur: * et omnis iníquitas oppilábit os suum.

Quis sápiens et custódiet hæc? * et intélliget misericórdias Dómini?

\antlabel{Ant.} Vidébunt recti, et lætabúntur, et intélligent misericórdias Dómini.

\variantheading{In feriali Officio in Officio sanctæ Mariæ in Sabbato et in Festis:}

\V\ Exáltent Dóminum in ecclésia plebis. \R\ Et in cáthedra seniórum laudent eum.

\variantheading{In Officio novem Lectionum}

\absolutiolabel{Absolutio.} A vínculis.

\benedlabel{Bened.} 7 Evangélica léctio. 8 \rubric{In Festis Domini:} Divínum auxílium. \rubric{In Festis B. M. V.:} Cujus festum cólimus, ipsa Virgo. \rubric{In Festis Sanctorum:} Cujus (\textit{vel} Quorum \textit{aut} Quarum) festum cólimus. 9 Ad societátem. (Per evangélica dicta.)

\variantheading{In Officio trium Lectionum}

\absolutiolabel{Absolutio.} A vínculis.

\variantheading{In Feriis cum Homilia:}

\benedlabel{Bened.} 1 Evangélica. 2 Divínum. 3 Ad societátem.

\variantheading{In Octava simplici:}

\benedlabel{Bened.} 1 Ille nos benedícat. 2 Cujus (\textit{vel} Quorum \textit{aut} Quarum). 3 Ad societátem.

\variantheading{In Officio sanctæ Mariæ in Sabbato}

\absolutiolabel{Absol.} Précibus et méritis beátæ Maríæ semper Vírginis et ómnium Sanctórum, perdúcat nos Dóminus ad regna cælórum. Amen.

\benedlabel{Bened.} 1 Nos cum prole pia benedícat Virgo María. Amen.

2 Ipsa Virgo vírginum intercédat pro nobis ad Dóminum. Amen.

3 Per Vírginem Matrem concédat nobis Dóminus salútem et pacem. Amen.

% ===== END OF BATCH 36 -- Covers autumnalis.md lines 11780-12200,
% opening "Sabbato" (Saturday) and completing its "Ad Matutinum" in
% full: the ferial Invitatorium "Pópulus Dómini, et oves páscuæ
% ejus," the hymn "Summæ Parens cleméntiæ," and the three
% Nocturns -- Psalm 104 in the first, Psalm 105 in the second,
% Psalm 106 in the third -- with their versicle, Absolutio, and
% Benedictiones apparatus. The ninth-lesson "Bened." line (line
% 12176) marks its three conditional clauses ("In Festis Domini:",
% "In Festis B. M. V.:", "In Festis Sanctorum:") as actual red
% rubric text in the source this time ({RUBRIC: ...}), unlike the
% plain black-text embedding of the identical clauses at every
% prior occurrence of this line this volume; this is rendered with
% \rubric{} on each of the three labels, preserving the genuine
% embedded-rubric distinction exactly as the source marks it here.
% Saturday Matins closes with a Sabbath-specific apparatus unique
% to this day: "In Officio sanctæ Mariæ in Sabbato," with its own
% full-sentence Absolutio ("Précibus et méritis beátæ Maríæ...")
% and a three-part Marian Benedictio given as three separate
% paragraphs in the source (each ending its own "Amen."), reproduced
% here with the same paragraph breaks. Several more partial-caps
% OCR defects ("SUmmæ," "COnfitémini"/"COnfiteántur" repeatedly,
% "ET irritavérunt" twice, "MIsit ténebras") are silently normalized
% to full caps. The document continues with "Ad Laudes I"
% (autumnalis.md line 12202).
% =====
\hourheading{Ad Laudes I}

\rubric{In Sabbatis extra Vigílias commúnes occurréntibus, in Festis novem Lectiónum, in ómnibus Officiis de Octáva non privilegiáta, atque in Officio sanctæ Maríæ in Sabbato.}

(Pater noster et Ave.)

\V\ Deus, in adjutórium.

\antlabel{Ant.} Fílii Sion * exsúltent in Rege suo.

\psalmheading{Psalmus 149}

\lettrine[lines=2]{C}{ANTATE} Dómino cánticum novum: * laus ejus in ecclésia sanctórum.

Lætétur Israël in eo, qui fecit eum: * et fílii Sion exsúltent in rege suo.

Laudent nomen ejus in choro: * in týmpano, et psaltério psallant ei:

Quia beneplácitum est Dómino in pópulo suo: * et exaltábit mansuétos in salútem.

Exsultábunt sancti in glória: * lætabúntur in cubílibus suis.

Exaltatiónes Dei in gútture eórum: * et gládii ancípites in mánibus eórum.

Ad faciéndam vindíctam in natiónibus: * increpatiónes in pópulis.

Ad alligándos reges eórum in compédibus: * et nóbiles eórum in mánicis férreis.

Ut fáciant in eis judícium conscríptum: * glória hæc est ómnibus sanctis ejus.

\antlabel{Ant.} Fílii Sion exsúltent in Rege suo.

\antlabel{Ant.} Quam magnificáta * sunt ópera tua, Dómine!

\psalmheading{Psalmus 91}

\lettrine[lines=2]{B}{ONUM} est confitéri Dómino: * et psállere nómini tuo, Altíssime.

Ad annuntiándum mane misericórdiam tuam: * et veritátem tuam per noctem.

In decachórdo, psaltério: * cum cántico, in cíthara.

Quia delectásti me, Dómine, in factúra tua: * et in opéribus mánuum tuárum exsultábo.

Quam magnificáta sunt ópera tua, Dómine! * nimis profúndæ factæ sunt cogitatiónes tuæ.

Vir insípiens non cognóscet: * et stultus non intélliget hæc.

Cum exórti fúerint peccatóres sicut fœnum: * et apparúerint omnes, qui operántur iniquitátem:

Ut intéreant in sæculum sæculi: * tu autem Altíssimus in ætérnum, Dómine.

Quóniam ecce inimíci tui, Dómine, quóniam ecce inimíci tui períbunt: * et dispergéntur omnes, qui operántur iniquitátem.

Et exaltábitur sicut unicórnis cornu meum: * et senéctus mea in misericórdia úberi.

Et despéxit óculus meus inimícos meos: * et in insurgéntibus in me malignántibus áudiet auris mea.

Justus ut palma florébit: * sicut cedrus Líbani multiplicábitur.

Plantáti in domo Dómini, * in átriis domus Dei nostri florébunt.

Adhuc multiplicabúntur in senécta úberi: * et patiéntes erunt, ut annúntient:

Quóniam rectus Dóminus Deus noster: * et non est iníquitas in eo.

\antlabel{Ant.} Quam magnificáta sunt ópera tua, Dómine!

\antlabel{Ant.} Lætábitur justus * in Dómino, et sperábit in eo.

\psalmheading{Psalmus 63}

\lettrine[lines=2]{E}{XAUDI}, Deus, oratiónem meam cum déprecor: * a timóre inimíci éripe ánimam meam.

Protexísti me a convéntu malignántium: * a multitúdine operántium iniquitátem.

Quia exacuérunt ut gládium linguas suas: * intendérunt arcum rem amáram, ut sagíttent in occúltis immaculátum.

Súbito sagittábunt eum, et non timébunt: * firmavérunt sibi sermónem nequam.

Narravérunt ut abscónderent láqueos: * dixérunt: Quis vidébit eos?

Scrutáti sunt iniquitátes: * defecérunt scrutántes scrutínio.

Accédet homo ad cor altum: * et exaltábitur Deus.

Sagíttæ parvulórum factæ sunt plagæ eórum: * et infirmátæ sunt contra eos linguæ eórum.

Conturbáti sunt omnes qui vidébant eos: * et tímuit omnis homo.

Et annuntiavérunt ópera Dei, * et facta ejus intellexérunt.

Lætábitur justus in Dómino, et sperábit in eo, * et laudabúntur omnes recti corde.

\antlabel{Ant.} Lætábitur justus in Dómino, et sperábit in eo.

\antlabel{Ant.} Osténde nobis, Dómine, * lumen miseratiónum tuárum.

\canticumheading{Canticum Ecclesiastici}{Eccli. 36, 1-16}

\lettrine[lines=2]{M}{ISERERE} nostri, Deus ómnium, et réspice nos, * et osténde nobis lucem miseratiónum tuárum:

Et immítte timórem tuum super gentes, * quæ non exquisiérunt te,

Ut cognóscant quia non est Deus nisi tu, * et enárrent magnália tua.

Alleva manum tuam super gentes aliénas, * ut vídeant poténtiam tuam.

Sicut enim in conspéctu eórum sanctificátus es in nobis, * sic in conspéctu nostro magnificáberis in eis,

Ut cognóscant te, sicut et nos cognóvimus, * quóniam non est Deus præter te, Dómine.

Innova signa, et immúta mirabília. * Glorífica manum, et bráchium déxtrum.

Excita furórem, et effúnde iram. * Tolle adversárium, et afflíge inimícum.

Festína tempus, et meménto finis, * ut enárrent mirabília tua.

In ira flammæ devorétur qui salvátur: * et qui péssimant plebem tuam, invéniant perditiónem.

Cóntere caput príncipum inimicórum, * dicéntium: Non est álius præter nos.

Cóngrega omnes tribus Jacob: ut cognóscant quia non est Deus nisi tu, * et enárrent magnália tua:

Et hereditábis eos, * sicut ab inítio.

Miserére plebi tuæ, super quam invocátum est nomen tuum: * et Israël, quem coæquásti primogénito tuo.

Miserére civitáti sanctificatiónis tuæ Jerúsalem, * civitáti requiéi tuæ.

Reple Sion inenarrabílibus verbis tuis, * et glória tua pópulum tuum.

\antlabel{Ant.} Osténde nobis, Dómine, lucem miseratiónum tuárum.

\antlabel{Ant.} Omnis spíritus * laudet Dóminum.

\psalmheading{Psalmus 150}

\lettrine[lines=2]{L}{AUDATE} Dóminum in sanctis ejus: * laudáte eum in firmaménto virtútis ejus.

Laudáte eum in virtútibus ejus: * laudáte eum secúndum multitúdinem magnitúdinis ejus.

Laudáte eum in sono tubæ: * laudáte eum in psaltério, et cíthara.

Laudáte eum in týmpano, et choro: * laudáte eum in chordis, et órgano.

Laudáte eum in cýmbalis benesonántibus: laudáte eum in cýmbalis jubilatiónis: * omnis spíritus laudet Dóminum.

\antlabel{Ant.} Omnis spíritus laudet Dóminum.

\rubric{In Officio Festi vel Octávæ ac sanctæ Maríæ in Sabbato Capítulum, Hymnus, Versus, Ant. ad Benedictus et Orátio ut in Próprio vel Communi.}

\rubric{Réliqua ut in Ordinário. 19}

% ===== END OF BATCH 37 -- Covers autumnalis.md lines 12202-12357,
% "Ad Laudes I" for Sabbato in full: Psalms 149, 91, 63, the
% Canticum Ecclesiastici ("Miserére nostri, Deus ómnium"), and
% Psalm 150. Unlike every other feria's Ad Laudes I this volume,
% Saturday's carries no ferial Capitulum/Hymnus of its own -- the
% source instead gives a standalone rubric directing that these,
% along with the Versus, Ant. ad Benedictus, and Oratio, be taken
% "ut in Próprio vel Communi" (from the Proper or Common of the
% occurring feast or the Officium of Our Lady on Saturday), since
% Saturday's Lauds is always completed from elsewhere rather than
% having its own fixed ferial conclusion. Several more psalm
% incipits carry the partial-caps OCR defect ("LAudáte," "MIserére
% nostri," "BOnum est") and are silently normalized to full caps.
% The document continues with "Ad Laudes II" (autumnalis.md line
% 12360).
% =====
\hourheading{Ad Laudes II}

\rubric{In Vigíliis commúnibus, quando fit Officium de Feria.}

(Pater noster et Ave.)

\V\ Deus, in adjutórium.

\antlabel{Ant.} Benígne fac, Dómine.

\psalmheading{Psalmus 50}

\lettrine[lines=2]{M}{ISERERE} mei, Deus, * secúndum magnam misericórdiam tuam.

Et secúndum multitúdinem miseratiónum tuárum, * dele iniquitátem meam.

Amplius lava me ab iniquitáte mea: * et a peccáto meo munda me.

Quóniam iniquitátem meam ego cognósco: * et peccátum meum contra me est semper.

Tibi soli peccávi, et malum coram te feci: * ut justificéris in sermónibus tuis, et vincas cum judicáris.

Ecce enim, in iniquitátibus concéptus sum: * et in peccátis concépit me mater mea.

Ecce enim, veritátem dilexísti: * incérta et occúlta sapiéntiæ tuæ manifestásti mihi.

Aspérges me hyssópo, et mundábor: * lavábis me, et super nivem dealbábor.

Audítui meo dabis gáudium et lætítiam: * et exsultábunt ossa humiliáta.

Avérte fáciem tuam a peccátis meis: * et omnes iniquitátes meas dele.

Cor mundum crea in me, Deus: * et spíritum rectum ínnova in viscéribus meis.

Ne projícias me a fácie tua: * et spíritum sanctum tuum ne áuferas a me.

Redde mihi lætítiam salutáris tui: * et spíritu principáli confírma me.

Docébo iníquos vias tuas: * et ímpii ad te converténtur.

Líbera me de sanguínibus, Deus, Deus salútis meæ: * et exsultábit lingua mea justítiam tuam.

Dómine, lábia mea apéries: * et os meum annuntiábit laudem tuam.

Quóniam si voluísses sacrifícium, dedíssem útique: * holocáustis non delectáberis.

Sacrifícium Deo spíritus contribulátus: * cor contrítum, et humiliátum, Deus, non despícies.

Benígne fac, Dómine, in bona voluntáte tua Sion: * ut ædificéntur muri Jerúsalem.

Tunc acceptábis sacrifícium justítiæ, oblatiónes, et holocáusta: * tunc impónent super altáre tuum vítulos.

\antlabel{Ant.} Benígne fac, Dómine, in bona voluntáte tua Sion.

\antlabel{Ant.} Rectus Dóminus.

\psalmheading{Psalmus 91}

\lettrine[lines=2]{B}{ONUM} est confitéri Dómino: * et psállere nómini tuo, Altíssime.

Ad annuntiándum mane misericórdiam tuam: * et veritátem tuam per noctem.

In decachórdo, psaltério: * cum cántico, in cíthara.

Quia delectásti me, Dómine, in factúra tua: * et in opéribus mánuum tuárum exsultábo.

Quam magnificáta sunt ópera tua, Dómine! * nimis profúndæ factæ sunt cogitatiónes tuæ.

Vir insípiens non cognóscet: * et stultus non intélliget hæc.

Cum exórti fúerint peccatóres sicut fœnum: * et apparúerint omnes, qui operántur iniquitátem:

Ut intéreant in sæculum sæculi: * tu autem Altíssimus in ætérnum, Dómine.

Quóniam ecce inimíci tui, Dómine, quóniam ecce inimíci tui períbunt: * et dispergéntur omnes, qui operántur iniquitátem.

Et exaltábitur sicut unicórnis cornu meum: * et senéctus mea in misericórdia úberi.

Et despéxit óculus meus inimícos meos: * et in insurgéntibus in me malignántibus áudiet auris mea.

Justus ut palma florébit: * sicut cedrus Líbani multiplicábitur.

Plantáti in domo Dómini, * in átriis domus Dei nostri florébunt.

Adhuc multiplicabúntur in senécta úberi: * et patiéntes erunt, ut annúntient:

Quóniam rectus Dóminus, Deus noster: * et non est iníquitas in eo.

\antlabel{Ant.} Rectus Dóminus, Deus noster: et non est iníquitas in eo.

\antlabel{Ant.} A timóre inimíci.

\psalmheading{Psalmus 63}

\lettrine[lines=2]{E}{XAUDI}, Deus, oratiónem meam cum déprecor: * a timóre inimíci éripe ánimam meam.

Protexísti me a convéntu malignántium: * a multitúdine operántium iniquitátem.

Quia exacuérunt ut gládium linguas suas: * intendérunt arcum rem amáram, ut sagíttent in occúltis immaculátum.

Súbito sagittábunt eum, et non timébunt: * firmavérunt sibi sermónem nequam.

Narravérunt ut abscónderent láqueos: * dixérunt: Quis vidébit eos?

Scrutáti sunt iniquitátes: * defecérunt scrutántes scrutínio.

Accédet homo ad cor altum: * et exaltábitur Deus.

Sagíttæ parvulórum factæ sunt plagæ eórum: * et infirmátæ sunt contra eos linguæ eórum.

Conturbáti sunt omnes qui vidébant eos: * et tímuit omnis homo.

Et annuntiavérunt ópera Dei, * et facta ejus intellexérunt.

Lætábitur justus in Dómino, et sperábit in eo, * et laudabúntur omnes recti corde.

\antlabel{Ant.} A timóre inimíci éripe, Dómine, ánimam meam.

\antlabel{Ant.} In servis suis.

\canticumheading{Canticum Moysis}{Deut. 32, 1-43}

\lettrine[lines=2]{A}{UDITE}, cæli, quæ loquor: * áudiat terra verba oris mei.

Concréscat ut plúvia doctrína mea, * fluat ut ros elóquium meum,

Quasi imber super herbam, et quasi stillæ super grámina. * Quia nomen Dómini invocábo:

Date magnificéntiam Deo nostro. * Dei perfécta sunt ópera, et omnes viæ ejus judícia:

Deus fidélis, et absque ulla iniquitáte, justus et rectus. * Peccavérunt ei, et non fílii ejus in sórdibus:

Generátio prava atque pervérsa. * Hæccine réddis Dómino, pópule stulte et insípiens?

Numquid non ipse est pater tuus, * qui possédit te, et fecit, et creávit te?

Meménto diérum antiquórum, * cógita generatiónes síngulas:

Intérroga patrem tuum, et annuntiábit tibi: * majóres tuos, et dicent tibi.

Quando dividébat Altíssimus gentes: * quando separábat fílios Adam,

Constítuit términos populórum * juxta númerum filiórum Israël.

Pars autem Dómini, pópulus ejus: * Jacob funículus hereditátis ejus.

Invénit eum in terra desérta, * in loco horróris, et vastæ solitúdinis:

Circumdúxit eum, et dócuit: * et custodívit quasi pupíllam óculi sui.

Sicut áquila próvocans ad volándum pullos suos, * et super eos vólitans,

Expándit alas suas, et assúmpsit eum, * atque portávit in húmeris suis.

Dóminus solus dux ejus fuit: * et non erat cum eo deus aliénus.

Constítuit eum super excélsam terram: * ut coméderet fructus agrórum,

Ut súgeret mel de petra, * oleúmque de saxo duríssimo.

Butýrum de arménto, et lac de óvibus * cum ádipe agnórum, et aríetum filiórum Basan:

Et hircos cum medúlla trítici, * et sánguinem uvæ bíberet meracíssimum.

Incrassátus est diléctus, et recalcitrávit: * incrassátus, impinguátus, dilatátus,

Derelíquit Deum, factórem suum, * et recéssit a Deo, salutári suo.

Provocavérunt eum in diis aliénis, * et in abominatiónibus ad iracúndiam concitavérunt.

Immolavérunt dæmóniis, et non Deo, * diis, quos ignorábant:

Novi recentésque venérunt, * quos non coluérunt patres eórum.

Deum qui te génuit dereliquísti, * et oblítus es Dómini, creatóris tui.

Vidit Dóminus, et ad iracúndiam concitátus est: * quia provocavérunt eum fílii sui et fíliæ.

Et ait: Abscóndam fáciem meam ab eis, * et considerábo novíssima eórum:

Generátio enim pervérsa est, * et infidéles fílii.

Ipsi me provocavérunt in eo, qui non erat Deus, * et irritavérunt in vanitátibus suis:

Et ego provocábo eos in eo, qui non est pópulus, * et in gente stulta irritábo illos.

Ignis succénsus est in furóre meo, * et ardébit usque ad inférni novíssima:

Devorabítque terram cum gérmine suo, * et móntium fundaménta combúret.

Congregábo super eos mala, * et sagíttas meas complébo in eis.

Consuméntur fame, * et devorábunt eos aves morsu amaríssimo:

Dentes bestiárum immíttam in eos, * cum furóre trahéntium super terram, atque serpéntium.

Foris vastábit eos gládius, et intus pavor, * júvenem simul ac vírginem, lacténtem cum hómine sene.

Dixi: Ubinam sunt? * cessáre fáciam ex homínibus memóriam eórum.

Sed propter iram inimicórum dístuli: * ne forte superbírent hostes eórum,

Et dícerent: * Manus nostra excélsa, et non

Dóminus, * fecit hæc ómnia.

Gens absque consílio est, et sine prudéntia. * Utinam sáperent, et intellígerent, ac novíssima providérent.

Quómodo persequátur unus mille, * et duo fugent decem míllia?

Nonne ídeo, quia Deus suus véndidit eos, * et Dóminus conclúsit illos?

Non enim est Deus noster ut dii eórum: * et inimíci nostri sunt júdices.

De vínea Sodomórum, vínea eórum, * et de suburbánis Gomórrhæ:

Uva eórum uva fellis, * et botri amaríssimi.

Fel dracónum vinum eórum, * et venénum áspidum insanábile.

Nonne hæc cóndita sunt apud me, * et signáta in thesáuris meis?

Mea est últio, et ego retríbuam in témpore, * ut labátur pes eórum:

Juxta est dies perditiónis, * et adésse festínant témpora.

Judicábit Dóminus pópulum suum, * et in servis suis miserébitur:

Vidébit quod infirmáta sit manus, * et clausi quoque defecérunt, residuíque consúmpti sunt.

Et dicet: Ubi sunt dii eórum, * in quibus habébant fidúciam?

De quorum víctimis comedébant ádipes, * et bibébant vinum libáminum:

Surgant, et opituléntur vobis, * et in necessitáte vos prótegant.

Vidéte quod ego sim solus, * et non sit álius Deus præter me:

Ego occídam, et ego vívere fáciam: percútiam, et ego sanábo, * et non est qui de manu mea possit erúere.

Levábo ad cælum manum meam, et dicam: * Vivo ego in ætérnum.

Si acúero ut fulgur gládium meum, * et arripúerit judícium manus mea:

Reddam ultiónem hóstibus meis, * et his qui odérunt me retríbuam.

Inebriábo sagíttas meas sánguine, * et gládius meus devorábit carnes,

De cruóre occisórum, * et de captivitáte, nudáti inimicórum cápitis.

Laudáte, gentes, pópulum ejus, * quia sánguinem servórum suórum ulciscétur:

Et vindíctam retríbuet in hostes eórum, * et propítius erit terræ pópuli sui.

\antlabel{Ant.} In servis suis miserébitur Dóminus: et propítius erit terræ pópuli sui.

\antlabel{Ant.} Laudáte. \textdagger

\psalmheading{Psalmus 150}

\lettrine[lines=2]{L}{AUDATE} \textdagger\ Dóminum in sanctis ejus: * laudáte eum in firmaménto virtútis ejus.

Laudáte eum in virtútibus ejus: * laudáte eum secúndum multitúdinem magnitúdinis ejus.

Laudáte eum in sono tubæ: * laudáte eum in psaltério, et cíthara.

Laudáte eum in týmpano, et choro: * laudáte eum in chordis, et órgano.

Laudáte eum in cýmbalis benesonántibus: laudáte eum in cýmbalis jubilatiónis: * omnis spíritus laudet Dóminum.

\antlabel{Ant.} Laudáte Dóminum secúndum multitúdinem magnitúdinis ejus.

\capitulumlabel{Capit.}{Rom. 13, 12-13}

\lettrine[lines=2]{N}{OX} præcéssit, dies autem appropinquávit. Abjiciámus ergo ópera tenebrárum, et induámur arma lucis. Sicut in die honéste ambulémus.

\hymnusheading{Hymnus}

\lettrine[lines=2]{A}{URORA} jam spargit polum:\\
Terris dies illábitur:\\
Lucis resúltat spículum:\\
Discédat omne lúbricum.

Phantásma noctis éxsulet:\\
Mentis reátus córruat:\\
Quidquid tenébris hórridum\\
Nox áttulit culpæ, cadat.

Ut mane, quod nos últimum\\
Hic deprecámur cérnui,\\
Cum luce nobis éffluat,\\
Hoc dum canóre cóncrepat.

Deo Patri sit glória,\\
Ejúsque soli Fílio,\\
Cum Spíritu Paráclito,\\
Nunc et per omne sæculum. Amen.

\V\ Repléti sumus mane misericórdia tua. \R\ Exsultávimus, et delectáti sumus.

\antlabel{Ad Bened. Ant.} Illúmina, Dómine, * sedéntes in ténebris et umbra mortis, et dírige pedes nostros in viam pacis.

\rubric{Expléto Cántico ac repetíta Antíphona, dicúntur sequéntes}

\variantheading{Preces feriales}

KYRIE, eléison. Christe, eléison. Kýrie, eléison.

Pater noster, \rubric{quod dícitur a solo Hebdomadário totum clara voce, usque ad} \V\ Et ne nos indúcas in tentatiónem. \R\ Sed líbera nos a malo. \V\ Ego dixi: Dómine, miserére mei. \R\ Sana ánimam meam, quia peccávi tibi. \V\ Convértere, Dómine, úsquequo? \R\ Et deprecábilis esto super servos tuos. \V\ Fiat misericórdia tua, Dómine, super nos. \R\ Quemádmodum sperávimus in te. \V\ Sacerdótes tui induántur justítiam. \R\ Et sancti tui exsúltent. \V\ Orémus pro beatíssimo Papa nostro N. \R\ Dóminus consérvet eum, et vivíficet eum, et beátum fáciat eum in terra, et non tradat eum in ánimam inimicórum ejus. \V\ Orémus et pro Antístite nostro N. \R\ Stet et pascat in fortitúdine tua, Dómine, in sublimitáte nóminis tui. \V\ Dómine, salvum fac regem. \R\ Et exáudi nos in die, qua invocavérimus te. \V\ Salvum fac pópulum tuum, Dómine, et bénedic hereditáti tuæ. \R\ Et rege eos, et extólle illos usque in ætérnum. \V\ Meménto Congregatiónis tuæ. \R\ Quam possedísti ab inítio. \V\ Fiat pax in virtúte tua. \R\ Et abundántia in túrribus tuis. \V\ Orémus pro benefactóribus nostris. \R\ Retribúere dignáre, Dómine, ómnibus, nobis bona faciéntibus propter nomen tuum, vitam ætérnam. Amen. \V\ Orémus pro fidélibus defúnctis. \R\ Réquiem ætérnam dona eis, Dómine, et lux perpétua lúceat eis. \V\ Requiéscant in pace. \R\ Amen. \V\ Pro frátribus nostris abséntibus. \R\ Salvos fac servos tuos, Deus meus, sperántes in te. \V\ Pro afflíctis et captívis. \R\ Líbera eos, Deus Israël, ex ómnibus tribulatiónibus suis. \V\ Mitte eis, Dómine, auxílium de sancto. \R\ Et de Sion tuére eos. \V\ Dómine, Deus virtútum, convérte nos. \R\ Et osténde fáciem tuam, et salvi érimus. \V\ Exsúrge, Christe, ádjuva nos. \R\ Et líbera nos propter nomen tuum. \V\ Dómine, exáudi oratiónem meam. \R\ Et clamor meus ad te véniat.

\rubric{Réliqua ut in Ordinário. 19}

% ===== END OF BATCH 38 -- Covers autumnalis.md lines 12360-12673,
% "Ad Laudes II" for Sabbato in full: Psalms 50, 91, 63, the
% Canticum Moysis ("Audíte, cæli, quæ loquor," the long canticle of
% Deuteronomy 32), and Psalm 150, with its own Capitulum "Nox
% præcéssit" and hymn "Auróra jam spargit polum" (unlike Ad Laudes
% I, which borrowed these from the Proper/Common), plus the "Preces
% feriales" litany. The source breaks one canticle verse oddly
% across two short lines ("Et dícerent: * Manus nostra excélsa, et
% non" / "Dóminus, * fecit hæc ómnia."), reproduced exactly as the
% source paragraphs it. The document continues with "Ad Primam"
% (autumnalis.md line 12676).
% =====
\hourheading{Ad Primam}

Pater, Ave et Credo.

\V\ Deus, in adjutórium.

\hymnusheading{Hymnus}

\lettrine[lines=2]{J}{AM} lucis orto sídere,\\
Deum precémur súpplices,\\
Ut in diúrnis áctibus\\
Nos servet a nocéntibus.

Linguam refrénans témperet,\\
Ne litis horror ínsonet:\\
Visum fovéndo cóntegat,\\
Ne vanitátes háuriat.

Sint pura cordis íntima,\\
Absístat et vecórdia:\\
Carnis terat supérbiam\\
Potus cibíque párcitas.

Ut, cum dies abscésserit\\
Noctémque sors redúxerit,\\
Mundi per abstinéntiam\\
Ipsi canámus glóriam.

Deo Patri sit glória,\\
Ejúsque soli Fílio,\\
Cum Spíritu Paráclito,\\
Nunc et per omne sæculum.\\
Amen.

\antlabel{Ant.} Exaltáre, Dómine.

\psalmheading{Psalmus 93, i}

\lettrine[lines=2]{D}{EUS} ultiónum Dóminus: * Deus ultiónum líbere egit.

Exaltáre, qui júdicas terram: * redde retributiónem supérbis.

Usquequo peccatóres, Dómine, * úsquequo peccatóres gloriabúntur:

Effabúntur, et loquéntur iniquitátem: * loquéntur omnes, qui operántur injustítiam?

Pópulum tuum, Dómine, humiliavérunt: * et hereditátem tuam vexavérunt.

Víduam, et ádvenam interfecérunt: * et pupíllos occidérunt.

Et dixérunt: Non vidébit Dóminus, * nec intélliget Deus Jacob.

Intellígite, insipiéntes in pópulo: * et stulti, áliquándo sápite.

Qui plantávit aurem, non áudiet? * aut qui finxit óculum, non consíderat?

Qui córripit gentes, non árguet: * qui docet hóminem sciéntiam?

Dóminus scit cogitatiónes hóminum, * quóniam vanæ sunt.

\psalmheading{Psalmus 93, ii}

\lettrine[lines=2]{B}{EATUS} homo, quem tu erudíeris, Dómine: * et de lege tua docúeris eum,

Ut mítiges ei a diébus malis: * donec fodiátur peccatóri fóvea.

Quia non repéllet Dóminus plebem suam: * et hereditátem suam non derelínquet.

Quoadúsque justítia convertátur in judícium: * et qui juxta illam omnes qui recto sunt corde.

Quis consúrget mihi advérsus malignántes? * aut quis stabit mecum advérsus operántes iniquitátem?

Nisi quia Dóminus adjúvit me: * paulo minus habitásset in inférno ánima mea.

Si dicébam: Motus est pes meus: * misericórdia tua, Dómine, adjuvábat me.

Secúndum multitúdinem dolórum meórum in corde meo: * consolatiónes tuæ lætificavérunt ánimam meam.

Numquid adhæret tibi sedes iniquitátis: * qui fingis labórem in præcépto?

Captábunt in ánimam justi: * et sánguinem innocéntem condemnábunt.

Et factus est mihi Dóminus in refúgium: * et Deus meus in adjutórium spei meæ.

Et reddet illis iniquitátem ipsórum: et in malítia eórum dispérdet eos: * dispérdet illos Dóminus, Deus noster.

\psalmheading{Psalmus 107}

\lettrine[lines=2]{P}{ARATUM} cor meum, Deus, parátum cor meum: * cantábo, et psallam in glória mea.

Exsúrge, glória mea, exsúrge, psaltérium et cíthara: * exsúrgam dilúculo.

Confitébor tibi in pópulis, Dómine: * et psallam tibi in natiónibus.

Quia magna est super cælos misericórdia tua: * et usque ad nubes véritas tua:

Exaltáre super cælos, Deus, et super omnem terram glória tua: * ut liberéntur dilécti tui.

Salvum fac déxtera tua, et exáudi me: * Deus locútus est in sancto suo:

Exsultábo, et dívidam Síchimam, * et convállem tabernaculórum dimétiar.

Meus est Gálaad, et meus est Manásses: * et Ephraïm suscéptio cápitis mei.

Juda rex meus: * Moab lebes spei meæ.

In Idumæam exténdam calceaméntum meum: * mihi alienígenæ amíci facti sunt.

Quis dedúcet me in civitátem munítam? * quis dedúcet me usque in Idumæam?

Nonne tu, Deus, qui repulísti nos, * et non exíbis, Deus, in virtútibus nostris?

Da nobis auxílium de tribulatióne: * quia vana salus hóminis.

In Deo faciémus virtútem: * et ipse ad níhilum dedúcet inimícos nostros.

\rubric{¶ Quando ad Laudes dictus fúerit Ps. 50 Miserére, hic subjúngitur sequens Ps. 149, qui in Láudibus locum non hábuit.}

\psalmheading{Psalmus 149}

\lettrine[lines=2]{C}{ANTATE} Dómino cánticum novum: * laus ejus in ecclésia sanctórum.

Lætétur Israël in eo, qui fecit eum: * et fílii Sion exsúltent in rege suo.

Laudent nomen ejus in choro: * in týmpano, et psaltério psallant ei:

Quia beneplácitum est Dómino in pópulo suo: * et exaltábit mansuétos in salútem.

Exsultábunt sancti in glória: * lætabúntur in cubílibus suis.

Exaltatiónes Dei in gútture eórum: * et gládii ancípites in mánibus eórum.

Ad faciéndam vindíctam in natiónibus: * increpatiónes in pópulis.

Ad alligándos reges eórum in compédibus: * et nóbiles eórum in mánicis férreis.

Ut fáciant in eis judícium conscríptum: * glória hæc est ómnibus sanctis ejus.

\rubric{Secus vero, tribus Psalmis explétis, statim dícitur Ant.}

\antlabel{Ant.} Exaltáre, Dómine, qui júdicas terram: redde retributiónem supérbis.

\rubric{In Officio Festi vel Octávæ, Dominicæ anticipátæ ac sanctæ Maríæ in Sabbato:}

\capitulumlabel{Capitulum}{1 Tim. 1, 17}

\lettrine[lines=2]{R}{EGI} sæculórum immortáli et invisíbili, soli Deo honor et glória in sæcula sæculórum. Amen.

\rubric{In ómnibus autem Fériis et in Vigíliis commúnibus:}

\capitulumlabel{Capitulum}{Zach. 8, 19}

\lettrine[lines=2]{P}{ACEM} et veritátem dilígite, ait Dóminus omnípotens.

\R\ br. Christe, Fili Dei vivi, * Miserére nobis. Christe. \V\ \textit{Qui sedes ad déxteram Patris.} Miserére nobis. Glória Patri. Christe. \V\ Exsúrge, Christe, ádjuva nos. \R\ Et líbera nos propter nomen tuum.

\rubric{In Officio sanctæ Maríæ in Sabbato dícitur sequens}

\variantheading{Preces dominicales et feriales (si dicendæ sint)}

KYRIE, eléison. Christe, eléison. Kýrie, eléison.

Pater noster \rubric{secreto usque ad} \V\ Et ne nos indúcas in tentatiónem. \R\ Sed líbera nos a malo.

Credo in Deum \rubric{secreto usque ad} \V\ Carnis resurrectiónem. \R\ Vitam ætérnam. Amen. \V\ Et ego ad te, Dómine, clamávi. \R\ Et mane orátio mea prævéniet te. \V\ Repleátur os meum laude. \R\ Ut cantem glóriam tuam, tota die magnitúdinem tuam. \V\ Dómine, avérte fáciem tuam a peccátis meis. \R\ Et omnes iniquitátes meas dele. \V\ Cor mundum crea in me, Deus. \R\ Et spíritum rectum ínnova in viscéribus meis. \V\ Ne projícias me a fácie tua. \R\ Et spíritum sanctum tuum ne áuferas a me. \V\ Redde mihi lætítiam salutáris tui. \R\ Et spíritu principáli confírma me.

\variantheading{In Précibus feriálibus tantum}

\V\ Eripe me, Dómine, ab hómine malo. \R\ A viro iníquo éripe me. \V\ Eripe me de inimícis meis, Deus meus. \R\ Et ab insurgéntibus in me líbera me. \V\ Eripe me de operántibus iniquitátem. \R\ Et de viris sánguinum salva me. \V\ Sic psalmum dicam nómini tuo in sæculum sæculi. \R\ Ut reddam vota mea de die in diem. \V\ Exáudi nos, Deus, salutáris noster. \R\ Spes ómnium fínium terræ, et in mari longe. \V\ Deus, in adjutórium meum inténde. \R\ Dómine, ad adjuvándum me festína. \V\ Sanctus Deus, Sanctus fortis, Sanctus immortális. \R\ Miserére nobis. \V\ Bénedic, ánima mea, Dómino. \R\ Et ómnia, quæ intra me sunt, nómini sancto ejus. \V\ Bénedic, ánima mea, Dómino. \R\ Et noli oblivísci omnes retributiónes ejus. \V\ Qui propitiátur ómnibus iniquitátibus tuis. \R\ Qui sanat omnes infirmitátes tuas. \V\ Qui rédimit de intéritu vitam tuam. \R\ Qui corónat te in misericórdia et miseratiónibus. \V\ Qui replet in bonis desidérium tuum. \R\ Renovábitur ut áquilæ juvéntus tua.

\variantheading{In utrisque Précibus}

\V\ Adjutórium nostrum in nómine Dómini. \R\ Qui fecit cælum et terram.

\formulalabel{Confíteor. Misereátur. Indulgéntiam.}

\V\ Dignáre, Dómine, die isto. \R\ Sine peccáto nos custodíre. \V\ Miserére nostri, Dómine. \R\ Miserére nostri. \V\ Fiat misericórdia tua, Dómine, super nos. \R\ Quemádmodum sperávimus in te. \V\ Dómine, exáudi oratiónem meam. \R\ Et clamor meus ad te véniat.

\V\ Dóminus vobíscum.

Orémus. \oratiolabel{Oratio}

\lettrine[lines=2]{D}{OMINE}, Deus omnípotens, qui ad princípium hujus diéi nos perveníre fecísti: tua nos hódie salva virtúte; ut in hac die ad nullum declinémus peccátum, sed semper ad tuam justítiam faciéndam nostra procédant elóquia, dirigántur cogitatiónes et ópera. Per Dóminum. \R\ Amen.

\V\ Dóminus vobíscum. \V\ Benedicámus Dómino.

\rubric{Deinde in Choro legitur Martyrológium. Postea Hebdomadárius dicit:}

\V\ Pretiósa in conspéctu Dómini. \R\ Mors Sanctórum ejus.

\rubric{Et absolúte, sine} Orémus:

\oratiolabel{Oratio}

\lettrine[lines=2]{S}{ANCTA} María, et ómnes Sancti intercédant pro nobis ad Dóminum, ut nos mereámur ab eo adjuvári et salvári, qui vivit et regnat in sæcula sæculórum. \R\ Amen.

\V\ Deus, in adjutórium meum inténde. \R\ Dómine, ad adjuvándum me festína. \rubric{Et dícitur ter; último additur:} Glória Patri. Sicut erat.

Kýrie, eléison. Christe, eléison. Kýrie, eléison.

Pater noster \rubric{secreto usque ad} \V\ Et ne nos indúcas in tentatiónem. \R\ Sed líbera nos a malo. \V\ Réspice in servos tuos, Dómine, et in ópera tua, et dírige fílios eórum. \R\ Et sit splendor Dómini, Dei nostri, super nos, et ópera mánuum nostrárum dírige super nos, et opus mánuum nostrárum dírige. \V\ Glória Patri. \R\ Sicut erat.

Orémus. \oratiolabel{Oratio}

\lettrine[lines=2]{D}{IRIGERE} et sanctificáre, régere et gubernáre dignáre, Dómine Deus, Rex cæli et terræ, hódie corda et córpora nostra, sensus, sermónes et actus nostros in lege tua, et in opéribus mandatórum tuórum: ut hic, et in ætérnum, te auxiliánte, salvi et líberi esse mereámur, Salvátor mundi: Qui vivis et regnas in sæcula sæculórum. \R\ Amen.

\V\ Jube, domne.

\benedlabel{Bened.} Dies et actus nostros in sua pace dispónat Dóminus omnípotens. \R\ Amen.

\variantheading{In Officio feriali}

\lettrine[lines=2]{D}{OMINUS} autem dírigat corda et córpora nostra in caritáte Dei, et patiéntia Christi. Tu autem.

\variantheading{E Proprio Sanctorum}

\variantheading{Pro Oct. Omn. Sanctorum}

\lettrine[lines=2]{B}{ENEDÍCTIO} et cláritas et sapiéntia et gratiárum áctio, honor, virtus et fortitúdo Deo nostro in sǽcula sæculórum. Amen.

\variantheading{E Communi Sanctorum}

\variantheading{Pro Apostolis}

\lettrine[lines=2]{I}{BANT} Apóstoli gaudéntes a conspéctu concílii, quóniam digni hábiti sunt pro nómine Jesu contuméliam pati.

\variantheading{Pro Martyre}

\lettrine[lines=2]{J}{USTUS} cor suum trádidit ad vigilándum dilúculo ad Dóminum, qui fecit illum, et in conspéctu Altíssimi deprecábitur.

\variantheading{Pro pluribus Martyribus}

\lettrine[lines=2]{F}{ULGEBUNT} justi, et tamquam scintíllæ in arundinéto discúrrent. Judicábunt natiónes, et dominabúntur pópulis: et regnábit Dóminus illórum in perpétuum.

\variantheading{Pro Conf. Pont.}

\lettrine[lines=2]{F}{UNGI} sacerdótio, et habére laudem in nómine ipsíus, et offérre illi incénsum dignum in odórem suavitátis.

\variantheading{Pro Conf. non Pont.}

\lettrine[lines=2]{J}{USTUM} dedúxit Dóminus per vias rectas, et osténdit illi regnum Dei, et dedit illi sciéntiam sanctórum; honestávit illum in labóribus, et complévit labóres illíus.

\variantheading{Pro Virgine Martyre}

\lettrine[lines=2]{D}{OMINE}, Deus meus, exaltásti super terram habitatiónem meam, et pro morte defluénte deprecáta sum. Invocávi Dóminum, Patrem Dómini mei, ut non derelínquat me in die tribulatiónis meæ, et in témpore superbórum sine adjutório.

\variantheading{Pro Virgine non Martyre}

\lettrine[lines=2]{O}{} \textsc{quam} pulchra est casta generátio cum claritáte! immortális est enim memória illíus: quóniam apud Deum nota est, et apud hómines.

\variantheading{Pro Martyre tantum}

\lettrine[lines=2]{L}{AUDABIT} usque ad mortem ánima mea Dóminum, quóniam éruis sustinéntes te, et líberas eos de manu angústiæ, Dómine, Deus noster.

\variantheading{Pro nec Virg. nec Mart.}

\lettrine[lines=2]{M}{ULTAE} fíliæ congregavérunt divítias: tu supergréssa es univérsas. Fallax grátia, et vana est pulchritúdo: múlier timens Dóminum, ipsa laudábitur.

\variantheading{Pro Dedicatione Ecclesiæ}

\lettrine[lines=2]{E}{T} abstérget Deus omnem lácrimam ab óculis eórum: et mors ultra non erit, neque luctus, neque clamor, neque dolor erit ultra, quia prima abiérunt. Et dixit qui sedébat in throno: Ecce, nova fácio ómnia.

\variantheading{In Officiis B. Mariæ V.}

\lettrine[lines=2]{I}{N} platéis sicut cinnamómum et bálsamum aromatízans odórem dedi: quasi myrrha elécta dedi suavitátem odóris.

\V\ Adjutórium nostrum in nómine Dómini. \R\ Qui fecit cælum et terram.

\V\ Benedícite. \R\ Deus.

\benedictiolabel{Benedictio} Dóminus nos benedícat, et ab omni malo deféndat, et ad vitam perdúcat ætérnam. Et fidélium ánimæ per misericórdiam Dei requiéscant in pace. \R\ Amen.

\rubric{Deinde dícitur tantum} Pater noster \rubric{secréto, nisi sequátur ália Hora.}

% ===== END OF BATCH 39 -- Covers autumnalis.md lines 12676-12954,
% completing "Ad Primam" for Sabbato in full: the hymn "Jam lucis
% orto sídere," Psalms 93/i, 93/ii, 107, the conditional Psalm 149
% insertion, both Capitula (the first now naming "Dominicæ
% anticipátæ" -- the anticipated Sunday -- alongside the usual
% Festum/Octave/Marian-Saturday conditions), the brief Responsory,
% a new standalone rubric specific to this day directing that the
% Preces which follow are said "In Officio sanctæ Maríæ in
% Sabbato," the full Preces apparatus, the closing Orationes, and
% the Lectio-brevis apparatus. Several more psalm incipits and
% Oratio openings carry the partial-caps OCR defect ("DEus
% ultiónum," "BEatus homo," "PAratum cor meum," "CAntáte," "REgi"
% [Capitulum], "DOmine, Deus omnípotens" [Oratio], "SAncta María"
% [Oratio], "DIrígere" [Oratio]) and are silently normalized to
% full caps. The document continues with "Ad Tertiam" (autumnalis.md
% line 12958).
% =====
\hourheading{Ad Tertiam}

Pater noster et Ave.

\V\ Deus, in adjutórium.

\hymnusheading{Hymnus}

\lettrine[lines=2]{N}{UNC}, Sancte, nobis, Spíritus,\\
Unum Patri cum Fílio,\\
Dignáre promptus íngeri\\
Nostro refúsus péctori.

Os, lingua, mens, sensus, vigor\\
Confessiónem pérsonent,\\
Flamméscat igne cáritas,\\
Accéndat ardor próximos.

Præsta, Pater piíssime,\\
Patríque compar Unice,\\
Cum Spíritu Paráclito\\
Regnans per omne sæculum.\\
Amen.

\antlabel{Ant.} Clamor meus.

\psalmheading{Psalmus 101, i}

\lettrine[lines=2]{D}{OMINE}, exáudi oratiónem meam: * et clamor meus ad te véniat.

Non avértas fáciem tuam a me: * in quacúmque die tríbulor, inclína ad me aurem tuam.

In quacúmque die invocávero te, * velóciter exáudi me.

Quia defecérunt sicut fumus dies mei: * et ossa mea sicut crémium aruérunt.

Percússus sum ut fœnum, et áruit cor meum: * quia oblítus sum comédere panem meum.

A voce gémitus mei * adhæsit os meum carni meæ.

Símilis factus sum pellicáno solitúdinis: * factus sum sicut nyctícorax in domicílio.

Vigiláví, * et factus sum sicut passer solitárius in tecto.

Tota die exprobrábant mihi inimíci mei: * et qui laudábant me, advérsum me jurábant.

Quia cínerem tamquam panem manducábam, * et potum meum cum fletu miscébam.

A fácie iræ et indignatiónis tuæ: * quia élevans allisísti me.

Dies mei sicut umbra declinavérunt: * et ego sicut fœnum árui.

Tu autem, Dómine, in ætérnum pérmanes: * et memoriále tuum in generatiónem et generatiónem.

\psalmheading{Psalmus 101, ii}

\lettrine[lines=2]{T}{U} exsúrgens miseréberis Sion: * quia tempus miseréndi ejus, quia venit tempus.

Quóniam placuérunt servis tuis lápides ejus: * et terræ ejus miserebúntur.

Et timébunt gentes nomen tuum, Dómine, * et omnes reges terræ glóriam tuam.

Quia ædificávit Dóminus Sion: * et vidébitur in glória sua.

Respéxit in oratiónem humílium: * et non sprevit precem eórum.

Scribántur hæc in generatióne áltera: * et pópulus, qui creábitur, laudábit Dóminum:

Quia prospéxit de excélso sancto suo: * Dóminus de cælo in terram aspéxit:

Ut audíret gémitus compeditórum: * ut sólveret fílios interemptórum:

Ut annúntient in Sion nomen Dómini: * et laudem ejus in Jerúsalem.

In conveniéndo pópulos in unum, * et reges ut sérviant Dómino.

\psalmheading{Psalmus 101, iii}

\lettrine[lines=2]{R}{ESPONDIT} ei in via virtútis suæ: * Paucitátem diérum meórum núntia mihi.

Ne révoces me in dimídio diérum meórum: * in generatiónem et generatiónem anni tui.

Inítio tu, Dómine, terram fundásti: * et ópera mánuum tuárum sunt cæli.

Ipsi períbunt, tu autem pérmanes: * et omnes sicut vestiméntum veteráscent.

Et sicut opertórium mutábis eos, et mutabúntur: * tu autem idem ipse es, et anni tui non defícient.

Fílii servórum tuórum habitábunt: * et semen eórum in sæculum dirigétur.

\antlabel{Ant.} Clamor meus, Dómine, ad te pervéniat: * non avértas fáciem tuam a me.

\variantheading{In Officio feriali}

\capitulumlabel{Capitulum}{Jer. 17, 14}

\lettrine[lines=2]{S}{ANA} me, Dómine, et sanábor: salvum me fac, et salvus ero: quóniam laus mea tu es.

\R\ br. Sana ánimam meam, * Quia peccávi tibi. Sana. \V\ Ego dixi: Dómine, miserére mei. Quia. Glória Patri. Sana. \V\ Adjútor meus esto, ne derelínquas me. \R\ Neque despícias me, Deus, salutáris meus.

\variantheading{In Vigiliis communibus:}

\variantheading{Preces feriales}

KYRIE, eléison. Christe, eléison. Kýrie, eléison.

Pater noster \rubric{secreto usque ad} \V\ Et ne nos indúcas in tentatiónem. \R\ Sed líbera nos a malo. \V\ Dómine, Deus virtútum, convérte nos. \R\ Et osténde fáciem tuam, et salvi érimus. \V\ Exsúrge, Christe, ádjuva nos. \R\ Et líbera nos propter nomen tuum. \V\ Dómine, exáudi oratiónem meam. \R\ Et clamor meus ad te véniat.

\rubric{Réliqua ut in Ordinário. 19}

\hourheading{Ad Sextam}

Pater noster et Ave.

\V\ Deus, in adjutórium.

\hymnusheading{Hymnus}

\lettrine[lines=2]{R}{ECTOR} potens, verax Deus,\\
Qui témperas rerum vices,\\
Splendóre mane illúminas,\\
Et ígnibus merídiem:

Exstíngue flammas lítium,\\
Aufer calórem nóxium,\\
Confer salútem córporum,\\
Verámque pacem córdium.

Præsta, Pater piíssime,\\
Patríque compar Unice,\\
Cum Spíritu Paráclito\\
Regnans per omne sæculum.\\
Amen.

\antlabel{Ant.} Dómine, Deus meus.

\psalmheading{Psalmus 103, i}

\lettrine[lines=2]{B}{ENEDIC}, ánima mea, Dómino: * Dómine, Deus meus, magnificátus es vehementer.

Confessiónem, et decórem induísti: * amíctus lúmine sicut vestiménto:

Exténdens cælum sicut pellem: * qui tegis aquis superióra ejus.

Qui ponis nubem ascénsum tuum: * qui ámbulas super pennas ventórum.

Qui facis ángelos tuos, spíritus: * et minístros tuos ignem uréntem.

Qui fundásti terram super stabilitátem suam: * non inclinábitur in sæculum sæculi.

Abýssus, sicut vestiméntum, amíctus ejus: * super montes stabunt aquæ.

Ab increpatióne tua fúgient: * a voce tonítrui tui formidábunt.

Ascéndunt montes: et descéndunt campi * in locum, quem fundásti eis.

Términum posuísti, quem non transgrediéntur: * neque converténtur operíre terram.

Qui emíttis fontes in convállibus: * inter médium móntium pertransíbunt aquæ.

Potábunt omnes béstiæ agri: * exspectábunt ónagri in siti sua.

Super ea volúcres cæli habitábunt: * de médio petrárum dabunt voces.

\psalmheading{Psalmus 103, ii}

\lettrine[lines=2]{R}{IGANS} montes de superióribus suis: * de fructu óperum tuórum satiábitur terra:

Prodúcens fœnum juméntis, * et herbam servitúti hóminum:

Ut edúcas panem de terra: * et vinum lætíficet cor hóminis:

Ut exhílaret fáciem in óleo: * et panis cor hóminis confírmet.

Saturabúntur ligna campi, et cedri Líbani, quas plantávit: * illic pásseres nidificábunt.

Heródii domus dux est eórum: * montes excélsi cervis: petra refúgium herináciis.

Fecit lunam in témpora: * sol cognóvit occásum suum.

Posuísti ténebras, et facta est nox: * in ipsa pertransíbunt omnes béstiæ silvæ.

Cátuli leónum rugiéntes, ut rápiant, * et quærant a Deo escam sibi.

Ortus est sol, et congregáti sunt: * et in cubílibus suis collocabúntur.

Exíbit homo ad opus suum: * et ad operatiónem suam usque ad vésperum.

\psalmheading{Psalmus 103, iii}

\lettrine[lines=2]{Q}{UAM} magnificáta sunt ópera tua, Dómine! * ómnia in sapiéntia fecísti: impléta est terra possessióne tua.

Hoc mare magnum, et spatiósum mánibus: * illic reptília, quorum non est númerus.

Animália pusílla cum magnis: * illic naves pertransíbunt.

Draco iste, quem formásti ad illudéndum ei: * ómnia a te exspéctant ut des illis escam in témpore.

Dante te illis, cólligent: * aperiénte te manum tuam, ómnia implebúntur bonitáte.

Avertánte autem te fáciem, turbabúntur: * áuferes spíritum eórum, et defícient, et in púlverem suum reverténtur.

Emíttes spíritum tuum, et creabúntur: * et renovábis fáciem terræ.

Sit glória Dómini in sæculum: * lætábitur Dóminus in opéribus suis:

Qui réspicit terram, et facit eam trémere: * qui tangit montes, et fúmigant.

Cantábo Dómino in vita mea: * psallam Deo meo, quámdiu sum.

Jucúndum sit ei elóquium meum: * ego vero delectábor in Dómino.

Defíciant peccatóres a terra, et iníqui ita ut non sint: * bénedic, ánima mea, Dómino.

\antlabel{Ant.} Dómine, Deus meus, magnificátus es veheménter.

\variantheading{In Officio feriali}

\capitulumlabel{Capitulum}{Rom. 13, 8}

\lettrine[lines=2]{N}{EMINI} quidquam debeátis, nisi ut ínvicem diligátis: qui enim díligit próximum, legem implévit.

\R\ br. Benedícam Dóminum * In omni témpore. Benedícam. \V\ Semper laus ejus in ore meo. In. Glória Patri. Benedícam. \V\ Dóminus regit me, et nihil mihi déerit. \R\ In loco páscuæ ibi me collocávit.

\variantheading{In Vigiliis communibus:}

\variantheading{Preces feriales}

KYRIE, eléison. Christe, eléison. Kýrie, eléison.

Pater noster \rubric{secreto usque ad} \V\ Et ne nos indúcas in tentatiónem. \R\ Sed líbera nos a malo. \V\ Dómine, Deus virtútum, convérte nos. \R\ Et osténde fáciem tuam, et salvi érimus. \V\ Exsúrge, Christe, ádjuva nos. \R\ Et líbera nos propter nomen tuum. \V\ Dómine, exáudi oratiónem meam. \R\ Et clamor meus ad te véniat.

\rubric{Réliqua ut in Ordinário. 20}

\hourheading{Ad Nonam}

Pater noster et Ave.

\V\ Deus, in adjutórium.

\hymnusheading{Hymnus}

\lettrine[lines=2]{R}{ERUM}, Deus, tenax vigor,\\
Immótus in te pérmanens,\\
Lucis diúrnæ témpora\\
Succéssibus detérminans:

Largíre lumen véspere,\\
Quo vita nusquam décidat,\\
Sed præmium mortis sacræ\\
Perénnis instet glória.

Præsta, Pater piíssime,\\
Patríque compar Unice,\\
Cum Spíritu Paráclito\\
Regnans per omne sæculum.\\
Amen.

\antlabel{Ant.} Ne tacúeris, Deus.

\psalmheading{Psalmus 108, i}

\lettrine[lines=2]{D}{EUS}, laudem meam ne tacúeris: * quia os peccatóris, et os dolósi super me apértum est.

Locúti sunt advérsum me lingua dolósa, et sermónibus ódii circumdedérunt me: * et expugnavérunt me gratis.

Pro eo ut me dilígerent, detrahébant mihi: * ego autem orábam.

Et posuérunt advérsum me mala pro bonis: * et ódium pro dilectióne mea.

Constítue super eum peccatórem: * et diábolus stet a dextris ejus.

Cum judicátur, éxeat condemnátus: * et orátio ejus fiat in peccátum.

Fiant dies ejus pauci: * et episcopátum ejus accípiat alter.

Fiant fílii ejus órphani: * et uxor ejus vídua.

Nutántes transferántur fílii ejus, et mendícent: * et ejiciántur de habitatiónibus suis.

Scrutétur fænerátor omnem substántiam ejus: * et dirípiant aliéni labóres ejus.

Non sit illi adjútor: * nec sit qui misereátur pupíllis ejus.

Fiant nati ejus in intéritum: * in generatióne una deleátur nomen ejus.

\psalmheading{Psalmus 108, ii}

\lettrine[lines=2]{I}{N} memóriam rédeat iníquitas patrum ejus in conspéctu Dómini: * et peccátum matris ejus non deleátur.

Fiant contra Dóminum semper, et dispéreat de terra memória eórum: * pro eo quod non est recordátus fácere misericórdiam.

Et persecútus est hóminem ínopem, et mendícum, * et compúnctum corde mortificáre.

Et diléxit maledictiónem, et véniet ei: * et nóluit benedictiónem, et elongábitur ab eo.

Et índuit maledictiónem sicut vestiméntum, * et intrávit sicut aqua in interióra ejus, et sicut óleum in óssibus ejus.

Fiat ei sicut vestiméntum, quo operítur: * et sicut zona, qua semper præcíngitur.

Hoc opus eórum, qui détrahunt mihi apud Dóminum: * et qui loquúntur mala advérsum ánimam meam.

Et tu, Dómine, Dómine, fac mecum propter nomen tuum: * quia suávis est misericórdia tua.

\psalmheading{Psalmus 108, iii}

\lettrine[lines=2]{L}{IBERA} me quia egénus, et pauper ego sum: * et cor meum conturbátum est intra me.

Sicut umbra cum declínat, ablátus sum: * et excússus sum sicut locústæ.

Génua mea infirmáta sunt a jejúnio: * et caro mea immutáta est propter óleum.

Et ego factus sum oppróbrium illis: * vidérunt me, et movérunt cápita sua.

Adjuva me, Dómine, Deus meus: * salvum me fac secúndum misericórdiam tuam.

Et sciant quia manus tua hæc: * et tu, Dómine, fecísti eam.

Maledícent illi, et tu benedíces: * qui insúrgunt in me, confundántur: servus autem tuus lætábitur.

Induántur qui détrahunt mihi, pudóre: * et operiántur sicut diplóide confusióne sua.

Confitébor Dómino nimis in ore meo: * et in médio multórum laudábo eum.

Quia ástitit a dextris páuperis, * ut salvam fáciat a persequéntibus ánimam meam.

\antlabel{Ant.} Ne tacúeris, Deus, quia sermónibus ódii circumdedérunt me.

\variantheading{In Officio feriali}

\capitulumlabel{Capitulum}{1 Petri 1, 17-19}

\lettrine[lines=2]{I}{N} timóre incolátus vestri témpore conversámini: sciéntes quod non corruptibílibus auro vel argénto redémpti estis, sed pretióso sánguine quasi Agni immaculáti Christi.

\R\ br. Rédime me, Dómine, * Et miserére mei. Rédime. \V\ Pes enim meus stetit in via recta. Et. Glória Patri. Rédime.

\V\ Ab occúltis meis munda me, Dómine. \R\ Et ab aliénis parce servo tuo.

\variantheading{In Vigiliis communibus:}

\variantheading{Preces feriales}

KYRIE, eléison. Christe, eléison. Kýrie, eléison.

Pater noster \rubric{secreto usque ad} \V\ Et ne nos indúcas in tentatiónem. \R\ Sed líbera nos a malo. \V\ Dómine, Deus virtútum, convérte nos. \R\ Et osténde fáciem tuam, et salvi érimus. \V\ Exsúrge, Christe, ádjuva nos. \R\ Et líbera nos propter nomen tuum. \V\ Dómine, exáudi oratiónem meam. \R\ Et clamor meus ad te véniat.

\rubric{Réliqua ut in Ordinário. 22}

% ===== END OF BATCH 40 -- Covers autumnalis.md lines 12958-13309,
% completing the Little Hours of Sabbato in full: "Ad Tertiam"
% (Psalms 101/i-iii), "Ad Sextam" (Psalms 103/i-iii), and "Ad
% Nonam" (Psalms 108/i-iii), each with its own hymn, Capitulum,
% brief Responsory, and "In Vigiliis communibus: Preces feriales"
% apparatus. Several more partial-caps OCR defects ("REspóndit,"
% "BEnedic," "RIgans," "QUam magnificáta," "DEus, laudem," "IN
% memóriam," "LIbera me") are silently normalized to full caps.
% The document continues with "Ad Vesperas" (autumnalis.md line
% 13312), which carries a special rubric about this Saturday's
% Vespers serving as First Vespers of the following Sunday.
% =====
\hourheading{Ad Vesperas}

Pater noster et Ave.

\V\ Deus, in adjutórium.

\antlabel{Ant.} Benedíctus Dóminus. \textdagger

\rubric{¶ Antíphonæ cum Psalmis assignátæ semper in Sábbato dicúntur pro I Vésperis Domínicæ sequéntis, quando agéndum est Offícium de ipsa Domínica, in Domínicis minóribus, extra Octávas Sanctíssimi Córporis Christi ac Sacratíssimi Cordis Jesu.}

\psalmheading{Psalmus 143, i}

\lettrine[lines=2]{B}{ENEDICTUS} Dóminus, \textdagger\ Deus meus, qui docet manus meas ad prœlium, * et dígitos meos ad bellum.

Misericórdia mea, et refúgium meum: * suscéptor meus, et liberátor meus:

Protéctor meus, et in ipso sperávi: * qui subdit pópulum meum sub me.

Dómine, quid est homo, quia innotuísti ei? * aut fílius hóminis, quia réputas eum?

Homo vanitáti símilis factus est: * dies ejus sicut umbra prætéreunt.

Dómine, inclína cælos tuos, et descénde: * tange montes, et fumigábunt.

Fúlgura coruscatiónem, et dissipábis eos: * emítte sagíttas tuas, et conturbábis eos.

Emítte manum tuam de alto, éripe me, et líbera me de aquis multis: * de manu filiórum alienórum.

Quorum os locútum est vanitátem: * et déxtera eórum, déxtera iniquitátis.

\antlabel{Ant.} Benedíctus Dóminus, suscéptor meus, et liberátor meus.

\antlabel{Ant.} Beátus pópulus.

\psalmheading{Psalmus 143, ii}

\lettrine[lines=2]{D}{EUS}, cánticum novum cantábo tibi: * in psaltério decachórdo psallam tibi.

Qui das salútem régibus: * qui redemísti David, servum tuum, de gládio malígno: éripe me.

Et érue me de manu filiórum alienórum, quorum os locútum est vanitátem: * et déxtera eórum, déxtera iniquitátis:

Quorum fílii, sicut novéllæ plantatiónes * in juventúte sua.

Fíliæ eórum compósitæ: * circumornátæ ut similitúdo templi.

Promptuária eórum plena, * eructántia ex hoc in illud.

Oves eórum fetósæ, abundántes in egréssibus suis: * boves eórum crassæ.

Non est ruína macériæ, neque tránsitus: * neque clamor in platéis eórum.

Beátum dixérunt pópulum, cui hæc sunt: * beátus pópulus, cujus Dóminus Deus ejus.

\antlabel{Ant.} Beátus pópulus, cujus Dóminus Deus ejus.

\antlabel{Ant.} Magnus Dóminus.

\psalmheading{Psalmus 144, i}

\lettrine[lines=2]{E}{XALTABO} te, Deus meus, rex: * et benedícam nómini tuo in sæculum, et in sæculum sæculi.

Per síngulos dies benedícam tibi: * et laudábo nomen tuum in sæculum, et in sæculum sæculi.

Magnus Dóminus, et laudábilis nimis: * et magnitúdinis ejus non est finis.

Generátio et generátio laudábit ópera tua: * et poténtiam tuam pronuntiábunt.

Magnificéntiam glóriæ sanctitátis tuæ loquéntur: * et mirabília tua narrábunt.

Et virtútem terribílium tuórum dicent: * et magnitúdinem tuam narrábunt.

Memóriam abundántiæ suavitátis tuæ eructábunt: * et justítia tua exsultábunt.

\antlabel{Ant.} Magnus Dóminus, et laudábilis nimis: et magnitúdinis ejus non est finis.

\antlabel{Ant.} Suávis Dóminus.

\psalmheading{Psalmus 144, ii}

\lettrine[lines=2]{M}{ISERATOR}, et miséricors Dóminus: * pátiens, et multum miséricors.

Suávis Dóminus univérsis: * et miseratiónes ejus super ómnia ópera ejus.

Confiteántur tibi, Dómine, ómnia ópera tua: * et sancti tui benedícant tibi.

Glóriam regni tui dicent: * et poténtiam tuam loquéntur:

Ut notam fáciant fíliis hóminum poténtiam tuam: * et glóriam magnificéntiæ regni tui.

Regnum tuum regnum ómnium sæculórum: * et dominátio tua in omni generatióne et generatiónem.

\antlabel{Ant.} Suávis Dóminus univérsis: et miseratiónes ejus super ómnia ópera ejus.

\antlabel{Ant.} Fidélis Dóminus. \textdagger

\psalmheading{Psalmus 144, iii}

\lettrine[lines=2]{F}{IDELIS} Dóminus \textdagger\ in ómnibus verbis suis: * et sanctus in ómnibus opéribus suis.

Allevat Dóminus omnes qui córruunt: * et érigit omnes elísos.

Óculi ómnium in te sperant, Dómine: * et tu das escam illórum in témpore opportúno.

Áperis tu manum tuam: * et imples omne ánimal benedictióne.

Justus Dóminus in ómnibus viis suis: * et sanctus in ómnibus opéribus suis.

Prope est Dóminus ómnibus invocántibus eum: * ómnibus invocántibus eum in veritáte.

Voluntátem timéntium se fáciet: * et deprecatiónem eórum exáudiet: et salvos fáciet eos.

Custódit Dóminus omnes diligéntes se: * et omnes peccatóres dispérdet.

Laudatiónem Dómini loquétur os meum: * et benedícat omnis caro nómini sancto ejus in sæculum, et in sæculum sæculi.

\antlabel{Ant.} Fidélis Dóminus in ómnibus verbis suis: et sanctus in ómnibus opéribus suis.

\rubric{In Sabbatis extra Octávas Ssmi Córporis Christi ac Sacratíssimi Cordis Jesu dícitur}

\variantheading{In Officio feriali}

\capitulumlabel{Capitulum}{Rom. 11, 33}

\lettrine[lines=2]{O}{} \textsc{altitudo} divitiárum sapiéntiæ et sciéntiæ Dei: quam incomprehensíbilia sunt judícia ejus, et investigábiles viæ ejus!

\hymnusheading{Hymnus}

\lettrine[lines=2]{J}{AM} sol recédit ígneus:\\
Tu, lux perénnis, Únitas,\\
Nostris, beáta Trínitas,\\
Infúnde lumen córdibus.

Te mane laudum cármine,\\
Te deprecámur véspere;\\
Dignéris ut te súpplices\\
Laudémus inter Cælites.

Patri, simúlque Fílio,\\
Tibíque, Sancte Spíritus,\\
Sicut fuit, sit júgiter\\
Sæclum per omne glória.\\
Amen.

\V\ Vespertína orátio ascéndat ad te, Dómine. \R\ Et descéndat super nos misericórdia tua.

\rubric{Ad Magníficat Antíphona ut in Próprio de Témpore.}

\rubric{Expléto Cántico ac repetíta Antíphona, dícitur Orátio Domínicæ sequéntis própria, quæ per totam hebdómadam suam in feriáli Officio adhibétur, nisi speciális Orátio habeátur.}

\rubric{Réliqua ut in Ordinário. 22}

\hourheading{Ad Completorium}

\V\ Jube, domne.

\benedlabel{Bened.} Noctem quiétam, et finem perféctum concédat nobis Dóminus omnípotens. \R\ Amen.

\lectioheading{Lectio br.}{1 Petri 5, 8-9}

\lettrine[lines=2]{F}{RATRES}: Sóbrii estóte, et vigiláte: quia adversárius vester diábolus tamquam leo rúgiens círcuit, quærens quem dévoret: cui resístite fortes in fide. Tu autem.

\R\ Deo grátias.

\V\ Adjutórium nostrum in nómine Dómini. \R\ Qui fecit cælum et terram.

Pater noster, \rubric{quod dícitur totum secréto.}

\formulalabel{Confíteor.} Misereátur. Indulgéntiam.

\V\ Convérte nos, Deus, salutáris noster. \R\ Et avérte iram tuam a nobis.

\V\ Deus, in adjutórium.

\antlabel{Ant.} Intret orátio mea.

\psalmheading{Psalmus 87}

\lettrine[lines=2]{D}{OMINE}, Deus salútis meæ: * in die clamávi, et nocte coram te.

Intret in conspéctu tuo orátio mea: * inclína aurem tuam ad precem meam:

Quia repléta est malis ánima mea: * et vita mea inférno appropinquávit.

Æstimátus sum cum descendéntibus in lacum: * factus sum sicut homo sine adjutório, inter mórtuos liber.

Sicut vulneráti dormiéntes in sepúlcris, quorum non es memor ámplius: * et ipsi de manu tua repúlsi sunt.

Posuérunt me in lacu inferióri: * in tenebrósis, et in umbra mortis.

Super me confirmátus est furor tuus: * et omnes fluctus tuos induxísti super me.

Longe fecísti notos meos a me: * posuérunt me abominatiónem sibi.

Tráditus sum, et non egrediébar: * óculi mei languérunt præ inópia.

Clamávi ad te, Dómine, tota die: * expándi ad te manus meas.

Numquid mórtuis fácies mirabília: * aut médici suscitábunt, et confitebúntur tibi?

Numquid narrábit áliquis in sepúlcro misericórdiam tuam, * et veritátem tuam in perditióne?

Numquid cognoscéntur in ténebris mirabília tua, * et justítia tua in terra obliviónis?

Et ego ad te, Dómine, clamávi: * et mane orátio mea prævéniet te.

Ut quid, Dómine, repéllis oratiónem meam: * avértis fáciem tuam a me?

Pauper sum ego, et in labóribus a juventúte mea: * exaltátus autem, humiliátus sum et conturbátus.

In me transiérunt iræ tuæ: * et terróres tui conturbavérunt me.

Circumdedérunt me sicut aqua tota die: * circumdedérunt me simul.

Elongásti a me amícum et próximum: * et notos meos a miséria.

\psalmheading{Psalmus 102, i}

\lettrine[lines=2]{B}{ENEDIC}, ánima mea, Dómino: * et ómnia, quæ intra me sunt, nómini sancto ejus.

Bénedic, ánima mea, Dómino: * et noli oblivísci omnes retributiónes ejus.

Qui propitiátur ómnibus iniquitátibus tuis: * qui sanat omnes infirmitátes tuas.

Qui rédimit de intéritu vitam tuam: * qui corónat te in misericórdia et miseratiónibus.

Qui replet in bonis desidérium tuum: * renovábitur ut áquilæ juvéntus tua:

Fáciens misericórdias Dóminus: * et judícium ómnibus injúriam patiéntibus.

Notas fecit vias suas Móysi, * fíliis Israël voluntátes suas.

Miserátor, et miséricors Dóminus: * longánimis, et multum miséricors.

Non in perpétuum irascétur: * neque in ætérnum comminábitur.

Non secúndum peccáta nostra fecit nobis: * neque secúndum iniquitátes nostras retríbuit nobis.

Quóniam secúndum altitúdinem cæli a terra: * corroborávit misericórdiam suam super timéntes se.

Quantum distat ortus ab occidénte: * longe fecit a nobis iniquitátes nostras.

\psalmheading{Psalmus 102, ii}

\lettrine[lines=2]{Q}{UOMODO} miserétur pater filiórum, misértus est Dóminus timéntibus se: * quóniam ipse cognóvit figméntum nostrum.

Recordátus est quóniam pulvis sumus: * homo, sicut fœnum dies ejus, tamquam flos agri sic efflorébit.

Quóniam spíritus pertransíbit in illo, et non subsístet: * et non cognóscet ámplius locum suum.

Misericórdia autem Dómini ab ætérno, * et usque in ætérnum super timéntes eum.

Et justítia illíus in fílios filiórum, * his qui servant testaméntum ejus:

Et mémores sunt mandatórum ipsíus, * ad faciéndum ea.

Dóminus in cælo parávit sedem suam: * et regnum ipsíus ómnibus dominábitur.

Benedícite Dómino, omnes Ángeli ejus: * poténtes virtúte, faciéntes verbum illíus, ad audiéndam vocem sermónum ejus.

Benedícite Dómino, omnes virtútes ejus: * minístri ejus, qui fácitis voluntátem ejus.

Benedícite Dómino, ómnia ópera ejus: * in omni loco dominatiónis ejus, benedíc, ánima mea, Dómino.

\antlabel{Ant.} Intret orátio mea in conspéctu tuo, Dómine.

\hymnusheading{Hymnus}

\lettrine[lines=2]{T}{E} lucis ante términum,\\
Rerum Creátor, póscimus,\\
Ut pro tua cleméntia\\
Sis præsul et custódia.

Procul recédant sómnia,\\
Et nóctium phantásmata;\\
Hostémque nostrum cómprime,\\
Ne polluántur córpora.

Præsta, Pater piíssime,\\
Patríque compar Unice,\\
Cum Spíritu Paráclito\\
Regnans per omne sæculum.\\
Amen.

\capitulumlabel{Capitulum}{Jer. 14, 9}

\lettrine[lines=2]{T}{U} autem in nobis es, Dómine, et nomen sanctum tuum invocátum est super nos: ne derelínquas nos, Dómine, Deus noster.

\R\ br. In manus tuas, Dómine, * Comméndo spíritum meum. In manus. \V\ Redemísti nos, Dómine, Deus veritátis. Comméndo. Glória Patri. In manus.

\V\ Custódi nos, Dómine, ut pupíllam óculi. \R\ Sub umbra alárum tuárum prótege nos.

\antlabel{Ant.} Salva nos.

\canticumheading{Canticum Simeonis}{Luc. 2, 29-32}

\lettrine[lines=2]{N}{UNC} dimíttis servum tuum, Dómine, * secúndum verbum tuum in pace:

Quia vidérunt óculi mei * salutáre tuum,

Quod parásti * ante fáciem ómnium populórum,

Lumen ad revelatiónem géntium, * et glóriam plebis tuæ Israël.

\antlabel{Ant.} Salva nos, Dómine, vigilántes, custódi nos dormiéntes; ut vigilémus cum Christo, et requiescámus in pace.

\variantheading{Preces (si dicendæ sint)}

KYRIE, eléison. Christe, eléison. Kýrie, eléison.

Pater noster \rubric{secreto usque ad} \V\ Et ne nos indúcas in tentatiónem. \R\ Sed líbera nos a malo.

Credo in Deum \rubric{secreto usque ad} \V\ Carnis resurrectiónem. \R\ Vitam ætérnam. Amen. \V\ Benedíctus es, Dómine, Deus patrum nostrórum. \R\ Et laudábilis et gloriósus in sæcula. \V\ Benedicámus Patrem et Fílium cum Sancto Spíritu. \R\ Laudémus, et superexaltémus eum in sæcula. \V\ Benedíctus es, Dómine, in firmaménto cæli. \R\ Et laudábilis, et gloriósus, et superexaltátus in sæcula. \V\ Benedícat et custódiat nos omnípotens et miséricors Dóminus. \R\ Amen. \V\ Dignáre, Dómine, nocte ista. \R\ Sine peccáto nos custodíre. \V\ Miserére nostri, Dómine. \R\ Miserére nostri. \V\ Fiat misericórdia tua, Dómine, super nos. \R\ Quemádmodum sperávimus in te. \V\ Dómine, exáudi oratiónem meam. \R\ Et clamor meus ad te véniat.

\V\ Dóminus vobíscum.

Orémus. \oratiolabel{Oratio}

\lettrine[lines=2]{V}{ISITA}, quæsumus, Dómine, habitatiónem istam, et omnes insídias inimíci ab ea longe repélle: Angeli tui sancti hábitent in ea, qui nos in pace custódiant; et benedíctio tua sit super nos semper. Per Dóminum nostrum. \R\ Amen.

\V\ Dóminus vobíscum.

\V\ Benedicámus Dómino.

\benedictiolabel{Benedictio} Benedícat et custódiat nos omnípotens et miséricors Dóminus, Pater, et Fílius, et Spíritus Sanctus. \R\ Amen.

\variantheading{Antiphona finalis B. M. V.}

\V\ Divínum auxílium.

Pater, Ave María et Credo.

% ===== END OF BATCH 41 -- Covers autumnalis.md lines 13312-13647,
% completing "Ad Vesperas" (Psalms 143/i, 143/ii, 144/i-iii, with
% the special rubric noting these serve as First Vespers of the
% following Sunday) and "Ad Completorium" (Psalms 87, 102/i,
% 102/ii, hymn "Te lucis ante términum," Canticum Simeonis, Preces),
% and with them the ENTIRE "Sabbato" Office and the complete
% seven-day ferial Psalterium cycle of the Pars Autumnalis (Dominica
% through Sabbato). Saturday Vespers here carries its own fixed
% "In Officio feriali" Capitulum ("O Altitúdo divitiárum," rendered
% with the single-letter-word lettrine convention since the source
% gives only "O" capitalized before the normally-cased "Altitúdo")
% and hymn ("Jam sol recédit ígneus," with the recurring stray OCR
% "\*" before its doxology silently dropped), unlike the weekday
% Vespers pattern, since this Capitulum/Hymn pair is fixed for every
% Saturday rather than varying with the day. Two rubrics direct that
% the Magnificat antiphon and the following week's Oratio are drawn
% "ut in Próprio de Témpore" -- from the Proper of the Time for the
% upcoming Sunday -- confirming this Vespers functions as First
% Vespers of Sunday. The Canticum Simeonis here omits the printed
% "Glória Patri" line that appears at every other occurrence of this
% canticle in this volume; transcribed exactly as the source omits
% it. The document's Psalterium section ends here; autumnalis.md
% continues at line 13650 with an "[IMAGE]" placeholder (omitted, a
% non-textual OCR marker, preceded by a "---" scan-artifact rule
% also omitted) followed immediately by the heading "Proprium de
% Tempore" at line 13652, opening with "Sabbato proximo Kalendis
% Septembris" (the Saturday nearest the Kalends of September) --
% this volume's own starting point in the liturgical year, genuinely
% different from aestiva.tex's parallel transition into "In Festo
% Ssmæ Trinitatis" (Trinity Sunday), confirming the Proprium de
% Tempore itself is NOT shared boilerplate and must be transcribed
% fresh from here on -- a major structural transition from the
% ferial Psalterium into the Proper of the Time.
% =====
\partheading{Proprium de Tempore}

\feriaheading{Sabbato proximo Kalendis Septembris}

\hourheading{Ad Vesperas}

\V\ Vespertína orátio ascéndat ad te, Dómine. \R\ Et descéndat super nos misericórdia tua.

\antlabel{Ad Magnif. Ant.} Cum audísset Job * nuntiórum verba, sustínuit patiénter, et ait: Si bona suscépimus de manu Dómini, mala autem quare non sustineámus? In ómnibus his non peccávit Job lábiis suis, neque stultum áliquid contra Deum locútus est.

\rubric{Orátio Domínicæ occurréntis, quæ habétur infra post Lectiónes últimæ hebdómadæ Novémbris cum áliis Oratiónibus usque ad Domínicam XXIV post Pentecósten.}

\feriaheading{Dominica I Septembris}

\rubric{Semiduplex}

\hourheading{Ad Matutinum}

\rubric{Invitatórium} Dóminum, qui fecit, \rubric{et Hymnus} Nocte surgéntes, \rubric{ut in Psaltério. 31; quæ dicúntur in ómnibus Domínicis usque ad diem 27 Septémbris inclusíve occurréntibus.}

\nocturnoheading{In I Nocturno}

\rubric{Incipit liber Job}

\lectioheading{Lectio i}{Cap. 1, 1-11}

\lettrine[lines=2]{V}{IR} erat in terra Hus nómine Job, et erat vir ille simplex et rectus ac timens Deum et recédens a malo. Natíque sunt ei septem fílii et tres fíliæ. Et fuit posséssio ejus septem míllia óvium et tria míllia camelórum, quingénta quoque juga boum et quingéntæ ásinæ ac família multa nimis: erátque vir ille magnus inter omnes Orientáles.

\R\ Si bona suscépimus de manu Dei, mala autem quare non sustineámus? * Dóminus dedit, Dóminus ábstulit; sicut Dómino plácuit, ita factum est: sit nomen Dómini benedíctum. \V\ Nudus egréssus sum de útero matris meæ et nudus revértar illuc. — Dóminus.

\lectioheading{Lectio ii}{}

\lettrine[lines=2]{E}{T} ibant fílii ejus et faciébant convívium per domos, unusquísque in die suo. Et mitténtes vocábant tres soróres suas, ut coméderent et bíberent cum eis. Cumque in orbem transíssent dies convívii, mittébat ad eos Job et sanctificábat illos, consurgénsque dilúculo offerébat holocáusta pro síngulis; dicébat enim: Ne forte peccáverint fílii mei et benedíxerint Deo in córdibus suis. Sic faciébat Job cunctis diébus.

\R\ Antequam cómedam suspíro, et tamquam inundántes aquæ sic rugítus meus; quia timor, quem timébam, evénit mihi, et quod verébar áccidit. Nonne dissimulávi? nonne sílui? nonne quiévi? * Et venit super me indignátio. \V\ Ecce non est auxílium mihi in me, et necessárii quoque mei recessérunt a me. — Et.

\lectioheading{Lectio iii}{}

\lettrine[lines=2]{Q}{UADAM} autem die, cum veníssent fílii Dei ut assísterent coram Dómino, affúit inter eos étiam Satan. Cui dixit Dóminus: Unde venis? Qui respóndens ait: Circuívi terram et perambulávi eam. Dixítque Dóminus ad eum: Numquid considerásti servum meum Job, quod non sit ei símilis in terra, homo símplex et rectus ac timens Deum et recédens a malo? Cui respóndens Satan ait: Numquid Job frustra timet Deum? nonne tu vallásti eum ac domum ejus, universámque substántiam per circúitum opéribus mánuum ejus benedixísti, et posséssio ejus crevit in terra? Sed exténde páululum manum tuam et tange cuncta quæ póssidet, nisi in fáciem benedíxerit tibi.

\R\ Quare detraxístis sermónibus veritátis? ad increpándum verba compónitis et subvértere nitímini amícum vestrum: * Verúmtamen quæ cogitástis, expléte. \V\ Quod justum est, judicáte; et non inveniétis in lingua mea iniquitátem. — Verúmtamen. Glória Patri. Verúmtamen.

\nocturnoheading{In II Nocturno}

\rubric{Ex libro Morálium sancti Gregórii Papæ}

\lectioheading{Lectio iv}{Lib. 2, cap. 1}

\lettrine[lines=2]{S}{CRIPTÚRA} sacra mentis óculis quasi quoddam spéculum oppónitur, ut intérna nostra fácies in ipsa videátur. Ibi étenim foeda, ibi pulchra nostra cognóscimus: ibi sentímus quantum profícimus, ibi a proféctu quam longe distámus. Narrat autem gesta Sanctórum, et ad imitatiónem corda próvocat infirmórum; dumque illórum victrícia facta commémorat, contra vitiórum prǽlia, debília nostra confírmat: fitque verbis illíus, ut eo mens minus inter certámina trépidet, quo ante se pósitos tot virórum fórtium triúmphos videt.

\R\ Indúta est caro mea putrédine, et sórdibus púlveris cutis mea áruit et contrácta est: * Meménto mei, Dómine, quóniam ventus est vita mea. \V\ Dies mei velócius transiérunt quam a texénte tela succíditur, et consúmpti sunt absque ulla spe. — Meménto.

\lectioheading{Lectio v}{}

\lettrine[lines=2]{N}{ONNÚMQUAM} vero non solum nobis eórum virtútes ásserit, sed étiam casus innotéscit; ut et in victória fórtium, quod imitándo arrípere, et rursum videámus in lápsibus quid debeámus timére. Ecce enim Job describítur tentatióne auctus, sed David tentatióne prostrátus; ut et majórum virtus spem nostram fóveat, et majórum casus ad cautélam nos humilitátis accíngat: quátenus dum illa gaudéntes súblevant, ista metuéntes premant; et audiéntis ánimus illinc spei fidúcia, hinc humilitáte erudítus, nec temeritáte supérbiat, quia ad spei fidúciam virtútis exémplo roborátur.

\R\ Páucitas diérum meórum finiétur brevi; dimítte me, Dómine, ut plangam páululum dolórem meum, * Antequam vadam et non revértar, ad terram tenebrósam et opértam mortis calígine. \V\ Manus tuæ, Dómine, fecérunt me, et plasmavérunt me totum in circúitu; et sic repénte præcípitas me? — Antequam.

\lectioheading{Lectio vi}{Lib. 1, cap. 1}

\lettrine[lines=2]{V}{IR} erat in terra Hus, nómine Job. Idcírco sanctus vir ubi habitáverit dícitur, ut ejus méritum virtútis exprimátur. Hus namque quis nésciat quod sit in terra Gentílium? Gentílitas autem eo obligáta vítiis exstítit, quo cognitiónem sui Condítoris ignorávit. Dicátur ítaque ubi habitáverit; ut hoc ejus láudibus profíciat, quod bonus inter malos fuit. Neque enim valde laudábile est, bonum esse cum bonis, sed bonum esse cum malis. Sicut enim gravióris culpæ est, inter bonos bonum non esse; ita imménsi est præcónii, bonum étiam inter malos exstitísse.

\R\ Non abscóndas me, Dómine, a fácie tua: manum tuam longe fac a me, * Et formído tua non me térreat. \V\ Córripe me, Dómine, in misericórdia, non in furóre tuo, ne forte ad níhilum rédigas me. — Et. Glória Patri. Et.

\rubric{Lectiónes III Nocturni de Homilía Domínicæ occurréntis post Pentecósten habéntur inférius circa finem hujus Partis de Témpore, in órdine aliárum Dominicárum usque ad Domínicam XXIV post Pentecósten. Simíliter Antíphonæ ad Benedíctus et Magníficat, et Oratiónes.}

\R\ vii Quis mihi tríbuat, ut in inférno prótegas me et abscóndas me, donec pertránseat furor tuus, Dómine, nisi tu, qui solus es Deus? * Et constítuas mihi tempus, in quo recordéris mei? \V\ Numquid sicut dies hóminis dies tui, ut quæras iniquitátem meam; cum sit nemo, qui de manu tua possit erúere? — Et.

\R\ viii Duo Séraphim, \rubric{ut infra in Homilíis.}

\rubric{Ad Laudes, in hac et sequéntibus Domínicis usque ad diem 27 Septémbris inclusíve occurréntibus, dícitur Hymnus} Ecce jam noctis, \rubric{ut in Psaltério. 40}

% ===== BATCH 1 of the Proprium de Tempore (autumnalis.md lines
% 13652-13730) -- genuinely new Autumn-specific content, transcribed
% fresh from autumnalis.md since the Proprium de Tempore is NOT
% shared boilerplate across Pars volumes (confirmed by the previous
% batch's diff work showing aestiva.tex's parallel section opens
% with Trinity Sunday rather than this volume's own starting point).
% Covers the brief "Sabbato proximo Kalendis Septembris" (First
% Vespers only, with its Magnificat antiphon from the Book of Job
% and a forward-pointing rubric to the proper Sunday Oratio given
% later in this Part) and the complete "Dominica I Septembris"
% Matins: the Invitatorium/Hymnus cross-reference (rendered as
% alternating \rubric{} labels and live incipit text per the
% established compact-citation convention), the six Lessons and six
% Responsories of Nocturns I-II (Job's affliction and the opening of
% St. Gregory's Moralia on Job), and the final two numbered
% Responsories (vii, viii) that stand in for Nocturn III, whose own
% Lessons, antiphons, and Orationes are deferred by the source's own
% rubric to the Homiliae section near the end of this Part (a
% structural forward-reference, not a page/footnote number, so no
% \pageref{} resolution applies; it will be transcribed fresh when
% that later section is reached). The em-dash "—" repeat-cue
% notation in each Responsory (e.g. "— Dóminus.") is preserved
% literally as the source's own shorthand for repeating the
% refrain's cue word, a notation not previously encountered in this
% project's established responsory-repetition conventions (which
% elsewhere spell the repeated word out in full); since it is
% genuinely how autumnalis.md itself abbreviates these responsories,
% it is kept exactly as printed rather than silently expanded. The
% "---" horizontal rule before "Dominica I Septembris" is omitted as
% a scan artifact. The document continues with "Feria Secunda infra
% Hebdomadam I Septembris" (autumnalis.md line 13731).
% =====
\feriaheading{Feria Secunda infra Hebdomadam I Septembris}

\rubric{De libro Job}

\lectioheading{Lectio i}{Cap. 1, 13-22}

\lettrine[lines=2]{C}{UM} autem quadam die fílii et fíliæ ejus coméderent et bíberent vinum in domo fratris sui primogéniti, núntius venit ad Job, qui díceret: Boves arábant et ásinæ pascebántur juxta eos, et irruérunt Sabæi tulerúntque ómnia et púeros percussérunt gládio, et evási ego solus ut nuntiárem tibi. Cumque adhuc ille loquerétur, venit alter et dixit: Ignis Dei cécidit e cælo et tactas oves puerósque consúmpsit, et effúgi ego solus ut nuntiárem tibi.

\R\ Versa est in luctum cíthara mea, et órganum meum in vocem fléntium: * Parce mihi, Dómine, nihil enim sunt dies mei. \V\ Cutis mea denigráta est super me, et ossa mea aruérunt. — Parce.

\lectioheading{Lectio ii}{}

\lettrine[lines=2]{S}{ED}, et illo adhuc loquénte, venit álius et dixit: Chaldǽi fecérunt tres turmas et invasérunt camélos et tulérunt eos, necnon et púeros percussérunt gládio, et ego fugi solus ut nuntiárem tibi. Adhuc loquebátur ille, et ecce álius intrávit et dixit: Fíliis tuis et filiábus vescéntibus et bibéntibus vinum in domo fratris sui primogéniti, repénte ventus veheméns írruit a regióne desérti et concússit quátuor ángulos domus, quæ córruens oppréssit líberos tuos, et mórtui sunt, et effúgi ego solus ut nuntiárem tibi.

\R\ Utinam appenderéntur peccáta mea, quibus iram mérui, * Et calamítas, quam pátior, in statéra. \V\ Quasi aréna maris hæc grávior apparéret, unde et verba mea dolóre sunt plena. — Et.

\rubric{In Octávis et Festis ritus simplícis: Glória Patri. Et.}

\lectioheading{Lectio iii}{}

\lettrine[lines=2]{T}{UNC} surréxit Job et scidit vestiménta sua et, tonso cápite, córruens in terram adorávit et dixit: Nudus egréssus sum de útero matris meæ et nudus revértar illuc. Dóminus dedit, Dóminus ábstulit; sicut Dómino plácuit, ita factum est: sit nomen Dómini benedíctum. In ómnibus his non peccávit Job lábiis suis neque stultum quid contra Deum locútus est.

\R\ Quare detraxístis sermónibus veritátis? ad increpándum verba compónitis et subvértere nitímini amícum vestrum: * Verúmtamen quæ cogitástis, expléte. \V\ Quod justum est, judicáte; et non inveniétis in lingua mea iniquitátem. — Verúmtamen. Glória Patri. Verúmtamen.

\feriaheading{Feria Tertia infra Hebdomadam I Septembris}

\rubric{De libro Job}

\lectioheading{Lectio i}{Cap. 2, 1-13}

\lettrine[lines=2]{F}{ACTUM} est autem, cum quadam die veníssent fílii Dei et starent coram Dómino, venísset quoque Satan inter eos et staret in conspéctu ejus, ut díceret Dóminus ad Satan: Unde venis? Qui respóndens ait: Circuívi terram et perambulávi eam. Et dixit Dóminus ad Satan: Numquid considerásti servum meum Job, quod non sit ei símilis in terra, vir símplex et rectus ac timens Deum et recédens a malo et adhuc rétinens innocéntiam? Tu autem commovísti me advérsus eum, ut afflígerem eum frustra. Cui respóndens Satan ait: Pellem pro pelle, et cuncta, quæ habet homo, dabit pro ánima sua. Alióqui mitte manum tuam et tange os ejus et carnem; et tunc vidébis quod in fáciem benedícat tibi.

\R\ Indúta est caro mea putrédine, et sórdibus púlveris cutis mea áruit et contrácta est: * Meménto mei, Dómine, quóniam ventus est vita mea. \V\ Dies mei velócius transiérunt quam a texénte tela succíditur, et consúmpti sunt absque ulla spe. — Meménto.

\lectioheading{Lectio ii}{}

\lettrine[lines=2]{D}{IXIT} ergo Dóminus ad Satan: Ecce in manu tua est, verúmtamen ánimam illíus serva. Egréssus ígitur Satan a fácie Dómini percússit Job úlcere péssimo a planta pedis usque ad vérticem ejus; qui testa sániem radébat sedens in sterquilínio. Dixit autem illi uxor sua: Adhuc tu pérmanes in simplicitáte tua? Bénedic Deo et mórere. Qui ait ad illam: Quasi una de stultis muliéribus locúta es: si bona suscépimus de manu Dei, mala quare non suscipiámus? In ómnibus his non peccávit Job lábiis suis.

\R\ Páucitas diérum meórum finiétur brevi; dimítte me, Dómine, ut plangam páululum dolórem meum, * Antequam vadam ad terram tenebrósam et opértam mortis calígine. \V\ Manus tuæ, Dómine, fecérunt me, et plasmavérunt me totum in circúitu; et sic repénte præcípitas me? — Antequam.

\rubric{In Octávis et Festis ritus simplícis: Glória Patri. Antequam.}

\lectioheading{Lectio iii}{}

\lettrine[lines=2]{I}{GITUR} audiéntes tres amíci Job omne malum quod accidísset ei, venérunt sínguli de loco suo, Elíphaz Themanítes et Báldad Suhítes et Sophar Naamathítes. Condíxerant enim ut páriter veniéntes visitárent eum et consolaréntur. Cumque elevássent procul óculos suos, non cognovérunt eum et exclamántes ploravérunt, scissísque véstibus, spargérunt púlverem super caput suum in cælum. Et sedérunt cum eo in terra septem diébus et septem nóctibus, et nemo loquebátur ei verbum; vidébant enim dolórem esse veheméntem.

\R\ Non abscóndas me, Dómine, a fácie tua: manum tuam longe fac a me, * Et formído tua non me térreat. \V\ Córripe me, Dómine, in misericórdia, non in furóre tuo, ne forte ad níhilum rédigas me. — Et. Glória Patri. Et.

\feriaheading{Feria Quarta infra Hebdomadam I Septembris}

\rubric{De libro Job}

\lectioheading{Lectio i}{Cap. 3, 1-16}

\lettrine[lines=2]{P}{OST} hæc apéruit Job os suum et maledíxit diéi suo et locútus est: Péreat dies in qua natus sum, et nox in qua dictum est: Concéptus est homo. Dies ille vertátur in ténebras: non requírat eum Deus désuper, et non illustrétur lúmine; obscúrent eum ténebræ et umbra mortis; óccupet eum calígo, et involvátur amaritúdine.

\R\ Quis mihi tríbuat, ut in inférno prótegas me et abscóndas me, donec pertránseat furor tuus, Dómine, nisi tu, qui solus es Deus? * Et constítuas mihi tempus, in quo recordéris mei? \V\ Numquid sicut dies hóminis dies tui, ut quæras iniquitátem meam; cum sit nemo, qui de manu tua possit erúere? — Et.

\lectioheading{Lectio ii}{}

\lettrine[lines=2]{N}{OCTEM} illam tenebrósus turbo possídeat; non computétur in diébus anni, nec numerétur in ménsibus. Sit nox illa solitária nec laude digna; maledícant ei qui maledícunt diéi, qui paráti sunt suscitáre leviáthan. Obtenebréntur stellæ calígine ejus, exspéctet lucem et non vídeat, nec ortum surgéntis auróræ; quia non conclúsit óstia ventris, qui portávit me, nec ábstulit mala ab óculis meis.

\R\ Utinam appenderéntur peccáta mea, quibus iram mérui, * Et calamítas, quam pátior, in statéra. \V\ Quasi aréna maris hæc grávior apparéret, unde et verba mea dolóre sunt plena. — Et.

\rubric{In Octávis et Festis ritus simplícis: Glória Patri. Et.}

\lectioheading{Lectio iii}{}

\lettrine[lines=2]{Q}{UARE} non in vulva mórtuus sum? egréssus ex útero non statim périi? quare excéptus génibus? cur lactátus ubéribus? Nunc enim dórmiens silérem et somno meo requiéscerem cum régibus et consúlibus terræ, qui ædíficant sibi solitúdines, aut cum princípibus, qui póssident aurum et replent domos suas argénto; aut sicut abortívum abscónditum non subsísterem, vel qui concépti non vidérunt lucem.

\R\ Quare detraxístis sermónibus veritátis? ad increpándum verba compónitis et subvértere nitímini amícum vestrum: * Verúmtamen quæ cogitástis, expléte. \V\ Quod justum est, judicáte; et non inveniétis in lingua mea iniquitátem. — Verúmtamen. Glória Patri. Verúmtamen.

\feriaheading{Feria Quinta infra Hebdomadam I Septembris}

\rubric{De libro Job}

\lectioheading{Lectio i}{Cap. 4, 1-18}

\lettrine[lines=2]{R}{ESPÓNDENS} autem Elíphaz Themanítes dixit: Si cœpérimus loqui tibi, fórsitan moléste accípies; sed concéptum sermónem tenére quis póterit? Ecce docuísti multos et manus lassas roborásti, vacillántes confirmavérunt sermónes tui et génua treméntia confortásti. Nunc autem venit super te plaga, et defecísti; tétigit te, et conturbátus es. Ubi est timor tuus, fortitúdo tua, patiéntia tua et perféctio viárum tuárum?

\R\ Si bona suscépimus de manu Dei, mala autem quare non sustineámus? * Dóminus dedit, Dóminus ábstulit; sicut Dómino plácuit, ita factum est: sit nomen Dómini benedíctum. \V\ Nudus egréssus sum de útero matris meæ et nudus revértar illuc. — Dóminus.

\lectioheading{Lectio ii}{}

\lettrine[lines=2]{R}{ECORDÁRE}, óbsecro te, quis umquam ínnocens périit? aut quando recti deléti sunt? Quin pótius vidi eos, qui operántur iniquitátem et séminant dolóres et metunt eos, flante Deo perísse et spíritu iræ ejus esse consúmptos. Rugítus leónis et vox leǽnæ et dentes catulórum leónum contríti sunt; tigris périit, eo quod non habéret prædam, et cátuli leónis dissipáti sunt.

\R\ Antequam cómedam suspíro, et tamquam inundántes aquæ sic rugítus meus; quia timor, quem timébam, evénit mihi, et quod verébar áccidit. Nonne dissimulávi? nonne sílui? nonne quiévi? * Et venit super me indignátio. \V\ Ecce non est auxílium mihi in me, et necessárii quoque mei recessérunt a me. — Et.

\rubric{In Octávis et Festis ritus simplícis: Glória Patri. Et.}

\lectioheading{Lectio iii}{}

\lettrine[lines=2]{P}{ORRO} ad me dictum est verbum abscónditum, et quasi furtíve suscépit auris mea venas susúrri ejus. In horróre visiónis noctúrnæ, quando solet sopor occupáre hómines, pavor ténuit me et tremor, et ómnia ossa mea pertérrita sunt, et cum spíritus, me præsénte, transíret, inhorruérunt pili carnis meæ. Stetit quidam, cujus non agnoscébam vultum, imágo coram óculis meis, et vocem quasi auræ lenis audívi: Numquid homo Dei comparatióne justificábitur? aut Factóre suo púrior erit vir? Ecce qui sérviunt ei non sunt stábiles, et in Angelis suis réperit pravitátem.

\R\ Quare detraxístis sermónibus veritátis? ad increpándum verba compónitis et subvértere nitímini amícum vestrum: * Verúmtamen quæ cogitástis, expléte. \V\ Quod justum est, judicáte; et non inveniétis in lingua mea iniquitátem. — Verúmtamen. Glória Patri. Verúmtamen.

\feriaheading{Feria Sexta infra Hebdomadam I Septembris}

\rubric{De libro Job}

\lectioheading{Lectio i}{Cap. 6, 1-13}

\lettrine[lines=2]{R}{ESPÓNDENS} autem Job dixit: Utinam appenderéntur peccáta mea, quibus iram mérui, et calamítas, quam pátior, in statéra. Quasi aréna maris hæc grávior apparéret, unde et verba mea dolóre sunt plena; quia sagíttæ Dómini in me sunt, quarum indignátio ébibit spíritum meum, et terróres Dómini mílitant contra me.

\R\ Indúta est caro mea putrédine, et sórdibus púlveris cutis mea áruit et contrácta est: * Meménto mei, Dómine, quóniam ventus est vita mea. \V\ Dies mei velócius transiérunt quam a texénte tela succíditur, et consúmpti sunt absque ulla spe. — Meménto.

\lectioheading{Lectio ii}{}

\lettrine[lines=2]{N}{UMQUID} rúgiet ónager cum habúerit herbam? aut mugíet bos, cum ante præsépe plenum stéterit? Aut póterit comédi insúlsum, quod non est sale condítum? aut potest áliquis gustáre, quod gustátum affert mortem? Quæ prius nolébat tángere ánima mea, nunc præ angústia cibi mei sunt.

\R\ Páucitas diérum meórum finiétur brevi; dimítte me, Dómine, ut plangam páululum dolórem meum, * Antequam vadam ad terram tenebrósam et opértam mortis calígine. \V\ Manus tuæ, Dómine, fecérunt me, et plasmavérunt me totum in circúitu; et sic repénte præcípitas me? — Antequam.

\rubric{In Octávis et Festis ritus simplícis: Glória Patri. Antequam.}

\lectioheading{Lectio iii}{}

\lettrine[lines=2]{Q}{UIS} det ut véniat petítio mea, et quod exspécto tríbuat mihi Deus? Et qui cœpit ipse me cónterat, solvat manum suam et succídat me? Et hæc mihi sit consolátio, ut afflígens me dolóre non parcat, nec contradícam sermónibus Sancti. Quæ est enim fortitúdo mea, ut sustíneam? aut quis finis meus, ut patiénter agam? Nec fortitúdo lápidum fortitúdo mea, nec caro mea ǽnea est. Ecce non est auxílium mihi in me, et necessárii quoque mei recessérunt a me.

\R\ Non abscóndas me, Dómine, a fácie tua: manum tuam longe fac a me, * Et formído tua non me térreat. \V\ Córripe me, Dómine, in misericórdia, non in furóre tuo, ne forte ad níhilum rédigas me. — Et. Glória Patri. Et.

\rubric{Ad Vésperas, nisi sequénti die persolvéndum sit Offícium de Festo novem Lectiónum, aut de qualíbet Octáva vel de Vigília, a Capítulo fit de sancta María, et de ea dícitur Offícium in Sábbato sequénti. Quod ítem servátur in ómnibus Sábbatis usque ad Sábbatum ante Domínicam I Advéntus inclusíve, præterquam in Sábbato Quátuor Témporum Septémbris, et in Sábbato, in quo anticipánda sit Domínica XXIII post Pentecósten.}

\feriaheading{Sabbato infra Hebdomadam I Septembris}

\rubric{De libro Job}

\lectioheading{Lectio i}{Cap. 7, 1-12}

\lettrine[lines=2]{M}{ILITIA} est vita hóminis super terram, et sicut dies mercenárii, dies ejus. Sicut servus desíderat umbram, et sicut mercenárius præstolátur finem óperis sui, sic et ego hábui menses vácuos et noctes laboriósas enumerávi mihi. Si dormíero, dicam: Quando consúrgam? Et rursum exspectábo vésperam, et replébor dolóribus usque ad ténebras.

\R\ Quis mihi tríbuat, ut in inférno prótegas me et abscóndas me, donec pertránseat furor tuus, Dómine, nisi tu, qui solus es Deus? * Et constítuas mihi tempus, in quo recordéris mei? \V\ Numquid sicut dies hóminis dies tui, ut quæras iniquitátem meam; cum sit nemo, qui de manu tua possit erúere? — Et.

\lectioheading{Lectio ii}{}

\lettrine[lines=2]{I}{NDÚTA} est caro mea putrédine, et sórdibus púlveris cutis mea áruit et contrácta est. Dies mei velócius transiérunt quam a texénte tela succíditur, et consúmpti sunt absque ulla spe. Meménto quia ventus est vita mea, et non revertétur óculus meus ut vídeat bona. Nec aspíciet me visus hóminis; óculi tui in me, et non subsístam.

\R\ Utinam appenderéntur peccáta mea, quibus iram mérui, * Et calamítas, quam pátior, in statéra. \V\ Quasi aréna maris hæc grávior apparéret, unde et verba mea dolóre sunt plena. — Et.

\rubric{In Octávis simplícibus et in Offício sanctæ Maríæ in Sábbato: Glória Patri. Et.}

\lectioheading{Lectio iii}{}

\lettrine[lines=2]{S}{ICUT} consúmitur nubes et pertránsit, sic qui descénderit ad ínferos non ascéndet; nec revertétur ultra in domum suam, neque cognóscet eum ámplius locus ejus. Quaprópter et ego non parcam ori meo; loquar in tribulatióne spíritus mei, confabulábor cum amaritúdine ánimæ meæ. Numquid mare ego sum aut cetus, quia circumdedísti me cárcere?

\R\ Quare detraxístis sermónibus veritátis? ad increpándum verba compónitis et subvértere nitímini amícum vestrum: * Verúmtamen quæ cogitástis, expléte. \V\ Quod justum est, judicáte; et non inveniétis in lingua mea iniquitátem. — Verúmtamen. Glória Patri. Verúmtamen.

\hourheading{Ad Vesperas}

\V\ Vespertína orátio ascéndat ad te, Dómine. \R\ Et descéndat super nos misericórdia tua.

\antlabel{Ad Magnif. Ant.} In ómnibus his * non peccávit Job lábiis suis neque stultum áliquid contra Deum locútus est.

\rubric{Orátio Domínicæ occurréntis.}

% ===== BATCH 2 of the Proprium de Tempore (autumnalis.md lines
% 13731-13902) -- completes the first full week of September: the
% six ferial days' Matins lessons (three Lessons and three
% Responsories apiece) continuing the Book of Job cycle through
% chapters 1-7, and the closing Vespers for "Sabbato infra
% Hebdomadam I Septembris" with its own Magnificat antiphon and a
% forward-pointing rubric to the proper Sunday Oratio. The lengthy
% governing rubric on Saturday Vespers switching to the Office of
% the Blessed Virgin Mary (except on the September Ember Saturday
% and certain anticipated Sundays) is transcribed in full as a
% single \rubric{} block. Three "In Octávis..." rubrics mark the
% alternate "Glória Patri" cue for feast-day/Marian-Saturday ranks;
% each is transcribed exactly as the source gives it, including the
% third day's distinct wording ("et in Officio sanctae Mariae in
% Sabbato" rather than "et Festis ritus simplicis"). Every
% Responsory's em-dash repeat-cue is preserved literally, consistent
% with Batch 1. The "---" horizontal rules before "Sabbato infra
% Hebdomadam I Septembris" and before "Dominica II Septembris" are
% omitted as scan artifacts. The document continues with "Dominica
% II Septembris" (autumnalis.md line 13903).
% =====
\feriaheading{Dominica II Septembris}

\rubric{Semiduplex}

\nocturnoheading{In I Nocturno}

\rubric{De libro Job}

\lectioheading{Lectio i}{Cap. 9, 1-17}

\lettrine[lines=2]{E}{T} respóndens Job ait: Vere scio quod ita sit, et quod non justificétur homo compósitus Deo. Si volúerit conténdere cum eo, non póterit ei respondére unum pro mille. Sápiens corde est et fortis róbore. Quis réstitit ei, et pacem hábuit? Qui tránstulit montes, et nesciérunt hi quos subvértit in furóre suo.

\R\ Si bona suscépimus de manu Dei, mala autem quare non sustineámus? * Dóminus dedit, Dóminus ábstulit; sicut Dómino plácuit, ita factum est: sit nomen Dómini benedíctum. \V\ Nudus egréssus sum de útero matris meæ et nudus revértar illuc. — Dóminus.

\lectioheading{Lectio ii}{}

\lettrine[lines=2]{Q}{UI} cómmovet terram de loco suo, et colúmnæ ejus concutiúntur; qui præcípit soli, et non óritur, et stellas claudit quasi sub signáculo; qui exténdit cælos solus, et gráditur super fluctus maris; qui facit Arctúrum et Oríona et Hýadas et interióra Austri; qui facit magna et incomprehensibília et mirabília, quorum non est númerus.

\R\ Antequam cómedam suspíro, et tamquam inundántes aquæ sic rugítus meus; quia timor, quem timébam, evénit mihi, et quod verébar áccidit. Nonne dissimulávi? nonne sílui? nonne quiévi? * Et venit super me indignátio. \V\ Ecce non est auxílium mihi in me, et necessárii quoque mei recessérunt a me. — Et.

\lectioheading{Lectio iii}{}

\lettrine[lines=2]{S}{I} vénerit ad me, non vidébo eum; si abíerit, non intélligam; si repénte intérroget, quis respondébit ei vel quis dícere potest: Cur ita facis? Deus, cujus iræ nemo resístere potest, et sub quo curvántur qui portant orbem; quantus ergo sum ego, ut respóndeam ei, et loquar verbis meis cum eo? Qui, étiam si habúero quíppiam justum, non respondébo, sed meum júdicem deprecábor. Et, cum invocántem exaudíerit me, non credo quod audíerit vocem meam. In túrbine enim cónteret me, et multiplicábit vúlnera mea, étiam sine causa.

\R\ Quare detraxístis sermónibus veritátis? ad increpándum verba compónitis et subvértere nitímini amícum vestrum: * Verúmtamen quæ cogitástis, expléte. \V\ Quod justum est, judicáte; et non inveniétis in lingua mea iniquitátem. — Verúmtamen. Glória Patri. Verúmtamen.

\nocturnoheading{In II Nocturno}

\rubric{Ex libro Morálium sancti Gregórii Papæ}

\lectioheading{Lectio iv}{Liber 9, cap. 2}

\lettrine[lines=2]{V}{ERE} scio quod ita sit, et quod non justificábitur homo compósitus Deo. Homo quippe Deo non compósitus justítiam pércipit, compósitus amíttit; quia quisquis se auctóri bonórum cómparat, bono se, quod accéperat, privat. Qui enim accépta bona sibi árrogat, suis contra Deum donis pugnat. Unde ergo despéctus erígitur, dignum est ut eréctus inde destruátur. Sanctus autem vir, quia omne virtútis nostræ méritum esse vítium cónspicit, si ab intérno árbitro distrícte judicétur, recte subjúngit: Si volúerit conténdere cum eo, non póterit respondére ei unum pro mille.

\R\ Indúta est caro mea putrédine, et sórdibus púlveris cutis mea áruit et contrácta est: * Meménto mei, Dómine, quóniam ventus est vita mea. \V\ Dies mei velócius transiérunt quam a texénte tela succíditur, et consúmpti sunt absque ulla spe. — Meménto.

\lectioheading{Lectio v}{}

\lettrine[lines=2]{I}{N} Scriptúra sancta millenárius númerus pro universitáte solet intélligi. Hinc étenim Psalmísta ait: Verbi, quod mandávit in mille generatiónes: cum profécto constet, quod ab ipso mundi exórdio usque ad Redemptóris advéntum per Evangelístam non ámplius quam septuagínta et septem propágines numeréntur. Quid ergo in millenário número nisi ad proferéndam novam sóbolem perfécta universitas præscítæ generatiónis exprímitur? Hinc et per Joánnem dícitur: Et regnábunt cum eo mille annis; quia vidélicet regnum sanctæ Ecclésiæ universitátis perfectióne solidátur.

\R\ Páucitas diérum meórum finiétur brevi; dimítte me, Dómine, ut plangam páululum dolórem meum, * Antequam vadam ad terram tenebrósam et opértam mortis calígine. \V\ Manus tuæ, Dómine, fecérunt me, et plasmavérunt me totum in circúitu; et sic repénte præcípitas me? — Antequam.

\lectioheading{Lectio vi}{}

\lettrine[lines=2]{Q}{UIA} vero monas décies multiplicáta in denárium dúcitur, denárius per semetípsum ductus in centenárium dilatátur, qui rursus per denárium ductus in millenárium ténditur; cum ab uno incípimus, ut ad millenárium veniámus, quid hoc loco uníus appellatióne, nisi bene vivéndi inítium? quid millenárii númeri amplitúdine, nisi ejúsdem bonæ vitæ perféctio designátur? Cum Deo autem conténdere, est, non ei tribúere, sed sibi glóriam suæ virtútis arrogáre. Sed sanctus vir conspíciat, quia et qui summa jam dona percépit, si de accéptis extóllitur, cuncta quæ accéperat, amíttit.

\R\ Non abscóndas me, Dómine, a fácie tua: manum tuam longe fac a me, * Et formído tua non me térreat. \V\ Córripe me, Dómine, in misericórdia, non in furóre tuo, ne forte ad níhilum rédigas me. — Et. Glória Patri. Et.

\R\ vii Quis mihi tríbuat, ut in inférno prótegas me et abscóndas me, donec pertránseat furor tuus, Dómine, nisi tu, qui solus es Deus? * Et constítuas mihi tempus, in quo recordéris mei? \V\ Numquid sicut dies hóminis dies tui, ut quæras iniquitátem meam; cum sit nemo, qui de manu tua possit erúere? — Et.

\R\ viii Duo Séraphim, \rubric{ut infra in Homilíis.}

\feriaheading{Feria Secunda infra Hebdomadam II Septembris}

\rubric{De libro Job}

\lectioheading{Lectio i}{Cap. 27, 1-15}

\lettrine[lines=2]{A}{DDIDIT} quoque Job, assúmens parábolam suam, et dixit: Vivit Deus, qui ábstulit judícium meum, et Omnípotens, qui ad amaritúdinem addúxit ánimam meam. Quia, donec súperest hálitus in me, et spíritus Dei in náribus meis, non loquéntur lábia mea iniquitátem, nec lingua mea meditábitur mendácium. Absit a me, ut justos vos esse júdicem; donec defíciam, non recédam ab innocéntia mea.

\R\ Versa est in luctum cíthara mea, et órganum meum in vocem fléntium: * Parce mihi, Dómine, nihil enim sunt dies mei. \V\ Cutis mea denigráta est super me, et ossa mea aruérunt. — Parce.

\lectioheading{Lectio ii}{}

\lettrine[lines=2]{J}{USTIFICATIÓNEM} meam, quam cœpi tenére, non déseram; neque enim reprehéndit me cor meum in omni vita mea. Sit ut ímpius inimícus meus, et adversárius meus quasi iníquus. Quæ est enim spes hypócritæ, si aváre rápiat, et non líberet Deus ánimam ejus? Numquid Deus áudiet clamórem ejus, cum vénerit super eum angústia? Aut póterit in Omnipoténte delectári, et invocáre Deum omni témpore?

\R\ Utinam appenderéntur peccáta mea, quibus iram mérui, * Et calamítas, quam pátior, in statéra. \V\ Quasi aréna maris hæc grávior apparéret, unde et verba mea dolóre sunt plena. — Et.

\rubric{In Octávis et Festis ritus simplícis: Glória Patri. Et.}

\lectioheading{Lectio iii}{}

\lettrine[lines=2]{D}{OCÉBO} vos per manum Dei quæ Omnípotens hábeat, nec abscóndam. Ecce vos omnes nostis, et quid sine causa vana loquímini? Hæc est pars hóminis ímpii apud Deum, et heréditas violentórum, quam ab Omnipoténte suscípient: si multiplicáti fúerint fílii ejus, in gládio erunt, et nepótes ejus non saturabúntur pane; qui réliqui fúerint ex eo sepeliéntur in intéritu, et víduæ illíus non plorábunt.

\R\ Quare detraxístis sermónibus veritátis? ad increpándum verba compónitis et subvértere nitímini amícum vestrum: * Verúmtamen quæ cogitástis, expléte. \V\ Quod justum est, judicáte; et non inveniétis in lingua mea iniquitátem. — Verúmtamen. Glória Patri. Verúmtamen.

\feriaheading{Feria Tertia infra Hebdomadam II Septembris}

\rubric{De libro Job}

\lectioheading{Lectio i}{Cap. 28, 12-28}

\lettrine[lines=2]{S}{APIÉNTIA} ubi invenítur? et quis est locus intelligéntiæ? Nescit homo prétium ejus, nec invenítur in terra suáviter vivéntium. Abýssus dicit: Non est in me, et mare lóquitur: Non est mecum. Non dábitur aurum obrízum pro ea, nec appendétur argéntum in commutatióne ejus; non conferétur tinctis Indiæ colóribus, nec lápidi sardónycho pretiosíssimo vel sapphíro.

\R\ Indúta est caro mea putrédine, et sórdibus púlveris cutis mea áruit et contrácta est: * Meménto mei, Dómine, quóniam ventus est vita mea. \V\ Dies mei velócius transiérunt quam a texénte tela succíditur, et consúmpti sunt absque ulla spe. — Meménto.

\lectioheading{Lectio ii}{}

\lettrine[lines=2]{N}{ON} adæquábitur ei aurum vel vitrum, nec commutabúntur pro ea vasa auri. Excélsa et eminéntia non memorabúntur comparatióne ejus; tráhitur autem sapiéntia de occúltis. Non adæquábitur ei topázius de Æthiópia, nec tinctúræ mundíssimæ componétur. Unde ergo sapiéntia venit? et quis est locus intelligéntiæ? Abscóndita est ab óculis ómnium vivéntium; vólucres quoque cæli latet. Perdítio et mors dixérunt: Áuribus nostris audívimus famam ejus.

\R\ Páucitas diérum meórum finiétur brevi; dimítte me, Dómine, ut plangam páululum dolórem meum, * Antequam vadam ad terram tenebrósam et opértam mortis calígine. \V\ Manus tuæ, Dómine, fecérunt me, et plasmavérunt me totum in circúitu; et sic repénte præcípitas me? — Antequam.

\rubric{In Octávis et Festis ritus simplícis: Glória Patri. Antequam.}

\lectioheading{Lectio iii}{}

\lettrine[lines=2]{D}{EUS} intélligit viam ejus, et ipse novit locum illíus; ipse enim fines mundi intuétur, et ómnia, quæ sub cælo sunt, réspicit. Qui fecit ventis pondus, et aquas appéndit in mensúra. Quando ponébat plúviis legem et viam procéllis sonántibus, tunc vidit illam et enarrávit et præparávit et investigávit. Et dixit hómini: Ecce timor Dómini, ipsa est sapiéntia, et recédere a malo, intelligéntia.

\R\ Non abscóndas me, Dómine, a fácie tua: manum tuam longe fac a me, * Et formído tua non me térreat. \V\ Córripe me, Dómine, in misericórdia, non in furóre tuo, ne forte ad níhilum rédigas me. — Et. Glória Patri. Et.

\feriaheading{Feria Quarta infra Hebdomadam II Septembris}

\rubric{De libro Job}

\lectioheading{Lectio i}{Cap. 31, 1-18}

\lettrine[lines=2]{P}{EPIGI} fœdus cum óculis meis, ut ne cogitárem quidem de vírgine. Quam enim partem habéret in me Deus désuper, et hereditátem Omnípotens de excélsis? Numquid non perdítio est iníquo, et alienátio operántibus injustítiam? Nonne ipse considérat vias meas, et cunctos gressus meos dinúmerat? Si ambulávi in vanitáte, et festinávit in dolo pes meus, appéndat me in statéra justa, et sciat Deus simplicitátem meam.

\R\ Quis mihi tríbuat, ut in inférno prótegas me et abscóndas me, donec pertránseat furor tuus, Dómine, nisi tu, qui solus es Deus? * Et constítuas mihi tempus, in quo recordéris mei? \V\ Numquid sicut dies hóminis dies tui, ut quæras iniquitátem meam; cum sit nemo, qui de manu tua possit erúere? — Et.

\lectioheading{Lectio ii}{}

\lettrine[lines=2]{S}{I} declinávit gressus meus de via, et si secútum est óculos meos cor meum, et si mánibus meis adhǽsit mácula, seram, et álius cómedat, et progénies mea eradicétur. Si decéptum est cor meum super mulíere, et si ad óstium amíci mei insidiátus sum, scortum altérius sit uxor mea, et super illam incurvéntur álii. Hoc enim nefas est et iníquitas máxima; ignis est usque ad perditiónem dévorans et ómnia eradícans genímina.

\R\ Utinam appenderéntur peccáta mea, quibus iram mérui, * Et calamítas, quam pátior, in statéra. \V\ Quasi aréna maris hæc grávior apparéret, unde et verba mea dolóre sunt plena. — Et.

\rubric{In Octávis et Festis ritus simplícis: Glória Patri. Et.}

\lectioheading{Lectio iii}{}

\lettrine[lines=2]{S}{I} contémpsi subíre judícium cum servo meo et ancílla mea, cum disceptárent advérsum me: quid enim fáciam, cum surréxerit ad judicándum Deus? et, cum quæsíerit, quid respondébo illi? Numquid non in útero fecit me, qui et illum operátus est? et formávit me in vulva unus? Si negávi, quod volébant, paupéribus et óculos víduæ exspectáre feci; si comédi buccéllam meam solus, et non comédit pupíllus ex ea (quia ab infántia mea crevit mecum miserátio, et de útero matris meæ egréssa est mecum).

\R\ Quare detraxístis sermónibus veritátis? ad increpándum verba compónitis et subvértere nitímini amícum vestrum: * Verúmtamen quæ cogitástis, expléte. \V\ Quod justum est, judicáte; et non inveniétis in lingua mea iniquitátem. — Verúmtamen. Glória Patri. Verúmtamen.

\feriaheading{Feria Quinta infra Hebdomadam II Septembris}

\rubric{De libro Job}

\lectioheading{Lectio i}{Cap. 38, 1-20}

\lettrine[lines=2]{R}{ESPÓNDENS} autem Dóminus Job de túrbine, dixit: Quis est iste invólvens senténtias sermónibus imperítis? Accínge, sicut vir, lumbos tuos; interrogábo te et respónde mihi. Ubi eras quando ponébam fundaménta terræ? Índica mihi, si habes intelligéntiam. Quis pósuit mensúras ejus, si nosti? vel quis teténdit super eam líneam? Super quo bases illíus solidátæ sunt? aut quis demísit lápidem angulárem ejus, cum me laudárent simul astra matutína, et jubilárent omnes fílii Dei?

\R\ Si bona suscépimus de manu Dei, mala autem quare non sustineámus? * Dóminus dedit, Dóminus ábstulit; sicut Dómino plácuit, ita factum est: sit nomen Dómini benedíctum. \V\ Nudus egréssus sum de útero matris meæ et nudus revértar illuc. — Dóminus.

\lectioheading{Lectio ii}{}

\lettrine[lines=2]{Q}{UIS} conclúsit óstiis mare, quando erumpébat quasi de vulva procédens, cum ponérem nubem vestiméntum ejus, et calígine illud, quasi pannis infántiæ, obvólverem? Circúmdedi illud términis meis, et pósui vectem et óstia, et dixi: Usque huc vénies, et non procédes ámplius, et hic confrínges tuméntes fluctus tuos. Numquid post ortum tuum præcepísti dilúculo, et ostendísti aurórae locum suum? Et tenuísti concútiens extréma terræ, et excussísti ímpios ex ea?

\R\ Antequam cómedam suspíro, et tamquam inundántes aquæ sic rugítus meus; quia timor, quem timébam, evénit mihi, et quod verébar áccidit. Nonne dissimulávi? nonne sílui? nonne quiévi? * Et venit super me indignátio. \V\ Ecce non est auxílium mihi in me et necessárii quoque mei recessérunt a me. — Et.

\rubric{In Octávis et Festis ritus simplícis: Glória Patri. Et.}

\lectioheading{Lectio iii}{Cap. 42, 1-6}

\lettrine[lines=2]{R}{ESTITUÉTUR} ut lutum signáculum, et stabit sicut vestiméntum: auferétur ab ímpiis lux sua, et bráchium excélsum confringétur? Numquid ingréssus es profúnda maris, et in novíssimis abýssi deambulásti? Numquid apértæ sunt tibi portæ mortis, et óstia tenebrósa vidísti? Numquid considerásti latitúdinem terræ? Índica mihi, si nosti, ómnia: in qua via lux hábitet, et tenebrárum quis locus sit; ut ducas unumquódque ad términos suos, et intélligas sémitas domus ejus.

\R\ Quare detraxístis sermónibus veritátis? ad increpándum verba compónitis et subvértere nitímini amícum vestrum: * Verúmtamen quæ cogitástis, expléte. \V\ Quod justum est, judicáte; et non inveniétis in lingua mea iniquitátem. — Verúmtamen. Glória Patri. Verúmtamen.

\feriaheading{Feria Sexta infra Hebdomadam II Septembris}

\rubric{De libro Job}

\lectioheading{Lectio i}{Cap. 40, 1-11}

\lettrine[lines=2]{R}{ESPÓNDENS} autem Dóminus Job de túrbine, dixit: Accínge sicut vir lumbos tuos; interrogábo te, et índica mihi. Numquid írritum fácies judícium meum, et condemnábis me, ut tu justificéris? Et, si habes bráchium sicut Deus, et si voce símili tonas, circúmda tibi decórem et in sublíme erígere, et esto gloriósus et speciósis indúere véstibus.

\R\ Indúta est caro mea putrédine, et sórdibus púlveris cutis mea áruit et contrácta est: * Meménto mei, Dómine, quóniam ventus est vita mea. \V\ Dies mei velócius transiérunt quam a texénte tela succíditur, et consúmpti sunt absque ulla spe. — Meménto.

\lectioheading{Lectio ii}{}

\lettrine[lines=2]{D}{ISPÉRGE} supérbos in furóre tuo, et respíciens omnem arrogántem humília; réspice cunctos supérbos et confúnde eos, cóntere ímpios in loco suo; abscónde eos in púlvere simul, et fácies eórum demérge in fóveam; et ego confitébor quod salváre te possit déxtera tua. Ecce béhemoth, quem feci tecum, fœnum, quasi bos, cómedet. Fortitúdo ejus in lumbis ejus, et virtus illíus in umbilíco ventris ejus.

\R\ Páucitas diérum meórum finiétur brevi; dimítte me, Dómine, ut plangam páululum dolórem meum, * Antequam vadam ad terram tenebrósam et opértam mortis calígine. \V\ Manus tuæ, Dómine, fecérunt me, et plasmavérunt me totum in circúitu; et sic repénte præcípitas me? — Antequam.

\rubric{In Octávis et Festis ritus simplícis: Glória Patri. Antequam.}

\lectioheading{Lectio iii}{Cap. 42, 1-6}

\lettrine[lines=2]{R}{ESPÓNDENS} autem Job Dómino dixit: Scio quia ómnia potes, et nulla te latet cogitátio. Quis est iste qui celat consílium absque sciéntia? Ideo insipiénter locútus sum, et quæ ultra modum excéderent sciéntiam meam. Audi, et ego loquar; interrogábo te, et respónde mihi. Audítu auris audívi te; nunc autem óculus meus videt te. Idcírco ipse me reprehéndo, et ago pœniténtiam in favílla et cínere.

\R\ Non abscóndas me, Dómine, a fácie tua: manum tuam longe fac a me, * Et formído tua non me térreat. \V\ Córripe me, Dómine, in misericórdia, non in furóre tuo, ne forte ad níhilum rédigas me. — Et. Glória Patri. Et.

\feriaheading{Sabbato infra Hebdomadam II Septembris}

\rubric{De libro Job}

\lectioheading{Lectio i}{Cap. 42, 7-16}

\lettrine[lines=2]{P}{OSTQUAM} autem locútus est Dóminus verba hæc ad Job, dixit ad Elíphaz Themanítem: Irátus est furor meus in te et in duos amícos tuos, quóniam non estis locúti coram me rectum, sicut servus meus Job. Súmite ergo vobis septem tauros et septem aríetes et ite ad servum meum Job et offérte holocáustum pro vobis; Job autem servus meus orábit pro vobis. Fáciem ejus suscípiam, ut non vobis imputétur stultítia; neque enim locúti estis ad me recta, sicut servus meus Job.

\R\ Quis mihi tríbuat, ut in inférno prótegas me et abscóndas me, donec pertránseat furor tuus, Dómine, nisi tu, qui solus es Deus? * Et constítuas mihi tempus, in quo recordéris mei? \V\ Numquid sicut dies hóminis dies tui, ut quæras iniquitátem meam; cum sit nemo, qui de manu tua possit erúere? — Et constítuas mihi tempus, in quo recordéris mei?

\lectioheading{Lectio ii}{}

\lettrine[lines=2]{A}{BIÉRUNT} ergo Elíphaz Themanítes et Báldad Suhítes et Sophar Naamathítes, et fecérunt sicut locútus fúerat Dóminus ad eos, et suscépit Dóminus fáciem Job. Dóminus quoque convérsus est ad pœniténtiam Job, cum oráret ille pro amícis suis; et áddidit Dóminus ómnia, quæcúmque fúerant Job, duplícia. Venérunt autem ad eum omnes fratres sui et univérsæ soróres suæ cuncti, qui nóverant eum prius, et comedérunt cum eo panem in domo ejus; et movérunt super eum caput, et consoláti sunt eum super omni malo, quod intúlerat Dóminus super eum, et dedérunt ei unusquísque ovem unam et ináurem áuream unam.

\R\ Utinam appenderéntur peccáta mea, quibus iram mérui, * Et calamítas, quam pátior, in statéra. \V\ Quasi aréna maris hæc grávior apparéret, unde et verba mea dolóre sunt plena. — Et.

\rubric{In Octávis simplícibus et in Offício sanctæ Maríæ in Sábbato: Glória Patri. Et.}

\lectioheading{Lectio iii}{}

\lettrine[lines=2]{D}{OMINUS} autem benedíxit novíssimis Job magis quam princípio ejus; et facta sunt ei quatuórdecim míllia óvium et sex míllia camelórum et mille juga boum et mille ásinæ, et fuérunt ei septem fílii et tres fíliæ, et vocávit nomen uníus Diem et nomen secúndæ Cássiam et nomen tértiæ Cornustíbii. Non sunt autem invéntæ mulíeres speciósæ sicut fíliæ Job in univérsa terra; dedítque eis pater suus hereditátem inter fratres eárum. Vixit autem Job post hæc centum quadragínta annis, et vidit fílios suos et fílios filiórum suórum usque ad quartam generatiónem; et mórtuus est senex et plenus diérum.

\R\ Quare detraxístis sermónibus veritátis? ad increpándum verba compónitis et subvértere nitímini amícum vestrum: * Verúmtamen quæ cogitástis, expléte. \V\ Quod justum est, judicáte; et non inveniétis in lingua mea iniquitátem. — Verúmtamen. Glória Patri. Verúmtamen.

\hourheading{Ad Vesperas}

\V\ Vespertína orátio ascéndat ad te, Dómine. \R\ Et descéndat super nos misericórdia tua.

\antlabel{Ad Magnif. Ant.} Ne reminiscáris, * Dómine, delícta mea vel paréntum meórum, neque vindíctam sumas de peccátis meis.

\rubric{Oratio Dominicæ occurrentis.}

% ===== BATCH 3 of the Proprium de Tempore (autumnalis.md lines
% 13903-14129) -- the complete second week of September: "Dominica
% II Septembris" Matins (Nocturns I-II continuing Job through
% chapter 9 and St. Gregory's Moralia Book 9, with the same deferred
% Nocturn III and numbered Responsories vii-viii pattern as Dominica
% I), the six ferial days' Matins lessons carrying the Job cycle
% through its conclusion (chapters 27-42, ending with Job's
% restoration), and the closing Saturday Vespers. Two genuine source
% variants from the established responsory pattern are preserved
% exactly: Feria Quinta's third lesson and Feria Sexta's third
% lesson both cite "Cap. 42, 1-6" as a fresh citation (rather than
% continuing the prior lesson's citation with an empty second
% argument), since the source itself restates the chapter reference
% at each occurrence; and the Sabbato's first Responsory repeats its
% FULL verse text after the em-dash ("— Et constítuas mihi tempus,
% in quo recordéris mei?") rather than the usual single-word cue,
% transcribed exactly as the source gives it. The "---" horizontal
% rules between days are omitted as scan artifacts throughout. The
% document continues with "Dominica III Septembris" (autumnalis.md
% line 14130).
% =====
\feriaheading{Dominica III Septembris}

\rubric{Semiduplex}

\nocturnoheading{In I Nocturno}

\rubric{Incipit liber Tobiæ}

\lectioheading{Lectio i}{Cap. 1, 1-15}

\lettrine[lines=2]{T}{OBIAS} ex tribu et civitáte Néphthali, quæ est in superióribus Galilǽæ supra Naásson, post viam quæ ducit ad Occidéntem in sinístro habens civitátem Sephet, cum captus esset in diébus Salmánasar regis Assyriórum, in captivitáte tamen pósitus viam veritátis non deséruit; ita ut ómnia, quæ habére póterat, quotídie concaptívis frátribus, qui erant ex ejus génere, impertíret. Cumque esset júnior ómnibus in tribu Néphthali, nihil tamen puerile gessit in ópere.

\R\ Peto, Dómine, ut de vínculo improperii hujus absólvas me, aut certe désuper terram erípias me: * Ne reminiscáris delícta mea vel paréntum meórum, neque vindíctam sumas de peccátis meis: quia éruis sustinéntes te, Dómine. \V\ Ómnia enim judícia tua justa sunt, et omnes viæ tuæ misericórdia et véritas: et nunc, Dómine, meménto mei. — Ne.

\lectioheading{Lectio ii}{}

\lettrine[lines=2]{D}{ENIQUE}, cum irent omnes ad vítulos áureos, quos Jeróboam fécerat rex Israël, hic solus fugiébat consórtia ómnium; sed pergébat in Jerúsalem ad templum Dómini et ibi adorábat Dóminum Deum Israël, ómnia primitíva sua et décimas suas fidéliter ófferens; ita ut in tértio anno prosélytis et ádvenis ministráret omnem decimatiónem. Hæc et his simília secúndum legem Dei puérulus observábat. Cum vero factus esset vir, accépit uxórem Annam de tribu sua genuítque ex ea fílium, nomen suum impónens ei. Quem ab infántia timére Deum dócuit et abstinére ab omni peccáto.

\R\ Omni témpore bénedic Deum, et pete ab eo ut vias tuas dírigat, * Et in omni témpore consília tua in ipso permáneant. \V\ Inquíre ut fácias quæ plácita sunt illi in veritáte, et in tota virtúte tua. — Et.

\lectioheading{Lectio iii}{}

\lettrine[lines=2]{I}{GITUR} cum per captivitátem devenísset cum uxóre sua et fílio in civitátem Níniven cum omni tribu sua, cum omnes éderent ex cibis gentílium, iste custodívit ánimam suam et numquam contamínátus est in escis eórum. Et quóniam memor fuit Dómini in toto corde suo, dedit illi Deus grátiam in conspéctu Salmánasar regis, et dedit illi potestátem quocúmque vellet ire, habens libertátem quæcúmque fácere voluísset. Pergébat ergo ad omnes, qui erant in captivitáte, et mónita salútis dabat eis.

\R\ Memor esto, fili, quóniam páuperem vitam gérimus: * Habébis multa bona, si timúeris Deum. \V\ In mente habéto eum, et cave nequándo prætermíttas præcépta ejus. — Habébis. Glória Patri. Habébis.

\nocturnoheading{In II Nocturno}

\rubric{Sermo sancti Leónis Papæ}

\rubric{Sermo 9 de jejúnio 7 mensis}

\lectioheading{Lectio iv}{}

\lettrine[lines=2]{S}{CIO} quidem, dilectíssimi, plúrimos vestrum ita in iis, quæ ad observántiam christiánæ fídei pértinent, esse devótos, ut nostris cohortatiónibus non indígeant admonéri. Quod enim dudum et tradítio decrévit, et consuetúdo firmávit; nec erudítio ignórat, nec pietas prætermíttit. Sed quia sacerdotális offícii est, erga omnes Ecclésiæ fílios curam habére commúnem, in id quod et rúdibus prosit et doctis, quos simul dilígimus, páriter incitámus; ut jejúnium, quod nobis séptimi mensis recúrsus indícit, fide álacri per castigatiónem ánimi et córporis celebrémus.

\R\ Sufficiébat nobis paupértas nostra, ut divítiæ computaréntur: numquam fuísset pecúnia ipsa, pro qua misísti fílium nostrum, * Báculum senectútis nostræ! \V\ Heu me, fili mi, ut quid te mísimus peregrinári, lumen oculórum nostrórum? — Báculum.

\lectioheading{Lectio v}{}

\lettrine[lines=2]{I}{DEO} enim ipsa continéntiæ observántia quátuor est assignáta tempóribus, ut in idípsum totíus anni redeúnte decúrsu, cognoscerémus nos indesinénter purificatiónibus indigére; semperque esse niténdum, dum hujus vitæ varietáte jactámur, ut peccátum, quod fragilitáte carnis et cupiditátum pollutióne contráhitur, jejúniis atque eleemósynis deleátur. Esuriámus páululum, dilectíssimi, et aliquántulum, quod juvándis possit prodésse paupéribus, nostræ consuetúdini subtrahámus.

\R\ Benedícite Deum cæli et coram ómnibus vivéntibus confitémini ei, * Quia fecit vobíscum misericórdiam suam. \V\ Ipsum benedícite et cantáte illi: et enarráte ómnia mirabília ejus. — Quia.

\lectioheading{Lectio vi}{}

\lettrine[lines=2]{E}{LECTÉTUR} consciéntia benignórum frúctibus largitátis: et gáudia tríbuens, quo es lætificándus, accípies. Diléctio próximi, diléctio Dei est, qui plenitúdinem legis et prophetárum in hac gémina caritátis unitáte constítuit; ut nemo ambígeret, Deo se offérre, quod hómini contulísset, dicénte Dómino Salvatóre, cum de aléndis juvandísque paupéribus loquerétur: Quod uni eórum fecísti, mihi fecísti. Quarta ígitur et sexta féria jejunémus; sábbato vero apud beátum Petrum Apóstolum vigílias celebrémus: cujus nos méritis et oratiónibus crédimus adjuvándos, ut misericórdi Deo jejúnio nostro et devotióne placeámus.

\R\ Tempus est ut revértar ad eum qui misit me; * Vos autem benedícite Deum et enarráte ómnia mirabília ejus. \V\ Confitémini ei coram ómnibus vivéntibus, quia fecit vobíscum misericórdiam suam. — Vos. Glória Patri. Vos.

\R\ vii Tribulatiónes civitátum audívimus, quas passæ sunt, et defécimus: timor et hebetúdo mentis cécidit super nos et super líberos nostros: ipsi montes nolunt recípere fugam nostram: * Dómine, miserére. \V\ Peccávimus cum pátribus nostris, injúste égimus, iniquitátem fécimus. — Dómine.

\R\ viii Duo Séraphim, \rubric{ut infra in Homilíis.}

\feriaheading{Feria Secunda infra Hebdomadam III Septembris}

\rubric{Si Lectiónes de Scriptúra hujus Fériæ II aut Fériæ III sequéntis, sua die dici non potúerint, et in Fériis Quátuor Témporum fíeri contíngat Offícium, in quo Lectiónes de Scriptúra occurrénti sint dicéndæ, neque in eis reponéndum sit Inítium libri Tobíæ Domínicæ præcedénti assignátum, Lectiónes ut supra impedítæ, et si opus sit, étiam Lectiónes insequéntium Feriárum, servató tamen eárum órdine, ultérius transferúntur, ita ut in Offício in Quátuor Témporibus occurrénti Lectiónes de Scriptúra valeant adhibéri. Quod ítem servátur de Lectiónibus Fériæ V ut supra impedítis, quæ simíliter in casu vel ultérius transferúntur, vel étiam anticipántur. Responsória tamen semper dicúntur de Féria occurrénti.}

\rubric{De libro Tobíæ}

\lectioheading{Lectio i}{Cap. 2, 1-4}

\lettrine[lines=2]{P}{OST} hæc vero, cum esset dies festus Dómini, et factum esset prándium bonum in domo Tobíæ, dixit fílio suo: Vade et adduc áliquos de tribu nostra timéntes Deum, ut epuléntur nobíscum. Cumque abíisset, revérsus nuntiávit ei unum ex fíliis Israël jugulátum jacére in platéa. Statímque exsíliens de accúbitu suo relínquens prándium, jejúnus pervénit ad corpus, tollénsque illud portávit ad domum suam occúlte, ut, dum sol occubuísset, caute sepelíret eum.

\R\ Peto, Dómine, ut de vínculo improperii hujus absólvas me, aut certe désuper terram erípias me: * Ne reminiscáris delícta mea vel paréntum meórum, neque vindíctam sumas de peccátis meis: quia éruis sustinéntes te, Dómine. \V\ Omnia enim judícia tua justa sunt, et omnes viæ tuæ misericórdia et véritas: et nunc, Dómine, meménto mei. — Ne.

\lectioheading{Lectio ii}{Cap. 2, 8-18}

\lettrine[lines=2]{A}{RGUÉBANT} autem eum omnes próximi ejus dicéntes: Jam hujus rei causa intérfici jussus es et vix effugísti mortis impérium; et íterum sépelis mórtuos? Sed Tobías plus timens Deum quam regem, rapiébat córpora occísórum et occultábat in domo sua et médiis nóctibus sepeliébat ea. Cóntigit autem ut quadam die fatigátus a sepultúra, véniens in domum suam jactásset se juxta paríetem et obdormísset, et ex nido hirúndinum dormiénti illi cálida stércora incíderent super óculos ejus fierétque cæcus. Hanc autem tentatiónem ídeo permísit Dóminus eveníre illi, ut postéris darétur exémplum patiéntiæ ejus, sicut et sancti Job.

\R\ Omni témpore bénedic Deum, et pete ab eo ut vias tuas dírigat, * Et in omni témpore consília tua in ipso permáneant. \V\ Inquíre ut fácias quæ plácita sunt illi in veritáte, et in tota virtúte tua. — Et.

\rubric{In Octávis et Festis ritus simplícis: Glória Patri. Et.}

\lectioheading{Lectio iii}{}

\lettrine[lines=2]{N}{AM}, cum ab infántia sua semper Deum timúerit et mandáta ejus custodíerit, non est contristátus contra Deum quod plaga cæcitátis evénerit ei; sed immóbilis in Dei timóre permánsit agens grátias Deo ómnibus diébus vitæ suæ. Nam sicut beáto Job insultábant reges, ita isti paréntes et cognáti ejus irridébant vitam ejus dicéntes: Ubi est spes tua, pro qua eleemósynas et sepultúras faciébas? Tobías vero increpábat eos dicens: Nolíte ita loqui: quóniam fílii sanctórum sumus et vitam illam exspectámus, quam Deus datúrus est his qui fidem suam numquam mutant ab eo.

\R\ Memor esto, fili, quóniam páuperem vitam gérimus: * Habébis multa bona, si timúeris Deum. \V\ In mente habéto eum, et cave nequándo prætermíttas præcépta ejus. — Habébis. Glória Patri. Habébis.

\feriaheading{Feria Tertia infra Hebdomadam III Septembris}

\rubric{De libro Tobíæ}

\lectioheading{Lectio i}{Cap. 2, 19-23}

\lettrine[lines=2]{A}{NNA} vero uxor ejus ibat ad opus textrínum quotídie, et de labóre mánuum suárum victum, quem cónsequi póterat, deferébat. Unde factum est ut hædum caprárum accípiens detulísset domi; cujus cum vocem baléntis vir ejus audísset, dixit: Vidéte, ne forte furtívus sit; réddite eum dóminis suis, quia non licet nobis aut édere ex furto áliquid aut contíngere.

\R\ Sufficiébat nobis paupértas nostra, ut divítiæ computaréntur: numquam fuísset pecúnia ipsa, pro qua misísti fílium nostrum, * Báculum senectútis nostræ! \V\ Heu me, fili mi, ut quid te mísimus peregrinári, lumen oculórum nostrórum? — Báculum.

\lectioheading{Lectio ii}{}

\lettrine[lines=2]{A}{D} hæc uxor ejus iráta respóndit: Manifésté vana facta est spes tua et eleemósynæ tuæ modo apparuérunt. Atque his et áliis hujuscémodi verbis exprobrábat ei. Tunc Tobías ingémuit et cœpit oráre cum lácrimis dicens: Justus es, Dómine, et ómnia judícia tua justa sunt, et omnes viæ tuæ misericórdia et véritas et judícium. Et nunc, Dómine, memor esto mei et ne vindíctam sumas de peccátis meis neque reminiscáris delícta mea vel paréntum meórum.

\R\ Benedícite Deum cæli et coram ómnibus vivéntibus confitémini ei, * Quia fecit vobíscum misericórdiam suam. \V\ Ipsum benedícite et cantáte illi: et enarráte ómnia mirabília ejus. — Quia.

\rubric{In Octávis et Festis ritus simplícis: Glória Patri. Quia.}

\lectioheading{Lectio iii}{}

\lettrine[lines=2]{Q}{UÓNIAM} non obedívimus præcéptis tuis, ídeo tráditi sumus in direptiónem et captivitátem et mortem et in fábulam et in impropérium ómnibus natiónibus, in quibus dispersísti nos. Et nunc, Dómine, magna judícia tua, quia non égimus secúndum præcépta tua et non ambulávimus sincériter coram te. Et nunc, Dómine, secúndum voluntátem tuam fac mecum et præcipe in pace récipi spíritum meum; éxpedit enim mihi mori magis quam vívere.

\R\ Tempus est ut revértar ad eum qui misit me; * Vos autem benedícite Deum et enarráte ómnia mirabília ejus. \V\ Confitémini ei coram ómnibus vivéntibus, quia fecit vobíscum misericórdiam suam. — Vos. Glória Patri. Vos.

\rubric{¶ Fériæ IV et VI ac Sábbato Quátuor Témporum sunt Fériæ majóres non privilegiátæ.}

% ===== BATCH 4 of the Proprium de Tempore (autumnalis.md lines
% 14130-14245) -- "Dominica III Septembris" Matins, opening the Book
% of Tobias cycle in Nocturn I and St. Leo's sermon on the September
% Ember fast in Nocturn II (again with the deferred Nocturn III and
% numbered Responsories vii-viii pattern), followed by "Feria
% Secunda" and "Feria Tertia infra Hebdomadam III Septembris"
% continuing Tobias. Feria Secunda's lesson heading carries an
% unusually long governing rubric on how Ember Day lesson transfers
% interact with the Tobias readings when a conflict arises; this is
% transcribed in full as a single \rubric{} block. The closing
% rubric "Fériæ IV et VI ac Sábbato Quátuor Témporum sunt Fériæ
% majóres non privilegiátæ" signals the upcoming September Ember
% Days. The document continues with "Feria Quarta, Quatuor Temporum
% Septembris" (autumnalis.md line 14246).
% =====
\feriaheading{Feria Quarta, Quatuor Temporum Septembris}

\rubric{Hódie, Féria VI et Sábbato sequéntibus, Antíphonæ et Psalmi ad Laudes, et hac Féria étiam tres últimæ Antíphonæ cum Psalmis ad Matutínum, sumúntur de 2 loco; ad Primam autem addítur quartus Psalmus, ut in Psaltério notátur. Et dicúntur Preces feriáles ad omnes Horas, ut in Ordinário, præterquam in Vésperis et ad Complétórium Sábbati.}

\rubric{Léctio sancti Evangélii secúndum Marcum}

\lectioheading{Lectio i}{Cap. 9, 16-28}

\lettrine[lines=2]{I}{N} illo témpore: Respóndens unus de turba dixit ad Jesum: Magíster, áttuli fílium meum ad te habéntem spíritum mutum. Et réliqua.

\rubric{Homilía sancti Bedæ Venerábilis Presbýteri}

\rubric{Liber 3, cap. 38 in Marcum 9}

\lettrine[lines=2]{D}{ÆMONÍACUM} hunc, quem descéndens de monte Dóminus sanávit, Marcus quidem surdum mutúmque, Matthǽus vero lunáticum fuísse commémorat. Significat autem eos, de quibus scriptum est: Stultus ut luna mutátur. Qui numquam in eódem statu permanéntes, nunc ad hæc, nunc ad illa vítia mutáti, crescunt atque decréscunt. Qui muti sunt, non confiténdo fidem; surdi, nec ipsum aliquátenus veritátis audiéndo sermónem. Spumant autem, cum stultítia tabéscunt; stultórum namque et languéntium atque hébetum est, spumas salivárum ex ore dimíttere. Strident déntibus, cum iracúndiæ furóre flamméscunt: aréscunt, cum ótio torpénte languéscunt, et nulla virtútis indústria confortáti enérviter vivunt.

\R\ Tribulatiónes civitátum audívimus, quas passæ sunt, et defécimus: timor et hebetúdo mentis cécidit super nos et super líberos nostros: ipsi montes nolunt recípere fugam nostram: * Dómine, miserére. \V\ Peccávimus cum pátribus nostris, injúste égimus, iniquitátem fécimus. — Dómine.

\lectioheading{Lectio ii}{}

\lettrine[lines=2]{Q}{UOD} autem ait: Et dixi discípulis tuis ut ejícerent illum, et non potuérunt; laténter Apóstolos accúsat, cum impossibílitas curándi, intérdum non ad imbecillitátem curántium, sed ad eórum, qui curándi sunt, fidem referátur, dicénte Dómino: Fiat tibi secúndum fidem tuam. Qui respóndens eis, dixit: O generátio incrédula, quámdiu apud vos ero? quámdiu vos pátiar? Non quod tædio superátus sit mansuétus ac mitis, qui non apéruit sicut agnus coram tondénte os suum, nec in verba furóris erúpit; sed quo in similitúdinem médici, si ægrótum vídeat contra sua præcépta se gérere, dicat: Úsquequo accédam ad domum tuam? quoúsque artis meæ perdam indústriam, me áliud jubénte, et te áliud perpetránte?

\R\ Omni témpore bénedic Deum, et pete ab eo ut vias tuas dírigat, * Et in omni témpore consília tua in ipso permáneant. \V\ Inquíre ut fácias quæ plácita sunt illi in veritáte, et in tota virtúte tua. — Et.

\lectioheading{Lectio iii}{}

\lettrine[lines=2]{E}{T} dixit illis: Hoc genus in nullo potest exíre, nisi in oratióne et jejúnio. Dum docet Apóstolos, quómodo dæmon nequíssimus débeat expélli, omnes instítuit ad vitam; ut scílicet novérimus, gravíssima quæque vel immundórum spirítuum, vel hóminum tentaménta jejúniis et oratiónibus esse superánda: iram quoque Dómini, cum in ultiónem nostrórum scélerum fúerit accénsa, hoc remédio singulári posse placári. Jejúnium autem generále est, non solum ab escis, sed et a cunctis illécebris abstinére carnálibus, immo ab ómnibus vitiórum continére se passiónibus. Sic et orátio generális non in verbis solum est, quibus divínam cleméntiam invocámus, verum étiam in ómnibus, quæ in obséquium nostri Condítoris fídei devotióne gérimus.

\R\ Memor esto, fili, quóniam páuperem vitam gérimus: * Habébis multa bona, si timúeris Deum. \V\ In mente habéto eum, et cave nequándo prætermíttas præcépta ejus. — Habébis. Glória Patri. Habébis.

\hourheading{Ad Laudes}

\V\ Repléti sumus mane misericórdia tua. \R\ Exsultávimus, et delectáti sumus.

\antlabel{Ad Bened. Ant.} Hoc genus * dæmoniórum in nullo potest exíre, nisi in oratióne et jejúnio.

\oratioheading{Oratio}

\lettrine[lines=2]{M}{ISERICÓRDIÆ} tuæ remédiis, quǽsumus, Dómine, fragílitas nostra subsístat: ut, quæ sua conditióne attéritur, tua cleméntia reparétur. Per Dóminum.

\rubric{Et dícitur usque ad Nonam. Ad Vésperas vero Orátio Domínicæ præcedéntis. Si autem Vésperæ dicéndæ non sint de Féria, de ea nec fit Commemorátio.}

\feriaheading{Feria Quinta infra Hebdomadam III Septembris}

\rubric{De libro Tobíæ}

\lectioheading{Lectio i}{Cap. 12, 1-17}

\lettrine[lines=2]{V}{OCÁVIT} ad se Tobías fílium suum dixítque ei: Quid póssumus dare viro isti sancto, qui venit tecum? Respóndens Tobías dixit patri suo: Pater, quam mercédem dábimus ei? aut quid dignum póterit esse benefíciis ejus? Me duxit et redúxit sanum, pecúniam a Gabélo ipse recépit, uxórem ipse me habére fecit et dæmónium ab ea ipse compéscuit, gáudium paréntibus ejus fecit, meípsum a devoratióne piscis erípuit, te quoque vidére fecit lumen cæli, et bonis ómnibus per eum repléti sumus. Quid illi ad hæc potérimus dignum dare? Sed peto te, pater mi, ut roges eum, si forte dignábitur medietátem de ómnibus quæ alláta sunt sibi assúmere.

\R\ Peto, Dómine, ut de vínculo improperii hujus absólvas me, aut certe désuper terram erípias me: * Ne reminiscáris delícta mea vel paréntum meórum, neque vindíctam sumas de peccátis meis: quia éruis sustinéntes te, Dómine. \V\ Omnia enim judícia tua justa sunt, et omnes viæ tuæ misericórdia et véritas: et nunc, Dómine, meménto mei. — Ne.

\lectioheading{Lectio ii}{}

\lettrine[lines=2]{E}{T} vocántes eum, pater scílicet et fílius, tulérunt eum in partem; et rogáre cœpérunt ut dignarétur dimídiam partem ómnium quæ attúlerant, accéptam habére. Tunc dixit eis occúlte: Benedícite Deum cæli et coram ómnibus vivéntibus confitémini ei, quia fecit vobíscum misericórdiam suam. Etenim sacraméntum regis abscóndere bonum est, ópera autem Dei reveláre et confitéri honoríficum est. Bona est orátio cum jejúnio, et eleemósyna magis quam thesáuros auri recóndere; quóniam eleemósyna a morte líberat, et ipsa est quæ purgat peccáta, et facit inveníre misericórdiam et vitam ætérnam. Qui autem fáciunt peccátum et iniquitátem hostes sunt ánimæ suæ.

\R\ Omni témpore bénedic Deum, et pete ab eo ut vias tuas dírigat, * Et in omni témpore consília tua in ipso permáneant. \V\ Inquíre ut fácias quæ plácita sunt illi in veritáte, et in tota virtúte tua. — Et.

\rubric{In Octávis et Festis ritus simplícis: Glória Patri. Et.}

\lectioheading{Lectio iii}{}

\lettrine[lines=2]{M}{ANIFÉSTO} ergo vobis veritátem et non abscóndam a vobis occúltum sermónem. Quando orábas cum lácrimis et sepeliébas mórtuos et derelinquébas prándium tuum et mórtuos abscondébas per diem in domo tua et nocte sepeliébas eos, ego óbtuli oratiónem tuam Dómino. Et quia accéptus eras Deo, necésse fuit ut tentátio probáret te. Et nunc misit me Dóminus ut curárem te et Saram uxórem fílii tui a dæmónio liberárem; ego enim sum Ráphaël Angelus, unus ex septem qui astámus ante Dóminum. Cumque hæc audíssent, turbáti sunt et treméntes cecidérunt super terram in fáciem suam. Dixítque eis Angelus: Pax vobis, nolíte timére.

\R\ Memor esto, fili, quóniam páuperem vitam gérimus: * Habébis multa bona, si timúeris Deum. \V\ In mente habéto eum, et cave nequándo prætermíttas præcépta ejus. — Habébis. Glória Patri. Habébis.

\feriaheading{Feria Sexta, Quatuor Temporum Septembris}

\rubric{Léctio sancti Evangélii secúndum Lucam}

\lectioheading{Lectio i}{Cap. 7, 36-50}

\lettrine[lines=2]{I}{N} illo témpore: Rogábat Jesum quidam de pharisǽis, ut manducáret cum illo. Et ingréssus domum pharisǽi discúbuit. Et réliqua.

\rubric{Homilía sancti Gregórii Papæ}

\rubric{Homilía 33 in Evang., post inítium}

\lettrine[lines=2]{Q}{UEM} pharisǽus de falsa justítia præsúmens, nisi Judáicum pópulum; quem peccátrix múlier ad vestígia Dómini véniens et plorans, nisi convérsam gentilitátem desígnat? Quæ cum alabástro venit, unguéntum fudit, retro secus pedes Dómini stetit, lácrimis pedes rigávit, capíllis tersit, eosdémque quos infundébat et tergébat pedes osculári non désiit. Nos ergo, nos illa múlier expréssit, si toto corde ad Dóminum post peccáta redeámus, si ejus pœniténtiæ luctus imitémur. Quid namque unguénto, nisi bonæ odor opiniónis exprímitur? Unde et Paulus dicit: Christi bonus odor sumus Deo in omni loco.

\R\ Sufficiébat nobis paupértas nostra, ut divítiæ computaréntur: numquam fuísset pecúnia ipsa, pro qua misísti fílium nostrum, * Báculum senectútis nostræ! \V\ Heu me, fili mi, ut quid te mísimus peregrinári, lumen oculórum nostrórum? — Báculum.

\lectioheading{Lectio ii}{}

\lettrine[lines=2]{S}{I} ergo recta ópera ágimus, quibus opiniónis bone odóre Ecclésiam respergámus, quid in Dómini córpore, nisi unguéntum fúndimus? Sed secus pedes Jesu múlier stetit: contra pedes enim Dómini stétimus, cum in peccátis pósiti ejus itinéribus renitebámur. Sed si ad veram pœniténtiam post peccáta convértimur, jam retro secus pedes stamus; quia ejus vestígia séquimur, quem impugnabámus. Lácrimis múlier pedes ejus rigat: quod nos quoque veráciter ágimus, si quibúslibet últimis membris Dómini per compassiónis afféctum inclinámur, si sanctis ejus in tribulatióne compátimur, si eórum tristítiam, nostram putámus.

\R\ Benedícite Deum cæli et coram ómnibus vivéntibus confitémini ei, * Quia fecit vobíscum misericórdiam suam. \V\ Ipsum benedícite et cantáte illi: et enarráte ómnia mirabília ejus. — Quia.

\lectioheading{Lectio iii}{}

\lettrine[lines=2]{C}{APILLIS} ergo pedes Dómini tergimus, quando sanctis ejus, quibus ex caritáte compátimur, étiam ex his quæ nobis supérfluunt, miserémur; quátenus sic mens per compassiónem dóleat, ut étiam larga manus afféctum dolóris osténdat. Rigat namque lácrimis Redemptóris pedes, sed capíllis suis non tergit, qui utcúmque proximórum dolóri compátitur, sed tamen eis ex his quæ sibi supérfluunt, non miserétur. Plorat, et non tergit, qui verba quidem dolóris tríbuit, sed non ministrándo quæ desunt, vim dolóris mínime abscíndit. Osculátur múlier pedes, quos tergit: quod nos quoque plene ágimus, si studióse dilígimus, quos ex largitáte continémus; ne gravis nobis sit necéssitas próximi; ne ipsa nobis ejus indigéntia, quæ sustentátur, fiat onerósa; et, cum manus necessária tríbuit, ánimus a dilectióne torpéscat.

\R\ Tempus est ut revértar ad eum qui misit me; * Vos autem benedícite Deum et enarráte ómnia mirabília ejus. \V\ Confitémini ei coram ómnibus vivéntibus, quia fecit vobíscum misericórdiam suam. — Vos. Glória Patri. Vos.

\hourheading{Ad Laudes}

\V\ Repléti sumus mane misericórdia tua. \R\ Exsultávimus, et delectáti sumus.

\antlabel{Ad Bened. Ant.} Múlier * quæ erat in civitáte peccátrix, stans retro secus pedes Dómini, lácrimis cœpit rigáre pedes ejus et capíllis cápitis sui tergébat, et deosculabátur pedes ejus et unguénto ungébat.

\oratioheading{Oratio}

\lettrine[lines=2]{P}{RÆSTA}, quǽsumus, omnípotens Deus: ut observatiónes sacras ánnua devotióne recoléntes, et córpore tibi placeámus, et mente. Per Dóminum.

\rubric{Et dícitur usque ad Nonam. Ad Vésperas vero Orátio Domínicæ præcedéntis. Si autem Vésperæ dicéndæ non sint de Féria, de ea nec fit Commemorátio.}

\feriaheading{Sabbato, Quatuor Temporum Septembris}

\rubric{Léctio sancti Evangélii secúndum Lucam}

\lectioheading{Lectio i}{Cap. 13, 6-17}

\lettrine[lines=2]{I}{N} illo témpore: Dicébat Jesus turbis hanc similitúdinem: Arbórem fici habébat quidam plantátam in vínea sua et venit quærens fructum in illa et non invénit. Et réliqua.

\rubric{Homilía sancti Gregórii Papæ}

\rubric{Homilía 31 in Evangelia}

\lettrine[lines=2]{D}{OMINUS} ac Redémptor noster per Evangélium suum aliquándo verbis, aliquándo rebus lóquitur: aliquándo áliud verbis, atque áliud rebus; aliquándo autem hoc verbis, quod rebus. Duas étenim res ex Evangélio, fratres, audístis: ficúlneam infructuósam, et mulíerem curvam; et utríque rei est pietas impénsa. Illud autem dixit per similitúdinem: istud egit per exhibitiónem. Sed hoc significat ficúlnea infructuósa, quod múlier inclináta; et hoc ficúlnea reserváta, quod múlier erécta.

\R\ Tribulatiónes civitátum audívimus, quas passæ sunt, et defécimus: timor et hebetúdo mentis cécidit super nos et super líberos nostros: ipsi montes nolunt recípere fugam nostram: * Dómine, miserére. \V\ Peccávimus cum pátribus nostris, injúste égimus, iniquitátem fécimus. — Dómine.

\lectioheading{Lectio ii}{}

\lettrine[lines=2]{Q}{UID} arbor fici, nisi humánam natúram desígnat? Quid múlier inclináta, nisi eámdem natúram significándo denúntiat? Quæ et bene plantáta est sicut ficus, et bene creáta sicut múlier; sed in culpam própria sponte lapsa, neque fructum servat operatiónis, neque statum rectitúdinis. Ad peccátum quippe ex voluntáte córruens, quia fructum obediéntiæ ferre nóluit, statum rectitúdinis amísit. Quæ ad Dei similitúdinem cóndita, dum in sua dignitáte non pérstitit, quod plantáta vel creáta fúerat, serváre contémpsit. Tértio dóminus víneæ ad ficúlneam venit, quia natúram géneris humáni ante legem, sub lege, sub grátia, exspectándo, admonéndo, visitándo requisívit.

\R\ Omni témpore bénedic Deum, et pete ab eo ut vias tuas dírigat, * Et in omni témpore consília tua in ipso permáneant. \V\ Inquíre ut fácias quæ plácita sunt illi in veritáte, et in tota virtúte tua. — Et.

\lectioheading{Lectio iii}{}

\lettrine[lines=2]{V}{ENIT} ante legem, quia per naturálem intelléctum quid unusquísque exémplo sui, quáliter erga próximum ágere debuísset, innótuit. Venit in lege, quia præcipiéndo dócuit. Venit post legem per grátiam, quia pietátis suæ præséntiam exhibéndo monstrávit. Sed tamen in tribus annis fructum se non invenísse conquéritur: quia quorúmdam pravórum mentes nec inspiráta lex naturális córrigit, nec præcépta erúdiunt, nec incarnatiónis ejus miracula convértunt. Quid vero per cultórem víneæ, nisi præpositórum ordo exprímitur? Qui dum præsunt Ecclésiæ, nimírum Domínicæ víneæ curam gerunt.

\R\ Memor esto, fili, quóniam páuperem vitam gérimus: * Habébis multa bona, si timúeris Deum. \V\ In mente habéto eum, et cave nequándo prætermíttas præcépta ejus. — Habébis. Glória Patri. Habébis.

\hourheading{Ad Laudes}

\V\ Repléti sumus mane misericórdia tua. \R\ Exsultávimus, et delectáti sumus.

\antlabel{Ad Bened. Ant.} Illúmina, Dómine, * sedéntes in ténebris et umbra mortis, et dírige pedes nostros in viam pacis.

\oratioheading{Oratio}

\lettrine[lines=2]{O}{MNÍPOTENS} sempitérne Deus, qui per continéntiam salutárem córporibus medéris et méntibus: majestátem tuam súpplices exorámus; ut pia jejunántium deprecatióne placátus, et præséntia nobis subsídia tríbuas et futúra. Per Dóminum.

\rubric{Et dícitur usque ad Nonam inclusíve.}

\hourheading{Ad Vesperas}

\V\ Vespertína orátio ascéndat ad te, Dómine. \R\ Et descéndat super nos misericórdia tua.

\antlabel{Ad Magnif. Ant.} Adonái, Dómine, * Deus magne et mirábilis, qui dedísti salútem in manu féminæ, exáudi preces servórum tuórum.

\rubric{Oratio Dominicæ occurrentis.}

% ===== BATCH 5 of the Proprium de Tempore (autumnalis.md lines
% 14246-14411) -- the complete September Ember Days (Feria Quarta,
% Feria Sexta, and Sabbato, Quatuor Temporum Septembris), each
% following the Sunday-Gospel-homily pattern: a brief Gospel
% pericope incipit ending "Et réliqua" under \lectioheading{Lectio
% i}, immediately followed within the same Lesson by a Patristic
% homily on that Gospel (Bede on Mark 9, Gregory on Luke 7:36-50,
% Gregory again on Luke 13:6-17), each homily's author and source
% given as separate \rubric{} labels before its own fresh lettrine,
% matching the established precedent from aestiva.tex's parallel
% "Lectio vii" Trinity Sunday construction. Feria Quinta infra
% Hebdomadam III Septembris (an ordinary day, not an Ember day,
% falling between the Wednesday and Friday Ember days) continues
% the Tobias cycle with the Archangel Raphael's self-revelation.
% Each Ember day's all-caps "AD LAUDES" (and, for the Saturday, "AD
% VESPERAS") sub-heading is normalized to the standard mixed-case
% \hourheading{} form per the established incidental-formatting-
% variation rule, each followed by its own brief antiphon, Oratio,
% and governing rubric on how long the Oratio is said. The document
% continues with "Dominica IV Septembris" (autumnalis.md line
% 14412).
% =====
\feriaheading{Dominica IV Septembris}

\rubric{Semiduplex}

\nocturnoheading{In I Nocturno}

\rubric{Incipit liber Judith}

\lectioheading{Lectio i}{Cap. 1, 1-10; 2, 1-3}

\lettrine[lines=2]{A}{RPHÁXAD} ítaque rex Medórum subjugáverat multas gentes império suo, et ipse ædificávit civitátem potentíssimam, quam appellávit Ecbátanis, ex lapídibus quadrátis et sectis. Fecit muros ejus in altitúdinem cubitórum septuagínta et in latitúdinem cubitórum trigínta; turres vero ejus pósuit in altitúdinem cubitórum centum. Per quadrum vero eárum latus utrúmque vicenórum pedum spátio tendebátur; posuítque portas ejus in altitúdinem túrrium; et gloriabátur quasi potens in poténtia exércitus sui et in glória quadrigárum suárum.

\R\ Adonái, Dómine, Deus magne et mirábilis, qui dedísti salútem in manu féminæ, * Exáudi preces servórum tuórum. \V\ Benedíctus es, Dómine, qui non derelínquis præsuméntes de te, et de sua virtúte gloriántes humílias. — Exáudi.

\lectioheading{Lectio ii}{}

\lettrine[lines=2]{A}{NNO} ígitur duodécimo regni sui Nabuchodónosor rex Assyriórum, qui regnábat in Nínive civitáte magna, pugnávit contra Arpháxad et obtínuit eum in campo magno qui appellátur Rágau, circa Euphráten et Tigrin et Jádason, in campo Erioch regis Elicórum. Tunc exaltátum est regnum Nabuchodónosor, et cor ejus elevátum est; et misit ad omnes qui habitábant in Cilícia et Damásco et Líbano, et ad gentes, quæ sunt in Carmélo et Cedar et inhabitántes Galilǽam in campo magno Esdrelon, et ad omnes qui erant in Samaría et trans flumen Jordánem usque ad Jerúsalem, et omnem terram Jesse, quoúsque perveniátur ad términos Æthiópiæ.

\R\ Tribulatiónes civitátum audívimus, quas passæ sunt, et defécimus: timor et hebetúdo mentis cécidit super nos et super líberos nostros: ipsi montes nolunt recípere fugam nostram: * Dómine, miserére. \V\ Peccávimus cum pátribus nostris, injúste égimus, iniquitátem fécimus. — Dómine.

\lectioheading{Lectio iii}{}

\lettrine[lines=2]{A}{D} hos omnes misit núntios Nabuchodónosor rex Assyriórum; qui omnes uno ánimo contradixérunt et remisérunt eos vácuos et sine honóre abjecérunt. Tunc indignátus Nabuchodónosor rex advérsus omnem terram illam, jurávit per thronum et regnum suum, quod defénderet se de ómnibus regiónibus his. Anno tértio décimo Nabuchodónosor regis, vigésima et secúnda die mensis primi factum est verbum in domo Nabuchodónosor regis Assyriórum, ut defénderet se. Vocavítque omnes majóres natu, omnésque duces et bellatóres suos, et hábuit cum eis mystérium consílii sui dixítque cogitatiónem suam in eo esse, ut omnem terram suo subjugáret império.

\R\ Benedícat te Dóminus in virtúte sua, qui per te ad níhilum redégit inimícos nostros: * Ut non defíciat laus tua de ore hóminum. \V\ Benedíctus Dóminus qui creávit cælum et terram; quia hódie nomen tuum ita magnificávit. — Ut. Glória Patri. Ut.

\nocturnoheading{In II Nocturno}

\rubric{Ex libro sancti Ambrósii Epíscopi de Elía et jejúnio}

\lectioheading{Lectio iv}{Cap. 9}

\lettrine[lines=2]{P}{OTÉNTES} vinum prohibéntur bíbere, ne, cum bíberint, obliviscántur sapiéntiam. Dénique bibébant vinum in ebrietáte poténtes, qui Holoférni príncipi milítiæ regis Assyriórum se trádere gestiébant; sed non bibébat fémina Judith, jejúnans ómnibus diébus viduitátis suæ, præter festórum diérum solemnitátes. His armis muníta procéssit et omnem Assyriórum circumvénit exércitum. Sóbrii vigóre consílii ábstulit Holoférnis caput, servávit pudicítiam, victóriam reportávit.

\R\ Nos álium Deum nescímus præter Dóminum, in quo sperámus: * Qui non déspicit nos, nec ámovet salútem suam a génere nostro. \V\ Indulgéntiam ipsíus fusis lácrimis postulémus, et humiliémus illi ánimas nostras. — Qui.

\lectioheading{Lectio v}{}

\lettrine[lines=2]{H}{ÆC} enim succíncta jejúnio in castris prætendébat aliénis; ille vino sepúltus jacébat, ut ictum vúlneris sentíre non posset. Itaque uníus mulíeris jejúnium innúmeros strávit exércitus Assyriórum. Esther quoque púlchrior facta est jejúnio; Dóminus enim grátiam sóbriæ mentis augébat. Omne genus suum, id est, totum pópulum Judæórum a persecutiónis acerbitáte liberávit, ita ut regem sibi fáceret esse subjéctum.

\R\ Dominátor, Dómine cælórum et terræ, Creátor aquárum, Rex univérsæ creatúræ: * Exáudi oratiónem servórum tuórum. \V\ Tu, Dómine, cui humílium semper et mansuetórum plácuit deprecátio. — Exáudi.

\lectioheading{Lectio vi}{}

\lettrine[lines=2]{I}{TAQUE} illa, quæ tríduo jejunávit contínuo, et corpus suum aqua lavit, plus plácuit, et vindíctam rétulit. Aman autem, dum se regáli jactat convívio, inter ipsa vina pœnam suæ ebrietátis exsólvit. Est ergo jejúnium reconciliatiónis sacrifícium, virtútis increméntum, quod fecit étiam féminas fortióres augménto grátiæ. Jejúnium nescit fænerátorem, non sortem fæneris novit: non rédolet usúras mensa jejunántium. Etiam ipsis jejúnium convívíis dat grátiam: dulcióres post famem épulæ fiunt, quæ assiduitáte fastídio sunt, et diutúrna continuatióne viléscunt. Condiméntum cibi jejúnium est: quanto avídior appeténtia, tanto esca jucúndior.

\R\ Dómine Deus, qui cónteris bella ab inítio, éleva bráchium tuum super Gentes, quæ cógitant servis tuis mala: * Et déxtera tua glorificétur in nobis. \V\ Allíde virtútem eórum in virtúte tua; cadat robur eórum in iracúndia tua. — Et. Glória Patri. Et.

\R\ vii Confórta me, Rex Sanctórum principátum tenens: * Et da sermónem rectum et bene sonántem in os meum. \V\ Dómine, Rex univérsæ potestátis, convérte consílium eórum super eos. — Et.

\R\ viii Duo Séraphim, \rubric{ut infra in Homilíis.}

\feriaheading{Feria Secunda infra Hebdomadam IV Septembris}

\rubric{De libro Judith}

\lectioheading{Lectio i}{Cap. 4, 1-12}

\lettrine[lines=2]{F}{ILII} Israël, qui habitábant in terra Juda, timuérunt valde a fácie Holoférnis; tremor et horror invásit sensus eórum, ne hoc fáceret Jerúsalem et templo Dómini, quod fécerat céteris civitátibus et templis eárum. Et misérunt in omnem Samaríam per circúitum usque Jéricho et præoccupavérunt omnes vértices móntium et muris circumdedérunt vicos suos et congregavérunt fruménta in præparatiónem pugnæ.

\R\ Adonái, Dómine, Deus magne et mirábilis, qui dedísti salútem in manu féminæ, * Exáudi preces servórum tuórum. \V\ Benedíctus es, Dómine, qui non derelínquis præsuméntes de te, et de sua virtúte gloriántes humílias. — Exáudi.

\lectioheading{Lectio ii}{}

\lettrine[lines=2]{S}{ACÉRDOS} étiam Elíachim scripsit ad univérsos qui erant contra Esdrelon, quæ est contra fáciem campi magni juxta Dóthain, et univérsos per quos viæ tránsitus esse póterat, ut obtinérent ascénsus móntium, per quos via esse póterat ad Jerúsalem, et illic custodírent ubi angústum iter esse póterat inter montes. Et fecérunt fílii Israël secúndum quod constitúerat eis sacérdos Dómini Elíachim. Et clamávit omnis pópulus ad Dóminum instántia magna, et humiliavérunt ánimas suas in jejúniis et oratiónibus ipsi et mulíeres eórum.

\R\ Tribulatiónes civitátum audívimus, quas passæ sunt, et defécimus: timor et hebetúdo mentis cécidit super nos et super líberos nostros: ipsi montes nolunt recípere fugam nostram: * Dómine, miserére. \V\ Peccávimus cum pátribus nostris, injúste égimus, iniquitátem fécimus. — Dómine.

\rubric{In Octávis et Festis ritus simplícis: Glória Patri. Dómine.}

\lectioheading{Lectio iii}{}

\lettrine[lines=2]{E}{T} induérunt se sacerdótes cilíciis, et infántes prostravérunt contra fáciem templi Dómini, et altáre Dómini operuérunt cilício; et clamavérunt ad Dóminum Deum Israël unanímiter, ne daréntur in prædam infántes eórum, et uxóres eórum in divisiónem, et civitátes eórum in extermínium, et sancta eórum in pollutiónem, et fíerent oppróbrium Géntibus. Tunc Elíachim sacérdos Dómini magnus circuívit omnem Israël allocutúsque est eos dicens: Scitóte quóniam exáudiet Dóminus preces vestras, si manéntes permanséritis in jejúniis et oratiónibus in conspéctu Dómini.

\R\ Benedícat te Dóminus in virtúte sua, qui per te ad níhilum redégit inimícos nostros: * Ut non defíciat laus tua de ore hóminum. \V\ Benedíctus Dóminus qui creávit cælum et terram; quia hódie nomen tuum ita magnificávit. — Ut. Glória Patri. Ut.

\feriaheading{Feria Tertia infra Hebdomadam IV Septembris}

\rubric{De libro Judith}

\lectioheading{Lectio i}{Cap. 8, 1-11}

\lettrine[lines=2]{E}{T} factum est, cum audísset hæc verba Judith vídua, quæ erat fília Merári, fílii Idox, fílii Joseph, fílii Oziæ, fílii Elai, fílii Jammor, fílii Gédeon, fílii Ráphaïm, fílii Achitob, fílii Melchíæ, fílii Enan, fílii Nathaníæ, fílii Saláthiel, fílii Símeon, fílii Ruben, et vir ejus fuit Manásses, qui mórtuus est in diébus messis hordeáceæ; instábat enim super alligántes manípulos in campo, et venit æstus super caput ejus, et mórtuus est in Bethúlia civitáte sua et sepúltus est illic cum pátribus suis. Erat autem Judith relícta ejus vídua jam annis tribus et ménsibus sex.

\R\ Nos álium Deum nescímus præter Dóminum, in quo sperámus: * Qui non déspicit nos, nec ámovet salútem suam a génere nostro. \V\ Indulgéntiam ipsíus fusis lácrimis postulémus, et humiliémus illi ánimas nostras. — Qui.

\lectioheading{Lectio ii}{}

\lettrine[lines=2]{E}{T} in superióribus domus suæ fecit sibi secrétum cubículum, in quo cum puéllis suis clausa morabátur. Et habens super lumbos suos cilícium, jejunábat ómnibus diébus vitæ suæ præter sábbata et neoménias et festa domus Israël. Erat autem elegánti aspéctu nimis; cui vir suus reliquérat divítias multas et famíliam copiósam ac possessiónes arméntis boum et grégibus óvium plenas. Et erat hæc in ómnibus famosíssima, quóniam timébat Dóminum valde, nec erat qui loquerétur de illa verbum malum.

\R\ Dominátor, Dómine cælórum et terræ, Creátor aquárum, Rex univérsæ creatúræ: * Exáudi oratiónem servórum tuórum. \V\ Tu, Dómine, cui humílium semper et mansuetórum plácuit deprecátio. — Exáudi.

\rubric{In Octávis et Festis ritus simplícis: Glória Patri. Exáudi.}

\lectioheading{Lectio iii}{}

\lettrine[lines=2]{H}{ÆC} ítaque, cum audísset quóniam Ozías promisísset quod, transácto quinto die, tráderet civitátem, misit ad presbýteros Chabri et Charmi. Et venérunt ad illam, et dixit illis: Quod est hoc verbum, in quo consénsit Ozías, ut tradat civitátem Assýriis, si intra quinque dies non vénerit vobis adjutórium? Et qui estis vos, qui tentátis Dóminum?

\R\ Dómine Deus, qui cónteris bella ab inítio, éleva bráchium tuum super gentes, quæ cógitant servis tuis mala: * Et déxtera tua glorificétur in nobis. \V\ Allíde virtútem eórum in virtúte tua; cadat robur eórum in iracúndia tua. — Et. Glória Patri. Et.

\feriaheading{Feria Quarta infra Hebdomadam IV Septembris}

\rubric{De libro Judith}

\lectioheading{Lectio i}{Cap. 10, 1-4}

\lettrine[lines=2]{F}{ACTUM} est autem, cum cessásset clamáre ad Dóminum surréxit de loco in quo jacúerat prostráta ad Dóminum; vocavítque abram suam, et descéndens in domum suam ábstulit a se cilícium et éxuit se vestiméntis viduitátis suæ et lavit corpus suum et unxit se myro óptimo et discriminávit crinem cápitis sui et impósuit mitram super caput suum et índuit se vestiméntis jucunditátis suæ induítque sandália pédibus suis assumpsítque dextralíola et lília et ináures et ánulos et ómnibus ornaméntis suis ornávit se. Cui étiam Dóminus cóntulit splendórem.

\R\ Confórta me, Rex Sanctórum principátum tenens: * Et da sermónem rectum et bene sonántem in os meum. \V\ Dómine, Rex univérsæ potestátis, convérte consílium eórum super eos. — Et.

\lectioheading{Lectio ii}{Cap. 10, 11-12}

\lettrine[lines=2]{F}{ACTUM} est autem, cum descénderet montem circa ortum diéi, occurrérunt ei exploratóres Assyriórum et tenuérunt eam dicéntes: Unde venis aut quo vadis? Quæ respóndit: Fília sum Hebræórum; ídeo ego fugi a fácie eórum, quóniam futúrum agnóvi quod dentur vobis in deprædatiónem, pro eo quod contemnéntes vos noluérunt ultro trádere seípsos, ut invenírent misericórdiam in conspéctu vestro.

\R\ Tribulatiónes civitátum audívimus, quas passæ sunt, et defécimus: timor et hebetúdo mentis cécidit super nos et super líberos nostros: ipsi montes nolunt recípere fugam nostram: * Dómine, miserére. \V\ Peccávimus cum pátribus nostris, injúste égimus, iniquitátem fécimus. — Dómine.

\rubric{In Octávis et Festis ritus simplícis: Glória Patri. Dómine.}

\lectioheading{Lectio iii}{Cap. 10, 16-20}

\lettrine[lines=2]{D}{UXERÚNTQUE} illam ad tabernáculum Holoférnis annuntiántes eam. Cumque intrásset ante fáciem ejus, statim captus est in suis óculis Holoférnes. Dixerúntque ad eum satéllites ejus: Quis contémnat pópulum Hebræórum, qui tam decóras mulíeres habent, ut non pro his mérito pugnáre contra eos debeámus? Videns ítaque Judith Holoférnem sedéntem in conopéo, quod erat ex púrpura et auro et smarágdo et lapídibus pretiósis intéxtum, et, cum in fáciem ejus intendísset, adorávit eum prostérnens se super terram.

\R\ Benedícat te Dóminus in virtúte sua, qui per te ad níhilum redégit inimícos nostros: * Ut non defíciat laus tua de ore hóminum. \V\ Benedíctus Dóminus qui creávit cælum et terram; quia hódie nomen tuum ita magnificávit. — Ut. Glória Patri. Ut.

\feriaheading{Feria Quinta infra Hebdomadam IV Septembris}

\rubric{Si hæc Féria occúrrat post diem 24 Septémbris, ita ut hæc sit última hebdómada mensis, hac nocte pónitur liber Esther cum suis Responsóriis, ut infra in Dominica V, 321: et légitur de eo étiam duóbus sequéntibus diébus, ita tamen ut, lecto Inítio Domínicæ assignáto, sequéntibus Fériis dicántur Lectiónes de respondénti Féria VI aut Sábbato ejúsdem hebdómadæ V.}

\rubric{De libro Judith}

\lectioheading{Lectio i}{Cap. 12, 10-13}

\lettrine[lines=2]{E}{T} factum est, in quarto die Holoférnes fecit cœnam servis suis et dixit ad Vágao eunúchum suum: Vade et suáde Hebrǽam illam ut sponte conséntiat habitáre mecum. Foedum est enim apud Assýrios, si fémina irrídeat virum agéndo ut immúnis ab eo tránseat. Tunc introívit Vágao ad Judith et dixit: Non vereátur bona puélla introíre ad dóminum meum, ut honorificétur ante fáciem ejus, ut mandúcet cum eo et bibat vinum in jucunditáte. Cui Judith respóndit: Quæ ego sum, ut contradícam dómino meo?

\R\ Adonái, Dómine, Deus magne et mirábilis, qui dedísti salútem in manu féminæ, * Exáudi preces servórum tuórum. \V\ Benedíctus es, Dómine, qui non derelínquis præsuméntes de te, et de sua virtúte gloriántes humílias. — Exáudi.

\lectioheading{Lectio ii}{Cap. 13, 1-11}

\lettrine[lines=2]{U}{T} autem sero factum est, festinavérunt servi illíus ad hospítia sua; et conclúsit Vágao óstia cubículi et ábiit. Erant autem omnes fatigáti a vino; erátque Judith sola in cubículo. Porro Holoférnes jacébat in lecto nímia ebrietáte sopítus. Dixítque Judith puéllæ suæ ut staret foris ante cubículum et observáret. Stetítque Judith ante lectum orans cum lácrimis et labiórum motu in siléntio dicens: Confírma me, Dómine, Deus Israël, et réspice in hac hora ad ópera mánuum meárum, ut, sicut promisísti, Jerúsalem civitátem tuam érigas, et hoc quod credens per te posse fíeri cogitávi, perfíciam.

\R\ Tribulatiónes civitátum audívimus, quas passæ sunt, et defécimus: timor et hebetúdo mentis cécidit super nos et super líberos nostros: ipsi montes nolunt recípere fugam nostram: * Dómine, miserére. \V\ Peccávimus cum pátribus nostris, injúste égimus, iniquitátem fécimus. — Dómine.

\rubric{In Octávis et Festis ritus simplícis: Glória Patri. Dómine.}

\lectioheading{Lectio iii}{}

\lettrine[lines=2]{E}{T}, cum hæc dixísset, accéssit ad colúmnam, quæ erat ad caput léctuli ejus, et pugiónem ejus, qui in ea ligátus pendébat, exsólvit. Cumque evaginásset illum, apprehéndit comam cápitis ejus et ait: Confírma me, Dómine Deus, in hac hora. Et percússit bis in cervícem ejus et abscídit caput ejus et ábstulit conopéum ejus a colúmnis et evólvit corpus ejus truncum. Et post pusíllum exívit et trádidit caput Holoférnis ancíllæ suæ et jussit ut mítteret illud in peram suam.

\R\ Benedícat te Dóminus in virtúte sua, qui per te ad níhilum redégit inimícos nostros: * Ut non defíciat laus tua de ore hóminum. \V\ Benedíctus Dóminus qui creávit cælum et terram; quia hódie nomen tuum ita magnificávit. — Ut. Glória Patri. Ut.

\feriaheading{Feria Sexta infra Hebdomadam IV Septembris}

\rubric{De libro Judith}

\lectioheading{Lectio i}{Cap. 15, 1-3}

\lettrine[lines=2]{C}{UMQUE} omnis exércitus decollátum Holoférnem audísset, fugit mens et consílium ab eis, et solo tremóre et metu agitáti fugæ præsídium sumunt, ita ut nullus loquerétur cum próximo suo, sed inclináto cápite, relíctis ómnibus, evádere festinábant Hebrǽos, quos armátos super se veníre audiébant, fugiéntes per vias campórum et sémitas cóllium. Vidéntes ítaque fílii Israël fugiéntes, secúti sunt illos, descenderúntque clangéntes tubis et ululántes post ipsos.

\R\ Nos álium Deum nescímus præter Dóminum, in quo sperámus: * Qui non déspicit nos, nec ámovet salútem suam a génere nostro. \V\ Indulgéntiam ipsíus fusis lácrimis postulémus, et humiliémus illi ánimas nostras. — Qui.

\lectioheading{Lectio ii}{Cap. 15, 5-7}

\lettrine[lines=2]{M}{ISIT} ítaque Ozías núntios per omnes civitátes et regiónes Israël. Omnis ítaque régio omnísque urbs eléctam juventútem armátam misit post eos; et persecúti sunt eos in ore gládii, quoúsque pervenírent ad extremitátem fínium suórum. Réliqui autem qui erant in Bethúlia, ingréssi sunt castra Assyriórum et prædam, quam fugiéntes Assýrii relíquerant, abstulérunt et onustáti sunt valde.

\R\ Dominátor, Dómine cælórum et terræ, Creátor aquárum, Rex univérsæ creatúræ: * Exáudi oratiónem servórum tuórum. \V\ Tu, Dómine, cui humílium semper et mansuetórum plácuit deprecátio. — Exáudi.

\rubric{In Octávis et Festis ritus simplícis: Glória Patri. Exáudi.}

\lectioheading{Lectio iii}{Cap. 15, 9-12}

\lettrine[lines=2]{J}{OACHIM} autem summus póntifex de Jerúsalem venit in Bethúliam cum univérsis presbýteris suis, ut vidéret Judith. Quæ cum exísset ad illum, benedixérunt eam omnes una voce dicéntes: Tu glória Jerúsalem, tu lætítia Israël, tu honorificéntia pópuli nostri. Quia fecísti viríliter, et confortátum est cor tuum, eo quod castitátem amáveris et, post virum tuum, álterum nescíeris, ídeo et manus Dómini confortávit te, et ídeo eris benedícta in ætérnum. Et dixit omnis pópulus: Fiat, fiat.

\R\ Dómine Deus, qui cónteris bella ab inítio, éleva bráchium tuum super gentes, quæ cógitant servis tuis mala: * Et déxtera tua glorificétur in nobis. \V\ Allíde virtútem eórum in virtúte tua; cadat robur eórum in iracúndia tua. — Et. Glória Patri. Et.

\feriaheading{Sabbato infra Hebdomadam IV Septembris}

\rubric{De libro Judith}

\lectioheading{Lectio i}{Cap. 16, 22-31}

\lettrine[lines=2]{E}{T} factum est post hæc, omnis pópulus post victóriam venit in Jerúsalem adoráre Dóminum; et mox ut purificáti sunt, obtulérunt omnes holocáusta et vota et repromissiónes suas. Porro Judith univérsa vasa béllica Holoférnis, quæ dedit illi pópulus, et conopéum, quod ipsa sustúlerat de cúbili ipsíus, óbtulit in anáthema oblivíónis.

\R\ Confórta me, Rex, Sanctórum principátum tenens: * Et da sermónem rectum et bene sonántem in os meum. \V\ Dómine, Rex univérsæ potestátis, convérte consílium eórum super eos. — Et.

\lectioheading{Lectio ii}{}

\lettrine[lines=2]{E}{RAT} autem pópulus jucúndus secúndum fáciem sanctórum; et per tres menses gáudium hujus victóriæ celebrátum est cum Judith. Post dies autem illos unusquísque rédiit in domum suam, et Judith magna facta est in Bethúlia et præclárior erat univérsæ terræ Israël. Erat étiam virtúti cástitas adjúncta, ita ut non cognósceret virum ómnibus diébus vitæ suæ, ex quo defúnctus est Manásses vir ejus. Erat autem, diébus festis, procédens cum magna glória.

\R\ Tribulatiónes civitátum audívimus, quas passæ sunt, et defécimus: timor et hebetúdo mentis cécidit super nos et super líberos nostros: ipsi montes nolunt recípere fugam nostram: * Dómine, miserére. \V\ Peccávimus cum pátribus nostris, injúste égimus, iniquitátem fécimus. — Dómine.

\rubric{In Octávis simplícibus et in Offício sanctæ Maríæ in Sábbato: Glória Patri. Dómine.}

\lectioheading{Lectio iii}{}

\lettrine[lines=2]{M}{ANSIT} autem in domo viri sui annos centum quinque, et dimísit abram suam líberam, et defúncta est ac sepúlta cum viro suo in Bethúlia. Luxítque illam omnis pópulus diébus septem. In omni autem spátio vitæ ejus non fuit qui perturbáret Israël, et post mortem ejus annis multis. Dies autem victóriæ hujus festivitátis ab Hebrǽis in número sanctórum diérum accípitur et cólitur a Judǽis ex illo témpore usque in præséntem diem.

\R\ Benedícat te Dóminus in virtúte sua, qui per te ad níhilum redégit inimícos nostros: * Ut non defíciat laus tua de ore hóminum. \V\ Benedíctus Dóminus qui creávit cælum et terram; quia hódie nomen tuum ita magnificávit. — Ut. Glória Patri. Ut.

\feriaheading{Sabbato ante Dominicam V Septembris}

\rubric{nisi hæc occúrrat a die 28 ad 30 Septémbris: quia tunc, sequénti Antíphona omíssa, pónitur Ant. Adapériat, quæ assignátur Sábbato ante Dominicam I Octóbris.}

\hourheading{Ad Vesperas}

\V\ Vespertína orátio ascéndat ad te, Dómine. \R\ Et descéndat super nos misericórdia tua.

\antlabel{Ad Magnif. Ant.} Dómine, Rex omnípotens, * in dicióne tua cuncta sunt pósita, et non est qui possit resístere voluntáti.

\rubric{Orátio Dominicæ occurréntis.}

% ===== BATCH 6 of the Proprium de Tempore (autumnalis.md lines
% 14412-14650) -- the complete fourth week of September: "Dominica
% IV Septembris" Matins opening the Book of Judith in Nocturn I and
% St. Ambrose's "De Elia et jejunio" in Nocturn II (same deferred
% Nocturn III pattern), the six ferial days completing the Judith
% cycle through Holofernes's death and Judith's long widowhood, and
% the brief "Sabbato ante Dominicam V Septembris" (Vespers only,
% with a rubric noting its Magnificat antiphon is replaced by
% "Adapériat" -- assigned to the Saturday before the First Sunday of
% October -- whenever this particular Saturday falls between
% September 28 and 30). Five genuine partial-caps OCR artifacts are
% silently normalized to full caps per the established rule
% (CUmque, MIsit, JOachim, ERat, MAnsit), each confirmed by direct
% inspection of autumnalis.md to be a true partial-caps defect (only
% the first one or two letters capitalized, the remainder lowercase)
% rather than a deliberate small-caps or proper-name styling; none
% carries an accent on its newly-capitalized remainder, so no accent
% is dropped in the normalization. A separate, more pervasive
% formatting quirk also appears starting partway through this
% batch (confirmed via direct grep of autumnalis.md to begin exactly
% at "Feria Tertia infra Hebdomadam IV Septembris" and continue
% through "Sabbato infra Hebdomadam IV Septembris," with zero
% occurrences before that point): many colons in the "In Octávis..."
% rubrics and the Responsories' verse markers are preceded by a
% space ("ritus simplícis :" rather than "ritus simplícis:"). Since
% this reflects only an incidental OCR/typesetting spacing variant,
% not a textual difference, and the surrounding document (including
% this very same batch's earlier material) consistently uses no
% space before a colon, this spacing is silently normalized to match
% the rest of the document, consistent with the established
% treatment of other incidental formatting variations (all-caps
% headings, lowercase Roman numerals) elsewhere in this project. The
% document continues with "Dominica V Septembris" (autumnalis.md
% line 14651), opening the Book of Esther.
% =====
\feriaheading{Dominica V Septembris}

\rubric{Semiduplex}

\nocturnoheading{In I Nocturno}

\rubric{Incipit liber Esther}

\lectioheading{Lectio i}{Cap. 1, 1-9}

\lettrine[lines=2]{I}{N} diébus Assuéri, qui regnávit ab India usque Æthiópiam super centum vigínti septem províncias, quando sedit in sólio regni sui, Susan cívitas regni ejus exórdium fuit. Tértio ígitur anno impérii sui, fecit grande convívium cunctis princípibus et púeris suis fortíssimis Persárum et Medórum ínclytis et præféctis provinciárum coram se, ut osténderet divítias glóriæ regni sui, ac magnitúdinem atque jactántiam poténtiæ suæ, multo témpore, centum vidélicet et octogínta diébus.

\R\ Dómine, mi Rex omnípotens, in dicióne tua cuncta sunt pósita, et non est qui possit resístere voluntáti tuæ: * Líbera nos propter nomen tuum. \V\ Exáudi oratiónem nostram, et convérte luctum nostrum in gáudium. — Líbera.

\lectioheading{Lectio ii}{}

\lettrine[lines=2]{C}{UMQUE} implerántur dies convívii, invitávit omnem pópulum, qui invéntus est in Susan, a máximo usque ad mínimum; et jussit septem diébus convívium præparári in vestíbulo horti et némoris, quod régio cultu et manu cónsitum erat. Et pendébant ex omni parte tentória aérii colóris et carbásini ac hyacínthini sustentáta fúnibus býssinis atque purpúreis, qui ebúrneis círculis insérti erant et colúmnis marmóreis fulciebántur. Léctuli quoque áurei et argéntei super paviméntum smarágdino et pário stratum lápide dispósiti erant, quod mira varietáte pictúra decorábat.

\R\ Confórta me, Rex, Sanctórum principátum tenens: * Et da sermónem rectum et bene sonántem in os meum. \V\ Dómine, Rex univérsæ potestátis, convérte consílium eórum super eos. — Et.

\lectioheading{Lectio iii}{}

\lettrine[lines=2]{B}{IBEBANT} autem, qui invitáti erant, áureis póculis et áliis atque áliis vasis cibi inferebántur. Vinum quoque, ut magnificéntia régia dignum erat, abúndans et præcípuum ponebátur. Nec erat qui noléntes cógeret ad bibéndum, sed, sicut rex statúerat præpónens mensis síngulos de princípibus suis, ut súmeret unusquísque quod vellet. Vasthi quoque regína fecit convívium feminárum in palátio, ubi rex Assuérus manére consuéverat.

\R\ Spem in álium numquam hábui, prætérquam in te, Deus Israël: * Qui irásceris, et propítius eris: et ómnia peccáta hóminum in tribulatióne dimíttis. \V\ Dómine Deus, Creátor cæli et terræ, réspice ad humilitátem nostram. — Qui. Glória Patri. Qui.

\nocturnoheading{In II Nocturno}

\rubric{Ex libro Officiórum sancti Ambrósii Epíscopi}

\lectioheading{Lectio iv}{Lib. 3, cap. 15}

\lettrine[lines=2]{Q}{UID} Esther regína, nonne ut pópulum suum perículo erúeret (quod erat decórum atque honéstum) morti se óbtulit, nec immítis regis trepidávit furórem? Ipse quoque rex Persárum, ferox atque túmido corde, tamen decórum judicávit índici insidiárum, quæ sibi parátæ forent, grátiam repræsentáre, populúmque líberum a servitúte erípere, erúere neci, nec párcere neci ejus, qui tam indecóra suasísset. Dénique quem secúndum a se, ac præcípuum inter omnes amícos habéret, cruci trádidit, quod dehonestátum se ejus frauduléntis consíliis animadvertísset.

\R\ Meménto mei, Dómine Deus, in bonum: * Et ne déleas miseratiónes meas quas feci in domo Dei mei et in cæremóniis ejus. \V\ Recordáre mei, Dómine, Deus meus. — Et.

\lectioheading{Lectio v}{Cap. 16}

\lettrine[lines=2]{E}{A} enim amicítia probábilis, quæ honestátem tuétur, præferénda sane ópibus, honóribus, potestátibus; honestáti vero præférri non solet, sed honestátem sequi. Qualis fuit Jónathæ, qui pro pietáte nec offénsam patris, nec salútis perículum refugiébat. Qualis fuit Achímelech, qui pro hospitális grátiæ offíciis necem pótius sibi, quam proditiónem fugiéntis amíci, subeúndam arbitrabátur. Nihil ígitur præferéndum honestáti; quæ tamen ne amicítiæ stúdio prætereátur, étiam hoc Scriptúra ádmonet.

\R\ Tribulatiónes civitátum audívimus, quas passæ sunt, et defécimus: timor et hebetúdo mentis cécidit super nos et super líberos nostros: ipsi montes nolunt recípere fugam nostram: * Dómine, miserére. \V\ Peccávimus cum pátribus nostris, injúste égimus, iniquitátem fécimus. — Dómine.

\lectioheading{Lectio vi}{}

\lettrine[lines=2]{S}{UNT} enim pleræque philosophórum quæstiónes: Utrum amíci causa quisquam contra pátriam sentíre necne débeat, ut amíco obédiat: utrum opórteat ut fidem déserat, dum indúlget atque inténdit amíci commoditátibus. Et Scriptúra quidem ait: Clava, et gládius, et sagítta ferráta, sic homo est testimónium dans falsum advérsus amícum suum. Sed consídera quid ástruat. Non testimónium reprehéndit dictum in amícum, sed falsum testimónium. Quid enim si Dei causa, quid si pátriæ, cogátur áliquis dícere testimónium? Numquid præponderáre debet amicítia religióni, præponderáre caritáti vítium?

\R\ Benedícat te Dóminus in virtúte sua, qui per te ad níhilum redégit inimícos nostros: * Ut non defíciat laus tua de ore hóminum. \V\ Benedíctus Dóminus qui creávit cælum et terram; quia hódie nomen tuum ita magnificávit. — Ut. Glória Patri. Ut.

\R\ Nos álium Deum nescímus præter Dóminum, in quo sperámus: * Qui non déspicit nos, nec ámovet salútem suam a génere nostro. \V\ Indulgéntiam ipsíus fusis lácrimis postulémus, et humiliémus illi ánimas nostras. — Qui.

\R\ Duo Séraphim, \rubric{ut infra in Homilíis.}

\feriaheading{Feria Secunda infra Hebdomadam V Septembris}

\rubric{De libro Esther}

\lectioheading{Lectio i}{Cap. 2, 5-11}

\lettrine[lines=2]{E}{RAT} vir Judǽus in Susan civitáte, vocábulo Mardochǽus, fílius Jair, fílii Sémei, fílii Cis, de stirpe Jémini, qui translátus fúerat de Jerúsalem eo témpore quo Jechoníam regem Juda Nabuchodónosor rex Babylónis transtúlerat. Qui fuit nutrícius fíliæ fratris sui Edíssæ, quæ áltero nómine vocabátur Esther, et utrúmque paréntem amíserat: pulchra nimis et decóra fácie. Mortuísque patre ejus ac matre, Mardochǽus sibi eam adoptávit in fíliam.

\R\ Dómine, mi Rex omnípotens, in dicióne tua cuncta sunt pósita, et non est qui possit resístere voluntáti tuæ: * Líbera nos propter nomen tuum. \V\ Exáudi oratiónem nostram, et convérte luctum nostrum in gáudium. — Líbera.

\lectioheading{Lectio ii}{}

\lettrine[lines=2]{C}{UMQUE} percrebuísset regis impérium, et juxta mandátum illíus multæ pulchræ vírgines adduceréntur Susan et Egéo traderéntur eunúcho; Esther quoque inter céteras puéllas ei trádita est, ut servarétur in número feminárum. Quæ plácuit ei et invénit grátiam in conspéctu illíus. Et præcépit eunúcho ut acceleráret mundum muliébrem et tráderet ei partes suas et septem puéllas speciosíssimas de domo regis, et tam ipsam quam pedíssequas ejus ornáret atque excóleret. Quæ nóluit indicáre ei pópulum et pátriam suam; Mardochǽus enim præcéperat ei, ut de hac re omníno reticéret. Qui deambulábat quotídie ante vestíbulum domus, in qua eléctæ vírgines servabántur, curam agens salútis Esther.

\R\ Confórta me, Rex, Sanctórum principátum tenens: * Et da sermónem rectum et bene sonántem in os meum. \V\ Dómine, Rex univérsæ potestátis, convérte consílium eórum super eos. — Et.

\rubric{In Octávis et Festis ritus simplícis: Glória Patri. Et.}

\lectioheading{Lectio iii}{Cap. 2, 15-17}

\lettrine[lines=2]{E}{VOLUTO} autem témpore per órdinem, instábat dies quo Esther fília Abíhail fratris Mardochǽi, quam sibi adoptáverat in fíliam, debéret intráre ad regem. Quæ non quæsívit muliébrem cultum, sed quæcúmque vóluit Egéus eunúchus custos vírginum, hæc ei ad ornátum dedit; erat enim formósa valde et incredíbili pulchritúdine: ómnium óculis gratiósa et amábilis videbátur. Ducta est ítaque ad cubículum regis Assuéri mense décimo, qui vocátur Tebeth, séptimo anno regni ejus. Et adamávit eam rex plus quam omnes mulíeres, habuítque grátiam et misericórdiam coram eo super omnes mulíeres, et pósuit diadéma regni in cápite ejus fecítque eam regnáre in loco Vasthi.

\R\ Spem in álium numquam hábui, prætérquam in te, Deus Israël: * Qui irásceris, et propítius eris: et ómnia peccáta hóminum in tribulatióne dimíttis. \V\ Dómine Deus, Creátor cæli et terræ, réspice ad humilitátem nostram. — Qui. Glória Patri. Qui.

\feriaheading{Feria Tertia infra Hebdomadam V Septembris}

\rubric{De libro Esther}

\lectioheading{Lectio i}{Cap. 3, 1-7}

\lettrine[lines=2]{P}{OST} hæc rex Assuérus exaltávit Aman fílium Amadáthi, qui erat de stirpe Agag, et pósuit sólium ejus super omnes príncipes quos habébat. Cunctíque servi regis, qui in fóribus palátii versabántur, flectébant génua et adorábant Aman; sic enim præcéperat eis imperátor. Solus Mardochǽus non flectébat genu neque adorábat eum. Cui dixérunt púeri regis, qui ad fores palátii præsidébant: Cur præter céteros non servas mandátum regis?

\R\ Meménto mei, Dómine Deus, in bonum: * Et ne déleas miseratiónes meas quas feci in domo Dei mei et in cæremóniis ejus. \V\ Recordáre mei, Dómine, Deus meus. — Et.

\lectioheading{Lectio ii}{}

\lettrine[lines=2]{C}{UMQUE} hoc crébrius dícerent, et ille nollet audíre, nuntiavérunt Aman, scire cupiéntes utrum perseváret in senténtia; díxerat enim eis se esse Judǽum. Quod cum audísset Aman et experiménto probásset quod Mardochǽus non fléctere sibi genu nec se adoráret, irátus est valde et pro níhilo duxit in unum Mardochǽum míttere manus suas.

\R\ Tribulatiónes civitátum audívimus, quas passæ sunt, et defécimus: timor et hebetúdo mentis cécidit super nos et super líberos nostros: ipsi montes nolunt recípere fugam nostram: * Dómine, miserére. \V\ Peccávimus cum pátribus nostris, injúste égimus, iniquitátem fécimus. — Dómine.

\rubric{In Octávis et Festis ritus simplícis: Glória Patri. Dómine.}

\lectioheading{Lectio iii}{}

\lettrine[lines=2]{A}{UDIERAT} enim quod esset gentis Judǽæ, magísque vóluit omnem Judæórum rum, qui erant in regno Assuéri, pérdere natiónem. Mense primo, cujus vocábulum est Nisan, anno duodécimo regni Assuéri, missa est sors in urnam, quæ Hebráice dícitur phur, coram Aman, quo die et quo mense gens Judæórum debéret intérfici; et exívit mensis duodécimus qui vocátur Adar.

\R\ Benedícat te Dóminus in virtúte sua, qui per te ad níhilum redégit inimícos nostros: * Ut non defíciat laus tua de ore hóminum. \V\ Benedíctus Dóminus qui creávit cælum et terram; quia hódie nomen tuum ita magnificávit. — Ut. Glória Patri. Ut.

\feriaheading{Feria Quarta infra Hebdomadam V Septembris}

\rubric{De libro Esther}

\lectioheading{Lectio i}{Cap. 4, 1-17}

\lettrine[lines=2]{Q}{UÆ} cum audísset Mardochǽus, scidit vestiménta sua et indútus est sacco spargens cínerem cápiti; et in platéa médiæ civitátis voce magna clamábat osténdens amaritúdinem ánimi sui, et hoc ejulátu usque ad fores palátii grádiens; non enim erat lícitum indútum sacco aulam regis intráre. In ómnibus quoque províncíis, óppidis ac locis, ad quæ crudéle regis dogma pervénerat, planctus ingens erat apud Judǽos, jejúnium, ululátus et fletus, sacco et cínere multis pro strato uténtibus. Ingréssæ autem sunt puéllæ Esther et eunúchi nuntiaverúntque ei. Quod áudiens consternáta est et vestem misit, ut, abláto sacco, indúerent eum: quam accípere nóluit. Accitóque Athach eunúcho, quem rex minístrum ei déderat, præcépit ei ut iret ad Mardochǽum et discéret ab eo cur hoc fáceret.

\R\ Nos álium Deum nescímus præter Dóminum, in quo sperámus: * Qui non déspicit nos, nec ámovet salútem suam a génere nostro. \V\ Indulgéntiam ipsíus fusis lácrimis postulémus, et humiliémus illi ánimas nostras. — Qui.

\lectioheading{Lectio ii}{}

\lettrine[lines=2]{E}{GRESSUSQUE} Athach ivit ad Mardochǽum stantem in platéa civitátis ante óstium palátii. Qui indicávit ei ómnia quæ accíderant, quómodo Aman promisísset, ut in thesáuros regis pro Judæórum nece inférret argéntum; exémplar quoque edícti, quod pendébat in Susan, dedit ei ut regínæ osténderet et monéret eam ut intráret ad regem et deprecarétur eum pro pópulo suo. Regréssus Athach nuntiávit Esther ómnia quæ Mardochǽus díxerat. Quæ respóndit ei et jussit ut díceret Mardochǽo: Omnes servi regis et cunctæ quæ sub dicióne ejus sunt norunt províncíæ, quod sive vir sive múlier non vocátus intérius átrium regis intráverit, absque ulla cunctatióne statim interficiátur, nisi forte rex áuream virgam ad eum teténderit pro signo cleméntiæ, atque ita possit vívere. Ego ígitur quómodo ad regem intráre pótero, quæ trigínta jam diébus non sum vocáta ad eum?

\R\ Confórta me, Rex, Sanctórum principátum tenens: * Et da sermónem rectum et bene sonántem in os meum. \V\ Dómine, Rex univérsæ potestátis, convérte consílium eórum super eos. — Et.

\rubric{In Octávis et Festis ritus simplícis: Glória Patri. Et.}

\lectioheading{Lectio iii}{}

\lettrine[lines=2]{Q}{UOD} cum audísset Mardochǽus, rursum mandávit Esther dicens: Ne putes quod ánimam tuam tantum líberes, quia in domo regis es, præ cunctis Judǽis; si enim nunc silúeris, per áliam occasiónem liberabúntur Judǽi, et tu et domus patris tui períbitis. Et quis novit utrum idcírco ad regnum véneris, ut in tali témpore pararéris? Rursúmque Esther hæc Mardochǽo verba mandávit: Vade et cóngrega omnes Judǽos, quos in Susan repéreris, et oráte pro me. Non comedátis et non bibátis tribus diébus et tribus nóctibus, et ego cum ancíllis meis simíliter jejunábo; et tunc ingrédiar ad regem contra legem fáciens non vocáta tradénsque me morti et perículo. Ivit ítaque Mardochǽus, et fecit ómnia quæ ei Esther præcéperat.

\R\ Spem in álium numquam hábui, prætérquam in te, Deus Israël: * Qui irásceris, et propítius eris: et ómnia peccáta hóminum in tribulatióne dimíttis. \V\ Dómine Deus, Creátor cæli et terræ, réspice ad humilitátem nostram. — Qui. Glória Patri. Qui.

\feriaheading{Feria Quinta infra Hebdomadam V Septembris}

\rubric{De libro Esther}

\lectioheading{Lectio i}{Cap. 5, 1-5}

\lettrine[lines=2]{D}{IE} autem tértio indúta est Esther regálibus vestiméntis et stetit in átrio domus régiæ, quod erat intérius contra basílicam regis. At ille sedébat super sólium suum in consistório palátii contra óstium domus. Cumque vidísset Esther regínam stantem, plácuit óculis ejus et exténdit contra eam virgam áuream quam tenébat manu. Quæ accédens osculáta est summitátem virgæ ejus, dixítque ad eam rex: Quid vis, Esther regína? quæ est petítio tua? Etiam si dimídiam partem regni petíeris, dábitur tibi. At illa respóndit: Si regi placet, óbsecro ut vénias ad me hódie et Aman tecum ad convívium quod parávi. Statímque rex, Vocáte, inquit, cito Aman ut Esther obédiat voluntáti.

\R\ Dómine, mi Rex omnípotens, in dicióne tua cuncta sunt pósita, et non est qui possit resístere voluntáti tuæ: * Líbera nos propter nomen tuum. \V\ Exáudi oratiónem nostram, et convérte luctum nostrum in gáudium. — Líbera.

\lectioheading{Lectio ii}{Cap. 5, 9-14}

\lettrine[lines=2]{E}{GRESSUS} est ítaque illo die Aman lætus et álacer. Cumque vidísset Mardochǽum sedéntem ante fores palátii, et non solum non assurrexísse sibi, sed nec motum quidem de loco sessiónis suæ, indignátus est valde; et, dissimuláta ira, revérsus in domum suam convocávit ad se amícos suos et Zares uxórem suam et expósuit illis magnitúdinem divitiárum suárum filiorúmque turbam et quanta eum glória super omnes príncipes et servos suos rex elevásset. Et post hæc ait: Regína quoque Esther nullum álium vocávit ad convívium cum rege, præter me, apud quam étiam cras cum rege pransúrus sum; et, cum hæc ómnia hábeam, nihil me habére puto, quámdiu vídero Mardochǽum Judǽum sedéntem ante fores régias.

\R\ Confórta me, Rex, Sanctórum principátum tenens: * Et da sermónem rectum et bene sonántem in os meum. \V\ Dómine, Rex univérsæ potestátis, convérte consílium eórum super eos. — Et.

\rubric{In Octávis et Festis ritus simplícis: Glória Patri. Et.}

\lectioheading{Lectio iii}{}

\lettrine[lines=2]{R}{ESPONDERUNTQUE} ei Zares uxor ejus et céteri amíci: Jube parári excélsam trabem habéntem altitúdinis quinquagínta cúbitos, et dic mane regi ut appendátur super eam Mardochǽus; et sic ibis cum rege lætus ad convívium. Plácuit ei consílium et jussit excélsam paráti crucem.

\R\ Spem in álium numquam hábui, prætérquam in te, Deus Israël: * Qui irásceris, et propítius eris: et ómnia peccáta hóminum in tribulatióne dimíttis. \V\ Dómine Deus, Creátor cæli et terræ, réspice ad humilitátem nostram. — Qui. Glória Patri. Qui.

\feriaheading{Feria Sexta infra Hebdomadam V Septembris}

\rubric{De libro Esther}

\lectioheading{Lectio i}{Cap. 6, 1-13}

\lettrine[lines=2]{N}{OCTEM} illam duxit rex insómnem jussítque sibi afférri histórias et annáles priórum témporum. Quæ cum illo præsénte legeréntur, ventum est ad illum locum, ubi scriptum erat quómodo nuntiásset Mardochǽus insídias Bágathan et Thares eunuchórum regem Assuérum juguláre cupiéntium. Quod cum audísset rex ait: Quid pro hac fide honóris ac præmii Mardochǽus consecútus est? Dixérunt ei servi illíus ac minístri: Níhil omníno mercédis accépit. Statímque rex, Quis est, inquit, in átrio? Aman quippe intérius átrium domus régiæ intráverat, ut suggéreret regi et jubéret Mardochǽum affígi patíbulo, quod ei fúerat præparátum. Respondérunt púeri: Aman stat in átrio. Dixítque rex: Ingrediátur.

\R\ Meménto mei, Dómine Deus, in bonum: * Et ne déleas miseratiónes meas quas feci in domo Dei mei et in cæremóniis ejus. \V\ Recordáre mei, Dómine, Deus meus. — Et.

\lectioheading{Lectio ii}{}

\lettrine[lines=2]{C}{UMQUE} esset ingréssus ait illi: Quid debet fíeri viro, quem rex honoráre desíderat? Cógitans autem in corde suo Aman et réputans quod nullum álium rex nisi se vellet honoráre, respóndit: Homo quem rex honoráre cupit debet índui véstibus régiis et impóni super equum qui de sella regis est, et accípere régium diadéma super caput suum, et primus de régiis princípibus ac tyránnis téneat equum ejus et per platéam civitátis incédens clamet et dicat: Sic honorábitur quemcúmque volúerit rex honoráre.

\R\ Tribulatiónes civitátum audívimus, quas passæ sunt, et defécimus: timor et hebetúdo mentis cécidit super nos et super líberos nostros: ipsi montes nolunt recípere fugam nostram: * Dómine, miserére. \V\ Peccávimus cum pátribus nostris, injúste égimus, iniquitátem fécimus. — Dómine, miserére.

\rubric{In Octávis et Festis ritus simplícis: Glória Patri. Dómine.}

\lectioheading{Lectio iii}{}

\lettrine[lines=2]{D}{IXITQUE} ei rex: Festína et, sumpta stola et equo, fac, ut locútus es, Mardochǽo Judǽo, qui sedet ante fores palátii. Cave ne quidquam de his quæ locútus es prætermíttas. Tulit ítaque Aman stolam et equum, indutúmque Mardochǽum in platéa civitátis et impósitum equo præcedébat atque clamábat: Hoc honóre condígnus est quemcúmque rex volúerit honoráre. Reversúsque est Mardochǽus ad jánuam palátii, et Aman festinávit ire in domum suam lugens et opérto cápite. Narravítque Zares uxóri suæ et amícis ómnia quæ eveníssent sibi.

\R\ Benedícat te Dóminus in virtúte sua, qui per te ad níhilum redégit inimícos nostros: * Ut non defíciat laus tua de ore hóminum. \V\ Benedíctus Dóminus qui creávit cælum et terram; quia hódie nomen tuum ita magnificávit. — Ut. Glória Patri. Ut.

\feriaheading{Sabbato infra Hebdomadam V Septembris}

\rubric{De libro Esther}

\lectioheading{Lectio i}{Cap. 7, 1-10}

\lettrine[lines=2]{I}{NTRAVIT} ítaque rex et Aman, ut bíberent cum regína. Dixítque ei rex étiam secúnda die, postquam vino incalúerat: Quæ est petítio tua, Esther, ut detur tibi, et quid vis fíeri? Etiam si dimídiam partem regni petíeris, impetrábis. Ad quem illa respóndit: Si invéni grátiam in óculis tuis, o rex, et si tibi placet, dona mihi ánimam meam, pro qua rogo, et pópulum meum, pro quo óbsecro. Tráditi enim sumus ego et pópulus meus, ut conterámur, jugulémur et pereámus. Atque, útinam in servos et fámulas venderémur: esset tolerábile malum et gemens tacérem; nunc autem hostis noster est cujus crudélitas redúndat in regem.

\R\ Nos álium Deum nescímus præter Dóminum, in quo sperámus: * Qui non déspicit nos, nec ámovet salútem suam a génere nostro. \V\ Indulgéntiam ipsíus fusis lácrimis postulémus, et humiliémus illi ánimas nostras. — Qui.

\lectioheading{Lectio ii}{}

\lettrine[lines=2]{R}{ESPONDENSQUE} rex Assuérus ait: Quis est iste et cujus poténtiæ, ut hæc áudeat fácere? Dixítque Esther: Hostis et inimícus noster péssimus iste est Aman. Quod ille áudiens íllico obstúpuit vultum regis ac regínæ ferre non sústinens. Rex autem irátus surréxit et de loco convívii intrávit in hortum arbóribus cónsitum. Aman quoque surréxit ut rogáret Esther regínam pro ánima sua; intelléxit enim a rege sibi parátum malum.

\R\ Confórta me, Rex, Sanctórum principátum tenens: * Et da sermónem rectum et bene sonántem in os meum. \V\ Dómine, Rex univérsæ potestátis, convérte consílium eórum super eos. — Et.

\rubric{In Octávis simplícibus et in Offício Sanctæ Maríæ in Sábbato: Glória Patri. Et.}

\lectioheading{Lectio iii}{}

\lettrine[lines=2]{Q}{UI} cum revérsus esset de horto nemóribus cónsito et intrásset convívii locum, réperit Aman super léctulum corruísse in quo jacébat Esther et ait: Etiam regínam vult opprímere, me præsénte, in domo mea. Necdum verbum de ore regis exíerat, et statim operuérunt fáciem ejus. Dixítque Harbóna unus de eunúchis qui stabant in ministério regis: En lignum, quod paráverat Mardochǽo qui locútus est pro rege, stat in domo Aman habens altitúdinis quinquagínta cúbitos. Cui dixit rex: Appéndite eum in eo. Suspénsus est ítaque Aman in patíbulo, quod paráverat Mardochǽo, et regis ira quiévit.

\R\ Spem in álium numquam hábui, prætérquam in te, Deus Israël: * Qui irásceris, et propítius eris: et ómnia peccáta hóminum in tribulatióne dimíttis. \V\ Dómine Deus, Creátor cæli et terræ, réspice ad humilitátem nostram. — Qui. Glória Patri. Qui.

\feriaheading{Sabbato ante Dominicam I Octobris}

\hourheading{Ad Vesperas}

\V\ Vespertína orátio ascéndat ad te, Dómine. \R\ Et descéndat super nos misericórdia tua.

\antlabel{Ad Magnif. Ant.} Adapériat Dóminus * cor vestrum in lege sua et in præcéptis suis et fáciat pacem Dóminus Deus noster.

\rubric{Oratio Dominicæ occurréntis.}

% ===== BATCH 7 of the Proprium de Tempore (autumnalis.md lines
% 14651-14883) -- the complete fifth and final week of September:
% "Dominica V Septembris" Matins opening the Book of Esther in
% Nocturn I and St. Ambrose's "De Officiis" in Nocturn II (this
% Sunday's Nocturn III responsories given without the "vii"/"viii"
% numbering used at the parallel point in Dominicas I-IV, a genuine
% source variation transcribed exactly as given), the six ferial
% days completing the Esther cycle through Haman's downfall, and
% the brief "Sabbato ante Dominicam I Octobris" Vespers with the
% "Adapériat Dóminus" antiphon that the governing rubric several
% batches ago had flagged as proper to this point. This closes out
% the entire month of September in the Proprium de Tempore. A
% dozen-plus genuine partial-caps OCR artifacts are silently
% normalized to full caps per the established rule (ERat, CUmque
% [several occurrences], EVolúto, POst, AUdíerat, EGressúsque, QUod,
% EGréssus, RESponderúntque, NOctem, DIxítque, INtrávit,
% RESpondénsque), each confirmed by direct inspection of
% autumnalis.md; accents on each word's newly-capitalized remainder
% are dropped per the standing rule, while complete short words
% already fully capped in the source (QUÆ, QUI, DIE) are left
% untouched. The pervasive space-before-colon formatting quirk
% identified in the previous batch continues throughout this one and
% is likewise silently normalized to standard colon spacing. The
% document continues with "Dominica I Octobris" (autumnalis.md line
% 14884), opening the Books of Machabees and the month of October.
% =====
\feriaheading{Dominica I Octobris}

\rubric{Semiduplex}

\hourheading{Ad Matutinum}

\rubric{Invitatorium Adorémus Dóminum et Hymnus Primo die, ut in Psaltério 32: quæ dicúntur étiam in réliquis Dominícis usque ad V Novémbris inclusíve.}

\nocturnoheading{In I Nocturno}

\rubric{Incipit liber primus Machabæórum}

\lectioheading{Lectio i}{Cap. 1, 1-16}

\lettrine[lines=2]{E}{T} factum est, postquam percússit Alexánder Philíppi Mácedo, qui primus regnávit in Grǽcia, egréssus de terra Cethim, Dárium regem Persárum et Medórum, constítuit prǽlia multa et obtínuit ómnium munitiónes et interfécit reges terræ et pertránsiit usque ad fines terræ et accépit spólia multitúdinis géntium, et síluit terra in conspéctu ejus. Et congregávit virtútem et exércitum fortem nimis, et exaltátum est et elevátum cor ejus, et obtínuit regiónes géntium et tyránnos, et facti sunt illi in tribútum. Et post hæc décidit in lectum et cognóvit quia moerétur et vocávit púeros suos nóbiles, qui secum erant nutríti a juventúte, et divísit illis regnum suum, cum adhuc víveret.

\R\ Adapériat Dóminus cor vestrum in lege sua et in præcéptis suis et fáciat pacem in diébus vestris: * Concédat vobis salútem, et rédimat vos a malis. \V\ Exáudiat Dóminus oratiónes vestras et reconciliétur vobis nec vos déserat in témpore malo. — Concédat.

\lectioheading{Lectio ii}{}

\lettrine[lines=2]{E}{T} regnávit Alexánder annis duódecim et mórtuus est. Et obtinuérunt púeri ejus regnum, unusquísque in loco suo, et imposuérunt omnes sibi diadémata post mortem ejus, et fílii eórum post eos annis multis. Et multiplicáta sunt mala in terra. Et éxiit ex eis radix peccátrix, Antíochus illústris, fílius Antíochi regis qui fúerat Romæ obses, et regnávit in anno centésimo trigésimo séptimo regni Grræcórum.

\R\ Exáudiat Dóminus oratiónes vestras, et reconciliétur vobis nec vos déserat in témpore malo * Dóminus, Deus noster. \V\ Det vobis cor ómnibus, ut colátis eum et faciátis ejus voluntátem. — Dóminus.

\lectioheading{Lectio iii}{}

\lettrine[lines=2]{I}{N} diébus illis exiérunt ex Israël fílii iníqui et suasérunt multis dicéntes: Eámus et disponámus testaméntum cum géntibus, quæ circa nos sunt, quia, ex quo recéssimus ab eis, invenérunt nos multa mala. Et bonus visus est sermo in óculis eórum. Et destinavérunt áliqui de pópulo et abiérunt ad regem, et dedit illis potestátem ut fácerent justítiam géntium; et ædificavérunt gymnásium in Jerosólymis secúndum leges natiónum et fecérunt sibi præpútia et recessérunt a testaménto sancto et juncti sunt natiónibus et venúmdati sunt ut fácerent malum.

\R\ Congregáti sunt inimíci nostri, et gloriántur in virtúte sua: cóntere fortitúdinem illórum, Dómine, et dispérge illos: * Ut cognóscant quia non est álius qui pugnet pro nobis, nisi tu, Deus noster. \V\ Dispérge illos in virtúte tua, et déstrue eos, protéctor noster, Dómine. — Ut. Glória Patri. Ut.

\nocturnoheading{In II Nocturno}

\rubric{Ex libro Officiórum sancti Ambrósii Epíscopi}

\lectioheading{Lectio iv}{Lib. 1, cap. 40}

\lettrine[lines=2]{F}{ORTASSE} áliquos béllica defíxos glória tenet, ut putent solam esse præliárem fortitúdinem; et ídeo me ad ista deflexísse, quia illa nostris déforet. Quam fortis Jesus Nave, ut uno prǽlio quinque reges captos stérneret cum pópulis suis! Deínde cum advérsum Gabaonítas urgéret prǽlium, et vererétur ne nox impedíret victóriam, magnitúdine mentis et fídei clamávit: Stet sol, et stetit, donec victória consummarétur. Gédeon in trecéntis viris de ingénti pópulo et acérbo hoste revéxit triúmphum. Jónathas adoléscens virtútem magnam fecit in prǽlio.

\R\ Impetum inimicórum ne timuéritis: mémores estóte quómodo salvi facti sunt patres nostri: * Et nunc clamémus in cælum et miserébitur nostri Deus noster. \V\ Mementóte mirabílium ejus, quæ fecit pharaóni et exercítui ejus in Mari Rubro. — Et.

\lectioheading{Lectio v}{}

\lettrine[lines=2]{Q}{UID} de Machabǽis loquar? Sed prius de pópulo dicam patrum; qui cum essent paráti ad repugnándum pro templo Dei et pro legítimis suis, dolo hóstium die lacessíti sábbati, maluérunt vulnéribus offérre nuda córpora, quam repugnáre, ne violárent sábbatum. Itaque omnes læti se obtulérunt morti. Sed Machabǽi considerántes quod hoc exémplo gens omnis posset períre, sábbato étiam, cum ipsi in bellum provocaréntur, ulti sunt innocéntium necem fratrum suórum. Unde póstea stimulátus rex Antíochus, cum bellum accénderet per duces suos Lýsiam, Nicánorem, Górgiam, ita cum Orientálibus suis et Assýriis attrítus est cópiis, ut quadragínta et octo míllia in médio campi a tribus míllibus prosterneréntur.

\R\ Congregátæ sunt gentes in multitúdine, ut dímicent contra nos, et ignorámus quid ágere debeámus: * Dómine Deus, ad te sunt óculi nostri, ne pereámus. \V\ Tu scis quæ cógitant in nos: quómodo potérimus subsístere ante fáciem illórum, nisi tu ádjuves nos? — Dómine.

\lectioheading{Lectio vi}{}

\lettrine[lines=2]{V}{IRTUTEM} ducis Judæ Machabǽi de uno ejus mílite consideráte. Namque Eleázarus, cum supereminéntem céteris elephántem lorica vestítum régia advérteret, arbitrátus quod in eo esset rex, cursu cóncito in médium legiónis se prorúpit: et, abjécto clýpeo, utráque manu interficiébat, donec perveníret ad béstiam, atque intrávit sub eam, et subjécto gládio interémit eam. Itaque cadens béstia oppréssit Eleázarum, atque ita mórtuus est. Quanta ígitur virtus ánimi! primo, ut mortem non timéret; deínde, ut circumfúsus legiónibus inimicórum, in confértos raperétur hostes, médium penetráret agmen, et contémpta morte ferócior, abjécto clýpeo, utráque manu vulnerátæ molem béstiæ subíret ac sustinéret: post infra ipsam succéderet, quo plenióri feríret ictu; cujus ruína inclúsus magis quam oppréssus, suo est sepúltus triúmpho.

\R\ Tua est poténtia, tuum regnum, Dómine: tu es super omnes gentes: * Da pacem, Dómine, in diébus nostris. \V\ Creátor ómnium, Deus, terríbilis et fortis, justus, et miséricors. — Da. Glória Patri. Da.

\R\ Refúlsit sol in clýpeos áureos, et resplenduérunt montes ab eis: * Et fortitúdo géntium dissipáta est. \V\ Erat enim exércitus magnus valde et fortis: et appropiávit Judas et exércitus ejus in prǽlio. — Et.

\R\ Duo Séraphim, \rubric{ut infra in Homilíis.}

\rubric{Ad Laudes, in ómnibus Domínicis usque ad V Novémbris inclusíve, Hymnus Ætérne rerum, ut in Psaltério.}

\feriaheading{Feria Secunda infra Hebdomadam I Octobris}

\rubric{De libro primo Machabæórum}

\lectioheading{Lectio i}{Cap. 1, 17-29}

\lettrine[lines=2]{E}{T} parátum est regnum in conspéctu Antíochi, et cœpit regnáre in terra Ægýpti, ut regnáret super duo regna, et intrávit in Ægýptum in multitúdine gravi, in cúrribus et elephántis et equítibus et copiósa návium multitúdine, et constítuit bellum advérsus Ptolemǽum regem Ægýpti, et véritus est Ptolemǽus a fácie ejus et fugit, et cecidérunt vulneráti multi, et comprehéndit civitátes munítas in terra Ægýpti et accépit spólia terræ Ægýpti.

\R\ Dixit Judas Simóni fratri suo: Elige tibi viros et vade, líbera fratres tuos in Galilǽam; ego autem et Jónathas frater tuus íbimus in Galaadítim: * Sicut fúerit volúntas in cælo, sic fiat. \V\ Accingímini, fílii poténtes, et estóte paráti: quóniam mélius est nobis mori in bello, quam vidére mala gentis nostræ et sanctórum. — Sicut.

\lectioheading{Lectio ii}{}

\lettrine[lines=2]{E}{T} convértit Antíochus, postquam percússit Ægýptum in centésimo et quadragésimo tértio anno, et ascéndit ad Israël. Et ascéndit Jerosólymam in multitúdine gravi et intrávit in sanctificatiónem cum supérbia et accépit altáre áureum et candelábrum lúminis et univérsa vasa ejus et mensam propositiónis et libatória et phíalas et mortaríola áurea et velum et corónas et ornaméntum áureum, quod in fácie templi erat, et commínuit ómnia.

\R\ Ornavérunt fáciem templi corónis áureis, et dedicavérunt altáre Dómino: * Et facta est lætítia magna in pópulo. \V\ In hymnis et confessiónibus benedicébant Dóminum. — Et.

\rubric{In Octávis et Festis ritus simplícis: Glória Patri. Et.}

\lectioheading{Lectio iii}{}

\lettrine[lines=2]{E}{T} accépit argéntum et aurum et vasa concupiscibília et accépit thesáuros occúltos, quos invénit, et, sublátis ómnibus, ábiit in terram suam et fecit cædem hóminum et locútus est in supérbia magna et factus est planctus magnus in Israël et in omni loco eórum, et ingemuérunt príncipes et senióres, vírgines et júvenes infirmáti sunt, et speciósitas mulíerum immutáta est. Omnis marítus sumpsit laméntum, et quæ sedébant in thoro maritáli lugébant, et commóta est terra super habitántes in ea, et univérsa domus Jacob índuit confusiónem.

\R\ In hymnis et confessiónibus benedicébant Dóminum, * Qui magna fecit in Israël, et victóriam dedit illis Dóminus omnípotens. \V\ Ornavérunt fáciem templi corónis áureis, et dedicavérunt altáre Dómino. — Qui. Glória Patri. Qui.

\feriaheading{Feria Tertia infra Hebdomadam I Octobris}

\rubric{De libro primo Machabæórum}

\lectioheading{Lectio i}{Cap. 2, 1-10}

\lettrine[lines=2]{I}{N} diébus illis surréxit Mathathías fílius Joánnis, fílii Simeónis, sacérdos ex fíliis Jóarib ab Jerúsalem et consédit in monte Modin. Et habébat fílios quinque, Joánnem, qui cognominabátur Gaddis, et Simónem, qui cognominabátur Thasi, et Judam, qui vocabátur Machabǽus, et Eleázarum, qui cognominabátur Abaron, et Jónathan, qui cognominabátur Apphus. Hi vidérunt mala quæ fiébant in pópulo Juda et in Jerúsalem.

\R\ Hic est fratrum amátor et pópuli Israël: * Hic est, qui multum orat pro pópulo et univérsa sancta civitáte Jerúsalem. \V\ Vir iste in pópulo suo mitíssimus appáruit. — Hic est, qui.

\lectioheading{Lectio ii}{}

\lettrine[lines=2]{E}{T} dixit Mathathías: Væ mihi! ut quid natus sum vidére contritiónem pópuli mei et contritiónem civitátis sanctæ et sedére illic, cum datur in mánibus inimicórum? Sancta in manu extraneórum facta sunt, templum ejus sicut homo ignóbilis, vasa glóriæ ejus captíva abdúcta sunt, trucidáti sunt senes ejus in platéis, et júvenes ejus cecidérunt in gládio inimicórum. Quæ gens non hereditávit regnum ejus et non obtínuit spólia ejus?

\R\ Tu, Dómine universórum, qui nullam habes indigéntiam, voluísti templum tuum fíeri in nobis; * Consérva domum istam immaculátam in ætérnum, Dómine. \V\ Tu elegísti, Dómine, domum istam ad invocándum nomen tuum in ea, ut esset domus oratiónis et obsecratiónis pópulo tuo. — Consérva.

\rubric{In Octávis et Festis ritus simplícis: Glória Patri. Consérva.}

\lectioheading{Lectio iii}{Cap. 2, 14-16}

\lettrine[lines=2]{E}{T} scidit vestiménta sua Mathathías et fílii ejus et operuérunt se cilíciis et planxérunt valde. Et venérunt illuc qui missi erant a rege Antíocho, ut cógerent eos qui confúgerant in civitátem Modin, immoláre et accéndere thura et a lege Dei discédere, et multi de pópulo Israël consentiéntes accessérunt ad eos, sed Mathathías et fílii ejus constánter stetérunt.

\R\ Aperi óculos tuos, Dómine, et vide afflictiónem nostram: circumdedérunt nos gentes ad puniéndum nos: * Sed tu, Dómine, exténde bráchium tuum, et líbera ánimas nostras. \V\ Afflíge oppriméntes nos et contuméliam faciéntes in supérbiam; et custódi partem tuam. — Sed. Glória Patri. Sed.

\feriaheading{Feria Quarta infra Hebdomadam I Octobris}

\rubric{De libro primo Machabæórum}

\lectioheading{Lectio i}{Cap. 2, 19-30}

\lettrine[lines=2]{E}{T} respóndit Mathathías et dixit magna voce: Etsi omnes gentes regi Antíocho obédiunt, ut discédat unusquísque a servitúte legis patrum suórum et conséntiat mandátis ejus, ego et fílii mei et fratres mei obediémus legi patrum nostrórum. Propítius sit nobis Deus, non est nobis útile relínquere legem et justítias Dei: non audiémus verba regis Antíochi nec sacrificábimus transgrediéntes legis nostræ mandáta, ut eámus áltera via.

\R\ Refúlsit sol in clýpeos áureos, et resplenduérunt montes ab eis: * Et fortitúdo géntium dissipáta est. \V\ Erat enim exércitus magnus valde et fortis: et appropiávit Judas et exércitus ejus in prǽlio. — Et.

\lectioheading{Lectio ii}{}

\lettrine[lines=2]{E}{T}, ut cessávit loqui verba hæc, accéssit quidam Judǽus in ómnium óculis sacrificáre idólis super aram in civitáte Modin secúndum jussum regis. Et vidit Mathathías et dóluit, et contremuérunt renes ejus, et accénsus est furor ejus secúndum judícium legis, et insíliens trucidávit eum super aram. Sed et virum quem rex Antíochus míserat, qui cogébat immoláre, occídit in ipso témpore et aram destrúxit et zelátus est legem, sicut fecit Phínees Zamri fílio Salómi.

\R\ Ornavérunt fáciem templi corónis áureis, et dedicavérunt altáre Dómino: * Et facta est lætítia magna in pópulo. \V\ In hymnis et confessiónibus benedicébant Dóminum. — Et.

\rubric{In Octávis et Festis ritus simplícis: Glória Patri. Et.}

\lectioheading{Lectio iii}{}

\lettrine[lines=2]{E}{T} exclamávit Mathathías voce magna in civitáte dicens: Omnis qui zelum habet legis státuens testaméntum éxeat post me. Et fugit ipse et fílii ejus in montes, et reliquérunt quæcúmque habébant in civitáte. Tunc descendérunt multi quæréntes judícium et justítiam in desértum. Et sedérunt ibi ipsi et fílii eórum et mulíeres eórum et pécora eórum, quóniam inundavérunt super eos mala.

\R\ In hymnis et confessiónibus benedicébant Dóminum, * Qui magna fecit in Israël, et victóriam dedit illis Dóminus omnípotens. \V\ Ornavérunt fáciem templi corónis áureis, et dedicavérunt altáre Dómino. — Qui. Glória Patri. Qui.

\feriaheading{Feria Quinta infra Hebdomadam I Octobris}

\rubric{De libro primo Machabæórum}

\lectioheading{Lectio i}{Cap. 2, 49-69}

\lettrine[lines=2]{E}{T} appropinquavérunt dies Mathathíæ moriéndi, et dixit fíliis suis: Nunc confortáta est supérbia et castigátio et tempus eversiónis et ira indignatiónis. Nunc ergo, o fílii, æmulatóres estóte legis et date ánimas vestras pro testaménto patrum vestrórum et mementóte óperum patrum, quæ fecérunt in generatiónibus suis, et accipiétis glóriam magnam et nomen ætérnum. Abraham, nonne in tentatióne invéntus est fidélis, et reputátum est ei ad justítiam? Joseph in témpore angústiæ suæ custodívit mandátum, et factus est dóminus Ægýpti; Phínees pater noster zelándo zelum Dei accépit testaméntum sacerdótii ætérni.

\R\ Adapériat Dóminus cor vestrum in lege sua et in præcéptis suis et fáciat pacem in diébus vestris: * Concédat vobis salútem, et rédimat vos a malis. \V\ Exáudiat Dóminus oratiónes vestras et reconciliétur vobis nec vos déserat in témpore malo. — Concédat.

\lectioheading{Lectio ii}{}

\lettrine[lines=2]{J}{ESUS}, dum implévit verbum, factus est dux in Israël; Caleb, dum testificátur in ecclésia, accépit hereditátem; David in sua misericórdia consecútus est sedem regni in sǽcula; Elías, dum zelat zelum legis, recéptus est in cælum; Ananías et Azarías et Mísael credéntes liberáti sunt de flamma; Dániel in sua simplicitáte liberátus est de ore leónum; et ita cogitáte per generatiónem et generatiónem, quia omnes qui sperant in eum non infirmántur. Et a verbis viri peccatóris ne timuéritis, quia glória ejus stercus et vermis est: hódie extóllitur et cras non inveniétur, quia convérsus est in terram suam, et cogitátio ejus périit.

\R\ Exáudiat Dóminus oratiónes vestras, et reconciliétur vobis nec vos déserat in témpore malo * Dóminus, Deus noster. \V\ Det vobis cor ómnibus, ut colátis eum et faciátis ejus voluntátem. — Dóminus.

\rubric{In Octávis et Festis ritus simplícis: Glória Patri. Dóminus.}

\lectioheading{Lectio iii}{}

\lettrine[lines=2]{V}{OS} ergo, fílii, confortámini et viríliter ágite in lege, quia in ipsa gloriósi éritis. Et ecce Simon frater vester, scio quod vir consílii est: ipsum audíte semper, et ipse erit vobis pater; et Judas Machabǽus fortis víribus a juventúte sua sit vobis princeps militíæ et ipse aget bellum pópuli, et adducétis ad vos omnes factóres legis et vindicáte vindíctam pópuli vestri, retribúite retributiónem géntibus et inténdite in præcéptum legis. Et benedíxit eos et appósitus est ad patres suos.

\R\ Congregáti sunt inimíci nostri, et gloriántur in virtúte sua: cóntere fortitúdinem illórum, Dómine, et dispérge illos: * Ut cognóscant quia non est álius qui pugnet pro nobis, nisi tu, Deus noster. \V\ Dispérge illos in virtúte tua, et déstrue eos, protéctor noster, Dómine. — Ut. Glória Patri. Ut.

\feriaheading{Feria Sexta infra Hebdomadam I Octobris}

\rubric{De libro primo Machabæórum}

\lectioheading{Lectio i}{Cap. 2, 70; 3, 1-3, 5-12, 25-28}

\lettrine[lines=2]{D}{EFUNCTUS} est Mathathías anno centésimo et quadragésimo sexto et sepúltus est a fíliis suis in sepúlcris patrum suórum in Modin, et planxérunt eum omnis Israël planctu magno. Et surréxit Judas, qui vocabátur Machabǽus, fílius ejus pro eo et adjuvábant eum omnes fratres ejus et univérsi qui se conjúnxerant patri ejus et præliabántur prǽlium Israël cum lætítia. Et dilatávit glóriam pópulo suo, et persecútus est iníquos perscrútans eos et, qui conturbábant pópulum suum, eos succéndit flammis, et repúlsi sunt inimíci ejus præ timóre ejus, et omnes operárii iniquitátis conturbáti sunt, et dirécta est salus in manu ejus.

\R\ Impetum inimicórum ne timuéritis: mémores estóte, quómodo salvi facti sunt patres nostri: * Et nunc clamémus in cælum, et miserébitur nostri Deus noster. \V\ Mementóte mirabílium ejus, quæ fecit Pharaóni et exercítui ejus in Mari Rubro. — Et.

\lectioheading{Lectio ii}{}

\lettrine[lines=2]{E}{T} exacerbábat reges multos et lætificábat Jacob in opéribus suis, et in sǽculum memória ejus in benedictióne, et perambulávit civitátes Juda et pérdidit ímpios ex eis et avértit iram ab Israël et nominátus est usque ad novíssimum terræ et congregávit pereúntes. Et congregávit Apollónius gentes et a Samaría virtútem multam et magnam ad bellándum contra Israël. Et cognóvit Judas et éxiit óbviam illi et percússit et occídit illum, et cecidérunt vulneráti multi, et réliqui fugérunt. Et accépit spólia eórum et gládium Apollónii ábstulit Judas et erat pugnans in eo ómnibus diébus.

\R\ Congregátæ sunt gentes in multitúdine, ut dímicent contra nos, et ignorámus quid ágere debeámus: * Dómine Deus, ad te sunt óculi nostri, ne pereámus. \V\ Tu scis quæ cógitant in nos: quómodo potérimus subsístere ante fáciem illórum, nisi tu ádjuves nos? — Dómine.

\rubric{In Octávis et Festis ritus simplícis: Glória Patri. Dómine.}

\lectioheading{Lectio iii}{}

\lettrine[lines=2]{E}{T} cécidit timor Judæ ac fratrum ejus et formído super omnes gentes in circúitu eórum; et pervénit ad regem nomen ejus, et de prǽliis Judæ narrábant omnes gentes. Ut audívit autem rex Antíochus sermónes istos, irátus est ánimo et misit et congregávit exércitum univérsi regni sui, castra fórtia valde, et apéruit ærárium suum et dedit stipéndia exercítui in annum et mandávit illis, ut essent paráti ad ómnia.

\R\ Tua est poténtia, tuum regnum, Dómine: tu es super omnes gentes: * Da pacem, Dómine, in diébus nostris. \V\ Creátor ómnium, Deus, terríbilis et fortis, justus et miséricors. — Da. Glória Patri. Da.

\feriaheading{Sabbato infra Hebdomadam I Octobris}

\rubric{De libro primo Machabæórum}

\lectioheading{Lectio i}{Cap. 3, 42-60}

\lettrine[lines=2]{E}{T} vidit Judas et fratres ejus quia multiplicáta sunt mala et exércitus applicábant ad fines eórum, et cognovérunt verba regis, quæ mandávit pópulo fácere in intéritum et consummatiónem, et dixérunt unusquísque ad próximum suum: Erigámus dejectiónem pópuli nostri et pugnémus pro pópulo nostro et sanctis nostris. Et congregátus est convéntus, ut essent paráti in prǽlium et ut orárent et péterent misericórdiam et miseratiónes. Et Jerúsalem non habitabátur, sed erat sicut desértum: non erat qui ingrederétur et egrederétur de natis ejus, et sanctum conculcabátur, et fílii alienigenárum erant in arce: ibi erat habitátio géntium, et abláta est volúptas a Jacob, et defécit ibi tíbia et cíthara.

\R\ Refúlsit sol in clýpeos áureos, et resplenduérunt montes ab eis: * Et fortitúdo géntium dissipáta est. \V\ Erat enim exércitus magnus valde et fortis: et appropiávit Judas et exércitus ejus in prǽlio. — Et.

\lectioheading{Lectio ii}{}

\lettrine[lines=2]{E}{T} congregáti sunt et venérunt in Maspha contra Jerúsalem, quia locus oratiónis erat in Maspha ante in Israël, et jejunavérunt illa die et induérunt se cilíciis et cínerem imposuérunt cápiti suo et discidérunt vestiménta sua et expandérunt libros legis, de quibus scrutabántur gentes similitúdinem simulacrórum suórum, et attulérunt ornaménta sacerdotália et primítias et décimas et suscitavérunt Nazarǽos, qui impléverant dies, et clamavérunt voce magna in cælum dicéntes: Quid faciémus istis et quo eos ducémus? Et sancta tua conculcáta sunt et contamináta sunt, et sacerdótes tui facti sunt in luctum et in humilitátem, et ecce natiónes convenérunt advérsum nos, ut nos dispérdant: tu scis quæ cógitant in nos. Quómodo potérimus subsístere ante fáciem eórum, nisi tu, Deus, ádjuves nos?

\R\ Ornavérunt fáciem templi corónis áureis, et dedicavérunt altáre Dómino: * Et facta est lætítia magna in pópulo. \V\ In hymnis et confessiónibus benedicébant Dóminum. — Et.

\rubric{In Octávis simplícibus et in Offício sanctæ Maríæ in Sábbato: Glória Patri. Et.}

\lectioheading{Lectio iii}{}

\lettrine[lines=2]{E}{T} tubis exclamavérunt voce magna. Et post hæc constítuit Judas duces pópuli, tribúnos et centuriónes et pentacontárchos et decuriónes, et dixit his, qui ædificábant domos et sponsábant uxóres et plantábant víneas et formidolósis, ut redírent unusquísque in domum suam secúndum legem. Et movérunt castra et collocavérunt ad Austrum Emmaum. Et ait Judas: Accingímini et estóte fílii poténtes et estóte paráti in mane, ut pugnétis advérsus natiónes has, quæ convenérunt advérsus nos dispérdere nos et sancta nostra; quóniam mélius est nos mori in bello, quam vidére mala gentis nostræ et sanctórum. Sicut autem fúerit volúntas in cælo, sic fiat.

\R\ In hymnis et confessiónibus benedicébant Dóminum, * Qui magna fecit in Israël, et victóriam dedit illis Dóminus omnípotens. \V\ Ornavérunt fáciem templi corónis áureis, et dedicavérunt altáre Dómino. — Qui. Glória Patri. Qui.

\hourheading{Ad Vesperas}

\V\ Vespertína orátio ascéndat ad te, Dómine. \R\ Et descéndat super nos misericórdia tua.

\antlabel{Ad Magnif. Ant.} Refúlsit sol * in clýpeos áureos, et resplenduérunt montes ab eis: et fortitúdo géntium dissipáta est.

\rubric{Oratio Dominicæ occurréntis.}

% ===== BATCH 8 of the Proprium de Tempore (autumnalis.md lines
% 14884-15116) -- the complete first week of October: "Dominica I
% Octobris" Matins opening the First Book of Machabees in Nocturn I
% and St. Ambrose's "De Officiis" in Nocturn II, the six ferial days
% continuing the Machabees cycle through Mathathias's death and
% Judas's early victories, and the closing Saturday Vespers. Both
% Nocturn III's responsories (as at the two immediately preceding
% Sundays) are given without the "vii"/"viii" numbering used at
% Dominicas I-IV Septembris. A new governing rubric fixes the hymn
% "Ætérne rerum" at Lauds for every Sunday through November 5. One
% genuine partial-caps OCR artifact is silently normalized
% (DEfúnctus to DEFUNCTUS, accent dropped on its newly-capitalized
% remainder), while complete short words already fully capped in
% the source (VOS, JESUS) are left untouched. Feria Sexta's citation
% ("Cap. 2, 70; 3, 1-3, 5-12, 25-28") appears in the source one line
% above its "Lectio i" header rather than below it as usual; it is
% nonetheless given as that Lesson's citation via
% \lectioheading{Lectio i}{...}, consistent with its evident
% function regardless of its exact placement in the source. The
% pervasive space-before-colon formatting quirk from the previous
% two batches continues here and is likewise silently normalized.
% The document continues with "Dominica II Octobris" (autumnalis.md
% line 15117).
% =====
       % or hiemalis.tex / verna.tex / aestiva.tex
%   \end{document}
%
% Since autumnalis.tex, aestiva.tex, hiemalis.tex, and verna.tex all
% now embed the versicle/response marks (℣ U+2123, ℟ U+211F) and the
% blessing cross (☩ U+2629) as literal Unicode characters rather than
% as macros, and the main body font (DejaVu Serif) does not contain
% glyphs for them, \newunicodechar is used below to substitute a font
% that does (DejaVu Sans) automatically wherever those three
% characters occur, without having to touch the body text itself.
% =====================================================================

\usepackage{fontspec}
\usepackage{textcomp}
\usepackage{lettrine}
\usepackage[dvipsnames]{xcolor}
\usepackage[margin=2cm]{geometry}
\usepackage{parskip}
\usepackage{newunicodechar}

\setmainfont{DejaVu Serif}
\newfontfamily\liturgicalsymbolfont{DejaVu Sans}

\newunicodechar{℣}{{\liturgicalsymbolfont\char"2123}}  % VERSICLE
\newunicodechar{℟}{{\liturgicalsymbolfont\char"211F}}  % RESPONSE
\newunicodechar{☩}{{\liturgicalsymbolfont\char"2629}}  % cross (blessing)

% --- Headings (one required argument each) ---
\newcommand{\partheading}[1]{{\bigskip\noindent\LARGE\bfseries #1\par\medskip}}
\newcommand{\partsubtitle}[1]{{\noindent\large\itshape #1\par\medskip}}
\newcommand{\sectionheading}[1]{{\bigskip\noindent\Large\bfseries #1\par\medskip}}
\newcommand{\feriaheading}[1]{{\bigskip\noindent\large\bfseries #1\par\smallskip}}
\newcommand{\hourheading}[1]{{\bigskip\noindent\large\bfseries #1\par\smallskip}}
\newcommand{\nocturnoheading}[1]{{\medskip\noindent\bfseries #1\par\smallskip}}
\newcommand{\psalmheading}[1]{{\medskip\noindent\itshape #1\par\smallskip}}
\newcommand{\hymnusheading}[1]{{\medskip\noindent\bfseries #1\par\smallskip}}
\newcommand{\oratioheading}[1]{{\medskip\noindent\itshape #1\par\smallskip}}
\newcommand{\variantheading}[1]{{\medskip\noindent\itshape #1\par\smallskip}}
\newcommand{\absolutioheading}[1]{{\medskip\noindent\itshape #1\par\smallskip}}
\newcommand{\suffragiumheading}[1]{{\medskip\noindent\bfseries #1\par\smallskip}}
\newcommand{\symbolumheading}[1]{{\medskip\noindent\bfseries #1\par\smallskip}}
\newcommand{\antiphonaheading}[1]{{\medskip\noindent\bfseries #1\par\smallskip}}
\newcommand{\feastheading}[1]{{\bigskip\noindent\large\bfseries #1\par\smallskip}}
\newcommand{\sanctusheading}[1]{{\bigskip\noindent\large\bfseries #1\par\smallskip}}
\newcommand{\diesheading}[1]{{\medskip\noindent\bfseries #1\par\smallskip}}
\newcommand{\monthheading}[1]{{\bigskip\noindent\Large\bfseries #1\par\medskip}}

% --- Purely decorative divider (no argument) ---
\newcommand{\sectionornament}{{\par\noindent\centering\rule{0.3\linewidth}{0.4pt}\par\medskip}}

% --- Headings (two required arguments: title + citation) ---
\newcommand{\lectioheading}[2]{{\medskip\noindent\itshape #1\quad\textit{#2}\par\smallskip}}
\newcommand{\capitulumlabel}[2]{{\medskip\noindent\itshape #1\quad\textit{#2}\par\smallskip}}
\newcommand{\canticumheading}[2]{{\medskip\noindent\itshape #1\quad\textit{#2}\par\smallskip}}
\newcommand{\lectiolabel}[2]{{\medskip\noindent\itshape #1\quad\textit{#2}\par\smallskip}}

% --- Labels (one required argument) ---
\newcommand{\antlabel}[1]{\textbf{#1}}
\newcommand{\invitatlabel}[1]{\textbf{#1}}
\newcommand{\psalmlabel}[1]{\textbf{#1}}
\newcommand{\benedlabel}[1]{\textbf{#1}}
\newcommand{\benedictiolabel}[1]{\textbf{#1}}
\newcommand{\absolutiolabel}[1]{\textbf{#1}}
\newcommand{\formulalabel}[1]{\textbf{#1}}
\newcommand{\oratiolabel}[1]{\textbf{#1}}

% --- Rubric (red text; embedded rubrics may nest \textit, etc.) ---
\newcommand{\rubric}[1]{\textcolor{red}{#1}}
% Black text embedded within a red \rubric{} span (governing instruction:
% "embedded rubrics (black text within red-colored rubrics) are preserved").
\newcommand{\rubricblack}[1]{\textcolor{black}{#1}}

% \P (pilcrow), \dag, \ddag, \label, \pageref, \hspace, \textbf,
% \textsuperscript, \textit, \textsc, \textdagger, \textdaggerdbl are all
% standard LaTeX/textcomp commands and need no definition here.
  (or copy its contents into a wrapper)
%   \begin{document}
%   % =====================================================================
% autumnalis.tex -- LaTeX conversion of autumnalis.md (Breviarium
% Romanum, Pars Autumnalis), BODY CONTENT ONLY.
%
% No preamble, \documentclass, or \newcommand definitions are included
% here, per instruction -- the command vocabulary below follows the
% same legend established for hiemalis.tex, verna.tex, and aestiva.tex
% (one-column vs. two-column, red vs. black, etc.).
%
% BATCH 1 of N: Ordinarium Divini Officii, Ad Matutinum and Ad Laudes
% (autumnalis.md lines 1-463)
%
% NOTE: autumnalis.md lines 1-1141 (the entire Ordinarium Divini
% Officii) were verified via direct diff to be BYTE-IDENTICAL to the
% corresponding lines of aestiva.md -- both being non-Lenten/
% non-Paschal volumes, each omits the same "Tempore Quadragesimæ" /
% "Tempore Passionis" Invitatorium-Hymnus variants and the Paschaltide
% "Commemoratio de Cruce" that verna.tex's Ordinarium carries. Given
% this confirmed byte-for-byte identity (not a parallel or merely
% similar occurrence, but the literal same source text), the verified
% LaTeX already produced for aestiva.tex's Batches 1-3 is reused
% directly below rather than re-keyed by hand, to avoid reintroducing
% errors (such as the accent-dropping mistakes) already caught and
% fixed in that earlier work. Every batch from "Psalterium Breviarii
% Romani" onward is fresh transcription from autumnalis.md's own text,
% as usual.
% =====================================================================

% [ornamental image beginning Pars Autumnalis]
\partheading{Ordinarium Divini Officii}
\partsubtitle{Juxta Ritum Romanum Persolvendi}

\rubric{Clericis in sacris constitútis, qui íntegrum divínum Offícium, quamvis in partes distribútum, coram Ssmo Sacraménto, sive públicæ adoratióni expósito sive in tabernáculo adserváto, devóte recitáverint, concéditur: Indulgéntia plenária, si prætérea sacramentálem confessiónem instítuerint, ad eucharísticum Convívium accésserint et ad mentem Summi Pontíficis oráverint. \textit{Pius Pp. XI, 23 Octóbris 1930.}}

\rubric{Iis vero, qui divínum Offícium coram Ssmo Sacraménto, ut supra, ex parte tantum recitáverint, concéditur: Indulgéntia quingentórum diérum pro unaquáque hora canónica. \textit{Pius Pp. XI, 18 Maji 1933.}}

\rubric{Iisdem autem cléricis, qui divínum Offícium in álias preces commutátum rite obtinúerint, si preces ejúsmodi coram Ssmo Sacraménto, uti supra, devóte recitáverint, concéditur: Indulgéntia plenária, addítis sacramentáli confessióne, sacra Communióne et oratióne ad mentem Summi Pontíficis. \textit{Pius Pp. XI, 7 Novémbris 1932.}}

\sectionheading{Ante Divinum Officium}

\rubric{Antequam inchoëtur Offícium laudabíliter dícitur, sub singulári semper número, séquens Orátio; pro qua Summus Póntifex Pius Pp. XI sub die 17 Novémbris 1933 Indulgéntiam trium annórum concéssit.}

\oratioheading{Oratio}

\lettrine[lines=2]{A}{PERI}, Dómine, os meum ad benedicéndum nomen sanctum tuum: munda quoque cor meum ab ómnibus vanis, pervérsis et aliénis cogitatiónibus; intelléctum illúmina, afféctum inflámma, ut digne, atténte ac devóte hoc Offícium recitáre váleam, et exaudíri mérear ante conspéctum divínæ Majestátis tuæ. Per Christum, Dóminum nostrum. \R\ Amen.

Dómine, in unióne illíus divínæ intentiónis, qua ipse in terris laudes Deo persolvísti, has tibi Horas (\textit{vel} hanc tibi Horam) persólvo.

\sectionheading{Ante singulas Horas}

\rubric{Ante Matutínum, sicut et ante Laudes a Matutíno extra Chorum separátas, et ante omnes Horas, præter quam ad Completórium, dícitur secréto:}

\lettrine[lines=2]{P}{ATER} noster, qui es in cælis, sanctificétur nomen tuum. Advéniat regnum tuum. Fiat volúntas tua, sicut in cælo et in terra. Panem nostrum quotidiánum da nobis hódie. Et dimítte nobis débita nostra, sicut et nos dimíttimus debitóribus nostris. Et ne nos indúcas in tentatiónem: sed líbera nos a malo. Amen.

\lettrine[lines=2]{A}{VE}, María, grátia plena; Dóminus tecum: benedícta tu in muliéribus, et benedíctus fructus ventris tui Jesus. Sancta María, Mater Dei, ora pro nobis peccatóribus, nunc et in hora mortis nostræ. Amen.

\rubric{In princípio autem Matutíni ac Primæ additur secréto Symbolum Apostolórum:}

\lettrine[lines=2]{C}{REDO} in Deum, Patrem omnipoténtem, Creatórem cæli et terræ. Et in Jesum Christum Fílium ejus únicum, Dóminum nostrum: qui concéptus est de Spíritu Sancto, natus ex María Vírgine, passus sub Póntio Piláto, crucifíxus, mórtuus, et sepúltus: descéndit ad ínferos; tértia die resurréxit a mórtuis; ascéndit ad cælos; sedet ad déxteram Dei Patris omnipoténtis: inde ventúrus est judicáre vivos et mórtuos. Credo in Spíritum Sanctum, sanctam Ecclésiam cathólicam, Sanctórum communiónem, remissiónem peccatórum, carnis resurrectiónem, vitam ætérnam. Amen.

\rubric{Officium ad Matutinum et ad reliquas Horas, nisi aliter suo loco indicetur, per totum annum persolvitur juxta formam, quæ in Rubricis hujus Ordinarii exhibetur.}

\hourheading{Ad Matutinum}

Pater noster, Ave María et Credo.

\rubric{Deinde, clara voce, dicitur Versus:}

\V\ \lettrine[lines=2]{D}{OMINE}, lábia mea apéries. \R\ Et os meum annuntiábit laudem tuam. \V\ Deus, in adjutórium meum inténde. \R\ Dómine, ad adjuvándum me festína.

Glória Patri, et Fílio, et Spirítui Sancto. Sicut erat in princípio, et nunc, et semper, et in sæcula sæculórum. Amen. Allelúja.

\rubric{Postea dicitur conveniens Invitatorium, quod ante Psalmum bis integrum recitatur, et ad singulos ejusdem Psalmi versus vel integrum vel dimidiatum ab Asterisco * alternis vicibus, ut infra, repetitur.}

\psalmheading{Psalmus 94}

\lettrine[lines=2]{V}{ENITE}, exsultémus Dómino, jubilémus Deo, salutári nostro: præoccupémus fáciem ejus in confessióne, et in psalmis jubilémus ei.

\rubric{Repetitur integrum Invitatorium.}

Quóniam Deus magnus Dóminus, et Rex magnus super omnes deos: quóniam non repéllet Dóminus plebem suam: quia in manu ejus sunt omnes fines terræ, et altitúdines móntium ipse cónspicit.

\rubric{Repetitur altera pars Invitatorii.}

\rubric{In sequenti Psalmi versu, ad verba veníte, adorémus, et procidámus ante Deum, genuflectitur.}

Quóniam ipsíus est mare, et ipse fecit illud, et áridam fundavérunt manus ejus: veníte, adorémus, et procidámus ante Deum: plorémus coram Dómino, qui fecit nos, quia ipse est Dóminus, Deus noster; nos autem pópulus ejus, et oves páscuæ ejus.

\rubric{Repetitur integrum Invitatorium.}

Hódie, si vocem ejus audíeritis, nolíte obduráre corda vestra, sicut in exacerbatióne secúndum diem tentatiónis in desérto: ubi tentavérunt me patres vestri, probavérunt et vidérunt ópera mea.

\rubric{Repetitur altera pars Invitatorii.}

Quadragínta annis próximus fui generatióni huic, et dixi: Semper hi errant corde; ipsi vero non cognovérunt vias meas: quibus jurávi in ira mea: Si introíbunt in réquiem meam.

\rubric{Repetitur integrum Invitatorium.}

Glória Patri, et Fílio, et Spirítui Sancto. Sicut erat in princípio, et nunc, et semper, et in sæcula sæculórum. Amen.

\rubric{Repetitur altera pars Invitatorii.}

\rubric{Denique repetitur integrum Invitatorium.}

\rubric{Expleto Psalmo, dicitur Hymnus Invitatorio respondens.}

\variantheading{In Officio dominicali}

\rubric{In omnibus Dominicis, extra Octavas Ssmi Corporis Christi ac Sacrat. Cordis Jesu, Invitatorium et Hymnus ut in Psalterio.}

\variantheading{In Officio feriali}

\rubric{In omnibus Feriis, Invitatorium et Hymnus de Feria currenti, ut in Psalterio.}

\variantheading{In Festis}

\rubric{In Officio cujuslibet Festi vel Octavæ, ac S. Mariæ in Sabbato, Invitatorium et Hymnus ut in Proprio vel Communi; in Dominicis autem infra Octavas Ssmi Corporis Christi ac Sacratissimi Cordis Jesu, si de eis fiat Officium, Invitatorium et Hymnus ut in Proprio de Tempore. Quæ norma ad reliquas item Horas erit observanda pro illis Officii partibus, quæ juxta Rubricas ex Ordinario vel Psalterio non fuerint in Festis desumendæ.}

\variantheading{In Officio novem Lectionum}

\rubric{Expleto Hymno, dicuntur Antiphonæ convenientes, quæ in Officiis ritus duplicis ante et post Psalmos integræ recitantur; in Officiis autem ritus semiduplicis initio Psalmi inchoantur tantum et usque ad Asteriscum * perducuntur, atque in fine integræ pronuntiantur; servata alterutra norma etiam in omnibus Antiphonis, quæ ad Laudes ac Vesperas adhibentur.}

\nocturnoheading{In I Nocturno}

\rubric{Sub congruentibus Antiphonis dicuntur tres Psalmi, ac deinde subjungitur Versus, prouti Officium occurrens requirit.}

\rubric{Post Versum cujuslibet Nocturni dicitur}

Pater noster \rubric{secreto usque ad} \V\ Et ne nos indúcas in tentatiónem. \R\ Sed líbera nos a malo.

\absolutioheading{Absolutio}

\lettrine[lines=2]{E}{XÁUDI}, Dómine Jesu Christe, preces servórum tuórum, et miserére nobis: Qui cum Patre et Spíritu Sancto vivis et regnas in sæcula sæculórum. \R\ Amen.

\V\ Jube, domne, benedícere.

\rubric{¶ Extra Chorum, quando ab uno tantum recitatur Officium, ante singulas Lectiones Matutini atque ad Lectionem brevem Primæ et Completorii, dicitur: Jube, Dómine, benedícere; et subjungitur congruens Benedictio. Ab Episcopo autem, ultimam Matutini Lectionem cantaturo, item dicitur: Jube, Dómine, benedícere; et respondetur a Choro: Amen.}

\variantheading{Pro I Lectione}

\benedlabel{Bened.} Benedictióne perpétua benedícat nos Pater ætérnus. \R\ Amen.

\rubric{Deinde dicuntur in unoquoque Nocturno Lectiones, prouti Officium occurrens requirit, et in fine cujuslibet Lectionis additur:}

\V\ Tu autem, Dómine, miserére nobis. \R\ Deo grátias.

\rubric{Post quamlibet vero Lectionem, quæ Hymnum Te Deum immediate non præcedat, congruens dicitur Responsorium, et in fine ultimi Responsorii cujusque Nocturni additur Versus: Glória Patri, et Fílio, et Spirítui Sancto, et Responsorium ab Asterisco *, et quidem a secundo, si duo fuerint, repetitur.}

\variantheading{Pro II Lectione}

\benedlabel{Bened.} Unigénitus Dei Fílius nos benedícere et adjuváre dignétur. \R\ Amen.

\variantheading{Pro III Lectione}

\benedlabel{Bened.} Spíritus Sancti grátia illúminet sensus et corda nostra. \R\ Amen.

\nocturnoheading{In II Nocturno}

\rubric{Sub congruentibus item Antiphonis dicuntur tres Psalmi et Versus, sicut in I Nocturno.}

\rubric{Post Versum dicitur.}

Pater noster \rubric{secreto usque ad} \V\ Et ne nos indúcas in tentatiónem. \R\ Sed líbera nos a malo.

\absolutioheading{Absolutio}

\lettrine[lines=2]{I}{PSIUS} píetas et misericórdia nos ádjuvet, qui cum Patre et Spíritu Sancto vivit et regnat in sæcula sæculórum. \R\ Amen.

\V\ Jube, domne, benedícere.

\variantheading{Pro IV Lectione}

\benedlabel{Bened.} Deus Pater omnípotens sit nobis propítius et clemens. \R\ Amen.

\variantheading{Pro V Lectione}

\benedlabel{Bened.} Christus perpétuæ det nobis gáudia vitæ. \R\ Amen.

\variantheading{Pro VI Lectione}

\benedlabel{Bened.} Ignem sui amóris accéndat Deus in córdibus nostris. \R\ Amen.

\nocturnoheading{In III Nocturno}

\rubric{Sub congruentibus denique Antiphonis dicuntur tres Psalmi et Versus, sicut in I et II Nocturno.}

\rubric{Post Versum dicitur}

Pater noster \rubric{secreto usque ad} \V\ Et ne nos indúcas in tentatiónem. \R\ Sed líbera nos a malo.

\absolutioheading{Absolutio}

\lettrine[lines=2]{A}{} \textsc{vinculis} peccatórum nostrórum absólvat nos omnípotens et miséricors Dóminus. \R\ Amen.

\V\ Jube, domne, benedícere.

\variantheading{Pro VII Lectione}

\benedlabel{Bened.} Evangélica léctio sit nobis salus et protéctio. \R\ Amen.

\variantheading{Pro VIII Lectione}

\rubric{In Festis Domini et in Dominicis:}

\benedlabel{Bened.} Divínum auxílium máneat semper nobíscum. \R\ Amen.

\rubric{In Festis beatæ Mariæ Virginis:}

\benedlabel{Bened.} Cujus festum cólimus, ipsa Virgo virgínum intercédat pro nobis ad Dóminum. \R\ Amen.

\rubric{In Festis Sanctorum:}

\benedlabel{Bened.} Cujus (\textit{vel} Quorum \textit{aut} Quarum) festum cólimus, ipse (\textit{vel} ipsa \textit{aut} ipsi \textit{vel} ipsæ) intercédat (\textit{vel} intercédant) pro nobis ad Dóminum. \R\ Amen.

\variantheading{Pro IX Lectione}

\benedlabel{Bened.} Ad societátem cívium supernórum perdúcat nos Rex Angelórum. \R\ Amen.

\rubric{Si autem legenda sit Lectio ix de Homilia cum Evangelio Dominicæ, vel Feriæ, aut Vigiliæ:}

\benedlabel{Bened.} Per evangélica dicta deleántur nostra delícta. \R\ Amen.

\variantheading{In Officio trium Lectionum}

\nocturnoheading{In Nocturno}

\rubric{In Officiis trium Lectionum, post Hymnum dicuntur Antiphonæ convenientes, quæ initio Psalmi inchoantur tantum et usque ad Asteriscum * perducuntur, ac deinde in fine integræ pronuntiantur.}

\rubric{Sub eisdem vero Antiphonis dicuntur novem Psalmi Feriæ currentis, quibus subjungitur Versus in III Nocturno positus, omissis Versibus pro I et II Nocturno assignatis.}

\rubric{Post Versum dicitur}

Pater noster \rubric{secreto usque ad} \V\ Et ne nos indúcas in tentatiónem. \R\ Sed líbera nos a malo.

\absolutioheading{Absolutio}

\rubric{Feria II et V}

Exáudi, Dómine, \rubric{ut in I Nocturno.}

\rubric{Feria III et VI}

Ipsíus píetas, \rubric{ut in II Nocturno.}

\rubric{Feria IV et Sabbato}

A vínculis, \rubric{ut in III Nocturno Officii novem Lectionum.}

\rubric{Deinde leguntur Lectiones cum Responsoriis, prouti Officium occurrens requirit; et ante eas dicuntur sequentes}

\variantheading{Benedictiones}

\rubric{In Feriis, quando legitur Homilia cum Evangelio, erit:}

\benedlabel{1 Bened.} Evangélica léctio sit nobis salus et protéctio. \R\ Amen.

\benedlabel{2 Bened.} Divínum auxílium máneat semper nobíscum. \R\ Amen.

\benedlabel{3 Bened.} Ad societátem cívium supernórum perdúcat nos Rex Angelórum. \R\ Amen.

\rubric{In Feriis, quando non legitur Homilia cum Evangelio, Feria II et V Benedictiones ut in I Nocturno, Feria III et VI ut in II Nocturno Officii novem Lectionum; Feria autem IV et Sabbato erit:}

\benedlabel{1 Bened.} Ille nos benedícat, qui sine fine vivit et regnat. \R\ Amen.

\benedlabel{2 Bened.} Divínum auxílium máneat semper nobíscum. \R\ Amen.

\benedlabel{3 Bened.} Ad societátem cívium supernórum perdúcat nos Rex Angelórum. \R\ Amen.

\rubric{In Festis Sanctorum erit:}

\benedlabel{1 Bened.} Ille nos benedícat, qui sine fine vivit et regnat. \R\ Amen.

\benedlabel{2 Bened.} Cujus (\textit{vel} Quorum \textit{aut} Quarum) festum cólimus, ipse (\textit{vel} ipsa \textit{aut} ipsi \textit{vel} ipsæ) intercédat (\textit{vel} intercédant) pro nobis ad Dóminum. \R\ Amen.

\benedlabel{3 Bened.} Ad societátem cívium supernórum perdúcat nos Rex Angelórum. \R\ Amen.

\rubric{In Officio de sancta María in Sábbato Absolútio et Benedictiónes, quæ ibídem ponúntur.}

\rubric{Post últimam Lectiónem, in ómnibus Domínicis per Annum minóribus, in Festis cujúsvis ritus, per omnes Octávas, et in Offício sanctæ Maríæ in Sábbato, dícitur Hymnus Ambrosiánus. Et numquam dícitur in feriáli Offício. Quando vero Hymnus prædíctus omíttitur, ejus loco dícitur IX aut III Responsórium.}

\hymnusheading{Hymnus Ambrosianus}

\lettrine[lines=2]{T}{E} Deum laudámus: * te Dóminum confitémur.

Te ætérnum Patrem * omnis terra venerátur.

Tibi omnes Angeli, * tibi Cæli, et univérsæ Potestátes:

Tibi Chérubim et Séraphim * incessábili voce proclámant:

Sanctus, * Sanctus, * Sanctus * Dóminus, Deus Sábaoth.

Pleni sunt cæli et terra * majestátis glóriæ tuæ.

Te gloriósus * Apostolórum chorus,

Te Prophetárum * laudábilis númerus,

Te Mártyrum candidátus * laudat exércitus.

Te per orbem terrárum * sancta confitétur Ecclésia,

Patrem * imménsæ majestátis;

Venerándum tuum verum * et únicum Fílium;

Sanctum quoque * Paráclitum Spíritum.

Tu Rex glóriæ, * Christe.

Tu Patris * sempitérnus es Fílius.

Tu, ad liberándum susceptúrus hóminem, * non horruísti Vírginis úterum.

Tu, devícto mortis acúleo, * aperuísti credéntibus regna cælórum.

Tu ad déxteram Dei sedes, * in glória Patris.

Judex créderis * esse ventúrus.

\rubric{¶ Sequens versus dicitur flexis genibus.}

Te ergo quæsumus, tuis fámulis súbveni, * quos pretióso sánguine redemísti.

Ætérna fac cum Sanctis tuis * in glória numerári.

Salvum fac pópulum tuum, Dómine, * et bénedic hereditáti tuæ.

Et rege eos, * et extólle illos usque in ætérnum.

Per síngulos dies * benedícimus te;

Et laudámus nomen tuum in séculum, * et in séculum séculi.

Dignáre, Dómine, die isto * sine peccáto nos custodíre.

Miserére nostri, Dómine, * miserére nostri.

Fiat misericórdia tua, Dómine, super nos, * quemádmodum sperávimus in te.

In te, Dómine, sperávi: * non confúndar in ætérnum.

\rubric{Dicto Te Deum, aut ultimo Responsorio, statim incipiuntur Laudes a Versu Deus, in adjutórium, ut infra.}

\rubric{In privata vero recitatione Matutinum a Laudibus separari potest; quo in casu, post Hymnum Te Deum, vel post ultimum Responsorium, dicitur:}

\V\ Dóminus vobíscum. \R\ Et cum spíritu tuo.

\rubric{¶ Hic Versus Dóminus vobíscum non dicitur ab eo, qui non est saltem in ordine Diaconatus; sed ejus loco substituitur, etiam in aliis Horis, quoties in antecedentibus Precibus dictus non fuerit:}

\V\ Dómine, exáudi oratiónem meam. \R\ Et clamor meus ad te véniat.

Orémus.

\rubric{Et dicitur Oratio ut ad Laudes.}

\rubric{Post Orationem additur:}

\V\ Dóminus vobíscum. \R\ Et cum spíritu tuo.

\V\ Benedicámus Dómino. \R\ Deo grátias.

\V\ Fidélium ánimæ per misericórdiam Dei requiéscant in pace. \R\ Amen.

\rubric{Deinde dicitur Pater noster totum secreto.}

\rubric{In publica vero recitatione Matutinum in Choro a Laudibus separari non potest.}

\hourheading{Ad Laudes}

\rubric{Si Laudes extra Chorum separentur a Matutino, ante eas dicitur secreto Pater noster et Ave María; secus absolute incipiuntur, ut sequitur:}

\V\ Deus, in adjutórium.

\rubric{Deinde, sub congruentibus Antiphonis, dicuntur Psalmi et Canticum, prouti Officium occurrens requirit.}

\rubric{Repetita vero post ultimum Psalmum Antiphona, dicuntur Capitulum, Hymnus et Versus, prouti Officium occurrens exigit.}

\variantheading{In Officio dominicali}

\rubric{In omnibus Dominicis, extra Octavas Ssmi Corporis Christi ac Sacratissimi Cordis Jesu, Capitulum, Hymnus et Versus ut in Psalterio.}

\variantheading{In Officio feriali}

\rubric{In omnibus Feriis, Capitulum, Hymnus et Versus de Feria currenti, ut in Psalterio.}

\rubric{Recitato Versu, dicitur cum Antiphona convenienti sequens}

\canticumheading{Canticum Zachariæ}{Luc. 1, 68-79}

\lettrine[lines=2]{B}{ENEDICTUS} \textdagger\ Dóminus, Deus Israël, * quia visitávit, et fecit redemptiónem plebis suæ:

Et eréxit cornu salútis nobis: * in domo David, púeri sui.

Sicut locútus est per os sanctórum, * qui a sæculo sunt, Prophetárum ejus:

Salútem ex inimícis nostris, * et de manu ómnium, qui odérunt nos:

Ad faciéndam misericórdiam cum pátribus nostris: * et memorári testaménti sui sancti.

Jusjurándum, quod jurávit ad Abraham, patrem nostrum, * datúrum se nobis:

Ut sine timóre, de manu inimicórum nostrórum liberáti, * serviámus illi.

In sanctitáte, et justítia coram ipso, * ómnibus diébus nostris.

Et tu, puer, Prophéta Altíssimi vocáberis: * præíbis enim ante fáciem Dómini paráre vias ejus:

Ad dandam sciéntiam salútis plebi ejus: * in remissiónem peccatórum eórum:

Per víscera misericórdiæ Dei nostri: * in quibus visitávit nos, óriens ex alto:

Illumináre his, qui in ténebris, et in umbra mortis sedent: * ad dirigéndos pedes nostros in viam pacis.

Glória Patri, et Fílio, * et Spirítui Sancto.

Sicut erat in princípio, et nunc, et semper, * et in sæcula sæculórum. Amen.

\rubric{¶ Hic Versus Glória Patri dicitur in fine omnium Canticorum, nisi aliter adnotetur.}

\rubric{Expleto Cantico ac repetita Antiphona, si Preces non fuerint recitandæ, statim dicitur Versus Dóminus vobíscum cum Oratione et reliquis omnibus usque ad finem Horæ, ut infra.}

\rubric{In omnibus vero Vigiliis, si fiat Officium de Vigilia currenti, post repetitam ad Benedictus Antiphonam dicuntur flexis genibus sequentes Preces, quæ aliis temporibus omittuntur:}

KYRIE, eléison. Christe, eléison. Kýrie, eléison.

\rubric{Pater noster, quod dicitur a solo Hebdomadario totum clara voce, usque ad} \V\ Et ne nos indúcas in tentatiónem. \R\ Sed líbera nos a malo.

\V\ Ego dixi: Dómine, miserére mei. \R\ Sana ánimam meam, quia peccávi tibi. \V\ Convértere, Dómine, úsquequo? \R\ Et deprecábilis esto super servos tuos. \V\ Fiat misericórdia tua, Dómine, super nos. \R\ Quemádmodum sperávimus in te. \V\ Sacerdótes tui induántur justítiam. \R\ Et sancti tui exsúltent.

\V\ Orémus pro beatíssimo Papa nostro N. \R\ Dóminus consérvet eum, et vivíficet eum, et beátum fáciat eum in terra, et non tradat eum in ánimam inimicórum ejus.

\V\ Orémus et pro Antístite nostro N. \R\ Stet et pascat in fortitúdine tua, Dómine, in sublimitáte nóminis tui.

\rubric{¶ Romæ præcedens Versus cum suo Responsorio omittitur; alibi vero, ad litteram N., ab omnibus nomen diœcesani Episcopi exprimatur. Vacante Apostolica vel Episcopali Sede, alteruter vel uterque respectivus Versus cum suo Responsorio prætéritur.}

\V\ Dómine, salvum fac regem. \R\ Et exáudi nos in die, qua invocavérimus te. \V\ Salvum fac pópulum tuum, Dómine, et bénedic hereditáti tuæ. \R\ Et rege eos, et extólle illos usque in ætérnum. \V\ Meménto Congregatiónis tuæ. \R\ Quam possedísti ab inítio. \V\ Fiat pax in virtúte tua. \R\ Et abundántia in túrribus tuis. \V\ Orémus pro benefactóribus nostris. \R\ Retribúere dignáre, Dómine, ómnibus, nobis bona faciéntibus propter nomen tuum, vitam ætérnam. Amen. \V\ Orémus pro fidélibus defúnctis. \R\ Réquiem ætérnam dona eis, Dómine, et lux perpétua lúceat eis. \V\ Requiéscant in pace. \R\ Amen. \V\ Pro frátribus nostris abséntibus. \R\ Salvos fac servos tuos, Deus meus, sperántes in te. \V\ Pro afflíctis et captívis. \R\ Líbera eos, Deus Israël, ex ómnibus tribulatiónibus suis. \V\ Mitte eis, Dómine, auxílium de sancto. \R\ Et de Sion tuére eos. \V\ Dómine, Deus virtútum, convérte nos. \R\ Et osténde fáciem tuam, et salvi érimus. \V\ Exsúrge, Christe, ádjuva nos. \R\ Et líbera nos propter nomen tuum. \V\ Dómine, exáudi oratiónem meam. \R\ Et clamor meus ad te véniat.

\rubric{Absolutis Precibus, vel, si Preces locum non habeant, repetita Antiphona post Benedictus, statim dicitur:}

\V\ Dóminus vobíscum. \R\ Et cum spíritu tuo.

Orémus.

\rubric{Et dicitur Oratio conveniens. Postea fiunt Commemorationes, si occurrant.}

\rubric{Deinde in Officio semiduplici aut simplici tam de Tempore quam de Sanctis, præter omnes dies, in quibus habita sit Commemoratio Duplicis, vel factum fuerit Officium aut Commemoratio cujusvis Octavæ, fit sequens}

\suffragiumheading{Suffragium de omnibus Sanctis}

\antlabel{Ant.} Beáta Dei Génitrix Virgo María, Sanctíque omnes intercédant pro nobis ad Dóminum.

\V\ Mirificávit Dóminus Sanctos suos. \R\ Et exaudívit eos clamántes ad se.

Orémus. \oratiolabel{Oratio}

\lettrine[lines=2]{A}{} \textsc{cunctis} nos, quæsumus, Dómine, mentis et córporis defénde perículis: et, intercedénte beáta et gloriósa semper Vírgine Dei Genitríce María, cum beáto Joseph, beátis Apóstolis tuis Petro et Paulo, atque beáto N. et ómnibus Sanctis, salútem nobis tríbue benígnus et pacem; ut, destrúctis adversitátibus et erróribus univérsis, Ecclésia tua secúra tibi sérviat libertáte. Per eúmdem Dóminum nostrum. \R\ Amen.

\rubric{¶ In hac Oratione, ad litteram N., exprimitur nomen Titularis propriæ Ecclesiæ, dummodo Titulus non sit Persona divina vel Mysterium Domini, aut de ipso ratione Vigiliæ vel Festi factum non fuerit Officium vel Commemoratio, aut illius nomen non sit in eadem Suffragii Oratione expressum; ac nomina sanctorum Angelorum et sancti Joannis Baptistæ, si Titulares fuerint, præponuntur nomini sancti Joseph. In hisce autem omnibus casibus omittuntur verba atque beáto N.}

\rubric{Si vero Officium vel Commemoratio fuerit de beáta María Vírgine, tunc dicitur ut infra; atque insuper, sicubi Titularis juxta superiorem Rubricam ante omnes fuerit nominandus, ad Orationem particula cum omittitur ante nomen sanctorum Petri et Pauli ac præponitur nomini sancti Joseph.}

\antlabel{Ant.} Sancti omnes intercédant pro nobis ad Dóminum.

\V\ Mirificávit Dóminus Sanctos suos. \R\ Et exaudívit eos clamántes ad se.

Orémus. \oratiolabel{Oratio}

\lettrine[lines=2]{A}{} \textsc{cunctis} nos, quæsumus, Dómine, mentis et córporis defénde perículis: et, intercedénte beáto Joseph, cum beátis Apóstolis tuis Petro et Paulo, atque beáto N. et ómnibus Sanctis, salútem nobis tríbue benígnus et pacem; ut, destrúctis adversitátibus et erróribus univérsis, Ecclésia tua secúra tibi sérviat libertáte. Per Dóminum. \R\ Amen.

\rubric{Post ultimam Orationem additur:}

\V\ Dóminus vobíscum. \R\ Et cum spíritu tuo.

\V\ Benedicámus Dómino. \R\ Deo grátias.

\V\ Fidélium ánimæ per misericórdiam Dei requiéscant in pace. \R\ Amen.

\rubric{Deinde, in publica Officii recitatione, si discedendum sit a Choro, itemque in recitatione privata, si tunc terminetur Officium; alioquin, in fine ultimæ Horæ dicitur Pater noster, totum secreto; eoque recitato, subjungitur:}

\V\ Dóminus det nobis suam pacem. \R\ Et vitam ætérnam. Amen.

\rubric{Et immediate dicitur, cum suis Versu et Oratione, finalis beatæ Mariæ Virginis Antiphona, quæ inferius post Completorium assignatur.}

\rubric{Postea concluditur:}

\V\ Divínum auxílium máneat semper nobíscum. \R\ Amen.

% ===== END OF BATCH 1 -- Covers autumnalis.md lines 1-463 (reused
% verbatim from aestiva.tex's Batch 1, per the byte-identity noted
% above): the opening of the Ordinarium Divini Officii (its
% introductory indulgence rubrics and the "Aperi, Dómine" Oratio),
% "Ante singulas Horas" (the Pater noster, Ave María, and Credo
% spelled out in full for the first and only time -- each given a
% lettrine since the source presents them in full capitals as
% complete standalone prayers, unlike their thousands of later
% plain-text abbreviated recurrences throughout the volume), "Ad
% Matutinum" (the Invitatorium Psalm 94, the three-Nocturn apparatus
% with its Absolutios, Benedictiones, and "Pro N Lectione" headings,
% the Officium trium Lectionum variant, and the Te Deum under its
% customary heading "Hymnus Ambrosianus"), and "Ad Laudes" through
% the Suffragium de omnibus Sanctis (in both its Marian-titular and
% Joseph-titular forms) and the closing versicles. This Pars
% Autumnalis Ordinarium matches its Pars Aestiva counterpart exactly
% (both being non-Lenten/non-Paschal volumes) and is notably simpler
% than Pars Verna's (Batch 1 of verna.tex), which carries the extra
% "Tempore Quadragesimæ"/"Tempore Passionis" Invitatorium-Hymnus
% variants and the Paschaltide "Commemoratio de Cruce". Roman
% numerals in the "Pro N Lectione" headings are capitalized (e.g.
% \variantheading{Pro I Lectione}) though the source itself gives them
% lowercase ("Pro i Lectione"), matching the established
% normalization already present throughout hiemalis.tex, verna.tex,
% and aestiva.tex for this specific heading type. The document
% continues with "Ad Primam" (autumnalis.md line 465) =====
\hourheading{Ad Primam}

Pater, Ave et Credo.

\V\ Deus, in adjutórium.

\hymnusheading{Hymnus}

\lettrine[lines=2]{J}{AM} lucis orto sídere,\\
Deum precémur súpplices,\\
Ut in diúrnis áctibus\\
Nos servet a nocéntibus.

Linguam refrénans témperet,\\
Ne litis horror ínsonet:\\
Visum fovéndo cóntegat,\\
Ne vanitátes háuriat.

Sint pura cordis íntima,\\
Absístat et vecórdia:\\
Carnis terat supérbiam\\
Potus cibíque párcitas.

Ut, cum dies abscésserit\\
Noctémque sors redúxerit,\\
Mundi per abstinéntiam\\
Ipsi canámus glóriam.

Deo Patri sit glória,\\
Ejúsque soli Fílio,\\
Cum Spíritu Paráclito\\
Nunc et per omne sæculum. Amen.

\rubric{¶ Conclusio communis, in Hymno præcedenti, et in aliis idem metrum habentibus, semper omittitur, quando specialis in omnibus Horis adhibenda præscribitur; et, si plures Conclusiones propriæ occurrant, sumitur Conclusio Officii currentis, aut secus Officii ipsa die et primo quidem loco inter cetera propriam Conclusionem habentia commemorandi, aut demum de occurrenti Octava communi.}

\rubric{Expleto Hymno, conveniens dicitur Antiphona usque ad Asteriscum *, prouti occurrens Officium requirit.}

\variantheading{In Officio dominicali}

\rubric{In omnibus Dominicis, extra Octavas Ssmi Corporis Christi et Sacrat. Cordis Jesu, Antiphona ut in Psalterio.}

\variantheading{In Officio feriali}

\rubric{In omnibus Feriis, Antiphona de Feria currenti ut in Psalterio.}

\rubric{Post enuntiatam Antiphonam dicuntur tres Psalmi Officii diei congruentes.}

\rubric{¶ Quando autem adhibetur posterius Laudum schema, Psalmus, in priori Laudum ferialium schemate omissus, ad Primam Feriæ currentis post alios Psalmos resumitur.}

\rubric{Repetita post ultimum Psalmum integre Antiphona, conveniens dicitur Capitulum; nimirum:}

\rubric{In omnibus Dominicis, in Officio cujuslibet Festi vel Octavæ ac sanctæ Mariæ in Sabbato:}

\capitulumlabel{Capitulum}{1 Tim. 1, 17}

\lettrine[lines=2]{R}{EGI} sæculórum immortáli et invisíbili, soli Deo honor et glória in sæcula sæculórum. Amen. \R\ Deo grátias.

\rubric{¶ Sic respondetur in fine omnium Capitulorum.}

\rubric{In omnibus autem Feriis et in Vigiliis communibus:}

\capitulumlabel{Capitulum}{Zach. 8, 19}

\lettrine[lines=2]{P}{ACEM} et veritátem dilígite, ait Dóminus omnípotens.

\rubric{Dicto Capitulo, subjungitur Responsorium breve, ut sequitur. Versus tamen Qui sedes ad dexteram Patris semper omittitur, quando specialis Versus adhibendus præscribitur, Officii nempe currentis, aut secus Officii ad Laudes et primo quidem loco inter cetera specialem Versum habentia commemorati, aut demum de occurrenti Octava communi proprius.}

\R\ br. Christe, Fili Dei vivi, * Miserére nobis. \rubric{Et repetitur:} Christe, Fili Dei vivi, miserére nobis. \V\ Qui sedes ad déxteram Patris. \R\ Miserére nobis. \V\ Glória Patri, et Fílio, et Spirítui Sancto. \R\ Christe, Fili Dei vivi, miserére nobis. \V\ Exsúrge, Christe, ádjuva nos. \R\ Et líbera nos propter nomen tuum.

\rubric{¶ Ita dicuntur semper Responsoria brevia, etiam ad Tertiam, Sextam, Nonam et Completorium.}

\rubric{Expleto Responsorio brevi, si Preces non fuerint recitandæ, statim dicitur Versus Dóminus vobíscum cum subsequenti Oratione et reliquis omnibus usque ad finem Horæ ut infra.}

\rubric{In quolibet Officio semiduplici, in Officio sanctæ Mariæ in Sabbato, in Festis simplicibus, atque in Feriis per Annum communibus, post Responsorium breve dicuntur sequentes Preces dominicales; quæ tamen omittuntur omnibus diebus, in quibus ad Laudes habita sit Commemoratio Duplicis, vel factum fuerit Officium aut Commemoratio cujusvis Octavæ:}

KYRIE, eléison. Christe, eléison. Kýrie, eléison.

\rubric{Pater noster secreto usque ad} \V\ Et ne nos indúcas in tentatiónem. \R\ Sed líbera nos a malo.

\rubric{Credo in Deum secreto usque ad} \V\ Carnis resurrectiónem. \R\ Vitam ætérnam. Amen.

\V\ Et ego ad te, Dómine, clamávi. \R\ Et mane orátio mea prævéniet te. \V\ Repleátur os meum laude. \R\ Ut cantem glóriam tuam, tota die magnitúdinem tuam. \V\ Dómine, avérte fáciem tuam a peccátis meis. \R\ Et omnes iniquitátes meas dele. \V\ Cor mundum crea in me, Deus. \R\ Et spíritum rectum ínnova in viscéribus meis. \V\ Ne projícias me a fácie tua. \R\ Et spíritum sanctum tuum ne áuferas a me. \V\ Redde mihi lætítiam salutáris tui. \R\ Et spíritu principáli confírma me. \V\ Adjutórium nostrum in nómine Dómini. \R\ Qui fecit cælum et terram.

\rubric{Deinde Hebdomadarius facit Confessionem cum Absolutione, eamque Chorus repetit, ut infra post feriales Preces.}

\rubric{In omnibus vero Vigiliis, si ad Laudes recitatæ fuerint preces, post Responsorium breve dicuntur flexis genibus sequentes Preces feriales:}

KYRIE, eléison. Christe, eléison. Kýrie, eléison.

\rubric{Pater noster secreto usque ad} \V\ Et ne nos indúcas in tentatiónem. \R\ Sed líbera nos a malo.

\rubric{Credo in Deum secreto usque ad} \V\ Carnis resurrectiónem. \R\ Vitam ætérnam. Amen.

\V\ Et ego ad te, Dómine, clamávi. \R\ Et mane orátio mea prævéniet te. \V\ Repleátur os meum laude. \R\ Ut cantem glóriam tuam, tota die magnitúdinem tuam. \V\ Dómine, avérte fáciem tuam a peccátis meis. \R\ Et omnes iniquitátes meas dele. \V\ Cor mundum crea in me, Deus. \R\ Et spíritum rectum ínnova in viscéribus meis. \V\ Ne projícias me a fácie tua. \R\ Et spíritum sanctum tuum ne áuferas a me. \V\ Redde mihi lætítiam salutáris tui. \R\ Et spíritu principáli confírma me. \V\ Eripe me, Dómine, ab hómine malo. \R\ A viro iníquo éripe me. \V\ Eripe me de inimícis meis, Deus meus. \R\ Et ab insurgéntibus in me líbera me. \V\ Eripe me de operántibus iniquitátem. \R\ Et de viris sánguinum salva me. \V\ Sic psalmum dicam nómini tuo in sæculum sæculi. \R\ Ut reddam vota mea de die in diem. \V\ Exáudi nos, Deus, salutáris noster. \R\ Spes ómnium fínium terræ, et in mari longe. \V\ Deus, in adjutórium meum inténde. \R\ Dómine, ad adjuvándum me festína. \V\ Sanctus Deus, Sanctus fortis, Sanctus immortális. \R\ Miserére nobis. \V\ Bénedic, ánima mea, Dómino. \R\ Et ómnia, quæ intra me sunt, nómini sancto ejus. \V\ Bénedic, ánima mea, Dómino. \R\ Et noli oblivísci omnes retributiónes ejus. \V\ Qui propitiátur ómnibus iniquitátibus tuis. \R\ Qui sanat omnes infirmitátes tuas. \V\ Qui rédimit de intéritu vitam tuam. \R\ Qui corónat te in misericórdia et miseratiónibus. \V\ Qui replet in bonis desidérium tuum. \R\ Renovábitur ut áquilæ juvéntus tua. \V\ Adjutórium nostrum in nómine Dómini. \R\ Qui fecit cælum et terram.

\rubric{Deinde Hebdomadarius facit Confessionem:}

\lettrine[lines=2]{C}{ONFITEOR} Deo omnipoténti, beátæ Maríæ semper Vírgini, beáto Michaéli Archángelo, beáto Joánni Baptístæ, sanctis Apóstolis Petro et Paulo, ómnibus Sanctis, et vobis, fratres, quia peccávi nimis cogitatióne, verbo et ópere: mea culpa, mea culpa, mea máxima culpa. Ideo precor beátam Maríam semper Vírginem, beátum Michaélem Archángelum, beátum Joánnem Baptístam, sanctos Apóstolos Petrum et Paulum, omnes Sanctos, et vos, fratres, oráre pro me ad Dóminum, Deum nostrum.

\rubric{Chorus respondet:}

\lettrine[lines=2]{M}{ISEREATUR} tui omnípotens Deus, et, dimíssis peccátis tuis, perdúcat te ad vitam ætérnam. \R\ Amen.

\rubric{Deinde repetit Confessionem, et ubi dicitur: vobis, fratres, et vos, fratres, dicat: tibi, pater, et te, pater.}

\rubric{Facta Confessione a Choro, Hebdomadarius dicit:}

\lettrine[lines=2]{M}{ISEREATUR} vestri omnípotens Deus, et, dimíssis peccátis vestris, perdúcat vos ad vitam ætérnam. \R\ Amen.

\lettrine[lines=2]{I}{NDULGENTIAM}, absolutiónem, et remissiónem peccatórum nostrórum tríbuat nobis omnípotens et miséricors Dóminus. \R\ Amen.

\rubric{Extra Chorum, si unus vel duo tantum recitent Officium et in Choro Monialium, semel tantum ac simul ab omnibus ita fit Confessio:}

\lettrine[lines=2]{C}{ONFITEOR} Deo omnipoténti, beátæ Maríæ semper Vírgini, beáto Michaéli Archángelo, beáto Joánni Baptístæ, sanctis Apóstolis Petro et Paulo, et ómnibus Sanctis, quia peccávi nimis cogitatióne, verbo et ópere: mea culpa, mea culpa, mea máxima culpa. Ideo precor beátam Maríam semper Vírginem, beátum Michaélem Archángelum, beátum Joánnem Baptístam, sanctos Apóstolos Petrum et Paulum, et omnes Sanctos, oráre pro me ad Dóminum, Deum nostrum.

\rubric{Deinde dicitur:}

\lettrine[lines=2]{M}{ISEREATUR} nostri omnípotens Deus, et, dimíssis peccátis nostris, perdúcat nos ad vitam ætérnam. Amen.

\lettrine[lines=2]{I}{NDULGENTIAM}, absolutiónem, et remissiónem peccatórum nostrórum tríbuat nobis omnípotens et miséricors Dóminus. Amen.

\rubric{Facta Absolutione, subjungitur:}

\V\ Dignáre, Dómine, die isto. \R\ Sine peccáto nos custodíre. \V\ Miserére nostri, Dómine. \R\ Miserére nostri. \V\ Fiat misericórdia tua, Dómine, super nos. \R\ Quemádmodum sperávimus in te. \V\ Dómine, exáudi oratiónem meam. \R\ Et clamor meus ad te véniat.

\rubric{Absolutis Precibus, vel, si Preces locum non habeant, post Responsorium breve statim dicitur:}

\V\ Dóminus vobíscum. \R\ Et cum spíritu tuo.

Orémus. \oratiolabel{Oratio}

\lettrine[lines=2]{D}{OMINE}, Deus omnípotens, qui ad princípium hujus diéi nos perveníre fecísti: tua nos hódie salva virtúte; ut in hac die ad nullum declinémus peccátum, sed semper ad tuam justítiam faciéndam nostra procédant elóquia, dirigántur cogitatiónes et ópera. Per Dóminum. \R\ Amen.

\V\ Dóminus vobíscum. \R\ Et cum spíritu tuo. \V\ Benedicámus Dómino. \R\ Deo grátias.

\rubric{Deinde in Choro legitur Martyrologium, quod laudabiliter fit etiam extra Chorum. Postea Hebdomadarius dicit:}

\V\ Pretiósa in conspéctu Dómini. \R\ Mors Sanctórum ejus.

\rubric{Et absolute, sine} Orémus:

\oratiolabel{Oratio}

\lettrine[lines=2]{S}{ANCTA} María, et omnes Sancti intercédant pro nobis ad Dóminum, ut nos mereámur ab eo adjuvári et salvári, qui vivit et regnat in sæcula sæculórum. \R\ Amen.

\V\ Deus, in adjutórium meum inténde. \R\ Dómine, ad adjuvándum me festína. \rubric{Et dicitur ter; ultimo additur:} Glória Patri, et Fílio, et Spirítui Sancto. Sicut erat in princípio, et nunc, et semper, et in sæcula sæculórum. Amen.

Kýrie, eléison. Christe, eléison. Kýrie, eléison.

Pater noster \rubric{secreto usque ad} \V\ Et ne nos indúcas in tentatiónem. \R\ Sed líbera nos a malo.

\V\ Réspice in servos tuos, Dómine, et in ópera tua, et dírige fílios eórum. \R\ Et sit splendor Dómini, Dei nostri, super nos, et ópera mánuum nostrárum dírige super nos, et opus mánuum nostrárum dírige. \V\ Glória Patri, et Fílio, et Spirítui Sancto. \R\ Sicut erat in princípio, et nunc, et semper, et in sæcula sæculórum. Amen.

Orémus. \oratiolabel{Oratio}

\lettrine[lines=2]{D}{IRIGERE} et sanctificáre, régere et gubernáre dignáre, Dómine Deus, Rex cæli et terræ, hódie corda et córpora nostra, sensus, sermónes et actus nostros in lege tua, et in opéribus mandatórum tuórum: ut hic, et in ætérnum, te auxiliánte, salvi et líberi esse mereámur, Salvátor mundi: Qui vivis et regnas in sæcula sæculórum. \R\ Amen.

\rubric{Deinde dicitur Lectio brevis, ad absolutionem Capituli, prouti Officium occurrens requirit.}

\V\ Jube, domne, benedícere. \benedlabel{Bened.} Dies et actus nostros in sua pace dispónat Dóminus omnípotens. \R\ Amen.

\variantheading{In Officio dominicali et feriali}

\rubric{A Feria II usque ad Feriam IV post Festum sanctissimæ Trinitatis, et a Sabbato post Octavam Sacratissimi Cordis Jesu et deinceps.}

\lectioheading{Lectio brevis}{2 Thess. 3, 5}

\lettrine[lines=2]{D}{OMINUS} autem dírigat corda et córpora nostra in caritáte Dei, et patiéntia Christi. Tu autem, Dómine, miserére nobis. \R\ Deo grátias.

\rubric{¶ Ita semper concluduntur Lectiones breves ad Primam, et ita in earum fine respondetur.}

\rubric{In Officio autem cujuslibet Festi vel Octavæ, atque S. Mariæ in Sabbato; itemque in Dominicis infra Octavas Ssmi Corporis Christi ac Sacrat. Cordis Jesu, si de eis fiat Officium:}

\rubric{Pro Lectione brevi sumitur Capitulum, quod ad Nonam ponitur ipsius Officii currentis, et in fine dicitur: Tu autem, Dómine, ut supra.}

\rubric{Post Lectionem brevem dicitur:}

\V\ Adjutórium nostrum in nómine Dómini. \R\ Qui fecit cælum et terram.

\V\ Benedícite. \R\ Deus.

\benedlabel{Bened.} Dóminus nos benedícat, et ab omni malo deféndat, et ad vitam perdúcat ætérnam. Et fidélium ánimæ per misericórdiam Dei requiéscant in pace. \R\ Amen.

\rubric{Deinde dicitur tantum Pater noster secreto, nisi sequatur alia Hora.}

\hourheading{Ad Tertiam}

Pater noster et Ave.

\V\ Deus, in adjutórium.

\hymnusheading{Hymnus}

\lettrine[lines=2]{N}{UNC}, Sancte, nobis, Spíritus,\\
Unum Patri cum Fílio,\\
Dignáre promptus íngeri\\
Nostro refúsus péctori.

Os, lingua, mens, sensus, vigor,\\
Confessiónem pérsonent,\\
Flamméscat igne cáritas,\\
Accéndat ardor próximos.

Præsta, Pater piíssime,\\
Patríque compar Unice,\\
Cum Spíritu Paráclito\\
Regnans per omne sæculum. Amen.

\rubric{Expleto Hymno, conveniens dicitur Antiphona usque ad Asteriscum *, et tres Psalmi, prouti Officium occurrens requirit; atque Psalmis absolutis, integra repetitur Antiphona.}

\variantheading{In Officio dominicali}

\rubric{In omnibus Dominicis, extra Octavas Ssmi Corporis Christi ac Sacrat. Cordis Jesu, Aña ut in Psalterio, et dicitur sequens}

\capitulumlabel{Capitulum}{1 Joann. 4, 16}

\lettrine[lines=2]{D}{EUS} cáritas est: et qui manet in caritáte, in Deo manet, et Deus in eo.

\R\ br. Inclína cor meum, Deus, * In testimónia tua. Inclína. \V\ Avérte óculos meos, ne vídeant vanitátem: in via tua vivífica me. In testimónia tua. Glória Patri. Inclína.

\V\ Ego dixi: Dómine, miserére mei. \R\ Sana ánimam meam, quia peccávi tibi.

\variantheading{In Officio feriali}

\rubric{In omnibus Feriis, Antiphona de Feria currenti, ut in Psalterio; et dicitur quotidie sequens}

\capitulumlabel{Capitulum}{Jer. 17, 14}

\lettrine[lines=2]{S}{ANA} me, Dómine, et sanábor: salvum me fac, et salvus ero: quóniam laus mea tu es.

\R\ br. Sana ánimam meam, * Quia peccávi tibi. Sana. \V\ Ego dixi: Dómine, miserére mei. Quia. Glória Patri. Sana.

\V\ Adjútor meus esto, ne derelínquas me. \R\ Neque despícias me, Deus, salutáris meus.

\rubric{Expleto Responsorio brevi, si Preces non fuerint recitandæ, statim dicitur Versus Dóminus vobíscum cum Oratione et reliquis omnibus usque ad finem Horæ, ut infra.}

\rubric{Quando vero ad Laudes recitatæ sunt Preces, post Responsorium breve dicuntur sequentes flexis genibus:}

KYRIE, eléison. Christe, eléison. Kýrie, eléison.

Pater noster \rubric{secreto usque ad} \V\ Et ne nos indúcas in tentatiónem. \R\ Sed líbera nos a malo. \V\ Dómine, Deus virtútum, convérte nos. \R\ Et osténde fáciem tuam, et salvi érimus. \V\ Exsúrge, Christe, ádjuva nos. \R\ Et líbera nos propter nomen tuum. \V\ Dómine, exáudi oratiónem meam. \R\ Et clamor meus ad te véniat.

\rubric{Absolutis Precibus, vel, si Preces locum non habeant, post Responsorium breve statim dicitur:}

\V\ Dóminus vobíscum. \R\ Et cum spíritu tuo.

Orémus.

\rubric{Et dicitur Oratio conveniens.}

\rubric{Post Orationem additur:}

\V\ Dóminus vobíscum. \R\ Et cum spíritu tuo.

\V\ Benedicámus Dómino. \R\ Deo grátias.

\V\ Fidélium ánimæ per misericórdiam Dei requiéscant in pace. \R\ Amen.

\rubric{Deinde dicitur tantum Pater noster secreto, nisi sequatur alia Hora.}

\hourheading{Ad Sextam}

Pater noster et Ave.

\V\ Deus, in adjutórium.

\hymnusheading{Hymnus}

\lettrine[lines=2]{R}{ECTOR} potens, verax Deus,\\
Qui témperas rerum vices,\\
Splendóre mane illúminas,\\
Et ígnibus merídiem:

Exstíngue flammas lítium,\\
Aufer calórem nóxium,\\
Confer salútem córporum,\\
Verámque pacem córdium.

Præsta, Pater piíssime,\\
Patríque compar Unice,\\
Cum Spíritu Paráclito\\
Regnans per omne sæculum. Amen.

\rubric{Expleto Hymno, conveniens dicitur Antiphona usque ad Asteriscum *, et tres Psalmi prouti Officium occurrens requirit; atque, Psalmis absolutis, integra repetitur Antiphona.}

\variantheading{In Officio dominicali}

\rubric{In omnibus Dominicis, extra Octavas Ssmi Corporis Christi ac Sacrat. Cordis Jesu, Aña ut in Psalterio, et dicitur sequens}

\capitulumlabel{Capitulum}{Gal. 6, 2}

\lettrine[lines=2]{A}{LTER} altérius ónera portáte, et sic adimplébitis legem Christi.

\R\ br. In ætérnum, Dómine, * Pérmanet verbum tuum. In ætérnum. \V\ In sæculum sæculi véritas tua. Pérmanet. Glória Patri. In ætérnum.

\V\ Dóminus regit me, et nihil mihi déerit. \R\ In loco páscuæ ibi me collocávit.

\variantheading{In Officio feriali}

\rubric{In omnibus Feriis, Antiphona de Feria currenti, ut in Psalterio; et dicitur quotidie sequens}

\capitulumlabel{Capitulum}{Rom. 13, 8}

\lettrine[lines=2]{N}{EMINI} quidquam debeátis, nisi ut ínvicem diligátis: qui enim díligit próximum, legem implévit.

\R\ br. Benedícam Dóminum * In omni témpore. Benedícam. \V\ Semper laus ejus in ore meo. In. Glória Patri. Benedícam.

\V\ Dóminus regit me, et nihil mihi déerit. \R\ In loco páscuæ ibi me collocávit.

\rubric{Expleto Responsorio brevi, si Preces non fuerint recitandæ, statim dicitur Versus Dóminus vobíscum cum Oratione et reliquis omnibus usque ad finem Horæ, ut infra.}

\rubric{Quando vero ad Laudes recitatæ sunt Preces, post Responsorium breve dicuntur sequentes flexis genibus:}

KYRIE, eléison. Christe, eléison. Kýrie, eléison.

Pater noster \rubric{secreto usque ad} \V\ Et ne nos indúcas in tentatiónem. \R\ Sed líbera nos a malo. \V\ Dómine, Deus virtútum, convérte nos. \R\ Et osténde fáciem tuam, et salvi érimus. \V\ Exsúrge, Christe, ádjuva nos. \R\ Et líbera nos propter nomen tuum. \V\ Dómine, exáudi oratiónem meam. \R\ Et clamor meus ad te véniat.

\rubric{Absolutis Precibus, vel, si Preces locum non habeant, post Responsorium breve statim dicitur:}

\V\ Dóminus vobíscum. \R\ Et cum spíritu tuo.

Orémus.

\rubric{Et dicitur Oratio conveniens.}

\rubric{Post Orationem additur:}

\V\ Dóminus vobíscum. \R\ Et cum spíritu tuo.

\V\ Benedicámus Dómino. \R\ Deo grátias.

\V\ Fidélium ánimæ per misericórdiam Dei requiéscant in pace. \R\ Amen.

\rubric{Deinde dicitur tantum Pater noster secreto, nisi sequatur alia Hora.}

\hourheading{Ad Nonam}

Pater noster et Ave.

\V\ Deus, in adjutórium.

\hymnusheading{Hymnus}

\lettrine[lines=2]{R}{ERUM}, Deus, tenax vigor,\\
Immótus in te pérmanens,\\
Lucis diúrnæ témpora\\
Successíbus detérminans:

Largíre lumen véspere,\\
Quo vita nusquam décidat,\\
Sed præmium mortis sacræ\\
Perénnis instet glória.

Præsta, Pater piíssime,\\
Patríque compar Unice,\\
Cum Spíritu Paráclito\\
Regnans per omne sæculum. Amen.

\rubric{Expleto Hymno, conveniens dicitur Antiphona usque ad Asteriscum *, et tres Psalmi, prouti Officium occurrens requirit; atque, Psalmis absolutis, integra repetitur Antiphona.}

\variantheading{In Officio dominicali}

\rubric{In omnibus Dominicis, extra Octavas Ssmi Corporis Christi ac Sacrat. Cordis Jesu, Aña ut in Psalterio, et dicitur sequens}

\capitulumlabel{Capitulum}{1 Cor. 6, 20}

\lettrine[lines=2]{E}{MPTI} enim estis prétio magno. Glorificáte et portáte Deum in córpore vestro.

\R\ br. Clamávi in toto corde meo: * Exáudi me, Dómine. Clamávi. \V\ Justificatiónes tuas requíram. Exáudi. Glória Patri. Clamávi.

\V\ Ab occúltis meis munda me, Dómine. \R\ Et ab aliénis parce servo tuo.

\variantheading{In Officio feriali}

\rubric{In omnibus Feriis, Antiphona de Feria currenti, ut in Psalterio; et dicitur quotidie sequens}

\capitulumlabel{Capitulum}{1 Petr. 1, 17-19}

\lettrine[lines=2]{I}{N} timóre incolátus vestri témpore conversámini: sciéntes quod non corruptibílibus auro vel argénto redémpti estis, sed pretióso sánguine quasi Agni immaculáti Christi.

\R\ br. Rédime me, Dómine, * Et miserére mei. Rédime. \V\ Pes enim meus stetit in via recta. Et. Glória Patri. Rédime.

\V\ Ab occúltis meis munda me, Dómine. \R\ Et ab aliénis parce servo tuo.

\rubric{Expleto Responsorio brevi, si Preces non fuerint recitandæ, statim dicitur Versus Dóminus vobíscum cum Oratione et reliquis omnibus usque ad finem Horæ, ut infra.}

\rubric{Quando vero ad Laudes recitatæ sunt Preces, post Responsorium breve dicuntur sequentes flexis genibus:}

KYRIE, eléison. Christe, eléison. Kýrie, eléison.

Pater noster \rubric{secreto usque ad} \V\ Et ne nos indúcas in tentatiónem. \R\ Sed líbera nos a malo. \V\ Dómine, Deus virtútum, convérte nos. \R\ Et osténde fáciem tuam, et salvi érimus. \V\ Exsúrge, Christe, ádjuva nos. \R\ Et líbera nos propter nomen tuum. \V\ Dómine, exáudi oratiónem meam. \R\ Et clamor meus ad te véniat.

\rubric{Absolutis Precibus, vel, si Preces locum non habeant, post Responsorium breve statim dicitur:}

\V\ Dóminus vobíscum. \R\ Et cum spíritu tuo.

Orémus.

\rubric{Et dicitur Oratio conveniens.}

\rubric{Post Orationem additur:}

\V\ Dóminus vobíscum. \R\ Et cum spíritu tuo.

\V\ Benedicámus Dómino. \R\ Deo grátias.

\V\ Fidélium ánimæ per misericórdiam Dei requiéscant in pace. \R\ Amen.

\rubric{Deinde dicitur tantum Pater noster secreto, nisi sequatur alia Hora.}

% ===== END OF BATCH 2 -- Covers autumnalis.md lines 465-874 (reused
% verbatim from aestiva.tex's Batch 2, per the byte-identity noted
% above), completing "Ad Primam" (hymn "Jam lucis orto sidere," the
% full Confiteor/Misereatur/Indulgentiam apparatus, the Martyrologium
% cue, the "Dirigere et sanctificare" Oratio, and the Lectio brevis,
% here given in its single simplified form -- "In Officio dominicali
% et feriali" -- rather than the three-way Lenten/Paschal-season
% branching found in verna.tex's Prime, since this volume's Lectio
% brevis never varies by season), "Ad Tertiam," "Ad Sextam," and "Ad
% Nonam," each with its proper hymn and its Sunday/weekday
% Capitulum-and-brief-responsory pair. Two spellings of the same word
% recur within this single batch exactly as the source alternates
% them: "sæculum"/"sæculi" at Sexta's responsory (line 758) versus
% the different spelling pattern seen in the Te Deum of Batch 1 --
% each transcribed as its own occurrence gives it, per the standing
% rule against silent normalization. The document continues with
% "Ad Vesperas" (autumnalis.md line 875) =====
\V\ Dóminus vobíscum. \R\ Et cum spíritu tuo.

Orémus.

\rubric{Et dicitur Oratio conveniens.}

\rubric{Postea fiunt Commemorationes, si occurrant.}

\rubric{Deinde, in Officio semiduplici aut simplici tam de Tempore quam de Sanctis, præter omnes dies, in quibus habita sit Commemoratio Duplicis, etiam concurrentis, vel factum fuerit Officium aut Commemoratio cujusvis Octavæ, fit sequens}

\suffragiumheading{Suffragium de omnibus Sanctis}

\antlabel{Ant.} Beáta Dei Génitrix Virgo María, Sanctíque omnes intercédant pro nobis ad Dóminum.

\V\ Mirificávit Dóminus Sanctos suos. \R\ Et exaudívit eos clamántes ad se.

Orémus. \oratiolabel{Oratio}

\lettrine[lines=2]{A}{} \textsc{cunctis} nos, quæsumus, Dómine, mentis et córporis defénde perículis: et, intercedénte beáta et gloriósa semper Vírgine Dei Genitríce María, cum beáto Joseph, beátis Apóstolis tuis Petro et Paulo, atque beáto N. et ómnibus Sanctis, salútem nobis tríbue benígnus et pacem; ut, destrúctis adversitátibus et erróribus univérsis, Ecclésia tua secúra tibi sérviat libertáte. Per eúmdem Dóminum. \R\ Amen.

\rubric{¶ In hac Oratione, ad litteram N., exprimitur nomen Titularis propriæ Ecclesiæ, dummodo Titulus non sit Persona divina vel Mysterium Domini, aut de ipso ratione Festi factum non fuerit Officium vel Commemoratio, aut illius nomen non sit in eadem Suffragii Oratione expressum; ac nomina sanctorum Angelorum et sancti Joannis Baptistæ, si Titulares fuerint, præponuntur nomini sancti Joseph. In hisce autem omnibus casibus omittuntur verba atque beáto N.}

\rubric{Si vero Officium vel Commemoratio fuerit de beáta María Vírgine, tunc dicitur ut infra; atque insuper, sicubi Titularis juxta superiorem Rubricam ante omnes fuerit nominandus, ad Orationem particula cum omittitur ante nomen sanctorum Petri et Pauli ac præponitur nomini sancti Joseph.}

\antlabel{Ant.} Sancti omnes intercédant pro nobis ad Dóminum.

\V\ Mirificávit Dóminus Sanctos suos. \R\ Et exaudívit eos clamántes ad se.

Orémus. \oratiolabel{Oratio}

\lettrine[lines=2]{A}{} \textsc{cunctis} nos, quæsumus, Dómine, mentis et córporis defénde perículis: et, intercedénte beáto Joseph, cum beátis Apóstolis tuis Petro et Paulo, atque beáto N. et ómnibus Sanctis, salútem nobis tríbue benígnus et pacem; ut, destrúctis adversitátibus et erróribus univérsis, Ecclésia tua secúra tibi sérviat libertáte. Per Dóminum.

\rubric{Post ultimam Orationem additur:}

\V\ Dóminus vobíscum. \R\ Et cum spíritu tuo.

\V\ Benedicámus Dómino. \R\ Deo grátias.

\V\ Fidélium ánimæ per misericórdiam Dei requiéscant in pace. \R\ Amen.

\rubric{Si post Vesperas immediate sequatur Completorium, dicto Versu Fidelium animæ, statim incipitur Versus Jube, domne, benedicere, ut infra ad Completorium; secus autem, si tunc terminetur Officium, dicitur tantum Pater noster secreto.}

\hourheading{Ad Completorium}

\rubric{Ante Completorium non dicitur Pater noster et Ave María, sicut in aliis Horis, sed absolute Lector incipit:}

\V\ Jube, domne, benedícere. \benedlabel{Bened.} Noctem quiétam, et finem perféctum concédat nobis Dóminus omnípotens. \R\ Amen.

\lectioheading{Lectio brevis}{1 Petr. 5, 8-9}

\lettrine[lines=2]{F}{RATRES}: Sóbrii estóte, et vigiláte: quia adversárius vester diábolus tamquam leo rúgiens círcuit, quærens quem dévoret: cui resístite fortes in fide. Tu autem, Dómine, miserére nobis. \R\ Deo grátias.

\V\ Adjutórium nostrum in nómine Dómini. \R\ Qui fecit cælum et terram.

Pater noster \rubric{quod dicitur totum secreto.}

\rubric{Deinde Hebdomadarius facit Confessionem:}

\lettrine[lines=2]{C}{ONFITEOR} Deo omnipoténti, beátæ Maríæ semper Vírgini, beáto Michaéli Archángelo, beáto Joánni Baptístæ, sanctis Apóstolis Petro et Paulo, ómnibus Sanctis, et vobis, fratres, quia peccávi nimis cogitatióne, verbo et ópere: mea culpa, mea culpa, mea máxima culpa. Ideo precor beátam Maríam semper Vírginem, beátum Michaélem Archángelum, beátum Joánnem Baptístam, sanctos Apóstolos Petrum et Paulum, omnes Sanctos, et vos, fratres, oráre pro me ad Dóminum, Deum nostrum.

\rubric{Chorus respondet:}

\lettrine[lines=2]{M}{ISEREATUR} tui omnípotens Deus, et, dimíssis peccátis tuis, perdúcat te ad vitam ætérnam. \R\ Amen.

\rubric{Deinde repetit Confessionem, et ubi dicitur: vobis, fratres, et vos, fratres, dicat: tibi, pater, et te, pater.}

\rubric{Facta Confessione a Choro, Hebdomadarius dicit:}

\lettrine[lines=2]{M}{ISEREATUR} vestri omnípotens Deus, et, dimíssis peccátis vestris, perdúcat vos ad vitam ætérnam. \R\ Amen.

\lettrine[lines=2]{I}{NDULGENTIAM}, absolutiónem, et remissiónem peccatórum nostrórum tríbuat nobis omnípotens et miséricors Dóminus. \R\ Amen.

\rubric{Extra Chorum, si unus vel duo tantum recitent Officium, et in Choro Monialium, semel tantum ac simul ab omnibus ita fit Confessio:}

\lettrine[lines=2]{C}{ONFITEOR} Deo omnipoténti, beátæ Maríæ semper Vírgini, beáto Michaéli Archángelo, beáto Joánni Baptístæ, sanctis Apóstolis Petro et Paulo, et ómnibus Sanctis, quia peccávi nimis cogitatióne, verbo et ópere: mea culpa, mea culpa, mea máxima culpa. Ideo precor beátam Maríam semper Vírginem, beátum Michaélem Archángelum, beátum Joánnem Baptístam, sanctos Apóstolos Petrum et Paulum, et omnes Sanctos, oráre pro me ad Dóminum, Deum nostrum.

\rubric{Deinde dicitur:}

\lettrine[lines=2]{M}{ISEREATUR} nostri omnípotens Deus, et, dimíssis peccátis nostris, perdúcat nos ad vitam ætérnam. Amen.

\lettrine[lines=2]{I}{NDULGENTIAM}, absolutiónem, et remissiónem peccatórum nostrórum tríbuat nobis omnípotens et miséricors Dóminus. Amen.

\rubric{Facta Absolutione, dicitur:}

\V\ Convérte nos, Deus, salutáris noster. \R\ Et avérte iram tuam a nobis. \V\ Deus, in adjutórium.

\rubric{Deinde, inchoata convenienti Antiphona, dicuntur tres Psalmi, prouti Officium occurrens requirit.}

\rubric{Repetita integre post ultimum Psalmum Antiphona, dicitur sequens}

\hymnusheading{Hymnus}

\lettrine[lines=2]{T}{E} lucis ante términum,\\
Rerum Creátor, póscimus,\\
Ut pro tua cleméntia\\
Sis præsul et custódia.

Procul recédant sómnia,\\
Et nóctium phantásmata;\\
Hostémque nostrum cómprime,\\
Ne polluántur córpora.

Præsta, Pater piíssime,\\
Patríque compar Unice,\\
Cum Spíritu Paráclito\\
Regnans per omne sæculum. Amen.

\rubric{Absoluto Hymno, dicitur sequens Capitulum et Responsorium breve:}

\capitulumlabel{Capitulum}{Jer. 14, 9}

\lettrine[lines=2]{T}{U} autem in nobis es, Dómine, et nomen sanctum tuum invocátum est super nos: ne derelínquas nos, Dómine, Deus noster.

\R\ br. In manus tuas, Dómine, * Comméndo spíritum meum. In manus tuas, Dómine, comméndo spíritum meum. \V\ Redemísti nos, Dómine, Deus veritátis. Comméndo spíritum meum. Glória Patri, et Fílio, et Spirítui Sancto. In manus tuas, Dómine, comméndo spíritum meum.

\V\ Custódi nos, Dómine, ut pupíllam óculi. \R\ Sub umbra alárum tuárum prótege nos.

\rubric{Postea ad Canticum dicitur sequens}

\antlabel{Ant.} Salva nos.

\canticumheading{Canticum Simeonis}{Luc. 2, 29-32}

\lettrine[lines=2]{N}{UNC} dimíttis servum tuum, Dómine, * secúndum verbum tuum in pace:

Quia vidérunt óculi mei * salutáre tuum.

Quod parásti * ante fáciem ómnium populórum,

Lumen ad revelatiónem géntium, * et glóriam plebis tuæ Israël.

Glória Patri.

\antlabel{Ant.} Salva nos, Dómine, vigilántes, custódi nos dormiéntes; ut vigilémus cum Christo, et requiescámus in pace.

\rubric{Deinde, si Preces non fuerint recitandæ, statim dicitur Versus Dóminus vobíscum cum subsequenti Oratione et reliquis omnibus usque ad finem Completorii, ut infra.}

\rubric{In quolibet Officio semiduplici, in Officio sanctæ Mariæ in Sabbato, in Festis simplicibus, atque in Feriis per Annum communibus, post repetitam ad Nunc dimittis Antiphonam dicuntur sequentes Preces, quæ tamen omittuntur omnibus diebus, in quibus ad Vesperas habita sit Commemoratio Duplicis, vel factum fuerit Officium aut Commemoratio cujusvis Octavæ.}

KYRIE, eléison. Christe, eléison. Kýrie, eléison.

Pater noster \rubric{secreto usque ad} \V\ Et ne nos indúcas in tentatiónem. \R\ Sed líbera nos a malo.

Credo in Deum \rubric{secreto usque ad} \V\ Carnis resurrectiónem. \R\ Vitam ætérnam. Amen.

\V\ Benedíctus es, Dómine, Deus patrum nostrórum. \R\ Et laudábilis et gloriósus in sæcula. \V\ Benedicámus Patrem et Fílium cum Sancto Spíritu. \R\ Laudémus, et superexaltémus eum in sæcula. \V\ Benedíctus es, Dómine, in firmaménto cæli. \R\ Et laudábilis, et gloriósus, et superexaltátus in sæcula. \V\ Benedícat et custódiat nos omnípotens et miséricors Dóminus. \R\ Amen. \V\ Dignáre, Dómine, nocte ista. \R\ Sine peccáto nos custodíre. \V\ Miserére nostri, Dómine. \R\ Miserére nostri. \V\ Fiat misericórdia tua, Dómine, super nos. \R\ Quemádmodum sperávimus in te. \V\ Dómine, exáudi oratiónem meam. \R\ Et clamor meus ad te véniat.

\rubric{Absolutis Precibus, vel, si Preces locum non habeant, repetita Antiphona post Nunc dimittis, statim dicitur:}

\V\ Dóminus vobíscum. \R\ Et cum spíritu tuo.

Orémus. \oratiolabel{Oratio}

\lettrine[lines=2]{V}{ISITA}, quæsumus, Dómine, habitatiónem istam, et omnes insídias inimíci ab ea longe repélle: Angeli tui sancti hábitent in ea, qui nos in pace custódiant; et benedíctio tua sit super nos semper. Per Dóminum nostrum. \R\ Amen.

\rubric{Post Orationem additur:}

\V\ Dóminus vobíscum. \R\ Et cum spíritu tuo.

\V\ Benedicámus Dómino. \R\ Deo grátias.

\benedictiolabel{Benedictio} Benedícat et custódiat nos omnípotens et miséricors Dóminus, Pater, et Fílius, et Spíritus Sanctus. \R\ Amen.

\rubric{Et non dicitur Versus Fidelium animæ, sed immediate subjungitur finalis beatæ Mariæ Virginis Antiphona, quæ infra assignatur.}

\rubric{Postea subjungitur:}

\V\ Divínum auxílium máneat semper nobíscum. \R\ Amen.

\rubric{Deinde dicitur secreto Pater noster, Ave María et Credo.}

\rubric{Si post Completorium immediate sequatur Matutinum, tunc bis dicitur Pater noster, Ave María et Credo, pro fine videlicet Officii diei præcedentis, et pro initio Officii diei sequentis.}

\sectionheading{Post singulas Horas}

\rubric{In fine Laudum, vel, si post Laudes immediate subsequatur alia Hora, in fine ultimæ Horæ, sicuti et semper in fine Completorii, quamvis immediate sequatur Matutinum diei sequentis, tam in publica quam in privata divini Officii recitatione; atque in publica insuper recitatione post alias quascumque Horas, si tunc terminetur Officium et discedatur a Choro, dicitur sequens Antiphona finalis beatæ Mariæ Virginis.}

\rubric{¶ Post Laudes ac ceteras Horas, præter Completorium, dicto Versu Fidelium animæ et Oratione Dominica secreto recitata, subjungitur:}

\V\ Dóminus det nobis suam pacem. \R\ Et vitam ætérnam. Amen.

\rubric{Atque statim dicitur Antiphona finalis, ut infra.}

\sectionheading{Antiphona finalis beatæ Mariæ Virginis}

\rubric{A I Vesperis Festi sanctissimæ Trinitatis et deinceps:}

\antiphonaheading{Antiphona}

\lettrine[lines=2]{S}{ALVE}, Regína, mater misericórdiæ; vita, dulcédo et spes nostra, salve. Ad te clamámus éxsules fílii Hevæ. Ad te suspirámus geméntes et flentes in hac lacrimárum valle. Eja ergo, advocáta nostra, illos tuos misericórdes óculos ad nos convérte. Et Jesum, benedíctum fructum ventris tui, nobis post hoc exsílium osténde. O clemens, o pia, o dulcis Virgo María.

\V\ Ora pro nobis, sancta Dei Génitrix. \R\ Ut digni efficiámur promissiónibus Christi.

Orémus. \oratiolabel{Oratio}

\lettrine[lines=2]{O}{MNIPOTENS} sempitérne Deus, qui gloriósæ Vírginis Matris Maríæ corpus et ánimam, ut dignum Fílii tui habitáculum éffici mererétur, Spíritu Sancto cooperánte, præparásti: da, ut, cujus commemoratióne lætámur; ejus pia intercessióne, ab instántibus malis, et a morte perpétua liberémur. Per eúmdem Christum, Dóminum nostrum. \R\ Amen.

\rubric{Absoluta vero Antiphona finali cum suis Versu et Oratione, concluditur:}

\V\ Divínum auxílium máneat semper nobíscum. \R\ Amen.

\sectionheading{Post divinum Officium}

\rubric{Cum expletum fuerit Officium, laudabiliter dicitur sequens Oratio; pro qua Summus Pontifex Pius Papa XI sub die 1 Decembris 1933 indulgentiam trium annorum concessit. Præterea Clericis in sacris constitutis et Sacerdotibus eam devote post Officium recitantibus Summus Pontifex Leo X defectus et culpas in eo persolvendo ex humana fragilitate contractas, indulsit. Dicitur autem flexis semper genibus in privata etiam recitatione, præter quam ab iis, qui ob certam infirmitatem vel gravioris impedimenti causam nequeant genuflectere.}

\oratioheading{Oratio}

\lettrine[lines=2]{S}{ACROSANCTÆ} et indivíduæ Trinitáti, crucifíxi Dómini nostri Jesu Christi humanitáti, beatíssimæ et gloriosíssimæ sempérque Vírginis Maríæ fecúndæ integritáti, et ómnium Sanctórum universitáti sit sempitérna laus, honor, virtus et glória ab omni creatúra, nobísque remíssio ómnium peccatórum, per infiníta sæcula sæculórum. \R\ Amen.

\V\ Beáta víscera Maríæ Vírginis, quæ portavérunt ætérni Patris Fílium. \R\ Et beáta úbera, quæ lactavérunt Christum Dóminum.

\rubric{Et dicitur secreto} Pater noster et Ave María.

% [ornamental image ending Ordinarium Divini Officii]

% ===== END OF BATCH 3 -- Covers autumnalis.md lines 922-1140 (reused
% verbatim from aestiva.tex's Batch 3, per the byte-identity noted
% above), completing the Ordinarium Divini Officii in its entirety:
% the close of "Ad Vesperas" (the Suffragium de omnibus Sanctis in
% both its Marian and Josephine forms, the second occurrence ending
% without a printed "℟. Amen," transcribed literally), "Ad
% Completorium" (the Lectio brevis, the full Confiteor/Misereatur/
% Indulgentiam apparatus, the hymn "Te lucis ante terminum," the Nunc
% dimittis with its proper brief responsory "In manus tuas," and the
% closing "Benedicat et custodiat nos" under
% \benedictiolabel{Benedictio}, matching the established distinction
% from the abbreviated \benedlabel{Bened.} used earlier at the same
% Hour's opening blessing), "Post singulas Horas," and the single
% "Antiphona finalis beatæ Mariæ Virginis" -- "Salve, Regína" alone,
% assigned from First Vespers of Trinity Sunday onward with no other
% seasonal Marian antiphon given in this shared Ordinarium text --
% and finally "Post divinum Officium" with its closing Oratio
% "Sacrosanctæ et individuæ Trinitati." This completes the Ordinarium
% Divini Officii for the Pars Autumnalis, matching Pars Aestiva's
% exactly. The document continues with "Psalterium Breviarii Romani"
% (autumnalis.md line 1142) =====
\partheading{Psalterium Breviarii Romani}
\partsubtitle{Per omnes ac singulos hebdomadæ dies dispositum}

\feriaheading{Dominica}

\hourheading{Ad Matutinum}

Pater, Ave et Credo.

\V\ Dómine, lábia mea.

\V\ Deus, in adjutórium.

Glória Patri. Sicut erat. Allelúja.

\rubric{In Dominica IV et reliquis Dominicis post Pentecosten occurrentibus dicitur:}

\invitatlabel{Invitat.} Dóminum, qui fecit nos, * Veníte, adorémus. \psalmlabel{Ps. 94} Veníte.

\hymnusheading{Hymnus}

\lettrine[lines=2]{N}{OCTE} surgéntes vigilémus omnes,\\
Semper in psalmis meditémur, atque\\
Voce concórdi Dómino canámus\\
Dúlciter hymnos.

Ut pio Regi páriter canéntes,\\
Cum suis Sanctis mereámur aulam\\
Ingredi cæli, simul et perénnem\\
Dúcere vitam.

Præstet hoc nobis Déitas beáta\\
Patris, ac Nati, paritérque Sancti\\
Spíritus, cujus résonat per omnem\\
Glória mundum. Amen.

\nocturnoheading{In I Nocturno}

\antlabel{Ant.} Beátus vir. \textdagger

\rubric{¶ Et non repetitur in Psalmo; quod semper fit, quando Antiphona incipit a primis Psalmi vel Cantici verbis. Et Psalmus vel Canticum tunc incipitur ab eo verbo, ante quod vel inchoata vel integra desinit Antiphona; si tamen eadem sint verba, et continuetur Antiphona cum Psalmo vel Cantico, neque additur in fine ipsius Antiphonæ Allelúja.}

\rubric{Antiphonæ autem cum Psalmis et Versibus assignatæ, adhibentur semper, extra Octavas Ssmi Corporis Christi et Sacrat. Cordis Jesu, quando fit Officium de Dominica.}

\rubric{¶ Asteriscus *, in Psalmorum versibus et in aliis Officii partibus insertus, denotat pausam cantus et recitationis in Choro; præter quam in Antiphonis, in quibus nempe significat tum præintonationem Antiphonæ, tum Antiphonam ipsam ante Psalmum seu Canticum eousque producendam esse ad omnes et singulas Horas in Officio ritus semiduplicis ac simplicis, atque ad Horas tantum minores in Officio ritus duplicis.}

\psalmheading{Psalmus 1}

\lettrine[lines=2]{B}{EATUS} vir, \textdagger\ qui non ábiit in consílio impiórum, et in via peccatórum non stetit, * et in cáthedra pestiléntiæ non sedit:

Sed in lege Dómini volúntas ejus, * et in lege ejus meditábitur die ac nocte.

Et erit tamquam lignum, quod plantátum est secus decúrsus aquárum, * quod fructum suum dabit in témpore suo:

Et fólium ejus non défluet: * et ómnia quæcúmque fáciet, prosperabúntur.

Non sic ímpii, non sic: * sed tamquam pulvis, quem prójicit ventus a fácie terræ.

Ideo non resúrgent ímpii in judício: * neque peccatóres in concílio justórum.

Quóniam novit Dóminus viam justórum: * et iter impiórum períbit.

Glória Patri, et Fílio, * et Spirítui Sancto.

Sicut erat in princípio, et nunc, et semper, * et in sæcula sæculórum. Amen.

\rubric{¶ Hic versus Glória Patri semper dicitur in fine omnium Psalmorum et Canticorum, nisi aliter notetur.}

\antlabel{Ant.} Beátus vir, qui in lege Dómini meditátur.

\antlabel{Ant.} Servíte Dómino.

\psalmheading{Psalmus 2}

\lettrine[lines=2]{Q}{UARE} fremuérunt gentes, * et pópuli meditáti sunt inánia?

Astitérunt reges terræ, et príncipes convenérunt in unum * advérsus Dóminum, et advérsus Christum ejus.

Dirumpámus víncula eórum: * et projiciámus a nobis jugum ipsórum.

Qui hábitat in cælis, irridébit eos: * et Dóminus subsannábit eos.

Tunc loquétur ad eos in ira sua, * et in furóre suo conturbábit eos.

Ego autem constitútus sum Rex ab eo super Sion montem sanctum ejus, * prædicans præcéptum ejus.

Dóminus dixit ad me: * Fílius meus es tu, ego hódie génui te.

Póstula a me, et dabo tibi gentes hereditátem tuam, * et possessiónem tuam términos terræ.

Reges eos in virga férrea, * et tamquam vas fíguli confrínges eos.

Et nunc, reges, intellígite: * erudímini, qui judicátis terram.

Servíte Dómino in timóre: * et exsultáte ei cum tremóre.

Apprehéndite disciplínam, nequándo irascátur Dóminus, * et pereátis de via justa.

Cum exárserit in brevi ira ejus, * beáti omnes qui confídunt in eo.

\antlabel{Ant.} Servíte Dómino in timóre: et exsultáte ei cum tremóre.

\antlabel{Ant.} Exsúrge.

\psalmheading{Psalmus 3}

\lettrine[lines=2]{D}{OMINE}, quid multiplicáti sunt qui tríbulant me? * multi insúrgunt advérsum me.

Multi dicunt ánimæ meæ: * Non est salus ipsi in Deo ejus.

Tu autem, Dómine, suscéptor meus es, * glória mea, et exáltans caput meum.

Voce mea ad Dóminum clamávi: * et exaudívit me de monte sancto suo.

Ego dormívi, et soporátus sum: * et exsurréxi, quia Dóminus suscépit me.

Non timébo míllia pópuli circumdántis me: * exsúrge, Dómine, salvum me fac, Deus meus.

Quóniam tu percussísti omnes adversántes mihi sine causa: * dentes peccatórum contrivísti.

Dómini est salus: * et super pópulum tuum benedíctio tua.

\antlabel{Ant.} Exsúrge, Dómine, salvum me fac, Deus meus.

\V\ Memor fui nocte nóminis tui, Dómine. \R\ Et custodívi legem tuam.

\absolutiolabel{Absol.} Exáudi, Dómine.

\benedlabel{Bened.} 1 Benedictióne perpétua. 2 Unigénitus Dei Fílius. 3 Spíritus Sancti grátia.

\nocturnoheading{In II Nocturno}

\antlabel{Ant.} Quam admirábile.

\psalmheading{Psalmus 8}

\lettrine[lines=2]{D}{OMINE}, Dóminus noster, * quam admirábile est nomen tuum in univérsa terra!

Quóniam eleváta est magnificéntia tua, * super cælos.

Ex ore infántium et lacténtium perfecísti laudem propter inimícos tuos, * ut déstruas inimícum et ultórem.

Quóniam vidébo cælos tuos, ópera digitórum tuórum: * lunam et stellas, quæ tu fundásti.

Quid est homo, quod memor es ejus? * aut fílius hóminis, quóniam vísitas eum?

Minuísti eum paulo minus ab Angelis, glória et honóre coronásti eum: * et constituísti eum super ópera mánuum tuárum.

Omnia subjecísti sub pédibus ejus, * oves et boves univérsas: ínsuper et pécora campi.

Vólucres cæli, et pisces maris, * qui perámbulant sémitas maris.

Dómine, Dóminus noster, * quam admirábile est nomen tuum in univérsa terra!

\antlabel{Ant.} Quam admirábile est nomen tuum, Dómine, in univérsa terra!

\antlabel{Ant.} Sedísti super thronum.

\psalmheading{Psalmus 9, i}

\lettrine[lines=2]{C}{ONFITEBOR} tibi, Dómine, in toto corde meo: * narrábo ómnia mirabília tua.

Lætábor et exsultábo in te: * psallam nómini tuo, Altíssime,

In converténdo inimícum meum retrórsum: * infirmabúntur, et períbunt a fácie tua.

Quóniam fecísti judícium meum et causam meam: * sedísti super thronum, qui júdicas justítiam.

Increpásti gentes, et périit ímpius: * nomen eórum delésti in ætérnum, et in sæculum sæculi.

Inimíci defecérunt frámeæ in finem: * et civitátes eórum destruxísti.

Périit memória eórum cum sónitu: * et Dóminus in ætérnum pérmanet.

Parávit in judício thronum suum: * et ipse judicábit orbem terræ in æquitáte, judicábit pópulos in justítia.

Et factus est Dóminus refúgium páuperi: * adjútor in opportunitátibus, in tribulatióne.

Et sperent in te qui novérunt nomen tuum: * quóniam non dereliquísti quæréntes te, Dómine.

\antlabel{Ant.} Sedísti super thronum, qui júdicas justítiam.

\antlabel{Ant.} Exsúrge, Dómine.

\psalmheading{Psalmus 9, ii}

\lettrine[lines=2]{P}{SALLITE} Dómino, qui hábitat in Sion: * annuntiáte inter gentes stúdia ejus:

Quóniam requírens sánguinem eórum recordátus est: * non est oblítus clamórem páuperum.

Miserére mei, Dómine: * vide humilitátem meam de inimícis meis.

Qui exáltas me de portis mortis, * ut annúntiem omnes laudatiónes tuas in portis fíliæ Sion.

Exsultábo in salutári tuo: * infíxæ sunt gentes in intéritu, quem fecérunt.

In láqueo isto, quem abscondérunt, * comprehénsus est pes eórum.

Cognoscétur Dóminus judícia fáciens: * in opéribus mánuum suárum comprehénsus est peccátor.

Convertántur peccatóres in inférnum, * omnes gentes quæ obliviscúntur Deum.

Quóniam non in finem oblívio erit páuperis: * patiéntia páuperum non períbit in finem.

Exsúrge, Dómine, non confortétur homo: * judicéntur gentes in conspéctu tuo.

Constítue, Dómine, legislatórem super eos: * ut sciant gentes quóniam hómines sunt.

\antlabel{Ant.} Exsúrge, Dómine, non prævaleat homo.

\V\ Média nocte surgébam ad confiténdum tibi. \R\ Super judícia justificatiónis tuæ.

\absolutiolabel{Absol.} Ipsíus píetas et.

\benedlabel{Bened.} 4 Deus Pater omnípotens. 5 Christus perpétuæ. 6 Ignem sui amóris.

\nocturnoheading{In III Nocturno}

\antlabel{Ant.} Ut quid, Dómine. \textdagger

\psalmheading{Psalmus 9, iii}

\lettrine[lines=2]{U}{T} quid, Dómine, \textdagger\ recessísti longe, * déspicis in opportunitátibus, in tribulatióne?

Dum supérbit ímpius, incénditur pauper: * comprehendúntur in consíliis quibus cógitant.

Quóniam laudátur peccátor in desidériis ánimæ suæ: * et iníquus benedícitur.

Exacerbávit Dóminum peccátor, * secúndum multitúdinem iræ suæ non quæret.

Non est Deus in conspéctu ejus: * inquinátæ sunt viæ illíus in omni témpore.

Auferúntur judícia tua a fácie ejus: * ómnium inimicórum suórum dominábitur.

Dixit enim in corde suo: * Non movébor a generatióne in generatiónem sine malo.

Cujus maledictióne os plenum est, et amaritúdine, et dolo: * sub lingua ejus labor et dolor.

Sedet in insídiis cum divítibus in occúltis, * ut interfíciat innocéntem.

Oculi ejus in páuperem respíciunt: * insidiátur in abscóndito, quasi leo in spelúnca sua.

Insidiátur ut rápiat páuperem: * rápere páuperem, dum áttrahit eum.

In láqueo suo humiliábit eum: * inclinábit se, et cadet, cum dominátus fúerit páuperum.

Dixit enim in corde suo: * Oblítus est Deus, avértit fáciem suam ne vídeat in finem.

\antlabel{Ant.} Ut quid, Dómine, recessísti longe?

\antlabel{Ant.} Exsúrge. \textdagger

\psalmheading{Psalmus 9, iv}

\lettrine[lines=2]{E}{XSURGE}, \textdagger\ Dómine Deus, exaltétur manus tua: * ne obliviscáris páuperum.

Propter quid irritávit ímpius Deum? * dixit enim in corde suo: Non requíret.

Vides quóniam tu labórem et dolórem consíderas: * ut tradas eos in manus tuas.

Tibi derelíctus est pauper: * órphano tu eris adjútor.

Cóntere bráchium peccatóris et malígni: * quærétur peccátum illíus, et non inveniétur.

Dóminus regnábit in ætérnum, et in sæculum sæculi: * períbitis, gentes, de terra illíus.

Desidérium páuperum exaudívit Dóminus: * præparatiónem cordis eórum audívit auris tua.

Judicáre pupíllo et húmili, * ut non appónat ultra magnificáre se homo super terram.

\antlabel{Ant.} Exsúrge, Dómine Deus, exaltétur manus tua.

\antlabel{Ant.} Justus Dóminus.

\psalmheading{Psalmus 10}

\lettrine[lines=2]{I}{N} Dómino confído: quómodo dícitis ánimæ meæ: * Tránsmigra in montem sicut passer?

Quóniam ecce peccatóres intendérunt arcum, paravérunt sagíttas suas in pháretra, * ut sagíttent in obscúro rectos corde.

Quóniam quæ perfecísti, destruxérunt: * justus autem quid fecit?

Dóminus in templo sancto suo, * Dóminus in cælo sedes ejus:

Oculi ejus in páuperem respíciunt: * pálpebræ ejus intérrogant fílios hóminum.

Dóminus intérrogat justum et ímpium: * qui autem díligit iniquitátem, odit ánimam suam.

Pluet super peccatóres láqueos: * ignis, et sulphur, et spíritus procellárum pars cálicis eórum.

Quóniam justus Dóminus, et justítias diléxit: * æquitátem vidit vultus ejus.

\antlabel{Ant.} Justus Dóminus, et justítiam diléxit.

\V\ Prævenérunt óculi mei ad te dilúculo. \R\ Ut meditárer elóquia tua, Dómine.

\absolutiolabel{Absolutio} A vínculis.

\benedlabel{Bened.} 7 Evangélica léctio. 8 Divínum auxílium. 9 Ad societátem cívium.

% ===== END OF BATCH 4 -- Covers autumnalis.md lines 1142-1442, opening
% "Psalterium Breviarii Romani" with "Dominica" / "Ad Matutinum" in
% its ordinary-time form ("In Dominica IV et reliquis Dominicis post
% Pentecosten occurrentibus"): the Invitatorium "Dóminum, qui fecit
% nos," the hymn "Nocte surgentes," and the three Nocturns (Psalms 1,
% 2, 3; 8, 9/i, 9/ii; 9/iii, 9/iv, 10), each with its Absolutio and
% triple Benedictiones. Neither this Invitatorium nor this hymn
% appears anywhere in hiemalis.tex or verna.tex, confirming this is
% genuinely new ordinary-time-after-Pentecost material specific to
% the summer and autumn volumes. A structurally significant
% observation: after Psalm 1 prints "Glória Patri" in full together
% with the rubric establishing that this doxology is "semper dicitur
% in fine omnium Psalmorum et Canticorum, nisi aliter notetur," EVERY
% SUBSEQUENT PSALM IN THIS NOCTURNAL SET (2, 3, 8, 9/i, 9/ii, 9/iii,
% 9/iv, 10) OMITS THE PRINTED "Glória Patri" LINE ENTIRELY -- the
% source trusts the reader to supply it from the standing rubric
% rather than reprinting it each time. This transcription follows
% the source exactly: no "Glória Patri" line is added after any of
% these psalms, since none appears in autumnalis.md itself; only where
% the source itself prints it (as after Psalm 1) does the LaTeX
% include it. The embedded dagger marks (e.g. "Beátus vir, †",
% "Ut quid, Dómine, †", "Exsúrge, †") are rendered \textdagger\
% exactly where the source places them, including inside an
% antiphon incipit alone ("Ant. Ut quid, Dómine. †"). The second
% Nocturn's Absolutio is given as the abbreviated "Absol." but the
% third Nocturn's uses the unabbreviated "Absolutio" -- both render
% via the same flexible \absolutiolabel{} command, consistent with
% established practice. The document continues with "Ad Laudes"
% (autumnalis.md line 1443) =====
Sicut ádipe et pinguédine repleátur ánima mea: * et lábiis exsultatiónis laudábit os meum.

Si memor fui tui super stratum meum, in matutínis meditábor in te: * quia fuísti adjútor meus.

Et in velaménto alárum tuárum exsultábo, adhǽsit ánima mea post te: * me suscépit déxtera tua.

Ipsi vero in vanum quæsiérunt ánimam meam, introíbunt in inferióra terræ: * tradéntur in manus gládii, partes vúlpium erunt.

Rex vero lætábitur in Deo, laudabúntur omnes qui jurant in eo: * quia obstrúctum est os loquéntium iníqua.

\antlabel{Ant.} Benedícam te in vita mea, Dómine: et in nómine tuo levábo manus meas, allelúja.

\antlabel{Ant.} Tres púeri.

\canticumheading{Canticum trium Puerorum [4]}{Dan. 57-88 et 56}

\lettrine[lines=2]{B}{ENEDICITE}, ómnia ópera Dómini, Dómino: * laudáte et superexaltáte eum in sæcula.

Benedícite, Angeli Dómini, Dómino: * benedícite, cæli, Dómino.

Benedícite, aquæ omnes, quæ super cælos sunt, Dómino: * benedícite, omnes virtútes Dómini, Dómino.

Benedícite, sol et luna, Dómino: * benedícite, stellæ cæli, Dómino.

Benedícite, omnis imber et ros, Dómino: * benedícite, omnes spíritus Dei, Dómino.

Benedícite, ignis et æstus, Dómino: * benedícite, frigus et æstus, Dómino.

Benedícite, rores et pruína, Dómino: * benedícite, gelu et frigus, Dómino.

Benedícite, glácies et nives, Dómino: * benedícite, noctes et dies, Dómino.

Benedícite, lux et ténebræ, Dómino: * benedícite, fúlgura et nubes, Dómino.

Benedícat terra Dóminum: * laudet et superexáltet eum in sæcula.

Benedícite, montes et colles, Dómino: * benedícite, univérsa germinántia in terra, Dómino.

Benedícite, fontes, Dómino: * benedícite, mária et flúmina, Dómino.

Benedícite, cete, et ómnia, quæ movéntur in aquis, Dómino: * benedícite, omnes vólucres cæli, Dómino.

Benedícite, omnes béstiæ et pécora, Dómino: * benedícite, fílii hóminum, Dómino.

Benedícat Israël Dóminum: * laudet et superexáltet eum in sæcula.

Benedícite, sacerdótes Dómini, Dómino: * benedícite, servi Dómini, Dómino.

Benedícite, spíritus, et ánimæ justórum, Dómino: * benedícite, sancti, et húmiles corde, Dómino.

Benedícite, Ananía, Azaría, Mísaël, Dómino: * laudáte et superexaltáte eum in sæcula.

Benedicámus Patrem et Fílium cum Sancto Spíritu: * laudémus et superexaltémus eum in sæcula.

Benedíctus es, Dómine, in firmaménto cæli: * et laudábilis, et gloriósus, et superexaltátus in sæcula.

\rubric{¶ Hic non dicitur} Glória Patri, \rubric{neque} Amen.

\antlabel{Ant.} Tres púeri jussu regis in fornácem missi sunt, non timéntes flammam ignis, dicéntes: Benedíctus Deus, allelúja.

\antlabel{Ant.} Allelúja.

\psalmheading{Psalmus 148 [5]}

\lettrine[lines=2]{L}{AUDATE} Dóminum de cælis: * laudáte eum in excélsis.

Laudáte eum, omnes Angeli ejus: * laudáte eum, omnes virtútes ejus.

Laudáte eum, sol et luna: * laudáte eum, omnes stellæ et lumen.

Laudáte eum, cæli cælórum: * et aquæ omnes, quæ super cælos sunt, laudent nomen Dómini.

Quia ipse dixit, et facta sunt: * ipse mandávit, et creáta sunt.

Státuit ea in ætérnum, et in sæculum sæculi: * præcéptum pósuit, et non præteríbit.

Laudáte Dóminum de terra, * dracónes, et omnes abýssi.

Ignis, grando, nix, glácies, spíritus procellárum: * quæ fáciunt verbum ejus:

Montes, et omnes colles: * ligna fructífera, et omnes cedri.

Béstiæ, et univérsa pécora: * serpéntes, et vólucres pennátæ:

Reges terræ, et omnes pópuli: * príncipes, et omnes júdices terræ.

Júvenes, et vírgines: senes cum junióribus laudent nomen Dómini: * quia exaltátum est nomen ejus solíus.

Conféssio ejus super cælum et terram: * et exaltávit cornu pópuli sui.

Hymnus ómnibus sanctis ejus: * fíliis Israël, pópulo appropinquánti sibi.

\antlabel{Ant.} Allelúja, laudáte Dóminum de cælis, allelúja, allelúja.

\rubric{In omnibus Dominicis extra Octavas Ssmi Corporis Christi ac Sacratissimi Cordis Jesu, dicitur:}

\capitulumlabel{Capitulum}{Apoc. 7, 12}

\lettrine[lines=2]{B}{ENEDICTIO}, et cláritas, et sapiéntia, et gratiárum áctio, honor, virtus, et fortitúdo Deo nostro in sæcula sæculórum. Amen. \R\ Deo grátias.

\rubric{In Dominica IV et reliquis Dominicis post Pentecosten occurrentibus, dicitur}

\hymnusheading{Hymnus}

\lettrine[lines=2]{E}{CCE}, jam noctis tenuátur umbra,\\
Lux et auróræ rútilans corúscat:\\
Súpplices rerum Dóminum canóra\\
Voce precémur:

Ut reos culpæ miserátus omnem\\
Pellat angórem, tríbuat salútem,\\
Donet et nobis bona sempitérnæ\\
Múnera pacis.

Præstet hoc nobis Déitas beáta\\
Patris, ac Nati, paritérque Sancti\\
Spíritus, cujus résonat per omnem\\
Glória mundum. Amen.

\V\ Dóminus regnávit, decórem índuit. \R\ Induit Dóminus fortitúdinem, et præcínxit se virtúte.

\rubric{Ad Benedictus Antiphona ut in Proprio de Tempore.}

\canticumheading{Canticum Zachariæ}{Luc. 1, 68-79}

\lettrine[lines=2]{B}{ENEDICTUS} \textdagger\ Dóminus, Deus Israël, * quia visitávit, et fecit redemptiónem plebis suæ:

Et eréxit cornu salútis nobis: * in domo David, púeri sui.

Sicut locútus est per os sanctórum, * qui a sæculo sunt, prophetárum ejus:

Salútem ex inimícis nostris, * et de manu ómnium, qui odérunt nos:

Ad faciéndam misericórdiam cum pátribus nostris: * et memorári testaménti sui sancti.

Jusjurándum, quod jurávit ad Abraham, patrem nostrum, * datúrum se nobis:

Ut sine timóre, de manu inimicórum nostrórum liberáti, * serviámus illi.

In sanctitáte, et justítia coram ipso, * ómnibus diébus nostris.

Et tu, puer, Prophéta Altíssimi vocáberis: * præíbis enim ante fáciem Dómini paráre vias ejus:

Ad dandam sciéntiam salútis plebi ejus: * in remissiónem peccatórum eórum:

Per víscera misericórdiæ Dei nostri: * in quibus visitávit nos, óriens ex alto:

Illumináre his, qui in ténebris, et in umbra mortis sedent: * ad dirigéndos pedes nostros in viam pacis.

\rubric{Oratio conveniens. Reliqua ut in Ordinario. 11}

\rubric{¶ Psalmi cum Cantico de Dominica superius appositi dicuntur etiam, sub Antiphonis Officio diei congruentibus, in Dominicis infra Octavas Ssmi Corporis Christi ac Sacrat. Cordis Jesu, in Duplicibus I vel II classis, in Festis novem Lectionum Domini, B. Mariæ Virginis, Angelorum, S. Joannis Baptistæ, S. Joseph, Apostolorum, Evangelistarum et per Octavas privilegiatas; itemque in quolibet Festo novem Lectionum proprias Antiphonas ad Laudes habenti, quando de eis persolvitur Officium.}

\hourheading{Ad Primam}

Pater, Ave et Credo.

\V\ Deus, in adjutórium.

\hymnusheading{Hymnus}

\lettrine[lines=2]{J}{AM} lucis orto sídere,\\
Deum precémur súpplices,\\
Ut in diúrnis áctibus\\
Nos servet a nocéntibus.

Linguam refrénans témperet,\\
Ne litis horror ínsonet:\\
Visum fovéndo cóntegat,\\
Ne vanitátes háuriat.

Sint pura cordis íntima,\\
Absístat et vecórdia:\\
Carnis terat supérbiam\\
Potus cibíque párcitas.

Ut, cum dies abscésserit\\
Noctémque sors redúxerit,\\
Mundi per abstinéntiam\\
Ipsi canámus glóriam.

Deo Patri sit glória,\\
Ejúsque soli Fílio,\\
Cum Spíritu Paráclito,\\
Nunc et per omne sæculum. Amen.

\antlabel{Ant.} Allelúja.

\rubric{¶ Aña cum Psalmis, ad Horas minores assignata, adhibetur semper, extra Octavas Ssmi Corporis Christi ac Sacrat. Cordis Jesu, quando fit Officium de qualibet Dominica minori.}

\psalmheading{Psalmus 117}

\lettrine[lines=2]{C}{ONFITEMINI} Dómino quóniam bonus: * quóniam in sæculum misericórdia ejus.

Dicat nunc Israël quóniam bonus: * quóniam in sæculum misericórdia ejus.

Dicat nunc domus Aaron: * quóniam in sæculum misericórdia ejus.

Dicant nunc qui timent Dóminum: * quóniam in sæculum misericórdia ejus.

De tribulatióne invocávi Dóminum: * et exaudívit me in latitúdine Dóminus.

Dóminus mihi adjútor: * non timébo quid fáciat mihi homo.

Dóminus mihi adjútor: * et ego despíciam inimícos meos.

Bonum est confídere in Dómino, * quam confídere in hómine:

Bonum est speráre in Dómino, * quam speráre in princípibus.

Omnes gentes circuiérunt me: * et in nómine Dómini quia ultus sum in eos.

Circumdántes circumdedérunt me: * et in nómine Dómini quia ultus sum in eos.

Circumdedérunt me sicut apes, et exarsérunt sicut ignis in spinis: * et in nómine Dómini quia ultus sum in eos.

Impúlsus evérsus sum ut cáderem: * et Dóminus suscépit me.

Fortitúdo mea, et laus mea Dóminus: * et factus est mihi in salútem.

Vox exsultatiónis, et salútis * in tabernáculis justórum.

Déxtera Dómini fecit virtútem: déxtera Dómini exaltávit me, * déxtera Dómini fecit virtútem.

Non móriar, sed vivam: * et narrábo ópera Dómini.

Castígans castigávit me Dóminus: * et morti non trádidit me.

Aperíte mihi portas justítiæ, ingréssus in eas confitébor Dómino: * hæc porta Dómini, justi intrábunt in eam.

Confitébor tibi quóniam exaudísti me: * et factus es mihi in salútem.

Lápidem, quem reprobavérunt ædificántes: * hic factus est in caput ánguli.

A Dómino factum est istud: * et est mirábile in óculis nostris.

Hæc est dies, quam fecit Dóminus: * exsultémus, et lætémur in ea.

O Dómine, salvum me fac, o Dómine, bene prosperáre: * benedíctus qui venit in nómine Dómini.

Benedíximus vobis de domo Dómini: * Deus Dóminus, et illúxit nobis.

Constitúite diem solémnem in condénsis, * usque ad cornu altáris.

Deus meus es tu, et confitébor tibi: * Deus meus es tu, et exaltábo te.

Confitébor tibi quóniam exaudísti me: * et factus es mihi in salútem.

Confitémini Dómino quóniam bonus: * quóniam in sæculum misericórdia ejus.

\rubric{¶ In Dominicis infra Octavas Ssmi Corporis Christi ac Sacrat. Cordis Jesu, in Duplicibus I vel II classis, in Festis novem Lectionum Domini, B. Mariæ Virginis, Angelorum, S. Joannis Baptistæ, S. Joseph, Apostolorum, Evangelistarum, et per Octavas privilegiatas: si de eis fiat Officium, ad Primam, Tertiam, Sextam et Nonam dicuntur Psalmi de Dominica, sub Antiphonis Officio diei congruentibus, nisi aliter suo loco notetur; sed ad Primam, loco Psalmi 117 Confitemini, in ipsa etiam die Dominica semper dicitur}

\psalmheading{Psalmus 53}

\lettrine[lines=2]{D}{EUS}, in nómine tuo salvum me fac: * et in virtúte tua júdica me.

Deus, exáudi oratiónem meam: * áuribus pércipe verba oris mei.

Quóniam aliéni insurrexérunt advérsum me, et fortes quæsiérunt ánimam meam: * et non proposuérunt Deum ante conspéctum suum.

Ecce enim, Deus ádjuvat me: * et Dóminus suscéptor est ánimæ meæ.

Avérte mala inimícis meis: * et in veritáte tua dispérde illos.

Voluntárie sacrificábo tibi, * et confitébor nómini tuo, Dómine: quóniam bonum est:

Quóniam ex omni tribulatióne eripuísti me: * et super inimícos meos despéxit óculus meus.

\psalmheading{Psalmus 118, i}

\lettrine[lines=2]{B}{EATI} immaculáti in via: * qui ámbulant in lege Dómini.

Beáti, qui scrutántur testimónia ejus: * in toto corde exquírunt eum.

Non enim qui operántur iniquitátem: * in viis ejus ambulavérunt.

Tu mandásti * mandáta tua custodíri nimis.

Utinam dirigántur viæ meæ, * ad custodiéndas justificatiónes tuas!

Tunc non confúndar, * cum perspéxero in ómnibus mandátis tuis.

Confitébor tibi in directióne cordis: * in eo quod dídici judícia justítiæ tuæ.

Justificatiónes tuas custódiam: * non me derelínquas usquequáque.

In quo córrigit adolescéntior viam suam? * in custodiéndo sermónes tuos.

In toto corde meo exquisívi te: * ne repéllas me a mandátis tuis.

In corde meo abscóndi elóquia tua: * ut non peccem tibi.

Benedíctus es, Dómine: * doce me justificatiónes tuas.

In lábiis meis, * pronuntiávi ómnia judícia oris tui.

In via testimoniórum tuórum delectátus sum, * sicut in ómnibus divítiis.

In mandátis tuis exercébor: * et considerábo vias tuas.

In justificatiónibus tuis meditábor: * non oblivíscar sermónes tuos.

\psalmheading{Psalmus 118, ii}

\lettrine[lines=2]{R}{ETRIBUE} servo tuo, vivífica me: * et custódiam sermónes tuos.

Revéla óculos meos: * et considerábo mirabília de lege tua.

Íncola ego sum in terra: * non abscóndas a me mandáta tua.

Concupívit ánima mea desideráre justificatiónes tuas, * in omni témpore.

Increpásti supérbos: * maledícti qui declínant a mandátis tuis.

Aufer a me oppróbrium, et contémptum: * quia testimónia tua exquisívi.

Etenim sedérunt príncipes, et advérsum me loquebántur: * servus autem tuus exercebátur in justificatiónibus tuis.

Nam et testimónia tua meditátio mea est: * et consílium meum justificatiónes tuæ.

Adhǽsit paviménto ánima mea: * vivífica me secúndum verbum tuum.

Vias meas enuntiávi, et exaudísti me: * doce me justificatiónes tuas.

Viam justificatiónum tuárum ínstrue me: * et exercébor in mirabílibus tuis.

Dormitávit ánima mea præ tædio: * confírma me in verbis tuis.

Viam iniquitátis ámove a me: * et de lege tua miserére mei.

Viam veritátis elégi: * judícia tua non sum oblítus.

Adhǽsi testimóniis tuis, Dómine: * noli me confúndere.

Viam mandatórum tuórum cucúrri, * cum dilatásti cor meum.

\rubric{¶ Si Athanasianum Symbolum omittatur, expletis Psalmis præcedentibus, statim dicitur Antiphona, ut infra. 45}

\rubric{Sequens Symbolum dicitur in Dominicis post Pentecosten, quando fit Officium de Dominica ut in Psalterio, neque ulla Commemoratio Officii duplicis vel Octavæ recurrit, et in Festo sanctissimæ Trinitatis.}

\symbolumheading{Symbolum Athanasianum}

QUICUMQUE vult salvus esse, * ante ómnia opus est, ut téneat cathólicam fidem:

Quam nisi quisque íntegram inviolatámque serváverit, * absque dúbio in ætérnum períbit.

Fides autem cathólica hæc est: * ut unum Deum in Trinitáte, et Trinitátem in unitáte venerémur.

Neque confundéntes persónas, * neque substántiam separántes.

Ália est enim persóna Patris, ália Fílii, * ália Spíritus Sancti:

Sed Patris, et Fílii, et Spíritus Sancti una est divínitas, * æquális glória, coætérna majéstas.

Quális Pater, talis Fílius, * talis Spíritus Sanctus.

Increátus Pater, increátus Fílius, * increátus Spíritus Sanctus.

Imménsus Pater, imménsus Fílius, * imménsus Spíritus Sanctus.

Ætérnus Pater, ætérnus Fílius, * ætérnus Spíritus Sanctus.

Et tamen non tres ætérni, * sed unus ætérnus.

Sicut non tres increáti, nec tres imménsi, * sed unus increátus, et unus imménsus.

Simíliter omnípotens Pater, omnípotens Fílius, * omnípotens Spíritus Sanctus.

Et tamen non tres omnipoténtes, * sed unus omnípotens.

Ita Deus Pater, Deus Fílius, * Deus Spíritus Sanctus.

Et tamen non tres Dii, * sed unus est Deus.

Ita Dóminus Pater, Dóminus Fílius, * Dóminus Spíritus Sanctus.

Et tamen non tres Dómini, * sed unus est Dóminus.

Quia, sicut singillátim unamquámque persónam Deum ac Dóminum confitéri christiána veritáte compéllimur: * ita tres Deos aut Dóminos dícere cathólica religióne prohibémur.

Pater a nullo est factus: * nec creátus, nec génitus.

Fílius a Patre solo est: * non factus, nec creátus, sed génitus.

Spíritus Sanctus a Patre et Fílio: * non factus, nec creátus, nec génitus, sed procédens.

Unus ergo Pater, non tres Patres: unus Fílius, non tres Fílii: * unus Spíritus Sanctus, non tres Spíritus Sancti.

Et in hac Trinitáte nihil prius aut postérius, nihil majus aut minus: * sed totæ tres persónæ coætérnæ sibi sunt et coæquáles.

Ita ut per ómnia, sicut jam supra dictum est, * et únitas in Trinitáte, et Trínitas in unitáte veneránda sit.

Qui vult ergo salvus esse, * ita de Trinitáte séntiat.

Sed necessárium est ad ætérnam salútem, * ut Incarnatiónem quoque Dómini nostri Jesu Christi fidéliter credat.

Est ergo fides recta, ut credámus et confiteámur, * quia Dóminus noster Jesus Christus, Dei Fílius, Deus et homo est.

Deus est ex substántia Patris ante sæcula génitus: * et homo est ex substántia matris in sæculo natus.

Perféctus Deus, perféctus homo: * ex ánima rationáli et humána carne subsístens.

Æquális Patri secúndum divinitátem: * minor Patre secúndum humanitátem.

Qui, licet Deus sit et homo, * non duo tamen, sed unus est Christus.

Unus autem non conversióne divinitátis in carnem, * sed assumptióne humanitátis in Deum.

Unus omníno, non confusióne substántiæ, * sed unitáte persónæ.

Nam sicut ánima ratiónális et caro unus est homo: * ita Deus et homo unus est Christus.

Qui passus est pro salúte nostra: descéndit ad ínferos: * tértia die resurréxit a mórtuis.

Ascéndit ad cælos, sedet ad déxteram Dei Patris omnipoténtis: * inde ventúrus est judicáre vivos et mórtuos.

Ad cujus advéntum omnes hómines resúrgere habent cum corpóribus suis: * et redditúri sunt de factis própriis ratiónem.

Et qui bona egérunt, ibunt in vitam ætérnam: * qui vero mala, in ignem ætérnum.

Hæc est fides cathólica, * quam nisi quisque fidéliter firmitérque credíderit, salvus esse non póterit.

Glória Patri.

\rubric{Expleto Symbolo, vel, si hoc non habeat locum, Psalmis absolutis, statim dicitur Antiphona.}

\antlabel{Ant.} Allelúja, confitémini Dómino quóniam in sæculum misericórdia ejus, allelúja, allelúja.

\capitulumlabel{Capitulum}{1 Tim. 1, 17}

\lettrine[lines=2]{R}{EGI} sæculórum immortáli et invisíbili, soli Deo honor et glória in sæcula sæculórum. Amen.

\R\ br. Christe, Fili Dei vivi, * Miserére nobis. \rubric{Et repetitur:} Christe, Fili Dei vivi, miserére nobis. \V\ Qui sedes ad déxteram Patris. \R\ Miserére nobis. \V\ Glória Patri, et Fílio, et Spirítui Sancto. \R\ Christe, Fili Dei vivi, miserére nobis. \V\ Exsúrge, Christe, ádjuva nos. \R\ Et líbera nos propter nomen tuum.

\variantheading{Preces (si dicendæ sint)}

KYRIE, eléison. Christe, eléison. Kýrie, eléison.

Pater noster \rubric{secreto usque ad} \V\ Et ne nos indúcas in tentatiónem. \R\ Sed líbera nos a malo.

Credo in Deum \rubric{secreto usque ad} \V\ Carnis resurrectiónem. \R\ Vitam ætérnam. Amen.

\V\ Et ego ad te, Dómine, clamávi. \R\ Et mane orátio mea prævéniet te. \V\ Repleátur os meum laude. \R\ Ut cantem glóriam tuam, tota die magnitúdinem tuam. \V\ Dómine, avérte fáciem tuam a peccátis meis. \R\ Et omnes iniquitátes meas dele. \V\ Cor mundum crea in me, Deus. \R\ Et spíritum rectum ínnova in viscéribus meis. \V\ Ne projícias me a fácie tua. \R\ Et spíritum sanctum tuum ne áuferas a me. \V\ Redde mihi lætítiam salutáris tui. \R\ Et spíritu principáli confírma me. \V\ Adjutórium nostrum in nómine Dómini. \R\ Qui fecit cælum et terram.

\formulalabel{Confíteor. Misereátur. Indulgéntiam.}

\V\ Dignáre, Dómine, die isto. \R\ Sine peccáto nos custodíre. \V\ Miserére nostri, Dómine. \R\ Miserére nostri. \V\ Fiat misericórdia tua, Dómine, super nos. \R\ Quemádmodum sperávimus in te. \V\ Dómine, exáudi oratiónem meam. \R\ Et clamor meus ad te véniat.

\V\ Dóminus vobíscum.

Orémus. \oratiolabel{Oratio}

\lettrine[lines=2]{D}{OMINE}, Deus omnípotens, qui ad princípium hujus diéi nos perveníre fecísti: tua nos hódie salva virtúte; ut in hac die ad nullum declinémus peccátum, sed semper ad tuam justítiam faciéndam nostra procédant elóquia, dirigántur cogitatiónes et ópera. Per Dóminum. \R\ Amen.

\V\ Dóminus vobíscum. \V\ Benedicámus Dómino.

\rubric{In Choro Martyrologium.}

\V\ Pretiósa in conspéctu Dómini. \R\ Mors Sanctórum ejus.

\rubric{Et absolute, sine} Orémus:

\oratiolabel{Oratio}

\lettrine[lines=2]{S}{ANCTA} María, et omnes Sancti intercédant pro nobis ad Dóminum, ut nos mereámur ab eo adjuvári et salvári, qui vivit et regnat in sæcula sæculórum. \R\ Amen.

\V\ Deus, in adjutórium meum inténde. \R\ Dómine, ad adjuvándum me festína. \rubric{Et dicitur ter; ultimo additur:} Glória Patri. Sicut erat.

Kýrie, eléison. Christe, eléison. Kýrie, eléison.

Pater noster \rubric{secreto usque ad} \V\ Et ne nos indúcas in tentatiónem. \R\ Sed líbera nos a malo. \V\ Réspice in servos tuos, Dómine, et in ópera tua, et dírige fílios eórum. \R\ Et sit splendor Dómini, Dei nostri, super nos, et ópera mánuum nostrárum dírige super nos, et opus mánuum nostrárum dírige. \V\ Glória Patri. \R\ Sicut erat.

Orémus. \oratiolabel{Oratio}

\lettrine[lines=2]{D}{IRIGERE} et sanctificáre, régere et gubernáre dignáre, Dómine Deus, Rex cæli et terræ, hódie corda et córpora nostra, sensus, sermónes et actus nostros in lege tua, et in opéribus mandatórum tuórum: ut hic, et in ætérnum, te auxiliánte, salvi et líberi esse mereámur, Salvátor mundi: Qui vivis et regnas in sæcula sæculórum. \R\ Amen.

\V\ Jube, domne.

\benedlabel{Bened.} Dies et actus nostros in sua pace dispónat Dóminus omnípotens. \R\ Amen.

% ===== END OF BATCH 5 -- Covers autumnalis.md lines 1443-1984,
% completing "Ad Laudes" (Psalms 92, 99, 62, the Canticle of the
% Three Young Men "Benedicite," Psalm 148, the hymn "Ecce, jam noctis
% tenuatur umbra" for ordinary Sundays after Pentecost, and the
% Benedictus) and the body of "Ad Primam" (Psalm 117 "Confitemini,"
% the special Sunday substitute Psalm 53, Psalm 118 sections i-ii,
% the complete Athanasian Creed "Quicumque vult," the brief
% responsory, the Preces, the compact bold "Confíteor. Misereátur.
% Indulgéntiam." cross-reference rendered \formulalabel{} since the
% source bold-marks it here (matching the established rule that this
% command is used only when the source itself bolds the phrase,
% per precedent in verna.tex), the Martyrologium cue, and the closing
% Orationes). TWO SILENT CORRECTIONS of what appear to be markdown/
% OCR artifacts, per the standing rule for unambiguous mechanical
% errors: (1) a stray literal "\*" interrupting the hymn "Jam lucis
% orto sidere" at its final stanza ("\* Deo Patri sit gloria," verna.md
% line 1687) is dropped, since no hymn anywhere in this project uses
% chant-pause asterisks within its verse, and this same hymn appeared
% correctly without it in Batch 2; (2) a spurious paragraph break
% splitting the second stanza of "Ecce, jam noctis tenuatur umbra"
% into two two-line fragments (lines 1612-1616) is closed up into one
% continuous four-line stanza, matching the Sapphic-Adonic meter of
% every other stanza in the same hymn. The Athanasian Creed
% "Quicumque vult," though given in full capitals at its opening
% ("QUICUMQUE vult salvus esse"), receives no lettrine, matching the
% precedent already set in hiemalis.tex: the governing instruction's
% lettrine rule names oratio, psalm, canticle, hymn, and lectio
% specifically, and a Symbolum (Creed) is none of these, so
% capitalization alone does not trigger a lettrine here as it does
% for unheaded Oratio-like text elsewhere. Several versicle/response
% pairs in this batch print only the versicle with no response
% ("℣. Dóminus vobíscum." alone, "℣. Dóminus vobíscum. ℣.
% Benedicámus Dómino." with neither response given) -- transcribed
% exactly as the source abbreviates them, without supplying the
% missing responses. The document continues with the Lectio brevis
% variants for Prime (autumnalis.md line 1986), beginning with "In
% Officio dominicali" and continuing through several Proprio de
% Tempore feast-specific substitutions (Ssmæ Trinitatis, Ssmi
% Corporis Christi, etc.) =====
\variantheading{In Officio dominicali}

\lettrine[lines=2]{D}{OMINUS} autem dírigat corda et córpora nostra in caritáte Dei, et patiéntia Christi. Tu autem.

\variantheading{E Proprio de Tempore}

\variantheading{D. N. J. C. Regis}

\lettrine[lines=2]{I}{N} ipso complácuit omnem plenitúdinem inhabitáre, et per eum reconciliáre ómnia in ipsum, pacíficans per sánguinem crucis ejus sive quæ in terris, sive quæ in cælis sunt, in Christo Jesu Dómino nostro.

\variantheading{E Proprio Sanctorum}

\variantheading{S. Raphaëlis Archangeli}

\lettrine[lines=2]{T}{EMPUS} est, ut revértar ad eum, qui me misit: vos autem benedícite Deum, et narráte ómnia mirabília ejus.

\variantheading{Omnium Sanctorum}

\lettrine[lines=2]{B}{ENEDÍCTIO} et cláritas et sapiéntia et gratiárum áctio, honor, virtus et fortitúdo Deo nostro in sǽcula sæculórum. Amen.

\variantheading{S. Andreæ Apost.}

\lettrine[lines=2]{I}{SAÍAS} enim dicit: Dómine, quis crédidit audítui nostro? Ergo fides ex audítu, audítus autem per verbum Christi. Sed dico: Numquid non audiérunt? Et quidem in omnem terram exívit sonus eórum, et in fines orbis terræ verba eórum.

\variantheading{E Communi Sanctorum}

\variantheading{Pro Apostolis}

\lettrine[lines=2]{I}{BANT} Apóstoli gaudéntes a conspéctu concílii, quóniam digni hábiti sunt pro nómine Jesu contuméliam pati.

\variantheading{Pro Martyre}

\lettrine[lines=2]{J}{USTUS} cor suum trádidit ad vigilándum dilúculo ad Dóminum, qui fecit illum, et in conspéctu Altíssimi deprecábitur.

\variantheading{Pro pluribus Martyribus}

\lettrine[lines=2]{F}{ULGEBUNT} justi, et tamquam scintíllæ in arundinéto discúrrent. Judicábunt natiónes, et dominabúntur pópulis: et regnábit Dóminus illórum in perpétuum.

\variantheading{Pro Conf. Pont.}

\lettrine[lines=2]{F}{UNGI} sacerdótio, et habére laudem in nómine ipsíus, et offérre illi incénsum dignum in odórem suavitátis.

\variantheading{Pro Conf. non Pont.}

\lettrine[lines=2]{J}{USTUM} dedúxit Dóminus per vias rectas, et osténdit illi regnum Dei, et dedit illi sciéntiam sanctórum; honestávit illum in labóribus, et complévit labóres illíus.

\variantheading{Pro Virgine Martyre}

\lettrine[lines=2]{D}{OMINE}, Deus meus, exaltásti super terram habitatiónem meam, et pro morte defluénte deprecáta sum. Invocávi Dóminum, Patrem Dómini mei, ut non derelínquat me in die tribulatiónis meæ, et in témpore superbórum sine adjutório.

\variantheading{Pro Virgine non Martyre}

\lettrine[lines=2]{O}{} \textsc{quam} pulchra est casta generátio cum claritáte! immortális est enim memória illíus: quóniam apud Deum nota est, et apud hómines.

\variantheading{Pro Martyre tantum}

\lettrine[lines=2]{L}{AUDABIT} usque ad mortem ánima mea Dóminum, quóniam éruis sustinéntes te, et líberas eos de manu angústiæ, Dómine, Deus noster.

\variantheading{Pro nec Virg. nec Mart.}

\lettrine[lines=2]{M}{ULTAE} fíliæ congregavérunt divítias: tu supergréssa es univérsas. Fallax grátia, et vana est pulchritúdo: múlier timens Dóminum, ipsa laudábitur.

\variantheading{Pro Dedicatione Ecclesiæ}

\lettrine[lines=2]{E}{T} abstérget Deus omnem lácrimam ab óculis eórum: et mors ultra non erit, neque luctus, neque clamor, neque dolor erit ultra, quia prima abiérunt. Et dixit qui sedébat in throno: Ecce, nova fácio ómnia.

\variantheading{In Festis B. Mariæ V.}

\lettrine[lines=2]{I}{N} platéis sicut cinnamómum et bálsamum aromatízans odórem dedi: quasi myrrha elécta dedi suavitátem odóris.

\V\ Adjutórium nostrum in nómine Dómini. \R\ Qui fecit cælum et terram.

\V\ Benedícite. \R\ Deus.

\benedictiolabel{Benedictio} Dóminus nos benedícat, et ab omni malo deféndat, et ad vitam perdúcat ætérnam. Et fidélium ánimæ per misericórdiam Dei requiéscant in pace. \R\ Amen.

\rubric{Deinde dicitur tantum} Pater noster \rubric{secreto, nisi sequatur alia Hora.}

\hourheading{Ad Tertiam}

Pater et Ave.

\V\ Deus, in adjutórium.

\hymnusheading{Hymnus}

\lettrine[lines=2]{N}{UNC}, Sancte, nobis, Spíritus,\\
Unum Patri cum Fílio,\\
Dignáre promptus íngeri\\
Nostro refúsus péctori.

Os, lingua, mens, sensus, vigor,\\
Confessiónem pérsonent,\\
Flamméscat igne cáritas,\\
Accéndat ardor próximos.

Præsta, Pater piíssime,\\
Patríque compar Unice,\\
Cum Spíritu Paráclito\\
Regnans per omne sæculum. Amen.

\antlabel{Ant.} Allelúja.

\psalmheading{Psalmus 118, iii}

\lettrine[lines=2]{L}{EGEM} pone mihi, Dómine, viam justificatiónum tuárum: * et exquíram eam semper.

Da mihi intelléctum, et scrutábor legem tuam: * et custódiam illam in toto corde meo.

Deduc me in sémitam mandatórum tuórum: * quia ipsam vólui.

Inclína cor meum in testimónia tua: * et non in avarítiam.

Avérte óculos meos ne vídeant vanitátem: * in via tua vivífica me.

Státue servo tuo elóquium tuum, * in timóre tuo.

Ámputa oppróbrium meum quod suspicátus sum: * quia judícia tua jucúnda.

Ecce, concupívi mandáta tua: * in æquitáte tua vivífica me.

Et véniat super me misericórdia tua, Dómine: * salutáre tuum secúndum elóquium tuum.

Et respondébo exprobrántibus mihi verbum: * quia sperávi in sermónibus tuis.

Et ne áuferas de ore meo verbum veritátis usquequáque: * quia in judíciis tuis supersperávi.

Et custódiam legem tuam semper: * in sæculum et in sæculum sæculi.

Et ambulábam in latitúdine: * quia mandáta tua exquisívi.

Et loquébar in testimóniis tuis in conspéctu regum: * et non confundébar.

Et meditábar in mandátis tuis, * quæ diléxi.

Et levávi manus meas ad mandáta tua, quæ diléxi: * et exercébar in justificatiónibus tuis.

\psalmheading{Psalmus 118, iv}

\lettrine[lines=2]{M}{EMOR} esto verbi tui servo tuo, * in quo mihi spem dedísti.

Hæc me consoláta est in humilitáte mea: * quia elóquium tuum vivificábit me.

Supérbi iníque agébant usquequáque: * a lege autem tua non declinávi.

Memor fui judiciórum tuórum a sæculo, Dómine: * et consolátus sum.

Deféctio ténuit me, * pro peccatóribus derelinquéntibus legem tuam.

Cantábiles mihi erant justificatiónes tuæ, * in loco peregrinatiónis meæ.

Memor fui nocte nóminis tui, Dómine: * et custodívi legem tuam.

Hæc facta est mihi: * quia justificatiónes tuas exquisívi.

Pórtio mea, Dómine, * dixi custodíre legem tuam.

Deprecátus sum fáciem tuam in toto corde meo: * miserére mei secúndum elóquium tuum.

Cogitávi vias meas: * et convérti pedes meos in testimónia tua.

Parátus sum, et non sum turbátus: * ut custódiam mandáta tua.

Funes peccatórum circumpléxi sunt me: * et legem tuam non sum oblítus.

Média nocte surgébam ad confiténdum tibi, * super judícia justificatiónis tuæ.

Párticeps ego sum ómnium timéntium te: * et custodiéntium mandáta tua.

Misericórdia tua, Dómine, plena est terra: * justificatiónes tuas doce me.

\psalmheading{Psalmus 118, v}

\lettrine[lines=2]{B}{ONITATEM} fecísti cum servo tuo, Dómine, * secúndum verbum tuum.

Bonitátem, et disciplínam, et sciéntiam doce me: * quia mandátis tuis crédidi.

Priúsquam humiliárer ego delíqui: * proptérea elóquium tuum custodívi.

Bonus es tu: * et in bonitáte tua doce me justificatiónes tuas.

Multiplicáta est super me iníquitas superbórum: * ego autem in toto corde meo scrutábor mandáta tua.

Coagulátum est sicut lac cor eórum: * ego vero legem tuam meditátus sum.

Bonum mihi quia humiliásti me: * ut discam justificatiónes tuas.

Bonum mihi lex oris tui, * super míllia auri et argénti.

Manus tuæ fecérunt me, et plasmavérunt me: * da mihi intelléctum, et discam mandáta tua.

Qui timent te vidébunt me, et lætabúntur: * quia in verba tua supersperávi.

Cognóvi, Dómine, quia æquitas judícia tua: * et in veritáte tua humiliásti me.

Fiat misericórdia tua ut consolétur me, * secúndum elóquium tuum servo tuo.

Véniant mihi miseratiónes tuæ, et vivam: * quia lex tua meditátio mea est.

Confundántur supérbi, quia injúste iniquitátem fecérunt in me: * ego autem exercébor in mandátis tuis.

Convertántur mihi timéntes te: * et qui novérunt testimónia tua.

Fiat cor meum immaculátum in justificatiónibus tuis, * ut non confúndar.

\antlabel{Ant.} Allelúja, deduc me, Dómine, in sémitam mandatórum tuórum, allelúja, allelúja.

\capitulumlabel{Capitulum}{1 Joann. 4, 16}

\lettrine[lines=2]{D}{EUS} cáritas est: et qui manet in caritáte, in Deo manet, et Deus in eo.

\R\ br. Inclína cor meum, Deus, * In testimónia tua. Inclína. \V\ Avérte óculos meos, ne vídeant vanitátem: in via tua vivífica me. In testimónia tua. Glória Patri. Inclína. \V\ Ego dixi: Dómine, miserére mei. \R\ Sana ánimam meam, quia peccávi tibi.

\rubric{Oratio conveniens. Reliqua ut in Ordinario. 18}

% ===== END OF BATCH 6 -- Covers autumnalis.md lines 1986-2200,
% completing the Prime Lectio brevis variant list (substitute lessons
% across "In Officio dominicali," "E Proprio de Tempore" [here just
% Christ the King, the one movable feast of the Proper of Time this
% volume's Autumn scope reaches], "E Proprio Sanctorum" [St. Raphael
% the Archangel, All Saints, St. Andrew the Apostle -- autumnalis.md's
% own shorter three-feast list, verified via direct diff against
% aestiva.md's parallel eight-feast Summer list (Trinity Sunday
% through St. Lawrence) and substituted fresh from autumnalis.md's
% own text rather than carried over], and "E Communi Sanctorum" [the
% full set of twelve Common-of-Saints categories from Apostles
% through Marian feasts, confirmed byte-identical to aestiva.md's
% copy since the Common-of-Saints capitula do not vary by season])
% and the complete "Ad Tertiam" (hymn, Psalm 118 sections iii-v,
% Capitulum, and brief responsory). Each lesson, lacking its own
% repeated "Lectio brevis" label, is nonetheless lettrined on its
% full-caps opening word, consistent with its functional role as a
% Lectio substitute. "Pro Virgine non Martyre" opens "O QUAM
% pulchra..." with BOTH words in full caps (unlike the normal-case "O
% Mater"/"O Felícem" forms seen in Batch 166 of verna.tex), matching
% instead the all-caps "O GLORIOSA"/"O QUALE" pattern from that same
% project, so it receives the single-letter lettrine plus
% \textsc{Quam}. A second stray "\*" artifact, identical in kind to
% the one corrected in Batch 5, interrupts the final stanza of "Nunc,
% Sancte, nobis, Spiritus" at autumnalis.md line 2083 (within Ad
% Tertiam's hymn) and is silently dropped here for the same reason
% (no hymn anywhere in this project uses a chant-pause asterisk
% within its verse). The identical artifact recurs in Ad Sextam's
% hymn in the next batch and will be corrected the same way there.
% The document continues with "Ad Sextam" (autumnalis.md line 2202)
% =====
\hourheading{Ad Sextam}

Pater noster et Ave.

\V\ Deus, in adjutórium.

\hymnusheading{Hymnus}

\lettrine[lines=2]{R}{ECTOR} potens, verax Deus,\\
Qui témperas rerum vices,\\
Splendóre mane illúminas,\\
Et ígnibus merídiem:

Exstíngue flammas lítium,\\
Aufer calórem nóxium,\\
Confer salútem córporum,\\
Verámque pacem córdium.

Præsta, Pater piíssime,\\
Patríque compar Unice,\\
Cum Spíritu Paráclito\\
Regnans per omne sæculum. Amen.

\antlabel{Ant.} Allelúja.

\psalmheading{Psalmus 118, vi}

\lettrine[lines=2]{D}{EFECIT} in salutáre tuum ánima mea: * et in verbum tuum supersperávi.

Defecérunt óculi mei in elóquium tuum, * dicéntes: Quando consoláberis me?

Quia factus sum sicut uter in pruína: * justificatiónes tuas non sum oblítus.

Quot sunt dies servi tui? * quando fácies de persequéntibus me judícium?

Narravérunt mihi iníqui fabulatiónes: * sed non ut lex tua.

Omnia mandáta tua véritas: * iníque persecúti sunt me, ádjuva me.

Paulo minus consummavérunt me in terra: * ego autem non derelíqui mandáta tua.

Secúndum misericórdiam tuam vivífica me: * et custódiam testimónia oris tui.

In ætérnum, Dómine, * verbum tuum pérmanet in cælo.

In generatiónem et generatiónem véritas tua: * fundásti terram, et pérmanet.

Ordinatióne tua persevérat dies: * quóniam ómnia sérviunt tibi.

Nisi quod lex tua meditátio mea est: * tunc forte periíssem in humilitáte mea.

In ætérnum non oblivíscar justificatiónes tuas: * quia in ipsis vivificásti me.

Tuus sum ego, salvum me fac: * quóniam justificatiónes tuas exquisívi.

Me exspectavérunt peccatóres ut pérderent me: * testimónia tua intelléxi.

Omnis consummatiónis vidi finem: * latum mandátum tuum nimis.

\psalmheading{Psalmus 118, vii}

\lettrine[lines=2]{Q}{UOMODO} diléxi legem tuam, Dómine? * tota die meditátio mea est.

Super inimícos meos prudéntem me fecísti mandáto tuo: * quia in ætérnum mihi est.

Super omnes docéntes me intelléxi: * quia testimónia tua meditátio mea est.

Super senes intelléxi: * quia mandáta tua quæsívi.

Ab omni via mala prohíbui pedes meos: * ut custódiam verba tua.

A judíciis tuis non declinávi: * quia tu legem posuísti mihi.

Quam dúlcia fáucibus meis elóquia tua, * super mel ori meo!

A mandátis tuis intelléxi: * proptérea odívi omnem viam iniquitátis.

Lucérna pédibus meis verbum tuum, * et lumen sémitis meis.

Jurávi, et státui * custodíre judícia justítiæ tuæ.

Humiliátus sum usquequáque, Dómine: * vivífica me secúndum verbum tuum.

Voluntária oris mei beneplácita fac, Dómine: * et judícia tua doce me.

Anima mea in mánibus meis semper: * et legem tuam non sum oblítus.

Posuérunt peccatóres láqueum mihi: * et de mandátis tuis non errávi.

Hereditáte acquisívi testimónia tua in ætérnum: * quia exsultátio cordis mei sunt.

Inclinávi cor meum ad faciéndas justificatiónes tuas in ætérnum, * propter retributiónem.

\psalmheading{Psalmus 118, viii}

\lettrine[lines=2]{I}{NIQUOS} ódio hábui: * et legem tuam diléxi.

Adjútor et suscéptor meus es tu: * et in verbum tuum supersperávi.

Declináte a me, malígni: * et scrutábor mandáta Dei mei.

Súscipe me secúndum elóquium tuum, et vivam: * et non confúndas me ab exspectatióne mea.

Adjuva me, et salvus ero: * et meditábor in justificatiónibus tuis semper.

Sprevísti omnes discedéntes a judíciis tuis: * quia injústa cogitátio eórum.

Prævaricántes reputávi omnes peccatóres terræ: * ídeo diléxi testimónia tua.

Confíge timóre tuo carnes meas: * a judíciis enim tuis tímui.

Feci judícium et justítiam: * non tradas me calumniántibus me.

Súscipe servum tuum in bonum: * non calumniéntur me supérbi.

Óculi mei defecérunt in salutáre tuum: * et in elóquium justítiæ tuæ.

Fac cum servo tuo secúndum misericórdiam tuam: * et justificatiónes tuas doce me.

Servus tuus sum ego: * da mihi intelléctum, ut sciam testimónia tua.

Tempus faciéndi, Dómine: * dissipavérunt legem tuam.

Ideo diléxi mandáta tua, * super aurum et topázion.

Proptérea ad ómnia mandáta tua dirigébar: * omnem viam iníquam ódio hábui.

\antlabel{Ant.} Allelúja, tuus sum ego, salvum me fac, Dómine, allelúja, allelúja.

\capitulumlabel{Capitulum}{Gal. 6, 2}

\lettrine[lines=2]{A}{LTER} altérius ónera portáte, et sic adimplébitis legem Christi.

\R\ br. In ætérnum, Dómine, * Pérmanet verbum tuum. In ætérnum. \V\ In sæculum sæculi véritas tua. Pérmanet. Glória Patri. In ætérnum.

\V\ Dóminus regit me, et nihil mihi déerit. \R\ In loco páscuæ ibi me collocávit.

\rubric{Oratio conveniens. Reliqua ut in Ordinario. 20}

\hourheading{Ad Nonam}

Pater noster et Ave.

\V\ Deus, in adjutórium.

\hymnusheading{Hymnus}

\lettrine[lines=2]{R}{ERUM}, Deus, tenax vigor,\\
Immótus in te pérmanens,\\
Lucis diúrnæ témpora\\
Successíbus detérminans:

Largíre lumen véspere,\\
Quo vita nusquam décidat,\\
Sed præmium mortis sacræ\\
Perénnis instet glória.

Præsta, Pater piíssime,\\
Patríque compar Unice,\\
Cum Spíritu Paráclito\\
Regnans per omne sæculum. Amen.

\antlabel{Ant.} Allelúja.

\psalmheading{Psalmus 118, ix}

\lettrine[lines=2]{M}{IRABILIA} testimónia tua: * ídeo scrutáta est ea ánima mea.

Declarátio sermónum tuórum illúminat: * et intelléctum dat párvulis.

Os meum apérui, et attráxi spíritum: * quia mandáta tua desiderábam.

Aspice in me, et miserére mei, * secúndum judícium diligéntium nomen tuum.

Gressus meos dírige secúndum elóquium tuum: * et non dominétur mei omnis injustítia.

Rédime me a calúmniis hóminum: * ut custódiam mandáta tua.

Fáciem tuam illúmina super servum tuum: * et doce me justificatiónes tuas.

Exitus aquárum deduxérunt óculi mei: * quia non custodiérunt legem tuam.

Justus es, Dómine: * et rectum judícium tuum.

Mandásti justítiam testimónia tua: * et veritátem tuam nimis.

Tabéscere me fecit zelus meus: * quia oblíti sunt verba tua inimíci mei.

Ignítum elóquium tuum veheménter: * et servus tuus diléxit illud.

Adolescéntulus sum ego et contémptus: * justificatiónes tuas non sum oblítus.

Justítia tua, justítia in ætérnum: * et lex tua véritas.

Tribulátio, et angústia invenérunt me: * mandáta tua meditátio mea est.

Æquitas testimónia tua in ætérnum: * intelléctum da mihi, et vivam. Glória Patri.

\psalmheading{Psalmus 118, x}

\lettrine[lines=2]{C}{LAMAVI} in toto corde meo, exáudi me, Dómine: * justificatiónes tuas requíram.

Clamávi ad te, salvum me fac: * ut custódiam mandáta tua.

Prævéni in maturitáte, et clamávi: * quia in verba tua supersperávi.

Prævenérunt óculi mei ad te dilúculo: * ut meditárer elóquia tua.

Vocem meam audi secúndum misericórdiam tuam, Dómine: * et secúndum judícium tuum vivífica me.

Appropinquavérunt persequéntes me iniquitáti: * a lege autem tua longe facti sunt.

Prope es tu, Dómine: * et omnes viæ tuæ véritas.

Inítio cognóvi de testimóniis tuis: * quia in ætérnum fundásti ea.

Vide humilitátem meam, et éripe me: * quia legem tuam non sum oblítus.

Júdica judícium meum, et rédime me: * propter elóquium tuum vivífica me.

Longe a peccatóribus salus: * quia justificatiónes tuas non exquisiérunt.

Misericórdiæ tuæ multæ, Dómine: * secúndum judícium tuum vivífica me.

Multi qui persequúntur me, et tríbulant me: * a testimóniis tuis non declinávi.

Vidi prævaricántes, et tabescébam: * quia elóquia tua non custodiérunt.

Vide quóniam mandáta tua diléxi, Dómine: * in misericórdia tua vivífica me.

Princípium verbórum tuórum, véritas: * in ætérnum ómnia judícia justítiæ tuæ.

\psalmheading{Psalmus 118, xi}

\lettrine[lines=2]{P}{RINCIPES} persecúti sunt me gratis: * et a verbis tuis formidávit cor meum.

Lætábor ego super elóquia tua: * sicut qui invénit spólia multa.

Iniquitátem ódio hábui, et abominátus sum: * legem autem tuam diléxi.

Sépties in die laudem dixi tibi, * super judícia justítiæ tuæ.

Pax multa diligéntibus legem tuam: * et non est illis scándalum.

Exspectábam salutáre tuum, Dómine: * et mandáta tua diléxi.

Custodívit ánima mea testimónia tua: * et diléxit ea veheménter.

Servávi mandáta tua, et testimónia tua: * quia omnes viæ meæ in conspéctu tuo.

Appropínquet deprecátio mea in conspéctu tuo, Dómine: * juxta elóquium tuum da mihi intelléctum.

Intret postulátio mea in conspéctu tuo: * secúndum elóquium tuum éripe me.

Eructábunt lábia mea hymnum, * cum docúeris me justificatiónes tuas.

Pronuntiábit lingua mea elóquium tuum: * quia ómnia mandáta tua æquitas.

Fiat manus tua ut salvet me: * quóniam mandáta tua elégi.

Concupívi salutáre tuum, Dómine: * et lex tua meditátio mea est.

Vivet ánima mea, et laudábit te: * et judícia tua adjuvábunt me.

Errávi, sicut ovis, quæ périit: * quære servum tuum, quia mandáta tua non sum oblítus.

\antlabel{Ant.} Allelúja, fáciem tuam, Dómine, illúmina super servum tuum, allelúja, allelúja.

\capitulumlabel{Capitulum}{1 Cor. 6, 20}

\lettrine[lines=2]{E}{MPTI} enim estis prétio magno. Glorificáte et portáte Deum in córpore vestro.

\R\ br. Clamávi in toto corde meo: * Exáudi me, Dómine. Clamávi. \V\ Justificatiónes tuas requíram. Exáudi. Glória Patri. Clamávi. \V\ Ab occúltis meis munda me, Dómine. \R\ Et ab aliénis parce servo tuo.

\rubric{Oratio conveniens. Reliqua ut in Ordinario. 21}

% ===== END OF BATCH 7 -- Covers autumnalis.md lines 2202-2477,
% completing "Ad Sextam" (Psalm 118 sections vi-viii, Capitulum, and
% brief responsory) and "Ad Nonam" (Psalm 118 sections ix-xi,
% Capitulum, and brief responsory). A third occurrence of the stray
% "\*" artifact already corrected twice before (verna.md's hymn
% repetitions aside -- this is autumnalis.md's own recurring OCR/
% markdown glitch) interrupts Ad Sextam's hymn "Rector potens, verax
% Deus" at its final stanza (autumnalis.md line 2359) and is silently
% dropped, consistent with Batches 5 and 6. Psalm 118's ninth section
% (Mirabilia) is the first of this long psalm's many sections to
% print "Glória Patri." at its end (autumnalis.md line 2398) rather than
% omitting it per the usual pattern established since Batch 4 --
% transcribed exactly as the source gives it, neither adding nor
% removing a doxology relative to what each specific occurrence
% prints. The document continues with "Ad Vesperas" (autumnalis.md line
% 2515), opening the classic Sunday Vespers psalms (109 "Dixit
% Dominus," 110, etc.) =====
\hourheading{Ad Vesperas}

Pater noster et Ave.

\V\ Deus, in adjutórium.

\antlabel{Ant.} Dixit Dóminus. \textdagger

\rubric{¶ Antiphonæ cum Psalmis assignatæ adhibentur semper, quando fit Officium de Dominica, in Dominicis extra Octavas Ssmi Corporis Christi ac Sacratissimi Cordis Jesu.}

\psalmheading{Psalmus 109}

\lettrine[lines=2]{D}{IXIT} Dóminus \textdagger\ Dómino meo: * Sede a dextris meis:

Donec ponam inimícos tuos, * scabéllum pedum tuórum.

Virgam virtútis tuæ emíttet Dóminus ex Sion: * domináre in médio inimicórum tuórum.

Tecum princípium in die virtútis tuæ in splendóribus sanctórum: * ex útero ante lucíferum génui te.

Jurávit Dóminus, et non pænitébit eum: * Tu es sacérdos in ætérnum secúndum órdinem Melchísedech.

Dóminus a dextris tuis, * confrégit in die iræ suæ reges.

Judicábit in natiónibus, implébit ruínas: * conquassábit cápita in terra multórum.

De torrénte in via bibet: * proptérea exaltábit caput.

\antlabel{Ant.} Dixit Dóminus Dómino meo: Sede a dextris meis.

\antlabel{Ant.} Magna ópera Dómini.

\psalmheading{Psalmus 110}

\lettrine[lines=2]{C}{ONFITEBOR} tibi, Dómine, in toto corde meo: * in consílio justórum, et congregatióne.

Magna ópera Dómini: * exquisíta in omnes voluntátes ejus.

Conféssio et magnificéntia opus ejus: * et justítia ejus manet in sæculum sæculi.

Memóriam fecit mirabílium suórum, miséricors et miserátor Dóminus: * escam dedit timéntibus se.

Memor erit in sæculum testaménti sui: * virtútem óperum suórum annuntiábit pópulo suo:

Ut det illis hereditátem géntium: * ópera mánuum ejus véritas, et judícium.

Fidélia ómnia mandáta ejus: confirmáta in sæculum sæculi, * facta in veritáte et æquitáte.

Redemptiónem misit pópulo suo: * mandávit in ætérnum testaméntum suum.

Sanctum, et terríbile nomen ejus: * inítium sapiéntiæ timor Dómini.

Intelléctus bonus ómnibus faciéntibus eum: * laudátio ejus manet in sæculum sæculi.

\antlabel{Ant.} Magna ópera Dómini: exquisíta in omnes voluntátes ejus.

\antlabel{Ant.} Qui timet Dóminum.

\psalmheading{Psalmus 111}

\lettrine[lines=2]{B}{EATUS} vir, qui timet Dóminum: * in mandátis ejus volet nimis.

Potens in terra erit semen ejus: * generátio rectórum benedicétur.

Glória, et divítiæ in domo ejus: * et justítia ejus manet in sæculum sæculi.

Exórtum est in ténebris lumen rectis: * miséricors, et miserátor, et justus.

Jucúndus homo qui miserétur et cómmodat, dispónet sermónes suos in judício: * quia in ætérnum non commovébitur.

In memória ætérna erit justus: * ab auditióne mala non timébit.

Parátum cor ejus speráre in Dómino, confirmátum est cor ejus: * non commovébitur donec despíciat inimícos suos.

Dispérsit, dedit paupéribus: justítia ejus manet in sæculum sæculi, * cornu ejus exaltábitur in glória.

Peccátor vidébit, et irascétur, déntibus suis fremet et tabéscet: * desidérium peccatórum períbit.

\antlabel{Ant.} Qui timet Dóminum, in mandátis ejus cupit nimis.

\antlabel{Ant.} Sit nomen Dómini.

\psalmheading{Psalmus 112}

\lettrine[lines=2]{L}{AUDATE}, púeri, Dóminum: * laudáte nomen Dómini.

Sit nomen Dómini benedíctum, * ex hoc nunc, et usque in sæculum.

A solis ortu usque ad occásum, * laudábile nomen Dómini.

Excélsus super omnes gentes Dóminus, * et super cælos glória ejus.

Quis sicut Dóminus, Deus noster, qui in altis hábitat, * et humília réspicit in cælo et in terra?

Súscitans a terra ínopem, * et de stércore érigens páuperem:

Ut cóllocet eum cum princípibus, * cum princípibus pópuli sui.

Qui habitáre facit stérilem in domo, * matrem filiórum lætántem.

\antlabel{Ant.} Sit nomen Dómini benedíctum in sæcula.

\antlabel{Ant.} Deus autem noster.

\psalmheading{Psalmus 113}

\lettrine[lines=2]{I}{N} éxitu Israël de Ægýpto, * domus Jacob de pópulo bárbaro:

Facta est Judǽa sanctificátio ejus, * Israël potéstas ejus.

Mare vidit, et fugit: * Jordánis convérsus est retrórsum.

Montes exsultavérunt ut aríetes, * et colles sicut agni óvium.

Quid est tibi, mare, quod fugísti: * et tu, Jordánis, quia convérsus es retrórsum?

Montes, exsultástis sicut aríetes, * et colles, sicut agni óvium?

A fácie Dómini mota est terra, * a fácie Dei Jacob.

Qui convértit petram in stagna aquárum, * et rupem in fontes aquárum.

Non nobis, Dómine, non nobis: * sed nómini tuo da glóriam.

Super misericórdia tua, et veritáte tua: * nequándo dicant gentes: Ubi est Deus eórum?

Deus autem noster in cælo: * ómnia quæcúmque vóluit, fecit.

Simulácra géntium argéntum, et aurum, * ópera mánuum hóminum.

Os habent, et non loquéntur: * óculos habent, et non vidébunt.

Aures habent, et non áudient: * nares habent, et non odorábunt.

Manus habent, et non palpábunt: pedes habent, et non ambulábunt: * non clamábunt in gútture suo.

Símiles illis fiant qui fáciunt ea: * et omnes qui confídunt in eis.

Domus Israël sperávit in Dómino: * adjútor eórum et protéctor eórum est.

Domus Aaron sperávit in Dómino: * adjútor eórum et protéctor eórum est.

Qui timent Dóminum, speravérunt in Dómino: * adjútor eórum et protéctor eórum est.

Dóminus memor fuit nostri: * et benedíxit nobis:

Benedíxit dómui Israël: * benedíxit dómui Aaron.

Benedíxit ómnibus, qui timent Dóminum, * pusíllis cum majóribus.

Adjíciat Dóminus super vos: * super vos, et super fílios vestros.

Benedícti vos a Dómino, * qui fecit cælum, et terram.

Cælum cæli Dómino: * terram autem dedit fíliis hóminum.

Non mórtui laudábunt te, Dómine: * neque omnes, qui descéndunt in inférnum.

Sed nos qui vívimus, benedícimus Dómino, * ex hoc nunc et usque in sæculum.

\antlabel{Ant.} Deus autem noster in cælo: ómnia quæcúmque vóluit, fecit.

\capitulumlabel{Capitulum}{2 Cor. 1, 3-4}

\lettrine[lines=2]{B}{ENEDICTUS} Deus, et Pater Dómini nostri Jesu Christi, Pater misericordiárum, et Deus totíus consolatiónis, qui consolátur nos in omni tribulatióne nostra.

\hymnusheading{Hymnus}

\lettrine[lines=2]{L}{UCIS} Creátor óptime,\\
Lucem diérum próferens,\\
Primórdiis lucis novæ,\\
Mundi parans oríginem:

Qui mane junctum vésperi\\
Diem vocári præcipis:\\
Illábitur tetrum chaos,\\
Audi preces cum flétibus.

Ne mens graváta crímine,\\
Vitæ sit exsul múnere,\\
Dum nil perénne cógitat,\\
Seséque culpis ílligat.

Cæléste pulset óstium:\\
Vitále tollat præmium:\\
Vitémus omne nóxium:\\
Purgémus omne péssimum.

Præsta, Pater piíssime,\\
Patríque compar Unice,\\
Cum Spíritu Paráclito\\
Regnans per omne sæculum. Amen.

\V\ Dirigátur, Dómine, orátio mea. \R\ Sicut incénsum in conspéctu tuo.

\rubric{Ad Magnificat Antiphona ut in Proprio de Tempore.}

\canticumheading{Canticum B. Mariæ Virg.}{Luc. 1, 46-55}

\lettrine[lines=2]{M}{AGNIFICAT} \textdagger\ ánima mea Dóminum:

Et exsultávit spíritus meus * in Deo, salutári meo.

Quia respéxit humilitátem ancíllæ suæ: * ecce enim ex hoc beátam me dicent omnes generatiónes.

Quia fecit mihi magna qui potens est: * et sanctum nomen ejus.

Et misericórdia ejus a progénie in progénies * timéntibus eum.

Fecit poténtiam in bráchio suo: * dispérsit supérbos mente cordis sui.

Depósuit poténtes de sede, * et exaltávit húmiles.

Esuriéntes implévit bonis: * et dívites dimísit inánes.

Suscépit Israël, púerum suum, * recordátus misericórdiæ suæ.

Sicut locútus est ad patres nostros, * Abraham, et sémini ejus in sæcula.

\rubric{Oratio conveniens. Reliqua ut in Ordinario. 22}

\hourheading{Ad Completorium}

\V\ Jube, domne.

\benedlabel{Bened.} Noctem quiétam, et finem perféctum concédat nobis Dóminus omnípotens. \R\ Amen.

\lectioheading{Lectio br.}{1 Petri 5, 8-9}

\lettrine[lines=2]{F}{RATRES}: Sóbrii estóte, et vigiláte: quia adversárius vester diábolus tamquam leo rúgiens círcuit, quærens quem dévoret: cui resístite fortes in fide. Tu autem, Dómine, miserére nobis.

\R\ Deo grátias.

\V\ Adjutórium nostrum in nómine Dómini. \R\ Qui fecit cælum et terram.

\rubric{Pater noster, quod dicitur totum secreto.}

\formulalabel{Confíteor. Misereátur. Indulgéntiam.}

\V\ Convérte nos, Deus, salutáris noster. \R\ Et avérte iram tuam a nobis.

\V\ Deus, in adjutórium.

\antlabel{Ant.} Miserére.

\psalmheading{Psalmus 4}

\lettrine[lines=2]{C}{UM} invocárem exaudívit me Deus justítiæ meæ: * in tribulatióne dilatásti mihi.

Miserére mei, * et exáudi oratiónem meam.

Fílii hóminum, úsquequo gravi corde? * ut quid dilígitis vanitátem, et quæritis mendácium?

Et scitóte quóniam mirificávit Dóminus sanctum suum: * Dóminus exáudiet me cum clamávero ad eum.

Irascímini, et nolíte peccáre: * quæ dícitis in córdibus vestris, in cubílibus vestris compungímini.

Sacrificáte sacrifícium justítiæ, et speráte in Dómino. * Multi dicunt: Quis osténdit nobis bona?

Signátum est super nos lumen vultus tui, Dómine: * dedísti lætítiam in corde meo.

A fructu fruménti, vini, et ólei sui * multiplicáti sunt.

In pace in idípsum * dórmiam, et requiéscam;

Quóniam tu, Dómine, singuláriter in spe * constituísti me.

\psalmheading{Psalmus 90}

\lettrine[lines=2]{Q}{UI} hábitat in adjutório Altíssimi, * in protectióne Dei cæli commorábitur.

Dicet Dómino: Suscéptor meus es tu, et refúgium meum: * Deus meus, sperábo in eum.

Quóniam ipse liberávit me de láqueo venántium, * et a verbo áspero.

Scápulis suis obumbrábit tibi: * et sub pennis ejus sperábis.

Scuto circúmdabit te véritas ejus: * non timébis a timóre noctúrno.

A sagítta volánte in die, a negótio perambulánte in ténebris: * ab incúrsu, et dæmónio meridiáno.

Cadent a látere tuo mille, et decem míllia a dextris tuis: * ad te autem non appropinquábit.

Verúmtamen óculis tuis considerábis: * et retributiónem peccatórum vidébis.

Quóniam tu es, Dómine, spes mea: * Altíssimum posuísti refúgium tuum.

Non accédet ad te malum: * et flagéllum non appropinquábit tabernáculo tuo.

Quóniam Angelis suis mandávit de te: * ut custódiant te in ómnibus viis tuis.

In mánibus portábunt te: * ne forte offéndas ad lápidem pedem tuum.

Super áspidem, et basilíscum ambulábis: * et conculcábis leónem et dracónem.

Quóniam in me sperávit, liberábo eum: * prótegam eum, quóniam cognóvit nomen meum.

Clamábit ad me, et ego exáudiam eum: * cum ipso sum in tribulatióne: erípiam eum et glorificábo eum.

Longitúdine diérum replébo eum: * et osténdam illi salutáre meum.

\psalmheading{Psalmus 133}

\lettrine[lines=2]{E}{CCE}, nunc benedícite Dóminum, * omnes servi Dómini:

Qui statis in domo Dómini, * in átriis domus Dei nostri.

In nóctibus extóllite manus vestras in sancta, * et benedícite Dóminum.

Benedícat te Dóminus ex Sion, * qui fecit cælum et terram.

\antlabel{Ant.} Miserére mihi, Dómine, et exáudi oratiónem meam.

\hymnusheading{Hymnus}

\lettrine[lines=2]{T}{E} lucis ante términum,\\
Rerum Creátor, póscimus,\\
Ut pro tua cleméntia\\
Sis præsul et custódia.

Procul recédant sómnia,\\
Et nóctium phantásmata;\\
Hostémque nostrum cómprime,\\
Ne polluántur córpora.

Præsta, Pater piíssime,\\
Patríque compar Unice,\\
Cum Spíritu Paráclito\\
Regnans per omne sæculum. Amen.

\capitulumlabel{Capitulum}{Jer. 14, 9}

\lettrine[lines=2]{T}{U} autem in nobis es, Dómine, et nomen sanctum tuum invocátum est super nos: ne derelínquas nos, Dómine, Deus noster.

\R\ br. In manus tuas, Dómine, * Comméndo spíritum meum. In manus. \V\ Redemísti nos, Dómine, Deus veritátis. Comméndo. Glória Patri. In manus.

\V\ Custódi nos, Dómine, ut pupíllam óculi. \R\ Sub umbra alárum tuárum prótege nos.

\canticumheading{Canticum Simeonis}{Luc. 2, 29-32}

\lettrine[lines=2]{N}{UNC} dimíttis servum tuum, Dómine, * secúndum verbum tuum in pace:

Quia vidérunt óculi mei * salutáre tuum,

Quod parásti * ante fáciem ómnium populórum,

Lumen ad revelatiónem géntium, * et glóriam plebis tuæ Israël.

Glória Patri.

\antlabel{Ant.} Salva nos, Dómine, vigilántes, custódi nos dormiéntes; ut vigilémus cum Christo, et requiescámus in pace.

\variantheading{Preces (si dicendæ sint)}

KYRIE, eléison. Christe, eléison. Kýrie, eléison.

Pater noster \rubric{secreto usque ad} \V\ Et ne nos indúcas in tentatiónem. \R\ Sed líbera nos a malo.

Credo in Deum \rubric{secreto usque ad} \V\ Carnis resurrectiónem. \R\ Vitam ætérnam. Amen.

\V\ Benedíctus es, Dómine, Deus patrum nostrórum. \R\ Et laudábilis et gloriósus in sæcula. \V\ Benedicámus Patrem et Fílium cum Sancto Spíritu. \R\ Laudémus, et superexaltémus eum in sæcula. \V\ Benedíctus es, Dómine, in firmaménto cæli. \R\ Et laudábilis, et gloriósus, et superexaltátus in sæcula. \V\ Benedícat et custódiat nos omnípotens et miséricors Dóminus. \R\ Amen. \V\ Dignáre, Dómine, nocte ista. \R\ Sine peccáto nos custodíre. \V\ Miserére nostri, Dómine. \R\ Miserére nostri. \V\ Fiat misericórdia tua, Dómine, super nos. \R\ Quemádmodum sperávimus in te. \V\ Dómine, exáudi oratiónem meam. \R\ Et clamor meus ad te véniat.

\V\ Dóminus vobíscum.

Orémus. \oratiolabel{Oratio}

\lettrine[lines=2]{V}{ISITA}, quæsumus, Dómine, habitatiónem istam, et omnes insídias inimíci ab ea longe repélle: Angeli tui sancti hábitent in ea, qui nos in pace custódiant; et benedíctio tua sit super nos semper. Per Dóminum nostrum. \R\ Amen.

\V\ Dóminus vobíscum.

\V\ Benedicámus Dómino.

\benedictiolabel{Benedictio} Benedícat et custódiat nos omnípotens et miséricors Dóminus, Pater, et Fílius, et Spíritus Sanctus. \R\ Amen.

\variantheading{Antiphona finalis B. M. V.}

\V\ Divínum auxílium.

Pater, Ave et Credo.

\rubric{¶ Psalmi de Dominica superius appositi dicuntur ad Completorium, adhibita simul Antiphona, etiam post utrasque Vesperas in Dominicis infra Octavas Ssmi Corporis Christi ac Sacrat. Cordis Jesu, in Duplicibus I vel II classis, in Festis novem Lectionum Domini, B. Mariæ Virginis, Angelorum, S. Joannis Baptistæ, S. Joseph, Apostolorum, Evangelistarum, et per Octavas privilegiatas; si tamen præcedentes I Vesperæ de eis integræ vel saltem a Capitulo sint celebratæ, ac II Vesperæ de ipsis integræ fuerint persolutæ.}

% ===== END OF BATCH 8 -- Covers autumnalis.md lines 2478-2863,
% completing "Ad Vesperas" (Psalms 109-113, Capitulum, hymn "Lucis
% Creator optime," and the Magnificat) and "Ad Completorium" (Lectio
% brevis, Confiteor shorthand, Psalms 4, 90, 133, hymn "Te lucis ante
% terminum," Nunc dimittis, and the Preces), and with them the ENTIRE
% "Dominica" Office. Two more occurrences of the recurring stray "\*"
% artifact (in "Lucis Creator optime" at line 2703 and "Te lucis ante
% terminum" at line 2842, both at the hymn's final stanza) are
% silently dropped, consistent with every prior occurrence in this
% volume. The closing "Antiphona finalis B. M. V." is given only as a
% bare cross-reference label with the incipit "℣. Divínum auxílium."
% (not the full "...máneat semper nobíscum. ℟. Amen."), transcribed
% exactly as abbreviated rather than completed from the Ordinarium's
% fuller text. The document continues with "Feria Secunda" (autumnalis.md
% line 2864), the first weekday of the Psalterium cycle =====
\feriaheading{Feria Secunda}

\hourheading{Ad Matutinum}

Pater, Ave et Credo.

\V\ Dómine, lábia mea.

\V\ Deus, in adjutórium.

\variantheading{In Officio feriali}

\rubric{Invitatorium et Hymnus dicuntur quotidie de Feria currenti, ut in Psalterio hic et alibi apponuntur.}

\invitatlabel{Invitat.} Veníte, * Exsultémus Dómino. \psalmlabel{Ps. 94} Veníte, exsultémus.

\rubric{In primo Psalmi versu non dicitur: Veníte, exsultémus Dómino; sed post repetitum Invitatorium, statim subjungitur: Jubilémus Deo, salutári nostro:}

\hymnusheading{Hymnus}

\lettrine[lines=2]{S}{OMNO} reféctis ártubus,\\
Spreto cubíli, súrgimus:\\
Nobis, Pater, canéntibus\\
Adésse te depóscimus.

Te lingua primum cóncinat,\\
Te mentis ardor ámbiat:\\
Ut áctuum sequéntium\\
Tu, Sancte, sis exórdium.

Cedant ténebræ lúmini,\\
Et nox diúrno síderi,\\
Ut culpa, quam nox íntulit,\\
Lucis labáscat múnere.

Precámur iídem súpplices,\\
Noxas ut omnes ámputes,\\
Et ore te canéntium\\
Laudéris omni témpore.

Præsta, Pater piíssime,\\
Patríque compar Unice,\\
Cum Spíritu Paráclito\\
Regnans per omne sæculum.\\
Amen.

\nocturnoheading{In I Nocturno}

\antlabel{Ant.} Dóminus de cælo * prospéxit super fílios hóminum.

\rubric{Antiphonæ, in hac et aliis Feriis appositæ, dicuntur semper in Officio feriali et in Festis sive trium sive novem Lectionum, quando juxta Rubricas sumendi sint Psalmi de Feria, ut in Psalterio.}

\rubric{In Officio autem trium Lectionum, tam de Tempore quam de Sanctis, omittuntur Versus et reliqua omnia, quæ in Officio novem Lectionum solent recitari in fine I et II Nocturni, et continuantur Psalmi cum suis Antiphonis usque ad Versum III Nocturni.}

\psalmheading{Psalmus 13}

\lettrine[lines=2]{D}{IXIT} insípiens in corde suo: * Non est Deus.

Corrúpti sunt, et abominábiles facti sunt in stúdiis suis: * non est qui fáciat bonum, non est usque ad unum.

Dóminus de cælo prospéxit super fílios hóminum, * ut vídeat si est intélligens, aut requírens Deum.

Omnes declinavérunt, simul inútiles facti sunt: * non est qui fáciat bonum, non est usque ad unum.

Sepúlcrum patens est guttur eórum: linguis suis dolóse agébant, * venénum áspidum sub lábiis eórum.

Quorum os maledictióne et amaritúdine plenum est: * velóces pedes eórum ad effundéndum sánguinem.

Contrítio et infelícitas in viis eórum, et viam pacis non cognovérunt: * non est timor Dei ante óculos eórum.

Nonne cognóscent omnes qui operántur iniquitátem, * qui dévorant plebem meam sicut escam panis?

Dóminum non invocavérunt, * illic trepidavérunt timóre, ubi non erat timor.

Quóniam Dóminus in generatióne justa est, consílium ínopis confudístis: * quóniam Dóminus spes ejus est.

Quis dabit ex Sion salutáre Israël? * cum avérterit Dóminus captivitátem plebis suæ, exsultábit Jacob, et lætábitur Israël.

\antlabel{Ant.} Dóminus de cælo prospéxit super fílios hóminum.

\antlabel{Ant.} Qui operátur justítiam, * requiéscet in monte sancto tuo, Dómine.

\psalmheading{Psalmus 14}

\lettrine[lines=2]{D}{OMINE}, quis habitábit in tabernáculo tuo? * aut quis requiéscet in monte sancto tuo?

Qui ingréditur sine mácula, * et operátur justítiam:

Qui lóquitur veritátem in corde suo, * qui non egit dolum in lingua sua:

Nec fecit próximo suo malum, * et oppróbrium non accépit advérsus próximos suos.

Ad níhilum dedúctus est in conspéctu ejus malígnus: * timéntes autem Dóminum gloríficat:

Qui jurat próximo suo, et non décipit, * qui pecúniam suam non dedit ad usúram, et múnera super innocéntem non accépit.

Qui facit hæc, * non movébitur in ætérnum.

\antlabel{Ant.} Qui operátur justítiam, requiéscet in monte sancto tuo, Dómine.

\antlabel{Ant.} Inclína, Dómine, * aurem tuam mihi, et exáudi verba mea.

\psalmheading{Psalmus 16}

\lettrine[lines=2]{E}{XAUDI}, Dómine, justítiam meam: * inténde deprecatiónem meam.

Áuribus pércipe oratiónem meam, * non in lábiis dolósis.

De vultu tuo judícium meum pródeat: * óculi tui vídeant æquitátes.

Probásti cor meum, et visitásti nocte: * igne me examinásti, et non est invénta in me iníquitas.

Ut non loquátur os meum ópera hóminum: * propter verba labiórum tuórum ego custodívi vias duras.

Pérfice gressus meos in sémitis tuis: * ut non moveántur vestígia mea.

Ego clamávi, quóniam exaudísti me, Deus: * inclína aurem tuam mihi, et exáudi verba mea.

Mirífica misericórdias tuas, * qui salvos facis sperántes in te.

A resisténtibus déxteræ tuæ custódi me, * ut pupíllam óculi.

Sub umbra alárum tuárum prótege me: * a fácie impiórum qui me afflixérunt.

Inimíci mei ánimam meam circumdedérunt, ádipem suum conclusérunt: * os eórum locútum est supérbiam.

Projiciéntes me nunc circumdedérunt me: * óculos suos statuérunt declináre in terram.

Suscepérunt me sicut leo parátus ad prǽdam: * et sicut cátulus leónis hábitans in abdítis.

Exsúrge, Dómine, prǽveni eum, et supplánta eum: * éripe ánimam meam ab ímpio, fráméam tuam ab inimícis manus tuæ.

Dómine, a paucis de terra dívide eos in vita eórum: * de abscónditis tuis adimplétus est venter eórum.

Saturáti sunt fíliis: * et dimisérunt relíquias suas párvulis suis.

Ego autem in justítia apparébo conspéctui tuo: * satiábor cum apparúerit glória tua.

\antlabel{Ant.} Inclína, Dómine, aurem tuam mihi, et exáudi verba mea.

\variantheading{In Festis novem Lectionum:}

\V\ Mirífica, Dómine, misericórdias tuas. \R\ Qui salvos facis sperántes in te.

\absolutiolabel{Absol.} Exáudi, Dómine.

\benedlabel{Bened.} 1 Benedictióne perpétua. 2 Unigénitus Dei Fílius. 3 Spíritus Sancti grátia.

\nocturnoheading{In II Nocturno}

\antlabel{Ant.} Díligam te, * \textdagger\ Dómine, virtus mea.

\psalmheading{Psalmus 17, i}

\lettrine[lines=2]{D}{ILIGAM} te, \textdagger\ Dómine, fortitúdo mea: * Dóminus firmaméntum meum, et refúgium meum, et liberátor meus.

Deus meus adjútor meus, * et sperábo in eum.

Protéctor meus, et cornu salútis meæ, * et suscéptor meus.

Laudans invocábo Dóminum: * et ab inimícis meis salvus ero.

Circumdedérunt me dolóres mortis: * et torréntes iniquitátis conturbavérunt me.

Dolóres inférni circumdedérunt me: * præoccupavérunt me láquei mortis.

In tribulatióne mea invocávi Dóminum, * et ad Deum meum clamávi.

Et exaudívit de templo sancto suo vocem meam: * et clamor meus in conspéctu ejus, introívit in aures ejus.

Commóta est, et contrémuit terra: * fundaménta móntium conturbáta sunt, et commóta sunt, quóniam irátus est eis.

Ascéndit fumus in ira ejus: * et ignis a fácie ejus exársit: * carbónes succénsi sunt ab eo.

Inclinávit cælos, et descéndit: * et calígo sub pédibus ejus.

Et ascéndit super Chérubim, et volávit: * volávit super pennas ventórum.

Et pósuit ténebras latíbulum suum, in circúitu ejus tabernáculum ejus: * tenebrósa aqua in núbibus áëris.

Præ fulgóre in conspéctu ejus nubes transiérunt, * grando et carbónes ignis.

Et intónuit de cælo Dóminus, et Altíssimus dedit vocem suam: * grando et carbónes ignis.

Et misit sagíttas suas, et dissipávit eos: * fúlgura multiplicávit, et conturbávit eos.

Et apparuérunt fontes aquárum, * et reveláta sunt fundaménta orbis terrárum:

Ab increpatióne tua, Dómine, * ab inspiratióne spíritus iræ tuæ.

\antlabel{Ant.} Díligam te, Dómine, virtus mea.

\antlabel{Ant.} Retríbuet mihi Dóminus * secúndum justítiam meam.

\psalmheading{Psalmus 17, ii}

\lettrine[lines=2]{M}{ISIT} de summo, et accépit me: * et assúmpsit me de aquis multis.

Erípuit me de inimícis meis fortíssimis, et ab his qui odérunt me: * quóniam confortáti sunt super me.

Prævenérunt me in die afflictiónis meæ: * et factus est Dóminus protéctor meus.

Et edúxit me in latitúdinem: * salvum me fecit, quóniam vóluit me.

Et retríbuet mihi Dóminus secúndum justítiam meam, * et secúndum puritátem mánuum meárum retríbuet mihi:

Quia custodívi vias Dómini, * nec ímpie gessi a Deo meo.

Quóniam ómnia judícia ejus in conspéctu meo: * et justítias ejus non répuli a me.

Et ero immaculátus cum eo: * et observábo me ab iniquitáte mea.

Et retríbuet mihi Dóminus secúndum justítiam meam: * et secúndum puritátem mánuum meárum in conspéctu oculórum ejus.

Cum sancto sanctus eris, * et cum viro innocénte ínnocens eris:

Et cum elécto eléctus eris: * et cum pervérso pervertéris.

Quóniam tu pópulum húmilem salvum fácies: * et óculos superbórum humiliábis.

Quóniam tu illúminas lucérnam meam, Dómine: * Deus meus, illúmina ténebras meas.

Quóniam in te erípiar a tentatióne, * et in Deo meo transgrédiar murum.

Deus meus, impollúta via ejus: elóquia Dómini igne examináta: * protéctor est ómnium sperántium in se.

Quóniam quis Deus præter Dóminum? * aut quis Deus præter Deum nostrum?

Deus, qui præcínxit me virtúte: * et pósuit immaculátam viam meam.

Qui perfécit pedes meos tamquam cervórum, * et super excélsa státuens me.

Qui docet manus meas ad prǽlium: * et posuísti, ut arcum áëreum, bráchia mea.

\antlabel{Ant.} Retríbuet mihi Dóminus secúndum justítiam meam.

\antlabel{Ant.} Vivit Dóminus, * et benedíctus Deus salútis meæ.

\psalmheading{Psalmus 17, iii}

\lettrine[lines=2]{E}{T} dedísti mihi protectiónem salútis tuæ: * et déxtera tua suscépit me:

Et disciplína tua corréxit me in finem: * et disciplína tua ipsa me docébit.

Dilatásti gressus meos subtus me: * et non sunt infirmáta vestígia mea:

Pérsequar inimícos meos, et comprehéndam illos: * et non convértar, donec defíciant.

Confríngam illos, nec póterunt stare: * cadent subtus pedes meos.

Et præcinxísti me virtúte ad bellum: * et supplantásti insurgéntes in me subtus me.

Et inimícos meos dedísti mihi dorsum, * et odiéntes me disperdidísti.

Clamavérunt, nec erat qui salvos fáceret, ad Dóminum: * nec exaudívit eos.

Et commínuam eos, ut púlverem ante fáciem venti: * ut lutum plateárum delébo eos.

Erípies me de contradictiónibus pópuli: * constítues me in caput géntium.

Pópulus, quem non cognóvi, servívit mihi: * in audítu auris obedívit mihi.

Fílii aliéni mentíti sunt mihi: * fílii aliéni inveteráti sunt, et claudicavérunt a sémitis suis.

Vivit Dóminus, et benedíctus Deus meus, * et exaltétur Deus salútis meæ.

Deus, qui das vindíctas mihi, et subdis pópulos sub me, * liberátor meus de inimícis meis iracúndis.

Et ab insurgéntibus in me exaltábis me: * a viro iníquo éripies me.

Proptérea confitébor tibi in natiónibus, Dómine: * et nómini tuo psalmum dicam.

Magníficans salútes Regis ejus, et fáciens misericórdiam Christo suo David, * et sémini ejus usque in séculum.

\antlabel{Ant.} Vivit Dóminus, et benedíctus Deus salútis meæ.

\variantheading{In Festis novem Lectionum:}

\V\ Confitébor tibi in natiónibus, Dómine. \R\ Et nómini tuo psalmum dicam.

\absolutiolabel{Absol.} Ipsíus píetas et.

\benedlabel{Bened.} 4 Deus Pater omnípotens. 5 Christus perpétuæ. 6 Ignem sui amóris.

\nocturnoheading{In III Nocturno}

\antlabel{Ant.} Exáudiat te * \textdagger\ Dóminus in die tribulatiónis. \textdagger

\psalmheading{Psalmus 19}

\lettrine[lines=2]{E}{XAUDIAT} te \textdagger\ Dóminus in die tribulatiónis: * \textdagger\ prótegat te nomen Dei Jacob.

Mittat tibi auxílium de sancto: * et de Sion tueátur te.

Memor sit omnis sacrifícii tui: * et holocáustum tuum pingue fiat.

Tríbuat tibi secúndum cor tuum: * et omne consílium tuum confírmet.

Lætábimur in salutári tuo: * et in nómine Dei nostri magnificábimur.

Ímpleat Dóminus omnes petitiónes tuas: * nunc cognóvi quóniam salvum fecit Dóminus Christum suum.

Exáudiet illum de cælo sancto suo: * in potentátibus salus déxteræ ejus.

Hi in cúrribus, et hi in equis: * nos autem in nómine Dómini, Dei nostri, invocábimus.

Ipsi obligáti sunt, et cecidérunt: * nos autem surréximus et erécti sumus.

Dómine, salvum fac regem: * et exáudi nos in die, qua invocavérimus te.

\antlabel{Ant.} Exáudiat te Dóminus in die tribulatiónis.

\antlabel{Ant.} Dómine, * \textdagger\ in virtúte tua lætábitur rex. \textdaggerdbl

\psalmheading{Psalmus 20}

\lettrine[lines=2]{D}{OMINE}, \textdagger\ in virtúte tua lætábitur rex: * \textdaggerdbl\ et super salutáre tuum exsultábit veheménter.

Desidérium cordis ejus tribuísti ei: * et voluntáte labiórum ejus non fraudásti eum.

Quóniam prævenísti eum in benedictiónibus dulcédinis: * posuísti in cápite ejus corónam de lápide pretióso.

Vitam pétiit a te: * et tribuísti ei longitúdinem diérum in séculum, et in séculum séculi.

Magna est glória ejus in salutári tuo: * glóriam et magnum decórem impónes super eum.

Quóniam dabis eum in benedictiónem in séculum séculi: * lætificábis eum in gáudio cum vultu tuo.

Quóniam rex sperat in Dómino: * et in misericórdia Altíssimi non commovébitur.

Inveniátur manus tua ómnibus inimícis tuis: * déxtera tua invéniat omnes, qui te odérunt.

Pones eos ut clíbanum ignis in témpore vultus tui: * Dóminus in ira sua conturbábit eos, et devorábit eos ignis.

Fructum eórum de terra perdes: * et semen eórum a fíliis hóminum.

Quóniam declinavérunt in te mala: * cogitavérunt consília, quæ non potuérunt stabilíre.

Quóniam pones eos dorsum: * in relíquiis tuis præparábis vultum eórum.

Exaltáre, Dómine, in virtúte tua: * cantábimus et psallémus virtútes tuas.

\antlabel{Ant.} Dómine, in virtúte tua lætábitur rex.

\antlabel{Ant.} Exaltábo te, * \textdagger\ Dómine, quóniam suscepísti me. \textdaggerdbl

\psalmheading{Psalmus 29}

\lettrine[lines=2]{E}{XALTABO} te, \textdagger\ Dómine, quóniam suscepísti me: * \textdaggerdbl\ nec delectásti inimícos meos super me.

Dómine, Deus meus, clamávi ad te, * et sanásti me.

Dómine, eduxísti ab inférno ánimam meam: * salvásti me a descendéntibus in lacum.

Psállite Dómino, sancti ejus: * et confitémini memóriæ sanctitátis ejus.

Quóniam ira in indignatióne ejus: * et vita in voluntáte ejus.

Ad vésperum demorábitur fletus: * et ad matutínum lætítia.

Ego autem dixi in abundántia mea: * Non movébor in ætérnum.

Dómine, in voluntáte tua, * præstitísti decóri meo virtútem.

Avertísti fáciem tuam a me, * et factus sum conturbátus.

Ad te, Dómine, clamábo: * et ad Deum meum deprecábor.

Quæ utílitas in sánguine meo, * dum descéndo in corruptiónem?

Numquid confitébitur tibi pulvis, * aut annuntiábit veritátem tuam?

Audívit Dóminus, et misértus est mei: * Dóminus factus est adjútor meus.

Convertísti planctum meum in gáudium mihi: * conscidísti saccum meum, et circumdedísti me lætítia:

Ut cantet tibi glória mea, et non compúngar: * Dómine, Deus meus, in ætérnum confitébor tibi.

\antlabel{Ant.} Exaltábo te, Dómine, quóniam suscepísti me.

\variantheading{In feriali Officio et in Festis:}

\V\ Psállite Dómino, sancti ejus. \R\ Et confitémini memóriæ sanctitátis ejus.

\variantheading{In Officio novem Lectionum}

\absolutiolabel{Absolutio} A vínculis.

\benedlabel{Bened.} 7 Evangélica léctio. 8 In Festis Domini: Divínum auxílium. In Festis B. M. V.: Cujus festum cólimus, ipsa Virgo. In Festis Sanctorum: Cujus (\textit{vel} Quorum \textit{aut} Quarum) festum cólimus. 9 Ad societátem. (Per evangélica dicta.)

\variantheading{In Officio trium Lectionum}

\absolutiolabel{Absol.} Exáudi, Dómine.

\variantheading{In Feriis, cum Homilia:}

\benedlabel{Bened.} 1 Evangélica. 2 Divínum. 3 Ad societátem.

\variantheading{In Feriis, sine Homilia:}

\benedlabel{Bened.} 1 Benedictióne. 2 Unigénitus. 3 Spíritus Sancti.

\variantheading{In Festis Sanctorum:}

\benedlabel{Bened.} 1 Ille nos benedícat qui sine fine vivit et regnat. \R\ Amen. 2 Cujus (\textit{vel} Quorum \textit{aut} Quarum). 3 Ad societátem.

% ===== END OF BATCH 9 -- Covers autumnalis.md lines 2864-3269, opening
% "Feria Secunda" (Monday) and completing its "Ad Matutinum" in
% full: the Invitatorium "Veníte, Exsultémus Dómino," the hymn
% "Somno reféctis artubus," and the three Nocturns (Psalms 13, 14,
% 16; 17/i, 17/ii, 17/iii; 19, 20, 29), each followed by its
% condition-specific versicle/Absolutio/Benedictiones apparatus
% (the ferial psalter prints these only "In Festis novem
% Lectionum"/"In Officio novem Lectionum" etc., since on an ordinary
% ferial day with three lessons only, most of this apparatus is
% omitted per the rubric given at the start of this Nocturnal set).
% Several antiphons and psalm-incipits in the Second and Third
% Nocturns carry multiple chant-pointing marks in a single line --
% asterisk, dagger, and even double-dagger together (e.g. "Dómine, *
% † in virtúte tua lætábitur rex. ‡") -- each rendered in its exact
% source position as *, \textdagger\, and \textdaggerdbl\
% respectively, with no marks added, dropped, or reordered. The
% compact multi-option "Bened." line for the ninth lesson ("8 In
% Festis Domini: Divínum auxílium. In Festis B. M. V.: ... In Festis
% Sanctorum: ...") is transcribed as one continuous \benedlabel{}
% argument exactly as the source runs it together. The document
% continues with "Ad Laudes I" (autumnalis.md line 3270), the first of
% two alternate ferial Lauds schemes =====
\hourheading{Ad Laudes I}

\rubric{In Feriis extra Vigílias commúnes occurréntibus, atque in Festis.}

(Pater noster et Ave.)

\V\ Deus, in adjutórium.

\antlabel{Ant.} Jubiláte * Deo in voce exsultatiónis.

\rubric{¶ Antíphonæ, in hac et áliis Fériis apposítæ, dicúntur semper in Officio feriáli et in Festis sive trium sive novem Lectiónum, quando juxta Rúbricas suméndi sint Psalmi de Féria.}

\psalmheading{Psalmus 46}

\lettrine[lines=2]{O}{MNES} gentes, pláudite mánibus: * jubiláte Deo in voce exsultatiónis.

Quóniam Dóminus excélsus, terríbilis: * Rex magnus super omnem terram.

Subjécit pópulos nobis: * et gentes sub pédibus nostris.

Elégit nobis hereditátem suam: * spéciem Jacob, quam diléxit.

Ascéndit Deus in júbilo: * et Dóminus in voce tubæ.

Psállite Deo nostro, psállite: * psállite Regi nostro, psállite.

Quóniam Rex omnis terræ Deus: * psállite sapiénter.

Regnábit Deus super gentes: * Deus sedet super sedem sanctam suam.

Príncipes populórum congregáti sunt cum Deo Abraham: * quóniam dii fortes terræ veheménter eleváti sunt.

\antlabel{Ant.} Jubiláte Deo in voce exsultatiónis.

\antlabel{Ant.} Inténde * voci oratiónis meæ, Rex meus et Deus meus.

\psalmheading{Psalmus 5}

\lettrine[lines=2]{V}{ERBA} mea áuribus pércipe, Dómine, * intéllige clamórem meum.

Inténde voci oratiónis meæ, * Rex meus et Deus meus.

Quóniam ad te orábo: * Dómine, mane exáudies vocem meam.

Mane astábo tibi et vidébo: * quóniam non Deus volens iniquitátem tu es.

Neque habitábit juxta te malígnus: * neque permanébunt injústi ante óculos tuos.

Odísti omnes, qui operántur iniquitátem: * perdes omnes, qui loquúntur mendácium.

Virum sánguinum et dolósum abominábitur Dóminus: * ego autem in multitúdine misericórdiæ tuæ.

Introíbo in domum tuam: * adorábo ad templum sanctum tuum in timóre tuo.

Dómine, deduc me in justítia tua: * propter inimícos meos dírige in conspéctu tuo viam meam.

Quóniam non est in ore eórum véritas: * cor eórum vanum est.

Sepúlcrum patens est guttur eórum, linguis suis dolóse agébant, * júdica illos, Deus.

Décidant a cogitatiónibus suis, secúndum multitúdinem impietátum eórum expélle eos, * quóniam irritavérunt te, Dómine.

Et læténtur omnes, qui sperant in te, * in ætérnum exsultábunt: et habitábis in eis.

Et gloriabúntur in te omnes, qui díligunt nomen tuum, * quóniam tu benedíces justo.

Dómine, ut scuto bonæ voluntátis tuæ * coronásti nos.

\antlabel{Ant.} Inténde voci oratiónis meæ, Rex meus et Deus meus.

\antlabel{Ant.} Deus majestátis * intónuit: afférte glóriam nómini ejus.

\psalmheading{Psalmus 28}

\lettrine[lines=2]{A}{FFERTE} Dómino, fílii Dei: * afférte Dómino fílios aríetum.

Afférte Dómino glóriam et honórem, afférte Dómino glóriam nómini ejus: * adoráte Dóminum in átrio sancto ejus.

Vox Dómini super aquas, Deus majestátis intónuit: * Dóminus super aquas multas.

Vox Dómini in virtúte: * vox Dómini in magnificéntia.

Vox Dómini confringéntis cedros: * et confrínget Dóminus cedros Líbani:

Et commínuet eas tamquam vítulum Líbani: * et diléctus quemádmodum fílius unicórnium.

Vox Dómini intercidéntis flammam ignis: * vox Dómini concutiéntis desértum: et commovébit Dóminus desértum Cades.

Vox Dómini præparántis cervos, et revelábit condénsa: * et in templo ejus omnes dicent glóriam.

Dóminus dilúvium inhabitáre facit: * et sedébit Dóminus Rex in ætérnum.

Dóminus virtútem pópulo suo dabit: * Dóminus benedícet pópulo suo in pace.

\antlabel{Ant.} Deus majestátis intónuit: afférte glóriam nómini ejus.

\antlabel{Ant.} Laudámus nomen tuum * ínclitum, Deus noster.

\canticumheading{Canticum David}{1 Par. 29, 10-13}

\lettrine[lines=2]{B}{ENEDICTUS} es, Dómine, Deus Israël patris nostri, * ab ætérno in ætérnum.

Tua est, Dómine, magnificéntia, et poténtia, * et glória, atque victória:

Et tibi laus: * cuncta enim quæ in cælo sunt, et in terra, tua sunt:

Tuum, Dómine, regnum, * et tu es super omnes príncipes.

Tuæ divítiæ, et tua est glória: * tu domináris ómnium.

In manu tua virtus et poténtia: * in manu tua magnitúdo, et impérium ómnium.

Nunc ígitur, Deus noster, confitémur tibi, * et laudámus nomen tuum ínclitum.

\antlabel{Ant.} Laudámus nomen tuum ínclitum, Deus noster.

\antlabel{Ant.} Laudáte * \textdagger\ Dóminum, omnes gentes. \textdaggerdbl

\psalmheading{Psalmus 116}

\lettrine[lines=2]{L}{AUDATE} \textdagger\ Dóminum, omnes gentes: * \textdaggerdbl\ laudáte eum, omnes pópuli:

Quóniam confirmáta est super nos misericórdia ejus: * et véritas Dómini manet in ætérnum.

\antlabel{Ant.} Laudáte Dóminum, omnes gentes.

\variantheading{In Officio feriali}

\rubric{Capítulum, Hymnus et Versus dicúntur quotídie de Féria curránti, ut in Psaltério hic et álibi appónuntur; et inde súmitur étiam Ant. ad Benedíctus.}

\capitulumlabel{Capit.}{Rom. 13, 12-13}

\lettrine[lines=2]{N}{OX} præcéssit, dies autem appropinquávit. Abjiciámus ergo ópera tenebrárum, et induámur arma lucis. Sicut in die honéste ambulémus.

\hymnusheading{Hymnus}

\lettrine[lines=2]{S}{PLENDOR} Patérnæ glóriæ,\\
De luce lucem próferens,\\
Lux lucis, et fons lúminis,\\
Diem dies illúminans:

Verúsque sol illábere,\\
Micans nitóre pérpeti:\\
Jubárque Sancti Spíritus\\
Infúnde nostris sénsibus.

Votis vocémus et Patrem,\\
Patrem poténtis grátiæ,\\
Patrem perénnis glóriæ:\\
Culpam reléget lúbricam.

Confírmet actus strénuos:\\
Dentes retúndat ínvidi:\\
Casus secúndet ásperos:\\
Agénda recte dírigat.

Mentem gubérnet et regat:\\
Sit pura nobis cástitas:\\
Fides calóre férveat,\\
Fraudis venéna nésciat.

Christúsque nobis sit cibus,\\
Potúsque noster sit fides:\\
Læti bibámus sóbriam\\
Profusiónem Spíritus.

Lætus dies hic tránseat:\\
Pudor sit ut dilúculum:\\
Fides velut merídies:\\
Crepúsculum mens nésciat.

Auróra lucem próvehit,\\
Cum luce nobis pródeat\\
In Patre totus Fílius,\\
Et totus in Verbo Pater.

Deo Patri sit glória,\\
Ejúsque soli Fílio,\\
Cum Spíritu Paráclito,\\
Nunc et per omne sæculum. Amen.

\V\ Repléti sumus mane misericórdia tua. \R\ Exsultávimus, et delectáti sumus.

\antlabel{Ad Bened. Ant.} Benedíctus * \textdagger\ Dóminus, Deus Israël, quia visitávit et liberávit nos.

\rubric{Réliqua ut in Ordinário. 9}

\hourheading{Ad Laudes II}

\rubric{In Vigíliis commúnibus, quando fit Offícium de Féria.}

(Pater noster et Ave.)

\V\ Deus, in adjutórium.

\antlabel{Ant.} Miserére. \textdagger

\psalmheading{Psalmus 50}

\lettrine[lines=2]{M}{ISERERE} \textdagger\ mei, Deus: * secúndum magnam misericórdiam tuam.

Et secúndum multitúdinem miseratiónum tuárum, * dele iniquitátem meam.

Amplius lava me ab iniquitáte mea: * et a peccáto meo munda me.

Quóniam iniquitátem meam ego cognósco: * et peccátum meum contra me est semper.

Tibi soli peccávi, et malum coram te feci: * ut justificéris in sermónibus tuis, et vincas cum judicáris.

Ecce enim, in iniquitátibus concéptus sum: * et in peccátis concépit me mater mea.

Ecce enim, veritátem dilexísti: * incérta et occúlta sapiéntiæ tuæ manifestásti mihi.

Aspérges me hyssópo, et mundábor: * lavábis me, et super nivem dealbábor.

Audítui meo dabis gáudium et lætítiam: * et exsultábunt ossa humiliáta.

Avérte fáciem tuam a peccátis meis: * et omnes iniquitátes meas dele.

Cor mundum crea in me, Deus: * et spíritum rectum ínnova in viscéribus meis.

Ne projícias me a fácie tua: * et spíritum sanctum tuum ne áuferas a me.

Redde mihi lætítiam salutáris tui: * et spíritu principáli confírma me.

Docébo iníquos vias tuas: * et ímpii ad te converténtur.

Líbera me de sanguínibus, Deus, Deus salútis meæ: * et exsultábit lingua mea justítiam tuam.

Dómine, lábia mea apéries: * et os meum annuntiábit laudem tuam.

Quóniam si voluísses sacrifícium, dedíssem útique: * holocáustis non delectáberis.

Sacrifícium Deo spíritus contribulátus: * cor contrítum, et humiliátum, Deus, non despícies.

Benígne fac, Dómine, in bona voluntáte tua Sion: * ut ædificéntur muri Jerúsalem.

Tunc acceptábis sacrifícium justítiæ, oblatiónes, et holocáusta: * tunc impónent super altáre tuum vítulos.

\antlabel{Ant.} Miserére mei, Deus, secúndum magnam misericórdiam tuam.

\antlabel{Ant.} Deduc me.

\psalmheading{Psalmus 5}

\lettrine[lines=2]{V}{ERBA} mea áuribus pércipe, Dómine, * intéllige clamórem meum.

Inténde voci oratiónis meæ, * Rex meus et Deus meus.

Quóniam ad te orábo: * Dómine, mane exáudies vocem meam.

Mane astábo tibi et vidébo: * quóniam non Deus volens iniquitátem tu es.

Neque habitábit juxta te malígnus: * neque permanébunt injústi ante óculos tuos.

Odísti omnes, qui operántur iniquitátem: * perdes omnes, qui loquúntur mendácium.

Virum sánguinum et dolósum abominábitur Dóminus: * ego autem in multitúdine misericórdiæ tuæ.

Introíbo in domum tuam: * adorábo ad templum sanctum tuum in timóre tuo.

Dómine, deduc me in justítia tua: * propter inimícos meos dírige in conspéctu tuo viam meam.

Quóniam non est in ore eórum véritas: * cor eórum vanum est.

Sepúlcrum patens est guttur eórum, linguis suis dolóse agébant, * júdica illos, Deus.

Décidant a cogitatiónibus suis, secúndum multitúdinem impietátum eórum expélle eos, * quóniam irritavérunt te, Dómine.

Et læténtur omnes, qui sperant in te, * in ætérnum exsultábunt: et habitábis in eis.

Et gloriabúntur in te omnes, qui díligunt nomen tuum, * quóniam tu benedíces justo.

Dómine, ut scuto bonæ voluntátis tuæ * coronásti nos.

\antlabel{Ant.} Deduc me in justítia tua, Dómine.

\antlabel{Ant.} Dóminus dabit virtútem.

\psalmheading{Psalmus 28}

\lettrine[lines=2]{A}{FFERTE} Dómino, fílii Dei: * afférte Dómino fílios aríetum.

Afférte Dómino glóriam et honórem, afférte Dómino glóriam nómini ejus: * adoráte Dóminum in átrio sancto ejus.

Vox Dómini super aquas, Deus majestátis intónuit: * Dóminus super aquas multas.

Vox Dómini in virtúte: * vox Dómini in magnificéntia.

Vox Dómini confringéntis cedros: * et confrínget Dóminus cedros Líbani:

Et commínuet eas tamquam vítulum Líbani: * et diléctus quemádmodum fílius unicórnium.

Vox Dómini intercidéntis flammam ignis: * vox Dómini concutiéntis desértum: et commovébit Dóminus desértum Cades.

Vox Dómini præparántis cervos, et revelábit condénsa: * et in templo ejus omnes dicent glóriam.

Dóminus dilúvium inhabitáre facit: * et sedébit Dóminus Rex in ætérnum.

Dóminus virtútem pópulo suo dabit: * Dóminus benedícet pópulo suo in pace.

\antlabel{Ant.} Dóminus dabit virtútem et benedícet pópulo suo in pace.

\antlabel{Ant.} Convérsus est furor tuus.

\canticumheading{Canticum Isaiæ}{Is. 12, 1-6}

\lettrine[lines=2]{C}{ONFITEBOR} tibi, Dómine, quóniam irátus es mihi: * convérsus est furor tuus, et consolátus es me.

Ecce Deus salvátor meus, * fiduciáliter agam, et non timébo:

Quia fortitúdo mea, et laus mea Dóminus, * et factus est mihi in salútem.

Hauriétis aquas in gáudio de fóntibus Salvatóris: * et dicétis in die illa: Confitémini Dómino, et invocáte nomen ejus:

Notas fácite in pópulis adinventiónes ejus: * mementóte quóniam excélsum est nomen ejus.

Cantáte Dómino quóniam magnífice fecit: * annuntiáte hoc in univérsa terra.

Exsúlta, et lauda, habitátio Sion: * quia magnus in médio tui Sanctus Israël.

\antlabel{Ant.} Convérsus est furor tuus, Dómine, et consolátus es me.

\antlabel{Ant.} Laudáte. \textdagger

\psalmheading{Psalmus 116}

\lettrine[lines=2]{L}{AUDATE} \textdagger\ Dóminum, omnes gentes: * \textdaggerdbl\ laudáte eum, omnes pópuli:

Quóniam confirmáta est super nos misericórdia ejus: * et véritas Dómini manet in ætérnum.

\antlabel{Ant.} Laudáte Dóminum, quóniam confirmáta est super nos misericórdia ejus.

\capitulumlabel{Capit.}{Rom. 13, 12-13}

\lettrine[lines=2]{N}{OX} præcéssit, dies autem appropinquávit. Abjiciámus ergo ópera tenebrárum, et induámur arma lucis. Sicut in die honéste ambulémus.

\hymnusheading{Hymnus}

\lettrine[lines=2]{S}{PLENDOR} Patérnæ glóriæ,\\
De luce lucem próferens,\\
Lux lucis, et fons lúminis,\\
Diem dies illúminans:

Verúsque sol illábere,\\
Micans nitóre pérpeti:\\
Jubárque Sancti Spíritus\\
Infúnde nostris sénsibus.

Votis vocémus et Patrem,\\
Patrem poténtis grátiæ,\\
Patrem perénnis glóriæ:\\
Culpam reléget lúbricam.

Confírmet actus strénuos:\\
Dentes retúndat ínvidi:\\
Casus secúndet ásperos:\\
Agénda recte dírigat.

Mentem gubérnet et regat:\\
Sit pura nobis cástitas:\\
Fides calóre férveat,\\
Fraudis venéna nésciat.

Christúsque nobis sit cibus,\\
Potúsque noster sit fides:\\
Læti bibámus sóbriam\\
Profusiónem Spíritus.

Lætus dies hic tránseat:\\
Pudor sit ut dilúculum:\\
Fides velut merídies:\\
Crepúsculum mens nésciat.

Auróra lucem próvehit,\\
Cum luce nobis pródeat\\
In Patre totus Fílius,\\
Et totus in Verbo Pater.

Deo Patri sit glória,\\
Ejúsque soli Fílio,\\
Cum Spíritu Paráclito,\\
Nunc et per omne sæculum. Amen.

\V\ Repléti sumus mane misericórdia tua. \R\ Exsultávimus, et delectáti sumus.

\antlabel{Ad Bened. Ant.} Benedíctus * \textdagger\ Dóminus, Deus Israël, quia visitávit et liberávit nos.

\variantheading{Preces feriales}

KYRIE, eléison. Christe, eléison. Kýrie, eléison.

Pater noster \rubric{quod dícitur a solo Hebdomadário totum clara voce, usque ad} \V\ Et ne nos indúcas in tentatiónem. \R\ Sed líbera nos a malo. \V\ Ego dixi: Dómine, miserére mei. \R\ Sana ánimam meam, quia peccávi tibi. \V\ Convértere, Dómine, úsquequo? \R\ Et deprecábilis esto super servos tuos. \V\ Fiat misericórdia tua, Dómine, super nos. \R\ Quemádmodum sperávimus in te. \V\ Sacerdótes tui induántur justítiam. \R\ Et sancti tui exsúltent. \V\ Orémus pro beatíssimo Papa nostro N. \R\ Dóminus consérvet eum, et vivíficet eum, et beátum fáciat eum in terra, et non tradat eum in ánimam inimicórum ejus. \V\ Orémus et pro Antístite nostro N. \R\ Stet et pascat in fortitúdine tua, Dómine, in sublimitáte nóminis tui. \V\ Dómine, salvum fac regem. \R\ Et exáudi nos in die, qua invocavérimus te. \V\ Salvum fac pópulum tuum, Dómine, et bénedic hereditáti tuæ. \R\ Et rege eos, et extólle illos usque in ætérnum. \V\ Meménto Congregatiónis tuæ. \R\ Quam possedísti ab inítio. \V\ Fiat pax in virtúte tua. \R\ Et abundántia in túrribus tuis. \V\ Orémus pro benefactóribus nostris. \R\ Retribúere dignáre, Dómine, ómnibus, nobis bona faciéntibus propter nomen tuum, vitam ætérnam. Amen. \V\ Orémus pro fidélibus defúnctis. \R\ Réquiem ætérnam dona eis, Dómine, et lux perpétua lúceat eis. \V\ Requiéscant in pace. \R\ Amen. \V\ Pro frátribus nostris abséntibus. \R\ Salvos fac servos tuos, Deus meus, sperántes in te. \V\ Pro afflíctis et captívis. \R\ Líbera eos, Deus Israël, ex ómnibus tribulatiónibus suis. \V\ Mitte eis, Dómine, auxílium de sancto. \R\ Et de Sion tuére eos. \V\ Dómine, Deus virtútum, convérte nos. \R\ Et osténde fáciem tuam, et salvi érimus. \V\ Exsúrge, Christe, ádjuva nos. \R\ Et líbera nos propter nomen tuum. \V\ Dómine, exáudi oratiónem meam. \R\ Et clamor meus ad te véniat.

\rubric{Réliqua ut in Ordinário. 9}

% ===== END OF BATCH 10 -- Covers autumnalis.md lines 3270-3672, the
% two alternate ferial Lauds schemes, "Ad Laudes I" (for ordinary
% ferias and feasts) and "Ad Laudes II" (for common Vigils when the
% Office of the Feria is said). Both share the same five psalms
% (46, 5, 28, Canticum, 116) with different antiphon sets, the same
% Capitulum "Nox præcéssit," and the identical hymn "Splendor
% Paternæ glóriæ" -- all transcribed fresh in both places per the
% governing rule, even though the Hymnus and Capitulum text is
% word-for-word identical between the two schemes in autumnalis.md
% itself. Ad Laudes II additionally carries the "Preces feriales"
% litany, transcribed in full; "Kyrie, eléison" at its head is left
% as plain text with no lettrine, matching precedent in verna.tex
% and hiemalis.tex, since it is not one of the five lettrine genres.
% The document continues with "Ad Primam" (autumnalis.md line 3674).
% =====
\hourheading{Ad Primam}

Pater, Ave et Credo.

\V\ Deus, in adjutórium.

\hymnusheading{Hymnus}

\lettrine[lines=2]{J}{AM} lucis orto sídere,\\
Deum precémur súpplices,\\
Ut in diúrnis áctibus\\
Nos servet a nocéntibus.

Linguam refrénans témperet,\\
Ne litis horror ínsonet:\\
Visum fovéndo cóntegat,\\
Ne vanitátes háuriat.

Sint pura cordis íntima,\\
Absístat et vecórdia:\\
Carnis terat supérbiam\\
Potus cibíque párcitas.

Ut, cum dies abscésserit\\
Noctémque sors redúxerit,\\
Mundi per abstinéntiam\\
Ipsi canámus glóriam.

Deo Patri sit glória,\\
Ejúsque soli Fílio,\\
Cum Spíritu Paráclito,\\
Nunc et per omne sæculum.\\
Amen.

\antlabel{Ant.} Innocens mánibus.

\rubric{¶ Antiphonæ, in hac et áliis Fériis ad Primam, Tértiam, Sextam et Nonam apposítæ, dicúntur semper in Officio feriáli et in Festis sive trium sive novem Lectiónum, quando juxta Rúbricas suméndi sint Psalmi de Féria.}

\psalmheading{Psalmus 23}

\lettrine[lines=2]{D}{OMINI} est terra, et plenitúdo ejus: * orbis terrárum, et univérsi qui hábitant in eo.

Quia ipse super mária fundávit eum: * et super flúmina præparávit eum.

Quis ascéndet in montem Dómini? * aut quis stabit in loco sancto ejus?

Innocens mánibus et mundo corde, * qui non accépit in vano ánimam suam, nec jurávit in dolo próximo suo.

Hic accípiet benedictiónem a Dómino: * et misericórdiam a Deo, salutári suo.

Hæc est generátio quæréntium eum, * quæréntium fáciem Dei Jacob.

Attóllite portas, príncipes, vestras, et elevámini, portæ æternáles: * et introíbit Rex glóriæ.

Quis est iste Rex glóriæ? * Dóminus fortis et potens: Dóminus potens in prœlio.

Attóllite portas, príncipes, vestras, et elevámini, portæ æternáles: * et introíbit Rex glóriæ.

Quis est iste Rex glóriæ? * Dóminus virtútum ipse est Rex glóriæ.

\psalmheading{Psalmus 18, i}

\lettrine[lines=2]{C}{ÆLI} enárrant glóriam Dei, * et ópera mánuum ejus annúntiat firmaméntum.

Dies diéi erúctat verbum, * et nox nocti índicat sciéntiam.

Non sunt loquélæ, neque sermónes, * quorum non audiántur voces eórum.

In omnem terram exívit sonus eórum: * et in fines orbis terræ verba eórum.

In sole pósuit tabernáculum suum: * et ipse tamquam sponsus procédens de thálamo suo:

Exsultávit ut gigas ad curréndam viam, * a summo cælo egréssio ejus:

Et occúrsus ejus usque ad summum ejus: * nec est qui se abscóndat a calóre ejus.

\psalmheading{Psalmus 18, ii}

\lettrine[lines=2]{L}{EX} Dómini immaculáta, convértens ánimas: * testimónium Dómini fidéle, sapiéntiam præstans párvulis.

Justítiæ Dómini rectæ, lætificántes corda: * præcéptum Dómini lúcidum, illúminans óculos.

Timor Dómini sanctus, pérmanens in sæculum sæculi: * judícia Dómini vera, justificáta in semetípsa.

Desiderabília super aurum et lápidem pretiósum multum: * et dulcióra super mel et favum.

Étenim servus tuus custódit ea: * in custodiéndis illis retribútio multa.

Delícta quis intélligit? ab occúltis meis munda me: * et ab aliénis parce servo tuo.

Si mei non fúerint domináti, tunc immaculátus ero: * et emundábor a delícto máximo.

Et erunt ut compláceant elóquia oris mei: * et meditátio cordis mei in conspéctu tuo semper.

Dómine, adjútor meus, * et redémptor meus.

\rubric{¶ Quando ad Laudes dictus fúerit Ps. 50 Miserére, hic subjúngitur sequens Ps. 46, qui in Laudibus locum non hábuit.}

\psalmheading{Psalmus 46}

\lettrine[lines=2]{O}{MNES} gentes, pláudite mánibus: * jubiláte Deo in voce exsultatiónis.

Quóniam Dóminus excélsus, terríbilis: * Rex magnus super omnem terram.

Subjécit pópulos nobis: * et gentes sub pédibus nostris.

Elégit nobis hereditátem suam: * spéciem Jacob, quam diléxit.

Ascéndit Deus in júbilo: * et Dóminus in voce tubæ.

Psállite Deo nostro, psállite: * psállite Regi nostro, psállite.

Quóniam Rex omnis terræ Deus: * psállite sapiénter.

Regnábit Deus super gentes: * Deus sedet super sedem sanctam suam.

Príncipes populórum congregáti sunt cum Deo Abraham: * quóniam dii fortes terræ veheménter eleváti sunt.

\rubric{Secus vero, tribus Psalmis explétis, statim dícitur Ant.}

\antlabel{Ant.} Innocens mánibus et mundo corde ascéndet in montem Dómini.

\rubric{In Officio Festi vel Octávæ:}

\capitulumlabel{Capitulum}{1 Tim. 1, 17}

\lettrine[lines=2]{R}{EGI} sæculórum immortáli et invisíbili, soli Deo honor et glória in sæcula sæculórum. Amen.

\rubric{In ómnibus autem Fériis et in Vigíliis commúnibus:}

\capitulumlabel{Capitulum}{Zach. 8, 19}

\lettrine[lines=2]{P}{ACEM} et veritátem dilígite, ait Dóminus omnípotens.

\R\ br. Christe, Fili Dei vivi, * Miserére nobis. Christe. \V\ \textit{Qui sedes ad déxteram Patris.} Miserére nobis. Glória Patri. Christe. \V\ Exsúrge, Christe, ádjuva nos. \R\ Et líbera nos propter nomen tuum.

\variantheading{Preces dominicales et feriales (si dicendæ sint)}

KYRIE, eléison. Christe, eléison. Kýrie, eléison.

Pater noster \rubric{secreto usque ad} \V\ Et ne nos indúcas in tentatiónem. \R\ Sed líbera nos a malo.

Credo in Deum \rubric{secreto usque ad} \V\ Carnis resurrectiónem. \R\ Vitam ætérnam. Amen. \V\ Et ego ad te, Dómine, clamávi. \R\ Et mane orátio mea prævéniet te. \V\ Repleátur os meum laude. \R\ Ut cantem glóriam tuam, tota die magnitúdinem tuam. \V\ Dómine, avérte fáciem tuam a peccátis meis. \R\ Et omnes iniquitátes meas dele. \V\ Cor mundum crea in me, Deus. \R\ Et spíritum rectum ínnova in viscéribus meis. \V\ Ne projícias me a fácie tua. \R\ Et spíritum sanctum tuum ne áuferas a me. \V\ Redde mihi lætítiam salutáris tui. \R\ Et spíritu principáli confírma me.

\variantheading{In Précibus feriálibus tantum}

\V\ Eripe me, Dómine, ab hómine malo. \R\ A viro iníquo éripe me. \V\ Eripe me de inimícis meis, Deus meus. \R\ Et ab insurgéntibus in me líbera me. \V\ Eripe me de operántibus iniquitátem. \R\ Et de viris sánguinum salva me. \V\ Sic psalmum dicam nómini tuo in sæculum sæculi. \R\ Ut reddam vota mea de die in diem. \V\ Exáudi nos, Deus, salutáris noster. \R\ Spes ómnium fínium terræ, et in mari longe. \V\ Deus, in adjutórium meum inténde. \R\ Dómine, ad adjuvándum me festína. \V\ Sanctus Deus, Sanctus fortis, Sanctus immortális. \R\ Miserére nobis. \V\ Bénedic, ánima mea, Dómino. \R\ Et ómnia, quæ intra me sunt, nómini sancto ejus. \V\ Bénedic, ánima mea, Dómino. \R\ Et noli oblivísci omnes retributiónes ejus. \V\ Qui propitiátur ómnibus iniquitátibus tuis. \R\ Qui sanat omnes infirmitátes tuas. \V\ Qui rédimit de intéritu vitam tuam. \R\ Qui corónat te in misericórdia et miseratiónibus. \V\ Qui replet in bonis desidérium tuum. \R\ Renovábitur ut áquilæ juvéntus tua.

\variantheading{In utrisque Précibus}

\V\ Adjutórium nostrum in nómine Dómini. \R\ Qui fecit cælum et terram.

\formulalabel{Confíteor. Misereátur. Indulgéntiam.}

\V\ Dignáre, Dómine, die isto. \R\ Sine peccáto nos custodíre. \V\ Miserére nostri, Dómine. \R\ Miserére nostri. \V\ Fiat misericórdia tua, Dómine, super nos. \R\ Quemádmodum sperávimus in te. \V\ Dómine, exáudi oratiónem meam. \R\ Et clamor meus ad te véniat.

\V\ Dóminus vobíscum.

Orémus. \oratiolabel{Oratio}

\lettrine[lines=2]{D}{OMINE}, Deus omnípotens, qui ad princípium hujus diéi nos perveníre fecísti: tua nos hódie salva virtúte; ut in hac die ad nullum declinémus peccátum, sed semper ad tuam justítiam faciéndam nostra procédant elóquia, dirigántur cogitatiónes et ópera. Per Dóminum. \R\ Amen.

\V\ Dóminus vobíscum. \V\ Benedicámus Dómino.

\rubric{In Choro Martyrológium.}

\V\ Pretiósa in conspéctu Dómini. \R\ Mors Sanctórum ejus.

\rubric{Et absolúte, sine} Orémus:

\oratiolabel{Oratio}

\lettrine[lines=2]{S}{ANCTA} María, et omnes Sancti intercédant pro nobis ad Dóminum, ut nos mereámur ab eo adjuvári et salvári, qui vivit et regnat in sæcula sæculórum. \R\ Amen.

\V\ Deus, in adjutórium meum inténde. \R\ Dómine, ad adjuvándum me festína. \rubric{Et dícitur ter; último additur:} Glória Patri. Sicut erat.

Kýrie, eléison. Christe, eléison. Kýrie, eléison.

Pater noster \rubric{secreto usque ad} \V\ Et ne nos indúcas in tentatiónem. \R\ Sed líbera nos a malo. \V\ Réspice in servos tuos, Dómine, et in ópera tua, et dírige fílios eórum. \R\ Et sit splendor Dómini, Dei nostri, super nos, et ópera mánuum nostrárum dírige super nos, et opus mánuum nostrárum dírige. \V\ Glória Patri. \R\ Sicut erat.

Orémus. \oratiolabel{Oratio}

\lettrine[lines=2]{D}{IRIGERE} et sanctificáre, régere et gubernáre dignáre, Dómine Deus, Rex cæli et terræ, hódie corda et córpora nostra, sensus, sermónes et actus nostros in lege tua, et in opéribus mandatórum tuórum: ut hic, et in ætérnum, te auxiliánte, salvi et líberi esse mereámur, Salvátor mundi: Qui vivis et regnas in sæcula sæculórum. \R\ Amen.

\V\ Jube, domne.

\benedlabel{Bened.} Dies et actus nostros in sua pace dispónat Dóminus omnípotens. \R\ Amen.

\variantheading{In Officio feriali}

\lettrine[lines=2]{D}{OMINUS} autem dírigat corda et córpora nostra in caritáte Dei, et patiéntia Christi. Tu autem.

\variantheading{E Proprio Sanctorum}

\variantheading{Pro Oct. Omn. Sanctorum}

\lettrine[lines=2]{B}{ENEDÍCTIO} et cláritas et sapiéntia et gratiárum áctio, honor, virtus et fortitúdo Deo nostro in sǽcula sæculórum. Amen.

\variantheading{E Communi Sanctorum}

\variantheading{Pro Apostolis}

\lettrine[lines=2]{I}{BANT} Apóstoli gaudéntes a conspéctu concílii, quóniam digni hábiti sunt pro nómine Jesu contuméliam pati.

\variantheading{Pro Martyre}

\lettrine[lines=2]{J}{USTUS} cor suum trádidit ad vigilándum dilúculo ad Dóminum, qui fecit illum, et in conspéctu Altíssimi deprecábitur.

\variantheading{Pro pluribus Martyribus}

\lettrine[lines=2]{F}{ULGEBUNT} justi, et tamquam scintíllæ in arundinéto discúrrent. Judicábunt natiónes, et dominabúntur pópulis: et regnábit Dóminus illórum in perpétuum.

\variantheading{Pro Conf. Pont.}

\lettrine[lines=2]{F}{UNGI} sacerdótio, et habére laudem in nómine ipsíus, et offérre illi incénsum dignum in odórem suavitátis.

\variantheading{Pro Conf. non Pont.}

\lettrine[lines=2]{J}{USTUM} dedúxit Dóminus per vias rectas, et osténdit illi regnum Dei, et dedit illi sciéntiam sanctórum; honestávit illum in labóribus, et complévit labóres illíus.

\variantheading{Pro Virgine Martyre}

\lettrine[lines=2]{D}{OMINE}, Deus meus, exaltásti super terram habitatiónem meam, et pro morte defluénte deprecáta sum. Invocávi Dóminum, Patrem Dómini mei, ut non derelínquat me in die tribulatiónis meæ, et in témpore superbórum sine adjutório.

\variantheading{Pro Virgine non Martyre}

\lettrine[lines=2]{O}{} \textsc{quam} pulchra est casta generátio cum claritáte! immortális est enim memória illíus: quóniam apud Deum nota est, et apud hómines.

\variantheading{Pro Martyre tantum}

\lettrine[lines=2]{L}{AUDABIT} usque ad mortem ánima mea Dóminum, quóniam éruis sustinéntes te, et líberas eos de manu angústiæ, Dómine, Deus noster.

\variantheading{Pro nec Virg. nec Mart.}

\lettrine[lines=2]{M}{ULTAE} fíliæ congregavérunt divítias: tu supergréssa es univérsas. Fallax grátia, et vana est pulchritúdo: múlier timens Dóminum, ipsa laudábitur.

\variantheading{Pro Dedicatione Ecclesiæ}

\lettrine[lines=2]{E}{T} abstérget Deus omnem lácrimam ab óculis eórum: et mors ultra non erit, neque luctus, neque clamor, neque dolor erit ultra, quia prima abiérunt. Et dixit qui sedébat in throno: Ecce, nova fácio ómnia.

\variantheading{In Officiis B. Mariæ V.}

\lettrine[lines=2]{I}{N} platéis sicut cinnamómum et bálsamum aromatízans odórem dedi: quasi myrrha elécta dedi suavitátem odóris.

\V\ Adjutórium nostrum in nómine Dómini. \R\ Qui fecit cælum et terram.

\V\ Benedícite. \R\ Deus.

\benedictiolabel{Benedictio} Dóminus nos benedícat, et ab omni malo deféndat, et ad vitam perdúcat ætérnam. Et fidélium ánimæ per misericórdiam Dei requiéscant in pace. \R\ Amen.

\rubric{Deinde dícitur tantum} Pater noster \rubric{secréto, nisi sequátur ália Hora.}

% ===== END OF BATCH 11 -- Covers autumnalis.md lines 3674-3927,
% completing "Ad Primam" for Feria Secunda in full: the hymn "Jam
% lucis orto sídere" (with its own stray OCR "\*" before the final
% doxology stanza silently dropped, the same recurring artifact
% corrected in every prior occurrence this volume), Psalms 23, 18/i,
% 18/ii, and the conditional insertion of Psalm 46 (used when Ps. 50
% was said at Lauds instead, per the rubric), the two alternate
% Capitula ("In Officio Festi vel Octávæ" / "In ómnibus autem
% Fériis"), the brief Responsory "Christe, Fili Dei vivi," the
% "Preces dominicales et feriales (si dicendæ sint)" apparatus with
% its "In Précibus feriálibus tantum" and "In utrísque Précibus"
% subdivisions (the bold "Confíteor. Misereátur. Indulgéntiam."
% rendered with \formulalabel{} per the established rule for bold
% source markup), the two closing Orationes and the Martyrologium
% versicle, and finally the full "Lectio brevis" apparatus (the
% ferial variant "Dóminus autem dírigat," followed by "E Proprio
% Sanctorum" with its single "Pro Oct. Omn. Sanctorum" heading (this
% volume's one-item Autumn substitute for aestiva.tex's two-item
% Summer list "Pro Oct. S. Joannis Bapt." / "Pro Invent. S.
% Stephani," verified via direct diff against aestiva.md and
% transcribed fresh from autumnalis.md's own text) and "E Communi
% Sanctorum" with its full run of eleven Pro-headings and "In
% Officiis B. Mariæ V.," all reusing the \variantheading{} command
% established already in this volume's own Sunday Prime, each short
% text given a two-line lettrine as a Lectio-brevis genre text). The
% embedded
% rubric "Et absolúte, sine Orémus:" is preserved exactly as the
% source interleaves red rubric and black directive text. The
% document continues with "Ad Tertiam" (autumnalis.md line 3928).
% =====
\hourheading{Ad Tertiam}

Pater noster et Ave.

\V\ Deus, in adjutórium.

\hymnusheading{Hymnus}

\lettrine[lines=2]{N}{UNC}, Sancte, nobis, Spíritus,\\
Unum Patri cum Fílio,\\
Dignáre promptus íngeri\\
Nostro refúsus péctori.

Os, lingua, mens, sensus, vigor\\
Confessiónem pérsonent,\\
Flamméscat igne cáritas,\\
Accéndat ardor próximos.

Præsta, Pater piíssime,\\
Patríque compar Unice,\\
Cum Spíritu Paráclito\\
Regnans per omne sæculum.\\
Amen.

\antlabel{Ant.} Illuminátio mea.

\psalmheading{Psalmus 26, i}

\lettrine[lines=2]{D}{OMINUS} illuminátio mea, et salus mea, * quem timébo?

Dóminus protéctor vitæ meæ, * a quo trepidábo?

Dum apprópiant super me nocéntes, * ut edant carnes meas:

Qui tríbulant me inimíci mei, * ipsi infirmáti sunt, et cecidérunt.

Si consístant advérsum me castra, * non timébit cor meum.

Si exsúrgat advérsum me prælium, * in hoc ego sperábo.

Unam pétii a Dómino, hanc requíram, * ut inhábitem in domo Dómini ómnibus diébus vitæ meæ:

Ut vídeam voluptátem Dómini, * et vísitem templum ejus.

Quóniam abscóndit me in tabernáculo suo: * in die malórum protéxit me in abscóndito tabernáculi sui.

In petra exaltávit me: * et nunc exaltávit caput meum super inimícos meos.

Circuívi, et immolávi in tabernáculo ejus hóstiam vociferatiónis: * cantábo, et psalmum dicam Dómino.

\psalmheading{Psalmus 26, ii}

\lettrine[lines=2]{E}{XAUDI}, Dómine, vocem meam, qua clamávi ad te: * miserére mei, et exáudi me.

Tibi dixit cor meum, exquisívit te fácies mea: * fáciem tuam, Dómine, requíram.

Ne avértas fáciem tuam a me: * ne declínes in ira a servo tuo.

Adjútor meus esto: * ne derelínquas me, neque despícias me, Deus, salutáris meus.

Quóniam pater meus, et mater mea dereliquérunt me: * Dóminus autem assúmpsit me.

Legem pone mihi, Dómine, in via tua: * et dírige me in sémitam rectam propter inimícos meos.

Ne tradíderis me in ánimas tribulántium me: * quóniam insurrexérunt in me testes iníqui, et mentíta est iníquitas sibi.

Credo vidére bona Dómini * in terra vivéntium.

Exspécta Dóminum, viríliter age: * et confortétur cor tuum, et sústine Dóminum.

\psalmheading{Psalmus 27}

\lettrine[lines=2]{A}{D} te, Dómine, clamábo, Deus meus, ne síleas a me: * nequándo táceas a me, et assimilábor descendéntibus in lacum.

Exáudi, Dómine, vocem deprecatiónis meæ dum oro ad te: * dum extóllo manus meas ad templum sanctum tuum.

Ne simul trahas me cum peccatóribus: * et cum operántibus iniquitátem ne perdas me.

Qui loquúntur pacem cum próximo suo: * mala autem in córdibus eórum.

Da illis secúndum ópera eórum, * et secúndum nequítiam adinventiónum ipsórum.

Secúndum ópera mánuum eórum tríbue illis: * redde retributiónem eórum ipsis.

Quóniam non intellexérunt ópera Dómini, et in ópera mánuum ejus * déstrues illos, et non ædificábis eos.

Benedíctus Dóminus: * quóniam exaudívit vocem deprecatiónis meæ.

Dóminus adjútor meus, et protéctor meus: * in ipso sperávit cor meum, et adjútus sum.

Et reflóruit caro mea: * et ex voluntáte mea confitébor ei.

Dóminus fortitúdo plebis suæ: * et protéctor salvatiónum Christi sui est.

Salvum fac pópulum tuum, Dómine, et bénedic hereditáti tuæ: * et rege eos, et extólle illos usque in ætérnum.

\antlabel{Ant.} Illuminátio mea, et salus mea Dóminus.

\variantheading{In Officio feriali}

\capitulumlabel{Capitulum}{Jer. 17, 14}

\lettrine[lines=2]{S}{ANA} me, Dómine, et sanábor: salvum me fac, et salvus ero: quóniam laus mea tu es.

\R\ br. Sana ánimam meam, * Quia peccávi tibi. Sana. \V\ Ego dixi: Dómine, miserére mei. Quia. Glória Patri. Sana. \V\ Adjútor meus esto, ne derelínquas me. \R\ Neque despícias me, Deus, salutáris meus.

\variantheading{In Vigiliis communibus:}

\variantheading{Preces feriales}

KYRIE, eléison. Christe, eléison. Kýrie, eléison.

Pater noster \rubric{secreto usque ad} \V\ Et ne nos indúcas in tentatiónem. \R\ Sed líbera nos a malo. \V\ Dómine, Deus virtútum, convérte nos. \R\ Et osténde fáciem tuam, et salvi érimus. \V\ Exsúrge, Christe, ádjuva nos. \R\ Et líbera nos propter nomen tuum. \V\ Dómine, exáudi oratiónem meam. \R\ Et clamor meus ad te véniat.

\rubric{Réliqua ut in Ordinário. 19}

\hourheading{Ad Sextam}

Pater noster et Ave.

\V\ Deus, in adjutórium.

\hymnusheading{Hymnus}

\lettrine[lines=2]{R}{ECTOR} potens, verax Deus,\\
Qui témperas rerum vices,\\
Splendóre mane illúminas,\\
Et ígnibus merídiem:

Exstíngue flammas lítium,\\
Aufer calórem nóxium,\\
Confer salútem córporum,\\
Verámque pacem córdium.

Præsta, Pater piíssime,\\
Patríque compar Unice,\\
Cum Spíritu Paráclito\\
Regnans per omne sæculum.\\
Amen.

\antlabel{Ant.} In tua justítia.

\psalmheading{Psalmus 30, i}

\lettrine[lines=2]{I}{N} te, Dómine, sperávi, non confúndar in ætérnum: * in justítia tua líbera me.

Inclína ad me aurem tuam, * accélera ut éruas me.

Esto mihi in Deum protectórem, et in domum refúgii: * ut salvum me fácias.

Quóniam fortitúdo mea, et refúgium meum es tu: * et propter nomen tuum dedúces me, et enútries me.

Edúces me de láqueo hoc, quem abscondérunt mihi: * quóniam tu es protéctor meus.

In manus tuas comméndo spíritum meum: * redemísti me, Dómine, Deus veritátis.

Odísti observántes vanitátes, * supervácue.

Ego autem in Dómino sperávi: * exsultábo, et lætábor in misericórdia tua.

Quóniam respexísti humilitátem meam, * salvásti de necessitátibus ánimam meam.

Nec conclusísti me in mánibus inimíci: * statuísti in loco spatióso pedes meos.

\psalmheading{Psalmus 30, ii}

\lettrine[lines=2]{M}{ISERERE} mei, Dómine, quóniam tríbulor: * conturbátus est in ira óculus meus, ánima mea, et venter meus:

Quóniam defécit in dolóre vita mea: * et anni mei in gemítibus.

Infirmáta est in paupertáte virtus mea: * et ossa mea conturbáta sunt.

Super omnes inimícos meos factus sum oppróbrium et vicínis meis valde: * et timor notis meis.

Qui vidébant me, foras fugérunt a me: * oblivióni datus sum, tamquam mórtuus a corde.

Factus sum tamquam vas pérditum: * quóniam audívi vituperatiónem multórum commorántium in circúitu.

In eo dum convenírent simul advérsum me, * accípere ánimam meam consiliáti sunt.

Ego autem in te sperávi, Dómine: * dixi: Deus meus es tu: in mánibus tuis sortes meæ.

Eripe me de manu inimicórum meórum, * et a persequéntibus me.

Illústra fáciem tuam super servum tuum, salvum me fac in misericórdia tua: * Dómine, non confúndar, quóniam invocávi te.

Erubéscant ímpii, et deducántur in inférnum: * muta fiant lábia dolósa.

Quæ loquúntur advérsus justum iniquitátem, * in supérbia, et in abusióne.

\psalmheading{Psalmus 30, iii}

\lettrine[lines=2]{Q}{UAM} magna multitúdo dulcédinis tuæ, Dómine, * quam abscondísti timéntibus te.

Perfecísti eis, qui sperant in te, * in conspéctu filiórum hóminum.

Abscóndes eos in abscóndito faciéi tuæ, * a conturbatióne hóminum.

Próteges eos in tabernáculo tuo * a contradictióne linguárum.

Benedíctus Dóminus: * quóniam mirificávit misericórdiam suam mihi in civitáte muníta.

Ego autem dixi in excéssu mentis meæ: * Projéctus sum a fácie oculórum tuórum.

Ideo exaudísti vocem oratiónis meæ, * dum clamárem ad te.

Dilígite Dóminum, omnes sancti ejus: * quóniam veritátem requíret Dóminus, et retríbuet abundánter faciéntibus supérbiam.

Viríliter ágite, et confortétur cor vestrum, * omnes, qui sperátis in Dómino.

\antlabel{Ant.} In tua justítia líbera me, Dómine.

\variantheading{In Officio feriali}

\capitulumlabel{Capitulum}{Rom. 13, 8}

\lettrine[lines=2]{N}{EMINI} quidquam debeátis, nisi ut ínvicem diligátis: qui enim díligit próximum, legem implévit.

\R\ br. Benedícam Dóminum * In omni témpore. Benedícam. \V\ Semper laus ejus in ore meo. In. Glória Patri. Benedícam. \V\ Dóminus regit me, et nihil mihi déerit. \R\ In loco páscuæ ibi me collocávit.

\variantheading{In Vigiliis communibus:}

\variantheading{Preces feriales}

KYRIE, eléison. Christe, eléison. Kýrie, eléison.

Pater noster \rubric{secreto usque ad} \V\ Et ne nos indúcas in tentatiónem. \R\ Sed líbera nos a malo. \V\ Dómine, Deus virtútum, convérte nos. \R\ Et osténde fáciem tuam, et salvi érimus. \V\ Exsúrge, Christe, ádjuva nos. \R\ Et líbera nos propter nomen tuum. \V\ Dómine, exáudi oratiónem meam. \R\ Et clamor meus ad te véniat.

\rubric{Réliqua ut in Ordinário. 20}

\hourheading{Ad Nonam}

Pater noster et Ave.

\V\ Deus, in adjutórium.

\hymnusheading{Hymnus}

\lettrine[lines=2]{R}{ERUM}, Deus, tenax vigor,\\
Immótus in te pérmanens,\\
Lucis diúrnæ témpora\\
Succéssibus detérminans:

Largíre lumen véspere,\\
Quo vita nusquam décidat,\\
Sed prǽmium mortis sacræ\\
Perénnis instet glória.

Præsta, Pater piíssime,\\
Patríque compar Unice,\\
Cum Spíritu Paráclito\\
Regnans per omne sæculum.\\
Amen.

\antlabel{Ant.} Exsultáte, justi.

\psalmheading{Psalmus 31}

\lettrine[lines=2]{B}{EATI} quorum remíssæ sunt iniquitátes: * et quorum tecta sunt peccáta.

Beátus vir, cui non imputávit Dóminus peccátum, * nec est in spíritu ejus dolus.

Quóniam tácui, inveteravérunt ossa mea, * dum clamárem tota die.

Quóniam die ac nocte graváta est super me manus tua: * convérsus sum in ærúmna mea, dum confígitur spina.

Delíctum meum cógnitum tibi feci: * et injustítiam meam non abscóndi.

Dixi: Confitébor advérsum me injustítiam meam Dómino: * et tu remisísti impietátem peccáti mei.

Pro hac orábit ad te omnis sanctus, * in témpore opportúno.

Verúmtamen in dilúvio aquárum multárum, * ad eum non approximábunt.

Tu es refúgium meum a tribulatióne, quæ circúmdedit me: * exsultátio mea, érue me a circumdántibus me.

Intelléctum tibi dabo, et ínstruam te in via hac, qua gradiéris: * firmábo super te óculos meos.

Nolíte fíeri sicut equus et mulus, * quibus non est intelléctus.

In camo et freno maxíllas eórum constrínge, * qui non appróximant ad te.

Multa flagélla peccatóris, * sperántem autem in Dómino misericórdia circúmdabit.

Lætámini in Dómino et exsultáte, justi, * et gloriámini, omnes recti corde.

\psalmheading{Psalmus 32, i}

\lettrine[lines=2]{E}{XSULTATE}, justi, in Dómino: * rectos decet collaudátio.

Confitémini Dómino in cíthara: * in psaltério decem chordárum psállite illi.

Cantáte ei cánticum novum: * bene psállite ei in vociferatióne.

Quia rectum est verbum Dómini, * et ómnia ópera ejus in fide.

Díligit misericórdiam et judícium: * misericórdia Dómini plena est terra.

Verbo Dómini cæli firmáti sunt: * et spíritu oris ejus omnis virtus eórum.

Cóngregans sicut in utre aquas maris: * ponens in thesáuris abýssos.

Tímeat Dóminum omnis terra: * ab eo autem commoveántur omnes inhabitántes orbem.

Quóniam ipse dixit, et facta sunt: * ipse mandávit, et creáta sunt.

Dóminus díssipat consília géntium: * réprobat autem cogitatiónes populórum, et réprobat consília príncipum.

Consílium autem Dómini in ætérnum manet: * cogitatiónes cordis ejus in generatióne et generatiónem.

\psalmheading{Psalmus 32, ii}

\lettrine[lines=2]{B}{EATA} gens, cujus est Dóminus, Deus ejus: * pópulus, quem elégit in hereditátem sibi.

De cælo respéxit Dóminus: * vidit omnes fílios hóminum.

De præparáto habitáculo suo * respéxit super omnes, qui hábitant terram.

Qui finxit sigillátim corda eórum: * qui intélligit ómnia ópera eórum.

Non salvábitur rex per multam virtútem: * et gigas non salvábitur in multitúdine virtútis suæ.

Fallax equus ad salútem: * in abundántia autem virtútis suæ non salvábitur.

Ecce óculi Dómini super metuéntes eum: * et in eis, qui sperant super misericórdia ejus:

Ut éruat a morte ánimas eórum: * et alat eos in fame.

Anima nostra sústinet Dóminum: * quóniam adjútor et protéctor noster est.

Quia in eo lætábitur cor nostrum: * et in nómine sancto ejus sperávimus.

Fiat misericórdia tua, Dómine, super nos: * quemádmodum sperávimus in te.

\antlabel{Ant.} Exsultáte, justi, et gloriámini, omnes recti corde.

\variantheading{In Officio feriali}

\capitulumlabel{Capit.}{1 Petri 1, 17-19}

\lettrine[lines=2]{I}{N} timóre incolátus vestri témpore conversámini: sciéntes quod non corruptibílibus auro vel argénto redémpti estis, sed pretióso sánguine quasi Agni immaculáti Christi.

\R\ br. Rédime me, Dómine, * Et miserére mei. Rédime. \V\ Pes enim meus stetit in via recta. Et miserére mei. Glória Patri. Rédime. \V\ Ab occúltis meis munda me, Dómine. \R\ Et ab aliénis parce servo tuo.

\variantheading{In Vigiliis communibus:}

\variantheading{Preces feriales}

KYRIE, eléison. Christe, eléison. Kýrie, eléison.

Pater noster \rubric{secreto usque ad} \V\ Et ne nos indúcas in tentatiónem. \R\ Sed líbera nos a malo. \V\ Dómine, Deus virtútum, convérte nos. \R\ Et osténde fáciem tuam, et salvi érimus. \V\ Exsúrge, Christe, ádjuva nos. \R\ Et líbera nos propter nomen tuum. \V\ Dómine, exáudi oratiónem meam. \R\ Et clamor meus ad te véniat.

\rubric{Réliqua ut in Ordinário. 22}

% ===== END OF BATCH 12 -- Covers autumnalis.md lines 3928-4286,
% completing the Little Hours of Feria Secunda in full: "Ad
% Tertiam" (Psalms 26/i, 26/ii, 27), "Ad Sextam" (Psalms 30/i,
% 30/ii, 30/iii), and "Ad Nonam" (Psalms 31, 32/i, 32/ii), each with
% its own hymn, Capitulum, brief Responsory, and "In Vigiliis
% communibus: Preces feriales" apparatus. Three more occurrences of
% the recurring stray OCR "\*" artifact (in the hymns "Nunc, Sancte,
% nobis, Spíritus," "Rector potens, verax Deus," and "Rerum, Deus,
% tenax vigor," each immediately before the final "Præsta, Pater
% piíssime" doxology stanza) are silently dropped, consistent with
% every prior occurrence in this volume. The document continues
% with "Ad Vesperas" (autumnalis.md line 4288).
% =====
\hourheading{Ad Vesperas}

Pater noster et Ave.

\V\ Deus, in adjutórium.

\antlabel{Ant.} Inclinávit Dóminus * aurem suam mihi.

\rubric{¶ Antíphonæ, in hac et áliis Fériis apposítæ, dicúntur semper in Officio feriáli et in Festis sive trium sive novem Lectiónum, quando juxta Rúbricas suméndi sint Psalmi de Féria, ut in Psaltério.}

\psalmheading{Psalmus 114}

\lettrine[lines=2]{D}{ILEXI}, quóniam exáudiet Dóminus * vocem oratiónis meæ.

Quia inclinávit aurem suam mihi: * et in diébus meis invocábo.

Circumdedérunt me dolóres mortis: * et perícula inférni invenérunt me.

Tribulatiónem et dolórem invéni: * et nomen Dómini invocávi.

O Dómine, líbera ánimam meam: * miséricors Dóminus, et justus, et Deus noster miserétur.

Custódiens párvulos Dóminus: * humiliátus sum, et liberávit me.

Convértere, ánima mea, in réquiem tuam: * quia Dóminus benefécit tibi.

Quia erípuit ánimam meam de morte: * óculos meos a lácrimis, pedes meos a lapsu.

Placébo Dómino * in regióne vivórum.

\antlabel{Ant.} Inclinávit Dóminus aurem suam mihi.

\antlabel{Ant.} Vota mea * Dómino reddam coram omni pópulo ejus.

\psalmheading{Psalmus 115}

\lettrine[lines=2]{C}{REDIDI}, propter quod locútus sum: * ego autem humiliátus sum nimis.

Ego dixi in excéssu meo: * Omnis homo mendax.

Quid retríbuam Dómino, * pro ómnibus, quæ retríbuit mihi?

Cálicem salutáris accípiam: * et nomen Dómini invocábo.

Vota mea Dómino reddam coram omni pópulo ejus: * pretiósa in conspéctu Dómini mors sanctórum ejus:

O Dómine, quia ego servus tuus: * ego servus tuus, et fílius ancíllæ tuæ.

Dirupísti víncula mea: * tibi sacrificábo hóstiam laudis, et nomen Dómini invocábo.

Vota mea Dómino reddam in conspéctu omnis pópuli ejus: * in átriis domus Dómini, in médio tui, Jerúsalem.

\antlabel{Ant.} Vota mea Dómino reddam coram omni pópulo ejus.

\antlabel{Ant.} Clamávi, * et Dóminus exaudívit me.

\psalmheading{Psalmus 119}

\lettrine[lines=2]{A}{D} Dóminum cum tribulárer clamávi: * et exaudívit me.

Dómine, líbera ánimam meam a lábiis iníquis, * et a lingua dolósa.

Quid detur tibi, aut quid apponátur tibi * ad linguam dolósam?

Sagíttæ poténtis acútæ, * cum carbónibus desolatóriis.

Heu mihi, quia incolátus meus prolongátus est: * habitávi cum habitántibus Cedar: multum íncola fuit ánima mea.

Cum his, qui odérunt pacem, eram pacíficus: * cum loquébar illis, impugnábant me gratis.

\antlabel{Ant.} Clamávi, et Dóminus exaudívit me.

\antlabel{Ant.} Auxílium meum * a Dómino qui fecit cælum et terram.

\psalmheading{Psalmus 120}

\lettrine[lines=2]{L}{EVAVI} óculos meos in montes, * unde véniet auxílium mihi.

Auxílium meum a Dómino, * qui fecit cælum et terram.

Non det in commotiónem pedem tuum: * neque dormítet qui custódit te.

Ecce, non dormitábit neque dórmiet, * qui custódit Israël.

Dóminus custódit te, Dóminus protéctio tua, * super mánum déxteram tuam.

Per diem sol non uret te: * neque luna per noctem.

Dóminus custódit te ab omni malo: * custódiat ánimam tuam Dóminus.

Dóminus custódiat intróitum tuum, et éxitum tuum: * ex hoc nunc, et usque in sæculum.

\antlabel{Ant.} Auxílium meum a Dómino, qui fecit cælum et terram.

\antlabel{Ant.} Lætátus sum * \textdagger\ in his, quæ dicta sunt mihi. \textdaggerdbl

\psalmheading{Psalmus 121}

\lettrine[lines=2]{L}{ÆTATUS} sum \textdagger\ in his, quæ dicta sunt mihi: * \textdaggerdbl\ In domum Dómini íbimus.

Stantes erant pedes nostri, * in átriis tuis, Jerúsalem.

Jerúsalem, quæ ædificátur ut cívitas: * cujus participátio ejus in idípsum.

Illuc enim ascendérunt tribus, tribus Dómini: * testimónium Israël ad confiténdum nómini Dómini.

Quia illic sedérunt sedes in judício, * sedes super domum David.

Rogáte quæ ad pacem sunt Jerúsalem: * et abundántia diligéntibus te:

Fiat pax in virtúte tua: * et abundántia in túrribus tuis.

Propter fratres meos, et próximos meos, * loquébar pacem de te:

Propter domum Dómini, Dei nostri, * quæsívi bona tibi.

\antlabel{Ant.} Lætátus sum in his, quæ dicta sunt mihi.

\variantheading{In Officio feriali}

\rubric{Capítulum, Hymnus, Versus, et Ant. ad Magníficat dicúntur quotídie de Féria curránti, ut in Psaltério hic et álibi appónuntur.}

\capitulumlabel{Capitulum}{2 Cor. 1, 3-4}

\lettrine[lines=2]{B}{ENEDICTUS} Deus, et Pater Dómini nostri Jesu Christi, Pater misericordiárum, et Deus totíus consolatiónis, qui consolátur nos in omni tribulatióne nostra.

\hymnusheading{Hymnus}

\lettrine[lines=2]{I}{MMENSE} cæli Cónditor,\\
Qui mixta ne confúnderent,\\
Aquæ fluénta dívidens,\\
Cælum dedísti límitem.

Firmans locum cæléstibus,\\
Simúlque terræ rívulis;\\
Ut unda flammas témperet,\\
Terræ solum ne díssipent.

Infúnde nunc, piíssime,\\
Donum perénnis grátiæ:\\
Fraudis novæ ne cásibus\\
Nos error átterat vetus.

Lucem fides adáugeat:\\
Sic lúminis jubar ferat:\\
Hæc vana cuncta próterat:\\
Hanc falsa nulla cómprimant.

Præsta, Pater piíssime,\\
Patríque compar Unice,\\
Cum Spíritu Paráclito\\
Regnans per omne sæculum.\\
Amen.

\V\ Dirigátur, Dómine, orátio mea. \R\ Sicut incénsum in conspéctu tuo.

\antlabel{Ad Magnif. Ant.} Magníficat * \textdagger\ ánima mea Dóminum, quia respéxit Deus humilitátem meam.

\rubric{Réliqua ut in Ordinário. 22}

\hourheading{Ad Completorium}

\V\ Jube, domne.

\benedlabel{Bened.} Noctem quiétam, et finem perféctum concédat nobis Dóminus omnípotens. \R\ Amen.

\lectioheading{Lectio br.}{1 Petri 5, 8-9}

\lettrine[lines=2]{F}{RATRES}: Sóbrii estóte, et vigiláte: quia adversárius vester diábolus tamquam leo rúgiens círcuit, quærens quem dévoret: cui resístite fortes in fide. Tu autem, Dómine, miserére nobis. \R\ Deo grátias.

\V\ Adjutórium nostrum in nómine Dómini. \R\ Qui fecit cælum et terram.

Pater noster \rubric{quod dícitur totum secréto.}

\formulalabel{Confíteor.} Misereátur. Indulgéntiam.

\V\ Convérte nos, Deus, salutáris noster. \R\ Et avérte iram tuam a nobis.

\V\ Deus, in adjutórium.

\antlabel{Ant.} Salvum me fac.

\rubric{¶ Antíphonæ, in hac et áliis Fériis apposítæ, dicúntur semper ad Completórium tum de Témpore tum de Sanctis, quando juxta Rúbricas suméndi sint Psalmi de Féria.}

\psalmheading{Psalmus 6}

\lettrine[lines=2]{D}{OMINE}, ne in furóre tuo árguas me, * neque in ira tua corrípias me.

Miserére mei, Dómine, quóniam infírmus sum: * sana me, Dómine, quóniam conturbáta sunt ossa mea.

Et ánima mea turbáta est valde: * sed tu, Dómine, úsquequo?

Convértere, Dómine, et éripe ánimam meam: * salvum me fac propter misericórdiam tuam.

Quóniam non est in morte qui memor sit tui: * in inférno autem quis confitébitur tibi?

Laborávi in gémitu meo, lavábo per síngulas noctes lectum meum: * lácrimis meis stratum meum rigábo.

Turbátus est a furóre óculus meus: * inveterávi inter omnes inimícos meos.

Discédite a me, omnes, qui operámini iniquitátem: * quóniam exaudívit Dóminus vocem fletus mei.

Exaudívit Dóminus deprecatiónem meam, * Dóminus oratiónem meam suscépit.

Erubéscant, et conturbéntur veheménter omnes inimíci mei: * convertántur et erubéscant valde velóciter.

\psalmheading{Psalmus 7, i}

\lettrine[lines=2]{D}{OMINE}, Deus meus, in te sperávi: * salvum me fac ex ómnibus persequéntibus me, et líbera me.

Nequándo rápiat ut leo ánimam meam, * dum non est qui rédimat, neque qui salvum fáciat.

Dómine, Deus meus, si feci istud, * si est iníquitas in mánibus meis:

Si réddidi retribuéntibus mihi mala, * décidam mérito ab inimícis meis inánis.

Persequátur inimícus ánimam meam, et comprehéndat, et concúlcet in terra vitam meam, * et glóriam meam in púlverem dedúcat.

Exsúrge, Dómine, in ira tua: * et exaltáre in fínibus inimicórum meórum.

Et exsúrge, Dómine, Deus meus, in præcépto quod mandásti: * et synagóga populórum circúmdabit te.

Et propter hanc in altum regrédere: * Dóminus júdicat pópulos.

Júdica me, Dómine, secúndum justítiam meam, * et secúndum innocéntiam meam super me.

Consumétur nequítia peccatórum, et díriges justum, * scrutans corda et renes, Deus.

\psalmheading{Psalmus 7, ii}

\lettrine[lines=2]{J}{USTUM} adjutórium meum a Dómino, * qui salvos facit rectos corde.

Deus judex justus, fortis, et pátiens: * numquid iráscitur per síngulos dies?

Nisi convérsi fuéritis, gládium suum vibrábit: * arcum suum teténdit, et parávit illum.

Et in eo parávit vasa mortis: * sagíttas suas ardéntibus effécit.

Ecce, partúriit injustítiam: * concépit dolórem, et péperit iniquitátem.

Lacum apéruit, et effódit eum: * et íncidit in fóveam quam fecit.

Convertétur dolor ejus in caput ejus: * et in vérticem ipsíus iníquitas ejus descéndet.

Confitébor Dómino secúndum justítiam ejus: * et psallam nómini Dómini altíssimi.

\antlabel{Ant.} Salvum me fac, Dómine, propter misericórdiam tuam.

\hymnusheading{Hymnus}

\lettrine[lines=2]{T}{E} lucis ante términum,\\
Rerum Creátor, póscimus,\\
Ut pro tua cleméntia\\
Sis præsul et custódia.

Procul recédant sómnia,\\
Et nóctium phantásmata,\\
Hostémque nostrum cómprime,\\
Ne polluántur córpora.

Præsta, Pater piíssime,\\
Patríque compar Unice,\\
Cum Spíritu Paráclito\\
Regnans per omne sæculum.\\
Amen.

\capitulumlabel{Capitulum}{Jer. 14, 9}

\lettrine[lines=2]{T}{U} autem in nobis es, Dómine, et nomen sanctum tuum invocátum est super nos: ne derelínquas nos, Dómine, Deus noster.

\R\ br. In manus tuas, Dómine, * Comméndo spíritum meum. In manus. \V\ Redemísti nos, Dómine, Deus veritátis. Comméndo. Glória Patri. In manus.

\V\ Custódi nos, Dómine, ut pupíllam óculi. \R\ Sub umbra alárum tuárum prótege nos.

\antlabel{Ant.} Salva nos.

\canticumheading{Canticum Simeonis}{Luc. 2, 29-32}

\lettrine[lines=2]{N}{UNC} dimíttis servum tuum, Dómine, * secúndum verbum tuum in pace:

Quia vidérunt óculi mei * salutáre tuum,

Quod parásti * ante fáciem ómnium populórum,

Lumen ad revelatiónem géntium, * et glóriam plebis tuæ Israël.

Glória Patri.

\antlabel{Ant.} Salva nos, Dómine, vigilántes, custódi nos dormiéntes; ut vigilémus cum Christo, et requiescámus in pace.

\variantheading{Preces (si dicendæ sint)}

KYRIE, eléison. Christe, eléison. Kýrie, eléison.

Pater noster \rubric{secreto usque ad} \V\ Et ne nos indúcas in tentatiónem. \R\ Sed líbera nos a malo.

Credo in Deum \rubric{secreto usque ad} \V\ Carnis resurrectiónem. \R\ Vitam ætérnam. Amen. \V\ Benedíctus es, Dómine, Deus patrum nostrórum. \R\ Et laudábilis et gloriósus in sæcula. \V\ Benedicámus Patrem et Fílium cum Sancto Spíritu. \R\ Laudémus, et superexaltémus eum in sæcula. \V\ Benedíctus es, Dómine, in firmaménto cæli. \R\ Et laudábilis, et gloriósus, et superexaltátus in sæcula. \V\ Benedícat et custódiat nos omnípotens et miséricors Dóminus. \R\ Amen. \V\ Dignáre, Dómine, nocte ista. \R\ Sine peccáto nos custodíre. \V\ Miserére nostri, Dómine. \R\ Miserére nostri. \V\ Fiat misericórdia tua, Dómine, super nos. \R\ Quemádmodum sperávimus in te. \V\ Dómine, exáudi oratiónem meam. \R\ Et clamor meus ad te véniat.

\V\ Dóminus vobíscum.

Orémus. \oratiolabel{Oratio}

\lettrine[lines=2]{V}{ISITA}, quæsumus, Dómine, habitatiónem istam, et omnes insídias inimíci ab ea longe repélle: Angeli tui sancti hábitent in ea, qui nos in pace custódiant; et benedíctio tua sit super nos semper. Per Dóminum nostrum. \R\ Amen.

\V\ Dóminus vobíscum.

\V\ Benedicámus Dómino.

\benedictiolabel{Benedictio} Benedícat et custódiat nos omnípotens et miséricors Dóminus, Pater, et Fílius, et Spíritus Sanctus. \R\ Amen.

\variantheading{Antiphona finalis B. M. V.}

\V\ Divínum auxílium.

Pater, Ave et Credo.

% ===== END OF BATCH 13 -- Covers autumnalis.md lines 4288-4607,
% completing "Ad Vesperas" (Psalms 114, 115, 119, 120, 121, the
% Capitulum, the hymn "Immense cæli Cónditor," and the Magnificat
% antiphon) and "Ad Completorium" (Psalms 6, 7/i, 7/ii, the hymn
% "Te lucis ante términum," the Capitulum, the brief Responsory "In
% manus tuas," the Canticum Simeonis, and the Preces), and with
% them the ENTIRE "Feria Secunda" Office. Two more occurrences of
% the recurring stray OCR "\*" artifact (in "Immense cæli Cónditor"
% at autumnalis.md line 4436 and "Te lucis ante términum" at line 4550,
% both at the hymn's final doxology stanza) are silently dropped,
% consistent with every prior occurrence in this volume. The source
% at line 4464 bolds only "Confíteor." (not the full "Misereátur.
% Indulgéntiam."), unlike the fully-bold form seen at every other
% occurrence transcribed so far in this volume; this is a genuine,
% recurring alternation in autumnalis.md itself (confirmed by grep
% across the file) and is transcribed exactly as marked here:
% \formulalabel{Confíteor.} followed by plain-text "Misereátur.
% Indulgéntiam." The closing "Antiphona finalis B. M. V." is again
% given only as the bare cross-reference incipit "℣. Divínum
% auxílium," matching the identical abbreviated treatment already
% established for Sunday's Compline earlier in this volume. The
% document continues with "Feria Tertia" (autumnalis.md line 4610),
% the second weekday of the Psalterium cycle.
% =====
\feriaheading{Feria Tertia}

\hourheading{Ad Matutinum}

Pater, Ave et Credo.

\V\ Dómine, lábia mea.

\V\ Deus, in adjutórium.

\variantheading{In Officio feriali}

\invitatlabel{Invitat.} Jubilémus Deo, * Salutári nostro. \psalmlabel{Ps. 94} Veníte, exsultémus.

\hymnusheading{Hymnus}

\lettrine[lines=2]{C}{ONSORS} Patérni lúminis,\\
Lux ipse lucis, et dies,\\
Noctem canéndo rúmpimus:\\
Assíste postulántibus.

Aufer tenébras méntium,\\
Fuga catérvas dæmonum,\\
Expélle somnoléntiam,\\
Ne pigritántes óbruat.

Sic, Christe, nobis ómnibus\\
Indúlgeas credéntibus,\\
Ut prosit exorántibus\\
Quod præcinéntes psállimus.

Præsta, Pater piíssime,\\
Patríque compar Unice,\\
Cum Spíritu Paráclito\\
Regnans per omne sæculum.\\
Amen.

\nocturnoheading{In I Nocturno}

\antlabel{Ant.} Expúgna, Dómine, * impugnántes me.

\psalmheading{Psalmus 34, i}

\lettrine[lines=2]{J}{UDICA}, Dómine, nocéntes me, * expúgna impugnántes me.

Apprehénde arma et scutum: * exsúrge in adjutórium mihi.

Effúnde frámeam, et conclúde advérsus eos, qui persequúntur me: * dic ánimæ meæ: Salus tua ego sum.

Confundántur et revereántur, * quæréntes ánimam meam.

Avertántur retrórsum, et confundántur * cogitántes mihi mala.

Fiant tamquam pulvis ante fáciem venti: * et Angelus Dómini coárctans eos.

Fiat via illórum ténebræ et lúbricum: * et Angelus Dómini pérsequens eos.

Quóniam gratis abscondérunt mihi intéritum láquei sui: * supervácue exprobravérunt ánimam meam.

Véniat illi láqueus, quem ignórat: et cáptio, quam abscóndit, apprehéndat eum: * et in láqueum cadat in ipsum.

Anima autem mea exsultábit in Dómino: * et delectábitur super salutári suo.

Omnia ossa mea dicent: * Dómine, quis símilis tibi?

Erípiens ínopem de manu fortiórum ejus: * egénum et páuperem a diripiéntibus eum.

\antlabel{Ant.} Expúgna, Dómine, impugnántes me.

\antlabel{Ant.} Restítue ánimam meam * a malefáctis eórum, Dómine.

\psalmheading{Psalmus 34, ii}

\lettrine[lines=2]{S}{URGENTES} testes iníqui, * quæ ignorábam interrogábant me.

Retribuébant mihi mala pro bonis: * sterilitátem ánimæ meæ.

Ego autem cum mihi molésti essent, * induébar cilício.

Humiliábam in jejúnio ánimam meam: * et orátio mea in sinu meo convertétur.

Quasi próximum, et quasi fratrem nostrum, sic complacébam: * quasi lugens et contristátus, sic humiliábar.

Et advérsum me lætáti sunt, et convenérunt: * congregáta sunt super me flagélla, et ignorávi.

Dissipáti sunt, nec compúncti, tentavérunt me, subsannavérunt me subsannatióne: * frenduérunt super me déntibus suis.

Dómine, quando respícies? * restítue ánimam meam a malignitáte eórum, a leónibus únicam meam.

\antlabel{Ant.} Restítue ánimam meam a malefáctis eórum, Dómine.

\antlabel{Ant.} Exsúrge, Dómine, * et inténde judício meo.

\psalmheading{Psalmus 34, iii}

\lettrine[lines=2]{C}{ONFITEBOR} tibi in ecclésia magna, * in pópulo gravi laudábo te.

Non supergáudeant mihi qui adversántur mihi iníque: * qui odérunt me gratis et ánnuunt óculis.

Quóniam mihi quidem pacífice loquebántur: * et in iracúndia terræ loquéntes, dolos cogitábant.

Et dilatavérunt super me os suum: * dixérunt: Euge, euge, vidérunt óculi nostri.

Vidísti, Dómine, ne síleas: * Dómine, ne discédas a me.

Exsúrge et inténde judício meo: * Deus meus, et Dóminus meus in causam meam.

Júdica me secúndum justítiam tuam, Dómine, Deus meus, * et non supergáudeant mihi.

Non dicant in córdibus suis: Euge, euge, ánimæ nostræ: * nec dicant: Devorávimus eum.

Erubéscant et revereántur simul, * qui gratulántur malis meis.

Induántur confusióne et reveréntia * qui magna loquúntur super me.

Exsúltent et læténtur qui volunt justítiam meam: * et dicant semper: Magnificétur Dóminus, qui volunt pacem servi ejus.

Et lingua mea meditábitur justítiam tuam, * tota die laudem tuam.

\antlabel{Ant.} Exsúrge, Dómine, et inténde judício meo.

\variantheading{In Festis novem Lectionum:}

\V\ Lingua mea meditábitur justítiam tuam. \R\ Tota die laudem tuam, Dómine.

\absolutiolabel{Absol.} Exáudi, Dómine.

\benedlabel{Bened.} 1 Benedictióne perpétua. 2 Unigénitus Dei Fílius. 3 Spíritus Sancti grátia.

\nocturnoheading{In II Nocturno}

\antlabel{Ant.} Noli æmulári * \textdagger\ in eo, qui prosperátur, et facit iniquitátem.

\psalmheading{Psalmus 36, i}

\lettrine[lines=2]{N}{OLI} æmulári in malignántibus: * neque zeláveris faciéntes iniquitátem.

Quóniam tamquam fœnum velóciter aréscent: * et quemádmodum ólera herbárum cito décident.

Spera in Dómino, et fac bonitátem: * et inhábita terram, et pascéris in divítiis ejus.

Delectáre in Dómino: * et dabit tibi petitiónes cordis tui.

Revéla Dómino viam tuam, et spera in eo: * et ipse fáciet.

Et edúcet quasi lumen justítiam tuam: et judícium tuum tamquam merídiem: * súbditus esto Dómino, et ora eum.

Noli æmulári in eo, qui prosperátur in via sua: * in hómine faciénte injustítias.

Désine ab ira, et derelínque furórem: * noli æmulári ut malignéris.

Quóniam qui malignántur, exterminabúntur: * sustinéntes autem Dóminum, ipsi hereditábunt terram.

Et adhuc pusíllum, et non erit peccátor: * et quæres locum ejus, et non inviénies.

Mansuéti autem hereditábunt terram: * et delectabúntur in multitúdine pacis.

Observábit peccátor justum: * et stridébit super eum déntibus suis.

Dóminus autem irridébit eum: * quóniam próspicit quod véniet dies ejus.

Gládium evaginavérunt peccatóres: * intendérunt arcum suum.

Ut dejíciant páuperem et ínopem: * ut trucídent rectos corde.

Gládius eórum intret in corda ipsórum: * et arcus eórum confringátur.

\antlabel{Ant.} Noli æmulári in eo, qui prosperátur, et facit iniquitátem.

\antlabel{Ant.} Bráchia peccatórum * conteréntur: confírmat autem justos Dóminus.

\psalmheading{Psalmus 36, ii}

\lettrine[lines=2]{M}{ELIUS} est módicum justo, * super divítias peccatórum multas.

Quóniam bráchia peccatórum conteréntur: * confírmat autem justos Dóminus.

Novit Dóminus dies immaculatórum: * et heréditas eórum in ætérnum erit.

Non confundéntur in témpore malo, et in diébus famis saturabúntur: * quia peccatóres períbunt.

Inimíci vero Dómini mox ut honorificáti fúerint et exaltáti: * deficiéntes, quemádmodum fumus defícient.

Mutuábitur peccátor, et non solvet: * justus autem miserétur et tríbuet.

Quia benedicéntes ei hereditábunt terram: * maledicéntes autem ei períbunt.

Apud Dóminum gressus hóminis dirigéntur: * et viam ejus volet.

Cum cecíderit, non collidétur: * quia Dóminus suppónit manum suam.

Júnior fui, étenim sénui: * et non vidi justum derelíctum, nec semen ejus quærens panem.

Tota die miserétur et cómmodat: * et semen illíus in benedictióne erit.

Declína a malo, et fac bonum: * et inhábita in sæculum sæculi.

Quia Dóminus amat judícium, et non derelínquet sanctos suos: * in ætérnum conservabúntur.

Injústi puniéntur: * et semen impiórum períbit.

Justi autem hereditábunt terram: * et inhabitábunt in sæculum sæculi super eam.

\antlabel{Ant.} Bráchia peccatórum conteréntur: confírmat autem justos Dóminus.

\antlabel{Ant.} Custódi innocéntiam, * et vide æquitátem.

\psalmheading{Psalmus 36, iii}

\lettrine[lines=2]{O}{S} justi meditábitur sapiéntiam, * et lingua ejus loquétur judícium.

Lex Dei ejus in corde ipsíus: * et non supplantabúntur gressus ejus.

Consíderat peccátor justum: * et quærit mortificáre eum.

Dóminus autem non derelínquet eum in mánibus ejus: * nec damnábit eum, cum judicábitur illi.

Exspécta Dóminum, et custódi viam ejus: et exaltábit te ut hereditáte cápias terram: * cum períerint peccatóres vidébis.

Vidi ímpium superexaltátum, * et elevátum sicut cedros Líbani.

Et transívi, et ecce non erat: * et quæsívi eum, et non est invéntus locus ejus.

Custódi innocéntiam, et vide æquitátem: * quóniam sunt relíquiæ hómini pacífico.

Injústi autem disperíbunt simul: * relíquiæ impiórum interíbunt.

Salus autem justórum a Dómino: * et protéctor eórum in témpore tribulatiónis.

Et adjuvábit eos Dóminus et liberábit eos: * et éruet eos a peccatóribus, et salvábit eos: quia speravérunt in eo.

\antlabel{Ant.} Custódi innocéntiam, et vide æquitátem.

\variantheading{In Festis novem Lectionum:}

\V\ Exspécta Dóminum, et custódi viam ejus. \R\ Exaltábit te ut hereditáte cápias terram.

\absolutiolabel{Absolutio} Ipsíus píetas et.

\benedlabel{Bened.} 4 Deus Pater omnípotens. 5 Christus perpétuæ. 6 Ignem sui amóris.

\nocturnoheading{In III Nocturno}

\antlabel{Ant.} Ne in ira tua * corrípias me, Dómine.

\psalmheading{Psalmus 37, i}

\lettrine[lines=2]{D}{OMINE}, ne in furóre tuo árguas me, * neque in ira tua corrípias me.

Quóniam sagíttæ tuæ infíxæ sunt mihi: * et confirmásti super me manum tuam.

Non est sánitas in carne mea a fácie iræ tuæ: * non est pax óssibus meis a fácie peccatórum meórum.

Quóniam iniquitátes meæ supergréssæ sunt caput meum: * et sicut onus grave gravátæ sunt super me.

Putruérunt et corrúptæ sunt cicatríces meæ, * a fácie insipiéntiæ meæ.

Miser factus sum, et curvátus sum usque in finem: * tota die contristátus ingrediébar.

Quóniam lumbi mei impléti sunt illusiónibus: * et non est sánitas in carne mea.

Afflíctus sum, et humiliátus sum nimis: * rugiébam a gémitu cordis mei.

Dómine, ante te omne desidérium meum: * et gémitus meus a te non est abscónditus.

Cor meum conturbátum est, derelíquit me virtus mea: * et lumen oculórum meórum, et ipsum non est mecum.

\antlabel{Ant.} Ne in ira tua corrípias me, Dómine.

\antlabel{Ant.} Inténde in adjutórium meum, * Dómine, virtus salútis meæ.

\psalmheading{Psalmus 37, ii}

\lettrine[lines=2]{A}{MICI} mei, et próximi mei * advérsum me appropinquavérunt, et stetérunt.

Et qui juxta me erant, de longe stetérunt: * et vim faciébant qui quærébant ánimam meam.

Et qui inquirébant mala mihi, locúti sunt vanitátes: * et dolos tota die meditabántur.

Ego autem tamquam surdus non audiébam: * et sicut mutus non apériens os suum.

Et factus sum sicut homo non áudiens: * et non habens in ore suo redargutiónes.

Quóniam in te, Dómine, sperávi: * tu exáudies me, Dómine, Deus meus.

Quia dixi: Nequándo supergáudeant mihi inimíci mei: * et dum commovéntur pedes mei, super me magna locúti sunt.

Quóniam ego in flagélla parátus sum: * et dolor meus in conspéctu meo semper.

Quóniam iniquitátem meam annuntiábo: * et cogitábo pro peccáto meo.

Inimíci autem mei vivunt, et confirmáti sunt super me: * et multiplicáti sunt qui odérunt me iníque.

Qui retríbuunt mala pro bonis, detrahébant mihi: * quóniam sequébar bonitátem.

Ne derelínquas me, Dómine, Deus meus: * ne discésseris a me.

Inténde in adjutórium meum, * Dómine, Deus salútis meæ.

\antlabel{Ant.} Inténde in adjutórium meum, Dómine, virtus salútis meæ.

\antlabel{Ant.} Amove, Dómine, * a me plagas tuas.

\psalmheading{Psalmus 38}

\lettrine[lines=2]{D}{IXI}: Custódiam vias meas: * ut non delínquam in lingua mea.

Pósui ori meo custódiam, * cum consísteret peccátor advérsum me.

Obmútui, et humiliátus sum, et sílui a bonis: * et dolor meus renovátus est.

Concáluit cor meum intra me: * et in meditatióne mea exardéscet ignis.

Locútus sum in lingua mea: * Notum fac mihi, Dómine, finem meum.

Et númerum diérum meórum quis est: * ut sciam quid desit mihi.

Ecce, mensurábiles posuísti dies meos: * et substántia mea tamquam níhilum ante te.

Verúmtamen univérsa vánitas, * omnis homo vivens.

Verúmtamen in imágine pertránsit homo: * sed et frustra conturbátur.

Thesaurízat: * et ignórat cui congregábit ea.

Et nunc quæ est exspectátio mea? Nonne Dóminus? * Et substántia mea apud te est.

Ab ómnibus iniquitátibus meis érue me: * oppróbrium insipiénti dedísti me.

Obmútui, et non apérui os meum, quóniam tu fecísti: * ámove a me plagas tuas.

A fortitúdine manus tuæ ego defíci in increpatiónibus: * propter iniquitátem corripuísti hóminem.

Et tabéscere fecísti sicut aráneam ánimam ejus: * verúmtamen vane conturbátur omnis homo.

Exáudi oratiónem meam, Dómine, et deprecatiónem meam: * áuribus pércipe lácrimas meas.

Ne síleas: quóniam ádvena ego sum apud te, et peregrínus, * sicut omnes patres mei.

Remítte mihi, ut refrígerer priúsquam ábeam, * et ámplius non ero.

\antlabel{Ant.} Amove, Dómine, a me plagas tuas.

\variantheading{In feriali Officio et in Festis:}

\V\ Deus, ne síleas a me, remítte mihi. \R\ Quóniam íncola ego sum apud te, et peregrínus.

\variantheading{In Officio novem Lectionum}

\absolutiolabel{Absolutio.} A vínculis.

\benedlabel{Bened.} 7 Evangélica léctio. 8 In Festis Dómini: Divínum auxílium. In Festis B. M. V.: Cujus festum cólimus, ipsa Virgo. In Festis Sanctórum: Cujus (\textit{vel} Quorum \textit{aut} Quarum) festum cólimus. 9 Ad societátem. (Per evangélica dicta.)

\variantheading{In Officio trium Lectionum}

\absolutiolabel{Absolutio.} Ipsíus píetas et.

\variantheading{In Feriis, cum Homilia:}

\benedlabel{Bened.} 1 Evangélica. 2 Divínum. 3 Ad societátem.

\variantheading{In Feriis, sine Homilia:}

\benedlabel{Bened.} 1 Deus Pater. 2 Christus perpétuæ. 3 Ignem sui amóris.

\variantheading{In Festis Sanctorum:}

\benedlabel{Bened.} 1 Ille nos benedícat, qui sine fine vivit et regnat. \R\ Amen. 2 Cujus (\textit{vel} Quorum \textit{aut} Quarum). 3 Ad societátem.

% ===== END OF BATCH 14 -- Covers autumnalis.md lines 4610-4978,
% opening "Feria Tertia" (Tuesday) and completing its "Ad
% Matutinum" in full: the ferial Invitatorium "Jubilémus Deo,
% Salutári nostro," the hymn "Consors Patérni lúminis" (with the
% recurring stray OCR asterisk before its final doxology stanza
% silently dropped, consistent with every prior occurrence), and
% the three Nocturns (Psalms 34/i-iii; 36/i-iii; 37/i, 37/ii, 38),
% each with its versicle, Absolutio, and Benedictiones apparatus.
% "Absolutio" is rendered with the exact trailing period the source
% gives it at lines 4959 and 4965 ("Absolutio.") though not at line
% 4847 ("Absolutio"), preserving the source's own inconsistency.
% The compact multi-option "Bened." line for the ninth lesson is
% again transcribed as one continuous argument, matching the
% identical text already seen in Feria Secunda's Matins. The
% document continues with "Ad Laudes I" (autumnalis.md line 4981).
% =====
\hourheading{Ad Laudes I}

\rubric{In Feriis extra Vigílias commúnes occurréntibus, atque in Festis.}

(Pater noster et Ave.)

\V\ Deus, in adjutórium.

\antlabel{Ant.} Cantáte * \textdagger\ Dómino, et benedícite nómini ejus.

\psalmheading{Psalmus 95}

\lettrine[lines=2]{C}{ANTATE} \textdagger\ Dómino cánticum novum: * cantáte Dómino, omnis terra.

Cantáte Dómino, et benedícite nómini ejus: * annuntiáte de die in diem salutáre ejus.

Annuntiáte inter gentes glóriam ejus, * in ómnibus pópulis mirabília ejus.

Quóniam magnus Dóminus, et laudábilis nimis: * terríbilis est super omnes deos.

Quóniam omnes dii géntium dæmónia: * Dóminus autem cælos fecit.

Conféssio, et pulchritúdo in conspéctu ejus: * sanctimónia et magnificéntia in sanctificatióne ejus.

Afférte Dómino, pátriæ géntium, afférte Dómino glóriam et honórem: * afférte Dómino glóriam nómini ejus.

Tóllite hóstias, et introíte in átria ejus: * adoráte Dóminum in átrio sancto ejus.

Commoveátur a fácie ejus univérsa terra: * dícite in géntibus quia Dóminus regnávit.

Etenim corréxit orbem terræ qui non commovébitur: * judicábit pópulos in æquitáte.

Læténtur cæli, et exsúltet terra: commoveátur mare, et plenitúdo ejus: * gaudébunt campi, et ómnia quæ in eis sunt.

Tunc exsultábunt ómnia ligna silvárum a fácie Dómini, quia venit: * quóniam venit judicáre terram.

Judicábit orbem terræ in æquitáte, * et pópulos in veritáte tua.

\antlabel{Ant.} Cantáte Dómino, et benedícite nómini ejus.

\antlabel{Ant.} Salutáre vultus mei * Deus meus.

\psalmheading{Psalmus 42}

\lettrine[lines=2]{J}{UDICA} me, Deus, et discérne causam meam de gente non sancta, * ab hómine iníquo, et dolóso érue me.

Quia tu es, Deus, fortitúdo mea: * quare me repulísti? et quare tristis incédo, dum afflígit me inimícus?

Emítte lucem tuam et veritátem tuam: * ipsa me deduxérunt, et adduxérunt in montem sanctum tuum, et in tabernácula tua.

Et introíbo ad altáre Dei: * ad Deum, qui lætíficat juventútem meam.

Confitébor tibi in cíthara, Deus, Deus meus: * quare tristis es, ánima mea? et quare contúrbas me?

Spera in Deo, quóniam adhuc confitébor illi: * salutáre vultus mei, et Deus meus.

\antlabel{Ant.} Salutáre vultus mei Deus meus.

\antlabel{Ant.} Illúmina, Dómine, * vultum tuum super nos.

\psalmheading{Psalmus 66}

\lettrine[lines=2]{D}{EUS} misereátur nostri, et benedícat nobis: * illúminet vultum suum super nos, et misereátur nostri.

Ut cognoscámus in terra viam tuam, * in ómnibus géntibus salutáre tuum.

Confiteántur tibi pópuli, Deus: * confiteántur tibi pópuli omnes.

Læténtur et exsúltent gentes: * quóniam júdicas pópulos in æquitáte, et gentes in terra dírigis.

Confiteántur tibi pópuli, Deus, confiteántur tibi pópuli omnes: * terra dedit fructum suum.

Benedícat nos Deus, Deus noster, benedícat nos Deus: * et métuant eum omnes fines terræ.

\antlabel{Ant.} Illúmina, Dómine, vultum tuum super nos.

\antlabel{Ant.} Exaltáte * Regem sæculórum in opéribus vestris.

\canticumheading{Canticum Tobiæ}{Tob. 13, 1-10}

\lettrine[lines=2]{M}{AGNUS} es, Dómine, in ætérnum, * et in ómnia sæcula regnum tuum:

Quóniam tu flagéllas, et salvas: dedúcis ad ínferos, et redúcis: * et non est qui effúgiat manum tuam.

Confitémini Dómino, fílii Israël, * et in conspéctu géntium laudáte eum:

Quóniam ídeo dispérsit vos inter gentes, quæ ignórant eum, * ut vos enarrétis mirabília ejus.

Et faciátis scire eos, * quia non est álius Deus omnípotens præter eum.

Ipse castigávit nos propter iniquitátes nostras: * et ipse salvábit nos propter misericórdiam suam.

Aspícite ergo quæ fecit nobíscum, et cum timóre et tremóre confitémini illi: * Regémque sæculórum exaltáte in opéribus vestris.

Ego autem in terra captivitátis meæ confitébor illi: * quóniam osténdit majestátem suam in gentem peccatrícem.

Convertímini ítaque, peccatóres, et fácite justítiam coram Deo, * credéntes quod fáciat vobíscum misericórdiam suam:

Ego autem, et ánima mea * in eo lætábimur.

Benedícite Dóminum, omnes elécti ejus: * ágite dies lætítiæ, et confitémini illi.

\antlabel{Ant.} Exaltáte Regem sæculórum in opéribus vestris.

\antlabel{Ant.} Laudáte * \textdagger\ nomen Dómini, qui statis in domo Dómini.

\psalmheading{Psalmus 134}

\lettrine[lines=2]{L}{AUDATE} \textdagger\ nomen Dómini, * laudáte, servi, Dóminum.

Qui statis in domo Dómini, * in átriis domus Dei nostri.

Laudáte Dóminum, quia bonus Dóminus: * psállite nómini ejus, quóniam suáve.

Quóniam Jacob elégit sibi Dóminus, * Israël in possessiónem sibi.

Quia ego cognóvi quod magnus est Dóminus, * et Deus noster præ ómnibus diis.

Omnia quæcúmque vóluit, Dóminus fecit in cælo, et in terra, * in mari, et in ómnibus abýssis.

Edúcens nubes ab extrémo terræ: * fúlgura in plúviam fecit.

Qui prodúcit ventos de thesáuris suis: * qui percússit primogénita Ægýpti ab hómine usque ad pecus.

Et misit signa, et prodígia in médio tui, Ægýpte: * in Pharaónem, et in omnes servos ejus.

Qui percússit gentes multas: * et occídit reges fortes:

Sehon, regem Amorrhæórum, et Og, regem Basan, * et ómnia regna Chánaan.

Et dedit terram eórum hereditátem, * hereditátem Israël, pópulo suo.

Dómine, nomen tuum in ætérnum: * Dómine, memoriále tuum in generatiónem et generatiónem.

Quia judicábit Dóminus pópulum suum: * et in servis suis deprecábitur.

Simulácra géntium argéntum, et aurum, * ópera mánuum hóminum.

Os habent, et non loquéntur: * óculos habent, et non vidébunt.

Aures habent, et non áudient: * neque enim est spíritus in ore ipsórum.

Símiles illis fiant qui fáciunt ea: * et omnes qui confídunt in eis.

Domus Israël, benedícite Dómino: * domus Aaron, benedícite Dómino.

Domus Levi, benedícite Dómino: * qui timétis Dóminum, benedícite Dómino.

Benedíctus Dóminus ex Sion, * qui hábitat in Jerúsalem.

\antlabel{Ant.} Laudáte nomen Dómini, qui statis in domo Dómini.

\variantheading{In Officio feriali}

\capitulumlabel{Capit.}{Rom. 13, 12-13}

\lettrine[lines=2]{N}{OX} præcéssit, dies autem appropinquávit. Abjiciámus ergo ópera tenebrárum, et induámur arma lucis. Sicut in die honéste ambulémus.

\hymnusheading{Hymnus}

\lettrine[lines=2]{A}{LES} diéi núntius\\
Lucem propínquam præcínit:\\
Nos excitátor méntium\\
Jam Christus ad vitam vocat.

Auférte, clamat, léctulos,\\
Ægro sopóre desídes:\\
Castíque, recti, ac sóbrii\\
Vigiláte, jam sum próximus.

Jesum ciámus vócibus,\\
Flentes, precántes, sóbrii:\\
Inténta supplicátio\\
Dormíre cor mundum vetat.

Tu, Christe, somnum díscute:\\
Tu rumpe noctis víncula:\\
Tu solve peccátum vetus,\\
Novúmque lumen íngere.

Deo Patri sit glória,\\
Ejúsque soli Fílio,\\
Cum Spíritu Paráclito,\\
Nunc et per omne sæculum.\\
Amen.

\V\ Repléti sumus mane misericórdia tua. \R\ Exsultávimus, et delectáti sumus.

\antlabel{Ad Bened. Ant.} Eréxit nobis * Dóminus cornu salútis in domo David, púeri sui.

\rubric{Réliqua ut in Ordinário. 11}

\hourheading{Ad Laudes II}

\rubric{In Vigíliis commúnibus, quando fit Offícium de Féria.}

(Pater noster et Ave.)

\V\ Deus, in adjutórium.

\antlabel{Ant.} Dele iniquitátem meam.

\psalmheading{Psalmus 50}

\lettrine[lines=2]{M}{ISERERE} \textdagger\ mei, Deus: * secúndum magnam misericórdiam tuam.

Et secúndum multitúdinem miseratiónum tuárum, * dele iniquitátem meam.

Amplius lava me ab iniquitáte mea: * et a peccáto meo munda me.

Quóniam iniquitátem meam ego cognósco: * et peccátum meum contra me est semper.

Tibi soli peccávi, et malum coram te feci: * ut justificéris in sermónibus tuis, et vincas cum judicáris.

Ecce enim, in iniquitátibus concéptus sum: * et in peccátis concépit me mater mea.

Ecce enim, veritátem dilexísti: * incérta et occúlta sapiéntiæ tuæ manifestásti mihi.

Aspérges me hyssópo, et mundábor: * lavábis me, et super nivem dealbábor.

Audítui meo dabis gáudium et lætítiam: * et exsultábunt ossa humiliáta.

Avérte fáciem tuam a peccátis meis: * et omnes iniquitátes meas dele.

Cor mundum crea in me, Deus: * et spíritum rectum ínnova in viscéribus meis.

Ne projícias me a fácie tua: * et spíritum sanctum tuum ne áuferas a me.

Redde mihi lætítiam salutáris tui: * et spíritu principáli confírma me.

Docébo iníquos vias tuas: * et ímpii ad te converténtur.

Líbera me de sanguínibus, Deus, Deus salútis meæ: * et exsultábit lingua mea justítiam tuam.

Dómine, lábia mea apéries: * et os meum annuntiábit laudem tuam.

Quóniam si voluísses sacrifícium, dedíssem útique: * holocáustis non delectáberis.

Sacrifícium Deo spíritus contribulátus: * cor contrítum, et humiliátum, Deus, non despícies.

Benígne fac, Dómine, in bona voluntáte tua Sion: * ut ædificéntur muri Jerúsalem.

Tunc acceptábis sacrifícium justítiæ, oblatiónes, et holocáusta: * tunc impónent super altáre tuum vítulos.

\antlabel{Ant.} Dele iniquitátem meam, Dómine, secúndum multitúdinem miseratiónum tuárum.

\antlabel{Ant.} Discérne causam meam.

\psalmheading{Psalmus 42}

\lettrine[lines=2]{J}{UDICA} me, Deus, et discérne causam meam de gente non sancta, * ab hómine iníquo, et dolóso érue me.

Quia tu es, Deus, fortitúdo mea: * quare me repulísti? et quare tristis incédo, dum afflígit me inimícus?

Emítte lucem tuam et veritátem tuam: * ipsa me deduxérunt, et adduxérunt in montem sanctum tuum, et in tabernácula tua.

Et introíbo ad altáre Dei: * ad Deum, qui lætíficat juventútem meam.

Confitébor tibi in cíthara, Deus, Deus meus: * quare tristis es, ánima mea? et quare contúrbas me?

Spera in Deo, quóniam adhuc confitébor illi: * salutáre vultus mei, et Deus meus.

\antlabel{Ant.} Discérne causam meam, Deus, de gente non sancta.

\antlabel{Ant.} Deus misereátur. \textdagger

\psalmheading{Psalmus 66}

\lettrine[lines=2]{D}{EUS} misereátur \textdagger\ nostri, et benedícat nobis: * illúminet vultum suum super nos, et misereátur nostri.

Ut cognoscámus in terra viam tuam: * in ómnibus géntibus salutáre tuum.

Confiteántur tibi pópuli, Deus: * confiteántur tibi pópuli omnes.

Læténtur et exsúltent gentes: * quóniam júdicas pópulos in æquitáte, et gentes in terra dírigis.

Confiteántur tibi pópuli, Deus, confiteántur tibi pópuli omnes: * terra dedit fructum suum.

Benedícat nos Deus, Deus noster, benedícat nos Deus: * et métuant eum omnes fines terræ.

\antlabel{Ant.} Deus misereátur nostri, et benedícat nos.

\antlabel{Ant.} Corrípies me, Dómine.

\canticumheading{Canticum Ezechiæ}{Is. 38, 10-20}

\lettrine[lines=2]{E}{GO} dixi: In dimídio diérum meórum * vadam ad portas ínferi.

Quæsívi resíduum annórum meórum. * Dixi: Non vidébo Dóminum Deum in terra vivéntium.

Non aspíciam hóminem ultra, * et habitatórem quiétis.

Generátio mea abláta est, et convolúta est a me, * quasi tabernáculum pastórum.

Præcísa est, velut a texénte, vita mea: dum adhuc ordírer, succídit me: * de mane usque ad vésperam fínies me.

Sperábam usque ad mane, * quasi leo sic contrívit ómnia ossa mea:

De mane usque ad vésperam fínies me: * sicut pullus hirúndinis sic clamábo, meditábor ut colúmba:

Attenuáti sunt óculi mei, * suspiciéntes in excélsum:

Dómine, vim pátior, respónde pro me. * Quid dicam, aut quid respondébit mihi, cum ipse fécerit?

Recogitábo tibi omnes annos meos * in amaritúdine ánimæ meæ.

Dómine, si sic vívitur, et in tálibus vita spíritus mei, corrípies me, et vivificábis me. * Ecce, in pace amaritúdo mea amaríssima:

Tu autem eruísti ánimam meam ut non períret: * projecísti post tergum tuum ómnia peccáta mea.

Quia non inférnus confitébitur tibi, neque mors laudábit te: * non exspectábunt qui descéndunt in lacum, veritátem tuam.

Vivens vivens ipse confitébitur tibi, sicut et ego hódie: * pater fíliis notam fáciet veritátem tuam.

Dómine, salvum me fac, * et psalmos nostros cantábimus cunctis diébus vitæ nostræ in domo Dómini.

\antlabel{Ant.} Corrípies me, Dómine, et vivificábis me.

\antlabel{Ant.} Laudáte. \textdagger

\psalmheading{Psalmus 134}

\lettrine[lines=2]{L}{AUDATE} \textdagger\ nomen Dómini, * laudáte, servi, Dóminum.

Qui statis in domo Dómini, * in átriis domus Dei nostri.

Laudáte Dóminum, quia bonus Dóminus: * psállite nómini ejus, quóniam suáve.

Quóniam Jacob elégit sibi Dóminus, * Israël in possessiónem sibi.

Quia ego cognóvi quod magnus est Dóminus: * et Deus noster præ ómnibus diis.

Omnia quæcúmque vóluit, Dóminus fecit in cælo, et in terra, * in mari, et in ómnibus abýssis.

Edúcens nubes ab extrémo terræ: * fúlgura in plúviam fecit.

Qui prodúcit ventos de thesáuris suis: * qui percússit primogénita Ægýpti ab hómine usque ad pecus.

Et misit signa, et prodígia in médio tui, Ægýpte: * in Pharaónem, et in omnes servos ejus.

Qui percússit gentes multas: * et occídit reges fortes:

Sehon, regem Amorrhæórum, et Og, regem Basan, * et ómnia regna Chánaan.

Et dedit terram eórum hereditátem, * hereditátem Israël, pópulo suo.

Dómine, nomen tuum in ætérnum: * Dómine, memoriále tuum in generatiónem et generatiónem.

Quia judicábit Dóminus pópulum suum: * et in servis suis deprecábitur.

Simulácra géntium argéntum, et aurum, * ópera mánuum hóminum.

Os habent, et non loquéntur: * óculos habent, et non vidébunt.

Aures habent, et non áudient: * neque enim est spíritus in ore ipsórum.

Símiles illis fiant qui fáciunt ea: * et omnes qui confídunt in eis.

Domus Israël, benedícite Dómino: * domus Aaron, benedícite Dómino.

Domus Levi, benedícite Dómino: * qui timétis Dóminum, benedícite Dómino.

Benedíctus Dóminus ex Sion, * qui hábitat in Jerúsalem.

\antlabel{Ant.} Laudáte Dóminum, quia benígnus est, et in servis suis deprecábitur.

\capitulumlabel{Capit.}{Rom. 13, 12-13}

\lettrine[lines=2]{N}{OX} præcéssit, dies autem appropinquávit. Abjiciámus ergo ópera tenebrárum, et induámur arma lucis. Sicut in die honéste ambulémus.

\hymnusheading{Hymnus}

\lettrine[lines=2]{A}{LES} diéi núntius\\
Lucem propínquam præcínit:\\
Nos excitátor méntium\\
Jam Christus ad vitam vocat.

Auférte, clamat, léctulos,\\
Ægro sopóre desídes:\\
Castíque, recti, ac sóbrii\\
Vigiláte, jam sum próximus.

Jesum ciámus vócibus,\\
Flentes, precántes, sóbrii:\\
Inténta supplicátio\\
Dormíre cor mundum vetat.

Tu, Christe, somnum díscute:\\
Tu rumpe noctis víncula:\\
Tu solve peccátum vetus,\\
Novúmque lumen íngere.

Deo Patri sit glória,\\
Ejúsque soli Fílio,\\
Cum Spíritu Paráclito,\\
Nunc et per omne sæculum.\\
Amen.

\V\ Repléti sumus mane misericórdia tua. \R\ Exsultávimus, et delectáti sumus.

\antlabel{Ad Bened. Ant.} Eréxit nobis * Dóminus cornu salútis in domo David, púeri tui.

\rubric{Expléto Cántico ac repetíta Antíphona, dicúntur sequéntes}

\variantheading{Preces feriales}

KYRIE, eléison. Christe, eléison. Kýrie, eléison.

Pater noster \rubric{quod dícitur a solo Hebdomadário totum clara voce, usque ad} \V\ Et ne nos indúcas in tentatiónem. \R\ Sed líbera nos a malo. \V\ Ego dixi: Dómine, miserére mei. \R\ Sana ánimam meam, quia peccávi tibi. \V\ Convértere, Dómine, úsquequo? \R\ Et deprecábilis esto super servos tuos. \V\ Fiat misericórdia tua, Dómine, super nos. \R\ Quemádmodum sperávimus in te. \V\ Sacerdótes tui induántur justítiam. \R\ Et sancti tui exsúltent. \V\ Orémus pro beatíssimo Papa nostro N. \R\ Dóminus consérvet eum, et vivíficet eum, et beátum fáciat eum in terra, et non tradat eum in ánimam inimicórum ejus. \V\ Orémus et pro Antístite nostro N. \R\ Stet et pascat in fortitúdine tua, Dómine, in sublimitáte nóminis tui. \V\ Dómine, salvum fac regem. \R\ Et exáudi nos in die, qua invocavérimus te. \V\ Salvum fac pópulum tuum, Dómine, et bénedic hereditáti tuæ. \R\ Et rege eos, et extólle illos usque in ætérnum. \V\ Meménto Congregatiónis tuæ. \R\ Quam possedísti ab inítio. \V\ Fiat pax in virtúte tua. \R\ Et abundántia in túrribus tuis. \V\ Orémus pro benefactóribus nostris. \R\ Retribúere dignáre, Dómine, ómnibus, nobis bona faciéntibus propter nomen tuum, vitam ætérnam. Amen. \V\ Orémus pro fidélibus defúnctis. \R\ Réquiem ætérnam dona eis, Dómine, et lux perpétua lúceat eis. \V\ Requiéscant in pace. \R\ Amen. \V\ Pro frátribus nostris abséntibus. \R\ Salvos fac servos tuos, Deus meus, sperántes in te. \V\ Pro afflíctis et captívis. \R\ Líbera eos, Deus Israël, ex ómnibus tribulatiónibus suis. \V\ Mitte eis, Dómine, auxílium de sancto. \R\ Et de Sion tuére eos. \V\ Dómine, Deus virtútum, convérte nos. \R\ Et osténde fáciem tuam, et salvi érimus. \V\ Exsúrge, Christe, ádjuva nos. \R\ Et líbera nos propter nomen tuum. \V\ Dómine, exáudi oratiónem meam. \R\ Et clamor meus ad te véniat.

\rubric{Réliqua ut in Ordinário. 9}

% ===== END OF BATCH 15 -- Covers autumnalis.md lines 4981-5396, the
% two alternate ferial Lauds schemes for Feria Tertia, "Ad Laudes
% I" (Psalms 95, 42, 66, Canticum Tobiæ, 134) and "Ad Laudes II"
% (Psalms 50, 42, 66, Canticum Ezechiæ, 134), each with the
% Capitulum "Nox præcéssit" and the hymn "Ales diéi núntius." The
% recurring stray OCR "\*" artifact appears only in the Ad Laudes I
% copy of the hymn (autumnalis.md line 5427, before the final doxology
% stanza) and is silently dropped there; the Ad Laudes II copy of
% the identical hymn text (line 5336) has NO such artifact in the
% source and is transcribed exactly as printed, confirming the
% artifact is a per-occurrence OCR defect rather than a feature of
% the hymn text itself. The "Ad Bened. Ant." differs by a single
% word between the two schemes ("púeri sui" in Ad Laudes I vs.
% "púeri tui" in Ad Laudes II) and each is transcribed exactly as
% its own source gives it. Ad Laudes II additionally carries the
% "Preces feriales" litany in full. The document continues with
% "Ad Primam" (autumnalis.md line 5399).
% =====
\hourheading{Ad Primam}

Pater, Ave et Credo.

\V\ Deus, in adjutórium.

\hymnusheading{Hymnus}

\lettrine[lines=2]{J}{AM} lucis orto sídere,\\
Deum precémur súpplices,\\
Ut in diúrnis áctibus\\
Nos servet a nocéntibus.

Linguam refrénans témperet,\\
Ne litis horror ínsonet:\\
Visum fovéndo cóntegat,\\
Ne vanitátes háuriat.

Sint pura cordis íntima,\\
Absístat et vecórdia:\\
Carnis terat supérbiam\\
Potus cibíque párcitas.

Ut, cum dies abscésserit\\
Noctémque sors redúxerit,\\
Mundi per abstinéntiam\\
Ipsi canámus glóriam.

Deo Patri sit glória,\\
Ejúsque soli Fílio,\\
Cum Spíritu Paráclito,\\
Nunc et per omne sæculum.\\
Amen.

\antlabel{Ant.} Deus meus.

\psalmheading{Psalmus 24, i}

\lettrine[lines=2]{A}{D} te, Dómine, levávi ánimam meam: * Deus meus, in te confído, non erubéscam.

Neque irrídeant me inimíci mei: * étenim univérsi, qui sústinent te, non confundéntur.

Confundántur omnes iníqua agéntes * supervácue.

Vias tuas, Dómine, demónstra mihi: * et sémitas tuas édoce me.

Dírige me in veritáte tua, et doce me: * quia tu es, Deus, salvátor meus, et te sustínui tota die.

Reminíscere miseratiónum tuárum, Dómine, * et misericordiárum tuárum, quæ a sæculo sunt.

Delícta juventútis meæ, * et ignorántias meas ne memíneris.

Secúndum misericórdiam tuam meménto mei tu: * propter bonitátem tuam, Dómine.

\psalmheading{Psalmus 24, ii}

\lettrine[lines=2]{D}{ULCIS} et rectus Dóminus: * propter hoc legem dabit delinquéntibus in via.

Díriget mansuétos in judício: * docébit mites vias suas.

Univérsæ viæ Dómini, misericórdia et véritas, * requiréntibus testaméntum ejus et testimónia ejus.

Propter nomen tuum, Dómine, propitiáberis peccáto meo: * multum est enim.

Quis est homo qui timet Dóminum? * legem státuit ei in via, quam elégit.

Anima ejus in bonis demorábitur: * et semen ejus hereditábit terram.

Firmaméntum est Dóminus timéntibus eum: * et testaméntum ipsíus ut manifestétur illis.

\psalmheading{Psalmus 24, iii}

\lettrine[lines=2]{O}{CULI} mei semper ad Dóminum: * quóniam ipse evéllet de láqueo pedes meos.

Réspice in me, et miserére mei: * quia únicus et pauper sum ego.

Tribulatiónes cordis mei multiplicátæ sunt: * de necessitátibus meis érue me.

Vide humilitátem meam, et labórem meum: * et dimítte univérsa delícta mea.

Réspice inimícos meos quóniam multiplicáti sunt, * et ódio iníquo odérunt me.

Custódi ánimam meam, et érue me: * non erubéscam quóniam sperávi in te.

Innocéntes et recti adhæsérunt mihi: * quia sustínui te.

Líbera, Deus, Israël, * ex ómnibus tribulatiónibus suis.

\rubric{¶ Quando ad Laudes dictus fúerit Ps. 50 Miserére, hic subjúngitur seq. Ps. 95, qui in Láudibus locum non hábuit.}

\psalmheading{Psalmus 95}

\lettrine[lines=2]{C}{ANTATE} Dómino cánticum novum: * cantáte Dómino, omnis terra.

Cantáte Dómino, et benedícite nómini ejus: * annuntiáte de die in diem salutáre ejus.

Annuntiáte inter gentes glóriam ejus, * in ómnibus pópulis mirabília ejus.

Quóniam magnus Dóminus, et laudábilis nimis: * terríbilis est super omnes deos.

Quóniam omnes dii géntium dæmónia: * Dóminus autem cælos fecit.

Conféssio, et pulchritúdo in conspéctu ejus: * sanctimónia et magnificéntia in sanctificatióne ejus.

Afférte Dómino, pátriæ géntium, afférte Dómino glóriam et honórem: * afférte Dómino glóriam nómini ejus.

Tóllite hóstias, et introíte in átria ejus: * adoráte Dóminum in átrio sancto ejus.

Commoveátur a fácie ejus univérsa terra: * dícite in géntibus quia Dóminus regnávit.

Etenim corréxit orbem terræ qui non commovébitur: * judicábit pópulos in æquitáte.

Læténtur cæli, et exsúltet terra: commoveátur mare, et plenitúdo ejus: * gaudébunt campi, et ómnia quæ in eis sunt.

Tunc exsultábunt ómnia ligna silvárum a fácie Dómini, quia venit: * quóniam venit judicáre terram.

Judicábit orbem terræ in æquitáte, * et pópulos in veritáte sua.

\rubric{Secus vero, tribus Psalmis explétis, statim dícitur Ant.}

\antlabel{Ant.} Deus meus, in te confído, non erubéscam.

\rubric{In Officio Festi vel Octávæ:}

\capitulumlabel{Capitulum}{1 Tim. 1, 17}

\lettrine[lines=2]{R}{EGI} sæculórum immortáli et invisíbili, soli Deo honor et glória in sæcula sæculórum. Amen.

\rubric{In ómnibus autem Fériis et in Vigíliis commúnibus:}

\capitulumlabel{Capitulum}{Zach. 8, 19}

\lettrine[lines=2]{P}{ACEM} et veritátem dilígite, ait Dóminus omnípotens.

\R\ br. Christe, Fili Dei vivi, * Miserére nobis. Christe. \V\ \textit{Qui sedes ad déxteram Patris.} Miserére nobis. Glória Patri. Christe. \V\ Exsúrge, Christe, ádjuva nos. \R\ Et líbera nos propter nomen tuum.

\variantheading{Preces dominicales et feriales (si dicendæ sint)}

KYRIE, eléison. Christe, eléison. Kýrie, eléison.

Pater noster \rubric{secreto usque ad} \V\ Et ne nos indúcas in tentatiónem. \R\ Sed líbera nos a malo.

Credo in Deum \rubric{secreto usque ad} \V\ Carnis resurrectiónem. \R\ Vitam ætérnam. Amen. \V\ Et ego ad te, Dómine, clamávi. \R\ Et mane orátio mea prævéniet te. \V\ Repleátur os meum laude. \R\ Ut cantem glóriam tuam, tota die magnitúdinem tuam. \V\ Dómine, avérte fáciem tuam a peccátis meis. \R\ Et omnes iniquitátes meas dele. \V\ Cor mundum crea in me, Deus. \R\ Et spíritum rectum ínnova in viscéribus meis. \V\ Ne projícias me a fácie tua. \R\ Et spíritum sanctum tuum ne áuferas a me. \V\ Redde mihi lætítiam salutáris tui. \R\ Et spíritu principáli confírma me.

\variantheading{In Précibus feriálibus tantum}

\V\ Eripe me, Dómine, ab hómine malo. \R\ A viro iníquo éripe me. \V\ Eripe me de inimícis meis, Deus meus. \R\ Et ab insurgéntibus in me líbera me. \V\ Eripe me de operántibus iniquitátem. \R\ Et de viris sánguinum salva me. \V\ Sic psalmum dicam nómini tuo in sæculum sæculi. \R\ Ut reddam vota mea de die in diem. \V\ Exáudi nos, Deus, salutáris noster. \R\ Spes ómnium fínium terræ, et in mari longe. \V\ Deus, in adjutórium meum inténde. \R\ Dómine, ad adjuvándum me festína. \V\ Sanctus Deus, Sanctus fortis, Sanctus immortális. \R\ Miserére nobis. \V\ Bénedic, ánima mea, Dómino. \R\ Et ómnia, quæ intra me sunt, nómini sancto ejus. \V\ Bénedic, ánima mea, Dómino. \R\ Et noli oblivísci omnes retributiónes ejus. \V\ Qui propitiátur ómnibus iniquitátibus tuis. \R\ Qui sanat omnes infirmitátes tuas. \V\ Qui rédimit de intéritu vitam tuam. \R\ Qui corónat te in misericórdia et miseratiónibus. \V\ Qui replet in bonis desidérium tuum. \R\ Renovábitur ut áquilæ juvéntus tua.

\variantheading{In utrisque Précibus}

\V\ Adjutórium nostrum in nómine Dómini. \R\ Qui fecit cælum et terram.

\formulalabel{Confíteor. Misereátur. Indulgéntiam.}

\V\ Dignáre, Dómine, die isto. \R\ Sine peccáto nos custodíre. \V\ Miserére nostri, Dómine. \R\ Miserére nostri. \V\ Fiat misericórdia tua, Dómine, super nos. \R\ Quemádmodum sperávimus in te. \V\ Dómine, exáudi oratiónem meam. \R\ Et clamor meus ad te véniat.

\V\ Dóminus vobíscum.

Orémus. \oratiolabel{Oratio}

\lettrine[lines=2]{D}{OMINE}, Deus omnípotens, qui ad princípium hujus diéi nos perveníre fecísti: tua nos hódie salva virtúte; ut in hac die ad nullum declinémus peccátum, sed semper ad tuam justítiam faciéndam nostra procédant elóquia, dirigántur cogitatiónes et ópera. Per Dóminum. \R\ Amen.

\V\ Dóminus vobíscum. \V\ Benedicámus Dómino.

\rubric{In Choro Martyrológium.}

\V\ Pretiósa in conspéctu Dómini. \R\ Mors Sanctórum ejus.

\rubric{Et absolúte, sine} Orémus:

\oratiolabel{Oratio}

\lettrine[lines=2]{S}{ANCTA} María, et omnes Sancti intercédant pro nobis ad Dóminum, ut nos mereámur ab eo adjuvári et salvári, qui vivit et regnat in sæcula sæculórum. \R\ Amen.

\V\ Deus, in adjutórium meum inténde. \R\ Dómine, ad adjuvándum me festína. \rubric{Et dícitur ter; último additur:} Glória Patri. Sicut erat.

Kýrie, eléison. Christe, eléison. Kýrie, eléison.

Pater noster \rubric{secreto usque ad} \V\ Et ne nos indúcas in tentatiónem. \R\ Sed líbera nos a malo. \V\ Réspice in servos tuos, Dómine, et in ópera tua, et dírige fílios eórum. \R\ Et sit splendor Dómini, Dei nostri, super nos, et ópera mánuum nostrárum dírige super nos, et opus mánuum nostrárum dírige. \V\ Glória Patri. \R\ Sicut erat.

Orémus. \oratiolabel{Oratio}

\lettrine[lines=2]{D}{IRIGERE} et sanctificáre, régere et gubernáre dignáre, Dómine Deus, Rex cæli et terræ, hódie corda et córpora nostra, sensus, sermónes et actus nostros in lege tua, et in opéribus mandatórum tuórum: ut hic, et in ætérnum, te auxiliánte, salvi et líberi esse mereámur, Salvátor mundi: Qui vivis et regnas in sæcula sæculórum. \R\ Amen.

\V\ Jube, domne.

\benedlabel{Bened.} Dies et actus nostros in sua pace dispónat Dóminus omnípotens. \R\ Amen.

\variantheading{In Officio feriali}

\lettrine[lines=2]{D}{OMINUS} autem dírigat corda et córpora nostra in caritáte Dei, et patiéntia Christi. Tu autem.

\variantheading{E Proprio de Tempore}

\variantheading{Pro Oct. Omn. Sanctorum}

\lettrine[lines=2]{B}{ENEDÍCTIO} et cláritas et sapiéntia et gratiárum áctio, honor, virtus et fortitúdo Deo nostro in sǽcula sæculórum. Amen.

\variantheading{E Communi Sanctorum}

\variantheading{Pro Apostolis}

\lettrine[lines=2]{I}{BANT} Apóstoli gaudéntes a conspéctu concílii, quóniam digni hábiti sunt pro nómine Jesu contuméliam pati.

\variantheading{Pro Martyre}

\lettrine[lines=2]{J}{USTUS} cor suum trádidit ad vigilándum dilúculo ad Dóminum, qui fecit illum, et in conspéctu Altíssimi deprecábitur.

\variantheading{Pro pluribus Martyribus}

\lettrine[lines=2]{F}{ULGEBUNT} justi, et tamquam scintíllæ in arundinéto discúrrent. Judicábunt natiónes, et dominabúntur pópulis: et regnábit Dóminus illórum in perpétuum.

\variantheading{Pro Conf. Pont.}

\lettrine[lines=2]{F}{UNGI} sacerdótio, et habére laudem in nómine ipsíus, et offérre illi incénsum dignum in odórem suavitátis.

\variantheading{Pro Conf. non Pont.}

\lettrine[lines=2]{J}{USTUM} dedúxit Dóminus per vias rectas, et osténdit illi regnum Dei, et dedit illi sciéntiam sanctórum; honestávit illum in labóribus, et complévit labóres illíus.

\variantheading{Pro Virgine Martyre}

\lettrine[lines=2]{D}{OMINE}, Deus meus, exaltásti super terram habitatiónem meam, et pro morte defluénte deprecáta sum. Invocávi Dóminum, Patrem Dómini mei, ut non derelínquat me in die tribulatiónis meæ, et in témpore superbórum sine adjutório.

\variantheading{Pro Virgine non Martyre}

\lettrine[lines=2]{O}{} \textsc{quam} pulchra est casta generátio cum claritáte! immortális est enim memória illíus: quóniam apud Deum nota est, et apud hómines.

\variantheading{Pro Martyre tantum}

\lettrine[lines=2]{L}{AUDABIT} usque ad mortem ánima mea Dóminum, quóniam éruis sustinéntes te, et líberas eos de manu angústiæ, Dómine, Deus noster.

\variantheading{Pro nec Virg. nec Mart.}

\lettrine[lines=2]{M}{ULTAE} fíliæ congregavérunt divítias: tu supergréssa es univérsas. Fallax grátia, et vana est pulchritúdo: múlier timens Dóminum, ipsa laudábitur.

\variantheading{Pro Dedicatione Ecclesiæ}

\lettrine[lines=2]{E}{T} abstérget Deus omnem lácrimam ab óculis eórum: et mors ultra non erit, neque luctus, neque clamor, neque dolor erit ultra, quia prima abiérunt. Et dixit qui sedébat in throno: Ecce, nova fácio ómnia.

\variantheading{In Officiis B. Mariæ V.}

\lettrine[lines=2]{I}{N} platéis sicut cinnamómum et bálsamum aromatízans odórem dedi: quasi myrrha elécta dedi suavitátem odóris.

\V\ Adjutórium nostrum in nómine Dómini. \R\ Qui fecit cælum et terram.

\V\ Benedícite. \R\ Deus.

\benedictiolabel{Benedictio} Dóminus nos benedícat, et ab omni malo deféndat, et ad vitam perdúcat ætérnam. Et fidélium ánimæ per misericórdiam Dei requiéscant in pace. \R\ Amen.

\rubric{Deinde dícitur tantum} Pater noster \rubric{secréto, nisi sequátur ália Hora.}

% ===== END OF BATCH 16 -- Covers autumnalis.md lines 5399-5652,
% completing "Ad Primam" for Feria Tertia in full: the hymn "Jam
% lucis orto sídere" (stray OCR "\*" before the final doxology
% stanza silently dropped), Psalms 24/i-iii, the conditional Psalm
% 95 insertion (used when Ps. 50 was said at Lauds instead -- its
% final verse here reads "et pópulos in veritáte sua," genuinely
% differing from the "veritáte tua" reading seen in the same
% psalm's earlier Ad Laudes I occurrence this batch-cycle, and
% transcribed exactly as the source gives each), both Capitula, the
% brief Responsory, the full Preces apparatus, the closing
% Orationes, and the Lectio-brevis apparatus. The source heads this
% apparatus "E Proprio de Tempore" (autumnalis.md line 5593) rather
% than "E Proprio Sanctorum" as at the parallel point in Feria
% Secunda, despite containing the same single "Pro Oct. Omn.
% Sanctorum" heading used throughout this volume's substitute for
% aestiva.tex's two-item Summer list; this is transcribed exactly as
% the source heads it, confirmed via direct inspection of
% autumnalis.md at this point. The document continues with "Ad
% Tertiam" (autumnalis.md line 5655).
% =====
\hourheading{Ad Tertiam}

Pater noster et Ave.

\V\ Deus, in adjutórium.

\hymnusheading{Hymnus}

\lettrine[lines=2]{N}{UNC}, Sancte, nobis, Spíritus,\\
Unum Patri cum Fílio,\\
Dignáre promptus íngeri\\
Nostro refúsus péctori.

Os, lingua, mens, sensus, vigor\\
Confessiónem pérsonent,\\
Flamméscat igne cáritas,\\
Accéndat ardor próximos.

Præsta, Pater piíssime,\\
Patríque compar Unice,\\
Cum Spíritu Paráclito\\
Regnans per omne sæculum.\\
Amen.

\antlabel{Ant.} Respéxit me.

\psalmheading{Psalmus 39, i}

\lettrine[lines=2]{E}{XSPECTANS} exspectávi Dóminum, * et inténdit mihi.

Et exaudívit preces meas: * et edúxit me de lacu misériæ, et de luto fæcis.

Et státuit super petram pedes meos: * et diréxit gressus meos.

Et immísit in os meum cánticum novum, * carmen Deo nostro.

Vidébunt multi, et timébunt: * et sperábunt in Dómino.

Beátus vir, cujus est nomen Dómini spes ejus: * et non respéxit in vanitátes et insánias falsas.

Multa fecísti tu, Dómine, Deus meus, mirabília tua: * et cogitatiónibus tuis non est qui símilis sit tibi.

Annuntiávi et locútus sum: * multiplicáti sunt super númerum.

Sacrifícium et oblatiónem noluísti: * aures autem perfecísti mihi.

Holocáustum et pro peccáto non postulásti: * tunc dixi: Ecce, vénio.

In cápite libri scriptum est de me ut fácerem voluntátem tuam: * Deus meus, vólui, et legem tuam in médio cordis mei.

\psalmheading{Psalmus 39, ii}

\lettrine[lines=2]{A}{NNUNTIAVI} justítiam tuam in ecclésia magna, * ecce, lábia mea non prohibébo: Dómine, tu scisti.

Justítiam tuam non abscóndi in corde meo: * veritátem tuam et salutáre tuum dixi.

Non abscóndi misericórdiam tuam et veritátem tuam * a concílio multo.

Tu autem, Dómine, ne longe fácias miseratiónes tuas a me: * misericórdia tua et véritas tua semper suscepérunt me.

Quóniam circumdedérunt me mala, quorum non est númerus: * comprehendérunt me iniquitátes meæ, et non pótui ut vidérem.

Multiplicátæ sunt super capíllos cápitis mei: * et cor meum derelíquit me.

Compláceat tibi, Dómine, ut éruas me: * Dómine, ad adjuvándum me réspice.

\psalmheading{Psalmus 39, iii}

\lettrine[lines=2]{C}{ONFUNDANTUR} et revereántur simul, qui quærunt ánimam meam, * ut áuferant eam.

Convertántur retrórsum, et revereántur, * qui volunt mihi mala.

Ferant conféstim confusiónem suam, * qui dicunt mihi: Euge, euge.

Exsúltent et læténtur super te omnes quæréntes te: * et dicant semper: Magnificétur Dóminus: qui díligunt salutáre tuum.

Ego autem mendícus sum, et pauper: * Dóminus sollícitus est mei.

Adjútor meus, et protéctor meus tu es: * Deus meus, ne tardáveris.

\antlabel{Ant.} Respéxit me, et exaudívit deprecatiónem meam Dóminus.

\variantheading{In Officio feriali}

\capitulumlabel{Capitulum}{Jer. 17, 14}

\lettrine[lines=2]{S}{ANA} me, Dómine, et sanábor: salvum me fac, et salvus ero: quóniam laus mea tu es.

\R\ br. Sana ánimam meam, * Quia peccávi tibi. Sana. \V\ Ego dixi: Dómine, miserére mei. Quia. Glória Patri. Sana. \V\ Adjútor meus esto, ne derelínquas me. \R\ Neque despícias me, Deus, salutáris meus.

\variantheading{In Vigiliis communibus:}

\variantheading{Preces feriales}

KYRIE, eléison. Christe, eléison. Kýrie, eléison.

Pater noster \rubric{secreto usque ad} \V\ Et ne nos indúcas in tentatiónem. \R\ Sed líbera nos a malo. \V\ Dómine, Deus virtútum, convérte nos. \R\ Et osténde fáciem tuam, et salvi érimus. \V\ Exsúrge, Christe, ádjuva nos. \R\ Et líbera nos propter nomen tuum. \V\ Dómine, exáudi oratiónem meam. \R\ Et clamor meus ad te véniat.

\rubric{Réliqua ut in Ordinário. 19}

\hourheading{Ad Sextam}

Pater noster et Ave.

\V\ Deus, in adjutórium.

\hymnusheading{Hymnus}

\lettrine[lines=2]{R}{ECTOR} potens, verax Deus,\\
Qui témperas rerum vices,\\
Splendóre mane illúminas,\\
Et ígnibus merídiem:

Exstíngue flammas lítium,\\
Aufer calórem nóxium,\\
Confer salútem córporum,\\
Verámque pacem córdium.

Præsta, Pater piíssime,\\
Patríque compar Unice,\\
Cum Spíritu Paráclito\\
Regnans per omne sæculum.\\
Amen.

\antlabel{Ant.} Suscepísti me, Dómine.

\psalmheading{Psalmus 40}

\lettrine[lines=2]{B}{EATUS} qui intélligit super egénum, et páuperem: * in die mala liberábit eum Dóminus.

Dóminus consérvet eum, et vivíficet eum, et beátum fáciat eum in terra: * et non tradat eum in ánimam inimicórum ejus.

Dóminus opem ferat illi super lectum dolóris ejus: * univérsum stratum ejus versásti in infirmitáte ejus.

Ego dixi: Dómine, miserére mei: * sana ánimam meam, quia peccávi tibi.

Inimíci mei dixérunt mala mihi: * Quando moriétur, et períbit nomen ejus?

Et si ingrediebátur ut vidéret, vana loquebátur: * cor ejus congregávit iniquitátem sibi.

Egrediebátur foras, * et loquebátur in idípsum.

Advérsum me susurrábant omnes inimíci mei: * advérsum me cogitábant mala mihi.

Verbum iníquum constituérunt advérsum me: * Numquid qui dormit non adjíciet ut resúrgat?

Etenim homo pacis meæ, in quo sperávi: * qui edébat panes meos, magnificávit super me supplantatiónem.

Tu autem, Dómine, miserére mei, et resúscita me: * et retríbuam eis.

In hoc cognóvi quóniam voluísti me: * quóniam non gaudébit inimícus meus super me.

Me autem propter innocéntiam suscepísti: * et confirmásti me in conspéctu tuo in ætérnum.

Benedíctus Dóminus, Deus Israël, a sæculo et usque in sæculum: * fiat, fiat.

\psalmheading{Psalmus 41, i}

\lettrine[lines=2]{Q}{UEMADMODUM} desíderat cervus ad fontes aquárum: * ita desíderat ánima mea ad te, Deus.

Sitívit ánima mea ad Deum fortem vivum: * quando véniam, et apparébo ante fáciem Dei?

Fuérunt mihi lácrimæ meæ panes die ac nocte: * dum dícitur mihi quotídie: Ubi est Deus tuus?

Hæc recordátus sum, et effúdi in me ánimam meam: * quóniam transíbo in locum tabernáculi admirábilis, usque ad domum Dei.

In voce exsultatiónis, et confessiónis: * sonus epulántis.

Quare tristis es, ánima mea? * et quare contúrbas me?

Spera in Deo, quóniam adhuc confitébor illi: * salutáre vultus mei, et Deus meus.

\psalmheading{Psalmus 41, ii}

\lettrine[lines=2]{A}{D} meípsum ánima mea conturbáta est: * proptérea memor ero tui de terra Jordánis, et Hermóniim a monte módico.

Abýssus abýssum ínvocat, * in voce cataractárum tuárum.

Omnia excélsa tua, et fluctus tui * super me transiérunt.

In die mandávit Dóminus misericórdiam suam: * et nocte cánticum ejus.

Apud me orátio Deo vitæ meæ, * dicam Deo: Suscéptor meus es.

Quare oblítus es mei? * et quare contristátus incédo, dum afflígit me inimícus?

Dum confringúntur ossa mea, * exprobravérunt mihi qui tríbulant me inimíci mei.

Dum dicunt mihi per síngulos dies: Ubi est Deus tuus? * quare tristis es, ánima mea? et quare contúrbas me?

Spera in Deo, quóniam adhuc confitébor illi: * salutáre vultus mei, et Deus meus.

\antlabel{Ant.} Suscepísti me, Dómine: et confirmásti me in conspéctu tuo.

\variantheading{In Officio feriali}

\capitulumlabel{Capitulum}{Rom. 13, 8}

\lettrine[lines=2]{N}{EMINI} quidquam debeátis, nisi ut ínvicem diligátis: qui enim díligit próximum, legem implévit.

\R\ br. Benedícam Dóminum * In omni témpore. Benedícam. \V\ Semper laus ejus in ore meo. In. Glória Patri. Benedícam. \V\ Dóminus regit me, et nihil mihi déerit. \R\ In loco páscuæ ibi me collocávit.

\variantheading{In Vigiliis communibus:}

\variantheading{Preces feriales}

KYRIE, eléison. Christe, eléison. Kýrie, eléison.

Pater noster \rubric{secreto usque ad} \V\ Et ne nos indúcas in tentatiónem. \R\ Sed líbera nos a malo. \V\ Dómine, Deus virtútum, convérte nos. \R\ Et osténde fáciem tuam, et salvi érimus. \V\ Exsúrge, Christe, ádjuva nos. \R\ Et líbera nos propter nomen tuum. \V\ Dómine, exáudi oratiónem meam. \R\ Et clamor meus ad te véniat.

\rubric{Réliqua ut in Ordinário. 20}

\hourheading{Ad Nonam}

Pater noster et Ave.

\V\ Deus, in adjutórium.

\hymnusheading{Hymnus}

\lettrine[lines=2]{R}{ERUM}, Deus, tenax vigor,\\
Immótus in te pérmanens,\\
Lucis diúrnæ témpora\\
Succéssibus detérminans:

Largíre lumen véspere,\\
Quo vita nusquam décidat,\\
Sed prǽmium mortis sacræ\\
Perénnis instet glória.

Præsta, Pater piíssime,\\
Patríque compar Unice,\\
Cum Spíritu Paráclito\\
Regnans per omne sæculum.\\
Amen.

\antlabel{Ant.} Salvásti nos.

\psalmheading{Psalmus 43, i}

\lettrine[lines=2]{D}{EUS}, áuribus nostris audívimus: * patres nostri annuntiavérunt nobis.

Opus, quod operátus es in diébus eórum, * et in diébus antíquis.

Manus tua gentes dispérdidit, et plantásti eos: * afflixísti pópulos, et expulísti eos.

Nec enim in gládio suo possedérunt terram, * et bráchium eórum non salvávit eos:

Sed déxtera tua, et bráchium tuum, et illuminátio vultus tui: * quóniam complacuísti in eis.

Tu es ipse Rex meus et Deus meus: * qui mandas salútes Jacob.

In te inimícos nostros ventilábimus cornu: * et in nómine tuo spernémus insurgéntes in nobis.

Non enim in arcu meo sperábo: * et gládius meus non salvábit me.

Salvásti enim nos de affligéntibus nos: * et odiéntes nos confudísti.

In Deo laudábimur tota die: * et in nómine tuo confitébimur in sæculum.

\psalmheading{Psalmus 43, ii}

\lettrine[lines=2]{N}{UNC} autem repulísti et confudísti nos: * et non egrediéris, Deus, in virtútibus nostris.

Avertísti nos retrórsum post inimícos nostros: * et qui odérunt nos, diripiébant sibi.

Dedísti nos tamquam oves escárum: * et in géntibus dispersísti nos.

Vendidísti pópulum tuum sine prétio: * et non fuit multitúdo in commutatiónibus eórum.

Posuísti nos oppróbrium vicínis nostris, * subsannatiónem et derísum his, qui sunt in circúitu nostro.

Posuísti nos in similitúdinem géntibus: * commotiónem cápitis in pópulis.

Tota die verecúndia mea contra me est, * et confúsio faciéi meæ coopéruit me.

A voce exprobrántis, et obloquéntis: * a fácie inimíci, et persequéntis.

Hæc ómnia venérunt super nos, nec oblíti sumus te: * et iníque non égimus in testaménto tuo.

Et non recéssit retro cor nostrum: * et declinásti sémitas nostras a via tua:

Quóniam humiliásti nos in loco afflictiónis, * et coopéruit nos umbra mortis.

\psalmheading{Psalmus 43, iii}

\lettrine[lines=2]{S}{I} oblíti sumus nomen Dei nostri, * et si expándimus manus nostras ad deum aliénum:

Nonne Deus requíret ista? * ipse enim novit abscóndita cordis.

Quóniam propter te mortificámur tota die: * æstimáti sumus sicut oves occisiónis.

Exsúrge, quare obdórmis, Dómine? * exsúrge, et ne repéllas in finem.

Quare fáciem tuam avértis, * oblivísceris inópiæ nostræ, et tribulatiónis nostræ?

Quóniam humiliáta est in púlvere ánima nostra: * conglutinátus est in terra venter noster.

Exsúrge, Dómine, ádjuva nos: * et rédime nos propter nomen tuum.

\antlabel{Ant.} Salvásti nos, Dómine: et in nómine tuo confitébimur in sæcula.

\variantheading{In Officio feriali}

\capitulumlabel{Capit.}{1 Petri 1, 17-19}

\lettrine[lines=2]{I}{N} timóre incolátus vestri témpore conversámini: sciéntes quod non corruptibílibus auro vel argénto redémpti estis, sed pretióso sánguine quasi Agni immaculáti Christi.

\R\ br. Rédime me, Dómine, * Et miserére mei. Rédime. \V\ Pes enim meus stetit in via recta. Et miserére mei. Glória Patri. Rédime. \V\ Ab occúltis meis munda me, Dómine. \R\ Et ab aliénis parce servo tuo.

\variantheading{In Vigiliis communibus:}

\variantheading{Preces feriales}

KYRIE, eléison. Christe, eléison. Kýrie, eléison.

Pater noster \rubric{secreto usque ad} \V\ Et ne nos indúcas in tentatiónem. \R\ Sed líbera nos a malo. \V\ Dómine, Deus virtútum, convérte nos. \R\ Et osténde fáciem tuam, et salvi érimus. \V\ Exsúrge, Christe, ádjuva nos. \R\ Et líbera nos propter nomen tuum. \V\ Dómine, exáudi oratiónem meam. \R\ Et clamor meus ad te véniat.

\rubric{Réliqua ut in Ordinário. 22}

% ===== END OF BATCH 17 -- Covers autumnalis.md lines 5655-5978,
% completing the Little Hours of Feria Tertia in full: "Ad
% Tertiam" (Psalms 39/i-iii), "Ad Sextam" (Psalms 40, 41/i, 41/ii),
% and "Ad Nonam" (Psalms 43/i-iii), each with its own hymn,
% Capitulum, brief Responsory, and "In Vigiliis communibus: Preces
% feriales" apparatus. Three further occurrences of the recurring
% stray OCR "\*" artifact (in the hymns "Nunc, Sancte, nobis,
% Spíritus," "Rector potens, verax Deus," and "Rerum, Deus, tenax
% vigor," each before the final "Præsta, Pater piíssime" doxology
% stanza) are silently dropped, consistent with every prior
% occurrence in this volume. The document continues with "Ad
% Vesperas" (autumnalis.md line 5981).
% =====
\hourheading{Ad Vesperas}

Pater noster et Ave.

\V\ Deus, in adjutórium.

\antlabel{Ant.} Qui hábitas in cælis, * miserére nobis.

\psalmheading{Psalmus 122}

\lettrine[lines=2]{A}{D} te levávi óculos meos, * qui hábitas in cælis.

Ecce, sicut óculi servórum * in mánibus dominórum suórum.

Sicut óculi ancíllæ in mánibus dóminæ suæ: * ita óculi nostri ad Dóminum, Deum nostrum, donec misereátur nostri.

Miserére nostri, Dómine, miserére nostri: * quia multum repléti sumus despectióne:

Quia multum repléta est ánima nostra: * oppróbrium abundántibus, et despéctio supérbis.

\antlabel{Ant.} Qui hábitas in cælis, miserére nobis.

\antlabel{Ant.} Adjutórium nostrum * in nómine Dómini.

\psalmheading{Psalmus 123}

\lettrine[lines=2]{N}{ISI} quia Dóminus erat in nobis, dicat nunc Israël: * nisi quia Dóminus erat in nobis.

Cum exsúrgerent hómines in nos, * forte vivos deglutíssent nos:

Cum irascerétur furor eórum in nos, * fórsitan aqua absorbuísset nos.

Torréntem pertransívit ánima nostra: * fórsitan pertransísset ánima nostra aquam intolerábilem.

Benedíctus Dóminus * qui non dedit nos in captiónem déntibus eórum.

Anima nostra sicut passer erépta est * de láqueo venántium:

Láqueus contrítus est, * et nos liberáti sumus.

Adjutórium nostrum in nómine Dómini, * qui fecit cælum et terram.

\antlabel{Ant.} Adjutórium nostrum in nómine Dómini.

\antlabel{Ant.} In circúitu pópuli sui * Dóminus, ex hoc nunc et usque in sæculum.

\psalmheading{Psalmus 124}

\lettrine[lines=2]{Q}{UI} confídunt in Dómino, sicut mons Sion: * non commovébitur in ætérnum, qui hábitat in Jerúsalem.

Montes in circúitu ejus: * et Dóminus in circúitu pópuli sui, ex hoc nunc et usque in sæculum.

Quia non relínquet Dóminus virgam peccatórum super sortem justórum: * ut non exténdant justi ad iniquitátem manus suas.

Bénefac, Dómine, bonis, * et rectis corde.

Declinántes autem in obligatiónes addúcet Dóminus cum operántibus iniquitátem: * pax super Israël.

\antlabel{Ant.} In circúitu pópuli sui Dóminus, ex hoc nunc et usque in sæculum.

\antlabel{Ant.} Magnificávit Dóminus * fácere nobíscum: facti sumus lætántes.

\psalmheading{Psalmus 125}

\lettrine[lines=2]{I}{N} converténdo Dóminus captivitátem Sion: * facti sumus sicut consoláti:

Tunc replétum est gáudio os nostrum: * et lingua nostra exsultatióne.

Tunc dicent inter gentes: * Magnificávit Dóminus fácere cum eis.

Magnificávit Dóminus fácere nobíscum: * facti sumus lætántes.

Convérte, Dómine, captivitátem nostram, * sicut torrens in Austro.

Qui séminant in lácrimis, * in exsultatióne metent.

Eúntes ibant et flebant, * mitténtes sémina sua.

Veniéntes autem vénient cum exsultatióne, * portántes manípulos suos.

\antlabel{Ant.} Magnificávit Dóminus fácere nobíscum: facti sumus lætántes.

\antlabel{Ant.} Dóminus ædíficet * nobis domum, et custódiat civitátem.

\psalmheading{Psalmus 126}

\lettrine[lines=2]{N}{ISI} Dóminus ædificáverit domum, * in vanum laboravérunt qui ædíficant eam.

Nisi Dóminus custodíerit civitátem, * frustra vígilat qui custódit eam.

Vanum est vobis ante lucem súrgere: * súrgite postquam sedéritis, qui manducátis panem dolóris.

Cum déderit diléctis suis somnum: * ecce heréditas Dómini fílii: merces, fructus ventris.

Sicut sagíttæ in manu poténtis: * ita fílii excussórum.

Beátus vir, qui implévit desidérium suum ex ipsis: * non confundétur cum loquétur inimícis suis in porta.

\antlabel{Ant.} Dóminus ædíficet nobis domum, et custódiat civitátem.

\variantheading{In Officio feriali}

\capitulumlabel{Capitulum}{2 Cor. 1, 3-4}

\lettrine[lines=2]{B}{ENEDICTUS} Deus, et Pater Dómini nostri Jesu Christi, Pater misericordiárum, et Deus totíus consolatiónis, qui consolátur nos in omni tribulatióne nostra.

\hymnusheading{Hymnus}

\lettrine[lines=2]{T}{ELLURIS} alme Cónditor\\
Mundi solum qui séparans,\\
Pulsis aquæ moléstiis,\\
Terram dedísti immóbilem:

Ut germen aptum próferens,\\
Fulvis decóra flóribus,\\
Fecúnda fructu sísteret,\\
Pastúmque gratum rédderet.

Mentis perústæ vúlnera\\
Munda viróre grátiæ:\\
Ut facta fletu díluat,\\
Motúsque pravos átterat.

Jussis tuis obtémperet:\\
Nullis malis appróximet:\\
Bonis repléri gáudeat,\\
Et mortis ictum nésciat.

Præsta, Pater piíssime,\\
Patríque compar Unice,\\
Cum Spíritu Paráclito\\
Regnans per omne sæculum.\\
Amen.

\V\ Dirigátur, Dómine, orátio mea. \R\ Sicut incénsum in conspéctu tuo.

\antlabel{Ad Magnif. Ant.} Exsultávit * spíritus meus in Deo, salutári meo.

\rubric{Réliqua ut in Ordinário. 22}

\hourheading{Ad Completorium}

\V\ Jube, domne.

\benedlabel{Bened.} Noctem quiétam, et finem perféctum concédat nobis Dóminus omnípotens. \R\ Amen.

\lectioheading{Lectio br.}{1 Petri 5, 8-9}

\lettrine[lines=2]{F}{RATRES}: Sóbrii estóte, et vigiláte: quia adversárius vester diábolus tamquam leo rúgiens círcuit, quærens quem dévoret: cui resístite fortes in fide. Tu autem, Dómine, miserére nobis. \R\ Deo grátias.

\V\ Adjutórium nostrum in nómine Dómini. \R\ Qui fecit cælum et terram.

Pater noster \rubric{quod dícitur totum secréto.}

\formulalabel{Confíteor.} Misereátur. Indulgéntiam.

\V\ Convérte nos, Deus, salutáris noster. \R\ Et avérte iram tuam a nobis.

\V\ Deus, in adjutórium.

\antlabel{Ant.} Tu Dómine.

\psalmheading{Psalmus 11}

\lettrine[lines=2]{S}{ALVUM} me fac, Dómine, quóniam defécit sanctus: * quóniam diminútæ sunt veritátes a fíliis hóminum.

Vana locúti sunt unusquísque ad próximum suum: * lábia dolósa, in corde et corde locúti sunt.

Dispérdat Dóminus univérsa lábia dolósa, * et linguam magníloquam.

Qui dixérunt: Linguam nostram magnificábimus, lábia nostra a nobis sunt, * quis noster Dóminus est?

Propter misériam ínopum, et gémitum páuperum, * nunc exsúrgam, dicit Dóminus.

Ponam in salutári: * fiduciáliter agam in eo.

Elóquia Dómini, elóquia casta: * argéntum igne examinátum, probátum terræ purgátum séptuplum.

Tu, Dómine, servábis nos: et custódies nos * a generatióne hac in ætérnum.

In circúitu ímpii ámbulant: * secúndum altitúdinem tuam multiplicásti fílios hóminum.

\psalmheading{Psalmus 12}

\lettrine[lines=2]{U}{SQUEQUO}, Dómine, oblivascéris me in finem? * Usquequo avértis fáciem tuam a me?

Quámdiu ponam consília in ánima mea, * dolórem in corde meo per diem?

Usquequo exaltábitur inimícus meus super me? * réspice, et exáudi me, Dómine, Deus meus.

Illúmina óculos meos ne umquam obdórmiam in morte: * nequándo dicat inimícus meus: Præválui advérsus eum.

Qui tríbulant me, exsultábunt si motus fúero: * ego autem in misericórdia tua sperávi.

Exsultábit cor meum in salutári tuo: cantábo Dómino qui bona tríbuit mihi: * et psallam nómini Dómini altíssimi.

\psalmheading{Psalmus 15}

\lettrine[lines=2]{C}{ONSERVA} me, Dómine, quóniam sperávi in te. * Dixi Dómino: Deus meus es tu, quóniam bonórum meórum non eges.

Sanctis, qui sunt in terra ejus, * mirificávit omnes voluntátes meas in eis.

Multiplicátæ sunt infirmitátes eórum: * póstea acceleravérunt.

Non congregábo conventícula eórum de sanguínibus, * nec memor ero nóminum eórum per lábia mea.

Dóminus pars hereditátis meæ, et cálicis mei: * tu es, qui restítues hereditátem meam mihi.

Funes cecidérunt mihi in præcláris: * étenim heréditas mea præclára est mihi.

Benedícam Dóminum, qui tríbuit mihi intelléctum: * ínsuper et usque ad noctem increpuérunt me renes mei.

Providébam Dóminum in conspéctu meo semper: * quóniam a dextris est mihi, ne commóvear.

Propter hoc lætátum est cor meum, et exsultávit lingua mea: * ínsuper et caro mea requiéscet in spe.

Quóniam non derelínques ánimam meam in inférno: * nec dabis sanctum tuum vidére corruptiónem.

Notas mihi fecísti vias vitæ, adimplébis me lætítia cum vultu tuo: * delectatiónes in déxtera tua usque in finem.

\antlabel{Ant.} Tu, Dómine, servábis nos: et custódies nos in ætérnum.

\hymnusheading{Hymnus}

\lettrine[lines=2]{T}{E} lucis ante términum,\\
Rerum Creátor, póscimus,\\
Ut pro tua cleméntia\\
Sis præsul et custódia.

Procul recédant sómnia,\\
Et nóctium phantásmata,\\
Hostémque nostrum cómprime,\\
Ne polluántur córpora.

Præsta, Pater piíssime,\\
Patríque compar Unice,\\
Cum Spíritu Paráclito\\
Regnans per omne sæculum.\\
Amen.

\capitulumlabel{Capitulum}{Jer. 14, 9}

\lettrine[lines=2]{T}{U} autem in nobis es, Dómine, et nomen sanctum tuum invocátum est super nos: ne derelínquas nos, Dómine, Deus noster.

\R\ br. In manus tuas, Dómine, * Comméndo spíritum meum. In manus. \V\ Redemísti nos, Dómine, Deus veritátis. Comméndo. Glória Patri. In manus.

\V\ Custódi nos, Dómine, ut pupíllam óculi. \R\ Sub umbra alárum tuárum prótege nos.

\antlabel{Ant.} Salva nos.

\canticumheading{Canticum Simeonis}{Luc. 2, 29-32}

\lettrine[lines=2]{N}{UNC} dimíttis servum tuum, Dómine, * secúndum verbum tuum in pace:

Quia vidérunt óculi mei * salutáre tuum,

Quod parásti * ante fáciem ómnium populórum,

Lumen ad revelatiónem géntium, * et glóriam plebis tuæ Israël.

Glória Patri.

\antlabel{Ant.} Salva nos, Dómine, vigilántes, custódi nos dormiéntes; ut vigilémus cum Christo, et requiescámus in pace.

\variantheading{Preces (si dicendæ sint)}

KYRIE, eléison. Christe, eléison. Kýrie, eléison.

Pater noster \rubric{secreto usque ad} \V\ Et ne nos indúcas in tentatiónem. \R\ Sed líbera nos a malo.

Credo in Deum \rubric{secreto usque ad} \V\ Carnis resurrectiónem. \R\ Vitam ætérnam. Amen. \V\ Benedíctus es, Dómine, Deus patrum nostrórum. \R\ Et laudábilis et gloriósus in sæcula. \V\ Benedicámus Patrem et Fílium cum Sancto Spíritu. \R\ Laudémus, et superexaltémus eum in sæcula. \V\ Benedíctus es, Dómine, in firmaménto cæli. \R\ Et laudábilis, et gloriósus, et superexaltátus in sæcula. \V\ Benedícat et custódiat nos omnípotens et miséricors Dóminus. \R\ Amen. \V\ Dignáre, Dómine, nocte ista. \R\ Sine peccáto nos custodíre. \V\ Miserére nostri, Dómine. \R\ Miserére nostri. \V\ Fiat misericórdia tua, Dómine, super nos. \R\ Quemádmodum sperávimus in te. \V\ Dómine, exáudi oratiónem meam. \R\ Et clamor meus ad te véniat.

\V\ Dóminus vobíscum.

Orémus. \oratiolabel{Oratio}

\lettrine[lines=2]{V}{ISITA}, quæsumus, Dómine, habitatiónem istam, et omnes insídias inimíci ab ea longe repélle: Angeli tui sancti hábitent in ea, qui nos in pace custódiant; et benedíctio tua sit super nos semper. Per Dóminum nostrum. \R\ Amen.

\V\ Dóminus vobíscum.

\V\ Benedicámus Dómino.

\benedictiolabel{Benedictio} Benedícat et custódiat nos omnípotens et miséricors Dóminus, Pater, et Fílius, et Spíritus Sanctus. \R\ Amen.

\variantheading{Antiphona finalis B. M. V.}

\V\ Divínum auxílium.

Pater, Ave et Credo.

% ===== END OF BATCH 18 -- Covers autumnalis.md lines 5981-6274,
% completing "Ad Vesperas" (the Gradual Psalms 122-126, Capitulum,
% hymn "Telluris alme Cónditor") and "Ad Completorium" (Psalms 11,
% 12, 15, hymn "Te lucis ante términum," Canticum Simeonis, Preces),
% and with them the ENTIRE "Feria Tertia" Office. Two more
% occurrences of the recurring stray OCR "\*" artifact (in "Telluris
% alme Cónditor" at autumnalis.md line 6109 and "Te lucis ante
% términum" at line 6217, both at the hymn's final doxology stanza)
% are silently dropped. The partial-bold "Confíteor." (not the full
% "Misereátur. Indulgéntiam.") again appears at Compline (line
% 6137), matching the alternation already documented in Feria
% Secunda's Compline, and is rendered the same way:
% \formulalabel{Confíteor.} followed by plain text. The document
% continues with "Feria Quarta" (autumnalis.md line 6277), the third
% weekday of the Psalterium cycle.
% =====
\feriaheading{Feria Quarta}

\hourheading{Ad Matutinum}

Pater, Ave et Credo.

\V\ Dómine, lábia mea.

\V\ Deus, in adjutórium.

\variantheading{In Officio feriali}

\invitatlabel{Invitat.} Deum magnum Dóminum, * Veníte, adorémus. \psalmlabel{Ps. 94} Veníte.

\hymnusheading{Hymnus}

\lettrine[lines=2]{R}{ERUM} Creátor óptime,\\
Rectórque noster, áspice:\\
Nos a quiéte nóxia\\
Mersos sopóre líbera.

Te, sancte Christe, póscimus,\\
Ignósce culpis ómnibus:\\
Ad confiténdum súrgimus,\\
Morásque noctis rúmpimus.

Mentes manúsque tóllimus,\\
Prophéta sicut nóctibus\\
Nobis geréndum præcipit,\\
Paulúsque gestis cénsuit.

Vides malum, quod fécimus:\\
Occúlta nostra pándimus:\\
Preces geméntes fúndimus,\\
Dimítte quod peccávimus.

Præsta, Pater piíssime,\\
Patríque compar Unice,\\
Cum Spíritu Paráclito\\
Regnans per omne sæculum.\\
Amen.

\nocturnoheading{In I Nocturno}

\antlabel{Ant.} Speciósus forma * præ fíliis hóminum, diffúsa est grátia in lábiis tuis.

\psalmheading{Psalmus 44, i}

\lettrine[lines=2]{E}{RUCTAVIT} cor meum verbum bonum: * dico ego ópera mea Regi.

Lingua mea cálamus scribæ: * velóciter scribéntis.

Speciósus forma præ fíliis hóminum, diffúsa est grátia in lábiis tuis: * proptérea benedíxit te Deus in ætérnum.

Accíngere gládio tuo super femur tuum, * potentíssime.

Spécie tua et pulchritúdine tua * inténde, próspere procéde, et regna.

Propter veritátem, et mansuetúdinem, et justítiam: * et dedúcet te mirabíliter déxtera tua.

Sagíttæ tuæ acútæ, pópuli sub te cadent, * in corda inimicórum Regis.

Sedes tua, Deus, in sæculum sæculi: * virga directiónis virga regni tui.

Dilexísti justítiam, et odísti iniquitátem: * proptérea unxit te Deus, Deus tuus, óleo lætítiæ præ consórtibus tuis.

Myrrha, et gutta, et cásia a vestiméntis tuis, a dómibus ebúrneis: * ex quibus delectavérunt te fíliæ regum in honóre tuo.

Astitit regína a dextris tuis in vestítu deauráto: * circúmdata varietáte.

\antlabel{Ant.} Speciósus forma præ fíliis hóminum, diffúsa est grátia in lábiis tuis.

\antlabel{Ant.} Confitebúntur tibi * pópuli, Deus, in ætérnum.

\psalmheading{Psalmus 44, ii}

\lettrine[lines=2]{A}{UDI}, fília, et vide, et inclína aurem tuam: * et oblivíscere pópulum tuum, et domum patris tui.

Et concupíscet Rex decórem tuum: * quóniam ipse est Dóminus, Deus tuus, et adorábunt eum.

Et fíliæ Tyri in munéribus * vultum tuum deprecabúntur: omnes dívites plebis.

Omnis glória ejus fíliæ Regis ab intus, * in fímbriis áureis circumamícta varietátibus.

Adducéntur Regi vírgines post eam: * próximæ ejus afferéntur tibi.

Afferéntur in lætítia et exsultatióne: * adducéntur in templum Regis.

Pro pátribus tuis nati sunt tibi fílii: * constítues eos príncipes super omnem terram.

Mémores erunt nóminis tui: * in omni generatióne et generatiónem.

Proptérea pópuli confitebúntur tibi in ætérnum: * et in sæculum sæculi.

\antlabel{Ant.} Confitebúntur tibi pópuli, Deus, in ætérnum.

\antlabel{Ant.} Adjútor in tribulatiónibus * Deus noster.

\psalmheading{Psalmus 45}

\lettrine[lines=2]{D}{EUS} noster refúgium, et virtus: * adjútor in tribulatiónibus, quæ invenérunt nos nimis.

Proptérea non timébimus dum turbábitur terra: * et transferéntur montes in cor maris.

Sonuérunt, et turbátæ sunt aquæ eórum: * conturbáti sunt montes in fortitúdine ejus.

Flúminis ímpetus lætíficat civitátem Dei: * sanctificávit tabernáculum suum Altíssimus.

Deus in médio ejus, non commovébitur: * adjuvábit eam Deus mane dilúculo.

Conturbátæ sunt gentes, et inclináta sunt regna: * dedit vocem suam, mota est terra.

Dóminus virtútum nobíscum: * suscéptor noster Deus Jacob.

Veníte, et vidéte ópera Dómini, quæ pósuit prodígia super terram: * áuferens bella usque ad finem terræ.

Arcum cónteret, et confrínget arma: * et scuta combúret igni.

Vacáte, et vidéte quóniam ego sum Deus: * exaltábor in géntibus, et exaltábor in terra.

Dóminus virtútum nobíscum: * suscéptor noster Deus Jacob.

\antlabel{Ant.} Adjútor in tribulatiónibus Deus noster.

\variantheading{In Festis novem Lectionum:}

\V\ Dóminus virtútum nobíscum. \R\ Suscéptor noster Deus Jacob.

\absolutiolabel{Absolutio.} Exáudi, Dómine.

\benedlabel{Bened.} 1 Benedictióne perpétua. 2 Unigénitus Dei Fílius. 3 Spíritus Sancti grátia.

\nocturnoheading{In II Nocturno}

\antlabel{Ant.} Magnus Dóminus, * \textdagger\ et laudábilis nimis in civitáte Dei nostri. \textdaggerdbl

\psalmheading{Psalmus 47}

\lettrine[lines=2]{M}{AGNUS} Dóminus, \textdagger\ et laudábilis nimis * in civitáte Dei nostri, \textdaggerdbl\ in monte sancto ejus.

Fundátur exsultatióne univérsæ terræ mons Sion, * látera Aquilónis, cívitas Regis magni.

Deus in dómibus ejus cognoscétur, * cum suscípiet eam.

Quóniam ecce reges terræ congregáti sunt: * convenérunt in unum.

Ipsi vidéntes sic admiráti sunt, conturbáti sunt, commóti sunt: * tremor apprehéndit eos.

Ibi dolóres ut parturiéntis: * in spíritu vehemménti cónteres naves Tharsis.

Sicut audívimus, sic vídimus in civitáte Dómini virtútum, in civitáte Dei nostri: * Deus fundávit eam in ætérnum.

Suscépimus, Deus, misericórdiam tuam, * in médio templi tui.

Secúndum nomen tuum, Deus, sic et laus tua in fines terræ: * justítia plena est déxtera tua.

Lætétur mons Sion, et exsúltent fíliæ Judæ: * propter judícia tua, Dómine.

Circúmdate Sion, et complectímini eam: * narráte in túrribus ejus.

Pónite corda vestra in virtúte ejus: * et distribúite domos ejus, ut enarrétis in progénie áltera.

Quóniam hic est Deus, Deus noster in ætérnum, et in sæculum sæculi: * ipse reget nos in sæcula.

\antlabel{Ant.} Magnus Dóminus, et laudábilis nimis in civitáte Dei nostri.

\antlabel{Ant.} Os meum loquétur * sapiéntiam: et meditátio cordis mei prudéntiam.

\psalmheading{Psalmus 48, i}

\lettrine[lines=2]{A}{UDITE} hæc, omnes gentes: * áuribus percípite omnes, qui habitátis orbem:

Quique terrígenæ, et fílii hóminum: * simul in unum dives et pauper.

Os meum loquétur sapiéntiam: * et meditátio cordis mei prudéntiam.

Inclinábo in parábolam aurem meam: * apériam in psaltério propositiónem meam.

Cur timébo in die mala? * iníquitas calcánei mei circúmdabit me:

Qui confídunt in virtúte sua: * et in multitúdine divitiárum suárum gloriántur.

Frater non rédimit, rédimet homo: * non dabit Deo placatiónem suam.

Et prétium redemptiónis ánimæ suæ: * et laborábit in ætérnum, et vivet adhuc in finem.

Non vidébit intéritum, cum víderit sapiéntes moriéntes: * simul insípiens, et stultus períbunt.

Et relínquent aliénis divítias suas: * et sepúlcra eórum domus illórum in ætérnum.

Tabernácula eórum in progénie et progénie: * vocavérunt nómina sua in terris suis.

Et homo, cum in honóre esset, non intelléxit: * comparátus est juméntis insipiéntibus, et símilis factus est illis.

\antlabel{Ant.} Os meum loquétur sapiéntiam: et meditátio cordis mei prudéntiam.

\antlabel{Ant.} Ne timúeris: * cum dívite non descéndet in sepúlcrum glória ejus.

\psalmheading{Psalmus 48, ii}

\lettrine[lines=2]{H}{ÆC} via illórum scándalum ipsis: * et póstea in ore suo complacébunt.

Sicut oves in inférno pósiti sunt: * mors depáscet eos.

Et dominabúntur eórum justi in matutíno: * et auxílium eórum veteráscet in inférno a glória eórum.

Verúmtamen Deus rédimet ánimam meam de manu ínferi, * cum accéperit me.

Ne timúeris, cum dives factus fúerit homo: * et cum multiplicáta fúerit glória domus ejus.

Quóniam cum interíerit, non sumet ómnia: * neque descéndet cum eo glória ejus.

Quia ánima ejus in vita ipsíus benedicétur: * confitébitur tibi cum beneféceris ei.

Introíbit usque in progénies patrum suórum: * et usque in ætérnum non vidébit lumen.

Homo, cum in honóre esset, non intelléxit: * comparátus est juméntis insipiéntibus, et símilis factus est illis.

\antlabel{Ant.} Ne timúeris: cum dívite non descéndet in sepúlcrum glória ejus.

\variantheading{In Festis novem Lectionum:}

\V\ Deus rédimet ánimam meam de manu ínferi. \R\ Cum accéperit me.

\absolutiolabel{Absolutio.} Ipsíus píetas et.

\benedlabel{Bened.} 4 Deus Pater omnípotens. 5 Christus perpétuæ. 6 Ignem sui amóris.

\nocturnoheading{In III Nocturno}

\variantheading{I}

\rubric{In Fériis extra Vigílias commúnes occurréntibus, et in Festis, cum ad Laudes non dícitur Psalmus 50 Miserére.}

\antlabel{Ant.} Deus deórum, * \textdagger\ Dóminus locútus est. \textdaggerdbl

\psalmheading{Psalmus 49, i}

\lettrine[lines=2]{D}{EUS} deórum, \textdagger\ Dóminus locútus est: * \textdaggerdbl\ et vocávit terram.

A solis ortu usque ad occásum: * ex Sion spécies decóris ejus.

Deus maniféste véniet: * Deus noster et non silébit.

Ignis in conspéctu ejus exardéscet: * et in circúitu ejus tempéstas válida.

Advocábit cælum desúrsum: * et terram discérnere pópulum suum.

Congregáte illi sanctos ejus: * qui órdinant testaméntum ejus super sacrifícia.

Et annuntiábunt cæli justítiam ejus: * quóniam Deus judex est.

Audi, pópulus meus, et loquar: Israël, et testificábor tibi: * Deus, Deus tuus ego sum.

Non in sacrifíciis tuis árguam te: * holocáusta autem tua in conspéctu meo sunt semper.

Non accípiam de domo tua vítulos: * neque de grégibus tuis hircos.

Quóniam meæ sunt omnes feræ silvárum, * juménta in móntibus et boves.

Cognóvi ómnia volatília cæli: * et pulchritúdo agri mecum est.

Si esuríero, non dicam tibi: * meus est enim orbis terræ, et plenitúdo ejus.

Numquid manducábo carnes taurórum? * aut sánguinem hircórum potábo?

Immola Deo sacrifícium laudis: * et redde Altíssimo vota tua.

Et ínvoca me in die tribulatiónis: * éruam te, et honorificábis me.

\antlabel{Ant.} Deus deórum, Dóminus locútus est.

\antlabel{Ant.} Intellígite, * qui obliviscímini Deum.

\psalmheading{Psalmus 49, ii}

\lettrine[lines=2]{P}{ECCATORI} autem dixit Deus: * Quare tu enárras justítias meas, et assúmis testaméntum meum per os tuum?

Tu vero odísti disciplínam: * et projecísti sermónes meos retrórsum:

Si vidébas furem, currébas cum eo: * et cum adúlteris portiónem tuam ponébas.

Os tuum abundávit malítia: * et lingua tua concinnábat dolos.

Sedens advérsus fratrem tuum loquebáris, et advérsus fílium matris tuæ ponébas scándalum: * hæc fecísti, et tácui.

Existimásti, iníque, quod ero tui símilis: * árguam te, et státuam contra fáciem tuam.

Intellígite hæc, qui oblivíscimini Deum: * nequándo rápiat, et non sit qui erípiat.

Sacrifícium laudis honorificábit me: * et illic iter, quo osténdam illi salutáre Dei.

\antlabel{Ant.} Intellígite, qui obliviscímini Deum.

\antlabel{Ant.} Acceptábis sacrifícium * justítiæ super altáre tuum, Dómine.

\psalmheading{Psalmus 50}

\lettrine[lines=2]{M}{ISERERE} mei, Deus, * secúndum magnam misericórdiam tuam.

Et secúndum multitúdinem miseratiónum tuárum, * dele iniquitátem meam.

Amplius lava me ab iniquitáte mea: * et a peccáto meo munda me.

Quóniam iniquitátem meam ego cognósco: * et peccátum meum contra me est semper.

Tibi soli peccávi, et malum coram te feci: * ut justificéris in sermónibus tuis, et vincas cum judicáris.

Ecce enim, in iniquitátibus concéptus sum: * et in peccátis concépit me mater mea.

Ecce enim, veritátem dilexísti: * incérta et occúlta sapiéntiæ tuæ manifestásti mihi.

Aspérges me hyssópo, et mundábor: * lavábis me, et super nivem dealbábor.

Audítui meo dabis gáudium et lætítiam: * et exsultábunt ossa humiliáta.

Avérte fáciem tuam a peccátis meis: * et omnes iniquitátes meas dele.

Cor mundum crea in me, Deus: * et spíritum rectum ínnova in viscéribus meis.

Ne projícias me a fácie tua: * et spíritum sanctum tuum ne áuferas a me.

Redde mihi lætítiam salutáris tui: * et spíritu principáli confírma me.

Docébo iníquos vias tuas: * et ímpii ad te converténtur.

Líbera me de sanguínibus, Deus, Deus salútis meæ: * et exsultábit lingua mea justítiam tuam.

Dómine, lábia mea apéries: * et os meum annuntiábit laudem tuam.

Quóniam si voluísses sacrifícium, dedíssem útique: * holocáustis non delectáberis.

Sacrifícium Deo spíritus contribulátus: * cor contrítum, et humiliátum, Deus, non despícies.

Benígne fac, Dómine, in bona voluntáte tua Sion: * ut ædificéntur muri Jerúsalem.

Tunc acceptábis sacrifícium justítiæ, oblatiónes, et holocáusta: * tunc impónent super altáre tuum vítulos.

\antlabel{Ant.} Acceptábis sacrifícium justítiæ super altáre tuum, Dómine.

\variantheading{In feriali Officio et in Festis:}

\V\ Audítui meo dabis gáudium et lætítiam. \R\ Et exsultábunt ossa humiliáta.

\variantheading{In Officio novem Lectionum}

\absolutiolabel{Absolutio.} A vínculis.

\benedlabel{Bened.} 7 Evangélica léctio. 8 In Festis Dómini: Divínum auxílium. In Festis B. M. V.: Cujus festum cólimus, ipsa Virgo. In Festis Sanctórum: Cujus (\textit{vel} Quorum \textit{aut} Quarum) festum cólimus. 9 Ad societátem. (Per evangélica dicta.)

\variantheading{In Officio trium Lectionum}

\absolutiolabel{Absolutio.} A vínculis.

\variantheading{In Feriis cum Homilia:}

\benedlabel{Bened.} 1 Evangélica. 2 Divínum. 3 Ad societátem.

\variantheading{In Feriis sine Homilia:}

\benedlabel{Bened.} 1 Ille nos benedícat. 2 Divínum. 3 Ad societátem.

\variantheading{In Festis Sanctorum:}

\benedlabel{Bened.} 1 Ille nos benedícat, qui sine fine vivit et regnat. \R\ Amen. 2 Cujus (\textit{vel} Quorum \textit{aut} Quarum). 3 Ad societátem.

\nocturnoheading{In III Nocturno}

\variantheading{II}

\rubric{In Vigíliis commúnibus, cum ad Laudes dícitur Psalmus 50 Miserére.}

\antlabel{Ant.} Deus deórum. \textdagger

\psalmheading{Psalmus 49, i}

\lettrine[lines=2]{D}{EUS} deórum, \textdagger\ Dóminus locútus est: * \textdaggerdbl\ et vocávit terram.

A solis ortu usque ad occásum: * ex Sion spécies decóris ejus.

Deus maniféste véniet: * Deus noster et non silébit.

Ignis in conspéctu ejus exardéscet: * et in circúitu ejus tempéstas válida.

Advocábit cælum desúrsum: * et terram discérnere pópulum suum.

Congregáte illi sanctos ejus: * qui órdinant testaméntum ejus super sacrifícia.

Et annuntiábunt cæli justítiam ejus: * quóniam Deus judex est.

Audi, pópulus meus, et loquar: Israël, et testificábor tibi: * Deus, Deus tuus ego sum.

Non in sacrifíciis tuis árguam te: * holocáusta autem tua in conspéctu meo sunt semper.

Non accípiam de domo tua vítulos: * neque de grégibus tuis hircos.

Quóniam meæ sunt omnes feræ silvárum, * juménta in móntibus et boves.

Cognóvi ómnia volatília cæli: * et pulchritúdo agri mecum est.

Si esuríero, non dicam tibi: * meus est enim orbis terræ, et plenitúdo ejus.

Numquid manducábo carnes taurórum? * aut sánguinem hircórum potábo?

Immola Deo sacrifícium laudis: * et redde Altíssimo vota tua.

Et ínvoca me in die tribulatiónis: * éruam te, et honorificábis me.

\antlabel{Ant.} Deus deórum, Dóminus locútus est.

\antlabel{Ant.} Immola Deo.

\psalmheading{Psalmus 49, ii}

\lettrine[lines=2]{P}{ECCATORI} autem dixit Deus: * Quare tu enárras justítias meas, et assúmis testaméntum meum per os tuum?

Tu vero odísti disciplínam: * et projecísti sermónes meos retrórsum:

Si vidébas furem, currébas cum eo: * et cum adúlteris portiónem tuam ponébas.

Os tuum abundávit malítia: * et lingua tua concinnábat dolos.

Sedens advérsus fratrem tuum loquebáris, et advérsus fílium matris tuæ ponébas scándalum: * hæc fecísti, et tácui.

Existimásti, iníque, quod ero tui símilis: * árguam te, et státuam contra fáciem tuam.

Intellígite hæc, qui oblivíscimini Deum: * nequándo rápiat, et non sit qui erípiat.

Sacrifícium laudis honorificábit me: * et illic iter, quo osténdam illi salutáre Dei.

\antlabel{Ant.} Immola Deo sacrifícium laudis.

\antlabel{Ant.} Intellígite.

\psalmheading{Psalmus 49, iii}

\lettrine[lines=2]{P}{ECCATORI} autem dixit Deus: * Quare tu enárras justítias meas, et assúmis testaméntum meum per os tuum?

Tu vero odísti disciplínam: * et projecísti sermónes meos retrórsum:

Si vidébas furem, currébas cum eo: * et cum adúlteris portiónem tuam ponébas.

Os tuum abundávit malítia: * et lingua tua concinnábat dolos.

Sedens advérsus fratrem tuum loquebáris, et advérsus fílium matris tuæ ponébas scándalum: * hæc fecísti, et tácui.

Existimásti, iníque, quod ero tui símilis: * árguam te, et státuam contra fáciem tuam.

Intellígite hæc, qui oblivíscimini Deum: * nequándo rápiat, et non sit qui erípiat.

Sacrifícium laudis honorificábit me: * et illic iter, quo osténdam illi salutáre Dei.

\antlabel{Ant.} Intellígite, qui obliviscímini Deum.

\V\ Sacrifícium laudis honorificábit me. \R\ Et illic iter, quo osténdam illi salutáre Dei.

\absolutiolabel{Absolutio.} A vínculis.

\variantheading{In Feriis cum Homilia:}

\benedlabel{Bened.} 1 Evangélica. 2 Divínum. 3 Ad societátem.

\variantheading{In Feriis sine Homilia:}

\benedlabel{Bened.} 1 Ille nos benedícat. 2 Divínum. 3 Ad societátem.

% ===== END OF BATCH 19 -- Covers autumnalis.md lines 6277-6743,
% opening "Feria Quarta" (Wednesday) and completing its "Ad
% Matutinum" in full: the ferial Invitatorium "Deum magnum
% Dóminum," the hymn "Rerum Creátor óptime" (stray OCR "\*"
% silently dropped), and the three Nocturns (Psalms 44/i, 44/ii,
% 45; 47, 48/i, 48/ii). "In III Nocturno" here introduces, for the
% first time in this volume, two complete alternate forms labeled
% "I" and "II," each with its own italicized rubric describing when
% it applies (Form I when Ps. 50 was NOT said at Lauds, using
% Psalms 49/i, 49/ii, and the full Ps. 50; Form II when Ps. 50 WAS
% said at Lauds, using Psalms 49/i, 49/ii, and 49/iii instead) --
% rendered with \variantheading{I} and \variantheading{II}, matching
% the identical command already used for this exact structure in
% both verna.tex and hiemalis.tex. Notably, Form II's "Psalmus 49,
% ii" and "Psalmus 49, iii" give verbatim identical text in
% autumnalis.md itself (both "Peccátori autem dixit Deus..."); this is
% transcribed exactly as the source repeats it, without alteration,
% since it is a structural feature of the source and not a
% character-level OCR defect. The document continues with "Ad
% Laudes I" (autumnalis.md line 6745).
% =====
\hourheading{Ad Laudes I}

\rubric{In Fériis extra Vigílias commúnes occurréntibus, atque in Festis.}

(Pater noster et Ave.)

\V\ Deus, in adjutórium.

\antlabel{Ant.} Dóminus regnávit, * \textdagger\ exsúltet terra. \textdaggerdbl

\psalmheading{Psalmus 96}

\lettrine[lines=2]{D}{OMINUS} regnávit, \textdagger\ exsúltet terra: * \textdaggerdbl\ læténtur ínsulæ multæ.

Nubes, et calígo in circúitu ejus: * justítia, et judícium corréctio sedis ejus.

Ignis ante ipsum præcédet, * et inflammábit in circúitu inimícos ejus.

Illuxérunt fúlgura ejus orbi terræ: * vidit, et commóta est terra.

Montes, sicut cera fluxérunt a fácie Dómini: * a fácie Dómini omnis terra.

Annuntiavérunt cæli justítiam ejus: * et vidérunt omnes pópuli glóriam ejus.

Confundántur omnes, qui adórant sculptília: * et qui gloriántur in simulácris suis.

Adoráte eum, omnes Angeli ejus: * audívit, et lætáta est Sion.

Et exsultavérunt fíliæ Judæ, * propter judícia tua, Dómine:

Quóniam tu Dóminus Altíssimus super omnem terram: * nimis exaltátus es super omnes deos.

Qui dilígitis Dóminum, odíte malum: * custódit Dóminus ánimas sanctórum suórum, de manu peccatóris liberábit eos.

Lux orta est justo, * et rectis corde lætítia.

Lætámini, justi, in Dómino: * et confitémini memóriæ sanctificatiónis ejus.

\antlabel{Ant.} Dóminus regnávit, exsúltet terra.

\antlabel{Ant.} Te decet hymnus, * \textdagger\ Deus, in Sion. \textdaggerdbl

\psalmheading{Psalmus 64}

\lettrine[lines=2]{T}{E} decet hymnus, \textdagger\ Deus, in Sion: * \textdaggerdbl\ et tibi reddétur votum in Jerúsalem.

Exáudi oratiónem meam: * ad te omnis caro véniet.

Verba iniquórum prævaluérunt super nos: * et impietátibus nostris tu propitiáberis.

Beátus, quem elegísti, et assumpsísti: * inhabitábit in átriis tuis.

Replébimur in bonis domus tuæ: * sanctum est templum tuum, mirábile in æquitáte.

Exáudi nos, Deus, salutáris noster, * spes ómnium fínium terræ, et in mari longe.

Præparans montes in virtúte tua, accínctus poténtia: * qui contúrbas profúndum maris sonum flúctuum ejus.

Turbabúntur gentes, et timébunt qui hábitant términos a signis tuis: * éxitus matutíni, et véspere delectábis.

Visitásti terram, et inebriásti eam: * multiplicásti locupletáre eam.

Flumen Dei replétum est aquis, parásti cibum illórum: * quóniam ita est præparátio ejus.

Rivos ejus inébria, multíplica genímina ejus: * in stillicídiis ejus lætábitur gérminans.

Benedíces corónæ anni benignitátis tuæ: * et campi tui replebúntur ubertáte.

Pinguéscent speciósa desérti: * et exsultatióne colles accingéntur.

Indúti sunt aríetes óvium, et valles abundábunt fruménto: * clamábunt, étenim hymnum dicent.

\antlabel{Ant.} Te decet hymnus, Deus, in Sion.

\antlabel{Ant.} Tibi, Dómine, psallam, * et intélligam in via immaculáta.

\psalmheading{Psalmus 100}

\lettrine[lines=2]{M}{ISERICORDIAM}, et judícium * cantábo tibi, Dómine:

Psallam, et intélligam in via immaculáta, * quando vénies ad me.

Perambulábam in innocéntia cordis mei, * in médio domus meæ.

Non proponébam ante óculos meos rem injústam: * faciéntes prævaricatiónes odívi.

Non adhæsit mihi cor pravum: * declinántem a me malígnum non cognoscébam.

Detrahéntem secréto próximo suo, * hunc persequébar.

Supérbo óculo, et insatiábili corde, * cum hoc non edébam.

Oculi mei ad fidéles terræ ut sédeant mecum: * ámbulans in via immaculáta, hic mihi ministrábat.

Non habitábit in médio domus meæ qui facit supérbiam: * qui lóquitur iníqua, non diréxit in conspéctu oculórum meórum.

In matutíno interficiébam omnes peccatóres terræ: * ut dispérderem de civitáte Dómini omnes operántes iniquitátem.

\antlabel{Ant.} Tibi, Dómine, psallam, et intélligam in via immaculáta.

\antlabel{Ant.} Dómine, magnus es tu, * et præclárus in virtúte tua.

\canticumheading{Canticum Judith}{Judith. 16, 15-21}

\lettrine[lines=2]{H}{YMNUM} cantémus Dómino, * hymnum novum cantémus Deo nostro.

Adonái, Dómine, magnus es tu, et præclárus in virtúte tua, * et quem superáre nemo potest.

Tibi sérviat omnis creatúra tua: * quia dixísti, et facta sunt:

Misísti spíritum tuum, et creáta sunt: * et non est qui resístat voci tuæ.

Montes a fundaméntis movebúntur cum aquis: * petræ, sicut cera, liquéscent ante fáciem tuam.

Qui autem timent te, * magni erunt apud te per ómnia.

Væ genti insurgénti super genus meum: Dóminus enim omnípotens vindicábit in eis, * in die judícii visitábit illos.

Dabit enim ignem, et vermes in carnes eórum, * ut urántur, et séntiant usque in sempitérnum.

\antlabel{Ant.} Dómine, magnus es tu, et præclárus in virtúte tua.

\antlabel{Ant.} Laudábo Deum meum * in vita mea.

\psalmheading{Psalmus 145}

\lettrine[lines=2]{L}{AUDA}, ánima mea, Dóminum, laudábo Dóminum in vita mea: * psallam Deo meo quámdiu fúero.

Nolíte confídere in princípibus: * in fíliis hóminum, in quibus non est salus.

Exíbit spíritus ejus, et revertétur in terram suam: * in illa die períbunt omnes cogitatiónes eórum.

Beátus, cujus Deus Jacob adjútor ejus, spes ejus in Dómino, Deo ipsíus: * qui fecit cælum et terram, mare, et ómnia, quæ in eis sunt.

Qui custódit veritátem in sæculum, facit judícium injúriam patiéntibus: * dat escam esuriéntibus.

Dóminus solvit compedítos: * Dóminus illúminat cæcos.

Dóminus érigit elísos, * Dóminus díligit justos.

Dóminus custódit ádvenas, pupíllum et víduam suscípiet: * et vias peccatórum dispérdet.

Regnábit Dóminus in sæcula, Deus tuus, Sion, * in generatiónem et generatiónem.

\antlabel{Ant.} Laudábo Deum meum in vita mea.

\variantheading{In Officio feriali}

\capitulumlabel{Capit.}{Rom. 13, 12-13}

\lettrine[lines=2]{N}{OX} præcéssit, dies autem appropinquávit. Abjiciámus ergo ópera tenebrárum, et induámur arma lucis. Sicut in die honéste ambulémus.

\hymnusheading{Hymnus}

\lettrine[lines=2]{N}{OX}, et ténebræ, et núbila,\\
Confúsa mundi et túrbida:\\
Lux intrat, albéscit polus:\\
Christus venit: discédite.

Calígo terræ scínditur\\
Percússa solis spículo,\\
Rebúsque jam color redit,\\
Vultu niténtis síderis.

Te, Christe, solum nóvimus:\\
Te mente pura et símplici,\\
Flendo et canéndo quǽsumus,\\
Inténde nostris sénsibus.

Sunt multa fucis illíta,\\
Quæ luce purgéntur tua:\\
Tu, vera lux Cæléstium,\\
Vultu seréno illúmina.

Deo Patri sit glória,\\
Ejúsque soli Fílio,\\
Cum Spíritu Paráclito,\\
Nunc et per omne sæculum.\\
Amen.

\V\ Repléti sumus mane misericórdia tua. \R\ Exsultávimus, et delectáti sumus.

\antlabel{Ad Bened. Ant.} De manu ómnium, * qui odérunt nos, liberávit nos Dóminus.

\rubric{Réliqua ut in Ordinário. 11}

\hourheading{Ad Laudes II}

\rubric{In Vigíliis commúnibus, quando fit Offícium de Féria.}

(Pater noster et Ave.)

\V\ Deus, in adjutórium.

\antlabel{Ant.} Amplius lava me.

\psalmheading{Psalmus 50}

\lettrine[lines=2]{M}{ISERERE} mei, Deus, * secúndum magnam misericórdiam tuam.

Et secúndum multitúdinem miseratiónum tuárum, * dele iniquitátem meam.

Amplius lava me ab iniquitáte mea: * et a peccáto meo munda me.

Quóniam iniquitátem meam ego cognósco: * et peccátum meum contra me est semper.

Tibi soli peccávi, et malum coram te feci: * ut justificéris in sermónibus tuis, et vincas cum judicáris.

Ecce enim, in iniquitátibus concéptus sum: * et in peccátis concépit me mater mea.

Ecce enim, veritátem dilexísti: * incérta et occúlta sapiéntiæ tuæ manifestásti mihi.

Aspérges me hyssópo, et mundábor: * lavábis me, et super nivem dealbábor.

Audítui meo dabis gáudium et lætítiam: * et exsultábunt ossa humiliáta.

Avérte fáciem tuam a peccátis meis: * et omnes iniquitátes meas dele.

Cor mundum crea in me, Deus: * et spíritum rectum ínnova in viscéribus meis.

Ne projícias me a fácie tua: * et spíritum sanctum tuum ne áuferas a me.

Redde mihi lætítiam salutáris tui: * et spíritu principáli confírma me.

Docébo iníquos vias tuas: * et ímpii ad te converténtur.

Líbera me de sanguínibus, Deus, Deus salútis meæ: * et exsultábit lingua mea justítiam tuam.

Dómine, lábia mea apéries: * et os meum annuntiábit laudem tuam.

Quóniam si voluísses sacrifícium, dedíssem útique: * holocáustis non delectáberis.

Sacrifícium Deo spíritus contribulátus: * cor contrítum, et humiliátum, Deus, non despícies.

Benígne fac, Dómine, in bona voluntáte tua Sion: * ut ædificéntur muri Jerúsalem.

Tunc acceptábis sacrifícium justítiæ, oblatiónes, et holocáusta: * tunc impónent super altáre tuum vítulos.

\antlabel{Ant.} Amplius lava me, Dómine, ab injustítia mea.

\antlabel{Ant.} Impietátibus nostris.

\psalmheading{Psalmus 64}

\lettrine[lines=2]{T}{E} decet hymnus, Deus, in Sion: * et tibi reddétur votum in Jerúsalem.

Exáudi oratiónem meam: * ad te omnis caro véniet.

Verba iniquórum prævaluérunt super nos: * et impietátibus nostris tu propitiáberis.

Beátus, quem elegísti, et assumpsísti: * inhabitábit in átriis tuis.

Replébimur in bonis domus tuæ: * sanctum est templum tuum, mirábile in æquitáte.

Exáudi nos, Deus, salutáris noster, * spes ómnium fínium terræ, et in mari longe.

Præparans montes in virtúte tua, accínctus poténtia: * qui contúrbas profúndum maris sonum flúctuum ejus.

Turbabúntur gentes, et timébunt qui hábitant términos a signis tuis: * éxitus matutíni, et véspere delectábis.

Visitásti terram, et inebriásti eam: * multiplicásti locupletáre eam.

Flumen Dei replétum est aquis, parásti cibum illórum: * quóniam ita est præparátio ejus.

Rivos ejus inébria, multíplica genímina ejus: * in stillicídiis ejus lætábitur gérminans.

Benedíces corónæ anni benignitátis tuæ: * et campi tui replebúntur ubertáte.

Pinguéscent speciósa desérti: * et exsultatióne colles accingéntur.

Indúti sunt aríetes óvium, et valles abundábunt fruménto: * clamábunt, étenim hymnum dicent.

\antlabel{Ant.} Impietátibus nostris tu propitiáberis, Deus.

\antlabel{Ant.} In innocéntia.

\psalmheading{Psalmus 100}

\lettrine[lines=2]{M}{ISERICORDIAM}, et judícium * cantábo tibi, Dómine:

Psallam, et intélligam in via immaculáta, * quando vénies ad me.

Perambulábam in innocéntia cordis mei, * in médio domus meæ.

Non proponébam ante óculos meos rem injústam: * faciéntes prævaricatiónes odívi.

Non adhæsit mihi cor pravum: * declinántem a me malígnum non cognoscébam.

Detrahéntem secréto próximo suo, * hunc persequébar.

Supérbo óculo, et insatiábili corde, * cum hoc non edébam.

Oculi mei ad fidéles terræ ut sédeant mecum: * ámbulans in via immaculáta, hic mihi ministrábat.

Non habitábit in médio domus meæ qui facit supérbiam: * qui lóquitur iníqua, non diréxit in conspéctu oculórum meórum.

In matutíno interficiébam omnes peccatóres terræ: * ut dispérderem de civitáte Dómini omnes operántes iniquitátem.

\antlabel{Ant.} In innocéntia cordis mei perambulábo, Dómine.

\antlabel{Ant.} Exsultávit cor meum. \textdagger

\canticumheading{Canticum Annæ}{1 Reg. 2, 1-10}

\lettrine[lines=2]{E}{XSULTAVIT} cor meum \textdagger\ in Dómino, * et exaltátum est cornu meum in Deo meo.

Dilatátum est os meum super inimícos meos: * quia lætáta sum in salutári tuo.

Non est sanctus, ut est Dóminus: neque enim est álius extra te, * et non est fortis sicut Deus noster.

Nolíte multiplicáre loqui sublímia, * gloriántes:

Recédant vétera de ore vestro: quia Deus sciéntiárum, Dóminus est, * et ipsi præparántur cogitatiónes.

Arcus fórtium superátus est, * et infírmi accíncti sunt róbore.

Repléti prius, pro pánibus se locavérunt: * et famélici saturáti sunt.

Donec stérilis péperit plúrimos: * et quæ multos habébat fílios, infirmáta est.

Dóminus mortíficat et vivíficat, * dedúcit ad ínferos et redúcit.

Dóminus páuperem facit et ditat, * humíliat et súblevat.

Súscitat de púlvere egénum, * et de stércore élevat páuperem:

Ut sédeat cum princípibus, * et sólium glóriæ téneat.

Dómini enim sunt cárdines terræ, * et pósuit super eos orbem.

Pedes sanctórum suórum servábit, et ímpii in ténebris conticéscent: * quia non in fortitúdine sua roborábitur vir.

Dóminum formidábunt adversárii ejus: * et super ipsos in cælis tonábit:

Dóminus judicábit fines terræ, et dabit impérium regi suo, * et sublimábit cornu Christi sui.

\antlabel{Ant.} Exsultávit cor meum in Dómino, qui humíliat et súblevat.

\antlabel{Ant.} Lauda. \textdagger

\psalmheading{Psalmus 145}

\lettrine[lines=2]{L}{AUDA}, \textdagger\ ánima mea, Dóminum, laudábo Dóminum in vita mea: * psallam Deo meo quámdiu fúero.

Nolíte confídere in princípibus: * in fíliis hóminum, in quibus non est salus.

Exíbit spíritus ejus, et revertétur in terram suam: * in illa die períbunt omnes cogitatiónes eórum.

Beátus, cujus Deus Jacob adjútor ejus, spes ejus in Dómino, Deo ipsíus: * qui fecit cælum et terram, mare, et ómnia, quæ in eis sunt.

Qui custódit veritátem in sæculum, facit judícium injúriam patiéntibus: * dat escam esuriéntibus.

Dóminus solvit compedítos: * Dóminus illúminat cæcos.

Dóminus érigit elísos, * Dóminus díligit justos.

Dóminus custódit ádvenas, pupíllum et víduam suscípiet: * et vias peccatórum dispérdet.

Regnábit Dóminus in sæcula, Deus tuus, Sion, * in generatiónem et generatiónem.

\antlabel{Ant.} Lauda, ánima mea, Dóminum, qui érigit elísos, et díligit justos.

\capitulumlabel{Capit.}{Rom. 13, 12-13}

\lettrine[lines=2]{N}{OX} præcéssit, dies autem appropinquávit. Abjiciámus ergo ópera tenebrárum, et induámur arma lucis. Sicut in die honéste ambulémus.

\hymnusheading{Hymnus}

\lettrine[lines=2]{N}{OX}, et ténebræ, et núbila,\\
Confúsa mundi et túrbida:\\
Lux intrat, albéscit polus:\\
Christus venit: discédite.

Calígo terræ scínditur\\
Percússa solis spículo,\\
Rebúsque jam color redit,\\
Vultu niténtis síderis.

Te, Christe, solum nóvimus:\\
Te mente pura et símplici,\\
Flendo et canéndo quǽsumus,\\
Inténde nostris sénsibus.

Sunt multa fucis illíta,\\
Quæ luce purgéntur tua:\\
Tu, vera lux Cæléstium,\\
Vultu seréno illúmina.

Deo Patri sit glória,\\
Ejúsque soli Fílio,\\
Cum Spíritu Paráclito,\\
Nunc et per omne sæculum.\\
Amen.

\V\ Repléti sumus mane misericórdia tua. \R\ Exsultávimus, et delectáti sumus.

\antlabel{Ad Bened. Ant.} De manu ómnium, * qui odérunt nos, liberávit nos Dóminus.

\rubric{Expléto Cántico ac repetíta Antíphona, dicúntur sequéntes}

\variantheading{Preces feriales}

KYRIE, eléison. Christe, eléison. Kýrie, eléison.

Pater noster \rubric{quod dícitur a solo Hebdomadário totum clara voce, usque ad} \V\ Et ne nos indúcas in tentatiónem. \R\ Sed líbera nos a malo. \V\ Ego dixi: Dómine, miserére mei. \R\ Sana ánimam meam, quia peccávi tibi. \V\ Convértere, Dómine, úsquequo? \R\ Et deprecábilis esto super servos tuos. \V\ Fiat misericórdia tua, Dómine, super nos. \R\ Quemádmodum sperávimus in te. \V\ Sacerdótes tui induántur justítiam. \R\ Et sancti tui exsúltent. \V\ Orémus pro beatíssimo Papa nostro N. \R\ Dóminus consérvet eum, et vivíficet eum, et beátum fáciat eum in terra, et non tradat eum in ánimam inimicórum ejus. \V\ Orémus et pro Antístite nostro N. \R\ Stet et pascat in fortitúdine tua, Dómine, in sublimitáte nóminis tui. \V\ Dómine, salvum fac regem. \R\ Et exáudi nos in die, qua invocavérimus te. \V\ Salvum fac pópulum tuum, Dómine, et bénedic hereditáti tuæ. \R\ Et rege eos, et extólle illos usque in ætérnum. \V\ Meménto Congregatiónis tuæ. \R\ Quam possedísti ab inítio. \V\ Fiat pax in virtúte tua. \R\ Et abundántia in túrribus tuis. \V\ Orémus pro benefactóribus nostris. \R\ Retribúere dignáre, Dómine, ómnibus, nobis bona faciéntibus propter nomen tuum, vitam ætérnam. Amen. \V\ Orémus pro fidélibus defúnctis. \R\ Réquiem ætérnam dona eis, Dómine, et lux perpétua lúceat eis. \V\ Requiéscant in pace. \R\ Amen. \V\ Pro frátribus nostris abséntibus. \R\ Salvos fac servos tuos, Deus meus, sperántes in te. \V\ Pro afflíctis et captívis. \R\ Líbera eos, Deus Israël, ex ómnibus tribulatiónibus suis. \V\ Mitte eis, Dómine, auxílium de sancto. \R\ Et de Sion tuére eos. \V\ Dómine, Deus virtútum, convérte nos. \R\ Et osténde fáciem tuam, et salvi érimus. \V\ Exsúrge, Christe, ádjuva nos. \R\ Et líbera nos propter nomen tuum. \V\ Dómine, exáudi oratiónem meam. \R\ Et clamor meus ad te véniat.

\rubric{Réliqua ut in Ordinário. 11}

% ===== END OF BATCH 20 -- Covers autumnalis.md lines 6745-7156, the
% two alternate ferial Lauds schemes for Feria Quarta, "Ad Laudes
% I" (Psalms 96, 64, 100, Canticum Judith, 145) and "Ad Laudes II"
% (Psalms 50, 64, 100, Canticum Annæ, 145), each with the Capitulum
% "Nox præcéssit" and the hymn "Nox, et ténebræ, et núbila." Notably,
% Psalm 64's incipit carries the dagger and double-dagger
% chant-pointing marks in the Ad Laudes I copy (autumnalis.md line
% 6834: "Te decet hymnus, † Deus, in Sion: * ‡...") but NOT in the
% Ad Laudes II copy of the identical psalm (line 6991, plain "Te
% decet hymnus, Deus, in Sion: *..."); each is transcribed exactly
% as its own occurrence gives it, confirming this is a genuine
% source divergence rather than an OCR gap. Two antiphon incipits in
% Ad Laudes II ("Impietátibus nostris." and "In innocéntia.") carry
% no asterisk at all in the source and are transcribed as bare
% unmarked incipits. Ad Laudes II additionally carries the "Preces
% feriales" litany in full. The document continues with "Ad Primam"
% (autumnalis.md line 7159).
% =====
\hourheading{Ad Primam}

Pater, Ave et Credo.

\V\ Deus, in adjutórium.

\hymnusheading{Hymnus}

\lettrine[lines=2]{J}{AM} lucis orto sídere,\\
Deum precémur súpplices,\\
Ut in diúrnis áctibus\\
Nos servet a nocéntibus.

Linguam refrénans témperet,\\
Ne litis horror ínsonet:\\
Visum fovéndo cóntegat,\\
Ne vanitátes háuriat.

Sint pura cordis íntima,\\
Absístat et vecórdia:\\
Carnis terat supérbiam\\
Potus cibíque párcitas.

Ut, cum dies abscésserit\\
Noctémque sors redúxerit,\\
Mundi per abstinéntiam\\
Ipsi canámus glóriam.

Deo Patri sit glória,\\
Ejúsque soli Fílio,\\
Cum Spíritu Paráclito,\\
Nunc et per omne sæculum.\\
Amen.

\antlabel{Ant.} Misericórdia tua.

\psalmheading{Psalmus 25}

\lettrine[lines=2]{J}{UDICA} me, Dómine, quóniam ego in innocéntia mea ingréssus sum: * et in Dómino sperans non infirmábor.

Proba me, Dómine, et tenta me: * ure renes meos et cor meum.

Quóniam misericórdia tua ante óculos meos est: * et complácui in veritáte tua.

Non sedi cum concílio vanitátis: * et cum iníqua geréntibus non introíbo.

Odívi ecclésiam malignántium: * et cum ímpiis non sedébo.

Lavábo inter innocéntes manus meas: * et circúmdabo altáre tuum, Dómine:

Ut áudiam vocem laudis, * et enárrem univérsa mirabília tua.

Dómine, diléxi decórem domus tuæ, * et locum habitatiónis glóriæ tuæ.

Ne perdas cum ímpiis, Deus, ánimam meam, * et cum viris sánguinum vitam meam:

In quorum mánibus iniquitátes sunt: * déxtera eórum repléta est munéribus.

Ego autem in innocéntia mea ingréssus sum: * rédime me, et miserére mei.

Pes meus stetit in directo: * in ecclésiis benedícam te, Dómine.

\psalmheading{Psalmus 51}

\lettrine[lines=2]{Q}{UID} gloriáris in malítia, * qui potens es in iniquitáte?

Tota die injustítiam cogitávit lingua tua: * sicut novácula acúta fecísti dolum.

Dilexísti malítiam super benignitátem: * iniquitátem magis quam loqui æquitátem.

Dilexísti ómnia verba præcipitatiónis, * lingua dolósa.

Proptérea Deus déstruet te in finem, * evéllet te, et emigrábit te de tabernáculo tuo: et radícem tuam de terra vivéntium.

Vidébunt justi, et timébunt, et super eum ridébunt, et dicent: * Ecce homo, qui non pósuit Deum adjutórem suum:

Sed sperávit in multitúdine divitiárum suárum: * et prævalúit in vanitáte sua.

Ego autem, sicut olíva fructífera in domo Dei, * sperávi in misericórdia Dei in ætérnum: et in séculum séculi.

Confitébor tibi in séculum, quia fecísti: * et exspectábo nomen tuum, quóniam bonum est in conspéctu sanctórum tuórum.

\psalmheading{Psalmus 52}

\lettrine[lines=2]{D}{IXIT} insípiens in corde suo: * Non est Deus.

Corrúpti sunt, et abominábiles facti sunt in iniquitátibus: * non est qui fáciat bonum.

Deus de cælo prospéxit super fílios hóminum: * ut vídeas si est intélligens, aut requírens Deum.

Omnes declinavérunt, simul inútiles facti sunt: * non est qui fáciat bonum, non est usque ad unum.

Nonne scient omnes qui operántur iniquitátem, * qui dévorant plebem meam ut cibum panis?

Deum non invocavérunt: * illic trepidavérunt timóre, ubi non erat timor.

Quóniam Deus dissipávit ossa eórum qui homínibus placent: * confúsi sunt, quóniam Deus sprevit eos.

Quis dabit ex Sion salutáre Israël? * cum convérterit Deus captivitátem plebis suæ, exsultábit Jacob, et lætábitur Israël.

\rubric{¶ Quando ad Laudes dictus fuerit Ps. 50 Miserére, hic subjúngitur sequens Ps. 96, qui in Laudibus locum non habuit.}

\psalmheading{Psalmus 96}

\lettrine[lines=2]{D}{OMINUS} regnávit, exsúltet terra: * læténtur ínsulæ multæ.

Nubes, et calígo in circúitu ejus: * justítia, et judícium corréctio sedis ejus.

Ignis ante ipsum præcédet, * et inflammábit in circúitu inimícos ejus.

Illuxérunt fúlgura ejus orbi terræ: * vidit, et commóta est terra.

Montes, sicut cera fluxérunt a fácie Dómini: * a fácie Dómini omnis terra.

Annuntiavérunt cæli justítiam ejus: * et vidérunt omnes pópuli glóriam ejus.

Confundántur omnes, qui adórant sculptília: * et qui gloriántur in simulácris suis.

Adoráte eum, omnes Angeli ejus: * audívit, et lætáta est Sion.

Et exsultavérunt fíliæ Judæ, * propter judícia tua, Dómine:

Quóniam tu Dóminus Altíssimus super omnem terram: * nimis exaltátus es super omnes deos.

Qui dilígitis Dóminum, odíte malum: * custódit Dóminus ánimas sanctórum suórum, de manu peccatóris liberábit eos.

Lux orta est justo, * et rectis corde lætítia.

Lætámini, justi, in Dómino: * et confitémini memóriæ sanctificatiónis ejus.

\rubric{Secus vero, tribus Psalmis explétis, statim dícitur Ant.}

\antlabel{Ant.} Misericórdia tua, Dómine, ante óculos meos: et complácui in veritáte tua.

\rubric{In Officio Festi vel Octávæ:}

\capitulumlabel{Capitulum}{1 Tim. 1, 17}

\lettrine[lines=2]{R}{EGI} sæculórum immortáli et invisíbili, soli Deo honor et glória in sæcula sæculórum. Amen.

\rubric{In ómnibus autem Fériis et in Vigíliis commúnibus:}

\capitulumlabel{Capitulum}{Zach. 8, 19}

\lettrine[lines=2]{P}{ACEM} et veritátem dilígite, ait Dóminus omnípotens.

\R\ br. Christe, Fili Dei vivi, * Miserére nobis. Christe. \V\ \textit{Qui sedes ad déxteram Patris.} Miserére nobis. Glória Patri. Christe. \V\ Exsúrge, Christe, ádjuva nos. \R\ Et líbera nos propter nomen tuum.

\variantheading{Preces dominicales et feriales (si dicendæ sint)}

KYRIE, eléison. Christe, eléison. Kýrie, eléison.

Pater noster \rubric{secreto usque ad} \V\ Et ne nos indúcas in tentatiónem. \R\ Sed líbera nos a malo.

Credo in Deum \rubric{secreto usque ad} \V\ Carnis resurrectiónem. \R\ Vitam ætérnam. Amen. \V\ Et ego ad te, Dómine, clamávi. \R\ Et mane orátio mea prævéniet te. \V\ Repleátur os meum laude. \R\ Ut cantem glóriam tuam, tota die magnitúdinem tuam. \V\ Dómine, avérte fáciem tuam a peccátis meis. \R\ Et omnes iniquitátes meas dele. \V\ Cor mundum crea in me, Deus. \R\ Et spíritum rectum ínnova in viscéribus meis. \V\ Ne projícias me a fácie tua. \R\ Et spíritum sanctum tuum ne áuferas a me. \V\ Redde mihi lætítiam salutáris tui. \R\ Et spíritu principáli confírma me.

\variantheading{In Précibus feriálibus tantum}

\V\ Eripe me, Dómine, ab hómine malo. \R\ A viro iníquo éripe me. \V\ Eripe me de inimícis meis, Deus meus. \R\ Et ab insurgéntibus in me líbera me. \V\ Eripe me de operántibus iniquitátem. \R\ Et de viris sánguinum salva me. \V\ Sic psalmum dicam nómini tuo in sæculum sæculi. \R\ Ut reddam vota mea de die in diem. \V\ Exáudi nos, Deus, salutáris noster. \R\ Spes ómnium fínium terræ, et in mari longe. \V\ Deus, in adjutórium meum inténde. \R\ Dómine, ad adjuvándum me festína. \V\ Sanctus Deus, Sanctus fortis, Sanctus immortális. \R\ Miserére nobis. \V\ Bénedic, ánima mea, Dómino. \R\ Et ómnia, quæ intra me sunt, nómini sancto ejus. \V\ Bénedic, ánima mea, Dómino. \R\ Et noli oblivísci omnes retributiónes ejus. \V\ Qui propitiátur ómnibus iniquitátibus tuis. \R\ Qui sanat omnes infirmitátes tuas. \V\ Qui rédimit de intéritu vitam tuam. \R\ Qui corónat te in misericórdia et miseratiónibus. \V\ Qui replet in bonis desidérium tuum. \R\ Renovábitur ut áquilæ juvéntus tua.

\variantheading{In utrisque Précibus}

\V\ Adjutórium nostrum in nómine Dómini. \R\ Qui fecit cælum et terram.

\formulalabel{Confíteor. Misereátur. Indulgéntiam.}

\V\ Dignáre, Dómine, die isto. \R\ Sine peccáto nos custodíre. \V\ Miserére nostri, Dómine. \R\ Miserére nostri. \V\ Fiat misericórdia tua, Dómine, super nos. \R\ Quemádmodum sperávimus in te. \V\ Dómine, exáudi oratiónem meam. \R\ Et clamor meus ad te véniat.

\V\ Dóminus vobíscum.

Orémus. \oratiolabel{Oratio}

\lettrine[lines=2]{D}{OMINE}, Deus omnípotens, qui ad princípium hujus diéi nos perveníre fecísti: tua nos hódie salva virtúte; ut in hac die ad nullum declinémus peccátum, sed semper ad tuam justítiam faciéndam nostra procédant elóquia, dirigántur cogitatiónes et ópera. Per Dóminum. \R\ Amen.

\V\ Dóminus vobíscum. \V\ Benedicámus Dómino.

\rubric{In Choro Martyrológium.}

\V\ Pretiósa in conspéctu Dómini. \R\ Mors Sanctórum ejus.

\rubric{Et absolúte, sine} Orémus:

\oratiolabel{Oratio}

\lettrine[lines=2]{S}{ANCTA} María, et omnes Sancti intercédant pro nobis ad Dóminum, ut nos mereámur ab eo adjuvári et salvári, qui vivit et regnat in sæcula sæculórum. \R\ Amen.

\V\ Deus, in adjutórium meum inténde. \R\ Dómine, ad adjuvándum me festína. \rubric{Et dícitur ter; último additur:} Glória Patri. Sicut erat.

Kýrie, eléison. Christe, eléison. Kýrie, eléison.

Pater noster \rubric{secreto usque ad} \V\ Et ne nos indúcas in tentatiónem. \R\ Sed líbera nos a malo. \V\ Réspice in servos tuos, Dómine, et in ópera tua, et dírige fílios eórum. \R\ Et sit splendor Dómini, Dei nostri, super nos, et ópera mánuum nostrárum dírige super nos, et opus mánuum nostrárum dírige. \V\ Glória Patri. \R\ Sicut erat.

Orémus. \oratiolabel{Oratio}

\lettrine[lines=2]{D}{IRIGERE} et sanctificáre, régere et gubernáre dignáre, Dómine Deus, Rex cæli et terræ, hódie corda et córpora nostra, sensus, sermónes et actus nostros in lege tua, et in opéribus mandatórum tuórum: ut hic, et in ætérnum, te auxiliánte, salvi et líberi esse mereámur, Salvátor mundi: Qui vivis et regnas in sæcula sæculórum. \R\ Amen.

\V\ Jube, domne.

\benedlabel{Bened.} Dies et actus nostros in sua pace dispónat Dóminus omnípotens. \R\ Amen.

\variantheading{In Officio feriali}

\lettrine[lines=2]{D}{OMINUS} autem dírigat corda et córpora nostra in caritáte Dei, et patiéntia Christi. Tu autem.

\variantheading{E Proprio Sanctorum}

\variantheading{Pro Oct. Omn. Sanctorum}

\lettrine[lines=2]{B}{ENEDÍCTIO} et cláritas et sapiéntia et gratiárum áctio, honor, virtus et fortitúdo Deo nostro in sǽcula sæculórum. Amen.

\variantheading{E Communi Sanctorum}

\variantheading{Pro Apostolis}

\lettrine[lines=2]{I}{BANT} Apóstoli gaudéntes a conspéctu concílii, quóniam digni hábiti sunt pro nómine Jesu contuméliam pati.

\variantheading{Pro Martyre}

\lettrine[lines=2]{J}{USTUS} cor suum trádidit ad vigilándum dilúculo ad Dóminum, qui fecit illum, et in conspéctu Altíssimi deprecábitur.

\variantheading{Pro pluribus Martyribus}

\lettrine[lines=2]{F}{ULGEBUNT} justi, et tamquam scintíllæ in arundinéto discúrrent. Judicábunt natiónes, et dominabúntur pópulis: et regnábit Dóminus illórum in perpétuum.

\variantheading{Pro Conf. Pont.}

\lettrine[lines=2]{F}{UNGI} sacerdótio, et habére laudem in nómine ipsíus, et offérre illi incénsum dignum in odórem suavitátis.

\variantheading{Pro Conf. non Pont.}

\lettrine[lines=2]{J}{USTUM} dedúxit Dóminus per vias rectas, et osténdit illi regnum Dei, et dedit illi sciéntiam sanctórum; honestávit illum in labóribus, et complévit labóres illíus.

\variantheading{Pro Virgine Martyre}

\lettrine[lines=2]{D}{OMINE}, Deus meus, exaltásti super terram habitatiónem meam, et pro morte defluénte deprecáta sum. Invocávi Dóminum, Patrem Dómini mei, ut non derelínquat me in die tribulatiónis meæ, et in témpore superbórum sine adjutório.

\variantheading{Pro Virgine non Martyre}

\lettrine[lines=2]{O}{} \textsc{quam} pulchra est casta generátio cum claritáte! immortális est enim memória illíus: quóniam apud Deum nota est, et apud hómines.

\variantheading{Pro Martyre tantum}

\lettrine[lines=2]{L}{AUDABIT} usque ad mortem ánima mea Dóminum, quóniam éruis sustinéntes te, et líberas eos de manu angústiæ, Dómine, Deus noster.

\variantheading{Pro nec Virg. nec Mart.}

\lettrine[lines=2]{M}{ULTAE} fíliæ congregavérunt divítias: tu supergréssa es univérsas. Fallax grátia, et vana est pulchritúdo: múlier timens Dóminum, ipsa laudábitur.

\variantheading{Pro Dedicatione Ecclesiæ}

\lettrine[lines=2]{E}{T} abstérget Deus omnem lácrimam ab óculis eórum: et mors ultra non erit, neque luctus, neque clamor, neque dolor erit ultra, quia prima abiérunt. Et dixit qui sedébat in throno: Ecce, nova fácio ómnia.

\variantheading{In Officiis B. Mariæ V.}

\lettrine[lines=2]{I}{N} platéis sicut cinnamómum et bálsamum aromatízans odórem dedi: quasi myrrha elécta dedi suavitátem odóris.

\V\ Adjutórium nostrum in nómine Dómini. \R\ Qui fecit cælum et terram.

\V\ Benedícite. \R\ Deus.

\benedictiolabel{Benedictio} Dóminus nos benedícat, et ab omni malo deféndat, et ad vitam perdúcat ætérnam. Et fidélium ánimæ per misericórdiam Dei requiéscant in pace. \R\ Amen.

\rubric{Deinde dícitur tantum} Pater noster \rubric{secréto, nisi sequátur ália Hora.}

% ===== END OF BATCH 21 -- Covers autumnalis.md lines 7159-7424,
% completing "Ad Primam" for Feria Quarta in full: the hymn "Jam
% lucis orto sídere," Psalms 25, 51, 52, the conditional Psalm 96
% insertion, both Capitula, the brief Responsory, the full Preces
% apparatus, the closing Orationes, and the Lectio-brevis apparatus
% (E Proprio Sanctorum, E Communi Sanctorum). Four consecutive psalm
% incipits in autumnalis.md (Ps. 25 "JUdica," Ps. 51 "QUid," Ps. 52
% "DIxit," and the conditional Ps. 96 "DOminus") print only their
% first letter or two in full caps with the remainder of the first
% word in lowercase -- a mechanical OCR capitalization defect,
% confirmed for "DIxit" and "DOminus" by the many fully-capitalized
% correct occurrences of these exact psalm incipits already
% transcribed elsewhere in this volume (e.g. Psalm 13's "DIXIT
% insípiens" in Batch 9, and "DOMINUS regnávit" repeatedly in
% Batches 19-20) -- and silently normalized to full caps throughout,
% consistent with the uniform ALL-CAPS-opening-word convention
% applied to every psalm in this transcription. The document
% continues with "Ad Tertiam" (autumnalis.md line 7427).
% =====
\hourheading{Ad Tertiam}

Pater noster et Ave.

\V\ Deus, in adjutórium.

\hymnusheading{Hymnus}

\lettrine[lines=2]{N}{UNC}, Sancte, nobis, Spíritus,\\
Unum Patri cum Fílio,\\
Dignáre promptus íngeri\\
Nostro refúsus péctori.

Os, lingua, mens, sensus, vigor\\
Confessiónem pérsonent,\\
Flamméscat igne cáritas,\\
Accéndat ardor próximos.

Præsta, Pater piíssime,\\
Patríque compar Unice,\\
Cum Spíritu Paráclito\\
Regnans per omne sæculum.\\
Amen.

\antlabel{Ant.} Deus ádjuvat me.

\psalmheading{Psalmus 53}

\lettrine[lines=2]{D}{EUS}, in nómine tuo salvum me fac: * et in virtúte tua júdica me.

Deus, exáudi oratiónem meam: * áuribus pércipe verba oris mei.

Quóniam aliéni insurrexérunt advérsum me, et fortes quæsiérunt ánimam meam: * et non proposuérunt Deum ante conspéctum suum.

Ecce enim, Deus ádjuvat me: * et Dóminus suscéptor est ánimæ meæ.

Avérte mala inimícis meis: * et in veritáte tua dispérde illos.

Voluntárie sacrificábo tibi, * et confitébor nómini tuo, Dómine: quóniam bonum est:

Quóniam ex omni tribulatióne eripuísti me: * et super inimícos meos despéxit óculus meus.

\psalmheading{Psalmus 54, i}

\lettrine[lines=2]{E}{XAUDI}, Deus, oratiónem meam, et ne despéxeris deprecatiónem meam: * inténde mihi, et exáudi me.

Contristátus sum in exercitatióne mea: * et conturbátus sum a voce inimíci, et a tribulatióne peccatóris.

Quóniam declinavérunt in me iniquitátes: * et in ira molésti erant mihi.

Cor meum conturbátum est in me: * et formído mortis cécidit super me.

Timor et tremor venérunt super me: * et contexérunt me ténebræ.

Et dixi: Quis dabit mihi pennas sicut colúmbæ, * et volábo, et requiéscam?

Ecce, elongávi fúgiens: * et mansi in solitúdine.

Exspectábam eum, qui salvum me fecit * a pusillanimitáte spíritus et tempestáte.

Præcípita, Dómine, dívide linguas eórum: * quóniam vidi iniquitátem, et contradictiónem in civitáte.

Die ac nocte circúmdabit eam super muros ejus iníquitas: * et labor in médio ejus, et injustítia.

Et non defécit de platéis ejus * usúra, et dolus.

Quóniam si inimícus meus maledixísset mihi, * sustinuíssem útique.

Et si is, qui óderat me, super me magna locútus fuísset, * abscondíssem me fórsitan ab eo.

Tu vero, homo unánimis: * dux meus, et notus meus:

Qui simul mecum dulces capiébas cibos: * in domo Dei ambulávimus cum consénsu.

Véniat mors super illos: * et descéndant in inférnum vivéntes:

Quóniam nequítiæ in habitáculis eórum: * in médio eórum.

\psalmheading{Psalmus 54, ii}

\lettrine[lines=2]{E}{GO} autem ad Deum clamávi: * et Dóminus salvábit me.

Véspere, et mane, et merídie narrábo et annuntiábo: * et exáudiet vocem meam.

Rédimet in pace ánimam meam ab his, qui appropínquant mihi: * quóniam inter multos erant mecum.

Exáudiet Deus, et humiliábit illos, * qui est ante sæcula.

Non enim est illis commutátio, et non timuérunt Deum: * exténdit manum suam in retribuéndo.

Contaminavérunt testaméntum ejus, divísi sunt ab ira vultus ejus: * et appropinquávit cor illíus.

Mollíti sunt sermónes ejus super óleum: * et ipsi sunt jácula.

Jacta super Dóminum curam tuam, et ipse te enútriet: * non dabit in ætérnum fluctuatiónem justo.

Tu vero, Deus, dedúces eos, * in púteum intéritus.

Viri sánguinum, et dolósi non dimidiábunt dies suos: * ego autem sperábo in te, Dómine.

\antlabel{Ant.} Deus ádjuvat me: et Dóminus suscéptor est ánimæ meæ.

\variantheading{In Officio feriali}

\capitulumlabel{Capitulum}{Jer. 17, 14}

\lettrine[lines=2]{S}{ANA} me, Dómine, et sanábor: salvum me fac, et salvus ero: quóniam laus mea tu es.

\R\ br. Sana ánimam meam, * Quia peccávi tibi. Sana. \V\ Ego dixi: Dómine, miserére mei. Quia. Glória Patri. Sana. \V\ Adjútor meus esto, ne derelínquas me. \R\ Neque despícias me, Deus, salutáris meus.

\variantheading{In Vigiliis communibus:}

\variantheading{Preces feriales}

KYRIE, eléison. Christe, eléison. Kýrie, eléison.

Pater noster \rubric{secreto usque ad} \V\ Et ne nos indúcas in tentatiónem. \R\ Sed líbera nos a malo. \V\ Dómine, Deus virtútum, convérte nos. \R\ Et osténde fáciem tuam, et salvi érimus. \V\ Exsúrge, Christe, ádjuva nos. \R\ Et líbera nos propter nomen tuum. \V\ Dómine, exáudi oratiónem meam. \R\ Et clamor meus ad te véniat.

\rubric{Réliqua ut in Ordinário. 19}

\hourheading{Ad Sextam}

Pater noster et Ave.

\V\ Deus, in adjutórium.

\hymnusheading{Hymnus}

\lettrine[lines=2]{R}{ECTOR} potens, verax Deus,\\
Qui témperas rerum vices,\\
Splendóre mane illúminas,\\
Et ígnibus merídiem:

Exstíngue flammas lítium,\\
Aufer calórem nóxium,\\
Confer salútem córporum,\\
Verámque pacem córdium.

Præsta, Pater piíssime,\\
Patríque compar Unice,\\
Cum Spíritu Paráclito\\
Regnans per omne sæculum.\\
Amen.

\antlabel{Ant.} In Deo sperávi.

\psalmheading{Psalmus 55}

\lettrine[lines=2]{M}{ISERERE} mei, Deus, quóniam conculcávit me homo: * tota die impúgnans tribulávit me.

Conculcavérunt me inimíci mei tota die: * quóniam multi bellántes advérsum me.

Ab altitúdine diéi timébo: * ego vero in te sperábo.

In Deo laudábo sermónes meos, in Deo sperávi: * non timébo quid fáciat mihi caro.

Tota die verba mea exsecrabántur: * advérsum me omnes cogitatiónes eórum in malum.

Inhabitábunt et abscóndent: * ipsi calcáneum meum observábunt.

Sicut sustinuérunt ánimam meam, pro níhilo salvos fácies illos: * in ira pópulos confrínges.

Deus, vitam meam annuntiávi tibi: * posuísti lácrimas meas in conspéctu tuo.

Sicut et in promissióne tua: * tunc converténtur inimíci mei retrórsum:

In quacúmque die invocávero te: * ecce, cognóvi, quóniam Deus meus es.

In Deo laudábo verbum, in Dómino laudábo sermónem: * in Deo sperávi, non timébo quid fáciat mihi homo.

In me sunt, Deus, vota tua, * quæ reddam, laudatiónes tibi.

Quóniam eripuísti ánimam meam de morte, et pedes meos de lapsu: * ut pláceam coram Deo in lúmine vivéntium.

\psalmheading{Psalmus 56}

\lettrine[lines=2]{M}{ISERERE} mei, Deus, miserére mei: * quóniam in te confídit ánima mea.

Et in umbra alárum tuárum sperábo, * donec tránseat iníquitas.

Clamábo ad Deum altíssimum: * Deum qui benefécit mihi.

Misit de cælo, et liberávit me: * dedit in oppróbrium conculcántes me.

Misit Deus misericórdiam suam, et veritátem suam, * et erípuit ánimam meam de médio catulórum leónum: dormívi conturbátus.

Fílii hóminum dentes eórum arma et sagíttæ: * et lingua eórum gládius acútus.

Exaltáre super cælos, Deus, * et in omnem terram glória tua.

Láqueum paravérunt pédibus meis: * et incurvavérunt ánimam meam.

Fodérunt ante fáciem meam fóveam: * et incidérunt in eam.

Parátum cor meum, Deus, parátum cor meum: * cantábo, et psalmum dicam.

Exsúrge, glória mea, exsúrge, psaltérium et cíthara: * exsúrgam dilúculo.

Confitébor tibi in pópulis, Dómine: * et psalmum dicam tibi in géntibus:

Quóniam magnificáta est usque ad cælos misericórdia tua, * et usque ad nubes véritas tua.

Exaltáre super cælos, Deus: * et super omnem terram glória tua.

\psalmheading{Psalmus 57}

\lettrine[lines=2]{S}{I} vere útique justítiam loquímini: * recta judicáte, fílii hóminum.

Etenim in corde iniquitátes operámini: * in terra injustítias manus vestræ concínnant.

Alienáti sunt peccatóres a vulva, erravérunt ab útero: * locúti sunt falsa.

Furor illis secúndum similitúdinem serpéntis: * sicut áspidis surdæ, et obturántis aures suas.

Quæ non exáudiet vocem incantántium: * et venéfici incantántis sapiénter.

Deus cónteret dentes eórum in ore ipsórum: * molas leónum confrínget Dóminus.

Ad níhilum devénient tamquam aqua decúrrens: * inténdit arcum suum donec infirméntur.

Sicut cera, quæ fluit, auferéntur: * supercécidit ignis, et non vidérunt solem.

Priúsquam intellígerent spinæ vestræ rhamnum: * sicut vivéntes, sic in ira absórbet eos.

Lætábitur justus cum víderit vindíctam: * manus suas lavábit in sánguine peccatóris.

Et dicet homo: Si útique est fructus justo: * útique est Deus júdicans eos in terra.

\antlabel{Ant.} In Deo sperávi, non timébo quid fáciat mihi homo.

\variantheading{In Officio feriali}

\capitulumlabel{Capitulum}{Rom. 13, 8}

\lettrine[lines=2]{N}{EMINI} quidquam debeátis, nisi ut ínvicem diligátis: qui enim díligit próximum, legem implévit.

\R\ br. Benedícam Dóminum * In omni témpore. Benedícam. \V\ Semper laus ejus in ore meo. In. Glória Patri. Benedícam. \V\ Dóminus regit me, et nihil mihi déerit. \R\ In loco páscuæ ibi me collocávit.

\variantheading{In Vigiliis communibus:}

\variantheading{Preces feriales}

KYRIE, eléison. Christe, eléison. Kýrie, eléison.

Pater noster \rubric{secreto usque ad} \V\ Et ne nos indúcas in tentatiónem. \R\ Sed líbera nos a malo. \V\ Dómine, Deus virtútum, convérte nos. \R\ Et osténde fáciem tuam, et salvi érimus. \V\ Exsúrge, Christe, ádjuva nos. \R\ Et líbera nos propter nomen tuum. \V\ Dómine, exáudi oratiónem meam. \R\ Et clamor meus ad te véniat.

\rubric{Réliqua ut in Ordinário. 20}

\hourheading{Ad Nonam}

Pater noster et Ave.

\V\ Deus, in adjutórium.

\hymnusheading{Hymnus}

\lettrine[lines=2]{R}{ERUM}, Deus, tenax vigor,\\
Immótus in te pérmanens,\\
Lucis diúrnæ témpora\\
Succéssibus detérminans:

Largíre lumen véspere,\\
Quo vita nusquam décidat,\\
Sed prǽmium mortis sacræ\\
Perénnis instet glória.

Præsta, Pater piíssime,\\
Patríque compar Unice,\\
Cum Spíritu Paráclito\\
Regnans per omne sæculum.\\
Amen.

\antlabel{Ant.} Deus meus.

\psalmheading{Psalmus 58, i}

\lettrine[lines=2]{E}{RIPE} me de inimícis meis, Deus meus: * et ab insurgéntibus in me líbera me.

Eripe me de operántibus iniquitátem: * et de viris sánguinum salva me.

Quia ecce cepérunt ánimam meam: * irruérunt in me fortes.

Neque iníquitas mea, neque peccátum meum, Dómine: * sine iniquitáte cucúrri, et diréxi.

Exsúrge in occúrsum meum, et vide: * et tu, Dómine, Deus virtútum, Deus Israël,

Inténde ad visitándas omnes gentes: * non misereáris ómnibus, qui operántur iniquitátem.

Converténtur ad vésperam: et famem patiéntur ut canes, * et circuíbunt civitátem.

Ecce, loquéntur in ore suo, et gládius in lábiis eórum: * quóniam quis audívit?

Et tu, Dómine, deridébis eos: * ad níhilum dedúces omnes gentes.

Fortitúdinem meam ad te custódiam, quia, Deus, suscéptor meus es: * Deus meus, misericórdia ejus prævéniet me.

\psalmheading{Psalmus 58, ii}

\lettrine[lines=2]{D}{EUS} osténdet mihi super inimícos meos, ne occídas eos: * nequándo obliviscántur pópuli mei.

Dispérge illos in virtúte tua: * et depóne eos, protéctor meus, Dómine:

Delíctum oris eórum, sermónem labiórum ipsórum: * et comprehendántur in supérbia sua.

Et de exsecratióne et mendácio annuntiabúntur in consummatióne: * in ira consummatiónis, et non erunt.

Et scient quia Deus dominábitur Jacob: * et fínium terræ.

Converténtur ad vésperam: et famem patiéntur ut canes, * et circuíbunt civitátem.

Ipsi dispergéntur ad manducándum: * si vero non fúerint saturáti, et murmurábunt.

Ego autem cantábo fortitúdinem tuam: * et exsultábo mane misericórdiam tuam.

Quia factus es suscéptor meus, * et refúgium meum, in die tribulatiónis meæ.

Adjútor meus, tibi psallam, quia, Deus, suscéptor meus es: * Deus meus, misericórdia mea.

\psalmheading{Psalmus 59}

\lettrine[lines=2]{D}{EUS}, repulísti nos, et destruxísti nos: * irátus es, et misértus es nobis.

Commovísti terram, et conturbásti eam: * sana contritiónes ejus, quia commóta est.

Ostendísti pópulo tuo dura: * potásti nos vino compunctiónis.

Dedísti metuéntibus te significatiónem: * ut fúgiant a fácie arcus:

Ut liberéntur dilécti tui: * salvum fac déxtera tua, et exáudi me.

Deus locútus est in sancto suo: * Lætábor, et partíbor Síchimam: et convállem tabernaculórum metíbor.

Meus est Gálaad, et meus est Manásses: * et Ephraïm fortitúdo cápitis mei.

Juda rex meus: * Moab olla spei meæ.

In Idumǽam exténdam calceaméntum meum: * mihi aliénigenæ súbditi sunt.

Quis dedúcet me in civitátem munítam? * quis dedúcet me usque in Idumǽam?

Nonne tu, Deus, qui repulísti nos, * et non egrediéris, Deus, in virtútibus nostris?

Da nobis auxílium de tribulatióne: * quia vana salus hóminis.

In Deo faciémus virtútem: * et ipse ad níhilum dedúcet tribulántes nos.

\antlabel{Ant.} Deus meus, misericórdia tua prævéniet me.

\variantheading{In Officio feriali}

\capitulumlabel{Capit.}{1 Petri 1, 17-19}

\lettrine[lines=2]{I}{N} timóre incolátus vestri témpore conversámini: sciéntes quod non corruptibílibus auro vel argénto redémpti estis, sed pretióso sánguine quasi Agni immaculáti Christi.

\R\ br. Rédime me, Dómine, * Et miserére mei. Rédime. \V\ Pes enim meus stetit in via recta. Et miserére mei. Glória Patri. Rédime. \V\ Ab occúltis meis munda me, Dómine. \R\ Et ab aliénis parce servo tuo.

\variantheading{In Vigiliis communibus:}

\variantheading{Preces feriales}

KYRIE, eléison. Christe, eléison. Kýrie, eléison.

Pater noster \rubric{secreto usque ad} \V\ Et ne nos indúcas in tentatiónem. \R\ Sed líbera nos a malo. \V\ Dómine, Deus virtútum, convérte nos. \R\ Et osténde fáciem tuam, et salvi érimus. \V\ Exsúrge, Christe, ádjuva nos. \R\ Et líbera nos propter nomen tuum. \V\ Dómine, exáudi oratiónem meam. \R\ Et clamor meus ad te véniat.

\rubric{Réliqua ut in Ordinário. 22}

% ===== END OF BATCH 22 -- Covers autumnalis.md lines 7427-7796,
% completing the Little Hours of Feria Quarta in full: "Ad
% Tertiam" (Psalms 53, 54/i, 54/ii), "Ad Sextam" (Psalms 55, 56,
% 57), and "Ad Nonam" (Psalms 58/i, 58/ii, 59), each with its own
% hymn, Capitulum, brief Responsory, and "In Vigiliis communibus:
% Preces feriales" apparatus. Three further occurrences of the
% recurring stray OCR "\*" artifact in the hour-hymns are silently
% dropped, consistent with every prior occurrence. Several more
% psalm incipits in this batch carry the same partial-caps OCR
% defect already documented in Batch 21 ("DEus," "EXaudi," "EGo,"
% "SAna" [Capitulum], "MIserére" twice, "ERipe," "DEus" twice,
% "NEmini" [Capitulum]) and are silently normalized to full caps,
% consistent with the uniform convention and confirmed by the many
% correctly-capitalized occurrences of these same incipits elsewhere
% in this volume. The document continues with "Ad Vesperas"
% (autumnalis.md line 7799).
% =====
\hourheading{Ad Vesperas}

Pater noster et Ave.

\V\ Deus, in adjutórium.

\antlabel{Ant.} Beáti omnes, * \textdagger\ qui timent Dóminum. \textdaggerdbl

\psalmheading{Psalmus 127}

\lettrine[lines=2]{B}{EATI} omnes, \textdagger\ qui timent Dóminum, * \textdaggerdbl\ qui ámbulant in viis ejus.

Labóres mánuum tuárum quia manducábis: * beátus es, et bene tibi erit.

Uxor tua sicut vitis abúndans, * in latéribus domus tuæ.

Fílii tui sicut novéllæ olivárum, * in circúitu mensæ tuæ.

Ecce, sic benedicétur homo, * qui timet Dóminum.

Benedícat tibi Dóminus ex Sion: * et vídeas bona Jerúsalem ómnibus diébus vitæ tuæ.

Et vídeas fílios filiórum tuórum, * pacem super Israël.

\antlabel{Ant.} Beáti omnes, qui timent Dóminum.

\antlabel{Ant.} Confundántur omnes, * qui odérunt Sion.

\psalmheading{Psalmus 128}

\lettrine[lines=2]{S}{ÆPE} expugnavérunt me a juventúte mea, * dicat nunc Israël:

Sæpe expugnavérunt me a juventúte mea: * étenim non potuérunt mihi.

Supra dorsum meum fabricavérunt peccatóres: * prolongavérunt iniquitátem suam.

Dóminus justus concídit cervíces peccatórum: * confundántur et convertántur retrórsum omnes, qui odérunt Sion.

Fiant sicut fœnum tectórum: * quod priúsquam evellátur, exáruit:

De quo non implévit manum suam qui metit, * et sinum suum qui manípulos cólligit.

Et non dixérunt qui præteríbant: Benedíctio Dómini super vos: * benedíximus vobis in nómine Dómini.

\antlabel{Ant.} Confundántur omnes, qui odérunt Sion.

\antlabel{Ant.} De profúndis * \textdagger\ clamávi ad te, Dómine. \textdaggerdbl

\psalmheading{Psalmus 129}

\lettrine[lines=2]{D}{E} profúndis \textdagger\ clamávi ad te, Dómine: * \textdaggerdbl\ Dómine, exáudi vocem meam:

Fiant aures tuæ inténdentes, * in vocem deprecatiónis meæ.

Si iniquitátes observáveris, Dómine: * Dómine, quis sustinébit?

Quia apud te propitiátio est: * et propter legem tuam sustínui te, Dómine.

Sustínuit ánima mea in verbo ejus: * sperávit ánima mea in Dómino.

A custódia matutína usque ad noctem: * speret Israël in Dómino.

Quia apud Dóminum misericórdia: * et copiósa apud eum redémptio.

Et ipse rédimet Israël, * ex ómnibus iniquitátibus ejus.

\antlabel{Ant.} De profúndis clamávi ad te, Dómine.

\antlabel{Ant.} Dómine, * \textdagger\ non est exaltátum cor meum. \textdaggerdbl

\psalmheading{Psalmus 130}

\lettrine[lines=2]{D}{OMINE}, \textdagger\ non est exaltátum cor meum: * \textdaggerdbl\ neque eláti sunt óculi mei.

Neque ambulávi in magnis: * neque in mirabílibus super me.

Si non humíliter sentiébam: * sed exaltávi ánimam meam:

Sicut ablactátus est super matre sua, * ita retribútio in ánima mea.

Speret Israël in Dómino, * ex hoc nunc et usque in sæculum.

\antlabel{Ant.} Dómine, non est exaltátum cor meum.

\antlabel{Ant.} Elégit Dóminus * Sion in habitatiónem sibi.

\psalmheading{Psalmus 131}

\lettrine[lines=2]{M}{EMENTO}, Dómine, David, * et omnis mansuetúdinis ejus:

Sicut jurávit Dómino, * votum vovit Deo Jacob:

Si introíero in tabernáculum domus meæ, * si ascéndero in lectum strati mei:

Si dédero somnum óculis meis, * et pálpebris meis dormitatiónem:

Et réquiem tempóribus meis: donec invéniam locum Dómino, * tabernáculum Deo Jacob.

Ecce, audívimus eam in Ephráta: * invénimus eam in campis silvæ.

Introíbimus in tabernáculum ejus: * adorábimus in loco, ubi stetérunt pedes ejus.

Surge, Dómine, in réquiem tuam, * tu et arca sanctificatiónis tuæ.

Sacerdótes tui induántur justítiam: * et sancti tui exsúltent.

Propter David, servum tuum, * non avértas fáciem Christi tui.

Jurávit Dóminus David veritátem, et non frustrábitur eam: * De fructu ventris tui ponam super sedem tuam.

Si custodíerint fílii tui testaméntum meum, * et testimónia mea hæc, quæ docébo eos:

Et fílii eórum usque in sæculum, * sedébunt super sedem tuam.

Quóniam elégit Dóminus Sion: * elégit eam in habitatiónem sibi.

Hæc réquies mea in sæculum sæculi: * hic habitábo quóniam elégi eam.

Víduam ejus benedícens benedícam: * páuperes ejus saturábo pánibus.

Sacerdótes ejus índuam salutári: * et sancti ejus exsultatióne exsultábunt.

Illuc prodúcam cornu David, * parávi lucérnam Christo meo.

Inimícos ejus índuam confusióne: * super ipsum autem efflorébit sanctificátio mea.

\antlabel{Ant.} Elégit Dóminus Sion in habitatiónem sibi.

\variantheading{In Officio feriali}

\capitulumlabel{Capitulum}{2 Cor. 1, 3-4}

\lettrine[lines=2]{B}{ENEDICTUS} Deus, et Pater Dómini nostri Jesu Christi, Pater misericordiárum, et Deus totíus consolatiónis, qui consolátur nos in omni tribulatióne nostra.

\hymnusheading{Hymnus}

\lettrine[lines=2]{C}{ÆLI} Deus sanctíssime,\\
Qui lúcidas mundi plagas\\
Candóre pingis ígneo,\\
Augens decóro lúmine:

Quarto die qui flámmeam\\
Dum solis accéndis rotam,\\
Lunæ minístras órdinem,\\
Vagósque cursus síderum:

Ut nóctibus, vel lúmini\\
Diremptiónis términum,\\
Primórdiis et ménsium\\
Signum dares notíssimum;

Expélle noctem córdium:\\
Abstérge sordes méntium:\\
Resólve culpæ vínculum:\\
Evérte moles críminum.

Præsta, Pater piíssime,\\
Patríque compar Unice,\\
Cum Spíritu Paráclito\\
Regnans per omne sæculum.\\
Amen.

\V\ Dirigátur, Dómine, orátio mea. \R\ Sicut incénsum in conspéctu tuo.

\antlabel{Ad Magnif. Ant.} Respéxit Dóminus * humilitátem meam, et fecit in me magna, qui potens est.

\rubric{Réliqua ut in Ordinário. 22}

\hourheading{Ad Completorium}

\V\ Jube, domne.

\benedlabel{Bened.} Noctem quiétam, et finem perféctum concédat nobis Dóminus omnípotens. \R\ Amen.

\lectioheading{Lectio br.}{1 Petri 5, 8-9}

\lettrine[lines=2]{F}{RATRES}: Sóbrii estóte, et vigiláte: quia adversárius vester diábolus tamquam leo rúgiens círcuit, quærens quem dévoret: cui resístite fortes in fide. Tu autem, Dómine, miserére nobis. \R\ Deo grátias.

\V\ Adjutórium nostrum in nómine Dómini. \R\ Qui fecit cælum et terram.

Pater noster \rubric{quod dícitur totum secréto.}

\formulalabel{Confíteor.} Misereátur. Indulgéntiam.

\V\ Convérte nos, Deus, salutáris noster. \R\ Et avérte iram tuam a nobis.

\V\ Deus, in adjutórium.

\antlabel{Ant.} Immíttet Angelus Dómini.

\psalmheading{Psalmus 33, i}

\lettrine[lines=2]{B}{ENEDICAM} Dóminum in omni témpore: * semper laus ejus in ore meo.

In Dómino laudábitur ánima mea: * áudiant mansuéti, et læténtur.

Magnificáte Dóminum mecum: * et exaltémus nomen ejus in idípsum.

Exquisívi Dóminum, et exaudívit me: * et ex ómnibus tribulatiónibus meis erípuit me.

Accédite ad eum, et illuminámini: * et fácies vestræ non confundéntur.

Ipse pauper clamávit, et Dóminus exaudívit eum: * et de ómnibus tribulatiónibus ejus salvávit eum.

Immíttet Angelus Dómini in circúitu timéntium eum: * et erípiet eos.

Gustáte, et vidéte quóniam suávis est Dóminus: * beátus vir, qui sperat in eo.

Timéte Dóminum, omnes sancti ejus: * quóniam non est inópia timéntibus eum.

Dívites eguérunt et esuriérunt: * inquiréntes autem Dóminum non minuéntur omni bono.

\psalmheading{Psalmus 33, ii}

\lettrine[lines=2]{V}{ENITE}, fílii, audíte me: * timórem Dómini docébo vos.

Quis est homo qui vult vitam: * díligit dies vidére bonos?

Próhibe linguam tuam a malo: * et lábia tua ne loquántur dolum.

Divérte a malo, et fac bonum: * inquíre pacem, et perséquere eam.

Oculi Dómini super justos: * et aures ejus in preces eórum.

Vultus autem Dómini super faciéntes mala: * ut perdat de terra memóriam eórum.

Clamavérunt justi, et Dóminus exaudívit eos: * et ex ómnibus tribulatiónibus eórum liberávit eos.

Juxta est Dóminus iis, qui tribuláto sunt corde: * et húmiles spíritu salvábit.

Multæ tribulatiónes justórum: * et de ómnibus his liberábit eos Dóminus.

Custódit Dóminus ómnia ossa eórum: * unum ex his non conterétur.

Mors peccatórum péssima: * et qui odérunt justum, delínquent.

Rédimet Dóminus ánimas servórum suórum: * et non delínquent omnes qui sperant in eo.

\psalmheading{Psalmus 60}

\lettrine[lines=2]{E}{XAUDI}, Deus, deprecatiónem meam: * inténde oratióni meæ.

A fínibus terræ ad te clamávi: * dum anxiarétur cor meum, in petra exaltásti me.

Deduxísti me, quia factus es spes mea: * turris fortitúdinis a fácie inimíci.

Inhabitábo in tabernáculo tuo in sæcula: * prótegar in velaménto alárum tuárum.

Quóniam tu, Deus meus, exaudísti oratiónem meam: * dedísti hereditátem timéntibus nomen tuum.

Dies super dies regis adjícies: * annos ejus usque in diem generatiónis et generatiónis.

Pérmanet in ætérnum in conspéctu Dei: * misericórdiam et veritátem ejus quis requíret?

Sic psalmum dicam nómini tuo in sæculum sæculi: * ut reddam vota mea de die in diem.

\antlabel{Ant.} Immíttet Angelus Dómini in circúitu timéntium eum: et erípiet eos.

\hymnusheading{Hymnus}

\lettrine[lines=2]{T}{E} lucis ante términum,\\
Rerum Creátor, póscimus,\\
Ut pro tua cleméntia\\
Sis præsul et custódia.

Procul recédant sómnia,\\
Et nóctium phantásmata,\\
Hostémque nostrum cómprime,\\
Ne polluántur córpora.

Præsta, Pater piíssime,\\
Patríque compar Unice,\\
Cum Spíritu Paráclito\\
Regnans per omne sæculum.\\
Amen.

\capitulumlabel{Capitulum}{Jer. 14, 9}

\lettrine[lines=2]{T}{U} autem in nobis es, Dómine, et nomen sanctum tuum invocátum est super nos: ne derelínquas nos, Dómine, Deus noster.

\R\ br. In manus tuas, Dómine, * Comméndo spíritum meum. In manus. \V\ Redemísti nos, Dómine, Deus veritátis. Comméndo. Glória Patri. In manus.

\V\ Custódi nos, Dómine, ut pupíllam óculi. \R\ Sub umbra alárum tuárum prótege nos.

\antlabel{Ant.} Salva nos.

\canticumheading{Canticum Simeonis}{Luc. 2, 29-32}

\lettrine[lines=2]{N}{UNC} dimíttis servum tuum, Dómine, * secúndum verbum tuum in pace:

Quia vidérunt óculi mei * salutáre tuum,

Quod parásti * ante fáciem ómnium populórum,

Lumen ad revelatiónem géntium, * et glóriam plebis tuæ Israël.

Glória Patri.

\antlabel{Ant.} Salva nos, Dómine, vigilántes, custódi nos dormiéntes; ut vigilémus cum Christo, et requiescámus in pace.

\variantheading{Preces (si dicendæ sint)}

KYRIE, eléison. Christe, eléison. Kýrie, eléison.

Pater noster \rubric{secreto usque ad} \V\ Et ne nos indúcas in tentatiónem. \R\ Sed líbera nos a malo.

Credo in Deum \rubric{secreto usque ad} \V\ Carnis resurrectiónem. \R\ Vitam ætérnam. Amen. \V\ Benedíctus es, Dómine, Deus patrum nostrórum. \R\ Et laudábilis et gloriósus in sæcula. \V\ Benedicámus Patrem et Fílium cum Sancto Spíritu. \R\ Laudémus, et superexaltémus eum in sæcula. \V\ Benedíctus es, Dómine, in firmaménto cæli. \R\ Et laudábilis, et gloriósus, et superexaltátus in sæcula. \V\ Benedícat et custódiat nos omnípotens et miséricors Dóminus. \R\ Amen. \V\ Dignáre, Dómine, nocte ista. \R\ Sine peccáto nos custodíre. \V\ Miserére nostri, Dómine. \R\ Miserére nostri. \V\ Fiat misericórdia tua, Dómine, super nos. \R\ Quemádmodum sperávimus in te. \V\ Dómine, exáudi oratiónem meam. \R\ Et clamor meus ad te véniat.

\V\ Dóminus vobíscum.

Orémus. \oratiolabel{Oratio}

\lettrine[lines=2]{V}{ISITA}, quæsumus, Dómine, habitatiónem istam, et omnes insídias inimíci ab ea longe repélle: Angeli tui sancti hábitent in ea, qui nos in pace custódiant; et benedíctio tua sit super nos semper. Per Dóminum nostrum. \R\ Amen.

\V\ Dóminus vobíscum.

\V\ Benedicámus Dómino.

\benedictiolabel{Benedictio} Benedícat et custódiat nos omnípotens et miséricors Dóminus, Pater, et Fílius, et Spíritus Sanctus. \R\ Amen.

\variantheading{Antiphona finalis B. M. V.}

\V\ Divínum auxílium.

Pater, Ave et Credo.

% ===== END OF BATCH 23 -- Covers autumnalis.md lines 7799-8128,
% completing "Ad Vesperas" (Psalms 127, 128, 129, 130, 131, hymn
% "Cæli Deus sanctíssime") and "Ad Completorium" (Psalms 33/i,
% 33/ii, 60, hymn "Te lucis ante términum," Canticum Simeonis,
% Preces), and with them the ENTIRE "Feria Quarta" Office. More
% partial-caps OCR defects ("MEmento," "BEnedictus" [Capitulum],
% "DOmine," "DE profúndis," "BEnedicam," "VEnite," "EXaudi") are
% silently normalized to full caps, consistent with Batches 21-22.
% The partial-bold "Confíteor." (not the full "Misereátur.
% Indulgéntiam.") again appears at Compline, rendered the same way
% as in Feria Secunda and Feria Tertia. The document continues with
% "Feria Quinta" (autumnalis.md line 8131), the fourth weekday of the
% Psalterium cycle.
% =====
\feriaheading{Feria Quinta}

\hourheading{Ad Matutinum}

Pater, Ave et Credo.

\V\ Dómine, lábia mea.

\V\ Deus, in adjutórium.

\variantheading{In Officio feriali}

\invitatlabel{Invit.} Regem ventúrum Dóminum, * Veníte, adorémus. \psalmlabel{Ps. 94} Veníte.

\hymnusheading{Hymnus}

\lettrine[lines=2]{N}{OX} atra rerum cóntegit\\
Terræ colóres ómnium:\\
Nos confiténtes póscimus\\
Te, juste judex córdium:

Ut áuferas piácula,\\
Sordésque mentis ábluas:\\
Donésque, Christe, grátiam,\\
Ut arceántur crímina.

Mens ecce torpet ímpia,\\
Quam culpa mordet nóxia:\\
Obscúra gestit tóllere,\\
Et te, Redémptor, quærere.

Repélle tu calíginem\\
Intrínsecus quam máxime,\\
Ut in beáto gáudeat\\
Se collocári lúmine.

Præsta, Pater piíssime,\\
Patríque compar Unice,\\
Cum Spíritu Paráclito\\
Regnans per omne sæculum.\\
Amen.

\nocturnoheading{In I Nocturno}

\antlabel{Ant.} In Deo salutáre meum, * et glória mea: et spes mea in Deo est.

\psalmheading{Psalmus 61}

\lettrine[lines=2]{N}{ONNE} Deo subjécta erit ánima mea? * ab ipso enim salutáre meum.

Nam et ipse Deus meus, et salutáris meus: * suscéptor meus, non movébor ámplius.

Quoúsque irrúitis in hóminem? * interfícitis univérsi vos: tamquam paríeti inclináto et macériæ depúlsæ?

Verúmtamen prétium meum cogitavérunt repéllere, cucúrri in siti: * ore suo benedicébant, et corde suo maledicébant.

Verúmtamen Deo subjécta esto, ánima mea: * quóniam ab ipso patiéntia mea.

Quia ipse Deus meus, et salvátor meus: * adjútor meus, non emigrábo.

In Deo salutáre meum, et glória mea: * Deus auxílii mei, et spes mea in Deo est.

Speráte in eo, omnis congregátio pópuli, effúndite coram illo corda vestra: * Deus adjútor noster in ætérnum.

Verúmtamen vani fílii hóminum, mendáces fílii hóminum in statéris: * ut decípiant ipsi de vanitáte in idípsum.

Nolíte speráre in iniquitáte, et rapínas nolíte concupíscere: * divítiæ si áffluant, nolíte cor appónere.

Semel locútus est Deus, duo hæc audívi, quia potéstas Dei est, et tibi, Dómine, misericórdia: * quia tu reddes unicuíque juxta ópera sua.

\antlabel{Ant.} In Deo salutáre meum, et glória mea: et spes mea in Deo est.

\antlabel{Ant.} Vidéte ópera Dómini, * et audítam fácite vocem laudis ejus.

\psalmheading{Psalmus 65, i}

\lettrine[lines=2]{J}{UBILATE} Deo, omnis terra, psalmum dícite nómini ejus: * date glóriam laudi ejus.

Dícite Deo: Quam terribília sunt ópera tua, Dómine! * in multitúdine virtútis tuæ mentiéntur tibi inimíci tui.

Omnis terra adóret te, et psallat tibi: * psalmum dicat nómini tuo.

Veníte, et vidéte ópera Dei: * terríbilis in consíliis super fílios hóminum.

Qui convértit mare in áridam, in flúmine pertransíbunt pede: * ibi lætábimur in ipso.

Qui dominátur in virtúte sua in ætérnum, óculi ejus super gentes respíciunt: * qui exásperant non exalténtur in semetípsis.

Benedícite, gentes, Deum nostrum: * et audítam fácite vocem laudis ejus,

Qui pósuit ánimam meam ad vitam: * et non dedit in commotiónem pedes meos.

Quóniam probásti nos, Deus: * igne nos examinásti, sicut examinátur argéntum.

Induxísti nos in láqueum, posuísti tribulatiónes in dorso nostro: * imposuísti hómines super cápita nostra.

Transívimus per ignem et aquam: * et eduxísti nos in refrigérium.

\antlabel{Ant.} Vidéte ópera Dómini, et audítam fácite vocem laudis ejus.

\antlabel{Ant.} Audíte, omnes, * qui timétis Deum, quanta fecit ánimæ meæ.

\psalmheading{Psalmus 65, ii}

\lettrine[lines=2]{I}{NTROIBO} in domum tuam in holocáustis: * reddam tibi vota mea, quæ distinxérunt lábia mea.

Et locútum est os meum, * in tribulatióne mea.

Holocáusta medulláta ófferam tibi cum incénso aríetum: * ófferam tibi boves cum hircis.

Veníte, audíte, et narrábo, omnes, qui timétis Deum, * quanta fecit ánimæ meæ.

Ad ipsum ore meo clamávi, * et exaltávi sub lingua mea.

Iniquitátem si aspéxi in corde meo, * non exáudiet Dóminus.

Proptérea exaudívit Deus, * et atténdit voci deprecatiónis meæ.

Benedíctus Deus, * qui non amóvit oratiónem meam, et misericórdiam suam a me.

\antlabel{Ant.} Audíte, omnes, qui timétis Deum, quanta fecit ánimæ meæ.

\variantheading{In Festis novem Lectionum:}

\V\ Non amóvit Dóminus oratiónem meam. \R\ Et misericórdiam suam a me.

\absolutiolabel{Absolutio.} Exáudi, Dómine.

\benedlabel{Bened.} 1 Benedictióne perpétua. 2 Unigénitus Dei Fílius. 3 Spíritus Sancti grátia.

\nocturnoheading{In II Nocturno}

\antlabel{Ant.} Exsúrgat Deus, * \textdagger\ et dissipéntur inimíci ejus. \textdaggerdbl

\psalmheading{Psalmus 67, i}

\lettrine[lines=2]{E}{XSURGAT} Deus, \textdagger\ et dissipéntur inimíci ejus, * \textdaggerdbl\ et fúgiant qui odérunt eum, a fácie ejus.

Sicut déficit fumus, defíciant: * sicut fluit cera a fácie ignis, sic péreant peccatóres a fácie Dei.

Et justi epuléntur, et exsúltent in conspéctu Dei: * et delecténtur in lætítia.

Cantáte Deo, psalmum dícite nómini ejus: * iter fácite ei, qui ascéndit super occásum: Dóminus nomen illi.

Exsultáte in conspéctu ejus: * turbabúntur a fácie ejus, patris orphanórum, et júdicis viduárum.

Deus in loco sancto suo: * Deus qui inhabitáre facit uníus moris in domo:

Qui edúcit vinctos in fortitúdine, * simíliter eos, qui exásperant, qui hábitant in sepúlcris.

Deus, cum egrederéris in conspéctu pópuli tui, * cum pertransíres in desérto:

Terra mota est, étenim cæli distillavérunt a fácie Dei Sínaï, * a fácie Dei Israël.

Plúviam voluntáriam segregábis, Deus, hereditáti tuæ: * et infirmáta est, tu vero perfecísti eam.

Animália tua habitábunt in ea: * parásti in dulcédine tua páuperi, Deus.

\antlabel{Ant.} Exsúrgat Deus, et dissipéntur inimíci ejus.

\antlabel{Ant.} Deus noster, * Deus salvos faciéndi: et Dómini sunt éxitus mortis.

\psalmheading{Psalmus 67, ii}

\lettrine[lines=2]{D}{OMINUS} dabit verbum evangelizántibus, * virtúte multa.

Rex virtútum dilécti dilécti: * et speciéi domus divídere spólia.

Si dormiátis inter médios cleros, pennæ colúmbæ deargentátæ, * et posterióra dorsi ejus in pallóre auri.

Dum discérnit cæléstis reges super eam, nive dealbabúntur in Selmon: * mons Dei, mons pinguis.

Mons coagulátus, mons pinguis: * ut quid suspicámini montes coagulátos?

Mons, in quo beneplácitum est Deo habitáre in eo: * étenim Dóminus habitábit in finem.

Currus Dei decem míllibus múltiplex, míllia lætántium: * Dóminus in eis in Sina in sancto.

Ascendísti in altum, cepísti captivitátem: * accepísti dona in homínibus:

Etenim non credéntes, * inhabitáre Dóminum Deum.

Benedíctus Dóminus die quotídie: * prósperum iter fáciet nobis Deus salutárium nostrórum.

Deus noster, Deus salvos faciéndi: * et Dómini Dómini éxitus mortis.

Verúmtamen Deus confrínget cápita inimicórum suórum: * vérticem capílli perambulántium in delíctis suis.

Dixit Dóminus: Ex Basan convértam, * convértam in profúndum maris:

Ut intingátur pes tuus in sánguine: * lingua canum tuórum ex inimícis, ab ipso.

\antlabel{Ant.} Deus noster, Deus salvos faciéndi: et Dómini sunt éxitus mortis.

\antlabel{Ant.} In ecclésiis * benedícite Dómino Deo.

\psalmheading{Psalmus 67, iii}

\lettrine[lines=2]{V}{IDERUNT} ingréssus tuos, Deus, * ingréssus Dei mei: regis mei qui est in sancto.

Prævenérunt príncipes conjúncti psalléntibus, * in médio juvenculárum tympanistriárum.

In ecclésiis benedícite Deo Dómino, * de fóntibus Israël.

Ibi Bénjamin adolescéntulus, * in mentis excéssu.

Príncipes Juda, duces eórum: * príncipes Zábulon, príncipes Néphthali.

Manda, Deus, virtúti tuæ: * confírma hoc, Deus, quod operátus es in nobis.

A templo tuo in Jerúsalem, * tibi ófferent reges múnera.

Íncrepa feras arúndinis, congregátio taurórum in vaccis populórum: * ut exclúdant eos, qui probáti sunt argénto.

Díssipa gentes, quæ bella volunt: vénient legáti ex Ægýpto: * Æthiópia prævéniet manus ejus Deo.

Regna terræ, cantáte Deo: * psállite Dómino.

Psállite Deo, qui ascéndit super cælum cæli, * ad Oriéntem.

Ecce, dabit voci suæ vocem virtútis, date glóriam Deo super Israël, * magnificéntia ejus, et virtus ejus in núbibus.

Mirábilis Deus in sanctis suis, Deus Israël ipse dabit virtútem, et fortitúdinem plebi suæ, * benedíctus Deus.

\antlabel{Ant.} In ecclésiis benedícite Dómino Deo.

\variantheading{In Festis novem Lectionum:}

\V\ Mirábilis Deus in sanctis suis. \R\ Deus Israël dabit virtútem plebi tuæ.

\absolutiolabel{Absolutio.} Ipsíus píetas et.

\benedlabel{Bened.} 4 Deus Pater omnípotens. 5 Christus perpétuæ. 6 Ignem sui amóris.

\nocturnoheading{In III Nocturno}

\antlabel{Ant.} Salvum me fac, * \textdagger\ Deus: quóniam intravérunt aquæ usque ad ánimam meam. \textdaggerdbl

\psalmheading{Psalmus 68, i}

\lettrine[lines=2]{S}{ALVUM} me fac, \textdagger\ Deus: * \textdaggerdbl\ quóniam intravérunt aquæ usque ad ánimam meam.

Infíxus sum in limo profúndi: * et non est substántia.

Veni in altitúdinem maris: * et tempéstas demérsit me.

Laborávi clamans, raucæ factæ sunt fauces meæ: * defecérunt óculi mei, dum spero in Deum meum.

Multiplicáti sunt super capíllos cápitis mei, * qui odérunt me gratis.

Confortáti sunt qui persecúti sunt me inimíci mei injúste: * quæ non rápui, tunc exsolvébam.

Deus, tu scis insipiéntiam meam: * et delícta mea a te non sunt abscóndita.

Non erubéscant in me qui exspéctant te, Dómine, * Dómine virtútum.

Non confundántur super me * qui quærunt te, Deus Israël.

Quóniam propter te sustínui oppróbrium: * opéruit confúsio fáciem meam.

Extráneus factus sum frátribus meis, * et peregrínus fíliis matris meæ.

Quóniam zelus domus tuæ comédit me: * et oppróbria exprobrántium tibi cecidérunt super me.

Et opérui in jejúnio ánimam meam: * et factum est in oppróbrium mihi.

Et pósui vestiméntum meum cilícium: * et factus sum illis in parábolam.

Advérsum me loquebántur, qui sedébant in porta: * et in me psallébant qui bibébant vinum.

Ego vero oratiónem meam ad te, Dómine: * tempus benepláciti, Deus.

\antlabel{Ant.} Salvum me fac, Deus: quóniam intravérunt aquæ usque ad ánimam meam.

\antlabel{Ant.} Propter inimícos meos * éripe me, Dómine.

\psalmheading{Psalmus 68, ii}

\lettrine[lines=2]{I}{N} multitúdine misericórdiæ tuæ exáudi me, * in veritáte salútis tuæ:

Eripe me de luto, ut non infígar: * líbera me ab iis, qui odérunt me, et de profúndis aquárum.

Non me demérgat tempéstas aquæ, neque absórbeat me profúndum: * neque úrgeat super me púteus os suum.

Exáudi me, Dómine, quóniam benígna est misericórdia tua: * secúndum multitúdinem miseratiónum tuárum réspice in me.

Et ne avértas fáciem tuam a púero tuo: * quóniam tríbulor, velóciter exáudi me.

Inténde ánimæ meæ, et líbera eam: * propter inimícos meos éripe me.

Tu scis impropérium meum, et confusiónem meam, * et reveréntiam meam.

In conspéctu tuo sunt omnes qui tríbulant me, * impropérium exspectávit cor meum, et misériam.

Et sustínui qui simul contristarétur, et non fuit: * et qui consolarétur, et non invéni.

Et dedérunt in escam meam fel: * et in siti mea potavérunt me acéto.

Fiat mensa eórum coram ipsis in láqueum, * et in retributiónes, et in scándalum.

Obscuréntur óculi eórum ne vídeant: * et dorsum eórum semper incúrva.

Effúnde super eos iram tuam: * et furor iræ tuæ comprehéndat eos.

Fiat habitátio eórum desérta: * et in tabernáculis eórum non sit qui inhábitet.

Quóniam quem tu percussísti, persecúti sunt: * et super dolórem vúlnerum meórum addidérunt.

Appóne iniquitátem super iniquitátem eórum: * et non intrent in justítiam tuam.

Deleántur de libro vivéntium: * et cum justis non scribántur.

\antlabel{Ant.} Propter inimícos meos éripe me, Dómine.

\antlabel{Ant.} Quærite Dóminum, * et vivet ánima vestra.

\psalmheading{Psalmus 68, iii}

\lettrine[lines=2]{E}{GO} sum pauper et dolens: * salus tua, Deus, suscépit me.

Laudábo nomen Dei cum cántico: * et magnificábo eum in laude:

Et placébit Deo super vítulum novéllum: * córnua producéntem et úngulas.

Vídeant páuperes et læténtur: * quærite Deum, et vivet ánima vestra.

Quóniam exaudívit páuperes Dóminus: * et vinctos suos non despéxit.

Laudent illum cæli et terra, * mare, et ómnia reptília in eis.

Quóniam Deus salvam fáciet Sion: * et ædificabúntur civitátes Juda.

Et inhabitábunt ibi, * et hereditáte acquírent eam.

Et semen servórum ejus possidébit eam, * et qui díligunt nomen ejus, habitábunt in ea.

\antlabel{Ant.} Quærite Dóminum, et vivet ánima vestra.

\variantheading{In feriali Officio et in Festis:}

\V\ Laudábo nomen Dei cum cántico. \R\ Et magnificábo eum in laude.

\variantheading{In Officio novem Lectionum}

\absolutiolabel{Absolutio.} A vínculis.

\benedlabel{Bened.} 7 Evangélica léctio. 8 In Festis Dómini: Divínum auxílium. In Festis B. M. V.: Cujus festum cólimus, ipsa Virgo. In Festis Sanctórum: Cujus (\textit{vel} Quorum \textit{aut} Quarum) festum cólimus. 9 Ad societátem. (Per evangélica dicta.)

\variantheading{In Officio trium Lectionum}

\absolutiolabel{Absolutio.} Exáudi, Dómine.

\variantheading{In Feriis cum Homilia:}

\benedlabel{Bened.} 1 Evangélica. 2 Divínum. 3 Ad societátem.

\variantheading{In Feriis sine Homilia:}

\benedlabel{Bened.} 1 Benedictióne. 2 Unigénitus. 3 Spíritus Sancti.

\variantheading{In Festis Sanctorum:}

\benedlabel{Bened.} 1 Ille nos benedícat, qui sine fine vivit et regnat. \R\ Amen. 2 Cujus (\textit{vel} Quorum \textit{aut} Quarum). 3 Ad societátem.

% ===== END OF BATCH 24 -- Covers autumnalis.md lines 8133-8494,
% opening "Feria Quinta" (Thursday) and completing its "Ad
% Matutinum" in full: the ferial Invitatorium "Regem ventúrum
% Dóminum," the hymn "Nox atra rerum cóntegit," and the three
% Nocturns (Psalms 61, 65/i, 65/ii; 67/i-iii; 68/i-iii) with their
% versicle, Absolutio, and Benedictiones apparatus. "In III
% Nocturno" here has only a single form (no "I"/"II" alternation as
% in Feria Quarta), confirming that split was specific to the
% Psalm-49/Psalm-50 overlap and does not recur every day. Several
% more psalm incipits carry the partial-caps OCR defect already
% documented in prior batches ("NOnne," "JUbiláte," "INtroíbo,"
% "EXsurgat," "DOminus," "VIderunt," "SAlvum," "IN" multitúdine,
% "EGo") and are silently normalized to full caps. The document
% continues with "Ad Laudes I" (autumnalis.md line 8497).
% =====
\hourheading{Ad Laudes I}

\rubric{In Fériis extra Vigílias commúnes occurréntibus, atque in Festis.}

(Pater noster et Ave.)

\V\ Deus, in adjutórium.

\antlabel{Ant.} Jubiláte * in conspéctu regis Dómini.

\psalmheading{Psalmus 97}

\lettrine[lines=2]{C}{ANTATE} Dómino cánticum novum: * quia mirabília fecit.

Salvávit sibi déxtera ejus: * et bráchium sanctum ejus.

Notum fecit Dóminus salutáre suum: * in conspéctu géntium revelávit justítiam suam.

Recordátus est misericórdiæ suæ, * et veritátis suæ dómui Israël.

Vidérunt omnes términi terræ * salutáre Dei nostri.

Jubiláte Deo, omnis terra: * cantáte, et exsultáte, et psállite.

Psállite Dómino in cíthara, in cíthara et voce psalmi: * in tubis ductílibus, et voce tubæ córneæ.

Jubiláte in conspéctu regis Dómini: * moveátur mare, et plenitúdo ejus: orbis terrárum, et qui hábitant in eo.

Flúmina plaudent manu, simul montes exsultábunt a conspéctu Dómini: * quóniam venit judicáre terram.

Judicábit orbem terrárum in justítia, * et pópulos in æquitáte.

\antlabel{Ant.} Jubiláte in conspéctu regis Dómini.

\antlabel{Ant.} Dómine, * \textdagger\ refúgium factus es nobis. \textdaggerdbl

\psalmheading{Psalmus 89}

\lettrine[lines=2]{D}{OMINE}, \textdagger\ refúgium factus es nobis: * \textdaggerdbl\ a generatióne in generatiónem.

Priúsquam montes fíerent, aut formarétur terra et orbis: * a sæculo et usque in sæculum tu es, Deus.

Ne avértas hóminem in humilitátem: * et dixísti: Convertímini, fílii hóminum.

Quóniam mille anni ante óculos tuos, * tamquam dies hestérna, quæ prætériit.

Et custódia in nocte, * quæ pro níhilo habéntur, eórum anni erunt.

Mane sicut herba tránseat, mane flóreat, et tránseat: * véspere décidat, indúret et aréscat.

Quia defécimus in ira tua, * et in furóre tuo turbáti sumus.

Posuísti iniquitátes nostras in conspéctu tuo: * sæculum nostrum in illuminatióne vultus tui.

Quóniam omnes dies nostri defecérunt: * et in ira tua defécimus.

Anni nostri sicut aránea meditabúntur: * dies annórum nostrórum in ipsis, septuagínta anni.

Si autem in potentátibus octogínta anni: * et ámplius eórum, labor et dolor.

Quóniam supervénit mansuetúdo: * et corripiémur.

Quis novit potestátem iræ tuæ: * et præ timóre tuo iram tuam dinumeráre?

Déxteram tuam sic notam fac: * et erudítos corde in sapiéntia.

Convértere, Dómine, úsquequo? * et deprecábilis esto super servos tuos.

Repléti sumus mane misericórdia tua: * et exsultávimus, et delectáti sumus ómnibus diébus nostris.

Lætáti sumus pro diébus, quibus nos humiliásti: * annis, quibus vídimus mala.

Réspice in servos tuos, et in ópera tua: * et dírige fílios eórum.

Et sit splendor Dómini, Dei nostri, super nos, et ópera mánuum nostrárum dírige super nos: * et opus mánuum nostrárum dírige.

\antlabel{Ant.} Dómine, refúgium factus es nobis.

\antlabel{Ant.} Dómine, * in cælo misericórdia tua.

\psalmheading{Psalmus 35}

\lettrine[lines=2]{D}{IXIT} injústus ut delínquat in semetípso: * non est timor Dei ante óculos ejus.

Quóniam dolóse egit in conspéctu ejus: * ut inveniátur iníquitas ejus ad ódium.

Verba oris ejus iníquitas, et dolus: * nóluit intellígere ut bene ágeret.

Iniquitátem meditátus est in cubíli suo: * ástitit omni viæ non bonæ, malítiam autem non odívit.

Dómine, in cælo misericórdia tua: * et véritas tua usque ad nubes.

Justítia tua sicut montes Dei: * judícia tua abýssus multa.

Hómines, et juménta salvábis, Dómine: * quemádmodum multiplicásti misericórdiam tuam, Deus.

Fílii autem hóminum, * in tégmine alárum tuárum sperábunt.

Inebriabúntur ab ubertáte domus tuæ: * et torrénte voluptátis tuæ potábis eos.

Quóniam apud te est fons vitæ: * et in lúmine tuo vidébimus lumen.

Præténde misericórdiam tuam sciéntibus te, * et justítiam tuam his, qui recto sunt corde.

Non véniat mihi pes supérbiæ: * et manus peccatóris non móveat me.

Ibi cecidérunt qui operántur iniquitátem: * expúlsi sunt, nec potuérunt stare.

\antlabel{Ant.} Dómine, in cælo misericórdia tua.

\antlabel{Ant.} Pópulus meus, * ait Dóminus, bonis meis adimplébitur.

\canticumheading{Canticum Jeremiæ}{Jer. 31, 10-14}

\lettrine[lines=2]{A}{UDITE} verbum Dómini, gentes, * et annuntiáte in ínsulis, quæ procul sunt,

Et dícite: Qui dispérsit Israël, congregábit eum: * et custódiet eum sicut pastor gregem suum.

Redémit enim Dóminus Jacob, * et liberábit eum de manu potentióris.

Et vénient, et laudábunt in monte Sion: * et confluent ad bona Dómini.

Super fruménto, et vino, et óleo, * et fetu pécorum et armentórum:

Erítque ánima eórum quasi hortus irríguus, * et ultra non esúrient.

Tunc lætábitur virgo in choro, * júvenes et senes simul:

Et convértam luctum eórum in gáudium, * et consolábor eos, et lætificábo a dolóre suo.

Et inebriábo ánimam sacerdótum pinguédine: * et pópulus meus bonis meis adimplébitur.

\antlabel{Ant.} Pópulus meus, ait Dóminus, bonis meis adimplébitur.

\antlabel{Ant.} Deo nostro * jucúnda sit laudátio.

\psalmheading{Psalmus 146}

\lettrine[lines=2]{L}{AUDATE} Dóminum quóniam bonus est psalmus: * Deo nostro sit jucúnda, decóraque laudátio.

Ædíficans Jerúsalem Dóminus: * dispersiónes Israélis congregábit.

Qui sanat contrítos corde: * et álligat contritiónes eórum.

Qui númerat multitúdinem stellárum: * et ómnibus eis nómina vocat.

Magnus Dóminus noster, et magna virtus ejus: * et sapiéntiæ ejus non est númerus.

Suscípiens mansuétos Dóminus: * humílians autem peccatóres usque ad terram.

Præcínite Dómino in confessióne: * psállite Deo nostro in cíthara.

Qui óperit cælum núbibus: * et parat terræ plúviam.

Qui prodúcit in móntibus fœnum: * et herbam servitúti hóminum.

Qui dat juméntis escam ipsórum: * et pullis corvórum invocántibus eum.

Non in fortitúdine equi voluntátem habébit: * nec in tíbiis viri beneplácitum erit ei.

Beneplácitum est Dómino super timéntes eum: * et in eis, qui sperant super misericórdia ejus.

\antlabel{Ant.} Deo nostro jucúnda sit laudátio.

\variantheading{In Officio feriali}

\capitulumlabel{Capit.}{Rom. 13, 12-13}

\lettrine[lines=2]{N}{OX} præcéssit, dies autem appropinquávit. Abjiciámus ergo ópera tenebrárum, et induámur arma lucis. Sicut in die honéste ambulémus.

\hymnusheading{Hymnus}

\lettrine[lines=2]{L}{UX} ecce surgit áurea,\\
Pallens facéssat cæcitas,\\
Quæ nosmet in præceps diu\\
Erróre traxit dévio.

Hæc lux serénum cónferat,\\
Purósque nos præstet sibi:\\
Nihil loquámur súbdolum:\\
Volvámus obscúrum nihil.

Sic tota decúrrat dies,\\
Ne lingua mendax, ne manus\\
Oculíve peccent lúbrici,\\
Ne noxa corpus ínquinet.

Speculátor astat désuper,\\
Qui nos diébus ómnibus,\\
Actúsque nostros próspicit\\
A luce prima in vésperum.

Deo Patri sit glória,\\
Ejúsque soli Fílio,\\
Cum Spíritu Paráclito,\\
Nunc et per omne sæculum.\\
Amen.

\V\ Repléti sumus mane misericórdia tua. \R\ Exsultávimus, et delectáti sumus.

\antlabel{Ad Bened. Ant.} In sanctitáte * serviámus Dómino, et liberábit nos ab inimícis nostris.

\rubric{Réliqua ut in Ordinário. 9}

% ===== END OF BATCH 25 -- Covers autumnalis.md lines 8497-8702, "Ad
% Laudes I" for Feria Quinta in full: Psalms 97, 89, 35, the
% Canticum Jeremiæ, and Psalm 146, followed by the Capitulum "Nox
% præcéssit" and the hymn "Lux ecce surgit áurea." Several more
% psalm incipits carry the partial-caps OCR defect already
% documented in prior batches ("DOmine," "DIxit," "AUdíte,"
% "LAudáte," "NOx" [Capitulum]) and are silently normalized to full
% caps. The document continues with "Ad Laudes II" (autumnalis.md
% line 8705).
% =====
\hourheading{Ad Laudes II}

\rubric{In Vigíliis commúnibus, quando fit Offícium de Féria.}

(Pater noster et Ave.)

\V\ Deus, in adjutórium.

\antlabel{Ant.} Tibi soli peccávi.

\psalmheading{Psalmus 50}

\lettrine[lines=2]{M}{ISERERE} mei, Deus, * secúndum magnam misericórdiam tuam.

Et secúndum multitúdinem miseratiónum tuárum, * dele iniquitátem meam.

Amplius lava me ab iniquitáte mea: * et a peccáto meo munda me.

Quóniam iniquitátem meam ego cognósco: * et peccátum meum contra me est semper.

Tibi soli peccávi, et malum coram te feci: * ut justificéris in sermónibus tuis, et vincas cum judicáris.

Ecce enim, in iniquitátibus concéptus sum: * et in peccátis concépit me mater mea.

Ecce enim, veritátem dilexísti: * incérta et occúlta sapiéntiæ tuæ manifestásti mihi.

Aspérges me hyssópo, et mundábor: * lavábis me, et super nivem dealbábor.

Audítui meo dabis gáudium et lætítiam: * et exsultábunt ossa humiliáta.

Avérte fáciem tuam a peccátis meis: * et omnes iniquitátes meas dele.

Cor mundum crea in me, Deus: * et spíritum rectum ínnova in viscéribus meis.

Ne projícias me a fácie tua: * et spíritum sanctum tuum ne áuferas a me.

Redde mihi lætítiam salutáris tui: * et spíritu principáli confírma me.

Docébo iníquos vias tuas: * et ímpii ad te converténtur.

Líbera me de sanguínibus, Deus, Deus salútis meæ: * et exsultábit lingua mea justítiam tuam.

Dómine, lábia mea apéries: * et os meum annuntiábit laudem tuam.

Quóniam si voluísses sacrifícium, dedíssem útique: * holocáustis non delectáberis.

Sacrifícium Deo spíritus contribulátus: * cor contrítum, et humiliátum, Deus, non despícies.

Benígne fac, Dómine, in bona voluntáte tua Sion: * ut ædificéntur muri Jerúsalem.

Tunc acceptábis sacrifícium justítiæ, oblatiónes, et holocáusta: * tunc impónent super altáre tuum vítulos.

\antlabel{Ant.} Tibi soli peccávi, Dómine: miserére mei.

\antlabel{Ant.} Convértere, Dómine.

\psalmheading{Psalmus 89}

\lettrine[lines=2]{D}{OMINE}, refúgium factus es nobis: * a generatióne in generatiónem.

Priúsquam montes fíerent, aut formarétur terra et orbis: * a sæculo et usque in sæculum tu es, Deus.

Ne avértas hóminem in humilitátem: * et dixísti: Convertímini, fílii hóminum.

Quóniam mille anni ante óculos tuos, * tamquam dies hestérna, quæ prætériit.

Et custódia in nocte, * quæ pro níhilo habéntur, eórum anni erunt.

Mane sicut herba tránseat, mane flóreat, et tránseat: * véspere décidat, indúret et aréscat.

Quia defécimus in ira tua, * et in furóre tuo turbáti sumus.

Posuísti iniquitátes nostras in conspéctu tuo: * sæculum nostrum in illuminatióne vultus tui.

Quóniam omnes dies nostri defecérunt: * et in ira tua defécimus.

Anni nostri sicut aránea meditabúntur: * dies annórum nostrórum in ipsis, septuagínta anni.

Si autem in potentátibus, octogínta anni: * et ámplius eórum, labor et dolor.

Quóniam supervénit mansuetúdo: * et corripiémur.

Quis novit potestátem iræ tuæ: * et præ timóre tuo iram tuam dinumeráre?

Déxteram tuam sic notam fac: * et erudítos corde in sapiéntia.

Convértere, Dómine, úsquequo? * et deprecábilis esto super servos tuos.

Repléti sumus mane misericórdia tua: * et exsultávimus, et delectáti sumus ómnibus diébus nostris.

Lætáti sumus pro diébus, quibus nos humiliásti: * annis, quibus vídimus mala.

Réspice in servos tuos, et in ópera tua: * et dírige fílios eórum.

Et sit splendor Dómini, Dei nostri, super nos, et ópera mánuum nostrárum dírige super nos: * et opus mánuum nostrárum dírige.

\antlabel{Ant.} Convértere, Dómine, et deprecábilis esto super servos tuos.

\antlabel{Ant.} Multiplicásti, Deus.

\psalmheading{Psalmus 35}

\lettrine[lines=2]{D}{IXIT} injústus ut delínquat in semetípso: * non est timor Dei ante óculos ejus.

Quóniam dolóse egit in conspéctu ejus: * ut inveniátur iníquitas ejus ad ódium.

Verba oris ejus iníquitas, et dolus: * nóluit intellígere ut bene ágeret.

Iniquitátem meditátus est in cubíli suo: * ástitit omni viæ non bonæ, malítiam autem non odívit.

Dómine, in cælo misericórdia tua: * et véritas tua usque ad nubes.

Justítia tua sicut montes Dei: * judícia tua abýssus multa.

Hómines, et juménta salvábis, Dómine: * quemádmodum multiplicásti misericórdiam tuam, Deus.

Fílii autem hóminum, * in tégmine alárum tuárum sperábunt.

Inebriabúntur ab ubertáte domus tuæ: * et torrénte voluptátis tuæ potábis eos.

Quóniam apud te est fons vitæ: * et in lúmine tuo vidébimus lumen.

Præténde misericórdiam tuam sciéntibus te, * et justítiam tuam his, qui recto sunt corde.

Non véniat mihi pes supérbiæ: * et manus peccatóris non móveat me.

Ibi cecidérunt qui operántur iniquitátem: * expúlsi sunt, nec potuérunt stare.

\antlabel{Ant.} Multiplicásti, Deus, misericórdiam tuam.

\antlabel{Ant.} Fortitúdo mea.

\canticumheading{Canticum Moysis}{Exodi 15, 1-19}

\lettrine[lines=2]{C}{ANTEMUS} Dómino: glorióse enim magnificátus est, * equum et ascensórem dejécit in mare.

Fortitúdo mea, et laus mea Dóminus, * et factus est mihi in salútem:

Iste Deus meus, et glorificábo eum: * Deus patris mei, et exaltábo eum.

Dóminus quasi vir pugnátor, Omnípotens nomen ejus. * Currus Pharaónis et exércitum ejus projécit in mare.

Elécti príncipes ejus submérsi sunt in Mari Rubro: * abýssi operuérunt eos, descendérunt in profúndum quasi lapis.

Déxtera tua, Dómine, magnificáta est in fortitúdine: déxtera tua, Dómine, percússit inimícum. * Et in multitúdine glóriæ tuæ deposuísti adversários tuos:

Misísti iram tuam, quæ devorávit eos sicut stípulam. * Et in spíritu furóris tui congregátæ sunt aquæ:

Stetit unda fluens, * congregátæ sunt abýssi in médio mari.

Dixit inimícus: Pérsequar et comprehéndam, * dívidam spólia, implébitur ánima mea:

Evaginábo gládium meum, * interfíciet eos manus mea.

Flavit spíritus tuus, et opéruit eos mare: * submérsi sunt quasi plumbum in aquis vehemméntibus.

Quis símilis tui in fórtibus, Dómine? * quis símilis tui, magníficus in sanctitáte, terríbilis atque laudábilis, fáciens mirabília?

Extendísti manum tuam, et devorávit eos terra. * Dux fuísti in misericórdia tua pópulo quem redemísti:

Et portásti eum in fortitúdine tua, * ad habitáculum sanctum tuum.

Ascendérunt pópuli, et iráti sunt: * dolóres obtinuérunt habitatóres Philísthiim.

Tunc conturbáti sunt príncipes Edom, robústos Moab obtínuit tremor: * obriguérunt omnes habitatóres Chánaan.

Irruat super eos formído et pavor, * in magnitúdine bráchii tui.

Fiant immóbiles quasi lapis, donec pertránseat pópulus tuus, Dómine, * donec pertránseat pópulus tuus iste, quem possedísti.

Introdúces eos, et plantábis in monte hereditátis tuæ, * firmíssimo habitáculo tuo quod operátus es, Dómine:

Sanctuárium tuum, Dómine, quod firmavérunt manus tuæ. * Dóminus regnábit in ætérnum et ultra.

Ingréssus est enim eques Phárao cum cúrribus et equítibus ejus in mare: * et redúxit super eos Dóminus aquas maris:

Fílii autem Israël ambulavérunt per siccum * in médio ejus.

\antlabel{Ant.} Fortitúdo mea, et laus mea Dóminus: et factus est mihi in salútem.

\antlabel{Ant.} Laudáte. \textdagger

\psalmheading{Psalmus 146}

\lettrine[lines=2]{L}{AUDATE} \textdagger\ Dóminum quóniam bonus est psalmus: * Deo nostro sit jucúnda, decóraque laudátio.

Ædíficans Jerúsalem Dóminus: * dispersiónes Israélis congregábit.

Qui sanat contrítos corde: * et álligat contritiónes eórum.

Qui númerat multitúdinem stellárum: * et ómnibus eis nómina vocat.

Magnus Dóminus noster, et magna virtus ejus: * et sapiéntiæ ejus non est númerus.

Suscípiens mansuétos Dóminus: * humílians autem peccatóres usque ad terram.

Præcínite Dómino in confessióne: * psállite Deo nostro in cíthara.

Qui óperit cælum núbibus: * et parat terræ plúviam.

Qui prodúcit in móntibus fœnum: * et herbam servitúti hóminum.

Qui dat juméntis escam ipsórum: * et pullis corvórum invocántibus eum.

Non in fortitúdine equi voluntátem habébit: * nec in tíbiis viri beneplácitum erit ei.

Beneplácitum est Dómino super timéntes eum: * et in eis, qui sperant super misericórdia ejus.

\antlabel{Ant.} Laudáte Dóminum, qui sanat contrítos corde, et álligat contritiónes eórum.

\capitulumlabel{Capit.}{Rom. 13, 12-13}

\lettrine[lines=2]{N}{OX} præcéssit, dies autem appropinquávit. Abjiciámus ergo ópera tenebrárum, et induámur arma lucis. Sicut in die honéste ambulémus.

\hymnusheading{Hymnus}

\lettrine[lines=2]{L}{UX} ecce surgit áurea,\\
Pallens facéssat cæcitas,\\
Quæ nosmet in præceps diu\\
Erróre traxit dévio.

Hæc lux serénum cónferat,\\
Purósque nos præstet sibi:\\
Nihil loquámur súbdolum:\\
Volvámus obscúrum nihil.

Sic tota decúrrat dies,\\
Ne lingua mendax, ne manus\\
Oculíve peccent lúbrici,\\
Ne noxa corpus ínquinet,

Speculátor astat désuper,\\
Qui nos diébus ómnibus,\\
Actúsque nostros próspicit\\
A luce prima in vésperum.

Deo Patri sit glória,\\
Ejúsque soli Fílio,\\
Cum Spíritu Paráclito,\\
Nunc et per omne sæculum.\\
Amen.

\V\ Repléti sumus mane misericórdia tua. \R\ Exsultávimus, et delectáti sumus.

\antlabel{Ad Bened. Ant.} In sanctitáte * serviámus Dómino, et liberábit nos ab inimícis nostris.

\rubric{Expleto Cántico ac repetíta Antíphona, dicúntur sequéntes}

\variantheading{Preces feriales}

KYRIE, eléison. Christe, eléison. Kýrie, eléison.

Pater noster \rubric{quod dícitur a solo Hebdomadário totum clara voce, usque ad} \V\ Et ne nos indúcas in tentatiónem. \R\ Sed líbera nos a malo. \V\ Ego dixi: Dómine, miserére mei. \R\ Sana ánimam meam, quia peccávi tibi. \V\ Convértere, Dómine, úsquequo? \R\ Et deprecábilis esto super servos tuos. \V\ Fiat misericórdia tua, Dómine, super nos. \R\ Quemádmodum sperávimus in te. \V\ Sacerdótes tui induántur justítiam. \R\ Et sancti tui exsúltent. \V\ Orémus pro beatíssimo Papa nostro N. \R\ Dóminus consérvet eum, et vivíficet eum, et beátum fáciat eum in terra, et non tradat eum in ánimam inimicórum ejus. \V\ Orémus et pro Antístite nostro N. \R\ Stet et pascat in fortitúdine tua, Dómine, in sublimitáte nóminis tui. \V\ Dómine, salvum fac regem. \R\ Et exáudi nos in die, qua invocavérimus te. \V\ Salvum fac pópulum tuum, Dómine, et bénedic hereditáti tuæ. \R\ Et rege eos, et extólle illos usque in ætérnum. \V\ Meménto Congregatiónis tuæ. \R\ Quam possedísti ab inítio. \V\ Fiat pax in virtúte tua. \R\ Et abundántia in túrribus tuis. \V\ Orémus pro benefactóribus nostris. \R\ Retribúere dignáre, Dómine, ómnibus, nobis bona faciéntibus propter nomen tuum, vitam ætérnam. Amen. \V\ Orémus pro fidélibus defúnctis. \R\ Réquiem ætérnam dona eis, Dómine, et lux perpétua lúceat eis. \V\ Requiéscant in pace. \R\ Amen. \V\ Pro frátribus nostris abséntibus. \R\ Salvos fac servos tuos, Deus meus, sperántes in te. \V\ Pro afflíctis et captívis. \R\ Líbera eos, Deus Israël, ex ómnibus tribulatiónibus suis. \V\ Mitte eis, Dómine, auxílium de sancto. \R\ Et de Sion tuére eos. \V\ Dómine, Deus virtútum, convérte nos. \R\ Et osténde fáciem tuam, et salvi érimus. \V\ Exsúrge, Christe, ádjuva nos. \R\ Et líbera nos propter nomen tuum. \V\ Dómine, exáudi oratiónem meam. \R\ Et clamor meus ad te véniat.

\rubric{Réliqua ut in Ordinário. 9}

% ===== END OF BATCH 26 -- Covers autumnalis.md lines 8705-8964, "Ad
% Laudes II" for Feria Quinta in full: Psalms 50, 89, 35, the
% Canticum Moysis ("Cantémus Dómino"), and Psalm 146, sharing the
% Capitulum "Nox præcéssit" and the hymn "Lux ecce surgit áurea"
% with Ad Laudes I, plus the "Preces feriales" litany. Psalm 89
% appears here without the dagger/double-dagger chant-pointing
% marks present in its Ad Laudes I occurrence (Batch 25), confirmed
% as a genuine source divergence rather than an omission. The
% document continues with "Ad Primam" (autumnalis.md line 8967).
% =====
\hourheading{Ad Primam}

Pater, Ave et Credo.

\V\ Deus, in adjutórium.

\hymnusheading{Hymnus}

\lettrine[lines=2]{J}{AM} lucis orto sídere,\\
Deum precémur súpplices,\\
Ut in diúrnis áctibus\\
Nos servet a nocéntibus.

Linguam refrénans témperet,\\
Ne litis horror ínsonet:\\
Visum fovéndo cóntegat,\\
Ne vanitátes háuriat.

Sint pura cordis íntima,\\
Absístat et vecórdia:\\
Carnis terat supérbiam\\
Potus cibíque párcitas.

Ut, cum dies abscésserit\\
Noctémque sors redúxerit,\\
Mundi per abstinéntiam\\
Ipsi canámus glóriam.

Deo Patri sit glória,\\
Ejúsque soli Fílio,\\
Cum Spíritu Paráclito,\\
Nunc et per omne sæculum.\\
Amen.

\antlabel{Ant.} In loco páscuæ.

\psalmheading{Psalmus 22}

\lettrine[lines=2]{D}{OMINUS} regit me, et nihil mihi déerit: * in loco páscuæ ibi me collocávit.

Super aquam refectiónis educávit me: * ánimam meam convértit.

Dedúxit me super sémitas justítiæ, * propter nomen suum.

Nam, et si ambulávero in médio umbræ mortis, non timébo mala: * quóniam tu mecum es.

Virga tua, et báculus tuus: * ipsa me consoláta sunt.

Parásti in conspéctu meo mensam, * advérsus eos, qui tríbulant me.

Impinguásti in óleo caput meum: * et calix meus inébrians quam præclárus est!

Et misericórdia tua subsequétur me: * ómnibus diébus vitæ meæ:

Et ut inhábitem in domo Dómini, * in longitúdinem diérum.

\psalmheading{Psalmus 71, i}

\lettrine[lines=2]{D}{EUS}, judícium tuum regi da: * et justítiam tuam fílio regis:

Judicáre pópulum tuum in justítia, * et páuperes tuos in judício.

Suscípiant montes pacem pópulo: * et colles justítiam.

Judicábit páuperes pópuli, et salvos fáciet fílios páuperum: * et humiliábit calumniatórem.

Et permanébit cum sole, et ante lunam, * in generatióne et generatiónem.

Descéndet sicut plúvia in vellus: * et sicut stillicídia stillántia super terram.

Oriétur in diébus ejus justítia, et abundántia pacis: * donec auferátur luna.

Et dominábitur a mari usque ad mare: * et a flúmine usque ad términos orbis terrárum.

\psalmheading{Psalmus 71, ii}

\lettrine[lines=2]{C}{ORAM} illo prócident Æthíopes: * et inimíci ejus terram lingent.

Reges Tharsis, et ínsulæ múnera ófferent: * reges Arabum et Saba dona addúcent.

Et adorábunt eum omnes reges terræ: * omnes gentes sérvient ei:

Quia liberábit páuperem a poténte: * et páuperem, cui non erat adjútor.

Parcet páuperi et ínopi: * et ánimas páuperum salvas fáciet.

Ex usúris et iniquitáte rédimet ánimas eórum: * et honorábile nomen eórum coram illo.

Et vivet, et dábitur ei de auro Arábiæ, et adorábunt de ipso semper: * tota die benedícent ei.

Et erit firmaméntum in terra in summis móntium, superextollétur super Líbanum fructus ejus: * et florébunt de civitáte sicut fœnum terræ.

Sit nomen ejus benedíctum in sæcula: * ante solem pérmanet nomen ejus.

Et benedicéntur in ipso omnes tribus terræ: * omnes gentes magnificábunt eum.

Benedíctus Dóminus, Deus Israël, * qui facit mirabília solus:

Et benedíctum nomen majestátis ejus in ætérnum: * et replébitur majestáte ejus omnis terra: fiat, fiat.

\rubric{¶ Quando ad Laudes dictus fúerit Ps. 50 Miserére, hic subjúngitur sequens Ps. 97, qui in Laudibus locum non hábuit.}

\psalmheading{Psalmus 97}

\lettrine[lines=2]{C}{ANTATE} Dómino cánticum novum: * quia mirabília fecit.

Salvávit sibi déxtera ejus: * et bráchium sanctum ejus.

Notum fecit Dóminus salutáre suum: * in conspéctu géntium revelávit justítiam suam.

Recordátus est misericórdiæ suæ, * et veritátis suæ dómui Israël.

Vidérunt omnes términi terræ * salutáre Dei nostri.

Jubiláte Deo, omnis terra: * cantáte, et exsultáte, et psállite.

Psállite Dómino in cíthara, in cíthara et voce psalmi: * in tubis ductílibus, et voce tubæ córneæ.

Jubiláte in conspéctu regis Dómini: * moveátur mare, et plenitúdo ejus: orbis terrárum, et qui hábitant in eo.

Flúmina plaudent manu, simul montes exsultábunt a conspéctu Dómini: * quóniam venit judicáre terram.

Judicábit orbem terrárum in justítia, * et pópulos in æquitáte.

\rubric{Secus vero, tribus Psalmis explétis, statim dícitur Ant.}

\antlabel{Ant.} In loco páscuæ, ibi Dóminus me collocávit.

\rubric{In Officio Festi vel Octávæ:}

\capitulumlabel{Capitulum}{1 Tim. 1, 17}

\lettrine[lines=2]{R}{EGI} sæculórum immortáli et invisíbili, soli Deo honor et glória in sæcula sæculórum. Amen.

\rubric{In ómnibus autem Fériis et in Vigíliis commúnibus:}

\capitulumlabel{Capitulum}{Zach. 8, 19}

\lettrine[lines=2]{P}{ACEM} et veritátem dilígite, ait Dóminus omnípotens.

\R\ br. Christe, Fili Dei vivi, * Miserére nobis. Christe. \V\ \textit{Qui sedes ad déxteram Patris.} Miserére nobis. Glória Patri. Christe. \V\ Exsúrge, Christe, ádjuva nos. \R\ Et líbera nos propter nomen tuum.

\variantheading{Preces dominicales et feriales (si dicendæ sint)}

KYRIE, eléison. Christe, eléison. Kýrie, eléison.

Pater noster \rubric{secreto usque ad} \V\ Et ne nos indúcas in tentatiónem. \R\ Sed líbera nos a malo.

Credo in Deum \rubric{secreto usque ad} \V\ Carnis resurrectiónem. \R\ Vitam ætérnam. Amen. \V\ Et ego ad te, Dómine, clamávi. \R\ Et mane orátio mea prævéniet te. \V\ Repleátur os meum laude. \R\ Ut cantem glóriam tuam, tota die magnitúdinem tuam. \V\ Dómine, avérte fáciem tuam a peccátis meis. \R\ Et omnes iniquitátes meas dele. \V\ Cor mundum crea in me, Deus. \R\ Et spíritum rectum ínnova in viscéribus meis. \V\ Ne projícias me a fácie tua. \R\ Et spíritum sanctum tuum ne áuferas a me. \V\ Redde mihi lætítiam salutáris tui. \R\ Et spíritu principáli confírma me.

\variantheading{In Précibus feriálibus tantum}

\V\ Eripe me, Dómine, ab hómine malo. \R\ A viro iníquo éripe me. \V\ Eripe me de inimícis meis, Deus meus. \R\ Et ab insurgéntibus in me líbera me. \V\ Eripe me de operántibus iniquitátem. \R\ Et de viris sánguinum salva me. \V\ Sic psalmum dicam nómini tuo in sæculum sæculi. \R\ Ut reddam vota mea de die in diem. \V\ Exáudi nos, Deus, salutáris noster. \R\ Spes ómnium fínium terræ, et in mari longe. \V\ Deus, in adjutórium meum inténde. \R\ Dómine, ad adjuvándum me festína. \V\ Sanctus Deus, Sanctus fortis, Sanctus immortális. \R\ Miserére nobis. \V\ Bénedic, ánima mea, Dómino. \R\ Et ómnia, quæ intra me sunt, nómini sancto ejus. \V\ Bénedic, ánima mea, Dómino. \R\ Et noli oblivísci omnes retributiónes ejus. \V\ Qui propitiátur ómnibus iniquitátibus tuis. \R\ Qui sanat omnes infirmitátes tuas. \V\ Qui rédimit de intéritu vitam tuam. \R\ Qui corónat te in misericórdia et miseratiónibus. \V\ Qui replet in bonis desidérium tuum. \R\ Renovábitur ut áquilæ juvéntus tua.

\variantheading{In utrisque Précibus}

\V\ Adjutórium nostrum in nómine Dómini. \R\ Qui fecit cælum et terram.

\formulalabel{Confíteor. Misereátur. Indulgéntiam.}

\V\ Dignáre, Dómine, die isto. \R\ Sine peccáto nos custodíre. \V\ Miserére nostri, Dómine. \R\ Miserére nostri. \V\ Fiat misericórdia tua, Dómine, super nos. \R\ Quemádmodum sperávimus in te. \V\ Dómine, exáudi oratiónem meam. \R\ Et clamor meus ad te véniat.

\V\ Dóminus vobíscum.

Orémus. \oratiolabel{Oratio}

\lettrine[lines=2]{D}{OMINE}, Deus omnípotens, qui ad princípium hujus diéi nos perveníre fecísti: tua nos hódie salva virtúte; ut in hac die ad nullum declinémus peccátum, sed semper ad tuam justítiam faciéndam nostra procédant elóquia, dirigántur cogitatiónes et ópera. Per Dóminum. \R\ Amen.

\V\ Dóminus vobíscum. \V\ Benedicámus Dómino.

\rubric{In Choro Martyrológium.}

\V\ Pretiósa in conspéctu Dómini. \R\ Mors Sanctórum ejus.

\rubric{Et absolúte, sine} Orémus:

\oratiolabel{Oratio}

\lettrine[lines=2]{S}{ANCTA} María, et omnes Sancti intercédant pro nobis ad Dóminum, ut nos mereámur ab eo adjuvári et salvári, qui vivit et regnat in sæcula sæculórum. \R\ Amen.

\V\ Deus, in adjutórium meum inténde. \R\ Dómine, ad adjuvándum me festína. \rubric{Et dícitur ter; último additur:} Glória Patri. Sicut erat.

Kýrie, eléison. Christe, eléison. Kýrie, eléison.

Pater noster \rubric{secreto usque ad} \V\ Et ne nos indúcas in tentatiónem. \R\ Sed líbera nos a malo. \V\ Réspice in servos tuos, Dómine, et in ópera tua, et dírige fílios eórum. \R\ Et sit splendor Dómini, Dei nostri, super nos, et ópera mánuum nostrárum dírige super nos, et opus mánuum nostrárum dírige. \V\ Glória Patri. \R\ Sicut erat.

Orémus. \oratiolabel{Oratio}

\lettrine[lines=2]{D}{IRIGERE} et sanctificáre, régere et gubernáre dignáre, Dómine Deus, Rex cæli et terræ, hódie corda et córpora nostra, sensus, sermónes et actus nostros in lege tua, et in opéribus mandatórum tuórum: ut hic, et in ætérnum, te auxiliánte, salvi et líberi esse mereámur, Salvátor mundi: Qui vivis et regnas in sæcula sæculórum. \R\ Amen.

\V\ Jube, domne.

\benedlabel{Bened.} Dies et actus nostros in sua pace dispónat Dóminus omnípotens. \R\ Amen.

\variantheading{In Officio feriali}

\lettrine[lines=2]{D}{OMINUS} autem dírigat corda et córpora nostra in caritáte Dei, et patiéntia Christi. Tu autem.

\variantheading{E Proprio Sanctorum}

\variantheading{Pro Oct. Omn. Sanctorum}

\lettrine[lines=2]{B}{ENEDÍCTIO} et cláritas et sapiéntia et gratiárum áctio, honor, virtus et fortitúdo Deo nostro in sǽcula sæculórum. Amen.

\variantheading{E Communi Sanctorum}

\variantheading{Pro Apostolis}

\lettrine[lines=2]{I}{BANT} Apóstoli gaudéntes a conspéctu concílii, quóniam digni hábiti sunt pro nómine Jesu contuméliam pati.

\variantheading{Pro Martyre}

\lettrine[lines=2]{J}{USTUS} cor suum trádidit ad vigilándum dilúculo ad Dóminum, qui fecit illum, et in conspéctu Altíssimi deprecábitur.

\variantheading{Pro pluribus Martyribus}

\lettrine[lines=2]{F}{ULGEBUNT} justi, et tamquam scintíllæ in arundinéto discúrrent. Judicábunt natiónes, et dominabúntur pópulis: et regnábit Dóminus illórum in perpétuum.

\variantheading{Pro Conf. Pont.}

\lettrine[lines=2]{F}{UNGI} sacerdótio, et habére laudem in nómine ipsíus, et offérre illi incénsum dignum in odórem suavitátis.

\variantheading{Pro Conf. non Pont.}

\lettrine[lines=2]{J}{USTUM} dedúxit Dóminus per vias rectas, et osténdit illi regnum Dei, et dedit illi sciéntiam sanctórum; honestávit illum in labóribus, et complévit labóres illíus.

\variantheading{Pro Virgine Martyre}

\lettrine[lines=2]{D}{OMINE}, Deus meus, exaltásti super terram habitatiónem meam, et pro morte defluénte deprecáta sum. Invocávi Dóminum, Patrem Dómini mei, ut non derelínquat me in die tribulatiónis meæ, et in témpore superbórum sine adjutório.

\variantheading{Pro Virgine non Martyre}

\lettrine[lines=2]{O}{} \textsc{quam} pulchra est casta generátio cum claritáte! immortális est enim memória illíus: quóniam apud Deum nota est, et apud hómines.

\variantheading{Pro Martyre tantum}

\lettrine[lines=2]{L}{AUDABIT} usque ad mortem ánima mea Dóminum, quóniam éruis sustinéntes te, et líberas eos de manu angústiæ, Dómine, Deus noster.

\variantheading{Pro nec Virg. nec Mart.}

\lettrine[lines=2]{M}{ULTAE} fíliæ congregavérunt divítias: tu supergréssa es univérsas. Fallax grátia, et vana est pulchritúdo: múlier timens Dóminum, ipsa laudábitur.

\variantheading{Pro Dedicatione Ecclesiæ}

\lettrine[lines=2]{E}{T} abstérget Deus omnem lácrimam ab óculis eórum: et mors ultra non erit, neque luctus, neque clamor, neque dolor erit ultra, quia prima abiérunt. Et dixit qui sedébat in throno: Ecce, nova fácio ómnia.

\variantheading{In Officiis B. Mariæ V.}

\lettrine[lines=2]{I}{N} platéis sicut cinnamómum et bálsamum aromatízans odórem dedi: quasi myrrha elécta dedi suavitátem odóris.

\V\ Adjutórium nostrum in nómine Dómini. \R\ Qui fecit cælum et terram.

\V\ Benedícite. \R\ Deus.

\benedictiolabel{Benedictio} Dóminus nos benedícat, et ab omni malo deféndat, et ad vitam perdúcat ætérnam. Et fidélium ánimæ per misericórdiam Dei requiéscant in pace. \R\ Amen.

\rubric{Deinde dícitur tantum} Pater noster \rubric{secréto, nisi sequátur ália Hora.}

% ===== END OF BATCH 27 -- Covers autumnalis.md lines 8967-9226,
% completing "Ad Primam" for Feria Quinta in full: the hymn "Jam
% lucis orto sídere," Psalms 22, 71/i, 71/ii, the conditional Psalm
% 97 insertion, both Capitula, the brief Responsory, the full
% Preces apparatus, the closing Orationes, and the Lectio-brevis
% apparatus (identical category list to the previous three ferias).
% Several more psalm incipits carry the partial-caps OCR defect
% ("DOminus," "DEus," "COram," "CAntáte," "PAcem" [Capitulum]) and
% are silently normalized to full caps. The document continues
% with "Ad Tertiam" (autumnalis.md line 9229).
% =====
\hourheading{Ad Tertiam}

Pater noster et Ave.

\V\ Deus, in adjutórium.

\hymnusheading{Hymnus}

\lettrine[lines=2]{N}{UNC}, Sancte, nobis, Spíritus,\\
Unum Patri cum Fílio,\\
Dignáre promptus íngeri\\
Nostro refúsus péctori.

Os, lingua, mens, sensus, vigor\\
Confessiónem pérsonent,\\
Flamméscat igne cáritas,\\
Accéndat ardor próximos.

Præsta, Pater piíssime,\\
Patríque compar Unice,\\
Cum Spíritu Paráclito\\
Regnans per omne sæculum.\\
Amen.

\antlabel{Ant.} Quam bonus. \textdagger

\psalmheading{Psalmus 72, i}

\lettrine[lines=2]{Q}{UAM} bonus \textdagger\ Israël Deus, * his qui recto sunt corde!

Mei autem pæne moti sunt pedes: * pæne effúsi sunt gressus mei.

Quia zelávi super iníquos, * pacem peccatórum videns.

Quia non est respéctus morti eórum: * et firmaméntum in plaga eórum.

In labóre hóminum non sunt, * et cum homínibus non flagellabúntur:

Ideo ténuit eos supérbia, * opérti sunt iniquitáte et impietáte sua.

Pródiit quasi ex ádipe iníquitas eórum: * transiérunt in afféctum cordis.

Cogitavérunt, et locúti sunt nequítiam: * iniquitátem in excélso locúti sunt.

Posuérunt in cælum os suum: * et lingua eórum transívit in terra.

\psalmheading{Psalmus 72, ii}

\lettrine[lines=2]{I}{DEO} convertétur pópulus meus hic: * et dies pleni inveniéntur in eis.

Et dixérunt: Quómodo scit Deus, * et si est sciéntia in excélso?

Ecce, ipsi peccatóres, et abundántes in sæculo, * obtinuérunt divítias.

Et dixi: Ergo sine causa justificávi cor meum, * et lavi inter innocéntes manus meas:

Et fui flagellátus tota die, * et castigátio mea in matutínis.

Si dicébam: Narrábo sic: * ecce, natiónem filiórum tuórum reprobávi.

Existimábam ut cognóscerem hoc, * labor est ante me:

Donec intrem in Sanctuárium Dei: * et intélligam in novíssimis eórum.

\psalmheading{Psalmus 72, iii}

\lettrine[lines=2]{V}{ERUMTAMEN} propter dolos posuísti eis: * dejecísti eos dum allevaréntur.

Quómodo facti sunt in desolatiónem, súbito defecérunt: * periérunt propter iniquitátem suam.

Velut sómnium surgéntium, Dómine, * in civitáte tua imáginem ipsórum ad níhilum rédiges.

Quia inflammátum est cor meum, et renes mei commutáti sunt: * et ego ad níhilum redáctus sum, et nescívi.

Ut juméntum factus sum apud te: * et ego semper tecum.

Tenuísti manum déxteram meam: et in voluntáte tua deduxísti me, * et cum glória suscepísti me.

Quid enim mihi est in cælo? * et a te quid vólui super terram?

Defécit caro mea, et cor meum: * Deus cordis mei, et pars mea Deus in ætérnum.

Quia ecce, qui elóngant se a te, períbunt: * perdidísti omnes, qui fornicántur abs te.

Mihi autem adhærére Deo bonum est: * pónere in Dómino Deo spem meam:

Ut annúntiem omnes prædicatiónes tuas, * in portis fíliæ Sion.

\antlabel{Ant.} Quam bonus Israël Deus, his, qui recto sunt corde?

\variantheading{In Officio feriali}

\capitulumlabel{Capitulum}{Jer. 17, 14}

\lettrine[lines=2]{S}{ANA} me, Dómine, et sanábor: salvum me fac, et salvus ero: quóniam laus mea tu es.

\R\ br. Sana ánimam meam, * Quia peccávi tibi. Sana. \V\ Ego dixi: Dómine, miserére mei. Quia. Glória Patri. Sana.

\V\ Adjútor meus esto, ne derelínquas me. \R\ Neque despícias me, Deus, salutáris meus.

\variantheading{In Vigiliis communibus:}

\variantheading{Preces feriales}

KYRIE, eléison. Christe, eléison. Kýrie, eléison.

Pater noster \rubric{secreto usque ad} \V\ Et ne nos indúcas in tentatiónem. \R\ Sed líbera nos a malo. \V\ Dómine, Deus virtútum, convérte nos. \R\ Et osténde fáciem tuam, et salvi érimus. \V\ Exsúrge, Christe, ádjuva nos. \R\ Et líbera nos propter nomen tuum. \V\ Dómine, exáudi oratiónem meam. \R\ Et clamor meus ad te véniat.

\rubric{Réliqua ut in Ordinário. 19}

\hourheading{Ad Sextam}

Pater noster et Ave.

\V\ Deus, in adjutórium.

\hymnusheading{Hymnus}

\lettrine[lines=2]{R}{ECTOR} potens, verax Deus,\\
Qui témperas rerum vices,\\
Splendóre mane illúminas,\\
Et ígnibus merídiem:

Exstíngue flammas lítium,\\
Aufer calórem nóxium,\\
Confer salútem córporum,\\
Verámque pacem córdium.

Præsta, Pater piíssime,\\
Patríque compar Unice,\\
Cum Spíritu Paráclito\\
Regnans per omne sæculum.\\
Amen.

\antlabel{Ant.} Memor esto.

\psalmheading{Psalmus 73, i}

\lettrine[lines=2]{U}{T} quid, Deus, repulísti in finem: * irátus est furor tuus super oves páscuæ tuæ?

Memor esto congregatiónis tuæ, * quam possedísti ab inítio.

Redemísti virgam hereditátis tuæ: * mons Sion, in quo habitásti in eo.

Leva manus tuas in supérbias eórum in finem: * quanta malignátus est inimícus in sancto!

Et gloriáti sunt qui odérunt te: * in médio solemnitátis tuæ.

Posuérunt signa sua, signa: * et non cognovérunt sicut in éxitu super summum.

Quasi in silva lignórum secúribus excidérunt jánuas ejus in idípsum: * in secúri et áscia dejecérunt eam.

Incendérunt igni Sanctuárium tuum: * in terra polluérunt tabernáculum nóminis tui.

Dixérunt in corde suo cognátio eórum simul: * Quiéscere faciámus omnes dies festos Dei a terra.

Signa nostra non vídimus, jam non est prophéta: * et nos non cognóscet ámplius.

\psalmheading{Psalmus 73, ii}

\lettrine[lines=2]{U}{SQUEQUO}, Deus, improperábit inimícus: * irrítat adversárius nomen tuum in finem?

Ut quid avértis manum tuam, et déxteram tuam, * de médio sinu tuo in finem?

Deus autem Rex noster ante sæcula: * operátus est salútem in médio terræ.

Tu confirmásti in virtúte tua mare: * contribulásti cápita dracónum in aquis.

Tu confregísti cápita dracónis: * dedísti eum escam pópulis Æthíopum.

Tu dirupísti fontes, et torréntes: * tu siccásti flúvios Ethan.

Tuus est dies, et tua est nox: * tu fabricátus es auróram et solem.

Tu fecísti omnes términos terræ: * æstátem et ver tu plasmásti ea.

\psalmheading{Psalmus 73, iii}

\lettrine[lines=2]{M}{EMOR} esto hujus, inimícus improperávit Dómino: * et pópulus insípiens incitávit nomen tuum.

Ne tradas béstiis ánimas confiténtes tibi, * et ánimas páuperum tuórum ne oblivíscáris in finem.

Réspice in testaméntum tuum: * quia repléti sunt, qui obscuráti sunt terræ dómibus iniquitátum.

Ne avertátur húmilis factus confúsus: * pauper et inops laudábunt nomen tuum.

Exsúrge, Deus, júdica causam tuam: * memor esto improperiórum tuórum, eórum quæ ab insipiénte sunt tota die.

Ne obliviscáris voces inimicórum tuórum: * supérbia eórum, qui te odérunt, ascéndit semper.

\antlabel{Ant.} Memor esto congregatiónis tuæ, Dómine, quam possedísti ab inítio.

\variantheading{In Officio feriali}

\capitulumlabel{Capitulum}{Rom. 13, 8}

\lettrine[lines=2]{N}{EMINI} quidquam debeátis, nisi ut ínvicem diligátis: qui enim díligit próximum, legem implévit.

\R\ br. Benedícam Dóminum * In omni témpore. Benedícam. \V\ Semper laus ejus in ore meo. In. Glória Patri. Benedícam. \V\ Dóminus regit me, et nihil mihi déerit. \R\ In loco páscuæ ibi me collocávit.

\variantheading{In Vigiliis communibus:}

\variantheading{Preces feriales}

KYRIE, eléison. Christe, eléison. Kýrie, eléison.

Pater noster \rubric{secreto usque ad} \V\ Et ne nos indúcas in tentatiónem. \R\ Sed líbera nos a malo. \V\ Dómine, Deus virtútum, convérte nos. \R\ Et osténde fáciem tuam, et salvi érimus. \V\ Exsúrge, Christe, ádjuva nos. \R\ Et líbera nos propter nomen tuum. \V\ Dómine, exáudi oratiónem meam. \R\ Et clamor meus ad te véniat.

\rubric{Réliqua ut in Ordinário. 20}

\hourheading{Ad Nonam}

Pater noster et Ave.

\V\ Deus, in adjutórium.

\hymnusheading{Hymnus}

\lettrine[lines=2]{R}{ERUM}, Deus, tenax vigor,\\
Immótus in te pérmanens,\\
Lucis diúrnæ témpora\\
Succéssibus detérminans:

Largíre lumen véspere,\\
Quo vita nusquam décidat,\\
Sed prǽmium mortis sacræ\\
Perénnis instet glória.

Præsta, Pater piíssime,\\
Patríque compar Unice,\\
Cum Spíritu Paráclito\\
Regnans per omne sæculum.\\
Amen.

\antlabel{Ant.} Invocábimus.

\psalmheading{Psalmus 74}

\lettrine[lines=2]{C}{ONFITEBIMUR} tibi, Deus: * confitébimur, et invocábimus nomen tuum.

Narrábimus mirabília tua: * cum accépero tempus, ego justítias judicábo.

Liquefácta est terra, et omnes qui hábitant in ea: * ego confirmávi colúmnas ejus.

Dixi iníquis: Nolíte iníque ágere: * et delinquéntibus: Nolíte exaltáre cornu:

Nolíte extóllere in altum cornu vestrum: * nolíte loqui advérsus Deum iniquitátem.

Quia neque ab Oriénte, neque ab Occidénte, neque a desértis móntibus: * quóniam Deus judex est.

Hunc humíliat, et hunc exáltat: * quia calix in manu Dómini vini meri plenus misto.

Et inclinávit ex hoc in hoc: verúmtamen fæx ejus non est exinaníta: * bibent omnes peccatóres terræ.

Ego autem annuntiábo in sæculum: * cantábo Deo Jacob.

Et ómnia córnua peccatórum confríngam: * et exaltabúntur córnua justi.

\psalmheading{Psalmus 75, i}

\lettrine[lines=2]{N}{OTUS} in Judǽa Deus: * in Israël magnum nomen ejus.

Et factus est in pace locus ejus: * et habitátio ejus in Sion.

Ibi confrégit poténtias árcuum, * scutum, gládium, et bellum.

Illúminans tu mirabíliter a móntibus ætérnis: * turbáti sunt omnes insipiéntes corde.

Dormiérunt somnum suum: * et nihil invenérunt omnes viri divitiárum in mánibus suis.

Ab increpatióne tua, Deus Jacob, * dormitavérunt qui ascendérunt equos.

\psalmheading{Psalmus 75, ii}

\lettrine[lines=2]{T}{U} terríbilis es, et quis resístet tibi? * ex tunc ira tua.

De cælo audítum fecísti judícium: * terra trémuit et quiévit.

Cum exsúrgeret in judícium Deus, * ut salvos fáceret omnes mansuétos terræ.

Quóniam cogitátio hóminis confitébitur tibi: * et relíquiæ cogitatiónis diem festum agent tibi.

Vovéte, et réddite Dómino, Deo vestro: * omnes, qui in circúitu ejus affértis múnera.

Terríbili et ei qui aufert spíritum príncipum, * terríbili apud reges terræ.

\antlabel{Ant.} Invocábimus nomen tuum, Dómine: narrábimus mirabília tua.

\variantheading{In Officio feriali}

\capitulumlabel{Capit.}{1 Petri 1, 17-19}

\lettrine[lines=2]{I}{N} timóre incolátus vestri témpore conversámini: sciéntes quod non corruptibílibus auro vel argénto redémpti estis, sed pretióso sánguine quasi Agni immaculáti Christi.

\R\ br. Rédime me, Dómine, * Et miserére mei. Rédime. \V\ Pes enim meus stetit in via recta. Et miserére mei. Glória Patri. Rédime. \V\ Ab occúltis meis munda me, Dómine. \R\ Et ab aliénis parce servo tuo.

\variantheading{In Vigiliis communibus:}

\variantheading{Preces feriales}

KYRIE, eléison. Christe, eléison. Kýrie, eléison.

Pater noster \rubric{secreto usque ad} \V\ Et ne nos indúcas in tentatiónem. \R\ Sed líbera nos a malo. \V\ Dómine, Deus virtútum, convérte nos. \R\ Et osténde fáciem tuam, et salvi érimus. \V\ Exsúrge, Christe, ádjuva nos. \R\ Et líbera nos propter nomen tuum. \V\ Dómine, exáudi oratiónem meam. \R\ Et clamor meus ad te véniat.

\rubric{Réliqua ut in Ordinário. 22}

% ===== END OF BATCH 28 -- Covers autumnalis.md lines 9229-9538,
% completing the Little Hours of Feria Quinta in full: "Ad
% Tertiam" (Psalms 72/i-iii), "Ad Sextam" (Psalms 73/i-iii), and
% "Ad Nonam" (Psalms 74, 75/i, 75/ii), each with its own hymn,
% Capitulum, brief Responsory, and "In Vigiliis communibus: Preces
% feriales" apparatus. Three further occurrences of the recurring
% stray OCR "\*" artifact in the hour-hymns are silently dropped.
% "Ut quid" (Psalm 73, i) correctly shows only its first short word
% in full caps ("UT"), consistent with the convention that only the
% single opening word takes full capitalization -- not a defect,
% unlike the partial-caps cases seen elsewhere. Several further
% genuine partial-caps OCR defects ("IDeo," "VErúmtamen," "USquequo,"
% "MEmor," "COnfitebimur," "NOtus") are silently normalized to full
% caps, consistent with prior batches. The document continues with
% "Ad Vesperas" (autumnalis.md line 9541).
% =====
\hourheading{Ad Vesperas}

Pater noster et Ave.

\V\ Deus, in adjutórium.

\antlabel{Ant.} Ecce, quam bonum * \textdagger\ et quam jucúndum, habitáre fratres in unum. \textdaggerdbl

\psalmheading{Psalmus 132}

\lettrine[lines=2]{E}{CCE}, quam bonum \textdagger\ et quam jucúndum, * habitáre fratres in unum: \textdaggerdbl

Sicut unguéntum in cápite, * quod descéndit in barbam, barbam Aaron.

Quod descéndit in oram vestiménti ejus: * sicut ros Hermon, qui descéndit in montem Sion.

Quóniam illic mandávit Dóminus benedictiónem, * et vitam usque in sæculum.

\antlabel{Ant.} Ecce, quam bonum et quam jucúndum, habitáre fratres in unum!

\antlabel{Ant.} Confitémini Dómino * \textdagger\ quóniam in ætérnum misericórdia ejus.

\psalmheading{Psalmus 135, i}

\lettrine[lines=2]{C}{ONFITEMINI} Dómino \textdagger\ quóniam bonus: * quóniam in ætérnum misericórdia ejus.

Confitémini Deo deórum: * quóniam in ætérnum misericórdia ejus.

Confitémini Dómino dominórum: * quóniam in ætérnum misericórdia ejus.

Qui facit mirabília magna solus: * quóniam in ætérnum misericórdia ejus.

Qui fecit cælos in intelléctu: * quóniam in ætérnum misericórdia ejus.

Qui firmávit terram super aquas: * quóniam in ætérnum misericórdia ejus.

Qui fecit luminária magna: * quóniam in ætérnum misericórdia ejus.

Solem in potestátem diéi: * quóniam in ætérnum misericórdia ejus.

Lunam, et stellas in potestátem noctis: * quóniam in ætérnum misericórdia ejus.

\antlabel{Ant.} Confitémini Dómino quóniam in ætérnum misericórdia ejus.

\antlabel{Ant.} Confitémini Dómino, * quia in humilitáte nostra memor fuit nostri.

\psalmheading{Psalmus 135, ii}

\lettrine[lines=2]{Q}{UI} percússit Ægýptum cum primogénitis eórum: * quóniam in ætérnum misericórdia ejus.

Qui edúxit Israël de médio eórum: * quóniam in ætérnum misericórdia ejus.

In manu poténti, et bráchio excélso: * quóniam in ætérnum misericórdia ejus.

Qui divísit Mare Rubrum in divisiónes: * quóniam in ætérnum misericórdia ejus.

Et edúxit Israël per médium ejus: * quóniam in ætérnum misericórdia ejus.

Et excússit Pharaónem, et virtútem ejus in Mari Rubro: * quóniam in ætérnum misericórdia ejus.

Qui tradúxit pópulum suum per desértum: * quóniam in ætérnum misericórdia ejus.

Qui percússit reges magnos: * quóniam in ætérnum misericórdia ejus.

Et occídit reges fortes: * quóniam in ætérnum misericórdia ejus.

Sehon, regem Amorrhæórum: * quóniam in ætérnum misericórdia ejus.

Et Og, regem Basan: * quóniam in ætérnum misericórdia ejus.

Et dedit terram eórum hereditátem: * quóniam in ætérnum misericórdia ejus.

Hereditátem Israël, servo suo: * quóniam in ætérnum misericórdia ejus.

Quia in humilitáte nostra memor fuit nostri: * quóniam in ætérnum misericórdia ejus.

Et redémit nos ab inimícis nostris: * quóniam in ætérnum misericórdia ejus.

Qui dat escam omni carni: * quóniam in ætérnum misericórdia ejus.

Confitémini Deo cæli: * quóniam in ætérnum misericórdia ejus.

Confitémini Dómino dominórum: * quóniam in ætérnum misericórdia ejus.

\antlabel{Ant.} Confitémini Dómino, quia in humilitáte nostra memor fuit nostri.

\antlabel{Ant.} Adhǽreat lingua mea * fáucibus meis, si non memínero tui, Jerúsalem.

\psalmheading{Psalmus 136}

\lettrine[lines=2]{S}{UPER} flúmina Babylónis, illic sédimus et flévimus: * cum recordarémur Sion:

In salícibus in médio ejus, * suspéndimus órgana nostra.

Quia illic interrogavérunt nos, qui captívos duxérunt nos, * verba cantiónum:

Et qui abduxérunt nos: * Hymnum cantáte nobis de cánticis Sion.

Quómodo cantábimus cánticum Dómini: * in terra aliéna?

Si oblítus fúero tui, Jerúsalem, * oblivióni detur déxtera mea.

Adhǽreat lingua mea fáucibus meis, * si non memínero tui:

Si non proposúero Jerúsalem, * in princípio lætítiæ meæ.

Memor esto, Dómine, filiórum Edom, * in die Jerúsalem:

Qui dicunt: Exinaníte, exinaníte: * usque ad fundaméntum in ea.

Fília Babylónis mísera: * beátus, qui retríbuet tibi retributiónem tuam, quam retribuísti nobis.

Beátus, qui tenébit, * et allídet párvulos tuos ad petram.

\antlabel{Ant.} Adhǽreat lingua mea fáucibus meis, si non memínero tui, Jerúsalem.

\antlabel{Ant.} Confitébor * \textdagger\ nómini tuo, Dómine, super misericórdia et veritáte tua.

\psalmheading{Psalmus 137}

\lettrine[lines=2]{C}{ONFITEBOR} \textdagger\ tibi, Dómine, in toto corde meo: * quóniam audísti verba oris mei.

In conspéctu Angelórum psallam tibi: * adorábo ad templum sanctum tuum, et confitébor nómini tuo.

Super misericórdia tua, et veritáte tua: * quóniam magnificásti super omne, nomen sanctum tuum.

In quacúmque die invocávero te, exáudi me: * multiplicábis in ánima mea virtútem.

Confiteántur tibi, Dómine, omnes reges terræ: * quia audiérunt ómnia verba oris tui:

Et cantent in viis Dómini: * quóniam magna est glória Dómini.

Quóniam excélsus Dóminus, et humília réspicit: * et alta a longe cognóscit.

Si ambulávero in médio tribulatiónis, vivificábis me: * et super iram inimicórum meórum extendísti manum tuam, et salvum me fecit déxtera tua.

Dóminus retríbuet pro me: * Dómine, misericórdia tua in sæculum: ópera mánuum tuárum ne despícias.

\antlabel{Ant.} Confitébor nómini tuo, Dómine, super misericórdia et veritáte tua.

\variantheading{In Officio feriali}

\capitulumlabel{Capitulum}{2 Cor. 1, 3-4}

\lettrine[lines=2]{B}{ENEDICTUS} Deus, et Pater Dómini nostri Jesu Christi, Pater misericordiárum, et Deus totíus consolatiónis, qui consolátur nos in omni tribulatióne nostra.

\hymnusheading{Hymnus}

\lettrine[lines=2]{M}{AGNAE} Deus poténtiæ,\\
Qui fértili natos aqua\\
Partim relínquis gúrgiti,\\
Partim levas in áëra.

Demérsa lymphis ímprimens,\\
Subvécta cælis érigens:\\
Ut stirpe ab una pródita,\\
Divérsa répleant loca:

Largíre cunctis sérvulis,\\
Quos mundat unda Sánguinis,\\
Nescíre lapsus críminum,\\
Nec ferre mortis tædium.

Ut culpa nullum déprimat:\\
Nullum éfferat jactántia:\\
Elísa mens ne cóncidat:\\
Eláta mens ne córruat.

Præsta, Pater piíssime,\\
Patríque compar Unice,\\
Cum Spíritu Paráclito\\
Regnans per omne sæculum.\\
Amen.

\V\ Dirigátur, Dómine, orátio mea. \R\ Sicut incénsum in conspéctu tuo.

\antlabel{Ad Magnif. Ant.} Fecit Deus * poténtiam in bráchio suo: dispérsit supérbos mente cordis sui.

\rubric{Réliqua ut in Ordinário. 22}

\hourheading{Ad Completorium}

\V\ Jube, domne.

\benedlabel{Bened.} Noctem quiétam, et finem perféctum concédat nobis Dóminus omnípotens. \R\ Amen.

\lectioheading{Lectio br.}{1 Petri 5, 8-9}

\lettrine[lines=2]{F}{RATRES}: Sóbrii estóte, et vigiláte: quia adversárius vester diábolus tamquam leo rúgiens círcuit, quærens quem dévoret: cui resístite fortes in fide. Tu autem, Dómine, miserére nobis. \R\ Deo grátias.

\V\ Adjutórium nostrum in nómine Dómini. \R\ Qui fecit cælum et terram.

Pater noster \rubric{quod dícitur totum secréto.}

\formulalabel{Confíteor.} Misereátur. Indulgéntiam.

\V\ Convérte nos, Deus, salutáris noster. \R\ Et avérte iram tuam a nobis.

\V\ Deus, in adjutórium.

\antlabel{Ant.} Adjútor meus.

\psalmheading{Psalmus 69}

\lettrine[lines=2]{D}{EUS}, in adjutórium meum inténde: * Dómine, ad adjuvándum me festína.

Confundántur et reverántur, * qui quærunt ánimam meam.

Avertántur retrórsum, et erubéscant, * qui volunt mihi mala.

Avertántur statim erubescéntes, * qui dicunt mihi: Euge, euge.

Exsúltent et læténtur in te omnes qui quærunt te, * et dicant semper: Magnificétur Dóminus: qui díligunt salutáre tuum.

Ego vero egénus, et pauper sum: * Deus, ádjuva me.

Adjútor meus, et liberátor meus es tu: * Dómine, ne moréris.

\psalmheading{Psalmus 70, i}

\lettrine[lines=2]{I}{N} te, Dómine, sperávi, non confúndar in ætérnum: * in justítia tua líbera me, et éripe me.

Inclína ad me aurem tuam, * et salva me.

Esto mihi in Deum protectórem, et in locum munítum: * ut salvum me fácias.

Quóniam firmaméntum meum, * et refúgium meum es tu.

Deus meus, éripe me de manu peccatóris, * et de manu contra legem agéntis et iníqui:

Quóniam tu es patiéntia mea, Dómine: * Dómine, spes mea a juventúte mea.

In te confirmátus sum ex útero: * de ventre matris meæ tu es protéctor meus.

In te cantátio mea semper: * tamquam prodígium factus sum multis: et tu adjútor fortis.

Repleátur os meum laude, ut cantem glóriam tuam: * tota die magnitúdinem tuam.

Ne projícias me in témpore senectútis: * cum defécerit virtus mea, ne derelínquas me.

Quia dixérunt inimíci mei mihi: * et qui custodiébant ánimam meam, consílium fecérunt in unum,

Dicéntes: Deus dereliquit eum, persequímini, et comprehéndite eum: * quia non est qui erípiat.

Deus, ne elongéris a me: * Deus meus, in auxílium meum réspice.

\psalmheading{Psalmus 70, ii}

\lettrine[lines=2]{C}{ONFUNDANTUR}, et defíciant detrahéntes ánimæ meæ: * operiántur confusióne, et pudóre qui quærunt mala mihi.

Ego autem semper sperábo: * et adjíciam super omnem laudem tuam.

Os meum annuntiábit justítiam tuam: * tota die salutáre tuum.

Quóniam non cognóvi litteratúram, introíbo in poténtias Dómini: * Dómine, memorábor justítiæ tuæ solíus.

Deus, docuísti me a juventúte mea: * et usque nunc pronuntiábo mirabília tua.

Et usque in senéctam et sénium: * Deus, ne derelínquas me.

Donec annúntiem bráchium tuum * generatióni omni, quæ ventúra est:

Poténtiam tuam, et justítiam tuam, Deus, usque in altíssima, quæ fecísti magnália: * Deus, quis símilis tibi?

Quantas ostendísti mihi tribulatiónes multas et malas: et convérsus vivificásti me: * et de abýssis terræ íterum reduxísti me:

Multiplicásti magnificéntiam tuam: * et convérsus consolátus es me.

Nam et ego confitébor tibi in vasis psalmi veritátem tuam: * Deus, psallam tibi in cíthara, Sanctus Israël.

Exsultábunt lábia mea cum cantávero tibi: * et ánima mea, quam redemísti.

Sed et lingua mea tota die meditábitur justítiam tuam: * cum confúsi et reveríti fúerint, qui quærunt mala mihi.

\antlabel{Ant.} Adjútor meus, et liberátor meus esto, Dómine.

\hymnusheading{Hymnus}

\lettrine[lines=2]{T}{E} lucis ante términum,\\
Rerum Creátor, póscimus,\\
Ut pro tua cleméntia\\
Sis præsul et custódia.

Procul recédant sómnia,\\
Et nóctium phantásmata,\\
Hostémque nostrum cómprime,\\
Ne polluántur córpora.

Præsta, Pater piíssime,\\
Patríque compar Unice,\\
Cum Spíritu Paráclito\\
Regnans per omne sæculum.\\
Amen.

\capitulumlabel{Capitulum}{Jer. 14, 9}

\lettrine[lines=2]{T}{U} autem in nobis es, Dómine, et nomen sanctum tuum invocátum est super nos: ne derelínquas nos, Dómine, Deus noster.

\R\ br. In manus tuas, Dómine, * Comméndo spíritum meum. In manus. \V\ Redemísti nos, Dómine, Deus veritátis. Comméndo. Glória Patri. In manus.

\V\ Custódi nos, Dómine, ut pupíllam óculi. \R\ Sub umbra alárum tuárum prótege nos.

\antlabel{Ant.} Salva nos.

\canticumheading{Canticum Simeonis}{Luc. 2, 29-32}

\lettrine[lines=2]{N}{UNC} dimíttis servum tuum, Dómine, * secúndum verbum tuum in pace:

Quia vidérunt óculi mei * salutáre tuum,

Quod parásti * ante fáciem ómnium populórum,

Lumen ad revelatiónem géntium, * et glóriam plebis tuæ Israël.

Glória Patri.

\antlabel{Ant.} Salva nos, Dómine, vigilántes, custódi nos dormiéntes; ut vigilémus cum Christo, et requiescámus in pace.

\variantheading{Preces (si dicendæ sint)}

KYRIE, eléison. Christe, eléison. Kýrie, eléison.

Pater noster \rubric{secreto usque ad} \V\ Et ne nos indúcas in tentatiónem. \R\ Sed líbera nos a malo.

Credo in Deum \rubric{secreto usque ad} \V\ Carnis resurrectiónem. \R\ Vitam ætérnam. Amen. \V\ Benedíctus es, Dómine, Deus patrum nostrórum. \R\ Et laudábilis et gloriósus in sæcula. \V\ Benedicámus Patrem et Fílium cum Sancto Spíritu. \R\ Laudémus, et superexaltémus eum in sæcula. \V\ Benedíctus es, Dómine, in firmaménto cæli. \R\ Et laudábilis, et gloriósus, et superexaltátus in sæcula. \V\ Benedícat et custódiat nos omnípotens et miséricors Dóminus. \R\ Amen. \V\ Dignáre, Dómine, nocte ista. \R\ Sine peccáto nos custodíre. \V\ Miserére nostri, Dómine. \R\ Miserére nostri. \V\ Fiat misericórdia tua, Dómine, super nos. \R\ Quemádmodum sperávimus in te. \V\ Dómine, exáudi oratiónem meam. \R\ Et clamor meus ad te véniat.

\V\ Dóminus vobíscum.

Orémus. \oratiolabel{Oratio}

\lettrine[lines=2]{V}{ISITA}, quæsumus, Dómine, habitatiónem istam, et omnes insídias inimíci ab ea longe repélle: Angeli tui sancti hábitent in ea, qui nos in pace custódiant; et benedíctio tua sit super nos semper. Per Dóminum nostrum. \R\ Amen.

\V\ Dóminus vobíscum.

\V\ Benedicámus Dómino.

\benedictiolabel{Benedictio} Benedícat et custódiat nos omnípotens et miséricors Dóminus, Pater, et Fílius, et Spíritus Sanctus. \R\ Amen.

\variantheading{Antiphona finalis beátæ Maríæ Vírginis.}

\V\ Divínum auxílium.

Pater noster, Ave María et Credo.

% ===== END OF BATCH 29 -- Covers autumnalis.md lines 9541-9888,
% completing "Ad Vesperas" (Psalms 132, 135/i, 135/ii, 136, 137,
% hymn "Magnæ Deus poténtiæ") and "Ad Completorium" (Psalms 69,
% 70/i, 70/ii, hymn "Te lucis ante términum," Canticum Simeonis,
% Preces), and with them the ENTIRE "Feria Quinta" Office. The
% closing cross-reference at the end of Compline is given in its
% fuller unabbreviated form here -- "Antiphona finalis beátæ Maríæ
% Vírginis." and "Pater noster, Ave María et Credo." -- rather than
% the abbreviated "Antiphona finalis B. M. V." and "Pater, Ave et
% Credo." seen at the end of every previous feria's Compline this
% volume; transcribed exactly as the source gives each. The
% document continues with "FERIA SEXTA" (autumnalis.md line 9891),
% printed in the source in full capitals as a heading (transcribed
% with \feriaheading{} regardless, consistent with the role-based
% convention already applied to every other day-heading in this
% volume). The Invitatorium at Friday Matins carries a bracketed
% footnote marker "[3]" after its incipit in the source, preserved
% verbatim as part of the transcribed text.
% =====
\feriaheading{Feria Sexta}

\hourheading{Ad Matutinum}

Pater, Ave et Credo.

\V\ Dómine, lábia mea.

\V\ Deus, in adjutórium.

\variantheading{In Officio feriali}

\invitatlabel{Invit.} Dóminum, Deum nostrum, * Veníte, adorémus. \psalmlabel{Ps. 94} Veníte. [3]

\hymnusheading{Hymnus}

\lettrine[lines=2]{T}{U}, Trinitátis Únitas,\\
Orbem poténter quæ regis,\\
Atténde laudis cánticum,\\
Quod excubántes psállimus.

Nam léctulo consúrgimus\\
Noctis quiéto témpore,\\
Ut flagitémus ómnium\\
A te medélam vúlnerum.

Quo fraude quidquid dæmonum\\
In nóctibus delíquimus,\\
Abstérgat illud cælitus\\
Tuæ potéstas glóriæ.

Ne corpus astet sórdidum,\\
Nec torpor instet córdium,\\
Ne críminis contágio\\
Tepéscat ardor spíritus.

Ob hoc, Redémptor, quǽsumus,\\
Reple tuo nos lúmine,\\
Per quod diérum círculis\\
Nullis ruámus áctibus.

Præsta, Pater piíssime,\\
Patríque compar Unice,\\
Cum Spíritu Paráclito\\
Regnans per omne sæculum.\\
Amen.

\nocturnoheading{In I Nocturno}

\antlabel{Ant.} Suscitávit Dóminus * testimónium in Jacob: et legem pósuit in Israël.

\psalmheading{Psalmus 77, i}

\lettrine[lines=2]{A}{TTENDITE}, pópule meus, legem meam: * inclináte aurem vestram in verba oris mei.

Apériam in parábolis os meum: * loquar propositiónes ab inítio.

Quanta audívimus et cognóvimus ea: * et patres nostri narravérunt nobis.

Non sunt occultáta a fíliis eórum, * in generatióne áltera.

Narrántes laudes Dómini, et virtútes ejus, * et mirabília ejus, quæ fecit.

Et suscitávit testimónium in Jacob: * et legem pósuit in Israël.

Quanta mandávit pátribus nostris nota fácere ea fíliis suis: * ut cognóscat generátio áltera.

Fílii qui nascéntur, et exsúrgent, * et narrábunt fíliis suis,

Ut ponant in Deo spem suam, et non obliviscántur óperum Dei: * et mandáta ejus exquírant.

Ne fiant sicut patres eórum: * generátio prava et exásperans.

Generátio, quæ non diréxit cor suum: * et non est créditus cum Deo spíritus ejus.

\antlabel{Ant.} Suscitávit Dóminus testimónium in Jacob: et legem pósuit in Israël.

\antlabel{Ant.} Coram pátribus eórum * fecit Deus mirabília.

\psalmheading{Psalmus 77, ii}

\lettrine[lines=2]{F}{ILII} Ephrem intendéntes et mitténtes arcum: * convérsi sunt in die belli.

Non custodiérunt testaméntum Dei: * et in lege ejus noluérunt ambuláre.

Et oblíti sunt benefactórum ejus, * et mirabílium ejus quæ osténdit eis.

Coram pátribus eórum fecit mirabília in terra Ægýpti, * in campo Táneos.

Interrúpit mare, et perdúxit eos: * et státuit aquas quasi in utre.

Et dedúxit eos in nube diéi: * et tota nocte in illuminatióne ignis.

Interrúpit petram in erémo: * et adaquávit eos velut in abýsso multa.

Et edúxit aquam de petra: * et dedúxit tamquam flúmina aquas.

\antlabel{Ant.} Coram pátribus eórum fecit Deus mirabília.

\antlabel{Ant.} Jánuas cæli apéruit * Dóminus, et pluit illis manna ad manducándum.

\psalmheading{Psalmus 77, iii}

\lettrine[lines=2]{E}{T} apposuérunt adhuc peccáre ei: * in iram excitavérunt Excélsum in inaquóso.

Et tentavérunt Deum in córdibus suis, * ut péterent escas animábus suis.

Et male locúti sunt de Deo: * dixérunt: Numquid póterit Deus paráre mensam in desérto?

Quóniam percússit petram, et fluxérunt aquæ, * et torréntes inundavérunt.

Numquid et panem póterit dare, * aut paráre mensam pópulo suo?

Ideo audívit Dóminus, et dístulit: * et ignis accénsus est in Jacob, et ira ascéndit in Israël:

Quia non credidérunt in Deo, * nec speravérunt in salutári ejus:

Et mandávit núbibus désuper, * et jánuas cæli apéruit.

Et pluit illis manna ad manducándum, * et panem cæli dedit eis.

Panem Angelórum manducávit homo, * cibária misit eis in abundántia.

Tránstulit Austrum de cælo: * et indúxit in virtúte sua Africum.

Et pluit super eos sicut púlverem carnes: * et sicut arénam maris volatília pennáta.

Et cecidérunt in médio castrórum eórum: * circa tabernácula eórum.

Et manducavérunt, et saturáti sunt nimis, et desidérium eórum áttulit eis: * non sunt fraudáti a desidério suo.

Adhuc escæ eórum erant in ore ipsórum: * et ira Dei ascéndit super eos.

Et occídit pingues eórum, * et eléctos Israël impedívit.

\antlabel{Ant.} Jánuas cæli apéruit Dóminus, et pluit illis manna ad manducándum.

\variantheading{In Festis novem Lectionum:}

\V\ Deus percússit petram, et fluxérunt aquæ. \R\ Et torréntes inundavérunt.

\absolutiolabel{Absolutio.} Exáudi, Dómine.

\benedlabel{Bened.} 1 Benedictióne perpétua. 2 Unigénitus Dei Fílius. 3 Spíritus Sancti grátia.

\nocturnoheading{In II Nocturno}

\antlabel{Ant.} Deus adjútor * est eórum: et Excélsus redémptor eórum est.

\psalmheading{Psalmus 77, iv}

\lettrine[lines=2]{I}{N} ómnibus his peccavérunt adhuc: * et non credidérunt in mirabílibus ejus.

Et defecérunt in vanitáte dies eórum: * et anni eórum cum festinatióne.

Cum occíderet eos, quærébant eum: * et revertebántur, et dilúculo veniébant ad eum.

Et rememoráti sunt quia Deus adjútor est eórum: * et Deus excélsus redémptor eórum est.

Et dilexérunt eum in ore suo, * et lingua sua mentíti sunt ei:

Cor autem eórum non erat rectum cum eo: * nec fidéles habíti sunt in testaménto ejus.

Ipse autem est miséricors, et propítius fiet peccátis eórum: * et non dispérdet eos.

Et abundávit ut avérteret iram suam: * et non accéndit omnem iram suam:

Et recordátus est quia caro sunt: * spíritus vadens, et non rédiens.

Quóties exacerbavérunt eum in desérto, * in iram concitavérunt eum in inaquóso?

Et convérsi sunt, et tentavérunt Deum: * et Sanctum Israël exacerbavérunt.

\antlabel{Ant.} Deus adjútor est eórum: et Excélsus redémptor eórum est.

\antlabel{Ant.} Redémit eos * Dóminus de manu tribulántis.

\psalmheading{Psalmus 77, v}

\lettrine[lines=2]{N}{ON} sunt recordáti manus ejus, * die qua redémit eos de manu tribulántis.

Sicut pósuit in Ægýpto signa sua, * et prodígia sua in campo Táneos.

Et convértit in sánguinem flúmina eórum, * et imbres eórum, ne bíberent.

Misit in eos cœnomyíam, et comédit eos: * et ranam, et dispérdidit eos.

Et dedit ærúgini fructus eórum, * et labóres eórum locústæ.

Et occídit in grándine víneas eórum, * et moros eórum in prúina.

Et trádidit grándini juménta eórum: * et possessiónem eórum igni.

Misit in eos iram indignatiónis suæ: * indignatiónem, et iram, et tribulatiónem: immissiónes per ángelos malos.

Viam fecit sémitæ iræ suæ, non pepércit a morte animábus eórum: * et juménta eórum in morte conclúsit.

Et percússit omne primogénitum in terra Ægýpti: * primítias omnis labóris eórum in tabernáculis Cham.

Et ábstulit sicut oves pópulum suum: * et perdúxit eos tamquam gregem in desérto.

Et dedúxit eos in spe, et non timuérunt: * et inimícos eórum opéruit mare.

Et indúxit eos in montem sanctificatiónis suæ, * montem, quem acquisívit déxtera ejus.

Et ejécit a fácie eórum gentes: * et sorte divísit eis terram in funículo distributiónis.

Et habitáre fecit in tabernáculis eórum * tribus Israël.

Et tentavérunt, et exacerbavérunt Deum excélsum: * et testimónia ejus non custodiérunt.

Et avertérunt se, et non servavérunt pactum: * quemádmodum patres eórum, convérsi sunt in arcum pravum.

In iram concitavérunt eum in cóllibus suis: * et in sculptílibus suis ad æmulatiónem eum provocavérunt.

\antlabel{Ant.} Redémit eos Dóminus de manu tribulántis.

\antlabel{Ant.} Ædificávit * Deus sanctifícium suum in terra.

\psalmheading{Psalmus 77, vi}

\lettrine[lines=2]{A}{UDIVIT} Deus, et sprevit: * et ad níhilum redégit valde Israël.

Et répulit tabernáculum Silo, * tabernáculum suum, ubi habitávit in homínibus.

Et trádidit in captivitátem virtútem eórum: * et pulchritúdinem eórum in manus inimíci.

Et conclúsit in gládio pópulum suum: * et hereditátem suam sprevit.

Júvenes eórum comédit ignis: * et vírgines eórum non sunt lamentátæ.

Sacerdótes eórum in gládio cecidérunt: * et víduæ eórum non plorabántur.

Et excitátus est tamquam dórmiens Dóminus, * tamquam potens crapulátus a vino.

Et percússit inimícos suos in posterióra: * oppróbrium sempitérnum dedit illis.

Et répulit tabernáculum Joseph: * et tribum Ephraim non elégit:

Sed elégit tribum Juda, * montem Sion quem diléxit.

Et ædificávit sicut unicórnium sanctifícium suum in terra, * quam fundávit in sæcula.

Et elégit David, servum suum, et sústulit eum de grégibus óvium: * de post fetántes accépit eum.

Páscere Jacob, servum suum, * et Israël, hereditátem suam:

Et pavit eos in innocéntia cordis sui: * et in intelléctibus mánuum suárum dedúxit eos.

\antlabel{Ant.} Ædificávit Deus sanctifícium suum in terra.

\variantheading{In Festis novem Lectionum:}

\V\ Elégit Dóminus tribum Juda. \R\ Montem Sion quem diléxit.

\absolutiolabel{Absolutio.} Ipsíus píetas et.

\benedlabel{Bened.} 4 Deus Pater omnípotens. 5 Christus perpétuæ. 6 Ignem sui amóris.

\nocturnoheading{In III Nocturno}

\antlabel{Ant.} Adjuva nos, * Deus, salutáris noster: et propítius esto peccátis nostris.

\psalmheading{Psalmus 78}

\lettrine[lines=2]{D}{EUS}, venérunt gentes in hereditátem tuam, polluérunt templum sanctum tuum: * posuérunt Jerúsalem in pomórum custódiam.

Posuérunt morticína servórum tuórum, escas volatílibus cæli: * carnes sanctórum tuórum béstiis terræ.

Effudérunt sánguinem eórum tamquam aquam in circúitu Jerúsalem: * et non erat qui sepelíret.

Facti sumus oppróbrium vicínis nostris: * subsannátio et illúsio his, qui in circúitu nostro sunt.

Usquequo, Dómine, irascéris in finem: * accendétur velut ignis zelus tuus?

Effúnde iram tuam in gentes, quæ te non novérunt: * et in regna, quæ nomen tuum non invocavérunt:

Quia comedérunt Jacob: * et locum ejus desolavérunt.

Ne memíneris iniquitátum nostrárum antiquárum, cito antícipent nos misericórdiæ tuæ: * quia páuperes facti sumus nimis.

Adjuva nos, Deus, salutáris noster: et propter glóriam nóminis tui, Dómine, líbera nos: * et propítius esto peccátis nostris, propter nomen tuum:

Ne forte dicant in géntibus: Ubi est Deus eórum? * et innotéscat in natiónibus coram óculis nostris.

Ultio sánguinis servórum tuórum, qui effúsus est: * intróeat in conspéctu tuo gémitus compeditórum.

Secúndum magnitúdinem bráchii tui, * pósside fílios mortificatórum.

Et redde vicínis nostris séptuplum in sinu eórum: * impropérium ipsórum, quod exprobravérunt tibi, Dómine.

Nos autem pópulus tuus, et oves páscuæ tuæ, * confitébimur tibi in sæculum.

In generatiónem et generatiónem * annuntiábimus laudem tuam.

\antlabel{Ant.} Adjuva nos, Deus, salutáris noster: et propítius esto peccátis nostris.

\antlabel{Ant.} Ego sum Dóminus, * Deus tuus, Israël, qui edúxi te de terra Ægýpti.

\psalmheading{Psalmus 80}

\lettrine[lines=2]{E}{XSULTATE} Deo, adjutóri nostro: * jubiláte Deo Jacob.

Súmite psalmum, et date týmpanum: * psaltérium jucúndum cum cíthara.

Buccináte in Noeménia tuba, * in insígni die solemnitátis vestræ.

Quia præcéptum in Israël est: * et judícium Deo Jacob.

Testimónium in Joseph pósuit illud, cum exíret de terra Ægýpti: * linguam, quam non nóverat, audívit.

Divértit ab onéribus dorsum ejus: * manus ejus in cóphino serviérunt.

In tribulatióne invocásti me, et liberávi te: * exaudívi te in abscóndito tempestátis: probávi te apud aquam contradictiónis.

Audi, pópulus meus, et contestábor te: * Israël, si audíeris me, non erit in te deus recens, neque adorábis deum aliénum.

Ego enim sum Dóminus, Deus tuus, qui edúxi te de terra Ægýpti: * diláta os tuum, et implébo illud.

Et non audívit pópulus meus vocem meam: * et Israël non inténdit mihi.

Et dimísi eos secúndum desidéria cordis eórum, * ibunt in adinventiónibus suis.

Si pópulus meus audísset me: * Israël si in viis meis ambulásset:

Pro níhilo fórsitan inimícos eórum humiliássem: * et super tribulántes eos misíssem manum meam.

Inimíci Dómini mentíti sunt ei: * et erit tempus eórum in sæcula.

Et cibávit eos ex ádipe fruménti: * et de petra, melle saturávit eos.

\antlabel{Ant.} Ego sum Dóminus, Deus tuus, Israël, qui edúxi te de terra Ægýpti.

\antlabel{Ant.} Ne táceas, Deus, * quóniam inimíci tui extulérunt caput.

\psalmheading{Psalmus 82}

\lettrine[lines=2]{D}{EUS}, quis símilis erit tibi? * ne táceas, neque compescáris, Deus:

Quóniam ecce inimíci tui sonuérunt: * et qui odérunt te, extulérunt caput.

Super pópulum tuum malignavérunt consílium: * et cogitavérunt advérsus sanctos tuos.

Dixérunt: Veníte, et disperdámus eos de gente: * et non memorétur nomen Israël ultra.

Quóniam cogitavérunt unanímiter: * simul advérsum te testaméntum disposuérunt, tabernácula Idumæórum et Ismahelítæ:

Moab, et Agaréni, Gebal, et Ammon, et Amalec: * alienígenæ cum habitántibus Tyrum.

Etenim Assur venit cum illis: * facti sunt in adjutórium fíliis Lot.

Fac illis sicut Mádian, et Sísaræ: * sicut Jabin in torrénte Cisson.

Disperiérunt in Endor: * facti sunt ut stercus terræ.

Pone príncipes eórum sicut Oreb, et Zeb, * et Zébee, et Sálmana:

Omnes príncipes eórum: * qui dixérunt: Hereditáte possideámus Sanctuárium Dei.

Deus meus, pone illos ut rotam: * et sicut stípulam ante fáciem venti.

Sicut ignis, qui combúrit silvam: * et sicut flamma combúrens montes:

Ita persequéris illos in tempestáte tua: * et in ira tua turbábis eos.

Imple fácies eórum ignomínia: * et quærent nomen tuum, Dómine.

Erubéscant, et conturbéntur in sæculum sæculi: * et confundántur et péreant.

Et cognóscant quia nomen tibi Dóminus: * tu solus Altíssimus in omni terra.

\antlabel{Ant.} Ne táceas, Deus, quóniam inimíci tui extulérunt caput.

\variantheading{In feriali Officio et in Festis:}

\V\ Cognóscant quia nomen tibi Dóminus. \R\ Tu solus Altíssimus super omnem terram.

\variantheading{In Officio novem Lectionum}

\absolutiolabel{Absolutio.} A vínculis.

\benedlabel{Bened.} 7 Evangélica léctio. 8 In Festis Dómini: Divínum auxílium. In Festis B. M. V.: Cujus festum cólimus, ipsa Virgo. In Festis Sanctórum: Cujus (\textit{vel} Quorum \textit{aut} Quarum) festum cólimus. 9 Ad societátem. (Per evangélica dicta.)

\variantheading{In Officio trium Lectionum}

\absolutiolabel{Absolutio.} Ipsíus píetas et.

\variantheading{In Feriis cum Homilia:}

\benedlabel{Bened.} 1 Evangélica. 2 Divínum. 3 Ad societátem.

\variantheading{In Feriis sine Homilia:}

\benedlabel{Bened.} 1 Deus Pater. 2 Christus perpétuæ. 3 Ignem sui amóris.

\variantheading{In Festis Sanctorum:}

\benedlabel{Bened.} 1 Ille nos benedícat, qui sine fine vivit et regnat. \R\ Amen. 2 Cujus (\textit{vel} Quorum \textit{aut} Quarum.) 3 Ad societátem.

% ===== END OF BATCH 30 -- Covers autumnalis.md lines 9891-10290,
% opening "Feria Sexta" (Friday) and completing its "Ad Matutinum"
% in full: the ferial Invitatorium "Dóminum, Deum nostrum" (with a
% bracketed footnote marker "[3]" in the source preserved verbatim
% after the incipit), the hymn "Tu, Trinitátis Únitas," and the
% three Nocturns -- Psalm 77 split across six sections (i-vi) over
% the first two Nocturns, and Psalms 78, 80, 82 in the third --
% with their versicle, Absolutio, and Benedictiones apparatus. The
% day-heading itself appears in the source in full capitals ("FERIA
% SEXTA"), transcribed with \feriaheading{} regardless, consistent
% with the role-based heading convention applied throughout this
% volume. Several more psalm incipits carry the partial-caps OCR
% defect already documented in prior batches ("ATtendite,"
% "FIlii," "ET apposuérunt" is correctly two full words so left
% as-is being a short conjunction, "AUdivit," "DEus" twice,
% "EXsultáte") and are silently normalized to full caps where the
% defect is genuine. The document continues with "Ad Laudes I"
% (autumnalis.md line 10292).
% =====
\hourheading{Ad Laudes I}

\rubric{In Fériis extra Vigílias commúnes occurréntibus, atque in Festis.}

(Pater noster et Ave.)

\V\ Deus, in adjutórium.

\antlabel{Ant.} Exaltáte * Dóminum, Deum nostrum, et adoráte in monte sancto ejus.

\psalmheading{Psalmus 98}

\lettrine[lines=2]{D}{OMINUS} regnávit, irascántur pópuli: * qui sedet super Chérubim, moveátur terra.

Dóminus in Sion magnus: * et excélsus super omnes pópulos.

Confiteántur nómini tuo magno: quóniam terríbile, et sanctum est: * et honor regis judícium díligit.

Tu parásti directiónes: * judícium et justítiam in Jacob tu fecísti.

Exaltáte Dóminum, Deum nostrum, et adoráte scabéllum pedum ejus: * quóniam sanctum est.

Móyses et Aaron in sacerdótibus ejus: * et Sámuel inter eos, qui ínvocant nomen ejus:

Invocábant Dóminum, et ipse exaudiébat eos: * in colúmna nubis loquebátur ad eos.

Custodiébant testimónia ejus, * et præcéptum quod dedit illis.

Dómine, Deus noster, tu exaudiébas eos: * Deus, tu propítius fuísti eis, et ulcíscens in omnes adinventiónes eórum.

Exaltáte Dóminum, Deum nostrum, et adoráte in monte sancto ejus: * quóniam sanctus Dóminus, Deus noster.

\antlabel{Ant.} Exaltáte Dóminum, Deum nostrum, et adoráte in monte sancto ejus.

\antlabel{Ant.} Eripe me * de inimícis meis, Dómine, ad te confúgi.

\psalmheading{Psalmus 142}

\lettrine[lines=2]{D}{OMINE}, exáudi oratiónem meam: áuribus pércipe obsecratiónem meam in veritáte tua: * exáudi me in tua justítia.

Et non intres in judícium cum servo tuo: * quia non justificábitur in conspéctu tuo omnis vivens.

Quia persecútus est inimícus ánimam meam: * humiliávit in terra vitam meam.

Collocávit me in obscúris sicut mórtuos sæculi: * et anxiátus est super me spíritus meus, in me turbátum est cor meum.

Memor fui diérum antiquórum, meditátus sum in ómnibus opéribus tuis: * in factis mánuum tuárum meditábar.

Expándi manus meas ad te: * ánima mea sicut terra sine aqua tibi.

Velóciter exáudi me, Dómine: * defécit spíritus meus.

Non avértas fáciem tuam a me: * et símilis ero descendéntibus in lacum.

Audítam fac mihi mane misericórdiam tuam: * quia in te sperávi.

Notam fac mihi viam, in qua ámbulem: * quia ad te levávi ánimam meam.

Eripe me de inimícis meis, Dómine, ad te confúgi: * doce me fácere voluntátem tuam, quia Deus meus es tu.

Spíritus tuus bonus dedúcet me in terram rectam: * propter nomen tuum, Dómine, vivificábis me, in æquitáte tua.

Edúces de tribulatióne ánimam meam: * et in misericórdia tua dispérdes inimícos meos.

Et perdes omnes, qui tríbulant ánimam meam: * quóniam ego servus tuus sum.

\antlabel{Ant.} Eripe me de inimícis meis, Dómine, ad te confúgi.

\antlabel{Ant.} Benedixísti, * \textdagger\ Dómine, terram tuam: remisísti iniquitátem plebis tuæ.

\psalmheading{Psalmus 84}

\lettrine[lines=2]{B}{ENEDIXISTI}, \textdagger\ Dómine, terram tuam: * avertísti captivitátem Jacob.

Remisísti iniquitátem plebis tuæ: * operuísti ómnia peccáta eórum.

Mitigásti omnem iram tuam: * avertísti ab ira indignatiónis tuæ.

Convérte nos, Deus, salutáris noster: * et avérte iram tuam a nobis.

Numquid in ætérnum irascéris nobis? * aut exténdes iram tuam a generatióne in generatiónem?

Deus, tu convérsus vivificábis nos: * et plebs tua lætábitur in te.

Osténde nobis, Dómine, misericórdiam tuam: * et salutáre tuum da nobis.

Audiam quid loquátur in me Dóminus Deus: * quóniam loquétur pacem in plebem suam.

Et super sanctos suos: * et in eos, qui convertúntur ad cor.

Verúmtamen prope timéntibus eum salutáre ipsíus: * ut inhábitet glória in terra nostra.

Misericórdia, et véritas obviavérunt sibi: * justítia, et pax osculátæ sunt.

Véritas de terra orta est: * et justítia de cælo prospéxit.

Etenim Dóminus dabit benignitátem: * et terra nostra dabit fructum suum.

Justítia ante eum ambulábit: * et ponet in via gressus suos.

\antlabel{Ant.} Benedixísti, Dómine, terram tuam: remisísti iniquitátem plebis tuæ.

\antlabel{Ant.} In Dómino justificábitur, * et laudábitur omne semen Israël.

\canticumheading{Canticum Isaiæ}{Is. 45, 15-26}

\lettrine[lines=2]{V}{ERE} tu es Deus abscónditus, * Deus Israël, Salvátor.

Confúsi sunt, et erubuérunt omnes: * simul abiérunt in confusiónem fabricatóres errórum.

Israël salvátus est in Dómino salúte ætérna: * non confundémini, et non erubescétis usque in sæculum sæculi.

Quia hæc dicit Dóminus creans cælos, * ipse Deus formans terram, et fáciens eam, ipse plastes ejus:

Non in vanum creávit eam, ut habitarétur formávit eam: * Ego Dóminus, et non est álius.

Non in abscóndito locútus sum, * in loco terræ tenebróso:

Non dixi sémini Jacob frustra: Quærite me: * ego Dóminus loquens justítiam, annúntians recta.

Congregámini, et veníte, et accédite simul * qui salváti estis ex géntibus:

Nesciérunt qui levant lignum sculptúræ suæ, * et rogant deum non salvántem.

Annuntiáte, et veníte, et consiliámini simul: * Quis audítum fecit hoc ab inítio, ex tunc prædíxit illud?

Numquid non ego Dóminus, et non est ultra Deus absque me? * Deus justus, et salvans non est præter me.

Convertímini ad me, et salvi éritis, omnes fines terræ: * quia ego Deus, et non est álius.

In memetípso jurávi, egrediétur de ore meo justítiæ verbum, * et non revertétur:

Quia mihi curvábitur omne genu, * et jurábit omnis lingua.

Ergo in Dómino, dicet, meæ sunt justítiæ et impérium: * ad eum vénient, et confundéntur omnes qui repúgnant ei.

In Dómino justificábitur, et laudábitur * omne semen Israël.

\antlabel{Ant.} In Dómino justificábitur, et laudábitur omne semen Israël.

\antlabel{Ant.} Lauda, * \textdagger\ Jerúsalem, Dóminum. \textdaggerdbl

\psalmheading{Psalmus 147}

\lettrine[lines=2]{L}{AUDA}, \textdagger\ Jerúsalem, Dóminum: * \textdaggerdbl\ lauda Deum tuum, Sion.

Quóniam confortávit seras portárum tuárum: * benedíxit fíliis tuis in te.

Qui pósuit fines tuos pacem: * et ádipe fruménti sátiat te.

Qui emíttit elóquium suum terræ: * velóciter currit sermo ejus.

Qui dat nivem sicut lanam: * nébulam sicut cínerem spargit.

Mittit crystállum suam sicut buccéllas: * ante fáciem frígoris ejus quis sustinébit?

Emíttet verbum suum, et liquefáciet ea: * flabit spíritus ejus, et fluent aquæ.

Qui annúntiat verbum suum Jacob: * justítias, et judícia sua Israël.

Non fecit táliter omni natióni: * et judícia sua non manifestávit eis.

\antlabel{Ant.} Lauda, Jerúsalem, Dóminum.

\variantheading{In Officio feriali}

\capitulumlabel{Capit.}{Rom. 13, 12-13}

\lettrine[lines=2]{N}{OX} præcéssit, dies autem appropinquávit. Abjiciámus ergo ópera tenebrárum, et induámur arma lucis. Sicut in die honéste ambulémus.

\hymnusheading{Hymnus}

\lettrine[lines=2]{Æ}{TERNA} cæli glória,\\
Beáta spes mortálium,\\
Summi Tonántis Únice,\\
Castǽque proles Vírginis:

Da déxteram surgéntibus,\\
Exsúrgat et mens sóbria,\\
Flagrans et in laudem Dei\\
Grates rependat débitas.

Ortus refúlget Lúcifer,\\
Præítque solem núntius:\\
Cadunt tenébræ nóctium:\\
Lux sancta nos illúminet.

Manénsque nostris sénsibus,\\
Noctem repéllat sæculi,\\
Omníque fine témporis\\
Purgáta servet péctora.

Quæsíta jam primum fides\\
In corde radíces agat:\\
Secúnda spes congáudeat,\\
Qua major exstat cáritas.

Deo Patri sit glória,\\
Ejúsque soli Fílio,\\
Cum Spíritu Paráclito,\\
Nunc et per omne sæculum.\\
Amen.

\V\ Repléti sumus mane misericórdia tua. \R\ Exsultávimus, et delectáti sumus.

\antlabel{Ad Bened. Ant.} Per víscera misericórdiæ * Dei nostri visitávit nos Oriens ex alto.

\rubric{Réliqua ut in Ordinário. 9}

% ===== END OF BATCH 31 -- Covers autumnalis.md lines 10292-10502,
% "Ad Laudes I" for Feria Sexta in full: Psalms 98, 142, 84, the
% Canticum Isaiæ, and Psalm 147, followed by the Capitulum "Nox
% præcéssit" and the hymn "Æterna cæli glória" (whose opening
% lettrine letter is the ligature "Æ," matching the established
% convention of lettrine-ing exactly the source's first letter).
% Several more psalm incipits carry the partial-caps OCR defect
% ("DOmine, exáudi," "BEnedixísti," "VEre tu es," "LAuda,
% Jerúsalem") and are silently normalized to full caps. The
% document continues with "Ad Laudes II" (autumnalis.md line 10505).
% =====
\hourheading{Ad Laudes II}

\rubric{In Vigíliis commúnibus, quando fit Offícium de Féria.}

(Pater noster et Ave.)

\V\ Deus, in adjutórium.

\antlabel{Ant.} Cor contrítum.

\psalmheading{Psalmus 50}

\lettrine[lines=2]{M}{ISERERE} mei, Deus, * secúndum magnam misericórdiam tuam.

Et secúndum multitúdinem miseratiónum tuárum, * dele iniquitátem meam.

Amplius lava me ab iniquitáte mea: * et a peccáto meo munda me.

Quóniam iniquitátem meam ego cognósco: * et peccátum meum contra me est semper.

Tibi soli peccávi, et malum coram te feci: * ut justificéris in sermónibus tuis, et vincas cum judicáris.

Ecce enim, in iniquitátibus concéptus sum: * et in peccátis concépit me mater mea.

Ecce enim, veritátem dilexísti: * incérta et occúlta sapiéntiæ tuæ manifestásti mihi.

Aspérges me hyssópo, et mundábor: * lavábis me, et super nivem dealbábor.

Audítui meo dabis gáudium et lætítiam: * et exsultábunt ossa humiliáta.

Avérte fáciem tuam a peccátis meis: * et omnes iniquitátes meas dele.

Cor mundum crea in me, Deus: * et spíritum rectum ínnova in viscéribus meis.

Ne projícias me a fácie tua: * et spíritum sanctum tuum ne áuferas a me.

Redde mihi lætítiam salutáris tui: * et spíritu principáli confírma me.

Docébo iníquos vias tuas: * et ímpii ad te converténtur.

Líbera me de sanguínibus, Deus, Deus salútis meæ: * et exsultábit lingua mea justítiam tuam.

Dómine, lábia mea apéries: * et os meum annuntiábit laudem tuam.

Quóniam si voluísses sacrifícium, dedíssem útique: * holocáustis non delectáberis.

Sacrifícium Deo spíritus contribulátus: * cor contrítum, et humiliátum, Deus, non despícies.

Benígne fac, Dómine, in bona voluntáte tua Sion: * ut ædificéntur muri Jerúsalem.

Tunc acceptábis sacrifícium justítiæ, oblatiónes, et holocáusta: * tunc impónent super altáre tuum vítulos.

\antlabel{Ant.} Cor contrítum, et humiliátum, Deus, non despícies.

\antlabel{Ant.} Propter nomen tuum.

\psalmheading{Psalmus 142}

\lettrine[lines=2]{D}{OMINE}, exáudi oratiónem meam: áuribus pércipe obsecratiónem meam in veritáte tua: * exáudi me in tua justítia.

Et non intres in judícium cum servo tuo: * quia non justificábitur in conspéctu tuo omnis vivens.

Quia persecútus est inimícus ánimam meam: * humiliávit in terra vitam meam.

Collocávit me in obscúris sicut mórtuos sæculi: * et anxiátus est super me spíritus meus, in me turbátum est cor meum.

Memor fui diérum antiquórum, meditátus sum in ómnibus opéribus tuis: * in factis mánuum tuárum meditábar.

Expándi manus meas ad te: * ánima mea sicut terra sine aqua tibi.

Velóciter exáudi me, Dómine: * defécit spíritus meus.

Non avértas fáciem tuam a me: * et símilis ero descendéntibus in lacum.

Audítam fac mihi mane misericórdiam tuam: * quia in te sperávi.

Notam fac mihi viam, in qua ámbulem: * quia ad te levávi ánimam meam.

Eripe me de inimícis meis, Dómine, ad te confúgi: * doce me fácere voluntátem tuam, quia Deus meus es tu.

Spíritus tuus bonus dedúcet me in terram rectam: * propter nomen tuum, Dómine, vivificábis me, in æquitáte tua.

Edúces de tribulatióne ánimam meam: * et in misericórdia tua dispérdes inimícos meos.

Et perdes omnes, qui tríbulant ánimam meam: * quóniam ego servus tuus sum.

\antlabel{Ant.} Propter nomen tuum, Dómine, vivificábis me, in æquitáte tua.

\antlabel{Ant.} Deus, tu convérsus.

\psalmheading{Psalmus 84}

\lettrine[lines=2]{B}{ENEDIXISTI}, Dómine, terram tuam: * avertísti captivitátem Jacob.

Remisísti iniquitátem plebis tuæ: * operuísti ómnia peccáta eórum.

Mitigásti omnem iram tuam: * avertísti ab ira indignatiónis tuæ.

Convérte nos, Deus, salutáris noster: * et avérte iram tuam a nobis.

Numquid in ætérnum irascéris nobis? * aut exténdes iram tuam a generatióne in generatiónem?

Deus, tu convérsus vivificábis nos: * et plebs tua lætábitur in te.

Osténde nobis, Dómine, misericórdiam tuam: * et salutáre tuum da nobis.

Audiam quid loquátur in me Dóminus Deus: * quóniam loquétur pacem in plebem suam.

Et super sanctos suos: * et in eos, qui convertúntur ad cor.

Verúmtamen prope timéntibus eum salutáre ipsíus: * ut inhábitet glória in terra nostra.

Misericórdia, et véritas obviavérunt sibi: * justítia, et pax osculátæ sunt.

Véritas de terra orta est: * et justítia de cælo prospéxit.

Etenim Dóminus dabit benignitátem: * et terra nostra dabit fructum suum.

Justítia ante eum ambulábit: * et ponet in via gressus suos.

\antlabel{Ant.} Deus, tu convérsus vivificábis nos: et plebs tua lætábitur in te.

\antlabel{Ant.} Cum irátus fúeris.

\canticumheading{Canticum Habacuc}{Hab. 3, 2-19}

\lettrine[lines=2]{D}{OMINE}, audívi auditiónem tuam, * et tímui.

Dómine, opus tuum, * in médio annórum vivífica illud:

In médio annórum notum fácies: * cum irátus fúeris, misericórdiæ recordáberis.

Deus ab Austro véniet, * et sanctus de monte Pharan:

Opéruit cælos glória ejus: * et laudis ejus plena est terra.

Splendor ejus ut lux erit: * córnua in mánibus ejus:

Ibi abscóndita est fortitúdo ejus: * ante fáciem ejus ibit mors.

Et egrediétur diábolus ante pedes ejus. * Stetit, et mensus est terram.

Aspéxit, et dissólvit gentes: * et contríti sunt montes sæculi.

Incurváti sunt colles mundi, * ab itinéribus æternitátis ejus.

Pro iniquitáte vidi tentória Æthiópiæ, * turbabúntur pelles terræ Mádian.

Numquid in flumínibus irátus es, Dómine? * aut in flumínibus furor tuus? vel in mari indignátio tua?

Qui ascéndes super equos tuos: * et quadrígæ tuæ salvátio.

Súscitans suscitábis arcum tuum: * juraménta tríbubus quæ locútus es.

Flúvios scindes terræ: * vidérunt te, et doluérunt montes: gurges aquárum tránsiit.

Dedit abýssus vocem suam: * altitúdo manus suas levávit.

Sol, et luna stetérunt in habitáculo suo, * in luce sagittárum tuárum, ibunt in splendóre fulgurántis hastæ tuæ.

In frémitu conculcábis terram: * et in furóre obstupefácies gentes.

Egréssus es in salútem pópuli tui: * in salútem cum Christo tuo.

Percussísti caput de domo ímpii: * denudásti fundaméntum ejus usque ad collum.

Maledixísti sceptris ejus, cápiti bellatórum ejus, * veniéntibus ut turbo ad dispergéndum me.

Exsultátio eórum * sicut ejus, qui dévorat páuperem in abscóndito.

Viam fecísti in mari equis tuis, * in luto aquárum multárum.

Audívi, et conturbátus est venter meus: * a voce contremuérunt lábia mea.

Ingrediátur putrédo in óssibus meis, * et subter me scáteat.

Ut requiéscam in die tribulatiónis: * ut ascéndam ad pópulum accínctum nostrum.

Ficus enim non florébit: * et non erit germen in víneis.

Mentiétur opus olívæ: * et arva non áfferent cibum.

Abscindétur de ovíli pecus: * et non erit arméntum in præsépibus.

Ego autem in Dómino gaudébo: * et exsultábo in Deo Jesu meo.

Deus Dóminus fortitúdo mea: * et ponet pedes meos quasi cervórum.

Et super excélsa mea dedúcet me victor * in psalmis canéntem.

\antlabel{Ant.} Cum irátus fúeris, Dómine, misericórdiæ recordáberis.

\antlabel{Ant.} Lauda. \textdagger

\psalmheading{Psalmus 147}

\lettrine[lines=2]{L}{AUDA}, \textdagger\ Jerúsalem, Dóminum: * lauda Deum tuum, Sion.

Quóniam confortávit seras portárum tuárum: * benedíxit fíliis tuis in te.

Qui pósuit fines tuos pacem: * et ádipe fruménti sátiat te.

Qui emíttit elóquium suum terræ: * velóciter currit sermo ejus.

Qui dat nivem sicut lanam: * nébulam sicut cínerem spargit.

Mittit crystállum suam sicut buccéllas: * ante fáciem frígoris ejus quis sustinébit?

Emíttet verbum suum, et liquefáciet ea: * flabit spíritus ejus, et fluent aquæ.

Qui annúntiat verbum suum Jacob: * justítias, et judícia sua Israël.

Non fecit táliter omni natióni: * et judícia sua non manifestávit eis.

\antlabel{Ant.} Lauda Deum tuum, Sion, qui annúntiat judícia sua Israël.

\capitulumlabel{Capit.}{Rom. 13, 12-13}

\lettrine[lines=2]{N}{OX} præcéssit, dies autem appropinquávit. Abjiciámus ergo ópera tenebrárum, et induámur arma lucis. Sicut in die honéste ambulémus.

\hymnusheading{Hymnus}

\lettrine[lines=2]{Æ}{TERNA} cæli glória,\\
Beáta spes mortálium,\\
Summi Tonántis Únice,\\
Castǽque proles Vírginis:

Da déxteram surgéntibus,\\
Exsúrgat et mens sóbria,\\
Flagrans et in laudem Dei\\
Grates rependat débitas.

Ortus refúlget Lúcifer,\\
Præítque solem núntius:\\
Cadunt tenébræ nóctium:\\
Lux sancta nos illúminet.

Manénsque nostris sénsibus,\\
Noctem repéllat sæculi,\\
Omníque fine témporis\\
Purgáta servet péctora.

Quæsíta jam primum fides\\
In corde radíces agat:\\
Secúnda spes congáudeat,\\
Qua major exstat cáritas.

Deo Patri sit glória,\\
Ejúsque soli Fílio,\\
Cum Spíritu Paráclito,\\
Nunc et per omne sæculum.\\
Amen.

\V\ Repléti sumus mane misericórdia tua. \R\ Exsultávimus, et delectáti sumus.

\antlabel{Ad Bened. Ant.} Per víscera misericórdiæ * Dei nostri visitávit nos Oriens ex alto.

\rubric{Expleto Cántico ac repetíta Antíphona, dicúntur sequéntes}

\variantheading{Preces feriales}

KYRIE, eléison. Christe, eléison. Kýrie, eléison.

Pater noster \rubric{quod dícitur a solo Hebdomadário totum clara voce, usque ad} \V\ Et ne nos indúcas in tentatiónem. \R\ Sed líbera nos a malo. \V\ Ego dixi: Dómine, miserére mei. \R\ Sana ánimam meam, quia peccávi tibi. \V\ Convértere, Dómine, úsquequo? \R\ Et deprecábilis esto super servos tuos. \V\ Fiat misericórdia tua, Dómine, super nos. \R\ Quemádmodum sperávimus in te. \V\ Sacerdótes tui induántur justítiam. \R\ Et sancti tui exsúltent. \V\ Orémus pro beatíssimo Papa nostro N. \R\ Dóminus consérvet eum, et vivíficet eum, et beátum fáciat eum in terra, et non tradat eum in ánimam inimicórum ejus. \V\ Orémus et pro Antístite nostro N. \R\ Stet et pascat in fortitúdine tua, Dómine, in sublimitáte nóminis tui. \V\ Dómine, salvum fac regem. \R\ Et exáudi nos in die, qua invocavérimus te. \V\ Salvum fac pópulum tuum, Dómine, et bénedic hereditáti tuæ. \R\ Et rege eos, et extólle illos usque in ætérnum. \V\ Meménto Congregatiónis tuæ. \R\ Quam possedísti ab inítio. \V\ Fiat pax in virtúte tua. \R\ Et abundántia in túrribus tuis. \V\ Orémus pro benefactóribus nostris. \R\ Retribúere dignáre, Dómine, ómnibus, nobis bona faciéntibus propter nomen tuum, vitam ætérnam. Amen. \V\ Orémus pro fidélibus defúnctis. \R\ Réquiem ætérnam dona eis, Dómine, et lux perpétua lúceat eis. \V\ Requiéscant in pace. \R\ Amen. \V\ Pro frátribus nostris abséntibus. \R\ Salvos fac servos tuos, Deus meus, sperántes in te. \V\ Pro afflíctis et captívis. \R\ Líbera eos, Deus Israël, ex ómnibus tribulatiónibus suis. \V\ Mitte eis, Dómine, auxílium de sancto. \R\ Et de Sion tuére eos. \V\ Dómine, Deus virtútum, convérte nos. \R\ Et osténde fáciem tuam, et salvi érimus. \V\ Exsúrge, Christe, ádjuva nos. \R\ Et líbera nos propter nomen tuum. \V\ Dómine, exáudi oratiónem meam. \R\ Et clamor meus ad te véniat.

\rubric{Réliqua ut in Ordinário. 9}

% ===== END OF BATCH 32 -- Covers autumnalis.md lines 10505-10775,
% "Ad Laudes II" for Feria Sexta in full: Psalms 50, 142, 84, the
% Canticum Habacuc ("Dómine, audívi auditiónem tuam"), and Psalm
% 147, sharing the Capitulum "Nox præcéssit" and the hymn "Æterna
% cæli glória" with Ad Laudes I, plus the "Preces feriales" litany.
% Psalms 142 and 84 appear here without the dagger chant-pointing
% mark present in their Ad Laudes I occurrences (Batch 31), a
% genuine source divergence transcribed exactly as given in each
% place. The document continues with "Ad Primam" (autumnalis.md line
% 10778).
% =====
\hourheading{Ad Primam}

Pater, Ave et Credo.

\V\ Deus, in adjutórium.

\hymnusheading{Hymnus}

\lettrine[lines=2]{J}{AM} lucis orto sídere,\\
Deum precémur súpplices,\\
Ut in diúrnis áctibus\\
Nos servet a nocéntibus.

Linguam refrénans témperet,\\
Ne litis horror ínsonet:\\
Visum fovéndo cóntegat,\\
Ne vanitátes háuriat.

Sint pura cordis íntima,\\
Absístat et vecórdia:\\
Carnis terat supérbiam\\
Potus cibíque párcitas.

Ut, cum dies abscésserit\\
Noctémque sors redúxerit,\\
Mundi per abstinéntiam\\
Ipsi canámus glóriam.

Deo Patri sit glória,\\
Ejúsque soli Fílio,\\
Cum Spíritu Paráclito,\\
Nunc et per omne sæculum.\\
Amen.

\antlabel{Ant.} Ne discédas a me.

\psalmheading{Psalmus 21, i}

\lettrine[lines=2]{D}{EUS}, Deus meus, réspice in me: quare me dereliquísti? * longe a salúte mea verba delictórum meórum.

Deus meus, clamábo per diem, et non exáudies: * et nocte, et non ad insipiéntiam mihi.

Tu autem in sancto hábitas, * laus Israël.

In te speravérunt patres nostri: * speravérunt, et liberásti eos.

Ad te clamavérunt, et salvi facti sunt: * in te speravérunt, et non sunt confúsi.

Ego autem sum vermis, et non homo: * oppróbrium hóminum, et abjéctio plebis.

Omnes vidéntes me, derisérunt me: * locúti sunt lábiis, et movérunt caput.

Sperávit in Dómino, éripiat eum: * salvum fáciat eum, quóniam vult eum.

Quóniam tu es, qui extraxísti me de ventre: * spes mea ab ubéribus matris meæ. In te projéctus sum ex útero:

De ventre matris meæ Deus meus es tu, * ne discésseris a me.

Quóniam tribulátio próxima est: * quóniam non est qui ádjuvet.

\psalmheading{Psalmus 21, ii}

\lettrine[lines=2]{C}{IRCUMDEDERUNT} me vítuli multi: * tauri pingues obsedérunt me.

Aperuérunt super me os suum, * sicut leo rápiens et rúgiens.

Sicut aqua effúsus sum: * et dispérsa sunt ómnia ossa mea.

Factum est cor meum tamquam cera liquéscens * in médio ventris mei.

Aruit tamquam testa virtus mea, et lingua mea adhæsit fáucibus meis: * et in púlverem mortis deduxísti me.

Quóniam circumdedérunt me canes multi: * concílium malignántium obsédit me.

Fodérunt manus meas et pedes meos: * dinumeravérunt ómnia ossa mea.

Ipsi vero consideravérunt et inspexérunt me: * divisérunt sibi vestiménta mea, et super vestem meam misérunt sortem.

Tu autem, Dómine, ne elongáveris auxílium tuum a me: * ad defensiónem meam cónspice.

Erue a frámea, Deus, ánimam meam: * et de manu canis únicam meam:

Salva me ex ore leónis: * et a córnibus unicórnium humilitátem meam.

Narrábo nomen tuum frátribus meis: * in médio ecclésiæ laudábo te.

\psalmheading{Psalmus 21, iii}

\lettrine[lines=2]{Q}{UI} timétis Dóminum, laudáte eum: * univérsum semen Jacob, glorificáte eum.

Tímeat eum omne semen Israël: * quóniam non sprevit, neque despéxit deprecatiónem páuperis:

Nec avértit fáciem suam a me: * et cum clamárem ad eum, exaudívit me.

Apud te laus mea in ecclésia magna: * vota mea reddam in conspéctu timéntium eum.

Edent páuperes, et saturabúntur: et laudábunt Dóminum qui requírunt eum: * vivent corda eórum in sæculum sæculi.

Reminiscéntur et converténtur ad Dóminum * univérsi fines terræ:

Et adorábunt in conspéctu ejus * univérsæ famíliæ géntium.

Quóniam Dómini est regnum: * et ipse dominábitur géntium.

Manducavérunt et adoravérunt omnes pingues terræ: * in conspéctu ejus cadent omnes qui descéndunt in terram.

Et ánima mea illi vivet: * et semen meum sérviet ipsi.

Annuntiábitur Dómino generátio ventúra: * et annuntiábunt cæli justítiam ejus pópulo qui nascétur, quem fecit Dóminus.

\rubric{¶ Quando ad Laudes dictus fúerit Ps. 50 Miserére, hic subjúngitur sequens Ps. 98, qui in Laudibus locum non hábuit.}

\psalmheading{Psalmus 98}

\lettrine[lines=2]{D}{OMINUS} regnávit, irascántur pópuli: * qui sedet super Chérubim, moveátur terra.

Dóminus in Sion magnus: * et excélsus super omnes pópulos.

Confiteántur nómini tuo magno: quóniam terríbile, et sanctum est: * et honor regis judícium díligit.

Tu parásti directiónes: * judícium et justítiam in Jacob tu fecísti.

Exaltáte Dóminum, Deum nostrum, et adoráte scabéllum pedum ejus: * quóniam sanctum est.

Móyses et Aaron in sacerdótibus ejus: * et Sámuel inter eos, qui ínvocant nomen ejus:

Invocábant Dóminum, et ipse exaudiébat eos: * in colúmna nubis loquebátur ad eos.

Custodiébant testimónia ejus, * et præcéptum quod dedit illis.

Dómine, Deus noster, tu exaudiébas eos: * Deus, tu propítius fuísti eis, et ulcíscens in omnes adinventiónes eórum.

Exaltáte Dóminum, Deum nostrum, et adoráte in monte sancto ejus: * quóniam sanctus Dóminus, Deus noster.

\rubric{Secus vero, tribus Psalmis explétis, statim dícitur Ant.}

\antlabel{Ant.} Ne discédas a me, Dómine: quóniam tribulátio próxima est: quóniam non est qui ádjuvet.

\rubric{In Officio Festi vel Octávæ:}

\capitulumlabel{Capitulum}{1 Tim. 1, 17}

\lettrine[lines=2]{R}{EGI} sæculórum immortáli et invisíbili, soli Deo honor et glória in sæcula sæculórum. Amen.

\rubric{In ómnibus autem Fériis et in Vigíliis commúnibus:}

\capitulumlabel{Capitulum}{Zach. 8, 19}

\lettrine[lines=2]{P}{ACEM} et veritátem dilígite, ait Dóminus omnípotens.

\R\ br. Christe, Fili Dei vivi, * Miserére nobis. Christe. \V\ \textit{Qui sedes ad déxteram Patris.} Miserére nobis. Glória Patri. Christe. \V\ Exsúrge, Christe, ádjuva nos. \R\ Et líbera nos propter nomen tuum.

\variantheading{Preces dominicales et feriales (si dicendæ sint)}

KYRIE, eléison. Christe, eléison. Kýrie, eléison.

Pater noster \rubric{secreto usque ad} \V\ Et ne nos indúcas in tentatiónem. \R\ Sed líbera nos a malo.

Credo in Deum \rubric{secreto usque ad} \V\ Carnis resurrectiónem. \R\ Vitam ætérnam. Amen. \V\ Et ego ad te, Dómine, clamávi. \R\ Et mane orátio mea prævéniet te. \V\ Repleátur os meum laude. \R\ Ut cantem glóriam tuam, tota die magnitúdinem tuam. \V\ Dómine, avérte fáciem tuam a peccátis meis. \R\ Et omnes iniquitátes meas dele. \V\ Cor mundum crea in me, Deus. \R\ Et spíritum rectum ínnova in viscéribus meis. \V\ Ne projícias me a fácie tua. \R\ Et spíritum sanctum tuum ne áuferas a me. \V\ Redde mihi lætítiam salutáris tui. \R\ Et spíritu principáli confírma me.

\variantheading{In Précibus feriálibus tantum}

\V\ Eripe me, Dómine, ab hómine malo. \R\ A viro iníquo éripe me. \V\ Eripe me de inimícis meis, Deus meus. \R\ Et ab insurgéntibus in me líbera me. \V\ Eripe me de operántibus iniquitátem. \R\ Et de viris sánguinum salva me. \V\ Sic psalmum dicam nómini tuo in sæculum sæculi. \R\ Ut reddam vota mea de die in diem. \V\ Exáudi nos, Deus, salutáris noster. \R\ Spes ómnium fínium terræ, et in mari longe. \V\ Deus, in adjutórium meum inténde. \R\ Dómine, ad adjuvándum me festína. \V\ Sanctus Deus, Sanctus fortis, Sanctus immortális. \R\ Miserére nobis. \V\ Bénedic, ánima mea, Dómino. \R\ Et ómnia, quæ intra me sunt, nómini sancto ejus. \V\ Bénedic, ánima mea, Dómino. \R\ Et noli oblivísci omnes retributiónes ejus. \V\ Qui propitiátur ómnibus iniquitátibus tuis. \R\ Qui sanat omnes infirmitátes tuas. \V\ Qui rédimit de intéritu vitam tuam. \R\ Qui corónat te in misericórdia et miseratiónibus. \V\ Qui replet in bonis desidérium tuum. \R\ Renovábitur ut áquilæ juvéntus tua.

\variantheading{In utrisque Précibus}

\V\ Adjutórium nostrum in nómine Dómini. \R\ Qui fecit cælum et terram.

\formulalabel{Confíteor. Misereátur. Indulgéntiam.}

\V\ Dignáre, Dómine, die isto. \R\ Sine peccáto nos custodíre. \V\ Miserére nostri, Dómine. \R\ Miserére nostri. \V\ Fiat misericórdia tua, Dómine, super nos. \R\ Quemádmodum sperávimus in te. \V\ Dómine, exáudi oratiónem meam. \R\ Et clamor meus ad te véniat.

\V\ Dóminus vobíscum.

Orémus. \oratiolabel{Oratio}

\lettrine[lines=2]{D}{OMINE}, Deus omnípotens, qui ad princípium hujus diéi nos perveníre fecísti: tua nos hódie salva virtúte; ut in hac die ad nullum declinémus peccátum, sed semper ad tuam justítiam faciéndam nostra procédant elóquia, dirigántur cogitatiónes et ópera. Per Dóminum. \R\ Amen.

\V\ Dóminus vobíscum. \V\ Benedicámus Dómino.

\rubric{In Choro Martyrológium.}

\V\ Pretiósa in conspéctu Dómini. \R\ Mors Sanctórum ejus.

\rubric{Et absolúte, sine} Orémus:

\oratiolabel{Oratio}

\lettrine[lines=2]{S}{ANCTA} María, et omnes Sancti intercédant pro nobis ad Dóminum, ut nos mereámur ab eo adjuvári et salvári, qui vivit et regnat in sæcula sæculórum. \R\ Amen.

\V\ Deus, in adjutórium meum inténde. \R\ Dómine, ad adjuvándum me festína. \rubric{Et dícitur ter; último additur:} Glória Patri. Sicut erat.

Kýrie, eléison. Christe, eléison. Kýrie, eléison.

Pater noster \rubric{secreto usque ad} \V\ Et ne nos indúcas in tentatiónem. \R\ Sed líbera nos a malo. \V\ Réspice in servos tuos, Dómine, et in ópera tua, et dírige fílios eórum. \R\ Et sit splendor Dómini, Dei nostri, super nos, et ópera mánuum nostrárum dírige super nos, et opus mánuum nostrárum dírige. \V\ Glória Patri. \R\ Sicut erat.

Orémus. \oratiolabel{Oratio}

\lettrine[lines=2]{D}{IRIGERE} et sanctificáre, régere et gubernáre dignáre, Dómine Deus, Rex cæli et terræ, hódie corda et córpora nostra, sensus, sermónes et actus nostros in lege tua, et in opéribus mandatórum tuórum: ut hic, et in ætérnum, te auxiliánte, salvi et líberi esse mereámur, Salvátor mundi: Qui vivis et regnas in sæcula sæculórum. \R\ Amen.

\V\ Jube, domne.

\benedlabel{Bened.} Dies et actus nostros in sua pace dispónat Dóminus omnípotens. \R\ Amen.

\variantheading{In Officio feriali}

\lettrine[lines=2]{D}{OMINUS} autem dírigat corda et córpora nostra in caritáte Dei, et patiéntia Christi. Tu autem.

\variantheading{E Proprio Sanctorum}

\variantheading{Pro Oct. Omn. Sanctorum}

\lettrine[lines=2]{B}{ENEDÍCTIO} et cláritas et sapiéntia et gratiárum áctio, honor, virtus et fortitúdo Deo nostro in sǽcula sæculórum. Amen.

\variantheading{E Communi Sanctorum}

\variantheading{Pro Apostolis}

\lettrine[lines=2]{I}{BANT} Apóstoli gaudéntes a conspéctu concílii, quóniam digni hábiti sunt pro nómine Jesu contuméliam pati.

\variantheading{Pro Martyre}

\lettrine[lines=2]{J}{USTUS} cor suum trádidit ad vigilándum dilúculo ad Dóminum, qui fecit illum, et in conspéctu Altíssimi deprecábitur.

\variantheading{Pro pluribus Martyribus}

\lettrine[lines=2]{F}{ULGEBUNT} justi, et tamquam scintíllæ in arundinéto discúrrent. Judicábunt natiónes, et dominabúntur pópulis: et regnábit Dóminus illórum in perpétuum.

\variantheading{Pro Conf. Pont.}

\lettrine[lines=2]{F}{UNGI} sacerdótio, et habére laudem in nómine ipsíus, et offérre illi incénsum dignum in odórem suavitátis.

\variantheading{Pro Conf. non Pont.}

\lettrine[lines=2]{J}{USTUM} dedúxit Dóminus per vias rectas, et osténdit illi regnum Dei, et dedit illi sciéntiam sanctórum; honestávit illum in labóribus, et complévit labóres illíus.

\variantheading{Pro Virgine Martyre}

\lettrine[lines=2]{D}{OMINE}, Deus meus, exaltásti super terram habitatiónem meam, et pro morte defluénte deprecáta sum. Invocávi Dóminum, Patrem Dómini mei, ut non derelínquat me in die tribulatiónis meæ, et in témpore superbórum sine adjutório.

\variantheading{Pro Virgine non Martyre}

\lettrine[lines=2]{O}{} \textsc{quam} pulchra est casta generátio cum claritáte! immortális est enim memória illíus: quóniam apud Deum nota est, et apud hómines.

\variantheading{Pro Martyre tantum}

\lettrine[lines=2]{L}{AUDABIT} usque ad mortem ánima mea Dóminum, quóniam éruis sustinéntes te, et líberas eos de manu angústiæ, Dómine, Deus noster.

\variantheading{Pro nec Virg. nec Mart.}

\lettrine[lines=2]{M}{ULTAE} fíliæ congregavérunt divítias: tu supergréssa es univérsas. Fallax grátia, et vana est pulchritúdo: múlier timens Dóminum, ipsa laudábitur.

\variantheading{Pro Dedicatione Ecclesiæ}

\lettrine[lines=2]{E}{T} abstérget Deus omnem lácrimam ab óculis eórum: et mors ultra non erit, neque luctus, neque clamor, neque dolor erit ultra, quia prima abiérunt. Et dixit qui sedébat in throno: Ecce, nova fácio ómnia.

\variantheading{In Officiis B. Mariæ V.}

\lettrine[lines=2]{I}{N} platéis sicut cinnamómum et bálsamum aromatízans odórem dedi: quasi myrrha elécta dedi suavitátem odóris.

\V\ Adjutórium nostrum in nómine Dómini. \R\ Qui fecit cælum et terram.

\V\ Benedícite. \R\ Deus.

\benedictiolabel{Benedictio} Dóminus nos benedícat, et ab omni malo deféndat, et ad vitam perdúcat ætérnam. Et fidélium ánimæ per misericórdiam Dei requiéscant in pace. \R\ Amen.

\rubric{Deinde dícitur tantum} Pater noster \rubric{secréto, nisi sequátur ália Hora.}

% ===== END OF BATCH 33 -- Covers autumnalis.md lines 10778-11047,
% completing "Ad Primam" for Feria Sexta in full: the hymn "Jam
% lucis orto sídere," Psalms 21/i-iii, the conditional Psalm 98
% insertion, both Capitula, the brief Responsory, the full Preces
% apparatus, the closing Orationes, and the Lectio-brevis apparatus
% (identical category list to the previous four ferias). Several
% more psalm incipits carry the partial-caps OCR defect ("DEus,
% Deus meus," "CIrcumdedérunt," "DOminus regnávit," "PAcem"
% [Capitulum]) and are silently normalized to full caps. The
% document continues with "Ad Tertiam" (autumnalis.md line 11050).
% =====
\hourheading{Ad Tertiam}

Pater noster et Ave.

\V\ Deus, in adjutórium.

\hymnusheading{Hymnus}

\lettrine[lines=2]{N}{UNC}, Sancte, nobis, Spíritus,\\
Unum Patri cum Fílio,\\
Dignáre promptus íngeri\\
Nostro refúsus péctori.

Os, lingua, mens, sensus, vigor\\
Confessiónem pérsonent,\\
Flamméscat igne cáritas,\\
Accéndat ardor próximos.

Præsta, Pater piíssime,\\
Patríque compar Unice,\\
Cum Spíritu Paráclito\\
Regnans per omne sæculum.\\
Amen.

\antlabel{Ant.} Excita, Dómine.

\psalmheading{Psalmus 79, i}

\lettrine[lines=2]{Q}{UI} regis Israël, inténde: * qui dedúcis velut ovem Joseph.

Qui sedes super Chérubim, * manifestáre coram Ephraïm, Bénjamin, et Manásse.

Excita poténtiam tuam, et veni, * ut salvos fácias nos.

Deus, convérte nos: * et osténde fáciem tuam, et salvi érimus.

Dómine, Deus virtútum, * quoúsque irascéris super oratiónem servi tui?

Cibábis nos pane lacrimárum: * et potum dabis nobis in lácrimis in ménsura?

Posuísti nos in contradictiónem vicínis nostris: * et inimíci nostri subsannavérunt nos.

Deus virtútum, convérte nos: * et osténde fáciem tuam, et salvi érimus.

\psalmheading{Psalmus 79, ii}

\lettrine[lines=2]{V}{INEAM} de Ægýpto transtulísti: * ejecísti gentes, et plantásti eam.

Dux itíneris fuísti in conspéctu ejus: * plantásti radíces ejus, et implévit terram.

Opéruit montes umbra ejus: * et arbústa ejus cedros Dei.

Exténdit pálmites suos usque ad mare: * et usque ad flumen propágines ejus.

Ut quid destruxísti macériam ejus: * et vindémiant eam omnes, qui prætergrediúntur viam?

Exterminávit eam aper de silva: * et singuláris ferus depástus est eam.

Deus virtútum, convértere: * réspice de cælo, et vide, et vísita víneam istam.

Et pérfice eam, quam plantávit déxtera tua: * et super fílium hóminis, quem confirmásti tibi.

Incénsa igni, et suffóssa * ab increpatióne vultus tui períbunt.

Fiat manus tua super virum déxteræ tuæ: * et super fílium hóminis, quem confirmásti tibi.

Et non discédimus a te, vivificábis nos: * et nomen tuum invocábimus.

Dómine, Deus virtútum, convérte nos: * et osténde fáciem tuam, et salvi érimus.

\psalmheading{Psalmus 81}

\lettrine[lines=2]{D}{EUS} stetit in synagóga deórum: * in médio autem deos dijúdicat.

Usquequo judicátis iniquitátem: * et fácies peccatórum súmitis?

Judicáte egéno, et pupíllo: * húmilem, et páuperem justificáte.

Erípite páuperem: * et egénum de manu peccatóris liberáte.

Nesciérunt, neque intellexérunt, in ténebris ámbulant: * movebúntur ómnia fundaménta terræ.

Ego dixi: Dii estis, * et fílii Excélsi omnes.

Vos autem sicut hómines moriémini: * et sicut unus de princípibus cadétis.

Surge, Deus, júdica terram: * quóniam tu hereditábis in ómnibus géntibus.

\antlabel{Ant.} Excita, Dómine, poténtiam tuam, ut salvos fácias nos.

\variantheading{In Officio feriali}

\capitulumlabel{Capitulum}{Jer. 17, 14}

\lettrine[lines=2]{S}{ANA} me, Dómine, et sanábor: salvum me fac, et salvus ero: quóniam laus mea tu es.

\R\ br. Sana ánimam meam, * Quia peccávi tibi. Sana. \V\ Ego dixi: Dómine, miserére mei. Quia. Glória Patri. Sana.

\V\ Adjútor meus esto, ne derelínquas me. \R\ Neque despícias me, Deus, salutáris meus.

\variantheading{In Vigiliis communibus:}

\variantheading{Preces feriales}

KYRIE, eléison. Christe, eléison. Kýrie, eléison.

Pater noster \rubric{secreto usque ad} \V\ Et ne nos indúcas in tentatiónem. \R\ Sed líbera nos a malo. \V\ Dómine, Deus virtútum, convérte nos. \R\ Et osténde fáciem tuam, et salvi érimus. \V\ Exsúrge, Christe, ádjuva nos. \R\ Et líbera nos propter nomen tuum. \V\ Dómine, exáudi oratiónem meam. \R\ Et clamor meus ad te véniat.

\rubric{Réliqua ut in Ordinário. 19}

\hourheading{Ad Sextam}

Pater noster et Ave.

\V\ Deus, in adjutórium.

\hymnusheading{Hymnus}

\lettrine[lines=2]{R}{ECTOR} potens, verax Deus,\\
Qui témperas rerum vices,\\
Splendóre mane illúminas,\\
Et ígnibus merídiem:

Exstíngue flammas lítium,\\
Aufer calórem nóxium,\\
Confer salútem córporum,\\
Verámque pacem córdium.

Præsta, Pater piíssime,\\
Patríque compar Unice,\\
Cum Spíritu Paráclito\\
Regnans per omne sæculum.\\
Amen.

\antlabel{Ant.} Beáti, qui hábitant.

\psalmheading{Psalmus 83, i}

\lettrine[lines=2]{Q}{UAM} dilécta tabernácula tua, Dómine virtútum: * concupíscit, et déficit ánima mea in átria Dómini.

Cor meum, et caro mea * exsultavérunt in Deum vivum.

Etenim passer invénit sibi domum: * et turtur nidum sibi, ubi ponat pullos suos.

Altária tua, Dómine virtútum: * Rex meus, et Deus meus.

Beáti, qui hábitant in domo tua, Dómine: * in sæcula sæculórum laudábunt te.

Beátus vir, cujus est auxílium abs te: * ascensiónes in corde suo dispósuit in valle lacrimárum in loco, quem pósuit.

Etenim benedictiónem dabit legislátor, ibunt de virtúte in virtútem: * vidébitur Deus deórum in Sion.

\psalmheading{Psalmus 83, ii}

\lettrine[lines=2]{D}{OMINE}, Deus virtútum, exáudi oratiónem meam: * áuribus pércipe, Deus Jacob.

Protéctor noster, áspice, Deus: * et réspice in fáciem Christi tui:

Quia mélior est dies una in átriis tuis, * super míllia.

Elégi abjéctus esse in domo Dei mei: * magis quam habitáre in tabernáculis peccatórum.

Quia misericórdiam, et veritátem díligit Deus: * grátiam et glóriam dabit Dóminus.

Non privábit bonis eos, qui ámbulant in innocéntia: * Dómine virtútum, beátus homo, qui sperat in te.

\psalmheading{Psalmus 86}

\lettrine[lines=2]{F}{UNDAMENTA} ejus in móntibus sanctis: * díligit Dóminus portas Sion super ómnia tabernácula Jacob.

Gloriósa dicta sunt de te, * cívitas Dei.

Memor ero Rahab, et Babylónis * sciéntium me.

Ecce, alienígenæ, et Tyrus, et pópulus Æthíopum, * hi fuérunt illic.

Numquid Sion dicet: Homo, et homo natus est in ea: * et ipse fundávit eam Altíssimus?

Dóminus narrábit in scriptúris populórum, et príncipum: * horum, qui fuérunt in ea.

Sicut lætántium ómnium * habitátio est in te.

\antlabel{Ant.} Beáti, qui hábitant in domo tua, Dómine.

\variantheading{In Officio feriali}

\capitulumlabel{Capitulum}{Rom. 13, 8}

\lettrine[lines=2]{N}{EMINI} quidquam debeátis, nisi ut ínvicem diligátis: qui enim díligit próximum, legem implévit.

\R\ br. Benedícam Dóminum * In omni témpore. Benedícam. \V\ Semper laus ejus in ore meo. In. Glória Patri. Benedícam. \V\ Dóminus regit me, et nihil mihi déerit. \R\ In loco páscuæ ibi me collocávit.

\variantheading{In Vigiliis communibus:}

\variantheading{Preces feriales}

KYRIE, eléison. Christe, eléison. Kýrie, eléison.

Pater noster \rubric{secreto usque ad} \V\ Et ne nos indúcas in tentatiónem. \R\ Sed líbera nos a malo. \V\ Dómine, Deus virtútum, convérte nos. \R\ Et osténde fáciem tuam, et salvi érimus. \V\ Exsúrge, Christe, ádjuva nos. \R\ Et líbera nos propter nomen tuum. \V\ Dómine, exáudi oratiónem meam. \R\ Et clamor meus ad te véniat.

\rubric{Réliqua ut in Ordinário. 20}

\hourheading{Ad Nonam}

Pater noster et Ave.

\V\ Deus, in adjutórium.

\hymnusheading{Hymnus}

\lettrine[lines=2]{R}{ERUM}, Deus, tenax vigor,\\
Immótus in te pérmanens,\\
Lucis diúrnæ témpora\\
Succéssibus detérminans:

Largíre lumen véspere,\\
Quo vita nusquam décidat,\\
Sed prǽmium mortis sacræ\\
Perénnis instet glória.

Præsta, Pater piíssime,\\
Patríque compar Unice,\\
Cum Spíritu Paráclito\\
Regnans per omne sæculum.\\
Amen.

\antlabel{Ant.} Misericórdia et véritas.

\psalmheading{Psalmus 88, i}

\lettrine[lines=2]{M}{ISERICORDIAS} Dómini in ætérnum cantábo.

In generatiónem et generatiónem * annuntiábo veritátem tuam in ore meo.

Quóniam dixísti: In ætérnum misericórdia ædificábitur in cælis: * præparábitur véritas tua in eis.

Dispósui testaméntum eléctis meis, * jurávi David, servo meo:

Usque in ætérnum præparábo semen tuum. * Et ædificábo in generatiónem et generatiónem sedem tuam.

Confitebúntur cæli mirabília tua, Dómine: * étenim veritátem tuam in ecclésia sanctórum.

Quóniam quis in núbibus æquábitur Dómino: * símilis erit Deo in fíliis Dei?

Deus, qui glorificátur in consílio sanctórum: * magnus et terríbilis super omnes qui in circúitu ejus sunt.

Dómine, Deus virtútum, quis símilis tibi? * potens es, Dómine, et véritas tua in circúitu tuo.

Tu domináris potestáti maris: * motum autem flúctuum ejus tu mítigas.

Tu humiliásti sicut vulnerátum, supérbum: * in bráchio virtútis tuæ dispersísti inimícos tuos.

Tui sunt cæli, et tua est terra, orbem terræ et plenitúdinem ejus tu fundásti: * aquilónem, et mare tu creásti.

Thabor et Hermon in nómine tuo exsultábunt: * tuum bráchium cum poténtia.

Firmétur manus tua, et exaltétur déxtera tua: * justítia et judícium præparátio sedis tuæ.

Misericórdia et véritas præcédent fáciem tuam: * beátus pópulus, qui scit jubilatiónem.

Dómine, in lúmine vultus tui ambulábunt, et in nómine tuo exsultábunt tota die: * et in justítia tua exaltabúntur.

Quóniam glória virtútis eórum tu es: * et in beneplácito tuo exaltábitur cornu nostrum.

Quia Dómini est assúmptio nostra, * et Sancti Israël, regis nostri.

\psalmheading{Psalmus 88, ii}

\lettrine[lines=2]{T}{UNC} locútus es in visióne sanctis tuis, et dixísti: * Pósui adjutórium in poténte: et exaltávi eléctum de plebe mea.

Invéni David, servum meum: * óleo sancto meo unxi eum.

Manus enim mea auxiliábitur ei: * et bráchium meum confortábit eum.

Nihil profíciet inimícus in eo, * et fílius iniquitátis non appónet nocére ei.

Et concídam a fácie ipsíus inimícos ejus: * et odiéntes eum in fugam convértam.

Et véritas mea, et misericórdia mea cum ipso: * et in nómine meo exaltábitur cornu ejus.

Et ponam in mari manum ejus: * et in flumínibus déxteram ejus.

Ipse invocábit me: Pater meus es tu: * Deus meus, et suscéptor salútis meæ.

Et ego primogénitum ponam illum * excélsum præ régibus terræ.

In ætérnum servábo illi misericórdiam meam: * et testaméntum meum fidéle ipsi.

Et ponam in sæculum sæculi semen ejus: * et thronum ejus sicut dies cæli.

Si autem derelíquerint fílii ejus legem meam: * et in judíciis meis non ambuláverint:

Si justítias meas profanáverint: * et mandáta mea non custodíerint:

Visitábo in virga iniquitátes eórum: * et in verbéribus peccáta eórum.

Misericórdiam autem meam non dispérgam ab eo: * neque nocébo in veritáte mea:

Neque profanábo testaméntum meum: * et quæ procédunt de lábiis meis, non fáciam írrita.

Semel jurávi in sancto meo: Si David méntiar: * semen ejus in ætérnum manébit.

Et thronus ejus sicut sol in conspéctu meo, * et sicut luna perfécta in ætérnum: et testis in cælo fidélis.

\psalmheading{Psalmus 88, iii}

\lettrine[lines=2]{T}{U} vero repulísti et despexísti: * distulísti Christum tuum.

Evertísti testaméntum servi tui: * profanásti in terra Sanctuárium ejus.

Destruxísti omnes sepes ejus: * posuísti firmaméntum ejus formídinem.

Diripuérunt eum omnes transeúntes viam: * factus est oppróbrium vicínis suis.

Exaltásti déxteram deprimméntium eum: * lætificásti omnes inimícos ejus.

Avertísti adjutórium gládii ejus: * et non es auxiliátus ei in bello.

Destruxísti eum ab emundatióne: * et sedem ejus in terram collisísti.

Minorásti dies témporis ejus: * perfudísti eum confusióne.

Usquequo, Dómine, avértis in finem: * exardéscet sicut ignis ira tua?

Memoráre quæ mea substántia: * numquid enim vane constituísti omnes fílios hóminum?

Quis est homo, qui vivet, et non vidébit mortem: * éruet ánimam suam de manu ínferi?

Ubi sunt misericórdiæ tuæ antíquæ, Dómine, * sicut jurásti David in veritáte tua?

Memor esto, Dómine, oppróbrii servórum tuórum (quod contínui in sinu meo) multárum géntium.

Quod exprobravérunt inimíci tui, Dómine, * quod exprobravérunt commutatiónem Christi tui.

Benedíctus Dóminus in ætérnum: * fiat, fiat.

\antlabel{Ant.} Misericórdia et véritas præcédent fáciem tuam, Dómine.

\variantheading{In Officio feriali}

\capitulumlabel{Capit.}{1 Petri 1, 17-19}

\lettrine[lines=2]{I}{N} timóre incolátus vestri témpore conversámini: sciéntes quod non corruptibílibus auro vel argénto redémpti estis, sed pretióso sánguine quasi Agni immaculáti Christi.

\R\ br. Rédime me, Dómine, * Et miserére mei. Rédime. \V\ Pes enim meus stetit in via recta. Et miserére mei. Glória Patri. Rédime. \V\ Ab occúltis meis munda me, Dómine. \R\ Et ab aliénis parce servo tuo.

\variantheading{In Vigiliis communibus:}

\variantheading{Preces feriales}

KYRIE, eléison. Christe, eléison. Kýrie, eléison.

Pater noster \rubric{secreto usque ad} \V\ Et ne nos indúcas in tentatiónem. \R\ Sed líbera nos a malo. \V\ Dómine, Deus virtútum, convérte nos. \R\ Et osténde fáciem tuam, et salvi érimus. \V\ Exsúrge, Christe, ádjuva nos. \R\ Et líbera nos propter nomen tuum. \V\ Dómine, exáudi oratiónem meam. \R\ Et clamor meus ad te véniat.

\rubric{Réliqua ut in Ordinário. 22}

% ===== END OF BATCH 34 -- Covers autumnalis.md lines 11050-11408,
% completing the Little Hours of Feria Sexta in full: "Ad Tertiam"
% (Psalms 79/i, 79/ii, 81), "Ad Sextam" (Psalms 83/i, 83/ii, 86),
% and "Ad Nonam" (Psalms 88/i-iii), each with its own hymn,
% Capitulum, brief Responsory, and "In Vigiliis communibus: Preces
% feriales" apparatus. Three further occurrences of the recurring
% stray OCR "\*" artifact in the hour-hymns are silently dropped.
% Several more partial-caps OCR defects ("QUi," "VIneam," "DEus
% stetit," "QUam dilécta," "DOmine, Deus virtútum," "FUndaménta,"
% "MIsericórdias," "TUnc locútus," "TU vero") are silently
% normalized to full caps. The document continues with "Ad
% Vesperas" (autumnalis.md line 11412).
% =====
\hourheading{Ad Vesperas}

Pater noster et Ave.

\V\ Deus, in adjutórium.

\antlabel{Ant.} Dómine, * \textdagger\ probásti me, et cognovísti me. \textdaggerdbl

\psalmheading{Psalmus 138, i}

\lettrine[lines=2]{D}{OMINE}, \textdagger\ probásti me, et cognovísti me: * \textdaggerdbl\ tu cognovísti sessiónem meam, et resurrectiónem meam.

Intellexísti cogitatiónes meas de longe: * sémitam meam, et funículum meum investigásti.

Et omnes vias meas prævidísti: * quia non est sermo in lingua mea.

Ecce, Dómine, tu cognovísti ómnia novíssima, et antíqua: * tu formásti me, et posuísti super me manum tuam.

Mirábilis facta est sciéntia tua ex me: * confortáta est, et non pótero ad eam.

Quo ibo a spíritu tuo? * et quo a fácie tua fúgiam?

Si ascéndero in cælum, tu illic es: * si descéndero in inférnum, ades.

Si súmpsero pennas meas dilúculo, * et habitávero in extrémis maris:

Etenim illuc manus tua dedúcet me: * et tenébit me déxtera tua.

Et dixi: Fórsitan ténebræ conculcábunt me: * et nox illuminátio mea in delíciis meis.

Quia ténebræ non obscurabúntur a te, et nox sicut dies illuminábitur: * sicut ténebræ ejus, ita et lumen ejus.

Quia tu possedísti renes meos: * suscepísti me de útero matris meæ.

\antlabel{Ant.} Dómine, probásti me, et cognovísti me.

\antlabel{Ant.} Mirabília ópera tua, * Dómine, et ánima mea cognóscit nimis.

\psalmheading{Psalmus 138, ii}

\lettrine[lines=2]{C}{ONFITEBOR} tibi quia terribíliter magnificátus es: * mirabília ópera tua, et ánima mea cognóscit nimis.

Non est occultátum os meum a te, quod fecísti in occúlto: * et substántia mea in inferióribus terræ.

Imperféctum meum vidérunt óculi tui, et in libro tuo omnes scribéntur: * dies formabúntur, et nemo in eis.

Mihi autem nimis honorificáti sunt amíci tui, Deus: * nimis confortátus est principátus eórum.

Dinumerábo eos, et super arénam multiplicabúntur: * exsurréxi, et adhuc sum tecum.

Si occíderis, Deus, peccatóres: * viri sánguinum, declináte a me:

Quia dícitis in cogitatióne: * Accípient in vanitáte civitátes tuas.

Nonne qui odérunt te, Dómine, óderam? * et super inimícos tuos tabescébam?

Perfécto ódio óderam illos: * et inimíci facti sunt mihi.

Proba me, Deus, et scito cor meum: * intérroga me, et cognósce sémitas meas.

Et vide, si via iniquitátis in me est: * et deduc me in via ætérna.

\antlabel{Ant.} Mirabília ópera tua, Dómine, et ánima mea cognóscit nimis.

\antlabel{Ant.} Ne derelínquas me, * Dómine, virtus salútis meæ.

\psalmheading{Psalmus 139}

\lettrine[lines=2]{E}{RIPE} me, Dómine, ab hómine malo: * a viro iníquo éripe me.

Qui cogitavérunt iniquitátes in corde: * tota die constituébant prǽlia.

Acuérunt línguas suas sicut serpéntis: * venénum áspidum sub lábiis eórum.

Custódi me, Dómine, de manu peccatóris: * et ab homínibus iníquis éripe me.

Qui cogitavérunt supplantáre gressus meos: * absondérunt supérbi láqueum mihi:

Et funes extendérunt in láqueum: * juxta iter scándalum posuérunt mihi.

Dixi Dómino: Deus meus es tu: * exáudi, Dómine, vocem deprecatiónis meæ.

Dómine, Dómine, virtus salútis meæ: * obumbrásti super caput meum in die belli.

Ne tradas me, Dómine, a desidério meo peccatóri: * cogitavérunt contra me, ne derelínquas me, ne forte exalténtur.

Caput circúitus eórum: * labor labiórum ipsórum opériet eos.

Cadent super eos carbónes, in ignem dejícies eos: * in misériis non subsístent.

Vir linguósus non dirigétur in terra: * virum injústum mala cápient in intéritu.

Cognóvi quia fáciet Dóminus judícium ínopis: * et vindíctam páuperum.

Verúmtamen justi confitebúntur nómini tuo: * et habitábunt recti cum vultu tuo.

\antlabel{Ant.} Ne derelínquas me, Dómine, virtus salútis meæ.

\antlabel{Ant.} Dómine, * \textdagger\ clamávi ad te, exáudi me. \textdaggerdbl

\psalmheading{Psalmus 140}

\lettrine[lines=2]{D}{OMINE}, \textdagger\ clamávi ad te, exáudi me: * \textdaggerdbl\ inténde voci meæ, cum clamávero ad te.

Dirigátur orátio mea sicut incénsum in conspéctu tuo: * elevátio mánuum meárum sacrifícium vespertínum.

Pone, Dómine, custódiam ori meo: * et óstium circumstántiæ lábiis meis.

Non declínes cor meum in verba malítiæ, * ad excusándas excusatiónes in peccátis.

Cum homínibus operántibus iniquitátem: * et non communicábo cum eléctis eórum.

Corrípiet me justus in misericórdia, et increpábit me: * óleum autem peccatóris non impínguet caput meum.

Quóniam adhuc et orátio mea in beneplácitis eórum: * absórpti sunt juncti petræ júdices eórum.

Audient verba mea quóniam potuérunt: * sicut crassitúdo terræ erúpta est super terram.

Dissipáta sunt ossa nostra secus inférnum: * quia ad te, Dómine, Dómine, óculi mei: in te sperávi, non áuferas ánimam meam.

Custódi me a láqueo, quem statuérunt mihi: * et a scándalis operántium iniquitátem.

Cadent in retiáculo ejus peccatóres: * singuláriter sum ego donec tránseam.

\antlabel{Ant.} Dómine, clamávi ad te, exáudi me.

\antlabel{Ant.} Educ de custódia * ánimam meam, Dómine, ad confiténdum nómini tuo.

\psalmheading{Psalmus 141}

\lettrine[lines=2]{V}{OCE} mea ad Dóminum clamávi: * voce mea ad Dóminum deprecátus sum:

Effúndo in conspéctu ejus oratiónem meam, * et tribulatiónem meam ante ipsum pronúntio.

In deficiéndo ex me spíritum meum, * et tu cognovísti sémitas meas.

In via hac, qua ambulábam, * abscondérunt láqueum mihi.

Considerábam ad déxteram, et vidébam: * et non erat qui cognósceret me.

Périit fuga a me, * et non est qui requírat ánimam meam.

Clamávi ad te, Dómine, * dixi: Tu es spes mea, pórtio mea in terra vivéntium.

Inténde ad deprecatiónem meam: * quia humiliátus sum nimis.

Líbera me a persequéntibus me: * quia confortáti sunt super me.

Educ de custódia ánimam meam ad confiténdum nómini tuo: * me exspéctant justi, donec retríbuas mihi.

\antlabel{Ant.} Educ de custódia ánimam meam, Dómine, ad confiténdum nómini tuo.

\variantheading{In Officio feriali}

\capitulumlabel{Capitulum}{2 Cor. 1, 3-4}

\lettrine[lines=2]{B}{ENEDICTUS} Deus, et Pater Dómini nostri Jesu Christi, Pater misericordiárum, et Deus totíus consolatiónis, qui consolátur nos in omni tribulatióne nostra.

\hymnusheading{Hymnus}

\lettrine[lines=2]{H}{OMINIS} supérne Cónditor,\\
Qui cuncta solus órdinans,\\
Humum jubes prodúcere\\
Reptántis et feræ genus:

Et magna rerum córpora,\\
Dictu jubéntis vívida,\\
Per témporum certas vices\\
Obtemperáre sérvulis:

Repélle, quod cupídinis\\
Ciénte vi nos ímpetit,\\
Aut móribus se súggerit,\\
Aut áctibus se intérserit.

Da gaudiórum præmia,\\
Da gratiárum múnera:\\
Dissólve litis víncula:\\
Astrínge pacis fœdera.

Præsta, Pater piíssime,\\
Patríque compar Unice,\\
Cum Spíritu Paráclito\\
Regnans per omne sæculum.\\
Amen.

\V\ Dirigátur, Dómine, orátio mea. \R\ Sicut incénsum in conspéctu tuo.

\antlabel{Ad Magnif. Ant.} Depósuit Dóminus * poténtes de sede, et exaltávit húmiles.

\rubric{Réliqua ut in Ordinário. 22}

\hourheading{Ad Completorium}

\V\ Jube, domne.

\benedlabel{Bened.} Noctem quiétam, et finem perféctum concédat nobis Dóminus omnípotens. \R\ Amen.

\lectioheading{Lectio br.}{1 Petri 5, 8-9}

\lettrine[lines=2]{F}{RATRES}: Sóbrii estóte, et vigiláte: quia adversárius vester diábolus tamquam leo rúgiens círcuit, quærens quem dévoret: cui resístite fortes in fide. Tu autem, Dómine, miserére nobis. \R\ Deo grátias.

\V\ Adjutórium nostrum in nómine Dómini. \R\ Qui fecit cælum et terram.

Pater noster \rubric{quod dícitur totum secréto.}

\formulalabel{Confíteor.} Misereátur. Indulgéntiam.

\V\ Convérte nos, Deus, salutáris noster. \R\ Et avérte iram tuam a nobis.

\V\ Deus, in adjutórium.

\antlabel{Ant.} Voce mea. \textdagger

\psalmheading{Psalmus 76, i}

\lettrine[lines=2]{V}{OCE} mea \textdagger\ ad Dóminum clamávi: * voce mea ad Deum, et inténdit mihi.

In die tribulatiónis meæ Deum exquisívi, mánibus meis nocte contra eum: * et non sum decéptus.

Rénuit consolári ánima mea, * memor fui Dei, et delectátus sum, et exercitátus sum: et defécit spíritus meus.

Anticipavérunt vigílias óculi mei: * turbátus sum, et non sum locútus.

Cogitávi dies antíquos: * et annos ætérnos in mente hábui.

Et meditátus sum nocte cum corde meo, * et exercitábar, et scopébam spíritum meum.

Numquid in ætérnum projíciet Deus: * aut non appónet ut complacítior sit adhuc?

Aut in finem misericórdiam suam abscíndet, * a generatióne in generatiónem?

Aut obliviscétur miseréri Deus? * aut continébit in ira sua misericórdias suas?

Et dixi: Nunc cœpi: * hæc mutátio déxteræ Excélsi.

Memor fui óperum Dómini: * quia memor ero ab inítio mirabílium tuórum.

Et meditábor in ómnibus opéribus tuis: * et in adinventiónibus tuis exercébor.

\psalmheading{Psalmus 76, ii}

\lettrine[lines=2]{D}{EUS}, in sancto via tua: * quis Deus magnus sicut Deus noster? tu es Deus qui facis mirabília.

Notam fecísti in pópulis virtútem tuam: * redemísti in bráchio tuo pópulum tuum, fílios Jacob et Joseph.

Vidérunt te aquæ, Deus, vidérunt te aquæ: * et timuérunt, et turbátæ sunt abýssi.

Multitúdo sónitus aquárum: * vocem dedérunt nubes.

Etenim sagíttæ tuæ tránseunt: * vox tonítrui tui in rota.

Illuxérunt coruscatiónes tuæ orbi terræ: * commóta est, et contrémuit terra.

In mari via tua, et sémitæ tuæ in aquis multis: * et vestígia tua non cognoscéntur.

Deduxísti sicut oves pópulum tuum, * in manu Móysi et Aaron.

\psalmheading{Psalmus 85}

\lettrine[lines=2]{I}{NCLINA}, Dómine, aurem tuam, et exáudi me: * quóniam inops, et pauper sum ego.

Custódi ánimam meam, quóniam sanctus sum: * salvum fac servum tuum, Deus meus, sperántem in te.

Miserére mei, Dómine, quóniam ad te clamávi tota die: * lætífica ánimam servi tui, quóniam ad te, Dómine, ánimam meam levávi.

Quóniam tu, Dómine, suávis, et mitis: * et multæ misericórdiæ ómnibus invocántibus te.

Auribus pércipe, Dómine, oratiónem meam: * et inténde voci deprecatiónis meæ.

In die tribulatiónis meæ clamávi ad te: * quia exaudísti me.

Non est símilis tui in diis, Dómine: * et non est secúndum ópera tua.

Omnes gentes quascúmque fecísti, vénient, et adorábunt coram te, Dómine: * et glorificábunt nomen tuum.

Quóniam magnus es tu, et fáciens mirabília: * tu es Deus solus.

Deduc me, Dómine, in via tua, et ingrédiar in veritáte tua: * lætétur cor meum ut tímeat nomen tuum.

Confitébor tibi, Dómine, Deus meus, in toto corde meo, * et glorificábo nomen tuum in ætérnum:

Quia misericórdia tua magna est super me: * et eruísti ánimam meam ex inférno inferióri.

Deus, iníqui insurrexérunt super me, et synagóga poténtium quæsiérunt ánimam meam: * et non proposuérunt te in conspéctu suo.

Et tu, Dómine, Deus miserátor et miséricors, * pátiens, et multæ misericórdiæ, et verax.

Réspice in me, et miserére mei, * da impérium tuum púero tuo: et salvum fac fílium ancíllæ tuæ.

Fac mecum signum in bonum, ut vídeant qui odérunt me, et confundántur: * quóniam tu, Dómine, adjuvísti me, et consolátus es me.

\antlabel{Ant.} Voce mea ad Dóminum clamávi: neque obliviscétur miseréri Deus.

\hymnusheading{Hymnus}

\lettrine[lines=2]{T}{E} lucis ante términum,\\
Rerum Creátor, póscimus,\\
Ut pro tua cleméntia\\
Sis præsul et custódia.

Procul recédant sómnia,\\
Et nóctium phantásmata,\\
Hostémque nostrum cómprime,\\
Ne polluántur córpora.

Præsta, Pater piíssime,\\
Patríque compar Unice,\\
Cum Spíritu Paráclito\\
Regnans per omne sæculum.\\
Amen.

\capitulumlabel{Capitulum}{Jer. 14, 9}

\lettrine[lines=2]{T}{U} autem in nobis es, Dómine, et nomen sanctum tuum invocátum est super nos: ne derelínquas nos, Dómine, Deus noster.

\R\ br. In manus tuas, Dómine, * Comméndo spíritum meum. In manus. \V\ Redemísti nos, Dómine, Deus veritátis. Comméndo. Glória Patri. In manus.

\V\ Custódi nos, Dómine, ut pupíllam óculi. \R\ Sub umbra alárum tuárum prótege nos.

\antlabel{Ant.} Salva nos.

\canticumheading{Canticum Simeonis}{Luc. 2, 29-32}

\lettrine[lines=2]{N}{UNC} dimíttis servum tuum, Dómine, * secúndum verbum tuum in pace:

Quia vidérunt óculi mei * salutáre tuum,

Quod parásti * ante fáciem ómnium populórum,

Lumen ad revelatiónem géntium, * et glóriam plebis tuæ Israël.

Glória Patri.

\antlabel{Ant.} Salva nos, Dómine, vigilántes, custódi nos dormiéntes; ut vigilémus cum Christo, et requiescámus in pace.

\variantheading{Preces (si dicendæ sint)}

KYRIE, eléison. Christe, eléison. Kýrie, eléison.

Pater noster \rubric{secreto usque ad} \V\ Et ne nos indúcas in tentatiónem. \R\ Sed líbera nos a malo.

Credo in Deum \rubric{secreto usque ad} \V\ Carnis resurrectiónem. \R\ Vitam ætérnam. Amen. \V\ Benedíctus es, Dómine, Deus patrum nostrórum. \R\ Et laudábilis et gloriósus in sæcula. \V\ Benedicámus Patrem et Fílium cum Sancto Spíritu. \R\ Laudémus, et superexaltémus eum in sæcula. \V\ Benedíctus es, Dómine, in firmaménto cæli. \R\ Et laudábilis, et gloriósus, et superexaltátus in sæcula. \V\ Benedícat et custódiat nos omnípotens et miséricors Dóminus. \R\ Amen. \V\ Dignáre, Dómine, nocte ista. \R\ Sine peccáto nos custodíre. \V\ Miserére nostri, Dómine. \R\ Miserére nostri. \V\ Fiat misericórdia tua, Dómine, super nos. \R\ Quemádmodum sperávimus in te. \V\ Dómine, exáudi oratiónem meam. \R\ Et clamor meus ad te véniat.

\V\ Dóminus vobíscum.

Orémus. \oratiolabel{Oratio}

\lettrine[lines=2]{V}{ISITA}, quæsumus, Dómine, habitatiónem istam, et omnes insídias inimíci ab ea longe repélle: Angeli tui sancti hábitent in ea, qui nos in pace custódiant; et benedíctio tua sit super nos semper. Per Dóminum nostrum. \R\ Amen.

\V\ Dóminus vobíscum.

\V\ Benedicámus Dómino.

\benedictiolabel{Benedictio} Benedícat et custódiat nos omnípotens et miséricors Dóminus, Pater, et Fílius, et Spíritus Sanctus. \R\ Amen.

\variantheading{Antiphona finalis B. M. V.}

\V\ Divínum auxílium.

Pater noster, Ave María et Credo.

% ===== END OF BATCH 35 -- Covers autumnalis.md lines 11412-11776,
% completing "Ad Vesperas" (Psalms 138/i, 138/ii, 139, 140, 141,
% hymn "Hominis supérne Cónditor") and "Ad Completorium" (Psalms
% 76/i, 76/ii, 85, hymn "Te lucis ante términum," Canticum
% Simeonis, Preces), and with them the ENTIRE "Feria Sexta" Office.
% Several more partial-caps OCR defects ("COnfitébor," "ERipe me,"
% "VOce mea" twice, "DEus, in sancto," "INclina, Dómine,"
% "BEnedíctus" [Capitulum], "HOminis supérne" [hymn]) are silently
% normalized to full caps. The closing cross-reference reverts to
% the abbreviated "Antiphona finalis B. M. V." but keeps the fuller
% "Pater noster, Ave María et Credo." seen since Feria Quinta; each
% transcribed exactly as its own occurrence gives it. The document
% continues with "Sabbato" (autumnalis.md line 11780), the final
% weekday of the Psalterium cycle.
% =====
\feriaheading{Sabbato}

\hourheading{Ad Matutinum}

Pater, Ave et Credo.

\V\ Dómine, lábia mea.

\V\ Deus, in adjutórium.

\variantheading{In Officio feriali}

\invitatlabel{Invit.} Pópulus Dómini, et oves páscuæ ejus: * Veníte, adorémus. \psalmlabel{Ps. 94} Veníte.

\hymnusheading{Hymnus}

\lettrine[lines=2]{S}{UMMAE} Parens cleméntiæ,\\
Mundi regis qui máchinam,\\
Uníus et substántiæ,\\
Trinúsque persónis Deus:

Nostros pius cum cánticis\\
Fletus benígne súscipe:\\
Ut corde puro sórdium\\
Te perfruámur lárgius.

Lumbos jecúrque mórbidum\\
Flammis adúre cóngruis,\\
Accíncti ut artus éxcubent,\\
Luxu remóto péssimo.

Quicúmque ut horas nóctium\\
Nunc concinéndo rúmpimus,\\
Ditémur omnes áffatim\\
Donis beátæ pátriæ.

Præsta, Pater piíssime,\\
Patríque compar Unice,\\
Cum Spíritu Paráclito\\
Regnans per omne sæculum.\\
Amen.

\nocturnoheading{In I Nocturno}

\antlabel{Ant.} Memor fuit in sæculum * testaménti sui Dóminus, Deus noster.

\psalmheading{Psalmus 104, i}

\lettrine[lines=2]{C}{ONFITEMINI} Dómino, et invocáte nomen ejus: * annuntiáte inter gentes ópera ejus.

Cantáte ei, et psállite ei: * narráte ómnia mirabília ejus.

Laudámini in nómine sancto ejus: * lætétur cor quæréntium Dóminum.

Quærite Dóminum, et confirmámini: * quærite fáciem ejus semper.

Mementóte mirabílium ejus, quæ fecit: * prodígia ejus, et judícia oris ejus.

Semen Abraham, servi ejus: * fílii Jacob, elécti ejus.

Ipse Dóminus Deus noster: * in univérsa terra judícia ejus.

Memor fuit in sæculum testaménti sui: * verbi, quod mandávit in mille generatiónes:

Quod dispósuit ad Abraham: * et juraménti sui ad Isaac:

Et státuit illud Jacob in præcéptum: * et Israël in testaméntum ætérnum:

Dicens: Tibi dabo terram Chánaan, * funículum hereditátis vestræ.

Cum essent número brevi, * paucíssimi et íncolæ ejus:

Et pertransiérunt de gente in gentem, * et de regno ad pópulum álterum.

Non relíquit hóminem nocére eis: * et corrípuit pro eis reges.

Nolíte tángere christos meos: * et in prophétis meis nolíte malignári.

\antlabel{Ant.} Memor fuit in sæculum testaménti sui Dóminus, Deus noster.

\antlabel{Ant.} Auxit Dóminus * pópulum suum: et firmávit eum super inimícos ejus.

\psalmheading{Psalmus 104, ii}

\lettrine[lines=2]{E}{T} vocávit famem super terram: * et omne firmaméntum panis contrívit.

Misit ante eos virum: * in servum venúmdatus est Joseph.

Humiliavérunt in compédibus pedes ejus, ferrum pertránsiit ánimam ejus, * donec veníret verbum ejus.

Elóquium Dómini inflammávit eum: * misit rex, et solvit eum; princeps populórum, et dimísit eum.

Constítuit eum dóminum domus suæ: * et príncipem omnis possessiónis suæ:

Ut erudíret príncipes ejus sicut semetípsum: * et senes ejus prudéntiam docéret.

Et intrávit Israël in Ægýptum: * et Jacob áccola fuit in terra Cham.

Et auxit pópulum suum veheménter: * et firmávit eum super inimícos ejus.

Convértit cor eórum ut odírent pópulum ejus: * et dolum fácerent in servos ejus.

Misit Móysen, servum suum: * Aaron, quem elégit ipsum.

Pósuit in eis verba signórum suórum, * et prodigiórum in terra Cham.

\antlabel{Ant.} Auxit Dóminus pópulum suum: et firmávit eum super inimícos ejus.

\antlabel{Ant.} Edúxit Deus * pópulum suum in exsultatióne, et eléctos suos in lætítia.

\psalmheading{Psalmus 104, iii}

\lettrine[lines=2]{M}{ISIT} ténebras, et obscurávit: * et non exacerbávit sermónes suos.

Convértit aquas eórum in sánguinem: * et occídit pisces eórum.

Edidit terra eórum ranas * in penetrálibus regum ipsórum.

Dixit, et venit cœnomýia: * et cínifes in ómnibus fínibus eórum.

Pósuit plúvias eórum grándinem: * ignem comburéntem in terra ipsórum.

Et percússit víneas eórum, et ficúlneas eórum: * et contrívit lignum fínium eórum.

Dixit, et venit locústa, et bruchus, * cujus non erat númerus:

Et comédit omne fœnum in terra eórum: * et comédit omnem fructum terræ eórum.

Et percússit omne primogénitum in terra eórum: * primítias omnis labóris eórum.

Et edúxit eos cum argénto et auro: * et non erat in tríbubus eórum infírmus.

Lætáta est Ægýptus in profectióne eórum: * quia incúbuit timor eórum super eos.

Expándit nubem in protectiónem eórum, * et ignem ut lucéret eis per noctem.

Petiérunt, et venit cotúrnix: * et pane cæli saturávit eos.

Dirúpit petram, et fluxérunt aquæ: * abiérunt in sicco flúmina;

Quóniam memor fuit verbi sancti sui, * quod hábuit ad Abraham, púerum suum.

Et edúxit pópulum suum in exsultatióne, * et eléctos suos in lætítia.

Et dedit illis regiónes géntium: * et labóres populórum possedérunt:

Ut custódiant justificatiónes ejus, * et legem ejus requírant.

\antlabel{Ant.} Edúxit Deus pópulum suum in exsultatióne, et eléctos suos in lætítia.

\variantheading{In Festis novem Lectionum:}

\V\ Memor fuit Dóminus verbi sancti tui. \R\ Quod locútus est ad Abraham, púerum suum.

\absolutiolabel{Absolutio.} Exáudi, Dómine.

\benedlabel{Bened.} 1 Benedictióne perpétua. 2 Unigénitus Dei Fílius. 3 Spíritus Sancti grátia.

\nocturnoheading{In II Nocturno}

\antlabel{Ant.} Salvávit eos Dóminus * propter nomen suum.

\psalmheading{Psalmus 105, i}

\lettrine[lines=2]{C}{ONFITEMINI} Dómino quóniam bonus: * quóniam in sæculum misericórdia ejus.

Quis loquétur poténtias Dómini, * audítas fáciet omnes laudes ejus?

Beáti, qui custódiunt judícium, * et fáciunt justítiam in omni témpore.

Meménto nostri, Dómine, in beneplácito pópuli tui: * vísita nos in salutári tuo:

Ad vidéndum in bonitáte electórum tuórum, ad lætándum in lætítia gentis tuæ: * ut laudéris cum hereditáte tua.

Peccávimus cum pátribus nostris: * injúste égimus, iniquitátem fécimus.

Patres nostri in Ægýpto non intellexérunt mirabília tua: * non fuérunt mémores multitúdinis misericórdiæ tuæ.

Et irritavérunt ascendéntes in mare, * Mare Rubrum.

Et salvávit eos propter nomen suum: * ut notam fáceret poténtiam suam.

Et incrépuit Mare Rubrum, et exsiccátum est: * et dedúxit eos in abýssis sicut in desérto.

Et salvávit eos de manu odiéntium: * et redémit eos de manu inimíci.

Et opéruit aqua tribulántes eos: * unus ex eis non remánsit.

Et credidérunt verbis ejus: * et laudavérunt laudem ejus.

Cito fecérunt, oblíti sunt óperum ejus: * et non sustinuérunt consílium ejus.

Et concupiérunt concupiscéntiam in desérto: * et tentavérunt Deum in inaquóso.

Et dedit eis petitiónem ipsórum: * et misit saturitátem in ánimas eórum.

\antlabel{Ant.} Salvávit eos Dóminus propter nomen suum.

\antlabel{Ant.} Oblíti sunt Deum, * qui salvávit eos.

\psalmheading{Psalmus 105, ii}

\lettrine[lines=2]{E}{T} irritavérunt Móysen in castris: * Aaron, sanctum Dómini.

Apérta est terra, et deglutívit Dathan: * et opéruit super congregatiónem Abíron.

Et exársit ignis in synagóga eórum: * flamma combússit peccatóres.

Et fecérunt vítulum in Horeb: * et adoravérunt scúlptile.

Et mutavérunt glóriam suam * in similitúdinem vítuli comedéntis fœnum.

Oblíti sunt Deum, qui salvávit eos, * qui fecit magnália in Ægýpto, mirabília in terra Cham: terribília in Mari Rubro.

Et dixit ut dispérderet eos: * si non Móyses, eléctus ejus, stetísset in confractióne in conspéctu ejus:

Ut avérteret iram ejus ne dispérderet eos: * et pro níhilo habuérunt terram desiderábilem:

Non credidérunt verbo ejus, et murmuravérunt in tabernáculis suis: * non exaudiérunt vocem Dómini.

Et elevávit manum suam super eos: * ut prostérneret eos in desérto:

Et ut dejíceret semen eórum in natiónibus: * et dispérgeret eos in regiónibus.

Et initiáti sunt Beélphegor: * et comedérunt sacrifícia mortuórum.

Et irritavérunt eum in adinventiónibus suis: * et multiplicáta est in eis ruína.

Et stetit Phínees, et placávit: * et cessávit quassátio.

Et reputátum est ei in justítiam, * in generatiónem et generatiónem usque in sempitérnum.

\antlabel{Ant.} Oblíti sunt Deum, qui salvávit eos.

\antlabel{Ant.} Cum tribularéntur, * vidit Dóminus: et audívit oratiónem eórum.

\psalmheading{Psalmus 105, iii}

\lettrine[lines=2]{E}{T} irritavérunt eum ad Aquas contradictiónis: * et vexátus est Móyses propter eos: quia exacerbavérunt spíritum ejus.

Et distínxit in lábiis suis: * non disperdidérunt gentes, quas dixit Dóminus illis.

Et commísti sunt inter gentes, et didicérunt ópera eórum: et serviérunt sculptílibus eórum: * et factum est illis in scándalum.

Et immolavérunt fílios suos, * et fílias suas dæmóniis.

Et effudérunt sánguinem innocéntem: * sánguinem filiórum suórum et filiárum suárum, quas sacrificavérunt sculptílibus Chánaan.

Et infécta est terra in sanguínibus, et contamináta est in opéribus eórum: * et fornicáti sunt in adinventiónibus suis.

Et irátus est furóre Dóminus in pópulum suum: * et abominátus est hereditátem suam.

Et trádidit eos in manus géntium: * et domináti sunt eórum qui odérunt eos.

Et tribulavérunt eos inimíci eórum, et humiliáti sunt sub mánibus eórum: * sæpe liberávit eos.

Ipsi autem exacerbavérunt eum in consílio suo: * et humiliáti sunt in iniquitátibus suis.

Et vidit, cum tribularéntur: * et audívit oratiónem eórum.

Et memor fuit testaménti sui: * et pœnítuit eum secúndum multitúdinem misericórdiæ suæ.

Et dedit eos in misericórdias * in conspéctu ómnium qui céperant eos.

Salvos nos fac, Dómine, Deus noster: * et cóngrega nos de natiónibus:

Ut confiteámur nómini sancto tuo: * et gloriémur in laude tua.

Benedíctus Dóminus Deus Israël, a sæculo et usque in sæculum: * et dicet omnis pópulus: Fiat, fiat.

\antlabel{Ant.} Cum tribularéntur, vidit Dóminus: et audívit oratiónem eórum.

\variantheading{In Officiis novem Lectionum:}

\V\ Salvos nos fac, Dómine, Deus noster. \R\ Ut confiteámur nómini sancto tuo.

\absolutiolabel{Absolutio.} Ipsíus píetas et.

\benedlabel{Bened.} 4 Deus Pater omnípotens. 5 Christus perpétuæ. 6 Ignem sui amóris.

\nocturnoheading{In III Nocturno}

\antlabel{Ant.} Clamavérunt ad Dóminum: * et de necessitátibus eórum liberávit eos.

\psalmheading{Psalmus 106, i}

\lettrine[lines=2]{C}{ONFITEMINI} Dómino quóniam bonus: * quóniam in sæculum misericórdia ejus.

Dicant qui redémpti sunt a Dómino, quos redémit de manu inimíci: * et de regiónibus congregávit eos:

A solis ortu, et occásu: * ab aquilóne, et mari.

Erravérunt in solitúdine in inaquóso: * viam civitátis habitáculi non invenérunt.

Esuriéntes, et sitiéntes: * ánima eórum in ipsis defécit.

Et clamavérunt ad Dóminum cum tribularéntur: * et de necessitátibus eórum erípuit eos.

Et dedúxit eos in viam rectam: * ut irent in civitátem habitatiónis.

Confiteántur Dómino misericórdiæ ejus: * et mirabília ejus fíliis hóminum.

Quia satiávit ánimam inánem: * et ánimam esuriéntem satiávit bonis.

Sedéntes in ténebris, et umbra mortis: * vinctos in mendicitáte et ferro.

Quia exacerbavérunt elóquia Dei: * et consílium Altíssimi irritavérunt.

Et humiliátum est in labóribus cor eórum: * infirmáti sunt, nec fuit qui adjuváret.

Et clamavérunt ad Dóminum cum tribularéntur: * et de necessitátibus eórum liberávit eos.

Et edúxit eos de ténebris, et umbra mortis: * et víncula eórum disrúpit.

\antlabel{Ant.} Clamavérunt ad Dóminum: et de necessitátibus eórum liberávit eos.

\antlabel{Ant.} Ipsi vidérunt * ópera Dei, et mirabília ejus.

\psalmheading{Psalmus 106, ii}

\lettrine[lines=2]{C}{ONFITEANTUR} Dómino misericórdiæ ejus: * et mirabília ejus fíliis hóminum.

Quia contrívit portas áereas: * et vectes férreos confrégit.

Suscépit eos de via iniquitátis eórum: * propter injustítias enim suas humiliáti sunt.

Omnem escam abomináta est ánima eórum: * et appropinquavérunt usque ad portas mortis.

Et clamavérunt ad Dóminum cum tribularéntur: * et de necessitátibus eórum liberávit eos.

Misit verbum suum, et sanávit eos: * et erípuit eos de interitiónibus eórum.

Confiteántur Dómino misericórdiæ ejus: * et mirabília ejus fíliis hóminum.

Et sacríficent sacrifícium laudis: * et annúntient ópera ejus in exsultatióne.

Qui descéndunt mare in návibus, * faciéntes operatiónem in aquis multis.

Ipsi vidérunt ópera Dómini, * et mirabília ejus in profúndo.

Dixit, et stetit spíritus procéllæ: * et exaltáti sunt fluctus ejus.

Ascéndunt usque ad cælos, et descéndunt usque ad abýssos: * ánima eórum in malis tabescébat.

Turbáti sunt, et moti sunt sicut ébrius: * et omnis sapiéntia eórum devoráta est.

Et clamavérunt ad Dóminum cum tribularéntur: * et de necessitátibus eórum edúxit eos.

Et státuit procéllam ejus in auram: * et siluérunt fluctus ejus.

Et lætáti sunt quia siluérunt: * et dedúxit eos in portum voluntátis eórum.

\antlabel{Ant.} Ipsi vidérunt ópera Dei, et mirabília ejus.

\antlabel{Ant.} Vidébunt recti, * et lætabúntur, et intélligent misericórdias Dómini.

\psalmheading{Psalmus 106, iii}

\lettrine[lines=2]{C}{ONFITEANTUR} Dómino misericórdiæ ejus: * et mirabília ejus fíliis hóminum.

Et exáltent eum in ecclésia plebis: * et in cáthedra seniórum laudent eum.

Pósuit flúmina in desértum: * et éxitus aquárum in sitim.

Terram fructíferam in salsúginem, * a malítia inhabitántium in ea.

Pósuit desértum in stagna aquárum: * et terram sine aqua in éxitus aquárum.

Et collocávit illic esuriéntes: * et constituérunt civitátem habitatiónis.

Et seminavérunt agros, et plantavérunt víneas: * et fecérunt fructum nativitátis.

Et benedíxit eis, et multiplicáti sunt nimis: * et juménta eórum non minorávit.

Et pauci facti sunt: * et vexáti sunt a tribulatióne malórum, et dolóre.

Effúsa est contémptio super príncipes: * et erráre fecit eos in ínvio, et non in via.

Et adjúvit páuperem de inópia: * et pósuit sicut oves famílias.

Vidébunt recti, et lætabúntur: * et omnis iníquitas oppilábit os suum.

Quis sápiens et custódiet hæc? * et intélliget misericórdias Dómini?

\antlabel{Ant.} Vidébunt recti, et lætabúntur, et intélligent misericórdias Dómini.

\variantheading{In feriali Officio in Officio sanctæ Mariæ in Sabbato et in Festis:}

\V\ Exáltent Dóminum in ecclésia plebis. \R\ Et in cáthedra seniórum laudent eum.

\variantheading{In Officio novem Lectionum}

\absolutiolabel{Absolutio.} A vínculis.

\benedlabel{Bened.} 7 Evangélica léctio. 8 \rubric{In Festis Domini:} Divínum auxílium. \rubric{In Festis B. M. V.:} Cujus festum cólimus, ipsa Virgo. \rubric{In Festis Sanctorum:} Cujus (\textit{vel} Quorum \textit{aut} Quarum) festum cólimus. 9 Ad societátem. (Per evangélica dicta.)

\variantheading{In Officio trium Lectionum}

\absolutiolabel{Absolutio.} A vínculis.

\variantheading{In Feriis cum Homilia:}

\benedlabel{Bened.} 1 Evangélica. 2 Divínum. 3 Ad societátem.

\variantheading{In Octava simplici:}

\benedlabel{Bened.} 1 Ille nos benedícat. 2 Cujus (\textit{vel} Quorum \textit{aut} Quarum). 3 Ad societátem.

\variantheading{In Officio sanctæ Mariæ in Sabbato}

\absolutiolabel{Absol.} Précibus et méritis beátæ Maríæ semper Vírginis et ómnium Sanctórum, perdúcat nos Dóminus ad regna cælórum. Amen.

\benedlabel{Bened.} 1 Nos cum prole pia benedícat Virgo María. Amen.

2 Ipsa Virgo vírginum intercédat pro nobis ad Dóminum. Amen.

3 Per Vírginem Matrem concédat nobis Dóminus salútem et pacem. Amen.

% ===== END OF BATCH 36 -- Covers autumnalis.md lines 11780-12200,
% opening "Sabbato" (Saturday) and completing its "Ad Matutinum" in
% full: the ferial Invitatorium "Pópulus Dómini, et oves páscuæ
% ejus," the hymn "Summæ Parens cleméntiæ," and the three
% Nocturns -- Psalm 104 in the first, Psalm 105 in the second,
% Psalm 106 in the third -- with their versicle, Absolutio, and
% Benedictiones apparatus. The ninth-lesson "Bened." line (line
% 12176) marks its three conditional clauses ("In Festis Domini:",
% "In Festis B. M. V.:", "In Festis Sanctorum:") as actual red
% rubric text in the source this time ({RUBRIC: ...}), unlike the
% plain black-text embedding of the identical clauses at every
% prior occurrence of this line this volume; this is rendered with
% \rubric{} on each of the three labels, preserving the genuine
% embedded-rubric distinction exactly as the source marks it here.
% Saturday Matins closes with a Sabbath-specific apparatus unique
% to this day: "In Officio sanctæ Mariæ in Sabbato," with its own
% full-sentence Absolutio ("Précibus et méritis beátæ Maríæ...")
% and a three-part Marian Benedictio given as three separate
% paragraphs in the source (each ending its own "Amen."), reproduced
% here with the same paragraph breaks. Several more partial-caps
% OCR defects ("SUmmæ," "COnfitémini"/"COnfiteántur" repeatedly,
% "ET irritavérunt" twice, "MIsit ténebras") are silently normalized
% to full caps. The document continues with "Ad Laudes I"
% (autumnalis.md line 12202).
% =====
\hourheading{Ad Laudes I}

\rubric{In Sabbatis extra Vigílias commúnes occurréntibus, in Festis novem Lectiónum, in ómnibus Officiis de Octáva non privilegiáta, atque in Officio sanctæ Maríæ in Sabbato.}

(Pater noster et Ave.)

\V\ Deus, in adjutórium.

\antlabel{Ant.} Fílii Sion * exsúltent in Rege suo.

\psalmheading{Psalmus 149}

\lettrine[lines=2]{C}{ANTATE} Dómino cánticum novum: * laus ejus in ecclésia sanctórum.

Lætétur Israël in eo, qui fecit eum: * et fílii Sion exsúltent in rege suo.

Laudent nomen ejus in choro: * in týmpano, et psaltério psallant ei:

Quia beneplácitum est Dómino in pópulo suo: * et exaltábit mansuétos in salútem.

Exsultábunt sancti in glória: * lætabúntur in cubílibus suis.

Exaltatiónes Dei in gútture eórum: * et gládii ancípites in mánibus eórum.

Ad faciéndam vindíctam in natiónibus: * increpatiónes in pópulis.

Ad alligándos reges eórum in compédibus: * et nóbiles eórum in mánicis férreis.

Ut fáciant in eis judícium conscríptum: * glória hæc est ómnibus sanctis ejus.

\antlabel{Ant.} Fílii Sion exsúltent in Rege suo.

\antlabel{Ant.} Quam magnificáta * sunt ópera tua, Dómine!

\psalmheading{Psalmus 91}

\lettrine[lines=2]{B}{ONUM} est confitéri Dómino: * et psállere nómini tuo, Altíssime.

Ad annuntiándum mane misericórdiam tuam: * et veritátem tuam per noctem.

In decachórdo, psaltério: * cum cántico, in cíthara.

Quia delectásti me, Dómine, in factúra tua: * et in opéribus mánuum tuárum exsultábo.

Quam magnificáta sunt ópera tua, Dómine! * nimis profúndæ factæ sunt cogitatiónes tuæ.

Vir insípiens non cognóscet: * et stultus non intélliget hæc.

Cum exórti fúerint peccatóres sicut fœnum: * et apparúerint omnes, qui operántur iniquitátem:

Ut intéreant in sæculum sæculi: * tu autem Altíssimus in ætérnum, Dómine.

Quóniam ecce inimíci tui, Dómine, quóniam ecce inimíci tui períbunt: * et dispergéntur omnes, qui operántur iniquitátem.

Et exaltábitur sicut unicórnis cornu meum: * et senéctus mea in misericórdia úberi.

Et despéxit óculus meus inimícos meos: * et in insurgéntibus in me malignántibus áudiet auris mea.

Justus ut palma florébit: * sicut cedrus Líbani multiplicábitur.

Plantáti in domo Dómini, * in átriis domus Dei nostri florébunt.

Adhuc multiplicabúntur in senécta úberi: * et patiéntes erunt, ut annúntient:

Quóniam rectus Dóminus Deus noster: * et non est iníquitas in eo.

\antlabel{Ant.} Quam magnificáta sunt ópera tua, Dómine!

\antlabel{Ant.} Lætábitur justus * in Dómino, et sperábit in eo.

\psalmheading{Psalmus 63}

\lettrine[lines=2]{E}{XAUDI}, Deus, oratiónem meam cum déprecor: * a timóre inimíci éripe ánimam meam.

Protexísti me a convéntu malignántium: * a multitúdine operántium iniquitátem.

Quia exacuérunt ut gládium linguas suas: * intendérunt arcum rem amáram, ut sagíttent in occúltis immaculátum.

Súbito sagittábunt eum, et non timébunt: * firmavérunt sibi sermónem nequam.

Narravérunt ut abscónderent láqueos: * dixérunt: Quis vidébit eos?

Scrutáti sunt iniquitátes: * defecérunt scrutántes scrutínio.

Accédet homo ad cor altum: * et exaltábitur Deus.

Sagíttæ parvulórum factæ sunt plagæ eórum: * et infirmátæ sunt contra eos linguæ eórum.

Conturbáti sunt omnes qui vidébant eos: * et tímuit omnis homo.

Et annuntiavérunt ópera Dei, * et facta ejus intellexérunt.

Lætábitur justus in Dómino, et sperábit in eo, * et laudabúntur omnes recti corde.

\antlabel{Ant.} Lætábitur justus in Dómino, et sperábit in eo.

\antlabel{Ant.} Osténde nobis, Dómine, * lumen miseratiónum tuárum.

\canticumheading{Canticum Ecclesiastici}{Eccli. 36, 1-16}

\lettrine[lines=2]{M}{ISERERE} nostri, Deus ómnium, et réspice nos, * et osténde nobis lucem miseratiónum tuárum:

Et immítte timórem tuum super gentes, * quæ non exquisiérunt te,

Ut cognóscant quia non est Deus nisi tu, * et enárrent magnália tua.

Alleva manum tuam super gentes aliénas, * ut vídeant poténtiam tuam.

Sicut enim in conspéctu eórum sanctificátus es in nobis, * sic in conspéctu nostro magnificáberis in eis,

Ut cognóscant te, sicut et nos cognóvimus, * quóniam non est Deus præter te, Dómine.

Innova signa, et immúta mirabília. * Glorífica manum, et bráchium déxtrum.

Excita furórem, et effúnde iram. * Tolle adversárium, et afflíge inimícum.

Festína tempus, et meménto finis, * ut enárrent mirabília tua.

In ira flammæ devorétur qui salvátur: * et qui péssimant plebem tuam, invéniant perditiónem.

Cóntere caput príncipum inimicórum, * dicéntium: Non est álius præter nos.

Cóngrega omnes tribus Jacob: ut cognóscant quia non est Deus nisi tu, * et enárrent magnália tua:

Et hereditábis eos, * sicut ab inítio.

Miserére plebi tuæ, super quam invocátum est nomen tuum: * et Israël, quem coæquásti primogénito tuo.

Miserére civitáti sanctificatiónis tuæ Jerúsalem, * civitáti requiéi tuæ.

Reple Sion inenarrabílibus verbis tuis, * et glória tua pópulum tuum.

\antlabel{Ant.} Osténde nobis, Dómine, lucem miseratiónum tuárum.

\antlabel{Ant.} Omnis spíritus * laudet Dóminum.

\psalmheading{Psalmus 150}

\lettrine[lines=2]{L}{AUDATE} Dóminum in sanctis ejus: * laudáte eum in firmaménto virtútis ejus.

Laudáte eum in virtútibus ejus: * laudáte eum secúndum multitúdinem magnitúdinis ejus.

Laudáte eum in sono tubæ: * laudáte eum in psaltério, et cíthara.

Laudáte eum in týmpano, et choro: * laudáte eum in chordis, et órgano.

Laudáte eum in cýmbalis benesonántibus: laudáte eum in cýmbalis jubilatiónis: * omnis spíritus laudet Dóminum.

\antlabel{Ant.} Omnis spíritus laudet Dóminum.

\rubric{In Officio Festi vel Octávæ ac sanctæ Maríæ in Sabbato Capítulum, Hymnus, Versus, Ant. ad Benedictus et Orátio ut in Próprio vel Communi.}

\rubric{Réliqua ut in Ordinário. 19}

% ===== END OF BATCH 37 -- Covers autumnalis.md lines 12202-12357,
% "Ad Laudes I" for Sabbato in full: Psalms 149, 91, 63, the
% Canticum Ecclesiastici ("Miserére nostri, Deus ómnium"), and
% Psalm 150. Unlike every other feria's Ad Laudes I this volume,
% Saturday's carries no ferial Capitulum/Hymnus of its own -- the
% source instead gives a standalone rubric directing that these,
% along with the Versus, Ant. ad Benedictus, and Oratio, be taken
% "ut in Próprio vel Communi" (from the Proper or Common of the
% occurring feast or the Officium of Our Lady on Saturday), since
% Saturday's Lauds is always completed from elsewhere rather than
% having its own fixed ferial conclusion. Several more psalm
% incipits carry the partial-caps OCR defect ("LAudáte," "MIserére
% nostri," "BOnum est") and are silently normalized to full caps.
% The document continues with "Ad Laudes II" (autumnalis.md line
% 12360).
% =====
\hourheading{Ad Laudes II}

\rubric{In Vigíliis commúnibus, quando fit Officium de Feria.}

(Pater noster et Ave.)

\V\ Deus, in adjutórium.

\antlabel{Ant.} Benígne fac, Dómine.

\psalmheading{Psalmus 50}

\lettrine[lines=2]{M}{ISERERE} mei, Deus, * secúndum magnam misericórdiam tuam.

Et secúndum multitúdinem miseratiónum tuárum, * dele iniquitátem meam.

Amplius lava me ab iniquitáte mea: * et a peccáto meo munda me.

Quóniam iniquitátem meam ego cognósco: * et peccátum meum contra me est semper.

Tibi soli peccávi, et malum coram te feci: * ut justificéris in sermónibus tuis, et vincas cum judicáris.

Ecce enim, in iniquitátibus concéptus sum: * et in peccátis concépit me mater mea.

Ecce enim, veritátem dilexísti: * incérta et occúlta sapiéntiæ tuæ manifestásti mihi.

Aspérges me hyssópo, et mundábor: * lavábis me, et super nivem dealbábor.

Audítui meo dabis gáudium et lætítiam: * et exsultábunt ossa humiliáta.

Avérte fáciem tuam a peccátis meis: * et omnes iniquitátes meas dele.

Cor mundum crea in me, Deus: * et spíritum rectum ínnova in viscéribus meis.

Ne projícias me a fácie tua: * et spíritum sanctum tuum ne áuferas a me.

Redde mihi lætítiam salutáris tui: * et spíritu principáli confírma me.

Docébo iníquos vias tuas: * et ímpii ad te converténtur.

Líbera me de sanguínibus, Deus, Deus salútis meæ: * et exsultábit lingua mea justítiam tuam.

Dómine, lábia mea apéries: * et os meum annuntiábit laudem tuam.

Quóniam si voluísses sacrifícium, dedíssem útique: * holocáustis non delectáberis.

Sacrifícium Deo spíritus contribulátus: * cor contrítum, et humiliátum, Deus, non despícies.

Benígne fac, Dómine, in bona voluntáte tua Sion: * ut ædificéntur muri Jerúsalem.

Tunc acceptábis sacrifícium justítiæ, oblatiónes, et holocáusta: * tunc impónent super altáre tuum vítulos.

\antlabel{Ant.} Benígne fac, Dómine, in bona voluntáte tua Sion.

\antlabel{Ant.} Rectus Dóminus.

\psalmheading{Psalmus 91}

\lettrine[lines=2]{B}{ONUM} est confitéri Dómino: * et psállere nómini tuo, Altíssime.

Ad annuntiándum mane misericórdiam tuam: * et veritátem tuam per noctem.

In decachórdo, psaltério: * cum cántico, in cíthara.

Quia delectásti me, Dómine, in factúra tua: * et in opéribus mánuum tuárum exsultábo.

Quam magnificáta sunt ópera tua, Dómine! * nimis profúndæ factæ sunt cogitatiónes tuæ.

Vir insípiens non cognóscet: * et stultus non intélliget hæc.

Cum exórti fúerint peccatóres sicut fœnum: * et apparúerint omnes, qui operántur iniquitátem:

Ut intéreant in sæculum sæculi: * tu autem Altíssimus in ætérnum, Dómine.

Quóniam ecce inimíci tui, Dómine, quóniam ecce inimíci tui períbunt: * et dispergéntur omnes, qui operántur iniquitátem.

Et exaltábitur sicut unicórnis cornu meum: * et senéctus mea in misericórdia úberi.

Et despéxit óculus meus inimícos meos: * et in insurgéntibus in me malignántibus áudiet auris mea.

Justus ut palma florébit: * sicut cedrus Líbani multiplicábitur.

Plantáti in domo Dómini, * in átriis domus Dei nostri florébunt.

Adhuc multiplicabúntur in senécta úberi: * et patiéntes erunt, ut annúntient:

Quóniam rectus Dóminus, Deus noster: * et non est iníquitas in eo.

\antlabel{Ant.} Rectus Dóminus, Deus noster: et non est iníquitas in eo.

\antlabel{Ant.} A timóre inimíci.

\psalmheading{Psalmus 63}

\lettrine[lines=2]{E}{XAUDI}, Deus, oratiónem meam cum déprecor: * a timóre inimíci éripe ánimam meam.

Protexísti me a convéntu malignántium: * a multitúdine operántium iniquitátem.

Quia exacuérunt ut gládium linguas suas: * intendérunt arcum rem amáram, ut sagíttent in occúltis immaculátum.

Súbito sagittábunt eum, et non timébunt: * firmavérunt sibi sermónem nequam.

Narravérunt ut abscónderent láqueos: * dixérunt: Quis vidébit eos?

Scrutáti sunt iniquitátes: * defecérunt scrutántes scrutínio.

Accédet homo ad cor altum: * et exaltábitur Deus.

Sagíttæ parvulórum factæ sunt plagæ eórum: * et infirmátæ sunt contra eos linguæ eórum.

Conturbáti sunt omnes qui vidébant eos: * et tímuit omnis homo.

Et annuntiavérunt ópera Dei, * et facta ejus intellexérunt.

Lætábitur justus in Dómino, et sperábit in eo, * et laudabúntur omnes recti corde.

\antlabel{Ant.} A timóre inimíci éripe, Dómine, ánimam meam.

\antlabel{Ant.} In servis suis.

\canticumheading{Canticum Moysis}{Deut. 32, 1-43}

\lettrine[lines=2]{A}{UDITE}, cæli, quæ loquor: * áudiat terra verba oris mei.

Concréscat ut plúvia doctrína mea, * fluat ut ros elóquium meum,

Quasi imber super herbam, et quasi stillæ super grámina. * Quia nomen Dómini invocábo:

Date magnificéntiam Deo nostro. * Dei perfécta sunt ópera, et omnes viæ ejus judícia:

Deus fidélis, et absque ulla iniquitáte, justus et rectus. * Peccavérunt ei, et non fílii ejus in sórdibus:

Generátio prava atque pervérsa. * Hæccine réddis Dómino, pópule stulte et insípiens?

Numquid non ipse est pater tuus, * qui possédit te, et fecit, et creávit te?

Meménto diérum antiquórum, * cógita generatiónes síngulas:

Intérroga patrem tuum, et annuntiábit tibi: * majóres tuos, et dicent tibi.

Quando dividébat Altíssimus gentes: * quando separábat fílios Adam,

Constítuit términos populórum * juxta númerum filiórum Israël.

Pars autem Dómini, pópulus ejus: * Jacob funículus hereditátis ejus.

Invénit eum in terra desérta, * in loco horróris, et vastæ solitúdinis:

Circumdúxit eum, et dócuit: * et custodívit quasi pupíllam óculi sui.

Sicut áquila próvocans ad volándum pullos suos, * et super eos vólitans,

Expándit alas suas, et assúmpsit eum, * atque portávit in húmeris suis.

Dóminus solus dux ejus fuit: * et non erat cum eo deus aliénus.

Constítuit eum super excélsam terram: * ut coméderet fructus agrórum,

Ut súgeret mel de petra, * oleúmque de saxo duríssimo.

Butýrum de arménto, et lac de óvibus * cum ádipe agnórum, et aríetum filiórum Basan:

Et hircos cum medúlla trítici, * et sánguinem uvæ bíberet meracíssimum.

Incrassátus est diléctus, et recalcitrávit: * incrassátus, impinguátus, dilatátus,

Derelíquit Deum, factórem suum, * et recéssit a Deo, salutári suo.

Provocavérunt eum in diis aliénis, * et in abominatiónibus ad iracúndiam concitavérunt.

Immolavérunt dæmóniis, et non Deo, * diis, quos ignorábant:

Novi recentésque venérunt, * quos non coluérunt patres eórum.

Deum qui te génuit dereliquísti, * et oblítus es Dómini, creatóris tui.

Vidit Dóminus, et ad iracúndiam concitátus est: * quia provocavérunt eum fílii sui et fíliæ.

Et ait: Abscóndam fáciem meam ab eis, * et considerábo novíssima eórum:

Generátio enim pervérsa est, * et infidéles fílii.

Ipsi me provocavérunt in eo, qui non erat Deus, * et irritavérunt in vanitátibus suis:

Et ego provocábo eos in eo, qui non est pópulus, * et in gente stulta irritábo illos.

Ignis succénsus est in furóre meo, * et ardébit usque ad inférni novíssima:

Devorabítque terram cum gérmine suo, * et móntium fundaménta combúret.

Congregábo super eos mala, * et sagíttas meas complébo in eis.

Consuméntur fame, * et devorábunt eos aves morsu amaríssimo:

Dentes bestiárum immíttam in eos, * cum furóre trahéntium super terram, atque serpéntium.

Foris vastábit eos gládius, et intus pavor, * júvenem simul ac vírginem, lacténtem cum hómine sene.

Dixi: Ubinam sunt? * cessáre fáciam ex homínibus memóriam eórum.

Sed propter iram inimicórum dístuli: * ne forte superbírent hostes eórum,

Et dícerent: * Manus nostra excélsa, et non

Dóminus, * fecit hæc ómnia.

Gens absque consílio est, et sine prudéntia. * Utinam sáperent, et intellígerent, ac novíssima providérent.

Quómodo persequátur unus mille, * et duo fugent decem míllia?

Nonne ídeo, quia Deus suus véndidit eos, * et Dóminus conclúsit illos?

Non enim est Deus noster ut dii eórum: * et inimíci nostri sunt júdices.

De vínea Sodomórum, vínea eórum, * et de suburbánis Gomórrhæ:

Uva eórum uva fellis, * et botri amaríssimi.

Fel dracónum vinum eórum, * et venénum áspidum insanábile.

Nonne hæc cóndita sunt apud me, * et signáta in thesáuris meis?

Mea est últio, et ego retríbuam in témpore, * ut labátur pes eórum:

Juxta est dies perditiónis, * et adésse festínant témpora.

Judicábit Dóminus pópulum suum, * et in servis suis miserébitur:

Vidébit quod infirmáta sit manus, * et clausi quoque defecérunt, residuíque consúmpti sunt.

Et dicet: Ubi sunt dii eórum, * in quibus habébant fidúciam?

De quorum víctimis comedébant ádipes, * et bibébant vinum libáminum:

Surgant, et opituléntur vobis, * et in necessitáte vos prótegant.

Vidéte quod ego sim solus, * et non sit álius Deus præter me:

Ego occídam, et ego vívere fáciam: percútiam, et ego sanábo, * et non est qui de manu mea possit erúere.

Levábo ad cælum manum meam, et dicam: * Vivo ego in ætérnum.

Si acúero ut fulgur gládium meum, * et arripúerit judícium manus mea:

Reddam ultiónem hóstibus meis, * et his qui odérunt me retríbuam.

Inebriábo sagíttas meas sánguine, * et gládius meus devorábit carnes,

De cruóre occisórum, * et de captivitáte, nudáti inimicórum cápitis.

Laudáte, gentes, pópulum ejus, * quia sánguinem servórum suórum ulciscétur:

Et vindíctam retríbuet in hostes eórum, * et propítius erit terræ pópuli sui.

\antlabel{Ant.} In servis suis miserébitur Dóminus: et propítius erit terræ pópuli sui.

\antlabel{Ant.} Laudáte. \textdagger

\psalmheading{Psalmus 150}

\lettrine[lines=2]{L}{AUDATE} \textdagger\ Dóminum in sanctis ejus: * laudáte eum in firmaménto virtútis ejus.

Laudáte eum in virtútibus ejus: * laudáte eum secúndum multitúdinem magnitúdinis ejus.

Laudáte eum in sono tubæ: * laudáte eum in psaltério, et cíthara.

Laudáte eum in týmpano, et choro: * laudáte eum in chordis, et órgano.

Laudáte eum in cýmbalis benesonántibus: laudáte eum in cýmbalis jubilatiónis: * omnis spíritus laudet Dóminum.

\antlabel{Ant.} Laudáte Dóminum secúndum multitúdinem magnitúdinis ejus.

\capitulumlabel{Capit.}{Rom. 13, 12-13}

\lettrine[lines=2]{N}{OX} præcéssit, dies autem appropinquávit. Abjiciámus ergo ópera tenebrárum, et induámur arma lucis. Sicut in die honéste ambulémus.

\hymnusheading{Hymnus}

\lettrine[lines=2]{A}{URORA} jam spargit polum:\\
Terris dies illábitur:\\
Lucis resúltat spículum:\\
Discédat omne lúbricum.

Phantásma noctis éxsulet:\\
Mentis reátus córruat:\\
Quidquid tenébris hórridum\\
Nox áttulit culpæ, cadat.

Ut mane, quod nos últimum\\
Hic deprecámur cérnui,\\
Cum luce nobis éffluat,\\
Hoc dum canóre cóncrepat.

Deo Patri sit glória,\\
Ejúsque soli Fílio,\\
Cum Spíritu Paráclito,\\
Nunc et per omne sæculum. Amen.

\V\ Repléti sumus mane misericórdia tua. \R\ Exsultávimus, et delectáti sumus.

\antlabel{Ad Bened. Ant.} Illúmina, Dómine, * sedéntes in ténebris et umbra mortis, et dírige pedes nostros in viam pacis.

\rubric{Expléto Cántico ac repetíta Antíphona, dicúntur sequéntes}

\variantheading{Preces feriales}

KYRIE, eléison. Christe, eléison. Kýrie, eléison.

Pater noster, \rubric{quod dícitur a solo Hebdomadário totum clara voce, usque ad} \V\ Et ne nos indúcas in tentatiónem. \R\ Sed líbera nos a malo. \V\ Ego dixi: Dómine, miserére mei. \R\ Sana ánimam meam, quia peccávi tibi. \V\ Convértere, Dómine, úsquequo? \R\ Et deprecábilis esto super servos tuos. \V\ Fiat misericórdia tua, Dómine, super nos. \R\ Quemádmodum sperávimus in te. \V\ Sacerdótes tui induántur justítiam. \R\ Et sancti tui exsúltent. \V\ Orémus pro beatíssimo Papa nostro N. \R\ Dóminus consérvet eum, et vivíficet eum, et beátum fáciat eum in terra, et non tradat eum in ánimam inimicórum ejus. \V\ Orémus et pro Antístite nostro N. \R\ Stet et pascat in fortitúdine tua, Dómine, in sublimitáte nóminis tui. \V\ Dómine, salvum fac regem. \R\ Et exáudi nos in die, qua invocavérimus te. \V\ Salvum fac pópulum tuum, Dómine, et bénedic hereditáti tuæ. \R\ Et rege eos, et extólle illos usque in ætérnum. \V\ Meménto Congregatiónis tuæ. \R\ Quam possedísti ab inítio. \V\ Fiat pax in virtúte tua. \R\ Et abundántia in túrribus tuis. \V\ Orémus pro benefactóribus nostris. \R\ Retribúere dignáre, Dómine, ómnibus, nobis bona faciéntibus propter nomen tuum, vitam ætérnam. Amen. \V\ Orémus pro fidélibus defúnctis. \R\ Réquiem ætérnam dona eis, Dómine, et lux perpétua lúceat eis. \V\ Requiéscant in pace. \R\ Amen. \V\ Pro frátribus nostris abséntibus. \R\ Salvos fac servos tuos, Deus meus, sperántes in te. \V\ Pro afflíctis et captívis. \R\ Líbera eos, Deus Israël, ex ómnibus tribulatiónibus suis. \V\ Mitte eis, Dómine, auxílium de sancto. \R\ Et de Sion tuére eos. \V\ Dómine, Deus virtútum, convérte nos. \R\ Et osténde fáciem tuam, et salvi érimus. \V\ Exsúrge, Christe, ádjuva nos. \R\ Et líbera nos propter nomen tuum. \V\ Dómine, exáudi oratiónem meam. \R\ Et clamor meus ad te véniat.

\rubric{Réliqua ut in Ordinário. 19}

% ===== END OF BATCH 38 -- Covers autumnalis.md lines 12360-12673,
% "Ad Laudes II" for Sabbato in full: Psalms 50, 91, 63, the
% Canticum Moysis ("Audíte, cæli, quæ loquor," the long canticle of
% Deuteronomy 32), and Psalm 150, with its own Capitulum "Nox
% præcéssit" and hymn "Auróra jam spargit polum" (unlike Ad Laudes
% I, which borrowed these from the Proper/Common), plus the "Preces
% feriales" litany. The source breaks one canticle verse oddly
% across two short lines ("Et dícerent: * Manus nostra excélsa, et
% non" / "Dóminus, * fecit hæc ómnia."), reproduced exactly as the
% source paragraphs it. The document continues with "Ad Primam"
% (autumnalis.md line 12676).
% =====
\hourheading{Ad Primam}

Pater, Ave et Credo.

\V\ Deus, in adjutórium.

\hymnusheading{Hymnus}

\lettrine[lines=2]{J}{AM} lucis orto sídere,\\
Deum precémur súpplices,\\
Ut in diúrnis áctibus\\
Nos servet a nocéntibus.

Linguam refrénans témperet,\\
Ne litis horror ínsonet:\\
Visum fovéndo cóntegat,\\
Ne vanitátes háuriat.

Sint pura cordis íntima,\\
Absístat et vecórdia:\\
Carnis terat supérbiam\\
Potus cibíque párcitas.

Ut, cum dies abscésserit\\
Noctémque sors redúxerit,\\
Mundi per abstinéntiam\\
Ipsi canámus glóriam.

Deo Patri sit glória,\\
Ejúsque soli Fílio,\\
Cum Spíritu Paráclito,\\
Nunc et per omne sæculum.\\
Amen.

\antlabel{Ant.} Exaltáre, Dómine.

\psalmheading{Psalmus 93, i}

\lettrine[lines=2]{D}{EUS} ultiónum Dóminus: * Deus ultiónum líbere egit.

Exaltáre, qui júdicas terram: * redde retributiónem supérbis.

Usquequo peccatóres, Dómine, * úsquequo peccatóres gloriabúntur:

Effabúntur, et loquéntur iniquitátem: * loquéntur omnes, qui operántur injustítiam?

Pópulum tuum, Dómine, humiliavérunt: * et hereditátem tuam vexavérunt.

Víduam, et ádvenam interfecérunt: * et pupíllos occidérunt.

Et dixérunt: Non vidébit Dóminus, * nec intélliget Deus Jacob.

Intellígite, insipiéntes in pópulo: * et stulti, áliquándo sápite.

Qui plantávit aurem, non áudiet? * aut qui finxit óculum, non consíderat?

Qui córripit gentes, non árguet: * qui docet hóminem sciéntiam?

Dóminus scit cogitatiónes hóminum, * quóniam vanæ sunt.

\psalmheading{Psalmus 93, ii}

\lettrine[lines=2]{B}{EATUS} homo, quem tu erudíeris, Dómine: * et de lege tua docúeris eum,

Ut mítiges ei a diébus malis: * donec fodiátur peccatóri fóvea.

Quia non repéllet Dóminus plebem suam: * et hereditátem suam non derelínquet.

Quoadúsque justítia convertátur in judícium: * et qui juxta illam omnes qui recto sunt corde.

Quis consúrget mihi advérsus malignántes? * aut quis stabit mecum advérsus operántes iniquitátem?

Nisi quia Dóminus adjúvit me: * paulo minus habitásset in inférno ánima mea.

Si dicébam: Motus est pes meus: * misericórdia tua, Dómine, adjuvábat me.

Secúndum multitúdinem dolórum meórum in corde meo: * consolatiónes tuæ lætificavérunt ánimam meam.

Numquid adhæret tibi sedes iniquitátis: * qui fingis labórem in præcépto?

Captábunt in ánimam justi: * et sánguinem innocéntem condemnábunt.

Et factus est mihi Dóminus in refúgium: * et Deus meus in adjutórium spei meæ.

Et reddet illis iniquitátem ipsórum: et in malítia eórum dispérdet eos: * dispérdet illos Dóminus, Deus noster.

\psalmheading{Psalmus 107}

\lettrine[lines=2]{P}{ARATUM} cor meum, Deus, parátum cor meum: * cantábo, et psallam in glória mea.

Exsúrge, glória mea, exsúrge, psaltérium et cíthara: * exsúrgam dilúculo.

Confitébor tibi in pópulis, Dómine: * et psallam tibi in natiónibus.

Quia magna est super cælos misericórdia tua: * et usque ad nubes véritas tua:

Exaltáre super cælos, Deus, et super omnem terram glória tua: * ut liberéntur dilécti tui.

Salvum fac déxtera tua, et exáudi me: * Deus locútus est in sancto suo:

Exsultábo, et dívidam Síchimam, * et convállem tabernaculórum dimétiar.

Meus est Gálaad, et meus est Manásses: * et Ephraïm suscéptio cápitis mei.

Juda rex meus: * Moab lebes spei meæ.

In Idumæam exténdam calceaméntum meum: * mihi alienígenæ amíci facti sunt.

Quis dedúcet me in civitátem munítam? * quis dedúcet me usque in Idumæam?

Nonne tu, Deus, qui repulísti nos, * et non exíbis, Deus, in virtútibus nostris?

Da nobis auxílium de tribulatióne: * quia vana salus hóminis.

In Deo faciémus virtútem: * et ipse ad níhilum dedúcet inimícos nostros.

\rubric{¶ Quando ad Laudes dictus fúerit Ps. 50 Miserére, hic subjúngitur sequens Ps. 149, qui in Láudibus locum non hábuit.}

\psalmheading{Psalmus 149}

\lettrine[lines=2]{C}{ANTATE} Dómino cánticum novum: * laus ejus in ecclésia sanctórum.

Lætétur Israël in eo, qui fecit eum: * et fílii Sion exsúltent in rege suo.

Laudent nomen ejus in choro: * in týmpano, et psaltério psallant ei:

Quia beneplácitum est Dómino in pópulo suo: * et exaltábit mansuétos in salútem.

Exsultábunt sancti in glória: * lætabúntur in cubílibus suis.

Exaltatiónes Dei in gútture eórum: * et gládii ancípites in mánibus eórum.

Ad faciéndam vindíctam in natiónibus: * increpatiónes in pópulis.

Ad alligándos reges eórum in compédibus: * et nóbiles eórum in mánicis férreis.

Ut fáciant in eis judícium conscríptum: * glória hæc est ómnibus sanctis ejus.

\rubric{Secus vero, tribus Psalmis explétis, statim dícitur Ant.}

\antlabel{Ant.} Exaltáre, Dómine, qui júdicas terram: redde retributiónem supérbis.

\rubric{In Officio Festi vel Octávæ, Dominicæ anticipátæ ac sanctæ Maríæ in Sabbato:}

\capitulumlabel{Capitulum}{1 Tim. 1, 17}

\lettrine[lines=2]{R}{EGI} sæculórum immortáli et invisíbili, soli Deo honor et glória in sæcula sæculórum. Amen.

\rubric{In ómnibus autem Fériis et in Vigíliis commúnibus:}

\capitulumlabel{Capitulum}{Zach. 8, 19}

\lettrine[lines=2]{P}{ACEM} et veritátem dilígite, ait Dóminus omnípotens.

\R\ br. Christe, Fili Dei vivi, * Miserére nobis. Christe. \V\ \textit{Qui sedes ad déxteram Patris.} Miserére nobis. Glória Patri. Christe. \V\ Exsúrge, Christe, ádjuva nos. \R\ Et líbera nos propter nomen tuum.

\rubric{In Officio sanctæ Maríæ in Sabbato dícitur sequens}

\variantheading{Preces dominicales et feriales (si dicendæ sint)}

KYRIE, eléison. Christe, eléison. Kýrie, eléison.

Pater noster \rubric{secreto usque ad} \V\ Et ne nos indúcas in tentatiónem. \R\ Sed líbera nos a malo.

Credo in Deum \rubric{secreto usque ad} \V\ Carnis resurrectiónem. \R\ Vitam ætérnam. Amen. \V\ Et ego ad te, Dómine, clamávi. \R\ Et mane orátio mea prævéniet te. \V\ Repleátur os meum laude. \R\ Ut cantem glóriam tuam, tota die magnitúdinem tuam. \V\ Dómine, avérte fáciem tuam a peccátis meis. \R\ Et omnes iniquitátes meas dele. \V\ Cor mundum crea in me, Deus. \R\ Et spíritum rectum ínnova in viscéribus meis. \V\ Ne projícias me a fácie tua. \R\ Et spíritum sanctum tuum ne áuferas a me. \V\ Redde mihi lætítiam salutáris tui. \R\ Et spíritu principáli confírma me.

\variantheading{In Précibus feriálibus tantum}

\V\ Eripe me, Dómine, ab hómine malo. \R\ A viro iníquo éripe me. \V\ Eripe me de inimícis meis, Deus meus. \R\ Et ab insurgéntibus in me líbera me. \V\ Eripe me de operántibus iniquitátem. \R\ Et de viris sánguinum salva me. \V\ Sic psalmum dicam nómini tuo in sæculum sæculi. \R\ Ut reddam vota mea de die in diem. \V\ Exáudi nos, Deus, salutáris noster. \R\ Spes ómnium fínium terræ, et in mari longe. \V\ Deus, in adjutórium meum inténde. \R\ Dómine, ad adjuvándum me festína. \V\ Sanctus Deus, Sanctus fortis, Sanctus immortális. \R\ Miserére nobis. \V\ Bénedic, ánima mea, Dómino. \R\ Et ómnia, quæ intra me sunt, nómini sancto ejus. \V\ Bénedic, ánima mea, Dómino. \R\ Et noli oblivísci omnes retributiónes ejus. \V\ Qui propitiátur ómnibus iniquitátibus tuis. \R\ Qui sanat omnes infirmitátes tuas. \V\ Qui rédimit de intéritu vitam tuam. \R\ Qui corónat te in misericórdia et miseratiónibus. \V\ Qui replet in bonis desidérium tuum. \R\ Renovábitur ut áquilæ juvéntus tua.

\variantheading{In utrisque Précibus}

\V\ Adjutórium nostrum in nómine Dómini. \R\ Qui fecit cælum et terram.

\formulalabel{Confíteor. Misereátur. Indulgéntiam.}

\V\ Dignáre, Dómine, die isto. \R\ Sine peccáto nos custodíre. \V\ Miserére nostri, Dómine. \R\ Miserére nostri. \V\ Fiat misericórdia tua, Dómine, super nos. \R\ Quemádmodum sperávimus in te. \V\ Dómine, exáudi oratiónem meam. \R\ Et clamor meus ad te véniat.

\V\ Dóminus vobíscum.

Orémus. \oratiolabel{Oratio}

\lettrine[lines=2]{D}{OMINE}, Deus omnípotens, qui ad princípium hujus diéi nos perveníre fecísti: tua nos hódie salva virtúte; ut in hac die ad nullum declinémus peccátum, sed semper ad tuam justítiam faciéndam nostra procédant elóquia, dirigántur cogitatiónes et ópera. Per Dóminum. \R\ Amen.

\V\ Dóminus vobíscum. \V\ Benedicámus Dómino.

\rubric{Deinde in Choro legitur Martyrológium. Postea Hebdomadárius dicit:}

\V\ Pretiósa in conspéctu Dómini. \R\ Mors Sanctórum ejus.

\rubric{Et absolúte, sine} Orémus:

\oratiolabel{Oratio}

\lettrine[lines=2]{S}{ANCTA} María, et ómnes Sancti intercédant pro nobis ad Dóminum, ut nos mereámur ab eo adjuvári et salvári, qui vivit et regnat in sæcula sæculórum. \R\ Amen.

\V\ Deus, in adjutórium meum inténde. \R\ Dómine, ad adjuvándum me festína. \rubric{Et dícitur ter; último additur:} Glória Patri. Sicut erat.

Kýrie, eléison. Christe, eléison. Kýrie, eléison.

Pater noster \rubric{secreto usque ad} \V\ Et ne nos indúcas in tentatiónem. \R\ Sed líbera nos a malo. \V\ Réspice in servos tuos, Dómine, et in ópera tua, et dírige fílios eórum. \R\ Et sit splendor Dómini, Dei nostri, super nos, et ópera mánuum nostrárum dírige super nos, et opus mánuum nostrárum dírige. \V\ Glória Patri. \R\ Sicut erat.

Orémus. \oratiolabel{Oratio}

\lettrine[lines=2]{D}{IRIGERE} et sanctificáre, régere et gubernáre dignáre, Dómine Deus, Rex cæli et terræ, hódie corda et córpora nostra, sensus, sermónes et actus nostros in lege tua, et in opéribus mandatórum tuórum: ut hic, et in ætérnum, te auxiliánte, salvi et líberi esse mereámur, Salvátor mundi: Qui vivis et regnas in sæcula sæculórum. \R\ Amen.

\V\ Jube, domne.

\benedlabel{Bened.} Dies et actus nostros in sua pace dispónat Dóminus omnípotens. \R\ Amen.

\variantheading{In Officio feriali}

\lettrine[lines=2]{D}{OMINUS} autem dírigat corda et córpora nostra in caritáte Dei, et patiéntia Christi. Tu autem.

\variantheading{E Proprio Sanctorum}

\variantheading{Pro Oct. Omn. Sanctorum}

\lettrine[lines=2]{B}{ENEDÍCTIO} et cláritas et sapiéntia et gratiárum áctio, honor, virtus et fortitúdo Deo nostro in sǽcula sæculórum. Amen.

\variantheading{E Communi Sanctorum}

\variantheading{Pro Apostolis}

\lettrine[lines=2]{I}{BANT} Apóstoli gaudéntes a conspéctu concílii, quóniam digni hábiti sunt pro nómine Jesu contuméliam pati.

\variantheading{Pro Martyre}

\lettrine[lines=2]{J}{USTUS} cor suum trádidit ad vigilándum dilúculo ad Dóminum, qui fecit illum, et in conspéctu Altíssimi deprecábitur.

\variantheading{Pro pluribus Martyribus}

\lettrine[lines=2]{F}{ULGEBUNT} justi, et tamquam scintíllæ in arundinéto discúrrent. Judicábunt natiónes, et dominabúntur pópulis: et regnábit Dóminus illórum in perpétuum.

\variantheading{Pro Conf. Pont.}

\lettrine[lines=2]{F}{UNGI} sacerdótio, et habére laudem in nómine ipsíus, et offérre illi incénsum dignum in odórem suavitátis.

\variantheading{Pro Conf. non Pont.}

\lettrine[lines=2]{J}{USTUM} dedúxit Dóminus per vias rectas, et osténdit illi regnum Dei, et dedit illi sciéntiam sanctórum; honestávit illum in labóribus, et complévit labóres illíus.

\variantheading{Pro Virgine Martyre}

\lettrine[lines=2]{D}{OMINE}, Deus meus, exaltásti super terram habitatiónem meam, et pro morte defluénte deprecáta sum. Invocávi Dóminum, Patrem Dómini mei, ut non derelínquat me in die tribulatiónis meæ, et in témpore superbórum sine adjutório.

\variantheading{Pro Virgine non Martyre}

\lettrine[lines=2]{O}{} \textsc{quam} pulchra est casta generátio cum claritáte! immortális est enim memória illíus: quóniam apud Deum nota est, et apud hómines.

\variantheading{Pro Martyre tantum}

\lettrine[lines=2]{L}{AUDABIT} usque ad mortem ánima mea Dóminum, quóniam éruis sustinéntes te, et líberas eos de manu angústiæ, Dómine, Deus noster.

\variantheading{Pro nec Virg. nec Mart.}

\lettrine[lines=2]{M}{ULTAE} fíliæ congregavérunt divítias: tu supergréssa es univérsas. Fallax grátia, et vana est pulchritúdo: múlier timens Dóminum, ipsa laudábitur.

\variantheading{Pro Dedicatione Ecclesiæ}

\lettrine[lines=2]{E}{T} abstérget Deus omnem lácrimam ab óculis eórum: et mors ultra non erit, neque luctus, neque clamor, neque dolor erit ultra, quia prima abiérunt. Et dixit qui sedébat in throno: Ecce, nova fácio ómnia.

\variantheading{In Officiis B. Mariæ V.}

\lettrine[lines=2]{I}{N} platéis sicut cinnamómum et bálsamum aromatízans odórem dedi: quasi myrrha elécta dedi suavitátem odóris.

\V\ Adjutórium nostrum in nómine Dómini. \R\ Qui fecit cælum et terram.

\V\ Benedícite. \R\ Deus.

\benedictiolabel{Benedictio} Dóminus nos benedícat, et ab omni malo deféndat, et ad vitam perdúcat ætérnam. Et fidélium ánimæ per misericórdiam Dei requiéscant in pace. \R\ Amen.

\rubric{Deinde dícitur tantum} Pater noster \rubric{secréto, nisi sequátur ália Hora.}

% ===== END OF BATCH 39 -- Covers autumnalis.md lines 12676-12954,
% completing "Ad Primam" for Sabbato in full: the hymn "Jam lucis
% orto sídere," Psalms 93/i, 93/ii, 107, the conditional Psalm 149
% insertion, both Capitula (the first now naming "Dominicæ
% anticipátæ" -- the anticipated Sunday -- alongside the usual
% Festum/Octave/Marian-Saturday conditions), the brief Responsory,
% a new standalone rubric specific to this day directing that the
% Preces which follow are said "In Officio sanctæ Maríæ in
% Sabbato," the full Preces apparatus, the closing Orationes, and
% the Lectio-brevis apparatus. Several more psalm incipits and
% Oratio openings carry the partial-caps OCR defect ("DEus
% ultiónum," "BEatus homo," "PAratum cor meum," "CAntáte," "REgi"
% [Capitulum], "DOmine, Deus omnípotens" [Oratio], "SAncta María"
% [Oratio], "DIrígere" [Oratio]) and are silently normalized to
% full caps. The document continues with "Ad Tertiam" (autumnalis.md
% line 12958).
% =====
\hourheading{Ad Tertiam}

Pater noster et Ave.

\V\ Deus, in adjutórium.

\hymnusheading{Hymnus}

\lettrine[lines=2]{N}{UNC}, Sancte, nobis, Spíritus,\\
Unum Patri cum Fílio,\\
Dignáre promptus íngeri\\
Nostro refúsus péctori.

Os, lingua, mens, sensus, vigor\\
Confessiónem pérsonent,\\
Flamméscat igne cáritas,\\
Accéndat ardor próximos.

Præsta, Pater piíssime,\\
Patríque compar Unice,\\
Cum Spíritu Paráclito\\
Regnans per omne sæculum.\\
Amen.

\antlabel{Ant.} Clamor meus.

\psalmheading{Psalmus 101, i}

\lettrine[lines=2]{D}{OMINE}, exáudi oratiónem meam: * et clamor meus ad te véniat.

Non avértas fáciem tuam a me: * in quacúmque die tríbulor, inclína ad me aurem tuam.

In quacúmque die invocávero te, * velóciter exáudi me.

Quia defecérunt sicut fumus dies mei: * et ossa mea sicut crémium aruérunt.

Percússus sum ut fœnum, et áruit cor meum: * quia oblítus sum comédere panem meum.

A voce gémitus mei * adhæsit os meum carni meæ.

Símilis factus sum pellicáno solitúdinis: * factus sum sicut nyctícorax in domicílio.

Vigiláví, * et factus sum sicut passer solitárius in tecto.

Tota die exprobrábant mihi inimíci mei: * et qui laudábant me, advérsum me jurábant.

Quia cínerem tamquam panem manducábam, * et potum meum cum fletu miscébam.

A fácie iræ et indignatiónis tuæ: * quia élevans allisísti me.

Dies mei sicut umbra declinavérunt: * et ego sicut fœnum árui.

Tu autem, Dómine, in ætérnum pérmanes: * et memoriále tuum in generatiónem et generatiónem.

\psalmheading{Psalmus 101, ii}

\lettrine[lines=2]{T}{U} exsúrgens miseréberis Sion: * quia tempus miseréndi ejus, quia venit tempus.

Quóniam placuérunt servis tuis lápides ejus: * et terræ ejus miserebúntur.

Et timébunt gentes nomen tuum, Dómine, * et omnes reges terræ glóriam tuam.

Quia ædificávit Dóminus Sion: * et vidébitur in glória sua.

Respéxit in oratiónem humílium: * et non sprevit precem eórum.

Scribántur hæc in generatióne áltera: * et pópulus, qui creábitur, laudábit Dóminum:

Quia prospéxit de excélso sancto suo: * Dóminus de cælo in terram aspéxit:

Ut audíret gémitus compeditórum: * ut sólveret fílios interemptórum:

Ut annúntient in Sion nomen Dómini: * et laudem ejus in Jerúsalem.

In conveniéndo pópulos in unum, * et reges ut sérviant Dómino.

\psalmheading{Psalmus 101, iii}

\lettrine[lines=2]{R}{ESPONDIT} ei in via virtútis suæ: * Paucitátem diérum meórum núntia mihi.

Ne révoces me in dimídio diérum meórum: * in generatiónem et generatiónem anni tui.

Inítio tu, Dómine, terram fundásti: * et ópera mánuum tuárum sunt cæli.

Ipsi períbunt, tu autem pérmanes: * et omnes sicut vestiméntum veteráscent.

Et sicut opertórium mutábis eos, et mutabúntur: * tu autem idem ipse es, et anni tui non defícient.

Fílii servórum tuórum habitábunt: * et semen eórum in sæculum dirigétur.

\antlabel{Ant.} Clamor meus, Dómine, ad te pervéniat: * non avértas fáciem tuam a me.

\variantheading{In Officio feriali}

\capitulumlabel{Capitulum}{Jer. 17, 14}

\lettrine[lines=2]{S}{ANA} me, Dómine, et sanábor: salvum me fac, et salvus ero: quóniam laus mea tu es.

\R\ br. Sana ánimam meam, * Quia peccávi tibi. Sana. \V\ Ego dixi: Dómine, miserére mei. Quia. Glória Patri. Sana. \V\ Adjútor meus esto, ne derelínquas me. \R\ Neque despícias me, Deus, salutáris meus.

\variantheading{In Vigiliis communibus:}

\variantheading{Preces feriales}

KYRIE, eléison. Christe, eléison. Kýrie, eléison.

Pater noster \rubric{secreto usque ad} \V\ Et ne nos indúcas in tentatiónem. \R\ Sed líbera nos a malo. \V\ Dómine, Deus virtútum, convérte nos. \R\ Et osténde fáciem tuam, et salvi érimus. \V\ Exsúrge, Christe, ádjuva nos. \R\ Et líbera nos propter nomen tuum. \V\ Dómine, exáudi oratiónem meam. \R\ Et clamor meus ad te véniat.

\rubric{Réliqua ut in Ordinário. 19}

\hourheading{Ad Sextam}

Pater noster et Ave.

\V\ Deus, in adjutórium.

\hymnusheading{Hymnus}

\lettrine[lines=2]{R}{ECTOR} potens, verax Deus,\\
Qui témperas rerum vices,\\
Splendóre mane illúminas,\\
Et ígnibus merídiem:

Exstíngue flammas lítium,\\
Aufer calórem nóxium,\\
Confer salútem córporum,\\
Verámque pacem córdium.

Præsta, Pater piíssime,\\
Patríque compar Unice,\\
Cum Spíritu Paráclito\\
Regnans per omne sæculum.\\
Amen.

\antlabel{Ant.} Dómine, Deus meus.

\psalmheading{Psalmus 103, i}

\lettrine[lines=2]{B}{ENEDIC}, ánima mea, Dómino: * Dómine, Deus meus, magnificátus es vehementer.

Confessiónem, et decórem induísti: * amíctus lúmine sicut vestiménto:

Exténdens cælum sicut pellem: * qui tegis aquis superióra ejus.

Qui ponis nubem ascénsum tuum: * qui ámbulas super pennas ventórum.

Qui facis ángelos tuos, spíritus: * et minístros tuos ignem uréntem.

Qui fundásti terram super stabilitátem suam: * non inclinábitur in sæculum sæculi.

Abýssus, sicut vestiméntum, amíctus ejus: * super montes stabunt aquæ.

Ab increpatióne tua fúgient: * a voce tonítrui tui formidábunt.

Ascéndunt montes: et descéndunt campi * in locum, quem fundásti eis.

Términum posuísti, quem non transgrediéntur: * neque converténtur operíre terram.

Qui emíttis fontes in convállibus: * inter médium móntium pertransíbunt aquæ.

Potábunt omnes béstiæ agri: * exspectábunt ónagri in siti sua.

Super ea volúcres cæli habitábunt: * de médio petrárum dabunt voces.

\psalmheading{Psalmus 103, ii}

\lettrine[lines=2]{R}{IGANS} montes de superióribus suis: * de fructu óperum tuórum satiábitur terra:

Prodúcens fœnum juméntis, * et herbam servitúti hóminum:

Ut edúcas panem de terra: * et vinum lætíficet cor hóminis:

Ut exhílaret fáciem in óleo: * et panis cor hóminis confírmet.

Saturabúntur ligna campi, et cedri Líbani, quas plantávit: * illic pásseres nidificábunt.

Heródii domus dux est eórum: * montes excélsi cervis: petra refúgium herináciis.

Fecit lunam in témpora: * sol cognóvit occásum suum.

Posuísti ténebras, et facta est nox: * in ipsa pertransíbunt omnes béstiæ silvæ.

Cátuli leónum rugiéntes, ut rápiant, * et quærant a Deo escam sibi.

Ortus est sol, et congregáti sunt: * et in cubílibus suis collocabúntur.

Exíbit homo ad opus suum: * et ad operatiónem suam usque ad vésperum.

\psalmheading{Psalmus 103, iii}

\lettrine[lines=2]{Q}{UAM} magnificáta sunt ópera tua, Dómine! * ómnia in sapiéntia fecísti: impléta est terra possessióne tua.

Hoc mare magnum, et spatiósum mánibus: * illic reptília, quorum non est númerus.

Animália pusílla cum magnis: * illic naves pertransíbunt.

Draco iste, quem formásti ad illudéndum ei: * ómnia a te exspéctant ut des illis escam in témpore.

Dante te illis, cólligent: * aperiénte te manum tuam, ómnia implebúntur bonitáte.

Avertánte autem te fáciem, turbabúntur: * áuferes spíritum eórum, et defícient, et in púlverem suum reverténtur.

Emíttes spíritum tuum, et creabúntur: * et renovábis fáciem terræ.

Sit glória Dómini in sæculum: * lætábitur Dóminus in opéribus suis:

Qui réspicit terram, et facit eam trémere: * qui tangit montes, et fúmigant.

Cantábo Dómino in vita mea: * psallam Deo meo, quámdiu sum.

Jucúndum sit ei elóquium meum: * ego vero delectábor in Dómino.

Defíciant peccatóres a terra, et iníqui ita ut non sint: * bénedic, ánima mea, Dómino.

\antlabel{Ant.} Dómine, Deus meus, magnificátus es veheménter.

\variantheading{In Officio feriali}

\capitulumlabel{Capitulum}{Rom. 13, 8}

\lettrine[lines=2]{N}{EMINI} quidquam debeátis, nisi ut ínvicem diligátis: qui enim díligit próximum, legem implévit.

\R\ br. Benedícam Dóminum * In omni témpore. Benedícam. \V\ Semper laus ejus in ore meo. In. Glória Patri. Benedícam. \V\ Dóminus regit me, et nihil mihi déerit. \R\ In loco páscuæ ibi me collocávit.

\variantheading{In Vigiliis communibus:}

\variantheading{Preces feriales}

KYRIE, eléison. Christe, eléison. Kýrie, eléison.

Pater noster \rubric{secreto usque ad} \V\ Et ne nos indúcas in tentatiónem. \R\ Sed líbera nos a malo. \V\ Dómine, Deus virtútum, convérte nos. \R\ Et osténde fáciem tuam, et salvi érimus. \V\ Exsúrge, Christe, ádjuva nos. \R\ Et líbera nos propter nomen tuum. \V\ Dómine, exáudi oratiónem meam. \R\ Et clamor meus ad te véniat.

\rubric{Réliqua ut in Ordinário. 20}

\hourheading{Ad Nonam}

Pater noster et Ave.

\V\ Deus, in adjutórium.

\hymnusheading{Hymnus}

\lettrine[lines=2]{R}{ERUM}, Deus, tenax vigor,\\
Immótus in te pérmanens,\\
Lucis diúrnæ témpora\\
Succéssibus detérminans:

Largíre lumen véspere,\\
Quo vita nusquam décidat,\\
Sed præmium mortis sacræ\\
Perénnis instet glória.

Præsta, Pater piíssime,\\
Patríque compar Unice,\\
Cum Spíritu Paráclito\\
Regnans per omne sæculum.\\
Amen.

\antlabel{Ant.} Ne tacúeris, Deus.

\psalmheading{Psalmus 108, i}

\lettrine[lines=2]{D}{EUS}, laudem meam ne tacúeris: * quia os peccatóris, et os dolósi super me apértum est.

Locúti sunt advérsum me lingua dolósa, et sermónibus ódii circumdedérunt me: * et expugnavérunt me gratis.

Pro eo ut me dilígerent, detrahébant mihi: * ego autem orábam.

Et posuérunt advérsum me mala pro bonis: * et ódium pro dilectióne mea.

Constítue super eum peccatórem: * et diábolus stet a dextris ejus.

Cum judicátur, éxeat condemnátus: * et orátio ejus fiat in peccátum.

Fiant dies ejus pauci: * et episcopátum ejus accípiat alter.

Fiant fílii ejus órphani: * et uxor ejus vídua.

Nutántes transferántur fílii ejus, et mendícent: * et ejiciántur de habitatiónibus suis.

Scrutétur fænerátor omnem substántiam ejus: * et dirípiant aliéni labóres ejus.

Non sit illi adjútor: * nec sit qui misereátur pupíllis ejus.

Fiant nati ejus in intéritum: * in generatióne una deleátur nomen ejus.

\psalmheading{Psalmus 108, ii}

\lettrine[lines=2]{I}{N} memóriam rédeat iníquitas patrum ejus in conspéctu Dómini: * et peccátum matris ejus non deleátur.

Fiant contra Dóminum semper, et dispéreat de terra memória eórum: * pro eo quod non est recordátus fácere misericórdiam.

Et persecútus est hóminem ínopem, et mendícum, * et compúnctum corde mortificáre.

Et diléxit maledictiónem, et véniet ei: * et nóluit benedictiónem, et elongábitur ab eo.

Et índuit maledictiónem sicut vestiméntum, * et intrávit sicut aqua in interióra ejus, et sicut óleum in óssibus ejus.

Fiat ei sicut vestiméntum, quo operítur: * et sicut zona, qua semper præcíngitur.

Hoc opus eórum, qui détrahunt mihi apud Dóminum: * et qui loquúntur mala advérsum ánimam meam.

Et tu, Dómine, Dómine, fac mecum propter nomen tuum: * quia suávis est misericórdia tua.

\psalmheading{Psalmus 108, iii}

\lettrine[lines=2]{L}{IBERA} me quia egénus, et pauper ego sum: * et cor meum conturbátum est intra me.

Sicut umbra cum declínat, ablátus sum: * et excússus sum sicut locústæ.

Génua mea infirmáta sunt a jejúnio: * et caro mea immutáta est propter óleum.

Et ego factus sum oppróbrium illis: * vidérunt me, et movérunt cápita sua.

Adjuva me, Dómine, Deus meus: * salvum me fac secúndum misericórdiam tuam.

Et sciant quia manus tua hæc: * et tu, Dómine, fecísti eam.

Maledícent illi, et tu benedíces: * qui insúrgunt in me, confundántur: servus autem tuus lætábitur.

Induántur qui détrahunt mihi, pudóre: * et operiántur sicut diplóide confusióne sua.

Confitébor Dómino nimis in ore meo: * et in médio multórum laudábo eum.

Quia ástitit a dextris páuperis, * ut salvam fáciat a persequéntibus ánimam meam.

\antlabel{Ant.} Ne tacúeris, Deus, quia sermónibus ódii circumdedérunt me.

\variantheading{In Officio feriali}

\capitulumlabel{Capitulum}{1 Petri 1, 17-19}

\lettrine[lines=2]{I}{N} timóre incolátus vestri témpore conversámini: sciéntes quod non corruptibílibus auro vel argénto redémpti estis, sed pretióso sánguine quasi Agni immaculáti Christi.

\R\ br. Rédime me, Dómine, * Et miserére mei. Rédime. \V\ Pes enim meus stetit in via recta. Et. Glória Patri. Rédime.

\V\ Ab occúltis meis munda me, Dómine. \R\ Et ab aliénis parce servo tuo.

\variantheading{In Vigiliis communibus:}

\variantheading{Preces feriales}

KYRIE, eléison. Christe, eléison. Kýrie, eléison.

Pater noster \rubric{secreto usque ad} \V\ Et ne nos indúcas in tentatiónem. \R\ Sed líbera nos a malo. \V\ Dómine, Deus virtútum, convérte nos. \R\ Et osténde fáciem tuam, et salvi érimus. \V\ Exsúrge, Christe, ádjuva nos. \R\ Et líbera nos propter nomen tuum. \V\ Dómine, exáudi oratiónem meam. \R\ Et clamor meus ad te véniat.

\rubric{Réliqua ut in Ordinário. 22}

% ===== END OF BATCH 40 -- Covers autumnalis.md lines 12958-13309,
% completing the Little Hours of Sabbato in full: "Ad Tertiam"
% (Psalms 101/i-iii), "Ad Sextam" (Psalms 103/i-iii), and "Ad
% Nonam" (Psalms 108/i-iii), each with its own hymn, Capitulum,
% brief Responsory, and "In Vigiliis communibus: Preces feriales"
% apparatus. Several more partial-caps OCR defects ("REspóndit,"
% "BEnedic," "RIgans," "QUam magnificáta," "DEus, laudem," "IN
% memóriam," "LIbera me") are silently normalized to full caps.
% The document continues with "Ad Vesperas" (autumnalis.md line
% 13312), which carries a special rubric about this Saturday's
% Vespers serving as First Vespers of the following Sunday.
% =====
\hourheading{Ad Vesperas}

Pater noster et Ave.

\V\ Deus, in adjutórium.

\antlabel{Ant.} Benedíctus Dóminus. \textdagger

\rubric{¶ Antíphonæ cum Psalmis assignátæ semper in Sábbato dicúntur pro I Vésperis Domínicæ sequéntis, quando agéndum est Offícium de ipsa Domínica, in Domínicis minóribus, extra Octávas Sanctíssimi Córporis Christi ac Sacratíssimi Cordis Jesu.}

\psalmheading{Psalmus 143, i}

\lettrine[lines=2]{B}{ENEDICTUS} Dóminus, \textdagger\ Deus meus, qui docet manus meas ad prœlium, * et dígitos meos ad bellum.

Misericórdia mea, et refúgium meum: * suscéptor meus, et liberátor meus:

Protéctor meus, et in ipso sperávi: * qui subdit pópulum meum sub me.

Dómine, quid est homo, quia innotuísti ei? * aut fílius hóminis, quia réputas eum?

Homo vanitáti símilis factus est: * dies ejus sicut umbra prætéreunt.

Dómine, inclína cælos tuos, et descénde: * tange montes, et fumigábunt.

Fúlgura coruscatiónem, et dissipábis eos: * emítte sagíttas tuas, et conturbábis eos.

Emítte manum tuam de alto, éripe me, et líbera me de aquis multis: * de manu filiórum alienórum.

Quorum os locútum est vanitátem: * et déxtera eórum, déxtera iniquitátis.

\antlabel{Ant.} Benedíctus Dóminus, suscéptor meus, et liberátor meus.

\antlabel{Ant.} Beátus pópulus.

\psalmheading{Psalmus 143, ii}

\lettrine[lines=2]{D}{EUS}, cánticum novum cantábo tibi: * in psaltério decachórdo psallam tibi.

Qui das salútem régibus: * qui redemísti David, servum tuum, de gládio malígno: éripe me.

Et érue me de manu filiórum alienórum, quorum os locútum est vanitátem: * et déxtera eórum, déxtera iniquitátis:

Quorum fílii, sicut novéllæ plantatiónes * in juventúte sua.

Fíliæ eórum compósitæ: * circumornátæ ut similitúdo templi.

Promptuária eórum plena, * eructántia ex hoc in illud.

Oves eórum fetósæ, abundántes in egréssibus suis: * boves eórum crassæ.

Non est ruína macériæ, neque tránsitus: * neque clamor in platéis eórum.

Beátum dixérunt pópulum, cui hæc sunt: * beátus pópulus, cujus Dóminus Deus ejus.

\antlabel{Ant.} Beátus pópulus, cujus Dóminus Deus ejus.

\antlabel{Ant.} Magnus Dóminus.

\psalmheading{Psalmus 144, i}

\lettrine[lines=2]{E}{XALTABO} te, Deus meus, rex: * et benedícam nómini tuo in sæculum, et in sæculum sæculi.

Per síngulos dies benedícam tibi: * et laudábo nomen tuum in sæculum, et in sæculum sæculi.

Magnus Dóminus, et laudábilis nimis: * et magnitúdinis ejus non est finis.

Generátio et generátio laudábit ópera tua: * et poténtiam tuam pronuntiábunt.

Magnificéntiam glóriæ sanctitátis tuæ loquéntur: * et mirabília tua narrábunt.

Et virtútem terribílium tuórum dicent: * et magnitúdinem tuam narrábunt.

Memóriam abundántiæ suavitátis tuæ eructábunt: * et justítia tua exsultábunt.

\antlabel{Ant.} Magnus Dóminus, et laudábilis nimis: et magnitúdinis ejus non est finis.

\antlabel{Ant.} Suávis Dóminus.

\psalmheading{Psalmus 144, ii}

\lettrine[lines=2]{M}{ISERATOR}, et miséricors Dóminus: * pátiens, et multum miséricors.

Suávis Dóminus univérsis: * et miseratiónes ejus super ómnia ópera ejus.

Confiteántur tibi, Dómine, ómnia ópera tua: * et sancti tui benedícant tibi.

Glóriam regni tui dicent: * et poténtiam tuam loquéntur:

Ut notam fáciant fíliis hóminum poténtiam tuam: * et glóriam magnificéntiæ regni tui.

Regnum tuum regnum ómnium sæculórum: * et dominátio tua in omni generatióne et generatiónem.

\antlabel{Ant.} Suávis Dóminus univérsis: et miseratiónes ejus super ómnia ópera ejus.

\antlabel{Ant.} Fidélis Dóminus. \textdagger

\psalmheading{Psalmus 144, iii}

\lettrine[lines=2]{F}{IDELIS} Dóminus \textdagger\ in ómnibus verbis suis: * et sanctus in ómnibus opéribus suis.

Allevat Dóminus omnes qui córruunt: * et érigit omnes elísos.

Óculi ómnium in te sperant, Dómine: * et tu das escam illórum in témpore opportúno.

Áperis tu manum tuam: * et imples omne ánimal benedictióne.

Justus Dóminus in ómnibus viis suis: * et sanctus in ómnibus opéribus suis.

Prope est Dóminus ómnibus invocántibus eum: * ómnibus invocántibus eum in veritáte.

Voluntátem timéntium se fáciet: * et deprecatiónem eórum exáudiet: et salvos fáciet eos.

Custódit Dóminus omnes diligéntes se: * et omnes peccatóres dispérdet.

Laudatiónem Dómini loquétur os meum: * et benedícat omnis caro nómini sancto ejus in sæculum, et in sæculum sæculi.

\antlabel{Ant.} Fidélis Dóminus in ómnibus verbis suis: et sanctus in ómnibus opéribus suis.

\rubric{In Sabbatis extra Octávas Ssmi Córporis Christi ac Sacratíssimi Cordis Jesu dícitur}

\variantheading{In Officio feriali}

\capitulumlabel{Capitulum}{Rom. 11, 33}

\lettrine[lines=2]{O}{} \textsc{altitudo} divitiárum sapiéntiæ et sciéntiæ Dei: quam incomprehensíbilia sunt judícia ejus, et investigábiles viæ ejus!

\hymnusheading{Hymnus}

\lettrine[lines=2]{J}{AM} sol recédit ígneus:\\
Tu, lux perénnis, Únitas,\\
Nostris, beáta Trínitas,\\
Infúnde lumen córdibus.

Te mane laudum cármine,\\
Te deprecámur véspere;\\
Dignéris ut te súpplices\\
Laudémus inter Cælites.

Patri, simúlque Fílio,\\
Tibíque, Sancte Spíritus,\\
Sicut fuit, sit júgiter\\
Sæclum per omne glória.\\
Amen.

\V\ Vespertína orátio ascéndat ad te, Dómine. \R\ Et descéndat super nos misericórdia tua.

\rubric{Ad Magníficat Antíphona ut in Próprio de Témpore.}

\rubric{Expléto Cántico ac repetíta Antíphona, dícitur Orátio Domínicæ sequéntis própria, quæ per totam hebdómadam suam in feriáli Officio adhibétur, nisi speciális Orátio habeátur.}

\rubric{Réliqua ut in Ordinário. 22}

\hourheading{Ad Completorium}

\V\ Jube, domne.

\benedlabel{Bened.} Noctem quiétam, et finem perféctum concédat nobis Dóminus omnípotens. \R\ Amen.

\lectioheading{Lectio br.}{1 Petri 5, 8-9}

\lettrine[lines=2]{F}{RATRES}: Sóbrii estóte, et vigiláte: quia adversárius vester diábolus tamquam leo rúgiens círcuit, quærens quem dévoret: cui resístite fortes in fide. Tu autem.

\R\ Deo grátias.

\V\ Adjutórium nostrum in nómine Dómini. \R\ Qui fecit cælum et terram.

Pater noster, \rubric{quod dícitur totum secréto.}

\formulalabel{Confíteor.} Misereátur. Indulgéntiam.

\V\ Convérte nos, Deus, salutáris noster. \R\ Et avérte iram tuam a nobis.

\V\ Deus, in adjutórium.

\antlabel{Ant.} Intret orátio mea.

\psalmheading{Psalmus 87}

\lettrine[lines=2]{D}{OMINE}, Deus salútis meæ: * in die clamávi, et nocte coram te.

Intret in conspéctu tuo orátio mea: * inclína aurem tuam ad precem meam:

Quia repléta est malis ánima mea: * et vita mea inférno appropinquávit.

Æstimátus sum cum descendéntibus in lacum: * factus sum sicut homo sine adjutório, inter mórtuos liber.

Sicut vulneráti dormiéntes in sepúlcris, quorum non es memor ámplius: * et ipsi de manu tua repúlsi sunt.

Posuérunt me in lacu inferióri: * in tenebrósis, et in umbra mortis.

Super me confirmátus est furor tuus: * et omnes fluctus tuos induxísti super me.

Longe fecísti notos meos a me: * posuérunt me abominatiónem sibi.

Tráditus sum, et non egrediébar: * óculi mei languérunt præ inópia.

Clamávi ad te, Dómine, tota die: * expándi ad te manus meas.

Numquid mórtuis fácies mirabília: * aut médici suscitábunt, et confitebúntur tibi?

Numquid narrábit áliquis in sepúlcro misericórdiam tuam, * et veritátem tuam in perditióne?

Numquid cognoscéntur in ténebris mirabília tua, * et justítia tua in terra obliviónis?

Et ego ad te, Dómine, clamávi: * et mane orátio mea prævéniet te.

Ut quid, Dómine, repéllis oratiónem meam: * avértis fáciem tuam a me?

Pauper sum ego, et in labóribus a juventúte mea: * exaltátus autem, humiliátus sum et conturbátus.

In me transiérunt iræ tuæ: * et terróres tui conturbavérunt me.

Circumdedérunt me sicut aqua tota die: * circumdedérunt me simul.

Elongásti a me amícum et próximum: * et notos meos a miséria.

\psalmheading{Psalmus 102, i}

\lettrine[lines=2]{B}{ENEDIC}, ánima mea, Dómino: * et ómnia, quæ intra me sunt, nómini sancto ejus.

Bénedic, ánima mea, Dómino: * et noli oblivísci omnes retributiónes ejus.

Qui propitiátur ómnibus iniquitátibus tuis: * qui sanat omnes infirmitátes tuas.

Qui rédimit de intéritu vitam tuam: * qui corónat te in misericórdia et miseratiónibus.

Qui replet in bonis desidérium tuum: * renovábitur ut áquilæ juvéntus tua:

Fáciens misericórdias Dóminus: * et judícium ómnibus injúriam patiéntibus.

Notas fecit vias suas Móysi, * fíliis Israël voluntátes suas.

Miserátor, et miséricors Dóminus: * longánimis, et multum miséricors.

Non in perpétuum irascétur: * neque in ætérnum comminábitur.

Non secúndum peccáta nostra fecit nobis: * neque secúndum iniquitátes nostras retríbuit nobis.

Quóniam secúndum altitúdinem cæli a terra: * corroborávit misericórdiam suam super timéntes se.

Quantum distat ortus ab occidénte: * longe fecit a nobis iniquitátes nostras.

\psalmheading{Psalmus 102, ii}

\lettrine[lines=2]{Q}{UOMODO} miserétur pater filiórum, misértus est Dóminus timéntibus se: * quóniam ipse cognóvit figméntum nostrum.

Recordátus est quóniam pulvis sumus: * homo, sicut fœnum dies ejus, tamquam flos agri sic efflorébit.

Quóniam spíritus pertransíbit in illo, et non subsístet: * et non cognóscet ámplius locum suum.

Misericórdia autem Dómini ab ætérno, * et usque in ætérnum super timéntes eum.

Et justítia illíus in fílios filiórum, * his qui servant testaméntum ejus:

Et mémores sunt mandatórum ipsíus, * ad faciéndum ea.

Dóminus in cælo parávit sedem suam: * et regnum ipsíus ómnibus dominábitur.

Benedícite Dómino, omnes Ángeli ejus: * poténtes virtúte, faciéntes verbum illíus, ad audiéndam vocem sermónum ejus.

Benedícite Dómino, omnes virtútes ejus: * minístri ejus, qui fácitis voluntátem ejus.

Benedícite Dómino, ómnia ópera ejus: * in omni loco dominatiónis ejus, benedíc, ánima mea, Dómino.

\antlabel{Ant.} Intret orátio mea in conspéctu tuo, Dómine.

\hymnusheading{Hymnus}

\lettrine[lines=2]{T}{E} lucis ante términum,\\
Rerum Creátor, póscimus,\\
Ut pro tua cleméntia\\
Sis præsul et custódia.

Procul recédant sómnia,\\
Et nóctium phantásmata;\\
Hostémque nostrum cómprime,\\
Ne polluántur córpora.

Præsta, Pater piíssime,\\
Patríque compar Unice,\\
Cum Spíritu Paráclito\\
Regnans per omne sæculum.\\
Amen.

\capitulumlabel{Capitulum}{Jer. 14, 9}

\lettrine[lines=2]{T}{U} autem in nobis es, Dómine, et nomen sanctum tuum invocátum est super nos: ne derelínquas nos, Dómine, Deus noster.

\R\ br. In manus tuas, Dómine, * Comméndo spíritum meum. In manus. \V\ Redemísti nos, Dómine, Deus veritátis. Comméndo. Glória Patri. In manus.

\V\ Custódi nos, Dómine, ut pupíllam óculi. \R\ Sub umbra alárum tuárum prótege nos.

\antlabel{Ant.} Salva nos.

\canticumheading{Canticum Simeonis}{Luc. 2, 29-32}

\lettrine[lines=2]{N}{UNC} dimíttis servum tuum, Dómine, * secúndum verbum tuum in pace:

Quia vidérunt óculi mei * salutáre tuum,

Quod parásti * ante fáciem ómnium populórum,

Lumen ad revelatiónem géntium, * et glóriam plebis tuæ Israël.

\antlabel{Ant.} Salva nos, Dómine, vigilántes, custódi nos dormiéntes; ut vigilémus cum Christo, et requiescámus in pace.

\variantheading{Preces (si dicendæ sint)}

KYRIE, eléison. Christe, eléison. Kýrie, eléison.

Pater noster \rubric{secreto usque ad} \V\ Et ne nos indúcas in tentatiónem. \R\ Sed líbera nos a malo.

Credo in Deum \rubric{secreto usque ad} \V\ Carnis resurrectiónem. \R\ Vitam ætérnam. Amen. \V\ Benedíctus es, Dómine, Deus patrum nostrórum. \R\ Et laudábilis et gloriósus in sæcula. \V\ Benedicámus Patrem et Fílium cum Sancto Spíritu. \R\ Laudémus, et superexaltémus eum in sæcula. \V\ Benedíctus es, Dómine, in firmaménto cæli. \R\ Et laudábilis, et gloriósus, et superexaltátus in sæcula. \V\ Benedícat et custódiat nos omnípotens et miséricors Dóminus. \R\ Amen. \V\ Dignáre, Dómine, nocte ista. \R\ Sine peccáto nos custodíre. \V\ Miserére nostri, Dómine. \R\ Miserére nostri. \V\ Fiat misericórdia tua, Dómine, super nos. \R\ Quemádmodum sperávimus in te. \V\ Dómine, exáudi oratiónem meam. \R\ Et clamor meus ad te véniat.

\V\ Dóminus vobíscum.

Orémus. \oratiolabel{Oratio}

\lettrine[lines=2]{V}{ISITA}, quæsumus, Dómine, habitatiónem istam, et omnes insídias inimíci ab ea longe repélle: Angeli tui sancti hábitent in ea, qui nos in pace custódiant; et benedíctio tua sit super nos semper. Per Dóminum nostrum. \R\ Amen.

\V\ Dóminus vobíscum.

\V\ Benedicámus Dómino.

\benedictiolabel{Benedictio} Benedícat et custódiat nos omnípotens et miséricors Dóminus, Pater, et Fílius, et Spíritus Sanctus. \R\ Amen.

\variantheading{Antiphona finalis B. M. V.}

\V\ Divínum auxílium.

Pater, Ave María et Credo.

% ===== END OF BATCH 41 -- Covers autumnalis.md lines 13312-13647,
% completing "Ad Vesperas" (Psalms 143/i, 143/ii, 144/i-iii, with
% the special rubric noting these serve as First Vespers of the
% following Sunday) and "Ad Completorium" (Psalms 87, 102/i,
% 102/ii, hymn "Te lucis ante términum," Canticum Simeonis, Preces),
% and with them the ENTIRE "Sabbato" Office and the complete
% seven-day ferial Psalterium cycle of the Pars Autumnalis (Dominica
% through Sabbato). Saturday Vespers here carries its own fixed
% "In Officio feriali" Capitulum ("O Altitúdo divitiárum," rendered
% with the single-letter-word lettrine convention since the source
% gives only "O" capitalized before the normally-cased "Altitúdo")
% and hymn ("Jam sol recédit ígneus," with the recurring stray OCR
% "\*" before its doxology silently dropped), unlike the weekday
% Vespers pattern, since this Capitulum/Hymn pair is fixed for every
% Saturday rather than varying with the day. Two rubrics direct that
% the Magnificat antiphon and the following week's Oratio are drawn
% "ut in Próprio de Témpore" -- from the Proper of the Time for the
% upcoming Sunday -- confirming this Vespers functions as First
% Vespers of Sunday. The Canticum Simeonis here omits the printed
% "Glória Patri" line that appears at every other occurrence of this
% canticle in this volume; transcribed exactly as the source omits
% it. The document's Psalterium section ends here; autumnalis.md
% continues at line 13650 with an "[IMAGE]" placeholder (omitted, a
% non-textual OCR marker, preceded by a "---" scan-artifact rule
% also omitted) followed immediately by the heading "Proprium de
% Tempore" at line 13652, opening with "Sabbato proximo Kalendis
% Septembris" (the Saturday nearest the Kalends of September) --
% this volume's own starting point in the liturgical year, genuinely
% different from aestiva.tex's parallel transition into "In Festo
% Ssmæ Trinitatis" (Trinity Sunday), confirming the Proprium de
% Tempore itself is NOT shared boilerplate and must be transcribed
% fresh from here on -- a major structural transition from the
% ferial Psalterium into the Proper of the Time.
% =====
\partheading{Proprium de Tempore}

\feriaheading{Sabbato proximo Kalendis Septembris}

\hourheading{Ad Vesperas}

\V\ Vespertína orátio ascéndat ad te, Dómine. \R\ Et descéndat super nos misericórdia tua.

\antlabel{Ad Magnif. Ant.} Cum audísset Job * nuntiórum verba, sustínuit patiénter, et ait: Si bona suscépimus de manu Dómini, mala autem quare non sustineámus? In ómnibus his non peccávit Job lábiis suis, neque stultum áliquid contra Deum locútus est.

\rubric{Orátio Domínicæ occurréntis, quæ habétur infra post Lectiónes últimæ hebdómadæ Novémbris cum áliis Oratiónibus usque ad Domínicam XXIV post Pentecósten.}

\feriaheading{Dominica I Septembris}

\rubric{Semiduplex}

\hourheading{Ad Matutinum}

\rubric{Invitatórium} Dóminum, qui fecit, \rubric{et Hymnus} Nocte surgéntes, \rubric{ut in Psaltério. 31; quæ dicúntur in ómnibus Domínicis usque ad diem 27 Septémbris inclusíve occurréntibus.}

\nocturnoheading{In I Nocturno}

\rubric{Incipit liber Job}

\lectioheading{Lectio i}{Cap. 1, 1-11}

\lettrine[lines=2]{V}{IR} erat in terra Hus nómine Job, et erat vir ille simplex et rectus ac timens Deum et recédens a malo. Natíque sunt ei septem fílii et tres fíliæ. Et fuit posséssio ejus septem míllia óvium et tria míllia camelórum, quingénta quoque juga boum et quingéntæ ásinæ ac família multa nimis: erátque vir ille magnus inter omnes Orientáles.

\R\ Si bona suscépimus de manu Dei, mala autem quare non sustineámus? * Dóminus dedit, Dóminus ábstulit; sicut Dómino plácuit, ita factum est: sit nomen Dómini benedíctum. \V\ Nudus egréssus sum de útero matris meæ et nudus revértar illuc. — Dóminus.

\lectioheading{Lectio ii}{}

\lettrine[lines=2]{E}{T} ibant fílii ejus et faciébant convívium per domos, unusquísque in die suo. Et mitténtes vocábant tres soróres suas, ut coméderent et bíberent cum eis. Cumque in orbem transíssent dies convívii, mittébat ad eos Job et sanctificábat illos, consurgénsque dilúculo offerébat holocáusta pro síngulis; dicébat enim: Ne forte peccáverint fílii mei et benedíxerint Deo in córdibus suis. Sic faciébat Job cunctis diébus.

\R\ Antequam cómedam suspíro, et tamquam inundántes aquæ sic rugítus meus; quia timor, quem timébam, evénit mihi, et quod verébar áccidit. Nonne dissimulávi? nonne sílui? nonne quiévi? * Et venit super me indignátio. \V\ Ecce non est auxílium mihi in me, et necessárii quoque mei recessérunt a me. — Et.

\lectioheading{Lectio iii}{}

\lettrine[lines=2]{Q}{UADAM} autem die, cum veníssent fílii Dei ut assísterent coram Dómino, affúit inter eos étiam Satan. Cui dixit Dóminus: Unde venis? Qui respóndens ait: Circuívi terram et perambulávi eam. Dixítque Dóminus ad eum: Numquid considerásti servum meum Job, quod non sit ei símilis in terra, homo símplex et rectus ac timens Deum et recédens a malo? Cui respóndens Satan ait: Numquid Job frustra timet Deum? nonne tu vallásti eum ac domum ejus, universámque substántiam per circúitum opéribus mánuum ejus benedixísti, et posséssio ejus crevit in terra? Sed exténde páululum manum tuam et tange cuncta quæ póssidet, nisi in fáciem benedíxerit tibi.

\R\ Quare detraxístis sermónibus veritátis? ad increpándum verba compónitis et subvértere nitímini amícum vestrum: * Verúmtamen quæ cogitástis, expléte. \V\ Quod justum est, judicáte; et non inveniétis in lingua mea iniquitátem. — Verúmtamen. Glória Patri. Verúmtamen.

\nocturnoheading{In II Nocturno}

\rubric{Ex libro Morálium sancti Gregórii Papæ}

\lectioheading{Lectio iv}{Lib. 2, cap. 1}

\lettrine[lines=2]{S}{CRIPTÚRA} sacra mentis óculis quasi quoddam spéculum oppónitur, ut intérna nostra fácies in ipsa videátur. Ibi étenim foeda, ibi pulchra nostra cognóscimus: ibi sentímus quantum profícimus, ibi a proféctu quam longe distámus. Narrat autem gesta Sanctórum, et ad imitatiónem corda próvocat infirmórum; dumque illórum victrícia facta commémorat, contra vitiórum prǽlia, debília nostra confírmat: fitque verbis illíus, ut eo mens minus inter certámina trépidet, quo ante se pósitos tot virórum fórtium triúmphos videt.

\R\ Indúta est caro mea putrédine, et sórdibus púlveris cutis mea áruit et contrácta est: * Meménto mei, Dómine, quóniam ventus est vita mea. \V\ Dies mei velócius transiérunt quam a texénte tela succíditur, et consúmpti sunt absque ulla spe. — Meménto.

\lectioheading{Lectio v}{}

\lettrine[lines=2]{N}{ONNÚMQUAM} vero non solum nobis eórum virtútes ásserit, sed étiam casus innotéscit; ut et in victória fórtium, quod imitándo arrípere, et rursum videámus in lápsibus quid debeámus timére. Ecce enim Job describítur tentatióne auctus, sed David tentatióne prostrátus; ut et majórum virtus spem nostram fóveat, et majórum casus ad cautélam nos humilitátis accíngat: quátenus dum illa gaudéntes súblevant, ista metuéntes premant; et audiéntis ánimus illinc spei fidúcia, hinc humilitáte erudítus, nec temeritáte supérbiat, quia ad spei fidúciam virtútis exémplo roborátur.

\R\ Páucitas diérum meórum finiétur brevi; dimítte me, Dómine, ut plangam páululum dolórem meum, * Antequam vadam et non revértar, ad terram tenebrósam et opértam mortis calígine. \V\ Manus tuæ, Dómine, fecérunt me, et plasmavérunt me totum in circúitu; et sic repénte præcípitas me? — Antequam.

\lectioheading{Lectio vi}{Lib. 1, cap. 1}

\lettrine[lines=2]{V}{IR} erat in terra Hus, nómine Job. Idcírco sanctus vir ubi habitáverit dícitur, ut ejus méritum virtútis exprimátur. Hus namque quis nésciat quod sit in terra Gentílium? Gentílitas autem eo obligáta vítiis exstítit, quo cognitiónem sui Condítoris ignorávit. Dicátur ítaque ubi habitáverit; ut hoc ejus láudibus profíciat, quod bonus inter malos fuit. Neque enim valde laudábile est, bonum esse cum bonis, sed bonum esse cum malis. Sicut enim gravióris culpæ est, inter bonos bonum non esse; ita imménsi est præcónii, bonum étiam inter malos exstitísse.

\R\ Non abscóndas me, Dómine, a fácie tua: manum tuam longe fac a me, * Et formído tua non me térreat. \V\ Córripe me, Dómine, in misericórdia, non in furóre tuo, ne forte ad níhilum rédigas me. — Et. Glória Patri. Et.

\rubric{Lectiónes III Nocturni de Homilía Domínicæ occurréntis post Pentecósten habéntur inférius circa finem hujus Partis de Témpore, in órdine aliárum Dominicárum usque ad Domínicam XXIV post Pentecósten. Simíliter Antíphonæ ad Benedíctus et Magníficat, et Oratiónes.}

\R\ vii Quis mihi tríbuat, ut in inférno prótegas me et abscóndas me, donec pertránseat furor tuus, Dómine, nisi tu, qui solus es Deus? * Et constítuas mihi tempus, in quo recordéris mei? \V\ Numquid sicut dies hóminis dies tui, ut quæras iniquitátem meam; cum sit nemo, qui de manu tua possit erúere? — Et.

\R\ viii Duo Séraphim, \rubric{ut infra in Homilíis.}

\rubric{Ad Laudes, in hac et sequéntibus Domínicis usque ad diem 27 Septémbris inclusíve occurréntibus, dícitur Hymnus} Ecce jam noctis, \rubric{ut in Psaltério. 40}

% ===== BATCH 1 of the Proprium de Tempore (autumnalis.md lines
% 13652-13730) -- genuinely new Autumn-specific content, transcribed
% fresh from autumnalis.md since the Proprium de Tempore is NOT
% shared boilerplate across Pars volumes (confirmed by the previous
% batch's diff work showing aestiva.tex's parallel section opens
% with Trinity Sunday rather than this volume's own starting point).
% Covers the brief "Sabbato proximo Kalendis Septembris" (First
% Vespers only, with its Magnificat antiphon from the Book of Job
% and a forward-pointing rubric to the proper Sunday Oratio given
% later in this Part) and the complete "Dominica I Septembris"
% Matins: the Invitatorium/Hymnus cross-reference (rendered as
% alternating \rubric{} labels and live incipit text per the
% established compact-citation convention), the six Lessons and six
% Responsories of Nocturns I-II (Job's affliction and the opening of
% St. Gregory's Moralia on Job), and the final two numbered
% Responsories (vii, viii) that stand in for Nocturn III, whose own
% Lessons, antiphons, and Orationes are deferred by the source's own
% rubric to the Homiliae section near the end of this Part (a
% structural forward-reference, not a page/footnote number, so no
% \pageref{} resolution applies; it will be transcribed fresh when
% that later section is reached). The em-dash "—" repeat-cue
% notation in each Responsory (e.g. "— Dóminus.") is preserved
% literally as the source's own shorthand for repeating the
% refrain's cue word, a notation not previously encountered in this
% project's established responsory-repetition conventions (which
% elsewhere spell the repeated word out in full); since it is
% genuinely how autumnalis.md itself abbreviates these responsories,
% it is kept exactly as printed rather than silently expanded. The
% "---" horizontal rule before "Dominica I Septembris" is omitted as
% a scan artifact. The document continues with "Feria Secunda infra
% Hebdomadam I Septembris" (autumnalis.md line 13731).
% =====
\feriaheading{Feria Secunda infra Hebdomadam I Septembris}

\rubric{De libro Job}

\lectioheading{Lectio i}{Cap. 1, 13-22}

\lettrine[lines=2]{C}{UM} autem quadam die fílii et fíliæ ejus coméderent et bíberent vinum in domo fratris sui primogéniti, núntius venit ad Job, qui díceret: Boves arábant et ásinæ pascebántur juxta eos, et irruérunt Sabæi tulerúntque ómnia et púeros percussérunt gládio, et evási ego solus ut nuntiárem tibi. Cumque adhuc ille loquerétur, venit alter et dixit: Ignis Dei cécidit e cælo et tactas oves puerósque consúmpsit, et effúgi ego solus ut nuntiárem tibi.

\R\ Versa est in luctum cíthara mea, et órganum meum in vocem fléntium: * Parce mihi, Dómine, nihil enim sunt dies mei. \V\ Cutis mea denigráta est super me, et ossa mea aruérunt. — Parce.

\lectioheading{Lectio ii}{}

\lettrine[lines=2]{S}{ED}, et illo adhuc loquénte, venit álius et dixit: Chaldǽi fecérunt tres turmas et invasérunt camélos et tulérunt eos, necnon et púeros percussérunt gládio, et ego fugi solus ut nuntiárem tibi. Adhuc loquebátur ille, et ecce álius intrávit et dixit: Fíliis tuis et filiábus vescéntibus et bibéntibus vinum in domo fratris sui primogéniti, repénte ventus veheméns írruit a regióne desérti et concússit quátuor ángulos domus, quæ córruens oppréssit líberos tuos, et mórtui sunt, et effúgi ego solus ut nuntiárem tibi.

\R\ Utinam appenderéntur peccáta mea, quibus iram mérui, * Et calamítas, quam pátior, in statéra. \V\ Quasi aréna maris hæc grávior apparéret, unde et verba mea dolóre sunt plena. — Et.

\rubric{In Octávis et Festis ritus simplícis: Glória Patri. Et.}

\lectioheading{Lectio iii}{}

\lettrine[lines=2]{T}{UNC} surréxit Job et scidit vestiménta sua et, tonso cápite, córruens in terram adorávit et dixit: Nudus egréssus sum de útero matris meæ et nudus revértar illuc. Dóminus dedit, Dóminus ábstulit; sicut Dómino plácuit, ita factum est: sit nomen Dómini benedíctum. In ómnibus his non peccávit Job lábiis suis neque stultum quid contra Deum locútus est.

\R\ Quare detraxístis sermónibus veritátis? ad increpándum verba compónitis et subvértere nitímini amícum vestrum: * Verúmtamen quæ cogitástis, expléte. \V\ Quod justum est, judicáte; et non inveniétis in lingua mea iniquitátem. — Verúmtamen. Glória Patri. Verúmtamen.

\feriaheading{Feria Tertia infra Hebdomadam I Septembris}

\rubric{De libro Job}

\lectioheading{Lectio i}{Cap. 2, 1-13}

\lettrine[lines=2]{F}{ACTUM} est autem, cum quadam die veníssent fílii Dei et starent coram Dómino, venísset quoque Satan inter eos et staret in conspéctu ejus, ut díceret Dóminus ad Satan: Unde venis? Qui respóndens ait: Circuívi terram et perambulávi eam. Et dixit Dóminus ad Satan: Numquid considerásti servum meum Job, quod non sit ei símilis in terra, vir símplex et rectus ac timens Deum et recédens a malo et adhuc rétinens innocéntiam? Tu autem commovísti me advérsus eum, ut afflígerem eum frustra. Cui respóndens Satan ait: Pellem pro pelle, et cuncta, quæ habet homo, dabit pro ánima sua. Alióqui mitte manum tuam et tange os ejus et carnem; et tunc vidébis quod in fáciem benedícat tibi.

\R\ Indúta est caro mea putrédine, et sórdibus púlveris cutis mea áruit et contrácta est: * Meménto mei, Dómine, quóniam ventus est vita mea. \V\ Dies mei velócius transiérunt quam a texénte tela succíditur, et consúmpti sunt absque ulla spe. — Meménto.

\lectioheading{Lectio ii}{}

\lettrine[lines=2]{D}{IXIT} ergo Dóminus ad Satan: Ecce in manu tua est, verúmtamen ánimam illíus serva. Egréssus ígitur Satan a fácie Dómini percússit Job úlcere péssimo a planta pedis usque ad vérticem ejus; qui testa sániem radébat sedens in sterquilínio. Dixit autem illi uxor sua: Adhuc tu pérmanes in simplicitáte tua? Bénedic Deo et mórere. Qui ait ad illam: Quasi una de stultis muliéribus locúta es: si bona suscépimus de manu Dei, mala quare non suscipiámus? In ómnibus his non peccávit Job lábiis suis.

\R\ Páucitas diérum meórum finiétur brevi; dimítte me, Dómine, ut plangam páululum dolórem meum, * Antequam vadam ad terram tenebrósam et opértam mortis calígine. \V\ Manus tuæ, Dómine, fecérunt me, et plasmavérunt me totum in circúitu; et sic repénte præcípitas me? — Antequam.

\rubric{In Octávis et Festis ritus simplícis: Glória Patri. Antequam.}

\lectioheading{Lectio iii}{}

\lettrine[lines=2]{I}{GITUR} audiéntes tres amíci Job omne malum quod accidísset ei, venérunt sínguli de loco suo, Elíphaz Themanítes et Báldad Suhítes et Sophar Naamathítes. Condíxerant enim ut páriter veniéntes visitárent eum et consolaréntur. Cumque elevássent procul óculos suos, non cognovérunt eum et exclamántes ploravérunt, scissísque véstibus, spargérunt púlverem super caput suum in cælum. Et sedérunt cum eo in terra septem diébus et septem nóctibus, et nemo loquebátur ei verbum; vidébant enim dolórem esse veheméntem.

\R\ Non abscóndas me, Dómine, a fácie tua: manum tuam longe fac a me, * Et formído tua non me térreat. \V\ Córripe me, Dómine, in misericórdia, non in furóre tuo, ne forte ad níhilum rédigas me. — Et. Glória Patri. Et.

\feriaheading{Feria Quarta infra Hebdomadam I Septembris}

\rubric{De libro Job}

\lectioheading{Lectio i}{Cap. 3, 1-16}

\lettrine[lines=2]{P}{OST} hæc apéruit Job os suum et maledíxit diéi suo et locútus est: Péreat dies in qua natus sum, et nox in qua dictum est: Concéptus est homo. Dies ille vertátur in ténebras: non requírat eum Deus désuper, et non illustrétur lúmine; obscúrent eum ténebræ et umbra mortis; óccupet eum calígo, et involvátur amaritúdine.

\R\ Quis mihi tríbuat, ut in inférno prótegas me et abscóndas me, donec pertránseat furor tuus, Dómine, nisi tu, qui solus es Deus? * Et constítuas mihi tempus, in quo recordéris mei? \V\ Numquid sicut dies hóminis dies tui, ut quæras iniquitátem meam; cum sit nemo, qui de manu tua possit erúere? — Et.

\lectioheading{Lectio ii}{}

\lettrine[lines=2]{N}{OCTEM} illam tenebrósus turbo possídeat; non computétur in diébus anni, nec numerétur in ménsibus. Sit nox illa solitária nec laude digna; maledícant ei qui maledícunt diéi, qui paráti sunt suscitáre leviáthan. Obtenebréntur stellæ calígine ejus, exspéctet lucem et non vídeat, nec ortum surgéntis auróræ; quia non conclúsit óstia ventris, qui portávit me, nec ábstulit mala ab óculis meis.

\R\ Utinam appenderéntur peccáta mea, quibus iram mérui, * Et calamítas, quam pátior, in statéra. \V\ Quasi aréna maris hæc grávior apparéret, unde et verba mea dolóre sunt plena. — Et.

\rubric{In Octávis et Festis ritus simplícis: Glória Patri. Et.}

\lectioheading{Lectio iii}{}

\lettrine[lines=2]{Q}{UARE} non in vulva mórtuus sum? egréssus ex útero non statim périi? quare excéptus génibus? cur lactátus ubéribus? Nunc enim dórmiens silérem et somno meo requiéscerem cum régibus et consúlibus terræ, qui ædíficant sibi solitúdines, aut cum princípibus, qui póssident aurum et replent domos suas argénto; aut sicut abortívum abscónditum non subsísterem, vel qui concépti non vidérunt lucem.

\R\ Quare detraxístis sermónibus veritátis? ad increpándum verba compónitis et subvértere nitímini amícum vestrum: * Verúmtamen quæ cogitástis, expléte. \V\ Quod justum est, judicáte; et non inveniétis in lingua mea iniquitátem. — Verúmtamen. Glória Patri. Verúmtamen.

\feriaheading{Feria Quinta infra Hebdomadam I Septembris}

\rubric{De libro Job}

\lectioheading{Lectio i}{Cap. 4, 1-18}

\lettrine[lines=2]{R}{ESPÓNDENS} autem Elíphaz Themanítes dixit: Si cœpérimus loqui tibi, fórsitan moléste accípies; sed concéptum sermónem tenére quis póterit? Ecce docuísti multos et manus lassas roborásti, vacillántes confirmavérunt sermónes tui et génua treméntia confortásti. Nunc autem venit super te plaga, et defecísti; tétigit te, et conturbátus es. Ubi est timor tuus, fortitúdo tua, patiéntia tua et perféctio viárum tuárum?

\R\ Si bona suscépimus de manu Dei, mala autem quare non sustineámus? * Dóminus dedit, Dóminus ábstulit; sicut Dómino plácuit, ita factum est: sit nomen Dómini benedíctum. \V\ Nudus egréssus sum de útero matris meæ et nudus revértar illuc. — Dóminus.

\lectioheading{Lectio ii}{}

\lettrine[lines=2]{R}{ECORDÁRE}, óbsecro te, quis umquam ínnocens périit? aut quando recti deléti sunt? Quin pótius vidi eos, qui operántur iniquitátem et séminant dolóres et metunt eos, flante Deo perísse et spíritu iræ ejus esse consúmptos. Rugítus leónis et vox leǽnæ et dentes catulórum leónum contríti sunt; tigris périit, eo quod non habéret prædam, et cátuli leónis dissipáti sunt.

\R\ Antequam cómedam suspíro, et tamquam inundántes aquæ sic rugítus meus; quia timor, quem timébam, evénit mihi, et quod verébar áccidit. Nonne dissimulávi? nonne sílui? nonne quiévi? * Et venit super me indignátio. \V\ Ecce non est auxílium mihi in me, et necessárii quoque mei recessérunt a me. — Et.

\rubric{In Octávis et Festis ritus simplícis: Glória Patri. Et.}

\lectioheading{Lectio iii}{}

\lettrine[lines=2]{P}{ORRO} ad me dictum est verbum abscónditum, et quasi furtíve suscépit auris mea venas susúrri ejus. In horróre visiónis noctúrnæ, quando solet sopor occupáre hómines, pavor ténuit me et tremor, et ómnia ossa mea pertérrita sunt, et cum spíritus, me præsénte, transíret, inhorruérunt pili carnis meæ. Stetit quidam, cujus non agnoscébam vultum, imágo coram óculis meis, et vocem quasi auræ lenis audívi: Numquid homo Dei comparatióne justificábitur? aut Factóre suo púrior erit vir? Ecce qui sérviunt ei non sunt stábiles, et in Angelis suis réperit pravitátem.

\R\ Quare detraxístis sermónibus veritátis? ad increpándum verba compónitis et subvértere nitímini amícum vestrum: * Verúmtamen quæ cogitástis, expléte. \V\ Quod justum est, judicáte; et non inveniétis in lingua mea iniquitátem. — Verúmtamen. Glória Patri. Verúmtamen.

\feriaheading{Feria Sexta infra Hebdomadam I Septembris}

\rubric{De libro Job}

\lectioheading{Lectio i}{Cap. 6, 1-13}

\lettrine[lines=2]{R}{ESPÓNDENS} autem Job dixit: Utinam appenderéntur peccáta mea, quibus iram mérui, et calamítas, quam pátior, in statéra. Quasi aréna maris hæc grávior apparéret, unde et verba mea dolóre sunt plena; quia sagíttæ Dómini in me sunt, quarum indignátio ébibit spíritum meum, et terróres Dómini mílitant contra me.

\R\ Indúta est caro mea putrédine, et sórdibus púlveris cutis mea áruit et contrácta est: * Meménto mei, Dómine, quóniam ventus est vita mea. \V\ Dies mei velócius transiérunt quam a texénte tela succíditur, et consúmpti sunt absque ulla spe. — Meménto.

\lectioheading{Lectio ii}{}

\lettrine[lines=2]{N}{UMQUID} rúgiet ónager cum habúerit herbam? aut mugíet bos, cum ante præsépe plenum stéterit? Aut póterit comédi insúlsum, quod non est sale condítum? aut potest áliquis gustáre, quod gustátum affert mortem? Quæ prius nolébat tángere ánima mea, nunc præ angústia cibi mei sunt.

\R\ Páucitas diérum meórum finiétur brevi; dimítte me, Dómine, ut plangam páululum dolórem meum, * Antequam vadam ad terram tenebrósam et opértam mortis calígine. \V\ Manus tuæ, Dómine, fecérunt me, et plasmavérunt me totum in circúitu; et sic repénte præcípitas me? — Antequam.

\rubric{In Octávis et Festis ritus simplícis: Glória Patri. Antequam.}

\lectioheading{Lectio iii}{}

\lettrine[lines=2]{Q}{UIS} det ut véniat petítio mea, et quod exspécto tríbuat mihi Deus? Et qui cœpit ipse me cónterat, solvat manum suam et succídat me? Et hæc mihi sit consolátio, ut afflígens me dolóre non parcat, nec contradícam sermónibus Sancti. Quæ est enim fortitúdo mea, ut sustíneam? aut quis finis meus, ut patiénter agam? Nec fortitúdo lápidum fortitúdo mea, nec caro mea ǽnea est. Ecce non est auxílium mihi in me, et necessárii quoque mei recessérunt a me.

\R\ Non abscóndas me, Dómine, a fácie tua: manum tuam longe fac a me, * Et formído tua non me térreat. \V\ Córripe me, Dómine, in misericórdia, non in furóre tuo, ne forte ad níhilum rédigas me. — Et. Glória Patri. Et.

\rubric{Ad Vésperas, nisi sequénti die persolvéndum sit Offícium de Festo novem Lectiónum, aut de qualíbet Octáva vel de Vigília, a Capítulo fit de sancta María, et de ea dícitur Offícium in Sábbato sequénti. Quod ítem servátur in ómnibus Sábbatis usque ad Sábbatum ante Domínicam I Advéntus inclusíve, præterquam in Sábbato Quátuor Témporum Septémbris, et in Sábbato, in quo anticipánda sit Domínica XXIII post Pentecósten.}

\feriaheading{Sabbato infra Hebdomadam I Septembris}

\rubric{De libro Job}

\lectioheading{Lectio i}{Cap. 7, 1-12}

\lettrine[lines=2]{M}{ILITIA} est vita hóminis super terram, et sicut dies mercenárii, dies ejus. Sicut servus desíderat umbram, et sicut mercenárius præstolátur finem óperis sui, sic et ego hábui menses vácuos et noctes laboriósas enumerávi mihi. Si dormíero, dicam: Quando consúrgam? Et rursum exspectábo vésperam, et replébor dolóribus usque ad ténebras.

\R\ Quis mihi tríbuat, ut in inférno prótegas me et abscóndas me, donec pertránseat furor tuus, Dómine, nisi tu, qui solus es Deus? * Et constítuas mihi tempus, in quo recordéris mei? \V\ Numquid sicut dies hóminis dies tui, ut quæras iniquitátem meam; cum sit nemo, qui de manu tua possit erúere? — Et.

\lectioheading{Lectio ii}{}

\lettrine[lines=2]{I}{NDÚTA} est caro mea putrédine, et sórdibus púlveris cutis mea áruit et contrácta est. Dies mei velócius transiérunt quam a texénte tela succíditur, et consúmpti sunt absque ulla spe. Meménto quia ventus est vita mea, et non revertétur óculus meus ut vídeat bona. Nec aspíciet me visus hóminis; óculi tui in me, et non subsístam.

\R\ Utinam appenderéntur peccáta mea, quibus iram mérui, * Et calamítas, quam pátior, in statéra. \V\ Quasi aréna maris hæc grávior apparéret, unde et verba mea dolóre sunt plena. — Et.

\rubric{In Octávis simplícibus et in Offício sanctæ Maríæ in Sábbato: Glória Patri. Et.}

\lectioheading{Lectio iii}{}

\lettrine[lines=2]{S}{ICUT} consúmitur nubes et pertránsit, sic qui descénderit ad ínferos non ascéndet; nec revertétur ultra in domum suam, neque cognóscet eum ámplius locus ejus. Quaprópter et ego non parcam ori meo; loquar in tribulatióne spíritus mei, confabulábor cum amaritúdine ánimæ meæ. Numquid mare ego sum aut cetus, quia circumdedísti me cárcere?

\R\ Quare detraxístis sermónibus veritátis? ad increpándum verba compónitis et subvértere nitímini amícum vestrum: * Verúmtamen quæ cogitástis, expléte. \V\ Quod justum est, judicáte; et non inveniétis in lingua mea iniquitátem. — Verúmtamen. Glória Patri. Verúmtamen.

\hourheading{Ad Vesperas}

\V\ Vespertína orátio ascéndat ad te, Dómine. \R\ Et descéndat super nos misericórdia tua.

\antlabel{Ad Magnif. Ant.} In ómnibus his * non peccávit Job lábiis suis neque stultum áliquid contra Deum locútus est.

\rubric{Orátio Domínicæ occurréntis.}

% ===== BATCH 2 of the Proprium de Tempore (autumnalis.md lines
% 13731-13902) -- completes the first full week of September: the
% six ferial days' Matins lessons (three Lessons and three
% Responsories apiece) continuing the Book of Job cycle through
% chapters 1-7, and the closing Vespers for "Sabbato infra
% Hebdomadam I Septembris" with its own Magnificat antiphon and a
% forward-pointing rubric to the proper Sunday Oratio. The lengthy
% governing rubric on Saturday Vespers switching to the Office of
% the Blessed Virgin Mary (except on the September Ember Saturday
% and certain anticipated Sundays) is transcribed in full as a
% single \rubric{} block. Three "In Octávis..." rubrics mark the
% alternate "Glória Patri" cue for feast-day/Marian-Saturday ranks;
% each is transcribed exactly as the source gives it, including the
% third day's distinct wording ("et in Officio sanctae Mariae in
% Sabbato" rather than "et Festis ritus simplicis"). Every
% Responsory's em-dash repeat-cue is preserved literally, consistent
% with Batch 1. The "---" horizontal rules before "Sabbato infra
% Hebdomadam I Septembris" and before "Dominica II Septembris" are
% omitted as scan artifacts. The document continues with "Dominica
% II Septembris" (autumnalis.md line 13903).
% =====
\feriaheading{Dominica II Septembris}

\rubric{Semiduplex}

\nocturnoheading{In I Nocturno}

\rubric{De libro Job}

\lectioheading{Lectio i}{Cap. 9, 1-17}

\lettrine[lines=2]{E}{T} respóndens Job ait: Vere scio quod ita sit, et quod non justificétur homo compósitus Deo. Si volúerit conténdere cum eo, non póterit ei respondére unum pro mille. Sápiens corde est et fortis róbore. Quis réstitit ei, et pacem hábuit? Qui tránstulit montes, et nesciérunt hi quos subvértit in furóre suo.

\R\ Si bona suscépimus de manu Dei, mala autem quare non sustineámus? * Dóminus dedit, Dóminus ábstulit; sicut Dómino plácuit, ita factum est: sit nomen Dómini benedíctum. \V\ Nudus egréssus sum de útero matris meæ et nudus revértar illuc. — Dóminus.

\lectioheading{Lectio ii}{}

\lettrine[lines=2]{Q}{UI} cómmovet terram de loco suo, et colúmnæ ejus concutiúntur; qui præcípit soli, et non óritur, et stellas claudit quasi sub signáculo; qui exténdit cælos solus, et gráditur super fluctus maris; qui facit Arctúrum et Oríona et Hýadas et interióra Austri; qui facit magna et incomprehensibília et mirabília, quorum non est númerus.

\R\ Antequam cómedam suspíro, et tamquam inundántes aquæ sic rugítus meus; quia timor, quem timébam, evénit mihi, et quod verébar áccidit. Nonne dissimulávi? nonne sílui? nonne quiévi? * Et venit super me indignátio. \V\ Ecce non est auxílium mihi in me, et necessárii quoque mei recessérunt a me. — Et.

\lectioheading{Lectio iii}{}

\lettrine[lines=2]{S}{I} vénerit ad me, non vidébo eum; si abíerit, non intélligam; si repénte intérroget, quis respondébit ei vel quis dícere potest: Cur ita facis? Deus, cujus iræ nemo resístere potest, et sub quo curvántur qui portant orbem; quantus ergo sum ego, ut respóndeam ei, et loquar verbis meis cum eo? Qui, étiam si habúero quíppiam justum, non respondébo, sed meum júdicem deprecábor. Et, cum invocántem exaudíerit me, non credo quod audíerit vocem meam. In túrbine enim cónteret me, et multiplicábit vúlnera mea, étiam sine causa.

\R\ Quare detraxístis sermónibus veritátis? ad increpándum verba compónitis et subvértere nitímini amícum vestrum: * Verúmtamen quæ cogitástis, expléte. \V\ Quod justum est, judicáte; et non inveniétis in lingua mea iniquitátem. — Verúmtamen. Glória Patri. Verúmtamen.

\nocturnoheading{In II Nocturno}

\rubric{Ex libro Morálium sancti Gregórii Papæ}

\lectioheading{Lectio iv}{Liber 9, cap. 2}

\lettrine[lines=2]{V}{ERE} scio quod ita sit, et quod non justificábitur homo compósitus Deo. Homo quippe Deo non compósitus justítiam pércipit, compósitus amíttit; quia quisquis se auctóri bonórum cómparat, bono se, quod accéperat, privat. Qui enim accépta bona sibi árrogat, suis contra Deum donis pugnat. Unde ergo despéctus erígitur, dignum est ut eréctus inde destruátur. Sanctus autem vir, quia omne virtútis nostræ méritum esse vítium cónspicit, si ab intérno árbitro distrícte judicétur, recte subjúngit: Si volúerit conténdere cum eo, non póterit respondére ei unum pro mille.

\R\ Indúta est caro mea putrédine, et sórdibus púlveris cutis mea áruit et contrácta est: * Meménto mei, Dómine, quóniam ventus est vita mea. \V\ Dies mei velócius transiérunt quam a texénte tela succíditur, et consúmpti sunt absque ulla spe. — Meménto.

\lectioheading{Lectio v}{}

\lettrine[lines=2]{I}{N} Scriptúra sancta millenárius númerus pro universitáte solet intélligi. Hinc étenim Psalmísta ait: Verbi, quod mandávit in mille generatiónes: cum profécto constet, quod ab ipso mundi exórdio usque ad Redemptóris advéntum per Evangelístam non ámplius quam septuagínta et septem propágines numeréntur. Quid ergo in millenário número nisi ad proferéndam novam sóbolem perfécta universitas præscítæ generatiónis exprímitur? Hinc et per Joánnem dícitur: Et regnábunt cum eo mille annis; quia vidélicet regnum sanctæ Ecclésiæ universitátis perfectióne solidátur.

\R\ Páucitas diérum meórum finiétur brevi; dimítte me, Dómine, ut plangam páululum dolórem meum, * Antequam vadam ad terram tenebrósam et opértam mortis calígine. \V\ Manus tuæ, Dómine, fecérunt me, et plasmavérunt me totum in circúitu; et sic repénte præcípitas me? — Antequam.

\lectioheading{Lectio vi}{}

\lettrine[lines=2]{Q}{UIA} vero monas décies multiplicáta in denárium dúcitur, denárius per semetípsum ductus in centenárium dilatátur, qui rursus per denárium ductus in millenárium ténditur; cum ab uno incípimus, ut ad millenárium veniámus, quid hoc loco uníus appellatióne, nisi bene vivéndi inítium? quid millenárii númeri amplitúdine, nisi ejúsdem bonæ vitæ perféctio designátur? Cum Deo autem conténdere, est, non ei tribúere, sed sibi glóriam suæ virtútis arrogáre. Sed sanctus vir conspíciat, quia et qui summa jam dona percépit, si de accéptis extóllitur, cuncta quæ accéperat, amíttit.

\R\ Non abscóndas me, Dómine, a fácie tua: manum tuam longe fac a me, * Et formído tua non me térreat. \V\ Córripe me, Dómine, in misericórdia, non in furóre tuo, ne forte ad níhilum rédigas me. — Et. Glória Patri. Et.

\R\ vii Quis mihi tríbuat, ut in inférno prótegas me et abscóndas me, donec pertránseat furor tuus, Dómine, nisi tu, qui solus es Deus? * Et constítuas mihi tempus, in quo recordéris mei? \V\ Numquid sicut dies hóminis dies tui, ut quæras iniquitátem meam; cum sit nemo, qui de manu tua possit erúere? — Et.

\R\ viii Duo Séraphim, \rubric{ut infra in Homilíis.}

\feriaheading{Feria Secunda infra Hebdomadam II Septembris}

\rubric{De libro Job}

\lectioheading{Lectio i}{Cap. 27, 1-15}

\lettrine[lines=2]{A}{DDIDIT} quoque Job, assúmens parábolam suam, et dixit: Vivit Deus, qui ábstulit judícium meum, et Omnípotens, qui ad amaritúdinem addúxit ánimam meam. Quia, donec súperest hálitus in me, et spíritus Dei in náribus meis, non loquéntur lábia mea iniquitátem, nec lingua mea meditábitur mendácium. Absit a me, ut justos vos esse júdicem; donec defíciam, non recédam ab innocéntia mea.

\R\ Versa est in luctum cíthara mea, et órganum meum in vocem fléntium: * Parce mihi, Dómine, nihil enim sunt dies mei. \V\ Cutis mea denigráta est super me, et ossa mea aruérunt. — Parce.

\lectioheading{Lectio ii}{}

\lettrine[lines=2]{J}{USTIFICATIÓNEM} meam, quam cœpi tenére, non déseram; neque enim reprehéndit me cor meum in omni vita mea. Sit ut ímpius inimícus meus, et adversárius meus quasi iníquus. Quæ est enim spes hypócritæ, si aváre rápiat, et non líberet Deus ánimam ejus? Numquid Deus áudiet clamórem ejus, cum vénerit super eum angústia? Aut póterit in Omnipoténte delectári, et invocáre Deum omni témpore?

\R\ Utinam appenderéntur peccáta mea, quibus iram mérui, * Et calamítas, quam pátior, in statéra. \V\ Quasi aréna maris hæc grávior apparéret, unde et verba mea dolóre sunt plena. — Et.

\rubric{In Octávis et Festis ritus simplícis: Glória Patri. Et.}

\lectioheading{Lectio iii}{}

\lettrine[lines=2]{D}{OCÉBO} vos per manum Dei quæ Omnípotens hábeat, nec abscóndam. Ecce vos omnes nostis, et quid sine causa vana loquímini? Hæc est pars hóminis ímpii apud Deum, et heréditas violentórum, quam ab Omnipoténte suscípient: si multiplicáti fúerint fílii ejus, in gládio erunt, et nepótes ejus non saturabúntur pane; qui réliqui fúerint ex eo sepeliéntur in intéritu, et víduæ illíus non plorábunt.

\R\ Quare detraxístis sermónibus veritátis? ad increpándum verba compónitis et subvértere nitímini amícum vestrum: * Verúmtamen quæ cogitástis, expléte. \V\ Quod justum est, judicáte; et non inveniétis in lingua mea iniquitátem. — Verúmtamen. Glória Patri. Verúmtamen.

\feriaheading{Feria Tertia infra Hebdomadam II Septembris}

\rubric{De libro Job}

\lectioheading{Lectio i}{Cap. 28, 12-28}

\lettrine[lines=2]{S}{APIÉNTIA} ubi invenítur? et quis est locus intelligéntiæ? Nescit homo prétium ejus, nec invenítur in terra suáviter vivéntium. Abýssus dicit: Non est in me, et mare lóquitur: Non est mecum. Non dábitur aurum obrízum pro ea, nec appendétur argéntum in commutatióne ejus; non conferétur tinctis Indiæ colóribus, nec lápidi sardónycho pretiosíssimo vel sapphíro.

\R\ Indúta est caro mea putrédine, et sórdibus púlveris cutis mea áruit et contrácta est: * Meménto mei, Dómine, quóniam ventus est vita mea. \V\ Dies mei velócius transiérunt quam a texénte tela succíditur, et consúmpti sunt absque ulla spe. — Meménto.

\lectioheading{Lectio ii}{}

\lettrine[lines=2]{N}{ON} adæquábitur ei aurum vel vitrum, nec commutabúntur pro ea vasa auri. Excélsa et eminéntia non memorabúntur comparatióne ejus; tráhitur autem sapiéntia de occúltis. Non adæquábitur ei topázius de Æthiópia, nec tinctúræ mundíssimæ componétur. Unde ergo sapiéntia venit? et quis est locus intelligéntiæ? Abscóndita est ab óculis ómnium vivéntium; vólucres quoque cæli latet. Perdítio et mors dixérunt: Áuribus nostris audívimus famam ejus.

\R\ Páucitas diérum meórum finiétur brevi; dimítte me, Dómine, ut plangam páululum dolórem meum, * Antequam vadam ad terram tenebrósam et opértam mortis calígine. \V\ Manus tuæ, Dómine, fecérunt me, et plasmavérunt me totum in circúitu; et sic repénte præcípitas me? — Antequam.

\rubric{In Octávis et Festis ritus simplícis: Glória Patri. Antequam.}

\lectioheading{Lectio iii}{}

\lettrine[lines=2]{D}{EUS} intélligit viam ejus, et ipse novit locum illíus; ipse enim fines mundi intuétur, et ómnia, quæ sub cælo sunt, réspicit. Qui fecit ventis pondus, et aquas appéndit in mensúra. Quando ponébat plúviis legem et viam procéllis sonántibus, tunc vidit illam et enarrávit et præparávit et investigávit. Et dixit hómini: Ecce timor Dómini, ipsa est sapiéntia, et recédere a malo, intelligéntia.

\R\ Non abscóndas me, Dómine, a fácie tua: manum tuam longe fac a me, * Et formído tua non me térreat. \V\ Córripe me, Dómine, in misericórdia, non in furóre tuo, ne forte ad níhilum rédigas me. — Et. Glória Patri. Et.

\feriaheading{Feria Quarta infra Hebdomadam II Septembris}

\rubric{De libro Job}

\lectioheading{Lectio i}{Cap. 31, 1-18}

\lettrine[lines=2]{P}{EPIGI} fœdus cum óculis meis, ut ne cogitárem quidem de vírgine. Quam enim partem habéret in me Deus désuper, et hereditátem Omnípotens de excélsis? Numquid non perdítio est iníquo, et alienátio operántibus injustítiam? Nonne ipse considérat vias meas, et cunctos gressus meos dinúmerat? Si ambulávi in vanitáte, et festinávit in dolo pes meus, appéndat me in statéra justa, et sciat Deus simplicitátem meam.

\R\ Quis mihi tríbuat, ut in inférno prótegas me et abscóndas me, donec pertránseat furor tuus, Dómine, nisi tu, qui solus es Deus? * Et constítuas mihi tempus, in quo recordéris mei? \V\ Numquid sicut dies hóminis dies tui, ut quæras iniquitátem meam; cum sit nemo, qui de manu tua possit erúere? — Et.

\lectioheading{Lectio ii}{}

\lettrine[lines=2]{S}{I} declinávit gressus meus de via, et si secútum est óculos meos cor meum, et si mánibus meis adhǽsit mácula, seram, et álius cómedat, et progénies mea eradicétur. Si decéptum est cor meum super mulíere, et si ad óstium amíci mei insidiátus sum, scortum altérius sit uxor mea, et super illam incurvéntur álii. Hoc enim nefas est et iníquitas máxima; ignis est usque ad perditiónem dévorans et ómnia eradícans genímina.

\R\ Utinam appenderéntur peccáta mea, quibus iram mérui, * Et calamítas, quam pátior, in statéra. \V\ Quasi aréna maris hæc grávior apparéret, unde et verba mea dolóre sunt plena. — Et.

\rubric{In Octávis et Festis ritus simplícis: Glória Patri. Et.}

\lectioheading{Lectio iii}{}

\lettrine[lines=2]{S}{I} contémpsi subíre judícium cum servo meo et ancílla mea, cum disceptárent advérsum me: quid enim fáciam, cum surréxerit ad judicándum Deus? et, cum quæsíerit, quid respondébo illi? Numquid non in útero fecit me, qui et illum operátus est? et formávit me in vulva unus? Si negávi, quod volébant, paupéribus et óculos víduæ exspectáre feci; si comédi buccéllam meam solus, et non comédit pupíllus ex ea (quia ab infántia mea crevit mecum miserátio, et de útero matris meæ egréssa est mecum).

\R\ Quare detraxístis sermónibus veritátis? ad increpándum verba compónitis et subvértere nitímini amícum vestrum: * Verúmtamen quæ cogitástis, expléte. \V\ Quod justum est, judicáte; et non inveniétis in lingua mea iniquitátem. — Verúmtamen. Glória Patri. Verúmtamen.

\feriaheading{Feria Quinta infra Hebdomadam II Septembris}

\rubric{De libro Job}

\lectioheading{Lectio i}{Cap. 38, 1-20}

\lettrine[lines=2]{R}{ESPÓNDENS} autem Dóminus Job de túrbine, dixit: Quis est iste invólvens senténtias sermónibus imperítis? Accínge, sicut vir, lumbos tuos; interrogábo te et respónde mihi. Ubi eras quando ponébam fundaménta terræ? Índica mihi, si habes intelligéntiam. Quis pósuit mensúras ejus, si nosti? vel quis teténdit super eam líneam? Super quo bases illíus solidátæ sunt? aut quis demísit lápidem angulárem ejus, cum me laudárent simul astra matutína, et jubilárent omnes fílii Dei?

\R\ Si bona suscépimus de manu Dei, mala autem quare non sustineámus? * Dóminus dedit, Dóminus ábstulit; sicut Dómino plácuit, ita factum est: sit nomen Dómini benedíctum. \V\ Nudus egréssus sum de útero matris meæ et nudus revértar illuc. — Dóminus.

\lectioheading{Lectio ii}{}

\lettrine[lines=2]{Q}{UIS} conclúsit óstiis mare, quando erumpébat quasi de vulva procédens, cum ponérem nubem vestiméntum ejus, et calígine illud, quasi pannis infántiæ, obvólverem? Circúmdedi illud términis meis, et pósui vectem et óstia, et dixi: Usque huc vénies, et non procédes ámplius, et hic confrínges tuméntes fluctus tuos. Numquid post ortum tuum præcepísti dilúculo, et ostendísti aurórae locum suum? Et tenuísti concútiens extréma terræ, et excussísti ímpios ex ea?

\R\ Antequam cómedam suspíro, et tamquam inundántes aquæ sic rugítus meus; quia timor, quem timébam, evénit mihi, et quod verébar áccidit. Nonne dissimulávi? nonne sílui? nonne quiévi? * Et venit super me indignátio. \V\ Ecce non est auxílium mihi in me et necessárii quoque mei recessérunt a me. — Et.

\rubric{In Octávis et Festis ritus simplícis: Glória Patri. Et.}

\lectioheading{Lectio iii}{Cap. 42, 1-6}

\lettrine[lines=2]{R}{ESTITUÉTUR} ut lutum signáculum, et stabit sicut vestiméntum: auferétur ab ímpiis lux sua, et bráchium excélsum confringétur? Numquid ingréssus es profúnda maris, et in novíssimis abýssi deambulásti? Numquid apértæ sunt tibi portæ mortis, et óstia tenebrósa vidísti? Numquid considerásti latitúdinem terræ? Índica mihi, si nosti, ómnia: in qua via lux hábitet, et tenebrárum quis locus sit; ut ducas unumquódque ad términos suos, et intélligas sémitas domus ejus.

\R\ Quare detraxístis sermónibus veritátis? ad increpándum verba compónitis et subvértere nitímini amícum vestrum: * Verúmtamen quæ cogitástis, expléte. \V\ Quod justum est, judicáte; et non inveniétis in lingua mea iniquitátem. — Verúmtamen. Glória Patri. Verúmtamen.

\feriaheading{Feria Sexta infra Hebdomadam II Septembris}

\rubric{De libro Job}

\lectioheading{Lectio i}{Cap. 40, 1-11}

\lettrine[lines=2]{R}{ESPÓNDENS} autem Dóminus Job de túrbine, dixit: Accínge sicut vir lumbos tuos; interrogábo te, et índica mihi. Numquid írritum fácies judícium meum, et condemnábis me, ut tu justificéris? Et, si habes bráchium sicut Deus, et si voce símili tonas, circúmda tibi decórem et in sublíme erígere, et esto gloriósus et speciósis indúere véstibus.

\R\ Indúta est caro mea putrédine, et sórdibus púlveris cutis mea áruit et contrácta est: * Meménto mei, Dómine, quóniam ventus est vita mea. \V\ Dies mei velócius transiérunt quam a texénte tela succíditur, et consúmpti sunt absque ulla spe. — Meménto.

\lectioheading{Lectio ii}{}

\lettrine[lines=2]{D}{ISPÉRGE} supérbos in furóre tuo, et respíciens omnem arrogántem humília; réspice cunctos supérbos et confúnde eos, cóntere ímpios in loco suo; abscónde eos in púlvere simul, et fácies eórum demérge in fóveam; et ego confitébor quod salváre te possit déxtera tua. Ecce béhemoth, quem feci tecum, fœnum, quasi bos, cómedet. Fortitúdo ejus in lumbis ejus, et virtus illíus in umbilíco ventris ejus.

\R\ Páucitas diérum meórum finiétur brevi; dimítte me, Dómine, ut plangam páululum dolórem meum, * Antequam vadam ad terram tenebrósam et opértam mortis calígine. \V\ Manus tuæ, Dómine, fecérunt me, et plasmavérunt me totum in circúitu; et sic repénte præcípitas me? — Antequam.

\rubric{In Octávis et Festis ritus simplícis: Glória Patri. Antequam.}

\lectioheading{Lectio iii}{Cap. 42, 1-6}

\lettrine[lines=2]{R}{ESPÓNDENS} autem Job Dómino dixit: Scio quia ómnia potes, et nulla te latet cogitátio. Quis est iste qui celat consílium absque sciéntia? Ideo insipiénter locútus sum, et quæ ultra modum excéderent sciéntiam meam. Audi, et ego loquar; interrogábo te, et respónde mihi. Audítu auris audívi te; nunc autem óculus meus videt te. Idcírco ipse me reprehéndo, et ago pœniténtiam in favílla et cínere.

\R\ Non abscóndas me, Dómine, a fácie tua: manum tuam longe fac a me, * Et formído tua non me térreat. \V\ Córripe me, Dómine, in misericórdia, non in furóre tuo, ne forte ad níhilum rédigas me. — Et. Glória Patri. Et.

\feriaheading{Sabbato infra Hebdomadam II Septembris}

\rubric{De libro Job}

\lectioheading{Lectio i}{Cap. 42, 7-16}

\lettrine[lines=2]{P}{OSTQUAM} autem locútus est Dóminus verba hæc ad Job, dixit ad Elíphaz Themanítem: Irátus est furor meus in te et in duos amícos tuos, quóniam non estis locúti coram me rectum, sicut servus meus Job. Súmite ergo vobis septem tauros et septem aríetes et ite ad servum meum Job et offérte holocáustum pro vobis; Job autem servus meus orábit pro vobis. Fáciem ejus suscípiam, ut non vobis imputétur stultítia; neque enim locúti estis ad me recta, sicut servus meus Job.

\R\ Quis mihi tríbuat, ut in inférno prótegas me et abscóndas me, donec pertránseat furor tuus, Dómine, nisi tu, qui solus es Deus? * Et constítuas mihi tempus, in quo recordéris mei? \V\ Numquid sicut dies hóminis dies tui, ut quæras iniquitátem meam; cum sit nemo, qui de manu tua possit erúere? — Et constítuas mihi tempus, in quo recordéris mei?

\lectioheading{Lectio ii}{}

\lettrine[lines=2]{A}{BIÉRUNT} ergo Elíphaz Themanítes et Báldad Suhítes et Sophar Naamathítes, et fecérunt sicut locútus fúerat Dóminus ad eos, et suscépit Dóminus fáciem Job. Dóminus quoque convérsus est ad pœniténtiam Job, cum oráret ille pro amícis suis; et áddidit Dóminus ómnia, quæcúmque fúerant Job, duplícia. Venérunt autem ad eum omnes fratres sui et univérsæ soróres suæ cuncti, qui nóverant eum prius, et comedérunt cum eo panem in domo ejus; et movérunt super eum caput, et consoláti sunt eum super omni malo, quod intúlerat Dóminus super eum, et dedérunt ei unusquísque ovem unam et ináurem áuream unam.

\R\ Utinam appenderéntur peccáta mea, quibus iram mérui, * Et calamítas, quam pátior, in statéra. \V\ Quasi aréna maris hæc grávior apparéret, unde et verba mea dolóre sunt plena. — Et.

\rubric{In Octávis simplícibus et in Offício sanctæ Maríæ in Sábbato: Glória Patri. Et.}

\lectioheading{Lectio iii}{}

\lettrine[lines=2]{D}{OMINUS} autem benedíxit novíssimis Job magis quam princípio ejus; et facta sunt ei quatuórdecim míllia óvium et sex míllia camelórum et mille juga boum et mille ásinæ, et fuérunt ei septem fílii et tres fíliæ, et vocávit nomen uníus Diem et nomen secúndæ Cássiam et nomen tértiæ Cornustíbii. Non sunt autem invéntæ mulíeres speciósæ sicut fíliæ Job in univérsa terra; dedítque eis pater suus hereditátem inter fratres eárum. Vixit autem Job post hæc centum quadragínta annis, et vidit fílios suos et fílios filiórum suórum usque ad quartam generatiónem; et mórtuus est senex et plenus diérum.

\R\ Quare detraxístis sermónibus veritátis? ad increpándum verba compónitis et subvértere nitímini amícum vestrum: * Verúmtamen quæ cogitástis, expléte. \V\ Quod justum est, judicáte; et non inveniétis in lingua mea iniquitátem. — Verúmtamen. Glória Patri. Verúmtamen.

\hourheading{Ad Vesperas}

\V\ Vespertína orátio ascéndat ad te, Dómine. \R\ Et descéndat super nos misericórdia tua.

\antlabel{Ad Magnif. Ant.} Ne reminiscáris, * Dómine, delícta mea vel paréntum meórum, neque vindíctam sumas de peccátis meis.

\rubric{Oratio Dominicæ occurrentis.}

% ===== BATCH 3 of the Proprium de Tempore (autumnalis.md lines
% 13903-14129) -- the complete second week of September: "Dominica
% II Septembris" Matins (Nocturns I-II continuing Job through
% chapter 9 and St. Gregory's Moralia Book 9, with the same deferred
% Nocturn III and numbered Responsories vii-viii pattern as Dominica
% I), the six ferial days' Matins lessons carrying the Job cycle
% through its conclusion (chapters 27-42, ending with Job's
% restoration), and the closing Saturday Vespers. Two genuine source
% variants from the established responsory pattern are preserved
% exactly: Feria Quinta's third lesson and Feria Sexta's third
% lesson both cite "Cap. 42, 1-6" as a fresh citation (rather than
% continuing the prior lesson's citation with an empty second
% argument), since the source itself restates the chapter reference
% at each occurrence; and the Sabbato's first Responsory repeats its
% FULL verse text after the em-dash ("— Et constítuas mihi tempus,
% in quo recordéris mei?") rather than the usual single-word cue,
% transcribed exactly as the source gives it. The "---" horizontal
% rules between days are omitted as scan artifacts throughout. The
% document continues with "Dominica III Septembris" (autumnalis.md
% line 14130).
% =====
\feriaheading{Dominica III Septembris}

\rubric{Semiduplex}

\nocturnoheading{In I Nocturno}

\rubric{Incipit liber Tobiæ}

\lectioheading{Lectio i}{Cap. 1, 1-15}

\lettrine[lines=2]{T}{OBIAS} ex tribu et civitáte Néphthali, quæ est in superióribus Galilǽæ supra Naásson, post viam quæ ducit ad Occidéntem in sinístro habens civitátem Sephet, cum captus esset in diébus Salmánasar regis Assyriórum, in captivitáte tamen pósitus viam veritátis non deséruit; ita ut ómnia, quæ habére póterat, quotídie concaptívis frátribus, qui erant ex ejus génere, impertíret. Cumque esset júnior ómnibus in tribu Néphthali, nihil tamen puerile gessit in ópere.

\R\ Peto, Dómine, ut de vínculo improperii hujus absólvas me, aut certe désuper terram erípias me: * Ne reminiscáris delícta mea vel paréntum meórum, neque vindíctam sumas de peccátis meis: quia éruis sustinéntes te, Dómine. \V\ Ómnia enim judícia tua justa sunt, et omnes viæ tuæ misericórdia et véritas: et nunc, Dómine, meménto mei. — Ne.

\lectioheading{Lectio ii}{}

\lettrine[lines=2]{D}{ENIQUE}, cum irent omnes ad vítulos áureos, quos Jeróboam fécerat rex Israël, hic solus fugiébat consórtia ómnium; sed pergébat in Jerúsalem ad templum Dómini et ibi adorábat Dóminum Deum Israël, ómnia primitíva sua et décimas suas fidéliter ófferens; ita ut in tértio anno prosélytis et ádvenis ministráret omnem decimatiónem. Hæc et his simília secúndum legem Dei puérulus observábat. Cum vero factus esset vir, accépit uxórem Annam de tribu sua genuítque ex ea fílium, nomen suum impónens ei. Quem ab infántia timére Deum dócuit et abstinére ab omni peccáto.

\R\ Omni témpore bénedic Deum, et pete ab eo ut vias tuas dírigat, * Et in omni témpore consília tua in ipso permáneant. \V\ Inquíre ut fácias quæ plácita sunt illi in veritáte, et in tota virtúte tua. — Et.

\lectioheading{Lectio iii}{}

\lettrine[lines=2]{I}{GITUR} cum per captivitátem devenísset cum uxóre sua et fílio in civitátem Níniven cum omni tribu sua, cum omnes éderent ex cibis gentílium, iste custodívit ánimam suam et numquam contamínátus est in escis eórum. Et quóniam memor fuit Dómini in toto corde suo, dedit illi Deus grátiam in conspéctu Salmánasar regis, et dedit illi potestátem quocúmque vellet ire, habens libertátem quæcúmque fácere voluísset. Pergébat ergo ad omnes, qui erant in captivitáte, et mónita salútis dabat eis.

\R\ Memor esto, fili, quóniam páuperem vitam gérimus: * Habébis multa bona, si timúeris Deum. \V\ In mente habéto eum, et cave nequándo prætermíttas præcépta ejus. — Habébis. Glória Patri. Habébis.

\nocturnoheading{In II Nocturno}

\rubric{Sermo sancti Leónis Papæ}

\rubric{Sermo 9 de jejúnio 7 mensis}

\lectioheading{Lectio iv}{}

\lettrine[lines=2]{S}{CIO} quidem, dilectíssimi, plúrimos vestrum ita in iis, quæ ad observántiam christiánæ fídei pértinent, esse devótos, ut nostris cohortatiónibus non indígeant admonéri. Quod enim dudum et tradítio decrévit, et consuetúdo firmávit; nec erudítio ignórat, nec pietas prætermíttit. Sed quia sacerdotális offícii est, erga omnes Ecclésiæ fílios curam habére commúnem, in id quod et rúdibus prosit et doctis, quos simul dilígimus, páriter incitámus; ut jejúnium, quod nobis séptimi mensis recúrsus indícit, fide álacri per castigatiónem ánimi et córporis celebrémus.

\R\ Sufficiébat nobis paupértas nostra, ut divítiæ computaréntur: numquam fuísset pecúnia ipsa, pro qua misísti fílium nostrum, * Báculum senectútis nostræ! \V\ Heu me, fili mi, ut quid te mísimus peregrinári, lumen oculórum nostrórum? — Báculum.

\lectioheading{Lectio v}{}

\lettrine[lines=2]{I}{DEO} enim ipsa continéntiæ observántia quátuor est assignáta tempóribus, ut in idípsum totíus anni redeúnte decúrsu, cognoscerémus nos indesinénter purificatiónibus indigére; semperque esse niténdum, dum hujus vitæ varietáte jactámur, ut peccátum, quod fragilitáte carnis et cupiditátum pollutióne contráhitur, jejúniis atque eleemósynis deleátur. Esuriámus páululum, dilectíssimi, et aliquántulum, quod juvándis possit prodésse paupéribus, nostræ consuetúdini subtrahámus.

\R\ Benedícite Deum cæli et coram ómnibus vivéntibus confitémini ei, * Quia fecit vobíscum misericórdiam suam. \V\ Ipsum benedícite et cantáte illi: et enarráte ómnia mirabília ejus. — Quia.

\lectioheading{Lectio vi}{}

\lettrine[lines=2]{E}{LECTÉTUR} consciéntia benignórum frúctibus largitátis: et gáudia tríbuens, quo es lætificándus, accípies. Diléctio próximi, diléctio Dei est, qui plenitúdinem legis et prophetárum in hac gémina caritátis unitáte constítuit; ut nemo ambígeret, Deo se offérre, quod hómini contulísset, dicénte Dómino Salvatóre, cum de aléndis juvandísque paupéribus loquerétur: Quod uni eórum fecísti, mihi fecísti. Quarta ígitur et sexta féria jejunémus; sábbato vero apud beátum Petrum Apóstolum vigílias celebrémus: cujus nos méritis et oratiónibus crédimus adjuvándos, ut misericórdi Deo jejúnio nostro et devotióne placeámus.

\R\ Tempus est ut revértar ad eum qui misit me; * Vos autem benedícite Deum et enarráte ómnia mirabília ejus. \V\ Confitémini ei coram ómnibus vivéntibus, quia fecit vobíscum misericórdiam suam. — Vos. Glória Patri. Vos.

\R\ vii Tribulatiónes civitátum audívimus, quas passæ sunt, et defécimus: timor et hebetúdo mentis cécidit super nos et super líberos nostros: ipsi montes nolunt recípere fugam nostram: * Dómine, miserére. \V\ Peccávimus cum pátribus nostris, injúste égimus, iniquitátem fécimus. — Dómine.

\R\ viii Duo Séraphim, \rubric{ut infra in Homilíis.}

\feriaheading{Feria Secunda infra Hebdomadam III Septembris}

\rubric{Si Lectiónes de Scriptúra hujus Fériæ II aut Fériæ III sequéntis, sua die dici non potúerint, et in Fériis Quátuor Témporum fíeri contíngat Offícium, in quo Lectiónes de Scriptúra occurrénti sint dicéndæ, neque in eis reponéndum sit Inítium libri Tobíæ Domínicæ præcedénti assignátum, Lectiónes ut supra impedítæ, et si opus sit, étiam Lectiónes insequéntium Feriárum, servató tamen eárum órdine, ultérius transferúntur, ita ut in Offício in Quátuor Témporibus occurrénti Lectiónes de Scriptúra valeant adhibéri. Quod ítem servátur de Lectiónibus Fériæ V ut supra impedítis, quæ simíliter in casu vel ultérius transferúntur, vel étiam anticipántur. Responsória tamen semper dicúntur de Féria occurrénti.}

\rubric{De libro Tobíæ}

\lectioheading{Lectio i}{Cap. 2, 1-4}

\lettrine[lines=2]{P}{OST} hæc vero, cum esset dies festus Dómini, et factum esset prándium bonum in domo Tobíæ, dixit fílio suo: Vade et adduc áliquos de tribu nostra timéntes Deum, ut epuléntur nobíscum. Cumque abíisset, revérsus nuntiávit ei unum ex fíliis Israël jugulátum jacére in platéa. Statímque exsíliens de accúbitu suo relínquens prándium, jejúnus pervénit ad corpus, tollénsque illud portávit ad domum suam occúlte, ut, dum sol occubuísset, caute sepelíret eum.

\R\ Peto, Dómine, ut de vínculo improperii hujus absólvas me, aut certe désuper terram erípias me: * Ne reminiscáris delícta mea vel paréntum meórum, neque vindíctam sumas de peccátis meis: quia éruis sustinéntes te, Dómine. \V\ Omnia enim judícia tua justa sunt, et omnes viæ tuæ misericórdia et véritas: et nunc, Dómine, meménto mei. — Ne.

\lectioheading{Lectio ii}{Cap. 2, 8-18}

\lettrine[lines=2]{A}{RGUÉBANT} autem eum omnes próximi ejus dicéntes: Jam hujus rei causa intérfici jussus es et vix effugísti mortis impérium; et íterum sépelis mórtuos? Sed Tobías plus timens Deum quam regem, rapiébat córpora occísórum et occultábat in domo sua et médiis nóctibus sepeliébat ea. Cóntigit autem ut quadam die fatigátus a sepultúra, véniens in domum suam jactásset se juxta paríetem et obdormísset, et ex nido hirúndinum dormiénti illi cálida stércora incíderent super óculos ejus fierétque cæcus. Hanc autem tentatiónem ídeo permísit Dóminus eveníre illi, ut postéris darétur exémplum patiéntiæ ejus, sicut et sancti Job.

\R\ Omni témpore bénedic Deum, et pete ab eo ut vias tuas dírigat, * Et in omni témpore consília tua in ipso permáneant. \V\ Inquíre ut fácias quæ plácita sunt illi in veritáte, et in tota virtúte tua. — Et.

\rubric{In Octávis et Festis ritus simplícis: Glória Patri. Et.}

\lectioheading{Lectio iii}{}

\lettrine[lines=2]{N}{AM}, cum ab infántia sua semper Deum timúerit et mandáta ejus custodíerit, non est contristátus contra Deum quod plaga cæcitátis evénerit ei; sed immóbilis in Dei timóre permánsit agens grátias Deo ómnibus diébus vitæ suæ. Nam sicut beáto Job insultábant reges, ita isti paréntes et cognáti ejus irridébant vitam ejus dicéntes: Ubi est spes tua, pro qua eleemósynas et sepultúras faciébas? Tobías vero increpábat eos dicens: Nolíte ita loqui: quóniam fílii sanctórum sumus et vitam illam exspectámus, quam Deus datúrus est his qui fidem suam numquam mutant ab eo.

\R\ Memor esto, fili, quóniam páuperem vitam gérimus: * Habébis multa bona, si timúeris Deum. \V\ In mente habéto eum, et cave nequándo prætermíttas præcépta ejus. — Habébis. Glória Patri. Habébis.

\feriaheading{Feria Tertia infra Hebdomadam III Septembris}

\rubric{De libro Tobíæ}

\lectioheading{Lectio i}{Cap. 2, 19-23}

\lettrine[lines=2]{A}{NNA} vero uxor ejus ibat ad opus textrínum quotídie, et de labóre mánuum suárum victum, quem cónsequi póterat, deferébat. Unde factum est ut hædum caprárum accípiens detulísset domi; cujus cum vocem baléntis vir ejus audísset, dixit: Vidéte, ne forte furtívus sit; réddite eum dóminis suis, quia non licet nobis aut édere ex furto áliquid aut contíngere.

\R\ Sufficiébat nobis paupértas nostra, ut divítiæ computaréntur: numquam fuísset pecúnia ipsa, pro qua misísti fílium nostrum, * Báculum senectútis nostræ! \V\ Heu me, fili mi, ut quid te mísimus peregrinári, lumen oculórum nostrórum? — Báculum.

\lectioheading{Lectio ii}{}

\lettrine[lines=2]{A}{D} hæc uxor ejus iráta respóndit: Manifésté vana facta est spes tua et eleemósynæ tuæ modo apparuérunt. Atque his et áliis hujuscémodi verbis exprobrábat ei. Tunc Tobías ingémuit et cœpit oráre cum lácrimis dicens: Justus es, Dómine, et ómnia judícia tua justa sunt, et omnes viæ tuæ misericórdia et véritas et judícium. Et nunc, Dómine, memor esto mei et ne vindíctam sumas de peccátis meis neque reminiscáris delícta mea vel paréntum meórum.

\R\ Benedícite Deum cæli et coram ómnibus vivéntibus confitémini ei, * Quia fecit vobíscum misericórdiam suam. \V\ Ipsum benedícite et cantáte illi: et enarráte ómnia mirabília ejus. — Quia.

\rubric{In Octávis et Festis ritus simplícis: Glória Patri. Quia.}

\lectioheading{Lectio iii}{}

\lettrine[lines=2]{Q}{UÓNIAM} non obedívimus præcéptis tuis, ídeo tráditi sumus in direptiónem et captivitátem et mortem et in fábulam et in impropérium ómnibus natiónibus, in quibus dispersísti nos. Et nunc, Dómine, magna judícia tua, quia non égimus secúndum præcépta tua et non ambulávimus sincériter coram te. Et nunc, Dómine, secúndum voluntátem tuam fac mecum et præcipe in pace récipi spíritum meum; éxpedit enim mihi mori magis quam vívere.

\R\ Tempus est ut revértar ad eum qui misit me; * Vos autem benedícite Deum et enarráte ómnia mirabília ejus. \V\ Confitémini ei coram ómnibus vivéntibus, quia fecit vobíscum misericórdiam suam. — Vos. Glória Patri. Vos.

\rubric{¶ Fériæ IV et VI ac Sábbato Quátuor Témporum sunt Fériæ majóres non privilegiátæ.}

% ===== BATCH 4 of the Proprium de Tempore (autumnalis.md lines
% 14130-14245) -- "Dominica III Septembris" Matins, opening the Book
% of Tobias cycle in Nocturn I and St. Leo's sermon on the September
% Ember fast in Nocturn II (again with the deferred Nocturn III and
% numbered Responsories vii-viii pattern), followed by "Feria
% Secunda" and "Feria Tertia infra Hebdomadam III Septembris"
% continuing Tobias. Feria Secunda's lesson heading carries an
% unusually long governing rubric on how Ember Day lesson transfers
% interact with the Tobias readings when a conflict arises; this is
% transcribed in full as a single \rubric{} block. The closing
% rubric "Fériæ IV et VI ac Sábbato Quátuor Témporum sunt Fériæ
% majóres non privilegiátæ" signals the upcoming September Ember
% Days. The document continues with "Feria Quarta, Quatuor Temporum
% Septembris" (autumnalis.md line 14246).
% =====
\feriaheading{Feria Quarta, Quatuor Temporum Septembris}

\rubric{Hódie, Féria VI et Sábbato sequéntibus, Antíphonæ et Psalmi ad Laudes, et hac Féria étiam tres últimæ Antíphonæ cum Psalmis ad Matutínum, sumúntur de 2 loco; ad Primam autem addítur quartus Psalmus, ut in Psaltério notátur. Et dicúntur Preces feriáles ad omnes Horas, ut in Ordinário, præterquam in Vésperis et ad Complétórium Sábbati.}

\rubric{Léctio sancti Evangélii secúndum Marcum}

\lectioheading{Lectio i}{Cap. 9, 16-28}

\lettrine[lines=2]{I}{N} illo témpore: Respóndens unus de turba dixit ad Jesum: Magíster, áttuli fílium meum ad te habéntem spíritum mutum. Et réliqua.

\rubric{Homilía sancti Bedæ Venerábilis Presbýteri}

\rubric{Liber 3, cap. 38 in Marcum 9}

\lettrine[lines=2]{D}{ÆMONÍACUM} hunc, quem descéndens de monte Dóminus sanávit, Marcus quidem surdum mutúmque, Matthǽus vero lunáticum fuísse commémorat. Significat autem eos, de quibus scriptum est: Stultus ut luna mutátur. Qui numquam in eódem statu permanéntes, nunc ad hæc, nunc ad illa vítia mutáti, crescunt atque decréscunt. Qui muti sunt, non confiténdo fidem; surdi, nec ipsum aliquátenus veritátis audiéndo sermónem. Spumant autem, cum stultítia tabéscunt; stultórum namque et languéntium atque hébetum est, spumas salivárum ex ore dimíttere. Strident déntibus, cum iracúndiæ furóre flamméscunt: aréscunt, cum ótio torpénte languéscunt, et nulla virtútis indústria confortáti enérviter vivunt.

\R\ Tribulatiónes civitátum audívimus, quas passæ sunt, et defécimus: timor et hebetúdo mentis cécidit super nos et super líberos nostros: ipsi montes nolunt recípere fugam nostram: * Dómine, miserére. \V\ Peccávimus cum pátribus nostris, injúste égimus, iniquitátem fécimus. — Dómine.

\lectioheading{Lectio ii}{}

\lettrine[lines=2]{Q}{UOD} autem ait: Et dixi discípulis tuis ut ejícerent illum, et non potuérunt; laténter Apóstolos accúsat, cum impossibílitas curándi, intérdum non ad imbecillitátem curántium, sed ad eórum, qui curándi sunt, fidem referátur, dicénte Dómino: Fiat tibi secúndum fidem tuam. Qui respóndens eis, dixit: O generátio incrédula, quámdiu apud vos ero? quámdiu vos pátiar? Non quod tædio superátus sit mansuétus ac mitis, qui non apéruit sicut agnus coram tondénte os suum, nec in verba furóris erúpit; sed quo in similitúdinem médici, si ægrótum vídeat contra sua præcépta se gérere, dicat: Úsquequo accédam ad domum tuam? quoúsque artis meæ perdam indústriam, me áliud jubénte, et te áliud perpetránte?

\R\ Omni témpore bénedic Deum, et pete ab eo ut vias tuas dírigat, * Et in omni témpore consília tua in ipso permáneant. \V\ Inquíre ut fácias quæ plácita sunt illi in veritáte, et in tota virtúte tua. — Et.

\lectioheading{Lectio iii}{}

\lettrine[lines=2]{E}{T} dixit illis: Hoc genus in nullo potest exíre, nisi in oratióne et jejúnio. Dum docet Apóstolos, quómodo dæmon nequíssimus débeat expélli, omnes instítuit ad vitam; ut scílicet novérimus, gravíssima quæque vel immundórum spirítuum, vel hóminum tentaménta jejúniis et oratiónibus esse superánda: iram quoque Dómini, cum in ultiónem nostrórum scélerum fúerit accénsa, hoc remédio singulári posse placári. Jejúnium autem generále est, non solum ab escis, sed et a cunctis illécebris abstinére carnálibus, immo ab ómnibus vitiórum continére se passiónibus. Sic et orátio generális non in verbis solum est, quibus divínam cleméntiam invocámus, verum étiam in ómnibus, quæ in obséquium nostri Condítoris fídei devotióne gérimus.

\R\ Memor esto, fili, quóniam páuperem vitam gérimus: * Habébis multa bona, si timúeris Deum. \V\ In mente habéto eum, et cave nequándo prætermíttas præcépta ejus. — Habébis. Glória Patri. Habébis.

\hourheading{Ad Laudes}

\V\ Repléti sumus mane misericórdia tua. \R\ Exsultávimus, et delectáti sumus.

\antlabel{Ad Bened. Ant.} Hoc genus * dæmoniórum in nullo potest exíre, nisi in oratióne et jejúnio.

\oratioheading{Oratio}

\lettrine[lines=2]{M}{ISERICÓRDIÆ} tuæ remédiis, quǽsumus, Dómine, fragílitas nostra subsístat: ut, quæ sua conditióne attéritur, tua cleméntia reparétur. Per Dóminum.

\rubric{Et dícitur usque ad Nonam. Ad Vésperas vero Orátio Domínicæ præcedéntis. Si autem Vésperæ dicéndæ non sint de Féria, de ea nec fit Commemorátio.}

\feriaheading{Feria Quinta infra Hebdomadam III Septembris}

\rubric{De libro Tobíæ}

\lectioheading{Lectio i}{Cap. 12, 1-17}

\lettrine[lines=2]{V}{OCÁVIT} ad se Tobías fílium suum dixítque ei: Quid póssumus dare viro isti sancto, qui venit tecum? Respóndens Tobías dixit patri suo: Pater, quam mercédem dábimus ei? aut quid dignum póterit esse benefíciis ejus? Me duxit et redúxit sanum, pecúniam a Gabélo ipse recépit, uxórem ipse me habére fecit et dæmónium ab ea ipse compéscuit, gáudium paréntibus ejus fecit, meípsum a devoratióne piscis erípuit, te quoque vidére fecit lumen cæli, et bonis ómnibus per eum repléti sumus. Quid illi ad hæc potérimus dignum dare? Sed peto te, pater mi, ut roges eum, si forte dignábitur medietátem de ómnibus quæ alláta sunt sibi assúmere.

\R\ Peto, Dómine, ut de vínculo improperii hujus absólvas me, aut certe désuper terram erípias me: * Ne reminiscáris delícta mea vel paréntum meórum, neque vindíctam sumas de peccátis meis: quia éruis sustinéntes te, Dómine. \V\ Omnia enim judícia tua justa sunt, et omnes viæ tuæ misericórdia et véritas: et nunc, Dómine, meménto mei. — Ne.

\lectioheading{Lectio ii}{}

\lettrine[lines=2]{E}{T} vocántes eum, pater scílicet et fílius, tulérunt eum in partem; et rogáre cœpérunt ut dignarétur dimídiam partem ómnium quæ attúlerant, accéptam habére. Tunc dixit eis occúlte: Benedícite Deum cæli et coram ómnibus vivéntibus confitémini ei, quia fecit vobíscum misericórdiam suam. Etenim sacraméntum regis abscóndere bonum est, ópera autem Dei reveláre et confitéri honoríficum est. Bona est orátio cum jejúnio, et eleemósyna magis quam thesáuros auri recóndere; quóniam eleemósyna a morte líberat, et ipsa est quæ purgat peccáta, et facit inveníre misericórdiam et vitam ætérnam. Qui autem fáciunt peccátum et iniquitátem hostes sunt ánimæ suæ.

\R\ Omni témpore bénedic Deum, et pete ab eo ut vias tuas dírigat, * Et in omni témpore consília tua in ipso permáneant. \V\ Inquíre ut fácias quæ plácita sunt illi in veritáte, et in tota virtúte tua. — Et.

\rubric{In Octávis et Festis ritus simplícis: Glória Patri. Et.}

\lectioheading{Lectio iii}{}

\lettrine[lines=2]{M}{ANIFÉSTO} ergo vobis veritátem et non abscóndam a vobis occúltum sermónem. Quando orábas cum lácrimis et sepeliébas mórtuos et derelinquébas prándium tuum et mórtuos abscondébas per diem in domo tua et nocte sepeliébas eos, ego óbtuli oratiónem tuam Dómino. Et quia accéptus eras Deo, necésse fuit ut tentátio probáret te. Et nunc misit me Dóminus ut curárem te et Saram uxórem fílii tui a dæmónio liberárem; ego enim sum Ráphaël Angelus, unus ex septem qui astámus ante Dóminum. Cumque hæc audíssent, turbáti sunt et treméntes cecidérunt super terram in fáciem suam. Dixítque eis Angelus: Pax vobis, nolíte timére.

\R\ Memor esto, fili, quóniam páuperem vitam gérimus: * Habébis multa bona, si timúeris Deum. \V\ In mente habéto eum, et cave nequándo prætermíttas præcépta ejus. — Habébis. Glória Patri. Habébis.

\feriaheading{Feria Sexta, Quatuor Temporum Septembris}

\rubric{Léctio sancti Evangélii secúndum Lucam}

\lectioheading{Lectio i}{Cap. 7, 36-50}

\lettrine[lines=2]{I}{N} illo témpore: Rogábat Jesum quidam de pharisǽis, ut manducáret cum illo. Et ingréssus domum pharisǽi discúbuit. Et réliqua.

\rubric{Homilía sancti Gregórii Papæ}

\rubric{Homilía 33 in Evang., post inítium}

\lettrine[lines=2]{Q}{UEM} pharisǽus de falsa justítia præsúmens, nisi Judáicum pópulum; quem peccátrix múlier ad vestígia Dómini véniens et plorans, nisi convérsam gentilitátem desígnat? Quæ cum alabástro venit, unguéntum fudit, retro secus pedes Dómini stetit, lácrimis pedes rigávit, capíllis tersit, eosdémque quos infundébat et tergébat pedes osculári non désiit. Nos ergo, nos illa múlier expréssit, si toto corde ad Dóminum post peccáta redeámus, si ejus pœniténtiæ luctus imitémur. Quid namque unguénto, nisi bonæ odor opiniónis exprímitur? Unde et Paulus dicit: Christi bonus odor sumus Deo in omni loco.

\R\ Sufficiébat nobis paupértas nostra, ut divítiæ computaréntur: numquam fuísset pecúnia ipsa, pro qua misísti fílium nostrum, * Báculum senectútis nostræ! \V\ Heu me, fili mi, ut quid te mísimus peregrinári, lumen oculórum nostrórum? — Báculum.

\lectioheading{Lectio ii}{}

\lettrine[lines=2]{S}{I} ergo recta ópera ágimus, quibus opiniónis bone odóre Ecclésiam respergámus, quid in Dómini córpore, nisi unguéntum fúndimus? Sed secus pedes Jesu múlier stetit: contra pedes enim Dómini stétimus, cum in peccátis pósiti ejus itinéribus renitebámur. Sed si ad veram pœniténtiam post peccáta convértimur, jam retro secus pedes stamus; quia ejus vestígia séquimur, quem impugnabámus. Lácrimis múlier pedes ejus rigat: quod nos quoque veráciter ágimus, si quibúslibet últimis membris Dómini per compassiónis afféctum inclinámur, si sanctis ejus in tribulatióne compátimur, si eórum tristítiam, nostram putámus.

\R\ Benedícite Deum cæli et coram ómnibus vivéntibus confitémini ei, * Quia fecit vobíscum misericórdiam suam. \V\ Ipsum benedícite et cantáte illi: et enarráte ómnia mirabília ejus. — Quia.

\lectioheading{Lectio iii}{}

\lettrine[lines=2]{C}{APILLIS} ergo pedes Dómini tergimus, quando sanctis ejus, quibus ex caritáte compátimur, étiam ex his quæ nobis supérfluunt, miserémur; quátenus sic mens per compassiónem dóleat, ut étiam larga manus afféctum dolóris osténdat. Rigat namque lácrimis Redemptóris pedes, sed capíllis suis non tergit, qui utcúmque proximórum dolóri compátitur, sed tamen eis ex his quæ sibi supérfluunt, non miserétur. Plorat, et non tergit, qui verba quidem dolóris tríbuit, sed non ministrándo quæ desunt, vim dolóris mínime abscíndit. Osculátur múlier pedes, quos tergit: quod nos quoque plene ágimus, si studióse dilígimus, quos ex largitáte continémus; ne gravis nobis sit necéssitas próximi; ne ipsa nobis ejus indigéntia, quæ sustentátur, fiat onerósa; et, cum manus necessária tríbuit, ánimus a dilectióne torpéscat.

\R\ Tempus est ut revértar ad eum qui misit me; * Vos autem benedícite Deum et enarráte ómnia mirabília ejus. \V\ Confitémini ei coram ómnibus vivéntibus, quia fecit vobíscum misericórdiam suam. — Vos. Glória Patri. Vos.

\hourheading{Ad Laudes}

\V\ Repléti sumus mane misericórdia tua. \R\ Exsultávimus, et delectáti sumus.

\antlabel{Ad Bened. Ant.} Múlier * quæ erat in civitáte peccátrix, stans retro secus pedes Dómini, lácrimis cœpit rigáre pedes ejus et capíllis cápitis sui tergébat, et deosculabátur pedes ejus et unguénto ungébat.

\oratioheading{Oratio}

\lettrine[lines=2]{P}{RÆSTA}, quǽsumus, omnípotens Deus: ut observatiónes sacras ánnua devotióne recoléntes, et córpore tibi placeámus, et mente. Per Dóminum.

\rubric{Et dícitur usque ad Nonam. Ad Vésperas vero Orátio Domínicæ præcedéntis. Si autem Vésperæ dicéndæ non sint de Féria, de ea nec fit Commemorátio.}

\feriaheading{Sabbato, Quatuor Temporum Septembris}

\rubric{Léctio sancti Evangélii secúndum Lucam}

\lectioheading{Lectio i}{Cap. 13, 6-17}

\lettrine[lines=2]{I}{N} illo témpore: Dicébat Jesus turbis hanc similitúdinem: Arbórem fici habébat quidam plantátam in vínea sua et venit quærens fructum in illa et non invénit. Et réliqua.

\rubric{Homilía sancti Gregórii Papæ}

\rubric{Homilía 31 in Evangelia}

\lettrine[lines=2]{D}{OMINUS} ac Redémptor noster per Evangélium suum aliquándo verbis, aliquándo rebus lóquitur: aliquándo áliud verbis, atque áliud rebus; aliquándo autem hoc verbis, quod rebus. Duas étenim res ex Evangélio, fratres, audístis: ficúlneam infructuósam, et mulíerem curvam; et utríque rei est pietas impénsa. Illud autem dixit per similitúdinem: istud egit per exhibitiónem. Sed hoc significat ficúlnea infructuósa, quod múlier inclináta; et hoc ficúlnea reserváta, quod múlier erécta.

\R\ Tribulatiónes civitátum audívimus, quas passæ sunt, et defécimus: timor et hebetúdo mentis cécidit super nos et super líberos nostros: ipsi montes nolunt recípere fugam nostram: * Dómine, miserére. \V\ Peccávimus cum pátribus nostris, injúste égimus, iniquitátem fécimus. — Dómine.

\lectioheading{Lectio ii}{}

\lettrine[lines=2]{Q}{UID} arbor fici, nisi humánam natúram desígnat? Quid múlier inclináta, nisi eámdem natúram significándo denúntiat? Quæ et bene plantáta est sicut ficus, et bene creáta sicut múlier; sed in culpam própria sponte lapsa, neque fructum servat operatiónis, neque statum rectitúdinis. Ad peccátum quippe ex voluntáte córruens, quia fructum obediéntiæ ferre nóluit, statum rectitúdinis amísit. Quæ ad Dei similitúdinem cóndita, dum in sua dignitáte non pérstitit, quod plantáta vel creáta fúerat, serváre contémpsit. Tértio dóminus víneæ ad ficúlneam venit, quia natúram géneris humáni ante legem, sub lege, sub grátia, exspectándo, admonéndo, visitándo requisívit.

\R\ Omni témpore bénedic Deum, et pete ab eo ut vias tuas dírigat, * Et in omni témpore consília tua in ipso permáneant. \V\ Inquíre ut fácias quæ plácita sunt illi in veritáte, et in tota virtúte tua. — Et.

\lectioheading{Lectio iii}{}

\lettrine[lines=2]{V}{ENIT} ante legem, quia per naturálem intelléctum quid unusquísque exémplo sui, quáliter erga próximum ágere debuísset, innótuit. Venit in lege, quia præcipiéndo dócuit. Venit post legem per grátiam, quia pietátis suæ præséntiam exhibéndo monstrávit. Sed tamen in tribus annis fructum se non invenísse conquéritur: quia quorúmdam pravórum mentes nec inspiráta lex naturális córrigit, nec præcépta erúdiunt, nec incarnatiónis ejus miracula convértunt. Quid vero per cultórem víneæ, nisi præpositórum ordo exprímitur? Qui dum præsunt Ecclésiæ, nimírum Domínicæ víneæ curam gerunt.

\R\ Memor esto, fili, quóniam páuperem vitam gérimus: * Habébis multa bona, si timúeris Deum. \V\ In mente habéto eum, et cave nequándo prætermíttas præcépta ejus. — Habébis. Glória Patri. Habébis.

\hourheading{Ad Laudes}

\V\ Repléti sumus mane misericórdia tua. \R\ Exsultávimus, et delectáti sumus.

\antlabel{Ad Bened. Ant.} Illúmina, Dómine, * sedéntes in ténebris et umbra mortis, et dírige pedes nostros in viam pacis.

\oratioheading{Oratio}

\lettrine[lines=2]{O}{MNÍPOTENS} sempitérne Deus, qui per continéntiam salutárem córporibus medéris et méntibus: majestátem tuam súpplices exorámus; ut pia jejunántium deprecatióne placátus, et præséntia nobis subsídia tríbuas et futúra. Per Dóminum.

\rubric{Et dícitur usque ad Nonam inclusíve.}

\hourheading{Ad Vesperas}

\V\ Vespertína orátio ascéndat ad te, Dómine. \R\ Et descéndat super nos misericórdia tua.

\antlabel{Ad Magnif. Ant.} Adonái, Dómine, * Deus magne et mirábilis, qui dedísti salútem in manu féminæ, exáudi preces servórum tuórum.

\rubric{Oratio Dominicæ occurrentis.}

% ===== BATCH 5 of the Proprium de Tempore (autumnalis.md lines
% 14246-14411) -- the complete September Ember Days (Feria Quarta,
% Feria Sexta, and Sabbato, Quatuor Temporum Septembris), each
% following the Sunday-Gospel-homily pattern: a brief Gospel
% pericope incipit ending "Et réliqua" under \lectioheading{Lectio
% i}, immediately followed within the same Lesson by a Patristic
% homily on that Gospel (Bede on Mark 9, Gregory on Luke 7:36-50,
% Gregory again on Luke 13:6-17), each homily's author and source
% given as separate \rubric{} labels before its own fresh lettrine,
% matching the established precedent from aestiva.tex's parallel
% "Lectio vii" Trinity Sunday construction. Feria Quinta infra
% Hebdomadam III Septembris (an ordinary day, not an Ember day,
% falling between the Wednesday and Friday Ember days) continues
% the Tobias cycle with the Archangel Raphael's self-revelation.
% Each Ember day's all-caps "AD LAUDES" (and, for the Saturday, "AD
% VESPERAS") sub-heading is normalized to the standard mixed-case
% \hourheading{} form per the established incidental-formatting-
% variation rule, each followed by its own brief antiphon, Oratio,
% and governing rubric on how long the Oratio is said. The document
% continues with "Dominica IV Septembris" (autumnalis.md line
% 14412).
% =====
\feriaheading{Dominica IV Septembris}

\rubric{Semiduplex}

\nocturnoheading{In I Nocturno}

\rubric{Incipit liber Judith}

\lectioheading{Lectio i}{Cap. 1, 1-10; 2, 1-3}

\lettrine[lines=2]{A}{RPHÁXAD} ítaque rex Medórum subjugáverat multas gentes império suo, et ipse ædificávit civitátem potentíssimam, quam appellávit Ecbátanis, ex lapídibus quadrátis et sectis. Fecit muros ejus in altitúdinem cubitórum septuagínta et in latitúdinem cubitórum trigínta; turres vero ejus pósuit in altitúdinem cubitórum centum. Per quadrum vero eárum latus utrúmque vicenórum pedum spátio tendebátur; posuítque portas ejus in altitúdinem túrrium; et gloriabátur quasi potens in poténtia exércitus sui et in glória quadrigárum suárum.

\R\ Adonái, Dómine, Deus magne et mirábilis, qui dedísti salútem in manu féminæ, * Exáudi preces servórum tuórum. \V\ Benedíctus es, Dómine, qui non derelínquis præsuméntes de te, et de sua virtúte gloriántes humílias. — Exáudi.

\lectioheading{Lectio ii}{}

\lettrine[lines=2]{A}{NNO} ígitur duodécimo regni sui Nabuchodónosor rex Assyriórum, qui regnábat in Nínive civitáte magna, pugnávit contra Arpháxad et obtínuit eum in campo magno qui appellátur Rágau, circa Euphráten et Tigrin et Jádason, in campo Erioch regis Elicórum. Tunc exaltátum est regnum Nabuchodónosor, et cor ejus elevátum est; et misit ad omnes qui habitábant in Cilícia et Damásco et Líbano, et ad gentes, quæ sunt in Carmélo et Cedar et inhabitántes Galilǽam in campo magno Esdrelon, et ad omnes qui erant in Samaría et trans flumen Jordánem usque ad Jerúsalem, et omnem terram Jesse, quoúsque perveniátur ad términos Æthiópiæ.

\R\ Tribulatiónes civitátum audívimus, quas passæ sunt, et defécimus: timor et hebetúdo mentis cécidit super nos et super líberos nostros: ipsi montes nolunt recípere fugam nostram: * Dómine, miserére. \V\ Peccávimus cum pátribus nostris, injúste égimus, iniquitátem fécimus. — Dómine.

\lectioheading{Lectio iii}{}

\lettrine[lines=2]{A}{D} hos omnes misit núntios Nabuchodónosor rex Assyriórum; qui omnes uno ánimo contradixérunt et remisérunt eos vácuos et sine honóre abjecérunt. Tunc indignátus Nabuchodónosor rex advérsus omnem terram illam, jurávit per thronum et regnum suum, quod defénderet se de ómnibus regiónibus his. Anno tértio décimo Nabuchodónosor regis, vigésima et secúnda die mensis primi factum est verbum in domo Nabuchodónosor regis Assyriórum, ut defénderet se. Vocavítque omnes majóres natu, omnésque duces et bellatóres suos, et hábuit cum eis mystérium consílii sui dixítque cogitatiónem suam in eo esse, ut omnem terram suo subjugáret império.

\R\ Benedícat te Dóminus in virtúte sua, qui per te ad níhilum redégit inimícos nostros: * Ut non defíciat laus tua de ore hóminum. \V\ Benedíctus Dóminus qui creávit cælum et terram; quia hódie nomen tuum ita magnificávit. — Ut. Glória Patri. Ut.

\nocturnoheading{In II Nocturno}

\rubric{Ex libro sancti Ambrósii Epíscopi de Elía et jejúnio}

\lectioheading{Lectio iv}{Cap. 9}

\lettrine[lines=2]{P}{OTÉNTES} vinum prohibéntur bíbere, ne, cum bíberint, obliviscántur sapiéntiam. Dénique bibébant vinum in ebrietáte poténtes, qui Holoférni príncipi milítiæ regis Assyriórum se trádere gestiébant; sed non bibébat fémina Judith, jejúnans ómnibus diébus viduitátis suæ, præter festórum diérum solemnitátes. His armis muníta procéssit et omnem Assyriórum circumvénit exércitum. Sóbrii vigóre consílii ábstulit Holoférnis caput, servávit pudicítiam, victóriam reportávit.

\R\ Nos álium Deum nescímus præter Dóminum, in quo sperámus: * Qui non déspicit nos, nec ámovet salútem suam a génere nostro. \V\ Indulgéntiam ipsíus fusis lácrimis postulémus, et humiliémus illi ánimas nostras. — Qui.

\lectioheading{Lectio v}{}

\lettrine[lines=2]{H}{ÆC} enim succíncta jejúnio in castris prætendébat aliénis; ille vino sepúltus jacébat, ut ictum vúlneris sentíre non posset. Itaque uníus mulíeris jejúnium innúmeros strávit exércitus Assyriórum. Esther quoque púlchrior facta est jejúnio; Dóminus enim grátiam sóbriæ mentis augébat. Omne genus suum, id est, totum pópulum Judæórum a persecutiónis acerbitáte liberávit, ita ut regem sibi fáceret esse subjéctum.

\R\ Dominátor, Dómine cælórum et terræ, Creátor aquárum, Rex univérsæ creatúræ: * Exáudi oratiónem servórum tuórum. \V\ Tu, Dómine, cui humílium semper et mansuetórum plácuit deprecátio. — Exáudi.

\lectioheading{Lectio vi}{}

\lettrine[lines=2]{I}{TAQUE} illa, quæ tríduo jejunávit contínuo, et corpus suum aqua lavit, plus plácuit, et vindíctam rétulit. Aman autem, dum se regáli jactat convívio, inter ipsa vina pœnam suæ ebrietátis exsólvit. Est ergo jejúnium reconciliatiónis sacrifícium, virtútis increméntum, quod fecit étiam féminas fortióres augménto grátiæ. Jejúnium nescit fænerátorem, non sortem fæneris novit: non rédolet usúras mensa jejunántium. Etiam ipsis jejúnium convívíis dat grátiam: dulcióres post famem épulæ fiunt, quæ assiduitáte fastídio sunt, et diutúrna continuatióne viléscunt. Condiméntum cibi jejúnium est: quanto avídior appeténtia, tanto esca jucúndior.

\R\ Dómine Deus, qui cónteris bella ab inítio, éleva bráchium tuum super Gentes, quæ cógitant servis tuis mala: * Et déxtera tua glorificétur in nobis. \V\ Allíde virtútem eórum in virtúte tua; cadat robur eórum in iracúndia tua. — Et. Glória Patri. Et.

\R\ vii Confórta me, Rex Sanctórum principátum tenens: * Et da sermónem rectum et bene sonántem in os meum. \V\ Dómine, Rex univérsæ potestátis, convérte consílium eórum super eos. — Et.

\R\ viii Duo Séraphim, \rubric{ut infra in Homilíis.}

\feriaheading{Feria Secunda infra Hebdomadam IV Septembris}

\rubric{De libro Judith}

\lectioheading{Lectio i}{Cap. 4, 1-12}

\lettrine[lines=2]{F}{ILII} Israël, qui habitábant in terra Juda, timuérunt valde a fácie Holoférnis; tremor et horror invásit sensus eórum, ne hoc fáceret Jerúsalem et templo Dómini, quod fécerat céteris civitátibus et templis eárum. Et misérunt in omnem Samaríam per circúitum usque Jéricho et præoccupavérunt omnes vértices móntium et muris circumdedérunt vicos suos et congregavérunt fruménta in præparatiónem pugnæ.

\R\ Adonái, Dómine, Deus magne et mirábilis, qui dedísti salútem in manu féminæ, * Exáudi preces servórum tuórum. \V\ Benedíctus es, Dómine, qui non derelínquis præsuméntes de te, et de sua virtúte gloriántes humílias. — Exáudi.

\lectioheading{Lectio ii}{}

\lettrine[lines=2]{S}{ACÉRDOS} étiam Elíachim scripsit ad univérsos qui erant contra Esdrelon, quæ est contra fáciem campi magni juxta Dóthain, et univérsos per quos viæ tránsitus esse póterat, ut obtinérent ascénsus móntium, per quos via esse póterat ad Jerúsalem, et illic custodírent ubi angústum iter esse póterat inter montes. Et fecérunt fílii Israël secúndum quod constitúerat eis sacérdos Dómini Elíachim. Et clamávit omnis pópulus ad Dóminum instántia magna, et humiliavérunt ánimas suas in jejúniis et oratiónibus ipsi et mulíeres eórum.

\R\ Tribulatiónes civitátum audívimus, quas passæ sunt, et defécimus: timor et hebetúdo mentis cécidit super nos et super líberos nostros: ipsi montes nolunt recípere fugam nostram: * Dómine, miserére. \V\ Peccávimus cum pátribus nostris, injúste égimus, iniquitátem fécimus. — Dómine.

\rubric{In Octávis et Festis ritus simplícis: Glória Patri. Dómine.}

\lectioheading{Lectio iii}{}

\lettrine[lines=2]{E}{T} induérunt se sacerdótes cilíciis, et infántes prostravérunt contra fáciem templi Dómini, et altáre Dómini operuérunt cilício; et clamavérunt ad Dóminum Deum Israël unanímiter, ne daréntur in prædam infántes eórum, et uxóres eórum in divisiónem, et civitátes eórum in extermínium, et sancta eórum in pollutiónem, et fíerent oppróbrium Géntibus. Tunc Elíachim sacérdos Dómini magnus circuívit omnem Israël allocutúsque est eos dicens: Scitóte quóniam exáudiet Dóminus preces vestras, si manéntes permanséritis in jejúniis et oratiónibus in conspéctu Dómini.

\R\ Benedícat te Dóminus in virtúte sua, qui per te ad níhilum redégit inimícos nostros: * Ut non defíciat laus tua de ore hóminum. \V\ Benedíctus Dóminus qui creávit cælum et terram; quia hódie nomen tuum ita magnificávit. — Ut. Glória Patri. Ut.

\feriaheading{Feria Tertia infra Hebdomadam IV Septembris}

\rubric{De libro Judith}

\lectioheading{Lectio i}{Cap. 8, 1-11}

\lettrine[lines=2]{E}{T} factum est, cum audísset hæc verba Judith vídua, quæ erat fília Merári, fílii Idox, fílii Joseph, fílii Oziæ, fílii Elai, fílii Jammor, fílii Gédeon, fílii Ráphaïm, fílii Achitob, fílii Melchíæ, fílii Enan, fílii Nathaníæ, fílii Saláthiel, fílii Símeon, fílii Ruben, et vir ejus fuit Manásses, qui mórtuus est in diébus messis hordeáceæ; instábat enim super alligántes manípulos in campo, et venit æstus super caput ejus, et mórtuus est in Bethúlia civitáte sua et sepúltus est illic cum pátribus suis. Erat autem Judith relícta ejus vídua jam annis tribus et ménsibus sex.

\R\ Nos álium Deum nescímus præter Dóminum, in quo sperámus: * Qui non déspicit nos, nec ámovet salútem suam a génere nostro. \V\ Indulgéntiam ipsíus fusis lácrimis postulémus, et humiliémus illi ánimas nostras. — Qui.

\lectioheading{Lectio ii}{}

\lettrine[lines=2]{E}{T} in superióribus domus suæ fecit sibi secrétum cubículum, in quo cum puéllis suis clausa morabátur. Et habens super lumbos suos cilícium, jejunábat ómnibus diébus vitæ suæ præter sábbata et neoménias et festa domus Israël. Erat autem elegánti aspéctu nimis; cui vir suus reliquérat divítias multas et famíliam copiósam ac possessiónes arméntis boum et grégibus óvium plenas. Et erat hæc in ómnibus famosíssima, quóniam timébat Dóminum valde, nec erat qui loquerétur de illa verbum malum.

\R\ Dominátor, Dómine cælórum et terræ, Creátor aquárum, Rex univérsæ creatúræ: * Exáudi oratiónem servórum tuórum. \V\ Tu, Dómine, cui humílium semper et mansuetórum plácuit deprecátio. — Exáudi.

\rubric{In Octávis et Festis ritus simplícis: Glória Patri. Exáudi.}

\lectioheading{Lectio iii}{}

\lettrine[lines=2]{H}{ÆC} ítaque, cum audísset quóniam Ozías promisísset quod, transácto quinto die, tráderet civitátem, misit ad presbýteros Chabri et Charmi. Et venérunt ad illam, et dixit illis: Quod est hoc verbum, in quo consénsit Ozías, ut tradat civitátem Assýriis, si intra quinque dies non vénerit vobis adjutórium? Et qui estis vos, qui tentátis Dóminum?

\R\ Dómine Deus, qui cónteris bella ab inítio, éleva bráchium tuum super gentes, quæ cógitant servis tuis mala: * Et déxtera tua glorificétur in nobis. \V\ Allíde virtútem eórum in virtúte tua; cadat robur eórum in iracúndia tua. — Et. Glória Patri. Et.

\feriaheading{Feria Quarta infra Hebdomadam IV Septembris}

\rubric{De libro Judith}

\lectioheading{Lectio i}{Cap. 10, 1-4}

\lettrine[lines=2]{F}{ACTUM} est autem, cum cessásset clamáre ad Dóminum surréxit de loco in quo jacúerat prostráta ad Dóminum; vocavítque abram suam, et descéndens in domum suam ábstulit a se cilícium et éxuit se vestiméntis viduitátis suæ et lavit corpus suum et unxit se myro óptimo et discriminávit crinem cápitis sui et impósuit mitram super caput suum et índuit se vestiméntis jucunditátis suæ induítque sandália pédibus suis assumpsítque dextralíola et lília et ináures et ánulos et ómnibus ornaméntis suis ornávit se. Cui étiam Dóminus cóntulit splendórem.

\R\ Confórta me, Rex Sanctórum principátum tenens: * Et da sermónem rectum et bene sonántem in os meum. \V\ Dómine, Rex univérsæ potestátis, convérte consílium eórum super eos. — Et.

\lectioheading{Lectio ii}{Cap. 10, 11-12}

\lettrine[lines=2]{F}{ACTUM} est autem, cum descénderet montem circa ortum diéi, occurrérunt ei exploratóres Assyriórum et tenuérunt eam dicéntes: Unde venis aut quo vadis? Quæ respóndit: Fília sum Hebræórum; ídeo ego fugi a fácie eórum, quóniam futúrum agnóvi quod dentur vobis in deprædatiónem, pro eo quod contemnéntes vos noluérunt ultro trádere seípsos, ut invenírent misericórdiam in conspéctu vestro.

\R\ Tribulatiónes civitátum audívimus, quas passæ sunt, et defécimus: timor et hebetúdo mentis cécidit super nos et super líberos nostros: ipsi montes nolunt recípere fugam nostram: * Dómine, miserére. \V\ Peccávimus cum pátribus nostris, injúste égimus, iniquitátem fécimus. — Dómine.

\rubric{In Octávis et Festis ritus simplícis: Glória Patri. Dómine.}

\lectioheading{Lectio iii}{Cap. 10, 16-20}

\lettrine[lines=2]{D}{UXERÚNTQUE} illam ad tabernáculum Holoférnis annuntiántes eam. Cumque intrásset ante fáciem ejus, statim captus est in suis óculis Holoférnes. Dixerúntque ad eum satéllites ejus: Quis contémnat pópulum Hebræórum, qui tam decóras mulíeres habent, ut non pro his mérito pugnáre contra eos debeámus? Videns ítaque Judith Holoférnem sedéntem in conopéo, quod erat ex púrpura et auro et smarágdo et lapídibus pretiósis intéxtum, et, cum in fáciem ejus intendísset, adorávit eum prostérnens se super terram.

\R\ Benedícat te Dóminus in virtúte sua, qui per te ad níhilum redégit inimícos nostros: * Ut non defíciat laus tua de ore hóminum. \V\ Benedíctus Dóminus qui creávit cælum et terram; quia hódie nomen tuum ita magnificávit. — Ut. Glória Patri. Ut.

\feriaheading{Feria Quinta infra Hebdomadam IV Septembris}

\rubric{Si hæc Féria occúrrat post diem 24 Septémbris, ita ut hæc sit última hebdómada mensis, hac nocte pónitur liber Esther cum suis Responsóriis, ut infra in Dominica V, 321: et légitur de eo étiam duóbus sequéntibus diébus, ita tamen ut, lecto Inítio Domínicæ assignáto, sequéntibus Fériis dicántur Lectiónes de respondénti Féria VI aut Sábbato ejúsdem hebdómadæ V.}

\rubric{De libro Judith}

\lectioheading{Lectio i}{Cap. 12, 10-13}

\lettrine[lines=2]{E}{T} factum est, in quarto die Holoférnes fecit cœnam servis suis et dixit ad Vágao eunúchum suum: Vade et suáde Hebrǽam illam ut sponte conséntiat habitáre mecum. Foedum est enim apud Assýrios, si fémina irrídeat virum agéndo ut immúnis ab eo tránseat. Tunc introívit Vágao ad Judith et dixit: Non vereátur bona puélla introíre ad dóminum meum, ut honorificétur ante fáciem ejus, ut mandúcet cum eo et bibat vinum in jucunditáte. Cui Judith respóndit: Quæ ego sum, ut contradícam dómino meo?

\R\ Adonái, Dómine, Deus magne et mirábilis, qui dedísti salútem in manu féminæ, * Exáudi preces servórum tuórum. \V\ Benedíctus es, Dómine, qui non derelínquis præsuméntes de te, et de sua virtúte gloriántes humílias. — Exáudi.

\lectioheading{Lectio ii}{Cap. 13, 1-11}

\lettrine[lines=2]{U}{T} autem sero factum est, festinavérunt servi illíus ad hospítia sua; et conclúsit Vágao óstia cubículi et ábiit. Erant autem omnes fatigáti a vino; erátque Judith sola in cubículo. Porro Holoférnes jacébat in lecto nímia ebrietáte sopítus. Dixítque Judith puéllæ suæ ut staret foris ante cubículum et observáret. Stetítque Judith ante lectum orans cum lácrimis et labiórum motu in siléntio dicens: Confírma me, Dómine, Deus Israël, et réspice in hac hora ad ópera mánuum meárum, ut, sicut promisísti, Jerúsalem civitátem tuam érigas, et hoc quod credens per te posse fíeri cogitávi, perfíciam.

\R\ Tribulatiónes civitátum audívimus, quas passæ sunt, et defécimus: timor et hebetúdo mentis cécidit super nos et super líberos nostros: ipsi montes nolunt recípere fugam nostram: * Dómine, miserére. \V\ Peccávimus cum pátribus nostris, injúste égimus, iniquitátem fécimus. — Dómine.

\rubric{In Octávis et Festis ritus simplícis: Glória Patri. Dómine.}

\lectioheading{Lectio iii}{}

\lettrine[lines=2]{E}{T}, cum hæc dixísset, accéssit ad colúmnam, quæ erat ad caput léctuli ejus, et pugiónem ejus, qui in ea ligátus pendébat, exsólvit. Cumque evaginásset illum, apprehéndit comam cápitis ejus et ait: Confírma me, Dómine Deus, in hac hora. Et percússit bis in cervícem ejus et abscídit caput ejus et ábstulit conopéum ejus a colúmnis et evólvit corpus ejus truncum. Et post pusíllum exívit et trádidit caput Holoférnis ancíllæ suæ et jussit ut mítteret illud in peram suam.

\R\ Benedícat te Dóminus in virtúte sua, qui per te ad níhilum redégit inimícos nostros: * Ut non defíciat laus tua de ore hóminum. \V\ Benedíctus Dóminus qui creávit cælum et terram; quia hódie nomen tuum ita magnificávit. — Ut. Glória Patri. Ut.

\feriaheading{Feria Sexta infra Hebdomadam IV Septembris}

\rubric{De libro Judith}

\lectioheading{Lectio i}{Cap. 15, 1-3}

\lettrine[lines=2]{C}{UMQUE} omnis exércitus decollátum Holoférnem audísset, fugit mens et consílium ab eis, et solo tremóre et metu agitáti fugæ præsídium sumunt, ita ut nullus loquerétur cum próximo suo, sed inclináto cápite, relíctis ómnibus, evádere festinábant Hebrǽos, quos armátos super se veníre audiébant, fugiéntes per vias campórum et sémitas cóllium. Vidéntes ítaque fílii Israël fugiéntes, secúti sunt illos, descenderúntque clangéntes tubis et ululántes post ipsos.

\R\ Nos álium Deum nescímus præter Dóminum, in quo sperámus: * Qui non déspicit nos, nec ámovet salútem suam a génere nostro. \V\ Indulgéntiam ipsíus fusis lácrimis postulémus, et humiliémus illi ánimas nostras. — Qui.

\lectioheading{Lectio ii}{Cap. 15, 5-7}

\lettrine[lines=2]{M}{ISIT} ítaque Ozías núntios per omnes civitátes et regiónes Israël. Omnis ítaque régio omnísque urbs eléctam juventútem armátam misit post eos; et persecúti sunt eos in ore gládii, quoúsque pervenírent ad extremitátem fínium suórum. Réliqui autem qui erant in Bethúlia, ingréssi sunt castra Assyriórum et prædam, quam fugiéntes Assýrii relíquerant, abstulérunt et onustáti sunt valde.

\R\ Dominátor, Dómine cælórum et terræ, Creátor aquárum, Rex univérsæ creatúræ: * Exáudi oratiónem servórum tuórum. \V\ Tu, Dómine, cui humílium semper et mansuetórum plácuit deprecátio. — Exáudi.

\rubric{In Octávis et Festis ritus simplícis: Glória Patri. Exáudi.}

\lectioheading{Lectio iii}{Cap. 15, 9-12}

\lettrine[lines=2]{J}{OACHIM} autem summus póntifex de Jerúsalem venit in Bethúliam cum univérsis presbýteris suis, ut vidéret Judith. Quæ cum exísset ad illum, benedixérunt eam omnes una voce dicéntes: Tu glória Jerúsalem, tu lætítia Israël, tu honorificéntia pópuli nostri. Quia fecísti viríliter, et confortátum est cor tuum, eo quod castitátem amáveris et, post virum tuum, álterum nescíeris, ídeo et manus Dómini confortávit te, et ídeo eris benedícta in ætérnum. Et dixit omnis pópulus: Fiat, fiat.

\R\ Dómine Deus, qui cónteris bella ab inítio, éleva bráchium tuum super gentes, quæ cógitant servis tuis mala: * Et déxtera tua glorificétur in nobis. \V\ Allíde virtútem eórum in virtúte tua; cadat robur eórum in iracúndia tua. — Et. Glória Patri. Et.

\feriaheading{Sabbato infra Hebdomadam IV Septembris}

\rubric{De libro Judith}

\lectioheading{Lectio i}{Cap. 16, 22-31}

\lettrine[lines=2]{E}{T} factum est post hæc, omnis pópulus post victóriam venit in Jerúsalem adoráre Dóminum; et mox ut purificáti sunt, obtulérunt omnes holocáusta et vota et repromissiónes suas. Porro Judith univérsa vasa béllica Holoférnis, quæ dedit illi pópulus, et conopéum, quod ipsa sustúlerat de cúbili ipsíus, óbtulit in anáthema oblivíónis.

\R\ Confórta me, Rex, Sanctórum principátum tenens: * Et da sermónem rectum et bene sonántem in os meum. \V\ Dómine, Rex univérsæ potestátis, convérte consílium eórum super eos. — Et.

\lectioheading{Lectio ii}{}

\lettrine[lines=2]{E}{RAT} autem pópulus jucúndus secúndum fáciem sanctórum; et per tres menses gáudium hujus victóriæ celebrátum est cum Judith. Post dies autem illos unusquísque rédiit in domum suam, et Judith magna facta est in Bethúlia et præclárior erat univérsæ terræ Israël. Erat étiam virtúti cástitas adjúncta, ita ut non cognósceret virum ómnibus diébus vitæ suæ, ex quo defúnctus est Manásses vir ejus. Erat autem, diébus festis, procédens cum magna glória.

\R\ Tribulatiónes civitátum audívimus, quas passæ sunt, et defécimus: timor et hebetúdo mentis cécidit super nos et super líberos nostros: ipsi montes nolunt recípere fugam nostram: * Dómine, miserére. \V\ Peccávimus cum pátribus nostris, injúste égimus, iniquitátem fécimus. — Dómine.

\rubric{In Octávis simplícibus et in Offício sanctæ Maríæ in Sábbato: Glória Patri. Dómine.}

\lectioheading{Lectio iii}{}

\lettrine[lines=2]{M}{ANSIT} autem in domo viri sui annos centum quinque, et dimísit abram suam líberam, et defúncta est ac sepúlta cum viro suo in Bethúlia. Luxítque illam omnis pópulus diébus septem. In omni autem spátio vitæ ejus non fuit qui perturbáret Israël, et post mortem ejus annis multis. Dies autem victóriæ hujus festivitátis ab Hebrǽis in número sanctórum diérum accípitur et cólitur a Judǽis ex illo témpore usque in præséntem diem.

\R\ Benedícat te Dóminus in virtúte sua, qui per te ad níhilum redégit inimícos nostros: * Ut non defíciat laus tua de ore hóminum. \V\ Benedíctus Dóminus qui creávit cælum et terram; quia hódie nomen tuum ita magnificávit. — Ut. Glória Patri. Ut.

\feriaheading{Sabbato ante Dominicam V Septembris}

\rubric{nisi hæc occúrrat a die 28 ad 30 Septémbris: quia tunc, sequénti Antíphona omíssa, pónitur Ant. Adapériat, quæ assignátur Sábbato ante Dominicam I Octóbris.}

\hourheading{Ad Vesperas}

\V\ Vespertína orátio ascéndat ad te, Dómine. \R\ Et descéndat super nos misericórdia tua.

\antlabel{Ad Magnif. Ant.} Dómine, Rex omnípotens, * in dicióne tua cuncta sunt pósita, et non est qui possit resístere voluntáti.

\rubric{Orátio Dominicæ occurréntis.}

% ===== BATCH 6 of the Proprium de Tempore (autumnalis.md lines
% 14412-14650) -- the complete fourth week of September: "Dominica
% IV Septembris" Matins opening the Book of Judith in Nocturn I and
% St. Ambrose's "De Elia et jejunio" in Nocturn II (same deferred
% Nocturn III pattern), the six ferial days completing the Judith
% cycle through Holofernes's death and Judith's long widowhood, and
% the brief "Sabbato ante Dominicam V Septembris" (Vespers only,
% with a rubric noting its Magnificat antiphon is replaced by
% "Adapériat" -- assigned to the Saturday before the First Sunday of
% October -- whenever this particular Saturday falls between
% September 28 and 30). Five genuine partial-caps OCR artifacts are
% silently normalized to full caps per the established rule
% (CUmque, MIsit, JOachim, ERat, MAnsit), each confirmed by direct
% inspection of autumnalis.md to be a true partial-caps defect (only
% the first one or two letters capitalized, the remainder lowercase)
% rather than a deliberate small-caps or proper-name styling; none
% carries an accent on its newly-capitalized remainder, so no accent
% is dropped in the normalization. A separate, more pervasive
% formatting quirk also appears starting partway through this
% batch (confirmed via direct grep of autumnalis.md to begin exactly
% at "Feria Tertia infra Hebdomadam IV Septembris" and continue
% through "Sabbato infra Hebdomadam IV Septembris," with zero
% occurrences before that point): many colons in the "In Octávis..."
% rubrics and the Responsories' verse markers are preceded by a
% space ("ritus simplícis :" rather than "ritus simplícis:"). Since
% this reflects only an incidental OCR/typesetting spacing variant,
% not a textual difference, and the surrounding document (including
% this very same batch's earlier material) consistently uses no
% space before a colon, this spacing is silently normalized to match
% the rest of the document, consistent with the established
% treatment of other incidental formatting variations (all-caps
% headings, lowercase Roman numerals) elsewhere in this project. The
% document continues with "Dominica V Septembris" (autumnalis.md
% line 14651), opening the Book of Esther.
% =====
\feriaheading{Dominica V Septembris}

\rubric{Semiduplex}

\nocturnoheading{In I Nocturno}

\rubric{Incipit liber Esther}

\lectioheading{Lectio i}{Cap. 1, 1-9}

\lettrine[lines=2]{I}{N} diébus Assuéri, qui regnávit ab India usque Æthiópiam super centum vigínti septem províncias, quando sedit in sólio regni sui, Susan cívitas regni ejus exórdium fuit. Tértio ígitur anno impérii sui, fecit grande convívium cunctis princípibus et púeris suis fortíssimis Persárum et Medórum ínclytis et præféctis provinciárum coram se, ut osténderet divítias glóriæ regni sui, ac magnitúdinem atque jactántiam poténtiæ suæ, multo témpore, centum vidélicet et octogínta diébus.

\R\ Dómine, mi Rex omnípotens, in dicióne tua cuncta sunt pósita, et non est qui possit resístere voluntáti tuæ: * Líbera nos propter nomen tuum. \V\ Exáudi oratiónem nostram, et convérte luctum nostrum in gáudium. — Líbera.

\lectioheading{Lectio ii}{}

\lettrine[lines=2]{C}{UMQUE} implerántur dies convívii, invitávit omnem pópulum, qui invéntus est in Susan, a máximo usque ad mínimum; et jussit septem diébus convívium præparári in vestíbulo horti et némoris, quod régio cultu et manu cónsitum erat. Et pendébant ex omni parte tentória aérii colóris et carbásini ac hyacínthini sustentáta fúnibus býssinis atque purpúreis, qui ebúrneis círculis insérti erant et colúmnis marmóreis fulciebántur. Léctuli quoque áurei et argéntei super paviméntum smarágdino et pário stratum lápide dispósiti erant, quod mira varietáte pictúra decorábat.

\R\ Confórta me, Rex, Sanctórum principátum tenens: * Et da sermónem rectum et bene sonántem in os meum. \V\ Dómine, Rex univérsæ potestátis, convérte consílium eórum super eos. — Et.

\lectioheading{Lectio iii}{}

\lettrine[lines=2]{B}{IBEBANT} autem, qui invitáti erant, áureis póculis et áliis atque áliis vasis cibi inferebántur. Vinum quoque, ut magnificéntia régia dignum erat, abúndans et præcípuum ponebátur. Nec erat qui noléntes cógeret ad bibéndum, sed, sicut rex statúerat præpónens mensis síngulos de princípibus suis, ut súmeret unusquísque quod vellet. Vasthi quoque regína fecit convívium feminárum in palátio, ubi rex Assuérus manére consuéverat.

\R\ Spem in álium numquam hábui, prætérquam in te, Deus Israël: * Qui irásceris, et propítius eris: et ómnia peccáta hóminum in tribulatióne dimíttis. \V\ Dómine Deus, Creátor cæli et terræ, réspice ad humilitátem nostram. — Qui. Glória Patri. Qui.

\nocturnoheading{In II Nocturno}

\rubric{Ex libro Officiórum sancti Ambrósii Epíscopi}

\lectioheading{Lectio iv}{Lib. 3, cap. 15}

\lettrine[lines=2]{Q}{UID} Esther regína, nonne ut pópulum suum perículo erúeret (quod erat decórum atque honéstum) morti se óbtulit, nec immítis regis trepidávit furórem? Ipse quoque rex Persárum, ferox atque túmido corde, tamen decórum judicávit índici insidiárum, quæ sibi parátæ forent, grátiam repræsentáre, populúmque líberum a servitúte erípere, erúere neci, nec párcere neci ejus, qui tam indecóra suasísset. Dénique quem secúndum a se, ac præcípuum inter omnes amícos habéret, cruci trádidit, quod dehonestátum se ejus frauduléntis consíliis animadvertísset.

\R\ Meménto mei, Dómine Deus, in bonum: * Et ne déleas miseratiónes meas quas feci in domo Dei mei et in cæremóniis ejus. \V\ Recordáre mei, Dómine, Deus meus. — Et.

\lectioheading{Lectio v}{Cap. 16}

\lettrine[lines=2]{E}{A} enim amicítia probábilis, quæ honestátem tuétur, præferénda sane ópibus, honóribus, potestátibus; honestáti vero præférri non solet, sed honestátem sequi. Qualis fuit Jónathæ, qui pro pietáte nec offénsam patris, nec salútis perículum refugiébat. Qualis fuit Achímelech, qui pro hospitális grátiæ offíciis necem pótius sibi, quam proditiónem fugiéntis amíci, subeúndam arbitrabátur. Nihil ígitur præferéndum honestáti; quæ tamen ne amicítiæ stúdio prætereátur, étiam hoc Scriptúra ádmonet.

\R\ Tribulatiónes civitátum audívimus, quas passæ sunt, et defécimus: timor et hebetúdo mentis cécidit super nos et super líberos nostros: ipsi montes nolunt recípere fugam nostram: * Dómine, miserére. \V\ Peccávimus cum pátribus nostris, injúste égimus, iniquitátem fécimus. — Dómine.

\lectioheading{Lectio vi}{}

\lettrine[lines=2]{S}{UNT} enim pleræque philosophórum quæstiónes: Utrum amíci causa quisquam contra pátriam sentíre necne débeat, ut amíco obédiat: utrum opórteat ut fidem déserat, dum indúlget atque inténdit amíci commoditátibus. Et Scriptúra quidem ait: Clava, et gládius, et sagítta ferráta, sic homo est testimónium dans falsum advérsus amícum suum. Sed consídera quid ástruat. Non testimónium reprehéndit dictum in amícum, sed falsum testimónium. Quid enim si Dei causa, quid si pátriæ, cogátur áliquis dícere testimónium? Numquid præponderáre debet amicítia religióni, præponderáre caritáti vítium?

\R\ Benedícat te Dóminus in virtúte sua, qui per te ad níhilum redégit inimícos nostros: * Ut non defíciat laus tua de ore hóminum. \V\ Benedíctus Dóminus qui creávit cælum et terram; quia hódie nomen tuum ita magnificávit. — Ut. Glória Patri. Ut.

\R\ Nos álium Deum nescímus præter Dóminum, in quo sperámus: * Qui non déspicit nos, nec ámovet salútem suam a génere nostro. \V\ Indulgéntiam ipsíus fusis lácrimis postulémus, et humiliémus illi ánimas nostras. — Qui.

\R\ Duo Séraphim, \rubric{ut infra in Homilíis.}

\feriaheading{Feria Secunda infra Hebdomadam V Septembris}

\rubric{De libro Esther}

\lectioheading{Lectio i}{Cap. 2, 5-11}

\lettrine[lines=2]{E}{RAT} vir Judǽus in Susan civitáte, vocábulo Mardochǽus, fílius Jair, fílii Sémei, fílii Cis, de stirpe Jémini, qui translátus fúerat de Jerúsalem eo témpore quo Jechoníam regem Juda Nabuchodónosor rex Babylónis transtúlerat. Qui fuit nutrícius fíliæ fratris sui Edíssæ, quæ áltero nómine vocabátur Esther, et utrúmque paréntem amíserat: pulchra nimis et decóra fácie. Mortuísque patre ejus ac matre, Mardochǽus sibi eam adoptávit in fíliam.

\R\ Dómine, mi Rex omnípotens, in dicióne tua cuncta sunt pósita, et non est qui possit resístere voluntáti tuæ: * Líbera nos propter nomen tuum. \V\ Exáudi oratiónem nostram, et convérte luctum nostrum in gáudium. — Líbera.

\lectioheading{Lectio ii}{}

\lettrine[lines=2]{C}{UMQUE} percrebuísset regis impérium, et juxta mandátum illíus multæ pulchræ vírgines adduceréntur Susan et Egéo traderéntur eunúcho; Esther quoque inter céteras puéllas ei trádita est, ut servarétur in número feminárum. Quæ plácuit ei et invénit grátiam in conspéctu illíus. Et præcépit eunúcho ut acceleráret mundum muliébrem et tráderet ei partes suas et septem puéllas speciosíssimas de domo regis, et tam ipsam quam pedíssequas ejus ornáret atque excóleret. Quæ nóluit indicáre ei pópulum et pátriam suam; Mardochǽus enim præcéperat ei, ut de hac re omníno reticéret. Qui deambulábat quotídie ante vestíbulum domus, in qua eléctæ vírgines servabántur, curam agens salútis Esther.

\R\ Confórta me, Rex, Sanctórum principátum tenens: * Et da sermónem rectum et bene sonántem in os meum. \V\ Dómine, Rex univérsæ potestátis, convérte consílium eórum super eos. — Et.

\rubric{In Octávis et Festis ritus simplícis: Glória Patri. Et.}

\lectioheading{Lectio iii}{Cap. 2, 15-17}

\lettrine[lines=2]{E}{VOLUTO} autem témpore per órdinem, instábat dies quo Esther fília Abíhail fratris Mardochǽi, quam sibi adoptáverat in fíliam, debéret intráre ad regem. Quæ non quæsívit muliébrem cultum, sed quæcúmque vóluit Egéus eunúchus custos vírginum, hæc ei ad ornátum dedit; erat enim formósa valde et incredíbili pulchritúdine: ómnium óculis gratiósa et amábilis videbátur. Ducta est ítaque ad cubículum regis Assuéri mense décimo, qui vocátur Tebeth, séptimo anno regni ejus. Et adamávit eam rex plus quam omnes mulíeres, habuítque grátiam et misericórdiam coram eo super omnes mulíeres, et pósuit diadéma regni in cápite ejus fecítque eam regnáre in loco Vasthi.

\R\ Spem in álium numquam hábui, prætérquam in te, Deus Israël: * Qui irásceris, et propítius eris: et ómnia peccáta hóminum in tribulatióne dimíttis. \V\ Dómine Deus, Creátor cæli et terræ, réspice ad humilitátem nostram. — Qui. Glória Patri. Qui.

\feriaheading{Feria Tertia infra Hebdomadam V Septembris}

\rubric{De libro Esther}

\lectioheading{Lectio i}{Cap. 3, 1-7}

\lettrine[lines=2]{P}{OST} hæc rex Assuérus exaltávit Aman fílium Amadáthi, qui erat de stirpe Agag, et pósuit sólium ejus super omnes príncipes quos habébat. Cunctíque servi regis, qui in fóribus palátii versabántur, flectébant génua et adorábant Aman; sic enim præcéperat eis imperátor. Solus Mardochǽus non flectébat genu neque adorábat eum. Cui dixérunt púeri regis, qui ad fores palátii præsidébant: Cur præter céteros non servas mandátum regis?

\R\ Meménto mei, Dómine Deus, in bonum: * Et ne déleas miseratiónes meas quas feci in domo Dei mei et in cæremóniis ejus. \V\ Recordáre mei, Dómine, Deus meus. — Et.

\lectioheading{Lectio ii}{}

\lettrine[lines=2]{C}{UMQUE} hoc crébrius dícerent, et ille nollet audíre, nuntiavérunt Aman, scire cupiéntes utrum perseváret in senténtia; díxerat enim eis se esse Judǽum. Quod cum audísset Aman et experiménto probásset quod Mardochǽus non fléctere sibi genu nec se adoráret, irátus est valde et pro níhilo duxit in unum Mardochǽum míttere manus suas.

\R\ Tribulatiónes civitátum audívimus, quas passæ sunt, et defécimus: timor et hebetúdo mentis cécidit super nos et super líberos nostros: ipsi montes nolunt recípere fugam nostram: * Dómine, miserére. \V\ Peccávimus cum pátribus nostris, injúste égimus, iniquitátem fécimus. — Dómine.

\rubric{In Octávis et Festis ritus simplícis: Glória Patri. Dómine.}

\lectioheading{Lectio iii}{}

\lettrine[lines=2]{A}{UDIERAT} enim quod esset gentis Judǽæ, magísque vóluit omnem Judæórum rum, qui erant in regno Assuéri, pérdere natiónem. Mense primo, cujus vocábulum est Nisan, anno duodécimo regni Assuéri, missa est sors in urnam, quæ Hebráice dícitur phur, coram Aman, quo die et quo mense gens Judæórum debéret intérfici; et exívit mensis duodécimus qui vocátur Adar.

\R\ Benedícat te Dóminus in virtúte sua, qui per te ad níhilum redégit inimícos nostros: * Ut non defíciat laus tua de ore hóminum. \V\ Benedíctus Dóminus qui creávit cælum et terram; quia hódie nomen tuum ita magnificávit. — Ut. Glória Patri. Ut.

\feriaheading{Feria Quarta infra Hebdomadam V Septembris}

\rubric{De libro Esther}

\lectioheading{Lectio i}{Cap. 4, 1-17}

\lettrine[lines=2]{Q}{UÆ} cum audísset Mardochǽus, scidit vestiménta sua et indútus est sacco spargens cínerem cápiti; et in platéa médiæ civitátis voce magna clamábat osténdens amaritúdinem ánimi sui, et hoc ejulátu usque ad fores palátii grádiens; non enim erat lícitum indútum sacco aulam regis intráre. In ómnibus quoque províncíis, óppidis ac locis, ad quæ crudéle regis dogma pervénerat, planctus ingens erat apud Judǽos, jejúnium, ululátus et fletus, sacco et cínere multis pro strato uténtibus. Ingréssæ autem sunt puéllæ Esther et eunúchi nuntiaverúntque ei. Quod áudiens consternáta est et vestem misit, ut, abláto sacco, indúerent eum: quam accípere nóluit. Accitóque Athach eunúcho, quem rex minístrum ei déderat, præcépit ei ut iret ad Mardochǽum et discéret ab eo cur hoc fáceret.

\R\ Nos álium Deum nescímus præter Dóminum, in quo sperámus: * Qui non déspicit nos, nec ámovet salútem suam a génere nostro. \V\ Indulgéntiam ipsíus fusis lácrimis postulémus, et humiliémus illi ánimas nostras. — Qui.

\lectioheading{Lectio ii}{}

\lettrine[lines=2]{E}{GRESSUSQUE} Athach ivit ad Mardochǽum stantem in platéa civitátis ante óstium palátii. Qui indicávit ei ómnia quæ accíderant, quómodo Aman promisísset, ut in thesáuros regis pro Judæórum nece inférret argéntum; exémplar quoque edícti, quod pendébat in Susan, dedit ei ut regínæ osténderet et monéret eam ut intráret ad regem et deprecarétur eum pro pópulo suo. Regréssus Athach nuntiávit Esther ómnia quæ Mardochǽus díxerat. Quæ respóndit ei et jussit ut díceret Mardochǽo: Omnes servi regis et cunctæ quæ sub dicióne ejus sunt norunt províncíæ, quod sive vir sive múlier non vocátus intérius átrium regis intráverit, absque ulla cunctatióne statim interficiátur, nisi forte rex áuream virgam ad eum teténderit pro signo cleméntiæ, atque ita possit vívere. Ego ígitur quómodo ad regem intráre pótero, quæ trigínta jam diébus non sum vocáta ad eum?

\R\ Confórta me, Rex, Sanctórum principátum tenens: * Et da sermónem rectum et bene sonántem in os meum. \V\ Dómine, Rex univérsæ potestátis, convérte consílium eórum super eos. — Et.

\rubric{In Octávis et Festis ritus simplícis: Glória Patri. Et.}

\lectioheading{Lectio iii}{}

\lettrine[lines=2]{Q}{UOD} cum audísset Mardochǽus, rursum mandávit Esther dicens: Ne putes quod ánimam tuam tantum líberes, quia in domo regis es, præ cunctis Judǽis; si enim nunc silúeris, per áliam occasiónem liberabúntur Judǽi, et tu et domus patris tui períbitis. Et quis novit utrum idcírco ad regnum véneris, ut in tali témpore pararéris? Rursúmque Esther hæc Mardochǽo verba mandávit: Vade et cóngrega omnes Judǽos, quos in Susan repéreris, et oráte pro me. Non comedátis et non bibátis tribus diébus et tribus nóctibus, et ego cum ancíllis meis simíliter jejunábo; et tunc ingrédiar ad regem contra legem fáciens non vocáta tradénsque me morti et perículo. Ivit ítaque Mardochǽus, et fecit ómnia quæ ei Esther præcéperat.

\R\ Spem in álium numquam hábui, prætérquam in te, Deus Israël: * Qui irásceris, et propítius eris: et ómnia peccáta hóminum in tribulatióne dimíttis. \V\ Dómine Deus, Creátor cæli et terræ, réspice ad humilitátem nostram. — Qui. Glória Patri. Qui.

\feriaheading{Feria Quinta infra Hebdomadam V Septembris}

\rubric{De libro Esther}

\lectioheading{Lectio i}{Cap. 5, 1-5}

\lettrine[lines=2]{D}{IE} autem tértio indúta est Esther regálibus vestiméntis et stetit in átrio domus régiæ, quod erat intérius contra basílicam regis. At ille sedébat super sólium suum in consistório palátii contra óstium domus. Cumque vidísset Esther regínam stantem, plácuit óculis ejus et exténdit contra eam virgam áuream quam tenébat manu. Quæ accédens osculáta est summitátem virgæ ejus, dixítque ad eam rex: Quid vis, Esther regína? quæ est petítio tua? Etiam si dimídiam partem regni petíeris, dábitur tibi. At illa respóndit: Si regi placet, óbsecro ut vénias ad me hódie et Aman tecum ad convívium quod parávi. Statímque rex, Vocáte, inquit, cito Aman ut Esther obédiat voluntáti.

\R\ Dómine, mi Rex omnípotens, in dicióne tua cuncta sunt pósita, et non est qui possit resístere voluntáti tuæ: * Líbera nos propter nomen tuum. \V\ Exáudi oratiónem nostram, et convérte luctum nostrum in gáudium. — Líbera.

\lectioheading{Lectio ii}{Cap. 5, 9-14}

\lettrine[lines=2]{E}{GRESSUS} est ítaque illo die Aman lætus et álacer. Cumque vidísset Mardochǽum sedéntem ante fores palátii, et non solum non assurrexísse sibi, sed nec motum quidem de loco sessiónis suæ, indignátus est valde; et, dissimuláta ira, revérsus in domum suam convocávit ad se amícos suos et Zares uxórem suam et expósuit illis magnitúdinem divitiárum suárum filiorúmque turbam et quanta eum glória super omnes príncipes et servos suos rex elevásset. Et post hæc ait: Regína quoque Esther nullum álium vocávit ad convívium cum rege, præter me, apud quam étiam cras cum rege pransúrus sum; et, cum hæc ómnia hábeam, nihil me habére puto, quámdiu vídero Mardochǽum Judǽum sedéntem ante fores régias.

\R\ Confórta me, Rex, Sanctórum principátum tenens: * Et da sermónem rectum et bene sonántem in os meum. \V\ Dómine, Rex univérsæ potestátis, convérte consílium eórum super eos. — Et.

\rubric{In Octávis et Festis ritus simplícis: Glória Patri. Et.}

\lectioheading{Lectio iii}{}

\lettrine[lines=2]{R}{ESPONDERUNTQUE} ei Zares uxor ejus et céteri amíci: Jube parári excélsam trabem habéntem altitúdinis quinquagínta cúbitos, et dic mane regi ut appendátur super eam Mardochǽus; et sic ibis cum rege lætus ad convívium. Plácuit ei consílium et jussit excélsam paráti crucem.

\R\ Spem in álium numquam hábui, prætérquam in te, Deus Israël: * Qui irásceris, et propítius eris: et ómnia peccáta hóminum in tribulatióne dimíttis. \V\ Dómine Deus, Creátor cæli et terræ, réspice ad humilitátem nostram. — Qui. Glória Patri. Qui.

\feriaheading{Feria Sexta infra Hebdomadam V Septembris}

\rubric{De libro Esther}

\lectioheading{Lectio i}{Cap. 6, 1-13}

\lettrine[lines=2]{N}{OCTEM} illam duxit rex insómnem jussítque sibi afférri histórias et annáles priórum témporum. Quæ cum illo præsénte legeréntur, ventum est ad illum locum, ubi scriptum erat quómodo nuntiásset Mardochǽus insídias Bágathan et Thares eunuchórum regem Assuérum juguláre cupiéntium. Quod cum audísset rex ait: Quid pro hac fide honóris ac præmii Mardochǽus consecútus est? Dixérunt ei servi illíus ac minístri: Níhil omníno mercédis accépit. Statímque rex, Quis est, inquit, in átrio? Aman quippe intérius átrium domus régiæ intráverat, ut suggéreret regi et jubéret Mardochǽum affígi patíbulo, quod ei fúerat præparátum. Respondérunt púeri: Aman stat in átrio. Dixítque rex: Ingrediátur.

\R\ Meménto mei, Dómine Deus, in bonum: * Et ne déleas miseratiónes meas quas feci in domo Dei mei et in cæremóniis ejus. \V\ Recordáre mei, Dómine, Deus meus. — Et.

\lectioheading{Lectio ii}{}

\lettrine[lines=2]{C}{UMQUE} esset ingréssus ait illi: Quid debet fíeri viro, quem rex honoráre desíderat? Cógitans autem in corde suo Aman et réputans quod nullum álium rex nisi se vellet honoráre, respóndit: Homo quem rex honoráre cupit debet índui véstibus régiis et impóni super equum qui de sella regis est, et accípere régium diadéma super caput suum, et primus de régiis princípibus ac tyránnis téneat equum ejus et per platéam civitátis incédens clamet et dicat: Sic honorábitur quemcúmque volúerit rex honoráre.

\R\ Tribulatiónes civitátum audívimus, quas passæ sunt, et defécimus: timor et hebetúdo mentis cécidit super nos et super líberos nostros: ipsi montes nolunt recípere fugam nostram: * Dómine, miserére. \V\ Peccávimus cum pátribus nostris, injúste égimus, iniquitátem fécimus. — Dómine, miserére.

\rubric{In Octávis et Festis ritus simplícis: Glória Patri. Dómine.}

\lectioheading{Lectio iii}{}

\lettrine[lines=2]{D}{IXITQUE} ei rex: Festína et, sumpta stola et equo, fac, ut locútus es, Mardochǽo Judǽo, qui sedet ante fores palátii. Cave ne quidquam de his quæ locútus es prætermíttas. Tulit ítaque Aman stolam et equum, indutúmque Mardochǽum in platéa civitátis et impósitum equo præcedébat atque clamábat: Hoc honóre condígnus est quemcúmque rex volúerit honoráre. Reversúsque est Mardochǽus ad jánuam palátii, et Aman festinávit ire in domum suam lugens et opérto cápite. Narravítque Zares uxóri suæ et amícis ómnia quæ eveníssent sibi.

\R\ Benedícat te Dóminus in virtúte sua, qui per te ad níhilum redégit inimícos nostros: * Ut non defíciat laus tua de ore hóminum. \V\ Benedíctus Dóminus qui creávit cælum et terram; quia hódie nomen tuum ita magnificávit. — Ut. Glória Patri. Ut.

\feriaheading{Sabbato infra Hebdomadam V Septembris}

\rubric{De libro Esther}

\lectioheading{Lectio i}{Cap. 7, 1-10}

\lettrine[lines=2]{I}{NTRAVIT} ítaque rex et Aman, ut bíberent cum regína. Dixítque ei rex étiam secúnda die, postquam vino incalúerat: Quæ est petítio tua, Esther, ut detur tibi, et quid vis fíeri? Etiam si dimídiam partem regni petíeris, impetrábis. Ad quem illa respóndit: Si invéni grátiam in óculis tuis, o rex, et si tibi placet, dona mihi ánimam meam, pro qua rogo, et pópulum meum, pro quo óbsecro. Tráditi enim sumus ego et pópulus meus, ut conterámur, jugulémur et pereámus. Atque, útinam in servos et fámulas venderémur: esset tolerábile malum et gemens tacérem; nunc autem hostis noster est cujus crudélitas redúndat in regem.

\R\ Nos álium Deum nescímus præter Dóminum, in quo sperámus: * Qui non déspicit nos, nec ámovet salútem suam a génere nostro. \V\ Indulgéntiam ipsíus fusis lácrimis postulémus, et humiliémus illi ánimas nostras. — Qui.

\lectioheading{Lectio ii}{}

\lettrine[lines=2]{R}{ESPONDENSQUE} rex Assuérus ait: Quis est iste et cujus poténtiæ, ut hæc áudeat fácere? Dixítque Esther: Hostis et inimícus noster péssimus iste est Aman. Quod ille áudiens íllico obstúpuit vultum regis ac regínæ ferre non sústinens. Rex autem irátus surréxit et de loco convívii intrávit in hortum arbóribus cónsitum. Aman quoque surréxit ut rogáret Esther regínam pro ánima sua; intelléxit enim a rege sibi parátum malum.

\R\ Confórta me, Rex, Sanctórum principátum tenens: * Et da sermónem rectum et bene sonántem in os meum. \V\ Dómine, Rex univérsæ potestátis, convérte consílium eórum super eos. — Et.

\rubric{In Octávis simplícibus et in Offício Sanctæ Maríæ in Sábbato: Glória Patri. Et.}

\lectioheading{Lectio iii}{}

\lettrine[lines=2]{Q}{UI} cum revérsus esset de horto nemóribus cónsito et intrásset convívii locum, réperit Aman super léctulum corruísse in quo jacébat Esther et ait: Etiam regínam vult opprímere, me præsénte, in domo mea. Necdum verbum de ore regis exíerat, et statim operuérunt fáciem ejus. Dixítque Harbóna unus de eunúchis qui stabant in ministério regis: En lignum, quod paráverat Mardochǽo qui locútus est pro rege, stat in domo Aman habens altitúdinis quinquagínta cúbitos. Cui dixit rex: Appéndite eum in eo. Suspénsus est ítaque Aman in patíbulo, quod paráverat Mardochǽo, et regis ira quiévit.

\R\ Spem in álium numquam hábui, prætérquam in te, Deus Israël: * Qui irásceris, et propítius eris: et ómnia peccáta hóminum in tribulatióne dimíttis. \V\ Dómine Deus, Creátor cæli et terræ, réspice ad humilitátem nostram. — Qui. Glória Patri. Qui.

\feriaheading{Sabbato ante Dominicam I Octobris}

\hourheading{Ad Vesperas}

\V\ Vespertína orátio ascéndat ad te, Dómine. \R\ Et descéndat super nos misericórdia tua.

\antlabel{Ad Magnif. Ant.} Adapériat Dóminus * cor vestrum in lege sua et in præcéptis suis et fáciat pacem Dóminus Deus noster.

\rubric{Oratio Dominicæ occurréntis.}

% ===== BATCH 7 of the Proprium de Tempore (autumnalis.md lines
% 14651-14883) -- the complete fifth and final week of September:
% "Dominica V Septembris" Matins opening the Book of Esther in
% Nocturn I and St. Ambrose's "De Officiis" in Nocturn II (this
% Sunday's Nocturn III responsories given without the "vii"/"viii"
% numbering used at the parallel point in Dominicas I-IV, a genuine
% source variation transcribed exactly as given), the six ferial
% days completing the Esther cycle through Haman's downfall, and
% the brief "Sabbato ante Dominicam I Octobris" Vespers with the
% "Adapériat Dóminus" antiphon that the governing rubric several
% batches ago had flagged as proper to this point. This closes out
% the entire month of September in the Proprium de Tempore. A
% dozen-plus genuine partial-caps OCR artifacts are silently
% normalized to full caps per the established rule (ERat, CUmque
% [several occurrences], EVolúto, POst, AUdíerat, EGressúsque, QUod,
% EGréssus, RESponderúntque, NOctem, DIxítque, INtrávit,
% RESpondénsque), each confirmed by direct inspection of
% autumnalis.md; accents on each word's newly-capitalized remainder
% are dropped per the standing rule, while complete short words
% already fully capped in the source (QUÆ, QUI, DIE) are left
% untouched. The pervasive space-before-colon formatting quirk
% identified in the previous batch continues throughout this one and
% is likewise silently normalized to standard colon spacing. The
% document continues with "Dominica I Octobris" (autumnalis.md line
% 14884), opening the Books of Machabees and the month of October.
% =====
\feriaheading{Dominica I Octobris}

\rubric{Semiduplex}

\hourheading{Ad Matutinum}

\rubric{Invitatorium Adorémus Dóminum et Hymnus Primo die, ut in Psaltério 32: quæ dicúntur étiam in réliquis Dominícis usque ad V Novémbris inclusíve.}

\nocturnoheading{In I Nocturno}

\rubric{Incipit liber primus Machabæórum}

\lectioheading{Lectio i}{Cap. 1, 1-16}

\lettrine[lines=2]{E}{T} factum est, postquam percússit Alexánder Philíppi Mácedo, qui primus regnávit in Grǽcia, egréssus de terra Cethim, Dárium regem Persárum et Medórum, constítuit prǽlia multa et obtínuit ómnium munitiónes et interfécit reges terræ et pertránsiit usque ad fines terræ et accépit spólia multitúdinis géntium, et síluit terra in conspéctu ejus. Et congregávit virtútem et exércitum fortem nimis, et exaltátum est et elevátum cor ejus, et obtínuit regiónes géntium et tyránnos, et facti sunt illi in tribútum. Et post hæc décidit in lectum et cognóvit quia moerétur et vocávit púeros suos nóbiles, qui secum erant nutríti a juventúte, et divísit illis regnum suum, cum adhuc víveret.

\R\ Adapériat Dóminus cor vestrum in lege sua et in præcéptis suis et fáciat pacem in diébus vestris: * Concédat vobis salútem, et rédimat vos a malis. \V\ Exáudiat Dóminus oratiónes vestras et reconciliétur vobis nec vos déserat in témpore malo. — Concédat.

\lectioheading{Lectio ii}{}

\lettrine[lines=2]{E}{T} regnávit Alexánder annis duódecim et mórtuus est. Et obtinuérunt púeri ejus regnum, unusquísque in loco suo, et imposuérunt omnes sibi diadémata post mortem ejus, et fílii eórum post eos annis multis. Et multiplicáta sunt mala in terra. Et éxiit ex eis radix peccátrix, Antíochus illústris, fílius Antíochi regis qui fúerat Romæ obses, et regnávit in anno centésimo trigésimo séptimo regni Grræcórum.

\R\ Exáudiat Dóminus oratiónes vestras, et reconciliétur vobis nec vos déserat in témpore malo * Dóminus, Deus noster. \V\ Det vobis cor ómnibus, ut colátis eum et faciátis ejus voluntátem. — Dóminus.

\lectioheading{Lectio iii}{}

\lettrine[lines=2]{I}{N} diébus illis exiérunt ex Israël fílii iníqui et suasérunt multis dicéntes: Eámus et disponámus testaméntum cum géntibus, quæ circa nos sunt, quia, ex quo recéssimus ab eis, invenérunt nos multa mala. Et bonus visus est sermo in óculis eórum. Et destinavérunt áliqui de pópulo et abiérunt ad regem, et dedit illis potestátem ut fácerent justítiam géntium; et ædificavérunt gymnásium in Jerosólymis secúndum leges natiónum et fecérunt sibi præpútia et recessérunt a testaménto sancto et juncti sunt natiónibus et venúmdati sunt ut fácerent malum.

\R\ Congregáti sunt inimíci nostri, et gloriántur in virtúte sua: cóntere fortitúdinem illórum, Dómine, et dispérge illos: * Ut cognóscant quia non est álius qui pugnet pro nobis, nisi tu, Deus noster. \V\ Dispérge illos in virtúte tua, et déstrue eos, protéctor noster, Dómine. — Ut. Glória Patri. Ut.

\nocturnoheading{In II Nocturno}

\rubric{Ex libro Officiórum sancti Ambrósii Epíscopi}

\lectioheading{Lectio iv}{Lib. 1, cap. 40}

\lettrine[lines=2]{F}{ORTASSE} áliquos béllica defíxos glória tenet, ut putent solam esse præliárem fortitúdinem; et ídeo me ad ista deflexísse, quia illa nostris déforet. Quam fortis Jesus Nave, ut uno prǽlio quinque reges captos stérneret cum pópulis suis! Deínde cum advérsum Gabaonítas urgéret prǽlium, et vererétur ne nox impedíret victóriam, magnitúdine mentis et fídei clamávit: Stet sol, et stetit, donec victória consummarétur. Gédeon in trecéntis viris de ingénti pópulo et acérbo hoste revéxit triúmphum. Jónathas adoléscens virtútem magnam fecit in prǽlio.

\R\ Impetum inimicórum ne timuéritis: mémores estóte quómodo salvi facti sunt patres nostri: * Et nunc clamémus in cælum et miserébitur nostri Deus noster. \V\ Mementóte mirabílium ejus, quæ fecit pharaóni et exercítui ejus in Mari Rubro. — Et.

\lectioheading{Lectio v}{}

\lettrine[lines=2]{Q}{UID} de Machabǽis loquar? Sed prius de pópulo dicam patrum; qui cum essent paráti ad repugnándum pro templo Dei et pro legítimis suis, dolo hóstium die lacessíti sábbati, maluérunt vulnéribus offérre nuda córpora, quam repugnáre, ne violárent sábbatum. Itaque omnes læti se obtulérunt morti. Sed Machabǽi considerántes quod hoc exémplo gens omnis posset períre, sábbato étiam, cum ipsi in bellum provocaréntur, ulti sunt innocéntium necem fratrum suórum. Unde póstea stimulátus rex Antíochus, cum bellum accénderet per duces suos Lýsiam, Nicánorem, Górgiam, ita cum Orientálibus suis et Assýriis attrítus est cópiis, ut quadragínta et octo míllia in médio campi a tribus míllibus prosterneréntur.

\R\ Congregátæ sunt gentes in multitúdine, ut dímicent contra nos, et ignorámus quid ágere debeámus: * Dómine Deus, ad te sunt óculi nostri, ne pereámus. \V\ Tu scis quæ cógitant in nos: quómodo potérimus subsístere ante fáciem illórum, nisi tu ádjuves nos? — Dómine.

\lectioheading{Lectio vi}{}

\lettrine[lines=2]{V}{IRTUTEM} ducis Judæ Machabǽi de uno ejus mílite consideráte. Namque Eleázarus, cum supereminéntem céteris elephántem lorica vestítum régia advérteret, arbitrátus quod in eo esset rex, cursu cóncito in médium legiónis se prorúpit: et, abjécto clýpeo, utráque manu interficiébat, donec perveníret ad béstiam, atque intrávit sub eam, et subjécto gládio interémit eam. Itaque cadens béstia oppréssit Eleázarum, atque ita mórtuus est. Quanta ígitur virtus ánimi! primo, ut mortem non timéret; deínde, ut circumfúsus legiónibus inimicórum, in confértos raperétur hostes, médium penetráret agmen, et contémpta morte ferócior, abjécto clýpeo, utráque manu vulnerátæ molem béstiæ subíret ac sustinéret: post infra ipsam succéderet, quo plenióri feríret ictu; cujus ruína inclúsus magis quam oppréssus, suo est sepúltus triúmpho.

\R\ Tua est poténtia, tuum regnum, Dómine: tu es super omnes gentes: * Da pacem, Dómine, in diébus nostris. \V\ Creátor ómnium, Deus, terríbilis et fortis, justus, et miséricors. — Da. Glória Patri. Da.

\R\ Refúlsit sol in clýpeos áureos, et resplenduérunt montes ab eis: * Et fortitúdo géntium dissipáta est. \V\ Erat enim exércitus magnus valde et fortis: et appropiávit Judas et exércitus ejus in prǽlio. — Et.

\R\ Duo Séraphim, \rubric{ut infra in Homilíis.}

\rubric{Ad Laudes, in ómnibus Domínicis usque ad V Novémbris inclusíve, Hymnus Ætérne rerum, ut in Psaltério.}

\feriaheading{Feria Secunda infra Hebdomadam I Octobris}

\rubric{De libro primo Machabæórum}

\lectioheading{Lectio i}{Cap. 1, 17-29}

\lettrine[lines=2]{E}{T} parátum est regnum in conspéctu Antíochi, et cœpit regnáre in terra Ægýpti, ut regnáret super duo regna, et intrávit in Ægýptum in multitúdine gravi, in cúrribus et elephántis et equítibus et copiósa návium multitúdine, et constítuit bellum advérsus Ptolemǽum regem Ægýpti, et véritus est Ptolemǽus a fácie ejus et fugit, et cecidérunt vulneráti multi, et comprehéndit civitátes munítas in terra Ægýpti et accépit spólia terræ Ægýpti.

\R\ Dixit Judas Simóni fratri suo: Elige tibi viros et vade, líbera fratres tuos in Galilǽam; ego autem et Jónathas frater tuus íbimus in Galaadítim: * Sicut fúerit volúntas in cælo, sic fiat. \V\ Accingímini, fílii poténtes, et estóte paráti: quóniam mélius est nobis mori in bello, quam vidére mala gentis nostræ et sanctórum. — Sicut.

\lectioheading{Lectio ii}{}

\lettrine[lines=2]{E}{T} convértit Antíochus, postquam percússit Ægýptum in centésimo et quadragésimo tértio anno, et ascéndit ad Israël. Et ascéndit Jerosólymam in multitúdine gravi et intrávit in sanctificatiónem cum supérbia et accépit altáre áureum et candelábrum lúminis et univérsa vasa ejus et mensam propositiónis et libatória et phíalas et mortaríola áurea et velum et corónas et ornaméntum áureum, quod in fácie templi erat, et commínuit ómnia.

\R\ Ornavérunt fáciem templi corónis áureis, et dedicavérunt altáre Dómino: * Et facta est lætítia magna in pópulo. \V\ In hymnis et confessiónibus benedicébant Dóminum. — Et.

\rubric{In Octávis et Festis ritus simplícis: Glória Patri. Et.}

\lectioheading{Lectio iii}{}

\lettrine[lines=2]{E}{T} accépit argéntum et aurum et vasa concupiscibília et accépit thesáuros occúltos, quos invénit, et, sublátis ómnibus, ábiit in terram suam et fecit cædem hóminum et locútus est in supérbia magna et factus est planctus magnus in Israël et in omni loco eórum, et ingemuérunt príncipes et senióres, vírgines et júvenes infirmáti sunt, et speciósitas mulíerum immutáta est. Omnis marítus sumpsit laméntum, et quæ sedébant in thoro maritáli lugébant, et commóta est terra super habitántes in ea, et univérsa domus Jacob índuit confusiónem.

\R\ In hymnis et confessiónibus benedicébant Dóminum, * Qui magna fecit in Israël, et victóriam dedit illis Dóminus omnípotens. \V\ Ornavérunt fáciem templi corónis áureis, et dedicavérunt altáre Dómino. — Qui. Glória Patri. Qui.

\feriaheading{Feria Tertia infra Hebdomadam I Octobris}

\rubric{De libro primo Machabæórum}

\lectioheading{Lectio i}{Cap. 2, 1-10}

\lettrine[lines=2]{I}{N} diébus illis surréxit Mathathías fílius Joánnis, fílii Simeónis, sacérdos ex fíliis Jóarib ab Jerúsalem et consédit in monte Modin. Et habébat fílios quinque, Joánnem, qui cognominabátur Gaddis, et Simónem, qui cognominabátur Thasi, et Judam, qui vocabátur Machabǽus, et Eleázarum, qui cognominabátur Abaron, et Jónathan, qui cognominabátur Apphus. Hi vidérunt mala quæ fiébant in pópulo Juda et in Jerúsalem.

\R\ Hic est fratrum amátor et pópuli Israël: * Hic est, qui multum orat pro pópulo et univérsa sancta civitáte Jerúsalem. \V\ Vir iste in pópulo suo mitíssimus appáruit. — Hic est, qui.

\lectioheading{Lectio ii}{}

\lettrine[lines=2]{E}{T} dixit Mathathías: Væ mihi! ut quid natus sum vidére contritiónem pópuli mei et contritiónem civitátis sanctæ et sedére illic, cum datur in mánibus inimicórum? Sancta in manu extraneórum facta sunt, templum ejus sicut homo ignóbilis, vasa glóriæ ejus captíva abdúcta sunt, trucidáti sunt senes ejus in platéis, et júvenes ejus cecidérunt in gládio inimicórum. Quæ gens non hereditávit regnum ejus et non obtínuit spólia ejus?

\R\ Tu, Dómine universórum, qui nullam habes indigéntiam, voluísti templum tuum fíeri in nobis; * Consérva domum istam immaculátam in ætérnum, Dómine. \V\ Tu elegísti, Dómine, domum istam ad invocándum nomen tuum in ea, ut esset domus oratiónis et obsecratiónis pópulo tuo. — Consérva.

\rubric{In Octávis et Festis ritus simplícis: Glória Patri. Consérva.}

\lectioheading{Lectio iii}{Cap. 2, 14-16}

\lettrine[lines=2]{E}{T} scidit vestiménta sua Mathathías et fílii ejus et operuérunt se cilíciis et planxérunt valde. Et venérunt illuc qui missi erant a rege Antíocho, ut cógerent eos qui confúgerant in civitátem Modin, immoláre et accéndere thura et a lege Dei discédere, et multi de pópulo Israël consentiéntes accessérunt ad eos, sed Mathathías et fílii ejus constánter stetérunt.

\R\ Aperi óculos tuos, Dómine, et vide afflictiónem nostram: circumdedérunt nos gentes ad puniéndum nos: * Sed tu, Dómine, exténde bráchium tuum, et líbera ánimas nostras. \V\ Afflíge oppriméntes nos et contuméliam faciéntes in supérbiam; et custódi partem tuam. — Sed. Glória Patri. Sed.

\feriaheading{Feria Quarta infra Hebdomadam I Octobris}

\rubric{De libro primo Machabæórum}

\lectioheading{Lectio i}{Cap. 2, 19-30}

\lettrine[lines=2]{E}{T} respóndit Mathathías et dixit magna voce: Etsi omnes gentes regi Antíocho obédiunt, ut discédat unusquísque a servitúte legis patrum suórum et conséntiat mandátis ejus, ego et fílii mei et fratres mei obediémus legi patrum nostrórum. Propítius sit nobis Deus, non est nobis útile relínquere legem et justítias Dei: non audiémus verba regis Antíochi nec sacrificábimus transgrediéntes legis nostræ mandáta, ut eámus áltera via.

\R\ Refúlsit sol in clýpeos áureos, et resplenduérunt montes ab eis: * Et fortitúdo géntium dissipáta est. \V\ Erat enim exércitus magnus valde et fortis: et appropiávit Judas et exércitus ejus in prǽlio. — Et.

\lectioheading{Lectio ii}{}

\lettrine[lines=2]{E}{T}, ut cessávit loqui verba hæc, accéssit quidam Judǽus in ómnium óculis sacrificáre idólis super aram in civitáte Modin secúndum jussum regis. Et vidit Mathathías et dóluit, et contremuérunt renes ejus, et accénsus est furor ejus secúndum judícium legis, et insíliens trucidávit eum super aram. Sed et virum quem rex Antíochus míserat, qui cogébat immoláre, occídit in ipso témpore et aram destrúxit et zelátus est legem, sicut fecit Phínees Zamri fílio Salómi.

\R\ Ornavérunt fáciem templi corónis áureis, et dedicavérunt altáre Dómino: * Et facta est lætítia magna in pópulo. \V\ In hymnis et confessiónibus benedicébant Dóminum. — Et.

\rubric{In Octávis et Festis ritus simplícis: Glória Patri. Et.}

\lectioheading{Lectio iii}{}

\lettrine[lines=2]{E}{T} exclamávit Mathathías voce magna in civitáte dicens: Omnis qui zelum habet legis státuens testaméntum éxeat post me. Et fugit ipse et fílii ejus in montes, et reliquérunt quæcúmque habébant in civitáte. Tunc descendérunt multi quæréntes judícium et justítiam in desértum. Et sedérunt ibi ipsi et fílii eórum et mulíeres eórum et pécora eórum, quóniam inundavérunt super eos mala.

\R\ In hymnis et confessiónibus benedicébant Dóminum, * Qui magna fecit in Israël, et victóriam dedit illis Dóminus omnípotens. \V\ Ornavérunt fáciem templi corónis áureis, et dedicavérunt altáre Dómino. — Qui. Glória Patri. Qui.

\feriaheading{Feria Quinta infra Hebdomadam I Octobris}

\rubric{De libro primo Machabæórum}

\lectioheading{Lectio i}{Cap. 2, 49-69}

\lettrine[lines=2]{E}{T} appropinquavérunt dies Mathathíæ moriéndi, et dixit fíliis suis: Nunc confortáta est supérbia et castigátio et tempus eversiónis et ira indignatiónis. Nunc ergo, o fílii, æmulatóres estóte legis et date ánimas vestras pro testaménto patrum vestrórum et mementóte óperum patrum, quæ fecérunt in generatiónibus suis, et accipiétis glóriam magnam et nomen ætérnum. Abraham, nonne in tentatióne invéntus est fidélis, et reputátum est ei ad justítiam? Joseph in témpore angústiæ suæ custodívit mandátum, et factus est dóminus Ægýpti; Phínees pater noster zelándo zelum Dei accépit testaméntum sacerdótii ætérni.

\R\ Adapériat Dóminus cor vestrum in lege sua et in præcéptis suis et fáciat pacem in diébus vestris: * Concédat vobis salútem, et rédimat vos a malis. \V\ Exáudiat Dóminus oratiónes vestras et reconciliétur vobis nec vos déserat in témpore malo. — Concédat.

\lectioheading{Lectio ii}{}

\lettrine[lines=2]{J}{ESUS}, dum implévit verbum, factus est dux in Israël; Caleb, dum testificátur in ecclésia, accépit hereditátem; David in sua misericórdia consecútus est sedem regni in sǽcula; Elías, dum zelat zelum legis, recéptus est in cælum; Ananías et Azarías et Mísael credéntes liberáti sunt de flamma; Dániel in sua simplicitáte liberátus est de ore leónum; et ita cogitáte per generatiónem et generatiónem, quia omnes qui sperant in eum non infirmántur. Et a verbis viri peccatóris ne timuéritis, quia glória ejus stercus et vermis est: hódie extóllitur et cras non inveniétur, quia convérsus est in terram suam, et cogitátio ejus périit.

\R\ Exáudiat Dóminus oratiónes vestras, et reconciliétur vobis nec vos déserat in témpore malo * Dóminus, Deus noster. \V\ Det vobis cor ómnibus, ut colátis eum et faciátis ejus voluntátem. — Dóminus.

\rubric{In Octávis et Festis ritus simplícis: Glória Patri. Dóminus.}

\lectioheading{Lectio iii}{}

\lettrine[lines=2]{V}{OS} ergo, fílii, confortámini et viríliter ágite in lege, quia in ipsa gloriósi éritis. Et ecce Simon frater vester, scio quod vir consílii est: ipsum audíte semper, et ipse erit vobis pater; et Judas Machabǽus fortis víribus a juventúte sua sit vobis princeps militíæ et ipse aget bellum pópuli, et adducétis ad vos omnes factóres legis et vindicáte vindíctam pópuli vestri, retribúite retributiónem géntibus et inténdite in præcéptum legis. Et benedíxit eos et appósitus est ad patres suos.

\R\ Congregáti sunt inimíci nostri, et gloriántur in virtúte sua: cóntere fortitúdinem illórum, Dómine, et dispérge illos: * Ut cognóscant quia non est álius qui pugnet pro nobis, nisi tu, Deus noster. \V\ Dispérge illos in virtúte tua, et déstrue eos, protéctor noster, Dómine. — Ut. Glória Patri. Ut.

\feriaheading{Feria Sexta infra Hebdomadam I Octobris}

\rubric{De libro primo Machabæórum}

\lectioheading{Lectio i}{Cap. 2, 70; 3, 1-3, 5-12, 25-28}

\lettrine[lines=2]{D}{EFUNCTUS} est Mathathías anno centésimo et quadragésimo sexto et sepúltus est a fíliis suis in sepúlcris patrum suórum in Modin, et planxérunt eum omnis Israël planctu magno. Et surréxit Judas, qui vocabátur Machabǽus, fílius ejus pro eo et adjuvábant eum omnes fratres ejus et univérsi qui se conjúnxerant patri ejus et præliabántur prǽlium Israël cum lætítia. Et dilatávit glóriam pópulo suo, et persecútus est iníquos perscrútans eos et, qui conturbábant pópulum suum, eos succéndit flammis, et repúlsi sunt inimíci ejus præ timóre ejus, et omnes operárii iniquitátis conturbáti sunt, et dirécta est salus in manu ejus.

\R\ Impetum inimicórum ne timuéritis: mémores estóte, quómodo salvi facti sunt patres nostri: * Et nunc clamémus in cælum, et miserébitur nostri Deus noster. \V\ Mementóte mirabílium ejus, quæ fecit Pharaóni et exercítui ejus in Mari Rubro. — Et.

\lectioheading{Lectio ii}{}

\lettrine[lines=2]{E}{T} exacerbábat reges multos et lætificábat Jacob in opéribus suis, et in sǽculum memória ejus in benedictióne, et perambulávit civitátes Juda et pérdidit ímpios ex eis et avértit iram ab Israël et nominátus est usque ad novíssimum terræ et congregávit pereúntes. Et congregávit Apollónius gentes et a Samaría virtútem multam et magnam ad bellándum contra Israël. Et cognóvit Judas et éxiit óbviam illi et percússit et occídit illum, et cecidérunt vulneráti multi, et réliqui fugérunt. Et accépit spólia eórum et gládium Apollónii ábstulit Judas et erat pugnans in eo ómnibus diébus.

\R\ Congregátæ sunt gentes in multitúdine, ut dímicent contra nos, et ignorámus quid ágere debeámus: * Dómine Deus, ad te sunt óculi nostri, ne pereámus. \V\ Tu scis quæ cógitant in nos: quómodo potérimus subsístere ante fáciem illórum, nisi tu ádjuves nos? — Dómine.

\rubric{In Octávis et Festis ritus simplícis: Glória Patri. Dómine.}

\lectioheading{Lectio iii}{}

\lettrine[lines=2]{E}{T} cécidit timor Judæ ac fratrum ejus et formído super omnes gentes in circúitu eórum; et pervénit ad regem nomen ejus, et de prǽliis Judæ narrábant omnes gentes. Ut audívit autem rex Antíochus sermónes istos, irátus est ánimo et misit et congregávit exércitum univérsi regni sui, castra fórtia valde, et apéruit ærárium suum et dedit stipéndia exercítui in annum et mandávit illis, ut essent paráti ad ómnia.

\R\ Tua est poténtia, tuum regnum, Dómine: tu es super omnes gentes: * Da pacem, Dómine, in diébus nostris. \V\ Creátor ómnium, Deus, terríbilis et fortis, justus et miséricors. — Da. Glória Patri. Da.

\feriaheading{Sabbato infra Hebdomadam I Octobris}

\rubric{De libro primo Machabæórum}

\lectioheading{Lectio i}{Cap. 3, 42-60}

\lettrine[lines=2]{E}{T} vidit Judas et fratres ejus quia multiplicáta sunt mala et exércitus applicábant ad fines eórum, et cognovérunt verba regis, quæ mandávit pópulo fácere in intéritum et consummatiónem, et dixérunt unusquísque ad próximum suum: Erigámus dejectiónem pópuli nostri et pugnémus pro pópulo nostro et sanctis nostris. Et congregátus est convéntus, ut essent paráti in prǽlium et ut orárent et péterent misericórdiam et miseratiónes. Et Jerúsalem non habitabátur, sed erat sicut desértum: non erat qui ingrederétur et egrederétur de natis ejus, et sanctum conculcabátur, et fílii alienigenárum erant in arce: ibi erat habitátio géntium, et abláta est volúptas a Jacob, et defécit ibi tíbia et cíthara.

\R\ Refúlsit sol in clýpeos áureos, et resplenduérunt montes ab eis: * Et fortitúdo géntium dissipáta est. \V\ Erat enim exércitus magnus valde et fortis: et appropiávit Judas et exércitus ejus in prǽlio. — Et.

\lectioheading{Lectio ii}{}

\lettrine[lines=2]{E}{T} congregáti sunt et venérunt in Maspha contra Jerúsalem, quia locus oratiónis erat in Maspha ante in Israël, et jejunavérunt illa die et induérunt se cilíciis et cínerem imposuérunt cápiti suo et discidérunt vestiménta sua et expandérunt libros legis, de quibus scrutabántur gentes similitúdinem simulacrórum suórum, et attulérunt ornaménta sacerdotália et primítias et décimas et suscitavérunt Nazarǽos, qui impléverant dies, et clamavérunt voce magna in cælum dicéntes: Quid faciémus istis et quo eos ducémus? Et sancta tua conculcáta sunt et contamináta sunt, et sacerdótes tui facti sunt in luctum et in humilitátem, et ecce natiónes convenérunt advérsum nos, ut nos dispérdant: tu scis quæ cógitant in nos. Quómodo potérimus subsístere ante fáciem eórum, nisi tu, Deus, ádjuves nos?

\R\ Ornavérunt fáciem templi corónis áureis, et dedicavérunt altáre Dómino: * Et facta est lætítia magna in pópulo. \V\ In hymnis et confessiónibus benedicébant Dóminum. — Et.

\rubric{In Octávis simplícibus et in Offício sanctæ Maríæ in Sábbato: Glória Patri. Et.}

\lectioheading{Lectio iii}{}

\lettrine[lines=2]{E}{T} tubis exclamavérunt voce magna. Et post hæc constítuit Judas duces pópuli, tribúnos et centuriónes et pentacontárchos et decuriónes, et dixit his, qui ædificábant domos et sponsábant uxóres et plantábant víneas et formidolósis, ut redírent unusquísque in domum suam secúndum legem. Et movérunt castra et collocavérunt ad Austrum Emmaum. Et ait Judas: Accingímini et estóte fílii poténtes et estóte paráti in mane, ut pugnétis advérsus natiónes has, quæ convenérunt advérsus nos dispérdere nos et sancta nostra; quóniam mélius est nos mori in bello, quam vidére mala gentis nostræ et sanctórum. Sicut autem fúerit volúntas in cælo, sic fiat.

\R\ In hymnis et confessiónibus benedicébant Dóminum, * Qui magna fecit in Israël, et victóriam dedit illis Dóminus omnípotens. \V\ Ornavérunt fáciem templi corónis áureis, et dedicavérunt altáre Dómino. — Qui. Glória Patri. Qui.

\hourheading{Ad Vesperas}

\V\ Vespertína orátio ascéndat ad te, Dómine. \R\ Et descéndat super nos misericórdia tua.

\antlabel{Ad Magnif. Ant.} Refúlsit sol * in clýpeos áureos, et resplenduérunt montes ab eis: et fortitúdo géntium dissipáta est.

\rubric{Oratio Dominicæ occurréntis.}

% ===== BATCH 8 of the Proprium de Tempore (autumnalis.md lines
% 14884-15116) -- the complete first week of October: "Dominica I
% Octobris" Matins opening the First Book of Machabees in Nocturn I
% and St. Ambrose's "De Officiis" in Nocturn II, the six ferial days
% continuing the Machabees cycle through Mathathias's death and
% Judas's early victories, and the closing Saturday Vespers. Both
% Nocturn III's responsories (as at the two immediately preceding
% Sundays) are given without the "vii"/"viii" numbering used at
% Dominicas I-IV Septembris. A new governing rubric fixes the hymn
% "Ætérne rerum" at Lauds for every Sunday through November 5. One
% genuine partial-caps OCR artifact is silently normalized
% (DEfúnctus to DEFUNCTUS, accent dropped on its newly-capitalized
% remainder), while complete short words already fully capped in
% the source (VOS, JESUS) are left untouched. Feria Sexta's citation
% ("Cap. 2, 70; 3, 1-3, 5-12, 25-28") appears in the source one line
% above its "Lectio i" header rather than below it as usual; it is
% nonetheless given as that Lesson's citation via
% \lectioheading{Lectio i}{...}, consistent with its evident
% function regardless of its exact placement in the source. The
% pervasive space-before-colon formatting quirk from the previous
% two batches continues here and is likewise silently normalized.
% The document continues with "Dominica II Octobris" (autumnalis.md
% line 15117).
% =====
       % or hiemalis.tex / verna.tex / aestiva.tex
%   \end{document}
%
% Since autumnalis.tex, aestiva.tex, hiemalis.tex, and verna.tex all
% now embed the versicle/response marks (℣ U+2123, ℟ U+211F) and the
% blessing cross (☩ U+2629) as literal Unicode characters rather than
% as macros, and the main body font (DejaVu Serif) does not contain
% glyphs for them, \newunicodechar is used below to substitute a font
% that does (DejaVu Sans) automatically wherever those three
% characters occur, without having to touch the body text itself.
% =====================================================================

\usepackage{fontspec}
\usepackage{textcomp}
\usepackage{lettrine}
\usepackage[dvipsnames]{xcolor}
\usepackage[margin=2cm]{geometry}
\usepackage{parskip}
\usepackage{newunicodechar}

\setmainfont{DejaVu Serif}
\newfontfamily\liturgicalsymbolfont{DejaVu Sans}

\newunicodechar{℣}{{\liturgicalsymbolfont\char"2123}}  % VERSICLE
\newunicodechar{℟}{{\liturgicalsymbolfont\char"211F}}  % RESPONSE
\newunicodechar{☩}{{\liturgicalsymbolfont\char"2629}}  % cross (blessing)

% --- Headings (one required argument each) ---
\newcommand{\partheading}[1]{{\bigskip\noindent\LARGE\bfseries #1\par\medskip}}
\newcommand{\partsubtitle}[1]{{\noindent\large\itshape #1\par\medskip}}
\newcommand{\sectionheading}[1]{{\bigskip\noindent\Large\bfseries #1\par\medskip}}
\newcommand{\feriaheading}[1]{{\bigskip\noindent\large\bfseries #1\par\smallskip}}
\newcommand{\hourheading}[1]{{\bigskip\noindent\large\bfseries #1\par\smallskip}}
\newcommand{\nocturnoheading}[1]{{\medskip\noindent\bfseries #1\par\smallskip}}
\newcommand{\psalmheading}[1]{{\medskip\noindent\itshape #1\par\smallskip}}
\newcommand{\hymnusheading}[1]{{\medskip\noindent\bfseries #1\par\smallskip}}
\newcommand{\oratioheading}[1]{{\medskip\noindent\itshape #1\par\smallskip}}
\newcommand{\variantheading}[1]{{\medskip\noindent\itshape #1\par\smallskip}}
\newcommand{\absolutioheading}[1]{{\medskip\noindent\itshape #1\par\smallskip}}
\newcommand{\suffragiumheading}[1]{{\medskip\noindent\bfseries #1\par\smallskip}}
\newcommand{\symbolumheading}[1]{{\medskip\noindent\bfseries #1\par\smallskip}}
\newcommand{\antiphonaheading}[1]{{\medskip\noindent\bfseries #1\par\smallskip}}
\newcommand{\feastheading}[1]{{\bigskip\noindent\large\bfseries #1\par\smallskip}}
\newcommand{\sanctusheading}[1]{{\bigskip\noindent\large\bfseries #1\par\smallskip}}
\newcommand{\diesheading}[1]{{\medskip\noindent\bfseries #1\par\smallskip}}
\newcommand{\monthheading}[1]{{\bigskip\noindent\Large\bfseries #1\par\medskip}}

% --- Purely decorative divider (no argument) ---
\newcommand{\sectionornament}{{\par\noindent\centering\rule{0.3\linewidth}{0.4pt}\par\medskip}}

% --- Headings (two required arguments: title + citation) ---
\newcommand{\lectioheading}[2]{{\medskip\noindent\itshape #1\quad\textit{#2}\par\smallskip}}
\newcommand{\capitulumlabel}[2]{{\medskip\noindent\itshape #1\quad\textit{#2}\par\smallskip}}
\newcommand{\canticumheading}[2]{{\medskip\noindent\itshape #1\quad\textit{#2}\par\smallskip}}
\newcommand{\lectiolabel}[2]{{\medskip\noindent\itshape #1\quad\textit{#2}\par\smallskip}}

% --- Labels (one required argument) ---
\newcommand{\antlabel}[1]{\textbf{#1}}
\newcommand{\invitatlabel}[1]{\textbf{#1}}
\newcommand{\psalmlabel}[1]{\textbf{#1}}
\newcommand{\benedlabel}[1]{\textbf{#1}}
\newcommand{\benedictiolabel}[1]{\textbf{#1}}
\newcommand{\absolutiolabel}[1]{\textbf{#1}}
\newcommand{\formulalabel}[1]{\textbf{#1}}
\newcommand{\oratiolabel}[1]{\textbf{#1}}

% --- Rubric (red text; embedded rubrics may nest \textit, etc.) ---
\newcommand{\rubric}[1]{\textcolor{red}{#1}}
% Black text embedded within a red \rubric{} span (governing instruction:
% "embedded rubrics (black text within red-colored rubrics) are preserved").
\newcommand{\rubricblack}[1]{\textcolor{black}{#1}}

% \P (pilcrow), \dag, \ddag, \label, \pageref, \hspace, \textbf,
% \textsuperscript, \textit, \textsc, \textdagger, \textdaggerdbl are all
% standard LaTeX/textcomp commands and need no definition here.
  (or copy its contents into a wrapper)
%   \begin{document}
%   % =====================================================================
% autumnalis.tex -- LaTeX conversion of autumnalis.md (Breviarium
% Romanum, Pars Autumnalis), BODY CONTENT ONLY.
%
% No preamble, \documentclass, or \newcommand definitions are included
% here, per instruction -- the command vocabulary below follows the
% same legend established for hiemalis.tex, verna.tex, and aestiva.tex
% (one-column vs. two-column, red vs. black, etc.).
%
% BATCH 1 of N: Ordinarium Divini Officii, Ad Matutinum and Ad Laudes
% (autumnalis.md lines 1-463)
%
% NOTE: autumnalis.md lines 1-1141 (the entire Ordinarium Divini
% Officii) were verified via direct diff to be BYTE-IDENTICAL to the
% corresponding lines of aestiva.md -- both being non-Lenten/
% non-Paschal volumes, each omits the same "Tempore Quadragesimæ" /
% "Tempore Passionis" Invitatorium-Hymnus variants and the Paschaltide
% "Commemoratio de Cruce" that verna.tex's Ordinarium carries. Given
% this confirmed byte-for-byte identity (not a parallel or merely
% similar occurrence, but the literal same source text), the verified
% LaTeX already produced for aestiva.tex's Batches 1-3 is reused
% directly below rather than re-keyed by hand, to avoid reintroducing
% errors (such as the accent-dropping mistakes) already caught and
% fixed in that earlier work. Every batch from "Psalterium Breviarii
% Romani" onward is fresh transcription from autumnalis.md's own text,
% as usual.
% =====================================================================

% [ornamental image beginning Pars Autumnalis]
\partheading{Ordinarium Divini Officii}
\partsubtitle{Juxta Ritum Romanum Persolvendi}

\rubric{Clericis in sacris constitútis, qui íntegrum divínum Offícium, quamvis in partes distribútum, coram Ssmo Sacraménto, sive públicæ adoratióni expósito sive in tabernáculo adserváto, devóte recitáverint, concéditur: Indulgéntia plenária, si prætérea sacramentálem confessiónem instítuerint, ad eucharísticum Convívium accésserint et ad mentem Summi Pontíficis oráverint. \textit{Pius Pp. XI, 23 Octóbris 1930.}}

\rubric{Iis vero, qui divínum Offícium coram Ssmo Sacraménto, ut supra, ex parte tantum recitáverint, concéditur: Indulgéntia quingentórum diérum pro unaquáque hora canónica. \textit{Pius Pp. XI, 18 Maji 1933.}}

\rubric{Iisdem autem cléricis, qui divínum Offícium in álias preces commutátum rite obtinúerint, si preces ejúsmodi coram Ssmo Sacraménto, uti supra, devóte recitáverint, concéditur: Indulgéntia plenária, addítis sacramentáli confessióne, sacra Communióne et oratióne ad mentem Summi Pontíficis. \textit{Pius Pp. XI, 7 Novémbris 1932.}}

\sectionheading{Ante Divinum Officium}

\rubric{Antequam inchoëtur Offícium laudabíliter dícitur, sub singulári semper número, séquens Orátio; pro qua Summus Póntifex Pius Pp. XI sub die 17 Novémbris 1933 Indulgéntiam trium annórum concéssit.}

\oratioheading{Oratio}

\lettrine[lines=2]{A}{PERI}, Dómine, os meum ad benedicéndum nomen sanctum tuum: munda quoque cor meum ab ómnibus vanis, pervérsis et aliénis cogitatiónibus; intelléctum illúmina, afféctum inflámma, ut digne, atténte ac devóte hoc Offícium recitáre váleam, et exaudíri mérear ante conspéctum divínæ Majestátis tuæ. Per Christum, Dóminum nostrum. \R\ Amen.

Dómine, in unióne illíus divínæ intentiónis, qua ipse in terris laudes Deo persolvísti, has tibi Horas (\textit{vel} hanc tibi Horam) persólvo.

\sectionheading{Ante singulas Horas}

\rubric{Ante Matutínum, sicut et ante Laudes a Matutíno extra Chorum separátas, et ante omnes Horas, præter quam ad Completórium, dícitur secréto:}

\lettrine[lines=2]{P}{ATER} noster, qui es in cælis, sanctificétur nomen tuum. Advéniat regnum tuum. Fiat volúntas tua, sicut in cælo et in terra. Panem nostrum quotidiánum da nobis hódie. Et dimítte nobis débita nostra, sicut et nos dimíttimus debitóribus nostris. Et ne nos indúcas in tentatiónem: sed líbera nos a malo. Amen.

\lettrine[lines=2]{A}{VE}, María, grátia plena; Dóminus tecum: benedícta tu in muliéribus, et benedíctus fructus ventris tui Jesus. Sancta María, Mater Dei, ora pro nobis peccatóribus, nunc et in hora mortis nostræ. Amen.

\rubric{In princípio autem Matutíni ac Primæ additur secréto Symbolum Apostolórum:}

\lettrine[lines=2]{C}{REDO} in Deum, Patrem omnipoténtem, Creatórem cæli et terræ. Et in Jesum Christum Fílium ejus únicum, Dóminum nostrum: qui concéptus est de Spíritu Sancto, natus ex María Vírgine, passus sub Póntio Piláto, crucifíxus, mórtuus, et sepúltus: descéndit ad ínferos; tértia die resurréxit a mórtuis; ascéndit ad cælos; sedet ad déxteram Dei Patris omnipoténtis: inde ventúrus est judicáre vivos et mórtuos. Credo in Spíritum Sanctum, sanctam Ecclésiam cathólicam, Sanctórum communiónem, remissiónem peccatórum, carnis resurrectiónem, vitam ætérnam. Amen.

\rubric{Officium ad Matutinum et ad reliquas Horas, nisi aliter suo loco indicetur, per totum annum persolvitur juxta formam, quæ in Rubricis hujus Ordinarii exhibetur.}

\hourheading{Ad Matutinum}

Pater noster, Ave María et Credo.

\rubric{Deinde, clara voce, dicitur Versus:}

\V\ \lettrine[lines=2]{D}{OMINE}, lábia mea apéries. \R\ Et os meum annuntiábit laudem tuam. \V\ Deus, in adjutórium meum inténde. \R\ Dómine, ad adjuvándum me festína.

Glória Patri, et Fílio, et Spirítui Sancto. Sicut erat in princípio, et nunc, et semper, et in sæcula sæculórum. Amen. Allelúja.

\rubric{Postea dicitur conveniens Invitatorium, quod ante Psalmum bis integrum recitatur, et ad singulos ejusdem Psalmi versus vel integrum vel dimidiatum ab Asterisco * alternis vicibus, ut infra, repetitur.}

\psalmheading{Psalmus 94}

\lettrine[lines=2]{V}{ENITE}, exsultémus Dómino, jubilémus Deo, salutári nostro: præoccupémus fáciem ejus in confessióne, et in psalmis jubilémus ei.

\rubric{Repetitur integrum Invitatorium.}

Quóniam Deus magnus Dóminus, et Rex magnus super omnes deos: quóniam non repéllet Dóminus plebem suam: quia in manu ejus sunt omnes fines terræ, et altitúdines móntium ipse cónspicit.

\rubric{Repetitur altera pars Invitatorii.}

\rubric{In sequenti Psalmi versu, ad verba veníte, adorémus, et procidámus ante Deum, genuflectitur.}

Quóniam ipsíus est mare, et ipse fecit illud, et áridam fundavérunt manus ejus: veníte, adorémus, et procidámus ante Deum: plorémus coram Dómino, qui fecit nos, quia ipse est Dóminus, Deus noster; nos autem pópulus ejus, et oves páscuæ ejus.

\rubric{Repetitur integrum Invitatorium.}

Hódie, si vocem ejus audíeritis, nolíte obduráre corda vestra, sicut in exacerbatióne secúndum diem tentatiónis in desérto: ubi tentavérunt me patres vestri, probavérunt et vidérunt ópera mea.

\rubric{Repetitur altera pars Invitatorii.}

Quadragínta annis próximus fui generatióni huic, et dixi: Semper hi errant corde; ipsi vero non cognovérunt vias meas: quibus jurávi in ira mea: Si introíbunt in réquiem meam.

\rubric{Repetitur integrum Invitatorium.}

Glória Patri, et Fílio, et Spirítui Sancto. Sicut erat in princípio, et nunc, et semper, et in sæcula sæculórum. Amen.

\rubric{Repetitur altera pars Invitatorii.}

\rubric{Denique repetitur integrum Invitatorium.}

\rubric{Expleto Psalmo, dicitur Hymnus Invitatorio respondens.}

\variantheading{In Officio dominicali}

\rubric{In omnibus Dominicis, extra Octavas Ssmi Corporis Christi ac Sacrat. Cordis Jesu, Invitatorium et Hymnus ut in Psalterio.}

\variantheading{In Officio feriali}

\rubric{In omnibus Feriis, Invitatorium et Hymnus de Feria currenti, ut in Psalterio.}

\variantheading{In Festis}

\rubric{In Officio cujuslibet Festi vel Octavæ, ac S. Mariæ in Sabbato, Invitatorium et Hymnus ut in Proprio vel Communi; in Dominicis autem infra Octavas Ssmi Corporis Christi ac Sacratissimi Cordis Jesu, si de eis fiat Officium, Invitatorium et Hymnus ut in Proprio de Tempore. Quæ norma ad reliquas item Horas erit observanda pro illis Officii partibus, quæ juxta Rubricas ex Ordinario vel Psalterio non fuerint in Festis desumendæ.}

\variantheading{In Officio novem Lectionum}

\rubric{Expleto Hymno, dicuntur Antiphonæ convenientes, quæ in Officiis ritus duplicis ante et post Psalmos integræ recitantur; in Officiis autem ritus semiduplicis initio Psalmi inchoantur tantum et usque ad Asteriscum * perducuntur, atque in fine integræ pronuntiantur; servata alterutra norma etiam in omnibus Antiphonis, quæ ad Laudes ac Vesperas adhibentur.}

\nocturnoheading{In I Nocturno}

\rubric{Sub congruentibus Antiphonis dicuntur tres Psalmi, ac deinde subjungitur Versus, prouti Officium occurrens requirit.}

\rubric{Post Versum cujuslibet Nocturni dicitur}

Pater noster \rubric{secreto usque ad} \V\ Et ne nos indúcas in tentatiónem. \R\ Sed líbera nos a malo.

\absolutioheading{Absolutio}

\lettrine[lines=2]{E}{XÁUDI}, Dómine Jesu Christe, preces servórum tuórum, et miserére nobis: Qui cum Patre et Spíritu Sancto vivis et regnas in sæcula sæculórum. \R\ Amen.

\V\ Jube, domne, benedícere.

\rubric{¶ Extra Chorum, quando ab uno tantum recitatur Officium, ante singulas Lectiones Matutini atque ad Lectionem brevem Primæ et Completorii, dicitur: Jube, Dómine, benedícere; et subjungitur congruens Benedictio. Ab Episcopo autem, ultimam Matutini Lectionem cantaturo, item dicitur: Jube, Dómine, benedícere; et respondetur a Choro: Amen.}

\variantheading{Pro I Lectione}

\benedlabel{Bened.} Benedictióne perpétua benedícat nos Pater ætérnus. \R\ Amen.

\rubric{Deinde dicuntur in unoquoque Nocturno Lectiones, prouti Officium occurrens requirit, et in fine cujuslibet Lectionis additur:}

\V\ Tu autem, Dómine, miserére nobis. \R\ Deo grátias.

\rubric{Post quamlibet vero Lectionem, quæ Hymnum Te Deum immediate non præcedat, congruens dicitur Responsorium, et in fine ultimi Responsorii cujusque Nocturni additur Versus: Glória Patri, et Fílio, et Spirítui Sancto, et Responsorium ab Asterisco *, et quidem a secundo, si duo fuerint, repetitur.}

\variantheading{Pro II Lectione}

\benedlabel{Bened.} Unigénitus Dei Fílius nos benedícere et adjuváre dignétur. \R\ Amen.

\variantheading{Pro III Lectione}

\benedlabel{Bened.} Spíritus Sancti grátia illúminet sensus et corda nostra. \R\ Amen.

\nocturnoheading{In II Nocturno}

\rubric{Sub congruentibus item Antiphonis dicuntur tres Psalmi et Versus, sicut in I Nocturno.}

\rubric{Post Versum dicitur.}

Pater noster \rubric{secreto usque ad} \V\ Et ne nos indúcas in tentatiónem. \R\ Sed líbera nos a malo.

\absolutioheading{Absolutio}

\lettrine[lines=2]{I}{PSIUS} píetas et misericórdia nos ádjuvet, qui cum Patre et Spíritu Sancto vivit et regnat in sæcula sæculórum. \R\ Amen.

\V\ Jube, domne, benedícere.

\variantheading{Pro IV Lectione}

\benedlabel{Bened.} Deus Pater omnípotens sit nobis propítius et clemens. \R\ Amen.

\variantheading{Pro V Lectione}

\benedlabel{Bened.} Christus perpétuæ det nobis gáudia vitæ. \R\ Amen.

\variantheading{Pro VI Lectione}

\benedlabel{Bened.} Ignem sui amóris accéndat Deus in córdibus nostris. \R\ Amen.

\nocturnoheading{In III Nocturno}

\rubric{Sub congruentibus denique Antiphonis dicuntur tres Psalmi et Versus, sicut in I et II Nocturno.}

\rubric{Post Versum dicitur}

Pater noster \rubric{secreto usque ad} \V\ Et ne nos indúcas in tentatiónem. \R\ Sed líbera nos a malo.

\absolutioheading{Absolutio}

\lettrine[lines=2]{A}{} \textsc{vinculis} peccatórum nostrórum absólvat nos omnípotens et miséricors Dóminus. \R\ Amen.

\V\ Jube, domne, benedícere.

\variantheading{Pro VII Lectione}

\benedlabel{Bened.} Evangélica léctio sit nobis salus et protéctio. \R\ Amen.

\variantheading{Pro VIII Lectione}

\rubric{In Festis Domini et in Dominicis:}

\benedlabel{Bened.} Divínum auxílium máneat semper nobíscum. \R\ Amen.

\rubric{In Festis beatæ Mariæ Virginis:}

\benedlabel{Bened.} Cujus festum cólimus, ipsa Virgo virgínum intercédat pro nobis ad Dóminum. \R\ Amen.

\rubric{In Festis Sanctorum:}

\benedlabel{Bened.} Cujus (\textit{vel} Quorum \textit{aut} Quarum) festum cólimus, ipse (\textit{vel} ipsa \textit{aut} ipsi \textit{vel} ipsæ) intercédat (\textit{vel} intercédant) pro nobis ad Dóminum. \R\ Amen.

\variantheading{Pro IX Lectione}

\benedlabel{Bened.} Ad societátem cívium supernórum perdúcat nos Rex Angelórum. \R\ Amen.

\rubric{Si autem legenda sit Lectio ix de Homilia cum Evangelio Dominicæ, vel Feriæ, aut Vigiliæ:}

\benedlabel{Bened.} Per evangélica dicta deleántur nostra delícta. \R\ Amen.

\variantheading{In Officio trium Lectionum}

\nocturnoheading{In Nocturno}

\rubric{In Officiis trium Lectionum, post Hymnum dicuntur Antiphonæ convenientes, quæ initio Psalmi inchoantur tantum et usque ad Asteriscum * perducuntur, ac deinde in fine integræ pronuntiantur.}

\rubric{Sub eisdem vero Antiphonis dicuntur novem Psalmi Feriæ currentis, quibus subjungitur Versus in III Nocturno positus, omissis Versibus pro I et II Nocturno assignatis.}

\rubric{Post Versum dicitur}

Pater noster \rubric{secreto usque ad} \V\ Et ne nos indúcas in tentatiónem. \R\ Sed líbera nos a malo.

\absolutioheading{Absolutio}

\rubric{Feria II et V}

Exáudi, Dómine, \rubric{ut in I Nocturno.}

\rubric{Feria III et VI}

Ipsíus píetas, \rubric{ut in II Nocturno.}

\rubric{Feria IV et Sabbato}

A vínculis, \rubric{ut in III Nocturno Officii novem Lectionum.}

\rubric{Deinde leguntur Lectiones cum Responsoriis, prouti Officium occurrens requirit; et ante eas dicuntur sequentes}

\variantheading{Benedictiones}

\rubric{In Feriis, quando legitur Homilia cum Evangelio, erit:}

\benedlabel{1 Bened.} Evangélica léctio sit nobis salus et protéctio. \R\ Amen.

\benedlabel{2 Bened.} Divínum auxílium máneat semper nobíscum. \R\ Amen.

\benedlabel{3 Bened.} Ad societátem cívium supernórum perdúcat nos Rex Angelórum. \R\ Amen.

\rubric{In Feriis, quando non legitur Homilia cum Evangelio, Feria II et V Benedictiones ut in I Nocturno, Feria III et VI ut in II Nocturno Officii novem Lectionum; Feria autem IV et Sabbato erit:}

\benedlabel{1 Bened.} Ille nos benedícat, qui sine fine vivit et regnat. \R\ Amen.

\benedlabel{2 Bened.} Divínum auxílium máneat semper nobíscum. \R\ Amen.

\benedlabel{3 Bened.} Ad societátem cívium supernórum perdúcat nos Rex Angelórum. \R\ Amen.

\rubric{In Festis Sanctorum erit:}

\benedlabel{1 Bened.} Ille nos benedícat, qui sine fine vivit et regnat. \R\ Amen.

\benedlabel{2 Bened.} Cujus (\textit{vel} Quorum \textit{aut} Quarum) festum cólimus, ipse (\textit{vel} ipsa \textit{aut} ipsi \textit{vel} ipsæ) intercédat (\textit{vel} intercédant) pro nobis ad Dóminum. \R\ Amen.

\benedlabel{3 Bened.} Ad societátem cívium supernórum perdúcat nos Rex Angelórum. \R\ Amen.

\rubric{In Officio de sancta María in Sábbato Absolútio et Benedictiónes, quæ ibídem ponúntur.}

\rubric{Post últimam Lectiónem, in ómnibus Domínicis per Annum minóribus, in Festis cujúsvis ritus, per omnes Octávas, et in Offício sanctæ Maríæ in Sábbato, dícitur Hymnus Ambrosiánus. Et numquam dícitur in feriáli Offício. Quando vero Hymnus prædíctus omíttitur, ejus loco dícitur IX aut III Responsórium.}

\hymnusheading{Hymnus Ambrosianus}

\lettrine[lines=2]{T}{E} Deum laudámus: * te Dóminum confitémur.

Te ætérnum Patrem * omnis terra venerátur.

Tibi omnes Angeli, * tibi Cæli, et univérsæ Potestátes:

Tibi Chérubim et Séraphim * incessábili voce proclámant:

Sanctus, * Sanctus, * Sanctus * Dóminus, Deus Sábaoth.

Pleni sunt cæli et terra * majestátis glóriæ tuæ.

Te gloriósus * Apostolórum chorus,

Te Prophetárum * laudábilis númerus,

Te Mártyrum candidátus * laudat exércitus.

Te per orbem terrárum * sancta confitétur Ecclésia,

Patrem * imménsæ majestátis;

Venerándum tuum verum * et únicum Fílium;

Sanctum quoque * Paráclitum Spíritum.

Tu Rex glóriæ, * Christe.

Tu Patris * sempitérnus es Fílius.

Tu, ad liberándum susceptúrus hóminem, * non horruísti Vírginis úterum.

Tu, devícto mortis acúleo, * aperuísti credéntibus regna cælórum.

Tu ad déxteram Dei sedes, * in glória Patris.

Judex créderis * esse ventúrus.

\rubric{¶ Sequens versus dicitur flexis genibus.}

Te ergo quæsumus, tuis fámulis súbveni, * quos pretióso sánguine redemísti.

Ætérna fac cum Sanctis tuis * in glória numerári.

Salvum fac pópulum tuum, Dómine, * et bénedic hereditáti tuæ.

Et rege eos, * et extólle illos usque in ætérnum.

Per síngulos dies * benedícimus te;

Et laudámus nomen tuum in séculum, * et in séculum séculi.

Dignáre, Dómine, die isto * sine peccáto nos custodíre.

Miserére nostri, Dómine, * miserére nostri.

Fiat misericórdia tua, Dómine, super nos, * quemádmodum sperávimus in te.

In te, Dómine, sperávi: * non confúndar in ætérnum.

\rubric{Dicto Te Deum, aut ultimo Responsorio, statim incipiuntur Laudes a Versu Deus, in adjutórium, ut infra.}

\rubric{In privata vero recitatione Matutinum a Laudibus separari potest; quo in casu, post Hymnum Te Deum, vel post ultimum Responsorium, dicitur:}

\V\ Dóminus vobíscum. \R\ Et cum spíritu tuo.

\rubric{¶ Hic Versus Dóminus vobíscum non dicitur ab eo, qui non est saltem in ordine Diaconatus; sed ejus loco substituitur, etiam in aliis Horis, quoties in antecedentibus Precibus dictus non fuerit:}

\V\ Dómine, exáudi oratiónem meam. \R\ Et clamor meus ad te véniat.

Orémus.

\rubric{Et dicitur Oratio ut ad Laudes.}

\rubric{Post Orationem additur:}

\V\ Dóminus vobíscum. \R\ Et cum spíritu tuo.

\V\ Benedicámus Dómino. \R\ Deo grátias.

\V\ Fidélium ánimæ per misericórdiam Dei requiéscant in pace. \R\ Amen.

\rubric{Deinde dicitur Pater noster totum secreto.}

\rubric{In publica vero recitatione Matutinum in Choro a Laudibus separari non potest.}

\hourheading{Ad Laudes}

\rubric{Si Laudes extra Chorum separentur a Matutino, ante eas dicitur secreto Pater noster et Ave María; secus absolute incipiuntur, ut sequitur:}

\V\ Deus, in adjutórium.

\rubric{Deinde, sub congruentibus Antiphonis, dicuntur Psalmi et Canticum, prouti Officium occurrens requirit.}

\rubric{Repetita vero post ultimum Psalmum Antiphona, dicuntur Capitulum, Hymnus et Versus, prouti Officium occurrens exigit.}

\variantheading{In Officio dominicali}

\rubric{In omnibus Dominicis, extra Octavas Ssmi Corporis Christi ac Sacratissimi Cordis Jesu, Capitulum, Hymnus et Versus ut in Psalterio.}

\variantheading{In Officio feriali}

\rubric{In omnibus Feriis, Capitulum, Hymnus et Versus de Feria currenti, ut in Psalterio.}

\rubric{Recitato Versu, dicitur cum Antiphona convenienti sequens}

\canticumheading{Canticum Zachariæ}{Luc. 1, 68-79}

\lettrine[lines=2]{B}{ENEDICTUS} \textdagger\ Dóminus, Deus Israël, * quia visitávit, et fecit redemptiónem plebis suæ:

Et eréxit cornu salútis nobis: * in domo David, púeri sui.

Sicut locútus est per os sanctórum, * qui a sæculo sunt, Prophetárum ejus:

Salútem ex inimícis nostris, * et de manu ómnium, qui odérunt nos:

Ad faciéndam misericórdiam cum pátribus nostris: * et memorári testaménti sui sancti.

Jusjurándum, quod jurávit ad Abraham, patrem nostrum, * datúrum se nobis:

Ut sine timóre, de manu inimicórum nostrórum liberáti, * serviámus illi.

In sanctitáte, et justítia coram ipso, * ómnibus diébus nostris.

Et tu, puer, Prophéta Altíssimi vocáberis: * præíbis enim ante fáciem Dómini paráre vias ejus:

Ad dandam sciéntiam salútis plebi ejus: * in remissiónem peccatórum eórum:

Per víscera misericórdiæ Dei nostri: * in quibus visitávit nos, óriens ex alto:

Illumináre his, qui in ténebris, et in umbra mortis sedent: * ad dirigéndos pedes nostros in viam pacis.

Glória Patri, et Fílio, * et Spirítui Sancto.

Sicut erat in princípio, et nunc, et semper, * et in sæcula sæculórum. Amen.

\rubric{¶ Hic Versus Glória Patri dicitur in fine omnium Canticorum, nisi aliter adnotetur.}

\rubric{Expleto Cantico ac repetita Antiphona, si Preces non fuerint recitandæ, statim dicitur Versus Dóminus vobíscum cum Oratione et reliquis omnibus usque ad finem Horæ, ut infra.}

\rubric{In omnibus vero Vigiliis, si fiat Officium de Vigilia currenti, post repetitam ad Benedictus Antiphonam dicuntur flexis genibus sequentes Preces, quæ aliis temporibus omittuntur:}

KYRIE, eléison. Christe, eléison. Kýrie, eléison.

\rubric{Pater noster, quod dicitur a solo Hebdomadario totum clara voce, usque ad} \V\ Et ne nos indúcas in tentatiónem. \R\ Sed líbera nos a malo.

\V\ Ego dixi: Dómine, miserére mei. \R\ Sana ánimam meam, quia peccávi tibi. \V\ Convértere, Dómine, úsquequo? \R\ Et deprecábilis esto super servos tuos. \V\ Fiat misericórdia tua, Dómine, super nos. \R\ Quemádmodum sperávimus in te. \V\ Sacerdótes tui induántur justítiam. \R\ Et sancti tui exsúltent.

\V\ Orémus pro beatíssimo Papa nostro N. \R\ Dóminus consérvet eum, et vivíficet eum, et beátum fáciat eum in terra, et non tradat eum in ánimam inimicórum ejus.

\V\ Orémus et pro Antístite nostro N. \R\ Stet et pascat in fortitúdine tua, Dómine, in sublimitáte nóminis tui.

\rubric{¶ Romæ præcedens Versus cum suo Responsorio omittitur; alibi vero, ad litteram N., ab omnibus nomen diœcesani Episcopi exprimatur. Vacante Apostolica vel Episcopali Sede, alteruter vel uterque respectivus Versus cum suo Responsorio prætéritur.}

\V\ Dómine, salvum fac regem. \R\ Et exáudi nos in die, qua invocavérimus te. \V\ Salvum fac pópulum tuum, Dómine, et bénedic hereditáti tuæ. \R\ Et rege eos, et extólle illos usque in ætérnum. \V\ Meménto Congregatiónis tuæ. \R\ Quam possedísti ab inítio. \V\ Fiat pax in virtúte tua. \R\ Et abundántia in túrribus tuis. \V\ Orémus pro benefactóribus nostris. \R\ Retribúere dignáre, Dómine, ómnibus, nobis bona faciéntibus propter nomen tuum, vitam ætérnam. Amen. \V\ Orémus pro fidélibus defúnctis. \R\ Réquiem ætérnam dona eis, Dómine, et lux perpétua lúceat eis. \V\ Requiéscant in pace. \R\ Amen. \V\ Pro frátribus nostris abséntibus. \R\ Salvos fac servos tuos, Deus meus, sperántes in te. \V\ Pro afflíctis et captívis. \R\ Líbera eos, Deus Israël, ex ómnibus tribulatiónibus suis. \V\ Mitte eis, Dómine, auxílium de sancto. \R\ Et de Sion tuére eos. \V\ Dómine, Deus virtútum, convérte nos. \R\ Et osténde fáciem tuam, et salvi érimus. \V\ Exsúrge, Christe, ádjuva nos. \R\ Et líbera nos propter nomen tuum. \V\ Dómine, exáudi oratiónem meam. \R\ Et clamor meus ad te véniat.

\rubric{Absolutis Precibus, vel, si Preces locum non habeant, repetita Antiphona post Benedictus, statim dicitur:}

\V\ Dóminus vobíscum. \R\ Et cum spíritu tuo.

Orémus.

\rubric{Et dicitur Oratio conveniens. Postea fiunt Commemorationes, si occurrant.}

\rubric{Deinde in Officio semiduplici aut simplici tam de Tempore quam de Sanctis, præter omnes dies, in quibus habita sit Commemoratio Duplicis, vel factum fuerit Officium aut Commemoratio cujusvis Octavæ, fit sequens}

\suffragiumheading{Suffragium de omnibus Sanctis}

\antlabel{Ant.} Beáta Dei Génitrix Virgo María, Sanctíque omnes intercédant pro nobis ad Dóminum.

\V\ Mirificávit Dóminus Sanctos suos. \R\ Et exaudívit eos clamántes ad se.

Orémus. \oratiolabel{Oratio}

\lettrine[lines=2]{A}{} \textsc{cunctis} nos, quæsumus, Dómine, mentis et córporis defénde perículis: et, intercedénte beáta et gloriósa semper Vírgine Dei Genitríce María, cum beáto Joseph, beátis Apóstolis tuis Petro et Paulo, atque beáto N. et ómnibus Sanctis, salútem nobis tríbue benígnus et pacem; ut, destrúctis adversitátibus et erróribus univérsis, Ecclésia tua secúra tibi sérviat libertáte. Per eúmdem Dóminum nostrum. \R\ Amen.

\rubric{¶ In hac Oratione, ad litteram N., exprimitur nomen Titularis propriæ Ecclesiæ, dummodo Titulus non sit Persona divina vel Mysterium Domini, aut de ipso ratione Vigiliæ vel Festi factum non fuerit Officium vel Commemoratio, aut illius nomen non sit in eadem Suffragii Oratione expressum; ac nomina sanctorum Angelorum et sancti Joannis Baptistæ, si Titulares fuerint, præponuntur nomini sancti Joseph. In hisce autem omnibus casibus omittuntur verba atque beáto N.}

\rubric{Si vero Officium vel Commemoratio fuerit de beáta María Vírgine, tunc dicitur ut infra; atque insuper, sicubi Titularis juxta superiorem Rubricam ante omnes fuerit nominandus, ad Orationem particula cum omittitur ante nomen sanctorum Petri et Pauli ac præponitur nomini sancti Joseph.}

\antlabel{Ant.} Sancti omnes intercédant pro nobis ad Dóminum.

\V\ Mirificávit Dóminus Sanctos suos. \R\ Et exaudívit eos clamántes ad se.

Orémus. \oratiolabel{Oratio}

\lettrine[lines=2]{A}{} \textsc{cunctis} nos, quæsumus, Dómine, mentis et córporis defénde perículis: et, intercedénte beáto Joseph, cum beátis Apóstolis tuis Petro et Paulo, atque beáto N. et ómnibus Sanctis, salútem nobis tríbue benígnus et pacem; ut, destrúctis adversitátibus et erróribus univérsis, Ecclésia tua secúra tibi sérviat libertáte. Per Dóminum. \R\ Amen.

\rubric{Post ultimam Orationem additur:}

\V\ Dóminus vobíscum. \R\ Et cum spíritu tuo.

\V\ Benedicámus Dómino. \R\ Deo grátias.

\V\ Fidélium ánimæ per misericórdiam Dei requiéscant in pace. \R\ Amen.

\rubric{Deinde, in publica Officii recitatione, si discedendum sit a Choro, itemque in recitatione privata, si tunc terminetur Officium; alioquin, in fine ultimæ Horæ dicitur Pater noster, totum secreto; eoque recitato, subjungitur:}

\V\ Dóminus det nobis suam pacem. \R\ Et vitam ætérnam. Amen.

\rubric{Et immediate dicitur, cum suis Versu et Oratione, finalis beatæ Mariæ Virginis Antiphona, quæ inferius post Completorium assignatur.}

\rubric{Postea concluditur:}

\V\ Divínum auxílium máneat semper nobíscum. \R\ Amen.

% ===== END OF BATCH 1 -- Covers autumnalis.md lines 1-463 (reused
% verbatim from aestiva.tex's Batch 1, per the byte-identity noted
% above): the opening of the Ordinarium Divini Officii (its
% introductory indulgence rubrics and the "Aperi, Dómine" Oratio),
% "Ante singulas Horas" (the Pater noster, Ave María, and Credo
% spelled out in full for the first and only time -- each given a
% lettrine since the source presents them in full capitals as
% complete standalone prayers, unlike their thousands of later
% plain-text abbreviated recurrences throughout the volume), "Ad
% Matutinum" (the Invitatorium Psalm 94, the three-Nocturn apparatus
% with its Absolutios, Benedictiones, and "Pro N Lectione" headings,
% the Officium trium Lectionum variant, and the Te Deum under its
% customary heading "Hymnus Ambrosianus"), and "Ad Laudes" through
% the Suffragium de omnibus Sanctis (in both its Marian-titular and
% Joseph-titular forms) and the closing versicles. This Pars
% Autumnalis Ordinarium matches its Pars Aestiva counterpart exactly
% (both being non-Lenten/non-Paschal volumes) and is notably simpler
% than Pars Verna's (Batch 1 of verna.tex), which carries the extra
% "Tempore Quadragesimæ"/"Tempore Passionis" Invitatorium-Hymnus
% variants and the Paschaltide "Commemoratio de Cruce". Roman
% numerals in the "Pro N Lectione" headings are capitalized (e.g.
% \variantheading{Pro I Lectione}) though the source itself gives them
% lowercase ("Pro i Lectione"), matching the established
% normalization already present throughout hiemalis.tex, verna.tex,
% and aestiva.tex for this specific heading type. The document
% continues with "Ad Primam" (autumnalis.md line 465) =====
\hourheading{Ad Primam}

Pater, Ave et Credo.

\V\ Deus, in adjutórium.

\hymnusheading{Hymnus}

\lettrine[lines=2]{J}{AM} lucis orto sídere,\\
Deum precémur súpplices,\\
Ut in diúrnis áctibus\\
Nos servet a nocéntibus.

Linguam refrénans témperet,\\
Ne litis horror ínsonet:\\
Visum fovéndo cóntegat,\\
Ne vanitátes háuriat.

Sint pura cordis íntima,\\
Absístat et vecórdia:\\
Carnis terat supérbiam\\
Potus cibíque párcitas.

Ut, cum dies abscésserit\\
Noctémque sors redúxerit,\\
Mundi per abstinéntiam\\
Ipsi canámus glóriam.

Deo Patri sit glória,\\
Ejúsque soli Fílio,\\
Cum Spíritu Paráclito\\
Nunc et per omne sæculum. Amen.

\rubric{¶ Conclusio communis, in Hymno præcedenti, et in aliis idem metrum habentibus, semper omittitur, quando specialis in omnibus Horis adhibenda præscribitur; et, si plures Conclusiones propriæ occurrant, sumitur Conclusio Officii currentis, aut secus Officii ipsa die et primo quidem loco inter cetera propriam Conclusionem habentia commemorandi, aut demum de occurrenti Octava communi.}

\rubric{Expleto Hymno, conveniens dicitur Antiphona usque ad Asteriscum *, prouti occurrens Officium requirit.}

\variantheading{In Officio dominicali}

\rubric{In omnibus Dominicis, extra Octavas Ssmi Corporis Christi et Sacrat. Cordis Jesu, Antiphona ut in Psalterio.}

\variantheading{In Officio feriali}

\rubric{In omnibus Feriis, Antiphona de Feria currenti ut in Psalterio.}

\rubric{Post enuntiatam Antiphonam dicuntur tres Psalmi Officii diei congruentes.}

\rubric{¶ Quando autem adhibetur posterius Laudum schema, Psalmus, in priori Laudum ferialium schemate omissus, ad Primam Feriæ currentis post alios Psalmos resumitur.}

\rubric{Repetita post ultimum Psalmum integre Antiphona, conveniens dicitur Capitulum; nimirum:}

\rubric{In omnibus Dominicis, in Officio cujuslibet Festi vel Octavæ ac sanctæ Mariæ in Sabbato:}

\capitulumlabel{Capitulum}{1 Tim. 1, 17}

\lettrine[lines=2]{R}{EGI} sæculórum immortáli et invisíbili, soli Deo honor et glória in sæcula sæculórum. Amen. \R\ Deo grátias.

\rubric{¶ Sic respondetur in fine omnium Capitulorum.}

\rubric{In omnibus autem Feriis et in Vigiliis communibus:}

\capitulumlabel{Capitulum}{Zach. 8, 19}

\lettrine[lines=2]{P}{ACEM} et veritátem dilígite, ait Dóminus omnípotens.

\rubric{Dicto Capitulo, subjungitur Responsorium breve, ut sequitur. Versus tamen Qui sedes ad dexteram Patris semper omittitur, quando specialis Versus adhibendus præscribitur, Officii nempe currentis, aut secus Officii ad Laudes et primo quidem loco inter cetera specialem Versum habentia commemorati, aut demum de occurrenti Octava communi proprius.}

\R\ br. Christe, Fili Dei vivi, * Miserére nobis. \rubric{Et repetitur:} Christe, Fili Dei vivi, miserére nobis. \V\ Qui sedes ad déxteram Patris. \R\ Miserére nobis. \V\ Glória Patri, et Fílio, et Spirítui Sancto. \R\ Christe, Fili Dei vivi, miserére nobis. \V\ Exsúrge, Christe, ádjuva nos. \R\ Et líbera nos propter nomen tuum.

\rubric{¶ Ita dicuntur semper Responsoria brevia, etiam ad Tertiam, Sextam, Nonam et Completorium.}

\rubric{Expleto Responsorio brevi, si Preces non fuerint recitandæ, statim dicitur Versus Dóminus vobíscum cum subsequenti Oratione et reliquis omnibus usque ad finem Horæ ut infra.}

\rubric{In quolibet Officio semiduplici, in Officio sanctæ Mariæ in Sabbato, in Festis simplicibus, atque in Feriis per Annum communibus, post Responsorium breve dicuntur sequentes Preces dominicales; quæ tamen omittuntur omnibus diebus, in quibus ad Laudes habita sit Commemoratio Duplicis, vel factum fuerit Officium aut Commemoratio cujusvis Octavæ:}

KYRIE, eléison. Christe, eléison. Kýrie, eléison.

\rubric{Pater noster secreto usque ad} \V\ Et ne nos indúcas in tentatiónem. \R\ Sed líbera nos a malo.

\rubric{Credo in Deum secreto usque ad} \V\ Carnis resurrectiónem. \R\ Vitam ætérnam. Amen.

\V\ Et ego ad te, Dómine, clamávi. \R\ Et mane orátio mea prævéniet te. \V\ Repleátur os meum laude. \R\ Ut cantem glóriam tuam, tota die magnitúdinem tuam. \V\ Dómine, avérte fáciem tuam a peccátis meis. \R\ Et omnes iniquitátes meas dele. \V\ Cor mundum crea in me, Deus. \R\ Et spíritum rectum ínnova in viscéribus meis. \V\ Ne projícias me a fácie tua. \R\ Et spíritum sanctum tuum ne áuferas a me. \V\ Redde mihi lætítiam salutáris tui. \R\ Et spíritu principáli confírma me. \V\ Adjutórium nostrum in nómine Dómini. \R\ Qui fecit cælum et terram.

\rubric{Deinde Hebdomadarius facit Confessionem cum Absolutione, eamque Chorus repetit, ut infra post feriales Preces.}

\rubric{In omnibus vero Vigiliis, si ad Laudes recitatæ fuerint preces, post Responsorium breve dicuntur flexis genibus sequentes Preces feriales:}

KYRIE, eléison. Christe, eléison. Kýrie, eléison.

\rubric{Pater noster secreto usque ad} \V\ Et ne nos indúcas in tentatiónem. \R\ Sed líbera nos a malo.

\rubric{Credo in Deum secreto usque ad} \V\ Carnis resurrectiónem. \R\ Vitam ætérnam. Amen.

\V\ Et ego ad te, Dómine, clamávi. \R\ Et mane orátio mea prævéniet te. \V\ Repleátur os meum laude. \R\ Ut cantem glóriam tuam, tota die magnitúdinem tuam. \V\ Dómine, avérte fáciem tuam a peccátis meis. \R\ Et omnes iniquitátes meas dele. \V\ Cor mundum crea in me, Deus. \R\ Et spíritum rectum ínnova in viscéribus meis. \V\ Ne projícias me a fácie tua. \R\ Et spíritum sanctum tuum ne áuferas a me. \V\ Redde mihi lætítiam salutáris tui. \R\ Et spíritu principáli confírma me. \V\ Eripe me, Dómine, ab hómine malo. \R\ A viro iníquo éripe me. \V\ Eripe me de inimícis meis, Deus meus. \R\ Et ab insurgéntibus in me líbera me. \V\ Eripe me de operántibus iniquitátem. \R\ Et de viris sánguinum salva me. \V\ Sic psalmum dicam nómini tuo in sæculum sæculi. \R\ Ut reddam vota mea de die in diem. \V\ Exáudi nos, Deus, salutáris noster. \R\ Spes ómnium fínium terræ, et in mari longe. \V\ Deus, in adjutórium meum inténde. \R\ Dómine, ad adjuvándum me festína. \V\ Sanctus Deus, Sanctus fortis, Sanctus immortális. \R\ Miserére nobis. \V\ Bénedic, ánima mea, Dómino. \R\ Et ómnia, quæ intra me sunt, nómini sancto ejus. \V\ Bénedic, ánima mea, Dómino. \R\ Et noli oblivísci omnes retributiónes ejus. \V\ Qui propitiátur ómnibus iniquitátibus tuis. \R\ Qui sanat omnes infirmitátes tuas. \V\ Qui rédimit de intéritu vitam tuam. \R\ Qui corónat te in misericórdia et miseratiónibus. \V\ Qui replet in bonis desidérium tuum. \R\ Renovábitur ut áquilæ juvéntus tua. \V\ Adjutórium nostrum in nómine Dómini. \R\ Qui fecit cælum et terram.

\rubric{Deinde Hebdomadarius facit Confessionem:}

\lettrine[lines=2]{C}{ONFITEOR} Deo omnipoténti, beátæ Maríæ semper Vírgini, beáto Michaéli Archángelo, beáto Joánni Baptístæ, sanctis Apóstolis Petro et Paulo, ómnibus Sanctis, et vobis, fratres, quia peccávi nimis cogitatióne, verbo et ópere: mea culpa, mea culpa, mea máxima culpa. Ideo precor beátam Maríam semper Vírginem, beátum Michaélem Archángelum, beátum Joánnem Baptístam, sanctos Apóstolos Petrum et Paulum, omnes Sanctos, et vos, fratres, oráre pro me ad Dóminum, Deum nostrum.

\rubric{Chorus respondet:}

\lettrine[lines=2]{M}{ISEREATUR} tui omnípotens Deus, et, dimíssis peccátis tuis, perdúcat te ad vitam ætérnam. \R\ Amen.

\rubric{Deinde repetit Confessionem, et ubi dicitur: vobis, fratres, et vos, fratres, dicat: tibi, pater, et te, pater.}

\rubric{Facta Confessione a Choro, Hebdomadarius dicit:}

\lettrine[lines=2]{M}{ISEREATUR} vestri omnípotens Deus, et, dimíssis peccátis vestris, perdúcat vos ad vitam ætérnam. \R\ Amen.

\lettrine[lines=2]{I}{NDULGENTIAM}, absolutiónem, et remissiónem peccatórum nostrórum tríbuat nobis omnípotens et miséricors Dóminus. \R\ Amen.

\rubric{Extra Chorum, si unus vel duo tantum recitent Officium et in Choro Monialium, semel tantum ac simul ab omnibus ita fit Confessio:}

\lettrine[lines=2]{C}{ONFITEOR} Deo omnipoténti, beátæ Maríæ semper Vírgini, beáto Michaéli Archángelo, beáto Joánni Baptístæ, sanctis Apóstolis Petro et Paulo, et ómnibus Sanctis, quia peccávi nimis cogitatióne, verbo et ópere: mea culpa, mea culpa, mea máxima culpa. Ideo precor beátam Maríam semper Vírginem, beátum Michaélem Archángelum, beátum Joánnem Baptístam, sanctos Apóstolos Petrum et Paulum, et omnes Sanctos, oráre pro me ad Dóminum, Deum nostrum.

\rubric{Deinde dicitur:}

\lettrine[lines=2]{M}{ISEREATUR} nostri omnípotens Deus, et, dimíssis peccátis nostris, perdúcat nos ad vitam ætérnam. Amen.

\lettrine[lines=2]{I}{NDULGENTIAM}, absolutiónem, et remissiónem peccatórum nostrórum tríbuat nobis omnípotens et miséricors Dóminus. Amen.

\rubric{Facta Absolutione, subjungitur:}

\V\ Dignáre, Dómine, die isto. \R\ Sine peccáto nos custodíre. \V\ Miserére nostri, Dómine. \R\ Miserére nostri. \V\ Fiat misericórdia tua, Dómine, super nos. \R\ Quemádmodum sperávimus in te. \V\ Dómine, exáudi oratiónem meam. \R\ Et clamor meus ad te véniat.

\rubric{Absolutis Precibus, vel, si Preces locum non habeant, post Responsorium breve statim dicitur:}

\V\ Dóminus vobíscum. \R\ Et cum spíritu tuo.

Orémus. \oratiolabel{Oratio}

\lettrine[lines=2]{D}{OMINE}, Deus omnípotens, qui ad princípium hujus diéi nos perveníre fecísti: tua nos hódie salva virtúte; ut in hac die ad nullum declinémus peccátum, sed semper ad tuam justítiam faciéndam nostra procédant elóquia, dirigántur cogitatiónes et ópera. Per Dóminum. \R\ Amen.

\V\ Dóminus vobíscum. \R\ Et cum spíritu tuo. \V\ Benedicámus Dómino. \R\ Deo grátias.

\rubric{Deinde in Choro legitur Martyrologium, quod laudabiliter fit etiam extra Chorum. Postea Hebdomadarius dicit:}

\V\ Pretiósa in conspéctu Dómini. \R\ Mors Sanctórum ejus.

\rubric{Et absolute, sine} Orémus:

\oratiolabel{Oratio}

\lettrine[lines=2]{S}{ANCTA} María, et omnes Sancti intercédant pro nobis ad Dóminum, ut nos mereámur ab eo adjuvári et salvári, qui vivit et regnat in sæcula sæculórum. \R\ Amen.

\V\ Deus, in adjutórium meum inténde. \R\ Dómine, ad adjuvándum me festína. \rubric{Et dicitur ter; ultimo additur:} Glória Patri, et Fílio, et Spirítui Sancto. Sicut erat in princípio, et nunc, et semper, et in sæcula sæculórum. Amen.

Kýrie, eléison. Christe, eléison. Kýrie, eléison.

Pater noster \rubric{secreto usque ad} \V\ Et ne nos indúcas in tentatiónem. \R\ Sed líbera nos a malo.

\V\ Réspice in servos tuos, Dómine, et in ópera tua, et dírige fílios eórum. \R\ Et sit splendor Dómini, Dei nostri, super nos, et ópera mánuum nostrárum dírige super nos, et opus mánuum nostrárum dírige. \V\ Glória Patri, et Fílio, et Spirítui Sancto. \R\ Sicut erat in princípio, et nunc, et semper, et in sæcula sæculórum. Amen.

Orémus. \oratiolabel{Oratio}

\lettrine[lines=2]{D}{IRIGERE} et sanctificáre, régere et gubernáre dignáre, Dómine Deus, Rex cæli et terræ, hódie corda et córpora nostra, sensus, sermónes et actus nostros in lege tua, et in opéribus mandatórum tuórum: ut hic, et in ætérnum, te auxiliánte, salvi et líberi esse mereámur, Salvátor mundi: Qui vivis et regnas in sæcula sæculórum. \R\ Amen.

\rubric{Deinde dicitur Lectio brevis, ad absolutionem Capituli, prouti Officium occurrens requirit.}

\V\ Jube, domne, benedícere. \benedlabel{Bened.} Dies et actus nostros in sua pace dispónat Dóminus omnípotens. \R\ Amen.

\variantheading{In Officio dominicali et feriali}

\rubric{A Feria II usque ad Feriam IV post Festum sanctissimæ Trinitatis, et a Sabbato post Octavam Sacratissimi Cordis Jesu et deinceps.}

\lectioheading{Lectio brevis}{2 Thess. 3, 5}

\lettrine[lines=2]{D}{OMINUS} autem dírigat corda et córpora nostra in caritáte Dei, et patiéntia Christi. Tu autem, Dómine, miserére nobis. \R\ Deo grátias.

\rubric{¶ Ita semper concluduntur Lectiones breves ad Primam, et ita in earum fine respondetur.}

\rubric{In Officio autem cujuslibet Festi vel Octavæ, atque S. Mariæ in Sabbato; itemque in Dominicis infra Octavas Ssmi Corporis Christi ac Sacrat. Cordis Jesu, si de eis fiat Officium:}

\rubric{Pro Lectione brevi sumitur Capitulum, quod ad Nonam ponitur ipsius Officii currentis, et in fine dicitur: Tu autem, Dómine, ut supra.}

\rubric{Post Lectionem brevem dicitur:}

\V\ Adjutórium nostrum in nómine Dómini. \R\ Qui fecit cælum et terram.

\V\ Benedícite. \R\ Deus.

\benedlabel{Bened.} Dóminus nos benedícat, et ab omni malo deféndat, et ad vitam perdúcat ætérnam. Et fidélium ánimæ per misericórdiam Dei requiéscant in pace. \R\ Amen.

\rubric{Deinde dicitur tantum Pater noster secreto, nisi sequatur alia Hora.}

\hourheading{Ad Tertiam}

Pater noster et Ave.

\V\ Deus, in adjutórium.

\hymnusheading{Hymnus}

\lettrine[lines=2]{N}{UNC}, Sancte, nobis, Spíritus,\\
Unum Patri cum Fílio,\\
Dignáre promptus íngeri\\
Nostro refúsus péctori.

Os, lingua, mens, sensus, vigor,\\
Confessiónem pérsonent,\\
Flamméscat igne cáritas,\\
Accéndat ardor próximos.

Præsta, Pater piíssime,\\
Patríque compar Unice,\\
Cum Spíritu Paráclito\\
Regnans per omne sæculum. Amen.

\rubric{Expleto Hymno, conveniens dicitur Antiphona usque ad Asteriscum *, et tres Psalmi, prouti Officium occurrens requirit; atque Psalmis absolutis, integra repetitur Antiphona.}

\variantheading{In Officio dominicali}

\rubric{In omnibus Dominicis, extra Octavas Ssmi Corporis Christi ac Sacrat. Cordis Jesu, Aña ut in Psalterio, et dicitur sequens}

\capitulumlabel{Capitulum}{1 Joann. 4, 16}

\lettrine[lines=2]{D}{EUS} cáritas est: et qui manet in caritáte, in Deo manet, et Deus in eo.

\R\ br. Inclína cor meum, Deus, * In testimónia tua. Inclína. \V\ Avérte óculos meos, ne vídeant vanitátem: in via tua vivífica me. In testimónia tua. Glória Patri. Inclína.

\V\ Ego dixi: Dómine, miserére mei. \R\ Sana ánimam meam, quia peccávi tibi.

\variantheading{In Officio feriali}

\rubric{In omnibus Feriis, Antiphona de Feria currenti, ut in Psalterio; et dicitur quotidie sequens}

\capitulumlabel{Capitulum}{Jer. 17, 14}

\lettrine[lines=2]{S}{ANA} me, Dómine, et sanábor: salvum me fac, et salvus ero: quóniam laus mea tu es.

\R\ br. Sana ánimam meam, * Quia peccávi tibi. Sana. \V\ Ego dixi: Dómine, miserére mei. Quia. Glória Patri. Sana.

\V\ Adjútor meus esto, ne derelínquas me. \R\ Neque despícias me, Deus, salutáris meus.

\rubric{Expleto Responsorio brevi, si Preces non fuerint recitandæ, statim dicitur Versus Dóminus vobíscum cum Oratione et reliquis omnibus usque ad finem Horæ, ut infra.}

\rubric{Quando vero ad Laudes recitatæ sunt Preces, post Responsorium breve dicuntur sequentes flexis genibus:}

KYRIE, eléison. Christe, eléison. Kýrie, eléison.

Pater noster \rubric{secreto usque ad} \V\ Et ne nos indúcas in tentatiónem. \R\ Sed líbera nos a malo. \V\ Dómine, Deus virtútum, convérte nos. \R\ Et osténde fáciem tuam, et salvi érimus. \V\ Exsúrge, Christe, ádjuva nos. \R\ Et líbera nos propter nomen tuum. \V\ Dómine, exáudi oratiónem meam. \R\ Et clamor meus ad te véniat.

\rubric{Absolutis Precibus, vel, si Preces locum non habeant, post Responsorium breve statim dicitur:}

\V\ Dóminus vobíscum. \R\ Et cum spíritu tuo.

Orémus.

\rubric{Et dicitur Oratio conveniens.}

\rubric{Post Orationem additur:}

\V\ Dóminus vobíscum. \R\ Et cum spíritu tuo.

\V\ Benedicámus Dómino. \R\ Deo grátias.

\V\ Fidélium ánimæ per misericórdiam Dei requiéscant in pace. \R\ Amen.

\rubric{Deinde dicitur tantum Pater noster secreto, nisi sequatur alia Hora.}

\hourheading{Ad Sextam}

Pater noster et Ave.

\V\ Deus, in adjutórium.

\hymnusheading{Hymnus}

\lettrine[lines=2]{R}{ECTOR} potens, verax Deus,\\
Qui témperas rerum vices,\\
Splendóre mane illúminas,\\
Et ígnibus merídiem:

Exstíngue flammas lítium,\\
Aufer calórem nóxium,\\
Confer salútem córporum,\\
Verámque pacem córdium.

Præsta, Pater piíssime,\\
Patríque compar Unice,\\
Cum Spíritu Paráclito\\
Regnans per omne sæculum. Amen.

\rubric{Expleto Hymno, conveniens dicitur Antiphona usque ad Asteriscum *, et tres Psalmi prouti Officium occurrens requirit; atque, Psalmis absolutis, integra repetitur Antiphona.}

\variantheading{In Officio dominicali}

\rubric{In omnibus Dominicis, extra Octavas Ssmi Corporis Christi ac Sacrat. Cordis Jesu, Aña ut in Psalterio, et dicitur sequens}

\capitulumlabel{Capitulum}{Gal. 6, 2}

\lettrine[lines=2]{A}{LTER} altérius ónera portáte, et sic adimplébitis legem Christi.

\R\ br. In ætérnum, Dómine, * Pérmanet verbum tuum. In ætérnum. \V\ In sæculum sæculi véritas tua. Pérmanet. Glória Patri. In ætérnum.

\V\ Dóminus regit me, et nihil mihi déerit. \R\ In loco páscuæ ibi me collocávit.

\variantheading{In Officio feriali}

\rubric{In omnibus Feriis, Antiphona de Feria currenti, ut in Psalterio; et dicitur quotidie sequens}

\capitulumlabel{Capitulum}{Rom. 13, 8}

\lettrine[lines=2]{N}{EMINI} quidquam debeátis, nisi ut ínvicem diligátis: qui enim díligit próximum, legem implévit.

\R\ br. Benedícam Dóminum * In omni témpore. Benedícam. \V\ Semper laus ejus in ore meo. In. Glória Patri. Benedícam.

\V\ Dóminus regit me, et nihil mihi déerit. \R\ In loco páscuæ ibi me collocávit.

\rubric{Expleto Responsorio brevi, si Preces non fuerint recitandæ, statim dicitur Versus Dóminus vobíscum cum Oratione et reliquis omnibus usque ad finem Horæ, ut infra.}

\rubric{Quando vero ad Laudes recitatæ sunt Preces, post Responsorium breve dicuntur sequentes flexis genibus:}

KYRIE, eléison. Christe, eléison. Kýrie, eléison.

Pater noster \rubric{secreto usque ad} \V\ Et ne nos indúcas in tentatiónem. \R\ Sed líbera nos a malo. \V\ Dómine, Deus virtútum, convérte nos. \R\ Et osténde fáciem tuam, et salvi érimus. \V\ Exsúrge, Christe, ádjuva nos. \R\ Et líbera nos propter nomen tuum. \V\ Dómine, exáudi oratiónem meam. \R\ Et clamor meus ad te véniat.

\rubric{Absolutis Precibus, vel, si Preces locum non habeant, post Responsorium breve statim dicitur:}

\V\ Dóminus vobíscum. \R\ Et cum spíritu tuo.

Orémus.

\rubric{Et dicitur Oratio conveniens.}

\rubric{Post Orationem additur:}

\V\ Dóminus vobíscum. \R\ Et cum spíritu tuo.

\V\ Benedicámus Dómino. \R\ Deo grátias.

\V\ Fidélium ánimæ per misericórdiam Dei requiéscant in pace. \R\ Amen.

\rubric{Deinde dicitur tantum Pater noster secreto, nisi sequatur alia Hora.}

\hourheading{Ad Nonam}

Pater noster et Ave.

\V\ Deus, in adjutórium.

\hymnusheading{Hymnus}

\lettrine[lines=2]{R}{ERUM}, Deus, tenax vigor,\\
Immótus in te pérmanens,\\
Lucis diúrnæ témpora\\
Successíbus detérminans:

Largíre lumen véspere,\\
Quo vita nusquam décidat,\\
Sed præmium mortis sacræ\\
Perénnis instet glória.

Præsta, Pater piíssime,\\
Patríque compar Unice,\\
Cum Spíritu Paráclito\\
Regnans per omne sæculum. Amen.

\rubric{Expleto Hymno, conveniens dicitur Antiphona usque ad Asteriscum *, et tres Psalmi, prouti Officium occurrens requirit; atque, Psalmis absolutis, integra repetitur Antiphona.}

\variantheading{In Officio dominicali}

\rubric{In omnibus Dominicis, extra Octavas Ssmi Corporis Christi ac Sacrat. Cordis Jesu, Aña ut in Psalterio, et dicitur sequens}

\capitulumlabel{Capitulum}{1 Cor. 6, 20}

\lettrine[lines=2]{E}{MPTI} enim estis prétio magno. Glorificáte et portáte Deum in córpore vestro.

\R\ br. Clamávi in toto corde meo: * Exáudi me, Dómine. Clamávi. \V\ Justificatiónes tuas requíram. Exáudi. Glória Patri. Clamávi.

\V\ Ab occúltis meis munda me, Dómine. \R\ Et ab aliénis parce servo tuo.

\variantheading{In Officio feriali}

\rubric{In omnibus Feriis, Antiphona de Feria currenti, ut in Psalterio; et dicitur quotidie sequens}

\capitulumlabel{Capitulum}{1 Petr. 1, 17-19}

\lettrine[lines=2]{I}{N} timóre incolátus vestri témpore conversámini: sciéntes quod non corruptibílibus auro vel argénto redémpti estis, sed pretióso sánguine quasi Agni immaculáti Christi.

\R\ br. Rédime me, Dómine, * Et miserére mei. Rédime. \V\ Pes enim meus stetit in via recta. Et. Glória Patri. Rédime.

\V\ Ab occúltis meis munda me, Dómine. \R\ Et ab aliénis parce servo tuo.

\rubric{Expleto Responsorio brevi, si Preces non fuerint recitandæ, statim dicitur Versus Dóminus vobíscum cum Oratione et reliquis omnibus usque ad finem Horæ, ut infra.}

\rubric{Quando vero ad Laudes recitatæ sunt Preces, post Responsorium breve dicuntur sequentes flexis genibus:}

KYRIE, eléison. Christe, eléison. Kýrie, eléison.

Pater noster \rubric{secreto usque ad} \V\ Et ne nos indúcas in tentatiónem. \R\ Sed líbera nos a malo. \V\ Dómine, Deus virtútum, convérte nos. \R\ Et osténde fáciem tuam, et salvi érimus. \V\ Exsúrge, Christe, ádjuva nos. \R\ Et líbera nos propter nomen tuum. \V\ Dómine, exáudi oratiónem meam. \R\ Et clamor meus ad te véniat.

\rubric{Absolutis Precibus, vel, si Preces locum non habeant, post Responsorium breve statim dicitur:}

\V\ Dóminus vobíscum. \R\ Et cum spíritu tuo.

Orémus.

\rubric{Et dicitur Oratio conveniens.}

\rubric{Post Orationem additur:}

\V\ Dóminus vobíscum. \R\ Et cum spíritu tuo.

\V\ Benedicámus Dómino. \R\ Deo grátias.

\V\ Fidélium ánimæ per misericórdiam Dei requiéscant in pace. \R\ Amen.

\rubric{Deinde dicitur tantum Pater noster secreto, nisi sequatur alia Hora.}

% ===== END OF BATCH 2 -- Covers autumnalis.md lines 465-874 (reused
% verbatim from aestiva.tex's Batch 2, per the byte-identity noted
% above), completing "Ad Primam" (hymn "Jam lucis orto sidere," the
% full Confiteor/Misereatur/Indulgentiam apparatus, the Martyrologium
% cue, the "Dirigere et sanctificare" Oratio, and the Lectio brevis,
% here given in its single simplified form -- "In Officio dominicali
% et feriali" -- rather than the three-way Lenten/Paschal-season
% branching found in verna.tex's Prime, since this volume's Lectio
% brevis never varies by season), "Ad Tertiam," "Ad Sextam," and "Ad
% Nonam," each with its proper hymn and its Sunday/weekday
% Capitulum-and-brief-responsory pair. Two spellings of the same word
% recur within this single batch exactly as the source alternates
% them: "sæculum"/"sæculi" at Sexta's responsory (line 758) versus
% the different spelling pattern seen in the Te Deum of Batch 1 --
% each transcribed as its own occurrence gives it, per the standing
% rule against silent normalization. The document continues with
% "Ad Vesperas" (autumnalis.md line 875) =====
\V\ Dóminus vobíscum. \R\ Et cum spíritu tuo.

Orémus.

\rubric{Et dicitur Oratio conveniens.}

\rubric{Postea fiunt Commemorationes, si occurrant.}

\rubric{Deinde, in Officio semiduplici aut simplici tam de Tempore quam de Sanctis, præter omnes dies, in quibus habita sit Commemoratio Duplicis, etiam concurrentis, vel factum fuerit Officium aut Commemoratio cujusvis Octavæ, fit sequens}

\suffragiumheading{Suffragium de omnibus Sanctis}

\antlabel{Ant.} Beáta Dei Génitrix Virgo María, Sanctíque omnes intercédant pro nobis ad Dóminum.

\V\ Mirificávit Dóminus Sanctos suos. \R\ Et exaudívit eos clamántes ad se.

Orémus. \oratiolabel{Oratio}

\lettrine[lines=2]{A}{} \textsc{cunctis} nos, quæsumus, Dómine, mentis et córporis defénde perículis: et, intercedénte beáta et gloriósa semper Vírgine Dei Genitríce María, cum beáto Joseph, beátis Apóstolis tuis Petro et Paulo, atque beáto N. et ómnibus Sanctis, salútem nobis tríbue benígnus et pacem; ut, destrúctis adversitátibus et erróribus univérsis, Ecclésia tua secúra tibi sérviat libertáte. Per eúmdem Dóminum. \R\ Amen.

\rubric{¶ In hac Oratione, ad litteram N., exprimitur nomen Titularis propriæ Ecclesiæ, dummodo Titulus non sit Persona divina vel Mysterium Domini, aut de ipso ratione Festi factum non fuerit Officium vel Commemoratio, aut illius nomen non sit in eadem Suffragii Oratione expressum; ac nomina sanctorum Angelorum et sancti Joannis Baptistæ, si Titulares fuerint, præponuntur nomini sancti Joseph. In hisce autem omnibus casibus omittuntur verba atque beáto N.}

\rubric{Si vero Officium vel Commemoratio fuerit de beáta María Vírgine, tunc dicitur ut infra; atque insuper, sicubi Titularis juxta superiorem Rubricam ante omnes fuerit nominandus, ad Orationem particula cum omittitur ante nomen sanctorum Petri et Pauli ac præponitur nomini sancti Joseph.}

\antlabel{Ant.} Sancti omnes intercédant pro nobis ad Dóminum.

\V\ Mirificávit Dóminus Sanctos suos. \R\ Et exaudívit eos clamántes ad se.

Orémus. \oratiolabel{Oratio}

\lettrine[lines=2]{A}{} \textsc{cunctis} nos, quæsumus, Dómine, mentis et córporis defénde perículis: et, intercedénte beáto Joseph, cum beátis Apóstolis tuis Petro et Paulo, atque beáto N. et ómnibus Sanctis, salútem nobis tríbue benígnus et pacem; ut, destrúctis adversitátibus et erróribus univérsis, Ecclésia tua secúra tibi sérviat libertáte. Per Dóminum.

\rubric{Post ultimam Orationem additur:}

\V\ Dóminus vobíscum. \R\ Et cum spíritu tuo.

\V\ Benedicámus Dómino. \R\ Deo grátias.

\V\ Fidélium ánimæ per misericórdiam Dei requiéscant in pace. \R\ Amen.

\rubric{Si post Vesperas immediate sequatur Completorium, dicto Versu Fidelium animæ, statim incipitur Versus Jube, domne, benedicere, ut infra ad Completorium; secus autem, si tunc terminetur Officium, dicitur tantum Pater noster secreto.}

\hourheading{Ad Completorium}

\rubric{Ante Completorium non dicitur Pater noster et Ave María, sicut in aliis Horis, sed absolute Lector incipit:}

\V\ Jube, domne, benedícere. \benedlabel{Bened.} Noctem quiétam, et finem perféctum concédat nobis Dóminus omnípotens. \R\ Amen.

\lectioheading{Lectio brevis}{1 Petr. 5, 8-9}

\lettrine[lines=2]{F}{RATRES}: Sóbrii estóte, et vigiláte: quia adversárius vester diábolus tamquam leo rúgiens círcuit, quærens quem dévoret: cui resístite fortes in fide. Tu autem, Dómine, miserére nobis. \R\ Deo grátias.

\V\ Adjutórium nostrum in nómine Dómini. \R\ Qui fecit cælum et terram.

Pater noster \rubric{quod dicitur totum secreto.}

\rubric{Deinde Hebdomadarius facit Confessionem:}

\lettrine[lines=2]{C}{ONFITEOR} Deo omnipoténti, beátæ Maríæ semper Vírgini, beáto Michaéli Archángelo, beáto Joánni Baptístæ, sanctis Apóstolis Petro et Paulo, ómnibus Sanctis, et vobis, fratres, quia peccávi nimis cogitatióne, verbo et ópere: mea culpa, mea culpa, mea máxima culpa. Ideo precor beátam Maríam semper Vírginem, beátum Michaélem Archángelum, beátum Joánnem Baptístam, sanctos Apóstolos Petrum et Paulum, omnes Sanctos, et vos, fratres, oráre pro me ad Dóminum, Deum nostrum.

\rubric{Chorus respondet:}

\lettrine[lines=2]{M}{ISEREATUR} tui omnípotens Deus, et, dimíssis peccátis tuis, perdúcat te ad vitam ætérnam. \R\ Amen.

\rubric{Deinde repetit Confessionem, et ubi dicitur: vobis, fratres, et vos, fratres, dicat: tibi, pater, et te, pater.}

\rubric{Facta Confessione a Choro, Hebdomadarius dicit:}

\lettrine[lines=2]{M}{ISEREATUR} vestri omnípotens Deus, et, dimíssis peccátis vestris, perdúcat vos ad vitam ætérnam. \R\ Amen.

\lettrine[lines=2]{I}{NDULGENTIAM}, absolutiónem, et remissiónem peccatórum nostrórum tríbuat nobis omnípotens et miséricors Dóminus. \R\ Amen.

\rubric{Extra Chorum, si unus vel duo tantum recitent Officium, et in Choro Monialium, semel tantum ac simul ab omnibus ita fit Confessio:}

\lettrine[lines=2]{C}{ONFITEOR} Deo omnipoténti, beátæ Maríæ semper Vírgini, beáto Michaéli Archángelo, beáto Joánni Baptístæ, sanctis Apóstolis Petro et Paulo, et ómnibus Sanctis, quia peccávi nimis cogitatióne, verbo et ópere: mea culpa, mea culpa, mea máxima culpa. Ideo precor beátam Maríam semper Vírginem, beátum Michaélem Archángelum, beátum Joánnem Baptístam, sanctos Apóstolos Petrum et Paulum, et omnes Sanctos, oráre pro me ad Dóminum, Deum nostrum.

\rubric{Deinde dicitur:}

\lettrine[lines=2]{M}{ISEREATUR} nostri omnípotens Deus, et, dimíssis peccátis nostris, perdúcat nos ad vitam ætérnam. Amen.

\lettrine[lines=2]{I}{NDULGENTIAM}, absolutiónem, et remissiónem peccatórum nostrórum tríbuat nobis omnípotens et miséricors Dóminus. Amen.

\rubric{Facta Absolutione, dicitur:}

\V\ Convérte nos, Deus, salutáris noster. \R\ Et avérte iram tuam a nobis. \V\ Deus, in adjutórium.

\rubric{Deinde, inchoata convenienti Antiphona, dicuntur tres Psalmi, prouti Officium occurrens requirit.}

\rubric{Repetita integre post ultimum Psalmum Antiphona, dicitur sequens}

\hymnusheading{Hymnus}

\lettrine[lines=2]{T}{E} lucis ante términum,\\
Rerum Creátor, póscimus,\\
Ut pro tua cleméntia\\
Sis præsul et custódia.

Procul recédant sómnia,\\
Et nóctium phantásmata;\\
Hostémque nostrum cómprime,\\
Ne polluántur córpora.

Præsta, Pater piíssime,\\
Patríque compar Unice,\\
Cum Spíritu Paráclito\\
Regnans per omne sæculum. Amen.

\rubric{Absoluto Hymno, dicitur sequens Capitulum et Responsorium breve:}

\capitulumlabel{Capitulum}{Jer. 14, 9}

\lettrine[lines=2]{T}{U} autem in nobis es, Dómine, et nomen sanctum tuum invocátum est super nos: ne derelínquas nos, Dómine, Deus noster.

\R\ br. In manus tuas, Dómine, * Comméndo spíritum meum. In manus tuas, Dómine, comméndo spíritum meum. \V\ Redemísti nos, Dómine, Deus veritátis. Comméndo spíritum meum. Glória Patri, et Fílio, et Spirítui Sancto. In manus tuas, Dómine, comméndo spíritum meum.

\V\ Custódi nos, Dómine, ut pupíllam óculi. \R\ Sub umbra alárum tuárum prótege nos.

\rubric{Postea ad Canticum dicitur sequens}

\antlabel{Ant.} Salva nos.

\canticumheading{Canticum Simeonis}{Luc. 2, 29-32}

\lettrine[lines=2]{N}{UNC} dimíttis servum tuum, Dómine, * secúndum verbum tuum in pace:

Quia vidérunt óculi mei * salutáre tuum.

Quod parásti * ante fáciem ómnium populórum,

Lumen ad revelatiónem géntium, * et glóriam plebis tuæ Israël.

Glória Patri.

\antlabel{Ant.} Salva nos, Dómine, vigilántes, custódi nos dormiéntes; ut vigilémus cum Christo, et requiescámus in pace.

\rubric{Deinde, si Preces non fuerint recitandæ, statim dicitur Versus Dóminus vobíscum cum subsequenti Oratione et reliquis omnibus usque ad finem Completorii, ut infra.}

\rubric{In quolibet Officio semiduplici, in Officio sanctæ Mariæ in Sabbato, in Festis simplicibus, atque in Feriis per Annum communibus, post repetitam ad Nunc dimittis Antiphonam dicuntur sequentes Preces, quæ tamen omittuntur omnibus diebus, in quibus ad Vesperas habita sit Commemoratio Duplicis, vel factum fuerit Officium aut Commemoratio cujusvis Octavæ.}

KYRIE, eléison. Christe, eléison. Kýrie, eléison.

Pater noster \rubric{secreto usque ad} \V\ Et ne nos indúcas in tentatiónem. \R\ Sed líbera nos a malo.

Credo in Deum \rubric{secreto usque ad} \V\ Carnis resurrectiónem. \R\ Vitam ætérnam. Amen.

\V\ Benedíctus es, Dómine, Deus patrum nostrórum. \R\ Et laudábilis et gloriósus in sæcula. \V\ Benedicámus Patrem et Fílium cum Sancto Spíritu. \R\ Laudémus, et superexaltémus eum in sæcula. \V\ Benedíctus es, Dómine, in firmaménto cæli. \R\ Et laudábilis, et gloriósus, et superexaltátus in sæcula. \V\ Benedícat et custódiat nos omnípotens et miséricors Dóminus. \R\ Amen. \V\ Dignáre, Dómine, nocte ista. \R\ Sine peccáto nos custodíre. \V\ Miserére nostri, Dómine. \R\ Miserére nostri. \V\ Fiat misericórdia tua, Dómine, super nos. \R\ Quemádmodum sperávimus in te. \V\ Dómine, exáudi oratiónem meam. \R\ Et clamor meus ad te véniat.

\rubric{Absolutis Precibus, vel, si Preces locum non habeant, repetita Antiphona post Nunc dimittis, statim dicitur:}

\V\ Dóminus vobíscum. \R\ Et cum spíritu tuo.

Orémus. \oratiolabel{Oratio}

\lettrine[lines=2]{V}{ISITA}, quæsumus, Dómine, habitatiónem istam, et omnes insídias inimíci ab ea longe repélle: Angeli tui sancti hábitent in ea, qui nos in pace custódiant; et benedíctio tua sit super nos semper. Per Dóminum nostrum. \R\ Amen.

\rubric{Post Orationem additur:}

\V\ Dóminus vobíscum. \R\ Et cum spíritu tuo.

\V\ Benedicámus Dómino. \R\ Deo grátias.

\benedictiolabel{Benedictio} Benedícat et custódiat nos omnípotens et miséricors Dóminus, Pater, et Fílius, et Spíritus Sanctus. \R\ Amen.

\rubric{Et non dicitur Versus Fidelium animæ, sed immediate subjungitur finalis beatæ Mariæ Virginis Antiphona, quæ infra assignatur.}

\rubric{Postea subjungitur:}

\V\ Divínum auxílium máneat semper nobíscum. \R\ Amen.

\rubric{Deinde dicitur secreto Pater noster, Ave María et Credo.}

\rubric{Si post Completorium immediate sequatur Matutinum, tunc bis dicitur Pater noster, Ave María et Credo, pro fine videlicet Officii diei præcedentis, et pro initio Officii diei sequentis.}

\sectionheading{Post singulas Horas}

\rubric{In fine Laudum, vel, si post Laudes immediate subsequatur alia Hora, in fine ultimæ Horæ, sicuti et semper in fine Completorii, quamvis immediate sequatur Matutinum diei sequentis, tam in publica quam in privata divini Officii recitatione; atque in publica insuper recitatione post alias quascumque Horas, si tunc terminetur Officium et discedatur a Choro, dicitur sequens Antiphona finalis beatæ Mariæ Virginis.}

\rubric{¶ Post Laudes ac ceteras Horas, præter Completorium, dicto Versu Fidelium animæ et Oratione Dominica secreto recitata, subjungitur:}

\V\ Dóminus det nobis suam pacem. \R\ Et vitam ætérnam. Amen.

\rubric{Atque statim dicitur Antiphona finalis, ut infra.}

\sectionheading{Antiphona finalis beatæ Mariæ Virginis}

\rubric{A I Vesperis Festi sanctissimæ Trinitatis et deinceps:}

\antiphonaheading{Antiphona}

\lettrine[lines=2]{S}{ALVE}, Regína, mater misericórdiæ; vita, dulcédo et spes nostra, salve. Ad te clamámus éxsules fílii Hevæ. Ad te suspirámus geméntes et flentes in hac lacrimárum valle. Eja ergo, advocáta nostra, illos tuos misericórdes óculos ad nos convérte. Et Jesum, benedíctum fructum ventris tui, nobis post hoc exsílium osténde. O clemens, o pia, o dulcis Virgo María.

\V\ Ora pro nobis, sancta Dei Génitrix. \R\ Ut digni efficiámur promissiónibus Christi.

Orémus. \oratiolabel{Oratio}

\lettrine[lines=2]{O}{MNIPOTENS} sempitérne Deus, qui gloriósæ Vírginis Matris Maríæ corpus et ánimam, ut dignum Fílii tui habitáculum éffici mererétur, Spíritu Sancto cooperánte, præparásti: da, ut, cujus commemoratióne lætámur; ejus pia intercessióne, ab instántibus malis, et a morte perpétua liberémur. Per eúmdem Christum, Dóminum nostrum. \R\ Amen.

\rubric{Absoluta vero Antiphona finali cum suis Versu et Oratione, concluditur:}

\V\ Divínum auxílium máneat semper nobíscum. \R\ Amen.

\sectionheading{Post divinum Officium}

\rubric{Cum expletum fuerit Officium, laudabiliter dicitur sequens Oratio; pro qua Summus Pontifex Pius Papa XI sub die 1 Decembris 1933 indulgentiam trium annorum concessit. Præterea Clericis in sacris constitutis et Sacerdotibus eam devote post Officium recitantibus Summus Pontifex Leo X defectus et culpas in eo persolvendo ex humana fragilitate contractas, indulsit. Dicitur autem flexis semper genibus in privata etiam recitatione, præter quam ab iis, qui ob certam infirmitatem vel gravioris impedimenti causam nequeant genuflectere.}

\oratioheading{Oratio}

\lettrine[lines=2]{S}{ACROSANCTÆ} et indivíduæ Trinitáti, crucifíxi Dómini nostri Jesu Christi humanitáti, beatíssimæ et gloriosíssimæ sempérque Vírginis Maríæ fecúndæ integritáti, et ómnium Sanctórum universitáti sit sempitérna laus, honor, virtus et glória ab omni creatúra, nobísque remíssio ómnium peccatórum, per infiníta sæcula sæculórum. \R\ Amen.

\V\ Beáta víscera Maríæ Vírginis, quæ portavérunt ætérni Patris Fílium. \R\ Et beáta úbera, quæ lactavérunt Christum Dóminum.

\rubric{Et dicitur secreto} Pater noster et Ave María.

% [ornamental image ending Ordinarium Divini Officii]

% ===== END OF BATCH 3 -- Covers autumnalis.md lines 922-1140 (reused
% verbatim from aestiva.tex's Batch 3, per the byte-identity noted
% above), completing the Ordinarium Divini Officii in its entirety:
% the close of "Ad Vesperas" (the Suffragium de omnibus Sanctis in
% both its Marian and Josephine forms, the second occurrence ending
% without a printed "℟. Amen," transcribed literally), "Ad
% Completorium" (the Lectio brevis, the full Confiteor/Misereatur/
% Indulgentiam apparatus, the hymn "Te lucis ante terminum," the Nunc
% dimittis with its proper brief responsory "In manus tuas," and the
% closing "Benedicat et custodiat nos" under
% \benedictiolabel{Benedictio}, matching the established distinction
% from the abbreviated \benedlabel{Bened.} used earlier at the same
% Hour's opening blessing), "Post singulas Horas," and the single
% "Antiphona finalis beatæ Mariæ Virginis" -- "Salve, Regína" alone,
% assigned from First Vespers of Trinity Sunday onward with no other
% seasonal Marian antiphon given in this shared Ordinarium text --
% and finally "Post divinum Officium" with its closing Oratio
% "Sacrosanctæ et individuæ Trinitati." This completes the Ordinarium
% Divini Officii for the Pars Autumnalis, matching Pars Aestiva's
% exactly. The document continues with "Psalterium Breviarii Romani"
% (autumnalis.md line 1142) =====
\partheading{Psalterium Breviarii Romani}
\partsubtitle{Per omnes ac singulos hebdomadæ dies dispositum}

\feriaheading{Dominica}

\hourheading{Ad Matutinum}

Pater, Ave et Credo.

\V\ Dómine, lábia mea.

\V\ Deus, in adjutórium.

Glória Patri. Sicut erat. Allelúja.

\rubric{In Dominica IV et reliquis Dominicis post Pentecosten occurrentibus dicitur:}

\invitatlabel{Invitat.} Dóminum, qui fecit nos, * Veníte, adorémus. \psalmlabel{Ps. 94} Veníte.

\hymnusheading{Hymnus}

\lettrine[lines=2]{N}{OCTE} surgéntes vigilémus omnes,\\
Semper in psalmis meditémur, atque\\
Voce concórdi Dómino canámus\\
Dúlciter hymnos.

Ut pio Regi páriter canéntes,\\
Cum suis Sanctis mereámur aulam\\
Ingredi cæli, simul et perénnem\\
Dúcere vitam.

Præstet hoc nobis Déitas beáta\\
Patris, ac Nati, paritérque Sancti\\
Spíritus, cujus résonat per omnem\\
Glória mundum. Amen.

\nocturnoheading{In I Nocturno}

\antlabel{Ant.} Beátus vir. \textdagger

\rubric{¶ Et non repetitur in Psalmo; quod semper fit, quando Antiphona incipit a primis Psalmi vel Cantici verbis. Et Psalmus vel Canticum tunc incipitur ab eo verbo, ante quod vel inchoata vel integra desinit Antiphona; si tamen eadem sint verba, et continuetur Antiphona cum Psalmo vel Cantico, neque additur in fine ipsius Antiphonæ Allelúja.}

\rubric{Antiphonæ autem cum Psalmis et Versibus assignatæ, adhibentur semper, extra Octavas Ssmi Corporis Christi et Sacrat. Cordis Jesu, quando fit Officium de Dominica.}

\rubric{¶ Asteriscus *, in Psalmorum versibus et in aliis Officii partibus insertus, denotat pausam cantus et recitationis in Choro; præter quam in Antiphonis, in quibus nempe significat tum præintonationem Antiphonæ, tum Antiphonam ipsam ante Psalmum seu Canticum eousque producendam esse ad omnes et singulas Horas in Officio ritus semiduplicis ac simplicis, atque ad Horas tantum minores in Officio ritus duplicis.}

\psalmheading{Psalmus 1}

\lettrine[lines=2]{B}{EATUS} vir, \textdagger\ qui non ábiit in consílio impiórum, et in via peccatórum non stetit, * et in cáthedra pestiléntiæ non sedit:

Sed in lege Dómini volúntas ejus, * et in lege ejus meditábitur die ac nocte.

Et erit tamquam lignum, quod plantátum est secus decúrsus aquárum, * quod fructum suum dabit in témpore suo:

Et fólium ejus non défluet: * et ómnia quæcúmque fáciet, prosperabúntur.

Non sic ímpii, non sic: * sed tamquam pulvis, quem prójicit ventus a fácie terræ.

Ideo non resúrgent ímpii in judício: * neque peccatóres in concílio justórum.

Quóniam novit Dóminus viam justórum: * et iter impiórum períbit.

Glória Patri, et Fílio, * et Spirítui Sancto.

Sicut erat in princípio, et nunc, et semper, * et in sæcula sæculórum. Amen.

\rubric{¶ Hic versus Glória Patri semper dicitur in fine omnium Psalmorum et Canticorum, nisi aliter notetur.}

\antlabel{Ant.} Beátus vir, qui in lege Dómini meditátur.

\antlabel{Ant.} Servíte Dómino.

\psalmheading{Psalmus 2}

\lettrine[lines=2]{Q}{UARE} fremuérunt gentes, * et pópuli meditáti sunt inánia?

Astitérunt reges terræ, et príncipes convenérunt in unum * advérsus Dóminum, et advérsus Christum ejus.

Dirumpámus víncula eórum: * et projiciámus a nobis jugum ipsórum.

Qui hábitat in cælis, irridébit eos: * et Dóminus subsannábit eos.

Tunc loquétur ad eos in ira sua, * et in furóre suo conturbábit eos.

Ego autem constitútus sum Rex ab eo super Sion montem sanctum ejus, * prædicans præcéptum ejus.

Dóminus dixit ad me: * Fílius meus es tu, ego hódie génui te.

Póstula a me, et dabo tibi gentes hereditátem tuam, * et possessiónem tuam términos terræ.

Reges eos in virga férrea, * et tamquam vas fíguli confrínges eos.

Et nunc, reges, intellígite: * erudímini, qui judicátis terram.

Servíte Dómino in timóre: * et exsultáte ei cum tremóre.

Apprehéndite disciplínam, nequándo irascátur Dóminus, * et pereátis de via justa.

Cum exárserit in brevi ira ejus, * beáti omnes qui confídunt in eo.

\antlabel{Ant.} Servíte Dómino in timóre: et exsultáte ei cum tremóre.

\antlabel{Ant.} Exsúrge.

\psalmheading{Psalmus 3}

\lettrine[lines=2]{D}{OMINE}, quid multiplicáti sunt qui tríbulant me? * multi insúrgunt advérsum me.

Multi dicunt ánimæ meæ: * Non est salus ipsi in Deo ejus.

Tu autem, Dómine, suscéptor meus es, * glória mea, et exáltans caput meum.

Voce mea ad Dóminum clamávi: * et exaudívit me de monte sancto suo.

Ego dormívi, et soporátus sum: * et exsurréxi, quia Dóminus suscépit me.

Non timébo míllia pópuli circumdántis me: * exsúrge, Dómine, salvum me fac, Deus meus.

Quóniam tu percussísti omnes adversántes mihi sine causa: * dentes peccatórum contrivísti.

Dómini est salus: * et super pópulum tuum benedíctio tua.

\antlabel{Ant.} Exsúrge, Dómine, salvum me fac, Deus meus.

\V\ Memor fui nocte nóminis tui, Dómine. \R\ Et custodívi legem tuam.

\absolutiolabel{Absol.} Exáudi, Dómine.

\benedlabel{Bened.} 1 Benedictióne perpétua. 2 Unigénitus Dei Fílius. 3 Spíritus Sancti grátia.

\nocturnoheading{In II Nocturno}

\antlabel{Ant.} Quam admirábile.

\psalmheading{Psalmus 8}

\lettrine[lines=2]{D}{OMINE}, Dóminus noster, * quam admirábile est nomen tuum in univérsa terra!

Quóniam eleváta est magnificéntia tua, * super cælos.

Ex ore infántium et lacténtium perfecísti laudem propter inimícos tuos, * ut déstruas inimícum et ultórem.

Quóniam vidébo cælos tuos, ópera digitórum tuórum: * lunam et stellas, quæ tu fundásti.

Quid est homo, quod memor es ejus? * aut fílius hóminis, quóniam vísitas eum?

Minuísti eum paulo minus ab Angelis, glória et honóre coronásti eum: * et constituísti eum super ópera mánuum tuárum.

Omnia subjecísti sub pédibus ejus, * oves et boves univérsas: ínsuper et pécora campi.

Vólucres cæli, et pisces maris, * qui perámbulant sémitas maris.

Dómine, Dóminus noster, * quam admirábile est nomen tuum in univérsa terra!

\antlabel{Ant.} Quam admirábile est nomen tuum, Dómine, in univérsa terra!

\antlabel{Ant.} Sedísti super thronum.

\psalmheading{Psalmus 9, i}

\lettrine[lines=2]{C}{ONFITEBOR} tibi, Dómine, in toto corde meo: * narrábo ómnia mirabília tua.

Lætábor et exsultábo in te: * psallam nómini tuo, Altíssime,

In converténdo inimícum meum retrórsum: * infirmabúntur, et períbunt a fácie tua.

Quóniam fecísti judícium meum et causam meam: * sedísti super thronum, qui júdicas justítiam.

Increpásti gentes, et périit ímpius: * nomen eórum delésti in ætérnum, et in sæculum sæculi.

Inimíci defecérunt frámeæ in finem: * et civitátes eórum destruxísti.

Périit memória eórum cum sónitu: * et Dóminus in ætérnum pérmanet.

Parávit in judício thronum suum: * et ipse judicábit orbem terræ in æquitáte, judicábit pópulos in justítia.

Et factus est Dóminus refúgium páuperi: * adjútor in opportunitátibus, in tribulatióne.

Et sperent in te qui novérunt nomen tuum: * quóniam non dereliquísti quæréntes te, Dómine.

\antlabel{Ant.} Sedísti super thronum, qui júdicas justítiam.

\antlabel{Ant.} Exsúrge, Dómine.

\psalmheading{Psalmus 9, ii}

\lettrine[lines=2]{P}{SALLITE} Dómino, qui hábitat in Sion: * annuntiáte inter gentes stúdia ejus:

Quóniam requírens sánguinem eórum recordátus est: * non est oblítus clamórem páuperum.

Miserére mei, Dómine: * vide humilitátem meam de inimícis meis.

Qui exáltas me de portis mortis, * ut annúntiem omnes laudatiónes tuas in portis fíliæ Sion.

Exsultábo in salutári tuo: * infíxæ sunt gentes in intéritu, quem fecérunt.

In láqueo isto, quem abscondérunt, * comprehénsus est pes eórum.

Cognoscétur Dóminus judícia fáciens: * in opéribus mánuum suárum comprehénsus est peccátor.

Convertántur peccatóres in inférnum, * omnes gentes quæ obliviscúntur Deum.

Quóniam non in finem oblívio erit páuperis: * patiéntia páuperum non períbit in finem.

Exsúrge, Dómine, non confortétur homo: * judicéntur gentes in conspéctu tuo.

Constítue, Dómine, legislatórem super eos: * ut sciant gentes quóniam hómines sunt.

\antlabel{Ant.} Exsúrge, Dómine, non prævaleat homo.

\V\ Média nocte surgébam ad confiténdum tibi. \R\ Super judícia justificatiónis tuæ.

\absolutiolabel{Absol.} Ipsíus píetas et.

\benedlabel{Bened.} 4 Deus Pater omnípotens. 5 Christus perpétuæ. 6 Ignem sui amóris.

\nocturnoheading{In III Nocturno}

\antlabel{Ant.} Ut quid, Dómine. \textdagger

\psalmheading{Psalmus 9, iii}

\lettrine[lines=2]{U}{T} quid, Dómine, \textdagger\ recessísti longe, * déspicis in opportunitátibus, in tribulatióne?

Dum supérbit ímpius, incénditur pauper: * comprehendúntur in consíliis quibus cógitant.

Quóniam laudátur peccátor in desidériis ánimæ suæ: * et iníquus benedícitur.

Exacerbávit Dóminum peccátor, * secúndum multitúdinem iræ suæ non quæret.

Non est Deus in conspéctu ejus: * inquinátæ sunt viæ illíus in omni témpore.

Auferúntur judícia tua a fácie ejus: * ómnium inimicórum suórum dominábitur.

Dixit enim in corde suo: * Non movébor a generatióne in generatiónem sine malo.

Cujus maledictióne os plenum est, et amaritúdine, et dolo: * sub lingua ejus labor et dolor.

Sedet in insídiis cum divítibus in occúltis, * ut interfíciat innocéntem.

Oculi ejus in páuperem respíciunt: * insidiátur in abscóndito, quasi leo in spelúnca sua.

Insidiátur ut rápiat páuperem: * rápere páuperem, dum áttrahit eum.

In láqueo suo humiliábit eum: * inclinábit se, et cadet, cum dominátus fúerit páuperum.

Dixit enim in corde suo: * Oblítus est Deus, avértit fáciem suam ne vídeat in finem.

\antlabel{Ant.} Ut quid, Dómine, recessísti longe?

\antlabel{Ant.} Exsúrge. \textdagger

\psalmheading{Psalmus 9, iv}

\lettrine[lines=2]{E}{XSURGE}, \textdagger\ Dómine Deus, exaltétur manus tua: * ne obliviscáris páuperum.

Propter quid irritávit ímpius Deum? * dixit enim in corde suo: Non requíret.

Vides quóniam tu labórem et dolórem consíderas: * ut tradas eos in manus tuas.

Tibi derelíctus est pauper: * órphano tu eris adjútor.

Cóntere bráchium peccatóris et malígni: * quærétur peccátum illíus, et non inveniétur.

Dóminus regnábit in ætérnum, et in sæculum sæculi: * períbitis, gentes, de terra illíus.

Desidérium páuperum exaudívit Dóminus: * præparatiónem cordis eórum audívit auris tua.

Judicáre pupíllo et húmili, * ut non appónat ultra magnificáre se homo super terram.

\antlabel{Ant.} Exsúrge, Dómine Deus, exaltétur manus tua.

\antlabel{Ant.} Justus Dóminus.

\psalmheading{Psalmus 10}

\lettrine[lines=2]{I}{N} Dómino confído: quómodo dícitis ánimæ meæ: * Tránsmigra in montem sicut passer?

Quóniam ecce peccatóres intendérunt arcum, paravérunt sagíttas suas in pháretra, * ut sagíttent in obscúro rectos corde.

Quóniam quæ perfecísti, destruxérunt: * justus autem quid fecit?

Dóminus in templo sancto suo, * Dóminus in cælo sedes ejus:

Oculi ejus in páuperem respíciunt: * pálpebræ ejus intérrogant fílios hóminum.

Dóminus intérrogat justum et ímpium: * qui autem díligit iniquitátem, odit ánimam suam.

Pluet super peccatóres láqueos: * ignis, et sulphur, et spíritus procellárum pars cálicis eórum.

Quóniam justus Dóminus, et justítias diléxit: * æquitátem vidit vultus ejus.

\antlabel{Ant.} Justus Dóminus, et justítiam diléxit.

\V\ Prævenérunt óculi mei ad te dilúculo. \R\ Ut meditárer elóquia tua, Dómine.

\absolutiolabel{Absolutio} A vínculis.

\benedlabel{Bened.} 7 Evangélica léctio. 8 Divínum auxílium. 9 Ad societátem cívium.

% ===== END OF BATCH 4 -- Covers autumnalis.md lines 1142-1442, opening
% "Psalterium Breviarii Romani" with "Dominica" / "Ad Matutinum" in
% its ordinary-time form ("In Dominica IV et reliquis Dominicis post
% Pentecosten occurrentibus"): the Invitatorium "Dóminum, qui fecit
% nos," the hymn "Nocte surgentes," and the three Nocturns (Psalms 1,
% 2, 3; 8, 9/i, 9/ii; 9/iii, 9/iv, 10), each with its Absolutio and
% triple Benedictiones. Neither this Invitatorium nor this hymn
% appears anywhere in hiemalis.tex or verna.tex, confirming this is
% genuinely new ordinary-time-after-Pentecost material specific to
% the summer and autumn volumes. A structurally significant
% observation: after Psalm 1 prints "Glória Patri" in full together
% with the rubric establishing that this doxology is "semper dicitur
% in fine omnium Psalmorum et Canticorum, nisi aliter notetur," EVERY
% SUBSEQUENT PSALM IN THIS NOCTURNAL SET (2, 3, 8, 9/i, 9/ii, 9/iii,
% 9/iv, 10) OMITS THE PRINTED "Glória Patri" LINE ENTIRELY -- the
% source trusts the reader to supply it from the standing rubric
% rather than reprinting it each time. This transcription follows
% the source exactly: no "Glória Patri" line is added after any of
% these psalms, since none appears in autumnalis.md itself; only where
% the source itself prints it (as after Psalm 1) does the LaTeX
% include it. The embedded dagger marks (e.g. "Beátus vir, †",
% "Ut quid, Dómine, †", "Exsúrge, †") are rendered \textdagger\
% exactly where the source places them, including inside an
% antiphon incipit alone ("Ant. Ut quid, Dómine. †"). The second
% Nocturn's Absolutio is given as the abbreviated "Absol." but the
% third Nocturn's uses the unabbreviated "Absolutio" -- both render
% via the same flexible \absolutiolabel{} command, consistent with
% established practice. The document continues with "Ad Laudes"
% (autumnalis.md line 1443) =====
Sicut ádipe et pinguédine repleátur ánima mea: * et lábiis exsultatiónis laudábit os meum.

Si memor fui tui super stratum meum, in matutínis meditábor in te: * quia fuísti adjútor meus.

Et in velaménto alárum tuárum exsultábo, adhǽsit ánima mea post te: * me suscépit déxtera tua.

Ipsi vero in vanum quæsiérunt ánimam meam, introíbunt in inferióra terræ: * tradéntur in manus gládii, partes vúlpium erunt.

Rex vero lætábitur in Deo, laudabúntur omnes qui jurant in eo: * quia obstrúctum est os loquéntium iníqua.

\antlabel{Ant.} Benedícam te in vita mea, Dómine: et in nómine tuo levábo manus meas, allelúja.

\antlabel{Ant.} Tres púeri.

\canticumheading{Canticum trium Puerorum [4]}{Dan. 57-88 et 56}

\lettrine[lines=2]{B}{ENEDICITE}, ómnia ópera Dómini, Dómino: * laudáte et superexaltáte eum in sæcula.

Benedícite, Angeli Dómini, Dómino: * benedícite, cæli, Dómino.

Benedícite, aquæ omnes, quæ super cælos sunt, Dómino: * benedícite, omnes virtútes Dómini, Dómino.

Benedícite, sol et luna, Dómino: * benedícite, stellæ cæli, Dómino.

Benedícite, omnis imber et ros, Dómino: * benedícite, omnes spíritus Dei, Dómino.

Benedícite, ignis et æstus, Dómino: * benedícite, frigus et æstus, Dómino.

Benedícite, rores et pruína, Dómino: * benedícite, gelu et frigus, Dómino.

Benedícite, glácies et nives, Dómino: * benedícite, noctes et dies, Dómino.

Benedícite, lux et ténebræ, Dómino: * benedícite, fúlgura et nubes, Dómino.

Benedícat terra Dóminum: * laudet et superexáltet eum in sæcula.

Benedícite, montes et colles, Dómino: * benedícite, univérsa germinántia in terra, Dómino.

Benedícite, fontes, Dómino: * benedícite, mária et flúmina, Dómino.

Benedícite, cete, et ómnia, quæ movéntur in aquis, Dómino: * benedícite, omnes vólucres cæli, Dómino.

Benedícite, omnes béstiæ et pécora, Dómino: * benedícite, fílii hóminum, Dómino.

Benedícat Israël Dóminum: * laudet et superexáltet eum in sæcula.

Benedícite, sacerdótes Dómini, Dómino: * benedícite, servi Dómini, Dómino.

Benedícite, spíritus, et ánimæ justórum, Dómino: * benedícite, sancti, et húmiles corde, Dómino.

Benedícite, Ananía, Azaría, Mísaël, Dómino: * laudáte et superexaltáte eum in sæcula.

Benedicámus Patrem et Fílium cum Sancto Spíritu: * laudémus et superexaltémus eum in sæcula.

Benedíctus es, Dómine, in firmaménto cæli: * et laudábilis, et gloriósus, et superexaltátus in sæcula.

\rubric{¶ Hic non dicitur} Glória Patri, \rubric{neque} Amen.

\antlabel{Ant.} Tres púeri jussu regis in fornácem missi sunt, non timéntes flammam ignis, dicéntes: Benedíctus Deus, allelúja.

\antlabel{Ant.} Allelúja.

\psalmheading{Psalmus 148 [5]}

\lettrine[lines=2]{L}{AUDATE} Dóminum de cælis: * laudáte eum in excélsis.

Laudáte eum, omnes Angeli ejus: * laudáte eum, omnes virtútes ejus.

Laudáte eum, sol et luna: * laudáte eum, omnes stellæ et lumen.

Laudáte eum, cæli cælórum: * et aquæ omnes, quæ super cælos sunt, laudent nomen Dómini.

Quia ipse dixit, et facta sunt: * ipse mandávit, et creáta sunt.

Státuit ea in ætérnum, et in sæculum sæculi: * præcéptum pósuit, et non præteríbit.

Laudáte Dóminum de terra, * dracónes, et omnes abýssi.

Ignis, grando, nix, glácies, spíritus procellárum: * quæ fáciunt verbum ejus:

Montes, et omnes colles: * ligna fructífera, et omnes cedri.

Béstiæ, et univérsa pécora: * serpéntes, et vólucres pennátæ:

Reges terræ, et omnes pópuli: * príncipes, et omnes júdices terræ.

Júvenes, et vírgines: senes cum junióribus laudent nomen Dómini: * quia exaltátum est nomen ejus solíus.

Conféssio ejus super cælum et terram: * et exaltávit cornu pópuli sui.

Hymnus ómnibus sanctis ejus: * fíliis Israël, pópulo appropinquánti sibi.

\antlabel{Ant.} Allelúja, laudáte Dóminum de cælis, allelúja, allelúja.

\rubric{In omnibus Dominicis extra Octavas Ssmi Corporis Christi ac Sacratissimi Cordis Jesu, dicitur:}

\capitulumlabel{Capitulum}{Apoc. 7, 12}

\lettrine[lines=2]{B}{ENEDICTIO}, et cláritas, et sapiéntia, et gratiárum áctio, honor, virtus, et fortitúdo Deo nostro in sæcula sæculórum. Amen. \R\ Deo grátias.

\rubric{In Dominica IV et reliquis Dominicis post Pentecosten occurrentibus, dicitur}

\hymnusheading{Hymnus}

\lettrine[lines=2]{E}{CCE}, jam noctis tenuátur umbra,\\
Lux et auróræ rútilans corúscat:\\
Súpplices rerum Dóminum canóra\\
Voce precémur:

Ut reos culpæ miserátus omnem\\
Pellat angórem, tríbuat salútem,\\
Donet et nobis bona sempitérnæ\\
Múnera pacis.

Præstet hoc nobis Déitas beáta\\
Patris, ac Nati, paritérque Sancti\\
Spíritus, cujus résonat per omnem\\
Glória mundum. Amen.

\V\ Dóminus regnávit, decórem índuit. \R\ Induit Dóminus fortitúdinem, et præcínxit se virtúte.

\rubric{Ad Benedictus Antiphona ut in Proprio de Tempore.}

\canticumheading{Canticum Zachariæ}{Luc. 1, 68-79}

\lettrine[lines=2]{B}{ENEDICTUS} \textdagger\ Dóminus, Deus Israël, * quia visitávit, et fecit redemptiónem plebis suæ:

Et eréxit cornu salútis nobis: * in domo David, púeri sui.

Sicut locútus est per os sanctórum, * qui a sæculo sunt, prophetárum ejus:

Salútem ex inimícis nostris, * et de manu ómnium, qui odérunt nos:

Ad faciéndam misericórdiam cum pátribus nostris: * et memorári testaménti sui sancti.

Jusjurándum, quod jurávit ad Abraham, patrem nostrum, * datúrum se nobis:

Ut sine timóre, de manu inimicórum nostrórum liberáti, * serviámus illi.

In sanctitáte, et justítia coram ipso, * ómnibus diébus nostris.

Et tu, puer, Prophéta Altíssimi vocáberis: * præíbis enim ante fáciem Dómini paráre vias ejus:

Ad dandam sciéntiam salútis plebi ejus: * in remissiónem peccatórum eórum:

Per víscera misericórdiæ Dei nostri: * in quibus visitávit nos, óriens ex alto:

Illumináre his, qui in ténebris, et in umbra mortis sedent: * ad dirigéndos pedes nostros in viam pacis.

\rubric{Oratio conveniens. Reliqua ut in Ordinario. 11}

\rubric{¶ Psalmi cum Cantico de Dominica superius appositi dicuntur etiam, sub Antiphonis Officio diei congruentibus, in Dominicis infra Octavas Ssmi Corporis Christi ac Sacrat. Cordis Jesu, in Duplicibus I vel II classis, in Festis novem Lectionum Domini, B. Mariæ Virginis, Angelorum, S. Joannis Baptistæ, S. Joseph, Apostolorum, Evangelistarum et per Octavas privilegiatas; itemque in quolibet Festo novem Lectionum proprias Antiphonas ad Laudes habenti, quando de eis persolvitur Officium.}

\hourheading{Ad Primam}

Pater, Ave et Credo.

\V\ Deus, in adjutórium.

\hymnusheading{Hymnus}

\lettrine[lines=2]{J}{AM} lucis orto sídere,\\
Deum precémur súpplices,\\
Ut in diúrnis áctibus\\
Nos servet a nocéntibus.

Linguam refrénans témperet,\\
Ne litis horror ínsonet:\\
Visum fovéndo cóntegat,\\
Ne vanitátes háuriat.

Sint pura cordis íntima,\\
Absístat et vecórdia:\\
Carnis terat supérbiam\\
Potus cibíque párcitas.

Ut, cum dies abscésserit\\
Noctémque sors redúxerit,\\
Mundi per abstinéntiam\\
Ipsi canámus glóriam.

Deo Patri sit glória,\\
Ejúsque soli Fílio,\\
Cum Spíritu Paráclito,\\
Nunc et per omne sæculum. Amen.

\antlabel{Ant.} Allelúja.

\rubric{¶ Aña cum Psalmis, ad Horas minores assignata, adhibetur semper, extra Octavas Ssmi Corporis Christi ac Sacrat. Cordis Jesu, quando fit Officium de qualibet Dominica minori.}

\psalmheading{Psalmus 117}

\lettrine[lines=2]{C}{ONFITEMINI} Dómino quóniam bonus: * quóniam in sæculum misericórdia ejus.

Dicat nunc Israël quóniam bonus: * quóniam in sæculum misericórdia ejus.

Dicat nunc domus Aaron: * quóniam in sæculum misericórdia ejus.

Dicant nunc qui timent Dóminum: * quóniam in sæculum misericórdia ejus.

De tribulatióne invocávi Dóminum: * et exaudívit me in latitúdine Dóminus.

Dóminus mihi adjútor: * non timébo quid fáciat mihi homo.

Dóminus mihi adjútor: * et ego despíciam inimícos meos.

Bonum est confídere in Dómino, * quam confídere in hómine:

Bonum est speráre in Dómino, * quam speráre in princípibus.

Omnes gentes circuiérunt me: * et in nómine Dómini quia ultus sum in eos.

Circumdántes circumdedérunt me: * et in nómine Dómini quia ultus sum in eos.

Circumdedérunt me sicut apes, et exarsérunt sicut ignis in spinis: * et in nómine Dómini quia ultus sum in eos.

Impúlsus evérsus sum ut cáderem: * et Dóminus suscépit me.

Fortitúdo mea, et laus mea Dóminus: * et factus est mihi in salútem.

Vox exsultatiónis, et salútis * in tabernáculis justórum.

Déxtera Dómini fecit virtútem: déxtera Dómini exaltávit me, * déxtera Dómini fecit virtútem.

Non móriar, sed vivam: * et narrábo ópera Dómini.

Castígans castigávit me Dóminus: * et morti non trádidit me.

Aperíte mihi portas justítiæ, ingréssus in eas confitébor Dómino: * hæc porta Dómini, justi intrábunt in eam.

Confitébor tibi quóniam exaudísti me: * et factus es mihi in salútem.

Lápidem, quem reprobavérunt ædificántes: * hic factus est in caput ánguli.

A Dómino factum est istud: * et est mirábile in óculis nostris.

Hæc est dies, quam fecit Dóminus: * exsultémus, et lætémur in ea.

O Dómine, salvum me fac, o Dómine, bene prosperáre: * benedíctus qui venit in nómine Dómini.

Benedíximus vobis de domo Dómini: * Deus Dóminus, et illúxit nobis.

Constitúite diem solémnem in condénsis, * usque ad cornu altáris.

Deus meus es tu, et confitébor tibi: * Deus meus es tu, et exaltábo te.

Confitébor tibi quóniam exaudísti me: * et factus es mihi in salútem.

Confitémini Dómino quóniam bonus: * quóniam in sæculum misericórdia ejus.

\rubric{¶ In Dominicis infra Octavas Ssmi Corporis Christi ac Sacrat. Cordis Jesu, in Duplicibus I vel II classis, in Festis novem Lectionum Domini, B. Mariæ Virginis, Angelorum, S. Joannis Baptistæ, S. Joseph, Apostolorum, Evangelistarum, et per Octavas privilegiatas: si de eis fiat Officium, ad Primam, Tertiam, Sextam et Nonam dicuntur Psalmi de Dominica, sub Antiphonis Officio diei congruentibus, nisi aliter suo loco notetur; sed ad Primam, loco Psalmi 117 Confitemini, in ipsa etiam die Dominica semper dicitur}

\psalmheading{Psalmus 53}

\lettrine[lines=2]{D}{EUS}, in nómine tuo salvum me fac: * et in virtúte tua júdica me.

Deus, exáudi oratiónem meam: * áuribus pércipe verba oris mei.

Quóniam aliéni insurrexérunt advérsum me, et fortes quæsiérunt ánimam meam: * et non proposuérunt Deum ante conspéctum suum.

Ecce enim, Deus ádjuvat me: * et Dóminus suscéptor est ánimæ meæ.

Avérte mala inimícis meis: * et in veritáte tua dispérde illos.

Voluntárie sacrificábo tibi, * et confitébor nómini tuo, Dómine: quóniam bonum est:

Quóniam ex omni tribulatióne eripuísti me: * et super inimícos meos despéxit óculus meus.

\psalmheading{Psalmus 118, i}

\lettrine[lines=2]{B}{EATI} immaculáti in via: * qui ámbulant in lege Dómini.

Beáti, qui scrutántur testimónia ejus: * in toto corde exquírunt eum.

Non enim qui operántur iniquitátem: * in viis ejus ambulavérunt.

Tu mandásti * mandáta tua custodíri nimis.

Utinam dirigántur viæ meæ, * ad custodiéndas justificatiónes tuas!

Tunc non confúndar, * cum perspéxero in ómnibus mandátis tuis.

Confitébor tibi in directióne cordis: * in eo quod dídici judícia justítiæ tuæ.

Justificatiónes tuas custódiam: * non me derelínquas usquequáque.

In quo córrigit adolescéntior viam suam? * in custodiéndo sermónes tuos.

In toto corde meo exquisívi te: * ne repéllas me a mandátis tuis.

In corde meo abscóndi elóquia tua: * ut non peccem tibi.

Benedíctus es, Dómine: * doce me justificatiónes tuas.

In lábiis meis, * pronuntiávi ómnia judícia oris tui.

In via testimoniórum tuórum delectátus sum, * sicut in ómnibus divítiis.

In mandátis tuis exercébor: * et considerábo vias tuas.

In justificatiónibus tuis meditábor: * non oblivíscar sermónes tuos.

\psalmheading{Psalmus 118, ii}

\lettrine[lines=2]{R}{ETRIBUE} servo tuo, vivífica me: * et custódiam sermónes tuos.

Revéla óculos meos: * et considerábo mirabília de lege tua.

Íncola ego sum in terra: * non abscóndas a me mandáta tua.

Concupívit ánima mea desideráre justificatiónes tuas, * in omni témpore.

Increpásti supérbos: * maledícti qui declínant a mandátis tuis.

Aufer a me oppróbrium, et contémptum: * quia testimónia tua exquisívi.

Etenim sedérunt príncipes, et advérsum me loquebántur: * servus autem tuus exercebátur in justificatiónibus tuis.

Nam et testimónia tua meditátio mea est: * et consílium meum justificatiónes tuæ.

Adhǽsit paviménto ánima mea: * vivífica me secúndum verbum tuum.

Vias meas enuntiávi, et exaudísti me: * doce me justificatiónes tuas.

Viam justificatiónum tuárum ínstrue me: * et exercébor in mirabílibus tuis.

Dormitávit ánima mea præ tædio: * confírma me in verbis tuis.

Viam iniquitátis ámove a me: * et de lege tua miserére mei.

Viam veritátis elégi: * judícia tua non sum oblítus.

Adhǽsi testimóniis tuis, Dómine: * noli me confúndere.

Viam mandatórum tuórum cucúrri, * cum dilatásti cor meum.

\rubric{¶ Si Athanasianum Symbolum omittatur, expletis Psalmis præcedentibus, statim dicitur Antiphona, ut infra. 45}

\rubric{Sequens Symbolum dicitur in Dominicis post Pentecosten, quando fit Officium de Dominica ut in Psalterio, neque ulla Commemoratio Officii duplicis vel Octavæ recurrit, et in Festo sanctissimæ Trinitatis.}

\symbolumheading{Symbolum Athanasianum}

QUICUMQUE vult salvus esse, * ante ómnia opus est, ut téneat cathólicam fidem:

Quam nisi quisque íntegram inviolatámque serváverit, * absque dúbio in ætérnum períbit.

Fides autem cathólica hæc est: * ut unum Deum in Trinitáte, et Trinitátem in unitáte venerémur.

Neque confundéntes persónas, * neque substántiam separántes.

Ália est enim persóna Patris, ália Fílii, * ália Spíritus Sancti:

Sed Patris, et Fílii, et Spíritus Sancti una est divínitas, * æquális glória, coætérna majéstas.

Quális Pater, talis Fílius, * talis Spíritus Sanctus.

Increátus Pater, increátus Fílius, * increátus Spíritus Sanctus.

Imménsus Pater, imménsus Fílius, * imménsus Spíritus Sanctus.

Ætérnus Pater, ætérnus Fílius, * ætérnus Spíritus Sanctus.

Et tamen non tres ætérni, * sed unus ætérnus.

Sicut non tres increáti, nec tres imménsi, * sed unus increátus, et unus imménsus.

Simíliter omnípotens Pater, omnípotens Fílius, * omnípotens Spíritus Sanctus.

Et tamen non tres omnipoténtes, * sed unus omnípotens.

Ita Deus Pater, Deus Fílius, * Deus Spíritus Sanctus.

Et tamen non tres Dii, * sed unus est Deus.

Ita Dóminus Pater, Dóminus Fílius, * Dóminus Spíritus Sanctus.

Et tamen non tres Dómini, * sed unus est Dóminus.

Quia, sicut singillátim unamquámque persónam Deum ac Dóminum confitéri christiána veritáte compéllimur: * ita tres Deos aut Dóminos dícere cathólica religióne prohibémur.

Pater a nullo est factus: * nec creátus, nec génitus.

Fílius a Patre solo est: * non factus, nec creátus, sed génitus.

Spíritus Sanctus a Patre et Fílio: * non factus, nec creátus, nec génitus, sed procédens.

Unus ergo Pater, non tres Patres: unus Fílius, non tres Fílii: * unus Spíritus Sanctus, non tres Spíritus Sancti.

Et in hac Trinitáte nihil prius aut postérius, nihil majus aut minus: * sed totæ tres persónæ coætérnæ sibi sunt et coæquáles.

Ita ut per ómnia, sicut jam supra dictum est, * et únitas in Trinitáte, et Trínitas in unitáte veneránda sit.

Qui vult ergo salvus esse, * ita de Trinitáte séntiat.

Sed necessárium est ad ætérnam salútem, * ut Incarnatiónem quoque Dómini nostri Jesu Christi fidéliter credat.

Est ergo fides recta, ut credámus et confiteámur, * quia Dóminus noster Jesus Christus, Dei Fílius, Deus et homo est.

Deus est ex substántia Patris ante sæcula génitus: * et homo est ex substántia matris in sæculo natus.

Perféctus Deus, perféctus homo: * ex ánima rationáli et humána carne subsístens.

Æquális Patri secúndum divinitátem: * minor Patre secúndum humanitátem.

Qui, licet Deus sit et homo, * non duo tamen, sed unus est Christus.

Unus autem non conversióne divinitátis in carnem, * sed assumptióne humanitátis in Deum.

Unus omníno, non confusióne substántiæ, * sed unitáte persónæ.

Nam sicut ánima ratiónális et caro unus est homo: * ita Deus et homo unus est Christus.

Qui passus est pro salúte nostra: descéndit ad ínferos: * tértia die resurréxit a mórtuis.

Ascéndit ad cælos, sedet ad déxteram Dei Patris omnipoténtis: * inde ventúrus est judicáre vivos et mórtuos.

Ad cujus advéntum omnes hómines resúrgere habent cum corpóribus suis: * et redditúri sunt de factis própriis ratiónem.

Et qui bona egérunt, ibunt in vitam ætérnam: * qui vero mala, in ignem ætérnum.

Hæc est fides cathólica, * quam nisi quisque fidéliter firmitérque credíderit, salvus esse non póterit.

Glória Patri.

\rubric{Expleto Symbolo, vel, si hoc non habeat locum, Psalmis absolutis, statim dicitur Antiphona.}

\antlabel{Ant.} Allelúja, confitémini Dómino quóniam in sæculum misericórdia ejus, allelúja, allelúja.

\capitulumlabel{Capitulum}{1 Tim. 1, 17}

\lettrine[lines=2]{R}{EGI} sæculórum immortáli et invisíbili, soli Deo honor et glória in sæcula sæculórum. Amen.

\R\ br. Christe, Fili Dei vivi, * Miserére nobis. \rubric{Et repetitur:} Christe, Fili Dei vivi, miserére nobis. \V\ Qui sedes ad déxteram Patris. \R\ Miserére nobis. \V\ Glória Patri, et Fílio, et Spirítui Sancto. \R\ Christe, Fili Dei vivi, miserére nobis. \V\ Exsúrge, Christe, ádjuva nos. \R\ Et líbera nos propter nomen tuum.

\variantheading{Preces (si dicendæ sint)}

KYRIE, eléison. Christe, eléison. Kýrie, eléison.

Pater noster \rubric{secreto usque ad} \V\ Et ne nos indúcas in tentatiónem. \R\ Sed líbera nos a malo.

Credo in Deum \rubric{secreto usque ad} \V\ Carnis resurrectiónem. \R\ Vitam ætérnam. Amen.

\V\ Et ego ad te, Dómine, clamávi. \R\ Et mane orátio mea prævéniet te. \V\ Repleátur os meum laude. \R\ Ut cantem glóriam tuam, tota die magnitúdinem tuam. \V\ Dómine, avérte fáciem tuam a peccátis meis. \R\ Et omnes iniquitátes meas dele. \V\ Cor mundum crea in me, Deus. \R\ Et spíritum rectum ínnova in viscéribus meis. \V\ Ne projícias me a fácie tua. \R\ Et spíritum sanctum tuum ne áuferas a me. \V\ Redde mihi lætítiam salutáris tui. \R\ Et spíritu principáli confírma me. \V\ Adjutórium nostrum in nómine Dómini. \R\ Qui fecit cælum et terram.

\formulalabel{Confíteor. Misereátur. Indulgéntiam.}

\V\ Dignáre, Dómine, die isto. \R\ Sine peccáto nos custodíre. \V\ Miserére nostri, Dómine. \R\ Miserére nostri. \V\ Fiat misericórdia tua, Dómine, super nos. \R\ Quemádmodum sperávimus in te. \V\ Dómine, exáudi oratiónem meam. \R\ Et clamor meus ad te véniat.

\V\ Dóminus vobíscum.

Orémus. \oratiolabel{Oratio}

\lettrine[lines=2]{D}{OMINE}, Deus omnípotens, qui ad princípium hujus diéi nos perveníre fecísti: tua nos hódie salva virtúte; ut in hac die ad nullum declinémus peccátum, sed semper ad tuam justítiam faciéndam nostra procédant elóquia, dirigántur cogitatiónes et ópera. Per Dóminum. \R\ Amen.

\V\ Dóminus vobíscum. \V\ Benedicámus Dómino.

\rubric{In Choro Martyrologium.}

\V\ Pretiósa in conspéctu Dómini. \R\ Mors Sanctórum ejus.

\rubric{Et absolute, sine} Orémus:

\oratiolabel{Oratio}

\lettrine[lines=2]{S}{ANCTA} María, et omnes Sancti intercédant pro nobis ad Dóminum, ut nos mereámur ab eo adjuvári et salvári, qui vivit et regnat in sæcula sæculórum. \R\ Amen.

\V\ Deus, in adjutórium meum inténde. \R\ Dómine, ad adjuvándum me festína. \rubric{Et dicitur ter; ultimo additur:} Glória Patri. Sicut erat.

Kýrie, eléison. Christe, eléison. Kýrie, eléison.

Pater noster \rubric{secreto usque ad} \V\ Et ne nos indúcas in tentatiónem. \R\ Sed líbera nos a malo. \V\ Réspice in servos tuos, Dómine, et in ópera tua, et dírige fílios eórum. \R\ Et sit splendor Dómini, Dei nostri, super nos, et ópera mánuum nostrárum dírige super nos, et opus mánuum nostrárum dírige. \V\ Glória Patri. \R\ Sicut erat.

Orémus. \oratiolabel{Oratio}

\lettrine[lines=2]{D}{IRIGERE} et sanctificáre, régere et gubernáre dignáre, Dómine Deus, Rex cæli et terræ, hódie corda et córpora nostra, sensus, sermónes et actus nostros in lege tua, et in opéribus mandatórum tuórum: ut hic, et in ætérnum, te auxiliánte, salvi et líberi esse mereámur, Salvátor mundi: Qui vivis et regnas in sæcula sæculórum. \R\ Amen.

\V\ Jube, domne.

\benedlabel{Bened.} Dies et actus nostros in sua pace dispónat Dóminus omnípotens. \R\ Amen.

% ===== END OF BATCH 5 -- Covers autumnalis.md lines 1443-1984,
% completing "Ad Laudes" (Psalms 92, 99, 62, the Canticle of the
% Three Young Men "Benedicite," Psalm 148, the hymn "Ecce, jam noctis
% tenuatur umbra" for ordinary Sundays after Pentecost, and the
% Benedictus) and the body of "Ad Primam" (Psalm 117 "Confitemini,"
% the special Sunday substitute Psalm 53, Psalm 118 sections i-ii,
% the complete Athanasian Creed "Quicumque vult," the brief
% responsory, the Preces, the compact bold "Confíteor. Misereátur.
% Indulgéntiam." cross-reference rendered \formulalabel{} since the
% source bold-marks it here (matching the established rule that this
% command is used only when the source itself bolds the phrase,
% per precedent in verna.tex), the Martyrologium cue, and the closing
% Orationes). TWO SILENT CORRECTIONS of what appear to be markdown/
% OCR artifacts, per the standing rule for unambiguous mechanical
% errors: (1) a stray literal "\*" interrupting the hymn "Jam lucis
% orto sidere" at its final stanza ("\* Deo Patri sit gloria," verna.md
% line 1687) is dropped, since no hymn anywhere in this project uses
% chant-pause asterisks within its verse, and this same hymn appeared
% correctly without it in Batch 2; (2) a spurious paragraph break
% splitting the second stanza of "Ecce, jam noctis tenuatur umbra"
% into two two-line fragments (lines 1612-1616) is closed up into one
% continuous four-line stanza, matching the Sapphic-Adonic meter of
% every other stanza in the same hymn. The Athanasian Creed
% "Quicumque vult," though given in full capitals at its opening
% ("QUICUMQUE vult salvus esse"), receives no lettrine, matching the
% precedent already set in hiemalis.tex: the governing instruction's
% lettrine rule names oratio, psalm, canticle, hymn, and lectio
% specifically, and a Symbolum (Creed) is none of these, so
% capitalization alone does not trigger a lettrine here as it does
% for unheaded Oratio-like text elsewhere. Several versicle/response
% pairs in this batch print only the versicle with no response
% ("℣. Dóminus vobíscum." alone, "℣. Dóminus vobíscum. ℣.
% Benedicámus Dómino." with neither response given) -- transcribed
% exactly as the source abbreviates them, without supplying the
% missing responses. The document continues with the Lectio brevis
% variants for Prime (autumnalis.md line 1986), beginning with "In
% Officio dominicali" and continuing through several Proprio de
% Tempore feast-specific substitutions (Ssmæ Trinitatis, Ssmi
% Corporis Christi, etc.) =====
\variantheading{In Officio dominicali}

\lettrine[lines=2]{D}{OMINUS} autem dírigat corda et córpora nostra in caritáte Dei, et patiéntia Christi. Tu autem.

\variantheading{E Proprio de Tempore}

\variantheading{D. N. J. C. Regis}

\lettrine[lines=2]{I}{N} ipso complácuit omnem plenitúdinem inhabitáre, et per eum reconciliáre ómnia in ipsum, pacíficans per sánguinem crucis ejus sive quæ in terris, sive quæ in cælis sunt, in Christo Jesu Dómino nostro.

\variantheading{E Proprio Sanctorum}

\variantheading{S. Raphaëlis Archangeli}

\lettrine[lines=2]{T}{EMPUS} est, ut revértar ad eum, qui me misit: vos autem benedícite Deum, et narráte ómnia mirabília ejus.

\variantheading{Omnium Sanctorum}

\lettrine[lines=2]{B}{ENEDÍCTIO} et cláritas et sapiéntia et gratiárum áctio, honor, virtus et fortitúdo Deo nostro in sǽcula sæculórum. Amen.

\variantheading{S. Andreæ Apost.}

\lettrine[lines=2]{I}{SAÍAS} enim dicit: Dómine, quis crédidit audítui nostro? Ergo fides ex audítu, audítus autem per verbum Christi. Sed dico: Numquid non audiérunt? Et quidem in omnem terram exívit sonus eórum, et in fines orbis terræ verba eórum.

\variantheading{E Communi Sanctorum}

\variantheading{Pro Apostolis}

\lettrine[lines=2]{I}{BANT} Apóstoli gaudéntes a conspéctu concílii, quóniam digni hábiti sunt pro nómine Jesu contuméliam pati.

\variantheading{Pro Martyre}

\lettrine[lines=2]{J}{USTUS} cor suum trádidit ad vigilándum dilúculo ad Dóminum, qui fecit illum, et in conspéctu Altíssimi deprecábitur.

\variantheading{Pro pluribus Martyribus}

\lettrine[lines=2]{F}{ULGEBUNT} justi, et tamquam scintíllæ in arundinéto discúrrent. Judicábunt natiónes, et dominabúntur pópulis: et regnábit Dóminus illórum in perpétuum.

\variantheading{Pro Conf. Pont.}

\lettrine[lines=2]{F}{UNGI} sacerdótio, et habére laudem in nómine ipsíus, et offérre illi incénsum dignum in odórem suavitátis.

\variantheading{Pro Conf. non Pont.}

\lettrine[lines=2]{J}{USTUM} dedúxit Dóminus per vias rectas, et osténdit illi regnum Dei, et dedit illi sciéntiam sanctórum; honestávit illum in labóribus, et complévit labóres illíus.

\variantheading{Pro Virgine Martyre}

\lettrine[lines=2]{D}{OMINE}, Deus meus, exaltásti super terram habitatiónem meam, et pro morte defluénte deprecáta sum. Invocávi Dóminum, Patrem Dómini mei, ut non derelínquat me in die tribulatiónis meæ, et in témpore superbórum sine adjutório.

\variantheading{Pro Virgine non Martyre}

\lettrine[lines=2]{O}{} \textsc{quam} pulchra est casta generátio cum claritáte! immortális est enim memória illíus: quóniam apud Deum nota est, et apud hómines.

\variantheading{Pro Martyre tantum}

\lettrine[lines=2]{L}{AUDABIT} usque ad mortem ánima mea Dóminum, quóniam éruis sustinéntes te, et líberas eos de manu angústiæ, Dómine, Deus noster.

\variantheading{Pro nec Virg. nec Mart.}

\lettrine[lines=2]{M}{ULTAE} fíliæ congregavérunt divítias: tu supergréssa es univérsas. Fallax grátia, et vana est pulchritúdo: múlier timens Dóminum, ipsa laudábitur.

\variantheading{Pro Dedicatione Ecclesiæ}

\lettrine[lines=2]{E}{T} abstérget Deus omnem lácrimam ab óculis eórum: et mors ultra non erit, neque luctus, neque clamor, neque dolor erit ultra, quia prima abiérunt. Et dixit qui sedébat in throno: Ecce, nova fácio ómnia.

\variantheading{In Festis B. Mariæ V.}

\lettrine[lines=2]{I}{N} platéis sicut cinnamómum et bálsamum aromatízans odórem dedi: quasi myrrha elécta dedi suavitátem odóris.

\V\ Adjutórium nostrum in nómine Dómini. \R\ Qui fecit cælum et terram.

\V\ Benedícite. \R\ Deus.

\benedictiolabel{Benedictio} Dóminus nos benedícat, et ab omni malo deféndat, et ad vitam perdúcat ætérnam. Et fidélium ánimæ per misericórdiam Dei requiéscant in pace. \R\ Amen.

\rubric{Deinde dicitur tantum} Pater noster \rubric{secreto, nisi sequatur alia Hora.}

\hourheading{Ad Tertiam}

Pater et Ave.

\V\ Deus, in adjutórium.

\hymnusheading{Hymnus}

\lettrine[lines=2]{N}{UNC}, Sancte, nobis, Spíritus,\\
Unum Patri cum Fílio,\\
Dignáre promptus íngeri\\
Nostro refúsus péctori.

Os, lingua, mens, sensus, vigor,\\
Confessiónem pérsonent,\\
Flamméscat igne cáritas,\\
Accéndat ardor próximos.

Præsta, Pater piíssime,\\
Patríque compar Unice,\\
Cum Spíritu Paráclito\\
Regnans per omne sæculum. Amen.

\antlabel{Ant.} Allelúja.

\psalmheading{Psalmus 118, iii}

\lettrine[lines=2]{L}{EGEM} pone mihi, Dómine, viam justificatiónum tuárum: * et exquíram eam semper.

Da mihi intelléctum, et scrutábor legem tuam: * et custódiam illam in toto corde meo.

Deduc me in sémitam mandatórum tuórum: * quia ipsam vólui.

Inclína cor meum in testimónia tua: * et non in avarítiam.

Avérte óculos meos ne vídeant vanitátem: * in via tua vivífica me.

Státue servo tuo elóquium tuum, * in timóre tuo.

Ámputa oppróbrium meum quod suspicátus sum: * quia judícia tua jucúnda.

Ecce, concupívi mandáta tua: * in æquitáte tua vivífica me.

Et véniat super me misericórdia tua, Dómine: * salutáre tuum secúndum elóquium tuum.

Et respondébo exprobrántibus mihi verbum: * quia sperávi in sermónibus tuis.

Et ne áuferas de ore meo verbum veritátis usquequáque: * quia in judíciis tuis supersperávi.

Et custódiam legem tuam semper: * in sæculum et in sæculum sæculi.

Et ambulábam in latitúdine: * quia mandáta tua exquisívi.

Et loquébar in testimóniis tuis in conspéctu regum: * et non confundébar.

Et meditábar in mandátis tuis, * quæ diléxi.

Et levávi manus meas ad mandáta tua, quæ diléxi: * et exercébar in justificatiónibus tuis.

\psalmheading{Psalmus 118, iv}

\lettrine[lines=2]{M}{EMOR} esto verbi tui servo tuo, * in quo mihi spem dedísti.

Hæc me consoláta est in humilitáte mea: * quia elóquium tuum vivificábit me.

Supérbi iníque agébant usquequáque: * a lege autem tua non declinávi.

Memor fui judiciórum tuórum a sæculo, Dómine: * et consolátus sum.

Deféctio ténuit me, * pro peccatóribus derelinquéntibus legem tuam.

Cantábiles mihi erant justificatiónes tuæ, * in loco peregrinatiónis meæ.

Memor fui nocte nóminis tui, Dómine: * et custodívi legem tuam.

Hæc facta est mihi: * quia justificatiónes tuas exquisívi.

Pórtio mea, Dómine, * dixi custodíre legem tuam.

Deprecátus sum fáciem tuam in toto corde meo: * miserére mei secúndum elóquium tuum.

Cogitávi vias meas: * et convérti pedes meos in testimónia tua.

Parátus sum, et non sum turbátus: * ut custódiam mandáta tua.

Funes peccatórum circumpléxi sunt me: * et legem tuam non sum oblítus.

Média nocte surgébam ad confiténdum tibi, * super judícia justificatiónis tuæ.

Párticeps ego sum ómnium timéntium te: * et custodiéntium mandáta tua.

Misericórdia tua, Dómine, plena est terra: * justificatiónes tuas doce me.

\psalmheading{Psalmus 118, v}

\lettrine[lines=2]{B}{ONITATEM} fecísti cum servo tuo, Dómine, * secúndum verbum tuum.

Bonitátem, et disciplínam, et sciéntiam doce me: * quia mandátis tuis crédidi.

Priúsquam humiliárer ego delíqui: * proptérea elóquium tuum custodívi.

Bonus es tu: * et in bonitáte tua doce me justificatiónes tuas.

Multiplicáta est super me iníquitas superbórum: * ego autem in toto corde meo scrutábor mandáta tua.

Coagulátum est sicut lac cor eórum: * ego vero legem tuam meditátus sum.

Bonum mihi quia humiliásti me: * ut discam justificatiónes tuas.

Bonum mihi lex oris tui, * super míllia auri et argénti.

Manus tuæ fecérunt me, et plasmavérunt me: * da mihi intelléctum, et discam mandáta tua.

Qui timent te vidébunt me, et lætabúntur: * quia in verba tua supersperávi.

Cognóvi, Dómine, quia æquitas judícia tua: * et in veritáte tua humiliásti me.

Fiat misericórdia tua ut consolétur me, * secúndum elóquium tuum servo tuo.

Véniant mihi miseratiónes tuæ, et vivam: * quia lex tua meditátio mea est.

Confundántur supérbi, quia injúste iniquitátem fecérunt in me: * ego autem exercébor in mandátis tuis.

Convertántur mihi timéntes te: * et qui novérunt testimónia tua.

Fiat cor meum immaculátum in justificatiónibus tuis, * ut non confúndar.

\antlabel{Ant.} Allelúja, deduc me, Dómine, in sémitam mandatórum tuórum, allelúja, allelúja.

\capitulumlabel{Capitulum}{1 Joann. 4, 16}

\lettrine[lines=2]{D}{EUS} cáritas est: et qui manet in caritáte, in Deo manet, et Deus in eo.

\R\ br. Inclína cor meum, Deus, * In testimónia tua. Inclína. \V\ Avérte óculos meos, ne vídeant vanitátem: in via tua vivífica me. In testimónia tua. Glória Patri. Inclína. \V\ Ego dixi: Dómine, miserére mei. \R\ Sana ánimam meam, quia peccávi tibi.

\rubric{Oratio conveniens. Reliqua ut in Ordinario. 18}

% ===== END OF BATCH 6 -- Covers autumnalis.md lines 1986-2200,
% completing the Prime Lectio brevis variant list (substitute lessons
% across "In Officio dominicali," "E Proprio de Tempore" [here just
% Christ the King, the one movable feast of the Proper of Time this
% volume's Autumn scope reaches], "E Proprio Sanctorum" [St. Raphael
% the Archangel, All Saints, St. Andrew the Apostle -- autumnalis.md's
% own shorter three-feast list, verified via direct diff against
% aestiva.md's parallel eight-feast Summer list (Trinity Sunday
% through St. Lawrence) and substituted fresh from autumnalis.md's
% own text rather than carried over], and "E Communi Sanctorum" [the
% full set of twelve Common-of-Saints categories from Apostles
% through Marian feasts, confirmed byte-identical to aestiva.md's
% copy since the Common-of-Saints capitula do not vary by season])
% and the complete "Ad Tertiam" (hymn, Psalm 118 sections iii-v,
% Capitulum, and brief responsory). Each lesson, lacking its own
% repeated "Lectio brevis" label, is nonetheless lettrined on its
% full-caps opening word, consistent with its functional role as a
% Lectio substitute. "Pro Virgine non Martyre" opens "O QUAM
% pulchra..." with BOTH words in full caps (unlike the normal-case "O
% Mater"/"O Felícem" forms seen in Batch 166 of verna.tex), matching
% instead the all-caps "O GLORIOSA"/"O QUALE" pattern from that same
% project, so it receives the single-letter lettrine plus
% \textsc{Quam}. A second stray "\*" artifact, identical in kind to
% the one corrected in Batch 5, interrupts the final stanza of "Nunc,
% Sancte, nobis, Spiritus" at autumnalis.md line 2083 (within Ad
% Tertiam's hymn) and is silently dropped here for the same reason
% (no hymn anywhere in this project uses a chant-pause asterisk
% within its verse). The identical artifact recurs in Ad Sextam's
% hymn in the next batch and will be corrected the same way there.
% The document continues with "Ad Sextam" (autumnalis.md line 2202)
% =====
\hourheading{Ad Sextam}

Pater noster et Ave.

\V\ Deus, in adjutórium.

\hymnusheading{Hymnus}

\lettrine[lines=2]{R}{ECTOR} potens, verax Deus,\\
Qui témperas rerum vices,\\
Splendóre mane illúminas,\\
Et ígnibus merídiem:

Exstíngue flammas lítium,\\
Aufer calórem nóxium,\\
Confer salútem córporum,\\
Verámque pacem córdium.

Præsta, Pater piíssime,\\
Patríque compar Unice,\\
Cum Spíritu Paráclito\\
Regnans per omne sæculum. Amen.

\antlabel{Ant.} Allelúja.

\psalmheading{Psalmus 118, vi}

\lettrine[lines=2]{D}{EFECIT} in salutáre tuum ánima mea: * et in verbum tuum supersperávi.

Defecérunt óculi mei in elóquium tuum, * dicéntes: Quando consoláberis me?

Quia factus sum sicut uter in pruína: * justificatiónes tuas non sum oblítus.

Quot sunt dies servi tui? * quando fácies de persequéntibus me judícium?

Narravérunt mihi iníqui fabulatiónes: * sed non ut lex tua.

Omnia mandáta tua véritas: * iníque persecúti sunt me, ádjuva me.

Paulo minus consummavérunt me in terra: * ego autem non derelíqui mandáta tua.

Secúndum misericórdiam tuam vivífica me: * et custódiam testimónia oris tui.

In ætérnum, Dómine, * verbum tuum pérmanet in cælo.

In generatiónem et generatiónem véritas tua: * fundásti terram, et pérmanet.

Ordinatióne tua persevérat dies: * quóniam ómnia sérviunt tibi.

Nisi quod lex tua meditátio mea est: * tunc forte periíssem in humilitáte mea.

In ætérnum non oblivíscar justificatiónes tuas: * quia in ipsis vivificásti me.

Tuus sum ego, salvum me fac: * quóniam justificatiónes tuas exquisívi.

Me exspectavérunt peccatóres ut pérderent me: * testimónia tua intelléxi.

Omnis consummatiónis vidi finem: * latum mandátum tuum nimis.

\psalmheading{Psalmus 118, vii}

\lettrine[lines=2]{Q}{UOMODO} diléxi legem tuam, Dómine? * tota die meditátio mea est.

Super inimícos meos prudéntem me fecísti mandáto tuo: * quia in ætérnum mihi est.

Super omnes docéntes me intelléxi: * quia testimónia tua meditátio mea est.

Super senes intelléxi: * quia mandáta tua quæsívi.

Ab omni via mala prohíbui pedes meos: * ut custódiam verba tua.

A judíciis tuis non declinávi: * quia tu legem posuísti mihi.

Quam dúlcia fáucibus meis elóquia tua, * super mel ori meo!

A mandátis tuis intelléxi: * proptérea odívi omnem viam iniquitátis.

Lucérna pédibus meis verbum tuum, * et lumen sémitis meis.

Jurávi, et státui * custodíre judícia justítiæ tuæ.

Humiliátus sum usquequáque, Dómine: * vivífica me secúndum verbum tuum.

Voluntária oris mei beneplácita fac, Dómine: * et judícia tua doce me.

Anima mea in mánibus meis semper: * et legem tuam non sum oblítus.

Posuérunt peccatóres láqueum mihi: * et de mandátis tuis non errávi.

Hereditáte acquisívi testimónia tua in ætérnum: * quia exsultátio cordis mei sunt.

Inclinávi cor meum ad faciéndas justificatiónes tuas in ætérnum, * propter retributiónem.

\psalmheading{Psalmus 118, viii}

\lettrine[lines=2]{I}{NIQUOS} ódio hábui: * et legem tuam diléxi.

Adjútor et suscéptor meus es tu: * et in verbum tuum supersperávi.

Declináte a me, malígni: * et scrutábor mandáta Dei mei.

Súscipe me secúndum elóquium tuum, et vivam: * et non confúndas me ab exspectatióne mea.

Adjuva me, et salvus ero: * et meditábor in justificatiónibus tuis semper.

Sprevísti omnes discedéntes a judíciis tuis: * quia injústa cogitátio eórum.

Prævaricántes reputávi omnes peccatóres terræ: * ídeo diléxi testimónia tua.

Confíge timóre tuo carnes meas: * a judíciis enim tuis tímui.

Feci judícium et justítiam: * non tradas me calumniántibus me.

Súscipe servum tuum in bonum: * non calumniéntur me supérbi.

Óculi mei defecérunt in salutáre tuum: * et in elóquium justítiæ tuæ.

Fac cum servo tuo secúndum misericórdiam tuam: * et justificatiónes tuas doce me.

Servus tuus sum ego: * da mihi intelléctum, ut sciam testimónia tua.

Tempus faciéndi, Dómine: * dissipavérunt legem tuam.

Ideo diléxi mandáta tua, * super aurum et topázion.

Proptérea ad ómnia mandáta tua dirigébar: * omnem viam iníquam ódio hábui.

\antlabel{Ant.} Allelúja, tuus sum ego, salvum me fac, Dómine, allelúja, allelúja.

\capitulumlabel{Capitulum}{Gal. 6, 2}

\lettrine[lines=2]{A}{LTER} altérius ónera portáte, et sic adimplébitis legem Christi.

\R\ br. In ætérnum, Dómine, * Pérmanet verbum tuum. In ætérnum. \V\ In sæculum sæculi véritas tua. Pérmanet. Glória Patri. In ætérnum.

\V\ Dóminus regit me, et nihil mihi déerit. \R\ In loco páscuæ ibi me collocávit.

\rubric{Oratio conveniens. Reliqua ut in Ordinario. 20}

\hourheading{Ad Nonam}

Pater noster et Ave.

\V\ Deus, in adjutórium.

\hymnusheading{Hymnus}

\lettrine[lines=2]{R}{ERUM}, Deus, tenax vigor,\\
Immótus in te pérmanens,\\
Lucis diúrnæ témpora\\
Successíbus detérminans:

Largíre lumen véspere,\\
Quo vita nusquam décidat,\\
Sed præmium mortis sacræ\\
Perénnis instet glória.

Præsta, Pater piíssime,\\
Patríque compar Unice,\\
Cum Spíritu Paráclito\\
Regnans per omne sæculum. Amen.

\antlabel{Ant.} Allelúja.

\psalmheading{Psalmus 118, ix}

\lettrine[lines=2]{M}{IRABILIA} testimónia tua: * ídeo scrutáta est ea ánima mea.

Declarátio sermónum tuórum illúminat: * et intelléctum dat párvulis.

Os meum apérui, et attráxi spíritum: * quia mandáta tua desiderábam.

Aspice in me, et miserére mei, * secúndum judícium diligéntium nomen tuum.

Gressus meos dírige secúndum elóquium tuum: * et non dominétur mei omnis injustítia.

Rédime me a calúmniis hóminum: * ut custódiam mandáta tua.

Fáciem tuam illúmina super servum tuum: * et doce me justificatiónes tuas.

Exitus aquárum deduxérunt óculi mei: * quia non custodiérunt legem tuam.

Justus es, Dómine: * et rectum judícium tuum.

Mandásti justítiam testimónia tua: * et veritátem tuam nimis.

Tabéscere me fecit zelus meus: * quia oblíti sunt verba tua inimíci mei.

Ignítum elóquium tuum veheménter: * et servus tuus diléxit illud.

Adolescéntulus sum ego et contémptus: * justificatiónes tuas non sum oblítus.

Justítia tua, justítia in ætérnum: * et lex tua véritas.

Tribulátio, et angústia invenérunt me: * mandáta tua meditátio mea est.

Æquitas testimónia tua in ætérnum: * intelléctum da mihi, et vivam. Glória Patri.

\psalmheading{Psalmus 118, x}

\lettrine[lines=2]{C}{LAMAVI} in toto corde meo, exáudi me, Dómine: * justificatiónes tuas requíram.

Clamávi ad te, salvum me fac: * ut custódiam mandáta tua.

Prævéni in maturitáte, et clamávi: * quia in verba tua supersperávi.

Prævenérunt óculi mei ad te dilúculo: * ut meditárer elóquia tua.

Vocem meam audi secúndum misericórdiam tuam, Dómine: * et secúndum judícium tuum vivífica me.

Appropinquavérunt persequéntes me iniquitáti: * a lege autem tua longe facti sunt.

Prope es tu, Dómine: * et omnes viæ tuæ véritas.

Inítio cognóvi de testimóniis tuis: * quia in ætérnum fundásti ea.

Vide humilitátem meam, et éripe me: * quia legem tuam non sum oblítus.

Júdica judícium meum, et rédime me: * propter elóquium tuum vivífica me.

Longe a peccatóribus salus: * quia justificatiónes tuas non exquisiérunt.

Misericórdiæ tuæ multæ, Dómine: * secúndum judícium tuum vivífica me.

Multi qui persequúntur me, et tríbulant me: * a testimóniis tuis non declinávi.

Vidi prævaricántes, et tabescébam: * quia elóquia tua non custodiérunt.

Vide quóniam mandáta tua diléxi, Dómine: * in misericórdia tua vivífica me.

Princípium verbórum tuórum, véritas: * in ætérnum ómnia judícia justítiæ tuæ.

\psalmheading{Psalmus 118, xi}

\lettrine[lines=2]{P}{RINCIPES} persecúti sunt me gratis: * et a verbis tuis formidávit cor meum.

Lætábor ego super elóquia tua: * sicut qui invénit spólia multa.

Iniquitátem ódio hábui, et abominátus sum: * legem autem tuam diléxi.

Sépties in die laudem dixi tibi, * super judícia justítiæ tuæ.

Pax multa diligéntibus legem tuam: * et non est illis scándalum.

Exspectábam salutáre tuum, Dómine: * et mandáta tua diléxi.

Custodívit ánima mea testimónia tua: * et diléxit ea veheménter.

Servávi mandáta tua, et testimónia tua: * quia omnes viæ meæ in conspéctu tuo.

Appropínquet deprecátio mea in conspéctu tuo, Dómine: * juxta elóquium tuum da mihi intelléctum.

Intret postulátio mea in conspéctu tuo: * secúndum elóquium tuum éripe me.

Eructábunt lábia mea hymnum, * cum docúeris me justificatiónes tuas.

Pronuntiábit lingua mea elóquium tuum: * quia ómnia mandáta tua æquitas.

Fiat manus tua ut salvet me: * quóniam mandáta tua elégi.

Concupívi salutáre tuum, Dómine: * et lex tua meditátio mea est.

Vivet ánima mea, et laudábit te: * et judícia tua adjuvábunt me.

Errávi, sicut ovis, quæ périit: * quære servum tuum, quia mandáta tua non sum oblítus.

\antlabel{Ant.} Allelúja, fáciem tuam, Dómine, illúmina super servum tuum, allelúja, allelúja.

\capitulumlabel{Capitulum}{1 Cor. 6, 20}

\lettrine[lines=2]{E}{MPTI} enim estis prétio magno. Glorificáte et portáte Deum in córpore vestro.

\R\ br. Clamávi in toto corde meo: * Exáudi me, Dómine. Clamávi. \V\ Justificatiónes tuas requíram. Exáudi. Glória Patri. Clamávi. \V\ Ab occúltis meis munda me, Dómine. \R\ Et ab aliénis parce servo tuo.

\rubric{Oratio conveniens. Reliqua ut in Ordinario. 21}

% ===== END OF BATCH 7 -- Covers autumnalis.md lines 2202-2477,
% completing "Ad Sextam" (Psalm 118 sections vi-viii, Capitulum, and
% brief responsory) and "Ad Nonam" (Psalm 118 sections ix-xi,
% Capitulum, and brief responsory). A third occurrence of the stray
% "\*" artifact already corrected twice before (verna.md's hymn
% repetitions aside -- this is autumnalis.md's own recurring OCR/
% markdown glitch) interrupts Ad Sextam's hymn "Rector potens, verax
% Deus" at its final stanza (autumnalis.md line 2359) and is silently
% dropped, consistent with Batches 5 and 6. Psalm 118's ninth section
% (Mirabilia) is the first of this long psalm's many sections to
% print "Glória Patri." at its end (autumnalis.md line 2398) rather than
% omitting it per the usual pattern established since Batch 4 --
% transcribed exactly as the source gives it, neither adding nor
% removing a doxology relative to what each specific occurrence
% prints. The document continues with "Ad Vesperas" (autumnalis.md line
% 2515), opening the classic Sunday Vespers psalms (109 "Dixit
% Dominus," 110, etc.) =====
\hourheading{Ad Vesperas}

Pater noster et Ave.

\V\ Deus, in adjutórium.

\antlabel{Ant.} Dixit Dóminus. \textdagger

\rubric{¶ Antiphonæ cum Psalmis assignatæ adhibentur semper, quando fit Officium de Dominica, in Dominicis extra Octavas Ssmi Corporis Christi ac Sacratissimi Cordis Jesu.}

\psalmheading{Psalmus 109}

\lettrine[lines=2]{D}{IXIT} Dóminus \textdagger\ Dómino meo: * Sede a dextris meis:

Donec ponam inimícos tuos, * scabéllum pedum tuórum.

Virgam virtútis tuæ emíttet Dóminus ex Sion: * domináre in médio inimicórum tuórum.

Tecum princípium in die virtútis tuæ in splendóribus sanctórum: * ex útero ante lucíferum génui te.

Jurávit Dóminus, et non pænitébit eum: * Tu es sacérdos in ætérnum secúndum órdinem Melchísedech.

Dóminus a dextris tuis, * confrégit in die iræ suæ reges.

Judicábit in natiónibus, implébit ruínas: * conquassábit cápita in terra multórum.

De torrénte in via bibet: * proptérea exaltábit caput.

\antlabel{Ant.} Dixit Dóminus Dómino meo: Sede a dextris meis.

\antlabel{Ant.} Magna ópera Dómini.

\psalmheading{Psalmus 110}

\lettrine[lines=2]{C}{ONFITEBOR} tibi, Dómine, in toto corde meo: * in consílio justórum, et congregatióne.

Magna ópera Dómini: * exquisíta in omnes voluntátes ejus.

Conféssio et magnificéntia opus ejus: * et justítia ejus manet in sæculum sæculi.

Memóriam fecit mirabílium suórum, miséricors et miserátor Dóminus: * escam dedit timéntibus se.

Memor erit in sæculum testaménti sui: * virtútem óperum suórum annuntiábit pópulo suo:

Ut det illis hereditátem géntium: * ópera mánuum ejus véritas, et judícium.

Fidélia ómnia mandáta ejus: confirmáta in sæculum sæculi, * facta in veritáte et æquitáte.

Redemptiónem misit pópulo suo: * mandávit in ætérnum testaméntum suum.

Sanctum, et terríbile nomen ejus: * inítium sapiéntiæ timor Dómini.

Intelléctus bonus ómnibus faciéntibus eum: * laudátio ejus manet in sæculum sæculi.

\antlabel{Ant.} Magna ópera Dómini: exquisíta in omnes voluntátes ejus.

\antlabel{Ant.} Qui timet Dóminum.

\psalmheading{Psalmus 111}

\lettrine[lines=2]{B}{EATUS} vir, qui timet Dóminum: * in mandátis ejus volet nimis.

Potens in terra erit semen ejus: * generátio rectórum benedicétur.

Glória, et divítiæ in domo ejus: * et justítia ejus manet in sæculum sæculi.

Exórtum est in ténebris lumen rectis: * miséricors, et miserátor, et justus.

Jucúndus homo qui miserétur et cómmodat, dispónet sermónes suos in judício: * quia in ætérnum non commovébitur.

In memória ætérna erit justus: * ab auditióne mala non timébit.

Parátum cor ejus speráre in Dómino, confirmátum est cor ejus: * non commovébitur donec despíciat inimícos suos.

Dispérsit, dedit paupéribus: justítia ejus manet in sæculum sæculi, * cornu ejus exaltábitur in glória.

Peccátor vidébit, et irascétur, déntibus suis fremet et tabéscet: * desidérium peccatórum períbit.

\antlabel{Ant.} Qui timet Dóminum, in mandátis ejus cupit nimis.

\antlabel{Ant.} Sit nomen Dómini.

\psalmheading{Psalmus 112}

\lettrine[lines=2]{L}{AUDATE}, púeri, Dóminum: * laudáte nomen Dómini.

Sit nomen Dómini benedíctum, * ex hoc nunc, et usque in sæculum.

A solis ortu usque ad occásum, * laudábile nomen Dómini.

Excélsus super omnes gentes Dóminus, * et super cælos glória ejus.

Quis sicut Dóminus, Deus noster, qui in altis hábitat, * et humília réspicit in cælo et in terra?

Súscitans a terra ínopem, * et de stércore érigens páuperem:

Ut cóllocet eum cum princípibus, * cum princípibus pópuli sui.

Qui habitáre facit stérilem in domo, * matrem filiórum lætántem.

\antlabel{Ant.} Sit nomen Dómini benedíctum in sæcula.

\antlabel{Ant.} Deus autem noster.

\psalmheading{Psalmus 113}

\lettrine[lines=2]{I}{N} éxitu Israël de Ægýpto, * domus Jacob de pópulo bárbaro:

Facta est Judǽa sanctificátio ejus, * Israël potéstas ejus.

Mare vidit, et fugit: * Jordánis convérsus est retrórsum.

Montes exsultavérunt ut aríetes, * et colles sicut agni óvium.

Quid est tibi, mare, quod fugísti: * et tu, Jordánis, quia convérsus es retrórsum?

Montes, exsultástis sicut aríetes, * et colles, sicut agni óvium?

A fácie Dómini mota est terra, * a fácie Dei Jacob.

Qui convértit petram in stagna aquárum, * et rupem in fontes aquárum.

Non nobis, Dómine, non nobis: * sed nómini tuo da glóriam.

Super misericórdia tua, et veritáte tua: * nequándo dicant gentes: Ubi est Deus eórum?

Deus autem noster in cælo: * ómnia quæcúmque vóluit, fecit.

Simulácra géntium argéntum, et aurum, * ópera mánuum hóminum.

Os habent, et non loquéntur: * óculos habent, et non vidébunt.

Aures habent, et non áudient: * nares habent, et non odorábunt.

Manus habent, et non palpábunt: pedes habent, et non ambulábunt: * non clamábunt in gútture suo.

Símiles illis fiant qui fáciunt ea: * et omnes qui confídunt in eis.

Domus Israël sperávit in Dómino: * adjútor eórum et protéctor eórum est.

Domus Aaron sperávit in Dómino: * adjútor eórum et protéctor eórum est.

Qui timent Dóminum, speravérunt in Dómino: * adjútor eórum et protéctor eórum est.

Dóminus memor fuit nostri: * et benedíxit nobis:

Benedíxit dómui Israël: * benedíxit dómui Aaron.

Benedíxit ómnibus, qui timent Dóminum, * pusíllis cum majóribus.

Adjíciat Dóminus super vos: * super vos, et super fílios vestros.

Benedícti vos a Dómino, * qui fecit cælum, et terram.

Cælum cæli Dómino: * terram autem dedit fíliis hóminum.

Non mórtui laudábunt te, Dómine: * neque omnes, qui descéndunt in inférnum.

Sed nos qui vívimus, benedícimus Dómino, * ex hoc nunc et usque in sæculum.

\antlabel{Ant.} Deus autem noster in cælo: ómnia quæcúmque vóluit, fecit.

\capitulumlabel{Capitulum}{2 Cor. 1, 3-4}

\lettrine[lines=2]{B}{ENEDICTUS} Deus, et Pater Dómini nostri Jesu Christi, Pater misericordiárum, et Deus totíus consolatiónis, qui consolátur nos in omni tribulatióne nostra.

\hymnusheading{Hymnus}

\lettrine[lines=2]{L}{UCIS} Creátor óptime,\\
Lucem diérum próferens,\\
Primórdiis lucis novæ,\\
Mundi parans oríginem:

Qui mane junctum vésperi\\
Diem vocári præcipis:\\
Illábitur tetrum chaos,\\
Audi preces cum flétibus.

Ne mens graváta crímine,\\
Vitæ sit exsul múnere,\\
Dum nil perénne cógitat,\\
Seséque culpis ílligat.

Cæléste pulset óstium:\\
Vitále tollat præmium:\\
Vitémus omne nóxium:\\
Purgémus omne péssimum.

Præsta, Pater piíssime,\\
Patríque compar Unice,\\
Cum Spíritu Paráclito\\
Regnans per omne sæculum. Amen.

\V\ Dirigátur, Dómine, orátio mea. \R\ Sicut incénsum in conspéctu tuo.

\rubric{Ad Magnificat Antiphona ut in Proprio de Tempore.}

\canticumheading{Canticum B. Mariæ Virg.}{Luc. 1, 46-55}

\lettrine[lines=2]{M}{AGNIFICAT} \textdagger\ ánima mea Dóminum:

Et exsultávit spíritus meus * in Deo, salutári meo.

Quia respéxit humilitátem ancíllæ suæ: * ecce enim ex hoc beátam me dicent omnes generatiónes.

Quia fecit mihi magna qui potens est: * et sanctum nomen ejus.

Et misericórdia ejus a progénie in progénies * timéntibus eum.

Fecit poténtiam in bráchio suo: * dispérsit supérbos mente cordis sui.

Depósuit poténtes de sede, * et exaltávit húmiles.

Esuriéntes implévit bonis: * et dívites dimísit inánes.

Suscépit Israël, púerum suum, * recordátus misericórdiæ suæ.

Sicut locútus est ad patres nostros, * Abraham, et sémini ejus in sæcula.

\rubric{Oratio conveniens. Reliqua ut in Ordinario. 22}

\hourheading{Ad Completorium}

\V\ Jube, domne.

\benedlabel{Bened.} Noctem quiétam, et finem perféctum concédat nobis Dóminus omnípotens. \R\ Amen.

\lectioheading{Lectio br.}{1 Petri 5, 8-9}

\lettrine[lines=2]{F}{RATRES}: Sóbrii estóte, et vigiláte: quia adversárius vester diábolus tamquam leo rúgiens círcuit, quærens quem dévoret: cui resístite fortes in fide. Tu autem, Dómine, miserére nobis.

\R\ Deo grátias.

\V\ Adjutórium nostrum in nómine Dómini. \R\ Qui fecit cælum et terram.

\rubric{Pater noster, quod dicitur totum secreto.}

\formulalabel{Confíteor. Misereátur. Indulgéntiam.}

\V\ Convérte nos, Deus, salutáris noster. \R\ Et avérte iram tuam a nobis.

\V\ Deus, in adjutórium.

\antlabel{Ant.} Miserére.

\psalmheading{Psalmus 4}

\lettrine[lines=2]{C}{UM} invocárem exaudívit me Deus justítiæ meæ: * in tribulatióne dilatásti mihi.

Miserére mei, * et exáudi oratiónem meam.

Fílii hóminum, úsquequo gravi corde? * ut quid dilígitis vanitátem, et quæritis mendácium?

Et scitóte quóniam mirificávit Dóminus sanctum suum: * Dóminus exáudiet me cum clamávero ad eum.

Irascímini, et nolíte peccáre: * quæ dícitis in córdibus vestris, in cubílibus vestris compungímini.

Sacrificáte sacrifícium justítiæ, et speráte in Dómino. * Multi dicunt: Quis osténdit nobis bona?

Signátum est super nos lumen vultus tui, Dómine: * dedísti lætítiam in corde meo.

A fructu fruménti, vini, et ólei sui * multiplicáti sunt.

In pace in idípsum * dórmiam, et requiéscam;

Quóniam tu, Dómine, singuláriter in spe * constituísti me.

\psalmheading{Psalmus 90}

\lettrine[lines=2]{Q}{UI} hábitat in adjutório Altíssimi, * in protectióne Dei cæli commorábitur.

Dicet Dómino: Suscéptor meus es tu, et refúgium meum: * Deus meus, sperábo in eum.

Quóniam ipse liberávit me de láqueo venántium, * et a verbo áspero.

Scápulis suis obumbrábit tibi: * et sub pennis ejus sperábis.

Scuto circúmdabit te véritas ejus: * non timébis a timóre noctúrno.

A sagítta volánte in die, a negótio perambulánte in ténebris: * ab incúrsu, et dæmónio meridiáno.

Cadent a látere tuo mille, et decem míllia a dextris tuis: * ad te autem non appropinquábit.

Verúmtamen óculis tuis considerábis: * et retributiónem peccatórum vidébis.

Quóniam tu es, Dómine, spes mea: * Altíssimum posuísti refúgium tuum.

Non accédet ad te malum: * et flagéllum non appropinquábit tabernáculo tuo.

Quóniam Angelis suis mandávit de te: * ut custódiant te in ómnibus viis tuis.

In mánibus portábunt te: * ne forte offéndas ad lápidem pedem tuum.

Super áspidem, et basilíscum ambulábis: * et conculcábis leónem et dracónem.

Quóniam in me sperávit, liberábo eum: * prótegam eum, quóniam cognóvit nomen meum.

Clamábit ad me, et ego exáudiam eum: * cum ipso sum in tribulatióne: erípiam eum et glorificábo eum.

Longitúdine diérum replébo eum: * et osténdam illi salutáre meum.

\psalmheading{Psalmus 133}

\lettrine[lines=2]{E}{CCE}, nunc benedícite Dóminum, * omnes servi Dómini:

Qui statis in domo Dómini, * in átriis domus Dei nostri.

In nóctibus extóllite manus vestras in sancta, * et benedícite Dóminum.

Benedícat te Dóminus ex Sion, * qui fecit cælum et terram.

\antlabel{Ant.} Miserére mihi, Dómine, et exáudi oratiónem meam.

\hymnusheading{Hymnus}

\lettrine[lines=2]{T}{E} lucis ante términum,\\
Rerum Creátor, póscimus,\\
Ut pro tua cleméntia\\
Sis præsul et custódia.

Procul recédant sómnia,\\
Et nóctium phantásmata;\\
Hostémque nostrum cómprime,\\
Ne polluántur córpora.

Præsta, Pater piíssime,\\
Patríque compar Unice,\\
Cum Spíritu Paráclito\\
Regnans per omne sæculum. Amen.

\capitulumlabel{Capitulum}{Jer. 14, 9}

\lettrine[lines=2]{T}{U} autem in nobis es, Dómine, et nomen sanctum tuum invocátum est super nos: ne derelínquas nos, Dómine, Deus noster.

\R\ br. In manus tuas, Dómine, * Comméndo spíritum meum. In manus. \V\ Redemísti nos, Dómine, Deus veritátis. Comméndo. Glória Patri. In manus.

\V\ Custódi nos, Dómine, ut pupíllam óculi. \R\ Sub umbra alárum tuárum prótege nos.

\canticumheading{Canticum Simeonis}{Luc. 2, 29-32}

\lettrine[lines=2]{N}{UNC} dimíttis servum tuum, Dómine, * secúndum verbum tuum in pace:

Quia vidérunt óculi mei * salutáre tuum,

Quod parásti * ante fáciem ómnium populórum,

Lumen ad revelatiónem géntium, * et glóriam plebis tuæ Israël.

Glória Patri.

\antlabel{Ant.} Salva nos, Dómine, vigilántes, custódi nos dormiéntes; ut vigilémus cum Christo, et requiescámus in pace.

\variantheading{Preces (si dicendæ sint)}

KYRIE, eléison. Christe, eléison. Kýrie, eléison.

Pater noster \rubric{secreto usque ad} \V\ Et ne nos indúcas in tentatiónem. \R\ Sed líbera nos a malo.

Credo in Deum \rubric{secreto usque ad} \V\ Carnis resurrectiónem. \R\ Vitam ætérnam. Amen.

\V\ Benedíctus es, Dómine, Deus patrum nostrórum. \R\ Et laudábilis et gloriósus in sæcula. \V\ Benedicámus Patrem et Fílium cum Sancto Spíritu. \R\ Laudémus, et superexaltémus eum in sæcula. \V\ Benedíctus es, Dómine, in firmaménto cæli. \R\ Et laudábilis, et gloriósus, et superexaltátus in sæcula. \V\ Benedícat et custódiat nos omnípotens et miséricors Dóminus. \R\ Amen. \V\ Dignáre, Dómine, nocte ista. \R\ Sine peccáto nos custodíre. \V\ Miserére nostri, Dómine. \R\ Miserére nostri. \V\ Fiat misericórdia tua, Dómine, super nos. \R\ Quemádmodum sperávimus in te. \V\ Dómine, exáudi oratiónem meam. \R\ Et clamor meus ad te véniat.

\V\ Dóminus vobíscum.

Orémus. \oratiolabel{Oratio}

\lettrine[lines=2]{V}{ISITA}, quæsumus, Dómine, habitatiónem istam, et omnes insídias inimíci ab ea longe repélle: Angeli tui sancti hábitent in ea, qui nos in pace custódiant; et benedíctio tua sit super nos semper. Per Dóminum nostrum. \R\ Amen.

\V\ Dóminus vobíscum.

\V\ Benedicámus Dómino.

\benedictiolabel{Benedictio} Benedícat et custódiat nos omnípotens et miséricors Dóminus, Pater, et Fílius, et Spíritus Sanctus. \R\ Amen.

\variantheading{Antiphona finalis B. M. V.}

\V\ Divínum auxílium.

Pater, Ave et Credo.

\rubric{¶ Psalmi de Dominica superius appositi dicuntur ad Completorium, adhibita simul Antiphona, etiam post utrasque Vesperas in Dominicis infra Octavas Ssmi Corporis Christi ac Sacrat. Cordis Jesu, in Duplicibus I vel II classis, in Festis novem Lectionum Domini, B. Mariæ Virginis, Angelorum, S. Joannis Baptistæ, S. Joseph, Apostolorum, Evangelistarum, et per Octavas privilegiatas; si tamen præcedentes I Vesperæ de eis integræ vel saltem a Capitulo sint celebratæ, ac II Vesperæ de ipsis integræ fuerint persolutæ.}

% ===== END OF BATCH 8 -- Covers autumnalis.md lines 2478-2863,
% completing "Ad Vesperas" (Psalms 109-113, Capitulum, hymn "Lucis
% Creator optime," and the Magnificat) and "Ad Completorium" (Lectio
% brevis, Confiteor shorthand, Psalms 4, 90, 133, hymn "Te lucis ante
% terminum," Nunc dimittis, and the Preces), and with them the ENTIRE
% "Dominica" Office. Two more occurrences of the recurring stray "\*"
% artifact (in "Lucis Creator optime" at line 2703 and "Te lucis ante
% terminum" at line 2842, both at the hymn's final stanza) are
% silently dropped, consistent with every prior occurrence in this
% volume. The closing "Antiphona finalis B. M. V." is given only as a
% bare cross-reference label with the incipit "℣. Divínum auxílium."
% (not the full "...máneat semper nobíscum. ℟. Amen."), transcribed
% exactly as abbreviated rather than completed from the Ordinarium's
% fuller text. The document continues with "Feria Secunda" (autumnalis.md
% line 2864), the first weekday of the Psalterium cycle =====
\feriaheading{Feria Secunda}

\hourheading{Ad Matutinum}

Pater, Ave et Credo.

\V\ Dómine, lábia mea.

\V\ Deus, in adjutórium.

\variantheading{In Officio feriali}

\rubric{Invitatorium et Hymnus dicuntur quotidie de Feria currenti, ut in Psalterio hic et alibi apponuntur.}

\invitatlabel{Invitat.} Veníte, * Exsultémus Dómino. \psalmlabel{Ps. 94} Veníte, exsultémus.

\rubric{In primo Psalmi versu non dicitur: Veníte, exsultémus Dómino; sed post repetitum Invitatorium, statim subjungitur: Jubilémus Deo, salutári nostro:}

\hymnusheading{Hymnus}

\lettrine[lines=2]{S}{OMNO} reféctis ártubus,\\
Spreto cubíli, súrgimus:\\
Nobis, Pater, canéntibus\\
Adésse te depóscimus.

Te lingua primum cóncinat,\\
Te mentis ardor ámbiat:\\
Ut áctuum sequéntium\\
Tu, Sancte, sis exórdium.

Cedant ténebræ lúmini,\\
Et nox diúrno síderi,\\
Ut culpa, quam nox íntulit,\\
Lucis labáscat múnere.

Precámur iídem súpplices,\\
Noxas ut omnes ámputes,\\
Et ore te canéntium\\
Laudéris omni témpore.

Præsta, Pater piíssime,\\
Patríque compar Unice,\\
Cum Spíritu Paráclito\\
Regnans per omne sæculum.\\
Amen.

\nocturnoheading{In I Nocturno}

\antlabel{Ant.} Dóminus de cælo * prospéxit super fílios hóminum.

\rubric{Antiphonæ, in hac et aliis Feriis appositæ, dicuntur semper in Officio feriali et in Festis sive trium sive novem Lectionum, quando juxta Rubricas sumendi sint Psalmi de Feria, ut in Psalterio.}

\rubric{In Officio autem trium Lectionum, tam de Tempore quam de Sanctis, omittuntur Versus et reliqua omnia, quæ in Officio novem Lectionum solent recitari in fine I et II Nocturni, et continuantur Psalmi cum suis Antiphonis usque ad Versum III Nocturni.}

\psalmheading{Psalmus 13}

\lettrine[lines=2]{D}{IXIT} insípiens in corde suo: * Non est Deus.

Corrúpti sunt, et abominábiles facti sunt in stúdiis suis: * non est qui fáciat bonum, non est usque ad unum.

Dóminus de cælo prospéxit super fílios hóminum, * ut vídeat si est intélligens, aut requírens Deum.

Omnes declinavérunt, simul inútiles facti sunt: * non est qui fáciat bonum, non est usque ad unum.

Sepúlcrum patens est guttur eórum: linguis suis dolóse agébant, * venénum áspidum sub lábiis eórum.

Quorum os maledictióne et amaritúdine plenum est: * velóces pedes eórum ad effundéndum sánguinem.

Contrítio et infelícitas in viis eórum, et viam pacis non cognovérunt: * non est timor Dei ante óculos eórum.

Nonne cognóscent omnes qui operántur iniquitátem, * qui dévorant plebem meam sicut escam panis?

Dóminum non invocavérunt, * illic trepidavérunt timóre, ubi non erat timor.

Quóniam Dóminus in generatióne justa est, consílium ínopis confudístis: * quóniam Dóminus spes ejus est.

Quis dabit ex Sion salutáre Israël? * cum avérterit Dóminus captivitátem plebis suæ, exsultábit Jacob, et lætábitur Israël.

\antlabel{Ant.} Dóminus de cælo prospéxit super fílios hóminum.

\antlabel{Ant.} Qui operátur justítiam, * requiéscet in monte sancto tuo, Dómine.

\psalmheading{Psalmus 14}

\lettrine[lines=2]{D}{OMINE}, quis habitábit in tabernáculo tuo? * aut quis requiéscet in monte sancto tuo?

Qui ingréditur sine mácula, * et operátur justítiam:

Qui lóquitur veritátem in corde suo, * qui non egit dolum in lingua sua:

Nec fecit próximo suo malum, * et oppróbrium non accépit advérsus próximos suos.

Ad níhilum dedúctus est in conspéctu ejus malígnus: * timéntes autem Dóminum gloríficat:

Qui jurat próximo suo, et non décipit, * qui pecúniam suam non dedit ad usúram, et múnera super innocéntem non accépit.

Qui facit hæc, * non movébitur in ætérnum.

\antlabel{Ant.} Qui operátur justítiam, requiéscet in monte sancto tuo, Dómine.

\antlabel{Ant.} Inclína, Dómine, * aurem tuam mihi, et exáudi verba mea.

\psalmheading{Psalmus 16}

\lettrine[lines=2]{E}{XAUDI}, Dómine, justítiam meam: * inténde deprecatiónem meam.

Áuribus pércipe oratiónem meam, * non in lábiis dolósis.

De vultu tuo judícium meum pródeat: * óculi tui vídeant æquitátes.

Probásti cor meum, et visitásti nocte: * igne me examinásti, et non est invénta in me iníquitas.

Ut non loquátur os meum ópera hóminum: * propter verba labiórum tuórum ego custodívi vias duras.

Pérfice gressus meos in sémitis tuis: * ut non moveántur vestígia mea.

Ego clamávi, quóniam exaudísti me, Deus: * inclína aurem tuam mihi, et exáudi verba mea.

Mirífica misericórdias tuas, * qui salvos facis sperántes in te.

A resisténtibus déxteræ tuæ custódi me, * ut pupíllam óculi.

Sub umbra alárum tuárum prótege me: * a fácie impiórum qui me afflixérunt.

Inimíci mei ánimam meam circumdedérunt, ádipem suum conclusérunt: * os eórum locútum est supérbiam.

Projiciéntes me nunc circumdedérunt me: * óculos suos statuérunt declináre in terram.

Suscepérunt me sicut leo parátus ad prǽdam: * et sicut cátulus leónis hábitans in abdítis.

Exsúrge, Dómine, prǽveni eum, et supplánta eum: * éripe ánimam meam ab ímpio, fráméam tuam ab inimícis manus tuæ.

Dómine, a paucis de terra dívide eos in vita eórum: * de abscónditis tuis adimplétus est venter eórum.

Saturáti sunt fíliis: * et dimisérunt relíquias suas párvulis suis.

Ego autem in justítia apparébo conspéctui tuo: * satiábor cum apparúerit glória tua.

\antlabel{Ant.} Inclína, Dómine, aurem tuam mihi, et exáudi verba mea.

\variantheading{In Festis novem Lectionum:}

\V\ Mirífica, Dómine, misericórdias tuas. \R\ Qui salvos facis sperántes in te.

\absolutiolabel{Absol.} Exáudi, Dómine.

\benedlabel{Bened.} 1 Benedictióne perpétua. 2 Unigénitus Dei Fílius. 3 Spíritus Sancti grátia.

\nocturnoheading{In II Nocturno}

\antlabel{Ant.} Díligam te, * \textdagger\ Dómine, virtus mea.

\psalmheading{Psalmus 17, i}

\lettrine[lines=2]{D}{ILIGAM} te, \textdagger\ Dómine, fortitúdo mea: * Dóminus firmaméntum meum, et refúgium meum, et liberátor meus.

Deus meus adjútor meus, * et sperábo in eum.

Protéctor meus, et cornu salútis meæ, * et suscéptor meus.

Laudans invocábo Dóminum: * et ab inimícis meis salvus ero.

Circumdedérunt me dolóres mortis: * et torréntes iniquitátis conturbavérunt me.

Dolóres inférni circumdedérunt me: * præoccupavérunt me láquei mortis.

In tribulatióne mea invocávi Dóminum, * et ad Deum meum clamávi.

Et exaudívit de templo sancto suo vocem meam: * et clamor meus in conspéctu ejus, introívit in aures ejus.

Commóta est, et contrémuit terra: * fundaménta móntium conturbáta sunt, et commóta sunt, quóniam irátus est eis.

Ascéndit fumus in ira ejus: * et ignis a fácie ejus exársit: * carbónes succénsi sunt ab eo.

Inclinávit cælos, et descéndit: * et calígo sub pédibus ejus.

Et ascéndit super Chérubim, et volávit: * volávit super pennas ventórum.

Et pósuit ténebras latíbulum suum, in circúitu ejus tabernáculum ejus: * tenebrósa aqua in núbibus áëris.

Præ fulgóre in conspéctu ejus nubes transiérunt, * grando et carbónes ignis.

Et intónuit de cælo Dóminus, et Altíssimus dedit vocem suam: * grando et carbónes ignis.

Et misit sagíttas suas, et dissipávit eos: * fúlgura multiplicávit, et conturbávit eos.

Et apparuérunt fontes aquárum, * et reveláta sunt fundaménta orbis terrárum:

Ab increpatióne tua, Dómine, * ab inspiratióne spíritus iræ tuæ.

\antlabel{Ant.} Díligam te, Dómine, virtus mea.

\antlabel{Ant.} Retríbuet mihi Dóminus * secúndum justítiam meam.

\psalmheading{Psalmus 17, ii}

\lettrine[lines=2]{M}{ISIT} de summo, et accépit me: * et assúmpsit me de aquis multis.

Erípuit me de inimícis meis fortíssimis, et ab his qui odérunt me: * quóniam confortáti sunt super me.

Prævenérunt me in die afflictiónis meæ: * et factus est Dóminus protéctor meus.

Et edúxit me in latitúdinem: * salvum me fecit, quóniam vóluit me.

Et retríbuet mihi Dóminus secúndum justítiam meam, * et secúndum puritátem mánuum meárum retríbuet mihi:

Quia custodívi vias Dómini, * nec ímpie gessi a Deo meo.

Quóniam ómnia judícia ejus in conspéctu meo: * et justítias ejus non répuli a me.

Et ero immaculátus cum eo: * et observábo me ab iniquitáte mea.

Et retríbuet mihi Dóminus secúndum justítiam meam: * et secúndum puritátem mánuum meárum in conspéctu oculórum ejus.

Cum sancto sanctus eris, * et cum viro innocénte ínnocens eris:

Et cum elécto eléctus eris: * et cum pervérso pervertéris.

Quóniam tu pópulum húmilem salvum fácies: * et óculos superbórum humiliábis.

Quóniam tu illúminas lucérnam meam, Dómine: * Deus meus, illúmina ténebras meas.

Quóniam in te erípiar a tentatióne, * et in Deo meo transgrédiar murum.

Deus meus, impollúta via ejus: elóquia Dómini igne examináta: * protéctor est ómnium sperántium in se.

Quóniam quis Deus præter Dóminum? * aut quis Deus præter Deum nostrum?

Deus, qui præcínxit me virtúte: * et pósuit immaculátam viam meam.

Qui perfécit pedes meos tamquam cervórum, * et super excélsa státuens me.

Qui docet manus meas ad prǽlium: * et posuísti, ut arcum áëreum, bráchia mea.

\antlabel{Ant.} Retríbuet mihi Dóminus secúndum justítiam meam.

\antlabel{Ant.} Vivit Dóminus, * et benedíctus Deus salútis meæ.

\psalmheading{Psalmus 17, iii}

\lettrine[lines=2]{E}{T} dedísti mihi protectiónem salútis tuæ: * et déxtera tua suscépit me:

Et disciplína tua corréxit me in finem: * et disciplína tua ipsa me docébit.

Dilatásti gressus meos subtus me: * et non sunt infirmáta vestígia mea:

Pérsequar inimícos meos, et comprehéndam illos: * et non convértar, donec defíciant.

Confríngam illos, nec póterunt stare: * cadent subtus pedes meos.

Et præcinxísti me virtúte ad bellum: * et supplantásti insurgéntes in me subtus me.

Et inimícos meos dedísti mihi dorsum, * et odiéntes me disperdidísti.

Clamavérunt, nec erat qui salvos fáceret, ad Dóminum: * nec exaudívit eos.

Et commínuam eos, ut púlverem ante fáciem venti: * ut lutum plateárum delébo eos.

Erípies me de contradictiónibus pópuli: * constítues me in caput géntium.

Pópulus, quem non cognóvi, servívit mihi: * in audítu auris obedívit mihi.

Fílii aliéni mentíti sunt mihi: * fílii aliéni inveteráti sunt, et claudicavérunt a sémitis suis.

Vivit Dóminus, et benedíctus Deus meus, * et exaltétur Deus salútis meæ.

Deus, qui das vindíctas mihi, et subdis pópulos sub me, * liberátor meus de inimícis meis iracúndis.

Et ab insurgéntibus in me exaltábis me: * a viro iníquo éripies me.

Proptérea confitébor tibi in natiónibus, Dómine: * et nómini tuo psalmum dicam.

Magníficans salútes Regis ejus, et fáciens misericórdiam Christo suo David, * et sémini ejus usque in séculum.

\antlabel{Ant.} Vivit Dóminus, et benedíctus Deus salútis meæ.

\variantheading{In Festis novem Lectionum:}

\V\ Confitébor tibi in natiónibus, Dómine. \R\ Et nómini tuo psalmum dicam.

\absolutiolabel{Absol.} Ipsíus píetas et.

\benedlabel{Bened.} 4 Deus Pater omnípotens. 5 Christus perpétuæ. 6 Ignem sui amóris.

\nocturnoheading{In III Nocturno}

\antlabel{Ant.} Exáudiat te * \textdagger\ Dóminus in die tribulatiónis. \textdagger

\psalmheading{Psalmus 19}

\lettrine[lines=2]{E}{XAUDIAT} te \textdagger\ Dóminus in die tribulatiónis: * \textdagger\ prótegat te nomen Dei Jacob.

Mittat tibi auxílium de sancto: * et de Sion tueátur te.

Memor sit omnis sacrifícii tui: * et holocáustum tuum pingue fiat.

Tríbuat tibi secúndum cor tuum: * et omne consílium tuum confírmet.

Lætábimur in salutári tuo: * et in nómine Dei nostri magnificábimur.

Ímpleat Dóminus omnes petitiónes tuas: * nunc cognóvi quóniam salvum fecit Dóminus Christum suum.

Exáudiet illum de cælo sancto suo: * in potentátibus salus déxteræ ejus.

Hi in cúrribus, et hi in equis: * nos autem in nómine Dómini, Dei nostri, invocábimus.

Ipsi obligáti sunt, et cecidérunt: * nos autem surréximus et erécti sumus.

Dómine, salvum fac regem: * et exáudi nos in die, qua invocavérimus te.

\antlabel{Ant.} Exáudiat te Dóminus in die tribulatiónis.

\antlabel{Ant.} Dómine, * \textdagger\ in virtúte tua lætábitur rex. \textdaggerdbl

\psalmheading{Psalmus 20}

\lettrine[lines=2]{D}{OMINE}, \textdagger\ in virtúte tua lætábitur rex: * \textdaggerdbl\ et super salutáre tuum exsultábit veheménter.

Desidérium cordis ejus tribuísti ei: * et voluntáte labiórum ejus non fraudásti eum.

Quóniam prævenísti eum in benedictiónibus dulcédinis: * posuísti in cápite ejus corónam de lápide pretióso.

Vitam pétiit a te: * et tribuísti ei longitúdinem diérum in séculum, et in séculum séculi.

Magna est glória ejus in salutári tuo: * glóriam et magnum decórem impónes super eum.

Quóniam dabis eum in benedictiónem in séculum séculi: * lætificábis eum in gáudio cum vultu tuo.

Quóniam rex sperat in Dómino: * et in misericórdia Altíssimi non commovébitur.

Inveniátur manus tua ómnibus inimícis tuis: * déxtera tua invéniat omnes, qui te odérunt.

Pones eos ut clíbanum ignis in témpore vultus tui: * Dóminus in ira sua conturbábit eos, et devorábit eos ignis.

Fructum eórum de terra perdes: * et semen eórum a fíliis hóminum.

Quóniam declinavérunt in te mala: * cogitavérunt consília, quæ non potuérunt stabilíre.

Quóniam pones eos dorsum: * in relíquiis tuis præparábis vultum eórum.

Exaltáre, Dómine, in virtúte tua: * cantábimus et psallémus virtútes tuas.

\antlabel{Ant.} Dómine, in virtúte tua lætábitur rex.

\antlabel{Ant.} Exaltábo te, * \textdagger\ Dómine, quóniam suscepísti me. \textdaggerdbl

\psalmheading{Psalmus 29}

\lettrine[lines=2]{E}{XALTABO} te, \textdagger\ Dómine, quóniam suscepísti me: * \textdaggerdbl\ nec delectásti inimícos meos super me.

Dómine, Deus meus, clamávi ad te, * et sanásti me.

Dómine, eduxísti ab inférno ánimam meam: * salvásti me a descendéntibus in lacum.

Psállite Dómino, sancti ejus: * et confitémini memóriæ sanctitátis ejus.

Quóniam ira in indignatióne ejus: * et vita in voluntáte ejus.

Ad vésperum demorábitur fletus: * et ad matutínum lætítia.

Ego autem dixi in abundántia mea: * Non movébor in ætérnum.

Dómine, in voluntáte tua, * præstitísti decóri meo virtútem.

Avertísti fáciem tuam a me, * et factus sum conturbátus.

Ad te, Dómine, clamábo: * et ad Deum meum deprecábor.

Quæ utílitas in sánguine meo, * dum descéndo in corruptiónem?

Numquid confitébitur tibi pulvis, * aut annuntiábit veritátem tuam?

Audívit Dóminus, et misértus est mei: * Dóminus factus est adjútor meus.

Convertísti planctum meum in gáudium mihi: * conscidísti saccum meum, et circumdedísti me lætítia:

Ut cantet tibi glória mea, et non compúngar: * Dómine, Deus meus, in ætérnum confitébor tibi.

\antlabel{Ant.} Exaltábo te, Dómine, quóniam suscepísti me.

\variantheading{In feriali Officio et in Festis:}

\V\ Psállite Dómino, sancti ejus. \R\ Et confitémini memóriæ sanctitátis ejus.

\variantheading{In Officio novem Lectionum}

\absolutiolabel{Absolutio} A vínculis.

\benedlabel{Bened.} 7 Evangélica léctio. 8 In Festis Domini: Divínum auxílium. In Festis B. M. V.: Cujus festum cólimus, ipsa Virgo. In Festis Sanctorum: Cujus (\textit{vel} Quorum \textit{aut} Quarum) festum cólimus. 9 Ad societátem. (Per evangélica dicta.)

\variantheading{In Officio trium Lectionum}

\absolutiolabel{Absol.} Exáudi, Dómine.

\variantheading{In Feriis, cum Homilia:}

\benedlabel{Bened.} 1 Evangélica. 2 Divínum. 3 Ad societátem.

\variantheading{In Feriis, sine Homilia:}

\benedlabel{Bened.} 1 Benedictióne. 2 Unigénitus. 3 Spíritus Sancti.

\variantheading{In Festis Sanctorum:}

\benedlabel{Bened.} 1 Ille nos benedícat qui sine fine vivit et regnat. \R\ Amen. 2 Cujus (\textit{vel} Quorum \textit{aut} Quarum). 3 Ad societátem.

% ===== END OF BATCH 9 -- Covers autumnalis.md lines 2864-3269, opening
% "Feria Secunda" (Monday) and completing its "Ad Matutinum" in
% full: the Invitatorium "Veníte, Exsultémus Dómino," the hymn
% "Somno reféctis artubus," and the three Nocturns (Psalms 13, 14,
% 16; 17/i, 17/ii, 17/iii; 19, 20, 29), each followed by its
% condition-specific versicle/Absolutio/Benedictiones apparatus
% (the ferial psalter prints these only "In Festis novem
% Lectionum"/"In Officio novem Lectionum" etc., since on an ordinary
% ferial day with three lessons only, most of this apparatus is
% omitted per the rubric given at the start of this Nocturnal set).
% Several antiphons and psalm-incipits in the Second and Third
% Nocturns carry multiple chant-pointing marks in a single line --
% asterisk, dagger, and even double-dagger together (e.g. "Dómine, *
% † in virtúte tua lætábitur rex. ‡") -- each rendered in its exact
% source position as *, \textdagger\, and \textdaggerdbl\
% respectively, with no marks added, dropped, or reordered. The
% compact multi-option "Bened." line for the ninth lesson ("8 In
% Festis Domini: Divínum auxílium. In Festis B. M. V.: ... In Festis
% Sanctorum: ...") is transcribed as one continuous \benedlabel{}
% argument exactly as the source runs it together. The document
% continues with "Ad Laudes I" (autumnalis.md line 3270), the first of
% two alternate ferial Lauds schemes =====
\hourheading{Ad Laudes I}

\rubric{In Feriis extra Vigílias commúnes occurréntibus, atque in Festis.}

(Pater noster et Ave.)

\V\ Deus, in adjutórium.

\antlabel{Ant.} Jubiláte * Deo in voce exsultatiónis.

\rubric{¶ Antíphonæ, in hac et áliis Fériis apposítæ, dicúntur semper in Officio feriáli et in Festis sive trium sive novem Lectiónum, quando juxta Rúbricas suméndi sint Psalmi de Féria.}

\psalmheading{Psalmus 46}

\lettrine[lines=2]{O}{MNES} gentes, pláudite mánibus: * jubiláte Deo in voce exsultatiónis.

Quóniam Dóminus excélsus, terríbilis: * Rex magnus super omnem terram.

Subjécit pópulos nobis: * et gentes sub pédibus nostris.

Elégit nobis hereditátem suam: * spéciem Jacob, quam diléxit.

Ascéndit Deus in júbilo: * et Dóminus in voce tubæ.

Psállite Deo nostro, psállite: * psállite Regi nostro, psállite.

Quóniam Rex omnis terræ Deus: * psállite sapiénter.

Regnábit Deus super gentes: * Deus sedet super sedem sanctam suam.

Príncipes populórum congregáti sunt cum Deo Abraham: * quóniam dii fortes terræ veheménter eleváti sunt.

\antlabel{Ant.} Jubiláte Deo in voce exsultatiónis.

\antlabel{Ant.} Inténde * voci oratiónis meæ, Rex meus et Deus meus.

\psalmheading{Psalmus 5}

\lettrine[lines=2]{V}{ERBA} mea áuribus pércipe, Dómine, * intéllige clamórem meum.

Inténde voci oratiónis meæ, * Rex meus et Deus meus.

Quóniam ad te orábo: * Dómine, mane exáudies vocem meam.

Mane astábo tibi et vidébo: * quóniam non Deus volens iniquitátem tu es.

Neque habitábit juxta te malígnus: * neque permanébunt injústi ante óculos tuos.

Odísti omnes, qui operántur iniquitátem: * perdes omnes, qui loquúntur mendácium.

Virum sánguinum et dolósum abominábitur Dóminus: * ego autem in multitúdine misericórdiæ tuæ.

Introíbo in domum tuam: * adorábo ad templum sanctum tuum in timóre tuo.

Dómine, deduc me in justítia tua: * propter inimícos meos dírige in conspéctu tuo viam meam.

Quóniam non est in ore eórum véritas: * cor eórum vanum est.

Sepúlcrum patens est guttur eórum, linguis suis dolóse agébant, * júdica illos, Deus.

Décidant a cogitatiónibus suis, secúndum multitúdinem impietátum eórum expélle eos, * quóniam irritavérunt te, Dómine.

Et læténtur omnes, qui sperant in te, * in ætérnum exsultábunt: et habitábis in eis.

Et gloriabúntur in te omnes, qui díligunt nomen tuum, * quóniam tu benedíces justo.

Dómine, ut scuto bonæ voluntátis tuæ * coronásti nos.

\antlabel{Ant.} Inténde voci oratiónis meæ, Rex meus et Deus meus.

\antlabel{Ant.} Deus majestátis * intónuit: afférte glóriam nómini ejus.

\psalmheading{Psalmus 28}

\lettrine[lines=2]{A}{FFERTE} Dómino, fílii Dei: * afférte Dómino fílios aríetum.

Afférte Dómino glóriam et honórem, afférte Dómino glóriam nómini ejus: * adoráte Dóminum in átrio sancto ejus.

Vox Dómini super aquas, Deus majestátis intónuit: * Dóminus super aquas multas.

Vox Dómini in virtúte: * vox Dómini in magnificéntia.

Vox Dómini confringéntis cedros: * et confrínget Dóminus cedros Líbani:

Et commínuet eas tamquam vítulum Líbani: * et diléctus quemádmodum fílius unicórnium.

Vox Dómini intercidéntis flammam ignis: * vox Dómini concutiéntis desértum: et commovébit Dóminus desértum Cades.

Vox Dómini præparántis cervos, et revelábit condénsa: * et in templo ejus omnes dicent glóriam.

Dóminus dilúvium inhabitáre facit: * et sedébit Dóminus Rex in ætérnum.

Dóminus virtútem pópulo suo dabit: * Dóminus benedícet pópulo suo in pace.

\antlabel{Ant.} Deus majestátis intónuit: afférte glóriam nómini ejus.

\antlabel{Ant.} Laudámus nomen tuum * ínclitum, Deus noster.

\canticumheading{Canticum David}{1 Par. 29, 10-13}

\lettrine[lines=2]{B}{ENEDICTUS} es, Dómine, Deus Israël patris nostri, * ab ætérno in ætérnum.

Tua est, Dómine, magnificéntia, et poténtia, * et glória, atque victória:

Et tibi laus: * cuncta enim quæ in cælo sunt, et in terra, tua sunt:

Tuum, Dómine, regnum, * et tu es super omnes príncipes.

Tuæ divítiæ, et tua est glória: * tu domináris ómnium.

In manu tua virtus et poténtia: * in manu tua magnitúdo, et impérium ómnium.

Nunc ígitur, Deus noster, confitémur tibi, * et laudámus nomen tuum ínclitum.

\antlabel{Ant.} Laudámus nomen tuum ínclitum, Deus noster.

\antlabel{Ant.} Laudáte * \textdagger\ Dóminum, omnes gentes. \textdaggerdbl

\psalmheading{Psalmus 116}

\lettrine[lines=2]{L}{AUDATE} \textdagger\ Dóminum, omnes gentes: * \textdaggerdbl\ laudáte eum, omnes pópuli:

Quóniam confirmáta est super nos misericórdia ejus: * et véritas Dómini manet in ætérnum.

\antlabel{Ant.} Laudáte Dóminum, omnes gentes.

\variantheading{In Officio feriali}

\rubric{Capítulum, Hymnus et Versus dicúntur quotídie de Féria curránti, ut in Psaltério hic et álibi appónuntur; et inde súmitur étiam Ant. ad Benedíctus.}

\capitulumlabel{Capit.}{Rom. 13, 12-13}

\lettrine[lines=2]{N}{OX} præcéssit, dies autem appropinquávit. Abjiciámus ergo ópera tenebrárum, et induámur arma lucis. Sicut in die honéste ambulémus.

\hymnusheading{Hymnus}

\lettrine[lines=2]{S}{PLENDOR} Patérnæ glóriæ,\\
De luce lucem próferens,\\
Lux lucis, et fons lúminis,\\
Diem dies illúminans:

Verúsque sol illábere,\\
Micans nitóre pérpeti:\\
Jubárque Sancti Spíritus\\
Infúnde nostris sénsibus.

Votis vocémus et Patrem,\\
Patrem poténtis grátiæ,\\
Patrem perénnis glóriæ:\\
Culpam reléget lúbricam.

Confírmet actus strénuos:\\
Dentes retúndat ínvidi:\\
Casus secúndet ásperos:\\
Agénda recte dírigat.

Mentem gubérnet et regat:\\
Sit pura nobis cástitas:\\
Fides calóre férveat,\\
Fraudis venéna nésciat.

Christúsque nobis sit cibus,\\
Potúsque noster sit fides:\\
Læti bibámus sóbriam\\
Profusiónem Spíritus.

Lætus dies hic tránseat:\\
Pudor sit ut dilúculum:\\
Fides velut merídies:\\
Crepúsculum mens nésciat.

Auróra lucem próvehit,\\
Cum luce nobis pródeat\\
In Patre totus Fílius,\\
Et totus in Verbo Pater.

Deo Patri sit glória,\\
Ejúsque soli Fílio,\\
Cum Spíritu Paráclito,\\
Nunc et per omne sæculum. Amen.

\V\ Repléti sumus mane misericórdia tua. \R\ Exsultávimus, et delectáti sumus.

\antlabel{Ad Bened. Ant.} Benedíctus * \textdagger\ Dóminus, Deus Israël, quia visitávit et liberávit nos.

\rubric{Réliqua ut in Ordinário. 9}

\hourheading{Ad Laudes II}

\rubric{In Vigíliis commúnibus, quando fit Offícium de Féria.}

(Pater noster et Ave.)

\V\ Deus, in adjutórium.

\antlabel{Ant.} Miserére. \textdagger

\psalmheading{Psalmus 50}

\lettrine[lines=2]{M}{ISERERE} \textdagger\ mei, Deus: * secúndum magnam misericórdiam tuam.

Et secúndum multitúdinem miseratiónum tuárum, * dele iniquitátem meam.

Amplius lava me ab iniquitáte mea: * et a peccáto meo munda me.

Quóniam iniquitátem meam ego cognósco: * et peccátum meum contra me est semper.

Tibi soli peccávi, et malum coram te feci: * ut justificéris in sermónibus tuis, et vincas cum judicáris.

Ecce enim, in iniquitátibus concéptus sum: * et in peccátis concépit me mater mea.

Ecce enim, veritátem dilexísti: * incérta et occúlta sapiéntiæ tuæ manifestásti mihi.

Aspérges me hyssópo, et mundábor: * lavábis me, et super nivem dealbábor.

Audítui meo dabis gáudium et lætítiam: * et exsultábunt ossa humiliáta.

Avérte fáciem tuam a peccátis meis: * et omnes iniquitátes meas dele.

Cor mundum crea in me, Deus: * et spíritum rectum ínnova in viscéribus meis.

Ne projícias me a fácie tua: * et spíritum sanctum tuum ne áuferas a me.

Redde mihi lætítiam salutáris tui: * et spíritu principáli confírma me.

Docébo iníquos vias tuas: * et ímpii ad te converténtur.

Líbera me de sanguínibus, Deus, Deus salútis meæ: * et exsultábit lingua mea justítiam tuam.

Dómine, lábia mea apéries: * et os meum annuntiábit laudem tuam.

Quóniam si voluísses sacrifícium, dedíssem útique: * holocáustis non delectáberis.

Sacrifícium Deo spíritus contribulátus: * cor contrítum, et humiliátum, Deus, non despícies.

Benígne fac, Dómine, in bona voluntáte tua Sion: * ut ædificéntur muri Jerúsalem.

Tunc acceptábis sacrifícium justítiæ, oblatiónes, et holocáusta: * tunc impónent super altáre tuum vítulos.

\antlabel{Ant.} Miserére mei, Deus, secúndum magnam misericórdiam tuam.

\antlabel{Ant.} Deduc me.

\psalmheading{Psalmus 5}

\lettrine[lines=2]{V}{ERBA} mea áuribus pércipe, Dómine, * intéllige clamórem meum.

Inténde voci oratiónis meæ, * Rex meus et Deus meus.

Quóniam ad te orábo: * Dómine, mane exáudies vocem meam.

Mane astábo tibi et vidébo: * quóniam non Deus volens iniquitátem tu es.

Neque habitábit juxta te malígnus: * neque permanébunt injústi ante óculos tuos.

Odísti omnes, qui operántur iniquitátem: * perdes omnes, qui loquúntur mendácium.

Virum sánguinum et dolósum abominábitur Dóminus: * ego autem in multitúdine misericórdiæ tuæ.

Introíbo in domum tuam: * adorábo ad templum sanctum tuum in timóre tuo.

Dómine, deduc me in justítia tua: * propter inimícos meos dírige in conspéctu tuo viam meam.

Quóniam non est in ore eórum véritas: * cor eórum vanum est.

Sepúlcrum patens est guttur eórum, linguis suis dolóse agébant, * júdica illos, Deus.

Décidant a cogitatiónibus suis, secúndum multitúdinem impietátum eórum expélle eos, * quóniam irritavérunt te, Dómine.

Et læténtur omnes, qui sperant in te, * in ætérnum exsultábunt: et habitábis in eis.

Et gloriabúntur in te omnes, qui díligunt nomen tuum, * quóniam tu benedíces justo.

Dómine, ut scuto bonæ voluntátis tuæ * coronásti nos.

\antlabel{Ant.} Deduc me in justítia tua, Dómine.

\antlabel{Ant.} Dóminus dabit virtútem.

\psalmheading{Psalmus 28}

\lettrine[lines=2]{A}{FFERTE} Dómino, fílii Dei: * afférte Dómino fílios aríetum.

Afférte Dómino glóriam et honórem, afférte Dómino glóriam nómini ejus: * adoráte Dóminum in átrio sancto ejus.

Vox Dómini super aquas, Deus majestátis intónuit: * Dóminus super aquas multas.

Vox Dómini in virtúte: * vox Dómini in magnificéntia.

Vox Dómini confringéntis cedros: * et confrínget Dóminus cedros Líbani:

Et commínuet eas tamquam vítulum Líbani: * et diléctus quemádmodum fílius unicórnium.

Vox Dómini intercidéntis flammam ignis: * vox Dómini concutiéntis desértum: et commovébit Dóminus desértum Cades.

Vox Dómini præparántis cervos, et revelábit condénsa: * et in templo ejus omnes dicent glóriam.

Dóminus dilúvium inhabitáre facit: * et sedébit Dóminus Rex in ætérnum.

Dóminus virtútem pópulo suo dabit: * Dóminus benedícet pópulo suo in pace.

\antlabel{Ant.} Dóminus dabit virtútem et benedícet pópulo suo in pace.

\antlabel{Ant.} Convérsus est furor tuus.

\canticumheading{Canticum Isaiæ}{Is. 12, 1-6}

\lettrine[lines=2]{C}{ONFITEBOR} tibi, Dómine, quóniam irátus es mihi: * convérsus est furor tuus, et consolátus es me.

Ecce Deus salvátor meus, * fiduciáliter agam, et non timébo:

Quia fortitúdo mea, et laus mea Dóminus, * et factus est mihi in salútem.

Hauriétis aquas in gáudio de fóntibus Salvatóris: * et dicétis in die illa: Confitémini Dómino, et invocáte nomen ejus:

Notas fácite in pópulis adinventiónes ejus: * mementóte quóniam excélsum est nomen ejus.

Cantáte Dómino quóniam magnífice fecit: * annuntiáte hoc in univérsa terra.

Exsúlta, et lauda, habitátio Sion: * quia magnus in médio tui Sanctus Israël.

\antlabel{Ant.} Convérsus est furor tuus, Dómine, et consolátus es me.

\antlabel{Ant.} Laudáte. \textdagger

\psalmheading{Psalmus 116}

\lettrine[lines=2]{L}{AUDATE} \textdagger\ Dóminum, omnes gentes: * \textdaggerdbl\ laudáte eum, omnes pópuli:

Quóniam confirmáta est super nos misericórdia ejus: * et véritas Dómini manet in ætérnum.

\antlabel{Ant.} Laudáte Dóminum, quóniam confirmáta est super nos misericórdia ejus.

\capitulumlabel{Capit.}{Rom. 13, 12-13}

\lettrine[lines=2]{N}{OX} præcéssit, dies autem appropinquávit. Abjiciámus ergo ópera tenebrárum, et induámur arma lucis. Sicut in die honéste ambulémus.

\hymnusheading{Hymnus}

\lettrine[lines=2]{S}{PLENDOR} Patérnæ glóriæ,\\
De luce lucem próferens,\\
Lux lucis, et fons lúminis,\\
Diem dies illúminans:

Verúsque sol illábere,\\
Micans nitóre pérpeti:\\
Jubárque Sancti Spíritus\\
Infúnde nostris sénsibus.

Votis vocémus et Patrem,\\
Patrem poténtis grátiæ,\\
Patrem perénnis glóriæ:\\
Culpam reléget lúbricam.

Confírmet actus strénuos:\\
Dentes retúndat ínvidi:\\
Casus secúndet ásperos:\\
Agénda recte dírigat.

Mentem gubérnet et regat:\\
Sit pura nobis cástitas:\\
Fides calóre férveat,\\
Fraudis venéna nésciat.

Christúsque nobis sit cibus,\\
Potúsque noster sit fides:\\
Læti bibámus sóbriam\\
Profusiónem Spíritus.

Lætus dies hic tránseat:\\
Pudor sit ut dilúculum:\\
Fides velut merídies:\\
Crepúsculum mens nésciat.

Auróra lucem próvehit,\\
Cum luce nobis pródeat\\
In Patre totus Fílius,\\
Et totus in Verbo Pater.

Deo Patri sit glória,\\
Ejúsque soli Fílio,\\
Cum Spíritu Paráclito,\\
Nunc et per omne sæculum. Amen.

\V\ Repléti sumus mane misericórdia tua. \R\ Exsultávimus, et delectáti sumus.

\antlabel{Ad Bened. Ant.} Benedíctus * \textdagger\ Dóminus, Deus Israël, quia visitávit et liberávit nos.

\variantheading{Preces feriales}

KYRIE, eléison. Christe, eléison. Kýrie, eléison.

Pater noster \rubric{quod dícitur a solo Hebdomadário totum clara voce, usque ad} \V\ Et ne nos indúcas in tentatiónem. \R\ Sed líbera nos a malo. \V\ Ego dixi: Dómine, miserére mei. \R\ Sana ánimam meam, quia peccávi tibi. \V\ Convértere, Dómine, úsquequo? \R\ Et deprecábilis esto super servos tuos. \V\ Fiat misericórdia tua, Dómine, super nos. \R\ Quemádmodum sperávimus in te. \V\ Sacerdótes tui induántur justítiam. \R\ Et sancti tui exsúltent. \V\ Orémus pro beatíssimo Papa nostro N. \R\ Dóminus consérvet eum, et vivíficet eum, et beátum fáciat eum in terra, et non tradat eum in ánimam inimicórum ejus. \V\ Orémus et pro Antístite nostro N. \R\ Stet et pascat in fortitúdine tua, Dómine, in sublimitáte nóminis tui. \V\ Dómine, salvum fac regem. \R\ Et exáudi nos in die, qua invocavérimus te. \V\ Salvum fac pópulum tuum, Dómine, et bénedic hereditáti tuæ. \R\ Et rege eos, et extólle illos usque in ætérnum. \V\ Meménto Congregatiónis tuæ. \R\ Quam possedísti ab inítio. \V\ Fiat pax in virtúte tua. \R\ Et abundántia in túrribus tuis. \V\ Orémus pro benefactóribus nostris. \R\ Retribúere dignáre, Dómine, ómnibus, nobis bona faciéntibus propter nomen tuum, vitam ætérnam. Amen. \V\ Orémus pro fidélibus defúnctis. \R\ Réquiem ætérnam dona eis, Dómine, et lux perpétua lúceat eis. \V\ Requiéscant in pace. \R\ Amen. \V\ Pro frátribus nostris abséntibus. \R\ Salvos fac servos tuos, Deus meus, sperántes in te. \V\ Pro afflíctis et captívis. \R\ Líbera eos, Deus Israël, ex ómnibus tribulatiónibus suis. \V\ Mitte eis, Dómine, auxílium de sancto. \R\ Et de Sion tuére eos. \V\ Dómine, Deus virtútum, convérte nos. \R\ Et osténde fáciem tuam, et salvi érimus. \V\ Exsúrge, Christe, ádjuva nos. \R\ Et líbera nos propter nomen tuum. \V\ Dómine, exáudi oratiónem meam. \R\ Et clamor meus ad te véniat.

\rubric{Réliqua ut in Ordinário. 9}

% ===== END OF BATCH 10 -- Covers autumnalis.md lines 3270-3672, the
% two alternate ferial Lauds schemes, "Ad Laudes I" (for ordinary
% ferias and feasts) and "Ad Laudes II" (for common Vigils when the
% Office of the Feria is said). Both share the same five psalms
% (46, 5, 28, Canticum, 116) with different antiphon sets, the same
% Capitulum "Nox præcéssit," and the identical hymn "Splendor
% Paternæ glóriæ" -- all transcribed fresh in both places per the
% governing rule, even though the Hymnus and Capitulum text is
% word-for-word identical between the two schemes in autumnalis.md
% itself. Ad Laudes II additionally carries the "Preces feriales"
% litany, transcribed in full; "Kyrie, eléison" at its head is left
% as plain text with no lettrine, matching precedent in verna.tex
% and hiemalis.tex, since it is not one of the five lettrine genres.
% The document continues with "Ad Primam" (autumnalis.md line 3674).
% =====
\hourheading{Ad Primam}

Pater, Ave et Credo.

\V\ Deus, in adjutórium.

\hymnusheading{Hymnus}

\lettrine[lines=2]{J}{AM} lucis orto sídere,\\
Deum precémur súpplices,\\
Ut in diúrnis áctibus\\
Nos servet a nocéntibus.

Linguam refrénans témperet,\\
Ne litis horror ínsonet:\\
Visum fovéndo cóntegat,\\
Ne vanitátes háuriat.

Sint pura cordis íntima,\\
Absístat et vecórdia:\\
Carnis terat supérbiam\\
Potus cibíque párcitas.

Ut, cum dies abscésserit\\
Noctémque sors redúxerit,\\
Mundi per abstinéntiam\\
Ipsi canámus glóriam.

Deo Patri sit glória,\\
Ejúsque soli Fílio,\\
Cum Spíritu Paráclito,\\
Nunc et per omne sæculum.\\
Amen.

\antlabel{Ant.} Innocens mánibus.

\rubric{¶ Antiphonæ, in hac et áliis Fériis ad Primam, Tértiam, Sextam et Nonam apposítæ, dicúntur semper in Officio feriáli et in Festis sive trium sive novem Lectiónum, quando juxta Rúbricas suméndi sint Psalmi de Féria.}

\psalmheading{Psalmus 23}

\lettrine[lines=2]{D}{OMINI} est terra, et plenitúdo ejus: * orbis terrárum, et univérsi qui hábitant in eo.

Quia ipse super mária fundávit eum: * et super flúmina præparávit eum.

Quis ascéndet in montem Dómini? * aut quis stabit in loco sancto ejus?

Innocens mánibus et mundo corde, * qui non accépit in vano ánimam suam, nec jurávit in dolo próximo suo.

Hic accípiet benedictiónem a Dómino: * et misericórdiam a Deo, salutári suo.

Hæc est generátio quæréntium eum, * quæréntium fáciem Dei Jacob.

Attóllite portas, príncipes, vestras, et elevámini, portæ æternáles: * et introíbit Rex glóriæ.

Quis est iste Rex glóriæ? * Dóminus fortis et potens: Dóminus potens in prœlio.

Attóllite portas, príncipes, vestras, et elevámini, portæ æternáles: * et introíbit Rex glóriæ.

Quis est iste Rex glóriæ? * Dóminus virtútum ipse est Rex glóriæ.

\psalmheading{Psalmus 18, i}

\lettrine[lines=2]{C}{ÆLI} enárrant glóriam Dei, * et ópera mánuum ejus annúntiat firmaméntum.

Dies diéi erúctat verbum, * et nox nocti índicat sciéntiam.

Non sunt loquélæ, neque sermónes, * quorum non audiántur voces eórum.

In omnem terram exívit sonus eórum: * et in fines orbis terræ verba eórum.

In sole pósuit tabernáculum suum: * et ipse tamquam sponsus procédens de thálamo suo:

Exsultávit ut gigas ad curréndam viam, * a summo cælo egréssio ejus:

Et occúrsus ejus usque ad summum ejus: * nec est qui se abscóndat a calóre ejus.

\psalmheading{Psalmus 18, ii}

\lettrine[lines=2]{L}{EX} Dómini immaculáta, convértens ánimas: * testimónium Dómini fidéle, sapiéntiam præstans párvulis.

Justítiæ Dómini rectæ, lætificántes corda: * præcéptum Dómini lúcidum, illúminans óculos.

Timor Dómini sanctus, pérmanens in sæculum sæculi: * judícia Dómini vera, justificáta in semetípsa.

Desiderabília super aurum et lápidem pretiósum multum: * et dulcióra super mel et favum.

Étenim servus tuus custódit ea: * in custodiéndis illis retribútio multa.

Delícta quis intélligit? ab occúltis meis munda me: * et ab aliénis parce servo tuo.

Si mei non fúerint domináti, tunc immaculátus ero: * et emundábor a delícto máximo.

Et erunt ut compláceant elóquia oris mei: * et meditátio cordis mei in conspéctu tuo semper.

Dómine, adjútor meus, * et redémptor meus.

\rubric{¶ Quando ad Laudes dictus fúerit Ps. 50 Miserére, hic subjúngitur sequens Ps. 46, qui in Laudibus locum non hábuit.}

\psalmheading{Psalmus 46}

\lettrine[lines=2]{O}{MNES} gentes, pláudite mánibus: * jubiláte Deo in voce exsultatiónis.

Quóniam Dóminus excélsus, terríbilis: * Rex magnus super omnem terram.

Subjécit pópulos nobis: * et gentes sub pédibus nostris.

Elégit nobis hereditátem suam: * spéciem Jacob, quam diléxit.

Ascéndit Deus in júbilo: * et Dóminus in voce tubæ.

Psállite Deo nostro, psállite: * psállite Regi nostro, psállite.

Quóniam Rex omnis terræ Deus: * psállite sapiénter.

Regnábit Deus super gentes: * Deus sedet super sedem sanctam suam.

Príncipes populórum congregáti sunt cum Deo Abraham: * quóniam dii fortes terræ veheménter eleváti sunt.

\rubric{Secus vero, tribus Psalmis explétis, statim dícitur Ant.}

\antlabel{Ant.} Innocens mánibus et mundo corde ascéndet in montem Dómini.

\rubric{In Officio Festi vel Octávæ:}

\capitulumlabel{Capitulum}{1 Tim. 1, 17}

\lettrine[lines=2]{R}{EGI} sæculórum immortáli et invisíbili, soli Deo honor et glória in sæcula sæculórum. Amen.

\rubric{In ómnibus autem Fériis et in Vigíliis commúnibus:}

\capitulumlabel{Capitulum}{Zach. 8, 19}

\lettrine[lines=2]{P}{ACEM} et veritátem dilígite, ait Dóminus omnípotens.

\R\ br. Christe, Fili Dei vivi, * Miserére nobis. Christe. \V\ \textit{Qui sedes ad déxteram Patris.} Miserére nobis. Glória Patri. Christe. \V\ Exsúrge, Christe, ádjuva nos. \R\ Et líbera nos propter nomen tuum.

\variantheading{Preces dominicales et feriales (si dicendæ sint)}

KYRIE, eléison. Christe, eléison. Kýrie, eléison.

Pater noster \rubric{secreto usque ad} \V\ Et ne nos indúcas in tentatiónem. \R\ Sed líbera nos a malo.

Credo in Deum \rubric{secreto usque ad} \V\ Carnis resurrectiónem. \R\ Vitam ætérnam. Amen. \V\ Et ego ad te, Dómine, clamávi. \R\ Et mane orátio mea prævéniet te. \V\ Repleátur os meum laude. \R\ Ut cantem glóriam tuam, tota die magnitúdinem tuam. \V\ Dómine, avérte fáciem tuam a peccátis meis. \R\ Et omnes iniquitátes meas dele. \V\ Cor mundum crea in me, Deus. \R\ Et spíritum rectum ínnova in viscéribus meis. \V\ Ne projícias me a fácie tua. \R\ Et spíritum sanctum tuum ne áuferas a me. \V\ Redde mihi lætítiam salutáris tui. \R\ Et spíritu principáli confírma me.

\variantheading{In Précibus feriálibus tantum}

\V\ Eripe me, Dómine, ab hómine malo. \R\ A viro iníquo éripe me. \V\ Eripe me de inimícis meis, Deus meus. \R\ Et ab insurgéntibus in me líbera me. \V\ Eripe me de operántibus iniquitátem. \R\ Et de viris sánguinum salva me. \V\ Sic psalmum dicam nómini tuo in sæculum sæculi. \R\ Ut reddam vota mea de die in diem. \V\ Exáudi nos, Deus, salutáris noster. \R\ Spes ómnium fínium terræ, et in mari longe. \V\ Deus, in adjutórium meum inténde. \R\ Dómine, ad adjuvándum me festína. \V\ Sanctus Deus, Sanctus fortis, Sanctus immortális. \R\ Miserére nobis. \V\ Bénedic, ánima mea, Dómino. \R\ Et ómnia, quæ intra me sunt, nómini sancto ejus. \V\ Bénedic, ánima mea, Dómino. \R\ Et noli oblivísci omnes retributiónes ejus. \V\ Qui propitiátur ómnibus iniquitátibus tuis. \R\ Qui sanat omnes infirmitátes tuas. \V\ Qui rédimit de intéritu vitam tuam. \R\ Qui corónat te in misericórdia et miseratiónibus. \V\ Qui replet in bonis desidérium tuum. \R\ Renovábitur ut áquilæ juvéntus tua.

\variantheading{In utrisque Précibus}

\V\ Adjutórium nostrum in nómine Dómini. \R\ Qui fecit cælum et terram.

\formulalabel{Confíteor. Misereátur. Indulgéntiam.}

\V\ Dignáre, Dómine, die isto. \R\ Sine peccáto nos custodíre. \V\ Miserére nostri, Dómine. \R\ Miserére nostri. \V\ Fiat misericórdia tua, Dómine, super nos. \R\ Quemádmodum sperávimus in te. \V\ Dómine, exáudi oratiónem meam. \R\ Et clamor meus ad te véniat.

\V\ Dóminus vobíscum.

Orémus. \oratiolabel{Oratio}

\lettrine[lines=2]{D}{OMINE}, Deus omnípotens, qui ad princípium hujus diéi nos perveníre fecísti: tua nos hódie salva virtúte; ut in hac die ad nullum declinémus peccátum, sed semper ad tuam justítiam faciéndam nostra procédant elóquia, dirigántur cogitatiónes et ópera. Per Dóminum. \R\ Amen.

\V\ Dóminus vobíscum. \V\ Benedicámus Dómino.

\rubric{In Choro Martyrológium.}

\V\ Pretiósa in conspéctu Dómini. \R\ Mors Sanctórum ejus.

\rubric{Et absolúte, sine} Orémus:

\oratiolabel{Oratio}

\lettrine[lines=2]{S}{ANCTA} María, et omnes Sancti intercédant pro nobis ad Dóminum, ut nos mereámur ab eo adjuvári et salvári, qui vivit et regnat in sæcula sæculórum. \R\ Amen.

\V\ Deus, in adjutórium meum inténde. \R\ Dómine, ad adjuvándum me festína. \rubric{Et dícitur ter; último additur:} Glória Patri. Sicut erat.

Kýrie, eléison. Christe, eléison. Kýrie, eléison.

Pater noster \rubric{secreto usque ad} \V\ Et ne nos indúcas in tentatiónem. \R\ Sed líbera nos a malo. \V\ Réspice in servos tuos, Dómine, et in ópera tua, et dírige fílios eórum. \R\ Et sit splendor Dómini, Dei nostri, super nos, et ópera mánuum nostrárum dírige super nos, et opus mánuum nostrárum dírige. \V\ Glória Patri. \R\ Sicut erat.

Orémus. \oratiolabel{Oratio}

\lettrine[lines=2]{D}{IRIGERE} et sanctificáre, régere et gubernáre dignáre, Dómine Deus, Rex cæli et terræ, hódie corda et córpora nostra, sensus, sermónes et actus nostros in lege tua, et in opéribus mandatórum tuórum: ut hic, et in ætérnum, te auxiliánte, salvi et líberi esse mereámur, Salvátor mundi: Qui vivis et regnas in sæcula sæculórum. \R\ Amen.

\V\ Jube, domne.

\benedlabel{Bened.} Dies et actus nostros in sua pace dispónat Dóminus omnípotens. \R\ Amen.

\variantheading{In Officio feriali}

\lettrine[lines=2]{D}{OMINUS} autem dírigat corda et córpora nostra in caritáte Dei, et patiéntia Christi. Tu autem.

\variantheading{E Proprio Sanctorum}

\variantheading{Pro Oct. Omn. Sanctorum}

\lettrine[lines=2]{B}{ENEDÍCTIO} et cláritas et sapiéntia et gratiárum áctio, honor, virtus et fortitúdo Deo nostro in sǽcula sæculórum. Amen.

\variantheading{E Communi Sanctorum}

\variantheading{Pro Apostolis}

\lettrine[lines=2]{I}{BANT} Apóstoli gaudéntes a conspéctu concílii, quóniam digni hábiti sunt pro nómine Jesu contuméliam pati.

\variantheading{Pro Martyre}

\lettrine[lines=2]{J}{USTUS} cor suum trádidit ad vigilándum dilúculo ad Dóminum, qui fecit illum, et in conspéctu Altíssimi deprecábitur.

\variantheading{Pro pluribus Martyribus}

\lettrine[lines=2]{F}{ULGEBUNT} justi, et tamquam scintíllæ in arundinéto discúrrent. Judicábunt natiónes, et dominabúntur pópulis: et regnábit Dóminus illórum in perpétuum.

\variantheading{Pro Conf. Pont.}

\lettrine[lines=2]{F}{UNGI} sacerdótio, et habére laudem in nómine ipsíus, et offérre illi incénsum dignum in odórem suavitátis.

\variantheading{Pro Conf. non Pont.}

\lettrine[lines=2]{J}{USTUM} dedúxit Dóminus per vias rectas, et osténdit illi regnum Dei, et dedit illi sciéntiam sanctórum; honestávit illum in labóribus, et complévit labóres illíus.

\variantheading{Pro Virgine Martyre}

\lettrine[lines=2]{D}{OMINE}, Deus meus, exaltásti super terram habitatiónem meam, et pro morte defluénte deprecáta sum. Invocávi Dóminum, Patrem Dómini mei, ut non derelínquat me in die tribulatiónis meæ, et in témpore superbórum sine adjutório.

\variantheading{Pro Virgine non Martyre}

\lettrine[lines=2]{O}{} \textsc{quam} pulchra est casta generátio cum claritáte! immortális est enim memória illíus: quóniam apud Deum nota est, et apud hómines.

\variantheading{Pro Martyre tantum}

\lettrine[lines=2]{L}{AUDABIT} usque ad mortem ánima mea Dóminum, quóniam éruis sustinéntes te, et líberas eos de manu angústiæ, Dómine, Deus noster.

\variantheading{Pro nec Virg. nec Mart.}

\lettrine[lines=2]{M}{ULTAE} fíliæ congregavérunt divítias: tu supergréssa es univérsas. Fallax grátia, et vana est pulchritúdo: múlier timens Dóminum, ipsa laudábitur.

\variantheading{Pro Dedicatione Ecclesiæ}

\lettrine[lines=2]{E}{T} abstérget Deus omnem lácrimam ab óculis eórum: et mors ultra non erit, neque luctus, neque clamor, neque dolor erit ultra, quia prima abiérunt. Et dixit qui sedébat in throno: Ecce, nova fácio ómnia.

\variantheading{In Officiis B. Mariæ V.}

\lettrine[lines=2]{I}{N} platéis sicut cinnamómum et bálsamum aromatízans odórem dedi: quasi myrrha elécta dedi suavitátem odóris.

\V\ Adjutórium nostrum in nómine Dómini. \R\ Qui fecit cælum et terram.

\V\ Benedícite. \R\ Deus.

\benedictiolabel{Benedictio} Dóminus nos benedícat, et ab omni malo deféndat, et ad vitam perdúcat ætérnam. Et fidélium ánimæ per misericórdiam Dei requiéscant in pace. \R\ Amen.

\rubric{Deinde dícitur tantum} Pater noster \rubric{secréto, nisi sequátur ália Hora.}

% ===== END OF BATCH 11 -- Covers autumnalis.md lines 3674-3927,
% completing "Ad Primam" for Feria Secunda in full: the hymn "Jam
% lucis orto sídere" (with its own stray OCR "\*" before the final
% doxology stanza silently dropped, the same recurring artifact
% corrected in every prior occurrence this volume), Psalms 23, 18/i,
% 18/ii, and the conditional insertion of Psalm 46 (used when Ps. 50
% was said at Lauds instead, per the rubric), the two alternate
% Capitula ("In Officio Festi vel Octávæ" / "In ómnibus autem
% Fériis"), the brief Responsory "Christe, Fili Dei vivi," the
% "Preces dominicales et feriales (si dicendæ sint)" apparatus with
% its "In Précibus feriálibus tantum" and "In utrísque Précibus"
% subdivisions (the bold "Confíteor. Misereátur. Indulgéntiam."
% rendered with \formulalabel{} per the established rule for bold
% source markup), the two closing Orationes and the Martyrologium
% versicle, and finally the full "Lectio brevis" apparatus (the
% ferial variant "Dóminus autem dírigat," followed by "E Proprio
% Sanctorum" with its single "Pro Oct. Omn. Sanctorum" heading (this
% volume's one-item Autumn substitute for aestiva.tex's two-item
% Summer list "Pro Oct. S. Joannis Bapt." / "Pro Invent. S.
% Stephani," verified via direct diff against aestiva.md and
% transcribed fresh from autumnalis.md's own text) and "E Communi
% Sanctorum" with its full run of eleven Pro-headings and "In
% Officiis B. Mariæ V.," all reusing the \variantheading{} command
% established already in this volume's own Sunday Prime, each short
% text given a two-line lettrine as a Lectio-brevis genre text). The
% embedded
% rubric "Et absolúte, sine Orémus:" is preserved exactly as the
% source interleaves red rubric and black directive text. The
% document continues with "Ad Tertiam" (autumnalis.md line 3928).
% =====
\hourheading{Ad Tertiam}

Pater noster et Ave.

\V\ Deus, in adjutórium.

\hymnusheading{Hymnus}

\lettrine[lines=2]{N}{UNC}, Sancte, nobis, Spíritus,\\
Unum Patri cum Fílio,\\
Dignáre promptus íngeri\\
Nostro refúsus péctori.

Os, lingua, mens, sensus, vigor\\
Confessiónem pérsonent,\\
Flamméscat igne cáritas,\\
Accéndat ardor próximos.

Præsta, Pater piíssime,\\
Patríque compar Unice,\\
Cum Spíritu Paráclito\\
Regnans per omne sæculum.\\
Amen.

\antlabel{Ant.} Illuminátio mea.

\psalmheading{Psalmus 26, i}

\lettrine[lines=2]{D}{OMINUS} illuminátio mea, et salus mea, * quem timébo?

Dóminus protéctor vitæ meæ, * a quo trepidábo?

Dum apprópiant super me nocéntes, * ut edant carnes meas:

Qui tríbulant me inimíci mei, * ipsi infirmáti sunt, et cecidérunt.

Si consístant advérsum me castra, * non timébit cor meum.

Si exsúrgat advérsum me prælium, * in hoc ego sperábo.

Unam pétii a Dómino, hanc requíram, * ut inhábitem in domo Dómini ómnibus diébus vitæ meæ:

Ut vídeam voluptátem Dómini, * et vísitem templum ejus.

Quóniam abscóndit me in tabernáculo suo: * in die malórum protéxit me in abscóndito tabernáculi sui.

In petra exaltávit me: * et nunc exaltávit caput meum super inimícos meos.

Circuívi, et immolávi in tabernáculo ejus hóstiam vociferatiónis: * cantábo, et psalmum dicam Dómino.

\psalmheading{Psalmus 26, ii}

\lettrine[lines=2]{E}{XAUDI}, Dómine, vocem meam, qua clamávi ad te: * miserére mei, et exáudi me.

Tibi dixit cor meum, exquisívit te fácies mea: * fáciem tuam, Dómine, requíram.

Ne avértas fáciem tuam a me: * ne declínes in ira a servo tuo.

Adjútor meus esto: * ne derelínquas me, neque despícias me, Deus, salutáris meus.

Quóniam pater meus, et mater mea dereliquérunt me: * Dóminus autem assúmpsit me.

Legem pone mihi, Dómine, in via tua: * et dírige me in sémitam rectam propter inimícos meos.

Ne tradíderis me in ánimas tribulántium me: * quóniam insurrexérunt in me testes iníqui, et mentíta est iníquitas sibi.

Credo vidére bona Dómini * in terra vivéntium.

Exspécta Dóminum, viríliter age: * et confortétur cor tuum, et sústine Dóminum.

\psalmheading{Psalmus 27}

\lettrine[lines=2]{A}{D} te, Dómine, clamábo, Deus meus, ne síleas a me: * nequándo táceas a me, et assimilábor descendéntibus in lacum.

Exáudi, Dómine, vocem deprecatiónis meæ dum oro ad te: * dum extóllo manus meas ad templum sanctum tuum.

Ne simul trahas me cum peccatóribus: * et cum operántibus iniquitátem ne perdas me.

Qui loquúntur pacem cum próximo suo: * mala autem in córdibus eórum.

Da illis secúndum ópera eórum, * et secúndum nequítiam adinventiónum ipsórum.

Secúndum ópera mánuum eórum tríbue illis: * redde retributiónem eórum ipsis.

Quóniam non intellexérunt ópera Dómini, et in ópera mánuum ejus * déstrues illos, et non ædificábis eos.

Benedíctus Dóminus: * quóniam exaudívit vocem deprecatiónis meæ.

Dóminus adjútor meus, et protéctor meus: * in ipso sperávit cor meum, et adjútus sum.

Et reflóruit caro mea: * et ex voluntáte mea confitébor ei.

Dóminus fortitúdo plebis suæ: * et protéctor salvatiónum Christi sui est.

Salvum fac pópulum tuum, Dómine, et bénedic hereditáti tuæ: * et rege eos, et extólle illos usque in ætérnum.

\antlabel{Ant.} Illuminátio mea, et salus mea Dóminus.

\variantheading{In Officio feriali}

\capitulumlabel{Capitulum}{Jer. 17, 14}

\lettrine[lines=2]{S}{ANA} me, Dómine, et sanábor: salvum me fac, et salvus ero: quóniam laus mea tu es.

\R\ br. Sana ánimam meam, * Quia peccávi tibi. Sana. \V\ Ego dixi: Dómine, miserére mei. Quia. Glória Patri. Sana. \V\ Adjútor meus esto, ne derelínquas me. \R\ Neque despícias me, Deus, salutáris meus.

\variantheading{In Vigiliis communibus:}

\variantheading{Preces feriales}

KYRIE, eléison. Christe, eléison. Kýrie, eléison.

Pater noster \rubric{secreto usque ad} \V\ Et ne nos indúcas in tentatiónem. \R\ Sed líbera nos a malo. \V\ Dómine, Deus virtútum, convérte nos. \R\ Et osténde fáciem tuam, et salvi érimus. \V\ Exsúrge, Christe, ádjuva nos. \R\ Et líbera nos propter nomen tuum. \V\ Dómine, exáudi oratiónem meam. \R\ Et clamor meus ad te véniat.

\rubric{Réliqua ut in Ordinário. 19}

\hourheading{Ad Sextam}

Pater noster et Ave.

\V\ Deus, in adjutórium.

\hymnusheading{Hymnus}

\lettrine[lines=2]{R}{ECTOR} potens, verax Deus,\\
Qui témperas rerum vices,\\
Splendóre mane illúminas,\\
Et ígnibus merídiem:

Exstíngue flammas lítium,\\
Aufer calórem nóxium,\\
Confer salútem córporum,\\
Verámque pacem córdium.

Præsta, Pater piíssime,\\
Patríque compar Unice,\\
Cum Spíritu Paráclito\\
Regnans per omne sæculum.\\
Amen.

\antlabel{Ant.} In tua justítia.

\psalmheading{Psalmus 30, i}

\lettrine[lines=2]{I}{N} te, Dómine, sperávi, non confúndar in ætérnum: * in justítia tua líbera me.

Inclína ad me aurem tuam, * accélera ut éruas me.

Esto mihi in Deum protectórem, et in domum refúgii: * ut salvum me fácias.

Quóniam fortitúdo mea, et refúgium meum es tu: * et propter nomen tuum dedúces me, et enútries me.

Edúces me de láqueo hoc, quem abscondérunt mihi: * quóniam tu es protéctor meus.

In manus tuas comméndo spíritum meum: * redemísti me, Dómine, Deus veritátis.

Odísti observántes vanitátes, * supervácue.

Ego autem in Dómino sperávi: * exsultábo, et lætábor in misericórdia tua.

Quóniam respexísti humilitátem meam, * salvásti de necessitátibus ánimam meam.

Nec conclusísti me in mánibus inimíci: * statuísti in loco spatióso pedes meos.

\psalmheading{Psalmus 30, ii}

\lettrine[lines=2]{M}{ISERERE} mei, Dómine, quóniam tríbulor: * conturbátus est in ira óculus meus, ánima mea, et venter meus:

Quóniam defécit in dolóre vita mea: * et anni mei in gemítibus.

Infirmáta est in paupertáte virtus mea: * et ossa mea conturbáta sunt.

Super omnes inimícos meos factus sum oppróbrium et vicínis meis valde: * et timor notis meis.

Qui vidébant me, foras fugérunt a me: * oblivióni datus sum, tamquam mórtuus a corde.

Factus sum tamquam vas pérditum: * quóniam audívi vituperatiónem multórum commorántium in circúitu.

In eo dum convenírent simul advérsum me, * accípere ánimam meam consiliáti sunt.

Ego autem in te sperávi, Dómine: * dixi: Deus meus es tu: in mánibus tuis sortes meæ.

Eripe me de manu inimicórum meórum, * et a persequéntibus me.

Illústra fáciem tuam super servum tuum, salvum me fac in misericórdia tua: * Dómine, non confúndar, quóniam invocávi te.

Erubéscant ímpii, et deducántur in inférnum: * muta fiant lábia dolósa.

Quæ loquúntur advérsus justum iniquitátem, * in supérbia, et in abusióne.

\psalmheading{Psalmus 30, iii}

\lettrine[lines=2]{Q}{UAM} magna multitúdo dulcédinis tuæ, Dómine, * quam abscondísti timéntibus te.

Perfecísti eis, qui sperant in te, * in conspéctu filiórum hóminum.

Abscóndes eos in abscóndito faciéi tuæ, * a conturbatióne hóminum.

Próteges eos in tabernáculo tuo * a contradictióne linguárum.

Benedíctus Dóminus: * quóniam mirificávit misericórdiam suam mihi in civitáte muníta.

Ego autem dixi in excéssu mentis meæ: * Projéctus sum a fácie oculórum tuórum.

Ideo exaudísti vocem oratiónis meæ, * dum clamárem ad te.

Dilígite Dóminum, omnes sancti ejus: * quóniam veritátem requíret Dóminus, et retríbuet abundánter faciéntibus supérbiam.

Viríliter ágite, et confortétur cor vestrum, * omnes, qui sperátis in Dómino.

\antlabel{Ant.} In tua justítia líbera me, Dómine.

\variantheading{In Officio feriali}

\capitulumlabel{Capitulum}{Rom. 13, 8}

\lettrine[lines=2]{N}{EMINI} quidquam debeátis, nisi ut ínvicem diligátis: qui enim díligit próximum, legem implévit.

\R\ br. Benedícam Dóminum * In omni témpore. Benedícam. \V\ Semper laus ejus in ore meo. In. Glória Patri. Benedícam. \V\ Dóminus regit me, et nihil mihi déerit. \R\ In loco páscuæ ibi me collocávit.

\variantheading{In Vigiliis communibus:}

\variantheading{Preces feriales}

KYRIE, eléison. Christe, eléison. Kýrie, eléison.

Pater noster \rubric{secreto usque ad} \V\ Et ne nos indúcas in tentatiónem. \R\ Sed líbera nos a malo. \V\ Dómine, Deus virtútum, convérte nos. \R\ Et osténde fáciem tuam, et salvi érimus. \V\ Exsúrge, Christe, ádjuva nos. \R\ Et líbera nos propter nomen tuum. \V\ Dómine, exáudi oratiónem meam. \R\ Et clamor meus ad te véniat.

\rubric{Réliqua ut in Ordinário. 20}

\hourheading{Ad Nonam}

Pater noster et Ave.

\V\ Deus, in adjutórium.

\hymnusheading{Hymnus}

\lettrine[lines=2]{R}{ERUM}, Deus, tenax vigor,\\
Immótus in te pérmanens,\\
Lucis diúrnæ témpora\\
Succéssibus detérminans:

Largíre lumen véspere,\\
Quo vita nusquam décidat,\\
Sed prǽmium mortis sacræ\\
Perénnis instet glória.

Præsta, Pater piíssime,\\
Patríque compar Unice,\\
Cum Spíritu Paráclito\\
Regnans per omne sæculum.\\
Amen.

\antlabel{Ant.} Exsultáte, justi.

\psalmheading{Psalmus 31}

\lettrine[lines=2]{B}{EATI} quorum remíssæ sunt iniquitátes: * et quorum tecta sunt peccáta.

Beátus vir, cui non imputávit Dóminus peccátum, * nec est in spíritu ejus dolus.

Quóniam tácui, inveteravérunt ossa mea, * dum clamárem tota die.

Quóniam die ac nocte graváta est super me manus tua: * convérsus sum in ærúmna mea, dum confígitur spina.

Delíctum meum cógnitum tibi feci: * et injustítiam meam non abscóndi.

Dixi: Confitébor advérsum me injustítiam meam Dómino: * et tu remisísti impietátem peccáti mei.

Pro hac orábit ad te omnis sanctus, * in témpore opportúno.

Verúmtamen in dilúvio aquárum multárum, * ad eum non approximábunt.

Tu es refúgium meum a tribulatióne, quæ circúmdedit me: * exsultátio mea, érue me a circumdántibus me.

Intelléctum tibi dabo, et ínstruam te in via hac, qua gradiéris: * firmábo super te óculos meos.

Nolíte fíeri sicut equus et mulus, * quibus non est intelléctus.

In camo et freno maxíllas eórum constrínge, * qui non appróximant ad te.

Multa flagélla peccatóris, * sperántem autem in Dómino misericórdia circúmdabit.

Lætámini in Dómino et exsultáte, justi, * et gloriámini, omnes recti corde.

\psalmheading{Psalmus 32, i}

\lettrine[lines=2]{E}{XSULTATE}, justi, in Dómino: * rectos decet collaudátio.

Confitémini Dómino in cíthara: * in psaltério decem chordárum psállite illi.

Cantáte ei cánticum novum: * bene psállite ei in vociferatióne.

Quia rectum est verbum Dómini, * et ómnia ópera ejus in fide.

Díligit misericórdiam et judícium: * misericórdia Dómini plena est terra.

Verbo Dómini cæli firmáti sunt: * et spíritu oris ejus omnis virtus eórum.

Cóngregans sicut in utre aquas maris: * ponens in thesáuris abýssos.

Tímeat Dóminum omnis terra: * ab eo autem commoveántur omnes inhabitántes orbem.

Quóniam ipse dixit, et facta sunt: * ipse mandávit, et creáta sunt.

Dóminus díssipat consília géntium: * réprobat autem cogitatiónes populórum, et réprobat consília príncipum.

Consílium autem Dómini in ætérnum manet: * cogitatiónes cordis ejus in generatióne et generatiónem.

\psalmheading{Psalmus 32, ii}

\lettrine[lines=2]{B}{EATA} gens, cujus est Dóminus, Deus ejus: * pópulus, quem elégit in hereditátem sibi.

De cælo respéxit Dóminus: * vidit omnes fílios hóminum.

De præparáto habitáculo suo * respéxit super omnes, qui hábitant terram.

Qui finxit sigillátim corda eórum: * qui intélligit ómnia ópera eórum.

Non salvábitur rex per multam virtútem: * et gigas non salvábitur in multitúdine virtútis suæ.

Fallax equus ad salútem: * in abundántia autem virtútis suæ non salvábitur.

Ecce óculi Dómini super metuéntes eum: * et in eis, qui sperant super misericórdia ejus:

Ut éruat a morte ánimas eórum: * et alat eos in fame.

Anima nostra sústinet Dóminum: * quóniam adjútor et protéctor noster est.

Quia in eo lætábitur cor nostrum: * et in nómine sancto ejus sperávimus.

Fiat misericórdia tua, Dómine, super nos: * quemádmodum sperávimus in te.

\antlabel{Ant.} Exsultáte, justi, et gloriámini, omnes recti corde.

\variantheading{In Officio feriali}

\capitulumlabel{Capit.}{1 Petri 1, 17-19}

\lettrine[lines=2]{I}{N} timóre incolátus vestri témpore conversámini: sciéntes quod non corruptibílibus auro vel argénto redémpti estis, sed pretióso sánguine quasi Agni immaculáti Christi.

\R\ br. Rédime me, Dómine, * Et miserére mei. Rédime. \V\ Pes enim meus stetit in via recta. Et miserére mei. Glória Patri. Rédime. \V\ Ab occúltis meis munda me, Dómine. \R\ Et ab aliénis parce servo tuo.

\variantheading{In Vigiliis communibus:}

\variantheading{Preces feriales}

KYRIE, eléison. Christe, eléison. Kýrie, eléison.

Pater noster \rubric{secreto usque ad} \V\ Et ne nos indúcas in tentatiónem. \R\ Sed líbera nos a malo. \V\ Dómine, Deus virtútum, convérte nos. \R\ Et osténde fáciem tuam, et salvi érimus. \V\ Exsúrge, Christe, ádjuva nos. \R\ Et líbera nos propter nomen tuum. \V\ Dómine, exáudi oratiónem meam. \R\ Et clamor meus ad te véniat.

\rubric{Réliqua ut in Ordinário. 22}

% ===== END OF BATCH 12 -- Covers autumnalis.md lines 3928-4286,
% completing the Little Hours of Feria Secunda in full: "Ad
% Tertiam" (Psalms 26/i, 26/ii, 27), "Ad Sextam" (Psalms 30/i,
% 30/ii, 30/iii), and "Ad Nonam" (Psalms 31, 32/i, 32/ii), each with
% its own hymn, Capitulum, brief Responsory, and "In Vigiliis
% communibus: Preces feriales" apparatus. Three more occurrences of
% the recurring stray OCR "\*" artifact (in the hymns "Nunc, Sancte,
% nobis, Spíritus," "Rector potens, verax Deus," and "Rerum, Deus,
% tenax vigor," each immediately before the final "Præsta, Pater
% piíssime" doxology stanza) are silently dropped, consistent with
% every prior occurrence in this volume. The document continues
% with "Ad Vesperas" (autumnalis.md line 4288).
% =====
\hourheading{Ad Vesperas}

Pater noster et Ave.

\V\ Deus, in adjutórium.

\antlabel{Ant.} Inclinávit Dóminus * aurem suam mihi.

\rubric{¶ Antíphonæ, in hac et áliis Fériis apposítæ, dicúntur semper in Officio feriáli et in Festis sive trium sive novem Lectiónum, quando juxta Rúbricas suméndi sint Psalmi de Féria, ut in Psaltério.}

\psalmheading{Psalmus 114}

\lettrine[lines=2]{D}{ILEXI}, quóniam exáudiet Dóminus * vocem oratiónis meæ.

Quia inclinávit aurem suam mihi: * et in diébus meis invocábo.

Circumdedérunt me dolóres mortis: * et perícula inférni invenérunt me.

Tribulatiónem et dolórem invéni: * et nomen Dómini invocávi.

O Dómine, líbera ánimam meam: * miséricors Dóminus, et justus, et Deus noster miserétur.

Custódiens párvulos Dóminus: * humiliátus sum, et liberávit me.

Convértere, ánima mea, in réquiem tuam: * quia Dóminus benefécit tibi.

Quia erípuit ánimam meam de morte: * óculos meos a lácrimis, pedes meos a lapsu.

Placébo Dómino * in regióne vivórum.

\antlabel{Ant.} Inclinávit Dóminus aurem suam mihi.

\antlabel{Ant.} Vota mea * Dómino reddam coram omni pópulo ejus.

\psalmheading{Psalmus 115}

\lettrine[lines=2]{C}{REDIDI}, propter quod locútus sum: * ego autem humiliátus sum nimis.

Ego dixi in excéssu meo: * Omnis homo mendax.

Quid retríbuam Dómino, * pro ómnibus, quæ retríbuit mihi?

Cálicem salutáris accípiam: * et nomen Dómini invocábo.

Vota mea Dómino reddam coram omni pópulo ejus: * pretiósa in conspéctu Dómini mors sanctórum ejus:

O Dómine, quia ego servus tuus: * ego servus tuus, et fílius ancíllæ tuæ.

Dirupísti víncula mea: * tibi sacrificábo hóstiam laudis, et nomen Dómini invocábo.

Vota mea Dómino reddam in conspéctu omnis pópuli ejus: * in átriis domus Dómini, in médio tui, Jerúsalem.

\antlabel{Ant.} Vota mea Dómino reddam coram omni pópulo ejus.

\antlabel{Ant.} Clamávi, * et Dóminus exaudívit me.

\psalmheading{Psalmus 119}

\lettrine[lines=2]{A}{D} Dóminum cum tribulárer clamávi: * et exaudívit me.

Dómine, líbera ánimam meam a lábiis iníquis, * et a lingua dolósa.

Quid detur tibi, aut quid apponátur tibi * ad linguam dolósam?

Sagíttæ poténtis acútæ, * cum carbónibus desolatóriis.

Heu mihi, quia incolátus meus prolongátus est: * habitávi cum habitántibus Cedar: multum íncola fuit ánima mea.

Cum his, qui odérunt pacem, eram pacíficus: * cum loquébar illis, impugnábant me gratis.

\antlabel{Ant.} Clamávi, et Dóminus exaudívit me.

\antlabel{Ant.} Auxílium meum * a Dómino qui fecit cælum et terram.

\psalmheading{Psalmus 120}

\lettrine[lines=2]{L}{EVAVI} óculos meos in montes, * unde véniet auxílium mihi.

Auxílium meum a Dómino, * qui fecit cælum et terram.

Non det in commotiónem pedem tuum: * neque dormítet qui custódit te.

Ecce, non dormitábit neque dórmiet, * qui custódit Israël.

Dóminus custódit te, Dóminus protéctio tua, * super mánum déxteram tuam.

Per diem sol non uret te: * neque luna per noctem.

Dóminus custódit te ab omni malo: * custódiat ánimam tuam Dóminus.

Dóminus custódiat intróitum tuum, et éxitum tuum: * ex hoc nunc, et usque in sæculum.

\antlabel{Ant.} Auxílium meum a Dómino, qui fecit cælum et terram.

\antlabel{Ant.} Lætátus sum * \textdagger\ in his, quæ dicta sunt mihi. \textdaggerdbl

\psalmheading{Psalmus 121}

\lettrine[lines=2]{L}{ÆTATUS} sum \textdagger\ in his, quæ dicta sunt mihi: * \textdaggerdbl\ In domum Dómini íbimus.

Stantes erant pedes nostri, * in átriis tuis, Jerúsalem.

Jerúsalem, quæ ædificátur ut cívitas: * cujus participátio ejus in idípsum.

Illuc enim ascendérunt tribus, tribus Dómini: * testimónium Israël ad confiténdum nómini Dómini.

Quia illic sedérunt sedes in judício, * sedes super domum David.

Rogáte quæ ad pacem sunt Jerúsalem: * et abundántia diligéntibus te:

Fiat pax in virtúte tua: * et abundántia in túrribus tuis.

Propter fratres meos, et próximos meos, * loquébar pacem de te:

Propter domum Dómini, Dei nostri, * quæsívi bona tibi.

\antlabel{Ant.} Lætátus sum in his, quæ dicta sunt mihi.

\variantheading{In Officio feriali}

\rubric{Capítulum, Hymnus, Versus, et Ant. ad Magníficat dicúntur quotídie de Féria curránti, ut in Psaltério hic et álibi appónuntur.}

\capitulumlabel{Capitulum}{2 Cor. 1, 3-4}

\lettrine[lines=2]{B}{ENEDICTUS} Deus, et Pater Dómini nostri Jesu Christi, Pater misericordiárum, et Deus totíus consolatiónis, qui consolátur nos in omni tribulatióne nostra.

\hymnusheading{Hymnus}

\lettrine[lines=2]{I}{MMENSE} cæli Cónditor,\\
Qui mixta ne confúnderent,\\
Aquæ fluénta dívidens,\\
Cælum dedísti límitem.

Firmans locum cæléstibus,\\
Simúlque terræ rívulis;\\
Ut unda flammas témperet,\\
Terræ solum ne díssipent.

Infúnde nunc, piíssime,\\
Donum perénnis grátiæ:\\
Fraudis novæ ne cásibus\\
Nos error átterat vetus.

Lucem fides adáugeat:\\
Sic lúminis jubar ferat:\\
Hæc vana cuncta próterat:\\
Hanc falsa nulla cómprimant.

Præsta, Pater piíssime,\\
Patríque compar Unice,\\
Cum Spíritu Paráclito\\
Regnans per omne sæculum.\\
Amen.

\V\ Dirigátur, Dómine, orátio mea. \R\ Sicut incénsum in conspéctu tuo.

\antlabel{Ad Magnif. Ant.} Magníficat * \textdagger\ ánima mea Dóminum, quia respéxit Deus humilitátem meam.

\rubric{Réliqua ut in Ordinário. 22}

\hourheading{Ad Completorium}

\V\ Jube, domne.

\benedlabel{Bened.} Noctem quiétam, et finem perféctum concédat nobis Dóminus omnípotens. \R\ Amen.

\lectioheading{Lectio br.}{1 Petri 5, 8-9}

\lettrine[lines=2]{F}{RATRES}: Sóbrii estóte, et vigiláte: quia adversárius vester diábolus tamquam leo rúgiens círcuit, quærens quem dévoret: cui resístite fortes in fide. Tu autem, Dómine, miserére nobis. \R\ Deo grátias.

\V\ Adjutórium nostrum in nómine Dómini. \R\ Qui fecit cælum et terram.

Pater noster \rubric{quod dícitur totum secréto.}

\formulalabel{Confíteor.} Misereátur. Indulgéntiam.

\V\ Convérte nos, Deus, salutáris noster. \R\ Et avérte iram tuam a nobis.

\V\ Deus, in adjutórium.

\antlabel{Ant.} Salvum me fac.

\rubric{¶ Antíphonæ, in hac et áliis Fériis apposítæ, dicúntur semper ad Completórium tum de Témpore tum de Sanctis, quando juxta Rúbricas suméndi sint Psalmi de Féria.}

\psalmheading{Psalmus 6}

\lettrine[lines=2]{D}{OMINE}, ne in furóre tuo árguas me, * neque in ira tua corrípias me.

Miserére mei, Dómine, quóniam infírmus sum: * sana me, Dómine, quóniam conturbáta sunt ossa mea.

Et ánima mea turbáta est valde: * sed tu, Dómine, úsquequo?

Convértere, Dómine, et éripe ánimam meam: * salvum me fac propter misericórdiam tuam.

Quóniam non est in morte qui memor sit tui: * in inférno autem quis confitébitur tibi?

Laborávi in gémitu meo, lavábo per síngulas noctes lectum meum: * lácrimis meis stratum meum rigábo.

Turbátus est a furóre óculus meus: * inveterávi inter omnes inimícos meos.

Discédite a me, omnes, qui operámini iniquitátem: * quóniam exaudívit Dóminus vocem fletus mei.

Exaudívit Dóminus deprecatiónem meam, * Dóminus oratiónem meam suscépit.

Erubéscant, et conturbéntur veheménter omnes inimíci mei: * convertántur et erubéscant valde velóciter.

\psalmheading{Psalmus 7, i}

\lettrine[lines=2]{D}{OMINE}, Deus meus, in te sperávi: * salvum me fac ex ómnibus persequéntibus me, et líbera me.

Nequándo rápiat ut leo ánimam meam, * dum non est qui rédimat, neque qui salvum fáciat.

Dómine, Deus meus, si feci istud, * si est iníquitas in mánibus meis:

Si réddidi retribuéntibus mihi mala, * décidam mérito ab inimícis meis inánis.

Persequátur inimícus ánimam meam, et comprehéndat, et concúlcet in terra vitam meam, * et glóriam meam in púlverem dedúcat.

Exsúrge, Dómine, in ira tua: * et exaltáre in fínibus inimicórum meórum.

Et exsúrge, Dómine, Deus meus, in præcépto quod mandásti: * et synagóga populórum circúmdabit te.

Et propter hanc in altum regrédere: * Dóminus júdicat pópulos.

Júdica me, Dómine, secúndum justítiam meam, * et secúndum innocéntiam meam super me.

Consumétur nequítia peccatórum, et díriges justum, * scrutans corda et renes, Deus.

\psalmheading{Psalmus 7, ii}

\lettrine[lines=2]{J}{USTUM} adjutórium meum a Dómino, * qui salvos facit rectos corde.

Deus judex justus, fortis, et pátiens: * numquid iráscitur per síngulos dies?

Nisi convérsi fuéritis, gládium suum vibrábit: * arcum suum teténdit, et parávit illum.

Et in eo parávit vasa mortis: * sagíttas suas ardéntibus effécit.

Ecce, partúriit injustítiam: * concépit dolórem, et péperit iniquitátem.

Lacum apéruit, et effódit eum: * et íncidit in fóveam quam fecit.

Convertétur dolor ejus in caput ejus: * et in vérticem ipsíus iníquitas ejus descéndet.

Confitébor Dómino secúndum justítiam ejus: * et psallam nómini Dómini altíssimi.

\antlabel{Ant.} Salvum me fac, Dómine, propter misericórdiam tuam.

\hymnusheading{Hymnus}

\lettrine[lines=2]{T}{E} lucis ante términum,\\
Rerum Creátor, póscimus,\\
Ut pro tua cleméntia\\
Sis præsul et custódia.

Procul recédant sómnia,\\
Et nóctium phantásmata,\\
Hostémque nostrum cómprime,\\
Ne polluántur córpora.

Præsta, Pater piíssime,\\
Patríque compar Unice,\\
Cum Spíritu Paráclito\\
Regnans per omne sæculum.\\
Amen.

\capitulumlabel{Capitulum}{Jer. 14, 9}

\lettrine[lines=2]{T}{U} autem in nobis es, Dómine, et nomen sanctum tuum invocátum est super nos: ne derelínquas nos, Dómine, Deus noster.

\R\ br. In manus tuas, Dómine, * Comméndo spíritum meum. In manus. \V\ Redemísti nos, Dómine, Deus veritátis. Comméndo. Glória Patri. In manus.

\V\ Custódi nos, Dómine, ut pupíllam óculi. \R\ Sub umbra alárum tuárum prótege nos.

\antlabel{Ant.} Salva nos.

\canticumheading{Canticum Simeonis}{Luc. 2, 29-32}

\lettrine[lines=2]{N}{UNC} dimíttis servum tuum, Dómine, * secúndum verbum tuum in pace:

Quia vidérunt óculi mei * salutáre tuum,

Quod parásti * ante fáciem ómnium populórum,

Lumen ad revelatiónem géntium, * et glóriam plebis tuæ Israël.

Glória Patri.

\antlabel{Ant.} Salva nos, Dómine, vigilántes, custódi nos dormiéntes; ut vigilémus cum Christo, et requiescámus in pace.

\variantheading{Preces (si dicendæ sint)}

KYRIE, eléison. Christe, eléison. Kýrie, eléison.

Pater noster \rubric{secreto usque ad} \V\ Et ne nos indúcas in tentatiónem. \R\ Sed líbera nos a malo.

Credo in Deum \rubric{secreto usque ad} \V\ Carnis resurrectiónem. \R\ Vitam ætérnam. Amen. \V\ Benedíctus es, Dómine, Deus patrum nostrórum. \R\ Et laudábilis et gloriósus in sæcula. \V\ Benedicámus Patrem et Fílium cum Sancto Spíritu. \R\ Laudémus, et superexaltémus eum in sæcula. \V\ Benedíctus es, Dómine, in firmaménto cæli. \R\ Et laudábilis, et gloriósus, et superexaltátus in sæcula. \V\ Benedícat et custódiat nos omnípotens et miséricors Dóminus. \R\ Amen. \V\ Dignáre, Dómine, nocte ista. \R\ Sine peccáto nos custodíre. \V\ Miserére nostri, Dómine. \R\ Miserére nostri. \V\ Fiat misericórdia tua, Dómine, super nos. \R\ Quemádmodum sperávimus in te. \V\ Dómine, exáudi oratiónem meam. \R\ Et clamor meus ad te véniat.

\V\ Dóminus vobíscum.

Orémus. \oratiolabel{Oratio}

\lettrine[lines=2]{V}{ISITA}, quæsumus, Dómine, habitatiónem istam, et omnes insídias inimíci ab ea longe repélle: Angeli tui sancti hábitent in ea, qui nos in pace custódiant; et benedíctio tua sit super nos semper. Per Dóminum nostrum. \R\ Amen.

\V\ Dóminus vobíscum.

\V\ Benedicámus Dómino.

\benedictiolabel{Benedictio} Benedícat et custódiat nos omnípotens et miséricors Dóminus, Pater, et Fílius, et Spíritus Sanctus. \R\ Amen.

\variantheading{Antiphona finalis B. M. V.}

\V\ Divínum auxílium.

Pater, Ave et Credo.

% ===== END OF BATCH 13 -- Covers autumnalis.md lines 4288-4607,
% completing "Ad Vesperas" (Psalms 114, 115, 119, 120, 121, the
% Capitulum, the hymn "Immense cæli Cónditor," and the Magnificat
% antiphon) and "Ad Completorium" (Psalms 6, 7/i, 7/ii, the hymn
% "Te lucis ante términum," the Capitulum, the brief Responsory "In
% manus tuas," the Canticum Simeonis, and the Preces), and with
% them the ENTIRE "Feria Secunda" Office. Two more occurrences of
% the recurring stray OCR "\*" artifact (in "Immense cæli Cónditor"
% at autumnalis.md line 4436 and "Te lucis ante términum" at line 4550,
% both at the hymn's final doxology stanza) are silently dropped,
% consistent with every prior occurrence in this volume. The source
% at line 4464 bolds only "Confíteor." (not the full "Misereátur.
% Indulgéntiam."), unlike the fully-bold form seen at every other
% occurrence transcribed so far in this volume; this is a genuine,
% recurring alternation in autumnalis.md itself (confirmed by grep
% across the file) and is transcribed exactly as marked here:
% \formulalabel{Confíteor.} followed by plain-text "Misereátur.
% Indulgéntiam." The closing "Antiphona finalis B. M. V." is again
% given only as the bare cross-reference incipit "℣. Divínum
% auxílium," matching the identical abbreviated treatment already
% established for Sunday's Compline earlier in this volume. The
% document continues with "Feria Tertia" (autumnalis.md line 4610),
% the second weekday of the Psalterium cycle.
% =====
\feriaheading{Feria Tertia}

\hourheading{Ad Matutinum}

Pater, Ave et Credo.

\V\ Dómine, lábia mea.

\V\ Deus, in adjutórium.

\variantheading{In Officio feriali}

\invitatlabel{Invitat.} Jubilémus Deo, * Salutári nostro. \psalmlabel{Ps. 94} Veníte, exsultémus.

\hymnusheading{Hymnus}

\lettrine[lines=2]{C}{ONSORS} Patérni lúminis,\\
Lux ipse lucis, et dies,\\
Noctem canéndo rúmpimus:\\
Assíste postulántibus.

Aufer tenébras méntium,\\
Fuga catérvas dæmonum,\\
Expélle somnoléntiam,\\
Ne pigritántes óbruat.

Sic, Christe, nobis ómnibus\\
Indúlgeas credéntibus,\\
Ut prosit exorántibus\\
Quod præcinéntes psállimus.

Præsta, Pater piíssime,\\
Patríque compar Unice,\\
Cum Spíritu Paráclito\\
Regnans per omne sæculum.\\
Amen.

\nocturnoheading{In I Nocturno}

\antlabel{Ant.} Expúgna, Dómine, * impugnántes me.

\psalmheading{Psalmus 34, i}

\lettrine[lines=2]{J}{UDICA}, Dómine, nocéntes me, * expúgna impugnántes me.

Apprehénde arma et scutum: * exsúrge in adjutórium mihi.

Effúnde frámeam, et conclúde advérsus eos, qui persequúntur me: * dic ánimæ meæ: Salus tua ego sum.

Confundántur et revereántur, * quæréntes ánimam meam.

Avertántur retrórsum, et confundántur * cogitántes mihi mala.

Fiant tamquam pulvis ante fáciem venti: * et Angelus Dómini coárctans eos.

Fiat via illórum ténebræ et lúbricum: * et Angelus Dómini pérsequens eos.

Quóniam gratis abscondérunt mihi intéritum láquei sui: * supervácue exprobravérunt ánimam meam.

Véniat illi láqueus, quem ignórat: et cáptio, quam abscóndit, apprehéndat eum: * et in láqueum cadat in ipsum.

Anima autem mea exsultábit in Dómino: * et delectábitur super salutári suo.

Omnia ossa mea dicent: * Dómine, quis símilis tibi?

Erípiens ínopem de manu fortiórum ejus: * egénum et páuperem a diripiéntibus eum.

\antlabel{Ant.} Expúgna, Dómine, impugnántes me.

\antlabel{Ant.} Restítue ánimam meam * a malefáctis eórum, Dómine.

\psalmheading{Psalmus 34, ii}

\lettrine[lines=2]{S}{URGENTES} testes iníqui, * quæ ignorábam interrogábant me.

Retribuébant mihi mala pro bonis: * sterilitátem ánimæ meæ.

Ego autem cum mihi molésti essent, * induébar cilício.

Humiliábam in jejúnio ánimam meam: * et orátio mea in sinu meo convertétur.

Quasi próximum, et quasi fratrem nostrum, sic complacébam: * quasi lugens et contristátus, sic humiliábar.

Et advérsum me lætáti sunt, et convenérunt: * congregáta sunt super me flagélla, et ignorávi.

Dissipáti sunt, nec compúncti, tentavérunt me, subsannavérunt me subsannatióne: * frenduérunt super me déntibus suis.

Dómine, quando respícies? * restítue ánimam meam a malignitáte eórum, a leónibus únicam meam.

\antlabel{Ant.} Restítue ánimam meam a malefáctis eórum, Dómine.

\antlabel{Ant.} Exsúrge, Dómine, * et inténde judício meo.

\psalmheading{Psalmus 34, iii}

\lettrine[lines=2]{C}{ONFITEBOR} tibi in ecclésia magna, * in pópulo gravi laudábo te.

Non supergáudeant mihi qui adversántur mihi iníque: * qui odérunt me gratis et ánnuunt óculis.

Quóniam mihi quidem pacífice loquebántur: * et in iracúndia terræ loquéntes, dolos cogitábant.

Et dilatavérunt super me os suum: * dixérunt: Euge, euge, vidérunt óculi nostri.

Vidísti, Dómine, ne síleas: * Dómine, ne discédas a me.

Exsúrge et inténde judício meo: * Deus meus, et Dóminus meus in causam meam.

Júdica me secúndum justítiam tuam, Dómine, Deus meus, * et non supergáudeant mihi.

Non dicant in córdibus suis: Euge, euge, ánimæ nostræ: * nec dicant: Devorávimus eum.

Erubéscant et revereántur simul, * qui gratulántur malis meis.

Induántur confusióne et reveréntia * qui magna loquúntur super me.

Exsúltent et læténtur qui volunt justítiam meam: * et dicant semper: Magnificétur Dóminus, qui volunt pacem servi ejus.

Et lingua mea meditábitur justítiam tuam, * tota die laudem tuam.

\antlabel{Ant.} Exsúrge, Dómine, et inténde judício meo.

\variantheading{In Festis novem Lectionum:}

\V\ Lingua mea meditábitur justítiam tuam. \R\ Tota die laudem tuam, Dómine.

\absolutiolabel{Absol.} Exáudi, Dómine.

\benedlabel{Bened.} 1 Benedictióne perpétua. 2 Unigénitus Dei Fílius. 3 Spíritus Sancti grátia.

\nocturnoheading{In II Nocturno}

\antlabel{Ant.} Noli æmulári * \textdagger\ in eo, qui prosperátur, et facit iniquitátem.

\psalmheading{Psalmus 36, i}

\lettrine[lines=2]{N}{OLI} æmulári in malignántibus: * neque zeláveris faciéntes iniquitátem.

Quóniam tamquam fœnum velóciter aréscent: * et quemádmodum ólera herbárum cito décident.

Spera in Dómino, et fac bonitátem: * et inhábita terram, et pascéris in divítiis ejus.

Delectáre in Dómino: * et dabit tibi petitiónes cordis tui.

Revéla Dómino viam tuam, et spera in eo: * et ipse fáciet.

Et edúcet quasi lumen justítiam tuam: et judícium tuum tamquam merídiem: * súbditus esto Dómino, et ora eum.

Noli æmulári in eo, qui prosperátur in via sua: * in hómine faciénte injustítias.

Désine ab ira, et derelínque furórem: * noli æmulári ut malignéris.

Quóniam qui malignántur, exterminabúntur: * sustinéntes autem Dóminum, ipsi hereditábunt terram.

Et adhuc pusíllum, et non erit peccátor: * et quæres locum ejus, et non inviénies.

Mansuéti autem hereditábunt terram: * et delectabúntur in multitúdine pacis.

Observábit peccátor justum: * et stridébit super eum déntibus suis.

Dóminus autem irridébit eum: * quóniam próspicit quod véniet dies ejus.

Gládium evaginavérunt peccatóres: * intendérunt arcum suum.

Ut dejíciant páuperem et ínopem: * ut trucídent rectos corde.

Gládius eórum intret in corda ipsórum: * et arcus eórum confringátur.

\antlabel{Ant.} Noli æmulári in eo, qui prosperátur, et facit iniquitátem.

\antlabel{Ant.} Bráchia peccatórum * conteréntur: confírmat autem justos Dóminus.

\psalmheading{Psalmus 36, ii}

\lettrine[lines=2]{M}{ELIUS} est módicum justo, * super divítias peccatórum multas.

Quóniam bráchia peccatórum conteréntur: * confírmat autem justos Dóminus.

Novit Dóminus dies immaculatórum: * et heréditas eórum in ætérnum erit.

Non confundéntur in témpore malo, et in diébus famis saturabúntur: * quia peccatóres períbunt.

Inimíci vero Dómini mox ut honorificáti fúerint et exaltáti: * deficiéntes, quemádmodum fumus defícient.

Mutuábitur peccátor, et non solvet: * justus autem miserétur et tríbuet.

Quia benedicéntes ei hereditábunt terram: * maledicéntes autem ei períbunt.

Apud Dóminum gressus hóminis dirigéntur: * et viam ejus volet.

Cum cecíderit, non collidétur: * quia Dóminus suppónit manum suam.

Júnior fui, étenim sénui: * et non vidi justum derelíctum, nec semen ejus quærens panem.

Tota die miserétur et cómmodat: * et semen illíus in benedictióne erit.

Declína a malo, et fac bonum: * et inhábita in sæculum sæculi.

Quia Dóminus amat judícium, et non derelínquet sanctos suos: * in ætérnum conservabúntur.

Injústi puniéntur: * et semen impiórum períbit.

Justi autem hereditábunt terram: * et inhabitábunt in sæculum sæculi super eam.

\antlabel{Ant.} Bráchia peccatórum conteréntur: confírmat autem justos Dóminus.

\antlabel{Ant.} Custódi innocéntiam, * et vide æquitátem.

\psalmheading{Psalmus 36, iii}

\lettrine[lines=2]{O}{S} justi meditábitur sapiéntiam, * et lingua ejus loquétur judícium.

Lex Dei ejus in corde ipsíus: * et non supplantabúntur gressus ejus.

Consíderat peccátor justum: * et quærit mortificáre eum.

Dóminus autem non derelínquet eum in mánibus ejus: * nec damnábit eum, cum judicábitur illi.

Exspécta Dóminum, et custódi viam ejus: et exaltábit te ut hereditáte cápias terram: * cum períerint peccatóres vidébis.

Vidi ímpium superexaltátum, * et elevátum sicut cedros Líbani.

Et transívi, et ecce non erat: * et quæsívi eum, et non est invéntus locus ejus.

Custódi innocéntiam, et vide æquitátem: * quóniam sunt relíquiæ hómini pacífico.

Injústi autem disperíbunt simul: * relíquiæ impiórum interíbunt.

Salus autem justórum a Dómino: * et protéctor eórum in témpore tribulatiónis.

Et adjuvábit eos Dóminus et liberábit eos: * et éruet eos a peccatóribus, et salvábit eos: quia speravérunt in eo.

\antlabel{Ant.} Custódi innocéntiam, et vide æquitátem.

\variantheading{In Festis novem Lectionum:}

\V\ Exspécta Dóminum, et custódi viam ejus. \R\ Exaltábit te ut hereditáte cápias terram.

\absolutiolabel{Absolutio} Ipsíus píetas et.

\benedlabel{Bened.} 4 Deus Pater omnípotens. 5 Christus perpétuæ. 6 Ignem sui amóris.

\nocturnoheading{In III Nocturno}

\antlabel{Ant.} Ne in ira tua * corrípias me, Dómine.

\psalmheading{Psalmus 37, i}

\lettrine[lines=2]{D}{OMINE}, ne in furóre tuo árguas me, * neque in ira tua corrípias me.

Quóniam sagíttæ tuæ infíxæ sunt mihi: * et confirmásti super me manum tuam.

Non est sánitas in carne mea a fácie iræ tuæ: * non est pax óssibus meis a fácie peccatórum meórum.

Quóniam iniquitátes meæ supergréssæ sunt caput meum: * et sicut onus grave gravátæ sunt super me.

Putruérunt et corrúptæ sunt cicatríces meæ, * a fácie insipiéntiæ meæ.

Miser factus sum, et curvátus sum usque in finem: * tota die contristátus ingrediébar.

Quóniam lumbi mei impléti sunt illusiónibus: * et non est sánitas in carne mea.

Afflíctus sum, et humiliátus sum nimis: * rugiébam a gémitu cordis mei.

Dómine, ante te omne desidérium meum: * et gémitus meus a te non est abscónditus.

Cor meum conturbátum est, derelíquit me virtus mea: * et lumen oculórum meórum, et ipsum non est mecum.

\antlabel{Ant.} Ne in ira tua corrípias me, Dómine.

\antlabel{Ant.} Inténde in adjutórium meum, * Dómine, virtus salútis meæ.

\psalmheading{Psalmus 37, ii}

\lettrine[lines=2]{A}{MICI} mei, et próximi mei * advérsum me appropinquavérunt, et stetérunt.

Et qui juxta me erant, de longe stetérunt: * et vim faciébant qui quærébant ánimam meam.

Et qui inquirébant mala mihi, locúti sunt vanitátes: * et dolos tota die meditabántur.

Ego autem tamquam surdus non audiébam: * et sicut mutus non apériens os suum.

Et factus sum sicut homo non áudiens: * et non habens in ore suo redargutiónes.

Quóniam in te, Dómine, sperávi: * tu exáudies me, Dómine, Deus meus.

Quia dixi: Nequándo supergáudeant mihi inimíci mei: * et dum commovéntur pedes mei, super me magna locúti sunt.

Quóniam ego in flagélla parátus sum: * et dolor meus in conspéctu meo semper.

Quóniam iniquitátem meam annuntiábo: * et cogitábo pro peccáto meo.

Inimíci autem mei vivunt, et confirmáti sunt super me: * et multiplicáti sunt qui odérunt me iníque.

Qui retríbuunt mala pro bonis, detrahébant mihi: * quóniam sequébar bonitátem.

Ne derelínquas me, Dómine, Deus meus: * ne discésseris a me.

Inténde in adjutórium meum, * Dómine, Deus salútis meæ.

\antlabel{Ant.} Inténde in adjutórium meum, Dómine, virtus salútis meæ.

\antlabel{Ant.} Amove, Dómine, * a me plagas tuas.

\psalmheading{Psalmus 38}

\lettrine[lines=2]{D}{IXI}: Custódiam vias meas: * ut non delínquam in lingua mea.

Pósui ori meo custódiam, * cum consísteret peccátor advérsum me.

Obmútui, et humiliátus sum, et sílui a bonis: * et dolor meus renovátus est.

Concáluit cor meum intra me: * et in meditatióne mea exardéscet ignis.

Locútus sum in lingua mea: * Notum fac mihi, Dómine, finem meum.

Et númerum diérum meórum quis est: * ut sciam quid desit mihi.

Ecce, mensurábiles posuísti dies meos: * et substántia mea tamquam níhilum ante te.

Verúmtamen univérsa vánitas, * omnis homo vivens.

Verúmtamen in imágine pertránsit homo: * sed et frustra conturbátur.

Thesaurízat: * et ignórat cui congregábit ea.

Et nunc quæ est exspectátio mea? Nonne Dóminus? * Et substántia mea apud te est.

Ab ómnibus iniquitátibus meis érue me: * oppróbrium insipiénti dedísti me.

Obmútui, et non apérui os meum, quóniam tu fecísti: * ámove a me plagas tuas.

A fortitúdine manus tuæ ego defíci in increpatiónibus: * propter iniquitátem corripuísti hóminem.

Et tabéscere fecísti sicut aráneam ánimam ejus: * verúmtamen vane conturbátur omnis homo.

Exáudi oratiónem meam, Dómine, et deprecatiónem meam: * áuribus pércipe lácrimas meas.

Ne síleas: quóniam ádvena ego sum apud te, et peregrínus, * sicut omnes patres mei.

Remítte mihi, ut refrígerer priúsquam ábeam, * et ámplius non ero.

\antlabel{Ant.} Amove, Dómine, a me plagas tuas.

\variantheading{In feriali Officio et in Festis:}

\V\ Deus, ne síleas a me, remítte mihi. \R\ Quóniam íncola ego sum apud te, et peregrínus.

\variantheading{In Officio novem Lectionum}

\absolutiolabel{Absolutio.} A vínculis.

\benedlabel{Bened.} 7 Evangélica léctio. 8 In Festis Dómini: Divínum auxílium. In Festis B. M. V.: Cujus festum cólimus, ipsa Virgo. In Festis Sanctórum: Cujus (\textit{vel} Quorum \textit{aut} Quarum) festum cólimus. 9 Ad societátem. (Per evangélica dicta.)

\variantheading{In Officio trium Lectionum}

\absolutiolabel{Absolutio.} Ipsíus píetas et.

\variantheading{In Feriis, cum Homilia:}

\benedlabel{Bened.} 1 Evangélica. 2 Divínum. 3 Ad societátem.

\variantheading{In Feriis, sine Homilia:}

\benedlabel{Bened.} 1 Deus Pater. 2 Christus perpétuæ. 3 Ignem sui amóris.

\variantheading{In Festis Sanctorum:}

\benedlabel{Bened.} 1 Ille nos benedícat, qui sine fine vivit et regnat. \R\ Amen. 2 Cujus (\textit{vel} Quorum \textit{aut} Quarum). 3 Ad societátem.

% ===== END OF BATCH 14 -- Covers autumnalis.md lines 4610-4978,
% opening "Feria Tertia" (Tuesday) and completing its "Ad
% Matutinum" in full: the ferial Invitatorium "Jubilémus Deo,
% Salutári nostro," the hymn "Consors Patérni lúminis" (with the
% recurring stray OCR asterisk before its final doxology stanza
% silently dropped, consistent with every prior occurrence), and
% the three Nocturns (Psalms 34/i-iii; 36/i-iii; 37/i, 37/ii, 38),
% each with its versicle, Absolutio, and Benedictiones apparatus.
% "Absolutio" is rendered with the exact trailing period the source
% gives it at lines 4959 and 4965 ("Absolutio.") though not at line
% 4847 ("Absolutio"), preserving the source's own inconsistency.
% The compact multi-option "Bened." line for the ninth lesson is
% again transcribed as one continuous argument, matching the
% identical text already seen in Feria Secunda's Matins. The
% document continues with "Ad Laudes I" (autumnalis.md line 4981).
% =====
\hourheading{Ad Laudes I}

\rubric{In Feriis extra Vigílias commúnes occurréntibus, atque in Festis.}

(Pater noster et Ave.)

\V\ Deus, in adjutórium.

\antlabel{Ant.} Cantáte * \textdagger\ Dómino, et benedícite nómini ejus.

\psalmheading{Psalmus 95}

\lettrine[lines=2]{C}{ANTATE} \textdagger\ Dómino cánticum novum: * cantáte Dómino, omnis terra.

Cantáte Dómino, et benedícite nómini ejus: * annuntiáte de die in diem salutáre ejus.

Annuntiáte inter gentes glóriam ejus, * in ómnibus pópulis mirabília ejus.

Quóniam magnus Dóminus, et laudábilis nimis: * terríbilis est super omnes deos.

Quóniam omnes dii géntium dæmónia: * Dóminus autem cælos fecit.

Conféssio, et pulchritúdo in conspéctu ejus: * sanctimónia et magnificéntia in sanctificatióne ejus.

Afférte Dómino, pátriæ géntium, afférte Dómino glóriam et honórem: * afférte Dómino glóriam nómini ejus.

Tóllite hóstias, et introíte in átria ejus: * adoráte Dóminum in átrio sancto ejus.

Commoveátur a fácie ejus univérsa terra: * dícite in géntibus quia Dóminus regnávit.

Etenim corréxit orbem terræ qui non commovébitur: * judicábit pópulos in æquitáte.

Læténtur cæli, et exsúltet terra: commoveátur mare, et plenitúdo ejus: * gaudébunt campi, et ómnia quæ in eis sunt.

Tunc exsultábunt ómnia ligna silvárum a fácie Dómini, quia venit: * quóniam venit judicáre terram.

Judicábit orbem terræ in æquitáte, * et pópulos in veritáte tua.

\antlabel{Ant.} Cantáte Dómino, et benedícite nómini ejus.

\antlabel{Ant.} Salutáre vultus mei * Deus meus.

\psalmheading{Psalmus 42}

\lettrine[lines=2]{J}{UDICA} me, Deus, et discérne causam meam de gente non sancta, * ab hómine iníquo, et dolóso érue me.

Quia tu es, Deus, fortitúdo mea: * quare me repulísti? et quare tristis incédo, dum afflígit me inimícus?

Emítte lucem tuam et veritátem tuam: * ipsa me deduxérunt, et adduxérunt in montem sanctum tuum, et in tabernácula tua.

Et introíbo ad altáre Dei: * ad Deum, qui lætíficat juventútem meam.

Confitébor tibi in cíthara, Deus, Deus meus: * quare tristis es, ánima mea? et quare contúrbas me?

Spera in Deo, quóniam adhuc confitébor illi: * salutáre vultus mei, et Deus meus.

\antlabel{Ant.} Salutáre vultus mei Deus meus.

\antlabel{Ant.} Illúmina, Dómine, * vultum tuum super nos.

\psalmheading{Psalmus 66}

\lettrine[lines=2]{D}{EUS} misereátur nostri, et benedícat nobis: * illúminet vultum suum super nos, et misereátur nostri.

Ut cognoscámus in terra viam tuam, * in ómnibus géntibus salutáre tuum.

Confiteántur tibi pópuli, Deus: * confiteántur tibi pópuli omnes.

Læténtur et exsúltent gentes: * quóniam júdicas pópulos in æquitáte, et gentes in terra dírigis.

Confiteántur tibi pópuli, Deus, confiteántur tibi pópuli omnes: * terra dedit fructum suum.

Benedícat nos Deus, Deus noster, benedícat nos Deus: * et métuant eum omnes fines terræ.

\antlabel{Ant.} Illúmina, Dómine, vultum tuum super nos.

\antlabel{Ant.} Exaltáte * Regem sæculórum in opéribus vestris.

\canticumheading{Canticum Tobiæ}{Tob. 13, 1-10}

\lettrine[lines=2]{M}{AGNUS} es, Dómine, in ætérnum, * et in ómnia sæcula regnum tuum:

Quóniam tu flagéllas, et salvas: dedúcis ad ínferos, et redúcis: * et non est qui effúgiat manum tuam.

Confitémini Dómino, fílii Israël, * et in conspéctu géntium laudáte eum:

Quóniam ídeo dispérsit vos inter gentes, quæ ignórant eum, * ut vos enarrétis mirabília ejus.

Et faciátis scire eos, * quia non est álius Deus omnípotens præter eum.

Ipse castigávit nos propter iniquitátes nostras: * et ipse salvábit nos propter misericórdiam suam.

Aspícite ergo quæ fecit nobíscum, et cum timóre et tremóre confitémini illi: * Regémque sæculórum exaltáte in opéribus vestris.

Ego autem in terra captivitátis meæ confitébor illi: * quóniam osténdit majestátem suam in gentem peccatrícem.

Convertímini ítaque, peccatóres, et fácite justítiam coram Deo, * credéntes quod fáciat vobíscum misericórdiam suam:

Ego autem, et ánima mea * in eo lætábimur.

Benedícite Dóminum, omnes elécti ejus: * ágite dies lætítiæ, et confitémini illi.

\antlabel{Ant.} Exaltáte Regem sæculórum in opéribus vestris.

\antlabel{Ant.} Laudáte * \textdagger\ nomen Dómini, qui statis in domo Dómini.

\psalmheading{Psalmus 134}

\lettrine[lines=2]{L}{AUDATE} \textdagger\ nomen Dómini, * laudáte, servi, Dóminum.

Qui statis in domo Dómini, * in átriis domus Dei nostri.

Laudáte Dóminum, quia bonus Dóminus: * psállite nómini ejus, quóniam suáve.

Quóniam Jacob elégit sibi Dóminus, * Israël in possessiónem sibi.

Quia ego cognóvi quod magnus est Dóminus, * et Deus noster præ ómnibus diis.

Omnia quæcúmque vóluit, Dóminus fecit in cælo, et in terra, * in mari, et in ómnibus abýssis.

Edúcens nubes ab extrémo terræ: * fúlgura in plúviam fecit.

Qui prodúcit ventos de thesáuris suis: * qui percússit primogénita Ægýpti ab hómine usque ad pecus.

Et misit signa, et prodígia in médio tui, Ægýpte: * in Pharaónem, et in omnes servos ejus.

Qui percússit gentes multas: * et occídit reges fortes:

Sehon, regem Amorrhæórum, et Og, regem Basan, * et ómnia regna Chánaan.

Et dedit terram eórum hereditátem, * hereditátem Israël, pópulo suo.

Dómine, nomen tuum in ætérnum: * Dómine, memoriále tuum in generatiónem et generatiónem.

Quia judicábit Dóminus pópulum suum: * et in servis suis deprecábitur.

Simulácra géntium argéntum, et aurum, * ópera mánuum hóminum.

Os habent, et non loquéntur: * óculos habent, et non vidébunt.

Aures habent, et non áudient: * neque enim est spíritus in ore ipsórum.

Símiles illis fiant qui fáciunt ea: * et omnes qui confídunt in eis.

Domus Israël, benedícite Dómino: * domus Aaron, benedícite Dómino.

Domus Levi, benedícite Dómino: * qui timétis Dóminum, benedícite Dómino.

Benedíctus Dóminus ex Sion, * qui hábitat in Jerúsalem.

\antlabel{Ant.} Laudáte nomen Dómini, qui statis in domo Dómini.

\variantheading{In Officio feriali}

\capitulumlabel{Capit.}{Rom. 13, 12-13}

\lettrine[lines=2]{N}{OX} præcéssit, dies autem appropinquávit. Abjiciámus ergo ópera tenebrárum, et induámur arma lucis. Sicut in die honéste ambulémus.

\hymnusheading{Hymnus}

\lettrine[lines=2]{A}{LES} diéi núntius\\
Lucem propínquam præcínit:\\
Nos excitátor méntium\\
Jam Christus ad vitam vocat.

Auférte, clamat, léctulos,\\
Ægro sopóre desídes:\\
Castíque, recti, ac sóbrii\\
Vigiláte, jam sum próximus.

Jesum ciámus vócibus,\\
Flentes, precántes, sóbrii:\\
Inténta supplicátio\\
Dormíre cor mundum vetat.

Tu, Christe, somnum díscute:\\
Tu rumpe noctis víncula:\\
Tu solve peccátum vetus,\\
Novúmque lumen íngere.

Deo Patri sit glória,\\
Ejúsque soli Fílio,\\
Cum Spíritu Paráclito,\\
Nunc et per omne sæculum.\\
Amen.

\V\ Repléti sumus mane misericórdia tua. \R\ Exsultávimus, et delectáti sumus.

\antlabel{Ad Bened. Ant.} Eréxit nobis * Dóminus cornu salútis in domo David, púeri sui.

\rubric{Réliqua ut in Ordinário. 11}

\hourheading{Ad Laudes II}

\rubric{In Vigíliis commúnibus, quando fit Offícium de Féria.}

(Pater noster et Ave.)

\V\ Deus, in adjutórium.

\antlabel{Ant.} Dele iniquitátem meam.

\psalmheading{Psalmus 50}

\lettrine[lines=2]{M}{ISERERE} \textdagger\ mei, Deus: * secúndum magnam misericórdiam tuam.

Et secúndum multitúdinem miseratiónum tuárum, * dele iniquitátem meam.

Amplius lava me ab iniquitáte mea: * et a peccáto meo munda me.

Quóniam iniquitátem meam ego cognósco: * et peccátum meum contra me est semper.

Tibi soli peccávi, et malum coram te feci: * ut justificéris in sermónibus tuis, et vincas cum judicáris.

Ecce enim, in iniquitátibus concéptus sum: * et in peccátis concépit me mater mea.

Ecce enim, veritátem dilexísti: * incérta et occúlta sapiéntiæ tuæ manifestásti mihi.

Aspérges me hyssópo, et mundábor: * lavábis me, et super nivem dealbábor.

Audítui meo dabis gáudium et lætítiam: * et exsultábunt ossa humiliáta.

Avérte fáciem tuam a peccátis meis: * et omnes iniquitátes meas dele.

Cor mundum crea in me, Deus: * et spíritum rectum ínnova in viscéribus meis.

Ne projícias me a fácie tua: * et spíritum sanctum tuum ne áuferas a me.

Redde mihi lætítiam salutáris tui: * et spíritu principáli confírma me.

Docébo iníquos vias tuas: * et ímpii ad te converténtur.

Líbera me de sanguínibus, Deus, Deus salútis meæ: * et exsultábit lingua mea justítiam tuam.

Dómine, lábia mea apéries: * et os meum annuntiábit laudem tuam.

Quóniam si voluísses sacrifícium, dedíssem útique: * holocáustis non delectáberis.

Sacrifícium Deo spíritus contribulátus: * cor contrítum, et humiliátum, Deus, non despícies.

Benígne fac, Dómine, in bona voluntáte tua Sion: * ut ædificéntur muri Jerúsalem.

Tunc acceptábis sacrifícium justítiæ, oblatiónes, et holocáusta: * tunc impónent super altáre tuum vítulos.

\antlabel{Ant.} Dele iniquitátem meam, Dómine, secúndum multitúdinem miseratiónum tuárum.

\antlabel{Ant.} Discérne causam meam.

\psalmheading{Psalmus 42}

\lettrine[lines=2]{J}{UDICA} me, Deus, et discérne causam meam de gente non sancta, * ab hómine iníquo, et dolóso érue me.

Quia tu es, Deus, fortitúdo mea: * quare me repulísti? et quare tristis incédo, dum afflígit me inimícus?

Emítte lucem tuam et veritátem tuam: * ipsa me deduxérunt, et adduxérunt in montem sanctum tuum, et in tabernácula tua.

Et introíbo ad altáre Dei: * ad Deum, qui lætíficat juventútem meam.

Confitébor tibi in cíthara, Deus, Deus meus: * quare tristis es, ánima mea? et quare contúrbas me?

Spera in Deo, quóniam adhuc confitébor illi: * salutáre vultus mei, et Deus meus.

\antlabel{Ant.} Discérne causam meam, Deus, de gente non sancta.

\antlabel{Ant.} Deus misereátur. \textdagger

\psalmheading{Psalmus 66}

\lettrine[lines=2]{D}{EUS} misereátur \textdagger\ nostri, et benedícat nobis: * illúminet vultum suum super nos, et misereátur nostri.

Ut cognoscámus in terra viam tuam: * in ómnibus géntibus salutáre tuum.

Confiteántur tibi pópuli, Deus: * confiteántur tibi pópuli omnes.

Læténtur et exsúltent gentes: * quóniam júdicas pópulos in æquitáte, et gentes in terra dírigis.

Confiteántur tibi pópuli, Deus, confiteántur tibi pópuli omnes: * terra dedit fructum suum.

Benedícat nos Deus, Deus noster, benedícat nos Deus: * et métuant eum omnes fines terræ.

\antlabel{Ant.} Deus misereátur nostri, et benedícat nos.

\antlabel{Ant.} Corrípies me, Dómine.

\canticumheading{Canticum Ezechiæ}{Is. 38, 10-20}

\lettrine[lines=2]{E}{GO} dixi: In dimídio diérum meórum * vadam ad portas ínferi.

Quæsívi resíduum annórum meórum. * Dixi: Non vidébo Dóminum Deum in terra vivéntium.

Non aspíciam hóminem ultra, * et habitatórem quiétis.

Generátio mea abláta est, et convolúta est a me, * quasi tabernáculum pastórum.

Præcísa est, velut a texénte, vita mea: dum adhuc ordírer, succídit me: * de mane usque ad vésperam fínies me.

Sperábam usque ad mane, * quasi leo sic contrívit ómnia ossa mea:

De mane usque ad vésperam fínies me: * sicut pullus hirúndinis sic clamábo, meditábor ut colúmba:

Attenuáti sunt óculi mei, * suspiciéntes in excélsum:

Dómine, vim pátior, respónde pro me. * Quid dicam, aut quid respondébit mihi, cum ipse fécerit?

Recogitábo tibi omnes annos meos * in amaritúdine ánimæ meæ.

Dómine, si sic vívitur, et in tálibus vita spíritus mei, corrípies me, et vivificábis me. * Ecce, in pace amaritúdo mea amaríssima:

Tu autem eruísti ánimam meam ut non períret: * projecísti post tergum tuum ómnia peccáta mea.

Quia non inférnus confitébitur tibi, neque mors laudábit te: * non exspectábunt qui descéndunt in lacum, veritátem tuam.

Vivens vivens ipse confitébitur tibi, sicut et ego hódie: * pater fíliis notam fáciet veritátem tuam.

Dómine, salvum me fac, * et psalmos nostros cantábimus cunctis diébus vitæ nostræ in domo Dómini.

\antlabel{Ant.} Corrípies me, Dómine, et vivificábis me.

\antlabel{Ant.} Laudáte. \textdagger

\psalmheading{Psalmus 134}

\lettrine[lines=2]{L}{AUDATE} \textdagger\ nomen Dómini, * laudáte, servi, Dóminum.

Qui statis in domo Dómini, * in átriis domus Dei nostri.

Laudáte Dóminum, quia bonus Dóminus: * psállite nómini ejus, quóniam suáve.

Quóniam Jacob elégit sibi Dóminus, * Israël in possessiónem sibi.

Quia ego cognóvi quod magnus est Dóminus: * et Deus noster præ ómnibus diis.

Omnia quæcúmque vóluit, Dóminus fecit in cælo, et in terra, * in mari, et in ómnibus abýssis.

Edúcens nubes ab extrémo terræ: * fúlgura in plúviam fecit.

Qui prodúcit ventos de thesáuris suis: * qui percússit primogénita Ægýpti ab hómine usque ad pecus.

Et misit signa, et prodígia in médio tui, Ægýpte: * in Pharaónem, et in omnes servos ejus.

Qui percússit gentes multas: * et occídit reges fortes:

Sehon, regem Amorrhæórum, et Og, regem Basan, * et ómnia regna Chánaan.

Et dedit terram eórum hereditátem, * hereditátem Israël, pópulo suo.

Dómine, nomen tuum in ætérnum: * Dómine, memoriále tuum in generatiónem et generatiónem.

Quia judicábit Dóminus pópulum suum: * et in servis suis deprecábitur.

Simulácra géntium argéntum, et aurum, * ópera mánuum hóminum.

Os habent, et non loquéntur: * óculos habent, et non vidébunt.

Aures habent, et non áudient: * neque enim est spíritus in ore ipsórum.

Símiles illis fiant qui fáciunt ea: * et omnes qui confídunt in eis.

Domus Israël, benedícite Dómino: * domus Aaron, benedícite Dómino.

Domus Levi, benedícite Dómino: * qui timétis Dóminum, benedícite Dómino.

Benedíctus Dóminus ex Sion, * qui hábitat in Jerúsalem.

\antlabel{Ant.} Laudáte Dóminum, quia benígnus est, et in servis suis deprecábitur.

\capitulumlabel{Capit.}{Rom. 13, 12-13}

\lettrine[lines=2]{N}{OX} præcéssit, dies autem appropinquávit. Abjiciámus ergo ópera tenebrárum, et induámur arma lucis. Sicut in die honéste ambulémus.

\hymnusheading{Hymnus}

\lettrine[lines=2]{A}{LES} diéi núntius\\
Lucem propínquam præcínit:\\
Nos excitátor méntium\\
Jam Christus ad vitam vocat.

Auférte, clamat, léctulos,\\
Ægro sopóre desídes:\\
Castíque, recti, ac sóbrii\\
Vigiláte, jam sum próximus.

Jesum ciámus vócibus,\\
Flentes, precántes, sóbrii:\\
Inténta supplicátio\\
Dormíre cor mundum vetat.

Tu, Christe, somnum díscute:\\
Tu rumpe noctis víncula:\\
Tu solve peccátum vetus,\\
Novúmque lumen íngere.

Deo Patri sit glória,\\
Ejúsque soli Fílio,\\
Cum Spíritu Paráclito,\\
Nunc et per omne sæculum.\\
Amen.

\V\ Repléti sumus mane misericórdia tua. \R\ Exsultávimus, et delectáti sumus.

\antlabel{Ad Bened. Ant.} Eréxit nobis * Dóminus cornu salútis in domo David, púeri tui.

\rubric{Expléto Cántico ac repetíta Antíphona, dicúntur sequéntes}

\variantheading{Preces feriales}

KYRIE, eléison. Christe, eléison. Kýrie, eléison.

Pater noster \rubric{quod dícitur a solo Hebdomadário totum clara voce, usque ad} \V\ Et ne nos indúcas in tentatiónem. \R\ Sed líbera nos a malo. \V\ Ego dixi: Dómine, miserére mei. \R\ Sana ánimam meam, quia peccávi tibi. \V\ Convértere, Dómine, úsquequo? \R\ Et deprecábilis esto super servos tuos. \V\ Fiat misericórdia tua, Dómine, super nos. \R\ Quemádmodum sperávimus in te. \V\ Sacerdótes tui induántur justítiam. \R\ Et sancti tui exsúltent. \V\ Orémus pro beatíssimo Papa nostro N. \R\ Dóminus consérvet eum, et vivíficet eum, et beátum fáciat eum in terra, et non tradat eum in ánimam inimicórum ejus. \V\ Orémus et pro Antístite nostro N. \R\ Stet et pascat in fortitúdine tua, Dómine, in sublimitáte nóminis tui. \V\ Dómine, salvum fac regem. \R\ Et exáudi nos in die, qua invocavérimus te. \V\ Salvum fac pópulum tuum, Dómine, et bénedic hereditáti tuæ. \R\ Et rege eos, et extólle illos usque in ætérnum. \V\ Meménto Congregatiónis tuæ. \R\ Quam possedísti ab inítio. \V\ Fiat pax in virtúte tua. \R\ Et abundántia in túrribus tuis. \V\ Orémus pro benefactóribus nostris. \R\ Retribúere dignáre, Dómine, ómnibus, nobis bona faciéntibus propter nomen tuum, vitam ætérnam. Amen. \V\ Orémus pro fidélibus defúnctis. \R\ Réquiem ætérnam dona eis, Dómine, et lux perpétua lúceat eis. \V\ Requiéscant in pace. \R\ Amen. \V\ Pro frátribus nostris abséntibus. \R\ Salvos fac servos tuos, Deus meus, sperántes in te. \V\ Pro afflíctis et captívis. \R\ Líbera eos, Deus Israël, ex ómnibus tribulatiónibus suis. \V\ Mitte eis, Dómine, auxílium de sancto. \R\ Et de Sion tuére eos. \V\ Dómine, Deus virtútum, convérte nos. \R\ Et osténde fáciem tuam, et salvi érimus. \V\ Exsúrge, Christe, ádjuva nos. \R\ Et líbera nos propter nomen tuum. \V\ Dómine, exáudi oratiónem meam. \R\ Et clamor meus ad te véniat.

\rubric{Réliqua ut in Ordinário. 9}

% ===== END OF BATCH 15 -- Covers autumnalis.md lines 4981-5396, the
% two alternate ferial Lauds schemes for Feria Tertia, "Ad Laudes
% I" (Psalms 95, 42, 66, Canticum Tobiæ, 134) and "Ad Laudes II"
% (Psalms 50, 42, 66, Canticum Ezechiæ, 134), each with the
% Capitulum "Nox præcéssit" and the hymn "Ales diéi núntius." The
% recurring stray OCR "\*" artifact appears only in the Ad Laudes I
% copy of the hymn (autumnalis.md line 5427, before the final doxology
% stanza) and is silently dropped there; the Ad Laudes II copy of
% the identical hymn text (line 5336) has NO such artifact in the
% source and is transcribed exactly as printed, confirming the
% artifact is a per-occurrence OCR defect rather than a feature of
% the hymn text itself. The "Ad Bened. Ant." differs by a single
% word between the two schemes ("púeri sui" in Ad Laudes I vs.
% "púeri tui" in Ad Laudes II) and each is transcribed exactly as
% its own source gives it. Ad Laudes II additionally carries the
% "Preces feriales" litany in full. The document continues with
% "Ad Primam" (autumnalis.md line 5399).
% =====
\hourheading{Ad Primam}

Pater, Ave et Credo.

\V\ Deus, in adjutórium.

\hymnusheading{Hymnus}

\lettrine[lines=2]{J}{AM} lucis orto sídere,\\
Deum precémur súpplices,\\
Ut in diúrnis áctibus\\
Nos servet a nocéntibus.

Linguam refrénans témperet,\\
Ne litis horror ínsonet:\\
Visum fovéndo cóntegat,\\
Ne vanitátes háuriat.

Sint pura cordis íntima,\\
Absístat et vecórdia:\\
Carnis terat supérbiam\\
Potus cibíque párcitas.

Ut, cum dies abscésserit\\
Noctémque sors redúxerit,\\
Mundi per abstinéntiam\\
Ipsi canámus glóriam.

Deo Patri sit glória,\\
Ejúsque soli Fílio,\\
Cum Spíritu Paráclito,\\
Nunc et per omne sæculum.\\
Amen.

\antlabel{Ant.} Deus meus.

\psalmheading{Psalmus 24, i}

\lettrine[lines=2]{A}{D} te, Dómine, levávi ánimam meam: * Deus meus, in te confído, non erubéscam.

Neque irrídeant me inimíci mei: * étenim univérsi, qui sústinent te, non confundéntur.

Confundántur omnes iníqua agéntes * supervácue.

Vias tuas, Dómine, demónstra mihi: * et sémitas tuas édoce me.

Dírige me in veritáte tua, et doce me: * quia tu es, Deus, salvátor meus, et te sustínui tota die.

Reminíscere miseratiónum tuárum, Dómine, * et misericordiárum tuárum, quæ a sæculo sunt.

Delícta juventútis meæ, * et ignorántias meas ne memíneris.

Secúndum misericórdiam tuam meménto mei tu: * propter bonitátem tuam, Dómine.

\psalmheading{Psalmus 24, ii}

\lettrine[lines=2]{D}{ULCIS} et rectus Dóminus: * propter hoc legem dabit delinquéntibus in via.

Díriget mansuétos in judício: * docébit mites vias suas.

Univérsæ viæ Dómini, misericórdia et véritas, * requiréntibus testaméntum ejus et testimónia ejus.

Propter nomen tuum, Dómine, propitiáberis peccáto meo: * multum est enim.

Quis est homo qui timet Dóminum? * legem státuit ei in via, quam elégit.

Anima ejus in bonis demorábitur: * et semen ejus hereditábit terram.

Firmaméntum est Dóminus timéntibus eum: * et testaméntum ipsíus ut manifestétur illis.

\psalmheading{Psalmus 24, iii}

\lettrine[lines=2]{O}{CULI} mei semper ad Dóminum: * quóniam ipse evéllet de láqueo pedes meos.

Réspice in me, et miserére mei: * quia únicus et pauper sum ego.

Tribulatiónes cordis mei multiplicátæ sunt: * de necessitátibus meis érue me.

Vide humilitátem meam, et labórem meum: * et dimítte univérsa delícta mea.

Réspice inimícos meos quóniam multiplicáti sunt, * et ódio iníquo odérunt me.

Custódi ánimam meam, et érue me: * non erubéscam quóniam sperávi in te.

Innocéntes et recti adhæsérunt mihi: * quia sustínui te.

Líbera, Deus, Israël, * ex ómnibus tribulatiónibus suis.

\rubric{¶ Quando ad Laudes dictus fúerit Ps. 50 Miserére, hic subjúngitur seq. Ps. 95, qui in Láudibus locum non hábuit.}

\psalmheading{Psalmus 95}

\lettrine[lines=2]{C}{ANTATE} Dómino cánticum novum: * cantáte Dómino, omnis terra.

Cantáte Dómino, et benedícite nómini ejus: * annuntiáte de die in diem salutáre ejus.

Annuntiáte inter gentes glóriam ejus, * in ómnibus pópulis mirabília ejus.

Quóniam magnus Dóminus, et laudábilis nimis: * terríbilis est super omnes deos.

Quóniam omnes dii géntium dæmónia: * Dóminus autem cælos fecit.

Conféssio, et pulchritúdo in conspéctu ejus: * sanctimónia et magnificéntia in sanctificatióne ejus.

Afférte Dómino, pátriæ géntium, afférte Dómino glóriam et honórem: * afférte Dómino glóriam nómini ejus.

Tóllite hóstias, et introíte in átria ejus: * adoráte Dóminum in átrio sancto ejus.

Commoveátur a fácie ejus univérsa terra: * dícite in géntibus quia Dóminus regnávit.

Etenim corréxit orbem terræ qui non commovébitur: * judicábit pópulos in æquitáte.

Læténtur cæli, et exsúltet terra: commoveátur mare, et plenitúdo ejus: * gaudébunt campi, et ómnia quæ in eis sunt.

Tunc exsultábunt ómnia ligna silvárum a fácie Dómini, quia venit: * quóniam venit judicáre terram.

Judicábit orbem terræ in æquitáte, * et pópulos in veritáte sua.

\rubric{Secus vero, tribus Psalmis explétis, statim dícitur Ant.}

\antlabel{Ant.} Deus meus, in te confído, non erubéscam.

\rubric{In Officio Festi vel Octávæ:}

\capitulumlabel{Capitulum}{1 Tim. 1, 17}

\lettrine[lines=2]{R}{EGI} sæculórum immortáli et invisíbili, soli Deo honor et glória in sæcula sæculórum. Amen.

\rubric{In ómnibus autem Fériis et in Vigíliis commúnibus:}

\capitulumlabel{Capitulum}{Zach. 8, 19}

\lettrine[lines=2]{P}{ACEM} et veritátem dilígite, ait Dóminus omnípotens.

\R\ br. Christe, Fili Dei vivi, * Miserére nobis. Christe. \V\ \textit{Qui sedes ad déxteram Patris.} Miserére nobis. Glória Patri. Christe. \V\ Exsúrge, Christe, ádjuva nos. \R\ Et líbera nos propter nomen tuum.

\variantheading{Preces dominicales et feriales (si dicendæ sint)}

KYRIE, eléison. Christe, eléison. Kýrie, eléison.

Pater noster \rubric{secreto usque ad} \V\ Et ne nos indúcas in tentatiónem. \R\ Sed líbera nos a malo.

Credo in Deum \rubric{secreto usque ad} \V\ Carnis resurrectiónem. \R\ Vitam ætérnam. Amen. \V\ Et ego ad te, Dómine, clamávi. \R\ Et mane orátio mea prævéniet te. \V\ Repleátur os meum laude. \R\ Ut cantem glóriam tuam, tota die magnitúdinem tuam. \V\ Dómine, avérte fáciem tuam a peccátis meis. \R\ Et omnes iniquitátes meas dele. \V\ Cor mundum crea in me, Deus. \R\ Et spíritum rectum ínnova in viscéribus meis. \V\ Ne projícias me a fácie tua. \R\ Et spíritum sanctum tuum ne áuferas a me. \V\ Redde mihi lætítiam salutáris tui. \R\ Et spíritu principáli confírma me.

\variantheading{In Précibus feriálibus tantum}

\V\ Eripe me, Dómine, ab hómine malo. \R\ A viro iníquo éripe me. \V\ Eripe me de inimícis meis, Deus meus. \R\ Et ab insurgéntibus in me líbera me. \V\ Eripe me de operántibus iniquitátem. \R\ Et de viris sánguinum salva me. \V\ Sic psalmum dicam nómini tuo in sæculum sæculi. \R\ Ut reddam vota mea de die in diem. \V\ Exáudi nos, Deus, salutáris noster. \R\ Spes ómnium fínium terræ, et in mari longe. \V\ Deus, in adjutórium meum inténde. \R\ Dómine, ad adjuvándum me festína. \V\ Sanctus Deus, Sanctus fortis, Sanctus immortális. \R\ Miserére nobis. \V\ Bénedic, ánima mea, Dómino. \R\ Et ómnia, quæ intra me sunt, nómini sancto ejus. \V\ Bénedic, ánima mea, Dómino. \R\ Et noli oblivísci omnes retributiónes ejus. \V\ Qui propitiátur ómnibus iniquitátibus tuis. \R\ Qui sanat omnes infirmitátes tuas. \V\ Qui rédimit de intéritu vitam tuam. \R\ Qui corónat te in misericórdia et miseratiónibus. \V\ Qui replet in bonis desidérium tuum. \R\ Renovábitur ut áquilæ juvéntus tua.

\variantheading{In utrisque Précibus}

\V\ Adjutórium nostrum in nómine Dómini. \R\ Qui fecit cælum et terram.

\formulalabel{Confíteor. Misereátur. Indulgéntiam.}

\V\ Dignáre, Dómine, die isto. \R\ Sine peccáto nos custodíre. \V\ Miserére nostri, Dómine. \R\ Miserére nostri. \V\ Fiat misericórdia tua, Dómine, super nos. \R\ Quemádmodum sperávimus in te. \V\ Dómine, exáudi oratiónem meam. \R\ Et clamor meus ad te véniat.

\V\ Dóminus vobíscum.

Orémus. \oratiolabel{Oratio}

\lettrine[lines=2]{D}{OMINE}, Deus omnípotens, qui ad princípium hujus diéi nos perveníre fecísti: tua nos hódie salva virtúte; ut in hac die ad nullum declinémus peccátum, sed semper ad tuam justítiam faciéndam nostra procédant elóquia, dirigántur cogitatiónes et ópera. Per Dóminum. \R\ Amen.

\V\ Dóminus vobíscum. \V\ Benedicámus Dómino.

\rubric{In Choro Martyrológium.}

\V\ Pretiósa in conspéctu Dómini. \R\ Mors Sanctórum ejus.

\rubric{Et absolúte, sine} Orémus:

\oratiolabel{Oratio}

\lettrine[lines=2]{S}{ANCTA} María, et omnes Sancti intercédant pro nobis ad Dóminum, ut nos mereámur ab eo adjuvári et salvári, qui vivit et regnat in sæcula sæculórum. \R\ Amen.

\V\ Deus, in adjutórium meum inténde. \R\ Dómine, ad adjuvándum me festína. \rubric{Et dícitur ter; último additur:} Glória Patri. Sicut erat.

Kýrie, eléison. Christe, eléison. Kýrie, eléison.

Pater noster \rubric{secreto usque ad} \V\ Et ne nos indúcas in tentatiónem. \R\ Sed líbera nos a malo. \V\ Réspice in servos tuos, Dómine, et in ópera tua, et dírige fílios eórum. \R\ Et sit splendor Dómini, Dei nostri, super nos, et ópera mánuum nostrárum dírige super nos, et opus mánuum nostrárum dírige. \V\ Glória Patri. \R\ Sicut erat.

Orémus. \oratiolabel{Oratio}

\lettrine[lines=2]{D}{IRIGERE} et sanctificáre, régere et gubernáre dignáre, Dómine Deus, Rex cæli et terræ, hódie corda et córpora nostra, sensus, sermónes et actus nostros in lege tua, et in opéribus mandatórum tuórum: ut hic, et in ætérnum, te auxiliánte, salvi et líberi esse mereámur, Salvátor mundi: Qui vivis et regnas in sæcula sæculórum. \R\ Amen.

\V\ Jube, domne.

\benedlabel{Bened.} Dies et actus nostros in sua pace dispónat Dóminus omnípotens. \R\ Amen.

\variantheading{In Officio feriali}

\lettrine[lines=2]{D}{OMINUS} autem dírigat corda et córpora nostra in caritáte Dei, et patiéntia Christi. Tu autem.

\variantheading{E Proprio de Tempore}

\variantheading{Pro Oct. Omn. Sanctorum}

\lettrine[lines=2]{B}{ENEDÍCTIO} et cláritas et sapiéntia et gratiárum áctio, honor, virtus et fortitúdo Deo nostro in sǽcula sæculórum. Amen.

\variantheading{E Communi Sanctorum}

\variantheading{Pro Apostolis}

\lettrine[lines=2]{I}{BANT} Apóstoli gaudéntes a conspéctu concílii, quóniam digni hábiti sunt pro nómine Jesu contuméliam pati.

\variantheading{Pro Martyre}

\lettrine[lines=2]{J}{USTUS} cor suum trádidit ad vigilándum dilúculo ad Dóminum, qui fecit illum, et in conspéctu Altíssimi deprecábitur.

\variantheading{Pro pluribus Martyribus}

\lettrine[lines=2]{F}{ULGEBUNT} justi, et tamquam scintíllæ in arundinéto discúrrent. Judicábunt natiónes, et dominabúntur pópulis: et regnábit Dóminus illórum in perpétuum.

\variantheading{Pro Conf. Pont.}

\lettrine[lines=2]{F}{UNGI} sacerdótio, et habére laudem in nómine ipsíus, et offérre illi incénsum dignum in odórem suavitátis.

\variantheading{Pro Conf. non Pont.}

\lettrine[lines=2]{J}{USTUM} dedúxit Dóminus per vias rectas, et osténdit illi regnum Dei, et dedit illi sciéntiam sanctórum; honestávit illum in labóribus, et complévit labóres illíus.

\variantheading{Pro Virgine Martyre}

\lettrine[lines=2]{D}{OMINE}, Deus meus, exaltásti super terram habitatiónem meam, et pro morte defluénte deprecáta sum. Invocávi Dóminum, Patrem Dómini mei, ut non derelínquat me in die tribulatiónis meæ, et in témpore superbórum sine adjutório.

\variantheading{Pro Virgine non Martyre}

\lettrine[lines=2]{O}{} \textsc{quam} pulchra est casta generátio cum claritáte! immortális est enim memória illíus: quóniam apud Deum nota est, et apud hómines.

\variantheading{Pro Martyre tantum}

\lettrine[lines=2]{L}{AUDABIT} usque ad mortem ánima mea Dóminum, quóniam éruis sustinéntes te, et líberas eos de manu angústiæ, Dómine, Deus noster.

\variantheading{Pro nec Virg. nec Mart.}

\lettrine[lines=2]{M}{ULTAE} fíliæ congregavérunt divítias: tu supergréssa es univérsas. Fallax grátia, et vana est pulchritúdo: múlier timens Dóminum, ipsa laudábitur.

\variantheading{Pro Dedicatione Ecclesiæ}

\lettrine[lines=2]{E}{T} abstérget Deus omnem lácrimam ab óculis eórum: et mors ultra non erit, neque luctus, neque clamor, neque dolor erit ultra, quia prima abiérunt. Et dixit qui sedébat in throno: Ecce, nova fácio ómnia.

\variantheading{In Officiis B. Mariæ V.}

\lettrine[lines=2]{I}{N} platéis sicut cinnamómum et bálsamum aromatízans odórem dedi: quasi myrrha elécta dedi suavitátem odóris.

\V\ Adjutórium nostrum in nómine Dómini. \R\ Qui fecit cælum et terram.

\V\ Benedícite. \R\ Deus.

\benedictiolabel{Benedictio} Dóminus nos benedícat, et ab omni malo deféndat, et ad vitam perdúcat ætérnam. Et fidélium ánimæ per misericórdiam Dei requiéscant in pace. \R\ Amen.

\rubric{Deinde dícitur tantum} Pater noster \rubric{secréto, nisi sequátur ália Hora.}

% ===== END OF BATCH 16 -- Covers autumnalis.md lines 5399-5652,
% completing "Ad Primam" for Feria Tertia in full: the hymn "Jam
% lucis orto sídere" (stray OCR "\*" before the final doxology
% stanza silently dropped), Psalms 24/i-iii, the conditional Psalm
% 95 insertion (used when Ps. 50 was said at Lauds instead -- its
% final verse here reads "et pópulos in veritáte sua," genuinely
% differing from the "veritáte tua" reading seen in the same
% psalm's earlier Ad Laudes I occurrence this batch-cycle, and
% transcribed exactly as the source gives each), both Capitula, the
% brief Responsory, the full Preces apparatus, the closing
% Orationes, and the Lectio-brevis apparatus. The source heads this
% apparatus "E Proprio de Tempore" (autumnalis.md line 5593) rather
% than "E Proprio Sanctorum" as at the parallel point in Feria
% Secunda, despite containing the same single "Pro Oct. Omn.
% Sanctorum" heading used throughout this volume's substitute for
% aestiva.tex's two-item Summer list; this is transcribed exactly as
% the source heads it, confirmed via direct inspection of
% autumnalis.md at this point. The document continues with "Ad
% Tertiam" (autumnalis.md line 5655).
% =====
\hourheading{Ad Tertiam}

Pater noster et Ave.

\V\ Deus, in adjutórium.

\hymnusheading{Hymnus}

\lettrine[lines=2]{N}{UNC}, Sancte, nobis, Spíritus,\\
Unum Patri cum Fílio,\\
Dignáre promptus íngeri\\
Nostro refúsus péctori.

Os, lingua, mens, sensus, vigor\\
Confessiónem pérsonent,\\
Flamméscat igne cáritas,\\
Accéndat ardor próximos.

Præsta, Pater piíssime,\\
Patríque compar Unice,\\
Cum Spíritu Paráclito\\
Regnans per omne sæculum.\\
Amen.

\antlabel{Ant.} Respéxit me.

\psalmheading{Psalmus 39, i}

\lettrine[lines=2]{E}{XSPECTANS} exspectávi Dóminum, * et inténdit mihi.

Et exaudívit preces meas: * et edúxit me de lacu misériæ, et de luto fæcis.

Et státuit super petram pedes meos: * et diréxit gressus meos.

Et immísit in os meum cánticum novum, * carmen Deo nostro.

Vidébunt multi, et timébunt: * et sperábunt in Dómino.

Beátus vir, cujus est nomen Dómini spes ejus: * et non respéxit in vanitátes et insánias falsas.

Multa fecísti tu, Dómine, Deus meus, mirabília tua: * et cogitatiónibus tuis non est qui símilis sit tibi.

Annuntiávi et locútus sum: * multiplicáti sunt super númerum.

Sacrifícium et oblatiónem noluísti: * aures autem perfecísti mihi.

Holocáustum et pro peccáto non postulásti: * tunc dixi: Ecce, vénio.

In cápite libri scriptum est de me ut fácerem voluntátem tuam: * Deus meus, vólui, et legem tuam in médio cordis mei.

\psalmheading{Psalmus 39, ii}

\lettrine[lines=2]{A}{NNUNTIAVI} justítiam tuam in ecclésia magna, * ecce, lábia mea non prohibébo: Dómine, tu scisti.

Justítiam tuam non abscóndi in corde meo: * veritátem tuam et salutáre tuum dixi.

Non abscóndi misericórdiam tuam et veritátem tuam * a concílio multo.

Tu autem, Dómine, ne longe fácias miseratiónes tuas a me: * misericórdia tua et véritas tua semper suscepérunt me.

Quóniam circumdedérunt me mala, quorum non est númerus: * comprehendérunt me iniquitátes meæ, et non pótui ut vidérem.

Multiplicátæ sunt super capíllos cápitis mei: * et cor meum derelíquit me.

Compláceat tibi, Dómine, ut éruas me: * Dómine, ad adjuvándum me réspice.

\psalmheading{Psalmus 39, iii}

\lettrine[lines=2]{C}{ONFUNDANTUR} et revereántur simul, qui quærunt ánimam meam, * ut áuferant eam.

Convertántur retrórsum, et revereántur, * qui volunt mihi mala.

Ferant conféstim confusiónem suam, * qui dicunt mihi: Euge, euge.

Exsúltent et læténtur super te omnes quæréntes te: * et dicant semper: Magnificétur Dóminus: qui díligunt salutáre tuum.

Ego autem mendícus sum, et pauper: * Dóminus sollícitus est mei.

Adjútor meus, et protéctor meus tu es: * Deus meus, ne tardáveris.

\antlabel{Ant.} Respéxit me, et exaudívit deprecatiónem meam Dóminus.

\variantheading{In Officio feriali}

\capitulumlabel{Capitulum}{Jer. 17, 14}

\lettrine[lines=2]{S}{ANA} me, Dómine, et sanábor: salvum me fac, et salvus ero: quóniam laus mea tu es.

\R\ br. Sana ánimam meam, * Quia peccávi tibi. Sana. \V\ Ego dixi: Dómine, miserére mei. Quia. Glória Patri. Sana. \V\ Adjútor meus esto, ne derelínquas me. \R\ Neque despícias me, Deus, salutáris meus.

\variantheading{In Vigiliis communibus:}

\variantheading{Preces feriales}

KYRIE, eléison. Christe, eléison. Kýrie, eléison.

Pater noster \rubric{secreto usque ad} \V\ Et ne nos indúcas in tentatiónem. \R\ Sed líbera nos a malo. \V\ Dómine, Deus virtútum, convérte nos. \R\ Et osténde fáciem tuam, et salvi érimus. \V\ Exsúrge, Christe, ádjuva nos. \R\ Et líbera nos propter nomen tuum. \V\ Dómine, exáudi oratiónem meam. \R\ Et clamor meus ad te véniat.

\rubric{Réliqua ut in Ordinário. 19}

\hourheading{Ad Sextam}

Pater noster et Ave.

\V\ Deus, in adjutórium.

\hymnusheading{Hymnus}

\lettrine[lines=2]{R}{ECTOR} potens, verax Deus,\\
Qui témperas rerum vices,\\
Splendóre mane illúminas,\\
Et ígnibus merídiem:

Exstíngue flammas lítium,\\
Aufer calórem nóxium,\\
Confer salútem córporum,\\
Verámque pacem córdium.

Præsta, Pater piíssime,\\
Patríque compar Unice,\\
Cum Spíritu Paráclito\\
Regnans per omne sæculum.\\
Amen.

\antlabel{Ant.} Suscepísti me, Dómine.

\psalmheading{Psalmus 40}

\lettrine[lines=2]{B}{EATUS} qui intélligit super egénum, et páuperem: * in die mala liberábit eum Dóminus.

Dóminus consérvet eum, et vivíficet eum, et beátum fáciat eum in terra: * et non tradat eum in ánimam inimicórum ejus.

Dóminus opem ferat illi super lectum dolóris ejus: * univérsum stratum ejus versásti in infirmitáte ejus.

Ego dixi: Dómine, miserére mei: * sana ánimam meam, quia peccávi tibi.

Inimíci mei dixérunt mala mihi: * Quando moriétur, et períbit nomen ejus?

Et si ingrediebátur ut vidéret, vana loquebátur: * cor ejus congregávit iniquitátem sibi.

Egrediebátur foras, * et loquebátur in idípsum.

Advérsum me susurrábant omnes inimíci mei: * advérsum me cogitábant mala mihi.

Verbum iníquum constituérunt advérsum me: * Numquid qui dormit non adjíciet ut resúrgat?

Etenim homo pacis meæ, in quo sperávi: * qui edébat panes meos, magnificávit super me supplantatiónem.

Tu autem, Dómine, miserére mei, et resúscita me: * et retríbuam eis.

In hoc cognóvi quóniam voluísti me: * quóniam non gaudébit inimícus meus super me.

Me autem propter innocéntiam suscepísti: * et confirmásti me in conspéctu tuo in ætérnum.

Benedíctus Dóminus, Deus Israël, a sæculo et usque in sæculum: * fiat, fiat.

\psalmheading{Psalmus 41, i}

\lettrine[lines=2]{Q}{UEMADMODUM} desíderat cervus ad fontes aquárum: * ita desíderat ánima mea ad te, Deus.

Sitívit ánima mea ad Deum fortem vivum: * quando véniam, et apparébo ante fáciem Dei?

Fuérunt mihi lácrimæ meæ panes die ac nocte: * dum dícitur mihi quotídie: Ubi est Deus tuus?

Hæc recordátus sum, et effúdi in me ánimam meam: * quóniam transíbo in locum tabernáculi admirábilis, usque ad domum Dei.

In voce exsultatiónis, et confessiónis: * sonus epulántis.

Quare tristis es, ánima mea? * et quare contúrbas me?

Spera in Deo, quóniam adhuc confitébor illi: * salutáre vultus mei, et Deus meus.

\psalmheading{Psalmus 41, ii}

\lettrine[lines=2]{A}{D} meípsum ánima mea conturbáta est: * proptérea memor ero tui de terra Jordánis, et Hermóniim a monte módico.

Abýssus abýssum ínvocat, * in voce cataractárum tuárum.

Omnia excélsa tua, et fluctus tui * super me transiérunt.

In die mandávit Dóminus misericórdiam suam: * et nocte cánticum ejus.

Apud me orátio Deo vitæ meæ, * dicam Deo: Suscéptor meus es.

Quare oblítus es mei? * et quare contristátus incédo, dum afflígit me inimícus?

Dum confringúntur ossa mea, * exprobravérunt mihi qui tríbulant me inimíci mei.

Dum dicunt mihi per síngulos dies: Ubi est Deus tuus? * quare tristis es, ánima mea? et quare contúrbas me?

Spera in Deo, quóniam adhuc confitébor illi: * salutáre vultus mei, et Deus meus.

\antlabel{Ant.} Suscepísti me, Dómine: et confirmásti me in conspéctu tuo.

\variantheading{In Officio feriali}

\capitulumlabel{Capitulum}{Rom. 13, 8}

\lettrine[lines=2]{N}{EMINI} quidquam debeátis, nisi ut ínvicem diligátis: qui enim díligit próximum, legem implévit.

\R\ br. Benedícam Dóminum * In omni témpore. Benedícam. \V\ Semper laus ejus in ore meo. In. Glória Patri. Benedícam. \V\ Dóminus regit me, et nihil mihi déerit. \R\ In loco páscuæ ibi me collocávit.

\variantheading{In Vigiliis communibus:}

\variantheading{Preces feriales}

KYRIE, eléison. Christe, eléison. Kýrie, eléison.

Pater noster \rubric{secreto usque ad} \V\ Et ne nos indúcas in tentatiónem. \R\ Sed líbera nos a malo. \V\ Dómine, Deus virtútum, convérte nos. \R\ Et osténde fáciem tuam, et salvi érimus. \V\ Exsúrge, Christe, ádjuva nos. \R\ Et líbera nos propter nomen tuum. \V\ Dómine, exáudi oratiónem meam. \R\ Et clamor meus ad te véniat.

\rubric{Réliqua ut in Ordinário. 20}

\hourheading{Ad Nonam}

Pater noster et Ave.

\V\ Deus, in adjutórium.

\hymnusheading{Hymnus}

\lettrine[lines=2]{R}{ERUM}, Deus, tenax vigor,\\
Immótus in te pérmanens,\\
Lucis diúrnæ témpora\\
Succéssibus detérminans:

Largíre lumen véspere,\\
Quo vita nusquam décidat,\\
Sed prǽmium mortis sacræ\\
Perénnis instet glória.

Præsta, Pater piíssime,\\
Patríque compar Unice,\\
Cum Spíritu Paráclito\\
Regnans per omne sæculum.\\
Amen.

\antlabel{Ant.} Salvásti nos.

\psalmheading{Psalmus 43, i}

\lettrine[lines=2]{D}{EUS}, áuribus nostris audívimus: * patres nostri annuntiavérunt nobis.

Opus, quod operátus es in diébus eórum, * et in diébus antíquis.

Manus tua gentes dispérdidit, et plantásti eos: * afflixísti pópulos, et expulísti eos.

Nec enim in gládio suo possedérunt terram, * et bráchium eórum non salvávit eos:

Sed déxtera tua, et bráchium tuum, et illuminátio vultus tui: * quóniam complacuísti in eis.

Tu es ipse Rex meus et Deus meus: * qui mandas salútes Jacob.

In te inimícos nostros ventilábimus cornu: * et in nómine tuo spernémus insurgéntes in nobis.

Non enim in arcu meo sperábo: * et gládius meus non salvábit me.

Salvásti enim nos de affligéntibus nos: * et odiéntes nos confudísti.

In Deo laudábimur tota die: * et in nómine tuo confitébimur in sæculum.

\psalmheading{Psalmus 43, ii}

\lettrine[lines=2]{N}{UNC} autem repulísti et confudísti nos: * et non egrediéris, Deus, in virtútibus nostris.

Avertísti nos retrórsum post inimícos nostros: * et qui odérunt nos, diripiébant sibi.

Dedísti nos tamquam oves escárum: * et in géntibus dispersísti nos.

Vendidísti pópulum tuum sine prétio: * et non fuit multitúdo in commutatiónibus eórum.

Posuísti nos oppróbrium vicínis nostris, * subsannatiónem et derísum his, qui sunt in circúitu nostro.

Posuísti nos in similitúdinem géntibus: * commotiónem cápitis in pópulis.

Tota die verecúndia mea contra me est, * et confúsio faciéi meæ coopéruit me.

A voce exprobrántis, et obloquéntis: * a fácie inimíci, et persequéntis.

Hæc ómnia venérunt super nos, nec oblíti sumus te: * et iníque non égimus in testaménto tuo.

Et non recéssit retro cor nostrum: * et declinásti sémitas nostras a via tua:

Quóniam humiliásti nos in loco afflictiónis, * et coopéruit nos umbra mortis.

\psalmheading{Psalmus 43, iii}

\lettrine[lines=2]{S}{I} oblíti sumus nomen Dei nostri, * et si expándimus manus nostras ad deum aliénum:

Nonne Deus requíret ista? * ipse enim novit abscóndita cordis.

Quóniam propter te mortificámur tota die: * æstimáti sumus sicut oves occisiónis.

Exsúrge, quare obdórmis, Dómine? * exsúrge, et ne repéllas in finem.

Quare fáciem tuam avértis, * oblivísceris inópiæ nostræ, et tribulatiónis nostræ?

Quóniam humiliáta est in púlvere ánima nostra: * conglutinátus est in terra venter noster.

Exsúrge, Dómine, ádjuva nos: * et rédime nos propter nomen tuum.

\antlabel{Ant.} Salvásti nos, Dómine: et in nómine tuo confitébimur in sæcula.

\variantheading{In Officio feriali}

\capitulumlabel{Capit.}{1 Petri 1, 17-19}

\lettrine[lines=2]{I}{N} timóre incolátus vestri témpore conversámini: sciéntes quod non corruptibílibus auro vel argénto redémpti estis, sed pretióso sánguine quasi Agni immaculáti Christi.

\R\ br. Rédime me, Dómine, * Et miserére mei. Rédime. \V\ Pes enim meus stetit in via recta. Et miserére mei. Glória Patri. Rédime. \V\ Ab occúltis meis munda me, Dómine. \R\ Et ab aliénis parce servo tuo.

\variantheading{In Vigiliis communibus:}

\variantheading{Preces feriales}

KYRIE, eléison. Christe, eléison. Kýrie, eléison.

Pater noster \rubric{secreto usque ad} \V\ Et ne nos indúcas in tentatiónem. \R\ Sed líbera nos a malo. \V\ Dómine, Deus virtútum, convérte nos. \R\ Et osténde fáciem tuam, et salvi érimus. \V\ Exsúrge, Christe, ádjuva nos. \R\ Et líbera nos propter nomen tuum. \V\ Dómine, exáudi oratiónem meam. \R\ Et clamor meus ad te véniat.

\rubric{Réliqua ut in Ordinário. 22}

% ===== END OF BATCH 17 -- Covers autumnalis.md lines 5655-5978,
% completing the Little Hours of Feria Tertia in full: "Ad
% Tertiam" (Psalms 39/i-iii), "Ad Sextam" (Psalms 40, 41/i, 41/ii),
% and "Ad Nonam" (Psalms 43/i-iii), each with its own hymn,
% Capitulum, brief Responsory, and "In Vigiliis communibus: Preces
% feriales" apparatus. Three further occurrences of the recurring
% stray OCR "\*" artifact (in the hymns "Nunc, Sancte, nobis,
% Spíritus," "Rector potens, verax Deus," and "Rerum, Deus, tenax
% vigor," each before the final "Præsta, Pater piíssime" doxology
% stanza) are silently dropped, consistent with every prior
% occurrence in this volume. The document continues with "Ad
% Vesperas" (autumnalis.md line 5981).
% =====
\hourheading{Ad Vesperas}

Pater noster et Ave.

\V\ Deus, in adjutórium.

\antlabel{Ant.} Qui hábitas in cælis, * miserére nobis.

\psalmheading{Psalmus 122}

\lettrine[lines=2]{A}{D} te levávi óculos meos, * qui hábitas in cælis.

Ecce, sicut óculi servórum * in mánibus dominórum suórum.

Sicut óculi ancíllæ in mánibus dóminæ suæ: * ita óculi nostri ad Dóminum, Deum nostrum, donec misereátur nostri.

Miserére nostri, Dómine, miserére nostri: * quia multum repléti sumus despectióne:

Quia multum repléta est ánima nostra: * oppróbrium abundántibus, et despéctio supérbis.

\antlabel{Ant.} Qui hábitas in cælis, miserére nobis.

\antlabel{Ant.} Adjutórium nostrum * in nómine Dómini.

\psalmheading{Psalmus 123}

\lettrine[lines=2]{N}{ISI} quia Dóminus erat in nobis, dicat nunc Israël: * nisi quia Dóminus erat in nobis.

Cum exsúrgerent hómines in nos, * forte vivos deglutíssent nos:

Cum irascerétur furor eórum in nos, * fórsitan aqua absorbuísset nos.

Torréntem pertransívit ánima nostra: * fórsitan pertransísset ánima nostra aquam intolerábilem.

Benedíctus Dóminus * qui non dedit nos in captiónem déntibus eórum.

Anima nostra sicut passer erépta est * de láqueo venántium:

Láqueus contrítus est, * et nos liberáti sumus.

Adjutórium nostrum in nómine Dómini, * qui fecit cælum et terram.

\antlabel{Ant.} Adjutórium nostrum in nómine Dómini.

\antlabel{Ant.} In circúitu pópuli sui * Dóminus, ex hoc nunc et usque in sæculum.

\psalmheading{Psalmus 124}

\lettrine[lines=2]{Q}{UI} confídunt in Dómino, sicut mons Sion: * non commovébitur in ætérnum, qui hábitat in Jerúsalem.

Montes in circúitu ejus: * et Dóminus in circúitu pópuli sui, ex hoc nunc et usque in sæculum.

Quia non relínquet Dóminus virgam peccatórum super sortem justórum: * ut non exténdant justi ad iniquitátem manus suas.

Bénefac, Dómine, bonis, * et rectis corde.

Declinántes autem in obligatiónes addúcet Dóminus cum operántibus iniquitátem: * pax super Israël.

\antlabel{Ant.} In circúitu pópuli sui Dóminus, ex hoc nunc et usque in sæculum.

\antlabel{Ant.} Magnificávit Dóminus * fácere nobíscum: facti sumus lætántes.

\psalmheading{Psalmus 125}

\lettrine[lines=2]{I}{N} converténdo Dóminus captivitátem Sion: * facti sumus sicut consoláti:

Tunc replétum est gáudio os nostrum: * et lingua nostra exsultatióne.

Tunc dicent inter gentes: * Magnificávit Dóminus fácere cum eis.

Magnificávit Dóminus fácere nobíscum: * facti sumus lætántes.

Convérte, Dómine, captivitátem nostram, * sicut torrens in Austro.

Qui séminant in lácrimis, * in exsultatióne metent.

Eúntes ibant et flebant, * mitténtes sémina sua.

Veniéntes autem vénient cum exsultatióne, * portántes manípulos suos.

\antlabel{Ant.} Magnificávit Dóminus fácere nobíscum: facti sumus lætántes.

\antlabel{Ant.} Dóminus ædíficet * nobis domum, et custódiat civitátem.

\psalmheading{Psalmus 126}

\lettrine[lines=2]{N}{ISI} Dóminus ædificáverit domum, * in vanum laboravérunt qui ædíficant eam.

Nisi Dóminus custodíerit civitátem, * frustra vígilat qui custódit eam.

Vanum est vobis ante lucem súrgere: * súrgite postquam sedéritis, qui manducátis panem dolóris.

Cum déderit diléctis suis somnum: * ecce heréditas Dómini fílii: merces, fructus ventris.

Sicut sagíttæ in manu poténtis: * ita fílii excussórum.

Beátus vir, qui implévit desidérium suum ex ipsis: * non confundétur cum loquétur inimícis suis in porta.

\antlabel{Ant.} Dóminus ædíficet nobis domum, et custódiat civitátem.

\variantheading{In Officio feriali}

\capitulumlabel{Capitulum}{2 Cor. 1, 3-4}

\lettrine[lines=2]{B}{ENEDICTUS} Deus, et Pater Dómini nostri Jesu Christi, Pater misericordiárum, et Deus totíus consolatiónis, qui consolátur nos in omni tribulatióne nostra.

\hymnusheading{Hymnus}

\lettrine[lines=2]{T}{ELLURIS} alme Cónditor\\
Mundi solum qui séparans,\\
Pulsis aquæ moléstiis,\\
Terram dedísti immóbilem:

Ut germen aptum próferens,\\
Fulvis decóra flóribus,\\
Fecúnda fructu sísteret,\\
Pastúmque gratum rédderet.

Mentis perústæ vúlnera\\
Munda viróre grátiæ:\\
Ut facta fletu díluat,\\
Motúsque pravos átterat.

Jussis tuis obtémperet:\\
Nullis malis appróximet:\\
Bonis repléri gáudeat,\\
Et mortis ictum nésciat.

Præsta, Pater piíssime,\\
Patríque compar Unice,\\
Cum Spíritu Paráclito\\
Regnans per omne sæculum.\\
Amen.

\V\ Dirigátur, Dómine, orátio mea. \R\ Sicut incénsum in conspéctu tuo.

\antlabel{Ad Magnif. Ant.} Exsultávit * spíritus meus in Deo, salutári meo.

\rubric{Réliqua ut in Ordinário. 22}

\hourheading{Ad Completorium}

\V\ Jube, domne.

\benedlabel{Bened.} Noctem quiétam, et finem perféctum concédat nobis Dóminus omnípotens. \R\ Amen.

\lectioheading{Lectio br.}{1 Petri 5, 8-9}

\lettrine[lines=2]{F}{RATRES}: Sóbrii estóte, et vigiláte: quia adversárius vester diábolus tamquam leo rúgiens círcuit, quærens quem dévoret: cui resístite fortes in fide. Tu autem, Dómine, miserére nobis. \R\ Deo grátias.

\V\ Adjutórium nostrum in nómine Dómini. \R\ Qui fecit cælum et terram.

Pater noster \rubric{quod dícitur totum secréto.}

\formulalabel{Confíteor.} Misereátur. Indulgéntiam.

\V\ Convérte nos, Deus, salutáris noster. \R\ Et avérte iram tuam a nobis.

\V\ Deus, in adjutórium.

\antlabel{Ant.} Tu Dómine.

\psalmheading{Psalmus 11}

\lettrine[lines=2]{S}{ALVUM} me fac, Dómine, quóniam defécit sanctus: * quóniam diminútæ sunt veritátes a fíliis hóminum.

Vana locúti sunt unusquísque ad próximum suum: * lábia dolósa, in corde et corde locúti sunt.

Dispérdat Dóminus univérsa lábia dolósa, * et linguam magníloquam.

Qui dixérunt: Linguam nostram magnificábimus, lábia nostra a nobis sunt, * quis noster Dóminus est?

Propter misériam ínopum, et gémitum páuperum, * nunc exsúrgam, dicit Dóminus.

Ponam in salutári: * fiduciáliter agam in eo.

Elóquia Dómini, elóquia casta: * argéntum igne examinátum, probátum terræ purgátum séptuplum.

Tu, Dómine, servábis nos: et custódies nos * a generatióne hac in ætérnum.

In circúitu ímpii ámbulant: * secúndum altitúdinem tuam multiplicásti fílios hóminum.

\psalmheading{Psalmus 12}

\lettrine[lines=2]{U}{SQUEQUO}, Dómine, oblivascéris me in finem? * Usquequo avértis fáciem tuam a me?

Quámdiu ponam consília in ánima mea, * dolórem in corde meo per diem?

Usquequo exaltábitur inimícus meus super me? * réspice, et exáudi me, Dómine, Deus meus.

Illúmina óculos meos ne umquam obdórmiam in morte: * nequándo dicat inimícus meus: Præválui advérsus eum.

Qui tríbulant me, exsultábunt si motus fúero: * ego autem in misericórdia tua sperávi.

Exsultábit cor meum in salutári tuo: cantábo Dómino qui bona tríbuit mihi: * et psallam nómini Dómini altíssimi.

\psalmheading{Psalmus 15}

\lettrine[lines=2]{C}{ONSERVA} me, Dómine, quóniam sperávi in te. * Dixi Dómino: Deus meus es tu, quóniam bonórum meórum non eges.

Sanctis, qui sunt in terra ejus, * mirificávit omnes voluntátes meas in eis.

Multiplicátæ sunt infirmitátes eórum: * póstea acceleravérunt.

Non congregábo conventícula eórum de sanguínibus, * nec memor ero nóminum eórum per lábia mea.

Dóminus pars hereditátis meæ, et cálicis mei: * tu es, qui restítues hereditátem meam mihi.

Funes cecidérunt mihi in præcláris: * étenim heréditas mea præclára est mihi.

Benedícam Dóminum, qui tríbuit mihi intelléctum: * ínsuper et usque ad noctem increpuérunt me renes mei.

Providébam Dóminum in conspéctu meo semper: * quóniam a dextris est mihi, ne commóvear.

Propter hoc lætátum est cor meum, et exsultávit lingua mea: * ínsuper et caro mea requiéscet in spe.

Quóniam non derelínques ánimam meam in inférno: * nec dabis sanctum tuum vidére corruptiónem.

Notas mihi fecísti vias vitæ, adimplébis me lætítia cum vultu tuo: * delectatiónes in déxtera tua usque in finem.

\antlabel{Ant.} Tu, Dómine, servábis nos: et custódies nos in ætérnum.

\hymnusheading{Hymnus}

\lettrine[lines=2]{T}{E} lucis ante términum,\\
Rerum Creátor, póscimus,\\
Ut pro tua cleméntia\\
Sis præsul et custódia.

Procul recédant sómnia,\\
Et nóctium phantásmata,\\
Hostémque nostrum cómprime,\\
Ne polluántur córpora.

Præsta, Pater piíssime,\\
Patríque compar Unice,\\
Cum Spíritu Paráclito\\
Regnans per omne sæculum.\\
Amen.

\capitulumlabel{Capitulum}{Jer. 14, 9}

\lettrine[lines=2]{T}{U} autem in nobis es, Dómine, et nomen sanctum tuum invocátum est super nos: ne derelínquas nos, Dómine, Deus noster.

\R\ br. In manus tuas, Dómine, * Comméndo spíritum meum. In manus. \V\ Redemísti nos, Dómine, Deus veritátis. Comméndo. Glória Patri. In manus.

\V\ Custódi nos, Dómine, ut pupíllam óculi. \R\ Sub umbra alárum tuárum prótege nos.

\antlabel{Ant.} Salva nos.

\canticumheading{Canticum Simeonis}{Luc. 2, 29-32}

\lettrine[lines=2]{N}{UNC} dimíttis servum tuum, Dómine, * secúndum verbum tuum in pace:

Quia vidérunt óculi mei * salutáre tuum,

Quod parásti * ante fáciem ómnium populórum,

Lumen ad revelatiónem géntium, * et glóriam plebis tuæ Israël.

Glória Patri.

\antlabel{Ant.} Salva nos, Dómine, vigilántes, custódi nos dormiéntes; ut vigilémus cum Christo, et requiescámus in pace.

\variantheading{Preces (si dicendæ sint)}

KYRIE, eléison. Christe, eléison. Kýrie, eléison.

Pater noster \rubric{secreto usque ad} \V\ Et ne nos indúcas in tentatiónem. \R\ Sed líbera nos a malo.

Credo in Deum \rubric{secreto usque ad} \V\ Carnis resurrectiónem. \R\ Vitam ætérnam. Amen. \V\ Benedíctus es, Dómine, Deus patrum nostrórum. \R\ Et laudábilis et gloriósus in sæcula. \V\ Benedicámus Patrem et Fílium cum Sancto Spíritu. \R\ Laudémus, et superexaltémus eum in sæcula. \V\ Benedíctus es, Dómine, in firmaménto cæli. \R\ Et laudábilis, et gloriósus, et superexaltátus in sæcula. \V\ Benedícat et custódiat nos omnípotens et miséricors Dóminus. \R\ Amen. \V\ Dignáre, Dómine, nocte ista. \R\ Sine peccáto nos custodíre. \V\ Miserére nostri, Dómine. \R\ Miserére nostri. \V\ Fiat misericórdia tua, Dómine, super nos. \R\ Quemádmodum sperávimus in te. \V\ Dómine, exáudi oratiónem meam. \R\ Et clamor meus ad te véniat.

\V\ Dóminus vobíscum.

Orémus. \oratiolabel{Oratio}

\lettrine[lines=2]{V}{ISITA}, quæsumus, Dómine, habitatiónem istam, et omnes insídias inimíci ab ea longe repélle: Angeli tui sancti hábitent in ea, qui nos in pace custódiant; et benedíctio tua sit super nos semper. Per Dóminum nostrum. \R\ Amen.

\V\ Dóminus vobíscum.

\V\ Benedicámus Dómino.

\benedictiolabel{Benedictio} Benedícat et custódiat nos omnípotens et miséricors Dóminus, Pater, et Fílius, et Spíritus Sanctus. \R\ Amen.

\variantheading{Antiphona finalis B. M. V.}

\V\ Divínum auxílium.

Pater, Ave et Credo.

% ===== END OF BATCH 18 -- Covers autumnalis.md lines 5981-6274,
% completing "Ad Vesperas" (the Gradual Psalms 122-126, Capitulum,
% hymn "Telluris alme Cónditor") and "Ad Completorium" (Psalms 11,
% 12, 15, hymn "Te lucis ante términum," Canticum Simeonis, Preces),
% and with them the ENTIRE "Feria Tertia" Office. Two more
% occurrences of the recurring stray OCR "\*" artifact (in "Telluris
% alme Cónditor" at autumnalis.md line 6109 and "Te lucis ante
% términum" at line 6217, both at the hymn's final doxology stanza)
% are silently dropped. The partial-bold "Confíteor." (not the full
% "Misereátur. Indulgéntiam.") again appears at Compline (line
% 6137), matching the alternation already documented in Feria
% Secunda's Compline, and is rendered the same way:
% \formulalabel{Confíteor.} followed by plain text. The document
% continues with "Feria Quarta" (autumnalis.md line 6277), the third
% weekday of the Psalterium cycle.
% =====
\feriaheading{Feria Quarta}

\hourheading{Ad Matutinum}

Pater, Ave et Credo.

\V\ Dómine, lábia mea.

\V\ Deus, in adjutórium.

\variantheading{In Officio feriali}

\invitatlabel{Invitat.} Deum magnum Dóminum, * Veníte, adorémus. \psalmlabel{Ps. 94} Veníte.

\hymnusheading{Hymnus}

\lettrine[lines=2]{R}{ERUM} Creátor óptime,\\
Rectórque noster, áspice:\\
Nos a quiéte nóxia\\
Mersos sopóre líbera.

Te, sancte Christe, póscimus,\\
Ignósce culpis ómnibus:\\
Ad confiténdum súrgimus,\\
Morásque noctis rúmpimus.

Mentes manúsque tóllimus,\\
Prophéta sicut nóctibus\\
Nobis geréndum præcipit,\\
Paulúsque gestis cénsuit.

Vides malum, quod fécimus:\\
Occúlta nostra pándimus:\\
Preces geméntes fúndimus,\\
Dimítte quod peccávimus.

Præsta, Pater piíssime,\\
Patríque compar Unice,\\
Cum Spíritu Paráclito\\
Regnans per omne sæculum.\\
Amen.

\nocturnoheading{In I Nocturno}

\antlabel{Ant.} Speciósus forma * præ fíliis hóminum, diffúsa est grátia in lábiis tuis.

\psalmheading{Psalmus 44, i}

\lettrine[lines=2]{E}{RUCTAVIT} cor meum verbum bonum: * dico ego ópera mea Regi.

Lingua mea cálamus scribæ: * velóciter scribéntis.

Speciósus forma præ fíliis hóminum, diffúsa est grátia in lábiis tuis: * proptérea benedíxit te Deus in ætérnum.

Accíngere gládio tuo super femur tuum, * potentíssime.

Spécie tua et pulchritúdine tua * inténde, próspere procéde, et regna.

Propter veritátem, et mansuetúdinem, et justítiam: * et dedúcet te mirabíliter déxtera tua.

Sagíttæ tuæ acútæ, pópuli sub te cadent, * in corda inimicórum Regis.

Sedes tua, Deus, in sæculum sæculi: * virga directiónis virga regni tui.

Dilexísti justítiam, et odísti iniquitátem: * proptérea unxit te Deus, Deus tuus, óleo lætítiæ præ consórtibus tuis.

Myrrha, et gutta, et cásia a vestiméntis tuis, a dómibus ebúrneis: * ex quibus delectavérunt te fíliæ regum in honóre tuo.

Astitit regína a dextris tuis in vestítu deauráto: * circúmdata varietáte.

\antlabel{Ant.} Speciósus forma præ fíliis hóminum, diffúsa est grátia in lábiis tuis.

\antlabel{Ant.} Confitebúntur tibi * pópuli, Deus, in ætérnum.

\psalmheading{Psalmus 44, ii}

\lettrine[lines=2]{A}{UDI}, fília, et vide, et inclína aurem tuam: * et oblivíscere pópulum tuum, et domum patris tui.

Et concupíscet Rex decórem tuum: * quóniam ipse est Dóminus, Deus tuus, et adorábunt eum.

Et fíliæ Tyri in munéribus * vultum tuum deprecabúntur: omnes dívites plebis.

Omnis glória ejus fíliæ Regis ab intus, * in fímbriis áureis circumamícta varietátibus.

Adducéntur Regi vírgines post eam: * próximæ ejus afferéntur tibi.

Afferéntur in lætítia et exsultatióne: * adducéntur in templum Regis.

Pro pátribus tuis nati sunt tibi fílii: * constítues eos príncipes super omnem terram.

Mémores erunt nóminis tui: * in omni generatióne et generatiónem.

Proptérea pópuli confitebúntur tibi in ætérnum: * et in sæculum sæculi.

\antlabel{Ant.} Confitebúntur tibi pópuli, Deus, in ætérnum.

\antlabel{Ant.} Adjútor in tribulatiónibus * Deus noster.

\psalmheading{Psalmus 45}

\lettrine[lines=2]{D}{EUS} noster refúgium, et virtus: * adjútor in tribulatiónibus, quæ invenérunt nos nimis.

Proptérea non timébimus dum turbábitur terra: * et transferéntur montes in cor maris.

Sonuérunt, et turbátæ sunt aquæ eórum: * conturbáti sunt montes in fortitúdine ejus.

Flúminis ímpetus lætíficat civitátem Dei: * sanctificávit tabernáculum suum Altíssimus.

Deus in médio ejus, non commovébitur: * adjuvábit eam Deus mane dilúculo.

Conturbátæ sunt gentes, et inclináta sunt regna: * dedit vocem suam, mota est terra.

Dóminus virtútum nobíscum: * suscéptor noster Deus Jacob.

Veníte, et vidéte ópera Dómini, quæ pósuit prodígia super terram: * áuferens bella usque ad finem terræ.

Arcum cónteret, et confrínget arma: * et scuta combúret igni.

Vacáte, et vidéte quóniam ego sum Deus: * exaltábor in géntibus, et exaltábor in terra.

Dóminus virtútum nobíscum: * suscéptor noster Deus Jacob.

\antlabel{Ant.} Adjútor in tribulatiónibus Deus noster.

\variantheading{In Festis novem Lectionum:}

\V\ Dóminus virtútum nobíscum. \R\ Suscéptor noster Deus Jacob.

\absolutiolabel{Absolutio.} Exáudi, Dómine.

\benedlabel{Bened.} 1 Benedictióne perpétua. 2 Unigénitus Dei Fílius. 3 Spíritus Sancti grátia.

\nocturnoheading{In II Nocturno}

\antlabel{Ant.} Magnus Dóminus, * \textdagger\ et laudábilis nimis in civitáte Dei nostri. \textdaggerdbl

\psalmheading{Psalmus 47}

\lettrine[lines=2]{M}{AGNUS} Dóminus, \textdagger\ et laudábilis nimis * in civitáte Dei nostri, \textdaggerdbl\ in monte sancto ejus.

Fundátur exsultatióne univérsæ terræ mons Sion, * látera Aquilónis, cívitas Regis magni.

Deus in dómibus ejus cognoscétur, * cum suscípiet eam.

Quóniam ecce reges terræ congregáti sunt: * convenérunt in unum.

Ipsi vidéntes sic admiráti sunt, conturbáti sunt, commóti sunt: * tremor apprehéndit eos.

Ibi dolóres ut parturiéntis: * in spíritu vehemménti cónteres naves Tharsis.

Sicut audívimus, sic vídimus in civitáte Dómini virtútum, in civitáte Dei nostri: * Deus fundávit eam in ætérnum.

Suscépimus, Deus, misericórdiam tuam, * in médio templi tui.

Secúndum nomen tuum, Deus, sic et laus tua in fines terræ: * justítia plena est déxtera tua.

Lætétur mons Sion, et exsúltent fíliæ Judæ: * propter judícia tua, Dómine.

Circúmdate Sion, et complectímini eam: * narráte in túrribus ejus.

Pónite corda vestra in virtúte ejus: * et distribúite domos ejus, ut enarrétis in progénie áltera.

Quóniam hic est Deus, Deus noster in ætérnum, et in sæculum sæculi: * ipse reget nos in sæcula.

\antlabel{Ant.} Magnus Dóminus, et laudábilis nimis in civitáte Dei nostri.

\antlabel{Ant.} Os meum loquétur * sapiéntiam: et meditátio cordis mei prudéntiam.

\psalmheading{Psalmus 48, i}

\lettrine[lines=2]{A}{UDITE} hæc, omnes gentes: * áuribus percípite omnes, qui habitátis orbem:

Quique terrígenæ, et fílii hóminum: * simul in unum dives et pauper.

Os meum loquétur sapiéntiam: * et meditátio cordis mei prudéntiam.

Inclinábo in parábolam aurem meam: * apériam in psaltério propositiónem meam.

Cur timébo in die mala? * iníquitas calcánei mei circúmdabit me:

Qui confídunt in virtúte sua: * et in multitúdine divitiárum suárum gloriántur.

Frater non rédimit, rédimet homo: * non dabit Deo placatiónem suam.

Et prétium redemptiónis ánimæ suæ: * et laborábit in ætérnum, et vivet adhuc in finem.

Non vidébit intéritum, cum víderit sapiéntes moriéntes: * simul insípiens, et stultus períbunt.

Et relínquent aliénis divítias suas: * et sepúlcra eórum domus illórum in ætérnum.

Tabernácula eórum in progénie et progénie: * vocavérunt nómina sua in terris suis.

Et homo, cum in honóre esset, non intelléxit: * comparátus est juméntis insipiéntibus, et símilis factus est illis.

\antlabel{Ant.} Os meum loquétur sapiéntiam: et meditátio cordis mei prudéntiam.

\antlabel{Ant.} Ne timúeris: * cum dívite non descéndet in sepúlcrum glória ejus.

\psalmheading{Psalmus 48, ii}

\lettrine[lines=2]{H}{ÆC} via illórum scándalum ipsis: * et póstea in ore suo complacébunt.

Sicut oves in inférno pósiti sunt: * mors depáscet eos.

Et dominabúntur eórum justi in matutíno: * et auxílium eórum veteráscet in inférno a glória eórum.

Verúmtamen Deus rédimet ánimam meam de manu ínferi, * cum accéperit me.

Ne timúeris, cum dives factus fúerit homo: * et cum multiplicáta fúerit glória domus ejus.

Quóniam cum interíerit, non sumet ómnia: * neque descéndet cum eo glória ejus.

Quia ánima ejus in vita ipsíus benedicétur: * confitébitur tibi cum beneféceris ei.

Introíbit usque in progénies patrum suórum: * et usque in ætérnum non vidébit lumen.

Homo, cum in honóre esset, non intelléxit: * comparátus est juméntis insipiéntibus, et símilis factus est illis.

\antlabel{Ant.} Ne timúeris: cum dívite non descéndet in sepúlcrum glória ejus.

\variantheading{In Festis novem Lectionum:}

\V\ Deus rédimet ánimam meam de manu ínferi. \R\ Cum accéperit me.

\absolutiolabel{Absolutio.} Ipsíus píetas et.

\benedlabel{Bened.} 4 Deus Pater omnípotens. 5 Christus perpétuæ. 6 Ignem sui amóris.

\nocturnoheading{In III Nocturno}

\variantheading{I}

\rubric{In Fériis extra Vigílias commúnes occurréntibus, et in Festis, cum ad Laudes non dícitur Psalmus 50 Miserére.}

\antlabel{Ant.} Deus deórum, * \textdagger\ Dóminus locútus est. \textdaggerdbl

\psalmheading{Psalmus 49, i}

\lettrine[lines=2]{D}{EUS} deórum, \textdagger\ Dóminus locútus est: * \textdaggerdbl\ et vocávit terram.

A solis ortu usque ad occásum: * ex Sion spécies decóris ejus.

Deus maniféste véniet: * Deus noster et non silébit.

Ignis in conspéctu ejus exardéscet: * et in circúitu ejus tempéstas válida.

Advocábit cælum desúrsum: * et terram discérnere pópulum suum.

Congregáte illi sanctos ejus: * qui órdinant testaméntum ejus super sacrifícia.

Et annuntiábunt cæli justítiam ejus: * quóniam Deus judex est.

Audi, pópulus meus, et loquar: Israël, et testificábor tibi: * Deus, Deus tuus ego sum.

Non in sacrifíciis tuis árguam te: * holocáusta autem tua in conspéctu meo sunt semper.

Non accípiam de domo tua vítulos: * neque de grégibus tuis hircos.

Quóniam meæ sunt omnes feræ silvárum, * juménta in móntibus et boves.

Cognóvi ómnia volatília cæli: * et pulchritúdo agri mecum est.

Si esuríero, non dicam tibi: * meus est enim orbis terræ, et plenitúdo ejus.

Numquid manducábo carnes taurórum? * aut sánguinem hircórum potábo?

Immola Deo sacrifícium laudis: * et redde Altíssimo vota tua.

Et ínvoca me in die tribulatiónis: * éruam te, et honorificábis me.

\antlabel{Ant.} Deus deórum, Dóminus locútus est.

\antlabel{Ant.} Intellígite, * qui obliviscímini Deum.

\psalmheading{Psalmus 49, ii}

\lettrine[lines=2]{P}{ECCATORI} autem dixit Deus: * Quare tu enárras justítias meas, et assúmis testaméntum meum per os tuum?

Tu vero odísti disciplínam: * et projecísti sermónes meos retrórsum:

Si vidébas furem, currébas cum eo: * et cum adúlteris portiónem tuam ponébas.

Os tuum abundávit malítia: * et lingua tua concinnábat dolos.

Sedens advérsus fratrem tuum loquebáris, et advérsus fílium matris tuæ ponébas scándalum: * hæc fecísti, et tácui.

Existimásti, iníque, quod ero tui símilis: * árguam te, et státuam contra fáciem tuam.

Intellígite hæc, qui oblivíscimini Deum: * nequándo rápiat, et non sit qui erípiat.

Sacrifícium laudis honorificábit me: * et illic iter, quo osténdam illi salutáre Dei.

\antlabel{Ant.} Intellígite, qui obliviscímini Deum.

\antlabel{Ant.} Acceptábis sacrifícium * justítiæ super altáre tuum, Dómine.

\psalmheading{Psalmus 50}

\lettrine[lines=2]{M}{ISERERE} mei, Deus, * secúndum magnam misericórdiam tuam.

Et secúndum multitúdinem miseratiónum tuárum, * dele iniquitátem meam.

Amplius lava me ab iniquitáte mea: * et a peccáto meo munda me.

Quóniam iniquitátem meam ego cognósco: * et peccátum meum contra me est semper.

Tibi soli peccávi, et malum coram te feci: * ut justificéris in sermónibus tuis, et vincas cum judicáris.

Ecce enim, in iniquitátibus concéptus sum: * et in peccátis concépit me mater mea.

Ecce enim, veritátem dilexísti: * incérta et occúlta sapiéntiæ tuæ manifestásti mihi.

Aspérges me hyssópo, et mundábor: * lavábis me, et super nivem dealbábor.

Audítui meo dabis gáudium et lætítiam: * et exsultábunt ossa humiliáta.

Avérte fáciem tuam a peccátis meis: * et omnes iniquitátes meas dele.

Cor mundum crea in me, Deus: * et spíritum rectum ínnova in viscéribus meis.

Ne projícias me a fácie tua: * et spíritum sanctum tuum ne áuferas a me.

Redde mihi lætítiam salutáris tui: * et spíritu principáli confírma me.

Docébo iníquos vias tuas: * et ímpii ad te converténtur.

Líbera me de sanguínibus, Deus, Deus salútis meæ: * et exsultábit lingua mea justítiam tuam.

Dómine, lábia mea apéries: * et os meum annuntiábit laudem tuam.

Quóniam si voluísses sacrifícium, dedíssem útique: * holocáustis non delectáberis.

Sacrifícium Deo spíritus contribulátus: * cor contrítum, et humiliátum, Deus, non despícies.

Benígne fac, Dómine, in bona voluntáte tua Sion: * ut ædificéntur muri Jerúsalem.

Tunc acceptábis sacrifícium justítiæ, oblatiónes, et holocáusta: * tunc impónent super altáre tuum vítulos.

\antlabel{Ant.} Acceptábis sacrifícium justítiæ super altáre tuum, Dómine.

\variantheading{In feriali Officio et in Festis:}

\V\ Audítui meo dabis gáudium et lætítiam. \R\ Et exsultábunt ossa humiliáta.

\variantheading{In Officio novem Lectionum}

\absolutiolabel{Absolutio.} A vínculis.

\benedlabel{Bened.} 7 Evangélica léctio. 8 In Festis Dómini: Divínum auxílium. In Festis B. M. V.: Cujus festum cólimus, ipsa Virgo. In Festis Sanctórum: Cujus (\textit{vel} Quorum \textit{aut} Quarum) festum cólimus. 9 Ad societátem. (Per evangélica dicta.)

\variantheading{In Officio trium Lectionum}

\absolutiolabel{Absolutio.} A vínculis.

\variantheading{In Feriis cum Homilia:}

\benedlabel{Bened.} 1 Evangélica. 2 Divínum. 3 Ad societátem.

\variantheading{In Feriis sine Homilia:}

\benedlabel{Bened.} 1 Ille nos benedícat. 2 Divínum. 3 Ad societátem.

\variantheading{In Festis Sanctorum:}

\benedlabel{Bened.} 1 Ille nos benedícat, qui sine fine vivit et regnat. \R\ Amen. 2 Cujus (\textit{vel} Quorum \textit{aut} Quarum). 3 Ad societátem.

\nocturnoheading{In III Nocturno}

\variantheading{II}

\rubric{In Vigíliis commúnibus, cum ad Laudes dícitur Psalmus 50 Miserére.}

\antlabel{Ant.} Deus deórum. \textdagger

\psalmheading{Psalmus 49, i}

\lettrine[lines=2]{D}{EUS} deórum, \textdagger\ Dóminus locútus est: * \textdaggerdbl\ et vocávit terram.

A solis ortu usque ad occásum: * ex Sion spécies decóris ejus.

Deus maniféste véniet: * Deus noster et non silébit.

Ignis in conspéctu ejus exardéscet: * et in circúitu ejus tempéstas válida.

Advocábit cælum desúrsum: * et terram discérnere pópulum suum.

Congregáte illi sanctos ejus: * qui órdinant testaméntum ejus super sacrifícia.

Et annuntiábunt cæli justítiam ejus: * quóniam Deus judex est.

Audi, pópulus meus, et loquar: Israël, et testificábor tibi: * Deus, Deus tuus ego sum.

Non in sacrifíciis tuis árguam te: * holocáusta autem tua in conspéctu meo sunt semper.

Non accípiam de domo tua vítulos: * neque de grégibus tuis hircos.

Quóniam meæ sunt omnes feræ silvárum, * juménta in móntibus et boves.

Cognóvi ómnia volatília cæli: * et pulchritúdo agri mecum est.

Si esuríero, non dicam tibi: * meus est enim orbis terræ, et plenitúdo ejus.

Numquid manducábo carnes taurórum? * aut sánguinem hircórum potábo?

Immola Deo sacrifícium laudis: * et redde Altíssimo vota tua.

Et ínvoca me in die tribulatiónis: * éruam te, et honorificábis me.

\antlabel{Ant.} Deus deórum, Dóminus locútus est.

\antlabel{Ant.} Immola Deo.

\psalmheading{Psalmus 49, ii}

\lettrine[lines=2]{P}{ECCATORI} autem dixit Deus: * Quare tu enárras justítias meas, et assúmis testaméntum meum per os tuum?

Tu vero odísti disciplínam: * et projecísti sermónes meos retrórsum:

Si vidébas furem, currébas cum eo: * et cum adúlteris portiónem tuam ponébas.

Os tuum abundávit malítia: * et lingua tua concinnábat dolos.

Sedens advérsus fratrem tuum loquebáris, et advérsus fílium matris tuæ ponébas scándalum: * hæc fecísti, et tácui.

Existimásti, iníque, quod ero tui símilis: * árguam te, et státuam contra fáciem tuam.

Intellígite hæc, qui oblivíscimini Deum: * nequándo rápiat, et non sit qui erípiat.

Sacrifícium laudis honorificábit me: * et illic iter, quo osténdam illi salutáre Dei.

\antlabel{Ant.} Immola Deo sacrifícium laudis.

\antlabel{Ant.} Intellígite.

\psalmheading{Psalmus 49, iii}

\lettrine[lines=2]{P}{ECCATORI} autem dixit Deus: * Quare tu enárras justítias meas, et assúmis testaméntum meum per os tuum?

Tu vero odísti disciplínam: * et projecísti sermónes meos retrórsum:

Si vidébas furem, currébas cum eo: * et cum adúlteris portiónem tuam ponébas.

Os tuum abundávit malítia: * et lingua tua concinnábat dolos.

Sedens advérsus fratrem tuum loquebáris, et advérsus fílium matris tuæ ponébas scándalum: * hæc fecísti, et tácui.

Existimásti, iníque, quod ero tui símilis: * árguam te, et státuam contra fáciem tuam.

Intellígite hæc, qui oblivíscimini Deum: * nequándo rápiat, et non sit qui erípiat.

Sacrifícium laudis honorificábit me: * et illic iter, quo osténdam illi salutáre Dei.

\antlabel{Ant.} Intellígite, qui obliviscímini Deum.

\V\ Sacrifícium laudis honorificábit me. \R\ Et illic iter, quo osténdam illi salutáre Dei.

\absolutiolabel{Absolutio.} A vínculis.

\variantheading{In Feriis cum Homilia:}

\benedlabel{Bened.} 1 Evangélica. 2 Divínum. 3 Ad societátem.

\variantheading{In Feriis sine Homilia:}

\benedlabel{Bened.} 1 Ille nos benedícat. 2 Divínum. 3 Ad societátem.

% ===== END OF BATCH 19 -- Covers autumnalis.md lines 6277-6743,
% opening "Feria Quarta" (Wednesday) and completing its "Ad
% Matutinum" in full: the ferial Invitatorium "Deum magnum
% Dóminum," the hymn "Rerum Creátor óptime" (stray OCR "\*"
% silently dropped), and the three Nocturns (Psalms 44/i, 44/ii,
% 45; 47, 48/i, 48/ii). "In III Nocturno" here introduces, for the
% first time in this volume, two complete alternate forms labeled
% "I" and "II," each with its own italicized rubric describing when
% it applies (Form I when Ps. 50 was NOT said at Lauds, using
% Psalms 49/i, 49/ii, and the full Ps. 50; Form II when Ps. 50 WAS
% said at Lauds, using Psalms 49/i, 49/ii, and 49/iii instead) --
% rendered with \variantheading{I} and \variantheading{II}, matching
% the identical command already used for this exact structure in
% both verna.tex and hiemalis.tex. Notably, Form II's "Psalmus 49,
% ii" and "Psalmus 49, iii" give verbatim identical text in
% autumnalis.md itself (both "Peccátori autem dixit Deus..."); this is
% transcribed exactly as the source repeats it, without alteration,
% since it is a structural feature of the source and not a
% character-level OCR defect. The document continues with "Ad
% Laudes I" (autumnalis.md line 6745).
% =====
\hourheading{Ad Laudes I}

\rubric{In Fériis extra Vigílias commúnes occurréntibus, atque in Festis.}

(Pater noster et Ave.)

\V\ Deus, in adjutórium.

\antlabel{Ant.} Dóminus regnávit, * \textdagger\ exsúltet terra. \textdaggerdbl

\psalmheading{Psalmus 96}

\lettrine[lines=2]{D}{OMINUS} regnávit, \textdagger\ exsúltet terra: * \textdaggerdbl\ læténtur ínsulæ multæ.

Nubes, et calígo in circúitu ejus: * justítia, et judícium corréctio sedis ejus.

Ignis ante ipsum præcédet, * et inflammábit in circúitu inimícos ejus.

Illuxérunt fúlgura ejus orbi terræ: * vidit, et commóta est terra.

Montes, sicut cera fluxérunt a fácie Dómini: * a fácie Dómini omnis terra.

Annuntiavérunt cæli justítiam ejus: * et vidérunt omnes pópuli glóriam ejus.

Confundántur omnes, qui adórant sculptília: * et qui gloriántur in simulácris suis.

Adoráte eum, omnes Angeli ejus: * audívit, et lætáta est Sion.

Et exsultavérunt fíliæ Judæ, * propter judícia tua, Dómine:

Quóniam tu Dóminus Altíssimus super omnem terram: * nimis exaltátus es super omnes deos.

Qui dilígitis Dóminum, odíte malum: * custódit Dóminus ánimas sanctórum suórum, de manu peccatóris liberábit eos.

Lux orta est justo, * et rectis corde lætítia.

Lætámini, justi, in Dómino: * et confitémini memóriæ sanctificatiónis ejus.

\antlabel{Ant.} Dóminus regnávit, exsúltet terra.

\antlabel{Ant.} Te decet hymnus, * \textdagger\ Deus, in Sion. \textdaggerdbl

\psalmheading{Psalmus 64}

\lettrine[lines=2]{T}{E} decet hymnus, \textdagger\ Deus, in Sion: * \textdaggerdbl\ et tibi reddétur votum in Jerúsalem.

Exáudi oratiónem meam: * ad te omnis caro véniet.

Verba iniquórum prævaluérunt super nos: * et impietátibus nostris tu propitiáberis.

Beátus, quem elegísti, et assumpsísti: * inhabitábit in átriis tuis.

Replébimur in bonis domus tuæ: * sanctum est templum tuum, mirábile in æquitáte.

Exáudi nos, Deus, salutáris noster, * spes ómnium fínium terræ, et in mari longe.

Præparans montes in virtúte tua, accínctus poténtia: * qui contúrbas profúndum maris sonum flúctuum ejus.

Turbabúntur gentes, et timébunt qui hábitant términos a signis tuis: * éxitus matutíni, et véspere delectábis.

Visitásti terram, et inebriásti eam: * multiplicásti locupletáre eam.

Flumen Dei replétum est aquis, parásti cibum illórum: * quóniam ita est præparátio ejus.

Rivos ejus inébria, multíplica genímina ejus: * in stillicídiis ejus lætábitur gérminans.

Benedíces corónæ anni benignitátis tuæ: * et campi tui replebúntur ubertáte.

Pinguéscent speciósa desérti: * et exsultatióne colles accingéntur.

Indúti sunt aríetes óvium, et valles abundábunt fruménto: * clamábunt, étenim hymnum dicent.

\antlabel{Ant.} Te decet hymnus, Deus, in Sion.

\antlabel{Ant.} Tibi, Dómine, psallam, * et intélligam in via immaculáta.

\psalmheading{Psalmus 100}

\lettrine[lines=2]{M}{ISERICORDIAM}, et judícium * cantábo tibi, Dómine:

Psallam, et intélligam in via immaculáta, * quando vénies ad me.

Perambulábam in innocéntia cordis mei, * in médio domus meæ.

Non proponébam ante óculos meos rem injústam: * faciéntes prævaricatiónes odívi.

Non adhæsit mihi cor pravum: * declinántem a me malígnum non cognoscébam.

Detrahéntem secréto próximo suo, * hunc persequébar.

Supérbo óculo, et insatiábili corde, * cum hoc non edébam.

Oculi mei ad fidéles terræ ut sédeant mecum: * ámbulans in via immaculáta, hic mihi ministrábat.

Non habitábit in médio domus meæ qui facit supérbiam: * qui lóquitur iníqua, non diréxit in conspéctu oculórum meórum.

In matutíno interficiébam omnes peccatóres terræ: * ut dispérderem de civitáte Dómini omnes operántes iniquitátem.

\antlabel{Ant.} Tibi, Dómine, psallam, et intélligam in via immaculáta.

\antlabel{Ant.} Dómine, magnus es tu, * et præclárus in virtúte tua.

\canticumheading{Canticum Judith}{Judith. 16, 15-21}

\lettrine[lines=2]{H}{YMNUM} cantémus Dómino, * hymnum novum cantémus Deo nostro.

Adonái, Dómine, magnus es tu, et præclárus in virtúte tua, * et quem superáre nemo potest.

Tibi sérviat omnis creatúra tua: * quia dixísti, et facta sunt:

Misísti spíritum tuum, et creáta sunt: * et non est qui resístat voci tuæ.

Montes a fundaméntis movebúntur cum aquis: * petræ, sicut cera, liquéscent ante fáciem tuam.

Qui autem timent te, * magni erunt apud te per ómnia.

Væ genti insurgénti super genus meum: Dóminus enim omnípotens vindicábit in eis, * in die judícii visitábit illos.

Dabit enim ignem, et vermes in carnes eórum, * ut urántur, et séntiant usque in sempitérnum.

\antlabel{Ant.} Dómine, magnus es tu, et præclárus in virtúte tua.

\antlabel{Ant.} Laudábo Deum meum * in vita mea.

\psalmheading{Psalmus 145}

\lettrine[lines=2]{L}{AUDA}, ánima mea, Dóminum, laudábo Dóminum in vita mea: * psallam Deo meo quámdiu fúero.

Nolíte confídere in princípibus: * in fíliis hóminum, in quibus non est salus.

Exíbit spíritus ejus, et revertétur in terram suam: * in illa die períbunt omnes cogitatiónes eórum.

Beátus, cujus Deus Jacob adjútor ejus, spes ejus in Dómino, Deo ipsíus: * qui fecit cælum et terram, mare, et ómnia, quæ in eis sunt.

Qui custódit veritátem in sæculum, facit judícium injúriam patiéntibus: * dat escam esuriéntibus.

Dóminus solvit compedítos: * Dóminus illúminat cæcos.

Dóminus érigit elísos, * Dóminus díligit justos.

Dóminus custódit ádvenas, pupíllum et víduam suscípiet: * et vias peccatórum dispérdet.

Regnábit Dóminus in sæcula, Deus tuus, Sion, * in generatiónem et generatiónem.

\antlabel{Ant.} Laudábo Deum meum in vita mea.

\variantheading{In Officio feriali}

\capitulumlabel{Capit.}{Rom. 13, 12-13}

\lettrine[lines=2]{N}{OX} præcéssit, dies autem appropinquávit. Abjiciámus ergo ópera tenebrárum, et induámur arma lucis. Sicut in die honéste ambulémus.

\hymnusheading{Hymnus}

\lettrine[lines=2]{N}{OX}, et ténebræ, et núbila,\\
Confúsa mundi et túrbida:\\
Lux intrat, albéscit polus:\\
Christus venit: discédite.

Calígo terræ scínditur\\
Percússa solis spículo,\\
Rebúsque jam color redit,\\
Vultu niténtis síderis.

Te, Christe, solum nóvimus:\\
Te mente pura et símplici,\\
Flendo et canéndo quǽsumus,\\
Inténde nostris sénsibus.

Sunt multa fucis illíta,\\
Quæ luce purgéntur tua:\\
Tu, vera lux Cæléstium,\\
Vultu seréno illúmina.

Deo Patri sit glória,\\
Ejúsque soli Fílio,\\
Cum Spíritu Paráclito,\\
Nunc et per omne sæculum.\\
Amen.

\V\ Repléti sumus mane misericórdia tua. \R\ Exsultávimus, et delectáti sumus.

\antlabel{Ad Bened. Ant.} De manu ómnium, * qui odérunt nos, liberávit nos Dóminus.

\rubric{Réliqua ut in Ordinário. 11}

\hourheading{Ad Laudes II}

\rubric{In Vigíliis commúnibus, quando fit Offícium de Féria.}

(Pater noster et Ave.)

\V\ Deus, in adjutórium.

\antlabel{Ant.} Amplius lava me.

\psalmheading{Psalmus 50}

\lettrine[lines=2]{M}{ISERERE} mei, Deus, * secúndum magnam misericórdiam tuam.

Et secúndum multitúdinem miseratiónum tuárum, * dele iniquitátem meam.

Amplius lava me ab iniquitáte mea: * et a peccáto meo munda me.

Quóniam iniquitátem meam ego cognósco: * et peccátum meum contra me est semper.

Tibi soli peccávi, et malum coram te feci: * ut justificéris in sermónibus tuis, et vincas cum judicáris.

Ecce enim, in iniquitátibus concéptus sum: * et in peccátis concépit me mater mea.

Ecce enim, veritátem dilexísti: * incérta et occúlta sapiéntiæ tuæ manifestásti mihi.

Aspérges me hyssópo, et mundábor: * lavábis me, et super nivem dealbábor.

Audítui meo dabis gáudium et lætítiam: * et exsultábunt ossa humiliáta.

Avérte fáciem tuam a peccátis meis: * et omnes iniquitátes meas dele.

Cor mundum crea in me, Deus: * et spíritum rectum ínnova in viscéribus meis.

Ne projícias me a fácie tua: * et spíritum sanctum tuum ne áuferas a me.

Redde mihi lætítiam salutáris tui: * et spíritu principáli confírma me.

Docébo iníquos vias tuas: * et ímpii ad te converténtur.

Líbera me de sanguínibus, Deus, Deus salútis meæ: * et exsultábit lingua mea justítiam tuam.

Dómine, lábia mea apéries: * et os meum annuntiábit laudem tuam.

Quóniam si voluísses sacrifícium, dedíssem útique: * holocáustis non delectáberis.

Sacrifícium Deo spíritus contribulátus: * cor contrítum, et humiliátum, Deus, non despícies.

Benígne fac, Dómine, in bona voluntáte tua Sion: * ut ædificéntur muri Jerúsalem.

Tunc acceptábis sacrifícium justítiæ, oblatiónes, et holocáusta: * tunc impónent super altáre tuum vítulos.

\antlabel{Ant.} Amplius lava me, Dómine, ab injustítia mea.

\antlabel{Ant.} Impietátibus nostris.

\psalmheading{Psalmus 64}

\lettrine[lines=2]{T}{E} decet hymnus, Deus, in Sion: * et tibi reddétur votum in Jerúsalem.

Exáudi oratiónem meam: * ad te omnis caro véniet.

Verba iniquórum prævaluérunt super nos: * et impietátibus nostris tu propitiáberis.

Beátus, quem elegísti, et assumpsísti: * inhabitábit in átriis tuis.

Replébimur in bonis domus tuæ: * sanctum est templum tuum, mirábile in æquitáte.

Exáudi nos, Deus, salutáris noster, * spes ómnium fínium terræ, et in mari longe.

Præparans montes in virtúte tua, accínctus poténtia: * qui contúrbas profúndum maris sonum flúctuum ejus.

Turbabúntur gentes, et timébunt qui hábitant términos a signis tuis: * éxitus matutíni, et véspere delectábis.

Visitásti terram, et inebriásti eam: * multiplicásti locupletáre eam.

Flumen Dei replétum est aquis, parásti cibum illórum: * quóniam ita est præparátio ejus.

Rivos ejus inébria, multíplica genímina ejus: * in stillicídiis ejus lætábitur gérminans.

Benedíces corónæ anni benignitátis tuæ: * et campi tui replebúntur ubertáte.

Pinguéscent speciósa desérti: * et exsultatióne colles accingéntur.

Indúti sunt aríetes óvium, et valles abundábunt fruménto: * clamábunt, étenim hymnum dicent.

\antlabel{Ant.} Impietátibus nostris tu propitiáberis, Deus.

\antlabel{Ant.} In innocéntia.

\psalmheading{Psalmus 100}

\lettrine[lines=2]{M}{ISERICORDIAM}, et judícium * cantábo tibi, Dómine:

Psallam, et intélligam in via immaculáta, * quando vénies ad me.

Perambulábam in innocéntia cordis mei, * in médio domus meæ.

Non proponébam ante óculos meos rem injústam: * faciéntes prævaricatiónes odívi.

Non adhæsit mihi cor pravum: * declinántem a me malígnum non cognoscébam.

Detrahéntem secréto próximo suo, * hunc persequébar.

Supérbo óculo, et insatiábili corde, * cum hoc non edébam.

Oculi mei ad fidéles terræ ut sédeant mecum: * ámbulans in via immaculáta, hic mihi ministrábat.

Non habitábit in médio domus meæ qui facit supérbiam: * qui lóquitur iníqua, non diréxit in conspéctu oculórum meórum.

In matutíno interficiébam omnes peccatóres terræ: * ut dispérderem de civitáte Dómini omnes operántes iniquitátem.

\antlabel{Ant.} In innocéntia cordis mei perambulábo, Dómine.

\antlabel{Ant.} Exsultávit cor meum. \textdagger

\canticumheading{Canticum Annæ}{1 Reg. 2, 1-10}

\lettrine[lines=2]{E}{XSULTAVIT} cor meum \textdagger\ in Dómino, * et exaltátum est cornu meum in Deo meo.

Dilatátum est os meum super inimícos meos: * quia lætáta sum in salutári tuo.

Non est sanctus, ut est Dóminus: neque enim est álius extra te, * et non est fortis sicut Deus noster.

Nolíte multiplicáre loqui sublímia, * gloriántes:

Recédant vétera de ore vestro: quia Deus sciéntiárum, Dóminus est, * et ipsi præparántur cogitatiónes.

Arcus fórtium superátus est, * et infírmi accíncti sunt róbore.

Repléti prius, pro pánibus se locavérunt: * et famélici saturáti sunt.

Donec stérilis péperit plúrimos: * et quæ multos habébat fílios, infirmáta est.

Dóminus mortíficat et vivíficat, * dedúcit ad ínferos et redúcit.

Dóminus páuperem facit et ditat, * humíliat et súblevat.

Súscitat de púlvere egénum, * et de stércore élevat páuperem:

Ut sédeat cum princípibus, * et sólium glóriæ téneat.

Dómini enim sunt cárdines terræ, * et pósuit super eos orbem.

Pedes sanctórum suórum servábit, et ímpii in ténebris conticéscent: * quia non in fortitúdine sua roborábitur vir.

Dóminum formidábunt adversárii ejus: * et super ipsos in cælis tonábit:

Dóminus judicábit fines terræ, et dabit impérium regi suo, * et sublimábit cornu Christi sui.

\antlabel{Ant.} Exsultávit cor meum in Dómino, qui humíliat et súblevat.

\antlabel{Ant.} Lauda. \textdagger

\psalmheading{Psalmus 145}

\lettrine[lines=2]{L}{AUDA}, \textdagger\ ánima mea, Dóminum, laudábo Dóminum in vita mea: * psallam Deo meo quámdiu fúero.

Nolíte confídere in princípibus: * in fíliis hóminum, in quibus non est salus.

Exíbit spíritus ejus, et revertétur in terram suam: * in illa die períbunt omnes cogitatiónes eórum.

Beátus, cujus Deus Jacob adjútor ejus, spes ejus in Dómino, Deo ipsíus: * qui fecit cælum et terram, mare, et ómnia, quæ in eis sunt.

Qui custódit veritátem in sæculum, facit judícium injúriam patiéntibus: * dat escam esuriéntibus.

Dóminus solvit compedítos: * Dóminus illúminat cæcos.

Dóminus érigit elísos, * Dóminus díligit justos.

Dóminus custódit ádvenas, pupíllum et víduam suscípiet: * et vias peccatórum dispérdet.

Regnábit Dóminus in sæcula, Deus tuus, Sion, * in generatiónem et generatiónem.

\antlabel{Ant.} Lauda, ánima mea, Dóminum, qui érigit elísos, et díligit justos.

\capitulumlabel{Capit.}{Rom. 13, 12-13}

\lettrine[lines=2]{N}{OX} præcéssit, dies autem appropinquávit. Abjiciámus ergo ópera tenebrárum, et induámur arma lucis. Sicut in die honéste ambulémus.

\hymnusheading{Hymnus}

\lettrine[lines=2]{N}{OX}, et ténebræ, et núbila,\\
Confúsa mundi et túrbida:\\
Lux intrat, albéscit polus:\\
Christus venit: discédite.

Calígo terræ scínditur\\
Percússa solis spículo,\\
Rebúsque jam color redit,\\
Vultu niténtis síderis.

Te, Christe, solum nóvimus:\\
Te mente pura et símplici,\\
Flendo et canéndo quǽsumus,\\
Inténde nostris sénsibus.

Sunt multa fucis illíta,\\
Quæ luce purgéntur tua:\\
Tu, vera lux Cæléstium,\\
Vultu seréno illúmina.

Deo Patri sit glória,\\
Ejúsque soli Fílio,\\
Cum Spíritu Paráclito,\\
Nunc et per omne sæculum.\\
Amen.

\V\ Repléti sumus mane misericórdia tua. \R\ Exsultávimus, et delectáti sumus.

\antlabel{Ad Bened. Ant.} De manu ómnium, * qui odérunt nos, liberávit nos Dóminus.

\rubric{Expléto Cántico ac repetíta Antíphona, dicúntur sequéntes}

\variantheading{Preces feriales}

KYRIE, eléison. Christe, eléison. Kýrie, eléison.

Pater noster \rubric{quod dícitur a solo Hebdomadário totum clara voce, usque ad} \V\ Et ne nos indúcas in tentatiónem. \R\ Sed líbera nos a malo. \V\ Ego dixi: Dómine, miserére mei. \R\ Sana ánimam meam, quia peccávi tibi. \V\ Convértere, Dómine, úsquequo? \R\ Et deprecábilis esto super servos tuos. \V\ Fiat misericórdia tua, Dómine, super nos. \R\ Quemádmodum sperávimus in te. \V\ Sacerdótes tui induántur justítiam. \R\ Et sancti tui exsúltent. \V\ Orémus pro beatíssimo Papa nostro N. \R\ Dóminus consérvet eum, et vivíficet eum, et beátum fáciat eum in terra, et non tradat eum in ánimam inimicórum ejus. \V\ Orémus et pro Antístite nostro N. \R\ Stet et pascat in fortitúdine tua, Dómine, in sublimitáte nóminis tui. \V\ Dómine, salvum fac regem. \R\ Et exáudi nos in die, qua invocavérimus te. \V\ Salvum fac pópulum tuum, Dómine, et bénedic hereditáti tuæ. \R\ Et rege eos, et extólle illos usque in ætérnum. \V\ Meménto Congregatiónis tuæ. \R\ Quam possedísti ab inítio. \V\ Fiat pax in virtúte tua. \R\ Et abundántia in túrribus tuis. \V\ Orémus pro benefactóribus nostris. \R\ Retribúere dignáre, Dómine, ómnibus, nobis bona faciéntibus propter nomen tuum, vitam ætérnam. Amen. \V\ Orémus pro fidélibus defúnctis. \R\ Réquiem ætérnam dona eis, Dómine, et lux perpétua lúceat eis. \V\ Requiéscant in pace. \R\ Amen. \V\ Pro frátribus nostris abséntibus. \R\ Salvos fac servos tuos, Deus meus, sperántes in te. \V\ Pro afflíctis et captívis. \R\ Líbera eos, Deus Israël, ex ómnibus tribulatiónibus suis. \V\ Mitte eis, Dómine, auxílium de sancto. \R\ Et de Sion tuére eos. \V\ Dómine, Deus virtútum, convérte nos. \R\ Et osténde fáciem tuam, et salvi érimus. \V\ Exsúrge, Christe, ádjuva nos. \R\ Et líbera nos propter nomen tuum. \V\ Dómine, exáudi oratiónem meam. \R\ Et clamor meus ad te véniat.

\rubric{Réliqua ut in Ordinário. 11}

% ===== END OF BATCH 20 -- Covers autumnalis.md lines 6745-7156, the
% two alternate ferial Lauds schemes for Feria Quarta, "Ad Laudes
% I" (Psalms 96, 64, 100, Canticum Judith, 145) and "Ad Laudes II"
% (Psalms 50, 64, 100, Canticum Annæ, 145), each with the Capitulum
% "Nox præcéssit" and the hymn "Nox, et ténebræ, et núbila." Notably,
% Psalm 64's incipit carries the dagger and double-dagger
% chant-pointing marks in the Ad Laudes I copy (autumnalis.md line
% 6834: "Te decet hymnus, † Deus, in Sion: * ‡...") but NOT in the
% Ad Laudes II copy of the identical psalm (line 6991, plain "Te
% decet hymnus, Deus, in Sion: *..."); each is transcribed exactly
% as its own occurrence gives it, confirming this is a genuine
% source divergence rather than an OCR gap. Two antiphon incipits in
% Ad Laudes II ("Impietátibus nostris." and "In innocéntia.") carry
% no asterisk at all in the source and are transcribed as bare
% unmarked incipits. Ad Laudes II additionally carries the "Preces
% feriales" litany in full. The document continues with "Ad Primam"
% (autumnalis.md line 7159).
% =====
\hourheading{Ad Primam}

Pater, Ave et Credo.

\V\ Deus, in adjutórium.

\hymnusheading{Hymnus}

\lettrine[lines=2]{J}{AM} lucis orto sídere,\\
Deum precémur súpplices,\\
Ut in diúrnis áctibus\\
Nos servet a nocéntibus.

Linguam refrénans témperet,\\
Ne litis horror ínsonet:\\
Visum fovéndo cóntegat,\\
Ne vanitátes háuriat.

Sint pura cordis íntima,\\
Absístat et vecórdia:\\
Carnis terat supérbiam\\
Potus cibíque párcitas.

Ut, cum dies abscésserit\\
Noctémque sors redúxerit,\\
Mundi per abstinéntiam\\
Ipsi canámus glóriam.

Deo Patri sit glória,\\
Ejúsque soli Fílio,\\
Cum Spíritu Paráclito,\\
Nunc et per omne sæculum.\\
Amen.

\antlabel{Ant.} Misericórdia tua.

\psalmheading{Psalmus 25}

\lettrine[lines=2]{J}{UDICA} me, Dómine, quóniam ego in innocéntia mea ingréssus sum: * et in Dómino sperans non infirmábor.

Proba me, Dómine, et tenta me: * ure renes meos et cor meum.

Quóniam misericórdia tua ante óculos meos est: * et complácui in veritáte tua.

Non sedi cum concílio vanitátis: * et cum iníqua geréntibus non introíbo.

Odívi ecclésiam malignántium: * et cum ímpiis non sedébo.

Lavábo inter innocéntes manus meas: * et circúmdabo altáre tuum, Dómine:

Ut áudiam vocem laudis, * et enárrem univérsa mirabília tua.

Dómine, diléxi decórem domus tuæ, * et locum habitatiónis glóriæ tuæ.

Ne perdas cum ímpiis, Deus, ánimam meam, * et cum viris sánguinum vitam meam:

In quorum mánibus iniquitátes sunt: * déxtera eórum repléta est munéribus.

Ego autem in innocéntia mea ingréssus sum: * rédime me, et miserére mei.

Pes meus stetit in directo: * in ecclésiis benedícam te, Dómine.

\psalmheading{Psalmus 51}

\lettrine[lines=2]{Q}{UID} gloriáris in malítia, * qui potens es in iniquitáte?

Tota die injustítiam cogitávit lingua tua: * sicut novácula acúta fecísti dolum.

Dilexísti malítiam super benignitátem: * iniquitátem magis quam loqui æquitátem.

Dilexísti ómnia verba præcipitatiónis, * lingua dolósa.

Proptérea Deus déstruet te in finem, * evéllet te, et emigrábit te de tabernáculo tuo: et radícem tuam de terra vivéntium.

Vidébunt justi, et timébunt, et super eum ridébunt, et dicent: * Ecce homo, qui non pósuit Deum adjutórem suum:

Sed sperávit in multitúdine divitiárum suárum: * et prævalúit in vanitáte sua.

Ego autem, sicut olíva fructífera in domo Dei, * sperávi in misericórdia Dei in ætérnum: et in séculum séculi.

Confitébor tibi in séculum, quia fecísti: * et exspectábo nomen tuum, quóniam bonum est in conspéctu sanctórum tuórum.

\psalmheading{Psalmus 52}

\lettrine[lines=2]{D}{IXIT} insípiens in corde suo: * Non est Deus.

Corrúpti sunt, et abominábiles facti sunt in iniquitátibus: * non est qui fáciat bonum.

Deus de cælo prospéxit super fílios hóminum: * ut vídeas si est intélligens, aut requírens Deum.

Omnes declinavérunt, simul inútiles facti sunt: * non est qui fáciat bonum, non est usque ad unum.

Nonne scient omnes qui operántur iniquitátem, * qui dévorant plebem meam ut cibum panis?

Deum non invocavérunt: * illic trepidavérunt timóre, ubi non erat timor.

Quóniam Deus dissipávit ossa eórum qui homínibus placent: * confúsi sunt, quóniam Deus sprevit eos.

Quis dabit ex Sion salutáre Israël? * cum convérterit Deus captivitátem plebis suæ, exsultábit Jacob, et lætábitur Israël.

\rubric{¶ Quando ad Laudes dictus fuerit Ps. 50 Miserére, hic subjúngitur sequens Ps. 96, qui in Laudibus locum non habuit.}

\psalmheading{Psalmus 96}

\lettrine[lines=2]{D}{OMINUS} regnávit, exsúltet terra: * læténtur ínsulæ multæ.

Nubes, et calígo in circúitu ejus: * justítia, et judícium corréctio sedis ejus.

Ignis ante ipsum præcédet, * et inflammábit in circúitu inimícos ejus.

Illuxérunt fúlgura ejus orbi terræ: * vidit, et commóta est terra.

Montes, sicut cera fluxérunt a fácie Dómini: * a fácie Dómini omnis terra.

Annuntiavérunt cæli justítiam ejus: * et vidérunt omnes pópuli glóriam ejus.

Confundántur omnes, qui adórant sculptília: * et qui gloriántur in simulácris suis.

Adoráte eum, omnes Angeli ejus: * audívit, et lætáta est Sion.

Et exsultavérunt fíliæ Judæ, * propter judícia tua, Dómine:

Quóniam tu Dóminus Altíssimus super omnem terram: * nimis exaltátus es super omnes deos.

Qui dilígitis Dóminum, odíte malum: * custódit Dóminus ánimas sanctórum suórum, de manu peccatóris liberábit eos.

Lux orta est justo, * et rectis corde lætítia.

Lætámini, justi, in Dómino: * et confitémini memóriæ sanctificatiónis ejus.

\rubric{Secus vero, tribus Psalmis explétis, statim dícitur Ant.}

\antlabel{Ant.} Misericórdia tua, Dómine, ante óculos meos: et complácui in veritáte tua.

\rubric{In Officio Festi vel Octávæ:}

\capitulumlabel{Capitulum}{1 Tim. 1, 17}

\lettrine[lines=2]{R}{EGI} sæculórum immortáli et invisíbili, soli Deo honor et glória in sæcula sæculórum. Amen.

\rubric{In ómnibus autem Fériis et in Vigíliis commúnibus:}

\capitulumlabel{Capitulum}{Zach. 8, 19}

\lettrine[lines=2]{P}{ACEM} et veritátem dilígite, ait Dóminus omnípotens.

\R\ br. Christe, Fili Dei vivi, * Miserére nobis. Christe. \V\ \textit{Qui sedes ad déxteram Patris.} Miserére nobis. Glória Patri. Christe. \V\ Exsúrge, Christe, ádjuva nos. \R\ Et líbera nos propter nomen tuum.

\variantheading{Preces dominicales et feriales (si dicendæ sint)}

KYRIE, eléison. Christe, eléison. Kýrie, eléison.

Pater noster \rubric{secreto usque ad} \V\ Et ne nos indúcas in tentatiónem. \R\ Sed líbera nos a malo.

Credo in Deum \rubric{secreto usque ad} \V\ Carnis resurrectiónem. \R\ Vitam ætérnam. Amen. \V\ Et ego ad te, Dómine, clamávi. \R\ Et mane orátio mea prævéniet te. \V\ Repleátur os meum laude. \R\ Ut cantem glóriam tuam, tota die magnitúdinem tuam. \V\ Dómine, avérte fáciem tuam a peccátis meis. \R\ Et omnes iniquitátes meas dele. \V\ Cor mundum crea in me, Deus. \R\ Et spíritum rectum ínnova in viscéribus meis. \V\ Ne projícias me a fácie tua. \R\ Et spíritum sanctum tuum ne áuferas a me. \V\ Redde mihi lætítiam salutáris tui. \R\ Et spíritu principáli confírma me.

\variantheading{In Précibus feriálibus tantum}

\V\ Eripe me, Dómine, ab hómine malo. \R\ A viro iníquo éripe me. \V\ Eripe me de inimícis meis, Deus meus. \R\ Et ab insurgéntibus in me líbera me. \V\ Eripe me de operántibus iniquitátem. \R\ Et de viris sánguinum salva me. \V\ Sic psalmum dicam nómini tuo in sæculum sæculi. \R\ Ut reddam vota mea de die in diem. \V\ Exáudi nos, Deus, salutáris noster. \R\ Spes ómnium fínium terræ, et in mari longe. \V\ Deus, in adjutórium meum inténde. \R\ Dómine, ad adjuvándum me festína. \V\ Sanctus Deus, Sanctus fortis, Sanctus immortális. \R\ Miserére nobis. \V\ Bénedic, ánima mea, Dómino. \R\ Et ómnia, quæ intra me sunt, nómini sancto ejus. \V\ Bénedic, ánima mea, Dómino. \R\ Et noli oblivísci omnes retributiónes ejus. \V\ Qui propitiátur ómnibus iniquitátibus tuis. \R\ Qui sanat omnes infirmitátes tuas. \V\ Qui rédimit de intéritu vitam tuam. \R\ Qui corónat te in misericórdia et miseratiónibus. \V\ Qui replet in bonis desidérium tuum. \R\ Renovábitur ut áquilæ juvéntus tua.

\variantheading{In utrisque Précibus}

\V\ Adjutórium nostrum in nómine Dómini. \R\ Qui fecit cælum et terram.

\formulalabel{Confíteor. Misereátur. Indulgéntiam.}

\V\ Dignáre, Dómine, die isto. \R\ Sine peccáto nos custodíre. \V\ Miserére nostri, Dómine. \R\ Miserére nostri. \V\ Fiat misericórdia tua, Dómine, super nos. \R\ Quemádmodum sperávimus in te. \V\ Dómine, exáudi oratiónem meam. \R\ Et clamor meus ad te véniat.

\V\ Dóminus vobíscum.

Orémus. \oratiolabel{Oratio}

\lettrine[lines=2]{D}{OMINE}, Deus omnípotens, qui ad princípium hujus diéi nos perveníre fecísti: tua nos hódie salva virtúte; ut in hac die ad nullum declinémus peccátum, sed semper ad tuam justítiam faciéndam nostra procédant elóquia, dirigántur cogitatiónes et ópera. Per Dóminum. \R\ Amen.

\V\ Dóminus vobíscum. \V\ Benedicámus Dómino.

\rubric{In Choro Martyrológium.}

\V\ Pretiósa in conspéctu Dómini. \R\ Mors Sanctórum ejus.

\rubric{Et absolúte, sine} Orémus:

\oratiolabel{Oratio}

\lettrine[lines=2]{S}{ANCTA} María, et omnes Sancti intercédant pro nobis ad Dóminum, ut nos mereámur ab eo adjuvári et salvári, qui vivit et regnat in sæcula sæculórum. \R\ Amen.

\V\ Deus, in adjutórium meum inténde. \R\ Dómine, ad adjuvándum me festína. \rubric{Et dícitur ter; último additur:} Glória Patri. Sicut erat.

Kýrie, eléison. Christe, eléison. Kýrie, eléison.

Pater noster \rubric{secreto usque ad} \V\ Et ne nos indúcas in tentatiónem. \R\ Sed líbera nos a malo. \V\ Réspice in servos tuos, Dómine, et in ópera tua, et dírige fílios eórum. \R\ Et sit splendor Dómini, Dei nostri, super nos, et ópera mánuum nostrárum dírige super nos, et opus mánuum nostrárum dírige. \V\ Glória Patri. \R\ Sicut erat.

Orémus. \oratiolabel{Oratio}

\lettrine[lines=2]{D}{IRIGERE} et sanctificáre, régere et gubernáre dignáre, Dómine Deus, Rex cæli et terræ, hódie corda et córpora nostra, sensus, sermónes et actus nostros in lege tua, et in opéribus mandatórum tuórum: ut hic, et in ætérnum, te auxiliánte, salvi et líberi esse mereámur, Salvátor mundi: Qui vivis et regnas in sæcula sæculórum. \R\ Amen.

\V\ Jube, domne.

\benedlabel{Bened.} Dies et actus nostros in sua pace dispónat Dóminus omnípotens. \R\ Amen.

\variantheading{In Officio feriali}

\lettrine[lines=2]{D}{OMINUS} autem dírigat corda et córpora nostra in caritáte Dei, et patiéntia Christi. Tu autem.

\variantheading{E Proprio Sanctorum}

\variantheading{Pro Oct. Omn. Sanctorum}

\lettrine[lines=2]{B}{ENEDÍCTIO} et cláritas et sapiéntia et gratiárum áctio, honor, virtus et fortitúdo Deo nostro in sǽcula sæculórum. Amen.

\variantheading{E Communi Sanctorum}

\variantheading{Pro Apostolis}

\lettrine[lines=2]{I}{BANT} Apóstoli gaudéntes a conspéctu concílii, quóniam digni hábiti sunt pro nómine Jesu contuméliam pati.

\variantheading{Pro Martyre}

\lettrine[lines=2]{J}{USTUS} cor suum trádidit ad vigilándum dilúculo ad Dóminum, qui fecit illum, et in conspéctu Altíssimi deprecábitur.

\variantheading{Pro pluribus Martyribus}

\lettrine[lines=2]{F}{ULGEBUNT} justi, et tamquam scintíllæ in arundinéto discúrrent. Judicábunt natiónes, et dominabúntur pópulis: et regnábit Dóminus illórum in perpétuum.

\variantheading{Pro Conf. Pont.}

\lettrine[lines=2]{F}{UNGI} sacerdótio, et habére laudem in nómine ipsíus, et offérre illi incénsum dignum in odórem suavitátis.

\variantheading{Pro Conf. non Pont.}

\lettrine[lines=2]{J}{USTUM} dedúxit Dóminus per vias rectas, et osténdit illi regnum Dei, et dedit illi sciéntiam sanctórum; honestávit illum in labóribus, et complévit labóres illíus.

\variantheading{Pro Virgine Martyre}

\lettrine[lines=2]{D}{OMINE}, Deus meus, exaltásti super terram habitatiónem meam, et pro morte defluénte deprecáta sum. Invocávi Dóminum, Patrem Dómini mei, ut non derelínquat me in die tribulatiónis meæ, et in témpore superbórum sine adjutório.

\variantheading{Pro Virgine non Martyre}

\lettrine[lines=2]{O}{} \textsc{quam} pulchra est casta generátio cum claritáte! immortális est enim memória illíus: quóniam apud Deum nota est, et apud hómines.

\variantheading{Pro Martyre tantum}

\lettrine[lines=2]{L}{AUDABIT} usque ad mortem ánima mea Dóminum, quóniam éruis sustinéntes te, et líberas eos de manu angústiæ, Dómine, Deus noster.

\variantheading{Pro nec Virg. nec Mart.}

\lettrine[lines=2]{M}{ULTAE} fíliæ congregavérunt divítias: tu supergréssa es univérsas. Fallax grátia, et vana est pulchritúdo: múlier timens Dóminum, ipsa laudábitur.

\variantheading{Pro Dedicatione Ecclesiæ}

\lettrine[lines=2]{E}{T} abstérget Deus omnem lácrimam ab óculis eórum: et mors ultra non erit, neque luctus, neque clamor, neque dolor erit ultra, quia prima abiérunt. Et dixit qui sedébat in throno: Ecce, nova fácio ómnia.

\variantheading{In Officiis B. Mariæ V.}

\lettrine[lines=2]{I}{N} platéis sicut cinnamómum et bálsamum aromatízans odórem dedi: quasi myrrha elécta dedi suavitátem odóris.

\V\ Adjutórium nostrum in nómine Dómini. \R\ Qui fecit cælum et terram.

\V\ Benedícite. \R\ Deus.

\benedictiolabel{Benedictio} Dóminus nos benedícat, et ab omni malo deféndat, et ad vitam perdúcat ætérnam. Et fidélium ánimæ per misericórdiam Dei requiéscant in pace. \R\ Amen.

\rubric{Deinde dícitur tantum} Pater noster \rubric{secréto, nisi sequátur ália Hora.}

% ===== END OF BATCH 21 -- Covers autumnalis.md lines 7159-7424,
% completing "Ad Primam" for Feria Quarta in full: the hymn "Jam
% lucis orto sídere," Psalms 25, 51, 52, the conditional Psalm 96
% insertion, both Capitula, the brief Responsory, the full Preces
% apparatus, the closing Orationes, and the Lectio-brevis apparatus
% (E Proprio Sanctorum, E Communi Sanctorum). Four consecutive psalm
% incipits in autumnalis.md (Ps. 25 "JUdica," Ps. 51 "QUid," Ps. 52
% "DIxit," and the conditional Ps. 96 "DOminus") print only their
% first letter or two in full caps with the remainder of the first
% word in lowercase -- a mechanical OCR capitalization defect,
% confirmed for "DIxit" and "DOminus" by the many fully-capitalized
% correct occurrences of these exact psalm incipits already
% transcribed elsewhere in this volume (e.g. Psalm 13's "DIXIT
% insípiens" in Batch 9, and "DOMINUS regnávit" repeatedly in
% Batches 19-20) -- and silently normalized to full caps throughout,
% consistent with the uniform ALL-CAPS-opening-word convention
% applied to every psalm in this transcription. The document
% continues with "Ad Tertiam" (autumnalis.md line 7427).
% =====
\hourheading{Ad Tertiam}

Pater noster et Ave.

\V\ Deus, in adjutórium.

\hymnusheading{Hymnus}

\lettrine[lines=2]{N}{UNC}, Sancte, nobis, Spíritus,\\
Unum Patri cum Fílio,\\
Dignáre promptus íngeri\\
Nostro refúsus péctori.

Os, lingua, mens, sensus, vigor\\
Confessiónem pérsonent,\\
Flamméscat igne cáritas,\\
Accéndat ardor próximos.

Præsta, Pater piíssime,\\
Patríque compar Unice,\\
Cum Spíritu Paráclito\\
Regnans per omne sæculum.\\
Amen.

\antlabel{Ant.} Deus ádjuvat me.

\psalmheading{Psalmus 53}

\lettrine[lines=2]{D}{EUS}, in nómine tuo salvum me fac: * et in virtúte tua júdica me.

Deus, exáudi oratiónem meam: * áuribus pércipe verba oris mei.

Quóniam aliéni insurrexérunt advérsum me, et fortes quæsiérunt ánimam meam: * et non proposuérunt Deum ante conspéctum suum.

Ecce enim, Deus ádjuvat me: * et Dóminus suscéptor est ánimæ meæ.

Avérte mala inimícis meis: * et in veritáte tua dispérde illos.

Voluntárie sacrificábo tibi, * et confitébor nómini tuo, Dómine: quóniam bonum est:

Quóniam ex omni tribulatióne eripuísti me: * et super inimícos meos despéxit óculus meus.

\psalmheading{Psalmus 54, i}

\lettrine[lines=2]{E}{XAUDI}, Deus, oratiónem meam, et ne despéxeris deprecatiónem meam: * inténde mihi, et exáudi me.

Contristátus sum in exercitatióne mea: * et conturbátus sum a voce inimíci, et a tribulatióne peccatóris.

Quóniam declinavérunt in me iniquitátes: * et in ira molésti erant mihi.

Cor meum conturbátum est in me: * et formído mortis cécidit super me.

Timor et tremor venérunt super me: * et contexérunt me ténebræ.

Et dixi: Quis dabit mihi pennas sicut colúmbæ, * et volábo, et requiéscam?

Ecce, elongávi fúgiens: * et mansi in solitúdine.

Exspectábam eum, qui salvum me fecit * a pusillanimitáte spíritus et tempestáte.

Præcípita, Dómine, dívide linguas eórum: * quóniam vidi iniquitátem, et contradictiónem in civitáte.

Die ac nocte circúmdabit eam super muros ejus iníquitas: * et labor in médio ejus, et injustítia.

Et non defécit de platéis ejus * usúra, et dolus.

Quóniam si inimícus meus maledixísset mihi, * sustinuíssem útique.

Et si is, qui óderat me, super me magna locútus fuísset, * abscondíssem me fórsitan ab eo.

Tu vero, homo unánimis: * dux meus, et notus meus:

Qui simul mecum dulces capiébas cibos: * in domo Dei ambulávimus cum consénsu.

Véniat mors super illos: * et descéndant in inférnum vivéntes:

Quóniam nequítiæ in habitáculis eórum: * in médio eórum.

\psalmheading{Psalmus 54, ii}

\lettrine[lines=2]{E}{GO} autem ad Deum clamávi: * et Dóminus salvábit me.

Véspere, et mane, et merídie narrábo et annuntiábo: * et exáudiet vocem meam.

Rédimet in pace ánimam meam ab his, qui appropínquant mihi: * quóniam inter multos erant mecum.

Exáudiet Deus, et humiliábit illos, * qui est ante sæcula.

Non enim est illis commutátio, et non timuérunt Deum: * exténdit manum suam in retribuéndo.

Contaminavérunt testaméntum ejus, divísi sunt ab ira vultus ejus: * et appropinquávit cor illíus.

Mollíti sunt sermónes ejus super óleum: * et ipsi sunt jácula.

Jacta super Dóminum curam tuam, et ipse te enútriet: * non dabit in ætérnum fluctuatiónem justo.

Tu vero, Deus, dedúces eos, * in púteum intéritus.

Viri sánguinum, et dolósi non dimidiábunt dies suos: * ego autem sperábo in te, Dómine.

\antlabel{Ant.} Deus ádjuvat me: et Dóminus suscéptor est ánimæ meæ.

\variantheading{In Officio feriali}

\capitulumlabel{Capitulum}{Jer. 17, 14}

\lettrine[lines=2]{S}{ANA} me, Dómine, et sanábor: salvum me fac, et salvus ero: quóniam laus mea tu es.

\R\ br. Sana ánimam meam, * Quia peccávi tibi. Sana. \V\ Ego dixi: Dómine, miserére mei. Quia. Glória Patri. Sana. \V\ Adjútor meus esto, ne derelínquas me. \R\ Neque despícias me, Deus, salutáris meus.

\variantheading{In Vigiliis communibus:}

\variantheading{Preces feriales}

KYRIE, eléison. Christe, eléison. Kýrie, eléison.

Pater noster \rubric{secreto usque ad} \V\ Et ne nos indúcas in tentatiónem. \R\ Sed líbera nos a malo. \V\ Dómine, Deus virtútum, convérte nos. \R\ Et osténde fáciem tuam, et salvi érimus. \V\ Exsúrge, Christe, ádjuva nos. \R\ Et líbera nos propter nomen tuum. \V\ Dómine, exáudi oratiónem meam. \R\ Et clamor meus ad te véniat.

\rubric{Réliqua ut in Ordinário. 19}

\hourheading{Ad Sextam}

Pater noster et Ave.

\V\ Deus, in adjutórium.

\hymnusheading{Hymnus}

\lettrine[lines=2]{R}{ECTOR} potens, verax Deus,\\
Qui témperas rerum vices,\\
Splendóre mane illúminas,\\
Et ígnibus merídiem:

Exstíngue flammas lítium,\\
Aufer calórem nóxium,\\
Confer salútem córporum,\\
Verámque pacem córdium.

Præsta, Pater piíssime,\\
Patríque compar Unice,\\
Cum Spíritu Paráclito\\
Regnans per omne sæculum.\\
Amen.

\antlabel{Ant.} In Deo sperávi.

\psalmheading{Psalmus 55}

\lettrine[lines=2]{M}{ISERERE} mei, Deus, quóniam conculcávit me homo: * tota die impúgnans tribulávit me.

Conculcavérunt me inimíci mei tota die: * quóniam multi bellántes advérsum me.

Ab altitúdine diéi timébo: * ego vero in te sperábo.

In Deo laudábo sermónes meos, in Deo sperávi: * non timébo quid fáciat mihi caro.

Tota die verba mea exsecrabántur: * advérsum me omnes cogitatiónes eórum in malum.

Inhabitábunt et abscóndent: * ipsi calcáneum meum observábunt.

Sicut sustinuérunt ánimam meam, pro níhilo salvos fácies illos: * in ira pópulos confrínges.

Deus, vitam meam annuntiávi tibi: * posuísti lácrimas meas in conspéctu tuo.

Sicut et in promissióne tua: * tunc converténtur inimíci mei retrórsum:

In quacúmque die invocávero te: * ecce, cognóvi, quóniam Deus meus es.

In Deo laudábo verbum, in Dómino laudábo sermónem: * in Deo sperávi, non timébo quid fáciat mihi homo.

In me sunt, Deus, vota tua, * quæ reddam, laudatiónes tibi.

Quóniam eripuísti ánimam meam de morte, et pedes meos de lapsu: * ut pláceam coram Deo in lúmine vivéntium.

\psalmheading{Psalmus 56}

\lettrine[lines=2]{M}{ISERERE} mei, Deus, miserére mei: * quóniam in te confídit ánima mea.

Et in umbra alárum tuárum sperábo, * donec tránseat iníquitas.

Clamábo ad Deum altíssimum: * Deum qui benefécit mihi.

Misit de cælo, et liberávit me: * dedit in oppróbrium conculcántes me.

Misit Deus misericórdiam suam, et veritátem suam, * et erípuit ánimam meam de médio catulórum leónum: dormívi conturbátus.

Fílii hóminum dentes eórum arma et sagíttæ: * et lingua eórum gládius acútus.

Exaltáre super cælos, Deus, * et in omnem terram glória tua.

Láqueum paravérunt pédibus meis: * et incurvavérunt ánimam meam.

Fodérunt ante fáciem meam fóveam: * et incidérunt in eam.

Parátum cor meum, Deus, parátum cor meum: * cantábo, et psalmum dicam.

Exsúrge, glória mea, exsúrge, psaltérium et cíthara: * exsúrgam dilúculo.

Confitébor tibi in pópulis, Dómine: * et psalmum dicam tibi in géntibus:

Quóniam magnificáta est usque ad cælos misericórdia tua, * et usque ad nubes véritas tua.

Exaltáre super cælos, Deus: * et super omnem terram glória tua.

\psalmheading{Psalmus 57}

\lettrine[lines=2]{S}{I} vere útique justítiam loquímini: * recta judicáte, fílii hóminum.

Etenim in corde iniquitátes operámini: * in terra injustítias manus vestræ concínnant.

Alienáti sunt peccatóres a vulva, erravérunt ab útero: * locúti sunt falsa.

Furor illis secúndum similitúdinem serpéntis: * sicut áspidis surdæ, et obturántis aures suas.

Quæ non exáudiet vocem incantántium: * et venéfici incantántis sapiénter.

Deus cónteret dentes eórum in ore ipsórum: * molas leónum confrínget Dóminus.

Ad níhilum devénient tamquam aqua decúrrens: * inténdit arcum suum donec infirméntur.

Sicut cera, quæ fluit, auferéntur: * supercécidit ignis, et non vidérunt solem.

Priúsquam intellígerent spinæ vestræ rhamnum: * sicut vivéntes, sic in ira absórbet eos.

Lætábitur justus cum víderit vindíctam: * manus suas lavábit in sánguine peccatóris.

Et dicet homo: Si útique est fructus justo: * útique est Deus júdicans eos in terra.

\antlabel{Ant.} In Deo sperávi, non timébo quid fáciat mihi homo.

\variantheading{In Officio feriali}

\capitulumlabel{Capitulum}{Rom. 13, 8}

\lettrine[lines=2]{N}{EMINI} quidquam debeátis, nisi ut ínvicem diligátis: qui enim díligit próximum, legem implévit.

\R\ br. Benedícam Dóminum * In omni témpore. Benedícam. \V\ Semper laus ejus in ore meo. In. Glória Patri. Benedícam. \V\ Dóminus regit me, et nihil mihi déerit. \R\ In loco páscuæ ibi me collocávit.

\variantheading{In Vigiliis communibus:}

\variantheading{Preces feriales}

KYRIE, eléison. Christe, eléison. Kýrie, eléison.

Pater noster \rubric{secreto usque ad} \V\ Et ne nos indúcas in tentatiónem. \R\ Sed líbera nos a malo. \V\ Dómine, Deus virtútum, convérte nos. \R\ Et osténde fáciem tuam, et salvi érimus. \V\ Exsúrge, Christe, ádjuva nos. \R\ Et líbera nos propter nomen tuum. \V\ Dómine, exáudi oratiónem meam. \R\ Et clamor meus ad te véniat.

\rubric{Réliqua ut in Ordinário. 20}

\hourheading{Ad Nonam}

Pater noster et Ave.

\V\ Deus, in adjutórium.

\hymnusheading{Hymnus}

\lettrine[lines=2]{R}{ERUM}, Deus, tenax vigor,\\
Immótus in te pérmanens,\\
Lucis diúrnæ témpora\\
Succéssibus detérminans:

Largíre lumen véspere,\\
Quo vita nusquam décidat,\\
Sed prǽmium mortis sacræ\\
Perénnis instet glória.

Præsta, Pater piíssime,\\
Patríque compar Unice,\\
Cum Spíritu Paráclito\\
Regnans per omne sæculum.\\
Amen.

\antlabel{Ant.} Deus meus.

\psalmheading{Psalmus 58, i}

\lettrine[lines=2]{E}{RIPE} me de inimícis meis, Deus meus: * et ab insurgéntibus in me líbera me.

Eripe me de operántibus iniquitátem: * et de viris sánguinum salva me.

Quia ecce cepérunt ánimam meam: * irruérunt in me fortes.

Neque iníquitas mea, neque peccátum meum, Dómine: * sine iniquitáte cucúrri, et diréxi.

Exsúrge in occúrsum meum, et vide: * et tu, Dómine, Deus virtútum, Deus Israël,

Inténde ad visitándas omnes gentes: * non misereáris ómnibus, qui operántur iniquitátem.

Converténtur ad vésperam: et famem patiéntur ut canes, * et circuíbunt civitátem.

Ecce, loquéntur in ore suo, et gládius in lábiis eórum: * quóniam quis audívit?

Et tu, Dómine, deridébis eos: * ad níhilum dedúces omnes gentes.

Fortitúdinem meam ad te custódiam, quia, Deus, suscéptor meus es: * Deus meus, misericórdia ejus prævéniet me.

\psalmheading{Psalmus 58, ii}

\lettrine[lines=2]{D}{EUS} osténdet mihi super inimícos meos, ne occídas eos: * nequándo obliviscántur pópuli mei.

Dispérge illos in virtúte tua: * et depóne eos, protéctor meus, Dómine:

Delíctum oris eórum, sermónem labiórum ipsórum: * et comprehendántur in supérbia sua.

Et de exsecratióne et mendácio annuntiabúntur in consummatióne: * in ira consummatiónis, et non erunt.

Et scient quia Deus dominábitur Jacob: * et fínium terræ.

Converténtur ad vésperam: et famem patiéntur ut canes, * et circuíbunt civitátem.

Ipsi dispergéntur ad manducándum: * si vero non fúerint saturáti, et murmurábunt.

Ego autem cantábo fortitúdinem tuam: * et exsultábo mane misericórdiam tuam.

Quia factus es suscéptor meus, * et refúgium meum, in die tribulatiónis meæ.

Adjútor meus, tibi psallam, quia, Deus, suscéptor meus es: * Deus meus, misericórdia mea.

\psalmheading{Psalmus 59}

\lettrine[lines=2]{D}{EUS}, repulísti nos, et destruxísti nos: * irátus es, et misértus es nobis.

Commovísti terram, et conturbásti eam: * sana contritiónes ejus, quia commóta est.

Ostendísti pópulo tuo dura: * potásti nos vino compunctiónis.

Dedísti metuéntibus te significatiónem: * ut fúgiant a fácie arcus:

Ut liberéntur dilécti tui: * salvum fac déxtera tua, et exáudi me.

Deus locútus est in sancto suo: * Lætábor, et partíbor Síchimam: et convállem tabernaculórum metíbor.

Meus est Gálaad, et meus est Manásses: * et Ephraïm fortitúdo cápitis mei.

Juda rex meus: * Moab olla spei meæ.

In Idumǽam exténdam calceaméntum meum: * mihi aliénigenæ súbditi sunt.

Quis dedúcet me in civitátem munítam? * quis dedúcet me usque in Idumǽam?

Nonne tu, Deus, qui repulísti nos, * et non egrediéris, Deus, in virtútibus nostris?

Da nobis auxílium de tribulatióne: * quia vana salus hóminis.

In Deo faciémus virtútem: * et ipse ad níhilum dedúcet tribulántes nos.

\antlabel{Ant.} Deus meus, misericórdia tua prævéniet me.

\variantheading{In Officio feriali}

\capitulumlabel{Capit.}{1 Petri 1, 17-19}

\lettrine[lines=2]{I}{N} timóre incolátus vestri témpore conversámini: sciéntes quod non corruptibílibus auro vel argénto redémpti estis, sed pretióso sánguine quasi Agni immaculáti Christi.

\R\ br. Rédime me, Dómine, * Et miserére mei. Rédime. \V\ Pes enim meus stetit in via recta. Et miserére mei. Glória Patri. Rédime. \V\ Ab occúltis meis munda me, Dómine. \R\ Et ab aliénis parce servo tuo.

\variantheading{In Vigiliis communibus:}

\variantheading{Preces feriales}

KYRIE, eléison. Christe, eléison. Kýrie, eléison.

Pater noster \rubric{secreto usque ad} \V\ Et ne nos indúcas in tentatiónem. \R\ Sed líbera nos a malo. \V\ Dómine, Deus virtútum, convérte nos. \R\ Et osténde fáciem tuam, et salvi érimus. \V\ Exsúrge, Christe, ádjuva nos. \R\ Et líbera nos propter nomen tuum. \V\ Dómine, exáudi oratiónem meam. \R\ Et clamor meus ad te véniat.

\rubric{Réliqua ut in Ordinário. 22}

% ===== END OF BATCH 22 -- Covers autumnalis.md lines 7427-7796,
% completing the Little Hours of Feria Quarta in full: "Ad
% Tertiam" (Psalms 53, 54/i, 54/ii), "Ad Sextam" (Psalms 55, 56,
% 57), and "Ad Nonam" (Psalms 58/i, 58/ii, 59), each with its own
% hymn, Capitulum, brief Responsory, and "In Vigiliis communibus:
% Preces feriales" apparatus. Three further occurrences of the
% recurring stray OCR "\*" artifact in the hour-hymns are silently
% dropped, consistent with every prior occurrence. Several more
% psalm incipits in this batch carry the same partial-caps OCR
% defect already documented in Batch 21 ("DEus," "EXaudi," "EGo,"
% "SAna" [Capitulum], "MIserére" twice, "ERipe," "DEus" twice,
% "NEmini" [Capitulum]) and are silently normalized to full caps,
% consistent with the uniform convention and confirmed by the many
% correctly-capitalized occurrences of these same incipits elsewhere
% in this volume. The document continues with "Ad Vesperas"
% (autumnalis.md line 7799).
% =====
\hourheading{Ad Vesperas}

Pater noster et Ave.

\V\ Deus, in adjutórium.

\antlabel{Ant.} Beáti omnes, * \textdagger\ qui timent Dóminum. \textdaggerdbl

\psalmheading{Psalmus 127}

\lettrine[lines=2]{B}{EATI} omnes, \textdagger\ qui timent Dóminum, * \textdaggerdbl\ qui ámbulant in viis ejus.

Labóres mánuum tuárum quia manducábis: * beátus es, et bene tibi erit.

Uxor tua sicut vitis abúndans, * in latéribus domus tuæ.

Fílii tui sicut novéllæ olivárum, * in circúitu mensæ tuæ.

Ecce, sic benedicétur homo, * qui timet Dóminum.

Benedícat tibi Dóminus ex Sion: * et vídeas bona Jerúsalem ómnibus diébus vitæ tuæ.

Et vídeas fílios filiórum tuórum, * pacem super Israël.

\antlabel{Ant.} Beáti omnes, qui timent Dóminum.

\antlabel{Ant.} Confundántur omnes, * qui odérunt Sion.

\psalmheading{Psalmus 128}

\lettrine[lines=2]{S}{ÆPE} expugnavérunt me a juventúte mea, * dicat nunc Israël:

Sæpe expugnavérunt me a juventúte mea: * étenim non potuérunt mihi.

Supra dorsum meum fabricavérunt peccatóres: * prolongavérunt iniquitátem suam.

Dóminus justus concídit cervíces peccatórum: * confundántur et convertántur retrórsum omnes, qui odérunt Sion.

Fiant sicut fœnum tectórum: * quod priúsquam evellátur, exáruit:

De quo non implévit manum suam qui metit, * et sinum suum qui manípulos cólligit.

Et non dixérunt qui præteríbant: Benedíctio Dómini super vos: * benedíximus vobis in nómine Dómini.

\antlabel{Ant.} Confundántur omnes, qui odérunt Sion.

\antlabel{Ant.} De profúndis * \textdagger\ clamávi ad te, Dómine. \textdaggerdbl

\psalmheading{Psalmus 129}

\lettrine[lines=2]{D}{E} profúndis \textdagger\ clamávi ad te, Dómine: * \textdaggerdbl\ Dómine, exáudi vocem meam:

Fiant aures tuæ inténdentes, * in vocem deprecatiónis meæ.

Si iniquitátes observáveris, Dómine: * Dómine, quis sustinébit?

Quia apud te propitiátio est: * et propter legem tuam sustínui te, Dómine.

Sustínuit ánima mea in verbo ejus: * sperávit ánima mea in Dómino.

A custódia matutína usque ad noctem: * speret Israël in Dómino.

Quia apud Dóminum misericórdia: * et copiósa apud eum redémptio.

Et ipse rédimet Israël, * ex ómnibus iniquitátibus ejus.

\antlabel{Ant.} De profúndis clamávi ad te, Dómine.

\antlabel{Ant.} Dómine, * \textdagger\ non est exaltátum cor meum. \textdaggerdbl

\psalmheading{Psalmus 130}

\lettrine[lines=2]{D}{OMINE}, \textdagger\ non est exaltátum cor meum: * \textdaggerdbl\ neque eláti sunt óculi mei.

Neque ambulávi in magnis: * neque in mirabílibus super me.

Si non humíliter sentiébam: * sed exaltávi ánimam meam:

Sicut ablactátus est super matre sua, * ita retribútio in ánima mea.

Speret Israël in Dómino, * ex hoc nunc et usque in sæculum.

\antlabel{Ant.} Dómine, non est exaltátum cor meum.

\antlabel{Ant.} Elégit Dóminus * Sion in habitatiónem sibi.

\psalmheading{Psalmus 131}

\lettrine[lines=2]{M}{EMENTO}, Dómine, David, * et omnis mansuetúdinis ejus:

Sicut jurávit Dómino, * votum vovit Deo Jacob:

Si introíero in tabernáculum domus meæ, * si ascéndero in lectum strati mei:

Si dédero somnum óculis meis, * et pálpebris meis dormitatiónem:

Et réquiem tempóribus meis: donec invéniam locum Dómino, * tabernáculum Deo Jacob.

Ecce, audívimus eam in Ephráta: * invénimus eam in campis silvæ.

Introíbimus in tabernáculum ejus: * adorábimus in loco, ubi stetérunt pedes ejus.

Surge, Dómine, in réquiem tuam, * tu et arca sanctificatiónis tuæ.

Sacerdótes tui induántur justítiam: * et sancti tui exsúltent.

Propter David, servum tuum, * non avértas fáciem Christi tui.

Jurávit Dóminus David veritátem, et non frustrábitur eam: * De fructu ventris tui ponam super sedem tuam.

Si custodíerint fílii tui testaméntum meum, * et testimónia mea hæc, quæ docébo eos:

Et fílii eórum usque in sæculum, * sedébunt super sedem tuam.

Quóniam elégit Dóminus Sion: * elégit eam in habitatiónem sibi.

Hæc réquies mea in sæculum sæculi: * hic habitábo quóniam elégi eam.

Víduam ejus benedícens benedícam: * páuperes ejus saturábo pánibus.

Sacerdótes ejus índuam salutári: * et sancti ejus exsultatióne exsultábunt.

Illuc prodúcam cornu David, * parávi lucérnam Christo meo.

Inimícos ejus índuam confusióne: * super ipsum autem efflorébit sanctificátio mea.

\antlabel{Ant.} Elégit Dóminus Sion in habitatiónem sibi.

\variantheading{In Officio feriali}

\capitulumlabel{Capitulum}{2 Cor. 1, 3-4}

\lettrine[lines=2]{B}{ENEDICTUS} Deus, et Pater Dómini nostri Jesu Christi, Pater misericordiárum, et Deus totíus consolatiónis, qui consolátur nos in omni tribulatióne nostra.

\hymnusheading{Hymnus}

\lettrine[lines=2]{C}{ÆLI} Deus sanctíssime,\\
Qui lúcidas mundi plagas\\
Candóre pingis ígneo,\\
Augens decóro lúmine:

Quarto die qui flámmeam\\
Dum solis accéndis rotam,\\
Lunæ minístras órdinem,\\
Vagósque cursus síderum:

Ut nóctibus, vel lúmini\\
Diremptiónis términum,\\
Primórdiis et ménsium\\
Signum dares notíssimum;

Expélle noctem córdium:\\
Abstérge sordes méntium:\\
Resólve culpæ vínculum:\\
Evérte moles críminum.

Præsta, Pater piíssime,\\
Patríque compar Unice,\\
Cum Spíritu Paráclito\\
Regnans per omne sæculum.\\
Amen.

\V\ Dirigátur, Dómine, orátio mea. \R\ Sicut incénsum in conspéctu tuo.

\antlabel{Ad Magnif. Ant.} Respéxit Dóminus * humilitátem meam, et fecit in me magna, qui potens est.

\rubric{Réliqua ut in Ordinário. 22}

\hourheading{Ad Completorium}

\V\ Jube, domne.

\benedlabel{Bened.} Noctem quiétam, et finem perféctum concédat nobis Dóminus omnípotens. \R\ Amen.

\lectioheading{Lectio br.}{1 Petri 5, 8-9}

\lettrine[lines=2]{F}{RATRES}: Sóbrii estóte, et vigiláte: quia adversárius vester diábolus tamquam leo rúgiens círcuit, quærens quem dévoret: cui resístite fortes in fide. Tu autem, Dómine, miserére nobis. \R\ Deo grátias.

\V\ Adjutórium nostrum in nómine Dómini. \R\ Qui fecit cælum et terram.

Pater noster \rubric{quod dícitur totum secréto.}

\formulalabel{Confíteor.} Misereátur. Indulgéntiam.

\V\ Convérte nos, Deus, salutáris noster. \R\ Et avérte iram tuam a nobis.

\V\ Deus, in adjutórium.

\antlabel{Ant.} Immíttet Angelus Dómini.

\psalmheading{Psalmus 33, i}

\lettrine[lines=2]{B}{ENEDICAM} Dóminum in omni témpore: * semper laus ejus in ore meo.

In Dómino laudábitur ánima mea: * áudiant mansuéti, et læténtur.

Magnificáte Dóminum mecum: * et exaltémus nomen ejus in idípsum.

Exquisívi Dóminum, et exaudívit me: * et ex ómnibus tribulatiónibus meis erípuit me.

Accédite ad eum, et illuminámini: * et fácies vestræ non confundéntur.

Ipse pauper clamávit, et Dóminus exaudívit eum: * et de ómnibus tribulatiónibus ejus salvávit eum.

Immíttet Angelus Dómini in circúitu timéntium eum: * et erípiet eos.

Gustáte, et vidéte quóniam suávis est Dóminus: * beátus vir, qui sperat in eo.

Timéte Dóminum, omnes sancti ejus: * quóniam non est inópia timéntibus eum.

Dívites eguérunt et esuriérunt: * inquiréntes autem Dóminum non minuéntur omni bono.

\psalmheading{Psalmus 33, ii}

\lettrine[lines=2]{V}{ENITE}, fílii, audíte me: * timórem Dómini docébo vos.

Quis est homo qui vult vitam: * díligit dies vidére bonos?

Próhibe linguam tuam a malo: * et lábia tua ne loquántur dolum.

Divérte a malo, et fac bonum: * inquíre pacem, et perséquere eam.

Oculi Dómini super justos: * et aures ejus in preces eórum.

Vultus autem Dómini super faciéntes mala: * ut perdat de terra memóriam eórum.

Clamavérunt justi, et Dóminus exaudívit eos: * et ex ómnibus tribulatiónibus eórum liberávit eos.

Juxta est Dóminus iis, qui tribuláto sunt corde: * et húmiles spíritu salvábit.

Multæ tribulatiónes justórum: * et de ómnibus his liberábit eos Dóminus.

Custódit Dóminus ómnia ossa eórum: * unum ex his non conterétur.

Mors peccatórum péssima: * et qui odérunt justum, delínquent.

Rédimet Dóminus ánimas servórum suórum: * et non delínquent omnes qui sperant in eo.

\psalmheading{Psalmus 60}

\lettrine[lines=2]{E}{XAUDI}, Deus, deprecatiónem meam: * inténde oratióni meæ.

A fínibus terræ ad te clamávi: * dum anxiarétur cor meum, in petra exaltásti me.

Deduxísti me, quia factus es spes mea: * turris fortitúdinis a fácie inimíci.

Inhabitábo in tabernáculo tuo in sæcula: * prótegar in velaménto alárum tuárum.

Quóniam tu, Deus meus, exaudísti oratiónem meam: * dedísti hereditátem timéntibus nomen tuum.

Dies super dies regis adjícies: * annos ejus usque in diem generatiónis et generatiónis.

Pérmanet in ætérnum in conspéctu Dei: * misericórdiam et veritátem ejus quis requíret?

Sic psalmum dicam nómini tuo in sæculum sæculi: * ut reddam vota mea de die in diem.

\antlabel{Ant.} Immíttet Angelus Dómini in circúitu timéntium eum: et erípiet eos.

\hymnusheading{Hymnus}

\lettrine[lines=2]{T}{E} lucis ante términum,\\
Rerum Creátor, póscimus,\\
Ut pro tua cleméntia\\
Sis præsul et custódia.

Procul recédant sómnia,\\
Et nóctium phantásmata,\\
Hostémque nostrum cómprime,\\
Ne polluántur córpora.

Præsta, Pater piíssime,\\
Patríque compar Unice,\\
Cum Spíritu Paráclito\\
Regnans per omne sæculum.\\
Amen.

\capitulumlabel{Capitulum}{Jer. 14, 9}

\lettrine[lines=2]{T}{U} autem in nobis es, Dómine, et nomen sanctum tuum invocátum est super nos: ne derelínquas nos, Dómine, Deus noster.

\R\ br. In manus tuas, Dómine, * Comméndo spíritum meum. In manus. \V\ Redemísti nos, Dómine, Deus veritátis. Comméndo. Glória Patri. In manus.

\V\ Custódi nos, Dómine, ut pupíllam óculi. \R\ Sub umbra alárum tuárum prótege nos.

\antlabel{Ant.} Salva nos.

\canticumheading{Canticum Simeonis}{Luc. 2, 29-32}

\lettrine[lines=2]{N}{UNC} dimíttis servum tuum, Dómine, * secúndum verbum tuum in pace:

Quia vidérunt óculi mei * salutáre tuum,

Quod parásti * ante fáciem ómnium populórum,

Lumen ad revelatiónem géntium, * et glóriam plebis tuæ Israël.

Glória Patri.

\antlabel{Ant.} Salva nos, Dómine, vigilántes, custódi nos dormiéntes; ut vigilémus cum Christo, et requiescámus in pace.

\variantheading{Preces (si dicendæ sint)}

KYRIE, eléison. Christe, eléison. Kýrie, eléison.

Pater noster \rubric{secreto usque ad} \V\ Et ne nos indúcas in tentatiónem. \R\ Sed líbera nos a malo.

Credo in Deum \rubric{secreto usque ad} \V\ Carnis resurrectiónem. \R\ Vitam ætérnam. Amen. \V\ Benedíctus es, Dómine, Deus patrum nostrórum. \R\ Et laudábilis et gloriósus in sæcula. \V\ Benedicámus Patrem et Fílium cum Sancto Spíritu. \R\ Laudémus, et superexaltémus eum in sæcula. \V\ Benedíctus es, Dómine, in firmaménto cæli. \R\ Et laudábilis, et gloriósus, et superexaltátus in sæcula. \V\ Benedícat et custódiat nos omnípotens et miséricors Dóminus. \R\ Amen. \V\ Dignáre, Dómine, nocte ista. \R\ Sine peccáto nos custodíre. \V\ Miserére nostri, Dómine. \R\ Miserére nostri. \V\ Fiat misericórdia tua, Dómine, super nos. \R\ Quemádmodum sperávimus in te. \V\ Dómine, exáudi oratiónem meam. \R\ Et clamor meus ad te véniat.

\V\ Dóminus vobíscum.

Orémus. \oratiolabel{Oratio}

\lettrine[lines=2]{V}{ISITA}, quæsumus, Dómine, habitatiónem istam, et omnes insídias inimíci ab ea longe repélle: Angeli tui sancti hábitent in ea, qui nos in pace custódiant; et benedíctio tua sit super nos semper. Per Dóminum nostrum. \R\ Amen.

\V\ Dóminus vobíscum.

\V\ Benedicámus Dómino.

\benedictiolabel{Benedictio} Benedícat et custódiat nos omnípotens et miséricors Dóminus, Pater, et Fílius, et Spíritus Sanctus. \R\ Amen.

\variantheading{Antiphona finalis B. M. V.}

\V\ Divínum auxílium.

Pater, Ave et Credo.

% ===== END OF BATCH 23 -- Covers autumnalis.md lines 7799-8128,
% completing "Ad Vesperas" (Psalms 127, 128, 129, 130, 131, hymn
% "Cæli Deus sanctíssime") and "Ad Completorium" (Psalms 33/i,
% 33/ii, 60, hymn "Te lucis ante términum," Canticum Simeonis,
% Preces), and with them the ENTIRE "Feria Quarta" Office. More
% partial-caps OCR defects ("MEmento," "BEnedictus" [Capitulum],
% "DOmine," "DE profúndis," "BEnedicam," "VEnite," "EXaudi") are
% silently normalized to full caps, consistent with Batches 21-22.
% The partial-bold "Confíteor." (not the full "Misereátur.
% Indulgéntiam.") again appears at Compline, rendered the same way
% as in Feria Secunda and Feria Tertia. The document continues with
% "Feria Quinta" (autumnalis.md line 8131), the fourth weekday of the
% Psalterium cycle.
% =====
\feriaheading{Feria Quinta}

\hourheading{Ad Matutinum}

Pater, Ave et Credo.

\V\ Dómine, lábia mea.

\V\ Deus, in adjutórium.

\variantheading{In Officio feriali}

\invitatlabel{Invit.} Regem ventúrum Dóminum, * Veníte, adorémus. \psalmlabel{Ps. 94} Veníte.

\hymnusheading{Hymnus}

\lettrine[lines=2]{N}{OX} atra rerum cóntegit\\
Terræ colóres ómnium:\\
Nos confiténtes póscimus\\
Te, juste judex córdium:

Ut áuferas piácula,\\
Sordésque mentis ábluas:\\
Donésque, Christe, grátiam,\\
Ut arceántur crímina.

Mens ecce torpet ímpia,\\
Quam culpa mordet nóxia:\\
Obscúra gestit tóllere,\\
Et te, Redémptor, quærere.

Repélle tu calíginem\\
Intrínsecus quam máxime,\\
Ut in beáto gáudeat\\
Se collocári lúmine.

Præsta, Pater piíssime,\\
Patríque compar Unice,\\
Cum Spíritu Paráclito\\
Regnans per omne sæculum.\\
Amen.

\nocturnoheading{In I Nocturno}

\antlabel{Ant.} In Deo salutáre meum, * et glória mea: et spes mea in Deo est.

\psalmheading{Psalmus 61}

\lettrine[lines=2]{N}{ONNE} Deo subjécta erit ánima mea? * ab ipso enim salutáre meum.

Nam et ipse Deus meus, et salutáris meus: * suscéptor meus, non movébor ámplius.

Quoúsque irrúitis in hóminem? * interfícitis univérsi vos: tamquam paríeti inclináto et macériæ depúlsæ?

Verúmtamen prétium meum cogitavérunt repéllere, cucúrri in siti: * ore suo benedicébant, et corde suo maledicébant.

Verúmtamen Deo subjécta esto, ánima mea: * quóniam ab ipso patiéntia mea.

Quia ipse Deus meus, et salvátor meus: * adjútor meus, non emigrábo.

In Deo salutáre meum, et glória mea: * Deus auxílii mei, et spes mea in Deo est.

Speráte in eo, omnis congregátio pópuli, effúndite coram illo corda vestra: * Deus adjútor noster in ætérnum.

Verúmtamen vani fílii hóminum, mendáces fílii hóminum in statéris: * ut decípiant ipsi de vanitáte in idípsum.

Nolíte speráre in iniquitáte, et rapínas nolíte concupíscere: * divítiæ si áffluant, nolíte cor appónere.

Semel locútus est Deus, duo hæc audívi, quia potéstas Dei est, et tibi, Dómine, misericórdia: * quia tu reddes unicuíque juxta ópera sua.

\antlabel{Ant.} In Deo salutáre meum, et glória mea: et spes mea in Deo est.

\antlabel{Ant.} Vidéte ópera Dómini, * et audítam fácite vocem laudis ejus.

\psalmheading{Psalmus 65, i}

\lettrine[lines=2]{J}{UBILATE} Deo, omnis terra, psalmum dícite nómini ejus: * date glóriam laudi ejus.

Dícite Deo: Quam terribília sunt ópera tua, Dómine! * in multitúdine virtútis tuæ mentiéntur tibi inimíci tui.

Omnis terra adóret te, et psallat tibi: * psalmum dicat nómini tuo.

Veníte, et vidéte ópera Dei: * terríbilis in consíliis super fílios hóminum.

Qui convértit mare in áridam, in flúmine pertransíbunt pede: * ibi lætábimur in ipso.

Qui dominátur in virtúte sua in ætérnum, óculi ejus super gentes respíciunt: * qui exásperant non exalténtur in semetípsis.

Benedícite, gentes, Deum nostrum: * et audítam fácite vocem laudis ejus,

Qui pósuit ánimam meam ad vitam: * et non dedit in commotiónem pedes meos.

Quóniam probásti nos, Deus: * igne nos examinásti, sicut examinátur argéntum.

Induxísti nos in láqueum, posuísti tribulatiónes in dorso nostro: * imposuísti hómines super cápita nostra.

Transívimus per ignem et aquam: * et eduxísti nos in refrigérium.

\antlabel{Ant.} Vidéte ópera Dómini, et audítam fácite vocem laudis ejus.

\antlabel{Ant.} Audíte, omnes, * qui timétis Deum, quanta fecit ánimæ meæ.

\psalmheading{Psalmus 65, ii}

\lettrine[lines=2]{I}{NTROIBO} in domum tuam in holocáustis: * reddam tibi vota mea, quæ distinxérunt lábia mea.

Et locútum est os meum, * in tribulatióne mea.

Holocáusta medulláta ófferam tibi cum incénso aríetum: * ófferam tibi boves cum hircis.

Veníte, audíte, et narrábo, omnes, qui timétis Deum, * quanta fecit ánimæ meæ.

Ad ipsum ore meo clamávi, * et exaltávi sub lingua mea.

Iniquitátem si aspéxi in corde meo, * non exáudiet Dóminus.

Proptérea exaudívit Deus, * et atténdit voci deprecatiónis meæ.

Benedíctus Deus, * qui non amóvit oratiónem meam, et misericórdiam suam a me.

\antlabel{Ant.} Audíte, omnes, qui timétis Deum, quanta fecit ánimæ meæ.

\variantheading{In Festis novem Lectionum:}

\V\ Non amóvit Dóminus oratiónem meam. \R\ Et misericórdiam suam a me.

\absolutiolabel{Absolutio.} Exáudi, Dómine.

\benedlabel{Bened.} 1 Benedictióne perpétua. 2 Unigénitus Dei Fílius. 3 Spíritus Sancti grátia.

\nocturnoheading{In II Nocturno}

\antlabel{Ant.} Exsúrgat Deus, * \textdagger\ et dissipéntur inimíci ejus. \textdaggerdbl

\psalmheading{Psalmus 67, i}

\lettrine[lines=2]{E}{XSURGAT} Deus, \textdagger\ et dissipéntur inimíci ejus, * \textdaggerdbl\ et fúgiant qui odérunt eum, a fácie ejus.

Sicut déficit fumus, defíciant: * sicut fluit cera a fácie ignis, sic péreant peccatóres a fácie Dei.

Et justi epuléntur, et exsúltent in conspéctu Dei: * et delecténtur in lætítia.

Cantáte Deo, psalmum dícite nómini ejus: * iter fácite ei, qui ascéndit super occásum: Dóminus nomen illi.

Exsultáte in conspéctu ejus: * turbabúntur a fácie ejus, patris orphanórum, et júdicis viduárum.

Deus in loco sancto suo: * Deus qui inhabitáre facit uníus moris in domo:

Qui edúcit vinctos in fortitúdine, * simíliter eos, qui exásperant, qui hábitant in sepúlcris.

Deus, cum egrederéris in conspéctu pópuli tui, * cum pertransíres in desérto:

Terra mota est, étenim cæli distillavérunt a fácie Dei Sínaï, * a fácie Dei Israël.

Plúviam voluntáriam segregábis, Deus, hereditáti tuæ: * et infirmáta est, tu vero perfecísti eam.

Animália tua habitábunt in ea: * parásti in dulcédine tua páuperi, Deus.

\antlabel{Ant.} Exsúrgat Deus, et dissipéntur inimíci ejus.

\antlabel{Ant.} Deus noster, * Deus salvos faciéndi: et Dómini sunt éxitus mortis.

\psalmheading{Psalmus 67, ii}

\lettrine[lines=2]{D}{OMINUS} dabit verbum evangelizántibus, * virtúte multa.

Rex virtútum dilécti dilécti: * et speciéi domus divídere spólia.

Si dormiátis inter médios cleros, pennæ colúmbæ deargentátæ, * et posterióra dorsi ejus in pallóre auri.

Dum discérnit cæléstis reges super eam, nive dealbabúntur in Selmon: * mons Dei, mons pinguis.

Mons coagulátus, mons pinguis: * ut quid suspicámini montes coagulátos?

Mons, in quo beneplácitum est Deo habitáre in eo: * étenim Dóminus habitábit in finem.

Currus Dei decem míllibus múltiplex, míllia lætántium: * Dóminus in eis in Sina in sancto.

Ascendísti in altum, cepísti captivitátem: * accepísti dona in homínibus:

Etenim non credéntes, * inhabitáre Dóminum Deum.

Benedíctus Dóminus die quotídie: * prósperum iter fáciet nobis Deus salutárium nostrórum.

Deus noster, Deus salvos faciéndi: * et Dómini Dómini éxitus mortis.

Verúmtamen Deus confrínget cápita inimicórum suórum: * vérticem capílli perambulántium in delíctis suis.

Dixit Dóminus: Ex Basan convértam, * convértam in profúndum maris:

Ut intingátur pes tuus in sánguine: * lingua canum tuórum ex inimícis, ab ipso.

\antlabel{Ant.} Deus noster, Deus salvos faciéndi: et Dómini sunt éxitus mortis.

\antlabel{Ant.} In ecclésiis * benedícite Dómino Deo.

\psalmheading{Psalmus 67, iii}

\lettrine[lines=2]{V}{IDERUNT} ingréssus tuos, Deus, * ingréssus Dei mei: regis mei qui est in sancto.

Prævenérunt príncipes conjúncti psalléntibus, * in médio juvenculárum tympanistriárum.

In ecclésiis benedícite Deo Dómino, * de fóntibus Israël.

Ibi Bénjamin adolescéntulus, * in mentis excéssu.

Príncipes Juda, duces eórum: * príncipes Zábulon, príncipes Néphthali.

Manda, Deus, virtúti tuæ: * confírma hoc, Deus, quod operátus es in nobis.

A templo tuo in Jerúsalem, * tibi ófferent reges múnera.

Íncrepa feras arúndinis, congregátio taurórum in vaccis populórum: * ut exclúdant eos, qui probáti sunt argénto.

Díssipa gentes, quæ bella volunt: vénient legáti ex Ægýpto: * Æthiópia prævéniet manus ejus Deo.

Regna terræ, cantáte Deo: * psállite Dómino.

Psállite Deo, qui ascéndit super cælum cæli, * ad Oriéntem.

Ecce, dabit voci suæ vocem virtútis, date glóriam Deo super Israël, * magnificéntia ejus, et virtus ejus in núbibus.

Mirábilis Deus in sanctis suis, Deus Israël ipse dabit virtútem, et fortitúdinem plebi suæ, * benedíctus Deus.

\antlabel{Ant.} In ecclésiis benedícite Dómino Deo.

\variantheading{In Festis novem Lectionum:}

\V\ Mirábilis Deus in sanctis suis. \R\ Deus Israël dabit virtútem plebi tuæ.

\absolutiolabel{Absolutio.} Ipsíus píetas et.

\benedlabel{Bened.} 4 Deus Pater omnípotens. 5 Christus perpétuæ. 6 Ignem sui amóris.

\nocturnoheading{In III Nocturno}

\antlabel{Ant.} Salvum me fac, * \textdagger\ Deus: quóniam intravérunt aquæ usque ad ánimam meam. \textdaggerdbl

\psalmheading{Psalmus 68, i}

\lettrine[lines=2]{S}{ALVUM} me fac, \textdagger\ Deus: * \textdaggerdbl\ quóniam intravérunt aquæ usque ad ánimam meam.

Infíxus sum in limo profúndi: * et non est substántia.

Veni in altitúdinem maris: * et tempéstas demérsit me.

Laborávi clamans, raucæ factæ sunt fauces meæ: * defecérunt óculi mei, dum spero in Deum meum.

Multiplicáti sunt super capíllos cápitis mei, * qui odérunt me gratis.

Confortáti sunt qui persecúti sunt me inimíci mei injúste: * quæ non rápui, tunc exsolvébam.

Deus, tu scis insipiéntiam meam: * et delícta mea a te non sunt abscóndita.

Non erubéscant in me qui exspéctant te, Dómine, * Dómine virtútum.

Non confundántur super me * qui quærunt te, Deus Israël.

Quóniam propter te sustínui oppróbrium: * opéruit confúsio fáciem meam.

Extráneus factus sum frátribus meis, * et peregrínus fíliis matris meæ.

Quóniam zelus domus tuæ comédit me: * et oppróbria exprobrántium tibi cecidérunt super me.

Et opérui in jejúnio ánimam meam: * et factum est in oppróbrium mihi.

Et pósui vestiméntum meum cilícium: * et factus sum illis in parábolam.

Advérsum me loquebántur, qui sedébant in porta: * et in me psallébant qui bibébant vinum.

Ego vero oratiónem meam ad te, Dómine: * tempus benepláciti, Deus.

\antlabel{Ant.} Salvum me fac, Deus: quóniam intravérunt aquæ usque ad ánimam meam.

\antlabel{Ant.} Propter inimícos meos * éripe me, Dómine.

\psalmheading{Psalmus 68, ii}

\lettrine[lines=2]{I}{N} multitúdine misericórdiæ tuæ exáudi me, * in veritáte salútis tuæ:

Eripe me de luto, ut non infígar: * líbera me ab iis, qui odérunt me, et de profúndis aquárum.

Non me demérgat tempéstas aquæ, neque absórbeat me profúndum: * neque úrgeat super me púteus os suum.

Exáudi me, Dómine, quóniam benígna est misericórdia tua: * secúndum multitúdinem miseratiónum tuárum réspice in me.

Et ne avértas fáciem tuam a púero tuo: * quóniam tríbulor, velóciter exáudi me.

Inténde ánimæ meæ, et líbera eam: * propter inimícos meos éripe me.

Tu scis impropérium meum, et confusiónem meam, * et reveréntiam meam.

In conspéctu tuo sunt omnes qui tríbulant me, * impropérium exspectávit cor meum, et misériam.

Et sustínui qui simul contristarétur, et non fuit: * et qui consolarétur, et non invéni.

Et dedérunt in escam meam fel: * et in siti mea potavérunt me acéto.

Fiat mensa eórum coram ipsis in láqueum, * et in retributiónes, et in scándalum.

Obscuréntur óculi eórum ne vídeant: * et dorsum eórum semper incúrva.

Effúnde super eos iram tuam: * et furor iræ tuæ comprehéndat eos.

Fiat habitátio eórum desérta: * et in tabernáculis eórum non sit qui inhábitet.

Quóniam quem tu percussísti, persecúti sunt: * et super dolórem vúlnerum meórum addidérunt.

Appóne iniquitátem super iniquitátem eórum: * et non intrent in justítiam tuam.

Deleántur de libro vivéntium: * et cum justis non scribántur.

\antlabel{Ant.} Propter inimícos meos éripe me, Dómine.

\antlabel{Ant.} Quærite Dóminum, * et vivet ánima vestra.

\psalmheading{Psalmus 68, iii}

\lettrine[lines=2]{E}{GO} sum pauper et dolens: * salus tua, Deus, suscépit me.

Laudábo nomen Dei cum cántico: * et magnificábo eum in laude:

Et placébit Deo super vítulum novéllum: * córnua producéntem et úngulas.

Vídeant páuperes et læténtur: * quærite Deum, et vivet ánima vestra.

Quóniam exaudívit páuperes Dóminus: * et vinctos suos non despéxit.

Laudent illum cæli et terra, * mare, et ómnia reptília in eis.

Quóniam Deus salvam fáciet Sion: * et ædificabúntur civitátes Juda.

Et inhabitábunt ibi, * et hereditáte acquírent eam.

Et semen servórum ejus possidébit eam, * et qui díligunt nomen ejus, habitábunt in ea.

\antlabel{Ant.} Quærite Dóminum, et vivet ánima vestra.

\variantheading{In feriali Officio et in Festis:}

\V\ Laudábo nomen Dei cum cántico. \R\ Et magnificábo eum in laude.

\variantheading{In Officio novem Lectionum}

\absolutiolabel{Absolutio.} A vínculis.

\benedlabel{Bened.} 7 Evangélica léctio. 8 In Festis Dómini: Divínum auxílium. In Festis B. M. V.: Cujus festum cólimus, ipsa Virgo. In Festis Sanctórum: Cujus (\textit{vel} Quorum \textit{aut} Quarum) festum cólimus. 9 Ad societátem. (Per evangélica dicta.)

\variantheading{In Officio trium Lectionum}

\absolutiolabel{Absolutio.} Exáudi, Dómine.

\variantheading{In Feriis cum Homilia:}

\benedlabel{Bened.} 1 Evangélica. 2 Divínum. 3 Ad societátem.

\variantheading{In Feriis sine Homilia:}

\benedlabel{Bened.} 1 Benedictióne. 2 Unigénitus. 3 Spíritus Sancti.

\variantheading{In Festis Sanctorum:}

\benedlabel{Bened.} 1 Ille nos benedícat, qui sine fine vivit et regnat. \R\ Amen. 2 Cujus (\textit{vel} Quorum \textit{aut} Quarum). 3 Ad societátem.

% ===== END OF BATCH 24 -- Covers autumnalis.md lines 8133-8494,
% opening "Feria Quinta" (Thursday) and completing its "Ad
% Matutinum" in full: the ferial Invitatorium "Regem ventúrum
% Dóminum," the hymn "Nox atra rerum cóntegit," and the three
% Nocturns (Psalms 61, 65/i, 65/ii; 67/i-iii; 68/i-iii) with their
% versicle, Absolutio, and Benedictiones apparatus. "In III
% Nocturno" here has only a single form (no "I"/"II" alternation as
% in Feria Quarta), confirming that split was specific to the
% Psalm-49/Psalm-50 overlap and does not recur every day. Several
% more psalm incipits carry the partial-caps OCR defect already
% documented in prior batches ("NOnne," "JUbiláte," "INtroíbo,"
% "EXsurgat," "DOminus," "VIderunt," "SAlvum," "IN" multitúdine,
% "EGo") and are silently normalized to full caps. The document
% continues with "Ad Laudes I" (autumnalis.md line 8497).
% =====
\hourheading{Ad Laudes I}

\rubric{In Fériis extra Vigílias commúnes occurréntibus, atque in Festis.}

(Pater noster et Ave.)

\V\ Deus, in adjutórium.

\antlabel{Ant.} Jubiláte * in conspéctu regis Dómini.

\psalmheading{Psalmus 97}

\lettrine[lines=2]{C}{ANTATE} Dómino cánticum novum: * quia mirabília fecit.

Salvávit sibi déxtera ejus: * et bráchium sanctum ejus.

Notum fecit Dóminus salutáre suum: * in conspéctu géntium revelávit justítiam suam.

Recordátus est misericórdiæ suæ, * et veritátis suæ dómui Israël.

Vidérunt omnes términi terræ * salutáre Dei nostri.

Jubiláte Deo, omnis terra: * cantáte, et exsultáte, et psállite.

Psállite Dómino in cíthara, in cíthara et voce psalmi: * in tubis ductílibus, et voce tubæ córneæ.

Jubiláte in conspéctu regis Dómini: * moveátur mare, et plenitúdo ejus: orbis terrárum, et qui hábitant in eo.

Flúmina plaudent manu, simul montes exsultábunt a conspéctu Dómini: * quóniam venit judicáre terram.

Judicábit orbem terrárum in justítia, * et pópulos in æquitáte.

\antlabel{Ant.} Jubiláte in conspéctu regis Dómini.

\antlabel{Ant.} Dómine, * \textdagger\ refúgium factus es nobis. \textdaggerdbl

\psalmheading{Psalmus 89}

\lettrine[lines=2]{D}{OMINE}, \textdagger\ refúgium factus es nobis: * \textdaggerdbl\ a generatióne in generatiónem.

Priúsquam montes fíerent, aut formarétur terra et orbis: * a sæculo et usque in sæculum tu es, Deus.

Ne avértas hóminem in humilitátem: * et dixísti: Convertímini, fílii hóminum.

Quóniam mille anni ante óculos tuos, * tamquam dies hestérna, quæ prætériit.

Et custódia in nocte, * quæ pro níhilo habéntur, eórum anni erunt.

Mane sicut herba tránseat, mane flóreat, et tránseat: * véspere décidat, indúret et aréscat.

Quia defécimus in ira tua, * et in furóre tuo turbáti sumus.

Posuísti iniquitátes nostras in conspéctu tuo: * sæculum nostrum in illuminatióne vultus tui.

Quóniam omnes dies nostri defecérunt: * et in ira tua defécimus.

Anni nostri sicut aránea meditabúntur: * dies annórum nostrórum in ipsis, septuagínta anni.

Si autem in potentátibus octogínta anni: * et ámplius eórum, labor et dolor.

Quóniam supervénit mansuetúdo: * et corripiémur.

Quis novit potestátem iræ tuæ: * et præ timóre tuo iram tuam dinumeráre?

Déxteram tuam sic notam fac: * et erudítos corde in sapiéntia.

Convértere, Dómine, úsquequo? * et deprecábilis esto super servos tuos.

Repléti sumus mane misericórdia tua: * et exsultávimus, et delectáti sumus ómnibus diébus nostris.

Lætáti sumus pro diébus, quibus nos humiliásti: * annis, quibus vídimus mala.

Réspice in servos tuos, et in ópera tua: * et dírige fílios eórum.

Et sit splendor Dómini, Dei nostri, super nos, et ópera mánuum nostrárum dírige super nos: * et opus mánuum nostrárum dírige.

\antlabel{Ant.} Dómine, refúgium factus es nobis.

\antlabel{Ant.} Dómine, * in cælo misericórdia tua.

\psalmheading{Psalmus 35}

\lettrine[lines=2]{D}{IXIT} injústus ut delínquat in semetípso: * non est timor Dei ante óculos ejus.

Quóniam dolóse egit in conspéctu ejus: * ut inveniátur iníquitas ejus ad ódium.

Verba oris ejus iníquitas, et dolus: * nóluit intellígere ut bene ágeret.

Iniquitátem meditátus est in cubíli suo: * ástitit omni viæ non bonæ, malítiam autem non odívit.

Dómine, in cælo misericórdia tua: * et véritas tua usque ad nubes.

Justítia tua sicut montes Dei: * judícia tua abýssus multa.

Hómines, et juménta salvábis, Dómine: * quemádmodum multiplicásti misericórdiam tuam, Deus.

Fílii autem hóminum, * in tégmine alárum tuárum sperábunt.

Inebriabúntur ab ubertáte domus tuæ: * et torrénte voluptátis tuæ potábis eos.

Quóniam apud te est fons vitæ: * et in lúmine tuo vidébimus lumen.

Præténde misericórdiam tuam sciéntibus te, * et justítiam tuam his, qui recto sunt corde.

Non véniat mihi pes supérbiæ: * et manus peccatóris non móveat me.

Ibi cecidérunt qui operántur iniquitátem: * expúlsi sunt, nec potuérunt stare.

\antlabel{Ant.} Dómine, in cælo misericórdia tua.

\antlabel{Ant.} Pópulus meus, * ait Dóminus, bonis meis adimplébitur.

\canticumheading{Canticum Jeremiæ}{Jer. 31, 10-14}

\lettrine[lines=2]{A}{UDITE} verbum Dómini, gentes, * et annuntiáte in ínsulis, quæ procul sunt,

Et dícite: Qui dispérsit Israël, congregábit eum: * et custódiet eum sicut pastor gregem suum.

Redémit enim Dóminus Jacob, * et liberábit eum de manu potentióris.

Et vénient, et laudábunt in monte Sion: * et confluent ad bona Dómini.

Super fruménto, et vino, et óleo, * et fetu pécorum et armentórum:

Erítque ánima eórum quasi hortus irríguus, * et ultra non esúrient.

Tunc lætábitur virgo in choro, * júvenes et senes simul:

Et convértam luctum eórum in gáudium, * et consolábor eos, et lætificábo a dolóre suo.

Et inebriábo ánimam sacerdótum pinguédine: * et pópulus meus bonis meis adimplébitur.

\antlabel{Ant.} Pópulus meus, ait Dóminus, bonis meis adimplébitur.

\antlabel{Ant.} Deo nostro * jucúnda sit laudátio.

\psalmheading{Psalmus 146}

\lettrine[lines=2]{L}{AUDATE} Dóminum quóniam bonus est psalmus: * Deo nostro sit jucúnda, decóraque laudátio.

Ædíficans Jerúsalem Dóminus: * dispersiónes Israélis congregábit.

Qui sanat contrítos corde: * et álligat contritiónes eórum.

Qui númerat multitúdinem stellárum: * et ómnibus eis nómina vocat.

Magnus Dóminus noster, et magna virtus ejus: * et sapiéntiæ ejus non est númerus.

Suscípiens mansuétos Dóminus: * humílians autem peccatóres usque ad terram.

Præcínite Dómino in confessióne: * psállite Deo nostro in cíthara.

Qui óperit cælum núbibus: * et parat terræ plúviam.

Qui prodúcit in móntibus fœnum: * et herbam servitúti hóminum.

Qui dat juméntis escam ipsórum: * et pullis corvórum invocántibus eum.

Non in fortitúdine equi voluntátem habébit: * nec in tíbiis viri beneplácitum erit ei.

Beneplácitum est Dómino super timéntes eum: * et in eis, qui sperant super misericórdia ejus.

\antlabel{Ant.} Deo nostro jucúnda sit laudátio.

\variantheading{In Officio feriali}

\capitulumlabel{Capit.}{Rom. 13, 12-13}

\lettrine[lines=2]{N}{OX} præcéssit, dies autem appropinquávit. Abjiciámus ergo ópera tenebrárum, et induámur arma lucis. Sicut in die honéste ambulémus.

\hymnusheading{Hymnus}

\lettrine[lines=2]{L}{UX} ecce surgit áurea,\\
Pallens facéssat cæcitas,\\
Quæ nosmet in præceps diu\\
Erróre traxit dévio.

Hæc lux serénum cónferat,\\
Purósque nos præstet sibi:\\
Nihil loquámur súbdolum:\\
Volvámus obscúrum nihil.

Sic tota decúrrat dies,\\
Ne lingua mendax, ne manus\\
Oculíve peccent lúbrici,\\
Ne noxa corpus ínquinet.

Speculátor astat désuper,\\
Qui nos diébus ómnibus,\\
Actúsque nostros próspicit\\
A luce prima in vésperum.

Deo Patri sit glória,\\
Ejúsque soli Fílio,\\
Cum Spíritu Paráclito,\\
Nunc et per omne sæculum.\\
Amen.

\V\ Repléti sumus mane misericórdia tua. \R\ Exsultávimus, et delectáti sumus.

\antlabel{Ad Bened. Ant.} In sanctitáte * serviámus Dómino, et liberábit nos ab inimícis nostris.

\rubric{Réliqua ut in Ordinário. 9}

% ===== END OF BATCH 25 -- Covers autumnalis.md lines 8497-8702, "Ad
% Laudes I" for Feria Quinta in full: Psalms 97, 89, 35, the
% Canticum Jeremiæ, and Psalm 146, followed by the Capitulum "Nox
% præcéssit" and the hymn "Lux ecce surgit áurea." Several more
% psalm incipits carry the partial-caps OCR defect already
% documented in prior batches ("DOmine," "DIxit," "AUdíte,"
% "LAudáte," "NOx" [Capitulum]) and are silently normalized to full
% caps. The document continues with "Ad Laudes II" (autumnalis.md
% line 8705).
% =====
\hourheading{Ad Laudes II}

\rubric{In Vigíliis commúnibus, quando fit Offícium de Féria.}

(Pater noster et Ave.)

\V\ Deus, in adjutórium.

\antlabel{Ant.} Tibi soli peccávi.

\psalmheading{Psalmus 50}

\lettrine[lines=2]{M}{ISERERE} mei, Deus, * secúndum magnam misericórdiam tuam.

Et secúndum multitúdinem miseratiónum tuárum, * dele iniquitátem meam.

Amplius lava me ab iniquitáte mea: * et a peccáto meo munda me.

Quóniam iniquitátem meam ego cognósco: * et peccátum meum contra me est semper.

Tibi soli peccávi, et malum coram te feci: * ut justificéris in sermónibus tuis, et vincas cum judicáris.

Ecce enim, in iniquitátibus concéptus sum: * et in peccátis concépit me mater mea.

Ecce enim, veritátem dilexísti: * incérta et occúlta sapiéntiæ tuæ manifestásti mihi.

Aspérges me hyssópo, et mundábor: * lavábis me, et super nivem dealbábor.

Audítui meo dabis gáudium et lætítiam: * et exsultábunt ossa humiliáta.

Avérte fáciem tuam a peccátis meis: * et omnes iniquitátes meas dele.

Cor mundum crea in me, Deus: * et spíritum rectum ínnova in viscéribus meis.

Ne projícias me a fácie tua: * et spíritum sanctum tuum ne áuferas a me.

Redde mihi lætítiam salutáris tui: * et spíritu principáli confírma me.

Docébo iníquos vias tuas: * et ímpii ad te converténtur.

Líbera me de sanguínibus, Deus, Deus salútis meæ: * et exsultábit lingua mea justítiam tuam.

Dómine, lábia mea apéries: * et os meum annuntiábit laudem tuam.

Quóniam si voluísses sacrifícium, dedíssem útique: * holocáustis non delectáberis.

Sacrifícium Deo spíritus contribulátus: * cor contrítum, et humiliátum, Deus, non despícies.

Benígne fac, Dómine, in bona voluntáte tua Sion: * ut ædificéntur muri Jerúsalem.

Tunc acceptábis sacrifícium justítiæ, oblatiónes, et holocáusta: * tunc impónent super altáre tuum vítulos.

\antlabel{Ant.} Tibi soli peccávi, Dómine: miserére mei.

\antlabel{Ant.} Convértere, Dómine.

\psalmheading{Psalmus 89}

\lettrine[lines=2]{D}{OMINE}, refúgium factus es nobis: * a generatióne in generatiónem.

Priúsquam montes fíerent, aut formarétur terra et orbis: * a sæculo et usque in sæculum tu es, Deus.

Ne avértas hóminem in humilitátem: * et dixísti: Convertímini, fílii hóminum.

Quóniam mille anni ante óculos tuos, * tamquam dies hestérna, quæ prætériit.

Et custódia in nocte, * quæ pro níhilo habéntur, eórum anni erunt.

Mane sicut herba tránseat, mane flóreat, et tránseat: * véspere décidat, indúret et aréscat.

Quia defécimus in ira tua, * et in furóre tuo turbáti sumus.

Posuísti iniquitátes nostras in conspéctu tuo: * sæculum nostrum in illuminatióne vultus tui.

Quóniam omnes dies nostri defecérunt: * et in ira tua defécimus.

Anni nostri sicut aránea meditabúntur: * dies annórum nostrórum in ipsis, septuagínta anni.

Si autem in potentátibus, octogínta anni: * et ámplius eórum, labor et dolor.

Quóniam supervénit mansuetúdo: * et corripiémur.

Quis novit potestátem iræ tuæ: * et præ timóre tuo iram tuam dinumeráre?

Déxteram tuam sic notam fac: * et erudítos corde in sapiéntia.

Convértere, Dómine, úsquequo? * et deprecábilis esto super servos tuos.

Repléti sumus mane misericórdia tua: * et exsultávimus, et delectáti sumus ómnibus diébus nostris.

Lætáti sumus pro diébus, quibus nos humiliásti: * annis, quibus vídimus mala.

Réspice in servos tuos, et in ópera tua: * et dírige fílios eórum.

Et sit splendor Dómini, Dei nostri, super nos, et ópera mánuum nostrárum dírige super nos: * et opus mánuum nostrárum dírige.

\antlabel{Ant.} Convértere, Dómine, et deprecábilis esto super servos tuos.

\antlabel{Ant.} Multiplicásti, Deus.

\psalmheading{Psalmus 35}

\lettrine[lines=2]{D}{IXIT} injústus ut delínquat in semetípso: * non est timor Dei ante óculos ejus.

Quóniam dolóse egit in conspéctu ejus: * ut inveniátur iníquitas ejus ad ódium.

Verba oris ejus iníquitas, et dolus: * nóluit intellígere ut bene ágeret.

Iniquitátem meditátus est in cubíli suo: * ástitit omni viæ non bonæ, malítiam autem non odívit.

Dómine, in cælo misericórdia tua: * et véritas tua usque ad nubes.

Justítia tua sicut montes Dei: * judícia tua abýssus multa.

Hómines, et juménta salvábis, Dómine: * quemádmodum multiplicásti misericórdiam tuam, Deus.

Fílii autem hóminum, * in tégmine alárum tuárum sperábunt.

Inebriabúntur ab ubertáte domus tuæ: * et torrénte voluptátis tuæ potábis eos.

Quóniam apud te est fons vitæ: * et in lúmine tuo vidébimus lumen.

Præténde misericórdiam tuam sciéntibus te, * et justítiam tuam his, qui recto sunt corde.

Non véniat mihi pes supérbiæ: * et manus peccatóris non móveat me.

Ibi cecidérunt qui operántur iniquitátem: * expúlsi sunt, nec potuérunt stare.

\antlabel{Ant.} Multiplicásti, Deus, misericórdiam tuam.

\antlabel{Ant.} Fortitúdo mea.

\canticumheading{Canticum Moysis}{Exodi 15, 1-19}

\lettrine[lines=2]{C}{ANTEMUS} Dómino: glorióse enim magnificátus est, * equum et ascensórem dejécit in mare.

Fortitúdo mea, et laus mea Dóminus, * et factus est mihi in salútem:

Iste Deus meus, et glorificábo eum: * Deus patris mei, et exaltábo eum.

Dóminus quasi vir pugnátor, Omnípotens nomen ejus. * Currus Pharaónis et exércitum ejus projécit in mare.

Elécti príncipes ejus submérsi sunt in Mari Rubro: * abýssi operuérunt eos, descendérunt in profúndum quasi lapis.

Déxtera tua, Dómine, magnificáta est in fortitúdine: déxtera tua, Dómine, percússit inimícum. * Et in multitúdine glóriæ tuæ deposuísti adversários tuos:

Misísti iram tuam, quæ devorávit eos sicut stípulam. * Et in spíritu furóris tui congregátæ sunt aquæ:

Stetit unda fluens, * congregátæ sunt abýssi in médio mari.

Dixit inimícus: Pérsequar et comprehéndam, * dívidam spólia, implébitur ánima mea:

Evaginábo gládium meum, * interfíciet eos manus mea.

Flavit spíritus tuus, et opéruit eos mare: * submérsi sunt quasi plumbum in aquis vehemméntibus.

Quis símilis tui in fórtibus, Dómine? * quis símilis tui, magníficus in sanctitáte, terríbilis atque laudábilis, fáciens mirabília?

Extendísti manum tuam, et devorávit eos terra. * Dux fuísti in misericórdia tua pópulo quem redemísti:

Et portásti eum in fortitúdine tua, * ad habitáculum sanctum tuum.

Ascendérunt pópuli, et iráti sunt: * dolóres obtinuérunt habitatóres Philísthiim.

Tunc conturbáti sunt príncipes Edom, robústos Moab obtínuit tremor: * obriguérunt omnes habitatóres Chánaan.

Irruat super eos formído et pavor, * in magnitúdine bráchii tui.

Fiant immóbiles quasi lapis, donec pertránseat pópulus tuus, Dómine, * donec pertránseat pópulus tuus iste, quem possedísti.

Introdúces eos, et plantábis in monte hereditátis tuæ, * firmíssimo habitáculo tuo quod operátus es, Dómine:

Sanctuárium tuum, Dómine, quod firmavérunt manus tuæ. * Dóminus regnábit in ætérnum et ultra.

Ingréssus est enim eques Phárao cum cúrribus et equítibus ejus in mare: * et redúxit super eos Dóminus aquas maris:

Fílii autem Israël ambulavérunt per siccum * in médio ejus.

\antlabel{Ant.} Fortitúdo mea, et laus mea Dóminus: et factus est mihi in salútem.

\antlabel{Ant.} Laudáte. \textdagger

\psalmheading{Psalmus 146}

\lettrine[lines=2]{L}{AUDATE} \textdagger\ Dóminum quóniam bonus est psalmus: * Deo nostro sit jucúnda, decóraque laudátio.

Ædíficans Jerúsalem Dóminus: * dispersiónes Israélis congregábit.

Qui sanat contrítos corde: * et álligat contritiónes eórum.

Qui númerat multitúdinem stellárum: * et ómnibus eis nómina vocat.

Magnus Dóminus noster, et magna virtus ejus: * et sapiéntiæ ejus non est númerus.

Suscípiens mansuétos Dóminus: * humílians autem peccatóres usque ad terram.

Præcínite Dómino in confessióne: * psállite Deo nostro in cíthara.

Qui óperit cælum núbibus: * et parat terræ plúviam.

Qui prodúcit in móntibus fœnum: * et herbam servitúti hóminum.

Qui dat juméntis escam ipsórum: * et pullis corvórum invocántibus eum.

Non in fortitúdine equi voluntátem habébit: * nec in tíbiis viri beneplácitum erit ei.

Beneplácitum est Dómino super timéntes eum: * et in eis, qui sperant super misericórdia ejus.

\antlabel{Ant.} Laudáte Dóminum, qui sanat contrítos corde, et álligat contritiónes eórum.

\capitulumlabel{Capit.}{Rom. 13, 12-13}

\lettrine[lines=2]{N}{OX} præcéssit, dies autem appropinquávit. Abjiciámus ergo ópera tenebrárum, et induámur arma lucis. Sicut in die honéste ambulémus.

\hymnusheading{Hymnus}

\lettrine[lines=2]{L}{UX} ecce surgit áurea,\\
Pallens facéssat cæcitas,\\
Quæ nosmet in præceps diu\\
Erróre traxit dévio.

Hæc lux serénum cónferat,\\
Purósque nos præstet sibi:\\
Nihil loquámur súbdolum:\\
Volvámus obscúrum nihil.

Sic tota decúrrat dies,\\
Ne lingua mendax, ne manus\\
Oculíve peccent lúbrici,\\
Ne noxa corpus ínquinet,

Speculátor astat désuper,\\
Qui nos diébus ómnibus,\\
Actúsque nostros próspicit\\
A luce prima in vésperum.

Deo Patri sit glória,\\
Ejúsque soli Fílio,\\
Cum Spíritu Paráclito,\\
Nunc et per omne sæculum.\\
Amen.

\V\ Repléti sumus mane misericórdia tua. \R\ Exsultávimus, et delectáti sumus.

\antlabel{Ad Bened. Ant.} In sanctitáte * serviámus Dómino, et liberábit nos ab inimícis nostris.

\rubric{Expleto Cántico ac repetíta Antíphona, dicúntur sequéntes}

\variantheading{Preces feriales}

KYRIE, eléison. Christe, eléison. Kýrie, eléison.

Pater noster \rubric{quod dícitur a solo Hebdomadário totum clara voce, usque ad} \V\ Et ne nos indúcas in tentatiónem. \R\ Sed líbera nos a malo. \V\ Ego dixi: Dómine, miserére mei. \R\ Sana ánimam meam, quia peccávi tibi. \V\ Convértere, Dómine, úsquequo? \R\ Et deprecábilis esto super servos tuos. \V\ Fiat misericórdia tua, Dómine, super nos. \R\ Quemádmodum sperávimus in te. \V\ Sacerdótes tui induántur justítiam. \R\ Et sancti tui exsúltent. \V\ Orémus pro beatíssimo Papa nostro N. \R\ Dóminus consérvet eum, et vivíficet eum, et beátum fáciat eum in terra, et non tradat eum in ánimam inimicórum ejus. \V\ Orémus et pro Antístite nostro N. \R\ Stet et pascat in fortitúdine tua, Dómine, in sublimitáte nóminis tui. \V\ Dómine, salvum fac regem. \R\ Et exáudi nos in die, qua invocavérimus te. \V\ Salvum fac pópulum tuum, Dómine, et bénedic hereditáti tuæ. \R\ Et rege eos, et extólle illos usque in ætérnum. \V\ Meménto Congregatiónis tuæ. \R\ Quam possedísti ab inítio. \V\ Fiat pax in virtúte tua. \R\ Et abundántia in túrribus tuis. \V\ Orémus pro benefactóribus nostris. \R\ Retribúere dignáre, Dómine, ómnibus, nobis bona faciéntibus propter nomen tuum, vitam ætérnam. Amen. \V\ Orémus pro fidélibus defúnctis. \R\ Réquiem ætérnam dona eis, Dómine, et lux perpétua lúceat eis. \V\ Requiéscant in pace. \R\ Amen. \V\ Pro frátribus nostris abséntibus. \R\ Salvos fac servos tuos, Deus meus, sperántes in te. \V\ Pro afflíctis et captívis. \R\ Líbera eos, Deus Israël, ex ómnibus tribulatiónibus suis. \V\ Mitte eis, Dómine, auxílium de sancto. \R\ Et de Sion tuére eos. \V\ Dómine, Deus virtútum, convérte nos. \R\ Et osténde fáciem tuam, et salvi érimus. \V\ Exsúrge, Christe, ádjuva nos. \R\ Et líbera nos propter nomen tuum. \V\ Dómine, exáudi oratiónem meam. \R\ Et clamor meus ad te véniat.

\rubric{Réliqua ut in Ordinário. 9}

% ===== END OF BATCH 26 -- Covers autumnalis.md lines 8705-8964, "Ad
% Laudes II" for Feria Quinta in full: Psalms 50, 89, 35, the
% Canticum Moysis ("Cantémus Dómino"), and Psalm 146, sharing the
% Capitulum "Nox præcéssit" and the hymn "Lux ecce surgit áurea"
% with Ad Laudes I, plus the "Preces feriales" litany. Psalm 89
% appears here without the dagger/double-dagger chant-pointing
% marks present in its Ad Laudes I occurrence (Batch 25), confirmed
% as a genuine source divergence rather than an omission. The
% document continues with "Ad Primam" (autumnalis.md line 8967).
% =====
\hourheading{Ad Primam}

Pater, Ave et Credo.

\V\ Deus, in adjutórium.

\hymnusheading{Hymnus}

\lettrine[lines=2]{J}{AM} lucis orto sídere,\\
Deum precémur súpplices,\\
Ut in diúrnis áctibus\\
Nos servet a nocéntibus.

Linguam refrénans témperet,\\
Ne litis horror ínsonet:\\
Visum fovéndo cóntegat,\\
Ne vanitátes háuriat.

Sint pura cordis íntima,\\
Absístat et vecórdia:\\
Carnis terat supérbiam\\
Potus cibíque párcitas.

Ut, cum dies abscésserit\\
Noctémque sors redúxerit,\\
Mundi per abstinéntiam\\
Ipsi canámus glóriam.

Deo Patri sit glória,\\
Ejúsque soli Fílio,\\
Cum Spíritu Paráclito,\\
Nunc et per omne sæculum.\\
Amen.

\antlabel{Ant.} In loco páscuæ.

\psalmheading{Psalmus 22}

\lettrine[lines=2]{D}{OMINUS} regit me, et nihil mihi déerit: * in loco páscuæ ibi me collocávit.

Super aquam refectiónis educávit me: * ánimam meam convértit.

Dedúxit me super sémitas justítiæ, * propter nomen suum.

Nam, et si ambulávero in médio umbræ mortis, non timébo mala: * quóniam tu mecum es.

Virga tua, et báculus tuus: * ipsa me consoláta sunt.

Parásti in conspéctu meo mensam, * advérsus eos, qui tríbulant me.

Impinguásti in óleo caput meum: * et calix meus inébrians quam præclárus est!

Et misericórdia tua subsequétur me: * ómnibus diébus vitæ meæ:

Et ut inhábitem in domo Dómini, * in longitúdinem diérum.

\psalmheading{Psalmus 71, i}

\lettrine[lines=2]{D}{EUS}, judícium tuum regi da: * et justítiam tuam fílio regis:

Judicáre pópulum tuum in justítia, * et páuperes tuos in judício.

Suscípiant montes pacem pópulo: * et colles justítiam.

Judicábit páuperes pópuli, et salvos fáciet fílios páuperum: * et humiliábit calumniatórem.

Et permanébit cum sole, et ante lunam, * in generatióne et generatiónem.

Descéndet sicut plúvia in vellus: * et sicut stillicídia stillántia super terram.

Oriétur in diébus ejus justítia, et abundántia pacis: * donec auferátur luna.

Et dominábitur a mari usque ad mare: * et a flúmine usque ad términos orbis terrárum.

\psalmheading{Psalmus 71, ii}

\lettrine[lines=2]{C}{ORAM} illo prócident Æthíopes: * et inimíci ejus terram lingent.

Reges Tharsis, et ínsulæ múnera ófferent: * reges Arabum et Saba dona addúcent.

Et adorábunt eum omnes reges terræ: * omnes gentes sérvient ei:

Quia liberábit páuperem a poténte: * et páuperem, cui non erat adjútor.

Parcet páuperi et ínopi: * et ánimas páuperum salvas fáciet.

Ex usúris et iniquitáte rédimet ánimas eórum: * et honorábile nomen eórum coram illo.

Et vivet, et dábitur ei de auro Arábiæ, et adorábunt de ipso semper: * tota die benedícent ei.

Et erit firmaméntum in terra in summis móntium, superextollétur super Líbanum fructus ejus: * et florébunt de civitáte sicut fœnum terræ.

Sit nomen ejus benedíctum in sæcula: * ante solem pérmanet nomen ejus.

Et benedicéntur in ipso omnes tribus terræ: * omnes gentes magnificábunt eum.

Benedíctus Dóminus, Deus Israël, * qui facit mirabília solus:

Et benedíctum nomen majestátis ejus in ætérnum: * et replébitur majestáte ejus omnis terra: fiat, fiat.

\rubric{¶ Quando ad Laudes dictus fúerit Ps. 50 Miserére, hic subjúngitur sequens Ps. 97, qui in Laudibus locum non hábuit.}

\psalmheading{Psalmus 97}

\lettrine[lines=2]{C}{ANTATE} Dómino cánticum novum: * quia mirabília fecit.

Salvávit sibi déxtera ejus: * et bráchium sanctum ejus.

Notum fecit Dóminus salutáre suum: * in conspéctu géntium revelávit justítiam suam.

Recordátus est misericórdiæ suæ, * et veritátis suæ dómui Israël.

Vidérunt omnes términi terræ * salutáre Dei nostri.

Jubiláte Deo, omnis terra: * cantáte, et exsultáte, et psállite.

Psállite Dómino in cíthara, in cíthara et voce psalmi: * in tubis ductílibus, et voce tubæ córneæ.

Jubiláte in conspéctu regis Dómini: * moveátur mare, et plenitúdo ejus: orbis terrárum, et qui hábitant in eo.

Flúmina plaudent manu, simul montes exsultábunt a conspéctu Dómini: * quóniam venit judicáre terram.

Judicábit orbem terrárum in justítia, * et pópulos in æquitáte.

\rubric{Secus vero, tribus Psalmis explétis, statim dícitur Ant.}

\antlabel{Ant.} In loco páscuæ, ibi Dóminus me collocávit.

\rubric{In Officio Festi vel Octávæ:}

\capitulumlabel{Capitulum}{1 Tim. 1, 17}

\lettrine[lines=2]{R}{EGI} sæculórum immortáli et invisíbili, soli Deo honor et glória in sæcula sæculórum. Amen.

\rubric{In ómnibus autem Fériis et in Vigíliis commúnibus:}

\capitulumlabel{Capitulum}{Zach. 8, 19}

\lettrine[lines=2]{P}{ACEM} et veritátem dilígite, ait Dóminus omnípotens.

\R\ br. Christe, Fili Dei vivi, * Miserére nobis. Christe. \V\ \textit{Qui sedes ad déxteram Patris.} Miserére nobis. Glória Patri. Christe. \V\ Exsúrge, Christe, ádjuva nos. \R\ Et líbera nos propter nomen tuum.

\variantheading{Preces dominicales et feriales (si dicendæ sint)}

KYRIE, eléison. Christe, eléison. Kýrie, eléison.

Pater noster \rubric{secreto usque ad} \V\ Et ne nos indúcas in tentatiónem. \R\ Sed líbera nos a malo.

Credo in Deum \rubric{secreto usque ad} \V\ Carnis resurrectiónem. \R\ Vitam ætérnam. Amen. \V\ Et ego ad te, Dómine, clamávi. \R\ Et mane orátio mea prævéniet te. \V\ Repleátur os meum laude. \R\ Ut cantem glóriam tuam, tota die magnitúdinem tuam. \V\ Dómine, avérte fáciem tuam a peccátis meis. \R\ Et omnes iniquitátes meas dele. \V\ Cor mundum crea in me, Deus. \R\ Et spíritum rectum ínnova in viscéribus meis. \V\ Ne projícias me a fácie tua. \R\ Et spíritum sanctum tuum ne áuferas a me. \V\ Redde mihi lætítiam salutáris tui. \R\ Et spíritu principáli confírma me.

\variantheading{In Précibus feriálibus tantum}

\V\ Eripe me, Dómine, ab hómine malo. \R\ A viro iníquo éripe me. \V\ Eripe me de inimícis meis, Deus meus. \R\ Et ab insurgéntibus in me líbera me. \V\ Eripe me de operántibus iniquitátem. \R\ Et de viris sánguinum salva me. \V\ Sic psalmum dicam nómini tuo in sæculum sæculi. \R\ Ut reddam vota mea de die in diem. \V\ Exáudi nos, Deus, salutáris noster. \R\ Spes ómnium fínium terræ, et in mari longe. \V\ Deus, in adjutórium meum inténde. \R\ Dómine, ad adjuvándum me festína. \V\ Sanctus Deus, Sanctus fortis, Sanctus immortális. \R\ Miserére nobis. \V\ Bénedic, ánima mea, Dómino. \R\ Et ómnia, quæ intra me sunt, nómini sancto ejus. \V\ Bénedic, ánima mea, Dómino. \R\ Et noli oblivísci omnes retributiónes ejus. \V\ Qui propitiátur ómnibus iniquitátibus tuis. \R\ Qui sanat omnes infirmitátes tuas. \V\ Qui rédimit de intéritu vitam tuam. \R\ Qui corónat te in misericórdia et miseratiónibus. \V\ Qui replet in bonis desidérium tuum. \R\ Renovábitur ut áquilæ juvéntus tua.

\variantheading{In utrisque Précibus}

\V\ Adjutórium nostrum in nómine Dómini. \R\ Qui fecit cælum et terram.

\formulalabel{Confíteor. Misereátur. Indulgéntiam.}

\V\ Dignáre, Dómine, die isto. \R\ Sine peccáto nos custodíre. \V\ Miserére nostri, Dómine. \R\ Miserére nostri. \V\ Fiat misericórdia tua, Dómine, super nos. \R\ Quemádmodum sperávimus in te. \V\ Dómine, exáudi oratiónem meam. \R\ Et clamor meus ad te véniat.

\V\ Dóminus vobíscum.

Orémus. \oratiolabel{Oratio}

\lettrine[lines=2]{D}{OMINE}, Deus omnípotens, qui ad princípium hujus diéi nos perveníre fecísti: tua nos hódie salva virtúte; ut in hac die ad nullum declinémus peccátum, sed semper ad tuam justítiam faciéndam nostra procédant elóquia, dirigántur cogitatiónes et ópera. Per Dóminum. \R\ Amen.

\V\ Dóminus vobíscum. \V\ Benedicámus Dómino.

\rubric{In Choro Martyrológium.}

\V\ Pretiósa in conspéctu Dómini. \R\ Mors Sanctórum ejus.

\rubric{Et absolúte, sine} Orémus:

\oratiolabel{Oratio}

\lettrine[lines=2]{S}{ANCTA} María, et omnes Sancti intercédant pro nobis ad Dóminum, ut nos mereámur ab eo adjuvári et salvári, qui vivit et regnat in sæcula sæculórum. \R\ Amen.

\V\ Deus, in adjutórium meum inténde. \R\ Dómine, ad adjuvándum me festína. \rubric{Et dícitur ter; último additur:} Glória Patri. Sicut erat.

Kýrie, eléison. Christe, eléison. Kýrie, eléison.

Pater noster \rubric{secreto usque ad} \V\ Et ne nos indúcas in tentatiónem. \R\ Sed líbera nos a malo. \V\ Réspice in servos tuos, Dómine, et in ópera tua, et dírige fílios eórum. \R\ Et sit splendor Dómini, Dei nostri, super nos, et ópera mánuum nostrárum dírige super nos, et opus mánuum nostrárum dírige. \V\ Glória Patri. \R\ Sicut erat.

Orémus. \oratiolabel{Oratio}

\lettrine[lines=2]{D}{IRIGERE} et sanctificáre, régere et gubernáre dignáre, Dómine Deus, Rex cæli et terræ, hódie corda et córpora nostra, sensus, sermónes et actus nostros in lege tua, et in opéribus mandatórum tuórum: ut hic, et in ætérnum, te auxiliánte, salvi et líberi esse mereámur, Salvátor mundi: Qui vivis et regnas in sæcula sæculórum. \R\ Amen.

\V\ Jube, domne.

\benedlabel{Bened.} Dies et actus nostros in sua pace dispónat Dóminus omnípotens. \R\ Amen.

\variantheading{In Officio feriali}

\lettrine[lines=2]{D}{OMINUS} autem dírigat corda et córpora nostra in caritáte Dei, et patiéntia Christi. Tu autem.

\variantheading{E Proprio Sanctorum}

\variantheading{Pro Oct. Omn. Sanctorum}

\lettrine[lines=2]{B}{ENEDÍCTIO} et cláritas et sapiéntia et gratiárum áctio, honor, virtus et fortitúdo Deo nostro in sǽcula sæculórum. Amen.

\variantheading{E Communi Sanctorum}

\variantheading{Pro Apostolis}

\lettrine[lines=2]{I}{BANT} Apóstoli gaudéntes a conspéctu concílii, quóniam digni hábiti sunt pro nómine Jesu contuméliam pati.

\variantheading{Pro Martyre}

\lettrine[lines=2]{J}{USTUS} cor suum trádidit ad vigilándum dilúculo ad Dóminum, qui fecit illum, et in conspéctu Altíssimi deprecábitur.

\variantheading{Pro pluribus Martyribus}

\lettrine[lines=2]{F}{ULGEBUNT} justi, et tamquam scintíllæ in arundinéto discúrrent. Judicábunt natiónes, et dominabúntur pópulis: et regnábit Dóminus illórum in perpétuum.

\variantheading{Pro Conf. Pont.}

\lettrine[lines=2]{F}{UNGI} sacerdótio, et habére laudem in nómine ipsíus, et offérre illi incénsum dignum in odórem suavitátis.

\variantheading{Pro Conf. non Pont.}

\lettrine[lines=2]{J}{USTUM} dedúxit Dóminus per vias rectas, et osténdit illi regnum Dei, et dedit illi sciéntiam sanctórum; honestávit illum in labóribus, et complévit labóres illíus.

\variantheading{Pro Virgine Martyre}

\lettrine[lines=2]{D}{OMINE}, Deus meus, exaltásti super terram habitatiónem meam, et pro morte defluénte deprecáta sum. Invocávi Dóminum, Patrem Dómini mei, ut non derelínquat me in die tribulatiónis meæ, et in témpore superbórum sine adjutório.

\variantheading{Pro Virgine non Martyre}

\lettrine[lines=2]{O}{} \textsc{quam} pulchra est casta generátio cum claritáte! immortális est enim memória illíus: quóniam apud Deum nota est, et apud hómines.

\variantheading{Pro Martyre tantum}

\lettrine[lines=2]{L}{AUDABIT} usque ad mortem ánima mea Dóminum, quóniam éruis sustinéntes te, et líberas eos de manu angústiæ, Dómine, Deus noster.

\variantheading{Pro nec Virg. nec Mart.}

\lettrine[lines=2]{M}{ULTAE} fíliæ congregavérunt divítias: tu supergréssa es univérsas. Fallax grátia, et vana est pulchritúdo: múlier timens Dóminum, ipsa laudábitur.

\variantheading{Pro Dedicatione Ecclesiæ}

\lettrine[lines=2]{E}{T} abstérget Deus omnem lácrimam ab óculis eórum: et mors ultra non erit, neque luctus, neque clamor, neque dolor erit ultra, quia prima abiérunt. Et dixit qui sedébat in throno: Ecce, nova fácio ómnia.

\variantheading{In Officiis B. Mariæ V.}

\lettrine[lines=2]{I}{N} platéis sicut cinnamómum et bálsamum aromatízans odórem dedi: quasi myrrha elécta dedi suavitátem odóris.

\V\ Adjutórium nostrum in nómine Dómini. \R\ Qui fecit cælum et terram.

\V\ Benedícite. \R\ Deus.

\benedictiolabel{Benedictio} Dóminus nos benedícat, et ab omni malo deféndat, et ad vitam perdúcat ætérnam. Et fidélium ánimæ per misericórdiam Dei requiéscant in pace. \R\ Amen.

\rubric{Deinde dícitur tantum} Pater noster \rubric{secréto, nisi sequátur ália Hora.}

% ===== END OF BATCH 27 -- Covers autumnalis.md lines 8967-9226,
% completing "Ad Primam" for Feria Quinta in full: the hymn "Jam
% lucis orto sídere," Psalms 22, 71/i, 71/ii, the conditional Psalm
% 97 insertion, both Capitula, the brief Responsory, the full
% Preces apparatus, the closing Orationes, and the Lectio-brevis
% apparatus (identical category list to the previous three ferias).
% Several more psalm incipits carry the partial-caps OCR defect
% ("DOminus," "DEus," "COram," "CAntáte," "PAcem" [Capitulum]) and
% are silently normalized to full caps. The document continues
% with "Ad Tertiam" (autumnalis.md line 9229).
% =====
\hourheading{Ad Tertiam}

Pater noster et Ave.

\V\ Deus, in adjutórium.

\hymnusheading{Hymnus}

\lettrine[lines=2]{N}{UNC}, Sancte, nobis, Spíritus,\\
Unum Patri cum Fílio,\\
Dignáre promptus íngeri\\
Nostro refúsus péctori.

Os, lingua, mens, sensus, vigor\\
Confessiónem pérsonent,\\
Flamméscat igne cáritas,\\
Accéndat ardor próximos.

Præsta, Pater piíssime,\\
Patríque compar Unice,\\
Cum Spíritu Paráclito\\
Regnans per omne sæculum.\\
Amen.

\antlabel{Ant.} Quam bonus. \textdagger

\psalmheading{Psalmus 72, i}

\lettrine[lines=2]{Q}{UAM} bonus \textdagger\ Israël Deus, * his qui recto sunt corde!

Mei autem pæne moti sunt pedes: * pæne effúsi sunt gressus mei.

Quia zelávi super iníquos, * pacem peccatórum videns.

Quia non est respéctus morti eórum: * et firmaméntum in plaga eórum.

In labóre hóminum non sunt, * et cum homínibus non flagellabúntur:

Ideo ténuit eos supérbia, * opérti sunt iniquitáte et impietáte sua.

Pródiit quasi ex ádipe iníquitas eórum: * transiérunt in afféctum cordis.

Cogitavérunt, et locúti sunt nequítiam: * iniquitátem in excélso locúti sunt.

Posuérunt in cælum os suum: * et lingua eórum transívit in terra.

\psalmheading{Psalmus 72, ii}

\lettrine[lines=2]{I}{DEO} convertétur pópulus meus hic: * et dies pleni inveniéntur in eis.

Et dixérunt: Quómodo scit Deus, * et si est sciéntia in excélso?

Ecce, ipsi peccatóres, et abundántes in sæculo, * obtinuérunt divítias.

Et dixi: Ergo sine causa justificávi cor meum, * et lavi inter innocéntes manus meas:

Et fui flagellátus tota die, * et castigátio mea in matutínis.

Si dicébam: Narrábo sic: * ecce, natiónem filiórum tuórum reprobávi.

Existimábam ut cognóscerem hoc, * labor est ante me:

Donec intrem in Sanctuárium Dei: * et intélligam in novíssimis eórum.

\psalmheading{Psalmus 72, iii}

\lettrine[lines=2]{V}{ERUMTAMEN} propter dolos posuísti eis: * dejecísti eos dum allevaréntur.

Quómodo facti sunt in desolatiónem, súbito defecérunt: * periérunt propter iniquitátem suam.

Velut sómnium surgéntium, Dómine, * in civitáte tua imáginem ipsórum ad níhilum rédiges.

Quia inflammátum est cor meum, et renes mei commutáti sunt: * et ego ad níhilum redáctus sum, et nescívi.

Ut juméntum factus sum apud te: * et ego semper tecum.

Tenuísti manum déxteram meam: et in voluntáte tua deduxísti me, * et cum glória suscepísti me.

Quid enim mihi est in cælo? * et a te quid vólui super terram?

Defécit caro mea, et cor meum: * Deus cordis mei, et pars mea Deus in ætérnum.

Quia ecce, qui elóngant se a te, períbunt: * perdidísti omnes, qui fornicántur abs te.

Mihi autem adhærére Deo bonum est: * pónere in Dómino Deo spem meam:

Ut annúntiem omnes prædicatiónes tuas, * in portis fíliæ Sion.

\antlabel{Ant.} Quam bonus Israël Deus, his, qui recto sunt corde?

\variantheading{In Officio feriali}

\capitulumlabel{Capitulum}{Jer. 17, 14}

\lettrine[lines=2]{S}{ANA} me, Dómine, et sanábor: salvum me fac, et salvus ero: quóniam laus mea tu es.

\R\ br. Sana ánimam meam, * Quia peccávi tibi. Sana. \V\ Ego dixi: Dómine, miserére mei. Quia. Glória Patri. Sana.

\V\ Adjútor meus esto, ne derelínquas me. \R\ Neque despícias me, Deus, salutáris meus.

\variantheading{In Vigiliis communibus:}

\variantheading{Preces feriales}

KYRIE, eléison. Christe, eléison. Kýrie, eléison.

Pater noster \rubric{secreto usque ad} \V\ Et ne nos indúcas in tentatiónem. \R\ Sed líbera nos a malo. \V\ Dómine, Deus virtútum, convérte nos. \R\ Et osténde fáciem tuam, et salvi érimus. \V\ Exsúrge, Christe, ádjuva nos. \R\ Et líbera nos propter nomen tuum. \V\ Dómine, exáudi oratiónem meam. \R\ Et clamor meus ad te véniat.

\rubric{Réliqua ut in Ordinário. 19}

\hourheading{Ad Sextam}

Pater noster et Ave.

\V\ Deus, in adjutórium.

\hymnusheading{Hymnus}

\lettrine[lines=2]{R}{ECTOR} potens, verax Deus,\\
Qui témperas rerum vices,\\
Splendóre mane illúminas,\\
Et ígnibus merídiem:

Exstíngue flammas lítium,\\
Aufer calórem nóxium,\\
Confer salútem córporum,\\
Verámque pacem córdium.

Præsta, Pater piíssime,\\
Patríque compar Unice,\\
Cum Spíritu Paráclito\\
Regnans per omne sæculum.\\
Amen.

\antlabel{Ant.} Memor esto.

\psalmheading{Psalmus 73, i}

\lettrine[lines=2]{U}{T} quid, Deus, repulísti in finem: * irátus est furor tuus super oves páscuæ tuæ?

Memor esto congregatiónis tuæ, * quam possedísti ab inítio.

Redemísti virgam hereditátis tuæ: * mons Sion, in quo habitásti in eo.

Leva manus tuas in supérbias eórum in finem: * quanta malignátus est inimícus in sancto!

Et gloriáti sunt qui odérunt te: * in médio solemnitátis tuæ.

Posuérunt signa sua, signa: * et non cognovérunt sicut in éxitu super summum.

Quasi in silva lignórum secúribus excidérunt jánuas ejus in idípsum: * in secúri et áscia dejecérunt eam.

Incendérunt igni Sanctuárium tuum: * in terra polluérunt tabernáculum nóminis tui.

Dixérunt in corde suo cognátio eórum simul: * Quiéscere faciámus omnes dies festos Dei a terra.

Signa nostra non vídimus, jam non est prophéta: * et nos non cognóscet ámplius.

\psalmheading{Psalmus 73, ii}

\lettrine[lines=2]{U}{SQUEQUO}, Deus, improperábit inimícus: * irrítat adversárius nomen tuum in finem?

Ut quid avértis manum tuam, et déxteram tuam, * de médio sinu tuo in finem?

Deus autem Rex noster ante sæcula: * operátus est salútem in médio terræ.

Tu confirmásti in virtúte tua mare: * contribulásti cápita dracónum in aquis.

Tu confregísti cápita dracónis: * dedísti eum escam pópulis Æthíopum.

Tu dirupísti fontes, et torréntes: * tu siccásti flúvios Ethan.

Tuus est dies, et tua est nox: * tu fabricátus es auróram et solem.

Tu fecísti omnes términos terræ: * æstátem et ver tu plasmásti ea.

\psalmheading{Psalmus 73, iii}

\lettrine[lines=2]{M}{EMOR} esto hujus, inimícus improperávit Dómino: * et pópulus insípiens incitávit nomen tuum.

Ne tradas béstiis ánimas confiténtes tibi, * et ánimas páuperum tuórum ne oblivíscáris in finem.

Réspice in testaméntum tuum: * quia repléti sunt, qui obscuráti sunt terræ dómibus iniquitátum.

Ne avertátur húmilis factus confúsus: * pauper et inops laudábunt nomen tuum.

Exsúrge, Deus, júdica causam tuam: * memor esto improperiórum tuórum, eórum quæ ab insipiénte sunt tota die.

Ne obliviscáris voces inimicórum tuórum: * supérbia eórum, qui te odérunt, ascéndit semper.

\antlabel{Ant.} Memor esto congregatiónis tuæ, Dómine, quam possedísti ab inítio.

\variantheading{In Officio feriali}

\capitulumlabel{Capitulum}{Rom. 13, 8}

\lettrine[lines=2]{N}{EMINI} quidquam debeátis, nisi ut ínvicem diligátis: qui enim díligit próximum, legem implévit.

\R\ br. Benedícam Dóminum * In omni témpore. Benedícam. \V\ Semper laus ejus in ore meo. In. Glória Patri. Benedícam. \V\ Dóminus regit me, et nihil mihi déerit. \R\ In loco páscuæ ibi me collocávit.

\variantheading{In Vigiliis communibus:}

\variantheading{Preces feriales}

KYRIE, eléison. Christe, eléison. Kýrie, eléison.

Pater noster \rubric{secreto usque ad} \V\ Et ne nos indúcas in tentatiónem. \R\ Sed líbera nos a malo. \V\ Dómine, Deus virtútum, convérte nos. \R\ Et osténde fáciem tuam, et salvi érimus. \V\ Exsúrge, Christe, ádjuva nos. \R\ Et líbera nos propter nomen tuum. \V\ Dómine, exáudi oratiónem meam. \R\ Et clamor meus ad te véniat.

\rubric{Réliqua ut in Ordinário. 20}

\hourheading{Ad Nonam}

Pater noster et Ave.

\V\ Deus, in adjutórium.

\hymnusheading{Hymnus}

\lettrine[lines=2]{R}{ERUM}, Deus, tenax vigor,\\
Immótus in te pérmanens,\\
Lucis diúrnæ témpora\\
Succéssibus detérminans:

Largíre lumen véspere,\\
Quo vita nusquam décidat,\\
Sed prǽmium mortis sacræ\\
Perénnis instet glória.

Præsta, Pater piíssime,\\
Patríque compar Unice,\\
Cum Spíritu Paráclito\\
Regnans per omne sæculum.\\
Amen.

\antlabel{Ant.} Invocábimus.

\psalmheading{Psalmus 74}

\lettrine[lines=2]{C}{ONFITEBIMUR} tibi, Deus: * confitébimur, et invocábimus nomen tuum.

Narrábimus mirabília tua: * cum accépero tempus, ego justítias judicábo.

Liquefácta est terra, et omnes qui hábitant in ea: * ego confirmávi colúmnas ejus.

Dixi iníquis: Nolíte iníque ágere: * et delinquéntibus: Nolíte exaltáre cornu:

Nolíte extóllere in altum cornu vestrum: * nolíte loqui advérsus Deum iniquitátem.

Quia neque ab Oriénte, neque ab Occidénte, neque a desértis móntibus: * quóniam Deus judex est.

Hunc humíliat, et hunc exáltat: * quia calix in manu Dómini vini meri plenus misto.

Et inclinávit ex hoc in hoc: verúmtamen fæx ejus non est exinaníta: * bibent omnes peccatóres terræ.

Ego autem annuntiábo in sæculum: * cantábo Deo Jacob.

Et ómnia córnua peccatórum confríngam: * et exaltabúntur córnua justi.

\psalmheading{Psalmus 75, i}

\lettrine[lines=2]{N}{OTUS} in Judǽa Deus: * in Israël magnum nomen ejus.

Et factus est in pace locus ejus: * et habitátio ejus in Sion.

Ibi confrégit poténtias árcuum, * scutum, gládium, et bellum.

Illúminans tu mirabíliter a móntibus ætérnis: * turbáti sunt omnes insipiéntes corde.

Dormiérunt somnum suum: * et nihil invenérunt omnes viri divitiárum in mánibus suis.

Ab increpatióne tua, Deus Jacob, * dormitavérunt qui ascendérunt equos.

\psalmheading{Psalmus 75, ii}

\lettrine[lines=2]{T}{U} terríbilis es, et quis resístet tibi? * ex tunc ira tua.

De cælo audítum fecísti judícium: * terra trémuit et quiévit.

Cum exsúrgeret in judícium Deus, * ut salvos fáceret omnes mansuétos terræ.

Quóniam cogitátio hóminis confitébitur tibi: * et relíquiæ cogitatiónis diem festum agent tibi.

Vovéte, et réddite Dómino, Deo vestro: * omnes, qui in circúitu ejus affértis múnera.

Terríbili et ei qui aufert spíritum príncipum, * terríbili apud reges terræ.

\antlabel{Ant.} Invocábimus nomen tuum, Dómine: narrábimus mirabília tua.

\variantheading{In Officio feriali}

\capitulumlabel{Capit.}{1 Petri 1, 17-19}

\lettrine[lines=2]{I}{N} timóre incolátus vestri témpore conversámini: sciéntes quod non corruptibílibus auro vel argénto redémpti estis, sed pretióso sánguine quasi Agni immaculáti Christi.

\R\ br. Rédime me, Dómine, * Et miserére mei. Rédime. \V\ Pes enim meus stetit in via recta. Et miserére mei. Glória Patri. Rédime. \V\ Ab occúltis meis munda me, Dómine. \R\ Et ab aliénis parce servo tuo.

\variantheading{In Vigiliis communibus:}

\variantheading{Preces feriales}

KYRIE, eléison. Christe, eléison. Kýrie, eléison.

Pater noster \rubric{secreto usque ad} \V\ Et ne nos indúcas in tentatiónem. \R\ Sed líbera nos a malo. \V\ Dómine, Deus virtútum, convérte nos. \R\ Et osténde fáciem tuam, et salvi érimus. \V\ Exsúrge, Christe, ádjuva nos. \R\ Et líbera nos propter nomen tuum. \V\ Dómine, exáudi oratiónem meam. \R\ Et clamor meus ad te véniat.

\rubric{Réliqua ut in Ordinário. 22}

% ===== END OF BATCH 28 -- Covers autumnalis.md lines 9229-9538,
% completing the Little Hours of Feria Quinta in full: "Ad
% Tertiam" (Psalms 72/i-iii), "Ad Sextam" (Psalms 73/i-iii), and
% "Ad Nonam" (Psalms 74, 75/i, 75/ii), each with its own hymn,
% Capitulum, brief Responsory, and "In Vigiliis communibus: Preces
% feriales" apparatus. Three further occurrences of the recurring
% stray OCR "\*" artifact in the hour-hymns are silently dropped.
% "Ut quid" (Psalm 73, i) correctly shows only its first short word
% in full caps ("UT"), consistent with the convention that only the
% single opening word takes full capitalization -- not a defect,
% unlike the partial-caps cases seen elsewhere. Several further
% genuine partial-caps OCR defects ("IDeo," "VErúmtamen," "USquequo,"
% "MEmor," "COnfitebimur," "NOtus") are silently normalized to full
% caps, consistent with prior batches. The document continues with
% "Ad Vesperas" (autumnalis.md line 9541).
% =====
\hourheading{Ad Vesperas}

Pater noster et Ave.

\V\ Deus, in adjutórium.

\antlabel{Ant.} Ecce, quam bonum * \textdagger\ et quam jucúndum, habitáre fratres in unum. \textdaggerdbl

\psalmheading{Psalmus 132}

\lettrine[lines=2]{E}{CCE}, quam bonum \textdagger\ et quam jucúndum, * habitáre fratres in unum: \textdaggerdbl

Sicut unguéntum in cápite, * quod descéndit in barbam, barbam Aaron.

Quod descéndit in oram vestiménti ejus: * sicut ros Hermon, qui descéndit in montem Sion.

Quóniam illic mandávit Dóminus benedictiónem, * et vitam usque in sæculum.

\antlabel{Ant.} Ecce, quam bonum et quam jucúndum, habitáre fratres in unum!

\antlabel{Ant.} Confitémini Dómino * \textdagger\ quóniam in ætérnum misericórdia ejus.

\psalmheading{Psalmus 135, i}

\lettrine[lines=2]{C}{ONFITEMINI} Dómino \textdagger\ quóniam bonus: * quóniam in ætérnum misericórdia ejus.

Confitémini Deo deórum: * quóniam in ætérnum misericórdia ejus.

Confitémini Dómino dominórum: * quóniam in ætérnum misericórdia ejus.

Qui facit mirabília magna solus: * quóniam in ætérnum misericórdia ejus.

Qui fecit cælos in intelléctu: * quóniam in ætérnum misericórdia ejus.

Qui firmávit terram super aquas: * quóniam in ætérnum misericórdia ejus.

Qui fecit luminária magna: * quóniam in ætérnum misericórdia ejus.

Solem in potestátem diéi: * quóniam in ætérnum misericórdia ejus.

Lunam, et stellas in potestátem noctis: * quóniam in ætérnum misericórdia ejus.

\antlabel{Ant.} Confitémini Dómino quóniam in ætérnum misericórdia ejus.

\antlabel{Ant.} Confitémini Dómino, * quia in humilitáte nostra memor fuit nostri.

\psalmheading{Psalmus 135, ii}

\lettrine[lines=2]{Q}{UI} percússit Ægýptum cum primogénitis eórum: * quóniam in ætérnum misericórdia ejus.

Qui edúxit Israël de médio eórum: * quóniam in ætérnum misericórdia ejus.

In manu poténti, et bráchio excélso: * quóniam in ætérnum misericórdia ejus.

Qui divísit Mare Rubrum in divisiónes: * quóniam in ætérnum misericórdia ejus.

Et edúxit Israël per médium ejus: * quóniam in ætérnum misericórdia ejus.

Et excússit Pharaónem, et virtútem ejus in Mari Rubro: * quóniam in ætérnum misericórdia ejus.

Qui tradúxit pópulum suum per desértum: * quóniam in ætérnum misericórdia ejus.

Qui percússit reges magnos: * quóniam in ætérnum misericórdia ejus.

Et occídit reges fortes: * quóniam in ætérnum misericórdia ejus.

Sehon, regem Amorrhæórum: * quóniam in ætérnum misericórdia ejus.

Et Og, regem Basan: * quóniam in ætérnum misericórdia ejus.

Et dedit terram eórum hereditátem: * quóniam in ætérnum misericórdia ejus.

Hereditátem Israël, servo suo: * quóniam in ætérnum misericórdia ejus.

Quia in humilitáte nostra memor fuit nostri: * quóniam in ætérnum misericórdia ejus.

Et redémit nos ab inimícis nostris: * quóniam in ætérnum misericórdia ejus.

Qui dat escam omni carni: * quóniam in ætérnum misericórdia ejus.

Confitémini Deo cæli: * quóniam in ætérnum misericórdia ejus.

Confitémini Dómino dominórum: * quóniam in ætérnum misericórdia ejus.

\antlabel{Ant.} Confitémini Dómino, quia in humilitáte nostra memor fuit nostri.

\antlabel{Ant.} Adhǽreat lingua mea * fáucibus meis, si non memínero tui, Jerúsalem.

\psalmheading{Psalmus 136}

\lettrine[lines=2]{S}{UPER} flúmina Babylónis, illic sédimus et flévimus: * cum recordarémur Sion:

In salícibus in médio ejus, * suspéndimus órgana nostra.

Quia illic interrogavérunt nos, qui captívos duxérunt nos, * verba cantiónum:

Et qui abduxérunt nos: * Hymnum cantáte nobis de cánticis Sion.

Quómodo cantábimus cánticum Dómini: * in terra aliéna?

Si oblítus fúero tui, Jerúsalem, * oblivióni detur déxtera mea.

Adhǽreat lingua mea fáucibus meis, * si non memínero tui:

Si non proposúero Jerúsalem, * in princípio lætítiæ meæ.

Memor esto, Dómine, filiórum Edom, * in die Jerúsalem:

Qui dicunt: Exinaníte, exinaníte: * usque ad fundaméntum in ea.

Fília Babylónis mísera: * beátus, qui retríbuet tibi retributiónem tuam, quam retribuísti nobis.

Beátus, qui tenébit, * et allídet párvulos tuos ad petram.

\antlabel{Ant.} Adhǽreat lingua mea fáucibus meis, si non memínero tui, Jerúsalem.

\antlabel{Ant.} Confitébor * \textdagger\ nómini tuo, Dómine, super misericórdia et veritáte tua.

\psalmheading{Psalmus 137}

\lettrine[lines=2]{C}{ONFITEBOR} \textdagger\ tibi, Dómine, in toto corde meo: * quóniam audísti verba oris mei.

In conspéctu Angelórum psallam tibi: * adorábo ad templum sanctum tuum, et confitébor nómini tuo.

Super misericórdia tua, et veritáte tua: * quóniam magnificásti super omne, nomen sanctum tuum.

In quacúmque die invocávero te, exáudi me: * multiplicábis in ánima mea virtútem.

Confiteántur tibi, Dómine, omnes reges terræ: * quia audiérunt ómnia verba oris tui:

Et cantent in viis Dómini: * quóniam magna est glória Dómini.

Quóniam excélsus Dóminus, et humília réspicit: * et alta a longe cognóscit.

Si ambulávero in médio tribulatiónis, vivificábis me: * et super iram inimicórum meórum extendísti manum tuam, et salvum me fecit déxtera tua.

Dóminus retríbuet pro me: * Dómine, misericórdia tua in sæculum: ópera mánuum tuárum ne despícias.

\antlabel{Ant.} Confitébor nómini tuo, Dómine, super misericórdia et veritáte tua.

\variantheading{In Officio feriali}

\capitulumlabel{Capitulum}{2 Cor. 1, 3-4}

\lettrine[lines=2]{B}{ENEDICTUS} Deus, et Pater Dómini nostri Jesu Christi, Pater misericordiárum, et Deus totíus consolatiónis, qui consolátur nos in omni tribulatióne nostra.

\hymnusheading{Hymnus}

\lettrine[lines=2]{M}{AGNAE} Deus poténtiæ,\\
Qui fértili natos aqua\\
Partim relínquis gúrgiti,\\
Partim levas in áëra.

Demérsa lymphis ímprimens,\\
Subvécta cælis érigens:\\
Ut stirpe ab una pródita,\\
Divérsa répleant loca:

Largíre cunctis sérvulis,\\
Quos mundat unda Sánguinis,\\
Nescíre lapsus críminum,\\
Nec ferre mortis tædium.

Ut culpa nullum déprimat:\\
Nullum éfferat jactántia:\\
Elísa mens ne cóncidat:\\
Eláta mens ne córruat.

Præsta, Pater piíssime,\\
Patríque compar Unice,\\
Cum Spíritu Paráclito\\
Regnans per omne sæculum.\\
Amen.

\V\ Dirigátur, Dómine, orátio mea. \R\ Sicut incénsum in conspéctu tuo.

\antlabel{Ad Magnif. Ant.} Fecit Deus * poténtiam in bráchio suo: dispérsit supérbos mente cordis sui.

\rubric{Réliqua ut in Ordinário. 22}

\hourheading{Ad Completorium}

\V\ Jube, domne.

\benedlabel{Bened.} Noctem quiétam, et finem perféctum concédat nobis Dóminus omnípotens. \R\ Amen.

\lectioheading{Lectio br.}{1 Petri 5, 8-9}

\lettrine[lines=2]{F}{RATRES}: Sóbrii estóte, et vigiláte: quia adversárius vester diábolus tamquam leo rúgiens círcuit, quærens quem dévoret: cui resístite fortes in fide. Tu autem, Dómine, miserére nobis. \R\ Deo grátias.

\V\ Adjutórium nostrum in nómine Dómini. \R\ Qui fecit cælum et terram.

Pater noster \rubric{quod dícitur totum secréto.}

\formulalabel{Confíteor.} Misereátur. Indulgéntiam.

\V\ Convérte nos, Deus, salutáris noster. \R\ Et avérte iram tuam a nobis.

\V\ Deus, in adjutórium.

\antlabel{Ant.} Adjútor meus.

\psalmheading{Psalmus 69}

\lettrine[lines=2]{D}{EUS}, in adjutórium meum inténde: * Dómine, ad adjuvándum me festína.

Confundántur et reverántur, * qui quærunt ánimam meam.

Avertántur retrórsum, et erubéscant, * qui volunt mihi mala.

Avertántur statim erubescéntes, * qui dicunt mihi: Euge, euge.

Exsúltent et læténtur in te omnes qui quærunt te, * et dicant semper: Magnificétur Dóminus: qui díligunt salutáre tuum.

Ego vero egénus, et pauper sum: * Deus, ádjuva me.

Adjútor meus, et liberátor meus es tu: * Dómine, ne moréris.

\psalmheading{Psalmus 70, i}

\lettrine[lines=2]{I}{N} te, Dómine, sperávi, non confúndar in ætérnum: * in justítia tua líbera me, et éripe me.

Inclína ad me aurem tuam, * et salva me.

Esto mihi in Deum protectórem, et in locum munítum: * ut salvum me fácias.

Quóniam firmaméntum meum, * et refúgium meum es tu.

Deus meus, éripe me de manu peccatóris, * et de manu contra legem agéntis et iníqui:

Quóniam tu es patiéntia mea, Dómine: * Dómine, spes mea a juventúte mea.

In te confirmátus sum ex útero: * de ventre matris meæ tu es protéctor meus.

In te cantátio mea semper: * tamquam prodígium factus sum multis: et tu adjútor fortis.

Repleátur os meum laude, ut cantem glóriam tuam: * tota die magnitúdinem tuam.

Ne projícias me in témpore senectútis: * cum defécerit virtus mea, ne derelínquas me.

Quia dixérunt inimíci mei mihi: * et qui custodiébant ánimam meam, consílium fecérunt in unum,

Dicéntes: Deus dereliquit eum, persequímini, et comprehéndite eum: * quia non est qui erípiat.

Deus, ne elongéris a me: * Deus meus, in auxílium meum réspice.

\psalmheading{Psalmus 70, ii}

\lettrine[lines=2]{C}{ONFUNDANTUR}, et defíciant detrahéntes ánimæ meæ: * operiántur confusióne, et pudóre qui quærunt mala mihi.

Ego autem semper sperábo: * et adjíciam super omnem laudem tuam.

Os meum annuntiábit justítiam tuam: * tota die salutáre tuum.

Quóniam non cognóvi litteratúram, introíbo in poténtias Dómini: * Dómine, memorábor justítiæ tuæ solíus.

Deus, docuísti me a juventúte mea: * et usque nunc pronuntiábo mirabília tua.

Et usque in senéctam et sénium: * Deus, ne derelínquas me.

Donec annúntiem bráchium tuum * generatióni omni, quæ ventúra est:

Poténtiam tuam, et justítiam tuam, Deus, usque in altíssima, quæ fecísti magnália: * Deus, quis símilis tibi?

Quantas ostendísti mihi tribulatiónes multas et malas: et convérsus vivificásti me: * et de abýssis terræ íterum reduxísti me:

Multiplicásti magnificéntiam tuam: * et convérsus consolátus es me.

Nam et ego confitébor tibi in vasis psalmi veritátem tuam: * Deus, psallam tibi in cíthara, Sanctus Israël.

Exsultábunt lábia mea cum cantávero tibi: * et ánima mea, quam redemísti.

Sed et lingua mea tota die meditábitur justítiam tuam: * cum confúsi et reveríti fúerint, qui quærunt mala mihi.

\antlabel{Ant.} Adjútor meus, et liberátor meus esto, Dómine.

\hymnusheading{Hymnus}

\lettrine[lines=2]{T}{E} lucis ante términum,\\
Rerum Creátor, póscimus,\\
Ut pro tua cleméntia\\
Sis præsul et custódia.

Procul recédant sómnia,\\
Et nóctium phantásmata,\\
Hostémque nostrum cómprime,\\
Ne polluántur córpora.

Præsta, Pater piíssime,\\
Patríque compar Unice,\\
Cum Spíritu Paráclito\\
Regnans per omne sæculum.\\
Amen.

\capitulumlabel{Capitulum}{Jer. 14, 9}

\lettrine[lines=2]{T}{U} autem in nobis es, Dómine, et nomen sanctum tuum invocátum est super nos: ne derelínquas nos, Dómine, Deus noster.

\R\ br. In manus tuas, Dómine, * Comméndo spíritum meum. In manus. \V\ Redemísti nos, Dómine, Deus veritátis. Comméndo. Glória Patri. In manus.

\V\ Custódi nos, Dómine, ut pupíllam óculi. \R\ Sub umbra alárum tuárum prótege nos.

\antlabel{Ant.} Salva nos.

\canticumheading{Canticum Simeonis}{Luc. 2, 29-32}

\lettrine[lines=2]{N}{UNC} dimíttis servum tuum, Dómine, * secúndum verbum tuum in pace:

Quia vidérunt óculi mei * salutáre tuum,

Quod parásti * ante fáciem ómnium populórum,

Lumen ad revelatiónem géntium, * et glóriam plebis tuæ Israël.

Glória Patri.

\antlabel{Ant.} Salva nos, Dómine, vigilántes, custódi nos dormiéntes; ut vigilémus cum Christo, et requiescámus in pace.

\variantheading{Preces (si dicendæ sint)}

KYRIE, eléison. Christe, eléison. Kýrie, eléison.

Pater noster \rubric{secreto usque ad} \V\ Et ne nos indúcas in tentatiónem. \R\ Sed líbera nos a malo.

Credo in Deum \rubric{secreto usque ad} \V\ Carnis resurrectiónem. \R\ Vitam ætérnam. Amen. \V\ Benedíctus es, Dómine, Deus patrum nostrórum. \R\ Et laudábilis et gloriósus in sæcula. \V\ Benedicámus Patrem et Fílium cum Sancto Spíritu. \R\ Laudémus, et superexaltémus eum in sæcula. \V\ Benedíctus es, Dómine, in firmaménto cæli. \R\ Et laudábilis, et gloriósus, et superexaltátus in sæcula. \V\ Benedícat et custódiat nos omnípotens et miséricors Dóminus. \R\ Amen. \V\ Dignáre, Dómine, nocte ista. \R\ Sine peccáto nos custodíre. \V\ Miserére nostri, Dómine. \R\ Miserére nostri. \V\ Fiat misericórdia tua, Dómine, super nos. \R\ Quemádmodum sperávimus in te. \V\ Dómine, exáudi oratiónem meam. \R\ Et clamor meus ad te véniat.

\V\ Dóminus vobíscum.

Orémus. \oratiolabel{Oratio}

\lettrine[lines=2]{V}{ISITA}, quæsumus, Dómine, habitatiónem istam, et omnes insídias inimíci ab ea longe repélle: Angeli tui sancti hábitent in ea, qui nos in pace custódiant; et benedíctio tua sit super nos semper. Per Dóminum nostrum. \R\ Amen.

\V\ Dóminus vobíscum.

\V\ Benedicámus Dómino.

\benedictiolabel{Benedictio} Benedícat et custódiat nos omnípotens et miséricors Dóminus, Pater, et Fílius, et Spíritus Sanctus. \R\ Amen.

\variantheading{Antiphona finalis beátæ Maríæ Vírginis.}

\V\ Divínum auxílium.

Pater noster, Ave María et Credo.

% ===== END OF BATCH 29 -- Covers autumnalis.md lines 9541-9888,
% completing "Ad Vesperas" (Psalms 132, 135/i, 135/ii, 136, 137,
% hymn "Magnæ Deus poténtiæ") and "Ad Completorium" (Psalms 69,
% 70/i, 70/ii, hymn "Te lucis ante términum," Canticum Simeonis,
% Preces), and with them the ENTIRE "Feria Quinta" Office. The
% closing cross-reference at the end of Compline is given in its
% fuller unabbreviated form here -- "Antiphona finalis beátæ Maríæ
% Vírginis." and "Pater noster, Ave María et Credo." -- rather than
% the abbreviated "Antiphona finalis B. M. V." and "Pater, Ave et
% Credo." seen at the end of every previous feria's Compline this
% volume; transcribed exactly as the source gives each. The
% document continues with "FERIA SEXTA" (autumnalis.md line 9891),
% printed in the source in full capitals as a heading (transcribed
% with \feriaheading{} regardless, consistent with the role-based
% convention already applied to every other day-heading in this
% volume). The Invitatorium at Friday Matins carries a bracketed
% footnote marker "[3]" after its incipit in the source, preserved
% verbatim as part of the transcribed text.
% =====
\feriaheading{Feria Sexta}

\hourheading{Ad Matutinum}

Pater, Ave et Credo.

\V\ Dómine, lábia mea.

\V\ Deus, in adjutórium.

\variantheading{In Officio feriali}

\invitatlabel{Invit.} Dóminum, Deum nostrum, * Veníte, adorémus. \psalmlabel{Ps. 94} Veníte. [3]

\hymnusheading{Hymnus}

\lettrine[lines=2]{T}{U}, Trinitátis Únitas,\\
Orbem poténter quæ regis,\\
Atténde laudis cánticum,\\
Quod excubántes psállimus.

Nam léctulo consúrgimus\\
Noctis quiéto témpore,\\
Ut flagitémus ómnium\\
A te medélam vúlnerum.

Quo fraude quidquid dæmonum\\
In nóctibus delíquimus,\\
Abstérgat illud cælitus\\
Tuæ potéstas glóriæ.

Ne corpus astet sórdidum,\\
Nec torpor instet córdium,\\
Ne críminis contágio\\
Tepéscat ardor spíritus.

Ob hoc, Redémptor, quǽsumus,\\
Reple tuo nos lúmine,\\
Per quod diérum círculis\\
Nullis ruámus áctibus.

Præsta, Pater piíssime,\\
Patríque compar Unice,\\
Cum Spíritu Paráclito\\
Regnans per omne sæculum.\\
Amen.

\nocturnoheading{In I Nocturno}

\antlabel{Ant.} Suscitávit Dóminus * testimónium in Jacob: et legem pósuit in Israël.

\psalmheading{Psalmus 77, i}

\lettrine[lines=2]{A}{TTENDITE}, pópule meus, legem meam: * inclináte aurem vestram in verba oris mei.

Apériam in parábolis os meum: * loquar propositiónes ab inítio.

Quanta audívimus et cognóvimus ea: * et patres nostri narravérunt nobis.

Non sunt occultáta a fíliis eórum, * in generatióne áltera.

Narrántes laudes Dómini, et virtútes ejus, * et mirabília ejus, quæ fecit.

Et suscitávit testimónium in Jacob: * et legem pósuit in Israël.

Quanta mandávit pátribus nostris nota fácere ea fíliis suis: * ut cognóscat generátio áltera.

Fílii qui nascéntur, et exsúrgent, * et narrábunt fíliis suis,

Ut ponant in Deo spem suam, et non obliviscántur óperum Dei: * et mandáta ejus exquírant.

Ne fiant sicut patres eórum: * generátio prava et exásperans.

Generátio, quæ non diréxit cor suum: * et non est créditus cum Deo spíritus ejus.

\antlabel{Ant.} Suscitávit Dóminus testimónium in Jacob: et legem pósuit in Israël.

\antlabel{Ant.} Coram pátribus eórum * fecit Deus mirabília.

\psalmheading{Psalmus 77, ii}

\lettrine[lines=2]{F}{ILII} Ephrem intendéntes et mitténtes arcum: * convérsi sunt in die belli.

Non custodiérunt testaméntum Dei: * et in lege ejus noluérunt ambuláre.

Et oblíti sunt benefactórum ejus, * et mirabílium ejus quæ osténdit eis.

Coram pátribus eórum fecit mirabília in terra Ægýpti, * in campo Táneos.

Interrúpit mare, et perdúxit eos: * et státuit aquas quasi in utre.

Et dedúxit eos in nube diéi: * et tota nocte in illuminatióne ignis.

Interrúpit petram in erémo: * et adaquávit eos velut in abýsso multa.

Et edúxit aquam de petra: * et dedúxit tamquam flúmina aquas.

\antlabel{Ant.} Coram pátribus eórum fecit Deus mirabília.

\antlabel{Ant.} Jánuas cæli apéruit * Dóminus, et pluit illis manna ad manducándum.

\psalmheading{Psalmus 77, iii}

\lettrine[lines=2]{E}{T} apposuérunt adhuc peccáre ei: * in iram excitavérunt Excélsum in inaquóso.

Et tentavérunt Deum in córdibus suis, * ut péterent escas animábus suis.

Et male locúti sunt de Deo: * dixérunt: Numquid póterit Deus paráre mensam in desérto?

Quóniam percússit petram, et fluxérunt aquæ, * et torréntes inundavérunt.

Numquid et panem póterit dare, * aut paráre mensam pópulo suo?

Ideo audívit Dóminus, et dístulit: * et ignis accénsus est in Jacob, et ira ascéndit in Israël:

Quia non credidérunt in Deo, * nec speravérunt in salutári ejus:

Et mandávit núbibus désuper, * et jánuas cæli apéruit.

Et pluit illis manna ad manducándum, * et panem cæli dedit eis.

Panem Angelórum manducávit homo, * cibária misit eis in abundántia.

Tránstulit Austrum de cælo: * et indúxit in virtúte sua Africum.

Et pluit super eos sicut púlverem carnes: * et sicut arénam maris volatília pennáta.

Et cecidérunt in médio castrórum eórum: * circa tabernácula eórum.

Et manducavérunt, et saturáti sunt nimis, et desidérium eórum áttulit eis: * non sunt fraudáti a desidério suo.

Adhuc escæ eórum erant in ore ipsórum: * et ira Dei ascéndit super eos.

Et occídit pingues eórum, * et eléctos Israël impedívit.

\antlabel{Ant.} Jánuas cæli apéruit Dóminus, et pluit illis manna ad manducándum.

\variantheading{In Festis novem Lectionum:}

\V\ Deus percússit petram, et fluxérunt aquæ. \R\ Et torréntes inundavérunt.

\absolutiolabel{Absolutio.} Exáudi, Dómine.

\benedlabel{Bened.} 1 Benedictióne perpétua. 2 Unigénitus Dei Fílius. 3 Spíritus Sancti grátia.

\nocturnoheading{In II Nocturno}

\antlabel{Ant.} Deus adjútor * est eórum: et Excélsus redémptor eórum est.

\psalmheading{Psalmus 77, iv}

\lettrine[lines=2]{I}{N} ómnibus his peccavérunt adhuc: * et non credidérunt in mirabílibus ejus.

Et defecérunt in vanitáte dies eórum: * et anni eórum cum festinatióne.

Cum occíderet eos, quærébant eum: * et revertebántur, et dilúculo veniébant ad eum.

Et rememoráti sunt quia Deus adjútor est eórum: * et Deus excélsus redémptor eórum est.

Et dilexérunt eum in ore suo, * et lingua sua mentíti sunt ei:

Cor autem eórum non erat rectum cum eo: * nec fidéles habíti sunt in testaménto ejus.

Ipse autem est miséricors, et propítius fiet peccátis eórum: * et non dispérdet eos.

Et abundávit ut avérteret iram suam: * et non accéndit omnem iram suam:

Et recordátus est quia caro sunt: * spíritus vadens, et non rédiens.

Quóties exacerbavérunt eum in desérto, * in iram concitavérunt eum in inaquóso?

Et convérsi sunt, et tentavérunt Deum: * et Sanctum Israël exacerbavérunt.

\antlabel{Ant.} Deus adjútor est eórum: et Excélsus redémptor eórum est.

\antlabel{Ant.} Redémit eos * Dóminus de manu tribulántis.

\psalmheading{Psalmus 77, v}

\lettrine[lines=2]{N}{ON} sunt recordáti manus ejus, * die qua redémit eos de manu tribulántis.

Sicut pósuit in Ægýpto signa sua, * et prodígia sua in campo Táneos.

Et convértit in sánguinem flúmina eórum, * et imbres eórum, ne bíberent.

Misit in eos cœnomyíam, et comédit eos: * et ranam, et dispérdidit eos.

Et dedit ærúgini fructus eórum, * et labóres eórum locústæ.

Et occídit in grándine víneas eórum, * et moros eórum in prúina.

Et trádidit grándini juménta eórum: * et possessiónem eórum igni.

Misit in eos iram indignatiónis suæ: * indignatiónem, et iram, et tribulatiónem: immissiónes per ángelos malos.

Viam fecit sémitæ iræ suæ, non pepércit a morte animábus eórum: * et juménta eórum in morte conclúsit.

Et percússit omne primogénitum in terra Ægýpti: * primítias omnis labóris eórum in tabernáculis Cham.

Et ábstulit sicut oves pópulum suum: * et perdúxit eos tamquam gregem in desérto.

Et dedúxit eos in spe, et non timuérunt: * et inimícos eórum opéruit mare.

Et indúxit eos in montem sanctificatiónis suæ, * montem, quem acquisívit déxtera ejus.

Et ejécit a fácie eórum gentes: * et sorte divísit eis terram in funículo distributiónis.

Et habitáre fecit in tabernáculis eórum * tribus Israël.

Et tentavérunt, et exacerbavérunt Deum excélsum: * et testimónia ejus non custodiérunt.

Et avertérunt se, et non servavérunt pactum: * quemádmodum patres eórum, convérsi sunt in arcum pravum.

In iram concitavérunt eum in cóllibus suis: * et in sculptílibus suis ad æmulatiónem eum provocavérunt.

\antlabel{Ant.} Redémit eos Dóminus de manu tribulántis.

\antlabel{Ant.} Ædificávit * Deus sanctifícium suum in terra.

\psalmheading{Psalmus 77, vi}

\lettrine[lines=2]{A}{UDIVIT} Deus, et sprevit: * et ad níhilum redégit valde Israël.

Et répulit tabernáculum Silo, * tabernáculum suum, ubi habitávit in homínibus.

Et trádidit in captivitátem virtútem eórum: * et pulchritúdinem eórum in manus inimíci.

Et conclúsit in gládio pópulum suum: * et hereditátem suam sprevit.

Júvenes eórum comédit ignis: * et vírgines eórum non sunt lamentátæ.

Sacerdótes eórum in gládio cecidérunt: * et víduæ eórum non plorabántur.

Et excitátus est tamquam dórmiens Dóminus, * tamquam potens crapulátus a vino.

Et percússit inimícos suos in posterióra: * oppróbrium sempitérnum dedit illis.

Et répulit tabernáculum Joseph: * et tribum Ephraim non elégit:

Sed elégit tribum Juda, * montem Sion quem diléxit.

Et ædificávit sicut unicórnium sanctifícium suum in terra, * quam fundávit in sæcula.

Et elégit David, servum suum, et sústulit eum de grégibus óvium: * de post fetántes accépit eum.

Páscere Jacob, servum suum, * et Israël, hereditátem suam:

Et pavit eos in innocéntia cordis sui: * et in intelléctibus mánuum suárum dedúxit eos.

\antlabel{Ant.} Ædificávit Deus sanctifícium suum in terra.

\variantheading{In Festis novem Lectionum:}

\V\ Elégit Dóminus tribum Juda. \R\ Montem Sion quem diléxit.

\absolutiolabel{Absolutio.} Ipsíus píetas et.

\benedlabel{Bened.} 4 Deus Pater omnípotens. 5 Christus perpétuæ. 6 Ignem sui amóris.

\nocturnoheading{In III Nocturno}

\antlabel{Ant.} Adjuva nos, * Deus, salutáris noster: et propítius esto peccátis nostris.

\psalmheading{Psalmus 78}

\lettrine[lines=2]{D}{EUS}, venérunt gentes in hereditátem tuam, polluérunt templum sanctum tuum: * posuérunt Jerúsalem in pomórum custódiam.

Posuérunt morticína servórum tuórum, escas volatílibus cæli: * carnes sanctórum tuórum béstiis terræ.

Effudérunt sánguinem eórum tamquam aquam in circúitu Jerúsalem: * et non erat qui sepelíret.

Facti sumus oppróbrium vicínis nostris: * subsannátio et illúsio his, qui in circúitu nostro sunt.

Usquequo, Dómine, irascéris in finem: * accendétur velut ignis zelus tuus?

Effúnde iram tuam in gentes, quæ te non novérunt: * et in regna, quæ nomen tuum non invocavérunt:

Quia comedérunt Jacob: * et locum ejus desolavérunt.

Ne memíneris iniquitátum nostrárum antiquárum, cito antícipent nos misericórdiæ tuæ: * quia páuperes facti sumus nimis.

Adjuva nos, Deus, salutáris noster: et propter glóriam nóminis tui, Dómine, líbera nos: * et propítius esto peccátis nostris, propter nomen tuum:

Ne forte dicant in géntibus: Ubi est Deus eórum? * et innotéscat in natiónibus coram óculis nostris.

Ultio sánguinis servórum tuórum, qui effúsus est: * intróeat in conspéctu tuo gémitus compeditórum.

Secúndum magnitúdinem bráchii tui, * pósside fílios mortificatórum.

Et redde vicínis nostris séptuplum in sinu eórum: * impropérium ipsórum, quod exprobravérunt tibi, Dómine.

Nos autem pópulus tuus, et oves páscuæ tuæ, * confitébimur tibi in sæculum.

In generatiónem et generatiónem * annuntiábimus laudem tuam.

\antlabel{Ant.} Adjuva nos, Deus, salutáris noster: et propítius esto peccátis nostris.

\antlabel{Ant.} Ego sum Dóminus, * Deus tuus, Israël, qui edúxi te de terra Ægýpti.

\psalmheading{Psalmus 80}

\lettrine[lines=2]{E}{XSULTATE} Deo, adjutóri nostro: * jubiláte Deo Jacob.

Súmite psalmum, et date týmpanum: * psaltérium jucúndum cum cíthara.

Buccináte in Noeménia tuba, * in insígni die solemnitátis vestræ.

Quia præcéptum in Israël est: * et judícium Deo Jacob.

Testimónium in Joseph pósuit illud, cum exíret de terra Ægýpti: * linguam, quam non nóverat, audívit.

Divértit ab onéribus dorsum ejus: * manus ejus in cóphino serviérunt.

In tribulatióne invocásti me, et liberávi te: * exaudívi te in abscóndito tempestátis: probávi te apud aquam contradictiónis.

Audi, pópulus meus, et contestábor te: * Israël, si audíeris me, non erit in te deus recens, neque adorábis deum aliénum.

Ego enim sum Dóminus, Deus tuus, qui edúxi te de terra Ægýpti: * diláta os tuum, et implébo illud.

Et non audívit pópulus meus vocem meam: * et Israël non inténdit mihi.

Et dimísi eos secúndum desidéria cordis eórum, * ibunt in adinventiónibus suis.

Si pópulus meus audísset me: * Israël si in viis meis ambulásset:

Pro níhilo fórsitan inimícos eórum humiliássem: * et super tribulántes eos misíssem manum meam.

Inimíci Dómini mentíti sunt ei: * et erit tempus eórum in sæcula.

Et cibávit eos ex ádipe fruménti: * et de petra, melle saturávit eos.

\antlabel{Ant.} Ego sum Dóminus, Deus tuus, Israël, qui edúxi te de terra Ægýpti.

\antlabel{Ant.} Ne táceas, Deus, * quóniam inimíci tui extulérunt caput.

\psalmheading{Psalmus 82}

\lettrine[lines=2]{D}{EUS}, quis símilis erit tibi? * ne táceas, neque compescáris, Deus:

Quóniam ecce inimíci tui sonuérunt: * et qui odérunt te, extulérunt caput.

Super pópulum tuum malignavérunt consílium: * et cogitavérunt advérsus sanctos tuos.

Dixérunt: Veníte, et disperdámus eos de gente: * et non memorétur nomen Israël ultra.

Quóniam cogitavérunt unanímiter: * simul advérsum te testaméntum disposuérunt, tabernácula Idumæórum et Ismahelítæ:

Moab, et Agaréni, Gebal, et Ammon, et Amalec: * alienígenæ cum habitántibus Tyrum.

Etenim Assur venit cum illis: * facti sunt in adjutórium fíliis Lot.

Fac illis sicut Mádian, et Sísaræ: * sicut Jabin in torrénte Cisson.

Disperiérunt in Endor: * facti sunt ut stercus terræ.

Pone príncipes eórum sicut Oreb, et Zeb, * et Zébee, et Sálmana:

Omnes príncipes eórum: * qui dixérunt: Hereditáte possideámus Sanctuárium Dei.

Deus meus, pone illos ut rotam: * et sicut stípulam ante fáciem venti.

Sicut ignis, qui combúrit silvam: * et sicut flamma combúrens montes:

Ita persequéris illos in tempestáte tua: * et in ira tua turbábis eos.

Imple fácies eórum ignomínia: * et quærent nomen tuum, Dómine.

Erubéscant, et conturbéntur in sæculum sæculi: * et confundántur et péreant.

Et cognóscant quia nomen tibi Dóminus: * tu solus Altíssimus in omni terra.

\antlabel{Ant.} Ne táceas, Deus, quóniam inimíci tui extulérunt caput.

\variantheading{In feriali Officio et in Festis:}

\V\ Cognóscant quia nomen tibi Dóminus. \R\ Tu solus Altíssimus super omnem terram.

\variantheading{In Officio novem Lectionum}

\absolutiolabel{Absolutio.} A vínculis.

\benedlabel{Bened.} 7 Evangélica léctio. 8 In Festis Dómini: Divínum auxílium. In Festis B. M. V.: Cujus festum cólimus, ipsa Virgo. In Festis Sanctórum: Cujus (\textit{vel} Quorum \textit{aut} Quarum) festum cólimus. 9 Ad societátem. (Per evangélica dicta.)

\variantheading{In Officio trium Lectionum}

\absolutiolabel{Absolutio.} Ipsíus píetas et.

\variantheading{In Feriis cum Homilia:}

\benedlabel{Bened.} 1 Evangélica. 2 Divínum. 3 Ad societátem.

\variantheading{In Feriis sine Homilia:}

\benedlabel{Bened.} 1 Deus Pater. 2 Christus perpétuæ. 3 Ignem sui amóris.

\variantheading{In Festis Sanctorum:}

\benedlabel{Bened.} 1 Ille nos benedícat, qui sine fine vivit et regnat. \R\ Amen. 2 Cujus (\textit{vel} Quorum \textit{aut} Quarum.) 3 Ad societátem.

% ===== END OF BATCH 30 -- Covers autumnalis.md lines 9891-10290,
% opening "Feria Sexta" (Friday) and completing its "Ad Matutinum"
% in full: the ferial Invitatorium "Dóminum, Deum nostrum" (with a
% bracketed footnote marker "[3]" in the source preserved verbatim
% after the incipit), the hymn "Tu, Trinitátis Únitas," and the
% three Nocturns -- Psalm 77 split across six sections (i-vi) over
% the first two Nocturns, and Psalms 78, 80, 82 in the third --
% with their versicle, Absolutio, and Benedictiones apparatus. The
% day-heading itself appears in the source in full capitals ("FERIA
% SEXTA"), transcribed with \feriaheading{} regardless, consistent
% with the role-based heading convention applied throughout this
% volume. Several more psalm incipits carry the partial-caps OCR
% defect already documented in prior batches ("ATtendite,"
% "FIlii," "ET apposuérunt" is correctly two full words so left
% as-is being a short conjunction, "AUdivit," "DEus" twice,
% "EXsultáte") and are silently normalized to full caps where the
% defect is genuine. The document continues with "Ad Laudes I"
% (autumnalis.md line 10292).
% =====
\hourheading{Ad Laudes I}

\rubric{In Fériis extra Vigílias commúnes occurréntibus, atque in Festis.}

(Pater noster et Ave.)

\V\ Deus, in adjutórium.

\antlabel{Ant.} Exaltáte * Dóminum, Deum nostrum, et adoráte in monte sancto ejus.

\psalmheading{Psalmus 98}

\lettrine[lines=2]{D}{OMINUS} regnávit, irascántur pópuli: * qui sedet super Chérubim, moveátur terra.

Dóminus in Sion magnus: * et excélsus super omnes pópulos.

Confiteántur nómini tuo magno: quóniam terríbile, et sanctum est: * et honor regis judícium díligit.

Tu parásti directiónes: * judícium et justítiam in Jacob tu fecísti.

Exaltáte Dóminum, Deum nostrum, et adoráte scabéllum pedum ejus: * quóniam sanctum est.

Móyses et Aaron in sacerdótibus ejus: * et Sámuel inter eos, qui ínvocant nomen ejus:

Invocábant Dóminum, et ipse exaudiébat eos: * in colúmna nubis loquebátur ad eos.

Custodiébant testimónia ejus, * et præcéptum quod dedit illis.

Dómine, Deus noster, tu exaudiébas eos: * Deus, tu propítius fuísti eis, et ulcíscens in omnes adinventiónes eórum.

Exaltáte Dóminum, Deum nostrum, et adoráte in monte sancto ejus: * quóniam sanctus Dóminus, Deus noster.

\antlabel{Ant.} Exaltáte Dóminum, Deum nostrum, et adoráte in monte sancto ejus.

\antlabel{Ant.} Eripe me * de inimícis meis, Dómine, ad te confúgi.

\psalmheading{Psalmus 142}

\lettrine[lines=2]{D}{OMINE}, exáudi oratiónem meam: áuribus pércipe obsecratiónem meam in veritáte tua: * exáudi me in tua justítia.

Et non intres in judícium cum servo tuo: * quia non justificábitur in conspéctu tuo omnis vivens.

Quia persecútus est inimícus ánimam meam: * humiliávit in terra vitam meam.

Collocávit me in obscúris sicut mórtuos sæculi: * et anxiátus est super me spíritus meus, in me turbátum est cor meum.

Memor fui diérum antiquórum, meditátus sum in ómnibus opéribus tuis: * in factis mánuum tuárum meditábar.

Expándi manus meas ad te: * ánima mea sicut terra sine aqua tibi.

Velóciter exáudi me, Dómine: * defécit spíritus meus.

Non avértas fáciem tuam a me: * et símilis ero descendéntibus in lacum.

Audítam fac mihi mane misericórdiam tuam: * quia in te sperávi.

Notam fac mihi viam, in qua ámbulem: * quia ad te levávi ánimam meam.

Eripe me de inimícis meis, Dómine, ad te confúgi: * doce me fácere voluntátem tuam, quia Deus meus es tu.

Spíritus tuus bonus dedúcet me in terram rectam: * propter nomen tuum, Dómine, vivificábis me, in æquitáte tua.

Edúces de tribulatióne ánimam meam: * et in misericórdia tua dispérdes inimícos meos.

Et perdes omnes, qui tríbulant ánimam meam: * quóniam ego servus tuus sum.

\antlabel{Ant.} Eripe me de inimícis meis, Dómine, ad te confúgi.

\antlabel{Ant.} Benedixísti, * \textdagger\ Dómine, terram tuam: remisísti iniquitátem plebis tuæ.

\psalmheading{Psalmus 84}

\lettrine[lines=2]{B}{ENEDIXISTI}, \textdagger\ Dómine, terram tuam: * avertísti captivitátem Jacob.

Remisísti iniquitátem plebis tuæ: * operuísti ómnia peccáta eórum.

Mitigásti omnem iram tuam: * avertísti ab ira indignatiónis tuæ.

Convérte nos, Deus, salutáris noster: * et avérte iram tuam a nobis.

Numquid in ætérnum irascéris nobis? * aut exténdes iram tuam a generatióne in generatiónem?

Deus, tu convérsus vivificábis nos: * et plebs tua lætábitur in te.

Osténde nobis, Dómine, misericórdiam tuam: * et salutáre tuum da nobis.

Audiam quid loquátur in me Dóminus Deus: * quóniam loquétur pacem in plebem suam.

Et super sanctos suos: * et in eos, qui convertúntur ad cor.

Verúmtamen prope timéntibus eum salutáre ipsíus: * ut inhábitet glória in terra nostra.

Misericórdia, et véritas obviavérunt sibi: * justítia, et pax osculátæ sunt.

Véritas de terra orta est: * et justítia de cælo prospéxit.

Etenim Dóminus dabit benignitátem: * et terra nostra dabit fructum suum.

Justítia ante eum ambulábit: * et ponet in via gressus suos.

\antlabel{Ant.} Benedixísti, Dómine, terram tuam: remisísti iniquitátem plebis tuæ.

\antlabel{Ant.} In Dómino justificábitur, * et laudábitur omne semen Israël.

\canticumheading{Canticum Isaiæ}{Is. 45, 15-26}

\lettrine[lines=2]{V}{ERE} tu es Deus abscónditus, * Deus Israël, Salvátor.

Confúsi sunt, et erubuérunt omnes: * simul abiérunt in confusiónem fabricatóres errórum.

Israël salvátus est in Dómino salúte ætérna: * non confundémini, et non erubescétis usque in sæculum sæculi.

Quia hæc dicit Dóminus creans cælos, * ipse Deus formans terram, et fáciens eam, ipse plastes ejus:

Non in vanum creávit eam, ut habitarétur formávit eam: * Ego Dóminus, et non est álius.

Non in abscóndito locútus sum, * in loco terræ tenebróso:

Non dixi sémini Jacob frustra: Quærite me: * ego Dóminus loquens justítiam, annúntians recta.

Congregámini, et veníte, et accédite simul * qui salváti estis ex géntibus:

Nesciérunt qui levant lignum sculptúræ suæ, * et rogant deum non salvántem.

Annuntiáte, et veníte, et consiliámini simul: * Quis audítum fecit hoc ab inítio, ex tunc prædíxit illud?

Numquid non ego Dóminus, et non est ultra Deus absque me? * Deus justus, et salvans non est præter me.

Convertímini ad me, et salvi éritis, omnes fines terræ: * quia ego Deus, et non est álius.

In memetípso jurávi, egrediétur de ore meo justítiæ verbum, * et non revertétur:

Quia mihi curvábitur omne genu, * et jurábit omnis lingua.

Ergo in Dómino, dicet, meæ sunt justítiæ et impérium: * ad eum vénient, et confundéntur omnes qui repúgnant ei.

In Dómino justificábitur, et laudábitur * omne semen Israël.

\antlabel{Ant.} In Dómino justificábitur, et laudábitur omne semen Israël.

\antlabel{Ant.} Lauda, * \textdagger\ Jerúsalem, Dóminum. \textdaggerdbl

\psalmheading{Psalmus 147}

\lettrine[lines=2]{L}{AUDA}, \textdagger\ Jerúsalem, Dóminum: * \textdaggerdbl\ lauda Deum tuum, Sion.

Quóniam confortávit seras portárum tuárum: * benedíxit fíliis tuis in te.

Qui pósuit fines tuos pacem: * et ádipe fruménti sátiat te.

Qui emíttit elóquium suum terræ: * velóciter currit sermo ejus.

Qui dat nivem sicut lanam: * nébulam sicut cínerem spargit.

Mittit crystállum suam sicut buccéllas: * ante fáciem frígoris ejus quis sustinébit?

Emíttet verbum suum, et liquefáciet ea: * flabit spíritus ejus, et fluent aquæ.

Qui annúntiat verbum suum Jacob: * justítias, et judícia sua Israël.

Non fecit táliter omni natióni: * et judícia sua non manifestávit eis.

\antlabel{Ant.} Lauda, Jerúsalem, Dóminum.

\variantheading{In Officio feriali}

\capitulumlabel{Capit.}{Rom. 13, 12-13}

\lettrine[lines=2]{N}{OX} præcéssit, dies autem appropinquávit. Abjiciámus ergo ópera tenebrárum, et induámur arma lucis. Sicut in die honéste ambulémus.

\hymnusheading{Hymnus}

\lettrine[lines=2]{Æ}{TERNA} cæli glória,\\
Beáta spes mortálium,\\
Summi Tonántis Únice,\\
Castǽque proles Vírginis:

Da déxteram surgéntibus,\\
Exsúrgat et mens sóbria,\\
Flagrans et in laudem Dei\\
Grates rependat débitas.

Ortus refúlget Lúcifer,\\
Præítque solem núntius:\\
Cadunt tenébræ nóctium:\\
Lux sancta nos illúminet.

Manénsque nostris sénsibus,\\
Noctem repéllat sæculi,\\
Omníque fine témporis\\
Purgáta servet péctora.

Quæsíta jam primum fides\\
In corde radíces agat:\\
Secúnda spes congáudeat,\\
Qua major exstat cáritas.

Deo Patri sit glória,\\
Ejúsque soli Fílio,\\
Cum Spíritu Paráclito,\\
Nunc et per omne sæculum.\\
Amen.

\V\ Repléti sumus mane misericórdia tua. \R\ Exsultávimus, et delectáti sumus.

\antlabel{Ad Bened. Ant.} Per víscera misericórdiæ * Dei nostri visitávit nos Oriens ex alto.

\rubric{Réliqua ut in Ordinário. 9}

% ===== END OF BATCH 31 -- Covers autumnalis.md lines 10292-10502,
% "Ad Laudes I" for Feria Sexta in full: Psalms 98, 142, 84, the
% Canticum Isaiæ, and Psalm 147, followed by the Capitulum "Nox
% præcéssit" and the hymn "Æterna cæli glória" (whose opening
% lettrine letter is the ligature "Æ," matching the established
% convention of lettrine-ing exactly the source's first letter).
% Several more psalm incipits carry the partial-caps OCR defect
% ("DOmine, exáudi," "BEnedixísti," "VEre tu es," "LAuda,
% Jerúsalem") and are silently normalized to full caps. The
% document continues with "Ad Laudes II" (autumnalis.md line 10505).
% =====
\hourheading{Ad Laudes II}

\rubric{In Vigíliis commúnibus, quando fit Offícium de Féria.}

(Pater noster et Ave.)

\V\ Deus, in adjutórium.

\antlabel{Ant.} Cor contrítum.

\psalmheading{Psalmus 50}

\lettrine[lines=2]{M}{ISERERE} mei, Deus, * secúndum magnam misericórdiam tuam.

Et secúndum multitúdinem miseratiónum tuárum, * dele iniquitátem meam.

Amplius lava me ab iniquitáte mea: * et a peccáto meo munda me.

Quóniam iniquitátem meam ego cognósco: * et peccátum meum contra me est semper.

Tibi soli peccávi, et malum coram te feci: * ut justificéris in sermónibus tuis, et vincas cum judicáris.

Ecce enim, in iniquitátibus concéptus sum: * et in peccátis concépit me mater mea.

Ecce enim, veritátem dilexísti: * incérta et occúlta sapiéntiæ tuæ manifestásti mihi.

Aspérges me hyssópo, et mundábor: * lavábis me, et super nivem dealbábor.

Audítui meo dabis gáudium et lætítiam: * et exsultábunt ossa humiliáta.

Avérte fáciem tuam a peccátis meis: * et omnes iniquitátes meas dele.

Cor mundum crea in me, Deus: * et spíritum rectum ínnova in viscéribus meis.

Ne projícias me a fácie tua: * et spíritum sanctum tuum ne áuferas a me.

Redde mihi lætítiam salutáris tui: * et spíritu principáli confírma me.

Docébo iníquos vias tuas: * et ímpii ad te converténtur.

Líbera me de sanguínibus, Deus, Deus salútis meæ: * et exsultábit lingua mea justítiam tuam.

Dómine, lábia mea apéries: * et os meum annuntiábit laudem tuam.

Quóniam si voluísses sacrifícium, dedíssem útique: * holocáustis non delectáberis.

Sacrifícium Deo spíritus contribulátus: * cor contrítum, et humiliátum, Deus, non despícies.

Benígne fac, Dómine, in bona voluntáte tua Sion: * ut ædificéntur muri Jerúsalem.

Tunc acceptábis sacrifícium justítiæ, oblatiónes, et holocáusta: * tunc impónent super altáre tuum vítulos.

\antlabel{Ant.} Cor contrítum, et humiliátum, Deus, non despícies.

\antlabel{Ant.} Propter nomen tuum.

\psalmheading{Psalmus 142}

\lettrine[lines=2]{D}{OMINE}, exáudi oratiónem meam: áuribus pércipe obsecratiónem meam in veritáte tua: * exáudi me in tua justítia.

Et non intres in judícium cum servo tuo: * quia non justificábitur in conspéctu tuo omnis vivens.

Quia persecútus est inimícus ánimam meam: * humiliávit in terra vitam meam.

Collocávit me in obscúris sicut mórtuos sæculi: * et anxiátus est super me spíritus meus, in me turbátum est cor meum.

Memor fui diérum antiquórum, meditátus sum in ómnibus opéribus tuis: * in factis mánuum tuárum meditábar.

Expándi manus meas ad te: * ánima mea sicut terra sine aqua tibi.

Velóciter exáudi me, Dómine: * defécit spíritus meus.

Non avértas fáciem tuam a me: * et símilis ero descendéntibus in lacum.

Audítam fac mihi mane misericórdiam tuam: * quia in te sperávi.

Notam fac mihi viam, in qua ámbulem: * quia ad te levávi ánimam meam.

Eripe me de inimícis meis, Dómine, ad te confúgi: * doce me fácere voluntátem tuam, quia Deus meus es tu.

Spíritus tuus bonus dedúcet me in terram rectam: * propter nomen tuum, Dómine, vivificábis me, in æquitáte tua.

Edúces de tribulatióne ánimam meam: * et in misericórdia tua dispérdes inimícos meos.

Et perdes omnes, qui tríbulant ánimam meam: * quóniam ego servus tuus sum.

\antlabel{Ant.} Propter nomen tuum, Dómine, vivificábis me, in æquitáte tua.

\antlabel{Ant.} Deus, tu convérsus.

\psalmheading{Psalmus 84}

\lettrine[lines=2]{B}{ENEDIXISTI}, Dómine, terram tuam: * avertísti captivitátem Jacob.

Remisísti iniquitátem plebis tuæ: * operuísti ómnia peccáta eórum.

Mitigásti omnem iram tuam: * avertísti ab ira indignatiónis tuæ.

Convérte nos, Deus, salutáris noster: * et avérte iram tuam a nobis.

Numquid in ætérnum irascéris nobis? * aut exténdes iram tuam a generatióne in generatiónem?

Deus, tu convérsus vivificábis nos: * et plebs tua lætábitur in te.

Osténde nobis, Dómine, misericórdiam tuam: * et salutáre tuum da nobis.

Audiam quid loquátur in me Dóminus Deus: * quóniam loquétur pacem in plebem suam.

Et super sanctos suos: * et in eos, qui convertúntur ad cor.

Verúmtamen prope timéntibus eum salutáre ipsíus: * ut inhábitet glória in terra nostra.

Misericórdia, et véritas obviavérunt sibi: * justítia, et pax osculátæ sunt.

Véritas de terra orta est: * et justítia de cælo prospéxit.

Etenim Dóminus dabit benignitátem: * et terra nostra dabit fructum suum.

Justítia ante eum ambulábit: * et ponet in via gressus suos.

\antlabel{Ant.} Deus, tu convérsus vivificábis nos: et plebs tua lætábitur in te.

\antlabel{Ant.} Cum irátus fúeris.

\canticumheading{Canticum Habacuc}{Hab. 3, 2-19}

\lettrine[lines=2]{D}{OMINE}, audívi auditiónem tuam, * et tímui.

Dómine, opus tuum, * in médio annórum vivífica illud:

In médio annórum notum fácies: * cum irátus fúeris, misericórdiæ recordáberis.

Deus ab Austro véniet, * et sanctus de monte Pharan:

Opéruit cælos glória ejus: * et laudis ejus plena est terra.

Splendor ejus ut lux erit: * córnua in mánibus ejus:

Ibi abscóndita est fortitúdo ejus: * ante fáciem ejus ibit mors.

Et egrediétur diábolus ante pedes ejus. * Stetit, et mensus est terram.

Aspéxit, et dissólvit gentes: * et contríti sunt montes sæculi.

Incurváti sunt colles mundi, * ab itinéribus æternitátis ejus.

Pro iniquitáte vidi tentória Æthiópiæ, * turbabúntur pelles terræ Mádian.

Numquid in flumínibus irátus es, Dómine? * aut in flumínibus furor tuus? vel in mari indignátio tua?

Qui ascéndes super equos tuos: * et quadrígæ tuæ salvátio.

Súscitans suscitábis arcum tuum: * juraménta tríbubus quæ locútus es.

Flúvios scindes terræ: * vidérunt te, et doluérunt montes: gurges aquárum tránsiit.

Dedit abýssus vocem suam: * altitúdo manus suas levávit.

Sol, et luna stetérunt in habitáculo suo, * in luce sagittárum tuárum, ibunt in splendóre fulgurántis hastæ tuæ.

In frémitu conculcábis terram: * et in furóre obstupefácies gentes.

Egréssus es in salútem pópuli tui: * in salútem cum Christo tuo.

Percussísti caput de domo ímpii: * denudásti fundaméntum ejus usque ad collum.

Maledixísti sceptris ejus, cápiti bellatórum ejus, * veniéntibus ut turbo ad dispergéndum me.

Exsultátio eórum * sicut ejus, qui dévorat páuperem in abscóndito.

Viam fecísti in mari equis tuis, * in luto aquárum multárum.

Audívi, et conturbátus est venter meus: * a voce contremuérunt lábia mea.

Ingrediátur putrédo in óssibus meis, * et subter me scáteat.

Ut requiéscam in die tribulatiónis: * ut ascéndam ad pópulum accínctum nostrum.

Ficus enim non florébit: * et non erit germen in víneis.

Mentiétur opus olívæ: * et arva non áfferent cibum.

Abscindétur de ovíli pecus: * et non erit arméntum in præsépibus.

Ego autem in Dómino gaudébo: * et exsultábo in Deo Jesu meo.

Deus Dóminus fortitúdo mea: * et ponet pedes meos quasi cervórum.

Et super excélsa mea dedúcet me victor * in psalmis canéntem.

\antlabel{Ant.} Cum irátus fúeris, Dómine, misericórdiæ recordáberis.

\antlabel{Ant.} Lauda. \textdagger

\psalmheading{Psalmus 147}

\lettrine[lines=2]{L}{AUDA}, \textdagger\ Jerúsalem, Dóminum: * lauda Deum tuum, Sion.

Quóniam confortávit seras portárum tuárum: * benedíxit fíliis tuis in te.

Qui pósuit fines tuos pacem: * et ádipe fruménti sátiat te.

Qui emíttit elóquium suum terræ: * velóciter currit sermo ejus.

Qui dat nivem sicut lanam: * nébulam sicut cínerem spargit.

Mittit crystállum suam sicut buccéllas: * ante fáciem frígoris ejus quis sustinébit?

Emíttet verbum suum, et liquefáciet ea: * flabit spíritus ejus, et fluent aquæ.

Qui annúntiat verbum suum Jacob: * justítias, et judícia sua Israël.

Non fecit táliter omni natióni: * et judícia sua non manifestávit eis.

\antlabel{Ant.} Lauda Deum tuum, Sion, qui annúntiat judícia sua Israël.

\capitulumlabel{Capit.}{Rom. 13, 12-13}

\lettrine[lines=2]{N}{OX} præcéssit, dies autem appropinquávit. Abjiciámus ergo ópera tenebrárum, et induámur arma lucis. Sicut in die honéste ambulémus.

\hymnusheading{Hymnus}

\lettrine[lines=2]{Æ}{TERNA} cæli glória,\\
Beáta spes mortálium,\\
Summi Tonántis Únice,\\
Castǽque proles Vírginis:

Da déxteram surgéntibus,\\
Exsúrgat et mens sóbria,\\
Flagrans et in laudem Dei\\
Grates rependat débitas.

Ortus refúlget Lúcifer,\\
Præítque solem núntius:\\
Cadunt tenébræ nóctium:\\
Lux sancta nos illúminet.

Manénsque nostris sénsibus,\\
Noctem repéllat sæculi,\\
Omníque fine témporis\\
Purgáta servet péctora.

Quæsíta jam primum fides\\
In corde radíces agat:\\
Secúnda spes congáudeat,\\
Qua major exstat cáritas.

Deo Patri sit glória,\\
Ejúsque soli Fílio,\\
Cum Spíritu Paráclito,\\
Nunc et per omne sæculum.\\
Amen.

\V\ Repléti sumus mane misericórdia tua. \R\ Exsultávimus, et delectáti sumus.

\antlabel{Ad Bened. Ant.} Per víscera misericórdiæ * Dei nostri visitávit nos Oriens ex alto.

\rubric{Expleto Cántico ac repetíta Antíphona, dicúntur sequéntes}

\variantheading{Preces feriales}

KYRIE, eléison. Christe, eléison. Kýrie, eléison.

Pater noster \rubric{quod dícitur a solo Hebdomadário totum clara voce, usque ad} \V\ Et ne nos indúcas in tentatiónem. \R\ Sed líbera nos a malo. \V\ Ego dixi: Dómine, miserére mei. \R\ Sana ánimam meam, quia peccávi tibi. \V\ Convértere, Dómine, úsquequo? \R\ Et deprecábilis esto super servos tuos. \V\ Fiat misericórdia tua, Dómine, super nos. \R\ Quemádmodum sperávimus in te. \V\ Sacerdótes tui induántur justítiam. \R\ Et sancti tui exsúltent. \V\ Orémus pro beatíssimo Papa nostro N. \R\ Dóminus consérvet eum, et vivíficet eum, et beátum fáciat eum in terra, et non tradat eum in ánimam inimicórum ejus. \V\ Orémus et pro Antístite nostro N. \R\ Stet et pascat in fortitúdine tua, Dómine, in sublimitáte nóminis tui. \V\ Dómine, salvum fac regem. \R\ Et exáudi nos in die, qua invocavérimus te. \V\ Salvum fac pópulum tuum, Dómine, et bénedic hereditáti tuæ. \R\ Et rege eos, et extólle illos usque in ætérnum. \V\ Meménto Congregatiónis tuæ. \R\ Quam possedísti ab inítio. \V\ Fiat pax in virtúte tua. \R\ Et abundántia in túrribus tuis. \V\ Orémus pro benefactóribus nostris. \R\ Retribúere dignáre, Dómine, ómnibus, nobis bona faciéntibus propter nomen tuum, vitam ætérnam. Amen. \V\ Orémus pro fidélibus defúnctis. \R\ Réquiem ætérnam dona eis, Dómine, et lux perpétua lúceat eis. \V\ Requiéscant in pace. \R\ Amen. \V\ Pro frátribus nostris abséntibus. \R\ Salvos fac servos tuos, Deus meus, sperántes in te. \V\ Pro afflíctis et captívis. \R\ Líbera eos, Deus Israël, ex ómnibus tribulatiónibus suis. \V\ Mitte eis, Dómine, auxílium de sancto. \R\ Et de Sion tuére eos. \V\ Dómine, Deus virtútum, convérte nos. \R\ Et osténde fáciem tuam, et salvi érimus. \V\ Exsúrge, Christe, ádjuva nos. \R\ Et líbera nos propter nomen tuum. \V\ Dómine, exáudi oratiónem meam. \R\ Et clamor meus ad te véniat.

\rubric{Réliqua ut in Ordinário. 9}

% ===== END OF BATCH 32 -- Covers autumnalis.md lines 10505-10775,
% "Ad Laudes II" for Feria Sexta in full: Psalms 50, 142, 84, the
% Canticum Habacuc ("Dómine, audívi auditiónem tuam"), and Psalm
% 147, sharing the Capitulum "Nox præcéssit" and the hymn "Æterna
% cæli glória" with Ad Laudes I, plus the "Preces feriales" litany.
% Psalms 142 and 84 appear here without the dagger chant-pointing
% mark present in their Ad Laudes I occurrences (Batch 31), a
% genuine source divergence transcribed exactly as given in each
% place. The document continues with "Ad Primam" (autumnalis.md line
% 10778).
% =====
\hourheading{Ad Primam}

Pater, Ave et Credo.

\V\ Deus, in adjutórium.

\hymnusheading{Hymnus}

\lettrine[lines=2]{J}{AM} lucis orto sídere,\\
Deum precémur súpplices,\\
Ut in diúrnis áctibus\\
Nos servet a nocéntibus.

Linguam refrénans témperet,\\
Ne litis horror ínsonet:\\
Visum fovéndo cóntegat,\\
Ne vanitátes háuriat.

Sint pura cordis íntima,\\
Absístat et vecórdia:\\
Carnis terat supérbiam\\
Potus cibíque párcitas.

Ut, cum dies abscésserit\\
Noctémque sors redúxerit,\\
Mundi per abstinéntiam\\
Ipsi canámus glóriam.

Deo Patri sit glória,\\
Ejúsque soli Fílio,\\
Cum Spíritu Paráclito,\\
Nunc et per omne sæculum.\\
Amen.

\antlabel{Ant.} Ne discédas a me.

\psalmheading{Psalmus 21, i}

\lettrine[lines=2]{D}{EUS}, Deus meus, réspice in me: quare me dereliquísti? * longe a salúte mea verba delictórum meórum.

Deus meus, clamábo per diem, et non exáudies: * et nocte, et non ad insipiéntiam mihi.

Tu autem in sancto hábitas, * laus Israël.

In te speravérunt patres nostri: * speravérunt, et liberásti eos.

Ad te clamavérunt, et salvi facti sunt: * in te speravérunt, et non sunt confúsi.

Ego autem sum vermis, et non homo: * oppróbrium hóminum, et abjéctio plebis.

Omnes vidéntes me, derisérunt me: * locúti sunt lábiis, et movérunt caput.

Sperávit in Dómino, éripiat eum: * salvum fáciat eum, quóniam vult eum.

Quóniam tu es, qui extraxísti me de ventre: * spes mea ab ubéribus matris meæ. In te projéctus sum ex útero:

De ventre matris meæ Deus meus es tu, * ne discésseris a me.

Quóniam tribulátio próxima est: * quóniam non est qui ádjuvet.

\psalmheading{Psalmus 21, ii}

\lettrine[lines=2]{C}{IRCUMDEDERUNT} me vítuli multi: * tauri pingues obsedérunt me.

Aperuérunt super me os suum, * sicut leo rápiens et rúgiens.

Sicut aqua effúsus sum: * et dispérsa sunt ómnia ossa mea.

Factum est cor meum tamquam cera liquéscens * in médio ventris mei.

Aruit tamquam testa virtus mea, et lingua mea adhæsit fáucibus meis: * et in púlverem mortis deduxísti me.

Quóniam circumdedérunt me canes multi: * concílium malignántium obsédit me.

Fodérunt manus meas et pedes meos: * dinumeravérunt ómnia ossa mea.

Ipsi vero consideravérunt et inspexérunt me: * divisérunt sibi vestiménta mea, et super vestem meam misérunt sortem.

Tu autem, Dómine, ne elongáveris auxílium tuum a me: * ad defensiónem meam cónspice.

Erue a frámea, Deus, ánimam meam: * et de manu canis únicam meam:

Salva me ex ore leónis: * et a córnibus unicórnium humilitátem meam.

Narrábo nomen tuum frátribus meis: * in médio ecclésiæ laudábo te.

\psalmheading{Psalmus 21, iii}

\lettrine[lines=2]{Q}{UI} timétis Dóminum, laudáte eum: * univérsum semen Jacob, glorificáte eum.

Tímeat eum omne semen Israël: * quóniam non sprevit, neque despéxit deprecatiónem páuperis:

Nec avértit fáciem suam a me: * et cum clamárem ad eum, exaudívit me.

Apud te laus mea in ecclésia magna: * vota mea reddam in conspéctu timéntium eum.

Edent páuperes, et saturabúntur: et laudábunt Dóminum qui requírunt eum: * vivent corda eórum in sæculum sæculi.

Reminiscéntur et converténtur ad Dóminum * univérsi fines terræ:

Et adorábunt in conspéctu ejus * univérsæ famíliæ géntium.

Quóniam Dómini est regnum: * et ipse dominábitur géntium.

Manducavérunt et adoravérunt omnes pingues terræ: * in conspéctu ejus cadent omnes qui descéndunt in terram.

Et ánima mea illi vivet: * et semen meum sérviet ipsi.

Annuntiábitur Dómino generátio ventúra: * et annuntiábunt cæli justítiam ejus pópulo qui nascétur, quem fecit Dóminus.

\rubric{¶ Quando ad Laudes dictus fúerit Ps. 50 Miserére, hic subjúngitur sequens Ps. 98, qui in Laudibus locum non hábuit.}

\psalmheading{Psalmus 98}

\lettrine[lines=2]{D}{OMINUS} regnávit, irascántur pópuli: * qui sedet super Chérubim, moveátur terra.

Dóminus in Sion magnus: * et excélsus super omnes pópulos.

Confiteántur nómini tuo magno: quóniam terríbile, et sanctum est: * et honor regis judícium díligit.

Tu parásti directiónes: * judícium et justítiam in Jacob tu fecísti.

Exaltáte Dóminum, Deum nostrum, et adoráte scabéllum pedum ejus: * quóniam sanctum est.

Móyses et Aaron in sacerdótibus ejus: * et Sámuel inter eos, qui ínvocant nomen ejus:

Invocábant Dóminum, et ipse exaudiébat eos: * in colúmna nubis loquebátur ad eos.

Custodiébant testimónia ejus, * et præcéptum quod dedit illis.

Dómine, Deus noster, tu exaudiébas eos: * Deus, tu propítius fuísti eis, et ulcíscens in omnes adinventiónes eórum.

Exaltáte Dóminum, Deum nostrum, et adoráte in monte sancto ejus: * quóniam sanctus Dóminus, Deus noster.

\rubric{Secus vero, tribus Psalmis explétis, statim dícitur Ant.}

\antlabel{Ant.} Ne discédas a me, Dómine: quóniam tribulátio próxima est: quóniam non est qui ádjuvet.

\rubric{In Officio Festi vel Octávæ:}

\capitulumlabel{Capitulum}{1 Tim. 1, 17}

\lettrine[lines=2]{R}{EGI} sæculórum immortáli et invisíbili, soli Deo honor et glória in sæcula sæculórum. Amen.

\rubric{In ómnibus autem Fériis et in Vigíliis commúnibus:}

\capitulumlabel{Capitulum}{Zach. 8, 19}

\lettrine[lines=2]{P}{ACEM} et veritátem dilígite, ait Dóminus omnípotens.

\R\ br. Christe, Fili Dei vivi, * Miserére nobis. Christe. \V\ \textit{Qui sedes ad déxteram Patris.} Miserére nobis. Glória Patri. Christe. \V\ Exsúrge, Christe, ádjuva nos. \R\ Et líbera nos propter nomen tuum.

\variantheading{Preces dominicales et feriales (si dicendæ sint)}

KYRIE, eléison. Christe, eléison. Kýrie, eléison.

Pater noster \rubric{secreto usque ad} \V\ Et ne nos indúcas in tentatiónem. \R\ Sed líbera nos a malo.

Credo in Deum \rubric{secreto usque ad} \V\ Carnis resurrectiónem. \R\ Vitam ætérnam. Amen. \V\ Et ego ad te, Dómine, clamávi. \R\ Et mane orátio mea prævéniet te. \V\ Repleátur os meum laude. \R\ Ut cantem glóriam tuam, tota die magnitúdinem tuam. \V\ Dómine, avérte fáciem tuam a peccátis meis. \R\ Et omnes iniquitátes meas dele. \V\ Cor mundum crea in me, Deus. \R\ Et spíritum rectum ínnova in viscéribus meis. \V\ Ne projícias me a fácie tua. \R\ Et spíritum sanctum tuum ne áuferas a me. \V\ Redde mihi lætítiam salutáris tui. \R\ Et spíritu principáli confírma me.

\variantheading{In Précibus feriálibus tantum}

\V\ Eripe me, Dómine, ab hómine malo. \R\ A viro iníquo éripe me. \V\ Eripe me de inimícis meis, Deus meus. \R\ Et ab insurgéntibus in me líbera me. \V\ Eripe me de operántibus iniquitátem. \R\ Et de viris sánguinum salva me. \V\ Sic psalmum dicam nómini tuo in sæculum sæculi. \R\ Ut reddam vota mea de die in diem. \V\ Exáudi nos, Deus, salutáris noster. \R\ Spes ómnium fínium terræ, et in mari longe. \V\ Deus, in adjutórium meum inténde. \R\ Dómine, ad adjuvándum me festína. \V\ Sanctus Deus, Sanctus fortis, Sanctus immortális. \R\ Miserére nobis. \V\ Bénedic, ánima mea, Dómino. \R\ Et ómnia, quæ intra me sunt, nómini sancto ejus. \V\ Bénedic, ánima mea, Dómino. \R\ Et noli oblivísci omnes retributiónes ejus. \V\ Qui propitiátur ómnibus iniquitátibus tuis. \R\ Qui sanat omnes infirmitátes tuas. \V\ Qui rédimit de intéritu vitam tuam. \R\ Qui corónat te in misericórdia et miseratiónibus. \V\ Qui replet in bonis desidérium tuum. \R\ Renovábitur ut áquilæ juvéntus tua.

\variantheading{In utrisque Précibus}

\V\ Adjutórium nostrum in nómine Dómini. \R\ Qui fecit cælum et terram.

\formulalabel{Confíteor. Misereátur. Indulgéntiam.}

\V\ Dignáre, Dómine, die isto. \R\ Sine peccáto nos custodíre. \V\ Miserére nostri, Dómine. \R\ Miserére nostri. \V\ Fiat misericórdia tua, Dómine, super nos. \R\ Quemádmodum sperávimus in te. \V\ Dómine, exáudi oratiónem meam. \R\ Et clamor meus ad te véniat.

\V\ Dóminus vobíscum.

Orémus. \oratiolabel{Oratio}

\lettrine[lines=2]{D}{OMINE}, Deus omnípotens, qui ad princípium hujus diéi nos perveníre fecísti: tua nos hódie salva virtúte; ut in hac die ad nullum declinémus peccátum, sed semper ad tuam justítiam faciéndam nostra procédant elóquia, dirigántur cogitatiónes et ópera. Per Dóminum. \R\ Amen.

\V\ Dóminus vobíscum. \V\ Benedicámus Dómino.

\rubric{In Choro Martyrológium.}

\V\ Pretiósa in conspéctu Dómini. \R\ Mors Sanctórum ejus.

\rubric{Et absolúte, sine} Orémus:

\oratiolabel{Oratio}

\lettrine[lines=2]{S}{ANCTA} María, et omnes Sancti intercédant pro nobis ad Dóminum, ut nos mereámur ab eo adjuvári et salvári, qui vivit et regnat in sæcula sæculórum. \R\ Amen.

\V\ Deus, in adjutórium meum inténde. \R\ Dómine, ad adjuvándum me festína. \rubric{Et dícitur ter; último additur:} Glória Patri. Sicut erat.

Kýrie, eléison. Christe, eléison. Kýrie, eléison.

Pater noster \rubric{secreto usque ad} \V\ Et ne nos indúcas in tentatiónem. \R\ Sed líbera nos a malo. \V\ Réspice in servos tuos, Dómine, et in ópera tua, et dírige fílios eórum. \R\ Et sit splendor Dómini, Dei nostri, super nos, et ópera mánuum nostrárum dírige super nos, et opus mánuum nostrárum dírige. \V\ Glória Patri. \R\ Sicut erat.

Orémus. \oratiolabel{Oratio}

\lettrine[lines=2]{D}{IRIGERE} et sanctificáre, régere et gubernáre dignáre, Dómine Deus, Rex cæli et terræ, hódie corda et córpora nostra, sensus, sermónes et actus nostros in lege tua, et in opéribus mandatórum tuórum: ut hic, et in ætérnum, te auxiliánte, salvi et líberi esse mereámur, Salvátor mundi: Qui vivis et regnas in sæcula sæculórum. \R\ Amen.

\V\ Jube, domne.

\benedlabel{Bened.} Dies et actus nostros in sua pace dispónat Dóminus omnípotens. \R\ Amen.

\variantheading{In Officio feriali}

\lettrine[lines=2]{D}{OMINUS} autem dírigat corda et córpora nostra in caritáte Dei, et patiéntia Christi. Tu autem.

\variantheading{E Proprio Sanctorum}

\variantheading{Pro Oct. Omn. Sanctorum}

\lettrine[lines=2]{B}{ENEDÍCTIO} et cláritas et sapiéntia et gratiárum áctio, honor, virtus et fortitúdo Deo nostro in sǽcula sæculórum. Amen.

\variantheading{E Communi Sanctorum}

\variantheading{Pro Apostolis}

\lettrine[lines=2]{I}{BANT} Apóstoli gaudéntes a conspéctu concílii, quóniam digni hábiti sunt pro nómine Jesu contuméliam pati.

\variantheading{Pro Martyre}

\lettrine[lines=2]{J}{USTUS} cor suum trádidit ad vigilándum dilúculo ad Dóminum, qui fecit illum, et in conspéctu Altíssimi deprecábitur.

\variantheading{Pro pluribus Martyribus}

\lettrine[lines=2]{F}{ULGEBUNT} justi, et tamquam scintíllæ in arundinéto discúrrent. Judicábunt natiónes, et dominabúntur pópulis: et regnábit Dóminus illórum in perpétuum.

\variantheading{Pro Conf. Pont.}

\lettrine[lines=2]{F}{UNGI} sacerdótio, et habére laudem in nómine ipsíus, et offérre illi incénsum dignum in odórem suavitátis.

\variantheading{Pro Conf. non Pont.}

\lettrine[lines=2]{J}{USTUM} dedúxit Dóminus per vias rectas, et osténdit illi regnum Dei, et dedit illi sciéntiam sanctórum; honestávit illum in labóribus, et complévit labóres illíus.

\variantheading{Pro Virgine Martyre}

\lettrine[lines=2]{D}{OMINE}, Deus meus, exaltásti super terram habitatiónem meam, et pro morte defluénte deprecáta sum. Invocávi Dóminum, Patrem Dómini mei, ut non derelínquat me in die tribulatiónis meæ, et in témpore superbórum sine adjutório.

\variantheading{Pro Virgine non Martyre}

\lettrine[lines=2]{O}{} \textsc{quam} pulchra est casta generátio cum claritáte! immortális est enim memória illíus: quóniam apud Deum nota est, et apud hómines.

\variantheading{Pro Martyre tantum}

\lettrine[lines=2]{L}{AUDABIT} usque ad mortem ánima mea Dóminum, quóniam éruis sustinéntes te, et líberas eos de manu angústiæ, Dómine, Deus noster.

\variantheading{Pro nec Virg. nec Mart.}

\lettrine[lines=2]{M}{ULTAE} fíliæ congregavérunt divítias: tu supergréssa es univérsas. Fallax grátia, et vana est pulchritúdo: múlier timens Dóminum, ipsa laudábitur.

\variantheading{Pro Dedicatione Ecclesiæ}

\lettrine[lines=2]{E}{T} abstérget Deus omnem lácrimam ab óculis eórum: et mors ultra non erit, neque luctus, neque clamor, neque dolor erit ultra, quia prima abiérunt. Et dixit qui sedébat in throno: Ecce, nova fácio ómnia.

\variantheading{In Officiis B. Mariæ V.}

\lettrine[lines=2]{I}{N} platéis sicut cinnamómum et bálsamum aromatízans odórem dedi: quasi myrrha elécta dedi suavitátem odóris.

\V\ Adjutórium nostrum in nómine Dómini. \R\ Qui fecit cælum et terram.

\V\ Benedícite. \R\ Deus.

\benedictiolabel{Benedictio} Dóminus nos benedícat, et ab omni malo deféndat, et ad vitam perdúcat ætérnam. Et fidélium ánimæ per misericórdiam Dei requiéscant in pace. \R\ Amen.

\rubric{Deinde dícitur tantum} Pater noster \rubric{secréto, nisi sequátur ália Hora.}

% ===== END OF BATCH 33 -- Covers autumnalis.md lines 10778-11047,
% completing "Ad Primam" for Feria Sexta in full: the hymn "Jam
% lucis orto sídere," Psalms 21/i-iii, the conditional Psalm 98
% insertion, both Capitula, the brief Responsory, the full Preces
% apparatus, the closing Orationes, and the Lectio-brevis apparatus
% (identical category list to the previous four ferias). Several
% more psalm incipits carry the partial-caps OCR defect ("DEus,
% Deus meus," "CIrcumdedérunt," "DOminus regnávit," "PAcem"
% [Capitulum]) and are silently normalized to full caps. The
% document continues with "Ad Tertiam" (autumnalis.md line 11050).
% =====
\hourheading{Ad Tertiam}

Pater noster et Ave.

\V\ Deus, in adjutórium.

\hymnusheading{Hymnus}

\lettrine[lines=2]{N}{UNC}, Sancte, nobis, Spíritus,\\
Unum Patri cum Fílio,\\
Dignáre promptus íngeri\\
Nostro refúsus péctori.

Os, lingua, mens, sensus, vigor\\
Confessiónem pérsonent,\\
Flamméscat igne cáritas,\\
Accéndat ardor próximos.

Præsta, Pater piíssime,\\
Patríque compar Unice,\\
Cum Spíritu Paráclito\\
Regnans per omne sæculum.\\
Amen.

\antlabel{Ant.} Excita, Dómine.

\psalmheading{Psalmus 79, i}

\lettrine[lines=2]{Q}{UI} regis Israël, inténde: * qui dedúcis velut ovem Joseph.

Qui sedes super Chérubim, * manifestáre coram Ephraïm, Bénjamin, et Manásse.

Excita poténtiam tuam, et veni, * ut salvos fácias nos.

Deus, convérte nos: * et osténde fáciem tuam, et salvi érimus.

Dómine, Deus virtútum, * quoúsque irascéris super oratiónem servi tui?

Cibábis nos pane lacrimárum: * et potum dabis nobis in lácrimis in ménsura?

Posuísti nos in contradictiónem vicínis nostris: * et inimíci nostri subsannavérunt nos.

Deus virtútum, convérte nos: * et osténde fáciem tuam, et salvi érimus.

\psalmheading{Psalmus 79, ii}

\lettrine[lines=2]{V}{INEAM} de Ægýpto transtulísti: * ejecísti gentes, et plantásti eam.

Dux itíneris fuísti in conspéctu ejus: * plantásti radíces ejus, et implévit terram.

Opéruit montes umbra ejus: * et arbústa ejus cedros Dei.

Exténdit pálmites suos usque ad mare: * et usque ad flumen propágines ejus.

Ut quid destruxísti macériam ejus: * et vindémiant eam omnes, qui prætergrediúntur viam?

Exterminávit eam aper de silva: * et singuláris ferus depástus est eam.

Deus virtútum, convértere: * réspice de cælo, et vide, et vísita víneam istam.

Et pérfice eam, quam plantávit déxtera tua: * et super fílium hóminis, quem confirmásti tibi.

Incénsa igni, et suffóssa * ab increpatióne vultus tui períbunt.

Fiat manus tua super virum déxteræ tuæ: * et super fílium hóminis, quem confirmásti tibi.

Et non discédimus a te, vivificábis nos: * et nomen tuum invocábimus.

Dómine, Deus virtútum, convérte nos: * et osténde fáciem tuam, et salvi érimus.

\psalmheading{Psalmus 81}

\lettrine[lines=2]{D}{EUS} stetit in synagóga deórum: * in médio autem deos dijúdicat.

Usquequo judicátis iniquitátem: * et fácies peccatórum súmitis?

Judicáte egéno, et pupíllo: * húmilem, et páuperem justificáte.

Erípite páuperem: * et egénum de manu peccatóris liberáte.

Nesciérunt, neque intellexérunt, in ténebris ámbulant: * movebúntur ómnia fundaménta terræ.

Ego dixi: Dii estis, * et fílii Excélsi omnes.

Vos autem sicut hómines moriémini: * et sicut unus de princípibus cadétis.

Surge, Deus, júdica terram: * quóniam tu hereditábis in ómnibus géntibus.

\antlabel{Ant.} Excita, Dómine, poténtiam tuam, ut salvos fácias nos.

\variantheading{In Officio feriali}

\capitulumlabel{Capitulum}{Jer. 17, 14}

\lettrine[lines=2]{S}{ANA} me, Dómine, et sanábor: salvum me fac, et salvus ero: quóniam laus mea tu es.

\R\ br. Sana ánimam meam, * Quia peccávi tibi. Sana. \V\ Ego dixi: Dómine, miserére mei. Quia. Glória Patri. Sana.

\V\ Adjútor meus esto, ne derelínquas me. \R\ Neque despícias me, Deus, salutáris meus.

\variantheading{In Vigiliis communibus:}

\variantheading{Preces feriales}

KYRIE, eléison. Christe, eléison. Kýrie, eléison.

Pater noster \rubric{secreto usque ad} \V\ Et ne nos indúcas in tentatiónem. \R\ Sed líbera nos a malo. \V\ Dómine, Deus virtútum, convérte nos. \R\ Et osténde fáciem tuam, et salvi érimus. \V\ Exsúrge, Christe, ádjuva nos. \R\ Et líbera nos propter nomen tuum. \V\ Dómine, exáudi oratiónem meam. \R\ Et clamor meus ad te véniat.

\rubric{Réliqua ut in Ordinário. 19}

\hourheading{Ad Sextam}

Pater noster et Ave.

\V\ Deus, in adjutórium.

\hymnusheading{Hymnus}

\lettrine[lines=2]{R}{ECTOR} potens, verax Deus,\\
Qui témperas rerum vices,\\
Splendóre mane illúminas,\\
Et ígnibus merídiem:

Exstíngue flammas lítium,\\
Aufer calórem nóxium,\\
Confer salútem córporum,\\
Verámque pacem córdium.

Præsta, Pater piíssime,\\
Patríque compar Unice,\\
Cum Spíritu Paráclito\\
Regnans per omne sæculum.\\
Amen.

\antlabel{Ant.} Beáti, qui hábitant.

\psalmheading{Psalmus 83, i}

\lettrine[lines=2]{Q}{UAM} dilécta tabernácula tua, Dómine virtútum: * concupíscit, et déficit ánima mea in átria Dómini.

Cor meum, et caro mea * exsultavérunt in Deum vivum.

Etenim passer invénit sibi domum: * et turtur nidum sibi, ubi ponat pullos suos.

Altária tua, Dómine virtútum: * Rex meus, et Deus meus.

Beáti, qui hábitant in domo tua, Dómine: * in sæcula sæculórum laudábunt te.

Beátus vir, cujus est auxílium abs te: * ascensiónes in corde suo dispósuit in valle lacrimárum in loco, quem pósuit.

Etenim benedictiónem dabit legislátor, ibunt de virtúte in virtútem: * vidébitur Deus deórum in Sion.

\psalmheading{Psalmus 83, ii}

\lettrine[lines=2]{D}{OMINE}, Deus virtútum, exáudi oratiónem meam: * áuribus pércipe, Deus Jacob.

Protéctor noster, áspice, Deus: * et réspice in fáciem Christi tui:

Quia mélior est dies una in átriis tuis, * super míllia.

Elégi abjéctus esse in domo Dei mei: * magis quam habitáre in tabernáculis peccatórum.

Quia misericórdiam, et veritátem díligit Deus: * grátiam et glóriam dabit Dóminus.

Non privábit bonis eos, qui ámbulant in innocéntia: * Dómine virtútum, beátus homo, qui sperat in te.

\psalmheading{Psalmus 86}

\lettrine[lines=2]{F}{UNDAMENTA} ejus in móntibus sanctis: * díligit Dóminus portas Sion super ómnia tabernácula Jacob.

Gloriósa dicta sunt de te, * cívitas Dei.

Memor ero Rahab, et Babylónis * sciéntium me.

Ecce, alienígenæ, et Tyrus, et pópulus Æthíopum, * hi fuérunt illic.

Numquid Sion dicet: Homo, et homo natus est in ea: * et ipse fundávit eam Altíssimus?

Dóminus narrábit in scriptúris populórum, et príncipum: * horum, qui fuérunt in ea.

Sicut lætántium ómnium * habitátio est in te.

\antlabel{Ant.} Beáti, qui hábitant in domo tua, Dómine.

\variantheading{In Officio feriali}

\capitulumlabel{Capitulum}{Rom. 13, 8}

\lettrine[lines=2]{N}{EMINI} quidquam debeátis, nisi ut ínvicem diligátis: qui enim díligit próximum, legem implévit.

\R\ br. Benedícam Dóminum * In omni témpore. Benedícam. \V\ Semper laus ejus in ore meo. In. Glória Patri. Benedícam. \V\ Dóminus regit me, et nihil mihi déerit. \R\ In loco páscuæ ibi me collocávit.

\variantheading{In Vigiliis communibus:}

\variantheading{Preces feriales}

KYRIE, eléison. Christe, eléison. Kýrie, eléison.

Pater noster \rubric{secreto usque ad} \V\ Et ne nos indúcas in tentatiónem. \R\ Sed líbera nos a malo. \V\ Dómine, Deus virtútum, convérte nos. \R\ Et osténde fáciem tuam, et salvi érimus. \V\ Exsúrge, Christe, ádjuva nos. \R\ Et líbera nos propter nomen tuum. \V\ Dómine, exáudi oratiónem meam. \R\ Et clamor meus ad te véniat.

\rubric{Réliqua ut in Ordinário. 20}

\hourheading{Ad Nonam}

Pater noster et Ave.

\V\ Deus, in adjutórium.

\hymnusheading{Hymnus}

\lettrine[lines=2]{R}{ERUM}, Deus, tenax vigor,\\
Immótus in te pérmanens,\\
Lucis diúrnæ témpora\\
Succéssibus detérminans:

Largíre lumen véspere,\\
Quo vita nusquam décidat,\\
Sed prǽmium mortis sacræ\\
Perénnis instet glória.

Præsta, Pater piíssime,\\
Patríque compar Unice,\\
Cum Spíritu Paráclito\\
Regnans per omne sæculum.\\
Amen.

\antlabel{Ant.} Misericórdia et véritas.

\psalmheading{Psalmus 88, i}

\lettrine[lines=2]{M}{ISERICORDIAS} Dómini in ætérnum cantábo.

In generatiónem et generatiónem * annuntiábo veritátem tuam in ore meo.

Quóniam dixísti: In ætérnum misericórdia ædificábitur in cælis: * præparábitur véritas tua in eis.

Dispósui testaméntum eléctis meis, * jurávi David, servo meo:

Usque in ætérnum præparábo semen tuum. * Et ædificábo in generatiónem et generatiónem sedem tuam.

Confitebúntur cæli mirabília tua, Dómine: * étenim veritátem tuam in ecclésia sanctórum.

Quóniam quis in núbibus æquábitur Dómino: * símilis erit Deo in fíliis Dei?

Deus, qui glorificátur in consílio sanctórum: * magnus et terríbilis super omnes qui in circúitu ejus sunt.

Dómine, Deus virtútum, quis símilis tibi? * potens es, Dómine, et véritas tua in circúitu tuo.

Tu domináris potestáti maris: * motum autem flúctuum ejus tu mítigas.

Tu humiliásti sicut vulnerátum, supérbum: * in bráchio virtútis tuæ dispersísti inimícos tuos.

Tui sunt cæli, et tua est terra, orbem terræ et plenitúdinem ejus tu fundásti: * aquilónem, et mare tu creásti.

Thabor et Hermon in nómine tuo exsultábunt: * tuum bráchium cum poténtia.

Firmétur manus tua, et exaltétur déxtera tua: * justítia et judícium præparátio sedis tuæ.

Misericórdia et véritas præcédent fáciem tuam: * beátus pópulus, qui scit jubilatiónem.

Dómine, in lúmine vultus tui ambulábunt, et in nómine tuo exsultábunt tota die: * et in justítia tua exaltabúntur.

Quóniam glória virtútis eórum tu es: * et in beneplácito tuo exaltábitur cornu nostrum.

Quia Dómini est assúmptio nostra, * et Sancti Israël, regis nostri.

\psalmheading{Psalmus 88, ii}

\lettrine[lines=2]{T}{UNC} locútus es in visióne sanctis tuis, et dixísti: * Pósui adjutórium in poténte: et exaltávi eléctum de plebe mea.

Invéni David, servum meum: * óleo sancto meo unxi eum.

Manus enim mea auxiliábitur ei: * et bráchium meum confortábit eum.

Nihil profíciet inimícus in eo, * et fílius iniquitátis non appónet nocére ei.

Et concídam a fácie ipsíus inimícos ejus: * et odiéntes eum in fugam convértam.

Et véritas mea, et misericórdia mea cum ipso: * et in nómine meo exaltábitur cornu ejus.

Et ponam in mari manum ejus: * et in flumínibus déxteram ejus.

Ipse invocábit me: Pater meus es tu: * Deus meus, et suscéptor salútis meæ.

Et ego primogénitum ponam illum * excélsum præ régibus terræ.

In ætérnum servábo illi misericórdiam meam: * et testaméntum meum fidéle ipsi.

Et ponam in sæculum sæculi semen ejus: * et thronum ejus sicut dies cæli.

Si autem derelíquerint fílii ejus legem meam: * et in judíciis meis non ambuláverint:

Si justítias meas profanáverint: * et mandáta mea non custodíerint:

Visitábo in virga iniquitátes eórum: * et in verbéribus peccáta eórum.

Misericórdiam autem meam non dispérgam ab eo: * neque nocébo in veritáte mea:

Neque profanábo testaméntum meum: * et quæ procédunt de lábiis meis, non fáciam írrita.

Semel jurávi in sancto meo: Si David méntiar: * semen ejus in ætérnum manébit.

Et thronus ejus sicut sol in conspéctu meo, * et sicut luna perfécta in ætérnum: et testis in cælo fidélis.

\psalmheading{Psalmus 88, iii}

\lettrine[lines=2]{T}{U} vero repulísti et despexísti: * distulísti Christum tuum.

Evertísti testaméntum servi tui: * profanásti in terra Sanctuárium ejus.

Destruxísti omnes sepes ejus: * posuísti firmaméntum ejus formídinem.

Diripuérunt eum omnes transeúntes viam: * factus est oppróbrium vicínis suis.

Exaltásti déxteram deprimméntium eum: * lætificásti omnes inimícos ejus.

Avertísti adjutórium gládii ejus: * et non es auxiliátus ei in bello.

Destruxísti eum ab emundatióne: * et sedem ejus in terram collisísti.

Minorásti dies témporis ejus: * perfudísti eum confusióne.

Usquequo, Dómine, avértis in finem: * exardéscet sicut ignis ira tua?

Memoráre quæ mea substántia: * numquid enim vane constituísti omnes fílios hóminum?

Quis est homo, qui vivet, et non vidébit mortem: * éruet ánimam suam de manu ínferi?

Ubi sunt misericórdiæ tuæ antíquæ, Dómine, * sicut jurásti David in veritáte tua?

Memor esto, Dómine, oppróbrii servórum tuórum (quod contínui in sinu meo) multárum géntium.

Quod exprobravérunt inimíci tui, Dómine, * quod exprobravérunt commutatiónem Christi tui.

Benedíctus Dóminus in ætérnum: * fiat, fiat.

\antlabel{Ant.} Misericórdia et véritas præcédent fáciem tuam, Dómine.

\variantheading{In Officio feriali}

\capitulumlabel{Capit.}{1 Petri 1, 17-19}

\lettrine[lines=2]{I}{N} timóre incolátus vestri témpore conversámini: sciéntes quod non corruptibílibus auro vel argénto redémpti estis, sed pretióso sánguine quasi Agni immaculáti Christi.

\R\ br. Rédime me, Dómine, * Et miserére mei. Rédime. \V\ Pes enim meus stetit in via recta. Et miserére mei. Glória Patri. Rédime. \V\ Ab occúltis meis munda me, Dómine. \R\ Et ab aliénis parce servo tuo.

\variantheading{In Vigiliis communibus:}

\variantheading{Preces feriales}

KYRIE, eléison. Christe, eléison. Kýrie, eléison.

Pater noster \rubric{secreto usque ad} \V\ Et ne nos indúcas in tentatiónem. \R\ Sed líbera nos a malo. \V\ Dómine, Deus virtútum, convérte nos. \R\ Et osténde fáciem tuam, et salvi érimus. \V\ Exsúrge, Christe, ádjuva nos. \R\ Et líbera nos propter nomen tuum. \V\ Dómine, exáudi oratiónem meam. \R\ Et clamor meus ad te véniat.

\rubric{Réliqua ut in Ordinário. 22}

% ===== END OF BATCH 34 -- Covers autumnalis.md lines 11050-11408,
% completing the Little Hours of Feria Sexta in full: "Ad Tertiam"
% (Psalms 79/i, 79/ii, 81), "Ad Sextam" (Psalms 83/i, 83/ii, 86),
% and "Ad Nonam" (Psalms 88/i-iii), each with its own hymn,
% Capitulum, brief Responsory, and "In Vigiliis communibus: Preces
% feriales" apparatus. Three further occurrences of the recurring
% stray OCR "\*" artifact in the hour-hymns are silently dropped.
% Several more partial-caps OCR defects ("QUi," "VIneam," "DEus
% stetit," "QUam dilécta," "DOmine, Deus virtútum," "FUndaménta,"
% "MIsericórdias," "TUnc locútus," "TU vero") are silently
% normalized to full caps. The document continues with "Ad
% Vesperas" (autumnalis.md line 11412).
% =====
\hourheading{Ad Vesperas}

Pater noster et Ave.

\V\ Deus, in adjutórium.

\antlabel{Ant.} Dómine, * \textdagger\ probásti me, et cognovísti me. \textdaggerdbl

\psalmheading{Psalmus 138, i}

\lettrine[lines=2]{D}{OMINE}, \textdagger\ probásti me, et cognovísti me: * \textdaggerdbl\ tu cognovísti sessiónem meam, et resurrectiónem meam.

Intellexísti cogitatiónes meas de longe: * sémitam meam, et funículum meum investigásti.

Et omnes vias meas prævidísti: * quia non est sermo in lingua mea.

Ecce, Dómine, tu cognovísti ómnia novíssima, et antíqua: * tu formásti me, et posuísti super me manum tuam.

Mirábilis facta est sciéntia tua ex me: * confortáta est, et non pótero ad eam.

Quo ibo a spíritu tuo? * et quo a fácie tua fúgiam?

Si ascéndero in cælum, tu illic es: * si descéndero in inférnum, ades.

Si súmpsero pennas meas dilúculo, * et habitávero in extrémis maris:

Etenim illuc manus tua dedúcet me: * et tenébit me déxtera tua.

Et dixi: Fórsitan ténebræ conculcábunt me: * et nox illuminátio mea in delíciis meis.

Quia ténebræ non obscurabúntur a te, et nox sicut dies illuminábitur: * sicut ténebræ ejus, ita et lumen ejus.

Quia tu possedísti renes meos: * suscepísti me de útero matris meæ.

\antlabel{Ant.} Dómine, probásti me, et cognovísti me.

\antlabel{Ant.} Mirabília ópera tua, * Dómine, et ánima mea cognóscit nimis.

\psalmheading{Psalmus 138, ii}

\lettrine[lines=2]{C}{ONFITEBOR} tibi quia terribíliter magnificátus es: * mirabília ópera tua, et ánima mea cognóscit nimis.

Non est occultátum os meum a te, quod fecísti in occúlto: * et substántia mea in inferióribus terræ.

Imperféctum meum vidérunt óculi tui, et in libro tuo omnes scribéntur: * dies formabúntur, et nemo in eis.

Mihi autem nimis honorificáti sunt amíci tui, Deus: * nimis confortátus est principátus eórum.

Dinumerábo eos, et super arénam multiplicabúntur: * exsurréxi, et adhuc sum tecum.

Si occíderis, Deus, peccatóres: * viri sánguinum, declináte a me:

Quia dícitis in cogitatióne: * Accípient in vanitáte civitátes tuas.

Nonne qui odérunt te, Dómine, óderam? * et super inimícos tuos tabescébam?

Perfécto ódio óderam illos: * et inimíci facti sunt mihi.

Proba me, Deus, et scito cor meum: * intérroga me, et cognósce sémitas meas.

Et vide, si via iniquitátis in me est: * et deduc me in via ætérna.

\antlabel{Ant.} Mirabília ópera tua, Dómine, et ánima mea cognóscit nimis.

\antlabel{Ant.} Ne derelínquas me, * Dómine, virtus salútis meæ.

\psalmheading{Psalmus 139}

\lettrine[lines=2]{E}{RIPE} me, Dómine, ab hómine malo: * a viro iníquo éripe me.

Qui cogitavérunt iniquitátes in corde: * tota die constituébant prǽlia.

Acuérunt línguas suas sicut serpéntis: * venénum áspidum sub lábiis eórum.

Custódi me, Dómine, de manu peccatóris: * et ab homínibus iníquis éripe me.

Qui cogitavérunt supplantáre gressus meos: * absondérunt supérbi láqueum mihi:

Et funes extendérunt in láqueum: * juxta iter scándalum posuérunt mihi.

Dixi Dómino: Deus meus es tu: * exáudi, Dómine, vocem deprecatiónis meæ.

Dómine, Dómine, virtus salútis meæ: * obumbrásti super caput meum in die belli.

Ne tradas me, Dómine, a desidério meo peccatóri: * cogitavérunt contra me, ne derelínquas me, ne forte exalténtur.

Caput circúitus eórum: * labor labiórum ipsórum opériet eos.

Cadent super eos carbónes, in ignem dejícies eos: * in misériis non subsístent.

Vir linguósus non dirigétur in terra: * virum injústum mala cápient in intéritu.

Cognóvi quia fáciet Dóminus judícium ínopis: * et vindíctam páuperum.

Verúmtamen justi confitebúntur nómini tuo: * et habitábunt recti cum vultu tuo.

\antlabel{Ant.} Ne derelínquas me, Dómine, virtus salútis meæ.

\antlabel{Ant.} Dómine, * \textdagger\ clamávi ad te, exáudi me. \textdaggerdbl

\psalmheading{Psalmus 140}

\lettrine[lines=2]{D}{OMINE}, \textdagger\ clamávi ad te, exáudi me: * \textdaggerdbl\ inténde voci meæ, cum clamávero ad te.

Dirigátur orátio mea sicut incénsum in conspéctu tuo: * elevátio mánuum meárum sacrifícium vespertínum.

Pone, Dómine, custódiam ori meo: * et óstium circumstántiæ lábiis meis.

Non declínes cor meum in verba malítiæ, * ad excusándas excusatiónes in peccátis.

Cum homínibus operántibus iniquitátem: * et non communicábo cum eléctis eórum.

Corrípiet me justus in misericórdia, et increpábit me: * óleum autem peccatóris non impínguet caput meum.

Quóniam adhuc et orátio mea in beneplácitis eórum: * absórpti sunt juncti petræ júdices eórum.

Audient verba mea quóniam potuérunt: * sicut crassitúdo terræ erúpta est super terram.

Dissipáta sunt ossa nostra secus inférnum: * quia ad te, Dómine, Dómine, óculi mei: in te sperávi, non áuferas ánimam meam.

Custódi me a láqueo, quem statuérunt mihi: * et a scándalis operántium iniquitátem.

Cadent in retiáculo ejus peccatóres: * singuláriter sum ego donec tránseam.

\antlabel{Ant.} Dómine, clamávi ad te, exáudi me.

\antlabel{Ant.} Educ de custódia * ánimam meam, Dómine, ad confiténdum nómini tuo.

\psalmheading{Psalmus 141}

\lettrine[lines=2]{V}{OCE} mea ad Dóminum clamávi: * voce mea ad Dóminum deprecátus sum:

Effúndo in conspéctu ejus oratiónem meam, * et tribulatiónem meam ante ipsum pronúntio.

In deficiéndo ex me spíritum meum, * et tu cognovísti sémitas meas.

In via hac, qua ambulábam, * abscondérunt láqueum mihi.

Considerábam ad déxteram, et vidébam: * et non erat qui cognósceret me.

Périit fuga a me, * et non est qui requírat ánimam meam.

Clamávi ad te, Dómine, * dixi: Tu es spes mea, pórtio mea in terra vivéntium.

Inténde ad deprecatiónem meam: * quia humiliátus sum nimis.

Líbera me a persequéntibus me: * quia confortáti sunt super me.

Educ de custódia ánimam meam ad confiténdum nómini tuo: * me exspéctant justi, donec retríbuas mihi.

\antlabel{Ant.} Educ de custódia ánimam meam, Dómine, ad confiténdum nómini tuo.

\variantheading{In Officio feriali}

\capitulumlabel{Capitulum}{2 Cor. 1, 3-4}

\lettrine[lines=2]{B}{ENEDICTUS} Deus, et Pater Dómini nostri Jesu Christi, Pater misericordiárum, et Deus totíus consolatiónis, qui consolátur nos in omni tribulatióne nostra.

\hymnusheading{Hymnus}

\lettrine[lines=2]{H}{OMINIS} supérne Cónditor,\\
Qui cuncta solus órdinans,\\
Humum jubes prodúcere\\
Reptántis et feræ genus:

Et magna rerum córpora,\\
Dictu jubéntis vívida,\\
Per témporum certas vices\\
Obtemperáre sérvulis:

Repélle, quod cupídinis\\
Ciénte vi nos ímpetit,\\
Aut móribus se súggerit,\\
Aut áctibus se intérserit.

Da gaudiórum præmia,\\
Da gratiárum múnera:\\
Dissólve litis víncula:\\
Astrínge pacis fœdera.

Præsta, Pater piíssime,\\
Patríque compar Unice,\\
Cum Spíritu Paráclito\\
Regnans per omne sæculum.\\
Amen.

\V\ Dirigátur, Dómine, orátio mea. \R\ Sicut incénsum in conspéctu tuo.

\antlabel{Ad Magnif. Ant.} Depósuit Dóminus * poténtes de sede, et exaltávit húmiles.

\rubric{Réliqua ut in Ordinário. 22}

\hourheading{Ad Completorium}

\V\ Jube, domne.

\benedlabel{Bened.} Noctem quiétam, et finem perféctum concédat nobis Dóminus omnípotens. \R\ Amen.

\lectioheading{Lectio br.}{1 Petri 5, 8-9}

\lettrine[lines=2]{F}{RATRES}: Sóbrii estóte, et vigiláte: quia adversárius vester diábolus tamquam leo rúgiens círcuit, quærens quem dévoret: cui resístite fortes in fide. Tu autem, Dómine, miserére nobis. \R\ Deo grátias.

\V\ Adjutórium nostrum in nómine Dómini. \R\ Qui fecit cælum et terram.

Pater noster \rubric{quod dícitur totum secréto.}

\formulalabel{Confíteor.} Misereátur. Indulgéntiam.

\V\ Convérte nos, Deus, salutáris noster. \R\ Et avérte iram tuam a nobis.

\V\ Deus, in adjutórium.

\antlabel{Ant.} Voce mea. \textdagger

\psalmheading{Psalmus 76, i}

\lettrine[lines=2]{V}{OCE} mea \textdagger\ ad Dóminum clamávi: * voce mea ad Deum, et inténdit mihi.

In die tribulatiónis meæ Deum exquisívi, mánibus meis nocte contra eum: * et non sum decéptus.

Rénuit consolári ánima mea, * memor fui Dei, et delectátus sum, et exercitátus sum: et defécit spíritus meus.

Anticipavérunt vigílias óculi mei: * turbátus sum, et non sum locútus.

Cogitávi dies antíquos: * et annos ætérnos in mente hábui.

Et meditátus sum nocte cum corde meo, * et exercitábar, et scopébam spíritum meum.

Numquid in ætérnum projíciet Deus: * aut non appónet ut complacítior sit adhuc?

Aut in finem misericórdiam suam abscíndet, * a generatióne in generatiónem?

Aut obliviscétur miseréri Deus? * aut continébit in ira sua misericórdias suas?

Et dixi: Nunc cœpi: * hæc mutátio déxteræ Excélsi.

Memor fui óperum Dómini: * quia memor ero ab inítio mirabílium tuórum.

Et meditábor in ómnibus opéribus tuis: * et in adinventiónibus tuis exercébor.

\psalmheading{Psalmus 76, ii}

\lettrine[lines=2]{D}{EUS}, in sancto via tua: * quis Deus magnus sicut Deus noster? tu es Deus qui facis mirabília.

Notam fecísti in pópulis virtútem tuam: * redemísti in bráchio tuo pópulum tuum, fílios Jacob et Joseph.

Vidérunt te aquæ, Deus, vidérunt te aquæ: * et timuérunt, et turbátæ sunt abýssi.

Multitúdo sónitus aquárum: * vocem dedérunt nubes.

Etenim sagíttæ tuæ tránseunt: * vox tonítrui tui in rota.

Illuxérunt coruscatiónes tuæ orbi terræ: * commóta est, et contrémuit terra.

In mari via tua, et sémitæ tuæ in aquis multis: * et vestígia tua non cognoscéntur.

Deduxísti sicut oves pópulum tuum, * in manu Móysi et Aaron.

\psalmheading{Psalmus 85}

\lettrine[lines=2]{I}{NCLINA}, Dómine, aurem tuam, et exáudi me: * quóniam inops, et pauper sum ego.

Custódi ánimam meam, quóniam sanctus sum: * salvum fac servum tuum, Deus meus, sperántem in te.

Miserére mei, Dómine, quóniam ad te clamávi tota die: * lætífica ánimam servi tui, quóniam ad te, Dómine, ánimam meam levávi.

Quóniam tu, Dómine, suávis, et mitis: * et multæ misericórdiæ ómnibus invocántibus te.

Auribus pércipe, Dómine, oratiónem meam: * et inténde voci deprecatiónis meæ.

In die tribulatiónis meæ clamávi ad te: * quia exaudísti me.

Non est símilis tui in diis, Dómine: * et non est secúndum ópera tua.

Omnes gentes quascúmque fecísti, vénient, et adorábunt coram te, Dómine: * et glorificábunt nomen tuum.

Quóniam magnus es tu, et fáciens mirabília: * tu es Deus solus.

Deduc me, Dómine, in via tua, et ingrédiar in veritáte tua: * lætétur cor meum ut tímeat nomen tuum.

Confitébor tibi, Dómine, Deus meus, in toto corde meo, * et glorificábo nomen tuum in ætérnum:

Quia misericórdia tua magna est super me: * et eruísti ánimam meam ex inférno inferióri.

Deus, iníqui insurrexérunt super me, et synagóga poténtium quæsiérunt ánimam meam: * et non proposuérunt te in conspéctu suo.

Et tu, Dómine, Deus miserátor et miséricors, * pátiens, et multæ misericórdiæ, et verax.

Réspice in me, et miserére mei, * da impérium tuum púero tuo: et salvum fac fílium ancíllæ tuæ.

Fac mecum signum in bonum, ut vídeant qui odérunt me, et confundántur: * quóniam tu, Dómine, adjuvísti me, et consolátus es me.

\antlabel{Ant.} Voce mea ad Dóminum clamávi: neque obliviscétur miseréri Deus.

\hymnusheading{Hymnus}

\lettrine[lines=2]{T}{E} lucis ante términum,\\
Rerum Creátor, póscimus,\\
Ut pro tua cleméntia\\
Sis præsul et custódia.

Procul recédant sómnia,\\
Et nóctium phantásmata,\\
Hostémque nostrum cómprime,\\
Ne polluántur córpora.

Præsta, Pater piíssime,\\
Patríque compar Unice,\\
Cum Spíritu Paráclito\\
Regnans per omne sæculum.\\
Amen.

\capitulumlabel{Capitulum}{Jer. 14, 9}

\lettrine[lines=2]{T}{U} autem in nobis es, Dómine, et nomen sanctum tuum invocátum est super nos: ne derelínquas nos, Dómine, Deus noster.

\R\ br. In manus tuas, Dómine, * Comméndo spíritum meum. In manus. \V\ Redemísti nos, Dómine, Deus veritátis. Comméndo. Glória Patri. In manus.

\V\ Custódi nos, Dómine, ut pupíllam óculi. \R\ Sub umbra alárum tuárum prótege nos.

\antlabel{Ant.} Salva nos.

\canticumheading{Canticum Simeonis}{Luc. 2, 29-32}

\lettrine[lines=2]{N}{UNC} dimíttis servum tuum, Dómine, * secúndum verbum tuum in pace:

Quia vidérunt óculi mei * salutáre tuum,

Quod parásti * ante fáciem ómnium populórum,

Lumen ad revelatiónem géntium, * et glóriam plebis tuæ Israël.

Glória Patri.

\antlabel{Ant.} Salva nos, Dómine, vigilántes, custódi nos dormiéntes; ut vigilémus cum Christo, et requiescámus in pace.

\variantheading{Preces (si dicendæ sint)}

KYRIE, eléison. Christe, eléison. Kýrie, eléison.

Pater noster \rubric{secreto usque ad} \V\ Et ne nos indúcas in tentatiónem. \R\ Sed líbera nos a malo.

Credo in Deum \rubric{secreto usque ad} \V\ Carnis resurrectiónem. \R\ Vitam ætérnam. Amen. \V\ Benedíctus es, Dómine, Deus patrum nostrórum. \R\ Et laudábilis et gloriósus in sæcula. \V\ Benedicámus Patrem et Fílium cum Sancto Spíritu. \R\ Laudémus, et superexaltémus eum in sæcula. \V\ Benedíctus es, Dómine, in firmaménto cæli. \R\ Et laudábilis, et gloriósus, et superexaltátus in sæcula. \V\ Benedícat et custódiat nos omnípotens et miséricors Dóminus. \R\ Amen. \V\ Dignáre, Dómine, nocte ista. \R\ Sine peccáto nos custodíre. \V\ Miserére nostri, Dómine. \R\ Miserére nostri. \V\ Fiat misericórdia tua, Dómine, super nos. \R\ Quemádmodum sperávimus in te. \V\ Dómine, exáudi oratiónem meam. \R\ Et clamor meus ad te véniat.

\V\ Dóminus vobíscum.

Orémus. \oratiolabel{Oratio}

\lettrine[lines=2]{V}{ISITA}, quæsumus, Dómine, habitatiónem istam, et omnes insídias inimíci ab ea longe repélle: Angeli tui sancti hábitent in ea, qui nos in pace custódiant; et benedíctio tua sit super nos semper. Per Dóminum nostrum. \R\ Amen.

\V\ Dóminus vobíscum.

\V\ Benedicámus Dómino.

\benedictiolabel{Benedictio} Benedícat et custódiat nos omnípotens et miséricors Dóminus, Pater, et Fílius, et Spíritus Sanctus. \R\ Amen.

\variantheading{Antiphona finalis B. M. V.}

\V\ Divínum auxílium.

Pater noster, Ave María et Credo.

% ===== END OF BATCH 35 -- Covers autumnalis.md lines 11412-11776,
% completing "Ad Vesperas" (Psalms 138/i, 138/ii, 139, 140, 141,
% hymn "Hominis supérne Cónditor") and "Ad Completorium" (Psalms
% 76/i, 76/ii, 85, hymn "Te lucis ante términum," Canticum
% Simeonis, Preces), and with them the ENTIRE "Feria Sexta" Office.
% Several more partial-caps OCR defects ("COnfitébor," "ERipe me,"
% "VOce mea" twice, "DEus, in sancto," "INclina, Dómine,"
% "BEnedíctus" [Capitulum], "HOminis supérne" [hymn]) are silently
% normalized to full caps. The closing cross-reference reverts to
% the abbreviated "Antiphona finalis B. M. V." but keeps the fuller
% "Pater noster, Ave María et Credo." seen since Feria Quinta; each
% transcribed exactly as its own occurrence gives it. The document
% continues with "Sabbato" (autumnalis.md line 11780), the final
% weekday of the Psalterium cycle.
% =====
\feriaheading{Sabbato}

\hourheading{Ad Matutinum}

Pater, Ave et Credo.

\V\ Dómine, lábia mea.

\V\ Deus, in adjutórium.

\variantheading{In Officio feriali}

\invitatlabel{Invit.} Pópulus Dómini, et oves páscuæ ejus: * Veníte, adorémus. \psalmlabel{Ps. 94} Veníte.

\hymnusheading{Hymnus}

\lettrine[lines=2]{S}{UMMAE} Parens cleméntiæ,\\
Mundi regis qui máchinam,\\
Uníus et substántiæ,\\
Trinúsque persónis Deus:

Nostros pius cum cánticis\\
Fletus benígne súscipe:\\
Ut corde puro sórdium\\
Te perfruámur lárgius.

Lumbos jecúrque mórbidum\\
Flammis adúre cóngruis,\\
Accíncti ut artus éxcubent,\\
Luxu remóto péssimo.

Quicúmque ut horas nóctium\\
Nunc concinéndo rúmpimus,\\
Ditémur omnes áffatim\\
Donis beátæ pátriæ.

Præsta, Pater piíssime,\\
Patríque compar Unice,\\
Cum Spíritu Paráclito\\
Regnans per omne sæculum.\\
Amen.

\nocturnoheading{In I Nocturno}

\antlabel{Ant.} Memor fuit in sæculum * testaménti sui Dóminus, Deus noster.

\psalmheading{Psalmus 104, i}

\lettrine[lines=2]{C}{ONFITEMINI} Dómino, et invocáte nomen ejus: * annuntiáte inter gentes ópera ejus.

Cantáte ei, et psállite ei: * narráte ómnia mirabília ejus.

Laudámini in nómine sancto ejus: * lætétur cor quæréntium Dóminum.

Quærite Dóminum, et confirmámini: * quærite fáciem ejus semper.

Mementóte mirabílium ejus, quæ fecit: * prodígia ejus, et judícia oris ejus.

Semen Abraham, servi ejus: * fílii Jacob, elécti ejus.

Ipse Dóminus Deus noster: * in univérsa terra judícia ejus.

Memor fuit in sæculum testaménti sui: * verbi, quod mandávit in mille generatiónes:

Quod dispósuit ad Abraham: * et juraménti sui ad Isaac:

Et státuit illud Jacob in præcéptum: * et Israël in testaméntum ætérnum:

Dicens: Tibi dabo terram Chánaan, * funículum hereditátis vestræ.

Cum essent número brevi, * paucíssimi et íncolæ ejus:

Et pertransiérunt de gente in gentem, * et de regno ad pópulum álterum.

Non relíquit hóminem nocére eis: * et corrípuit pro eis reges.

Nolíte tángere christos meos: * et in prophétis meis nolíte malignári.

\antlabel{Ant.} Memor fuit in sæculum testaménti sui Dóminus, Deus noster.

\antlabel{Ant.} Auxit Dóminus * pópulum suum: et firmávit eum super inimícos ejus.

\psalmheading{Psalmus 104, ii}

\lettrine[lines=2]{E}{T} vocávit famem super terram: * et omne firmaméntum panis contrívit.

Misit ante eos virum: * in servum venúmdatus est Joseph.

Humiliavérunt in compédibus pedes ejus, ferrum pertránsiit ánimam ejus, * donec veníret verbum ejus.

Elóquium Dómini inflammávit eum: * misit rex, et solvit eum; princeps populórum, et dimísit eum.

Constítuit eum dóminum domus suæ: * et príncipem omnis possessiónis suæ:

Ut erudíret príncipes ejus sicut semetípsum: * et senes ejus prudéntiam docéret.

Et intrávit Israël in Ægýptum: * et Jacob áccola fuit in terra Cham.

Et auxit pópulum suum veheménter: * et firmávit eum super inimícos ejus.

Convértit cor eórum ut odírent pópulum ejus: * et dolum fácerent in servos ejus.

Misit Móysen, servum suum: * Aaron, quem elégit ipsum.

Pósuit in eis verba signórum suórum, * et prodigiórum in terra Cham.

\antlabel{Ant.} Auxit Dóminus pópulum suum: et firmávit eum super inimícos ejus.

\antlabel{Ant.} Edúxit Deus * pópulum suum in exsultatióne, et eléctos suos in lætítia.

\psalmheading{Psalmus 104, iii}

\lettrine[lines=2]{M}{ISIT} ténebras, et obscurávit: * et non exacerbávit sermónes suos.

Convértit aquas eórum in sánguinem: * et occídit pisces eórum.

Edidit terra eórum ranas * in penetrálibus regum ipsórum.

Dixit, et venit cœnomýia: * et cínifes in ómnibus fínibus eórum.

Pósuit plúvias eórum grándinem: * ignem comburéntem in terra ipsórum.

Et percússit víneas eórum, et ficúlneas eórum: * et contrívit lignum fínium eórum.

Dixit, et venit locústa, et bruchus, * cujus non erat númerus:

Et comédit omne fœnum in terra eórum: * et comédit omnem fructum terræ eórum.

Et percússit omne primogénitum in terra eórum: * primítias omnis labóris eórum.

Et edúxit eos cum argénto et auro: * et non erat in tríbubus eórum infírmus.

Lætáta est Ægýptus in profectióne eórum: * quia incúbuit timor eórum super eos.

Expándit nubem in protectiónem eórum, * et ignem ut lucéret eis per noctem.

Petiérunt, et venit cotúrnix: * et pane cæli saturávit eos.

Dirúpit petram, et fluxérunt aquæ: * abiérunt in sicco flúmina;

Quóniam memor fuit verbi sancti sui, * quod hábuit ad Abraham, púerum suum.

Et edúxit pópulum suum in exsultatióne, * et eléctos suos in lætítia.

Et dedit illis regiónes géntium: * et labóres populórum possedérunt:

Ut custódiant justificatiónes ejus, * et legem ejus requírant.

\antlabel{Ant.} Edúxit Deus pópulum suum in exsultatióne, et eléctos suos in lætítia.

\variantheading{In Festis novem Lectionum:}

\V\ Memor fuit Dóminus verbi sancti tui. \R\ Quod locútus est ad Abraham, púerum suum.

\absolutiolabel{Absolutio.} Exáudi, Dómine.

\benedlabel{Bened.} 1 Benedictióne perpétua. 2 Unigénitus Dei Fílius. 3 Spíritus Sancti grátia.

\nocturnoheading{In II Nocturno}

\antlabel{Ant.} Salvávit eos Dóminus * propter nomen suum.

\psalmheading{Psalmus 105, i}

\lettrine[lines=2]{C}{ONFITEMINI} Dómino quóniam bonus: * quóniam in sæculum misericórdia ejus.

Quis loquétur poténtias Dómini, * audítas fáciet omnes laudes ejus?

Beáti, qui custódiunt judícium, * et fáciunt justítiam in omni témpore.

Meménto nostri, Dómine, in beneplácito pópuli tui: * vísita nos in salutári tuo:

Ad vidéndum in bonitáte electórum tuórum, ad lætándum in lætítia gentis tuæ: * ut laudéris cum hereditáte tua.

Peccávimus cum pátribus nostris: * injúste égimus, iniquitátem fécimus.

Patres nostri in Ægýpto non intellexérunt mirabília tua: * non fuérunt mémores multitúdinis misericórdiæ tuæ.

Et irritavérunt ascendéntes in mare, * Mare Rubrum.

Et salvávit eos propter nomen suum: * ut notam fáceret poténtiam suam.

Et incrépuit Mare Rubrum, et exsiccátum est: * et dedúxit eos in abýssis sicut in desérto.

Et salvávit eos de manu odiéntium: * et redémit eos de manu inimíci.

Et opéruit aqua tribulántes eos: * unus ex eis non remánsit.

Et credidérunt verbis ejus: * et laudavérunt laudem ejus.

Cito fecérunt, oblíti sunt óperum ejus: * et non sustinuérunt consílium ejus.

Et concupiérunt concupiscéntiam in desérto: * et tentavérunt Deum in inaquóso.

Et dedit eis petitiónem ipsórum: * et misit saturitátem in ánimas eórum.

\antlabel{Ant.} Salvávit eos Dóminus propter nomen suum.

\antlabel{Ant.} Oblíti sunt Deum, * qui salvávit eos.

\psalmheading{Psalmus 105, ii}

\lettrine[lines=2]{E}{T} irritavérunt Móysen in castris: * Aaron, sanctum Dómini.

Apérta est terra, et deglutívit Dathan: * et opéruit super congregatiónem Abíron.

Et exársit ignis in synagóga eórum: * flamma combússit peccatóres.

Et fecérunt vítulum in Horeb: * et adoravérunt scúlptile.

Et mutavérunt glóriam suam * in similitúdinem vítuli comedéntis fœnum.

Oblíti sunt Deum, qui salvávit eos, * qui fecit magnália in Ægýpto, mirabília in terra Cham: terribília in Mari Rubro.

Et dixit ut dispérderet eos: * si non Móyses, eléctus ejus, stetísset in confractióne in conspéctu ejus:

Ut avérteret iram ejus ne dispérderet eos: * et pro níhilo habuérunt terram desiderábilem:

Non credidérunt verbo ejus, et murmuravérunt in tabernáculis suis: * non exaudiérunt vocem Dómini.

Et elevávit manum suam super eos: * ut prostérneret eos in desérto:

Et ut dejíceret semen eórum in natiónibus: * et dispérgeret eos in regiónibus.

Et initiáti sunt Beélphegor: * et comedérunt sacrifícia mortuórum.

Et irritavérunt eum in adinventiónibus suis: * et multiplicáta est in eis ruína.

Et stetit Phínees, et placávit: * et cessávit quassátio.

Et reputátum est ei in justítiam, * in generatiónem et generatiónem usque in sempitérnum.

\antlabel{Ant.} Oblíti sunt Deum, qui salvávit eos.

\antlabel{Ant.} Cum tribularéntur, * vidit Dóminus: et audívit oratiónem eórum.

\psalmheading{Psalmus 105, iii}

\lettrine[lines=2]{E}{T} irritavérunt eum ad Aquas contradictiónis: * et vexátus est Móyses propter eos: quia exacerbavérunt spíritum ejus.

Et distínxit in lábiis suis: * non disperdidérunt gentes, quas dixit Dóminus illis.

Et commísti sunt inter gentes, et didicérunt ópera eórum: et serviérunt sculptílibus eórum: * et factum est illis in scándalum.

Et immolavérunt fílios suos, * et fílias suas dæmóniis.

Et effudérunt sánguinem innocéntem: * sánguinem filiórum suórum et filiárum suárum, quas sacrificavérunt sculptílibus Chánaan.

Et infécta est terra in sanguínibus, et contamináta est in opéribus eórum: * et fornicáti sunt in adinventiónibus suis.

Et irátus est furóre Dóminus in pópulum suum: * et abominátus est hereditátem suam.

Et trádidit eos in manus géntium: * et domináti sunt eórum qui odérunt eos.

Et tribulavérunt eos inimíci eórum, et humiliáti sunt sub mánibus eórum: * sæpe liberávit eos.

Ipsi autem exacerbavérunt eum in consílio suo: * et humiliáti sunt in iniquitátibus suis.

Et vidit, cum tribularéntur: * et audívit oratiónem eórum.

Et memor fuit testaménti sui: * et pœnítuit eum secúndum multitúdinem misericórdiæ suæ.

Et dedit eos in misericórdias * in conspéctu ómnium qui céperant eos.

Salvos nos fac, Dómine, Deus noster: * et cóngrega nos de natiónibus:

Ut confiteámur nómini sancto tuo: * et gloriémur in laude tua.

Benedíctus Dóminus Deus Israël, a sæculo et usque in sæculum: * et dicet omnis pópulus: Fiat, fiat.

\antlabel{Ant.} Cum tribularéntur, vidit Dóminus: et audívit oratiónem eórum.

\variantheading{In Officiis novem Lectionum:}

\V\ Salvos nos fac, Dómine, Deus noster. \R\ Ut confiteámur nómini sancto tuo.

\absolutiolabel{Absolutio.} Ipsíus píetas et.

\benedlabel{Bened.} 4 Deus Pater omnípotens. 5 Christus perpétuæ. 6 Ignem sui amóris.

\nocturnoheading{In III Nocturno}

\antlabel{Ant.} Clamavérunt ad Dóminum: * et de necessitátibus eórum liberávit eos.

\psalmheading{Psalmus 106, i}

\lettrine[lines=2]{C}{ONFITEMINI} Dómino quóniam bonus: * quóniam in sæculum misericórdia ejus.

Dicant qui redémpti sunt a Dómino, quos redémit de manu inimíci: * et de regiónibus congregávit eos:

A solis ortu, et occásu: * ab aquilóne, et mari.

Erravérunt in solitúdine in inaquóso: * viam civitátis habitáculi non invenérunt.

Esuriéntes, et sitiéntes: * ánima eórum in ipsis defécit.

Et clamavérunt ad Dóminum cum tribularéntur: * et de necessitátibus eórum erípuit eos.

Et dedúxit eos in viam rectam: * ut irent in civitátem habitatiónis.

Confiteántur Dómino misericórdiæ ejus: * et mirabília ejus fíliis hóminum.

Quia satiávit ánimam inánem: * et ánimam esuriéntem satiávit bonis.

Sedéntes in ténebris, et umbra mortis: * vinctos in mendicitáte et ferro.

Quia exacerbavérunt elóquia Dei: * et consílium Altíssimi irritavérunt.

Et humiliátum est in labóribus cor eórum: * infirmáti sunt, nec fuit qui adjuváret.

Et clamavérunt ad Dóminum cum tribularéntur: * et de necessitátibus eórum liberávit eos.

Et edúxit eos de ténebris, et umbra mortis: * et víncula eórum disrúpit.

\antlabel{Ant.} Clamavérunt ad Dóminum: et de necessitátibus eórum liberávit eos.

\antlabel{Ant.} Ipsi vidérunt * ópera Dei, et mirabília ejus.

\psalmheading{Psalmus 106, ii}

\lettrine[lines=2]{C}{ONFITEANTUR} Dómino misericórdiæ ejus: * et mirabília ejus fíliis hóminum.

Quia contrívit portas áereas: * et vectes férreos confrégit.

Suscépit eos de via iniquitátis eórum: * propter injustítias enim suas humiliáti sunt.

Omnem escam abomináta est ánima eórum: * et appropinquavérunt usque ad portas mortis.

Et clamavérunt ad Dóminum cum tribularéntur: * et de necessitátibus eórum liberávit eos.

Misit verbum suum, et sanávit eos: * et erípuit eos de interitiónibus eórum.

Confiteántur Dómino misericórdiæ ejus: * et mirabília ejus fíliis hóminum.

Et sacríficent sacrifícium laudis: * et annúntient ópera ejus in exsultatióne.

Qui descéndunt mare in návibus, * faciéntes operatiónem in aquis multis.

Ipsi vidérunt ópera Dómini, * et mirabília ejus in profúndo.

Dixit, et stetit spíritus procéllæ: * et exaltáti sunt fluctus ejus.

Ascéndunt usque ad cælos, et descéndunt usque ad abýssos: * ánima eórum in malis tabescébat.

Turbáti sunt, et moti sunt sicut ébrius: * et omnis sapiéntia eórum devoráta est.

Et clamavérunt ad Dóminum cum tribularéntur: * et de necessitátibus eórum edúxit eos.

Et státuit procéllam ejus in auram: * et siluérunt fluctus ejus.

Et lætáti sunt quia siluérunt: * et dedúxit eos in portum voluntátis eórum.

\antlabel{Ant.} Ipsi vidérunt ópera Dei, et mirabília ejus.

\antlabel{Ant.} Vidébunt recti, * et lætabúntur, et intélligent misericórdias Dómini.

\psalmheading{Psalmus 106, iii}

\lettrine[lines=2]{C}{ONFITEANTUR} Dómino misericórdiæ ejus: * et mirabília ejus fíliis hóminum.

Et exáltent eum in ecclésia plebis: * et in cáthedra seniórum laudent eum.

Pósuit flúmina in desértum: * et éxitus aquárum in sitim.

Terram fructíferam in salsúginem, * a malítia inhabitántium in ea.

Pósuit desértum in stagna aquárum: * et terram sine aqua in éxitus aquárum.

Et collocávit illic esuriéntes: * et constituérunt civitátem habitatiónis.

Et seminavérunt agros, et plantavérunt víneas: * et fecérunt fructum nativitátis.

Et benedíxit eis, et multiplicáti sunt nimis: * et juménta eórum non minorávit.

Et pauci facti sunt: * et vexáti sunt a tribulatióne malórum, et dolóre.

Effúsa est contémptio super príncipes: * et erráre fecit eos in ínvio, et non in via.

Et adjúvit páuperem de inópia: * et pósuit sicut oves famílias.

Vidébunt recti, et lætabúntur: * et omnis iníquitas oppilábit os suum.

Quis sápiens et custódiet hæc? * et intélliget misericórdias Dómini?

\antlabel{Ant.} Vidébunt recti, et lætabúntur, et intélligent misericórdias Dómini.

\variantheading{In feriali Officio in Officio sanctæ Mariæ in Sabbato et in Festis:}

\V\ Exáltent Dóminum in ecclésia plebis. \R\ Et in cáthedra seniórum laudent eum.

\variantheading{In Officio novem Lectionum}

\absolutiolabel{Absolutio.} A vínculis.

\benedlabel{Bened.} 7 Evangélica léctio. 8 \rubric{In Festis Domini:} Divínum auxílium. \rubric{In Festis B. M. V.:} Cujus festum cólimus, ipsa Virgo. \rubric{In Festis Sanctorum:} Cujus (\textit{vel} Quorum \textit{aut} Quarum) festum cólimus. 9 Ad societátem. (Per evangélica dicta.)

\variantheading{In Officio trium Lectionum}

\absolutiolabel{Absolutio.} A vínculis.

\variantheading{In Feriis cum Homilia:}

\benedlabel{Bened.} 1 Evangélica. 2 Divínum. 3 Ad societátem.

\variantheading{In Octava simplici:}

\benedlabel{Bened.} 1 Ille nos benedícat. 2 Cujus (\textit{vel} Quorum \textit{aut} Quarum). 3 Ad societátem.

\variantheading{In Officio sanctæ Mariæ in Sabbato}

\absolutiolabel{Absol.} Précibus et méritis beátæ Maríæ semper Vírginis et ómnium Sanctórum, perdúcat nos Dóminus ad regna cælórum. Amen.

\benedlabel{Bened.} 1 Nos cum prole pia benedícat Virgo María. Amen.

2 Ipsa Virgo vírginum intercédat pro nobis ad Dóminum. Amen.

3 Per Vírginem Matrem concédat nobis Dóminus salútem et pacem. Amen.

% ===== END OF BATCH 36 -- Covers autumnalis.md lines 11780-12200,
% opening "Sabbato" (Saturday) and completing its "Ad Matutinum" in
% full: the ferial Invitatorium "Pópulus Dómini, et oves páscuæ
% ejus," the hymn "Summæ Parens cleméntiæ," and the three
% Nocturns -- Psalm 104 in the first, Psalm 105 in the second,
% Psalm 106 in the third -- with their versicle, Absolutio, and
% Benedictiones apparatus. The ninth-lesson "Bened." line (line
% 12176) marks its three conditional clauses ("In Festis Domini:",
% "In Festis B. M. V.:", "In Festis Sanctorum:") as actual red
% rubric text in the source this time ({RUBRIC: ...}), unlike the
% plain black-text embedding of the identical clauses at every
% prior occurrence of this line this volume; this is rendered with
% \rubric{} on each of the three labels, preserving the genuine
% embedded-rubric distinction exactly as the source marks it here.
% Saturday Matins closes with a Sabbath-specific apparatus unique
% to this day: "In Officio sanctæ Mariæ in Sabbato," with its own
% full-sentence Absolutio ("Précibus et méritis beátæ Maríæ...")
% and a three-part Marian Benedictio given as three separate
% paragraphs in the source (each ending its own "Amen."), reproduced
% here with the same paragraph breaks. Several more partial-caps
% OCR defects ("SUmmæ," "COnfitémini"/"COnfiteántur" repeatedly,
% "ET irritavérunt" twice, "MIsit ténebras") are silently normalized
% to full caps. The document continues with "Ad Laudes I"
% (autumnalis.md line 12202).
% =====
\hourheading{Ad Laudes I}

\rubric{In Sabbatis extra Vigílias commúnes occurréntibus, in Festis novem Lectiónum, in ómnibus Officiis de Octáva non privilegiáta, atque in Officio sanctæ Maríæ in Sabbato.}

(Pater noster et Ave.)

\V\ Deus, in adjutórium.

\antlabel{Ant.} Fílii Sion * exsúltent in Rege suo.

\psalmheading{Psalmus 149}

\lettrine[lines=2]{C}{ANTATE} Dómino cánticum novum: * laus ejus in ecclésia sanctórum.

Lætétur Israël in eo, qui fecit eum: * et fílii Sion exsúltent in rege suo.

Laudent nomen ejus in choro: * in týmpano, et psaltério psallant ei:

Quia beneplácitum est Dómino in pópulo suo: * et exaltábit mansuétos in salútem.

Exsultábunt sancti in glória: * lætabúntur in cubílibus suis.

Exaltatiónes Dei in gútture eórum: * et gládii ancípites in mánibus eórum.

Ad faciéndam vindíctam in natiónibus: * increpatiónes in pópulis.

Ad alligándos reges eórum in compédibus: * et nóbiles eórum in mánicis férreis.

Ut fáciant in eis judícium conscríptum: * glória hæc est ómnibus sanctis ejus.

\antlabel{Ant.} Fílii Sion exsúltent in Rege suo.

\antlabel{Ant.} Quam magnificáta * sunt ópera tua, Dómine!

\psalmheading{Psalmus 91}

\lettrine[lines=2]{B}{ONUM} est confitéri Dómino: * et psállere nómini tuo, Altíssime.

Ad annuntiándum mane misericórdiam tuam: * et veritátem tuam per noctem.

In decachórdo, psaltério: * cum cántico, in cíthara.

Quia delectásti me, Dómine, in factúra tua: * et in opéribus mánuum tuárum exsultábo.

Quam magnificáta sunt ópera tua, Dómine! * nimis profúndæ factæ sunt cogitatiónes tuæ.

Vir insípiens non cognóscet: * et stultus non intélliget hæc.

Cum exórti fúerint peccatóres sicut fœnum: * et apparúerint omnes, qui operántur iniquitátem:

Ut intéreant in sæculum sæculi: * tu autem Altíssimus in ætérnum, Dómine.

Quóniam ecce inimíci tui, Dómine, quóniam ecce inimíci tui períbunt: * et dispergéntur omnes, qui operántur iniquitátem.

Et exaltábitur sicut unicórnis cornu meum: * et senéctus mea in misericórdia úberi.

Et despéxit óculus meus inimícos meos: * et in insurgéntibus in me malignántibus áudiet auris mea.

Justus ut palma florébit: * sicut cedrus Líbani multiplicábitur.

Plantáti in domo Dómini, * in átriis domus Dei nostri florébunt.

Adhuc multiplicabúntur in senécta úberi: * et patiéntes erunt, ut annúntient:

Quóniam rectus Dóminus Deus noster: * et non est iníquitas in eo.

\antlabel{Ant.} Quam magnificáta sunt ópera tua, Dómine!

\antlabel{Ant.} Lætábitur justus * in Dómino, et sperábit in eo.

\psalmheading{Psalmus 63}

\lettrine[lines=2]{E}{XAUDI}, Deus, oratiónem meam cum déprecor: * a timóre inimíci éripe ánimam meam.

Protexísti me a convéntu malignántium: * a multitúdine operántium iniquitátem.

Quia exacuérunt ut gládium linguas suas: * intendérunt arcum rem amáram, ut sagíttent in occúltis immaculátum.

Súbito sagittábunt eum, et non timébunt: * firmavérunt sibi sermónem nequam.

Narravérunt ut abscónderent láqueos: * dixérunt: Quis vidébit eos?

Scrutáti sunt iniquitátes: * defecérunt scrutántes scrutínio.

Accédet homo ad cor altum: * et exaltábitur Deus.

Sagíttæ parvulórum factæ sunt plagæ eórum: * et infirmátæ sunt contra eos linguæ eórum.

Conturbáti sunt omnes qui vidébant eos: * et tímuit omnis homo.

Et annuntiavérunt ópera Dei, * et facta ejus intellexérunt.

Lætábitur justus in Dómino, et sperábit in eo, * et laudabúntur omnes recti corde.

\antlabel{Ant.} Lætábitur justus in Dómino, et sperábit in eo.

\antlabel{Ant.} Osténde nobis, Dómine, * lumen miseratiónum tuárum.

\canticumheading{Canticum Ecclesiastici}{Eccli. 36, 1-16}

\lettrine[lines=2]{M}{ISERERE} nostri, Deus ómnium, et réspice nos, * et osténde nobis lucem miseratiónum tuárum:

Et immítte timórem tuum super gentes, * quæ non exquisiérunt te,

Ut cognóscant quia non est Deus nisi tu, * et enárrent magnália tua.

Alleva manum tuam super gentes aliénas, * ut vídeant poténtiam tuam.

Sicut enim in conspéctu eórum sanctificátus es in nobis, * sic in conspéctu nostro magnificáberis in eis,

Ut cognóscant te, sicut et nos cognóvimus, * quóniam non est Deus præter te, Dómine.

Innova signa, et immúta mirabília. * Glorífica manum, et bráchium déxtrum.

Excita furórem, et effúnde iram. * Tolle adversárium, et afflíge inimícum.

Festína tempus, et meménto finis, * ut enárrent mirabília tua.

In ira flammæ devorétur qui salvátur: * et qui péssimant plebem tuam, invéniant perditiónem.

Cóntere caput príncipum inimicórum, * dicéntium: Non est álius præter nos.

Cóngrega omnes tribus Jacob: ut cognóscant quia non est Deus nisi tu, * et enárrent magnália tua:

Et hereditábis eos, * sicut ab inítio.

Miserére plebi tuæ, super quam invocátum est nomen tuum: * et Israël, quem coæquásti primogénito tuo.

Miserére civitáti sanctificatiónis tuæ Jerúsalem, * civitáti requiéi tuæ.

Reple Sion inenarrabílibus verbis tuis, * et glória tua pópulum tuum.

\antlabel{Ant.} Osténde nobis, Dómine, lucem miseratiónum tuárum.

\antlabel{Ant.} Omnis spíritus * laudet Dóminum.

\psalmheading{Psalmus 150}

\lettrine[lines=2]{L}{AUDATE} Dóminum in sanctis ejus: * laudáte eum in firmaménto virtútis ejus.

Laudáte eum in virtútibus ejus: * laudáte eum secúndum multitúdinem magnitúdinis ejus.

Laudáte eum in sono tubæ: * laudáte eum in psaltério, et cíthara.

Laudáte eum in týmpano, et choro: * laudáte eum in chordis, et órgano.

Laudáte eum in cýmbalis benesonántibus: laudáte eum in cýmbalis jubilatiónis: * omnis spíritus laudet Dóminum.

\antlabel{Ant.} Omnis spíritus laudet Dóminum.

\rubric{In Officio Festi vel Octávæ ac sanctæ Maríæ in Sabbato Capítulum, Hymnus, Versus, Ant. ad Benedictus et Orátio ut in Próprio vel Communi.}

\rubric{Réliqua ut in Ordinário. 19}

% ===== END OF BATCH 37 -- Covers autumnalis.md lines 12202-12357,
% "Ad Laudes I" for Sabbato in full: Psalms 149, 91, 63, the
% Canticum Ecclesiastici ("Miserére nostri, Deus ómnium"), and
% Psalm 150. Unlike every other feria's Ad Laudes I this volume,
% Saturday's carries no ferial Capitulum/Hymnus of its own -- the
% source instead gives a standalone rubric directing that these,
% along with the Versus, Ant. ad Benedictus, and Oratio, be taken
% "ut in Próprio vel Communi" (from the Proper or Common of the
% occurring feast or the Officium of Our Lady on Saturday), since
% Saturday's Lauds is always completed from elsewhere rather than
% having its own fixed ferial conclusion. Several more psalm
% incipits carry the partial-caps OCR defect ("LAudáte," "MIserére
% nostri," "BOnum est") and are silently normalized to full caps.
% The document continues with "Ad Laudes II" (autumnalis.md line
% 12360).
% =====
\hourheading{Ad Laudes II}

\rubric{In Vigíliis commúnibus, quando fit Officium de Feria.}

(Pater noster et Ave.)

\V\ Deus, in adjutórium.

\antlabel{Ant.} Benígne fac, Dómine.

\psalmheading{Psalmus 50}

\lettrine[lines=2]{M}{ISERERE} mei, Deus, * secúndum magnam misericórdiam tuam.

Et secúndum multitúdinem miseratiónum tuárum, * dele iniquitátem meam.

Amplius lava me ab iniquitáte mea: * et a peccáto meo munda me.

Quóniam iniquitátem meam ego cognósco: * et peccátum meum contra me est semper.

Tibi soli peccávi, et malum coram te feci: * ut justificéris in sermónibus tuis, et vincas cum judicáris.

Ecce enim, in iniquitátibus concéptus sum: * et in peccátis concépit me mater mea.

Ecce enim, veritátem dilexísti: * incérta et occúlta sapiéntiæ tuæ manifestásti mihi.

Aspérges me hyssópo, et mundábor: * lavábis me, et super nivem dealbábor.

Audítui meo dabis gáudium et lætítiam: * et exsultábunt ossa humiliáta.

Avérte fáciem tuam a peccátis meis: * et omnes iniquitátes meas dele.

Cor mundum crea in me, Deus: * et spíritum rectum ínnova in viscéribus meis.

Ne projícias me a fácie tua: * et spíritum sanctum tuum ne áuferas a me.

Redde mihi lætítiam salutáris tui: * et spíritu principáli confírma me.

Docébo iníquos vias tuas: * et ímpii ad te converténtur.

Líbera me de sanguínibus, Deus, Deus salútis meæ: * et exsultábit lingua mea justítiam tuam.

Dómine, lábia mea apéries: * et os meum annuntiábit laudem tuam.

Quóniam si voluísses sacrifícium, dedíssem útique: * holocáustis non delectáberis.

Sacrifícium Deo spíritus contribulátus: * cor contrítum, et humiliátum, Deus, non despícies.

Benígne fac, Dómine, in bona voluntáte tua Sion: * ut ædificéntur muri Jerúsalem.

Tunc acceptábis sacrifícium justítiæ, oblatiónes, et holocáusta: * tunc impónent super altáre tuum vítulos.

\antlabel{Ant.} Benígne fac, Dómine, in bona voluntáte tua Sion.

\antlabel{Ant.} Rectus Dóminus.

\psalmheading{Psalmus 91}

\lettrine[lines=2]{B}{ONUM} est confitéri Dómino: * et psállere nómini tuo, Altíssime.

Ad annuntiándum mane misericórdiam tuam: * et veritátem tuam per noctem.

In decachórdo, psaltério: * cum cántico, in cíthara.

Quia delectásti me, Dómine, in factúra tua: * et in opéribus mánuum tuárum exsultábo.

Quam magnificáta sunt ópera tua, Dómine! * nimis profúndæ factæ sunt cogitatiónes tuæ.

Vir insípiens non cognóscet: * et stultus non intélliget hæc.

Cum exórti fúerint peccatóres sicut fœnum: * et apparúerint omnes, qui operántur iniquitátem:

Ut intéreant in sæculum sæculi: * tu autem Altíssimus in ætérnum, Dómine.

Quóniam ecce inimíci tui, Dómine, quóniam ecce inimíci tui períbunt: * et dispergéntur omnes, qui operántur iniquitátem.

Et exaltábitur sicut unicórnis cornu meum: * et senéctus mea in misericórdia úberi.

Et despéxit óculus meus inimícos meos: * et in insurgéntibus in me malignántibus áudiet auris mea.

Justus ut palma florébit: * sicut cedrus Líbani multiplicábitur.

Plantáti in domo Dómini, * in átriis domus Dei nostri florébunt.

Adhuc multiplicabúntur in senécta úberi: * et patiéntes erunt, ut annúntient:

Quóniam rectus Dóminus, Deus noster: * et non est iníquitas in eo.

\antlabel{Ant.} Rectus Dóminus, Deus noster: et non est iníquitas in eo.

\antlabel{Ant.} A timóre inimíci.

\psalmheading{Psalmus 63}

\lettrine[lines=2]{E}{XAUDI}, Deus, oratiónem meam cum déprecor: * a timóre inimíci éripe ánimam meam.

Protexísti me a convéntu malignántium: * a multitúdine operántium iniquitátem.

Quia exacuérunt ut gládium linguas suas: * intendérunt arcum rem amáram, ut sagíttent in occúltis immaculátum.

Súbito sagittábunt eum, et non timébunt: * firmavérunt sibi sermónem nequam.

Narravérunt ut abscónderent láqueos: * dixérunt: Quis vidébit eos?

Scrutáti sunt iniquitátes: * defecérunt scrutántes scrutínio.

Accédet homo ad cor altum: * et exaltábitur Deus.

Sagíttæ parvulórum factæ sunt plagæ eórum: * et infirmátæ sunt contra eos linguæ eórum.

Conturbáti sunt omnes qui vidébant eos: * et tímuit omnis homo.

Et annuntiavérunt ópera Dei, * et facta ejus intellexérunt.

Lætábitur justus in Dómino, et sperábit in eo, * et laudabúntur omnes recti corde.

\antlabel{Ant.} A timóre inimíci éripe, Dómine, ánimam meam.

\antlabel{Ant.} In servis suis.

\canticumheading{Canticum Moysis}{Deut. 32, 1-43}

\lettrine[lines=2]{A}{UDITE}, cæli, quæ loquor: * áudiat terra verba oris mei.

Concréscat ut plúvia doctrína mea, * fluat ut ros elóquium meum,

Quasi imber super herbam, et quasi stillæ super grámina. * Quia nomen Dómini invocábo:

Date magnificéntiam Deo nostro. * Dei perfécta sunt ópera, et omnes viæ ejus judícia:

Deus fidélis, et absque ulla iniquitáte, justus et rectus. * Peccavérunt ei, et non fílii ejus in sórdibus:

Generátio prava atque pervérsa. * Hæccine réddis Dómino, pópule stulte et insípiens?

Numquid non ipse est pater tuus, * qui possédit te, et fecit, et creávit te?

Meménto diérum antiquórum, * cógita generatiónes síngulas:

Intérroga patrem tuum, et annuntiábit tibi: * majóres tuos, et dicent tibi.

Quando dividébat Altíssimus gentes: * quando separábat fílios Adam,

Constítuit términos populórum * juxta númerum filiórum Israël.

Pars autem Dómini, pópulus ejus: * Jacob funículus hereditátis ejus.

Invénit eum in terra desérta, * in loco horróris, et vastæ solitúdinis:

Circumdúxit eum, et dócuit: * et custodívit quasi pupíllam óculi sui.

Sicut áquila próvocans ad volándum pullos suos, * et super eos vólitans,

Expándit alas suas, et assúmpsit eum, * atque portávit in húmeris suis.

Dóminus solus dux ejus fuit: * et non erat cum eo deus aliénus.

Constítuit eum super excélsam terram: * ut coméderet fructus agrórum,

Ut súgeret mel de petra, * oleúmque de saxo duríssimo.

Butýrum de arménto, et lac de óvibus * cum ádipe agnórum, et aríetum filiórum Basan:

Et hircos cum medúlla trítici, * et sánguinem uvæ bíberet meracíssimum.

Incrassátus est diléctus, et recalcitrávit: * incrassátus, impinguátus, dilatátus,

Derelíquit Deum, factórem suum, * et recéssit a Deo, salutári suo.

Provocavérunt eum in diis aliénis, * et in abominatiónibus ad iracúndiam concitavérunt.

Immolavérunt dæmóniis, et non Deo, * diis, quos ignorábant:

Novi recentésque venérunt, * quos non coluérunt patres eórum.

Deum qui te génuit dereliquísti, * et oblítus es Dómini, creatóris tui.

Vidit Dóminus, et ad iracúndiam concitátus est: * quia provocavérunt eum fílii sui et fíliæ.

Et ait: Abscóndam fáciem meam ab eis, * et considerábo novíssima eórum:

Generátio enim pervérsa est, * et infidéles fílii.

Ipsi me provocavérunt in eo, qui non erat Deus, * et irritavérunt in vanitátibus suis:

Et ego provocábo eos in eo, qui non est pópulus, * et in gente stulta irritábo illos.

Ignis succénsus est in furóre meo, * et ardébit usque ad inférni novíssima:

Devorabítque terram cum gérmine suo, * et móntium fundaménta combúret.

Congregábo super eos mala, * et sagíttas meas complébo in eis.

Consuméntur fame, * et devorábunt eos aves morsu amaríssimo:

Dentes bestiárum immíttam in eos, * cum furóre trahéntium super terram, atque serpéntium.

Foris vastábit eos gládius, et intus pavor, * júvenem simul ac vírginem, lacténtem cum hómine sene.

Dixi: Ubinam sunt? * cessáre fáciam ex homínibus memóriam eórum.

Sed propter iram inimicórum dístuli: * ne forte superbírent hostes eórum,

Et dícerent: * Manus nostra excélsa, et non

Dóminus, * fecit hæc ómnia.

Gens absque consílio est, et sine prudéntia. * Utinam sáperent, et intellígerent, ac novíssima providérent.

Quómodo persequátur unus mille, * et duo fugent decem míllia?

Nonne ídeo, quia Deus suus véndidit eos, * et Dóminus conclúsit illos?

Non enim est Deus noster ut dii eórum: * et inimíci nostri sunt júdices.

De vínea Sodomórum, vínea eórum, * et de suburbánis Gomórrhæ:

Uva eórum uva fellis, * et botri amaríssimi.

Fel dracónum vinum eórum, * et venénum áspidum insanábile.

Nonne hæc cóndita sunt apud me, * et signáta in thesáuris meis?

Mea est últio, et ego retríbuam in témpore, * ut labátur pes eórum:

Juxta est dies perditiónis, * et adésse festínant témpora.

Judicábit Dóminus pópulum suum, * et in servis suis miserébitur:

Vidébit quod infirmáta sit manus, * et clausi quoque defecérunt, residuíque consúmpti sunt.

Et dicet: Ubi sunt dii eórum, * in quibus habébant fidúciam?

De quorum víctimis comedébant ádipes, * et bibébant vinum libáminum:

Surgant, et opituléntur vobis, * et in necessitáte vos prótegant.

Vidéte quod ego sim solus, * et non sit álius Deus præter me:

Ego occídam, et ego vívere fáciam: percútiam, et ego sanábo, * et non est qui de manu mea possit erúere.

Levábo ad cælum manum meam, et dicam: * Vivo ego in ætérnum.

Si acúero ut fulgur gládium meum, * et arripúerit judícium manus mea:

Reddam ultiónem hóstibus meis, * et his qui odérunt me retríbuam.

Inebriábo sagíttas meas sánguine, * et gládius meus devorábit carnes,

De cruóre occisórum, * et de captivitáte, nudáti inimicórum cápitis.

Laudáte, gentes, pópulum ejus, * quia sánguinem servórum suórum ulciscétur:

Et vindíctam retríbuet in hostes eórum, * et propítius erit terræ pópuli sui.

\antlabel{Ant.} In servis suis miserébitur Dóminus: et propítius erit terræ pópuli sui.

\antlabel{Ant.} Laudáte. \textdagger

\psalmheading{Psalmus 150}

\lettrine[lines=2]{L}{AUDATE} \textdagger\ Dóminum in sanctis ejus: * laudáte eum in firmaménto virtútis ejus.

Laudáte eum in virtútibus ejus: * laudáte eum secúndum multitúdinem magnitúdinis ejus.

Laudáte eum in sono tubæ: * laudáte eum in psaltério, et cíthara.

Laudáte eum in týmpano, et choro: * laudáte eum in chordis, et órgano.

Laudáte eum in cýmbalis benesonántibus: laudáte eum in cýmbalis jubilatiónis: * omnis spíritus laudet Dóminum.

\antlabel{Ant.} Laudáte Dóminum secúndum multitúdinem magnitúdinis ejus.

\capitulumlabel{Capit.}{Rom. 13, 12-13}

\lettrine[lines=2]{N}{OX} præcéssit, dies autem appropinquávit. Abjiciámus ergo ópera tenebrárum, et induámur arma lucis. Sicut in die honéste ambulémus.

\hymnusheading{Hymnus}

\lettrine[lines=2]{A}{URORA} jam spargit polum:\\
Terris dies illábitur:\\
Lucis resúltat spículum:\\
Discédat omne lúbricum.

Phantásma noctis éxsulet:\\
Mentis reátus córruat:\\
Quidquid tenébris hórridum\\
Nox áttulit culpæ, cadat.

Ut mane, quod nos últimum\\
Hic deprecámur cérnui,\\
Cum luce nobis éffluat,\\
Hoc dum canóre cóncrepat.

Deo Patri sit glória,\\
Ejúsque soli Fílio,\\
Cum Spíritu Paráclito,\\
Nunc et per omne sæculum. Amen.

\V\ Repléti sumus mane misericórdia tua. \R\ Exsultávimus, et delectáti sumus.

\antlabel{Ad Bened. Ant.} Illúmina, Dómine, * sedéntes in ténebris et umbra mortis, et dírige pedes nostros in viam pacis.

\rubric{Expléto Cántico ac repetíta Antíphona, dicúntur sequéntes}

\variantheading{Preces feriales}

KYRIE, eléison. Christe, eléison. Kýrie, eléison.

Pater noster, \rubric{quod dícitur a solo Hebdomadário totum clara voce, usque ad} \V\ Et ne nos indúcas in tentatiónem. \R\ Sed líbera nos a malo. \V\ Ego dixi: Dómine, miserére mei. \R\ Sana ánimam meam, quia peccávi tibi. \V\ Convértere, Dómine, úsquequo? \R\ Et deprecábilis esto super servos tuos. \V\ Fiat misericórdia tua, Dómine, super nos. \R\ Quemádmodum sperávimus in te. \V\ Sacerdótes tui induántur justítiam. \R\ Et sancti tui exsúltent. \V\ Orémus pro beatíssimo Papa nostro N. \R\ Dóminus consérvet eum, et vivíficet eum, et beátum fáciat eum in terra, et non tradat eum in ánimam inimicórum ejus. \V\ Orémus et pro Antístite nostro N. \R\ Stet et pascat in fortitúdine tua, Dómine, in sublimitáte nóminis tui. \V\ Dómine, salvum fac regem. \R\ Et exáudi nos in die, qua invocavérimus te. \V\ Salvum fac pópulum tuum, Dómine, et bénedic hereditáti tuæ. \R\ Et rege eos, et extólle illos usque in ætérnum. \V\ Meménto Congregatiónis tuæ. \R\ Quam possedísti ab inítio. \V\ Fiat pax in virtúte tua. \R\ Et abundántia in túrribus tuis. \V\ Orémus pro benefactóribus nostris. \R\ Retribúere dignáre, Dómine, ómnibus, nobis bona faciéntibus propter nomen tuum, vitam ætérnam. Amen. \V\ Orémus pro fidélibus defúnctis. \R\ Réquiem ætérnam dona eis, Dómine, et lux perpétua lúceat eis. \V\ Requiéscant in pace. \R\ Amen. \V\ Pro frátribus nostris abséntibus. \R\ Salvos fac servos tuos, Deus meus, sperántes in te. \V\ Pro afflíctis et captívis. \R\ Líbera eos, Deus Israël, ex ómnibus tribulatiónibus suis. \V\ Mitte eis, Dómine, auxílium de sancto. \R\ Et de Sion tuére eos. \V\ Dómine, Deus virtútum, convérte nos. \R\ Et osténde fáciem tuam, et salvi érimus. \V\ Exsúrge, Christe, ádjuva nos. \R\ Et líbera nos propter nomen tuum. \V\ Dómine, exáudi oratiónem meam. \R\ Et clamor meus ad te véniat.

\rubric{Réliqua ut in Ordinário. 19}

% ===== END OF BATCH 38 -- Covers autumnalis.md lines 12360-12673,
% "Ad Laudes II" for Sabbato in full: Psalms 50, 91, 63, the
% Canticum Moysis ("Audíte, cæli, quæ loquor," the long canticle of
% Deuteronomy 32), and Psalm 150, with its own Capitulum "Nox
% præcéssit" and hymn "Auróra jam spargit polum" (unlike Ad Laudes
% I, which borrowed these from the Proper/Common), plus the "Preces
% feriales" litany. The source breaks one canticle verse oddly
% across two short lines ("Et dícerent: * Manus nostra excélsa, et
% non" / "Dóminus, * fecit hæc ómnia."), reproduced exactly as the
% source paragraphs it. The document continues with "Ad Primam"
% (autumnalis.md line 12676).
% =====
\hourheading{Ad Primam}

Pater, Ave et Credo.

\V\ Deus, in adjutórium.

\hymnusheading{Hymnus}

\lettrine[lines=2]{J}{AM} lucis orto sídere,\\
Deum precémur súpplices,\\
Ut in diúrnis áctibus\\
Nos servet a nocéntibus.

Linguam refrénans témperet,\\
Ne litis horror ínsonet:\\
Visum fovéndo cóntegat,\\
Ne vanitátes háuriat.

Sint pura cordis íntima,\\
Absístat et vecórdia:\\
Carnis terat supérbiam\\
Potus cibíque párcitas.

Ut, cum dies abscésserit\\
Noctémque sors redúxerit,\\
Mundi per abstinéntiam\\
Ipsi canámus glóriam.

Deo Patri sit glória,\\
Ejúsque soli Fílio,\\
Cum Spíritu Paráclito,\\
Nunc et per omne sæculum.\\
Amen.

\antlabel{Ant.} Exaltáre, Dómine.

\psalmheading{Psalmus 93, i}

\lettrine[lines=2]{D}{EUS} ultiónum Dóminus: * Deus ultiónum líbere egit.

Exaltáre, qui júdicas terram: * redde retributiónem supérbis.

Usquequo peccatóres, Dómine, * úsquequo peccatóres gloriabúntur:

Effabúntur, et loquéntur iniquitátem: * loquéntur omnes, qui operántur injustítiam?

Pópulum tuum, Dómine, humiliavérunt: * et hereditátem tuam vexavérunt.

Víduam, et ádvenam interfecérunt: * et pupíllos occidérunt.

Et dixérunt: Non vidébit Dóminus, * nec intélliget Deus Jacob.

Intellígite, insipiéntes in pópulo: * et stulti, áliquándo sápite.

Qui plantávit aurem, non áudiet? * aut qui finxit óculum, non consíderat?

Qui córripit gentes, non árguet: * qui docet hóminem sciéntiam?

Dóminus scit cogitatiónes hóminum, * quóniam vanæ sunt.

\psalmheading{Psalmus 93, ii}

\lettrine[lines=2]{B}{EATUS} homo, quem tu erudíeris, Dómine: * et de lege tua docúeris eum,

Ut mítiges ei a diébus malis: * donec fodiátur peccatóri fóvea.

Quia non repéllet Dóminus plebem suam: * et hereditátem suam non derelínquet.

Quoadúsque justítia convertátur in judícium: * et qui juxta illam omnes qui recto sunt corde.

Quis consúrget mihi advérsus malignántes? * aut quis stabit mecum advérsus operántes iniquitátem?

Nisi quia Dóminus adjúvit me: * paulo minus habitásset in inférno ánima mea.

Si dicébam: Motus est pes meus: * misericórdia tua, Dómine, adjuvábat me.

Secúndum multitúdinem dolórum meórum in corde meo: * consolatiónes tuæ lætificavérunt ánimam meam.

Numquid adhæret tibi sedes iniquitátis: * qui fingis labórem in præcépto?

Captábunt in ánimam justi: * et sánguinem innocéntem condemnábunt.

Et factus est mihi Dóminus in refúgium: * et Deus meus in adjutórium spei meæ.

Et reddet illis iniquitátem ipsórum: et in malítia eórum dispérdet eos: * dispérdet illos Dóminus, Deus noster.

\psalmheading{Psalmus 107}

\lettrine[lines=2]{P}{ARATUM} cor meum, Deus, parátum cor meum: * cantábo, et psallam in glória mea.

Exsúrge, glória mea, exsúrge, psaltérium et cíthara: * exsúrgam dilúculo.

Confitébor tibi in pópulis, Dómine: * et psallam tibi in natiónibus.

Quia magna est super cælos misericórdia tua: * et usque ad nubes véritas tua:

Exaltáre super cælos, Deus, et super omnem terram glória tua: * ut liberéntur dilécti tui.

Salvum fac déxtera tua, et exáudi me: * Deus locútus est in sancto suo:

Exsultábo, et dívidam Síchimam, * et convállem tabernaculórum dimétiar.

Meus est Gálaad, et meus est Manásses: * et Ephraïm suscéptio cápitis mei.

Juda rex meus: * Moab lebes spei meæ.

In Idumæam exténdam calceaméntum meum: * mihi alienígenæ amíci facti sunt.

Quis dedúcet me in civitátem munítam? * quis dedúcet me usque in Idumæam?

Nonne tu, Deus, qui repulísti nos, * et non exíbis, Deus, in virtútibus nostris?

Da nobis auxílium de tribulatióne: * quia vana salus hóminis.

In Deo faciémus virtútem: * et ipse ad níhilum dedúcet inimícos nostros.

\rubric{¶ Quando ad Laudes dictus fúerit Ps. 50 Miserére, hic subjúngitur sequens Ps. 149, qui in Láudibus locum non hábuit.}

\psalmheading{Psalmus 149}

\lettrine[lines=2]{C}{ANTATE} Dómino cánticum novum: * laus ejus in ecclésia sanctórum.

Lætétur Israël in eo, qui fecit eum: * et fílii Sion exsúltent in rege suo.

Laudent nomen ejus in choro: * in týmpano, et psaltério psallant ei:

Quia beneplácitum est Dómino in pópulo suo: * et exaltábit mansuétos in salútem.

Exsultábunt sancti in glória: * lætabúntur in cubílibus suis.

Exaltatiónes Dei in gútture eórum: * et gládii ancípites in mánibus eórum.

Ad faciéndam vindíctam in natiónibus: * increpatiónes in pópulis.

Ad alligándos reges eórum in compédibus: * et nóbiles eórum in mánicis férreis.

Ut fáciant in eis judícium conscríptum: * glória hæc est ómnibus sanctis ejus.

\rubric{Secus vero, tribus Psalmis explétis, statim dícitur Ant.}

\antlabel{Ant.} Exaltáre, Dómine, qui júdicas terram: redde retributiónem supérbis.

\rubric{In Officio Festi vel Octávæ, Dominicæ anticipátæ ac sanctæ Maríæ in Sabbato:}

\capitulumlabel{Capitulum}{1 Tim. 1, 17}

\lettrine[lines=2]{R}{EGI} sæculórum immortáli et invisíbili, soli Deo honor et glória in sæcula sæculórum. Amen.

\rubric{In ómnibus autem Fériis et in Vigíliis commúnibus:}

\capitulumlabel{Capitulum}{Zach. 8, 19}

\lettrine[lines=2]{P}{ACEM} et veritátem dilígite, ait Dóminus omnípotens.

\R\ br. Christe, Fili Dei vivi, * Miserére nobis. Christe. \V\ \textit{Qui sedes ad déxteram Patris.} Miserére nobis. Glória Patri. Christe. \V\ Exsúrge, Christe, ádjuva nos. \R\ Et líbera nos propter nomen tuum.

\rubric{In Officio sanctæ Maríæ in Sabbato dícitur sequens}

\variantheading{Preces dominicales et feriales (si dicendæ sint)}

KYRIE, eléison. Christe, eléison. Kýrie, eléison.

Pater noster \rubric{secreto usque ad} \V\ Et ne nos indúcas in tentatiónem. \R\ Sed líbera nos a malo.

Credo in Deum \rubric{secreto usque ad} \V\ Carnis resurrectiónem. \R\ Vitam ætérnam. Amen. \V\ Et ego ad te, Dómine, clamávi. \R\ Et mane orátio mea prævéniet te. \V\ Repleátur os meum laude. \R\ Ut cantem glóriam tuam, tota die magnitúdinem tuam. \V\ Dómine, avérte fáciem tuam a peccátis meis. \R\ Et omnes iniquitátes meas dele. \V\ Cor mundum crea in me, Deus. \R\ Et spíritum rectum ínnova in viscéribus meis. \V\ Ne projícias me a fácie tua. \R\ Et spíritum sanctum tuum ne áuferas a me. \V\ Redde mihi lætítiam salutáris tui. \R\ Et spíritu principáli confírma me.

\variantheading{In Précibus feriálibus tantum}

\V\ Eripe me, Dómine, ab hómine malo. \R\ A viro iníquo éripe me. \V\ Eripe me de inimícis meis, Deus meus. \R\ Et ab insurgéntibus in me líbera me. \V\ Eripe me de operántibus iniquitátem. \R\ Et de viris sánguinum salva me. \V\ Sic psalmum dicam nómini tuo in sæculum sæculi. \R\ Ut reddam vota mea de die in diem. \V\ Exáudi nos, Deus, salutáris noster. \R\ Spes ómnium fínium terræ, et in mari longe. \V\ Deus, in adjutórium meum inténde. \R\ Dómine, ad adjuvándum me festína. \V\ Sanctus Deus, Sanctus fortis, Sanctus immortális. \R\ Miserére nobis. \V\ Bénedic, ánima mea, Dómino. \R\ Et ómnia, quæ intra me sunt, nómini sancto ejus. \V\ Bénedic, ánima mea, Dómino. \R\ Et noli oblivísci omnes retributiónes ejus. \V\ Qui propitiátur ómnibus iniquitátibus tuis. \R\ Qui sanat omnes infirmitátes tuas. \V\ Qui rédimit de intéritu vitam tuam. \R\ Qui corónat te in misericórdia et miseratiónibus. \V\ Qui replet in bonis desidérium tuum. \R\ Renovábitur ut áquilæ juvéntus tua.

\variantheading{In utrisque Précibus}

\V\ Adjutórium nostrum in nómine Dómini. \R\ Qui fecit cælum et terram.

\formulalabel{Confíteor. Misereátur. Indulgéntiam.}

\V\ Dignáre, Dómine, die isto. \R\ Sine peccáto nos custodíre. \V\ Miserére nostri, Dómine. \R\ Miserére nostri. \V\ Fiat misericórdia tua, Dómine, super nos. \R\ Quemádmodum sperávimus in te. \V\ Dómine, exáudi oratiónem meam. \R\ Et clamor meus ad te véniat.

\V\ Dóminus vobíscum.

Orémus. \oratiolabel{Oratio}

\lettrine[lines=2]{D}{OMINE}, Deus omnípotens, qui ad princípium hujus diéi nos perveníre fecísti: tua nos hódie salva virtúte; ut in hac die ad nullum declinémus peccátum, sed semper ad tuam justítiam faciéndam nostra procédant elóquia, dirigántur cogitatiónes et ópera. Per Dóminum. \R\ Amen.

\V\ Dóminus vobíscum. \V\ Benedicámus Dómino.

\rubric{Deinde in Choro legitur Martyrológium. Postea Hebdomadárius dicit:}

\V\ Pretiósa in conspéctu Dómini. \R\ Mors Sanctórum ejus.

\rubric{Et absolúte, sine} Orémus:

\oratiolabel{Oratio}

\lettrine[lines=2]{S}{ANCTA} María, et ómnes Sancti intercédant pro nobis ad Dóminum, ut nos mereámur ab eo adjuvári et salvári, qui vivit et regnat in sæcula sæculórum. \R\ Amen.

\V\ Deus, in adjutórium meum inténde. \R\ Dómine, ad adjuvándum me festína. \rubric{Et dícitur ter; último additur:} Glória Patri. Sicut erat.

Kýrie, eléison. Christe, eléison. Kýrie, eléison.

Pater noster \rubric{secreto usque ad} \V\ Et ne nos indúcas in tentatiónem. \R\ Sed líbera nos a malo. \V\ Réspice in servos tuos, Dómine, et in ópera tua, et dírige fílios eórum. \R\ Et sit splendor Dómini, Dei nostri, super nos, et ópera mánuum nostrárum dírige super nos, et opus mánuum nostrárum dírige. \V\ Glória Patri. \R\ Sicut erat.

Orémus. \oratiolabel{Oratio}

\lettrine[lines=2]{D}{IRIGERE} et sanctificáre, régere et gubernáre dignáre, Dómine Deus, Rex cæli et terræ, hódie corda et córpora nostra, sensus, sermónes et actus nostros in lege tua, et in opéribus mandatórum tuórum: ut hic, et in ætérnum, te auxiliánte, salvi et líberi esse mereámur, Salvátor mundi: Qui vivis et regnas in sæcula sæculórum. \R\ Amen.

\V\ Jube, domne.

\benedlabel{Bened.} Dies et actus nostros in sua pace dispónat Dóminus omnípotens. \R\ Amen.

\variantheading{In Officio feriali}

\lettrine[lines=2]{D}{OMINUS} autem dírigat corda et córpora nostra in caritáte Dei, et patiéntia Christi. Tu autem.

\variantheading{E Proprio Sanctorum}

\variantheading{Pro Oct. Omn. Sanctorum}

\lettrine[lines=2]{B}{ENEDÍCTIO} et cláritas et sapiéntia et gratiárum áctio, honor, virtus et fortitúdo Deo nostro in sǽcula sæculórum. Amen.

\variantheading{E Communi Sanctorum}

\variantheading{Pro Apostolis}

\lettrine[lines=2]{I}{BANT} Apóstoli gaudéntes a conspéctu concílii, quóniam digni hábiti sunt pro nómine Jesu contuméliam pati.

\variantheading{Pro Martyre}

\lettrine[lines=2]{J}{USTUS} cor suum trádidit ad vigilándum dilúculo ad Dóminum, qui fecit illum, et in conspéctu Altíssimi deprecábitur.

\variantheading{Pro pluribus Martyribus}

\lettrine[lines=2]{F}{ULGEBUNT} justi, et tamquam scintíllæ in arundinéto discúrrent. Judicábunt natiónes, et dominabúntur pópulis: et regnábit Dóminus illórum in perpétuum.

\variantheading{Pro Conf. Pont.}

\lettrine[lines=2]{F}{UNGI} sacerdótio, et habére laudem in nómine ipsíus, et offérre illi incénsum dignum in odórem suavitátis.

\variantheading{Pro Conf. non Pont.}

\lettrine[lines=2]{J}{USTUM} dedúxit Dóminus per vias rectas, et osténdit illi regnum Dei, et dedit illi sciéntiam sanctórum; honestávit illum in labóribus, et complévit labóres illíus.

\variantheading{Pro Virgine Martyre}

\lettrine[lines=2]{D}{OMINE}, Deus meus, exaltásti super terram habitatiónem meam, et pro morte defluénte deprecáta sum. Invocávi Dóminum, Patrem Dómini mei, ut non derelínquat me in die tribulatiónis meæ, et in témpore superbórum sine adjutório.

\variantheading{Pro Virgine non Martyre}

\lettrine[lines=2]{O}{} \textsc{quam} pulchra est casta generátio cum claritáte! immortális est enim memória illíus: quóniam apud Deum nota est, et apud hómines.

\variantheading{Pro Martyre tantum}

\lettrine[lines=2]{L}{AUDABIT} usque ad mortem ánima mea Dóminum, quóniam éruis sustinéntes te, et líberas eos de manu angústiæ, Dómine, Deus noster.

\variantheading{Pro nec Virg. nec Mart.}

\lettrine[lines=2]{M}{ULTAE} fíliæ congregavérunt divítias: tu supergréssa es univérsas. Fallax grátia, et vana est pulchritúdo: múlier timens Dóminum, ipsa laudábitur.

\variantheading{Pro Dedicatione Ecclesiæ}

\lettrine[lines=2]{E}{T} abstérget Deus omnem lácrimam ab óculis eórum: et mors ultra non erit, neque luctus, neque clamor, neque dolor erit ultra, quia prima abiérunt. Et dixit qui sedébat in throno: Ecce, nova fácio ómnia.

\variantheading{In Officiis B. Mariæ V.}

\lettrine[lines=2]{I}{N} platéis sicut cinnamómum et bálsamum aromatízans odórem dedi: quasi myrrha elécta dedi suavitátem odóris.

\V\ Adjutórium nostrum in nómine Dómini. \R\ Qui fecit cælum et terram.

\V\ Benedícite. \R\ Deus.

\benedictiolabel{Benedictio} Dóminus nos benedícat, et ab omni malo deféndat, et ad vitam perdúcat ætérnam. Et fidélium ánimæ per misericórdiam Dei requiéscant in pace. \R\ Amen.

\rubric{Deinde dícitur tantum} Pater noster \rubric{secréto, nisi sequátur ália Hora.}

% ===== END OF BATCH 39 -- Covers autumnalis.md lines 12676-12954,
% completing "Ad Primam" for Sabbato in full: the hymn "Jam lucis
% orto sídere," Psalms 93/i, 93/ii, 107, the conditional Psalm 149
% insertion, both Capitula (the first now naming "Dominicæ
% anticipátæ" -- the anticipated Sunday -- alongside the usual
% Festum/Octave/Marian-Saturday conditions), the brief Responsory,
% a new standalone rubric specific to this day directing that the
% Preces which follow are said "In Officio sanctæ Maríæ in
% Sabbato," the full Preces apparatus, the closing Orationes, and
% the Lectio-brevis apparatus. Several more psalm incipits and
% Oratio openings carry the partial-caps OCR defect ("DEus
% ultiónum," "BEatus homo," "PAratum cor meum," "CAntáte," "REgi"
% [Capitulum], "DOmine, Deus omnípotens" [Oratio], "SAncta María"
% [Oratio], "DIrígere" [Oratio]) and are silently normalized to
% full caps. The document continues with "Ad Tertiam" (autumnalis.md
% line 12958).
% =====
\hourheading{Ad Tertiam}

Pater noster et Ave.

\V\ Deus, in adjutórium.

\hymnusheading{Hymnus}

\lettrine[lines=2]{N}{UNC}, Sancte, nobis, Spíritus,\\
Unum Patri cum Fílio,\\
Dignáre promptus íngeri\\
Nostro refúsus péctori.

Os, lingua, mens, sensus, vigor\\
Confessiónem pérsonent,\\
Flamméscat igne cáritas,\\
Accéndat ardor próximos.

Præsta, Pater piíssime,\\
Patríque compar Unice,\\
Cum Spíritu Paráclito\\
Regnans per omne sæculum.\\
Amen.

\antlabel{Ant.} Clamor meus.

\psalmheading{Psalmus 101, i}

\lettrine[lines=2]{D}{OMINE}, exáudi oratiónem meam: * et clamor meus ad te véniat.

Non avértas fáciem tuam a me: * in quacúmque die tríbulor, inclína ad me aurem tuam.

In quacúmque die invocávero te, * velóciter exáudi me.

Quia defecérunt sicut fumus dies mei: * et ossa mea sicut crémium aruérunt.

Percússus sum ut fœnum, et áruit cor meum: * quia oblítus sum comédere panem meum.

A voce gémitus mei * adhæsit os meum carni meæ.

Símilis factus sum pellicáno solitúdinis: * factus sum sicut nyctícorax in domicílio.

Vigiláví, * et factus sum sicut passer solitárius in tecto.

Tota die exprobrábant mihi inimíci mei: * et qui laudábant me, advérsum me jurábant.

Quia cínerem tamquam panem manducábam, * et potum meum cum fletu miscébam.

A fácie iræ et indignatiónis tuæ: * quia élevans allisísti me.

Dies mei sicut umbra declinavérunt: * et ego sicut fœnum árui.

Tu autem, Dómine, in ætérnum pérmanes: * et memoriále tuum in generatiónem et generatiónem.

\psalmheading{Psalmus 101, ii}

\lettrine[lines=2]{T}{U} exsúrgens miseréberis Sion: * quia tempus miseréndi ejus, quia venit tempus.

Quóniam placuérunt servis tuis lápides ejus: * et terræ ejus miserebúntur.

Et timébunt gentes nomen tuum, Dómine, * et omnes reges terræ glóriam tuam.

Quia ædificávit Dóminus Sion: * et vidébitur in glória sua.

Respéxit in oratiónem humílium: * et non sprevit precem eórum.

Scribántur hæc in generatióne áltera: * et pópulus, qui creábitur, laudábit Dóminum:

Quia prospéxit de excélso sancto suo: * Dóminus de cælo in terram aspéxit:

Ut audíret gémitus compeditórum: * ut sólveret fílios interemptórum:

Ut annúntient in Sion nomen Dómini: * et laudem ejus in Jerúsalem.

In conveniéndo pópulos in unum, * et reges ut sérviant Dómino.

\psalmheading{Psalmus 101, iii}

\lettrine[lines=2]{R}{ESPONDIT} ei in via virtútis suæ: * Paucitátem diérum meórum núntia mihi.

Ne révoces me in dimídio diérum meórum: * in generatiónem et generatiónem anni tui.

Inítio tu, Dómine, terram fundásti: * et ópera mánuum tuárum sunt cæli.

Ipsi períbunt, tu autem pérmanes: * et omnes sicut vestiméntum veteráscent.

Et sicut opertórium mutábis eos, et mutabúntur: * tu autem idem ipse es, et anni tui non defícient.

Fílii servórum tuórum habitábunt: * et semen eórum in sæculum dirigétur.

\antlabel{Ant.} Clamor meus, Dómine, ad te pervéniat: * non avértas fáciem tuam a me.

\variantheading{In Officio feriali}

\capitulumlabel{Capitulum}{Jer. 17, 14}

\lettrine[lines=2]{S}{ANA} me, Dómine, et sanábor: salvum me fac, et salvus ero: quóniam laus mea tu es.

\R\ br. Sana ánimam meam, * Quia peccávi tibi. Sana. \V\ Ego dixi: Dómine, miserére mei. Quia. Glória Patri. Sana. \V\ Adjútor meus esto, ne derelínquas me. \R\ Neque despícias me, Deus, salutáris meus.

\variantheading{In Vigiliis communibus:}

\variantheading{Preces feriales}

KYRIE, eléison. Christe, eléison. Kýrie, eléison.

Pater noster \rubric{secreto usque ad} \V\ Et ne nos indúcas in tentatiónem. \R\ Sed líbera nos a malo. \V\ Dómine, Deus virtútum, convérte nos. \R\ Et osténde fáciem tuam, et salvi érimus. \V\ Exsúrge, Christe, ádjuva nos. \R\ Et líbera nos propter nomen tuum. \V\ Dómine, exáudi oratiónem meam. \R\ Et clamor meus ad te véniat.

\rubric{Réliqua ut in Ordinário. 19}

\hourheading{Ad Sextam}

Pater noster et Ave.

\V\ Deus, in adjutórium.

\hymnusheading{Hymnus}

\lettrine[lines=2]{R}{ECTOR} potens, verax Deus,\\
Qui témperas rerum vices,\\
Splendóre mane illúminas,\\
Et ígnibus merídiem:

Exstíngue flammas lítium,\\
Aufer calórem nóxium,\\
Confer salútem córporum,\\
Verámque pacem córdium.

Præsta, Pater piíssime,\\
Patríque compar Unice,\\
Cum Spíritu Paráclito\\
Regnans per omne sæculum.\\
Amen.

\antlabel{Ant.} Dómine, Deus meus.

\psalmheading{Psalmus 103, i}

\lettrine[lines=2]{B}{ENEDIC}, ánima mea, Dómino: * Dómine, Deus meus, magnificátus es vehementer.

Confessiónem, et decórem induísti: * amíctus lúmine sicut vestiménto:

Exténdens cælum sicut pellem: * qui tegis aquis superióra ejus.

Qui ponis nubem ascénsum tuum: * qui ámbulas super pennas ventórum.

Qui facis ángelos tuos, spíritus: * et minístros tuos ignem uréntem.

Qui fundásti terram super stabilitátem suam: * non inclinábitur in sæculum sæculi.

Abýssus, sicut vestiméntum, amíctus ejus: * super montes stabunt aquæ.

Ab increpatióne tua fúgient: * a voce tonítrui tui formidábunt.

Ascéndunt montes: et descéndunt campi * in locum, quem fundásti eis.

Términum posuísti, quem non transgrediéntur: * neque converténtur operíre terram.

Qui emíttis fontes in convállibus: * inter médium móntium pertransíbunt aquæ.

Potábunt omnes béstiæ agri: * exspectábunt ónagri in siti sua.

Super ea volúcres cæli habitábunt: * de médio petrárum dabunt voces.

\psalmheading{Psalmus 103, ii}

\lettrine[lines=2]{R}{IGANS} montes de superióribus suis: * de fructu óperum tuórum satiábitur terra:

Prodúcens fœnum juméntis, * et herbam servitúti hóminum:

Ut edúcas panem de terra: * et vinum lætíficet cor hóminis:

Ut exhílaret fáciem in óleo: * et panis cor hóminis confírmet.

Saturabúntur ligna campi, et cedri Líbani, quas plantávit: * illic pásseres nidificábunt.

Heródii domus dux est eórum: * montes excélsi cervis: petra refúgium herináciis.

Fecit lunam in témpora: * sol cognóvit occásum suum.

Posuísti ténebras, et facta est nox: * in ipsa pertransíbunt omnes béstiæ silvæ.

Cátuli leónum rugiéntes, ut rápiant, * et quærant a Deo escam sibi.

Ortus est sol, et congregáti sunt: * et in cubílibus suis collocabúntur.

Exíbit homo ad opus suum: * et ad operatiónem suam usque ad vésperum.

\psalmheading{Psalmus 103, iii}

\lettrine[lines=2]{Q}{UAM} magnificáta sunt ópera tua, Dómine! * ómnia in sapiéntia fecísti: impléta est terra possessióne tua.

Hoc mare magnum, et spatiósum mánibus: * illic reptília, quorum non est númerus.

Animália pusílla cum magnis: * illic naves pertransíbunt.

Draco iste, quem formásti ad illudéndum ei: * ómnia a te exspéctant ut des illis escam in témpore.

Dante te illis, cólligent: * aperiénte te manum tuam, ómnia implebúntur bonitáte.

Avertánte autem te fáciem, turbabúntur: * áuferes spíritum eórum, et defícient, et in púlverem suum reverténtur.

Emíttes spíritum tuum, et creabúntur: * et renovábis fáciem terræ.

Sit glória Dómini in sæculum: * lætábitur Dóminus in opéribus suis:

Qui réspicit terram, et facit eam trémere: * qui tangit montes, et fúmigant.

Cantábo Dómino in vita mea: * psallam Deo meo, quámdiu sum.

Jucúndum sit ei elóquium meum: * ego vero delectábor in Dómino.

Defíciant peccatóres a terra, et iníqui ita ut non sint: * bénedic, ánima mea, Dómino.

\antlabel{Ant.} Dómine, Deus meus, magnificátus es veheménter.

\variantheading{In Officio feriali}

\capitulumlabel{Capitulum}{Rom. 13, 8}

\lettrine[lines=2]{N}{EMINI} quidquam debeátis, nisi ut ínvicem diligátis: qui enim díligit próximum, legem implévit.

\R\ br. Benedícam Dóminum * In omni témpore. Benedícam. \V\ Semper laus ejus in ore meo. In. Glória Patri. Benedícam. \V\ Dóminus regit me, et nihil mihi déerit. \R\ In loco páscuæ ibi me collocávit.

\variantheading{In Vigiliis communibus:}

\variantheading{Preces feriales}

KYRIE, eléison. Christe, eléison. Kýrie, eléison.

Pater noster \rubric{secreto usque ad} \V\ Et ne nos indúcas in tentatiónem. \R\ Sed líbera nos a malo. \V\ Dómine, Deus virtútum, convérte nos. \R\ Et osténde fáciem tuam, et salvi érimus. \V\ Exsúrge, Christe, ádjuva nos. \R\ Et líbera nos propter nomen tuum. \V\ Dómine, exáudi oratiónem meam. \R\ Et clamor meus ad te véniat.

\rubric{Réliqua ut in Ordinário. 20}

\hourheading{Ad Nonam}

Pater noster et Ave.

\V\ Deus, in adjutórium.

\hymnusheading{Hymnus}

\lettrine[lines=2]{R}{ERUM}, Deus, tenax vigor,\\
Immótus in te pérmanens,\\
Lucis diúrnæ témpora\\
Succéssibus detérminans:

Largíre lumen véspere,\\
Quo vita nusquam décidat,\\
Sed præmium mortis sacræ\\
Perénnis instet glória.

Præsta, Pater piíssime,\\
Patríque compar Unice,\\
Cum Spíritu Paráclito\\
Regnans per omne sæculum.\\
Amen.

\antlabel{Ant.} Ne tacúeris, Deus.

\psalmheading{Psalmus 108, i}

\lettrine[lines=2]{D}{EUS}, laudem meam ne tacúeris: * quia os peccatóris, et os dolósi super me apértum est.

Locúti sunt advérsum me lingua dolósa, et sermónibus ódii circumdedérunt me: * et expugnavérunt me gratis.

Pro eo ut me dilígerent, detrahébant mihi: * ego autem orábam.

Et posuérunt advérsum me mala pro bonis: * et ódium pro dilectióne mea.

Constítue super eum peccatórem: * et diábolus stet a dextris ejus.

Cum judicátur, éxeat condemnátus: * et orátio ejus fiat in peccátum.

Fiant dies ejus pauci: * et episcopátum ejus accípiat alter.

Fiant fílii ejus órphani: * et uxor ejus vídua.

Nutántes transferántur fílii ejus, et mendícent: * et ejiciántur de habitatiónibus suis.

Scrutétur fænerátor omnem substántiam ejus: * et dirípiant aliéni labóres ejus.

Non sit illi adjútor: * nec sit qui misereátur pupíllis ejus.

Fiant nati ejus in intéritum: * in generatióne una deleátur nomen ejus.

\psalmheading{Psalmus 108, ii}

\lettrine[lines=2]{I}{N} memóriam rédeat iníquitas patrum ejus in conspéctu Dómini: * et peccátum matris ejus non deleátur.

Fiant contra Dóminum semper, et dispéreat de terra memória eórum: * pro eo quod non est recordátus fácere misericórdiam.

Et persecútus est hóminem ínopem, et mendícum, * et compúnctum corde mortificáre.

Et diléxit maledictiónem, et véniet ei: * et nóluit benedictiónem, et elongábitur ab eo.

Et índuit maledictiónem sicut vestiméntum, * et intrávit sicut aqua in interióra ejus, et sicut óleum in óssibus ejus.

Fiat ei sicut vestiméntum, quo operítur: * et sicut zona, qua semper præcíngitur.

Hoc opus eórum, qui détrahunt mihi apud Dóminum: * et qui loquúntur mala advérsum ánimam meam.

Et tu, Dómine, Dómine, fac mecum propter nomen tuum: * quia suávis est misericórdia tua.

\psalmheading{Psalmus 108, iii}

\lettrine[lines=2]{L}{IBERA} me quia egénus, et pauper ego sum: * et cor meum conturbátum est intra me.

Sicut umbra cum declínat, ablátus sum: * et excússus sum sicut locústæ.

Génua mea infirmáta sunt a jejúnio: * et caro mea immutáta est propter óleum.

Et ego factus sum oppróbrium illis: * vidérunt me, et movérunt cápita sua.

Adjuva me, Dómine, Deus meus: * salvum me fac secúndum misericórdiam tuam.

Et sciant quia manus tua hæc: * et tu, Dómine, fecísti eam.

Maledícent illi, et tu benedíces: * qui insúrgunt in me, confundántur: servus autem tuus lætábitur.

Induántur qui détrahunt mihi, pudóre: * et operiántur sicut diplóide confusióne sua.

Confitébor Dómino nimis in ore meo: * et in médio multórum laudábo eum.

Quia ástitit a dextris páuperis, * ut salvam fáciat a persequéntibus ánimam meam.

\antlabel{Ant.} Ne tacúeris, Deus, quia sermónibus ódii circumdedérunt me.

\variantheading{In Officio feriali}

\capitulumlabel{Capitulum}{1 Petri 1, 17-19}

\lettrine[lines=2]{I}{N} timóre incolátus vestri témpore conversámini: sciéntes quod non corruptibílibus auro vel argénto redémpti estis, sed pretióso sánguine quasi Agni immaculáti Christi.

\R\ br. Rédime me, Dómine, * Et miserére mei. Rédime. \V\ Pes enim meus stetit in via recta. Et. Glória Patri. Rédime.

\V\ Ab occúltis meis munda me, Dómine. \R\ Et ab aliénis parce servo tuo.

\variantheading{In Vigiliis communibus:}

\variantheading{Preces feriales}

KYRIE, eléison. Christe, eléison. Kýrie, eléison.

Pater noster \rubric{secreto usque ad} \V\ Et ne nos indúcas in tentatiónem. \R\ Sed líbera nos a malo. \V\ Dómine, Deus virtútum, convérte nos. \R\ Et osténde fáciem tuam, et salvi érimus. \V\ Exsúrge, Christe, ádjuva nos. \R\ Et líbera nos propter nomen tuum. \V\ Dómine, exáudi oratiónem meam. \R\ Et clamor meus ad te véniat.

\rubric{Réliqua ut in Ordinário. 22}

% ===== END OF BATCH 40 -- Covers autumnalis.md lines 12958-13309,
% completing the Little Hours of Sabbato in full: "Ad Tertiam"
% (Psalms 101/i-iii), "Ad Sextam" (Psalms 103/i-iii), and "Ad
% Nonam" (Psalms 108/i-iii), each with its own hymn, Capitulum,
% brief Responsory, and "In Vigiliis communibus: Preces feriales"
% apparatus. Several more partial-caps OCR defects ("REspóndit,"
% "BEnedic," "RIgans," "QUam magnificáta," "DEus, laudem," "IN
% memóriam," "LIbera me") are silently normalized to full caps.
% The document continues with "Ad Vesperas" (autumnalis.md line
% 13312), which carries a special rubric about this Saturday's
% Vespers serving as First Vespers of the following Sunday.
% =====
\hourheading{Ad Vesperas}

Pater noster et Ave.

\V\ Deus, in adjutórium.

\antlabel{Ant.} Benedíctus Dóminus. \textdagger

\rubric{¶ Antíphonæ cum Psalmis assignátæ semper in Sábbato dicúntur pro I Vésperis Domínicæ sequéntis, quando agéndum est Offícium de ipsa Domínica, in Domínicis minóribus, extra Octávas Sanctíssimi Córporis Christi ac Sacratíssimi Cordis Jesu.}

\psalmheading{Psalmus 143, i}

\lettrine[lines=2]{B}{ENEDICTUS} Dóminus, \textdagger\ Deus meus, qui docet manus meas ad prœlium, * et dígitos meos ad bellum.

Misericórdia mea, et refúgium meum: * suscéptor meus, et liberátor meus:

Protéctor meus, et in ipso sperávi: * qui subdit pópulum meum sub me.

Dómine, quid est homo, quia innotuísti ei? * aut fílius hóminis, quia réputas eum?

Homo vanitáti símilis factus est: * dies ejus sicut umbra prætéreunt.

Dómine, inclína cælos tuos, et descénde: * tange montes, et fumigábunt.

Fúlgura coruscatiónem, et dissipábis eos: * emítte sagíttas tuas, et conturbábis eos.

Emítte manum tuam de alto, éripe me, et líbera me de aquis multis: * de manu filiórum alienórum.

Quorum os locútum est vanitátem: * et déxtera eórum, déxtera iniquitátis.

\antlabel{Ant.} Benedíctus Dóminus, suscéptor meus, et liberátor meus.

\antlabel{Ant.} Beátus pópulus.

\psalmheading{Psalmus 143, ii}

\lettrine[lines=2]{D}{EUS}, cánticum novum cantábo tibi: * in psaltério decachórdo psallam tibi.

Qui das salútem régibus: * qui redemísti David, servum tuum, de gládio malígno: éripe me.

Et érue me de manu filiórum alienórum, quorum os locútum est vanitátem: * et déxtera eórum, déxtera iniquitátis:

Quorum fílii, sicut novéllæ plantatiónes * in juventúte sua.

Fíliæ eórum compósitæ: * circumornátæ ut similitúdo templi.

Promptuária eórum plena, * eructántia ex hoc in illud.

Oves eórum fetósæ, abundántes in egréssibus suis: * boves eórum crassæ.

Non est ruína macériæ, neque tránsitus: * neque clamor in platéis eórum.

Beátum dixérunt pópulum, cui hæc sunt: * beátus pópulus, cujus Dóminus Deus ejus.

\antlabel{Ant.} Beátus pópulus, cujus Dóminus Deus ejus.

\antlabel{Ant.} Magnus Dóminus.

\psalmheading{Psalmus 144, i}

\lettrine[lines=2]{E}{XALTABO} te, Deus meus, rex: * et benedícam nómini tuo in sæculum, et in sæculum sæculi.

Per síngulos dies benedícam tibi: * et laudábo nomen tuum in sæculum, et in sæculum sæculi.

Magnus Dóminus, et laudábilis nimis: * et magnitúdinis ejus non est finis.

Generátio et generátio laudábit ópera tua: * et poténtiam tuam pronuntiábunt.

Magnificéntiam glóriæ sanctitátis tuæ loquéntur: * et mirabília tua narrábunt.

Et virtútem terribílium tuórum dicent: * et magnitúdinem tuam narrábunt.

Memóriam abundántiæ suavitátis tuæ eructábunt: * et justítia tua exsultábunt.

\antlabel{Ant.} Magnus Dóminus, et laudábilis nimis: et magnitúdinis ejus non est finis.

\antlabel{Ant.} Suávis Dóminus.

\psalmheading{Psalmus 144, ii}

\lettrine[lines=2]{M}{ISERATOR}, et miséricors Dóminus: * pátiens, et multum miséricors.

Suávis Dóminus univérsis: * et miseratiónes ejus super ómnia ópera ejus.

Confiteántur tibi, Dómine, ómnia ópera tua: * et sancti tui benedícant tibi.

Glóriam regni tui dicent: * et poténtiam tuam loquéntur:

Ut notam fáciant fíliis hóminum poténtiam tuam: * et glóriam magnificéntiæ regni tui.

Regnum tuum regnum ómnium sæculórum: * et dominátio tua in omni generatióne et generatiónem.

\antlabel{Ant.} Suávis Dóminus univérsis: et miseratiónes ejus super ómnia ópera ejus.

\antlabel{Ant.} Fidélis Dóminus. \textdagger

\psalmheading{Psalmus 144, iii}

\lettrine[lines=2]{F}{IDELIS} Dóminus \textdagger\ in ómnibus verbis suis: * et sanctus in ómnibus opéribus suis.

Allevat Dóminus omnes qui córruunt: * et érigit omnes elísos.

Óculi ómnium in te sperant, Dómine: * et tu das escam illórum in témpore opportúno.

Áperis tu manum tuam: * et imples omne ánimal benedictióne.

Justus Dóminus in ómnibus viis suis: * et sanctus in ómnibus opéribus suis.

Prope est Dóminus ómnibus invocántibus eum: * ómnibus invocántibus eum in veritáte.

Voluntátem timéntium se fáciet: * et deprecatiónem eórum exáudiet: et salvos fáciet eos.

Custódit Dóminus omnes diligéntes se: * et omnes peccatóres dispérdet.

Laudatiónem Dómini loquétur os meum: * et benedícat omnis caro nómini sancto ejus in sæculum, et in sæculum sæculi.

\antlabel{Ant.} Fidélis Dóminus in ómnibus verbis suis: et sanctus in ómnibus opéribus suis.

\rubric{In Sabbatis extra Octávas Ssmi Córporis Christi ac Sacratíssimi Cordis Jesu dícitur}

\variantheading{In Officio feriali}

\capitulumlabel{Capitulum}{Rom. 11, 33}

\lettrine[lines=2]{O}{} \textsc{altitudo} divitiárum sapiéntiæ et sciéntiæ Dei: quam incomprehensíbilia sunt judícia ejus, et investigábiles viæ ejus!

\hymnusheading{Hymnus}

\lettrine[lines=2]{J}{AM} sol recédit ígneus:\\
Tu, lux perénnis, Únitas,\\
Nostris, beáta Trínitas,\\
Infúnde lumen córdibus.

Te mane laudum cármine,\\
Te deprecámur véspere;\\
Dignéris ut te súpplices\\
Laudémus inter Cælites.

Patri, simúlque Fílio,\\
Tibíque, Sancte Spíritus,\\
Sicut fuit, sit júgiter\\
Sæclum per omne glória.\\
Amen.

\V\ Vespertína orátio ascéndat ad te, Dómine. \R\ Et descéndat super nos misericórdia tua.

\rubric{Ad Magníficat Antíphona ut in Próprio de Témpore.}

\rubric{Expléto Cántico ac repetíta Antíphona, dícitur Orátio Domínicæ sequéntis própria, quæ per totam hebdómadam suam in feriáli Officio adhibétur, nisi speciális Orátio habeátur.}

\rubric{Réliqua ut in Ordinário. 22}

\hourheading{Ad Completorium}

\V\ Jube, domne.

\benedlabel{Bened.} Noctem quiétam, et finem perféctum concédat nobis Dóminus omnípotens. \R\ Amen.

\lectioheading{Lectio br.}{1 Petri 5, 8-9}

\lettrine[lines=2]{F}{RATRES}: Sóbrii estóte, et vigiláte: quia adversárius vester diábolus tamquam leo rúgiens círcuit, quærens quem dévoret: cui resístite fortes in fide. Tu autem.

\R\ Deo grátias.

\V\ Adjutórium nostrum in nómine Dómini. \R\ Qui fecit cælum et terram.

Pater noster, \rubric{quod dícitur totum secréto.}

\formulalabel{Confíteor.} Misereátur. Indulgéntiam.

\V\ Convérte nos, Deus, salutáris noster. \R\ Et avérte iram tuam a nobis.

\V\ Deus, in adjutórium.

\antlabel{Ant.} Intret orátio mea.

\psalmheading{Psalmus 87}

\lettrine[lines=2]{D}{OMINE}, Deus salútis meæ: * in die clamávi, et nocte coram te.

Intret in conspéctu tuo orátio mea: * inclína aurem tuam ad precem meam:

Quia repléta est malis ánima mea: * et vita mea inférno appropinquávit.

Æstimátus sum cum descendéntibus in lacum: * factus sum sicut homo sine adjutório, inter mórtuos liber.

Sicut vulneráti dormiéntes in sepúlcris, quorum non es memor ámplius: * et ipsi de manu tua repúlsi sunt.

Posuérunt me in lacu inferióri: * in tenebrósis, et in umbra mortis.

Super me confirmátus est furor tuus: * et omnes fluctus tuos induxísti super me.

Longe fecísti notos meos a me: * posuérunt me abominatiónem sibi.

Tráditus sum, et non egrediébar: * óculi mei languérunt præ inópia.

Clamávi ad te, Dómine, tota die: * expándi ad te manus meas.

Numquid mórtuis fácies mirabília: * aut médici suscitábunt, et confitebúntur tibi?

Numquid narrábit áliquis in sepúlcro misericórdiam tuam, * et veritátem tuam in perditióne?

Numquid cognoscéntur in ténebris mirabília tua, * et justítia tua in terra obliviónis?

Et ego ad te, Dómine, clamávi: * et mane orátio mea prævéniet te.

Ut quid, Dómine, repéllis oratiónem meam: * avértis fáciem tuam a me?

Pauper sum ego, et in labóribus a juventúte mea: * exaltátus autem, humiliátus sum et conturbátus.

In me transiérunt iræ tuæ: * et terróres tui conturbavérunt me.

Circumdedérunt me sicut aqua tota die: * circumdedérunt me simul.

Elongásti a me amícum et próximum: * et notos meos a miséria.

\psalmheading{Psalmus 102, i}

\lettrine[lines=2]{B}{ENEDIC}, ánima mea, Dómino: * et ómnia, quæ intra me sunt, nómini sancto ejus.

Bénedic, ánima mea, Dómino: * et noli oblivísci omnes retributiónes ejus.

Qui propitiátur ómnibus iniquitátibus tuis: * qui sanat omnes infirmitátes tuas.

Qui rédimit de intéritu vitam tuam: * qui corónat te in misericórdia et miseratiónibus.

Qui replet in bonis desidérium tuum: * renovábitur ut áquilæ juvéntus tua:

Fáciens misericórdias Dóminus: * et judícium ómnibus injúriam patiéntibus.

Notas fecit vias suas Móysi, * fíliis Israël voluntátes suas.

Miserátor, et miséricors Dóminus: * longánimis, et multum miséricors.

Non in perpétuum irascétur: * neque in ætérnum comminábitur.

Non secúndum peccáta nostra fecit nobis: * neque secúndum iniquitátes nostras retríbuit nobis.

Quóniam secúndum altitúdinem cæli a terra: * corroborávit misericórdiam suam super timéntes se.

Quantum distat ortus ab occidénte: * longe fecit a nobis iniquitátes nostras.

\psalmheading{Psalmus 102, ii}

\lettrine[lines=2]{Q}{UOMODO} miserétur pater filiórum, misértus est Dóminus timéntibus se: * quóniam ipse cognóvit figméntum nostrum.

Recordátus est quóniam pulvis sumus: * homo, sicut fœnum dies ejus, tamquam flos agri sic efflorébit.

Quóniam spíritus pertransíbit in illo, et non subsístet: * et non cognóscet ámplius locum suum.

Misericórdia autem Dómini ab ætérno, * et usque in ætérnum super timéntes eum.

Et justítia illíus in fílios filiórum, * his qui servant testaméntum ejus:

Et mémores sunt mandatórum ipsíus, * ad faciéndum ea.

Dóminus in cælo parávit sedem suam: * et regnum ipsíus ómnibus dominábitur.

Benedícite Dómino, omnes Ángeli ejus: * poténtes virtúte, faciéntes verbum illíus, ad audiéndam vocem sermónum ejus.

Benedícite Dómino, omnes virtútes ejus: * minístri ejus, qui fácitis voluntátem ejus.

Benedícite Dómino, ómnia ópera ejus: * in omni loco dominatiónis ejus, benedíc, ánima mea, Dómino.

\antlabel{Ant.} Intret orátio mea in conspéctu tuo, Dómine.

\hymnusheading{Hymnus}

\lettrine[lines=2]{T}{E} lucis ante términum,\\
Rerum Creátor, póscimus,\\
Ut pro tua cleméntia\\
Sis præsul et custódia.

Procul recédant sómnia,\\
Et nóctium phantásmata;\\
Hostémque nostrum cómprime,\\
Ne polluántur córpora.

Præsta, Pater piíssime,\\
Patríque compar Unice,\\
Cum Spíritu Paráclito\\
Regnans per omne sæculum.\\
Amen.

\capitulumlabel{Capitulum}{Jer. 14, 9}

\lettrine[lines=2]{T}{U} autem in nobis es, Dómine, et nomen sanctum tuum invocátum est super nos: ne derelínquas nos, Dómine, Deus noster.

\R\ br. In manus tuas, Dómine, * Comméndo spíritum meum. In manus. \V\ Redemísti nos, Dómine, Deus veritátis. Comméndo. Glória Patri. In manus.

\V\ Custódi nos, Dómine, ut pupíllam óculi. \R\ Sub umbra alárum tuárum prótege nos.

\antlabel{Ant.} Salva nos.

\canticumheading{Canticum Simeonis}{Luc. 2, 29-32}

\lettrine[lines=2]{N}{UNC} dimíttis servum tuum, Dómine, * secúndum verbum tuum in pace:

Quia vidérunt óculi mei * salutáre tuum,

Quod parásti * ante fáciem ómnium populórum,

Lumen ad revelatiónem géntium, * et glóriam plebis tuæ Israël.

\antlabel{Ant.} Salva nos, Dómine, vigilántes, custódi nos dormiéntes; ut vigilémus cum Christo, et requiescámus in pace.

\variantheading{Preces (si dicendæ sint)}

KYRIE, eléison. Christe, eléison. Kýrie, eléison.

Pater noster \rubric{secreto usque ad} \V\ Et ne nos indúcas in tentatiónem. \R\ Sed líbera nos a malo.

Credo in Deum \rubric{secreto usque ad} \V\ Carnis resurrectiónem. \R\ Vitam ætérnam. Amen. \V\ Benedíctus es, Dómine, Deus patrum nostrórum. \R\ Et laudábilis et gloriósus in sæcula. \V\ Benedicámus Patrem et Fílium cum Sancto Spíritu. \R\ Laudémus, et superexaltémus eum in sæcula. \V\ Benedíctus es, Dómine, in firmaménto cæli. \R\ Et laudábilis, et gloriósus, et superexaltátus in sæcula. \V\ Benedícat et custódiat nos omnípotens et miséricors Dóminus. \R\ Amen. \V\ Dignáre, Dómine, nocte ista. \R\ Sine peccáto nos custodíre. \V\ Miserére nostri, Dómine. \R\ Miserére nostri. \V\ Fiat misericórdia tua, Dómine, super nos. \R\ Quemádmodum sperávimus in te. \V\ Dómine, exáudi oratiónem meam. \R\ Et clamor meus ad te véniat.

\V\ Dóminus vobíscum.

Orémus. \oratiolabel{Oratio}

\lettrine[lines=2]{V}{ISITA}, quæsumus, Dómine, habitatiónem istam, et omnes insídias inimíci ab ea longe repélle: Angeli tui sancti hábitent in ea, qui nos in pace custódiant; et benedíctio tua sit super nos semper. Per Dóminum nostrum. \R\ Amen.

\V\ Dóminus vobíscum.

\V\ Benedicámus Dómino.

\benedictiolabel{Benedictio} Benedícat et custódiat nos omnípotens et miséricors Dóminus, Pater, et Fílius, et Spíritus Sanctus. \R\ Amen.

\variantheading{Antiphona finalis B. M. V.}

\V\ Divínum auxílium.

Pater, Ave María et Credo.

% ===== END OF BATCH 41 -- Covers autumnalis.md lines 13312-13647,
% completing "Ad Vesperas" (Psalms 143/i, 143/ii, 144/i-iii, with
% the special rubric noting these serve as First Vespers of the
% following Sunday) and "Ad Completorium" (Psalms 87, 102/i,
% 102/ii, hymn "Te lucis ante términum," Canticum Simeonis, Preces),
% and with them the ENTIRE "Sabbato" Office and the complete
% seven-day ferial Psalterium cycle of the Pars Autumnalis (Dominica
% through Sabbato). Saturday Vespers here carries its own fixed
% "In Officio feriali" Capitulum ("O Altitúdo divitiárum," rendered
% with the single-letter-word lettrine convention since the source
% gives only "O" capitalized before the normally-cased "Altitúdo")
% and hymn ("Jam sol recédit ígneus," with the recurring stray OCR
% "\*" before its doxology silently dropped), unlike the weekday
% Vespers pattern, since this Capitulum/Hymn pair is fixed for every
% Saturday rather than varying with the day. Two rubrics direct that
% the Magnificat antiphon and the following week's Oratio are drawn
% "ut in Próprio de Témpore" -- from the Proper of the Time for the
% upcoming Sunday -- confirming this Vespers functions as First
% Vespers of Sunday. The Canticum Simeonis here omits the printed
% "Glória Patri" line that appears at every other occurrence of this
% canticle in this volume; transcribed exactly as the source omits
% it. The document's Psalterium section ends here; autumnalis.md
% continues at line 13650 with an "[IMAGE]" placeholder (omitted, a
% non-textual OCR marker, preceded by a "---" scan-artifact rule
% also omitted) followed immediately by the heading "Proprium de
% Tempore" at line 13652, opening with "Sabbato proximo Kalendis
% Septembris" (the Saturday nearest the Kalends of September) --
% this volume's own starting point in the liturgical year, genuinely
% different from aestiva.tex's parallel transition into "In Festo
% Ssmæ Trinitatis" (Trinity Sunday), confirming the Proprium de
% Tempore itself is NOT shared boilerplate and must be transcribed
% fresh from here on -- a major structural transition from the
% ferial Psalterium into the Proper of the Time.
% =====
\partheading{Proprium de Tempore}

\feriaheading{Sabbato proximo Kalendis Septembris}

\hourheading{Ad Vesperas}

\V\ Vespertína orátio ascéndat ad te, Dómine. \R\ Et descéndat super nos misericórdia tua.

\antlabel{Ad Magnif. Ant.} Cum audísset Job * nuntiórum verba, sustínuit patiénter, et ait: Si bona suscépimus de manu Dómini, mala autem quare non sustineámus? In ómnibus his non peccávit Job lábiis suis, neque stultum áliquid contra Deum locútus est.

\rubric{Orátio Domínicæ occurréntis, quæ habétur infra post Lectiónes últimæ hebdómadæ Novémbris cum áliis Oratiónibus usque ad Domínicam XXIV post Pentecósten.}

\feriaheading{Dominica I Septembris}

\rubric{Semiduplex}

\hourheading{Ad Matutinum}

\rubric{Invitatórium} Dóminum, qui fecit, \rubric{et Hymnus} Nocte surgéntes, \rubric{ut in Psaltério. 31; quæ dicúntur in ómnibus Domínicis usque ad diem 27 Septémbris inclusíve occurréntibus.}

\nocturnoheading{In I Nocturno}

\rubric{Incipit liber Job}

\lectioheading{Lectio i}{Cap. 1, 1-11}

\lettrine[lines=2]{V}{IR} erat in terra Hus nómine Job, et erat vir ille simplex et rectus ac timens Deum et recédens a malo. Natíque sunt ei septem fílii et tres fíliæ. Et fuit posséssio ejus septem míllia óvium et tria míllia camelórum, quingénta quoque juga boum et quingéntæ ásinæ ac família multa nimis: erátque vir ille magnus inter omnes Orientáles.

\R\ Si bona suscépimus de manu Dei, mala autem quare non sustineámus? * Dóminus dedit, Dóminus ábstulit; sicut Dómino plácuit, ita factum est: sit nomen Dómini benedíctum. \V\ Nudus egréssus sum de útero matris meæ et nudus revértar illuc. — Dóminus.

\lectioheading{Lectio ii}{}

\lettrine[lines=2]{E}{T} ibant fílii ejus et faciébant convívium per domos, unusquísque in die suo. Et mitténtes vocábant tres soróres suas, ut coméderent et bíberent cum eis. Cumque in orbem transíssent dies convívii, mittébat ad eos Job et sanctificábat illos, consurgénsque dilúculo offerébat holocáusta pro síngulis; dicébat enim: Ne forte peccáverint fílii mei et benedíxerint Deo in córdibus suis. Sic faciébat Job cunctis diébus.

\R\ Antequam cómedam suspíro, et tamquam inundántes aquæ sic rugítus meus; quia timor, quem timébam, evénit mihi, et quod verébar áccidit. Nonne dissimulávi? nonne sílui? nonne quiévi? * Et venit super me indignátio. \V\ Ecce non est auxílium mihi in me, et necessárii quoque mei recessérunt a me. — Et.

\lectioheading{Lectio iii}{}

\lettrine[lines=2]{Q}{UADAM} autem die, cum veníssent fílii Dei ut assísterent coram Dómino, affúit inter eos étiam Satan. Cui dixit Dóminus: Unde venis? Qui respóndens ait: Circuívi terram et perambulávi eam. Dixítque Dóminus ad eum: Numquid considerásti servum meum Job, quod non sit ei símilis in terra, homo símplex et rectus ac timens Deum et recédens a malo? Cui respóndens Satan ait: Numquid Job frustra timet Deum? nonne tu vallásti eum ac domum ejus, universámque substántiam per circúitum opéribus mánuum ejus benedixísti, et posséssio ejus crevit in terra? Sed exténde páululum manum tuam et tange cuncta quæ póssidet, nisi in fáciem benedíxerit tibi.

\R\ Quare detraxístis sermónibus veritátis? ad increpándum verba compónitis et subvértere nitímini amícum vestrum: * Verúmtamen quæ cogitástis, expléte. \V\ Quod justum est, judicáte; et non inveniétis in lingua mea iniquitátem. — Verúmtamen. Glória Patri. Verúmtamen.

\nocturnoheading{In II Nocturno}

\rubric{Ex libro Morálium sancti Gregórii Papæ}

\lectioheading{Lectio iv}{Lib. 2, cap. 1}

\lettrine[lines=2]{S}{CRIPTÚRA} sacra mentis óculis quasi quoddam spéculum oppónitur, ut intérna nostra fácies in ipsa videátur. Ibi étenim foeda, ibi pulchra nostra cognóscimus: ibi sentímus quantum profícimus, ibi a proféctu quam longe distámus. Narrat autem gesta Sanctórum, et ad imitatiónem corda próvocat infirmórum; dumque illórum victrícia facta commémorat, contra vitiórum prǽlia, debília nostra confírmat: fitque verbis illíus, ut eo mens minus inter certámina trépidet, quo ante se pósitos tot virórum fórtium triúmphos videt.

\R\ Indúta est caro mea putrédine, et sórdibus púlveris cutis mea áruit et contrácta est: * Meménto mei, Dómine, quóniam ventus est vita mea. \V\ Dies mei velócius transiérunt quam a texénte tela succíditur, et consúmpti sunt absque ulla spe. — Meménto.

\lectioheading{Lectio v}{}

\lettrine[lines=2]{N}{ONNÚMQUAM} vero non solum nobis eórum virtútes ásserit, sed étiam casus innotéscit; ut et in victória fórtium, quod imitándo arrípere, et rursum videámus in lápsibus quid debeámus timére. Ecce enim Job describítur tentatióne auctus, sed David tentatióne prostrátus; ut et majórum virtus spem nostram fóveat, et majórum casus ad cautélam nos humilitátis accíngat: quátenus dum illa gaudéntes súblevant, ista metuéntes premant; et audiéntis ánimus illinc spei fidúcia, hinc humilitáte erudítus, nec temeritáte supérbiat, quia ad spei fidúciam virtútis exémplo roborátur.

\R\ Páucitas diérum meórum finiétur brevi; dimítte me, Dómine, ut plangam páululum dolórem meum, * Antequam vadam et non revértar, ad terram tenebrósam et opértam mortis calígine. \V\ Manus tuæ, Dómine, fecérunt me, et plasmavérunt me totum in circúitu; et sic repénte præcípitas me? — Antequam.

\lectioheading{Lectio vi}{Lib. 1, cap. 1}

\lettrine[lines=2]{V}{IR} erat in terra Hus, nómine Job. Idcírco sanctus vir ubi habitáverit dícitur, ut ejus méritum virtútis exprimátur. Hus namque quis nésciat quod sit in terra Gentílium? Gentílitas autem eo obligáta vítiis exstítit, quo cognitiónem sui Condítoris ignorávit. Dicátur ítaque ubi habitáverit; ut hoc ejus láudibus profíciat, quod bonus inter malos fuit. Neque enim valde laudábile est, bonum esse cum bonis, sed bonum esse cum malis. Sicut enim gravióris culpæ est, inter bonos bonum non esse; ita imménsi est præcónii, bonum étiam inter malos exstitísse.

\R\ Non abscóndas me, Dómine, a fácie tua: manum tuam longe fac a me, * Et formído tua non me térreat. \V\ Córripe me, Dómine, in misericórdia, non in furóre tuo, ne forte ad níhilum rédigas me. — Et. Glória Patri. Et.

\rubric{Lectiónes III Nocturni de Homilía Domínicæ occurréntis post Pentecósten habéntur inférius circa finem hujus Partis de Témpore, in órdine aliárum Dominicárum usque ad Domínicam XXIV post Pentecósten. Simíliter Antíphonæ ad Benedíctus et Magníficat, et Oratiónes.}

\R\ vii Quis mihi tríbuat, ut in inférno prótegas me et abscóndas me, donec pertránseat furor tuus, Dómine, nisi tu, qui solus es Deus? * Et constítuas mihi tempus, in quo recordéris mei? \V\ Numquid sicut dies hóminis dies tui, ut quæras iniquitátem meam; cum sit nemo, qui de manu tua possit erúere? — Et.

\R\ viii Duo Séraphim, \rubric{ut infra in Homilíis.}

\rubric{Ad Laudes, in hac et sequéntibus Domínicis usque ad diem 27 Septémbris inclusíve occurréntibus, dícitur Hymnus} Ecce jam noctis, \rubric{ut in Psaltério. 40}

% ===== BATCH 1 of the Proprium de Tempore (autumnalis.md lines
% 13652-13730) -- genuinely new Autumn-specific content, transcribed
% fresh from autumnalis.md since the Proprium de Tempore is NOT
% shared boilerplate across Pars volumes (confirmed by the previous
% batch's diff work showing aestiva.tex's parallel section opens
% with Trinity Sunday rather than this volume's own starting point).
% Covers the brief "Sabbato proximo Kalendis Septembris" (First
% Vespers only, with its Magnificat antiphon from the Book of Job
% and a forward-pointing rubric to the proper Sunday Oratio given
% later in this Part) and the complete "Dominica I Septembris"
% Matins: the Invitatorium/Hymnus cross-reference (rendered as
% alternating \rubric{} labels and live incipit text per the
% established compact-citation convention), the six Lessons and six
% Responsories of Nocturns I-II (Job's affliction and the opening of
% St. Gregory's Moralia on Job), and the final two numbered
% Responsories (vii, viii) that stand in for Nocturn III, whose own
% Lessons, antiphons, and Orationes are deferred by the source's own
% rubric to the Homiliae section near the end of this Part (a
% structural forward-reference, not a page/footnote number, so no
% \pageref{} resolution applies; it will be transcribed fresh when
% that later section is reached). The em-dash "—" repeat-cue
% notation in each Responsory (e.g. "— Dóminus.") is preserved
% literally as the source's own shorthand for repeating the
% refrain's cue word, a notation not previously encountered in this
% project's established responsory-repetition conventions (which
% elsewhere spell the repeated word out in full); since it is
% genuinely how autumnalis.md itself abbreviates these responsories,
% it is kept exactly as printed rather than silently expanded. The
% "---" horizontal rule before "Dominica I Septembris" is omitted as
% a scan artifact. The document continues with "Feria Secunda infra
% Hebdomadam I Septembris" (autumnalis.md line 13731).
% =====
\feriaheading{Feria Secunda infra Hebdomadam I Septembris}

\rubric{De libro Job}

\lectioheading{Lectio i}{Cap. 1, 13-22}

\lettrine[lines=2]{C}{UM} autem quadam die fílii et fíliæ ejus coméderent et bíberent vinum in domo fratris sui primogéniti, núntius venit ad Job, qui díceret: Boves arábant et ásinæ pascebántur juxta eos, et irruérunt Sabæi tulerúntque ómnia et púeros percussérunt gládio, et evási ego solus ut nuntiárem tibi. Cumque adhuc ille loquerétur, venit alter et dixit: Ignis Dei cécidit e cælo et tactas oves puerósque consúmpsit, et effúgi ego solus ut nuntiárem tibi.

\R\ Versa est in luctum cíthara mea, et órganum meum in vocem fléntium: * Parce mihi, Dómine, nihil enim sunt dies mei. \V\ Cutis mea denigráta est super me, et ossa mea aruérunt. — Parce.

\lectioheading{Lectio ii}{}

\lettrine[lines=2]{S}{ED}, et illo adhuc loquénte, venit álius et dixit: Chaldǽi fecérunt tres turmas et invasérunt camélos et tulérunt eos, necnon et púeros percussérunt gládio, et ego fugi solus ut nuntiárem tibi. Adhuc loquebátur ille, et ecce álius intrávit et dixit: Fíliis tuis et filiábus vescéntibus et bibéntibus vinum in domo fratris sui primogéniti, repénte ventus veheméns írruit a regióne desérti et concússit quátuor ángulos domus, quæ córruens oppréssit líberos tuos, et mórtui sunt, et effúgi ego solus ut nuntiárem tibi.

\R\ Utinam appenderéntur peccáta mea, quibus iram mérui, * Et calamítas, quam pátior, in statéra. \V\ Quasi aréna maris hæc grávior apparéret, unde et verba mea dolóre sunt plena. — Et.

\rubric{In Octávis et Festis ritus simplícis: Glória Patri. Et.}

\lectioheading{Lectio iii}{}

\lettrine[lines=2]{T}{UNC} surréxit Job et scidit vestiménta sua et, tonso cápite, córruens in terram adorávit et dixit: Nudus egréssus sum de útero matris meæ et nudus revértar illuc. Dóminus dedit, Dóminus ábstulit; sicut Dómino plácuit, ita factum est: sit nomen Dómini benedíctum. In ómnibus his non peccávit Job lábiis suis neque stultum quid contra Deum locútus est.

\R\ Quare detraxístis sermónibus veritátis? ad increpándum verba compónitis et subvértere nitímini amícum vestrum: * Verúmtamen quæ cogitástis, expléte. \V\ Quod justum est, judicáte; et non inveniétis in lingua mea iniquitátem. — Verúmtamen. Glória Patri. Verúmtamen.

\feriaheading{Feria Tertia infra Hebdomadam I Septembris}

\rubric{De libro Job}

\lectioheading{Lectio i}{Cap. 2, 1-13}

\lettrine[lines=2]{F}{ACTUM} est autem, cum quadam die veníssent fílii Dei et starent coram Dómino, venísset quoque Satan inter eos et staret in conspéctu ejus, ut díceret Dóminus ad Satan: Unde venis? Qui respóndens ait: Circuívi terram et perambulávi eam. Et dixit Dóminus ad Satan: Numquid considerásti servum meum Job, quod non sit ei símilis in terra, vir símplex et rectus ac timens Deum et recédens a malo et adhuc rétinens innocéntiam? Tu autem commovísti me advérsus eum, ut afflígerem eum frustra. Cui respóndens Satan ait: Pellem pro pelle, et cuncta, quæ habet homo, dabit pro ánima sua. Alióqui mitte manum tuam et tange os ejus et carnem; et tunc vidébis quod in fáciem benedícat tibi.

\R\ Indúta est caro mea putrédine, et sórdibus púlveris cutis mea áruit et contrácta est: * Meménto mei, Dómine, quóniam ventus est vita mea. \V\ Dies mei velócius transiérunt quam a texénte tela succíditur, et consúmpti sunt absque ulla spe. — Meménto.

\lectioheading{Lectio ii}{}

\lettrine[lines=2]{D}{IXIT} ergo Dóminus ad Satan: Ecce in manu tua est, verúmtamen ánimam illíus serva. Egréssus ígitur Satan a fácie Dómini percússit Job úlcere péssimo a planta pedis usque ad vérticem ejus; qui testa sániem radébat sedens in sterquilínio. Dixit autem illi uxor sua: Adhuc tu pérmanes in simplicitáte tua? Bénedic Deo et mórere. Qui ait ad illam: Quasi una de stultis muliéribus locúta es: si bona suscépimus de manu Dei, mala quare non suscipiámus? In ómnibus his non peccávit Job lábiis suis.

\R\ Páucitas diérum meórum finiétur brevi; dimítte me, Dómine, ut plangam páululum dolórem meum, * Antequam vadam ad terram tenebrósam et opértam mortis calígine. \V\ Manus tuæ, Dómine, fecérunt me, et plasmavérunt me totum in circúitu; et sic repénte præcípitas me? — Antequam.

\rubric{In Octávis et Festis ritus simplícis: Glória Patri. Antequam.}

\lectioheading{Lectio iii}{}

\lettrine[lines=2]{I}{GITUR} audiéntes tres amíci Job omne malum quod accidísset ei, venérunt sínguli de loco suo, Elíphaz Themanítes et Báldad Suhítes et Sophar Naamathítes. Condíxerant enim ut páriter veniéntes visitárent eum et consolaréntur. Cumque elevássent procul óculos suos, non cognovérunt eum et exclamántes ploravérunt, scissísque véstibus, spargérunt púlverem super caput suum in cælum. Et sedérunt cum eo in terra septem diébus et septem nóctibus, et nemo loquebátur ei verbum; vidébant enim dolórem esse veheméntem.

\R\ Non abscóndas me, Dómine, a fácie tua: manum tuam longe fac a me, * Et formído tua non me térreat. \V\ Córripe me, Dómine, in misericórdia, non in furóre tuo, ne forte ad níhilum rédigas me. — Et. Glória Patri. Et.

\feriaheading{Feria Quarta infra Hebdomadam I Septembris}

\rubric{De libro Job}

\lectioheading{Lectio i}{Cap. 3, 1-16}

\lettrine[lines=2]{P}{OST} hæc apéruit Job os suum et maledíxit diéi suo et locútus est: Péreat dies in qua natus sum, et nox in qua dictum est: Concéptus est homo. Dies ille vertátur in ténebras: non requírat eum Deus désuper, et non illustrétur lúmine; obscúrent eum ténebræ et umbra mortis; óccupet eum calígo, et involvátur amaritúdine.

\R\ Quis mihi tríbuat, ut in inférno prótegas me et abscóndas me, donec pertránseat furor tuus, Dómine, nisi tu, qui solus es Deus? * Et constítuas mihi tempus, in quo recordéris mei? \V\ Numquid sicut dies hóminis dies tui, ut quæras iniquitátem meam; cum sit nemo, qui de manu tua possit erúere? — Et.

\lectioheading{Lectio ii}{}

\lettrine[lines=2]{N}{OCTEM} illam tenebrósus turbo possídeat; non computétur in diébus anni, nec numerétur in ménsibus. Sit nox illa solitária nec laude digna; maledícant ei qui maledícunt diéi, qui paráti sunt suscitáre leviáthan. Obtenebréntur stellæ calígine ejus, exspéctet lucem et non vídeat, nec ortum surgéntis auróræ; quia non conclúsit óstia ventris, qui portávit me, nec ábstulit mala ab óculis meis.

\R\ Utinam appenderéntur peccáta mea, quibus iram mérui, * Et calamítas, quam pátior, in statéra. \V\ Quasi aréna maris hæc grávior apparéret, unde et verba mea dolóre sunt plena. — Et.

\rubric{In Octávis et Festis ritus simplícis: Glória Patri. Et.}

\lectioheading{Lectio iii}{}

\lettrine[lines=2]{Q}{UARE} non in vulva mórtuus sum? egréssus ex útero non statim périi? quare excéptus génibus? cur lactátus ubéribus? Nunc enim dórmiens silérem et somno meo requiéscerem cum régibus et consúlibus terræ, qui ædíficant sibi solitúdines, aut cum princípibus, qui póssident aurum et replent domos suas argénto; aut sicut abortívum abscónditum non subsísterem, vel qui concépti non vidérunt lucem.

\R\ Quare detraxístis sermónibus veritátis? ad increpándum verba compónitis et subvértere nitímini amícum vestrum: * Verúmtamen quæ cogitástis, expléte. \V\ Quod justum est, judicáte; et non inveniétis in lingua mea iniquitátem. — Verúmtamen. Glória Patri. Verúmtamen.

\feriaheading{Feria Quinta infra Hebdomadam I Septembris}

\rubric{De libro Job}

\lectioheading{Lectio i}{Cap. 4, 1-18}

\lettrine[lines=2]{R}{ESPÓNDENS} autem Elíphaz Themanítes dixit: Si cœpérimus loqui tibi, fórsitan moléste accípies; sed concéptum sermónem tenére quis póterit? Ecce docuísti multos et manus lassas roborásti, vacillántes confirmavérunt sermónes tui et génua treméntia confortásti. Nunc autem venit super te plaga, et defecísti; tétigit te, et conturbátus es. Ubi est timor tuus, fortitúdo tua, patiéntia tua et perféctio viárum tuárum?

\R\ Si bona suscépimus de manu Dei, mala autem quare non sustineámus? * Dóminus dedit, Dóminus ábstulit; sicut Dómino plácuit, ita factum est: sit nomen Dómini benedíctum. \V\ Nudus egréssus sum de útero matris meæ et nudus revértar illuc. — Dóminus.

\lectioheading{Lectio ii}{}

\lettrine[lines=2]{R}{ECORDÁRE}, óbsecro te, quis umquam ínnocens périit? aut quando recti deléti sunt? Quin pótius vidi eos, qui operántur iniquitátem et séminant dolóres et metunt eos, flante Deo perísse et spíritu iræ ejus esse consúmptos. Rugítus leónis et vox leǽnæ et dentes catulórum leónum contríti sunt; tigris périit, eo quod non habéret prædam, et cátuli leónis dissipáti sunt.

\R\ Antequam cómedam suspíro, et tamquam inundántes aquæ sic rugítus meus; quia timor, quem timébam, evénit mihi, et quod verébar áccidit. Nonne dissimulávi? nonne sílui? nonne quiévi? * Et venit super me indignátio. \V\ Ecce non est auxílium mihi in me, et necessárii quoque mei recessérunt a me. — Et.

\rubric{In Octávis et Festis ritus simplícis: Glória Patri. Et.}

\lectioheading{Lectio iii}{}

\lettrine[lines=2]{P}{ORRO} ad me dictum est verbum abscónditum, et quasi furtíve suscépit auris mea venas susúrri ejus. In horróre visiónis noctúrnæ, quando solet sopor occupáre hómines, pavor ténuit me et tremor, et ómnia ossa mea pertérrita sunt, et cum spíritus, me præsénte, transíret, inhorruérunt pili carnis meæ. Stetit quidam, cujus non agnoscébam vultum, imágo coram óculis meis, et vocem quasi auræ lenis audívi: Numquid homo Dei comparatióne justificábitur? aut Factóre suo púrior erit vir? Ecce qui sérviunt ei non sunt stábiles, et in Angelis suis réperit pravitátem.

\R\ Quare detraxístis sermónibus veritátis? ad increpándum verba compónitis et subvértere nitímini amícum vestrum: * Verúmtamen quæ cogitástis, expléte. \V\ Quod justum est, judicáte; et non inveniétis in lingua mea iniquitátem. — Verúmtamen. Glória Patri. Verúmtamen.

\feriaheading{Feria Sexta infra Hebdomadam I Septembris}

\rubric{De libro Job}

\lectioheading{Lectio i}{Cap. 6, 1-13}

\lettrine[lines=2]{R}{ESPÓNDENS} autem Job dixit: Utinam appenderéntur peccáta mea, quibus iram mérui, et calamítas, quam pátior, in statéra. Quasi aréna maris hæc grávior apparéret, unde et verba mea dolóre sunt plena; quia sagíttæ Dómini in me sunt, quarum indignátio ébibit spíritum meum, et terróres Dómini mílitant contra me.

\R\ Indúta est caro mea putrédine, et sórdibus púlveris cutis mea áruit et contrácta est: * Meménto mei, Dómine, quóniam ventus est vita mea. \V\ Dies mei velócius transiérunt quam a texénte tela succíditur, et consúmpti sunt absque ulla spe. — Meménto.

\lectioheading{Lectio ii}{}

\lettrine[lines=2]{N}{UMQUID} rúgiet ónager cum habúerit herbam? aut mugíet bos, cum ante præsépe plenum stéterit? Aut póterit comédi insúlsum, quod non est sale condítum? aut potest áliquis gustáre, quod gustátum affert mortem? Quæ prius nolébat tángere ánima mea, nunc præ angústia cibi mei sunt.

\R\ Páucitas diérum meórum finiétur brevi; dimítte me, Dómine, ut plangam páululum dolórem meum, * Antequam vadam ad terram tenebrósam et opértam mortis calígine. \V\ Manus tuæ, Dómine, fecérunt me, et plasmavérunt me totum in circúitu; et sic repénte præcípitas me? — Antequam.

\rubric{In Octávis et Festis ritus simplícis: Glória Patri. Antequam.}

\lectioheading{Lectio iii}{}

\lettrine[lines=2]{Q}{UIS} det ut véniat petítio mea, et quod exspécto tríbuat mihi Deus? Et qui cœpit ipse me cónterat, solvat manum suam et succídat me? Et hæc mihi sit consolátio, ut afflígens me dolóre non parcat, nec contradícam sermónibus Sancti. Quæ est enim fortitúdo mea, ut sustíneam? aut quis finis meus, ut patiénter agam? Nec fortitúdo lápidum fortitúdo mea, nec caro mea ǽnea est. Ecce non est auxílium mihi in me, et necessárii quoque mei recessérunt a me.

\R\ Non abscóndas me, Dómine, a fácie tua: manum tuam longe fac a me, * Et formído tua non me térreat. \V\ Córripe me, Dómine, in misericórdia, non in furóre tuo, ne forte ad níhilum rédigas me. — Et. Glória Patri. Et.

\rubric{Ad Vésperas, nisi sequénti die persolvéndum sit Offícium de Festo novem Lectiónum, aut de qualíbet Octáva vel de Vigília, a Capítulo fit de sancta María, et de ea dícitur Offícium in Sábbato sequénti. Quod ítem servátur in ómnibus Sábbatis usque ad Sábbatum ante Domínicam I Advéntus inclusíve, præterquam in Sábbato Quátuor Témporum Septémbris, et in Sábbato, in quo anticipánda sit Domínica XXIII post Pentecósten.}

\feriaheading{Sabbato infra Hebdomadam I Septembris}

\rubric{De libro Job}

\lectioheading{Lectio i}{Cap. 7, 1-12}

\lettrine[lines=2]{M}{ILITIA} est vita hóminis super terram, et sicut dies mercenárii, dies ejus. Sicut servus desíderat umbram, et sicut mercenárius præstolátur finem óperis sui, sic et ego hábui menses vácuos et noctes laboriósas enumerávi mihi. Si dormíero, dicam: Quando consúrgam? Et rursum exspectábo vésperam, et replébor dolóribus usque ad ténebras.

\R\ Quis mihi tríbuat, ut in inférno prótegas me et abscóndas me, donec pertránseat furor tuus, Dómine, nisi tu, qui solus es Deus? * Et constítuas mihi tempus, in quo recordéris mei? \V\ Numquid sicut dies hóminis dies tui, ut quæras iniquitátem meam; cum sit nemo, qui de manu tua possit erúere? — Et.

\lectioheading{Lectio ii}{}

\lettrine[lines=2]{I}{NDÚTA} est caro mea putrédine, et sórdibus púlveris cutis mea áruit et contrácta est. Dies mei velócius transiérunt quam a texénte tela succíditur, et consúmpti sunt absque ulla spe. Meménto quia ventus est vita mea, et non revertétur óculus meus ut vídeat bona. Nec aspíciet me visus hóminis; óculi tui in me, et non subsístam.

\R\ Utinam appenderéntur peccáta mea, quibus iram mérui, * Et calamítas, quam pátior, in statéra. \V\ Quasi aréna maris hæc grávior apparéret, unde et verba mea dolóre sunt plena. — Et.

\rubric{In Octávis simplícibus et in Offício sanctæ Maríæ in Sábbato: Glória Patri. Et.}

\lectioheading{Lectio iii}{}

\lettrine[lines=2]{S}{ICUT} consúmitur nubes et pertránsit, sic qui descénderit ad ínferos non ascéndet; nec revertétur ultra in domum suam, neque cognóscet eum ámplius locus ejus. Quaprópter et ego non parcam ori meo; loquar in tribulatióne spíritus mei, confabulábor cum amaritúdine ánimæ meæ. Numquid mare ego sum aut cetus, quia circumdedísti me cárcere?

\R\ Quare detraxístis sermónibus veritátis? ad increpándum verba compónitis et subvértere nitímini amícum vestrum: * Verúmtamen quæ cogitástis, expléte. \V\ Quod justum est, judicáte; et non inveniétis in lingua mea iniquitátem. — Verúmtamen. Glória Patri. Verúmtamen.

\hourheading{Ad Vesperas}

\V\ Vespertína orátio ascéndat ad te, Dómine. \R\ Et descéndat super nos misericórdia tua.

\antlabel{Ad Magnif. Ant.} In ómnibus his * non peccávit Job lábiis suis neque stultum áliquid contra Deum locútus est.

\rubric{Orátio Domínicæ occurréntis.}

% ===== BATCH 2 of the Proprium de Tempore (autumnalis.md lines
% 13731-13902) -- completes the first full week of September: the
% six ferial days' Matins lessons (three Lessons and three
% Responsories apiece) continuing the Book of Job cycle through
% chapters 1-7, and the closing Vespers for "Sabbato infra
% Hebdomadam I Septembris" with its own Magnificat antiphon and a
% forward-pointing rubric to the proper Sunday Oratio. The lengthy
% governing rubric on Saturday Vespers switching to the Office of
% the Blessed Virgin Mary (except on the September Ember Saturday
% and certain anticipated Sundays) is transcribed in full as a
% single \rubric{} block. Three "In Octávis..." rubrics mark the
% alternate "Glória Patri" cue for feast-day/Marian-Saturday ranks;
% each is transcribed exactly as the source gives it, including the
% third day's distinct wording ("et in Officio sanctae Mariae in
% Sabbato" rather than "et Festis ritus simplicis"). Every
% Responsory's em-dash repeat-cue is preserved literally, consistent
% with Batch 1. The "---" horizontal rules before "Sabbato infra
% Hebdomadam I Septembris" and before "Dominica II Septembris" are
% omitted as scan artifacts. The document continues with "Dominica
% II Septembris" (autumnalis.md line 13903).
% =====
\feriaheading{Dominica II Septembris}

\rubric{Semiduplex}

\nocturnoheading{In I Nocturno}

\rubric{De libro Job}

\lectioheading{Lectio i}{Cap. 9, 1-17}

\lettrine[lines=2]{E}{T} respóndens Job ait: Vere scio quod ita sit, et quod non justificétur homo compósitus Deo. Si volúerit conténdere cum eo, non póterit ei respondére unum pro mille. Sápiens corde est et fortis róbore. Quis réstitit ei, et pacem hábuit? Qui tránstulit montes, et nesciérunt hi quos subvértit in furóre suo.

\R\ Si bona suscépimus de manu Dei, mala autem quare non sustineámus? * Dóminus dedit, Dóminus ábstulit; sicut Dómino plácuit, ita factum est: sit nomen Dómini benedíctum. \V\ Nudus egréssus sum de útero matris meæ et nudus revértar illuc. — Dóminus.

\lectioheading{Lectio ii}{}

\lettrine[lines=2]{Q}{UI} cómmovet terram de loco suo, et colúmnæ ejus concutiúntur; qui præcípit soli, et non óritur, et stellas claudit quasi sub signáculo; qui exténdit cælos solus, et gráditur super fluctus maris; qui facit Arctúrum et Oríona et Hýadas et interióra Austri; qui facit magna et incomprehensibília et mirabília, quorum non est númerus.

\R\ Antequam cómedam suspíro, et tamquam inundántes aquæ sic rugítus meus; quia timor, quem timébam, evénit mihi, et quod verébar áccidit. Nonne dissimulávi? nonne sílui? nonne quiévi? * Et venit super me indignátio. \V\ Ecce non est auxílium mihi in me, et necessárii quoque mei recessérunt a me. — Et.

\lectioheading{Lectio iii}{}

\lettrine[lines=2]{S}{I} vénerit ad me, non vidébo eum; si abíerit, non intélligam; si repénte intérroget, quis respondébit ei vel quis dícere potest: Cur ita facis? Deus, cujus iræ nemo resístere potest, et sub quo curvántur qui portant orbem; quantus ergo sum ego, ut respóndeam ei, et loquar verbis meis cum eo? Qui, étiam si habúero quíppiam justum, non respondébo, sed meum júdicem deprecábor. Et, cum invocántem exaudíerit me, non credo quod audíerit vocem meam. In túrbine enim cónteret me, et multiplicábit vúlnera mea, étiam sine causa.

\R\ Quare detraxístis sermónibus veritátis? ad increpándum verba compónitis et subvértere nitímini amícum vestrum: * Verúmtamen quæ cogitástis, expléte. \V\ Quod justum est, judicáte; et non inveniétis in lingua mea iniquitátem. — Verúmtamen. Glória Patri. Verúmtamen.

\nocturnoheading{In II Nocturno}

\rubric{Ex libro Morálium sancti Gregórii Papæ}

\lectioheading{Lectio iv}{Liber 9, cap. 2}

\lettrine[lines=2]{V}{ERE} scio quod ita sit, et quod non justificábitur homo compósitus Deo. Homo quippe Deo non compósitus justítiam pércipit, compósitus amíttit; quia quisquis se auctóri bonórum cómparat, bono se, quod accéperat, privat. Qui enim accépta bona sibi árrogat, suis contra Deum donis pugnat. Unde ergo despéctus erígitur, dignum est ut eréctus inde destruátur. Sanctus autem vir, quia omne virtútis nostræ méritum esse vítium cónspicit, si ab intérno árbitro distrícte judicétur, recte subjúngit: Si volúerit conténdere cum eo, non póterit respondére ei unum pro mille.

\R\ Indúta est caro mea putrédine, et sórdibus púlveris cutis mea áruit et contrácta est: * Meménto mei, Dómine, quóniam ventus est vita mea. \V\ Dies mei velócius transiérunt quam a texénte tela succíditur, et consúmpti sunt absque ulla spe. — Meménto.

\lectioheading{Lectio v}{}

\lettrine[lines=2]{I}{N} Scriptúra sancta millenárius númerus pro universitáte solet intélligi. Hinc étenim Psalmísta ait: Verbi, quod mandávit in mille generatiónes: cum profécto constet, quod ab ipso mundi exórdio usque ad Redemptóris advéntum per Evangelístam non ámplius quam septuagínta et septem propágines numeréntur. Quid ergo in millenário número nisi ad proferéndam novam sóbolem perfécta universitas præscítæ generatiónis exprímitur? Hinc et per Joánnem dícitur: Et regnábunt cum eo mille annis; quia vidélicet regnum sanctæ Ecclésiæ universitátis perfectióne solidátur.

\R\ Páucitas diérum meórum finiétur brevi; dimítte me, Dómine, ut plangam páululum dolórem meum, * Antequam vadam ad terram tenebrósam et opértam mortis calígine. \V\ Manus tuæ, Dómine, fecérunt me, et plasmavérunt me totum in circúitu; et sic repénte præcípitas me? — Antequam.

\lectioheading{Lectio vi}{}

\lettrine[lines=2]{Q}{UIA} vero monas décies multiplicáta in denárium dúcitur, denárius per semetípsum ductus in centenárium dilatátur, qui rursus per denárium ductus in millenárium ténditur; cum ab uno incípimus, ut ad millenárium veniámus, quid hoc loco uníus appellatióne, nisi bene vivéndi inítium? quid millenárii númeri amplitúdine, nisi ejúsdem bonæ vitæ perféctio designátur? Cum Deo autem conténdere, est, non ei tribúere, sed sibi glóriam suæ virtútis arrogáre. Sed sanctus vir conspíciat, quia et qui summa jam dona percépit, si de accéptis extóllitur, cuncta quæ accéperat, amíttit.

\R\ Non abscóndas me, Dómine, a fácie tua: manum tuam longe fac a me, * Et formído tua non me térreat. \V\ Córripe me, Dómine, in misericórdia, non in furóre tuo, ne forte ad níhilum rédigas me. — Et. Glória Patri. Et.

\R\ vii Quis mihi tríbuat, ut in inférno prótegas me et abscóndas me, donec pertránseat furor tuus, Dómine, nisi tu, qui solus es Deus? * Et constítuas mihi tempus, in quo recordéris mei? \V\ Numquid sicut dies hóminis dies tui, ut quæras iniquitátem meam; cum sit nemo, qui de manu tua possit erúere? — Et.

\R\ viii Duo Séraphim, \rubric{ut infra in Homilíis.}

\feriaheading{Feria Secunda infra Hebdomadam II Septembris}

\rubric{De libro Job}

\lectioheading{Lectio i}{Cap. 27, 1-15}

\lettrine[lines=2]{A}{DDIDIT} quoque Job, assúmens parábolam suam, et dixit: Vivit Deus, qui ábstulit judícium meum, et Omnípotens, qui ad amaritúdinem addúxit ánimam meam. Quia, donec súperest hálitus in me, et spíritus Dei in náribus meis, non loquéntur lábia mea iniquitátem, nec lingua mea meditábitur mendácium. Absit a me, ut justos vos esse júdicem; donec defíciam, non recédam ab innocéntia mea.

\R\ Versa est in luctum cíthara mea, et órganum meum in vocem fléntium: * Parce mihi, Dómine, nihil enim sunt dies mei. \V\ Cutis mea denigráta est super me, et ossa mea aruérunt. — Parce.

\lectioheading{Lectio ii}{}

\lettrine[lines=2]{J}{USTIFICATIÓNEM} meam, quam cœpi tenére, non déseram; neque enim reprehéndit me cor meum in omni vita mea. Sit ut ímpius inimícus meus, et adversárius meus quasi iníquus. Quæ est enim spes hypócritæ, si aváre rápiat, et non líberet Deus ánimam ejus? Numquid Deus áudiet clamórem ejus, cum vénerit super eum angústia? Aut póterit in Omnipoténte delectári, et invocáre Deum omni témpore?

\R\ Utinam appenderéntur peccáta mea, quibus iram mérui, * Et calamítas, quam pátior, in statéra. \V\ Quasi aréna maris hæc grávior apparéret, unde et verba mea dolóre sunt plena. — Et.

\rubric{In Octávis et Festis ritus simplícis: Glória Patri. Et.}

\lectioheading{Lectio iii}{}

\lettrine[lines=2]{D}{OCÉBO} vos per manum Dei quæ Omnípotens hábeat, nec abscóndam. Ecce vos omnes nostis, et quid sine causa vana loquímini? Hæc est pars hóminis ímpii apud Deum, et heréditas violentórum, quam ab Omnipoténte suscípient: si multiplicáti fúerint fílii ejus, in gládio erunt, et nepótes ejus non saturabúntur pane; qui réliqui fúerint ex eo sepeliéntur in intéritu, et víduæ illíus non plorábunt.

\R\ Quare detraxístis sermónibus veritátis? ad increpándum verba compónitis et subvértere nitímini amícum vestrum: * Verúmtamen quæ cogitástis, expléte. \V\ Quod justum est, judicáte; et non inveniétis in lingua mea iniquitátem. — Verúmtamen. Glória Patri. Verúmtamen.

\feriaheading{Feria Tertia infra Hebdomadam II Septembris}

\rubric{De libro Job}

\lectioheading{Lectio i}{Cap. 28, 12-28}

\lettrine[lines=2]{S}{APIÉNTIA} ubi invenítur? et quis est locus intelligéntiæ? Nescit homo prétium ejus, nec invenítur in terra suáviter vivéntium. Abýssus dicit: Non est in me, et mare lóquitur: Non est mecum. Non dábitur aurum obrízum pro ea, nec appendétur argéntum in commutatióne ejus; non conferétur tinctis Indiæ colóribus, nec lápidi sardónycho pretiosíssimo vel sapphíro.

\R\ Indúta est caro mea putrédine, et sórdibus púlveris cutis mea áruit et contrácta est: * Meménto mei, Dómine, quóniam ventus est vita mea. \V\ Dies mei velócius transiérunt quam a texénte tela succíditur, et consúmpti sunt absque ulla spe. — Meménto.

\lectioheading{Lectio ii}{}

\lettrine[lines=2]{N}{ON} adæquábitur ei aurum vel vitrum, nec commutabúntur pro ea vasa auri. Excélsa et eminéntia non memorabúntur comparatióne ejus; tráhitur autem sapiéntia de occúltis. Non adæquábitur ei topázius de Æthiópia, nec tinctúræ mundíssimæ componétur. Unde ergo sapiéntia venit? et quis est locus intelligéntiæ? Abscóndita est ab óculis ómnium vivéntium; vólucres quoque cæli latet. Perdítio et mors dixérunt: Áuribus nostris audívimus famam ejus.

\R\ Páucitas diérum meórum finiétur brevi; dimítte me, Dómine, ut plangam páululum dolórem meum, * Antequam vadam ad terram tenebrósam et opértam mortis calígine. \V\ Manus tuæ, Dómine, fecérunt me, et plasmavérunt me totum in circúitu; et sic repénte præcípitas me? — Antequam.

\rubric{In Octávis et Festis ritus simplícis: Glória Patri. Antequam.}

\lectioheading{Lectio iii}{}

\lettrine[lines=2]{D}{EUS} intélligit viam ejus, et ipse novit locum illíus; ipse enim fines mundi intuétur, et ómnia, quæ sub cælo sunt, réspicit. Qui fecit ventis pondus, et aquas appéndit in mensúra. Quando ponébat plúviis legem et viam procéllis sonántibus, tunc vidit illam et enarrávit et præparávit et investigávit. Et dixit hómini: Ecce timor Dómini, ipsa est sapiéntia, et recédere a malo, intelligéntia.

\R\ Non abscóndas me, Dómine, a fácie tua: manum tuam longe fac a me, * Et formído tua non me térreat. \V\ Córripe me, Dómine, in misericórdia, non in furóre tuo, ne forte ad níhilum rédigas me. — Et. Glória Patri. Et.

\feriaheading{Feria Quarta infra Hebdomadam II Septembris}

\rubric{De libro Job}

\lectioheading{Lectio i}{Cap. 31, 1-18}

\lettrine[lines=2]{P}{EPIGI} fœdus cum óculis meis, ut ne cogitárem quidem de vírgine. Quam enim partem habéret in me Deus désuper, et hereditátem Omnípotens de excélsis? Numquid non perdítio est iníquo, et alienátio operántibus injustítiam? Nonne ipse considérat vias meas, et cunctos gressus meos dinúmerat? Si ambulávi in vanitáte, et festinávit in dolo pes meus, appéndat me in statéra justa, et sciat Deus simplicitátem meam.

\R\ Quis mihi tríbuat, ut in inférno prótegas me et abscóndas me, donec pertránseat furor tuus, Dómine, nisi tu, qui solus es Deus? * Et constítuas mihi tempus, in quo recordéris mei? \V\ Numquid sicut dies hóminis dies tui, ut quæras iniquitátem meam; cum sit nemo, qui de manu tua possit erúere? — Et.

\lectioheading{Lectio ii}{}

\lettrine[lines=2]{S}{I} declinávit gressus meus de via, et si secútum est óculos meos cor meum, et si mánibus meis adhǽsit mácula, seram, et álius cómedat, et progénies mea eradicétur. Si decéptum est cor meum super mulíere, et si ad óstium amíci mei insidiátus sum, scortum altérius sit uxor mea, et super illam incurvéntur álii. Hoc enim nefas est et iníquitas máxima; ignis est usque ad perditiónem dévorans et ómnia eradícans genímina.

\R\ Utinam appenderéntur peccáta mea, quibus iram mérui, * Et calamítas, quam pátior, in statéra. \V\ Quasi aréna maris hæc grávior apparéret, unde et verba mea dolóre sunt plena. — Et.

\rubric{In Octávis et Festis ritus simplícis: Glória Patri. Et.}

\lectioheading{Lectio iii}{}

\lettrine[lines=2]{S}{I} contémpsi subíre judícium cum servo meo et ancílla mea, cum disceptárent advérsum me: quid enim fáciam, cum surréxerit ad judicándum Deus? et, cum quæsíerit, quid respondébo illi? Numquid non in útero fecit me, qui et illum operátus est? et formávit me in vulva unus? Si negávi, quod volébant, paupéribus et óculos víduæ exspectáre feci; si comédi buccéllam meam solus, et non comédit pupíllus ex ea (quia ab infántia mea crevit mecum miserátio, et de útero matris meæ egréssa est mecum).

\R\ Quare detraxístis sermónibus veritátis? ad increpándum verba compónitis et subvértere nitímini amícum vestrum: * Verúmtamen quæ cogitástis, expléte. \V\ Quod justum est, judicáte; et non inveniétis in lingua mea iniquitátem. — Verúmtamen. Glória Patri. Verúmtamen.

\feriaheading{Feria Quinta infra Hebdomadam II Septembris}

\rubric{De libro Job}

\lectioheading{Lectio i}{Cap. 38, 1-20}

\lettrine[lines=2]{R}{ESPÓNDENS} autem Dóminus Job de túrbine, dixit: Quis est iste invólvens senténtias sermónibus imperítis? Accínge, sicut vir, lumbos tuos; interrogábo te et respónde mihi. Ubi eras quando ponébam fundaménta terræ? Índica mihi, si habes intelligéntiam. Quis pósuit mensúras ejus, si nosti? vel quis teténdit super eam líneam? Super quo bases illíus solidátæ sunt? aut quis demísit lápidem angulárem ejus, cum me laudárent simul astra matutína, et jubilárent omnes fílii Dei?

\R\ Si bona suscépimus de manu Dei, mala autem quare non sustineámus? * Dóminus dedit, Dóminus ábstulit; sicut Dómino plácuit, ita factum est: sit nomen Dómini benedíctum. \V\ Nudus egréssus sum de útero matris meæ et nudus revértar illuc. — Dóminus.

\lectioheading{Lectio ii}{}

\lettrine[lines=2]{Q}{UIS} conclúsit óstiis mare, quando erumpébat quasi de vulva procédens, cum ponérem nubem vestiméntum ejus, et calígine illud, quasi pannis infántiæ, obvólverem? Circúmdedi illud términis meis, et pósui vectem et óstia, et dixi: Usque huc vénies, et non procédes ámplius, et hic confrínges tuméntes fluctus tuos. Numquid post ortum tuum præcepísti dilúculo, et ostendísti aurórae locum suum? Et tenuísti concútiens extréma terræ, et excussísti ímpios ex ea?

\R\ Antequam cómedam suspíro, et tamquam inundántes aquæ sic rugítus meus; quia timor, quem timébam, evénit mihi, et quod verébar áccidit. Nonne dissimulávi? nonne sílui? nonne quiévi? * Et venit super me indignátio. \V\ Ecce non est auxílium mihi in me et necessárii quoque mei recessérunt a me. — Et.

\rubric{In Octávis et Festis ritus simplícis: Glória Patri. Et.}

\lectioheading{Lectio iii}{Cap. 42, 1-6}

\lettrine[lines=2]{R}{ESTITUÉTUR} ut lutum signáculum, et stabit sicut vestiméntum: auferétur ab ímpiis lux sua, et bráchium excélsum confringétur? Numquid ingréssus es profúnda maris, et in novíssimis abýssi deambulásti? Numquid apértæ sunt tibi portæ mortis, et óstia tenebrósa vidísti? Numquid considerásti latitúdinem terræ? Índica mihi, si nosti, ómnia: in qua via lux hábitet, et tenebrárum quis locus sit; ut ducas unumquódque ad términos suos, et intélligas sémitas domus ejus.

\R\ Quare detraxístis sermónibus veritátis? ad increpándum verba compónitis et subvértere nitímini amícum vestrum: * Verúmtamen quæ cogitástis, expléte. \V\ Quod justum est, judicáte; et non inveniétis in lingua mea iniquitátem. — Verúmtamen. Glória Patri. Verúmtamen.

\feriaheading{Feria Sexta infra Hebdomadam II Septembris}

\rubric{De libro Job}

\lectioheading{Lectio i}{Cap. 40, 1-11}

\lettrine[lines=2]{R}{ESPÓNDENS} autem Dóminus Job de túrbine, dixit: Accínge sicut vir lumbos tuos; interrogábo te, et índica mihi. Numquid írritum fácies judícium meum, et condemnábis me, ut tu justificéris? Et, si habes bráchium sicut Deus, et si voce símili tonas, circúmda tibi decórem et in sublíme erígere, et esto gloriósus et speciósis indúere véstibus.

\R\ Indúta est caro mea putrédine, et sórdibus púlveris cutis mea áruit et contrácta est: * Meménto mei, Dómine, quóniam ventus est vita mea. \V\ Dies mei velócius transiérunt quam a texénte tela succíditur, et consúmpti sunt absque ulla spe. — Meménto.

\lectioheading{Lectio ii}{}

\lettrine[lines=2]{D}{ISPÉRGE} supérbos in furóre tuo, et respíciens omnem arrogántem humília; réspice cunctos supérbos et confúnde eos, cóntere ímpios in loco suo; abscónde eos in púlvere simul, et fácies eórum demérge in fóveam; et ego confitébor quod salváre te possit déxtera tua. Ecce béhemoth, quem feci tecum, fœnum, quasi bos, cómedet. Fortitúdo ejus in lumbis ejus, et virtus illíus in umbilíco ventris ejus.

\R\ Páucitas diérum meórum finiétur brevi; dimítte me, Dómine, ut plangam páululum dolórem meum, * Antequam vadam ad terram tenebrósam et opértam mortis calígine. \V\ Manus tuæ, Dómine, fecérunt me, et plasmavérunt me totum in circúitu; et sic repénte præcípitas me? — Antequam.

\rubric{In Octávis et Festis ritus simplícis: Glória Patri. Antequam.}

\lectioheading{Lectio iii}{Cap. 42, 1-6}

\lettrine[lines=2]{R}{ESPÓNDENS} autem Job Dómino dixit: Scio quia ómnia potes, et nulla te latet cogitátio. Quis est iste qui celat consílium absque sciéntia? Ideo insipiénter locútus sum, et quæ ultra modum excéderent sciéntiam meam. Audi, et ego loquar; interrogábo te, et respónde mihi. Audítu auris audívi te; nunc autem óculus meus videt te. Idcírco ipse me reprehéndo, et ago pœniténtiam in favílla et cínere.

\R\ Non abscóndas me, Dómine, a fácie tua: manum tuam longe fac a me, * Et formído tua non me térreat. \V\ Córripe me, Dómine, in misericórdia, non in furóre tuo, ne forte ad níhilum rédigas me. — Et. Glória Patri. Et.

\feriaheading{Sabbato infra Hebdomadam II Septembris}

\rubric{De libro Job}

\lectioheading{Lectio i}{Cap. 42, 7-16}

\lettrine[lines=2]{P}{OSTQUAM} autem locútus est Dóminus verba hæc ad Job, dixit ad Elíphaz Themanítem: Irátus est furor meus in te et in duos amícos tuos, quóniam non estis locúti coram me rectum, sicut servus meus Job. Súmite ergo vobis septem tauros et septem aríetes et ite ad servum meum Job et offérte holocáustum pro vobis; Job autem servus meus orábit pro vobis. Fáciem ejus suscípiam, ut non vobis imputétur stultítia; neque enim locúti estis ad me recta, sicut servus meus Job.

\R\ Quis mihi tríbuat, ut in inférno prótegas me et abscóndas me, donec pertránseat furor tuus, Dómine, nisi tu, qui solus es Deus? * Et constítuas mihi tempus, in quo recordéris mei? \V\ Numquid sicut dies hóminis dies tui, ut quæras iniquitátem meam; cum sit nemo, qui de manu tua possit erúere? — Et constítuas mihi tempus, in quo recordéris mei?

\lectioheading{Lectio ii}{}

\lettrine[lines=2]{A}{BIÉRUNT} ergo Elíphaz Themanítes et Báldad Suhítes et Sophar Naamathítes, et fecérunt sicut locútus fúerat Dóminus ad eos, et suscépit Dóminus fáciem Job. Dóminus quoque convérsus est ad pœniténtiam Job, cum oráret ille pro amícis suis; et áddidit Dóminus ómnia, quæcúmque fúerant Job, duplícia. Venérunt autem ad eum omnes fratres sui et univérsæ soróres suæ cuncti, qui nóverant eum prius, et comedérunt cum eo panem in domo ejus; et movérunt super eum caput, et consoláti sunt eum super omni malo, quod intúlerat Dóminus super eum, et dedérunt ei unusquísque ovem unam et ináurem áuream unam.

\R\ Utinam appenderéntur peccáta mea, quibus iram mérui, * Et calamítas, quam pátior, in statéra. \V\ Quasi aréna maris hæc grávior apparéret, unde et verba mea dolóre sunt plena. — Et.

\rubric{In Octávis simplícibus et in Offício sanctæ Maríæ in Sábbato: Glória Patri. Et.}

\lectioheading{Lectio iii}{}

\lettrine[lines=2]{D}{OMINUS} autem benedíxit novíssimis Job magis quam princípio ejus; et facta sunt ei quatuórdecim míllia óvium et sex míllia camelórum et mille juga boum et mille ásinæ, et fuérunt ei septem fílii et tres fíliæ, et vocávit nomen uníus Diem et nomen secúndæ Cássiam et nomen tértiæ Cornustíbii. Non sunt autem invéntæ mulíeres speciósæ sicut fíliæ Job in univérsa terra; dedítque eis pater suus hereditátem inter fratres eárum. Vixit autem Job post hæc centum quadragínta annis, et vidit fílios suos et fílios filiórum suórum usque ad quartam generatiónem; et mórtuus est senex et plenus diérum.

\R\ Quare detraxístis sermónibus veritátis? ad increpándum verba compónitis et subvértere nitímini amícum vestrum: * Verúmtamen quæ cogitástis, expléte. \V\ Quod justum est, judicáte; et non inveniétis in lingua mea iniquitátem. — Verúmtamen. Glória Patri. Verúmtamen.

\hourheading{Ad Vesperas}

\V\ Vespertína orátio ascéndat ad te, Dómine. \R\ Et descéndat super nos misericórdia tua.

\antlabel{Ad Magnif. Ant.} Ne reminiscáris, * Dómine, delícta mea vel paréntum meórum, neque vindíctam sumas de peccátis meis.

\rubric{Oratio Dominicæ occurrentis.}

% ===== BATCH 3 of the Proprium de Tempore (autumnalis.md lines
% 13903-14129) -- the complete second week of September: "Dominica
% II Septembris" Matins (Nocturns I-II continuing Job through
% chapter 9 and St. Gregory's Moralia Book 9, with the same deferred
% Nocturn III and numbered Responsories vii-viii pattern as Dominica
% I), the six ferial days' Matins lessons carrying the Job cycle
% through its conclusion (chapters 27-42, ending with Job's
% restoration), and the closing Saturday Vespers. Two genuine source
% variants from the established responsory pattern are preserved
% exactly: Feria Quinta's third lesson and Feria Sexta's third
% lesson both cite "Cap. 42, 1-6" as a fresh citation (rather than
% continuing the prior lesson's citation with an empty second
% argument), since the source itself restates the chapter reference
% at each occurrence; and the Sabbato's first Responsory repeats its
% FULL verse text after the em-dash ("— Et constítuas mihi tempus,
% in quo recordéris mei?") rather than the usual single-word cue,
% transcribed exactly as the source gives it. The "---" horizontal
% rules between days are omitted as scan artifacts throughout. The
% document continues with "Dominica III Septembris" (autumnalis.md
% line 14130).
% =====
\feriaheading{Dominica III Septembris}

\rubric{Semiduplex}

\nocturnoheading{In I Nocturno}

\rubric{Incipit liber Tobiæ}

\lectioheading{Lectio i}{Cap. 1, 1-15}

\lettrine[lines=2]{T}{OBIAS} ex tribu et civitáte Néphthali, quæ est in superióribus Galilǽæ supra Naásson, post viam quæ ducit ad Occidéntem in sinístro habens civitátem Sephet, cum captus esset in diébus Salmánasar regis Assyriórum, in captivitáte tamen pósitus viam veritátis non deséruit; ita ut ómnia, quæ habére póterat, quotídie concaptívis frátribus, qui erant ex ejus génere, impertíret. Cumque esset júnior ómnibus in tribu Néphthali, nihil tamen puerile gessit in ópere.

\R\ Peto, Dómine, ut de vínculo improperii hujus absólvas me, aut certe désuper terram erípias me: * Ne reminiscáris delícta mea vel paréntum meórum, neque vindíctam sumas de peccátis meis: quia éruis sustinéntes te, Dómine. \V\ Ómnia enim judícia tua justa sunt, et omnes viæ tuæ misericórdia et véritas: et nunc, Dómine, meménto mei. — Ne.

\lectioheading{Lectio ii}{}

\lettrine[lines=2]{D}{ENIQUE}, cum irent omnes ad vítulos áureos, quos Jeróboam fécerat rex Israël, hic solus fugiébat consórtia ómnium; sed pergébat in Jerúsalem ad templum Dómini et ibi adorábat Dóminum Deum Israël, ómnia primitíva sua et décimas suas fidéliter ófferens; ita ut in tértio anno prosélytis et ádvenis ministráret omnem decimatiónem. Hæc et his simília secúndum legem Dei puérulus observábat. Cum vero factus esset vir, accépit uxórem Annam de tribu sua genuítque ex ea fílium, nomen suum impónens ei. Quem ab infántia timére Deum dócuit et abstinére ab omni peccáto.

\R\ Omni témpore bénedic Deum, et pete ab eo ut vias tuas dírigat, * Et in omni témpore consília tua in ipso permáneant. \V\ Inquíre ut fácias quæ plácita sunt illi in veritáte, et in tota virtúte tua. — Et.

\lectioheading{Lectio iii}{}

\lettrine[lines=2]{I}{GITUR} cum per captivitátem devenísset cum uxóre sua et fílio in civitátem Níniven cum omni tribu sua, cum omnes éderent ex cibis gentílium, iste custodívit ánimam suam et numquam contamínátus est in escis eórum. Et quóniam memor fuit Dómini in toto corde suo, dedit illi Deus grátiam in conspéctu Salmánasar regis, et dedit illi potestátem quocúmque vellet ire, habens libertátem quæcúmque fácere voluísset. Pergébat ergo ad omnes, qui erant in captivitáte, et mónita salútis dabat eis.

\R\ Memor esto, fili, quóniam páuperem vitam gérimus: * Habébis multa bona, si timúeris Deum. \V\ In mente habéto eum, et cave nequándo prætermíttas præcépta ejus. — Habébis. Glória Patri. Habébis.

\nocturnoheading{In II Nocturno}

\rubric{Sermo sancti Leónis Papæ}

\rubric{Sermo 9 de jejúnio 7 mensis}

\lectioheading{Lectio iv}{}

\lettrine[lines=2]{S}{CIO} quidem, dilectíssimi, plúrimos vestrum ita in iis, quæ ad observántiam christiánæ fídei pértinent, esse devótos, ut nostris cohortatiónibus non indígeant admonéri. Quod enim dudum et tradítio decrévit, et consuetúdo firmávit; nec erudítio ignórat, nec pietas prætermíttit. Sed quia sacerdotális offícii est, erga omnes Ecclésiæ fílios curam habére commúnem, in id quod et rúdibus prosit et doctis, quos simul dilígimus, páriter incitámus; ut jejúnium, quod nobis séptimi mensis recúrsus indícit, fide álacri per castigatiónem ánimi et córporis celebrémus.

\R\ Sufficiébat nobis paupértas nostra, ut divítiæ computaréntur: numquam fuísset pecúnia ipsa, pro qua misísti fílium nostrum, * Báculum senectútis nostræ! \V\ Heu me, fili mi, ut quid te mísimus peregrinári, lumen oculórum nostrórum? — Báculum.

\lectioheading{Lectio v}{}

\lettrine[lines=2]{I}{DEO} enim ipsa continéntiæ observántia quátuor est assignáta tempóribus, ut in idípsum totíus anni redeúnte decúrsu, cognoscerémus nos indesinénter purificatiónibus indigére; semperque esse niténdum, dum hujus vitæ varietáte jactámur, ut peccátum, quod fragilitáte carnis et cupiditátum pollutióne contráhitur, jejúniis atque eleemósynis deleátur. Esuriámus páululum, dilectíssimi, et aliquántulum, quod juvándis possit prodésse paupéribus, nostræ consuetúdini subtrahámus.

\R\ Benedícite Deum cæli et coram ómnibus vivéntibus confitémini ei, * Quia fecit vobíscum misericórdiam suam. \V\ Ipsum benedícite et cantáte illi: et enarráte ómnia mirabília ejus. — Quia.

\lectioheading{Lectio vi}{}

\lettrine[lines=2]{E}{LECTÉTUR} consciéntia benignórum frúctibus largitátis: et gáudia tríbuens, quo es lætificándus, accípies. Diléctio próximi, diléctio Dei est, qui plenitúdinem legis et prophetárum in hac gémina caritátis unitáte constítuit; ut nemo ambígeret, Deo se offérre, quod hómini contulísset, dicénte Dómino Salvatóre, cum de aléndis juvandísque paupéribus loquerétur: Quod uni eórum fecísti, mihi fecísti. Quarta ígitur et sexta féria jejunémus; sábbato vero apud beátum Petrum Apóstolum vigílias celebrémus: cujus nos méritis et oratiónibus crédimus adjuvándos, ut misericórdi Deo jejúnio nostro et devotióne placeámus.

\R\ Tempus est ut revértar ad eum qui misit me; * Vos autem benedícite Deum et enarráte ómnia mirabília ejus. \V\ Confitémini ei coram ómnibus vivéntibus, quia fecit vobíscum misericórdiam suam. — Vos. Glória Patri. Vos.

\R\ vii Tribulatiónes civitátum audívimus, quas passæ sunt, et defécimus: timor et hebetúdo mentis cécidit super nos et super líberos nostros: ipsi montes nolunt recípere fugam nostram: * Dómine, miserére. \V\ Peccávimus cum pátribus nostris, injúste égimus, iniquitátem fécimus. — Dómine.

\R\ viii Duo Séraphim, \rubric{ut infra in Homilíis.}

\feriaheading{Feria Secunda infra Hebdomadam III Septembris}

\rubric{Si Lectiónes de Scriptúra hujus Fériæ II aut Fériæ III sequéntis, sua die dici non potúerint, et in Fériis Quátuor Témporum fíeri contíngat Offícium, in quo Lectiónes de Scriptúra occurrénti sint dicéndæ, neque in eis reponéndum sit Inítium libri Tobíæ Domínicæ præcedénti assignátum, Lectiónes ut supra impedítæ, et si opus sit, étiam Lectiónes insequéntium Feriárum, servató tamen eárum órdine, ultérius transferúntur, ita ut in Offício in Quátuor Témporibus occurrénti Lectiónes de Scriptúra valeant adhibéri. Quod ítem servátur de Lectiónibus Fériæ V ut supra impedítis, quæ simíliter in casu vel ultérius transferúntur, vel étiam anticipántur. Responsória tamen semper dicúntur de Féria occurrénti.}

\rubric{De libro Tobíæ}

\lectioheading{Lectio i}{Cap. 2, 1-4}

\lettrine[lines=2]{P}{OST} hæc vero, cum esset dies festus Dómini, et factum esset prándium bonum in domo Tobíæ, dixit fílio suo: Vade et adduc áliquos de tribu nostra timéntes Deum, ut epuléntur nobíscum. Cumque abíisset, revérsus nuntiávit ei unum ex fíliis Israël jugulátum jacére in platéa. Statímque exsíliens de accúbitu suo relínquens prándium, jejúnus pervénit ad corpus, tollénsque illud portávit ad domum suam occúlte, ut, dum sol occubuísset, caute sepelíret eum.

\R\ Peto, Dómine, ut de vínculo improperii hujus absólvas me, aut certe désuper terram erípias me: * Ne reminiscáris delícta mea vel paréntum meórum, neque vindíctam sumas de peccátis meis: quia éruis sustinéntes te, Dómine. \V\ Omnia enim judícia tua justa sunt, et omnes viæ tuæ misericórdia et véritas: et nunc, Dómine, meménto mei. — Ne.

\lectioheading{Lectio ii}{Cap. 2, 8-18}

\lettrine[lines=2]{A}{RGUÉBANT} autem eum omnes próximi ejus dicéntes: Jam hujus rei causa intérfici jussus es et vix effugísti mortis impérium; et íterum sépelis mórtuos? Sed Tobías plus timens Deum quam regem, rapiébat córpora occísórum et occultábat in domo sua et médiis nóctibus sepeliébat ea. Cóntigit autem ut quadam die fatigátus a sepultúra, véniens in domum suam jactásset se juxta paríetem et obdormísset, et ex nido hirúndinum dormiénti illi cálida stércora incíderent super óculos ejus fierétque cæcus. Hanc autem tentatiónem ídeo permísit Dóminus eveníre illi, ut postéris darétur exémplum patiéntiæ ejus, sicut et sancti Job.

\R\ Omni témpore bénedic Deum, et pete ab eo ut vias tuas dírigat, * Et in omni témpore consília tua in ipso permáneant. \V\ Inquíre ut fácias quæ plácita sunt illi in veritáte, et in tota virtúte tua. — Et.

\rubric{In Octávis et Festis ritus simplícis: Glória Patri. Et.}

\lectioheading{Lectio iii}{}

\lettrine[lines=2]{N}{AM}, cum ab infántia sua semper Deum timúerit et mandáta ejus custodíerit, non est contristátus contra Deum quod plaga cæcitátis evénerit ei; sed immóbilis in Dei timóre permánsit agens grátias Deo ómnibus diébus vitæ suæ. Nam sicut beáto Job insultábant reges, ita isti paréntes et cognáti ejus irridébant vitam ejus dicéntes: Ubi est spes tua, pro qua eleemósynas et sepultúras faciébas? Tobías vero increpábat eos dicens: Nolíte ita loqui: quóniam fílii sanctórum sumus et vitam illam exspectámus, quam Deus datúrus est his qui fidem suam numquam mutant ab eo.

\R\ Memor esto, fili, quóniam páuperem vitam gérimus: * Habébis multa bona, si timúeris Deum. \V\ In mente habéto eum, et cave nequándo prætermíttas præcépta ejus. — Habébis. Glória Patri. Habébis.

\feriaheading{Feria Tertia infra Hebdomadam III Septembris}

\rubric{De libro Tobíæ}

\lectioheading{Lectio i}{Cap. 2, 19-23}

\lettrine[lines=2]{A}{NNA} vero uxor ejus ibat ad opus textrínum quotídie, et de labóre mánuum suárum victum, quem cónsequi póterat, deferébat. Unde factum est ut hædum caprárum accípiens detulísset domi; cujus cum vocem baléntis vir ejus audísset, dixit: Vidéte, ne forte furtívus sit; réddite eum dóminis suis, quia non licet nobis aut édere ex furto áliquid aut contíngere.

\R\ Sufficiébat nobis paupértas nostra, ut divítiæ computaréntur: numquam fuísset pecúnia ipsa, pro qua misísti fílium nostrum, * Báculum senectútis nostræ! \V\ Heu me, fili mi, ut quid te mísimus peregrinári, lumen oculórum nostrórum? — Báculum.

\lectioheading{Lectio ii}{}

\lettrine[lines=2]{A}{D} hæc uxor ejus iráta respóndit: Manifésté vana facta est spes tua et eleemósynæ tuæ modo apparuérunt. Atque his et áliis hujuscémodi verbis exprobrábat ei. Tunc Tobías ingémuit et cœpit oráre cum lácrimis dicens: Justus es, Dómine, et ómnia judícia tua justa sunt, et omnes viæ tuæ misericórdia et véritas et judícium. Et nunc, Dómine, memor esto mei et ne vindíctam sumas de peccátis meis neque reminiscáris delícta mea vel paréntum meórum.

\R\ Benedícite Deum cæli et coram ómnibus vivéntibus confitémini ei, * Quia fecit vobíscum misericórdiam suam. \V\ Ipsum benedícite et cantáte illi: et enarráte ómnia mirabília ejus. — Quia.

\rubric{In Octávis et Festis ritus simplícis: Glória Patri. Quia.}

\lectioheading{Lectio iii}{}

\lettrine[lines=2]{Q}{UÓNIAM} non obedívimus præcéptis tuis, ídeo tráditi sumus in direptiónem et captivitátem et mortem et in fábulam et in impropérium ómnibus natiónibus, in quibus dispersísti nos. Et nunc, Dómine, magna judícia tua, quia non égimus secúndum præcépta tua et non ambulávimus sincériter coram te. Et nunc, Dómine, secúndum voluntátem tuam fac mecum et præcipe in pace récipi spíritum meum; éxpedit enim mihi mori magis quam vívere.

\R\ Tempus est ut revértar ad eum qui misit me; * Vos autem benedícite Deum et enarráte ómnia mirabília ejus. \V\ Confitémini ei coram ómnibus vivéntibus, quia fecit vobíscum misericórdiam suam. — Vos. Glória Patri. Vos.

\rubric{¶ Fériæ IV et VI ac Sábbato Quátuor Témporum sunt Fériæ majóres non privilegiátæ.}

% ===== BATCH 4 of the Proprium de Tempore (autumnalis.md lines
% 14130-14245) -- "Dominica III Septembris" Matins, opening the Book
% of Tobias cycle in Nocturn I and St. Leo's sermon on the September
% Ember fast in Nocturn II (again with the deferred Nocturn III and
% numbered Responsories vii-viii pattern), followed by "Feria
% Secunda" and "Feria Tertia infra Hebdomadam III Septembris"
% continuing Tobias. Feria Secunda's lesson heading carries an
% unusually long governing rubric on how Ember Day lesson transfers
% interact with the Tobias readings when a conflict arises; this is
% transcribed in full as a single \rubric{} block. The closing
% rubric "Fériæ IV et VI ac Sábbato Quátuor Témporum sunt Fériæ
% majóres non privilegiátæ" signals the upcoming September Ember
% Days. The document continues with "Feria Quarta, Quatuor Temporum
% Septembris" (autumnalis.md line 14246).
% =====
\feriaheading{Feria Quarta, Quatuor Temporum Septembris}

\rubric{Hódie, Féria VI et Sábbato sequéntibus, Antíphonæ et Psalmi ad Laudes, et hac Féria étiam tres últimæ Antíphonæ cum Psalmis ad Matutínum, sumúntur de 2 loco; ad Primam autem addítur quartus Psalmus, ut in Psaltério notátur. Et dicúntur Preces feriáles ad omnes Horas, ut in Ordinário, præterquam in Vésperis et ad Complétórium Sábbati.}

\rubric{Léctio sancti Evangélii secúndum Marcum}

\lectioheading{Lectio i}{Cap. 9, 16-28}

\lettrine[lines=2]{I}{N} illo témpore: Respóndens unus de turba dixit ad Jesum: Magíster, áttuli fílium meum ad te habéntem spíritum mutum. Et réliqua.

\rubric{Homilía sancti Bedæ Venerábilis Presbýteri}

\rubric{Liber 3, cap. 38 in Marcum 9}

\lettrine[lines=2]{D}{ÆMONÍACUM} hunc, quem descéndens de monte Dóminus sanávit, Marcus quidem surdum mutúmque, Matthǽus vero lunáticum fuísse commémorat. Significat autem eos, de quibus scriptum est: Stultus ut luna mutátur. Qui numquam in eódem statu permanéntes, nunc ad hæc, nunc ad illa vítia mutáti, crescunt atque decréscunt. Qui muti sunt, non confiténdo fidem; surdi, nec ipsum aliquátenus veritátis audiéndo sermónem. Spumant autem, cum stultítia tabéscunt; stultórum namque et languéntium atque hébetum est, spumas salivárum ex ore dimíttere. Strident déntibus, cum iracúndiæ furóre flamméscunt: aréscunt, cum ótio torpénte languéscunt, et nulla virtútis indústria confortáti enérviter vivunt.

\R\ Tribulatiónes civitátum audívimus, quas passæ sunt, et defécimus: timor et hebetúdo mentis cécidit super nos et super líberos nostros: ipsi montes nolunt recípere fugam nostram: * Dómine, miserére. \V\ Peccávimus cum pátribus nostris, injúste égimus, iniquitátem fécimus. — Dómine.

\lectioheading{Lectio ii}{}

\lettrine[lines=2]{Q}{UOD} autem ait: Et dixi discípulis tuis ut ejícerent illum, et non potuérunt; laténter Apóstolos accúsat, cum impossibílitas curándi, intérdum non ad imbecillitátem curántium, sed ad eórum, qui curándi sunt, fidem referátur, dicénte Dómino: Fiat tibi secúndum fidem tuam. Qui respóndens eis, dixit: O generátio incrédula, quámdiu apud vos ero? quámdiu vos pátiar? Non quod tædio superátus sit mansuétus ac mitis, qui non apéruit sicut agnus coram tondénte os suum, nec in verba furóris erúpit; sed quo in similitúdinem médici, si ægrótum vídeat contra sua præcépta se gérere, dicat: Úsquequo accédam ad domum tuam? quoúsque artis meæ perdam indústriam, me áliud jubénte, et te áliud perpetránte?

\R\ Omni témpore bénedic Deum, et pete ab eo ut vias tuas dírigat, * Et in omni témpore consília tua in ipso permáneant. \V\ Inquíre ut fácias quæ plácita sunt illi in veritáte, et in tota virtúte tua. — Et.

\lectioheading{Lectio iii}{}

\lettrine[lines=2]{E}{T} dixit illis: Hoc genus in nullo potest exíre, nisi in oratióne et jejúnio. Dum docet Apóstolos, quómodo dæmon nequíssimus débeat expélli, omnes instítuit ad vitam; ut scílicet novérimus, gravíssima quæque vel immundórum spirítuum, vel hóminum tentaménta jejúniis et oratiónibus esse superánda: iram quoque Dómini, cum in ultiónem nostrórum scélerum fúerit accénsa, hoc remédio singulári posse placári. Jejúnium autem generále est, non solum ab escis, sed et a cunctis illécebris abstinére carnálibus, immo ab ómnibus vitiórum continére se passiónibus. Sic et orátio generális non in verbis solum est, quibus divínam cleméntiam invocámus, verum étiam in ómnibus, quæ in obséquium nostri Condítoris fídei devotióne gérimus.

\R\ Memor esto, fili, quóniam páuperem vitam gérimus: * Habébis multa bona, si timúeris Deum. \V\ In mente habéto eum, et cave nequándo prætermíttas præcépta ejus. — Habébis. Glória Patri. Habébis.

\hourheading{Ad Laudes}

\V\ Repléti sumus mane misericórdia tua. \R\ Exsultávimus, et delectáti sumus.

\antlabel{Ad Bened. Ant.} Hoc genus * dæmoniórum in nullo potest exíre, nisi in oratióne et jejúnio.

\oratioheading{Oratio}

\lettrine[lines=2]{M}{ISERICÓRDIÆ} tuæ remédiis, quǽsumus, Dómine, fragílitas nostra subsístat: ut, quæ sua conditióne attéritur, tua cleméntia reparétur. Per Dóminum.

\rubric{Et dícitur usque ad Nonam. Ad Vésperas vero Orátio Domínicæ præcedéntis. Si autem Vésperæ dicéndæ non sint de Féria, de ea nec fit Commemorátio.}

\feriaheading{Feria Quinta infra Hebdomadam III Septembris}

\rubric{De libro Tobíæ}

\lectioheading{Lectio i}{Cap. 12, 1-17}

\lettrine[lines=2]{V}{OCÁVIT} ad se Tobías fílium suum dixítque ei: Quid póssumus dare viro isti sancto, qui venit tecum? Respóndens Tobías dixit patri suo: Pater, quam mercédem dábimus ei? aut quid dignum póterit esse benefíciis ejus? Me duxit et redúxit sanum, pecúniam a Gabélo ipse recépit, uxórem ipse me habére fecit et dæmónium ab ea ipse compéscuit, gáudium paréntibus ejus fecit, meípsum a devoratióne piscis erípuit, te quoque vidére fecit lumen cæli, et bonis ómnibus per eum repléti sumus. Quid illi ad hæc potérimus dignum dare? Sed peto te, pater mi, ut roges eum, si forte dignábitur medietátem de ómnibus quæ alláta sunt sibi assúmere.

\R\ Peto, Dómine, ut de vínculo improperii hujus absólvas me, aut certe désuper terram erípias me: * Ne reminiscáris delícta mea vel paréntum meórum, neque vindíctam sumas de peccátis meis: quia éruis sustinéntes te, Dómine. \V\ Omnia enim judícia tua justa sunt, et omnes viæ tuæ misericórdia et véritas: et nunc, Dómine, meménto mei. — Ne.

\lectioheading{Lectio ii}{}

\lettrine[lines=2]{E}{T} vocántes eum, pater scílicet et fílius, tulérunt eum in partem; et rogáre cœpérunt ut dignarétur dimídiam partem ómnium quæ attúlerant, accéptam habére. Tunc dixit eis occúlte: Benedícite Deum cæli et coram ómnibus vivéntibus confitémini ei, quia fecit vobíscum misericórdiam suam. Etenim sacraméntum regis abscóndere bonum est, ópera autem Dei reveláre et confitéri honoríficum est. Bona est orátio cum jejúnio, et eleemósyna magis quam thesáuros auri recóndere; quóniam eleemósyna a morte líberat, et ipsa est quæ purgat peccáta, et facit inveníre misericórdiam et vitam ætérnam. Qui autem fáciunt peccátum et iniquitátem hostes sunt ánimæ suæ.

\R\ Omni témpore bénedic Deum, et pete ab eo ut vias tuas dírigat, * Et in omni témpore consília tua in ipso permáneant. \V\ Inquíre ut fácias quæ plácita sunt illi in veritáte, et in tota virtúte tua. — Et.

\rubric{In Octávis et Festis ritus simplícis: Glória Patri. Et.}

\lectioheading{Lectio iii}{}

\lettrine[lines=2]{M}{ANIFÉSTO} ergo vobis veritátem et non abscóndam a vobis occúltum sermónem. Quando orábas cum lácrimis et sepeliébas mórtuos et derelinquébas prándium tuum et mórtuos abscondébas per diem in domo tua et nocte sepeliébas eos, ego óbtuli oratiónem tuam Dómino. Et quia accéptus eras Deo, necésse fuit ut tentátio probáret te. Et nunc misit me Dóminus ut curárem te et Saram uxórem fílii tui a dæmónio liberárem; ego enim sum Ráphaël Angelus, unus ex septem qui astámus ante Dóminum. Cumque hæc audíssent, turbáti sunt et treméntes cecidérunt super terram in fáciem suam. Dixítque eis Angelus: Pax vobis, nolíte timére.

\R\ Memor esto, fili, quóniam páuperem vitam gérimus: * Habébis multa bona, si timúeris Deum. \V\ In mente habéto eum, et cave nequándo prætermíttas præcépta ejus. — Habébis. Glória Patri. Habébis.

\feriaheading{Feria Sexta, Quatuor Temporum Septembris}

\rubric{Léctio sancti Evangélii secúndum Lucam}

\lectioheading{Lectio i}{Cap. 7, 36-50}

\lettrine[lines=2]{I}{N} illo témpore: Rogábat Jesum quidam de pharisǽis, ut manducáret cum illo. Et ingréssus domum pharisǽi discúbuit. Et réliqua.

\rubric{Homilía sancti Gregórii Papæ}

\rubric{Homilía 33 in Evang., post inítium}

\lettrine[lines=2]{Q}{UEM} pharisǽus de falsa justítia præsúmens, nisi Judáicum pópulum; quem peccátrix múlier ad vestígia Dómini véniens et plorans, nisi convérsam gentilitátem desígnat? Quæ cum alabástro venit, unguéntum fudit, retro secus pedes Dómini stetit, lácrimis pedes rigávit, capíllis tersit, eosdémque quos infundébat et tergébat pedes osculári non désiit. Nos ergo, nos illa múlier expréssit, si toto corde ad Dóminum post peccáta redeámus, si ejus pœniténtiæ luctus imitémur. Quid namque unguénto, nisi bonæ odor opiniónis exprímitur? Unde et Paulus dicit: Christi bonus odor sumus Deo in omni loco.

\R\ Sufficiébat nobis paupértas nostra, ut divítiæ computaréntur: numquam fuísset pecúnia ipsa, pro qua misísti fílium nostrum, * Báculum senectútis nostræ! \V\ Heu me, fili mi, ut quid te mísimus peregrinári, lumen oculórum nostrórum? — Báculum.

\lectioheading{Lectio ii}{}

\lettrine[lines=2]{S}{I} ergo recta ópera ágimus, quibus opiniónis bone odóre Ecclésiam respergámus, quid in Dómini córpore, nisi unguéntum fúndimus? Sed secus pedes Jesu múlier stetit: contra pedes enim Dómini stétimus, cum in peccátis pósiti ejus itinéribus renitebámur. Sed si ad veram pœniténtiam post peccáta convértimur, jam retro secus pedes stamus; quia ejus vestígia séquimur, quem impugnabámus. Lácrimis múlier pedes ejus rigat: quod nos quoque veráciter ágimus, si quibúslibet últimis membris Dómini per compassiónis afféctum inclinámur, si sanctis ejus in tribulatióne compátimur, si eórum tristítiam, nostram putámus.

\R\ Benedícite Deum cæli et coram ómnibus vivéntibus confitémini ei, * Quia fecit vobíscum misericórdiam suam. \V\ Ipsum benedícite et cantáte illi: et enarráte ómnia mirabília ejus. — Quia.

\lectioheading{Lectio iii}{}

\lettrine[lines=2]{C}{APILLIS} ergo pedes Dómini tergimus, quando sanctis ejus, quibus ex caritáte compátimur, étiam ex his quæ nobis supérfluunt, miserémur; quátenus sic mens per compassiónem dóleat, ut étiam larga manus afféctum dolóris osténdat. Rigat namque lácrimis Redemptóris pedes, sed capíllis suis non tergit, qui utcúmque proximórum dolóri compátitur, sed tamen eis ex his quæ sibi supérfluunt, non miserétur. Plorat, et non tergit, qui verba quidem dolóris tríbuit, sed non ministrándo quæ desunt, vim dolóris mínime abscíndit. Osculátur múlier pedes, quos tergit: quod nos quoque plene ágimus, si studióse dilígimus, quos ex largitáte continémus; ne gravis nobis sit necéssitas próximi; ne ipsa nobis ejus indigéntia, quæ sustentátur, fiat onerósa; et, cum manus necessária tríbuit, ánimus a dilectióne torpéscat.

\R\ Tempus est ut revértar ad eum qui misit me; * Vos autem benedícite Deum et enarráte ómnia mirabília ejus. \V\ Confitémini ei coram ómnibus vivéntibus, quia fecit vobíscum misericórdiam suam. — Vos. Glória Patri. Vos.

\hourheading{Ad Laudes}

\V\ Repléti sumus mane misericórdia tua. \R\ Exsultávimus, et delectáti sumus.

\antlabel{Ad Bened. Ant.} Múlier * quæ erat in civitáte peccátrix, stans retro secus pedes Dómini, lácrimis cœpit rigáre pedes ejus et capíllis cápitis sui tergébat, et deosculabátur pedes ejus et unguénto ungébat.

\oratioheading{Oratio}

\lettrine[lines=2]{P}{RÆSTA}, quǽsumus, omnípotens Deus: ut observatiónes sacras ánnua devotióne recoléntes, et córpore tibi placeámus, et mente. Per Dóminum.

\rubric{Et dícitur usque ad Nonam. Ad Vésperas vero Orátio Domínicæ præcedéntis. Si autem Vésperæ dicéndæ non sint de Féria, de ea nec fit Commemorátio.}

\feriaheading{Sabbato, Quatuor Temporum Septembris}

\rubric{Léctio sancti Evangélii secúndum Lucam}

\lectioheading{Lectio i}{Cap. 13, 6-17}

\lettrine[lines=2]{I}{N} illo témpore: Dicébat Jesus turbis hanc similitúdinem: Arbórem fici habébat quidam plantátam in vínea sua et venit quærens fructum in illa et non invénit. Et réliqua.

\rubric{Homilía sancti Gregórii Papæ}

\rubric{Homilía 31 in Evangelia}

\lettrine[lines=2]{D}{OMINUS} ac Redémptor noster per Evangélium suum aliquándo verbis, aliquándo rebus lóquitur: aliquándo áliud verbis, atque áliud rebus; aliquándo autem hoc verbis, quod rebus. Duas étenim res ex Evangélio, fratres, audístis: ficúlneam infructuósam, et mulíerem curvam; et utríque rei est pietas impénsa. Illud autem dixit per similitúdinem: istud egit per exhibitiónem. Sed hoc significat ficúlnea infructuósa, quod múlier inclináta; et hoc ficúlnea reserváta, quod múlier erécta.

\R\ Tribulatiónes civitátum audívimus, quas passæ sunt, et defécimus: timor et hebetúdo mentis cécidit super nos et super líberos nostros: ipsi montes nolunt recípere fugam nostram: * Dómine, miserére. \V\ Peccávimus cum pátribus nostris, injúste égimus, iniquitátem fécimus. — Dómine.

\lectioheading{Lectio ii}{}

\lettrine[lines=2]{Q}{UID} arbor fici, nisi humánam natúram desígnat? Quid múlier inclináta, nisi eámdem natúram significándo denúntiat? Quæ et bene plantáta est sicut ficus, et bene creáta sicut múlier; sed in culpam própria sponte lapsa, neque fructum servat operatiónis, neque statum rectitúdinis. Ad peccátum quippe ex voluntáte córruens, quia fructum obediéntiæ ferre nóluit, statum rectitúdinis amísit. Quæ ad Dei similitúdinem cóndita, dum in sua dignitáte non pérstitit, quod plantáta vel creáta fúerat, serváre contémpsit. Tértio dóminus víneæ ad ficúlneam venit, quia natúram géneris humáni ante legem, sub lege, sub grátia, exspectándo, admonéndo, visitándo requisívit.

\R\ Omni témpore bénedic Deum, et pete ab eo ut vias tuas dírigat, * Et in omni témpore consília tua in ipso permáneant. \V\ Inquíre ut fácias quæ plácita sunt illi in veritáte, et in tota virtúte tua. — Et.

\lectioheading{Lectio iii}{}

\lettrine[lines=2]{V}{ENIT} ante legem, quia per naturálem intelléctum quid unusquísque exémplo sui, quáliter erga próximum ágere debuísset, innótuit. Venit in lege, quia præcipiéndo dócuit. Venit post legem per grátiam, quia pietátis suæ præséntiam exhibéndo monstrávit. Sed tamen in tribus annis fructum se non invenísse conquéritur: quia quorúmdam pravórum mentes nec inspiráta lex naturális córrigit, nec præcépta erúdiunt, nec incarnatiónis ejus miracula convértunt. Quid vero per cultórem víneæ, nisi præpositórum ordo exprímitur? Qui dum præsunt Ecclésiæ, nimírum Domínicæ víneæ curam gerunt.

\R\ Memor esto, fili, quóniam páuperem vitam gérimus: * Habébis multa bona, si timúeris Deum. \V\ In mente habéto eum, et cave nequándo prætermíttas præcépta ejus. — Habébis. Glória Patri. Habébis.

\hourheading{Ad Laudes}

\V\ Repléti sumus mane misericórdia tua. \R\ Exsultávimus, et delectáti sumus.

\antlabel{Ad Bened. Ant.} Illúmina, Dómine, * sedéntes in ténebris et umbra mortis, et dírige pedes nostros in viam pacis.

\oratioheading{Oratio}

\lettrine[lines=2]{O}{MNÍPOTENS} sempitérne Deus, qui per continéntiam salutárem córporibus medéris et méntibus: majestátem tuam súpplices exorámus; ut pia jejunántium deprecatióne placátus, et præséntia nobis subsídia tríbuas et futúra. Per Dóminum.

\rubric{Et dícitur usque ad Nonam inclusíve.}

\hourheading{Ad Vesperas}

\V\ Vespertína orátio ascéndat ad te, Dómine. \R\ Et descéndat super nos misericórdia tua.

\antlabel{Ad Magnif. Ant.} Adonái, Dómine, * Deus magne et mirábilis, qui dedísti salútem in manu féminæ, exáudi preces servórum tuórum.

\rubric{Oratio Dominicæ occurrentis.}

% ===== BATCH 5 of the Proprium de Tempore (autumnalis.md lines
% 14246-14411) -- the complete September Ember Days (Feria Quarta,
% Feria Sexta, and Sabbato, Quatuor Temporum Septembris), each
% following the Sunday-Gospel-homily pattern: a brief Gospel
% pericope incipit ending "Et réliqua" under \lectioheading{Lectio
% i}, immediately followed within the same Lesson by a Patristic
% homily on that Gospel (Bede on Mark 9, Gregory on Luke 7:36-50,
% Gregory again on Luke 13:6-17), each homily's author and source
% given as separate \rubric{} labels before its own fresh lettrine,
% matching the established precedent from aestiva.tex's parallel
% "Lectio vii" Trinity Sunday construction. Feria Quinta infra
% Hebdomadam III Septembris (an ordinary day, not an Ember day,
% falling between the Wednesday and Friday Ember days) continues
% the Tobias cycle with the Archangel Raphael's self-revelation.
% Each Ember day's all-caps "AD LAUDES" (and, for the Saturday, "AD
% VESPERAS") sub-heading is normalized to the standard mixed-case
% \hourheading{} form per the established incidental-formatting-
% variation rule, each followed by its own brief antiphon, Oratio,
% and governing rubric on how long the Oratio is said. The document
% continues with "Dominica IV Septembris" (autumnalis.md line
% 14412).
% =====
\feriaheading{Dominica IV Septembris}

\rubric{Semiduplex}

\nocturnoheading{In I Nocturno}

\rubric{Incipit liber Judith}

\lectioheading{Lectio i}{Cap. 1, 1-10; 2, 1-3}

\lettrine[lines=2]{A}{RPHÁXAD} ítaque rex Medórum subjugáverat multas gentes império suo, et ipse ædificávit civitátem potentíssimam, quam appellávit Ecbátanis, ex lapídibus quadrátis et sectis. Fecit muros ejus in altitúdinem cubitórum septuagínta et in latitúdinem cubitórum trigínta; turres vero ejus pósuit in altitúdinem cubitórum centum. Per quadrum vero eárum latus utrúmque vicenórum pedum spátio tendebátur; posuítque portas ejus in altitúdinem túrrium; et gloriabátur quasi potens in poténtia exércitus sui et in glória quadrigárum suárum.

\R\ Adonái, Dómine, Deus magne et mirábilis, qui dedísti salútem in manu féminæ, * Exáudi preces servórum tuórum. \V\ Benedíctus es, Dómine, qui non derelínquis præsuméntes de te, et de sua virtúte gloriántes humílias. — Exáudi.

\lectioheading{Lectio ii}{}

\lettrine[lines=2]{A}{NNO} ígitur duodécimo regni sui Nabuchodónosor rex Assyriórum, qui regnábat in Nínive civitáte magna, pugnávit contra Arpháxad et obtínuit eum in campo magno qui appellátur Rágau, circa Euphráten et Tigrin et Jádason, in campo Erioch regis Elicórum. Tunc exaltátum est regnum Nabuchodónosor, et cor ejus elevátum est; et misit ad omnes qui habitábant in Cilícia et Damásco et Líbano, et ad gentes, quæ sunt in Carmélo et Cedar et inhabitántes Galilǽam in campo magno Esdrelon, et ad omnes qui erant in Samaría et trans flumen Jordánem usque ad Jerúsalem, et omnem terram Jesse, quoúsque perveniátur ad términos Æthiópiæ.

\R\ Tribulatiónes civitátum audívimus, quas passæ sunt, et defécimus: timor et hebetúdo mentis cécidit super nos et super líberos nostros: ipsi montes nolunt recípere fugam nostram: * Dómine, miserére. \V\ Peccávimus cum pátribus nostris, injúste égimus, iniquitátem fécimus. — Dómine.

\lectioheading{Lectio iii}{}

\lettrine[lines=2]{A}{D} hos omnes misit núntios Nabuchodónosor rex Assyriórum; qui omnes uno ánimo contradixérunt et remisérunt eos vácuos et sine honóre abjecérunt. Tunc indignátus Nabuchodónosor rex advérsus omnem terram illam, jurávit per thronum et regnum suum, quod defénderet se de ómnibus regiónibus his. Anno tértio décimo Nabuchodónosor regis, vigésima et secúnda die mensis primi factum est verbum in domo Nabuchodónosor regis Assyriórum, ut defénderet se. Vocavítque omnes majóres natu, omnésque duces et bellatóres suos, et hábuit cum eis mystérium consílii sui dixítque cogitatiónem suam in eo esse, ut omnem terram suo subjugáret império.

\R\ Benedícat te Dóminus in virtúte sua, qui per te ad níhilum redégit inimícos nostros: * Ut non defíciat laus tua de ore hóminum. \V\ Benedíctus Dóminus qui creávit cælum et terram; quia hódie nomen tuum ita magnificávit. — Ut. Glória Patri. Ut.

\nocturnoheading{In II Nocturno}

\rubric{Ex libro sancti Ambrósii Epíscopi de Elía et jejúnio}

\lectioheading{Lectio iv}{Cap. 9}

\lettrine[lines=2]{P}{OTÉNTES} vinum prohibéntur bíbere, ne, cum bíberint, obliviscántur sapiéntiam. Dénique bibébant vinum in ebrietáte poténtes, qui Holoférni príncipi milítiæ regis Assyriórum se trádere gestiébant; sed non bibébat fémina Judith, jejúnans ómnibus diébus viduitátis suæ, præter festórum diérum solemnitátes. His armis muníta procéssit et omnem Assyriórum circumvénit exércitum. Sóbrii vigóre consílii ábstulit Holoférnis caput, servávit pudicítiam, victóriam reportávit.

\R\ Nos álium Deum nescímus præter Dóminum, in quo sperámus: * Qui non déspicit nos, nec ámovet salútem suam a génere nostro. \V\ Indulgéntiam ipsíus fusis lácrimis postulémus, et humiliémus illi ánimas nostras. — Qui.

\lectioheading{Lectio v}{}

\lettrine[lines=2]{H}{ÆC} enim succíncta jejúnio in castris prætendébat aliénis; ille vino sepúltus jacébat, ut ictum vúlneris sentíre non posset. Itaque uníus mulíeris jejúnium innúmeros strávit exércitus Assyriórum. Esther quoque púlchrior facta est jejúnio; Dóminus enim grátiam sóbriæ mentis augébat. Omne genus suum, id est, totum pópulum Judæórum a persecutiónis acerbitáte liberávit, ita ut regem sibi fáceret esse subjéctum.

\R\ Dominátor, Dómine cælórum et terræ, Creátor aquárum, Rex univérsæ creatúræ: * Exáudi oratiónem servórum tuórum. \V\ Tu, Dómine, cui humílium semper et mansuetórum plácuit deprecátio. — Exáudi.

\lectioheading{Lectio vi}{}

\lettrine[lines=2]{I}{TAQUE} illa, quæ tríduo jejunávit contínuo, et corpus suum aqua lavit, plus plácuit, et vindíctam rétulit. Aman autem, dum se regáli jactat convívio, inter ipsa vina pœnam suæ ebrietátis exsólvit. Est ergo jejúnium reconciliatiónis sacrifícium, virtútis increméntum, quod fecit étiam féminas fortióres augménto grátiæ. Jejúnium nescit fænerátorem, non sortem fæneris novit: non rédolet usúras mensa jejunántium. Etiam ipsis jejúnium convívíis dat grátiam: dulcióres post famem épulæ fiunt, quæ assiduitáte fastídio sunt, et diutúrna continuatióne viléscunt. Condiméntum cibi jejúnium est: quanto avídior appeténtia, tanto esca jucúndior.

\R\ Dómine Deus, qui cónteris bella ab inítio, éleva bráchium tuum super Gentes, quæ cógitant servis tuis mala: * Et déxtera tua glorificétur in nobis. \V\ Allíde virtútem eórum in virtúte tua; cadat robur eórum in iracúndia tua. — Et. Glória Patri. Et.

\R\ vii Confórta me, Rex Sanctórum principátum tenens: * Et da sermónem rectum et bene sonántem in os meum. \V\ Dómine, Rex univérsæ potestátis, convérte consílium eórum super eos. — Et.

\R\ viii Duo Séraphim, \rubric{ut infra in Homilíis.}

\feriaheading{Feria Secunda infra Hebdomadam IV Septembris}

\rubric{De libro Judith}

\lectioheading{Lectio i}{Cap. 4, 1-12}

\lettrine[lines=2]{F}{ILII} Israël, qui habitábant in terra Juda, timuérunt valde a fácie Holoférnis; tremor et horror invásit sensus eórum, ne hoc fáceret Jerúsalem et templo Dómini, quod fécerat céteris civitátibus et templis eárum. Et misérunt in omnem Samaríam per circúitum usque Jéricho et præoccupavérunt omnes vértices móntium et muris circumdedérunt vicos suos et congregavérunt fruménta in præparatiónem pugnæ.

\R\ Adonái, Dómine, Deus magne et mirábilis, qui dedísti salútem in manu féminæ, * Exáudi preces servórum tuórum. \V\ Benedíctus es, Dómine, qui non derelínquis præsuméntes de te, et de sua virtúte gloriántes humílias. — Exáudi.

\lectioheading{Lectio ii}{}

\lettrine[lines=2]{S}{ACÉRDOS} étiam Elíachim scripsit ad univérsos qui erant contra Esdrelon, quæ est contra fáciem campi magni juxta Dóthain, et univérsos per quos viæ tránsitus esse póterat, ut obtinérent ascénsus móntium, per quos via esse póterat ad Jerúsalem, et illic custodírent ubi angústum iter esse póterat inter montes. Et fecérunt fílii Israël secúndum quod constitúerat eis sacérdos Dómini Elíachim. Et clamávit omnis pópulus ad Dóminum instántia magna, et humiliavérunt ánimas suas in jejúniis et oratiónibus ipsi et mulíeres eórum.

\R\ Tribulatiónes civitátum audívimus, quas passæ sunt, et defécimus: timor et hebetúdo mentis cécidit super nos et super líberos nostros: ipsi montes nolunt recípere fugam nostram: * Dómine, miserére. \V\ Peccávimus cum pátribus nostris, injúste égimus, iniquitátem fécimus. — Dómine.

\rubric{In Octávis et Festis ritus simplícis: Glória Patri. Dómine.}

\lectioheading{Lectio iii}{}

\lettrine[lines=2]{E}{T} induérunt se sacerdótes cilíciis, et infántes prostravérunt contra fáciem templi Dómini, et altáre Dómini operuérunt cilício; et clamavérunt ad Dóminum Deum Israël unanímiter, ne daréntur in prædam infántes eórum, et uxóres eórum in divisiónem, et civitátes eórum in extermínium, et sancta eórum in pollutiónem, et fíerent oppróbrium Géntibus. Tunc Elíachim sacérdos Dómini magnus circuívit omnem Israël allocutúsque est eos dicens: Scitóte quóniam exáudiet Dóminus preces vestras, si manéntes permanséritis in jejúniis et oratiónibus in conspéctu Dómini.

\R\ Benedícat te Dóminus in virtúte sua, qui per te ad níhilum redégit inimícos nostros: * Ut non defíciat laus tua de ore hóminum. \V\ Benedíctus Dóminus qui creávit cælum et terram; quia hódie nomen tuum ita magnificávit. — Ut. Glória Patri. Ut.

\feriaheading{Feria Tertia infra Hebdomadam IV Septembris}

\rubric{De libro Judith}

\lectioheading{Lectio i}{Cap. 8, 1-11}

\lettrine[lines=2]{E}{T} factum est, cum audísset hæc verba Judith vídua, quæ erat fília Merári, fílii Idox, fílii Joseph, fílii Oziæ, fílii Elai, fílii Jammor, fílii Gédeon, fílii Ráphaïm, fílii Achitob, fílii Melchíæ, fílii Enan, fílii Nathaníæ, fílii Saláthiel, fílii Símeon, fílii Ruben, et vir ejus fuit Manásses, qui mórtuus est in diébus messis hordeáceæ; instábat enim super alligántes manípulos in campo, et venit æstus super caput ejus, et mórtuus est in Bethúlia civitáte sua et sepúltus est illic cum pátribus suis. Erat autem Judith relícta ejus vídua jam annis tribus et ménsibus sex.

\R\ Nos álium Deum nescímus præter Dóminum, in quo sperámus: * Qui non déspicit nos, nec ámovet salútem suam a génere nostro. \V\ Indulgéntiam ipsíus fusis lácrimis postulémus, et humiliémus illi ánimas nostras. — Qui.

\lectioheading{Lectio ii}{}

\lettrine[lines=2]{E}{T} in superióribus domus suæ fecit sibi secrétum cubículum, in quo cum puéllis suis clausa morabátur. Et habens super lumbos suos cilícium, jejunábat ómnibus diébus vitæ suæ præter sábbata et neoménias et festa domus Israël. Erat autem elegánti aspéctu nimis; cui vir suus reliquérat divítias multas et famíliam copiósam ac possessiónes arméntis boum et grégibus óvium plenas. Et erat hæc in ómnibus famosíssima, quóniam timébat Dóminum valde, nec erat qui loquerétur de illa verbum malum.

\R\ Dominátor, Dómine cælórum et terræ, Creátor aquárum, Rex univérsæ creatúræ: * Exáudi oratiónem servórum tuórum. \V\ Tu, Dómine, cui humílium semper et mansuetórum plácuit deprecátio. — Exáudi.

\rubric{In Octávis et Festis ritus simplícis: Glória Patri. Exáudi.}

\lectioheading{Lectio iii}{}

\lettrine[lines=2]{H}{ÆC} ítaque, cum audísset quóniam Ozías promisísset quod, transácto quinto die, tráderet civitátem, misit ad presbýteros Chabri et Charmi. Et venérunt ad illam, et dixit illis: Quod est hoc verbum, in quo consénsit Ozías, ut tradat civitátem Assýriis, si intra quinque dies non vénerit vobis adjutórium? Et qui estis vos, qui tentátis Dóminum?

\R\ Dómine Deus, qui cónteris bella ab inítio, éleva bráchium tuum super gentes, quæ cógitant servis tuis mala: * Et déxtera tua glorificétur in nobis. \V\ Allíde virtútem eórum in virtúte tua; cadat robur eórum in iracúndia tua. — Et. Glória Patri. Et.

\feriaheading{Feria Quarta infra Hebdomadam IV Septembris}

\rubric{De libro Judith}

\lectioheading{Lectio i}{Cap. 10, 1-4}

\lettrine[lines=2]{F}{ACTUM} est autem, cum cessásset clamáre ad Dóminum surréxit de loco in quo jacúerat prostráta ad Dóminum; vocavítque abram suam, et descéndens in domum suam ábstulit a se cilícium et éxuit se vestiméntis viduitátis suæ et lavit corpus suum et unxit se myro óptimo et discriminávit crinem cápitis sui et impósuit mitram super caput suum et índuit se vestiméntis jucunditátis suæ induítque sandália pédibus suis assumpsítque dextralíola et lília et ináures et ánulos et ómnibus ornaméntis suis ornávit se. Cui étiam Dóminus cóntulit splendórem.

\R\ Confórta me, Rex Sanctórum principátum tenens: * Et da sermónem rectum et bene sonántem in os meum. \V\ Dómine, Rex univérsæ potestátis, convérte consílium eórum super eos. — Et.

\lectioheading{Lectio ii}{Cap. 10, 11-12}

\lettrine[lines=2]{F}{ACTUM} est autem, cum descénderet montem circa ortum diéi, occurrérunt ei exploratóres Assyriórum et tenuérunt eam dicéntes: Unde venis aut quo vadis? Quæ respóndit: Fília sum Hebræórum; ídeo ego fugi a fácie eórum, quóniam futúrum agnóvi quod dentur vobis in deprædatiónem, pro eo quod contemnéntes vos noluérunt ultro trádere seípsos, ut invenírent misericórdiam in conspéctu vestro.

\R\ Tribulatiónes civitátum audívimus, quas passæ sunt, et defécimus: timor et hebetúdo mentis cécidit super nos et super líberos nostros: ipsi montes nolunt recípere fugam nostram: * Dómine, miserére. \V\ Peccávimus cum pátribus nostris, injúste égimus, iniquitátem fécimus. — Dómine.

\rubric{In Octávis et Festis ritus simplícis: Glória Patri. Dómine.}

\lectioheading{Lectio iii}{Cap. 10, 16-20}

\lettrine[lines=2]{D}{UXERÚNTQUE} illam ad tabernáculum Holoférnis annuntiántes eam. Cumque intrásset ante fáciem ejus, statim captus est in suis óculis Holoférnes. Dixerúntque ad eum satéllites ejus: Quis contémnat pópulum Hebræórum, qui tam decóras mulíeres habent, ut non pro his mérito pugnáre contra eos debeámus? Videns ítaque Judith Holoférnem sedéntem in conopéo, quod erat ex púrpura et auro et smarágdo et lapídibus pretiósis intéxtum, et, cum in fáciem ejus intendísset, adorávit eum prostérnens se super terram.

\R\ Benedícat te Dóminus in virtúte sua, qui per te ad níhilum redégit inimícos nostros: * Ut non defíciat laus tua de ore hóminum. \V\ Benedíctus Dóminus qui creávit cælum et terram; quia hódie nomen tuum ita magnificávit. — Ut. Glória Patri. Ut.

\feriaheading{Feria Quinta infra Hebdomadam IV Septembris}

\rubric{Si hæc Féria occúrrat post diem 24 Septémbris, ita ut hæc sit última hebdómada mensis, hac nocte pónitur liber Esther cum suis Responsóriis, ut infra in Dominica V, 321: et légitur de eo étiam duóbus sequéntibus diébus, ita tamen ut, lecto Inítio Domínicæ assignáto, sequéntibus Fériis dicántur Lectiónes de respondénti Féria VI aut Sábbato ejúsdem hebdómadæ V.}

\rubric{De libro Judith}

\lectioheading{Lectio i}{Cap. 12, 10-13}

\lettrine[lines=2]{E}{T} factum est, in quarto die Holoférnes fecit cœnam servis suis et dixit ad Vágao eunúchum suum: Vade et suáde Hebrǽam illam ut sponte conséntiat habitáre mecum. Foedum est enim apud Assýrios, si fémina irrídeat virum agéndo ut immúnis ab eo tránseat. Tunc introívit Vágao ad Judith et dixit: Non vereátur bona puélla introíre ad dóminum meum, ut honorificétur ante fáciem ejus, ut mandúcet cum eo et bibat vinum in jucunditáte. Cui Judith respóndit: Quæ ego sum, ut contradícam dómino meo?

\R\ Adonái, Dómine, Deus magne et mirábilis, qui dedísti salútem in manu féminæ, * Exáudi preces servórum tuórum. \V\ Benedíctus es, Dómine, qui non derelínquis præsuméntes de te, et de sua virtúte gloriántes humílias. — Exáudi.

\lectioheading{Lectio ii}{Cap. 13, 1-11}

\lettrine[lines=2]{U}{T} autem sero factum est, festinavérunt servi illíus ad hospítia sua; et conclúsit Vágao óstia cubículi et ábiit. Erant autem omnes fatigáti a vino; erátque Judith sola in cubículo. Porro Holoférnes jacébat in lecto nímia ebrietáte sopítus. Dixítque Judith puéllæ suæ ut staret foris ante cubículum et observáret. Stetítque Judith ante lectum orans cum lácrimis et labiórum motu in siléntio dicens: Confírma me, Dómine, Deus Israël, et réspice in hac hora ad ópera mánuum meárum, ut, sicut promisísti, Jerúsalem civitátem tuam érigas, et hoc quod credens per te posse fíeri cogitávi, perfíciam.

\R\ Tribulatiónes civitátum audívimus, quas passæ sunt, et defécimus: timor et hebetúdo mentis cécidit super nos et super líberos nostros: ipsi montes nolunt recípere fugam nostram: * Dómine, miserére. \V\ Peccávimus cum pátribus nostris, injúste égimus, iniquitátem fécimus. — Dómine.

\rubric{In Octávis et Festis ritus simplícis: Glória Patri. Dómine.}

\lectioheading{Lectio iii}{}

\lettrine[lines=2]{E}{T}, cum hæc dixísset, accéssit ad colúmnam, quæ erat ad caput léctuli ejus, et pugiónem ejus, qui in ea ligátus pendébat, exsólvit. Cumque evaginásset illum, apprehéndit comam cápitis ejus et ait: Confírma me, Dómine Deus, in hac hora. Et percússit bis in cervícem ejus et abscídit caput ejus et ábstulit conopéum ejus a colúmnis et evólvit corpus ejus truncum. Et post pusíllum exívit et trádidit caput Holoférnis ancíllæ suæ et jussit ut mítteret illud in peram suam.

\R\ Benedícat te Dóminus in virtúte sua, qui per te ad níhilum redégit inimícos nostros: * Ut non defíciat laus tua de ore hóminum. \V\ Benedíctus Dóminus qui creávit cælum et terram; quia hódie nomen tuum ita magnificávit. — Ut. Glória Patri. Ut.

\feriaheading{Feria Sexta infra Hebdomadam IV Septembris}

\rubric{De libro Judith}

\lectioheading{Lectio i}{Cap. 15, 1-3}

\lettrine[lines=2]{C}{UMQUE} omnis exércitus decollátum Holoférnem audísset, fugit mens et consílium ab eis, et solo tremóre et metu agitáti fugæ præsídium sumunt, ita ut nullus loquerétur cum próximo suo, sed inclináto cápite, relíctis ómnibus, evádere festinábant Hebrǽos, quos armátos super se veníre audiébant, fugiéntes per vias campórum et sémitas cóllium. Vidéntes ítaque fílii Israël fugiéntes, secúti sunt illos, descenderúntque clangéntes tubis et ululántes post ipsos.

\R\ Nos álium Deum nescímus præter Dóminum, in quo sperámus: * Qui non déspicit nos, nec ámovet salútem suam a génere nostro. \V\ Indulgéntiam ipsíus fusis lácrimis postulémus, et humiliémus illi ánimas nostras. — Qui.

\lectioheading{Lectio ii}{Cap. 15, 5-7}

\lettrine[lines=2]{M}{ISIT} ítaque Ozías núntios per omnes civitátes et regiónes Israël. Omnis ítaque régio omnísque urbs eléctam juventútem armátam misit post eos; et persecúti sunt eos in ore gládii, quoúsque pervenírent ad extremitátem fínium suórum. Réliqui autem qui erant in Bethúlia, ingréssi sunt castra Assyriórum et prædam, quam fugiéntes Assýrii relíquerant, abstulérunt et onustáti sunt valde.

\R\ Dominátor, Dómine cælórum et terræ, Creátor aquárum, Rex univérsæ creatúræ: * Exáudi oratiónem servórum tuórum. \V\ Tu, Dómine, cui humílium semper et mansuetórum plácuit deprecátio. — Exáudi.

\rubric{In Octávis et Festis ritus simplícis: Glória Patri. Exáudi.}

\lectioheading{Lectio iii}{Cap. 15, 9-12}

\lettrine[lines=2]{J}{OACHIM} autem summus póntifex de Jerúsalem venit in Bethúliam cum univérsis presbýteris suis, ut vidéret Judith. Quæ cum exísset ad illum, benedixérunt eam omnes una voce dicéntes: Tu glória Jerúsalem, tu lætítia Israël, tu honorificéntia pópuli nostri. Quia fecísti viríliter, et confortátum est cor tuum, eo quod castitátem amáveris et, post virum tuum, álterum nescíeris, ídeo et manus Dómini confortávit te, et ídeo eris benedícta in ætérnum. Et dixit omnis pópulus: Fiat, fiat.

\R\ Dómine Deus, qui cónteris bella ab inítio, éleva bráchium tuum super gentes, quæ cógitant servis tuis mala: * Et déxtera tua glorificétur in nobis. \V\ Allíde virtútem eórum in virtúte tua; cadat robur eórum in iracúndia tua. — Et. Glória Patri. Et.

\feriaheading{Sabbato infra Hebdomadam IV Septembris}

\rubric{De libro Judith}

\lectioheading{Lectio i}{Cap. 16, 22-31}

\lettrine[lines=2]{E}{T} factum est post hæc, omnis pópulus post victóriam venit in Jerúsalem adoráre Dóminum; et mox ut purificáti sunt, obtulérunt omnes holocáusta et vota et repromissiónes suas. Porro Judith univérsa vasa béllica Holoférnis, quæ dedit illi pópulus, et conopéum, quod ipsa sustúlerat de cúbili ipsíus, óbtulit in anáthema oblivíónis.

\R\ Confórta me, Rex, Sanctórum principátum tenens: * Et da sermónem rectum et bene sonántem in os meum. \V\ Dómine, Rex univérsæ potestátis, convérte consílium eórum super eos. — Et.

\lectioheading{Lectio ii}{}

\lettrine[lines=2]{E}{RAT} autem pópulus jucúndus secúndum fáciem sanctórum; et per tres menses gáudium hujus victóriæ celebrátum est cum Judith. Post dies autem illos unusquísque rédiit in domum suam, et Judith magna facta est in Bethúlia et præclárior erat univérsæ terræ Israël. Erat étiam virtúti cástitas adjúncta, ita ut non cognósceret virum ómnibus diébus vitæ suæ, ex quo defúnctus est Manásses vir ejus. Erat autem, diébus festis, procédens cum magna glória.

\R\ Tribulatiónes civitátum audívimus, quas passæ sunt, et defécimus: timor et hebetúdo mentis cécidit super nos et super líberos nostros: ipsi montes nolunt recípere fugam nostram: * Dómine, miserére. \V\ Peccávimus cum pátribus nostris, injúste égimus, iniquitátem fécimus. — Dómine.

\rubric{In Octávis simplícibus et in Offício sanctæ Maríæ in Sábbato: Glória Patri. Dómine.}

\lectioheading{Lectio iii}{}

\lettrine[lines=2]{M}{ANSIT} autem in domo viri sui annos centum quinque, et dimísit abram suam líberam, et defúncta est ac sepúlta cum viro suo in Bethúlia. Luxítque illam omnis pópulus diébus septem. In omni autem spátio vitæ ejus non fuit qui perturbáret Israël, et post mortem ejus annis multis. Dies autem victóriæ hujus festivitátis ab Hebrǽis in número sanctórum diérum accípitur et cólitur a Judǽis ex illo témpore usque in præséntem diem.

\R\ Benedícat te Dóminus in virtúte sua, qui per te ad níhilum redégit inimícos nostros: * Ut non defíciat laus tua de ore hóminum. \V\ Benedíctus Dóminus qui creávit cælum et terram; quia hódie nomen tuum ita magnificávit. — Ut. Glória Patri. Ut.

\feriaheading{Sabbato ante Dominicam V Septembris}

\rubric{nisi hæc occúrrat a die 28 ad 30 Septémbris: quia tunc, sequénti Antíphona omíssa, pónitur Ant. Adapériat, quæ assignátur Sábbato ante Dominicam I Octóbris.}

\hourheading{Ad Vesperas}

\V\ Vespertína orátio ascéndat ad te, Dómine. \R\ Et descéndat super nos misericórdia tua.

\antlabel{Ad Magnif. Ant.} Dómine, Rex omnípotens, * in dicióne tua cuncta sunt pósita, et non est qui possit resístere voluntáti.

\rubric{Orátio Dominicæ occurréntis.}

% ===== BATCH 6 of the Proprium de Tempore (autumnalis.md lines
% 14412-14650) -- the complete fourth week of September: "Dominica
% IV Septembris" Matins opening the Book of Judith in Nocturn I and
% St. Ambrose's "De Elia et jejunio" in Nocturn II (same deferred
% Nocturn III pattern), the six ferial days completing the Judith
% cycle through Holofernes's death and Judith's long widowhood, and
% the brief "Sabbato ante Dominicam V Septembris" (Vespers only,
% with a rubric noting its Magnificat antiphon is replaced by
% "Adapériat" -- assigned to the Saturday before the First Sunday of
% October -- whenever this particular Saturday falls between
% September 28 and 30). Five genuine partial-caps OCR artifacts are
% silently normalized to full caps per the established rule
% (CUmque, MIsit, JOachim, ERat, MAnsit), each confirmed by direct
% inspection of autumnalis.md to be a true partial-caps defect (only
% the first one or two letters capitalized, the remainder lowercase)
% rather than a deliberate small-caps or proper-name styling; none
% carries an accent on its newly-capitalized remainder, so no accent
% is dropped in the normalization. A separate, more pervasive
% formatting quirk also appears starting partway through this
% batch (confirmed via direct grep of autumnalis.md to begin exactly
% at "Feria Tertia infra Hebdomadam IV Septembris" and continue
% through "Sabbato infra Hebdomadam IV Septembris," with zero
% occurrences before that point): many colons in the "In Octávis..."
% rubrics and the Responsories' verse markers are preceded by a
% space ("ritus simplícis :" rather than "ritus simplícis:"). Since
% this reflects only an incidental OCR/typesetting spacing variant,
% not a textual difference, and the surrounding document (including
% this very same batch's earlier material) consistently uses no
% space before a colon, this spacing is silently normalized to match
% the rest of the document, consistent with the established
% treatment of other incidental formatting variations (all-caps
% headings, lowercase Roman numerals) elsewhere in this project. The
% document continues with "Dominica V Septembris" (autumnalis.md
% line 14651), opening the Book of Esther.
% =====
\feriaheading{Dominica V Septembris}

\rubric{Semiduplex}

\nocturnoheading{In I Nocturno}

\rubric{Incipit liber Esther}

\lectioheading{Lectio i}{Cap. 1, 1-9}

\lettrine[lines=2]{I}{N} diébus Assuéri, qui regnávit ab India usque Æthiópiam super centum vigínti septem províncias, quando sedit in sólio regni sui, Susan cívitas regni ejus exórdium fuit. Tértio ígitur anno impérii sui, fecit grande convívium cunctis princípibus et púeris suis fortíssimis Persárum et Medórum ínclytis et præféctis provinciárum coram se, ut osténderet divítias glóriæ regni sui, ac magnitúdinem atque jactántiam poténtiæ suæ, multo témpore, centum vidélicet et octogínta diébus.

\R\ Dómine, mi Rex omnípotens, in dicióne tua cuncta sunt pósita, et non est qui possit resístere voluntáti tuæ: * Líbera nos propter nomen tuum. \V\ Exáudi oratiónem nostram, et convérte luctum nostrum in gáudium. — Líbera.

\lectioheading{Lectio ii}{}

\lettrine[lines=2]{C}{UMQUE} implerántur dies convívii, invitávit omnem pópulum, qui invéntus est in Susan, a máximo usque ad mínimum; et jussit septem diébus convívium præparári in vestíbulo horti et némoris, quod régio cultu et manu cónsitum erat. Et pendébant ex omni parte tentória aérii colóris et carbásini ac hyacínthini sustentáta fúnibus býssinis atque purpúreis, qui ebúrneis círculis insérti erant et colúmnis marmóreis fulciebántur. Léctuli quoque áurei et argéntei super paviméntum smarágdino et pário stratum lápide dispósiti erant, quod mira varietáte pictúra decorábat.

\R\ Confórta me, Rex, Sanctórum principátum tenens: * Et da sermónem rectum et bene sonántem in os meum. \V\ Dómine, Rex univérsæ potestátis, convérte consílium eórum super eos. — Et.

\lectioheading{Lectio iii}{}

\lettrine[lines=2]{B}{IBEBANT} autem, qui invitáti erant, áureis póculis et áliis atque áliis vasis cibi inferebántur. Vinum quoque, ut magnificéntia régia dignum erat, abúndans et præcípuum ponebátur. Nec erat qui noléntes cógeret ad bibéndum, sed, sicut rex statúerat præpónens mensis síngulos de princípibus suis, ut súmeret unusquísque quod vellet. Vasthi quoque regína fecit convívium feminárum in palátio, ubi rex Assuérus manére consuéverat.

\R\ Spem in álium numquam hábui, prætérquam in te, Deus Israël: * Qui irásceris, et propítius eris: et ómnia peccáta hóminum in tribulatióne dimíttis. \V\ Dómine Deus, Creátor cæli et terræ, réspice ad humilitátem nostram. — Qui. Glória Patri. Qui.

\nocturnoheading{In II Nocturno}

\rubric{Ex libro Officiórum sancti Ambrósii Epíscopi}

\lectioheading{Lectio iv}{Lib. 3, cap. 15}

\lettrine[lines=2]{Q}{UID} Esther regína, nonne ut pópulum suum perículo erúeret (quod erat decórum atque honéstum) morti se óbtulit, nec immítis regis trepidávit furórem? Ipse quoque rex Persárum, ferox atque túmido corde, tamen decórum judicávit índici insidiárum, quæ sibi parátæ forent, grátiam repræsentáre, populúmque líberum a servitúte erípere, erúere neci, nec párcere neci ejus, qui tam indecóra suasísset. Dénique quem secúndum a se, ac præcípuum inter omnes amícos habéret, cruci trádidit, quod dehonestátum se ejus frauduléntis consíliis animadvertísset.

\R\ Meménto mei, Dómine Deus, in bonum: * Et ne déleas miseratiónes meas quas feci in domo Dei mei et in cæremóniis ejus. \V\ Recordáre mei, Dómine, Deus meus. — Et.

\lectioheading{Lectio v}{Cap. 16}

\lettrine[lines=2]{E}{A} enim amicítia probábilis, quæ honestátem tuétur, præferénda sane ópibus, honóribus, potestátibus; honestáti vero præférri non solet, sed honestátem sequi. Qualis fuit Jónathæ, qui pro pietáte nec offénsam patris, nec salútis perículum refugiébat. Qualis fuit Achímelech, qui pro hospitális grátiæ offíciis necem pótius sibi, quam proditiónem fugiéntis amíci, subeúndam arbitrabátur. Nihil ígitur præferéndum honestáti; quæ tamen ne amicítiæ stúdio prætereátur, étiam hoc Scriptúra ádmonet.

\R\ Tribulatiónes civitátum audívimus, quas passæ sunt, et defécimus: timor et hebetúdo mentis cécidit super nos et super líberos nostros: ipsi montes nolunt recípere fugam nostram: * Dómine, miserére. \V\ Peccávimus cum pátribus nostris, injúste égimus, iniquitátem fécimus. — Dómine.

\lectioheading{Lectio vi}{}

\lettrine[lines=2]{S}{UNT} enim pleræque philosophórum quæstiónes: Utrum amíci causa quisquam contra pátriam sentíre necne débeat, ut amíco obédiat: utrum opórteat ut fidem déserat, dum indúlget atque inténdit amíci commoditátibus. Et Scriptúra quidem ait: Clava, et gládius, et sagítta ferráta, sic homo est testimónium dans falsum advérsus amícum suum. Sed consídera quid ástruat. Non testimónium reprehéndit dictum in amícum, sed falsum testimónium. Quid enim si Dei causa, quid si pátriæ, cogátur áliquis dícere testimónium? Numquid præponderáre debet amicítia religióni, præponderáre caritáti vítium?

\R\ Benedícat te Dóminus in virtúte sua, qui per te ad níhilum redégit inimícos nostros: * Ut non defíciat laus tua de ore hóminum. \V\ Benedíctus Dóminus qui creávit cælum et terram; quia hódie nomen tuum ita magnificávit. — Ut. Glória Patri. Ut.

\R\ Nos álium Deum nescímus præter Dóminum, in quo sperámus: * Qui non déspicit nos, nec ámovet salútem suam a génere nostro. \V\ Indulgéntiam ipsíus fusis lácrimis postulémus, et humiliémus illi ánimas nostras. — Qui.

\R\ Duo Séraphim, \rubric{ut infra in Homilíis.}

\feriaheading{Feria Secunda infra Hebdomadam V Septembris}

\rubric{De libro Esther}

\lectioheading{Lectio i}{Cap. 2, 5-11}

\lettrine[lines=2]{E}{RAT} vir Judǽus in Susan civitáte, vocábulo Mardochǽus, fílius Jair, fílii Sémei, fílii Cis, de stirpe Jémini, qui translátus fúerat de Jerúsalem eo témpore quo Jechoníam regem Juda Nabuchodónosor rex Babylónis transtúlerat. Qui fuit nutrícius fíliæ fratris sui Edíssæ, quæ áltero nómine vocabátur Esther, et utrúmque paréntem amíserat: pulchra nimis et decóra fácie. Mortuísque patre ejus ac matre, Mardochǽus sibi eam adoptávit in fíliam.

\R\ Dómine, mi Rex omnípotens, in dicióne tua cuncta sunt pósita, et non est qui possit resístere voluntáti tuæ: * Líbera nos propter nomen tuum. \V\ Exáudi oratiónem nostram, et convérte luctum nostrum in gáudium. — Líbera.

\lectioheading{Lectio ii}{}

\lettrine[lines=2]{C}{UMQUE} percrebuísset regis impérium, et juxta mandátum illíus multæ pulchræ vírgines adduceréntur Susan et Egéo traderéntur eunúcho; Esther quoque inter céteras puéllas ei trádita est, ut servarétur in número feminárum. Quæ plácuit ei et invénit grátiam in conspéctu illíus. Et præcépit eunúcho ut acceleráret mundum muliébrem et tráderet ei partes suas et septem puéllas speciosíssimas de domo regis, et tam ipsam quam pedíssequas ejus ornáret atque excóleret. Quæ nóluit indicáre ei pópulum et pátriam suam; Mardochǽus enim præcéperat ei, ut de hac re omníno reticéret. Qui deambulábat quotídie ante vestíbulum domus, in qua eléctæ vírgines servabántur, curam agens salútis Esther.

\R\ Confórta me, Rex, Sanctórum principátum tenens: * Et da sermónem rectum et bene sonántem in os meum. \V\ Dómine, Rex univérsæ potestátis, convérte consílium eórum super eos. — Et.

\rubric{In Octávis et Festis ritus simplícis: Glória Patri. Et.}

\lectioheading{Lectio iii}{Cap. 2, 15-17}

\lettrine[lines=2]{E}{VOLUTO} autem témpore per órdinem, instábat dies quo Esther fília Abíhail fratris Mardochǽi, quam sibi adoptáverat in fíliam, debéret intráre ad regem. Quæ non quæsívit muliébrem cultum, sed quæcúmque vóluit Egéus eunúchus custos vírginum, hæc ei ad ornátum dedit; erat enim formósa valde et incredíbili pulchritúdine: ómnium óculis gratiósa et amábilis videbátur. Ducta est ítaque ad cubículum regis Assuéri mense décimo, qui vocátur Tebeth, séptimo anno regni ejus. Et adamávit eam rex plus quam omnes mulíeres, habuítque grátiam et misericórdiam coram eo super omnes mulíeres, et pósuit diadéma regni in cápite ejus fecítque eam regnáre in loco Vasthi.

\R\ Spem in álium numquam hábui, prætérquam in te, Deus Israël: * Qui irásceris, et propítius eris: et ómnia peccáta hóminum in tribulatióne dimíttis. \V\ Dómine Deus, Creátor cæli et terræ, réspice ad humilitátem nostram. — Qui. Glória Patri. Qui.

\feriaheading{Feria Tertia infra Hebdomadam V Septembris}

\rubric{De libro Esther}

\lectioheading{Lectio i}{Cap. 3, 1-7}

\lettrine[lines=2]{P}{OST} hæc rex Assuérus exaltávit Aman fílium Amadáthi, qui erat de stirpe Agag, et pósuit sólium ejus super omnes príncipes quos habébat. Cunctíque servi regis, qui in fóribus palátii versabántur, flectébant génua et adorábant Aman; sic enim præcéperat eis imperátor. Solus Mardochǽus non flectébat genu neque adorábat eum. Cui dixérunt púeri regis, qui ad fores palátii præsidébant: Cur præter céteros non servas mandátum regis?

\R\ Meménto mei, Dómine Deus, in bonum: * Et ne déleas miseratiónes meas quas feci in domo Dei mei et in cæremóniis ejus. \V\ Recordáre mei, Dómine, Deus meus. — Et.

\lectioheading{Lectio ii}{}

\lettrine[lines=2]{C}{UMQUE} hoc crébrius dícerent, et ille nollet audíre, nuntiavérunt Aman, scire cupiéntes utrum perseváret in senténtia; díxerat enim eis se esse Judǽum. Quod cum audísset Aman et experiménto probásset quod Mardochǽus non fléctere sibi genu nec se adoráret, irátus est valde et pro níhilo duxit in unum Mardochǽum míttere manus suas.

\R\ Tribulatiónes civitátum audívimus, quas passæ sunt, et defécimus: timor et hebetúdo mentis cécidit super nos et super líberos nostros: ipsi montes nolunt recípere fugam nostram: * Dómine, miserére. \V\ Peccávimus cum pátribus nostris, injúste égimus, iniquitátem fécimus. — Dómine.

\rubric{In Octávis et Festis ritus simplícis: Glória Patri. Dómine.}

\lectioheading{Lectio iii}{}

\lettrine[lines=2]{A}{UDIERAT} enim quod esset gentis Judǽæ, magísque vóluit omnem Judæórum rum, qui erant in regno Assuéri, pérdere natiónem. Mense primo, cujus vocábulum est Nisan, anno duodécimo regni Assuéri, missa est sors in urnam, quæ Hebráice dícitur phur, coram Aman, quo die et quo mense gens Judæórum debéret intérfici; et exívit mensis duodécimus qui vocátur Adar.

\R\ Benedícat te Dóminus in virtúte sua, qui per te ad níhilum redégit inimícos nostros: * Ut non defíciat laus tua de ore hóminum. \V\ Benedíctus Dóminus qui creávit cælum et terram; quia hódie nomen tuum ita magnificávit. — Ut. Glória Patri. Ut.

\feriaheading{Feria Quarta infra Hebdomadam V Septembris}

\rubric{De libro Esther}

\lectioheading{Lectio i}{Cap. 4, 1-17}

\lettrine[lines=2]{Q}{UÆ} cum audísset Mardochǽus, scidit vestiménta sua et indútus est sacco spargens cínerem cápiti; et in platéa médiæ civitátis voce magna clamábat osténdens amaritúdinem ánimi sui, et hoc ejulátu usque ad fores palátii grádiens; non enim erat lícitum indútum sacco aulam regis intráre. In ómnibus quoque províncíis, óppidis ac locis, ad quæ crudéle regis dogma pervénerat, planctus ingens erat apud Judǽos, jejúnium, ululátus et fletus, sacco et cínere multis pro strato uténtibus. Ingréssæ autem sunt puéllæ Esther et eunúchi nuntiaverúntque ei. Quod áudiens consternáta est et vestem misit, ut, abláto sacco, indúerent eum: quam accípere nóluit. Accitóque Athach eunúcho, quem rex minístrum ei déderat, præcépit ei ut iret ad Mardochǽum et discéret ab eo cur hoc fáceret.

\R\ Nos álium Deum nescímus præter Dóminum, in quo sperámus: * Qui non déspicit nos, nec ámovet salútem suam a génere nostro. \V\ Indulgéntiam ipsíus fusis lácrimis postulémus, et humiliémus illi ánimas nostras. — Qui.

\lectioheading{Lectio ii}{}

\lettrine[lines=2]{E}{GRESSUSQUE} Athach ivit ad Mardochǽum stantem in platéa civitátis ante óstium palátii. Qui indicávit ei ómnia quæ accíderant, quómodo Aman promisísset, ut in thesáuros regis pro Judæórum nece inférret argéntum; exémplar quoque edícti, quod pendébat in Susan, dedit ei ut regínæ osténderet et monéret eam ut intráret ad regem et deprecarétur eum pro pópulo suo. Regréssus Athach nuntiávit Esther ómnia quæ Mardochǽus díxerat. Quæ respóndit ei et jussit ut díceret Mardochǽo: Omnes servi regis et cunctæ quæ sub dicióne ejus sunt norunt províncíæ, quod sive vir sive múlier non vocátus intérius átrium regis intráverit, absque ulla cunctatióne statim interficiátur, nisi forte rex áuream virgam ad eum teténderit pro signo cleméntiæ, atque ita possit vívere. Ego ígitur quómodo ad regem intráre pótero, quæ trigínta jam diébus non sum vocáta ad eum?

\R\ Confórta me, Rex, Sanctórum principátum tenens: * Et da sermónem rectum et bene sonántem in os meum. \V\ Dómine, Rex univérsæ potestátis, convérte consílium eórum super eos. — Et.

\rubric{In Octávis et Festis ritus simplícis: Glória Patri. Et.}

\lectioheading{Lectio iii}{}

\lettrine[lines=2]{Q}{UOD} cum audísset Mardochǽus, rursum mandávit Esther dicens: Ne putes quod ánimam tuam tantum líberes, quia in domo regis es, præ cunctis Judǽis; si enim nunc silúeris, per áliam occasiónem liberabúntur Judǽi, et tu et domus patris tui períbitis. Et quis novit utrum idcírco ad regnum véneris, ut in tali témpore pararéris? Rursúmque Esther hæc Mardochǽo verba mandávit: Vade et cóngrega omnes Judǽos, quos in Susan repéreris, et oráte pro me. Non comedátis et non bibátis tribus diébus et tribus nóctibus, et ego cum ancíllis meis simíliter jejunábo; et tunc ingrédiar ad regem contra legem fáciens non vocáta tradénsque me morti et perículo. Ivit ítaque Mardochǽus, et fecit ómnia quæ ei Esther præcéperat.

\R\ Spem in álium numquam hábui, prætérquam in te, Deus Israël: * Qui irásceris, et propítius eris: et ómnia peccáta hóminum in tribulatióne dimíttis. \V\ Dómine Deus, Creátor cæli et terræ, réspice ad humilitátem nostram. — Qui. Glória Patri. Qui.

\feriaheading{Feria Quinta infra Hebdomadam V Septembris}

\rubric{De libro Esther}

\lectioheading{Lectio i}{Cap. 5, 1-5}

\lettrine[lines=2]{D}{IE} autem tértio indúta est Esther regálibus vestiméntis et stetit in átrio domus régiæ, quod erat intérius contra basílicam regis. At ille sedébat super sólium suum in consistório palátii contra óstium domus. Cumque vidísset Esther regínam stantem, plácuit óculis ejus et exténdit contra eam virgam áuream quam tenébat manu. Quæ accédens osculáta est summitátem virgæ ejus, dixítque ad eam rex: Quid vis, Esther regína? quæ est petítio tua? Etiam si dimídiam partem regni petíeris, dábitur tibi. At illa respóndit: Si regi placet, óbsecro ut vénias ad me hódie et Aman tecum ad convívium quod parávi. Statímque rex, Vocáte, inquit, cito Aman ut Esther obédiat voluntáti.

\R\ Dómine, mi Rex omnípotens, in dicióne tua cuncta sunt pósita, et non est qui possit resístere voluntáti tuæ: * Líbera nos propter nomen tuum. \V\ Exáudi oratiónem nostram, et convérte luctum nostrum in gáudium. — Líbera.

\lectioheading{Lectio ii}{Cap. 5, 9-14}

\lettrine[lines=2]{E}{GRESSUS} est ítaque illo die Aman lætus et álacer. Cumque vidísset Mardochǽum sedéntem ante fores palátii, et non solum non assurrexísse sibi, sed nec motum quidem de loco sessiónis suæ, indignátus est valde; et, dissimuláta ira, revérsus in domum suam convocávit ad se amícos suos et Zares uxórem suam et expósuit illis magnitúdinem divitiárum suárum filiorúmque turbam et quanta eum glória super omnes príncipes et servos suos rex elevásset. Et post hæc ait: Regína quoque Esther nullum álium vocávit ad convívium cum rege, præter me, apud quam étiam cras cum rege pransúrus sum; et, cum hæc ómnia hábeam, nihil me habére puto, quámdiu vídero Mardochǽum Judǽum sedéntem ante fores régias.

\R\ Confórta me, Rex, Sanctórum principátum tenens: * Et da sermónem rectum et bene sonántem in os meum. \V\ Dómine, Rex univérsæ potestátis, convérte consílium eórum super eos. — Et.

\rubric{In Octávis et Festis ritus simplícis: Glória Patri. Et.}

\lectioheading{Lectio iii}{}

\lettrine[lines=2]{R}{ESPONDERUNTQUE} ei Zares uxor ejus et céteri amíci: Jube parári excélsam trabem habéntem altitúdinis quinquagínta cúbitos, et dic mane regi ut appendátur super eam Mardochǽus; et sic ibis cum rege lætus ad convívium. Plácuit ei consílium et jussit excélsam paráti crucem.

\R\ Spem in álium numquam hábui, prætérquam in te, Deus Israël: * Qui irásceris, et propítius eris: et ómnia peccáta hóminum in tribulatióne dimíttis. \V\ Dómine Deus, Creátor cæli et terræ, réspice ad humilitátem nostram. — Qui. Glória Patri. Qui.

\feriaheading{Feria Sexta infra Hebdomadam V Septembris}

\rubric{De libro Esther}

\lectioheading{Lectio i}{Cap. 6, 1-13}

\lettrine[lines=2]{N}{OCTEM} illam duxit rex insómnem jussítque sibi afférri histórias et annáles priórum témporum. Quæ cum illo præsénte legeréntur, ventum est ad illum locum, ubi scriptum erat quómodo nuntiásset Mardochǽus insídias Bágathan et Thares eunuchórum regem Assuérum juguláre cupiéntium. Quod cum audísset rex ait: Quid pro hac fide honóris ac præmii Mardochǽus consecútus est? Dixérunt ei servi illíus ac minístri: Níhil omníno mercédis accépit. Statímque rex, Quis est, inquit, in átrio? Aman quippe intérius átrium domus régiæ intráverat, ut suggéreret regi et jubéret Mardochǽum affígi patíbulo, quod ei fúerat præparátum. Respondérunt púeri: Aman stat in átrio. Dixítque rex: Ingrediátur.

\R\ Meménto mei, Dómine Deus, in bonum: * Et ne déleas miseratiónes meas quas feci in domo Dei mei et in cæremóniis ejus. \V\ Recordáre mei, Dómine, Deus meus. — Et.

\lectioheading{Lectio ii}{}

\lettrine[lines=2]{C}{UMQUE} esset ingréssus ait illi: Quid debet fíeri viro, quem rex honoráre desíderat? Cógitans autem in corde suo Aman et réputans quod nullum álium rex nisi se vellet honoráre, respóndit: Homo quem rex honoráre cupit debet índui véstibus régiis et impóni super equum qui de sella regis est, et accípere régium diadéma super caput suum, et primus de régiis princípibus ac tyránnis téneat equum ejus et per platéam civitátis incédens clamet et dicat: Sic honorábitur quemcúmque volúerit rex honoráre.

\R\ Tribulatiónes civitátum audívimus, quas passæ sunt, et defécimus: timor et hebetúdo mentis cécidit super nos et super líberos nostros: ipsi montes nolunt recípere fugam nostram: * Dómine, miserére. \V\ Peccávimus cum pátribus nostris, injúste égimus, iniquitátem fécimus. — Dómine, miserére.

\rubric{In Octávis et Festis ritus simplícis: Glória Patri. Dómine.}

\lectioheading{Lectio iii}{}

\lettrine[lines=2]{D}{IXITQUE} ei rex: Festína et, sumpta stola et equo, fac, ut locútus es, Mardochǽo Judǽo, qui sedet ante fores palátii. Cave ne quidquam de his quæ locútus es prætermíttas. Tulit ítaque Aman stolam et equum, indutúmque Mardochǽum in platéa civitátis et impósitum equo præcedébat atque clamábat: Hoc honóre condígnus est quemcúmque rex volúerit honoráre. Reversúsque est Mardochǽus ad jánuam palátii, et Aman festinávit ire in domum suam lugens et opérto cápite. Narravítque Zares uxóri suæ et amícis ómnia quæ eveníssent sibi.

\R\ Benedícat te Dóminus in virtúte sua, qui per te ad níhilum redégit inimícos nostros: * Ut non defíciat laus tua de ore hóminum. \V\ Benedíctus Dóminus qui creávit cælum et terram; quia hódie nomen tuum ita magnificávit. — Ut. Glória Patri. Ut.

\feriaheading{Sabbato infra Hebdomadam V Septembris}

\rubric{De libro Esther}

\lectioheading{Lectio i}{Cap. 7, 1-10}

\lettrine[lines=2]{I}{NTRAVIT} ítaque rex et Aman, ut bíberent cum regína. Dixítque ei rex étiam secúnda die, postquam vino incalúerat: Quæ est petítio tua, Esther, ut detur tibi, et quid vis fíeri? Etiam si dimídiam partem regni petíeris, impetrábis. Ad quem illa respóndit: Si invéni grátiam in óculis tuis, o rex, et si tibi placet, dona mihi ánimam meam, pro qua rogo, et pópulum meum, pro quo óbsecro. Tráditi enim sumus ego et pópulus meus, ut conterámur, jugulémur et pereámus. Atque, útinam in servos et fámulas venderémur: esset tolerábile malum et gemens tacérem; nunc autem hostis noster est cujus crudélitas redúndat in regem.

\R\ Nos álium Deum nescímus præter Dóminum, in quo sperámus: * Qui non déspicit nos, nec ámovet salútem suam a génere nostro. \V\ Indulgéntiam ipsíus fusis lácrimis postulémus, et humiliémus illi ánimas nostras. — Qui.

\lectioheading{Lectio ii}{}

\lettrine[lines=2]{R}{ESPONDENSQUE} rex Assuérus ait: Quis est iste et cujus poténtiæ, ut hæc áudeat fácere? Dixítque Esther: Hostis et inimícus noster péssimus iste est Aman. Quod ille áudiens íllico obstúpuit vultum regis ac regínæ ferre non sústinens. Rex autem irátus surréxit et de loco convívii intrávit in hortum arbóribus cónsitum. Aman quoque surréxit ut rogáret Esther regínam pro ánima sua; intelléxit enim a rege sibi parátum malum.

\R\ Confórta me, Rex, Sanctórum principátum tenens: * Et da sermónem rectum et bene sonántem in os meum. \V\ Dómine, Rex univérsæ potestátis, convérte consílium eórum super eos. — Et.

\rubric{In Octávis simplícibus et in Offício Sanctæ Maríæ in Sábbato: Glória Patri. Et.}

\lectioheading{Lectio iii}{}

\lettrine[lines=2]{Q}{UI} cum revérsus esset de horto nemóribus cónsito et intrásset convívii locum, réperit Aman super léctulum corruísse in quo jacébat Esther et ait: Etiam regínam vult opprímere, me præsénte, in domo mea. Necdum verbum de ore regis exíerat, et statim operuérunt fáciem ejus. Dixítque Harbóna unus de eunúchis qui stabant in ministério regis: En lignum, quod paráverat Mardochǽo qui locútus est pro rege, stat in domo Aman habens altitúdinis quinquagínta cúbitos. Cui dixit rex: Appéndite eum in eo. Suspénsus est ítaque Aman in patíbulo, quod paráverat Mardochǽo, et regis ira quiévit.

\R\ Spem in álium numquam hábui, prætérquam in te, Deus Israël: * Qui irásceris, et propítius eris: et ómnia peccáta hóminum in tribulatióne dimíttis. \V\ Dómine Deus, Creátor cæli et terræ, réspice ad humilitátem nostram. — Qui. Glória Patri. Qui.

\feriaheading{Sabbato ante Dominicam I Octobris}

\hourheading{Ad Vesperas}

\V\ Vespertína orátio ascéndat ad te, Dómine. \R\ Et descéndat super nos misericórdia tua.

\antlabel{Ad Magnif. Ant.} Adapériat Dóminus * cor vestrum in lege sua et in præcéptis suis et fáciat pacem Dóminus Deus noster.

\rubric{Oratio Dominicæ occurréntis.}

% ===== BATCH 7 of the Proprium de Tempore (autumnalis.md lines
% 14651-14883) -- the complete fifth and final week of September:
% "Dominica V Septembris" Matins opening the Book of Esther in
% Nocturn I and St. Ambrose's "De Officiis" in Nocturn II (this
% Sunday's Nocturn III responsories given without the "vii"/"viii"
% numbering used at the parallel point in Dominicas I-IV, a genuine
% source variation transcribed exactly as given), the six ferial
% days completing the Esther cycle through Haman's downfall, and
% the brief "Sabbato ante Dominicam I Octobris" Vespers with the
% "Adapériat Dóminus" antiphon that the governing rubric several
% batches ago had flagged as proper to this point. This closes out
% the entire month of September in the Proprium de Tempore. A
% dozen-plus genuine partial-caps OCR artifacts are silently
% normalized to full caps per the established rule (ERat, CUmque
% [several occurrences], EVolúto, POst, AUdíerat, EGressúsque, QUod,
% EGréssus, RESponderúntque, NOctem, DIxítque, INtrávit,
% RESpondénsque), each confirmed by direct inspection of
% autumnalis.md; accents on each word's newly-capitalized remainder
% are dropped per the standing rule, while complete short words
% already fully capped in the source (QUÆ, QUI, DIE) are left
% untouched. The pervasive space-before-colon formatting quirk
% identified in the previous batch continues throughout this one and
% is likewise silently normalized to standard colon spacing. The
% document continues with "Dominica I Octobris" (autumnalis.md line
% 14884), opening the Books of Machabees and the month of October.
% =====
\feriaheading{Dominica I Octobris}

\rubric{Semiduplex}

\hourheading{Ad Matutinum}

\rubric{Invitatorium Adorémus Dóminum et Hymnus Primo die, ut in Psaltério 32: quæ dicúntur étiam in réliquis Dominícis usque ad V Novémbris inclusíve.}

\nocturnoheading{In I Nocturno}

\rubric{Incipit liber primus Machabæórum}

\lectioheading{Lectio i}{Cap. 1, 1-16}

\lettrine[lines=2]{E}{T} factum est, postquam percússit Alexánder Philíppi Mácedo, qui primus regnávit in Grǽcia, egréssus de terra Cethim, Dárium regem Persárum et Medórum, constítuit prǽlia multa et obtínuit ómnium munitiónes et interfécit reges terræ et pertránsiit usque ad fines terræ et accépit spólia multitúdinis géntium, et síluit terra in conspéctu ejus. Et congregávit virtútem et exércitum fortem nimis, et exaltátum est et elevátum cor ejus, et obtínuit regiónes géntium et tyránnos, et facti sunt illi in tribútum. Et post hæc décidit in lectum et cognóvit quia moerétur et vocávit púeros suos nóbiles, qui secum erant nutríti a juventúte, et divísit illis regnum suum, cum adhuc víveret.

\R\ Adapériat Dóminus cor vestrum in lege sua et in præcéptis suis et fáciat pacem in diébus vestris: * Concédat vobis salútem, et rédimat vos a malis. \V\ Exáudiat Dóminus oratiónes vestras et reconciliétur vobis nec vos déserat in témpore malo. — Concédat.

\lectioheading{Lectio ii}{}

\lettrine[lines=2]{E}{T} regnávit Alexánder annis duódecim et mórtuus est. Et obtinuérunt púeri ejus regnum, unusquísque in loco suo, et imposuérunt omnes sibi diadémata post mortem ejus, et fílii eórum post eos annis multis. Et multiplicáta sunt mala in terra. Et éxiit ex eis radix peccátrix, Antíochus illústris, fílius Antíochi regis qui fúerat Romæ obses, et regnávit in anno centésimo trigésimo séptimo regni Grræcórum.

\R\ Exáudiat Dóminus oratiónes vestras, et reconciliétur vobis nec vos déserat in témpore malo * Dóminus, Deus noster. \V\ Det vobis cor ómnibus, ut colátis eum et faciátis ejus voluntátem. — Dóminus.

\lectioheading{Lectio iii}{}

\lettrine[lines=2]{I}{N} diébus illis exiérunt ex Israël fílii iníqui et suasérunt multis dicéntes: Eámus et disponámus testaméntum cum géntibus, quæ circa nos sunt, quia, ex quo recéssimus ab eis, invenérunt nos multa mala. Et bonus visus est sermo in óculis eórum. Et destinavérunt áliqui de pópulo et abiérunt ad regem, et dedit illis potestátem ut fácerent justítiam géntium; et ædificavérunt gymnásium in Jerosólymis secúndum leges natiónum et fecérunt sibi præpútia et recessérunt a testaménto sancto et juncti sunt natiónibus et venúmdati sunt ut fácerent malum.

\R\ Congregáti sunt inimíci nostri, et gloriántur in virtúte sua: cóntere fortitúdinem illórum, Dómine, et dispérge illos: * Ut cognóscant quia non est álius qui pugnet pro nobis, nisi tu, Deus noster. \V\ Dispérge illos in virtúte tua, et déstrue eos, protéctor noster, Dómine. — Ut. Glória Patri. Ut.

\nocturnoheading{In II Nocturno}

\rubric{Ex libro Officiórum sancti Ambrósii Epíscopi}

\lectioheading{Lectio iv}{Lib. 1, cap. 40}

\lettrine[lines=2]{F}{ORTASSE} áliquos béllica defíxos glória tenet, ut putent solam esse præliárem fortitúdinem; et ídeo me ad ista deflexísse, quia illa nostris déforet. Quam fortis Jesus Nave, ut uno prǽlio quinque reges captos stérneret cum pópulis suis! Deínde cum advérsum Gabaonítas urgéret prǽlium, et vererétur ne nox impedíret victóriam, magnitúdine mentis et fídei clamávit: Stet sol, et stetit, donec victória consummarétur. Gédeon in trecéntis viris de ingénti pópulo et acérbo hoste revéxit triúmphum. Jónathas adoléscens virtútem magnam fecit in prǽlio.

\R\ Impetum inimicórum ne timuéritis: mémores estóte quómodo salvi facti sunt patres nostri: * Et nunc clamémus in cælum et miserébitur nostri Deus noster. \V\ Mementóte mirabílium ejus, quæ fecit pharaóni et exercítui ejus in Mari Rubro. — Et.

\lectioheading{Lectio v}{}

\lettrine[lines=2]{Q}{UID} de Machabǽis loquar? Sed prius de pópulo dicam patrum; qui cum essent paráti ad repugnándum pro templo Dei et pro legítimis suis, dolo hóstium die lacessíti sábbati, maluérunt vulnéribus offérre nuda córpora, quam repugnáre, ne violárent sábbatum. Itaque omnes læti se obtulérunt morti. Sed Machabǽi considerántes quod hoc exémplo gens omnis posset períre, sábbato étiam, cum ipsi in bellum provocaréntur, ulti sunt innocéntium necem fratrum suórum. Unde póstea stimulátus rex Antíochus, cum bellum accénderet per duces suos Lýsiam, Nicánorem, Górgiam, ita cum Orientálibus suis et Assýriis attrítus est cópiis, ut quadragínta et octo míllia in médio campi a tribus míllibus prosterneréntur.

\R\ Congregátæ sunt gentes in multitúdine, ut dímicent contra nos, et ignorámus quid ágere debeámus: * Dómine Deus, ad te sunt óculi nostri, ne pereámus. \V\ Tu scis quæ cógitant in nos: quómodo potérimus subsístere ante fáciem illórum, nisi tu ádjuves nos? — Dómine.

\lectioheading{Lectio vi}{}

\lettrine[lines=2]{V}{IRTUTEM} ducis Judæ Machabǽi de uno ejus mílite consideráte. Namque Eleázarus, cum supereminéntem céteris elephántem lorica vestítum régia advérteret, arbitrátus quod in eo esset rex, cursu cóncito in médium legiónis se prorúpit: et, abjécto clýpeo, utráque manu interficiébat, donec perveníret ad béstiam, atque intrávit sub eam, et subjécto gládio interémit eam. Itaque cadens béstia oppréssit Eleázarum, atque ita mórtuus est. Quanta ígitur virtus ánimi! primo, ut mortem non timéret; deínde, ut circumfúsus legiónibus inimicórum, in confértos raperétur hostes, médium penetráret agmen, et contémpta morte ferócior, abjécto clýpeo, utráque manu vulnerátæ molem béstiæ subíret ac sustinéret: post infra ipsam succéderet, quo plenióri feríret ictu; cujus ruína inclúsus magis quam oppréssus, suo est sepúltus triúmpho.

\R\ Tua est poténtia, tuum regnum, Dómine: tu es super omnes gentes: * Da pacem, Dómine, in diébus nostris. \V\ Creátor ómnium, Deus, terríbilis et fortis, justus, et miséricors. — Da. Glória Patri. Da.

\R\ Refúlsit sol in clýpeos áureos, et resplenduérunt montes ab eis: * Et fortitúdo géntium dissipáta est. \V\ Erat enim exércitus magnus valde et fortis: et appropiávit Judas et exércitus ejus in prǽlio. — Et.

\R\ Duo Séraphim, \rubric{ut infra in Homilíis.}

\rubric{Ad Laudes, in ómnibus Domínicis usque ad V Novémbris inclusíve, Hymnus Ætérne rerum, ut in Psaltério.}

\feriaheading{Feria Secunda infra Hebdomadam I Octobris}

\rubric{De libro primo Machabæórum}

\lectioheading{Lectio i}{Cap. 1, 17-29}

\lettrine[lines=2]{E}{T} parátum est regnum in conspéctu Antíochi, et cœpit regnáre in terra Ægýpti, ut regnáret super duo regna, et intrávit in Ægýptum in multitúdine gravi, in cúrribus et elephántis et equítibus et copiósa návium multitúdine, et constítuit bellum advérsus Ptolemǽum regem Ægýpti, et véritus est Ptolemǽus a fácie ejus et fugit, et cecidérunt vulneráti multi, et comprehéndit civitátes munítas in terra Ægýpti et accépit spólia terræ Ægýpti.

\R\ Dixit Judas Simóni fratri suo: Elige tibi viros et vade, líbera fratres tuos in Galilǽam; ego autem et Jónathas frater tuus íbimus in Galaadítim: * Sicut fúerit volúntas in cælo, sic fiat. \V\ Accingímini, fílii poténtes, et estóte paráti: quóniam mélius est nobis mori in bello, quam vidére mala gentis nostræ et sanctórum. — Sicut.

\lectioheading{Lectio ii}{}

\lettrine[lines=2]{E}{T} convértit Antíochus, postquam percússit Ægýptum in centésimo et quadragésimo tértio anno, et ascéndit ad Israël. Et ascéndit Jerosólymam in multitúdine gravi et intrávit in sanctificatiónem cum supérbia et accépit altáre áureum et candelábrum lúminis et univérsa vasa ejus et mensam propositiónis et libatória et phíalas et mortaríola áurea et velum et corónas et ornaméntum áureum, quod in fácie templi erat, et commínuit ómnia.

\R\ Ornavérunt fáciem templi corónis áureis, et dedicavérunt altáre Dómino: * Et facta est lætítia magna in pópulo. \V\ In hymnis et confessiónibus benedicébant Dóminum. — Et.

\rubric{In Octávis et Festis ritus simplícis: Glória Patri. Et.}

\lectioheading{Lectio iii}{}

\lettrine[lines=2]{E}{T} accépit argéntum et aurum et vasa concupiscibília et accépit thesáuros occúltos, quos invénit, et, sublátis ómnibus, ábiit in terram suam et fecit cædem hóminum et locútus est in supérbia magna et factus est planctus magnus in Israël et in omni loco eórum, et ingemuérunt príncipes et senióres, vírgines et júvenes infirmáti sunt, et speciósitas mulíerum immutáta est. Omnis marítus sumpsit laméntum, et quæ sedébant in thoro maritáli lugébant, et commóta est terra super habitántes in ea, et univérsa domus Jacob índuit confusiónem.

\R\ In hymnis et confessiónibus benedicébant Dóminum, * Qui magna fecit in Israël, et victóriam dedit illis Dóminus omnípotens. \V\ Ornavérunt fáciem templi corónis áureis, et dedicavérunt altáre Dómino. — Qui. Glória Patri. Qui.

\feriaheading{Feria Tertia infra Hebdomadam I Octobris}

\rubric{De libro primo Machabæórum}

\lectioheading{Lectio i}{Cap. 2, 1-10}

\lettrine[lines=2]{I}{N} diébus illis surréxit Mathathías fílius Joánnis, fílii Simeónis, sacérdos ex fíliis Jóarib ab Jerúsalem et consédit in monte Modin. Et habébat fílios quinque, Joánnem, qui cognominabátur Gaddis, et Simónem, qui cognominabátur Thasi, et Judam, qui vocabátur Machabǽus, et Eleázarum, qui cognominabátur Abaron, et Jónathan, qui cognominabátur Apphus. Hi vidérunt mala quæ fiébant in pópulo Juda et in Jerúsalem.

\R\ Hic est fratrum amátor et pópuli Israël: * Hic est, qui multum orat pro pópulo et univérsa sancta civitáte Jerúsalem. \V\ Vir iste in pópulo suo mitíssimus appáruit. — Hic est, qui.

\lectioheading{Lectio ii}{}

\lettrine[lines=2]{E}{T} dixit Mathathías: Væ mihi! ut quid natus sum vidére contritiónem pópuli mei et contritiónem civitátis sanctæ et sedére illic, cum datur in mánibus inimicórum? Sancta in manu extraneórum facta sunt, templum ejus sicut homo ignóbilis, vasa glóriæ ejus captíva abdúcta sunt, trucidáti sunt senes ejus in platéis, et júvenes ejus cecidérunt in gládio inimicórum. Quæ gens non hereditávit regnum ejus et non obtínuit spólia ejus?

\R\ Tu, Dómine universórum, qui nullam habes indigéntiam, voluísti templum tuum fíeri in nobis; * Consérva domum istam immaculátam in ætérnum, Dómine. \V\ Tu elegísti, Dómine, domum istam ad invocándum nomen tuum in ea, ut esset domus oratiónis et obsecratiónis pópulo tuo. — Consérva.

\rubric{In Octávis et Festis ritus simplícis: Glória Patri. Consérva.}

\lectioheading{Lectio iii}{Cap. 2, 14-16}

\lettrine[lines=2]{E}{T} scidit vestiménta sua Mathathías et fílii ejus et operuérunt se cilíciis et planxérunt valde. Et venérunt illuc qui missi erant a rege Antíocho, ut cógerent eos qui confúgerant in civitátem Modin, immoláre et accéndere thura et a lege Dei discédere, et multi de pópulo Israël consentiéntes accessérunt ad eos, sed Mathathías et fílii ejus constánter stetérunt.

\R\ Aperi óculos tuos, Dómine, et vide afflictiónem nostram: circumdedérunt nos gentes ad puniéndum nos: * Sed tu, Dómine, exténde bráchium tuum, et líbera ánimas nostras. \V\ Afflíge oppriméntes nos et contuméliam faciéntes in supérbiam; et custódi partem tuam. — Sed. Glória Patri. Sed.

\feriaheading{Feria Quarta infra Hebdomadam I Octobris}

\rubric{De libro primo Machabæórum}

\lectioheading{Lectio i}{Cap. 2, 19-30}

\lettrine[lines=2]{E}{T} respóndit Mathathías et dixit magna voce: Etsi omnes gentes regi Antíocho obédiunt, ut discédat unusquísque a servitúte legis patrum suórum et conséntiat mandátis ejus, ego et fílii mei et fratres mei obediémus legi patrum nostrórum. Propítius sit nobis Deus, non est nobis útile relínquere legem et justítias Dei: non audiémus verba regis Antíochi nec sacrificábimus transgrediéntes legis nostræ mandáta, ut eámus áltera via.

\R\ Refúlsit sol in clýpeos áureos, et resplenduérunt montes ab eis: * Et fortitúdo géntium dissipáta est. \V\ Erat enim exércitus magnus valde et fortis: et appropiávit Judas et exércitus ejus in prǽlio. — Et.

\lectioheading{Lectio ii}{}

\lettrine[lines=2]{E}{T}, ut cessávit loqui verba hæc, accéssit quidam Judǽus in ómnium óculis sacrificáre idólis super aram in civitáte Modin secúndum jussum regis. Et vidit Mathathías et dóluit, et contremuérunt renes ejus, et accénsus est furor ejus secúndum judícium legis, et insíliens trucidávit eum super aram. Sed et virum quem rex Antíochus míserat, qui cogébat immoláre, occídit in ipso témpore et aram destrúxit et zelátus est legem, sicut fecit Phínees Zamri fílio Salómi.

\R\ Ornavérunt fáciem templi corónis áureis, et dedicavérunt altáre Dómino: * Et facta est lætítia magna in pópulo. \V\ In hymnis et confessiónibus benedicébant Dóminum. — Et.

\rubric{In Octávis et Festis ritus simplícis: Glória Patri. Et.}

\lectioheading{Lectio iii}{}

\lettrine[lines=2]{E}{T} exclamávit Mathathías voce magna in civitáte dicens: Omnis qui zelum habet legis státuens testaméntum éxeat post me. Et fugit ipse et fílii ejus in montes, et reliquérunt quæcúmque habébant in civitáte. Tunc descendérunt multi quæréntes judícium et justítiam in desértum. Et sedérunt ibi ipsi et fílii eórum et mulíeres eórum et pécora eórum, quóniam inundavérunt super eos mala.

\R\ In hymnis et confessiónibus benedicébant Dóminum, * Qui magna fecit in Israël, et victóriam dedit illis Dóminus omnípotens. \V\ Ornavérunt fáciem templi corónis áureis, et dedicavérunt altáre Dómino. — Qui. Glória Patri. Qui.

\feriaheading{Feria Quinta infra Hebdomadam I Octobris}

\rubric{De libro primo Machabæórum}

\lectioheading{Lectio i}{Cap. 2, 49-69}

\lettrine[lines=2]{E}{T} appropinquavérunt dies Mathathíæ moriéndi, et dixit fíliis suis: Nunc confortáta est supérbia et castigátio et tempus eversiónis et ira indignatiónis. Nunc ergo, o fílii, æmulatóres estóte legis et date ánimas vestras pro testaménto patrum vestrórum et mementóte óperum patrum, quæ fecérunt in generatiónibus suis, et accipiétis glóriam magnam et nomen ætérnum. Abraham, nonne in tentatióne invéntus est fidélis, et reputátum est ei ad justítiam? Joseph in témpore angústiæ suæ custodívit mandátum, et factus est dóminus Ægýpti; Phínees pater noster zelándo zelum Dei accépit testaméntum sacerdótii ætérni.

\R\ Adapériat Dóminus cor vestrum in lege sua et in præcéptis suis et fáciat pacem in diébus vestris: * Concédat vobis salútem, et rédimat vos a malis. \V\ Exáudiat Dóminus oratiónes vestras et reconciliétur vobis nec vos déserat in témpore malo. — Concédat.

\lectioheading{Lectio ii}{}

\lettrine[lines=2]{J}{ESUS}, dum implévit verbum, factus est dux in Israël; Caleb, dum testificátur in ecclésia, accépit hereditátem; David in sua misericórdia consecútus est sedem regni in sǽcula; Elías, dum zelat zelum legis, recéptus est in cælum; Ananías et Azarías et Mísael credéntes liberáti sunt de flamma; Dániel in sua simplicitáte liberátus est de ore leónum; et ita cogitáte per generatiónem et generatiónem, quia omnes qui sperant in eum non infirmántur. Et a verbis viri peccatóris ne timuéritis, quia glória ejus stercus et vermis est: hódie extóllitur et cras non inveniétur, quia convérsus est in terram suam, et cogitátio ejus périit.

\R\ Exáudiat Dóminus oratiónes vestras, et reconciliétur vobis nec vos déserat in témpore malo * Dóminus, Deus noster. \V\ Det vobis cor ómnibus, ut colátis eum et faciátis ejus voluntátem. — Dóminus.

\rubric{In Octávis et Festis ritus simplícis: Glória Patri. Dóminus.}

\lectioheading{Lectio iii}{}

\lettrine[lines=2]{V}{OS} ergo, fílii, confortámini et viríliter ágite in lege, quia in ipsa gloriósi éritis. Et ecce Simon frater vester, scio quod vir consílii est: ipsum audíte semper, et ipse erit vobis pater; et Judas Machabǽus fortis víribus a juventúte sua sit vobis princeps militíæ et ipse aget bellum pópuli, et adducétis ad vos omnes factóres legis et vindicáte vindíctam pópuli vestri, retribúite retributiónem géntibus et inténdite in præcéptum legis. Et benedíxit eos et appósitus est ad patres suos.

\R\ Congregáti sunt inimíci nostri, et gloriántur in virtúte sua: cóntere fortitúdinem illórum, Dómine, et dispérge illos: * Ut cognóscant quia non est álius qui pugnet pro nobis, nisi tu, Deus noster. \V\ Dispérge illos in virtúte tua, et déstrue eos, protéctor noster, Dómine. — Ut. Glória Patri. Ut.

\feriaheading{Feria Sexta infra Hebdomadam I Octobris}

\rubric{De libro primo Machabæórum}

\lectioheading{Lectio i}{Cap. 2, 70; 3, 1-3, 5-12, 25-28}

\lettrine[lines=2]{D}{EFUNCTUS} est Mathathías anno centésimo et quadragésimo sexto et sepúltus est a fíliis suis in sepúlcris patrum suórum in Modin, et planxérunt eum omnis Israël planctu magno. Et surréxit Judas, qui vocabátur Machabǽus, fílius ejus pro eo et adjuvábant eum omnes fratres ejus et univérsi qui se conjúnxerant patri ejus et præliabántur prǽlium Israël cum lætítia. Et dilatávit glóriam pópulo suo, et persecútus est iníquos perscrútans eos et, qui conturbábant pópulum suum, eos succéndit flammis, et repúlsi sunt inimíci ejus præ timóre ejus, et omnes operárii iniquitátis conturbáti sunt, et dirécta est salus in manu ejus.

\R\ Impetum inimicórum ne timuéritis: mémores estóte, quómodo salvi facti sunt patres nostri: * Et nunc clamémus in cælum, et miserébitur nostri Deus noster. \V\ Mementóte mirabílium ejus, quæ fecit Pharaóni et exercítui ejus in Mari Rubro. — Et.

\lectioheading{Lectio ii}{}

\lettrine[lines=2]{E}{T} exacerbábat reges multos et lætificábat Jacob in opéribus suis, et in sǽculum memória ejus in benedictióne, et perambulávit civitátes Juda et pérdidit ímpios ex eis et avértit iram ab Israël et nominátus est usque ad novíssimum terræ et congregávit pereúntes. Et congregávit Apollónius gentes et a Samaría virtútem multam et magnam ad bellándum contra Israël. Et cognóvit Judas et éxiit óbviam illi et percússit et occídit illum, et cecidérunt vulneráti multi, et réliqui fugérunt. Et accépit spólia eórum et gládium Apollónii ábstulit Judas et erat pugnans in eo ómnibus diébus.

\R\ Congregátæ sunt gentes in multitúdine, ut dímicent contra nos, et ignorámus quid ágere debeámus: * Dómine Deus, ad te sunt óculi nostri, ne pereámus. \V\ Tu scis quæ cógitant in nos: quómodo potérimus subsístere ante fáciem illórum, nisi tu ádjuves nos? — Dómine.

\rubric{In Octávis et Festis ritus simplícis: Glória Patri. Dómine.}

\lectioheading{Lectio iii}{}

\lettrine[lines=2]{E}{T} cécidit timor Judæ ac fratrum ejus et formído super omnes gentes in circúitu eórum; et pervénit ad regem nomen ejus, et de prǽliis Judæ narrábant omnes gentes. Ut audívit autem rex Antíochus sermónes istos, irátus est ánimo et misit et congregávit exércitum univérsi regni sui, castra fórtia valde, et apéruit ærárium suum et dedit stipéndia exercítui in annum et mandávit illis, ut essent paráti ad ómnia.

\R\ Tua est poténtia, tuum regnum, Dómine: tu es super omnes gentes: * Da pacem, Dómine, in diébus nostris. \V\ Creátor ómnium, Deus, terríbilis et fortis, justus et miséricors. — Da. Glória Patri. Da.

\feriaheading{Sabbato infra Hebdomadam I Octobris}

\rubric{De libro primo Machabæórum}

\lectioheading{Lectio i}{Cap. 3, 42-60}

\lettrine[lines=2]{E}{T} vidit Judas et fratres ejus quia multiplicáta sunt mala et exércitus applicábant ad fines eórum, et cognovérunt verba regis, quæ mandávit pópulo fácere in intéritum et consummatiónem, et dixérunt unusquísque ad próximum suum: Erigámus dejectiónem pópuli nostri et pugnémus pro pópulo nostro et sanctis nostris. Et congregátus est convéntus, ut essent paráti in prǽlium et ut orárent et péterent misericórdiam et miseratiónes. Et Jerúsalem non habitabátur, sed erat sicut desértum: non erat qui ingrederétur et egrederétur de natis ejus, et sanctum conculcabátur, et fílii alienigenárum erant in arce: ibi erat habitátio géntium, et abláta est volúptas a Jacob, et defécit ibi tíbia et cíthara.

\R\ Refúlsit sol in clýpeos áureos, et resplenduérunt montes ab eis: * Et fortitúdo géntium dissipáta est. \V\ Erat enim exércitus magnus valde et fortis: et appropiávit Judas et exércitus ejus in prǽlio. — Et.

\lectioheading{Lectio ii}{}

\lettrine[lines=2]{E}{T} congregáti sunt et venérunt in Maspha contra Jerúsalem, quia locus oratiónis erat in Maspha ante in Israël, et jejunavérunt illa die et induérunt se cilíciis et cínerem imposuérunt cápiti suo et discidérunt vestiménta sua et expandérunt libros legis, de quibus scrutabántur gentes similitúdinem simulacrórum suórum, et attulérunt ornaménta sacerdotália et primítias et décimas et suscitavérunt Nazarǽos, qui impléverant dies, et clamavérunt voce magna in cælum dicéntes: Quid faciémus istis et quo eos ducémus? Et sancta tua conculcáta sunt et contamináta sunt, et sacerdótes tui facti sunt in luctum et in humilitátem, et ecce natiónes convenérunt advérsum nos, ut nos dispérdant: tu scis quæ cógitant in nos. Quómodo potérimus subsístere ante fáciem eórum, nisi tu, Deus, ádjuves nos?

\R\ Ornavérunt fáciem templi corónis áureis, et dedicavérunt altáre Dómino: * Et facta est lætítia magna in pópulo. \V\ In hymnis et confessiónibus benedicébant Dóminum. — Et.

\rubric{In Octávis simplícibus et in Offício sanctæ Maríæ in Sábbato: Glória Patri. Et.}

\lectioheading{Lectio iii}{}

\lettrine[lines=2]{E}{T} tubis exclamavérunt voce magna. Et post hæc constítuit Judas duces pópuli, tribúnos et centuriónes et pentacontárchos et decuriónes, et dixit his, qui ædificábant domos et sponsábant uxóres et plantábant víneas et formidolósis, ut redírent unusquísque in domum suam secúndum legem. Et movérunt castra et collocavérunt ad Austrum Emmaum. Et ait Judas: Accingímini et estóte fílii poténtes et estóte paráti in mane, ut pugnétis advérsus natiónes has, quæ convenérunt advérsus nos dispérdere nos et sancta nostra; quóniam mélius est nos mori in bello, quam vidére mala gentis nostræ et sanctórum. Sicut autem fúerit volúntas in cælo, sic fiat.

\R\ In hymnis et confessiónibus benedicébant Dóminum, * Qui magna fecit in Israël, et victóriam dedit illis Dóminus omnípotens. \V\ Ornavérunt fáciem templi corónis áureis, et dedicavérunt altáre Dómino. — Qui. Glória Patri. Qui.

\hourheading{Ad Vesperas}

\V\ Vespertína orátio ascéndat ad te, Dómine. \R\ Et descéndat super nos misericórdia tua.

\antlabel{Ad Magnif. Ant.} Refúlsit sol * in clýpeos áureos, et resplenduérunt montes ab eis: et fortitúdo géntium dissipáta est.

\rubric{Oratio Dominicæ occurréntis.}

% ===== BATCH 8 of the Proprium de Tempore (autumnalis.md lines
% 14884-15116) -- the complete first week of October: "Dominica I
% Octobris" Matins opening the First Book of Machabees in Nocturn I
% and St. Ambrose's "De Officiis" in Nocturn II, the six ferial days
% continuing the Machabees cycle through Mathathias's death and
% Judas's early victories, and the closing Saturday Vespers. Both
% Nocturn III's responsories (as at the two immediately preceding
% Sundays) are given without the "vii"/"viii" numbering used at
% Dominicas I-IV Septembris. A new governing rubric fixes the hymn
% "Ætérne rerum" at Lauds for every Sunday through November 5. One
% genuine partial-caps OCR artifact is silently normalized
% (DEfúnctus to DEFUNCTUS, accent dropped on its newly-capitalized
% remainder), while complete short words already fully capped in
% the source (VOS, JESUS) are left untouched. Feria Sexta's citation
% ("Cap. 2, 70; 3, 1-3, 5-12, 25-28") appears in the source one line
% above its "Lectio i" header rather than below it as usual; it is
% nonetheless given as that Lesson's citation via
% \lectioheading{Lectio i}{...}, consistent with its evident
% function regardless of its exact placement in the source. The
% pervasive space-before-colon formatting quirk from the previous
% two batches continues here and is likewise silently normalized.
% The document continues with "Dominica II Octobris" (autumnalis.md
% line 15117).
% =====
       % or hiemalis.tex / verna.tex / aestiva.tex
%   \end{document}
%
% Since autumnalis.tex, aestiva.tex, hiemalis.tex, and verna.tex all
% now embed the versicle/response marks (℣ U+2123, ℟ U+211F) and the
% blessing cross (☩ U+2629) as literal Unicode characters rather than
% as macros, and the main body font (DejaVu Serif) does not contain
% glyphs for them, \newunicodechar is used below to substitute a font
% that does (DejaVu Sans) automatically wherever those three
% characters occur, without having to touch the body text itself.
% =====================================================================

\usepackage{fontspec}
\usepackage{textcomp}
\usepackage{lettrine}
\usepackage[dvipsnames]{xcolor}
\usepackage[margin=2cm]{geometry}
\usepackage{parskip}
\usepackage{newunicodechar}

\setmainfont{DejaVu Serif}
\newfontfamily\liturgicalsymbolfont{DejaVu Sans}

\newunicodechar{℣}{{\liturgicalsymbolfont\char"2123}}  % VERSICLE
\newunicodechar{℟}{{\liturgicalsymbolfont\char"211F}}  % RESPONSE
\newunicodechar{☩}{{\liturgicalsymbolfont\char"2629}}  % cross (blessing)

% --- Headings (one required argument each) ---
\newcommand{\partheading}[1]{{\bigskip\noindent\LARGE\bfseries #1\par\medskip}}
\newcommand{\partsubtitle}[1]{{\noindent\large\itshape #1\par\medskip}}
\newcommand{\sectionheading}[1]{{\bigskip\noindent\Large\bfseries #1\par\medskip}}
\newcommand{\feriaheading}[1]{{\bigskip\noindent\large\bfseries #1\par\smallskip}}
\newcommand{\hourheading}[1]{{\bigskip\noindent\large\bfseries #1\par\smallskip}}
\newcommand{\nocturnoheading}[1]{{\medskip\noindent\bfseries #1\par\smallskip}}
\newcommand{\psalmheading}[1]{{\medskip\noindent\itshape #1\par\smallskip}}
\newcommand{\hymnusheading}[1]{{\medskip\noindent\bfseries #1\par\smallskip}}
\newcommand{\oratioheading}[1]{{\medskip\noindent\itshape #1\par\smallskip}}
\newcommand{\variantheading}[1]{{\medskip\noindent\itshape #1\par\smallskip}}
\newcommand{\absolutioheading}[1]{{\medskip\noindent\itshape #1\par\smallskip}}
\newcommand{\suffragiumheading}[1]{{\medskip\noindent\bfseries #1\par\smallskip}}
\newcommand{\symbolumheading}[1]{{\medskip\noindent\bfseries #1\par\smallskip}}
\newcommand{\antiphonaheading}[1]{{\medskip\noindent\bfseries #1\par\smallskip}}
\newcommand{\feastheading}[1]{{\bigskip\noindent\large\bfseries #1\par\smallskip}}
\newcommand{\sanctusheading}[1]{{\bigskip\noindent\large\bfseries #1\par\smallskip}}
\newcommand{\diesheading}[1]{{\medskip\noindent\bfseries #1\par\smallskip}}
\newcommand{\monthheading}[1]{{\bigskip\noindent\Large\bfseries #1\par\medskip}}

% --- Purely decorative divider (no argument) ---
\newcommand{\sectionornament}{{\par\noindent\centering\rule{0.3\linewidth}{0.4pt}\par\medskip}}

% --- Headings (two required arguments: title + citation) ---
\newcommand{\lectioheading}[2]{{\medskip\noindent\itshape #1\quad\textit{#2}\par\smallskip}}
\newcommand{\capitulumlabel}[2]{{\medskip\noindent\itshape #1\quad\textit{#2}\par\smallskip}}
\newcommand{\canticumheading}[2]{{\medskip\noindent\itshape #1\quad\textit{#2}\par\smallskip}}
\newcommand{\lectiolabel}[2]{{\medskip\noindent\itshape #1\quad\textit{#2}\par\smallskip}}

% --- Labels (one required argument) ---
\newcommand{\antlabel}[1]{\textbf{#1}}
\newcommand{\invitatlabel}[1]{\textbf{#1}}
\newcommand{\psalmlabel}[1]{\textbf{#1}}
\newcommand{\benedlabel}[1]{\textbf{#1}}
\newcommand{\benedictiolabel}[1]{\textbf{#1}}
\newcommand{\absolutiolabel}[1]{\textbf{#1}}
\newcommand{\formulalabel}[1]{\textbf{#1}}
\newcommand{\oratiolabel}[1]{\textbf{#1}}

% --- Rubric (red text; embedded rubrics may nest \textit, etc.) ---
\newcommand{\rubric}[1]{\textcolor{red}{#1}}
% Black text embedded within a red \rubric{} span (governing instruction:
% "embedded rubrics (black text within red-colored rubrics) are preserved").
\newcommand{\rubricblack}[1]{\textcolor{black}{#1}}

% \P (pilcrow), \dag, \ddag, \label, \pageref, \hspace, \textbf,
% \textsuperscript, \textit, \textsc, \textdagger, \textdaggerdbl are all
% standard LaTeX/textcomp commands and need no definition here.
